% !TeX program = xelatex
% Generated by tools/flatten.py; edit the native project and regenerate.
\begin{filecontents*}[overwrite]{SGA4-native-amsbook.cls}
%%
%% This is file `amsbook.cls',
%% generated with the docstrip utility.
%%
%% The original source files were:
%%
%% amsclass.dtx  (with options: `amsbook,classes')
%% 
%% This is a generated file.
%% 
%% Copyright 1995, 199, 2004, 2009 American Mathematical Society.
%% 
%% American Mathematical Society
%% Technical Support
%% Publications Technical Group
%% 201 Charles Street
%% Providence, RI 02904
%% USA
%% tel: (401) 455-4080
%%      (800) 321-4267 (USA and Canada only)
%% fax: (401) 331-3842
%% email: tech-support@ams.org
%% 
%% Unlimited copying and redistribution of this file are permitted as
%% long as this file is not modified.  Modifications, and distribution
%% of modified versions, are permitted, but only if the resulting file
%% is renamed.
%% 
\NeedsTeXFormat{LaTeX2e}% LaTeX 2.09 can't be used (nor non-LaTeX)
[1995/06/01]% LaTeX date must be June 1995 or later
\ProvidesClass{amsbook}[2009/07/02 v2.20.1]
\global\expandafter\let\csname ver@amsthm.sty\expandafter\endcsname
  \csname ver@\@currname.\@currext\endcsname
\let\@xp=\expandafter
\let\@nx=\noexpand
\def\@oparg#1[#2]{\@ifnextchar[{#1}{#1[#2]}}
\long\def\@ifempty#1{\@xifempty#1@@..\@nil}
\long\def\@xifempty#1#2@#3#4#5\@nil{%
  \ifx#3#4\@xp\@firstoftwo\else\@xp\@secondoftwo\fi}
\long\def\@ifnotempty#1{\@ifempty{#1}{}}
\def\setboxz@h{\setbox\z@\hbox}
\def\@addpunct#1{%
  \relax\ifhmode
    \ifnum\spacefactor>\@m \else#1\fi
  \fi}
\def\nopunct{\spacefactor 1007 }
\def\frenchspacing{\sfcode`\.1006\sfcode`\?1005\sfcode`\!1004%
  \sfcode`\:1003\sfcode`\;1002\sfcode`\,1001 }
\def\@tempa#1#2\@nil{\edef\@classname{#1}}
\expandafter\@tempa\@currnamestack{}{}{}\@nil
\ifx\@classname\@empty \edef\@classname{\@currname}\fi
\def\@True{00}
\def\@False{01}
\newcommand\newswitch[2][False]{%
  \expandafter\@ifdefinable\csname ?@#2\endcsname{%
    \global\expandafter\let\csname ?@#2\expandafter\endcsname
      \csname @#1\endcsname
  }%
}
\newcommand{\setFalse}[1]{%
  \expandafter\let\csname ?@#1\endcsname\@False
}
\newcommand{\setTrue}[1]{%
  \expandafter\let\csname ?@#1\endcsname\@True
}
\newswitch{}
\DeclareRobustCommand{\except}[1]{%
  \if\csname ?@#1\endcsname \expandafter\@gobble
  \else \expandafter\@firstofone
  \fi
}
\DeclareRobustCommand{\for}[1]{%
  \if\csname ?@#1\endcsname \expandafter\@firstofone
  \else \expandafter\@gobble
  \fi
}
\DeclareRobustCommand{\forany}[1]{%
  \csname for@any@01\endcsname#1,?,\@nil
}
\@namedef{for@any@\@False}#1,{%
  \csname for@any@%
    \csname ?@\zap@space#1 \@empty\endcsname
  \endcsname
}
\@namedef{?@?}{x}
\@namedef{for@any@\@True}#1\@nil#2{#2}
\def\for@any@x{\@car\@gobble}
\DeclareOption{a4paper}{\paperheight 297mm\paperwidth 210mm
  \textheight 54.5pc }
\DeclareOption{letterpaper}{\paperheight 11in\paperwidth 8.5in }
\DeclareOption{landscape}{\@tempdima\paperheight
  \paperheight\paperwidth \paperwidth\@tempdima}
\DeclareOption{portrait}{}
\DeclareOption{oneside}{\@twosidefalse \@mparswitchfalse}
\DeclareOption{twoside}{\@twosidetrue  \@mparswitchtrue}
\DeclareOption{draft}{\overfullrule5\p@
  \ClassWarningNoLine{\@classname}{%
    When the draft option is used, the
      \protect\includegraphics\MessageBreak
    command will print blank placeholder boxes\MessageBreak
    for the graphics}%
}
\DeclareOption{final}{\overfullrule\z@ }
\def\dateposted#1{\def\@dateposted{#1}}%
\let\@dateposted\@empty
\newif\if@titlepage
\DeclareOption{titlepage}{\@titlepagetrue}
\DeclareOption{notitlepage}{\@titlepagefalse}
\newif\if@openright
\DeclareOption{openright}{\@openrighttrue}
\DeclareOption{openany}{\@openrightfalse}
\@openrighttrue
\DeclareOption{onecolumn}{\@twocolumnfalse}
\DeclareOption{twocolumn}{\@twocolumntrue}
\DeclareOption{nomath}{}
\DeclareOption{noamsfonts}{}
\DeclareOption{psamsfonts}{%
  \PassOptionsToPackage{psamsfonts}{amsfonts}%
  \PassOptionsToPackage{cmex10}{amsmath}}
\newif\iftagsleft@
\DeclareOption{leqno}{%
  \tagsleft@true \PassOptionsToPackage{leqno}{amsmath}}
\DeclareOption{reqno}{%
  \tagsleft@false \PassOptionsToPackage{reqno}{amsmath}}
\newif\ifctagsplit@
\DeclareOption{centertags}{%
  \ctagsplit@true \PassOptionsToPackage{centertags}{amsmath}}
\DeclareOption{tbtags}{%
  \ctagsplit@false \PassOptionsToPackage{tbtags}{amsmath}}
\DeclareOption{fleqn}{}%
\newcommand{\@mainsize}{10}
\newcommand{\@ptsize}{0}
\newcommand{\larger}[1][1]{%
  \count@\@currsizeindex \advance\count@#1\relax
  \ifnum\count@<\z@ \count@\z@ \else\ifnum\count@>12 \count@12 \fi\fi
  \ifcase\count@
    \Tiny\or\Tiny\or\tiny\or\SMALL\or\Small\or\small
    \or\normalsize
    \or\large\or\Large\or\LARGE\or\huge\or\Huge\else\Huge
  \fi
}
\newcommand{\smaller}[1][1]{\larger[-#1]}
\def\@adjustvertspacing{%
  \bigskipamount.7\baselineskip plus.7\baselineskip
  \medskipamount\bigskipamount \divide\medskipamount\tw@
  \smallskipamount\medskipamount \divide\smallskipamount\tw@
  \abovedisplayskip\medskipamount
  \belowdisplayskip \abovedisplayskip
  \abovedisplayshortskip\abovedisplayskip
  \advance\abovedisplayshortskip-1\abovedisplayskip
  \belowdisplayshortskip\abovedisplayshortskip
  \advance\belowdisplayshortskip 1\smallskipamount
  \jot\baselineskip \divide\jot 4 \relax
}
\renewcommand\normalsize{\@xsetfontsize\normalsize 6%
  \@adjustvertspacing \let\@listi\@listI}
\DeclareRobustCommand{\Tiny}{\@xsetfontsize\Tiny 1}
\DeclareRobustCommand{\tiny}{\@xsetfontsize\tiny 2}
\DeclareRobustCommand{\SMALL}{\@xsetfontsize\SMALL 3}
\DeclareRobustCommand{\Small}{\@xsetfontsize\Small 4%
  \@adjustvertspacing
  \def\@listi{\topsep\smallskipamount \parsep\z@skip \itemsep\z@skip
    \leftmargin=\leftmargini
    \labelwidth=\leftmargini \advance\labelwidth-\labelsep
  }%
}
\DeclareRobustCommand{\small}{\@xsetfontsize\small 5\@adjustvertspacing}
\def\footnotesize{\Small}
\def\scriptsize{\SMALL}
\DeclareRobustCommand{\large}{\@xsetfontsize\large 7\@adjustvertspacing}
\DeclareRobustCommand{\Large}{\@xsetfontsize\Large 8\@adjustvertspacing}
\DeclareRobustCommand{\LARGE}{\@xsetfontsize\LARGE 9}
\DeclareRobustCommand{\huge}{\@xsetfontsize\huge{10}}
\DeclareRobustCommand{\Huge}{\@xsetfontsize\Huge{11}}
\def\@xsetfontsize#1#2{%
  \chardef\@currsizeindex#2\relax
  \edef\@tempa{\@nx\@setfontsize\@nx#1%
    \@xp\ifcase\@xp\@currsizeindex\@typesizes
      \else{99}{99}\fi}%
  \@tempa
}
\chardef\@currsizeindex=6
\widowpenalty=10000
\clubpenalty=10000
\brokenpenalty=10000
\newdimen\linespacing
\lineskip=1pt \lineskiplimit=1pt
\normallineskip=1pt \normallineskiplimit=1pt
\let\baselinestretch=\@empty
\headheight=8pt \headsep=14pt
\footskip=18pt
\textheight=50.5pc \topskip=10pt
\textwidth=30pc
\columnsep=10pt \columnseprule=0pt
\marginparwidth=90pt
\marginparsep=11pt
\marginparpush=5pt
\AtBeginDocument{\settoheight{\footnotesep}{\footnotesize M$^1$}}
\skip\footins=7pt plus11pt
\skip\@mpfootins=\skip\footins
\fboxsep=3pt \fboxrule=.4pt
\arrayrulewidth=.4pt    \doublerulesep=2pt
\labelsep=5pt           \arraycolsep=\labelsep
\tabcolsep=\labelsep    \tabbingsep=\labelsep
\floatsep=15pt plus 12pt    \dblfloatsep=15pt plus 12pt
\textfloatsep=\floatsep  \dbltextfloatsep=15pt plus 12pt
\intextsep=\floatsep
\@fptop=0pt plus1fil     \@dblfptop=0pt plus1fil
\@fpbot=0pt plus1fil     \@dblfpbot=0pt plus1fil
\@fpsep=8pt plus2fil     \@dblfpsep=8pt plus2fil\relax
\parskip=0pt \relax
\newdimen\normalparindent
\normalparindent=18pt
\parindent=\normalparindent
\partopsep=0pt \relax \parsep=0pt \relax \itemsep=0pt \relax
\@lowpenalty=51         \@medpenalty=151        \@highpenalty=301
\@beginparpenalty=-\@lowpenalty
\@endparpenalty=-\@lowpenalty
\@itempenalty=-\@lowpenalty
\DeclareOption{10pt}{\def\@mainsize{10}\def\@ptsize{0}%
  \def\@typesizes{%
    \or{5}{6}\or{6}{7}\or{7}{8}\or{8}{10}\or{9}{11}%
    \or{10}{12}% normalsize
    \or{\@xipt}{13}\or{\@xiipt}{14}\or{\@xivpt}{17}%
    \or{\@xviipt}{20}\or{\@xxpt}{24}}%
  \normalsize \linespacing=\baselineskip
}
\DeclareOption{11pt}{\def\@mainsize{11}\def\@ptsize{1}%
  \def\@typesizes{%
    \or{6}{7}\or{7}{8}\or{8}{10}\or{9}{11}\or{10}{12}%
    \or{\@xipt}{13}% normalsize
    \or{\@xiipt}{14}\or{\@xivpt}{17}\or{\@xviipt}{20}%
    \or{\@xxpt}{24}\or{\@xxvpt}{30}}%
  \normalsize \linespacing=\baselineskip
}
\DeclareOption{12pt}{\def\@mainsize{12}\def\@ptsize{2}%
  \def\@typesizes{%
    \or{7}{8}\or{8}{10}\or{9}{11}\or{10}{12}\or{\@xipt}{13}%
    \or{\@xiipt}{14}% normalsize
    \or{\@xivpt}{17}\or{\@xviipt}{20}\or{\@xxpt}{24}%
    \or{\@xxvpt}{30}\or{\@xxvpt}{30}}%
  \normalsize \linespacing=\baselineskip
}
\DeclareOption{8pt}{\def\@mainsize{8}\def\@ptsize{8}%
  \def\@typesizes{%
    \or{5}{6}\or{5}{6}\or{5}{6}\or{6}{7}\or{7}{8}%
    \or{8}{10}% normalsize
    \or{9}{11}\or{10}{12}\or{\@xipt}{13}%
    \or{\@xiipt}{14}\or{\@xivpt}{17}}%
  \normalsize \linespacing=\baselineskip
}
\DeclareOption{9pt}{\def\@mainsize{9}\def\@ptsize{9}%
  \def\@typesizes{%
    \or{5}{6}\or{5}{6}\or{6}{7}\or{7}{8}\or{8}{10}%
    \or{9}{11}% normalsize
    \or{10}{12}\or{\@xipt}{13}\or{\@xiipt}{14}%
    \or{\@xivpt}{17}\or{\@xviipt}{20}}%
  \normalsize \linespacing=\baselineskip
}
\def\ps@empty{\let\@mkboth\@gobbletwo
  \let\@oddhead\@empty \let\@evenhead\@empty
  \let\@oddfoot\@empty \let\@evenfoot\@empty
  \global\topskip\normaltopskip}
\def\ps@plain{\ps@empty
  \def\@oddfoot{\normalfont\scriptsize \hfil\thepage\hfil}%
  \let\@evenfoot\@oddfoot}
\newswitch{runhead}
\def\ps@headings{\ps@empty
  \def\@evenhead{%
    \setTrue{runhead}%
    \normalfont\scriptsize
    \rlap{\thepage}\hfil
    \def\thanks{\protect\thanks@warning}%
    \leftmark{}{}\hfil}%
  \def\@oddhead{%
    \setTrue{runhead}%
    \normalfont\scriptsize \hfil
    \def\thanks{\protect\thanks@warning}%
    \rightmark{}{}\hfil \llap{\thepage}}%
  \let\@mkboth\markboth
  \def\partmark{\@secmark\markboth\partrunhead\partname}%
  \def\chaptermark{%
    \@secmark\markboth\chapterrunhead{}}%
  \def\sectionmark{%
    \@secmark\markright\sectionrunhead\sectionname}%
}
\let\sectionname\@empty
\let\subsectionname\@empty
\let\subsubsectionname\@empty
\let\paragraphname\@empty
\let\subparagraphname\@empty
\def\leftmark{\expandafter\@firstoftwo\topmark{}{}}
\def\rightmark{\expandafter\@secondoftwo\botmark{}{}}
\long\def\@nilgobble#1\@nil{}
\def\@secmark#1#2#3#4{%
  \begingroup \let\protect\@unexpandable@protect
  \edef\@tempa{\endgroup \toks@{\protect#2{#3}{\@secnumber}}}%
  \@tempa
  \toks@\@xp{\the\toks@{#4}}%
  \afterassignment\@nilgobble\@temptokena\@themark{}\@nil
  \edef\@tempa{\@nx\@mkboth{%
    \ifx\markright#1\the\@temptokena\else\the\toks@\fi}{\the\toks@}}%
  \@tempa}
\let\@secnumber\@empty
\def\markboth#1#2{%
  \begingroup
    \@temptokena{{#1}{#2}}\xdef\@themark{\the\@temptokena}%
    \mark{\the\@temptokena}%
  \endgroup
  \if@nobreak\ifvmode\nobreak\fi\fi}
\def\ps@myheadings{\ps@headings \let\@mkboth\@gobbletwo}
\newskip\normaltopskip
\normaltopskip=10pt \relax
\let\sectionmark\@gobble
\let\subsectionmark\@gobble
\let\subsubsectionmark\@gobble
\let\paragraphmark\@gobble

\DeclareOption{makeidx}{}
\ExecuteOptions{leqno,centertags,letterpaper,portrait,%
  10pt,twoside,onecolumn,final}
\ProcessOptions\relax
\if@compatibility
  \def\@tempa{\RequirePackage{amstex}\relax}%
\else
  \@ifclasswith{\@classname}{nomath}{%
    \let\@tempa\relax
  }{%
    \def\@tempa{\RequirePackage{amsmath}\relax}%
  }%
\fi
\@tempa % load amstex.sty or amsmath.sty
\@ifundefined{numberwithin}{%
  \newcommand{\numberwithin}[3][\arabic]{%
    \@ifundefined{c@#2}{\@nocounterr{#2}}{%
      \@ifundefined{c@#3}{\@nocnterr{#3}}{%
        \@addtoreset{#2}{#3}%
        \@xp\xdef\csname the#2\endcsname{%
          \@xp\@nx\csname the#3\endcsname .\@nx#1{#2}}}}%
  }
  \csname newtoks\endcsname\@emptytoks
}{}
\if@compatibility
\else
  \@ifclasswith{\@classname}{noamsfonts}{%
    % amsfonts package is not wanted
  }{%
    % amsfonts package IS wanted; test whether a recent enough version
    % seems to be installed
    \begingroup \fontencoding{U}\fontfamily{msa}\try@load@fontshape\endgroup
    \global\@xp\let\csname U+msa\endcsname\relax % reset
    \@ifundefined{U/msa/m/n}{%
      \ClassError{\@classname}{%
  Package `amsfonts' not installed, or version too old?\MessageBreak
  Unable to get font info for the `msam' fonts in the expected form%
      }{%
  The amsfonts package will not be loaded, to avoid probable\MessageBreak
  incompatibility problems. You can (a) use the `noamsfonts'
      documentclass\MessageBreak
  option next time, or (b) check that the amsfonts package is
      installed\MessageBreak
  correctly, and is not too old to be compatible.%
      }%
    }{%
      \RequirePackage{amsfonts}[1995/01/01]\relax
    }%
  }
\fi % end yesamsfonts branch
\let\cleardouble@page\cleardoublepage
\AtBeginDocument{%
  \ifx\cleardouble@page\cleardoublepage
    \def\cleardoublepage{\clearpage{\pagestyle{empty}\cleardouble@page}}
  \fi
}
\newcommand{\uppercasenonmath}[1]{\toks@\@emptytoks
  \@xp\@skipmath\@xp\@empty#1$$%
  \edef#1{{\@nx\protect\@nx\@upprep\the\toks@}}%
}
\newcommand{\@upprep}{%
  \spaceskip1.3\fontdimen2\font plus1.3\fontdimen3\font
  \upchars@}
\newcommand{\upchars@}{%
  \def\ss{SS}\def\i{I}\def\j{J}\def\ae{\AE}\def\oe{\OE}%
  \def\o{\O}\def\aa{\AA}\def\l{\L}\def\Mc{M{\scshape c}}}
\providecommand{\Mc}{Mc}
\newcommand{\@skipmath}{}
\long\def\@skipmath#1$#2${%
  \@xskipmath#1\(\)%
  \@ifnotempty{#2}{\toks@\@xp{\the\toks@$#2$}\@skipmath\@empty}}%
\newcommand{\@xskipmath}{}
\long\def\@xskipmath#1\(#2\){%
  \uppercase{\toks@\@xp\@xp\@xp{\@xp\the\@xp\toks@#1}}%
  \@ifnotempty{#2}{\toks@\@xp{\the\toks@\(#2\)}\@xskipmath\@empty}}%
\def\altucnm#1{%
  \MakeTextUppercase{\toks@{#1}}%
  \edef#1{\the\toks@}%
}
\AtBeginDocument{%
  \@ifundefined{MakeTextUppercase}{}{\let\uppercasenonmath\altucnm}%
}
\@ifundefined{MakeUppercase}{\let\MakeUppercase\uppercase}{}%
\newcommand{\today}{%
  \relax\ifcase\month\or
  January\or February\or March\or April\or May\or June\or
  July\or August\or September\or October\or November\or December\fi
  \space\number\day, \number\year}
\DeclareOldFontCommand{\rm}{\normalfont\rmfamily}{\mathrm}
\DeclareOldFontCommand{\sf}{\normalfont\sffamily}{\mathsf}
\DeclareOldFontCommand{\tt}{\normalfont\ttfamily}{\mathtt}
\DeclareOldFontCommand{\bf}{\normalfont\bfseries}{\mathbf}
\DeclareOldFontCommand{\it}{\normalfont\itshape}{\mathit}
\DeclareOldFontCommand{\sl}{\normalfont\slshape}{\@nomath\sl}
\DeclareOldFontCommand{\sc}{\normalfont\scshape}{\@nomath\sc}
%%\if@compatibility
%%\else
%%  \def\@obsolete@fontswitch#1#2#3{%
%%    \@latex@warning@no@line{%
%%      Command {\string#1...}\on@line\space is obsolete;\MessageBreak
%%      the LaTeX2e equivalent is \string#3{...}}%
%%    \gdef#1{\@fontswitch\relax#3}%
%%  }
%%  \DeclareRobustCommand*\cal{%
%%    \@xp\@obsolete@fontswitch\csname cal \endcsname\relax\mathcal}
%%  \DeclareRobustCommand*\mit{%
%%    \@xp\@obsolete@fontswitch\csname mit \endcsname\relax\mathnormal}
%%\fi
\renewcommand*{\title}[2][]{\gdef\shorttitle{#1}\gdef\@title{#2}}
\edef\title{\@nx\@dblarg
  \@xp\@nx\csname\string\title\endcsname}
\renewcommand{\author}[2][]{%
  \ifx\@empty\authors
    \gdef\authors{#2}%
  \else
    \g@addto@macro\authors{\and#2}%
    \g@addto@macro\addresses{\author{}}%
  \fi
  \@ifnotempty{#1}{%
    \ifx\@empty\shortauthors
      \gdef\shortauthors{#1}%
    \else
      \g@addto@macro\shortauthors{\and#1}%
    \fi
  }%
}
\edef\author{\@nx\@dblarg
  \@xp\@nx\csname\string\author\endcsname}
\let\shortauthors\@empty   \let\authors\@empty
\let\addresses\@empty      \let\thankses\@empty
\newcommand{\address}[2][]{\g@addto@macro\addresses{\address{#1}{#2}}}
\newcommand{\curraddr}[2][]{\g@addto@macro\addresses{\curraddr{#1}{#2}}}
\newcommand{\email}[2][]{\g@addto@macro\addresses{\email{#1}{#2}}}
\newcommand{\urladdr}[2][]{\g@addto@macro\addresses{\urladdr{#1}{#2}}}
\long\def\thanks@warning#1{%
  \ClassError{\@classname}{%
    \protect\thanks\space should be given separately, not inside author name.%
  }\@ehb
}
\renewcommand{\thanks}[1]{%
  \@ifnotempty{#1}{\g@addto@macro\thankses{\thanks{#1}}}%
}
\def\curraddrname{{\itshape Current address}}
\def\emailaddrname{{\itshape E-mail address}}
\def\urladdrname{{\itshape URL}}
\def\@setaddresses{\par
  \nobreak \begingroup
  \def\author##1{\nobreak\addvspace\bigskipamount}%
  \def\\{\unskip, \ignorespaces}%
  \interlinepenalty\@M
  \def\address##1##2{\begingroup
    \par\addvspace\bigskipamount\indent
    \@ifnotempty{##1}{(\ignorespaces##1\unskip) }%
    {\scshape\ignorespaces##2}\par\endgroup}%
  \def\curraddr##1##2{\begingroup
    \@ifnotempty{##2}{\nobreak\indent\curraddrname
      \@ifnotempty{##1}{, \ignorespaces##1\unskip}\/:\space
      ##2\par}\endgroup}%
  \def\email##1##2{\begingroup
    \@ifnotempty{##2}{\nobreak\indent\emailaddrname
      \@ifnotempty{##1}{, \ignorespaces##1\unskip}\/:\space
      \ttfamily##2\par}\endgroup}%
  \def\urladdr##1##2{\begingroup
    \def~{\char`\~}%
    \@ifnotempty{##2}{\nobreak\indent\urladdrname
      \@ifnotempty{##1}{, \ignorespaces##1\unskip}\/:\space
      \ttfamily##2\par}\endgroup}%
  \addresses
  \endgroup
}
\let\@date\@empty
\def\dedicatory#1{\def\@dedicatory{#1}}
\let\@dedicatory=\@empty
\def\keywords#1{\def\@keywords{#1}}
\let\@keywords=\@empty
\newcommand*\subjclass[2][1991]{%
  \def\@subjclass{#2}%
  \@ifundefined{subjclassname@#1}{%
    \ClassWarning{\@classname}{Unknown edition (#1) of Mathematics
      Subject Classification; using '1991'.}%
  }{%
    \@xp\let\@xp\subjclassname\csname subjclassname@#1\endcsname
  }%
}
\let\@subjclass=\@empty
\def\translname{Translated by}
\def\translator#1{%
  \ifx\@empty\@translators \def\@translators{#1}%
  \else\g@addto@macro\@translators{\and#1}\fi}
\let\@translators=\@empty
\newcommand{\xandlist}[4]{\@andlista{{#1}{#2}{#3}}#4\and\and}
\def\@andlista#1#2\and#3\and{\@andlistc{#2}\@ifnotempty{#3}{%
  \@andlistb#1{#3}}}
\def\@andlistb#1#2#3#4#5\and{%
  \@ifempty{#5}{%
    \@andlistc{#2#4}%
  }{%
    \@andlistc{#1#4}\@andlistb{#1}{#3}{#3}{#5}%
  }}
\let\@andlistc\@iden
\newcommand{\nxandlist}[4]{%
  \def\@andlistc##1{\toks@\@xp{\the\toks@##1}}%
  \toks@{\toks@\@emptytoks \@andlista{{#1}{#2}{#3}}}%
  \the\@xp\toks@#4\and\and
  \edef#4{\the\toks@}%
  \let\@andlistc\@iden}
\def\@@and{and}
\newcommand{\andify}{%
  \nxandlist{\unskip, }{\unskip{} \@@and~}{\unskip, \@@and~}}
\def\and{\unskip{ }\@@and{ }\ignorespaces}
\newcommand{\copyrightinfo}[2]{%
  \def\copyrightyear{#1}%
  \@ifnotempty{#2}{\def\copyrightholder{#2}}%
}
\copyrightinfo{0000}{(copyright holder)}
\newinsert\copyins
\skip\copyins=1.5pc
\count\copyins=1000 % magnification factor, 1000 = 100%
\dimen\copyins=.5\textheight % maximum allowed per page
\g@addto@macro\@reinserts{%
  \ifvoid\copyins\else\insert\copyins{\unvbox\copyins}\fi
}
\def\@copyinsfontsize{\fontsize{6}{7\p@}\normalfont\upshape}
\newif\if@extracrline  \@extracrlinefalse
\let\@extracrline\@empty
\relax
\def\@setcopyright{%
  \ifx\@empty\@serieslogo
  \else\ifx\@empty\copyrightyear
  \else
    \insert\copyins{\hsize\textwidth
      \parfillskip\z@\relax
      \leftskip\z@\@plus.9\textwidth\relax \rightskip\z@\relax
      \@copyinsfontsize
      \everypar{}%
      \vskip-\skip\copyins
      \if@extracrline
        \vskip-6pt
      \fi
      \nointerlineskip
      \leavevmode\hfill\vrule\@width\z@\@height\skip\copyins
      \copyright\copyrightyear\ \copyrightholder\ignorespaces
      \if@extracrline \@extracrline \fi
      \par
      \kern\z@}%
  \fi\fi
}
\def\@combinefloats{%
  \ifx \@toplist\@empty \else \@cflt \fi
  \ifx \@botlist\@empty \else \@cflb \fi
  \ifvoid\copyins \else \@cflci \fi
}
\def\@cflci{%
  \setbox\@outputbox\vbox{%
    \unvbox\@outputbox
    \vskip\skip\copyins
    \if@twocolumn \else \vskip\z@ plus\p@ \fi
    \hbox to\columnwidth{%
      \hss\vbox to\z@{\vss
        \if@twocolumn
          \if@firstcolumn \else \unvbox\copyins \fi
        \else
          \unvbox\copyins
        \fi
        }}}%
  \if@twocolumn \if@firstcolumn
    \insert\copyins{\unvbox\copyins}%
  \fi\fi
  \global\count\copyins=999 \relax
}
\newcommand{\abstractname}{Abstract}
\newcommand{\keywordsname}{Key words and phrases}
\newcommand{\subjclassname}{%
  \textup{1991} Mathematics Subject Classification}
\@xp\let\csname subjclassname@1991\endcsname \subjclassname
\@namedef{subjclassname@2000}{%
  \textup{2000} Mathematics Subject Classification}
\@namedef{subjclassname@2010}{%
  \textup{2010} Mathematics Subject Classification}
\def\@tempb{amsbook}
\ifx\@classname\@tempb
  \newcommand{\datename}{\textit{Date}:}
\else
  \newcommand{\datename}{Received by the editors}
\fi
\def\@setthanks{\def\thanks##1{\par##1\@addpunct.}\thankses}
\newbox\abstractbox
\newenvironment{abstract}{%
  \ifx\maketitle\relax
    \ClassWarning{\@classname}{Abstract should precede
      \protect\maketitle\space in AMS document classes; reported}%
  \fi
  \global\setbox\abstractbox=\vtop \bgroup
    \normalfont\Small
    \list{}{\labelwidth\z@
      \leftmargin3pc \rightmargin\leftmargin
      \listparindent\normalparindent \itemindent\z@
      \parsep\z@ \@plus\p@
      \let\fullwidthdisplay\relax
    }%
    \item[\hskip\labelsep\scshape\abstractname.]%
}{%
  \endlist\egroup
  \ifx\@setabstract\relax \@setabstracta \fi
}
\def\@setabstract{\@setabstracta \global\let\@setabstract\relax}
\def\@setabstracta{%
  \ifvoid\abstractbox
  \else
    \skip@20\p@ \advance\skip@-\lastskip
    \advance\skip@-\baselineskip \vskip\skip@
    \box\abstractbox
    \prevdepth\z@ % because \abstractbox is a vtop
  \fi
}
\def\titlepage{%
  \cleardoublepage
  \thispagestyle{empty}\setcounter{page}{0}}
\def\endtitlepage{\newpage}
\def\labelenumi{(\theenumi)}
\def\theenumi{\@arabic\c@enumi}
\def\labelenumii{(\theenumii)}
\def\theenumii{\@alph\c@enumii}
\def\p@enumii{\theenumi}
\def\labelenumiii{(\theenumiii)}
\def\theenumiii{\@roman\c@enumiii}
\def\p@enumiii{\theenumi(\theenumii)}
\def\labelenumiv{(\theenumiv)}
\def\theenumiv{\@Alph\c@enumiv}
\def\p@enumiv{\p@enumiii\theenumiii}
\def\labelitemi{$\m@th\bullet$}
\def\labelitemii{\bfseries --}% \upshape already done by \itemize
\def\labelitemiii{$\m@th\ast$}
\def\labelitemiv{$\m@th\cdot$}
\newenvironment{verse}{\let\\\@centercr
  \list{}{\itemsep\z@ \itemindent -1.5em\listparindent\itemindent
  \rightmargin\leftmargin \advance\leftmargin 1.5em}\item[]%
}{%
  \endlist
}
\let\endverse=\endlist % for efficiency
\newenvironment{quotation}{\list{}{%
    \leftmargin3pc \listparindent\normalparindent
    \itemindent\z@
    \rightmargin\leftmargin \parsep\z@ \@plus\p@}%
  \item[]%
}{%
  \endlist
}
\let\endquotation=\endlist % for efficiency
\newenvironment{quote}{%
  \list{}{\rightmargin\leftmargin}\item[]%
}{%
  \endlist
}
\let\endquote=\endlist % for efficiency
\def\trivlist{\parsep\parskip\@nmbrlistfalse
  \@trivlist \labelwidth\z@ \leftmargin\z@
  \itemindent\z@
  \let\@itemlabel\@empty
  \def\makelabel##1{\upshape##1}}
\renewenvironment{enumerate}{%
  \ifnum \@enumdepth >3 \@toodeep\else
      \advance\@enumdepth \@ne
      \edef\@enumctr{enum\romannumeral\the\@enumdepth}\list
      {\csname label\@enumctr\endcsname}{\usecounter
        {\@enumctr}\def\makelabel##1{\hss\llap{\upshape##1}}}\fi
}{%
  \endlist
}
\let\endenumerate=\endlist % for efficiency
\renewenvironment{itemize}{%
  \ifnum\@itemdepth>3 \@toodeep
  \else \advance\@itemdepth\@ne
    \edef\@itemitem{labelitem\romannumeral\the\@itemdepth}%
    \list{\csname\@itemitem\endcsname}%
      {\def\makelabel##1{\hss\llap{\upshape##1}}}%
  \fi
}{%
  \endlist
}
\let\enditemize=\endlist % for efficiency
\newcommand{\descriptionlabel}[1]{\hspace\labelsep \upshape\bfseries #1:}
\newenvironment{description}{\list{}{%
  \advance\leftmargini6\p@ \itemindent-12\p@
  \labelwidth\z@ \let\makelabel\descriptionlabel}%
}{
  \endlist
}
\let\enddescription=\endlist % for efficiency
\let\upn=\textup
\AtBeginDocument{%
  \labelsep=5pt\relax
  \setcounter{enumi}{13}\setcounter{enumii}{13}%
  \setcounter{enumiii}{13}\setcounter{enumiv}{13}%
  \settowidth\leftmargini{\labelenumi\hskip\labelsep}%
  \advance\leftmargini by \normalparindent
  \settowidth\leftmarginii{\labelenumii\hskip\labelsep}%
  \settowidth\leftmarginiii{\labelenumiii\hskip\labelsep}%
  \settowidth\leftmarginiv{\labelenumiv\hskip\labelsep}%
  \setcounter{enumi}{0}\setcounter{enumii}{0}%
  \setcounter{enumiii}{0}\setcounter{enumiv}{0}%
  \leftmarginv=10pt  \leftmarginvi=\leftmarginv
  \leftmargin=\leftmargini
  \labelwidth=\leftmargini \advance\labelwidth-\labelsep
  \@listi}
\newskip\listisep
\listisep\smallskipamount
\def\@listI{\leftmargin\leftmargini \parsep\z@skip
  \topsep\listisep \itemsep\z@skip
  \listparindent\normalparindent}
\let\@listi\@listI
\def\@listii{\leftmargin\leftmarginii
  \labelwidth\leftmarginii \advance\labelwidth-\labelsep
  \topsep\z@skip \parsep\z@skip \partopsep\z@skip \itemsep\z@skip}
\def\@listiii{\leftmargin\leftmarginiii
  \labelwidth\leftmarginiii \advance\labelwidth-\labelsep}
\def\@listiv{\leftmargin\leftmarginiv
  \labelwidth\leftmarginiv \advance\labelwidth-\labelsep}
\def\@listv{\leftmargin\leftmarginv
  \labelwidth\leftmarginv \advance\labelwidth-\labelsep}
\def\@listvi{\leftmargin\leftmarginvi
  \labelwidth\leftmarginvi \advance\labelwidth-\labelsep}
\@ifclasswith{\@classname}{fleqn}{%
  \let\@tempa\@iden
  \AtBeginDocument{\mathindent\leftmargini}%
}{\let\@tempa\@gobble}%
\@ifpackageloaded{amsmath}{\let\@tempa\@gobble}{%
  \@ifpackageloaded{amstex}{\let\@tempa\@gobble}{}%
}
\@tempa{%
  \def\[{\relax
    \ifmmode\@badmath
    \else
      \begin{trivlist}%
      \@beginparpenalty\predisplaypenalty
      \@endparpenalty\postdisplaypenalty
      \item[]\leavevmode
      \hbox to\linewidth\bgroup$\displaystyle
      \hskip\mathindent\bgroup
    \fi}%
  \def\]{\relax
    \ifmmode
      \egroup \m@th$\hfil \egroup
      \end{trivlist}%
    \else \@badmath
    \fi}%
  \renewenvironment{equation}{%
    \@beginparpenalty\predisplaypenalty
    \@endparpenalty\postdisplaypenalty
    \refstepcounter{equation}%
    \@topsep\abovedisplayskip \trivlist
    \item[]\leavevmode
    \hbox to\linewidth\bgroup\hskip\mathindent$\displaystyle
  }{%
    \m@th$\hfil \displaywidth\linewidth \hbox{\@eqnnum}\egroup
    \endtrivlist
  }%
  \renewenvironment{eqnarray}{%
    \stepcounter{equation}\let\@currentlabel\theequation
    \global\@eqnswtrue \global\@eqcnt\z@ \tabskip\mathindent
    \let\\=\@eqncr \abovedisplayskip\topsep
    \ifvmode \advance\abovedisplayskip\partopsep \fi
    \belowdisplayskip\abovedisplayskip
    \belowdisplayshortskip\abovedisplayskip
    \abovedisplayshortskip\abovedisplayskip
    $$\everycr{}\halign to\linewidth\bgroup
      \hskip\@centering
      $\displaystyle\tabskip\z@skip####\m@th$&%
      \@eqnsel \global\@eqcnt\@ne
      \hfil${}####{}\m@th$\hfil&%
      \global\@eqcnt\tw@
      $\displaystyle ####\m@th$\hfil\tabskip\@centering&%
      \global\@eqcnt\thr@@
      \hbox to \z@\bgroup\hss####\egroup\tabskip\z@skip\cr
    }{%
      \@@eqncr \egroup \global\advance\c@equation\m@ne$$%
      \global\@ignoretrue
    }%
  \newdimen\mathindent
  \mathindent\leftmargini
}
\def\@startsection#1#2#3#4#5#6{%
 \if@noskipsec \leavevmode \fi
 \par \@tempskipa #4\relax
 \@afterindenttrue
 \ifdim \@tempskipa <\z@ \@tempskipa -\@tempskipa \@afterindentfalse\fi
 \if@nobreak \everypar{}\else
     \addpenalty\@secpenalty\addvspace\@tempskipa\fi
 \@ifstar{\@dblarg{\@sect{#1}{\@m}{#3}{#4}{#5}{#6}}}%
         {\@dblarg{\@sect{#1}{#2}{#3}{#4}{#5}{#6}}}%
}
\def\@seccntformat#1{%
  \protect\textup{\protect\@secnumfont
    \csname the#1\endcsname
    \protect\@secnumpunct
  }%
}
\let\@secnumfont\@empty
\def\@sect#1#2#3#4#5#6[#7]#8{%
  \edef\@toclevel{\ifnum#2=\@m 0\else\number#2\fi}%
  \ifnum #2>\c@secnumdepth \let\@secnumber\@empty
  \else \@xp\let\@xp\@secnumber\csname the#1\endcsname\fi
  \@tempskipa #5\relax
  \ifnum #2>\c@secnumdepth
    \let\@svsec\@empty
  \else
    \refstepcounter{#1}%
    \edef\@secnumpunct{%
      \ifdim\@tempskipa>\z@ % not a run-in section heading
        \@ifnotempty{#8}{.\@nx\enspace}%
      \else
        \@ifempty{#8}{.}{.\@nx\enspace}%
      \fi
    }%
    \protected@edef\@svsec{%
      \ifnum#2<\@m
        \@ifundefined{#1name}{}{%
          \ignorespaces\csname #1name\endcsname\space
        }%
      \fi
      \@seccntformat{#1}%
    }%
  \fi
  \ifdim \@tempskipa>\z@ % then this is not a run-in section heading
    \begingroup #6\relax
    \@hangfrom{\hskip #3\relax\@svsec}{\interlinepenalty\@M #8\par}%
    \endgroup
  \csname #1mark\endcsname{#7}%
    \ifnum#2>\@m \else \@tocwrite{#1}{#8}\fi
  \else
  \def\@svsechd{#6\hskip #3\@svsec
    \@ifnotempty{#8}{\ignorespaces#8\unskip
       \@addpunct.}%
    \ifnum#2>\@m \else \@tocwrite{#1}{#8}\fi
  }%
  \fi
  \global\@nobreaktrue
  \@xsect{#5}}
\let\@ssect\relax
\newcounter{part}
\newcounter{chapter}
\newcounter{section}[chapter]
\def\thesection{\arabic{section}}
\newcounter{subsection}[section]
\newcounter{subsubsection}[subsection]
\newcounter{paragraph}[subsubsection]
\newcounter{subparagraph}[paragraph]
\renewcommand\thepart          {\arabic{part}}
\renewcommand\thechapter       {\arabic{chapter}}
\renewcommand\thesection       {\arabic{section}}
\renewcommand\thesubsection    {\thesection.\arabic{subsection}}
\renewcommand\thesubsubsection {\thesubsection .\arabic{subsubsection}}
\renewcommand\theparagraph     {\thesubsubsection.\arabic{paragraph}}
\renewcommand\thesubparagraph  {\theparagraph.\arabic{subparagraph}}
\setcounter{secnumdepth}{3}
\def\partname{Part}
\def\specialsection{\@startsection{section}{1}%
  \z@{2\linespacing\@plus\linespacing}{.5\linespacing}%
  {\large\scshape\centering}}
\def\part{\cleardoublepage \thispagestyle{empty}%
  \null\vfil  \markboth{}{}\secdef\@part\@spart}
\def\@part[#1]#2{%
  \ifnum \c@secnumdepth >-2\relax \refstepcounter{part}%
    \addcontentsline{toc}{part}{\partname\ \thepart.
        \protect\enspace\protect\noindent#1}%
  \else
    \addcontentsline{toc}{part}{#1}\fi
  \begingroup\centering
  \ifnum \c@secnumdepth >-2\relax
       {\fontsize{\@xviipt}{22}\bfseries
         \partname\ \thepart} \vskip 20\p@ \fi
  \fontsize{\@xxpt}{25}\bfseries
      #1\vfil\vfil\endgroup \newpage\thispagestyle{empty}}

\def\@spart#1{\addcontentsline{toc}{part}{\protect\noindent#1}%
  \begingroup\centering
  \fontsize{\@xxpt}{25}\bfseries
     #1\vfil\vfil\endgroup \newpage\thispagestyle{empty}}
\def\partrunhead#1#2#3{%
  \@ifnotempty{#2}{\uppercase{\ignorespaces#1 #2\unskip}\@ifnotempty{#3}{. }}%
  \def\@tempa{#3}%
  \ifx\@empty\@tempa\else
    \begingroup \def\\{ \ignorespaces}% defend against questionable usage
    \uppercasenonmath\@tempa\@tempa
    \endgroup
  \fi
}
\let\chapterrunhead\partrunhead
\let\sectionrunhead\partrunhead
\def\section{\@startsection{section}{1}%
  \z@{.7\linespacing\@plus\linespacing}{.5\linespacing}%
  {\normalfont\bfseries\centering}}
\def\subsection{\@startsection{subsection}{2}%
  \normalparindent{.5\linespacing\@plus.7\linespacing}{-.5em}%
  {\normalfont\bfseries}}
\def\subsubsection{\@startsection{subsubsection}{3}%
  \normalparindent\z@{-.5em}%
  {\normalfont\itshape}}
\def\paragraph{\@startsection{paragraph}{4}%
  \normalparindent\z@{-\fontdimen2\font}%
  \normalfont}
\def\subparagraph{\@startsection{subparagraph}{5}%
  \z@\z@{-\fontdimen2\font}%
  \normalfont}
\def\@Roman#1{\@xp\@slowromancap
  \romannumeral#1@}%
\def\@slowromancap#1{\ifx @#1% then terminate
  \else
    \if i#1I\else\if v#1V\else\if x#1X\else\if l#1L\else\if
    c#1C\else\if m#1M\else#1\fi\fi\fi\fi\fi\fi
    \@xp\@slowromancap
  \fi
}
\def\frontmatter{\cleardoublepage\pagenumbering{roman}}
\def\mainmatter{\cleardoublepage\pagenumbering{arabic}}
\def\backmatter{%
  \if@openright\cleardoublepage\else\clearpage\fi
  \let\chaptername\relax}
\def\aufm#1{\par\vspace*{12pt}{\flushright #1\par}}
\newenvironment{xcb}{%
  \setcounter{enumi}{0}%
  \settowidth{\leftmargini}{\labelenumi\hskip\labelsep}%
  \setcounter{enumii}{4}% letter d
  \settowidth{\leftmarginii}{\labelenumii\hskip\labelsep}%
  \@startsection{section}% counter name; ignored because of the
                                % * below
  {1}% sectioning level
  {\z@}% indent to the left of the section title
  {18\p@\@plus2\p@}% vertical space above
  {1sp}% Space below of 13pt base-to-base, so none needs to be added
      % here; but \z@ would cause the following text to be run-in, so we
      % use 1sp instead.
  {\bfseries}% The font of the subsection title
  *% always unnumbered
}{%
  \par
}
\def\chapter{%
  \if@openright\cleardoublepage\else\clearpage\fi
  \thispagestyle{plain}\global\@topnum\z@
  \@afterindenttrue \secdef\@chapter\@schapter}
\def\@chapter[#1]#2{\refstepcounter{chapter}%
  \ifnum\c@secnumdepth<\z@ \let\@secnumber\@empty
  \else \let\@secnumber\thechapter \fi
  \typeout{\chaptername\space\@secnumber}%
  \def\@toclevel{0}%
  \ifx\chaptername\appendixname \@tocwriteb\tocappendix{chapter}{#2}%
  \else \@tocwriteb\tocchapter{chapter}{#2}\fi
  \chaptermark{#1}%
  \addtocontents{lof}{\protect\addvspace{10\p@}}%
  \addtocontents{lot}{\protect\addvspace{10\p@}}%
  \@makechapterhead{#2}\@afterheading}
\def\@schapter#1{\typeout{#1}%
  \let\@secnumber\@empty
  \def\@toclevel{0}%
  \ifx\chaptername\appendixname \@tocwriteb\tocappendix{chapter}{#1}%
  \else \@tocwriteb\tocchapter{chapter}{#1}\fi
  \chaptermark{#1}%
  \addtocontents{lof}{\protect\addvspace{10\p@}}%
  \addtocontents{lot}{\protect\addvspace{10\p@}}%
  \@makeschapterhead{#1}\@afterheading}
\newcommand\chaptername{Expos\'e}
\newcommand\appendixname{Appendix}
\def\@makechapterhead#1{\global\topskip 7.5pc\relax
  \begingroup
  \fontsize{\@xivpt}{18}\bfseries\centering
    \ifnum\c@secnumdepth>\m@ne
      \leavevmode \hskip-\leftskip
      \rlap{\vbox to\z@{\vss
          \centerline{\normalsize\mdseries
              \uppercase\@xp{\chaptername}\enspace\thechapter}
          \vskip 3pc}}\hskip\leftskip\fi
     #1\par \endgroup
  \skip@34\p@ \advance\skip@-\normalbaselineskip
  \vskip\skip@ }
\def\@makeschapterhead#1{\global\topskip 7.5pc\relax
  \begingroup
  \fontsize{\@xivpt}{18}\bfseries\centering
  #1\par \endgroup
  \skip@34\p@ \advance\skip@-\normalbaselineskip
  \vskip\skip@ }
\def\appendix{\par
  \c@chapter\z@ \c@section\z@
  \let\chaptername\appendixname
  \def\thechapter{\@Alph\c@chapter}}
\newcommand{\@pnumwidth}{1.6em}
\newcommand{\@tocrmarg}{2.6em}
\setcounter{tocdepth}{1}
\newswitch{toc}
\newswitch{lof}
\newswitch{lot}
\def\@starttoc#1#2{%
  \begingroup
  \setTrue{#1}%
  \let\secdef\@gobbletwo \chapter
  \let\@secnumber\@empty % for \@tocwrite and \chaptermark
  \ifx\contentsname#2%
  \else \@tocwrite{chapter}{#2}\fi
  \typeout{#2}\@xp\chaptermark\@xp{#2}%
  \@makeschapterhead{#2}\@afterheading
  \parskip\z@skip
  \makeatletter
  \@input{\jobname.#1}%
  \if@filesw
    \@xp\newwrite\csname tf@#1\endcsname
    \immediate\@xp\openout\csname tf@#1\endcsname \jobname.#1\relax
  \fi
  \global\@nobreakfalse \endgroup
  \newpage
}
\def\contentsname{Table des mati\`eres}
\def\listfigurename{List of Figures}
\def\listtablename{List of Tables}
\def\tableofcontents{%
  \@starttoc{toc}\contentsname
}
\def\listoffigures{\@starttoc{lof}\listfigurename}
\def\listoftables{\@starttoc{lot}\listtablename}
\AtBeginDocument{%
  \@for\@tempa:=-1,0,1,2,3\do{%
    \@ifundefined{r@tocindent\@tempa}{%
      \@xp\gdef\csname r@tocindent\@tempa\endcsname{0pt}}{}%
  }%
}
\def\@writetocindents{%
  \begingroup
  \@for\@tempa:=-1,0,1,2,3\do{%
    \immediate\write\@auxout{%
      \string\newlabel{tocindent\@tempa}{%
        \csname r@tocindent\@tempa\endcsname}}%
  }%
  \endgroup}
\AtEndDocument{\@writetocindents}

\let\indentlabel\@empty
\def\@tochangmeasure#1{\sbox\z@{#1}%
  \ifdim\wd\z@>\csname r@tocindent\@toclevel\endcsname\relax
    \@xp\xdef\csname r@tocindent\@toclevel\endcsname{\the\wd\z@}%
  \fi
}
\def\@toclevel{0}
\def\@tocline#1#2#3#4#5#6#7{\relax
  \ifnum #1>\c@tocdepth % then omit
  \else
    \par \addpenalty\@secpenalty\addvspace{#2}%
    \begingroup \hyphenpenalty\@M
    \@ifempty{#4}{%
      \@tempdima\csname r@tocindent\number#1\endcsname\relax
    }{%
      \@tempdima#4\relax
    }%
    \parindent\z@ \leftskip#3\relax \advance\leftskip\@tempdima\relax
    \rightskip\@pnumwidth plus4em \parfillskip-\@pnumwidth
    #5\leavevmode\hskip-\@tempdima #6\nobreak\relax
    \hfil\hbox to\@pnumwidth{\@tocpagenum{#7}}\par
    \nobreak
    \endgroup
  \fi}
\def\@tocpagenum#1{\hss{\mdseries #1}}
\def\@tocwrite#1{\@xp\@tocwriteb\csname toc#1\endcsname{#1}}
\def\@tocwriteb#1#2#3{%
  \begingroup
    \def\@tocline##1##2##3##4##5##6{%
      \ifnum##1>\c@tocdepth
      \else \sbox\z@{##5\let\indentlabel\@tochangmeasure##6}\fi}%
    \csname l@#2\endcsname{#1{\csname#2name\endcsname}{\@secnumber}{}}%
  \endgroup
  \addcontentsline{toc}{#2}%
    {\protect#1{\csname#2name\endcsname}{\@secnumber}{#3}}}
\def\l@section{\@tocline{1}{0pt}{1pc}{}{}}
\newcommand{\tocsection}[3]{%
  \indentlabel{\@ifnotempty{#2}{\ignorespaces#1 #2.\quad}}#3}
\def\l@subsection{\@tocline{2}{0pt}{1pc}{5pc}{}}
\let\tocsubsection\tocsection
\def\l@subsubsection{\@tocline{3}{0pt}{1pc}{7pc}{}}
\let\tocsubsubsection\tocsection
\let\l@paragraph\l@subsubsection
\let\tocparagraph\tocsection
\let\l@subparagraph\l@subsubsection
\let\tocsubparagraph\tocsection
\def\l@part{\@tocline{-1}{12pt plus2pt}{0pt}{}{\bfseries}}
\let\tocpart\tocsection
\def\l@chapter{\@tocline{0}{8pt plus1pt}{0pt}{}{}}
\let\tocchapter\tocsection
\newcommand{\tocappendix}[3]{%
  \indentlabel{#1\@ifnotempty{#2}{ #2}.\quad}#3}
\def\l@figure{\@tocline{0}{3pt plus2pt}{0pt}{1.5pc}{}}
\let\l@table=\l@figure
\def\refname{References}
\def\bibname{Bibliography}
\def\@defaultbiblabelstyle#1{#1.}
\def\bibliographystyle#1{%
    \if@filesw\immediate\write\@auxout{\string\bibstyle{#1}}\fi
    \def\@tempa{#1}%
    \def\@tempb{amsplain}%
    \def\@tempc{}%
    \ifx\@tempa\@tempb
        \def\@biblabel##1{\@defaultbiblabelstyle{##1}}%
        \def\bibsetup{}%
    \else
        \def\bibsetup{\labelsep6\p@}%
        \ifx\@tempa\@tempc
            \def\@biblabel##1{}%
            \def\bibsetup{\labelwidth\z@ \leftmargin24\p@
                \itemindent-\leftmargin
                \labelsep\z@ }%
        \fi
    \fi}
\newcommand{\bibliofont}{\footnotesize}
\newcommand{\@bibtitlestyle}{%
  \@xp\chapter\@xp*\@xp{\bibname}%
}
\newenvironment{thebibliography}[1]{%
  \@bibtitlestyle
  \normalfont\bibliofont\labelsep .5em\relax
  \renewcommand\theenumiv{\arabic{enumiv}}\let\p@enumiv\@empty
  \list{\@biblabel{\theenumiv}}{\settowidth\labelwidth{\@biblabel{#1}}%
    \leftmargin\labelwidth \advance\leftmargin\labelsep
    \usecounter{enumiv}}%
  \sloppy \clubpenalty\@M \widowpenalty\clubpenalty
  \sfcode`\.=\@m
}{%
  \def\@noitemerr{\@latex@warning{Empty `thebibliography' environment}}%
  \endlist
}
\def\bysame{\leavevmode\hbox to3em{\hrulefill}\thinspace}
\def\newblock{}
\newcommand\MR[1]{\relax\ifhmode\unskip\spacefactor3000 \space\fi
  MR~\MRhref{#1}{#1}}
\let\MRhref\@gobble
\newcommand\URL{\begingroup
  \def\@sverb##1{%
    \def\@tempa####1##1{\@URL{####1}\egroup\endgroup}%
    \@tempa}%
  \verb}
\let\URLhref\@gobble
\def\@URL#1{\URLhref{#1}#1}
\newif\if@restonecol
\def\indexchap#1{\global\topskip 7.5pc\relax
  \twocolumn[{\fontsize{\@xivpt}{18}\bfseries\centering
    \vskip\topskip\hbox{}\vskip-\baselineskip% adjust top space
    #1\par
   \global\topskip 34\p@\relax
   \ifx\@empty\indexintro
   \else
     \begingroup \normalsize
       \skip@\topskip \advance\skip@ -\baselineskip
       \vskip\skip@
       \parbox[t]{24pc}{\normalfont\indexintro\par}%
     \endgroup
     \global\topskip 24\p@\relax
   \fi
  }]%
}
\newcommand{\indexintro}{}
\newcommand{\@indextitlestyle}{%
  \let\@makeschapterhead\indexchap
  \@xp\chapter\@xp*\@xp{\indexname}%
}
\def\theindex{\@restonecoltrue\if@twocolumn\@restonecolfalse\fi
  \columnseprule\z@ \columnsep 35\p@
  \@indextitlestyle
  \thispagestyle{plain}%
  \let\item\@idxitem
  \parindent\z@  \parskip\z@\@plus.3\p@\relax
  \raggedright
  \hyphenpenalty\@M
  \footnotesize}
\def\indexname{Index}
\def\@idxitem{\par\hangindent 2em}
\def\subitem{\par\hangindent 2em\hspace*{1em}}
\def\subsubitem{\par\hangindent 3em\hspace*{2em}}
\def\endtheindex{\if@restonecol\onecolumn\else\clearpage\fi}
\def\indexspace{\par\bigskip}
\@addtoreset{footnote}{chapter}
\def\footnoterule{\kern-.4\p@
        \hrule\@width 5pc\kern11\p@\kern-\footnotesep}
\def\@makefnmark{%
  \leavevmode
  \raise.9ex\hbox{\fontsize\sf@size\z@\normalfont\@thefnmark}%
}
\def\@makefntext{\indent\@makefnmark}
\long\def\@footnotetext#1{%
  \insert\footins{%
    \normalfont\footnotesize
    \interlinepenalty\interfootnotelinepenalty
    \splittopskip\footnotesep \splitmaxdepth \dp\strutbox
    \floatingpenalty\@MM \hsize\columnwidth
    \@parboxrestore \parindent\normalparindent \sloppy
    \protected@edef\@currentlabel{%
      \csname p@footnote\endcsname\@thefnmark}%
    \@makefntext{%
      \rule\z@\footnotesep\ignorespaces#1\unskip\strut\par}}}
\hfuzz=1pt \vfuzz=\hfuzz
\def\sloppy{\tolerance9999 \emergencystretch 3em\relax}
\setcounter{topnumber}{4}
\setcounter{bottomnumber}{4}
\setcounter{totalnumber}{4}
\setcounter{dbltopnumber}{4}
\renewcommand{\topfraction}{.97}
\renewcommand{\bottomfraction}{.97}
\renewcommand{\textfraction}{.03}
\renewcommand{\floatpagefraction}{.9}
\renewcommand{\dbltopfraction}{.97}
\renewcommand{\dblfloatpagefraction}{.9}
\setlength{\floatsep}{12pt plus 6pt minus 4pt}
\setlength{\textfloatsep}{15pt plus 8pt minus 5pt}
\setlength{\intextsep}{12pt plus 6pt minus 4pt}
\setlength{\dblfloatsep}{12pt plus 6pt minus 4pt}
\setlength{\dbltextfloatsep}{15pt plus 8pt minus 5pt}
\setlength{\@fptop}{0pt}% removed "plus 1fil"
\setlength{\@fpsep}{8pt}% removed "plus 2fil"
\setlength{\@fpbot}{0pt plus 1fil}
\setlength{\@dblfptop}{0pt}% removed "plus 1fil"
\setlength{\@dblfpsep}{8pt}% removed "plus 2fil"
\setlength{\@dblfpbot}{0pt plus 1fil}
\newcommand{\fps@figure}{tbp}
\newcommand{\fps@table}{tbp}
\newcounter{figure}[chapter]
\def\@captionheadfont{\scshape}
\def\@captionfont{\normalfont}
\def\ftype@figure{1}
\def\ext@figure{lof}
\def\fnum@figure{\figurename\ \thefigure}
\def\figurename{Figure}
\newenvironment{figure}{%
  \@float{figure}%
}{%
  \end@float
}
\newenvironment{figure*}{%
  \@dblfloat{figure}%
}{%
  \end@dblfloat
}
\newcounter{table}[chapter]
\def\ftype@table{2}
\def\ext@table{lot}
\def\fnum@table{\tablename\ \thetable}
\def\tablename{Table}
\newenvironment{table}{%
  \@float{table}%
}{%
  \end@float
}
\newenvironment{table*}{%
  \@dblfloat{table}%
}{%
  \end@dblfloat
}
\def\@floatboxreset{\global\@minipagefalse \centering}
\long\def\@makecaption#1#2{%
  \setbox\@tempboxa\vbox{\color@setgroup
    \advance\hsize-2\captionindent\noindent
    \@captionfont\@captionheadfont#1\@xp\@ifnotempty\@xp
        {\@cdr#2\@nil}{.\@captionfont\upshape\enspace#2}%
    \unskip\kern-2\captionindent\par
    \global\setbox\@ne\lastbox\color@endgroup}%
  \ifhbox\@ne % the normal case
    \setbox\@ne\hbox{\unhbox\@ne\unskip\unskip\unpenalty\unkern}%
  \fi
  \ifdim\wd\@tempboxa=\z@ % this means caption will fit on one line
    \setbox\@ne\hbox to\columnwidth{\hss\kern-2\captionindent\box\@ne\hss}%
  \else % tempboxa contained more than one line
    \setbox\@ne\vbox{\unvbox\@tempboxa\parskip\z@skip
        \noindent\unhbox\@ne\advance\hsize-2\captionindent\par}%
  \fi
  \ifnum\@tempcnta<64 % if the float IS a figure...
    \addvspace\abovecaptionskip
    \hbox to\hsize{\kern\captionindent\box\@ne\hss}%
  \else % if the float IS NOT a figure...
    \hbox to\hsize{\kern\captionindent\box\@ne\hss}%
    \nobreak
    \vskip\belowcaptionskip
  \fi
\relax
}
\newskip\abovecaptionskip \abovecaptionskip=12pt \relax
\newskip\belowcaptionskip \belowcaptionskip=12pt \relax
\newdimen\captionindent \captionindent=3pc
\def\nonbreakingspace{\unskip\nobreak\ \ignorespaces}
\def~{\protect\nonbreakingspace}
\def\@biblabel#1{\@ifnotempty{#1}{[#1]}}
\def\@citestyle{\m@th\upshape\mdseries}
\def\citeform#1{{\bfseries#1}}
\def\@cite#1#2{{%
  \@citestyle[\citeform{#1}\if@tempswa, #2\fi]}}
\@ifundefined{cite }{%
  \expandafter\let\csname cite \endcsname\cite
  \edef\cite{\@nx\protect\@xp\@nx\csname cite \endcsname}%
}{}
\def\fullwidthdisplay{\displayindent\z@ \displaywidth\columnwidth}
\edef\@tempa{\noexpand\fullwidthdisplay\the\everydisplay}
\everydisplay\expandafter{\@tempa}
\newcommand*\seeonlyname{see}
\newcommand*\seename{see also}
\newcommand*\alsoname{see also}
\newcommand*\seeonly[2]{\emph{\seeonlyname} #1}
\newcommand*\see[2]{\emph{\seename} #1}
\newcommand*\seealso[2]{\emph{\alsoname} #1}
\newcommand\printindex{\@input{\jobname.ind}}
\DeclareRobustCommand\textprime{\leavevmode
  \raise.8ex\hbox{\check@mathfonts\the\scriptfont2 \char48 }}
\def\maketitle{\par
  \@topnum\z@ % this prevents figures from falling at the top of page 1
  \begingroup
  \@maketitle
  \endgroup
  \c@footnote\z@
  \def\do##1{\let##1\relax}%
  \do\maketitle \do\@maketitle \do\title \do\@xtitle \do\@title
  \do\author \do\@xauthor \do\address \do\@xaddress
  \do\email \do\@xemail \do\curraddr \do\@xcurraddr
  \do\dedicatory \do\@dedicatory \do\thanks \do\thankses
  \do\keywords \do\@keywords \do\subjclass \do\@subjclass
}
\def\@maketitle{%
  \cleardoublepage \thispagestyle{empty}%
  \begingroup \topskip\z@skip
  \null\vfil
  \begingroup
  \LARGE\bfseries \centering
  \openup\medskipamount
  \@title\par\vspace{24pt}%
  \def\and{\par\medskip}\centering
  \mdseries\authors\par\bigskip
  \endgroup
  \vfil
  \ifx\@empty\addresses \else \@setaddresses \fi
  \vfil
  \ifx\@empty\@dedicatory
  \else \begingroup
    \centering{\footnotesize\itshape\@dedicatory\@@par}%
    \endgroup
  \fi
  \vfill
  \newpage\thispagestyle{empty}
  \begin{center}
    \ifx\@empty\@subjclass\else\@setsubjclass\fi
    \ifx\@empty\@keywords\else\@setkeywords\fi
    \ifx\@empty\@translators\else\vfil\@settranslators\fi
    \ifx\@empty\thankses\else\vfil\@setthanks\fi
  \end{center}
  \vfil
  \@setabstract
  \endgroup}
\def\@settranslators{\par
  \begingroup
  \translname: \andify\@translators \uppercasenonmath\@translators
  \@translators \@@par
  \endgroup}
\def\@setdate{\par\smallskip\@date\par\smallskip}
\def\@setsubjclass{\par\smallskip
  {\itshape\subjclassname.}\enspace\@subjclass\par\smallskip}
\def\@setkeywords{\par\smallskip
  {\itshape \keywordsname.}\enspace \@keywords\par\smallskip}

\newcommand{\theoremstyle}[1]{%
  \@ifundefined{th@#1}{%
    \PackageWarning{amsthm}{Unknown theoremstyle `#1'}%
    \thm@style{plain}%
  }{%
    \thm@style{#1}%
  }%
}
\newtoks\thm@style
\thm@style{plain}
\newtoks\thm@bodyfont  \thm@bodyfont{\itshape}
\newtoks\thm@headfont  \thm@headfont{\bfseries}
\newtoks\thm@notefont  \thm@notefont{}
\newtoks\thm@headpunct \thm@headpunct{.}
\newskip\thm@preskip \newskip\thm@postskip
\def\thm@space@setup{%
  \thm@preskip=.5\baselineskip\@plus.2\baselineskip
                             \@minus.2\baselineskip
  \thm@postskip=\thm@preskip
}
\renewcommand{\newtheorem}{\@ifstar{\@xnthm *}{\@xnthm \relax}}
\def\@xnthm#1#2{%
  \let\@tempa\relax
  \@xp\@ifdefinable\csname #2\endcsname{%
    \global\@xp\let\csname end#2\endcsname\@endtheorem
    \ifx *#1% unnumbered, need to get one more mandatory arg
      \edef\@tempa##1{%
        \gdef\@xp\@nx\csname#2\endcsname{%
          \@nx\@thm{\@xp\@nx\csname th@\the\thm@style\endcsname}%
            {}{##1}}}%
    \else % numbered theorem, need to check for optional arg
      \def\@tempa{\@oparg{\@ynthm{#2}}[]}%
    \fi
  }%
  \@tempa
}
\def\@ynthm#1[#2]#3{%
  \ifx\relax#2\relax
    \def\@tempa{\@oparg{\@xthm{#1}{#3}}[]}%
  \else
    \@ifundefined{c@#2}{%
      \def\@tempa{\@nocounterr{#2}}%
    }{%
      \@xp\xdef\csname the#1\endcsname{\@xp\@nx\csname the#2\endcsname}%
      \toks@{#3}%
      \@xp\xdef\csname#1\endcsname{%
        \@nx\@thm{%
          \let\@nx\thm@swap
            \if S\thm@swap\@nx\@firstoftwo\else\@nx\@gobble\fi
          \@xp\@nx\csname th@\the\thm@style\endcsname}%
            {#2}{\the\toks@}}%
      \let\@tempa\relax
    }%
  \fi
  \@tempa
}
\def\@xthm#1#2[#3]{%
  \ifx\relax#3\relax
    \newcounter{#1}%
  \else
    \newcounter{#1}[#3]%
    \@xp\xdef\csname the#1\endcsname{\@xp\@nx\csname the#3\endcsname
      \@thmcountersep\@thmcounter{#1}}%
  \fi
  \toks@{#2}%
  \@xp\xdef\csname#1\endcsname{%
    \@nx\@thm{%
      \let\@nx\thm@swap
        \if S\thm@swap\@nx\@firstoftwo\else\@nx\@gobble\fi
      \@xp\@nx\csname th@\the\thm@style\endcsname}%
      {#1}{\the\toks@}}%
}
\def\@thm#1#2#3{%
  \ifhmode\unskip\unskip\par\fi
  \normalfont
  \trivlist
  \let\thmheadnl\relax
  \let\thm@swap\@gobble
  \let\thm@indent\indent % indent
  \thm@headfont{\scshape}% heading font small caps
  \thm@notefont{\fontseries\mddefault\upshape}%
  \thm@headpunct{.}% add period after heading
  \thm@headsep 5\p@ plus\p@ minus\p@\relax
  \thm@space@setup
  #1% style overrides
  \@topsep \thm@preskip               % used by thm head
  \@topsepadd \thm@postskip           % used by \@endparenv
  \def\@tempa{#2}\ifx\@empty\@tempa
    \def\@tempa{\@oparg{\@begintheorem{#3}{}}[]}%
  \else
    \refstepcounter{#2}%
    \def\@tempa{\@oparg{\@begintheorem{#3}{\csname the#2\endcsname}}[]}%
  \fi
  \@tempa
}
\def\@restorelabelsep{\relax}
\let\@ythm\relax
\let\thmname\@iden \let\thmnote\@iden \let\thmnumber\@iden
\providecommand\@upn{\textup}
\def\thmhead@plain#1#2#3{%
  \thmname{#1}\thmnumber{\@ifnotempty{#1}{ }\@upn{#2}}%
  \thmnote{ {\the\thm@notefont(#3)}}}
\let\thmhead\thmhead@plain
\def\swappedhead#1#2#3{%
  \thmnumber{\@upn{\@secnumfont#2\@ifnotempty{#1}{.~}}}%
  \thmname{#1}%
  \thmnote{ {\the\thm@notefont(#3)}}}
\let\swappedhead@plain=\swappedhead
\let\thmheadnl\relax
\let\thm@indent\noindent
\let\thm@swap\@gobble
\def\@begintheorem#1#2[#3]{%
  \deferred@thm@head{\the\thm@headfont \thm@indent
    \@ifempty{#1}{\let\thmname\@gobble}{\let\thmname\@iden}%
    \@ifempty{#2}{\let\thmnumber\@gobble}{\let\thmnumber\@iden}%
    \@ifempty{#3}{\let\thmnote\@gobble}{\let\thmnote\@iden}%
    \thm@swap\swappedhead\thmhead{#1}{#2}{#3}%
    \the\thm@headpunct
    \thmheadnl % possibly a newline.
    \hskip\thm@headsep
  }%
  \ignorespaces}
\newskip\thm@headsep
\thm@headsep=5pt plus1pt minus1pt\relax
\let\adjust@parskip@nobreak=\@nbitem
\newtoks\dth@everypar
\dth@everypar={%
  \@minipagefalse \global\@newlistfalse
  \@noparitemfalse
  \if@inlabel
    \global\@inlabelfalse
    \begingroup \setbox\z@\lastbox
     \ifvoid\z@ \kern-\itemindent \fi
    \endgroup
    \unhbox\@labels
  \fi
  \if@nobreak \@nobreakfalse \clubpenalty\@M
  \else \clubpenalty\@clubpenalty \everypar{}%
  \fi
}%
\def\deferred@thm@head#1{%
  \if@inlabel \indent \par \fi % eject a section head if one is pending
  \if@nobreak
    \adjust@parskip@nobreak
  \else
    \addpenalty\@beginparpenalty
    \addvspace\@topsep
    \addvspace{-\parskip}%
  \fi
  \global\@inlabeltrue
  \everypar\dth@everypar
  \sbox\@labels{\normalfont#1}%
  \ignorespaces
}
\def\nonslanted{\relax
  \@xp\let\@xp\@tempa\csname\f@shape shape\endcsname
  \ifx\@tempa\itshape\upshape
  \else\ifx\@tempa\slshape\upshape\fi\fi}
\def\swapnumbers{\edef\thm@swap{\if S\thm@swap N\else S\fi}}
\def\thm@swap{N}%
\let\@opargbegintheorem\relax
\def\th@plain{%
%%  \let\thm@indent\noindent % no indent
%%  \thm@headfont{\bfseries}% heading font is bold
%%  \thm@notefont{}% same as heading font
%%  \thm@headpunct{.}% add period after heading
%%  \let\thm@swap\@gobble
%%  \thm@preskip\topsep
%%  \thm@postskip\theorempreskipamount
  \itshape % body font
}
\def\th@definition{%
  \normalfont % body font
}
\def\th@remark{%
  \normalfont % body font
}
\def\@endtheorem{\endtrivlist\@endpefalse }
\newcommand{\newtheoremstyle}[9]{%
  \@ifempty{#5}{\dimen@\z@skip}{\dimen@#5\relax}%
  \ifdim\dimen@=\z@
    \toks@{#4\let\thm@indent\noindent}%
  \else
    \toks@{#4\def\thm@indent{\noindent\hbox to#5{}}}%
  \fi
  \def\@tempa{#8}\ifx\space\@tempa
    \toks@\@xp{\the\toks@ \thm@headsep\fontdimen\tw@\font\relax}%
  \else
    \def\@tempb{\newline}%
    \ifx\@tempb\@tempa
      \toks@\@xp{\the\toks@ \thm@headsep\z@skip
        \def\thmheadnl{\newline}}%
    \else
      \toks@\@xp{\the\toks@ \thm@headsep#8\relax}%
    \fi
  \fi
  \begingroup
  \thm@space@setup
  \@defaultunits\@tempskipa#2\thm@preskip\relax\@nnil
  \@defaultunits\@tempskipb#3\thm@postskip\relax\@nnil
  \xdef\@gtempa{\thm@preskip\the\@tempskipa
    \thm@postskip\the\@tempskipb\relax}%
  \endgroup
  \@temptokena\@xp{\@gtempa
    \thm@headfont{#6}\thm@headpunct{#7}%
  }%
  \@ifempty{#9}{%
    \let\thmhead\thmhead@plain
  }{%
    \@namedef{thmhead@#1}##1##2##3{#9}%
    \@temptokena\@xp{\the\@temptokena
      \@xp\let\@xp\thmhead\csname thmhead@#1\endcsname}%
  }%
  \@xp\xdef\csname th@#1\endcsname{\the\toks@ \the\@temptokena}%
}
\DeclareRobustCommand{\qed}{%
  \ifmmode \mathqed
  \else
    \leavevmode\unskip\penalty9999 \hbox{}\nobreak\hfill
    \quad\hbox{\qedsymbol}%
  \fi
}
\let\QED@stack\@empty
\let\qed@elt\relax
\newcommand{\pushQED}[1]{%
  \toks@{\qed@elt{#1}}\@temptokena\expandafter{\QED@stack}%
  \xdef\QED@stack{\the\toks@\the\@temptokena}%
}
\newcommand{\popQED}{%
  \begingroup\let\qed@elt\popQED@elt \QED@stack\relax\relax\endgroup
}
\def\popQED@elt#1#2\relax{#1\gdef\QED@stack{#2}}
\newcommand{\qedhere}{%
  \begingroup \let\mathqed\math@qedhere
    \let\qed@elt\setQED@elt \QED@stack\relax\relax \endgroup
}
\newif\ifmeasuring@
\newif\iffirstchoice@ \firstchoice@true
\def\setQED@elt#1#2\relax{%
  \ifmeasuring@
  \else \iffirstchoice@ \gdef\QED@stack{\qed@elt{}#2}\fi
  \fi
  #1%
}
\def\qed@warning{%
  \PackageWarning{amsthm}{The \@nx\qedhere command may not work
    correctly here}%
}
\newcommand{\mathqed}{\quad\hbox{\qedsymbol}}
\def\linebox@qed{\hfil\hbox{\qedsymbol}\hfilneg}
\@ifpackageloaded{amsmath}{%
  \def\math@qedhere{%
    \@ifundefined{\@currenvir @qed}{%
      \qed@warning\quad\hbox{\qedsymbol}%
    }{%
      \@xp\aftergroup\csname\@currenvir @qed\endcsname
    }%
  }
  \def\displaymath@qed{%
    \relax
    \ifmmode
      \ifinner \aftergroup\linebox@qed
      \else
        \eqno
        \let\eqno\relax \let\leqno\relax \let\veqno\relax
        \hbox{\qedsymbol}%
      \fi
    \else
       \aftergroup\linebox@qed
    \fi
  }
  \@xp\let\csname equation*@qed\endcsname\displaymath@qed
  \def\equation@qed{%
    \iftagsleft@
      \hbox{\phantom{\quad\qedsymbol}}%
      \gdef\alt@tag{%
        \rlap{\hbox to\displaywidth{\hfil\qedsymbol}}%
        \global\let\alt@tag\@empty
      }%
    \else
      \gdef\alt@tag{%
        \global\let\alt@tag\@empty
        \vtop{\ialign{\hfil####\cr
                \tagform@\theequation\cr
                \qedsymbol\cr}}%
        \setbox\z@
      }%
    \fi
  }
  \def\qed@tag{%
    \global\tag@true \nonumber
    &\omit\setboxz@h {\strut@ \qedsymbol}\tagsleft@false
    \place@tag@gather
    \kern-\tabskip
    \ifst@rred \else \global\@eqnswtrue \fi \global\advance\row@\@ne \cr
  }
  \def\split@qed{%
    \def\endsplit{\crcr\egroup \egroup \ctagsplit@false \rendsplit@
      \aftergroup\align@qed
    }%
  }
  \def\align@qed{%
    \ifmeasuring@ \tag*{\qedsymbol}%
    \else \let\math@cr@@@\qed@tag
    \fi
  }
  \@xp\let\csname align*@qed\endcsname\align@qed
  \@xp\let\csname gather*@qed\endcsname\align@qed
%% Needs some patching up for amsmath 1.2
}{% end of amsmath branch, start plain LaTeX branch
  \def\math@qedhere{%
    \@ifundefined{\@currenvir @qed}{%
      \qed@warning \aftergroup\displaymath@qed
    }{%
      \@xp\aftergroup\csname\@currenvir @qed\endcsname
    }%
  }
  \def\displaymath@qed{%
    \relax
    \ifmmode
      \ifinner \aftergroup\aftergroup\aftergroup\linebox@qed
      \else
        \eqno \def\@badmath{$$}%
        \let\eqno\relax \let\leqno\relax \let\veqno\relax
        \hbox{\qedsymbol}%
      \fi
    \else
       \aftergroup\linebox@qed
    \fi
  }
  \@ifundefined{ver@leqno.clo}{%
    \def\equation@qed{\displaymath@qed \quad}%
  }{%
    \def\equation@qed{\displaymath@qed}%
  }
  \def\@tempa#1$#2#3\@nil{%
    \def\[{#1$#2\def\@currenvir{displaymath}#3}%
  }%
  \expandafter\@tempa\[\@nil
}
\@ifpackageloaded{amstex}{%
  \def\@tempa{TT}%
}{%
  \@ifpackageloaded{amsmath}{%
    \def\@tempb#1 v#2.#3\@nil{#2}%
    \ifnum\@xp\@xp\@xp\@tempb\csname ver@amsmath.sty\endcsname v0.0\@nil
       <\tw@
      \def\@tempa{TT}%
    \else
      \def\@tempa{TF}%
    \fi
  }{%
    \def\@tempa{TF}
  }%
}
\if\@tempa
  \renewcommand{\math@qedhere}{\quad\hbox{\qedsymbol}}%
\fi
\newcommand{\openbox}{\leavevmode
  \hbox to.77778em{%
  \hfil\vrule
  \vbox to.675em{\hrule width.6em\vfil\hrule}%
  \vrule\hfil}}
\DeclareRobustCommand{\textsquare}{%
  \begingroup \usefont{U}{msa}{m}{n}\thr@@\endgroup
}
\@ifclasswith{\@classname}{noamsfonts}{%
  \providecommand{\qedsymbol}{\openbox}%
}{}
\providecommand{\qedsymbol}{\textsquare}
\newenvironment{proof}[1][\proofname]{\par
  \pushQED{\qed}%
  \normalfont \topsep6\p@\@plus6\p@\relax
  \trivlist
  \itemindent\normalparindent
  \item[\hskip\labelsep
        \scshape
    #1\@addpunct{.}]\ignorespaces
}{%
  \popQED\endtrivlist\@endpefalse
}
\providecommand{\proofname}{Proof}
\def\bb@skip#1{%
  \skip@#1\relax \advance\skip@-\prevdepth \advance\skip@-\baselineskip
  \vskip\skip@}
\def\markleft#1{{\let\protect\noexpand
    \let\label\relax \let\index\relax \let\glossary\relax
    \expandafter\@markleft\@themark{#1}%
    \mark{\@themark}}%
  \if@nobreak\ifvmode\nobreak\fi\fi}
\def\@markleft#1#2#3{\gdef\@themark{{#3}{#2}}}
\def\@tempa{}
\edef\@dh{%
  \noexpand\mathhexbox{\hexnumber@\symAMSb}67}
\DeclareTextCommand{\dh}{OT1}{%
  \edef\@tempb{\scdefault}%
  \ifx\f@shape\@tempb
    \leavevmode
    \raisebox{-.8ex}{\makebox[\z@][l]{\hskip-.08em\accent"16\hss}}d%
  \else
    \@dh
  \fi
}
\DeclareTextCommand{\DH}{OT1}{%
  \leavevmode\raisebox{-.5ex}{\makebox[\z@][l]{\hskip-.07em\accent"16\hss}}D}
\DeclareTextCommand{\DJ}{OT1}{%
  \leavevmode\raisebox{-.5ex}{\makebox[\z@][l]{\hskip-.07em\accent"16\hss}}D}
\DeclareTextCommand{\dj}{OT1}{%
  \edef\@tempa{\f@shape}\edef\@tempb{\scdefault}%
  \ifx\@tempa\@tempb
    \leavevmode
    \raisebox{-.75ex}{\makebox[\z@][l]{\hskip-.08em\accent"16\hss}}d%
  \else
    \leavevmode\raisebox{.02ex}{\makebox[\z@][l]{\hskip.1em\accent"16\hss}}d%
  \fi}
\hyphenation{acad-e-my acad-e-mies af-ter-thought anom-aly anom-alies
an-ti-deriv-a-tive an-tin-o-my an-tin-o-mies apoth-e-o-ses
apoth-e-o-sis ap-pen-dix ar-che-typ-al as-sign-a-ble as-sist-ant-ship
as-ymp-tot-ic asyn-chro-nous at-trib-uted at-trib-ut-able bank-rupt
bank-rupt-cy bi-dif-fer-en-tial blue-print busier busiest
cat-a-stroph-ic cat-a-stroph-i-cally con-gress cross-hatched data-base
de-fin-i-tive de-riv-a-tive dis-trib-ute dri-ver dri-vers eco-nom-ics
econ-o-mist elit-ist equi-vari-ant ex-quis-ite ex-tra-or-di-nary
flow-chart for-mi-da-ble forth-right friv-o-lous ge-o-des-ic
ge-o-det-ic geo-met-ric griev-ance griev-ous griev-ous-ly
hexa-dec-i-mal ho-lo-no-my ho-mo-thetic ideals idio-syn-crasy
in-fin-ite-ly in-fin-i-tes-i-mal ir-rev-o-ca-ble key-stroke
lam-en-ta-ble light-weight mal-a-prop-ism man-u-script mar-gin-al
meta-bol-ic me-tab-o-lism meta-lan-guage me-trop-o-lis
met-ro-pol-i-tan mi-nut-est mol-e-cule mono-chrome mono-pole
mo-nop-oly mono-spline mo-not-o-nous mul-ti-fac-eted mul-ti-plic-able
non-euclid-ean non-iso-mor-phic non-smooth par-a-digm par-a-bol-ic
pa-rab-o-loid pa-ram-e-trize para-mount pen-ta-gon phe-nom-e-non
post-script pre-am-ble pro-ce-dur-al pro-hib-i-tive pro-hib-i-tive-ly
pseu-do-dif-fer-en-tial pseu-do-fi-nite pseu-do-nym qua-drat-ic
quad-ra-ture qua-si-smooth qua-si-sta-tion-ary qua-si-tri-an-gu-lar
quin-tes-sence quin-tes-sen-tial re-arrange-ment rec-tan-gle
ret-ri-bu-tion retro-fit retro-fit-ted right-eous right-eous-ness
ro-bot ro-bot-ics sched-ul-ing se-mes-ter semi-def-i-nite
semi-ho-mo-thet-ic set-up se-vere-ly side-step sov-er-eign spe-cious
spher-oid spher-oid-al star-tling star-tling-ly sta-tis-tics
sto-chas-tic straight-est strange-ness strat-a-gem strong-hold
sum-ma-ble symp-to-matic syn-chro-nous topo-graph-i-cal tra-vers-a-ble
tra-ver-sal tra-ver-sals treach-ery turn-around un-at-tached
un-err-ing-ly white-space wide-spread wing-spread wretch-ed
wretch-ed-ly Eng-lish Euler-ian Feb-ru-ary Gauss-ian
Hamil-ton-ian Her-mit-ian Jan-u-ary Japan-ese Kor-te-weg
Le-gendre Mar-kov-ian Noe-ther-ian No-vem-ber Rie-mann-ian Sep-tem-ber}
\def\calclayout{\advance\textheight -\headheight
  \advance\textheight -\headsep
  \oddsidemargin\paperwidth
  \advance\oddsidemargin -\textwidth
  \divide\oddsidemargin\tw@
  \ifdim\oddsidemargin<.5truein \oddsidemargin.5truein \fi
  \advance\oddsidemargin -1truein
  \evensidemargin\oddsidemargin
  \topmargin\paperheight \advance\topmargin -\textheight
  \advance\topmargin -\headheight \advance\topmargin -\headsep
  \divide\topmargin\tw@
  \ifdim\topmargin<.5truein \topmargin.5truein \fi
  \advance\topmargin -1truein\relax
}
\InputIfFileExists{amsbook.cfg}{}{%
  \calclayout % initialize
  \pagenumbering{arabic}%
  \pagestyle{headings}%
  \thispagestyle{plain}%
}
\if@compatibility \else\endinput\fi
\def\tiny{\Tiny}
\def\defaultfont{\normalfont}
\def\rom{\textup}
\let\@newpf\proof \let\proof\relax \let\endproof\relax
\newenvironment{pf}{\@newpf[\proofname]}{\popQED\endtrivlist}
\newenvironment{pf*}[1]{\@newpf[#1]}{\popQED\endtrivlist}
\endinput
%%
%% End of file `amsbook.cls'.

\end{filecontents*}
\begin{filecontents*}[overwrite]{SGA4-native-smfthm.sty}
%%
%% This is file `smfthm.sty',
%% generated with the docstrip utility.
%%
%% The original source files were:
%%
%% smfclass.dtx  (with options: `smfthm')
%% 
%%
%%  %%%%%%%%%%%%%%%%%%%%%%%%%%%%%%%%%%%%%%%%%%%%%%%%%%%%%%%%%
%%  %% Classes, paquetages pour les publications de la SMF %%
%%  %% (c) 1999-2001 Societe mathematique de France        %%
%%  %%%%%%%%%%%%%%%%%%%%%%%%%%%%%%%%%%%%%%%%%%%%%%%%%%%%%%%%%
%%
%% \CharacterTable
%%  {Upper-case    \A\B\C\D\E\F\G\H\I\J\K\L\M\N\O\P\Q\R\S\T\U\V\W\X\Y\Z
%%   Lower-case    \a\b\c\d\e\f\g\h\i\j\k\l\m\n\o\p\q\r\s\t\u\v\w\x\y\z
%%   Digits        \0\1\2\3\4\5\6\7\8\9
%%   Exclamation   \!     Double quote  \"     Hash (number) \#
%%   Dollar        \$     Percent       \%     Ampersand     \&
%%   Acute accent  \'     Left paren    \(     Right paren   \)
%%   Asterisk      \*     Plus          \+     Comma         \,
%%   Minus         \-     Point         \.     Solidus       \/
%%   Colon         \:     Semicolon     \;     Less than     \<
%%   Equals        \=     Greater than  \>     Question mark \?
%%   Commercial at \@     Left bracket  \[     Backslash     \\
%%   Right bracket \]     Circumflex    \^     Underscore    \_
%%   Grave accent  \`     Left brace    \{     Vertical bar  \|
%%   Right brace   \}     Tilde         \~}
%%
\ProvidesPackage{smfthm}[1999/03/05 v1.0.h
  Environnements predefinis de theoremes]
\newif\ifsmf@thmsin\smf@thmsintrue
\def\NumberTheoremsIn#1{\@ifempty{#1}%
     {\smf@thmsinfalse}%
     {\@ifundefined{c@#1}{\@nocounterr{#1}}%
       {\smf@thmsintrue\def\@NumberTheorems{#1}}}}
  \NumberTheoremsIn{section}
\def\NumberTheoremsAs#1{\@ifempty{#1}%
     {\smf@thmsintrue}%
     {\@ifundefined{c@#1}{\@nocounterr{#1}}%
       {\smf@thmsinfalse\def\@NumberTheorems{#1}}}}
  \NumberTheoremsAs{}
\def\SwapTheoremNumbers{\def\thm@swap{S}}
\def\NoSwapTheoremNumbers{\def\thm@swap{N}}
  \NoSwapTheoremNumbers
\def\smf@thmdefs{%
  \theoremstyle{plain}
  \ifsmf@thmsin
    \newcounter{smfthm}[\@NumberTheorems]
    \xdef\thesmfthm
        {\expandafter\noexpand\csname the\@NumberTheorems\endcsname
        .\noexpand\arabic{smfthm}}%
    \gdef\smf@thm{smfthm}
  \else
    \xdef\smf@thm{\@NumberTheorems}
  \fi
  \newtheorem{theo}[\smf@thm]{\theoname}%
  \newtheorem{prop}[\smf@thm]{\propname}%
  \newtheorem{conj}[\smf@thm]{\conjname}%
  \newtheorem{coro}[\smf@thm]{\coroname}%
  \newtheorem{lemm}[\smf@thm]{\lemmname}%
  \theoremstyle{definition}%
  \newtheorem{defi}[\smf@thm]{\definame}%
  \theoremstyle{remark}%
  \newtheorem{rema}[\smf@thm]{\remaname}%
  \newtheorem{exem}[\smf@thm]{\exemname}%
  \theoremstyle{plain}}
\def\smf@enoncedef{%
 \newenvironment{enonce}[2][plain]%
    {\let\smfenonce\relax \theoremstyle{##1}%
     \newtheorem{smfenonce}[\smf@thm]{##2}%
     \begin{smfenonce}}%
    {\end{smfenonce}}%
 \newenvironment{enonce*}[2][plain]%
    {\let\smfenonce\relax \theoremstyle{##1}%
     \newtheorem*{smfenonce}{##2}%
     \begin{smfenonce}}%
    {\end{smfenonce}}%
}
\AtBeginDocument{%
     \smf@thmdefs\smf@enoncedef}
\endinput
%%
%% End of file `smfthm.sty'.

\end{filecontents*}
% !TeX program = xelatex
% BEGIN native input: style/preamble.tex
% SGA 4 source layout: original AMS/SMF support files, Chinese adaptation only.
\documentclass[12pt]{SGA4-native-amsbook}
\usepackage[a4paper,left=3cm,right=3cm,marginpar=0.5cm,marginparsep=0.8cm,vmargin=2.8cm]{geometry}
% BEGIN native input: style/paquets.tex
\usepackage{polyglossia}
%\usepackage[francais]{babel}
\usepackage{fontspec}
\usepackage[math-style=ISO,bold-style=upright,sans-style=upright]{unicode-math}
%\usepackage[mathletters]{ucs}
%\usepackage[utf8]{inputenc}
%\usepackage[T1]{fontenc}
\usepackage[unicode,hypertexnames=false,colorlinks=true,linkcolor=blue,citecolor=Sepia,urlcolor=magenta]{hyperref}
\setmainlanguage{french}
\setmainfont[Ligatures=TeX]{Linux Libertine O}
% variante possible : TeXGyrePagella ?
\setmathfont{XITS Math}

% maths
\usepackage{amsmath}
%\usepackage{amssymb}
%\usepackage{mathrsfs}
%\usepackage{stmaryrd}
%\usepackage{euscript}

% couleurs
\usepackage[usenames,dvipsnames]{xcolor}

% graphiques
\usepackage{graphicx}
\usepackage[all]{xy}
%\usepackage{tikz}
%\usetikzlibrary{matrix,arrows,calc}
%\tikzset{column sep=3em,row sep=2em,text height=1.5ex, text depth=0.25ex,>=angle 90}
%\tikzset{description/.style={fill=white,inner sep=1pt}}

% hyperliens, 'etiquettes
\usepackage{xr-hyper}
\usepackage{SGA4-native-smfthm}
%\ifx\cacheretiquettes\undefined
%\usepackage[inline]{showlabels} \else \fi % pour rendre visibles les 'etiquettes
%\usepackage{../fichiers-configuration-latex/makeidx} % pour faire un index
%\usepackage[pagebackref,unicode]{hyperref}  % pour voir o`u sont cit'ees les r'ef'erences
% hyperref is loaded once above with the cumulative-reader destination policy.
%\hypersetup{breaklinks=true} %permet le retour `a la ligne dans les liens trop longs. Incompatible avec pdflatex
%\renewcommand*{\backref}[1]{{\small [cit\'e page(s)~#1.]}}

\usepackage{xspace}
\usepackage{pifont} % pour expose XVII


% END native input: style/paquets.tex
% BEGIN native input: style/numerotation.tex
\renewcommand{\thefootnote}{\roman{footnote}} % pour mettre les notes en bas de page sous la forme (ii) et non (2), qui ressemble trop (?) a un exposant.

\newcounter{pageoriginal}
\newcommand{\pageoriginale}{\refstepcounter{pageoriginal}\marginpar{\footnotesize\textbf{\thepageoriginal}}\xspace}
\let\pageoriginaled\pageoriginale

\usepackage{ifthen}
\newboolean{original}
\def\sisi{\ifthenelse{\boolean{original}}}
\setboolean{original}{false}

\newcounter{cnote}
\def\nde#1{\ifthenelse{\boolean{original}}{}{\stepcounter{cnote}\def\thefootnote{\arabic{cnote}}\protect\footnote[\value{cnote}]{N.D.E.: #1}}}
\def\ndemark{\ifthenelse{\boolean{original}}{}{\stepcounter{cnote}\def\thefootnote{\arabic{cnote}}\protect\footnotemark[\value{cnote}]}}
\def\ndetext#1{\ifthenelse{\boolean{original}}{}{\def\thefootnote{\arabic{cnote}}\protect\footnotetext[\value{cnote}]{N.D.E.: #1}}}

% END native input: style/numerotation.tex
% BEGIN native input: style/macros.tex
\setmainlanguage{french}
\setotherlanguage{english}


%% Fixed, engine-independent release identifier.  The inherited dynamic Git
%% probe used Lua io.popen, exposed build-state-dependent output, and printed
%% a visible question mark outside a Git checkout.
\def\showgitstatus{English release r2, 2026-07-23}


% macros pour LaTeX
\renewcommand{\mathsf}{}

% Police différente pour U+2113 SCRIPT SMALL L
\setmathfont[range="2113]{latinmodern-math.otf}

% Give every inherited semantic label its own stable PDF destination instead
% of inheriting whichever automatic equation/page anchor happened to be live.
% The legacy label remains as a compatibility alias used by existing edges.
\makeatletter
\def\etiquette[#1]#2{%
  \Hy@raisedlink{\hyper@anchorstart{sga4.semantic.#2}\hyper@anchorend}%
  \def\@currentHref{sga4.semantic.#2}%
  \def\@currentlabel{#1}%
  \label{#2}%
}
\makeatother
\def\eqhyperref[#1]#2{(\hyperref[#1]{#2})}
\def\pagehyperref[#1]#2{(page \hyperref[#1]{#2})}
\def\hyperrefext[#1][#2]#3{\hyperref[#2]{#1 #3}}
\def\eqhyperrefext[#1][#2]#3{\hyperref[#2]{#1 #3}}     %cf. exposes 17
\def\neweqhyperrefext[#1][#2]#3{\hyperref[#2]{#1 #3}}  %cf. exposes 17

\def\shat#1{\protect\widehat{#1}}
\def\stilde#1{\protect\widetilde{#1}}
%%%%%%%%%%%%%%%%%%


\def\etcn{etc.\xspace}
\def\etco{etc...\xspace}
\def\etc{\sisi{\etco}{\etcn}}
\def\ldotsn{\sisi{...}{\etcn}}

\def\cf{{cf.}\kern.3em}
\def\Cf{{Cf.}\kern.3em}
\def\ie{{i.e.}\kern.3em}
\def\iev{{i.e.},\xspace}
\def\loccito{loc.\kern2pt cit.\xspace}
\def\loccitn{\textit{loc.\kern2pt cit.}\xspace}
\def\loccit{\sisi{\loccito}{\loccitn}}
\def\EGA{\textup{EGA}\kern.2em}
\def\SGA{\textup{SGA}\kern.3em}

\def\quoi{}

\def\quav{\quad,\quad}
\def\quaet{\quad \text{et}\quad}
\def\fibreproduct#1#2#3{#1 ×_{#3} #2}
\let\fprod\fibreproduct
\def\ccheck#1#2{\mbox{${\check{\check{\rm #1}}}^{\lower2pt\hbox{$\scriptstyle #2$}}$}}
\def\sheaf#1{\protect\mathcal{#1}}
\def\vieuxsheaf#1{{#1}} % cf. problème avec sheaf{a\Phi_i} par exemple...
\DeclareMathOperator{\Hom}{Hom}
\DeclareMathOperator{\can}{can}
\DeclareMathOperator{\Ind}{Ind}
\DeclareMathOperator{\ind}{ind}
%\DeclareMathOperator{\Hhomtop}{\mathcal{H}\!\mathit{omtop}}
\DeclareMathOperator{\Hhomtop}{\mathbf{Homtop}}
\DeclareMathOperator{\Point}{Point}
\DeclareMathOperator{\Pro}{Pro}
\DeclareMathOperator{\Top}{Top}
\DeclareMathOperator{\Mod}{Mod}
\DeclareMathOperator{\Cart}{Cart}
%\DeclareMathOperator{\Mmorsite}{\mathcal{M}\!\mathit{orsite}}
\DeclareMathOperator{\Mmorsite}{\mathbf{Morsite}}
\DeclareMathOperator{\loc}{loc}
\DeclareMathOperator{\Loc}{Loc}
\DeclareMathOperator{\Pic}{Pic}
\DeclareMathOperator{\Et}{Et}
\DeclareMathOperator{\card}{card}
\DeclareMathOperator{\Ens}{Ens}
\DeclareMathOperator{\h}{\underline{H}}
\DeclareMathOperator{\Fonct}{Fonct}
\DeclareMathOperator{\Ker}{Ker}
\DeclareMathOperator{\but}{but}
\DeclareMathOperator{\Ob}{Ob}
\DeclareMathOperator{\ob}{ob}
\DeclareMathOperator{\Lib}{Lib}
\DeclareMathOperator{\esj}{esj}
\DeclareMathOperator{\id}{id}
\DeclareMathOperator{\Id}{Id}
\DeclareMathOperator{\ad}{ad}
\DeclareMathOperator{\Iim}{Im}
\DeclareMathOperator{\pr}{pr}
\DeclareMathOperator{\Cat}{Cat}
\DeclareMathOperator{\car}{car}
\DeclareMathOperator{\SysInd}{SysInd}
\DeclareMathOperator{\Sup}{Sup}
\DeclareMathOperator{\Coker}{Coker}
\DeclareMathOperator{\Filt}{Filt}
\DeclareMathOperator{\acc}{acc}
\DeclareMathOperator{\sur}{sur}
\DeclareMathOperator{\cat}{cat}
\DeclareMathOperator{\Ord}{Ord}
\DeclareMathOperator{\Gr}{Gr}
\DeclareMathOperator{\Cov}{Cov}
\DeclareMathOperator{\Caf}{Caf}
\DeclareMathOperator{\coker}{coker}
\DeclareMathOperator{\coim}{coim}
\DeclareMathOperator{\im}{im}
\DeclareMathOperator{\Ouv}{Ouv}
\DeclareMathOperator{\diag}{diag}
\DeclareMathOperator{\Esp}{Esp}
\DeclareMathOperator{\TOP}{TOP}
\DeclareMathOperator{\Unif}{Unif}
\DeclareMathOperator{\Iind}{\mathbf{Ind}}
\DeclareMathOperator{\Hhomcart}{\mathbf{Homcart}}
\DeclareMathOperator{\Hhomtopess}{\mathbf{Homtopess}}
\DeclareMathOperator{\Ppoints}{\mathbf{Points}}
\DeclareMathOperator{\Ppoint}{\mathbf{Point}}
\DeclareMathOperator{\Ttors}{\mathbf{Tors}}
\DeclareMathOperator{\Ssex}{\mathbf{Sex}}
\DeclareMathOperator{\Ffib}{\mathbf{Fib}}
\DeclareMathOperator{\Vvois}{\mathbf{Vois}}
\DeclareMathOperator{\Ppointess}{\mathbf{Pointess}}
\DeclareMathOperator{\scrHom}{\mathbf{Hom}}
\DeclareMathOperator{\sheafPic}{\mathbf{Pic}}
\DeclareMathOperator{\sheafAut}{\mathbf{Aut}}
\DeclareMathOperator{\Hhomtoplocan}{\mathbf{Homtoplocan}}
\DeclareMathOperator{\Hhomtopan}{\mathbf{Homtopan}}
\DeclareMathOperator{\Pptgeom}{\mathbf{Ptgéom}}
\DeclareMathOperator{\Morsite}{Morsite}
\DeclareMathOperator{\Fib}{Fib}
\DeclareMathOperator{\Fais}{Fais}
\DeclareMathOperator{\kar}{kar}
\DeclareMathOperator{\Groupes}{Groupes}
\DeclareMathOperator{\Restr}{Restr}
\DeclareMathOperator{\Isom}{Isom}
\DeclareMathOperator{\Prol}{Prol}
\DeclareMathOperator{\Res}{Res}
\DeclareMathOperator{\Tors}{Tors}
\DeclareMathOperator{\esp}{esp}
\DeclareMathOperator{\es}{es}
\DeclareMathOperator{\od}{od}
\DeclareMathOperator{\gr}{gr}
\DeclareMathOperator{\sob}{sob}
\DeclareMathOperator{\dfn}{dfn}
\DeclareMathOperator{\Sch}{Sch}
\DeclareMathOperator{\ouv}{ouv}
\DeclareMathOperator{\cons}{cons}
\DeclareMathOperator{\Nerf}{Nerf}
\DeclareMathOperator{\Points}{Points}
\DeclareMathOperator{\Espsob}{Espsob}
\DeclareMathOperator{\Homtop}{Homtop}
\DeclareMathOperator{\Pt}{Pt}
\DeclareMathOperator{\Quot}{Quot}
\DeclareMathOperator{\glob}{glob}
\DeclareMathOperator{\Homfib}{Homfib}
\DeclareMathOperator{\Qouv}{Qouv}
\DeclareMathOperator{\Quotouv}{Quotouv}
\DeclareMathOperator{\ens}{ens}
\DeclareMathOperator{\ab}{ab}
\DeclareMathOperator{\an}{an}
\DeclareMathOperator{\fs}{fs}
\DeclareMathOperator{\red}{red}
\DeclareMathOperator{\Ab}{Ab}
\DeclareMathOperator{\cosq}{cosq}
\DeclareMathOperator{\Ann}{Ann}
\DeclareMathOperator{\Zar}{Zar}


%%%%%%%%%%%%
%%%%% SGA4-2
\DeclareMathOperator{\Sex}{Sex}
\DeclareMathOperator{\vois}{vois}
\DeclareMathOperator{\Vois}{Vois}
\DeclareMathOperator{\Ext}{Ext}
\DeclareMathOperator{\ext}{ext}
\DeclareMathOperator{\Fer}{Fer}
\DeclareMathOperator{\Add}{Add}
\DeclareMathOperator{\Cocart}{Cocart}
\DeclareMathOperator{\cosk}{cosk}
\DeclareMathOperator{\plat}{plat}

\DeclareMathOperator{\faisc}{faisc}
\DeclareMathOperator{\sk}{sk}
\DeclareMathOperator{\ch}{ch}

\def\HH{\mathbf{H}}
\def\qc{\mathrm{qc}}
\def\et{\mathrm{et}}
\def\coh{\mathrm{coh}}

\def\Hvr#1{%
\begin{array}[b]{l}
{\kern2pt\scriptstyle r}\\[-8pt]
{\kern1pt\scriptstyle \vee}\\[-5pt]
H^{#1}
\end{array}}

\def\Hvv#1{%
\begin{array}[b]{l}
\kern2pt\rotatebox{90}{${\scriptscriptstyle\ll}$}\\[-6pt]
H^{#1}
\end{array}}

\def\Hvi#1{%
\begin{array}[b]{l}
{\kern0.2pt\scriptscriptstyle \infty}\\[-8pt]
{\kern1pt\scriptstyle \vee}\\[-5pt]
H^{#1}
\end{array}}

\def\Hv#1{%
\begin{array}[b]{l}
{\kern1pt\scriptstyle \vee}\\[-5pt]
H^{#1}
\end{array}}

\DeclareMathOperator{\cart}{cart}
\DeclareMathOperator{\Ffaisc}{\mathbf{Faisc}}


\def\Ev#1{%
\begin{array}[b]{l}
{\kern1pt\scriptstyle \vee}\\[-5pt]
E'_{#1}
\end{array}}

\DeclareMathOperator{\Cont}{Cont}
\DeclareMathOperator{\Ccont}{\mathscr{C}\!\mathit{ont}}
\DeclareMathOperator{\Site}{Site}
\newcommand\MM{\mathbf{M}}
\def\limtop#1{{\displaystyle{\varprojlim_{#1}\!\mbox{\rm top }}}}
\def\Limtop{\mathop{\varprojlim\!\mbox{\rm top}}}
\DeclareMathOperator{\spec}{spec}
\def\ss{\scriptstyle{\mathscr{S}}}
\DeclareMathOperator{\Gal}{Gal}
\def\smallm{\scriptstyle{\mathscr{M}}}
\DeclareMathOperator{\End}{End}
\DeclareMathOperator{\ass}{ass}
\DeclareMathOperator{\zar}{zar}
\DeclareMathOperator{\Ogg}{Ogg}
\DeclareMathOperator{\res}{res}
\DeclareMathOperator{\tr}{tr}
\DeclareMathOperator{\Supp}{Supp}
\DeclareMathOperator{\codim}{codim}
\DeclareMathOperator{\rk}{rk}
\DeclareMathOperator{\cur}{cur}
\DeclareMathOperator{\cdqc}{cdqc}
\newcommand\DD{\mathbf{D}}
\DeclareMathOperator{\cl}{cl}
\def\JJ{\mathbf{J}}
\def\FF{\mathbf{F}}
\def\BB{\mathbf{B}}
\DeclareMathOperator{\Alget}{Alget}
\DeclareMathOperator{\Ass}{Ass}
\DeclareMathOperator{\supp}{supp}
\DeclareMathOperator{\rad}{rad}
\DeclareMathOperator{\prof}{prof}
\DeclareMathOperator{\cint}{int}
\DeclareMathOperator{\sic}{sic}
\DeclareMathOperator{\Tor}{Tor}
\DeclareMathOperator{\tors}{tors}
\DeclareMathOperator{\Sym}{Sym}
\DeclareMathOperator{\parf}{parf}
\DeclareMathOperator{\tor}{tor}
\DeclareMathOperator{\defn}{defn}
\DeclareMathOperator{\Div}{Div}
\DeclareMathOperator{\Pol}{Pol}
\DeclareMathOperator{\Fex}{Fex}
\DeclareMathOperator{\Lie}{Lie}
\DeclareMathOperator{\tot}{tot}
\DeclareMathOperator{\Aut}{Aut}
\DeclareMathOperator{\HOM}{HOM}
\DeclareMathOperator{\pch}{pch}
\DeclareMathOperator{\Ch}{Ch}
\DeclareMathOperator{\EXT}{EXT}
\DeclareMathOperator{\PIC}{PIC}
\DeclareMathOperator{\Rf}{Rf}
\DeclareMathOperator{\inj}{inj}
\newcommand\VV{\mathbf{V}}

\def\Xv{%
\begin{array}[b]{l}
{\kern1pt\scriptstyle \vee}\\[-5pt]
X
\end{array}}

%%%%%%%
\newcommand\TT{\mathbf{T}}
\newcommand\XX{\mathbf{X}}
\newcommand\EE{\mathbf{E}}
%%%%%%

\def\qedsymbol{\textup{C.Q.F.D.}}
\let\CQFD\qed
\let\cqfd\qed

%%%%%\bigdot
\font\bigsymb=cmsy10 at 4pt
\def\bigdot{{\kern1.2pt\raise 1.5pt\hbox{\bigsymb\char15}}}
%%%%%

\def\Lright{\mathop{\vcenter{\hbox{\hbox{$\longrightarrow$}\lower8pt \hbox{\kern-16pt\hbox{$\longrightarrow$}}}}}}

\def\Aa{{\mathcal A}}
\def\Bb{{\mathcal B}}
\def\Cc{{\mathcal C}}
\def\Dd{{\mathcal D}}
\def\Ee{{\mathcal E}}
\def\Ff{{\mathcal F}}
\def\Gg{{\mathcal G}}
\def\Hh{{\mathcal H}}
\def\Ii{{\mathcal I}}
\def\Jj{{\mathcal J}}
\def\Kk{{\mathcal K}}
\def\Ll{{\mathcal L}}
\def\Mm{{\mathcal M}}
\def\Nn{{\mathcal N}}
\def\Oo{{\mathcal O}}
\def\Pp{{\mathcal P}}
\def\Qq{{\mathcal Q}}
\def\Rr{{\mathcal R}}
\def\Ss{{\mathcal S}}
\def\Tt{{\mathcal T}}
\def\Uu{{\mathcal U}}
\def\Vv{{\mathcal V}}
\def\Ww{{\mathcal W}}
\def\Xx{{\mathcal X}}
\def\Yy{{\mathcal Y}}
\def\Zz{{\mathcal Z}}

\newcommand{\CC}{\mathbf{C}}
\newcommand{\GG}{\mathbf{G}}
\newcommand{\KK}{\mathbf{K}}
\newcommand{\LL}{\mathbf{L}}
\newcommand{\NN}{\mathbf{N}}
\newcommand{\PP}{\mathbf{P}}
\newcommand\QQ{\mathbf{Q}}
\newcommand\RR{\mathbf{R}}
\newcommand\ZZ{\mathbf{Z}}

\def\to{\mathchoice{\longrightarrow}{\rightarrow}{\rightarrow}{\rightarrow}}
\def\from{\mathchoice{\longleftarrow}{\leftarrow}{\leftarrow}{\leftarrow}}
\def\mto{\mathchoice{\longmapsto}{\mapsto}{\mapsto}{\mapsto}}
\def\mfrom{\mathrel{\from\joinrel\mapstochar\,}}
\def\hto{\mathrel{\lhook\joinrel\to}}
\def\hfrom{\mathrel{\from\joinrel\rhook}}

\def\implies{\mathchoice{\Longrightarrow}{\Rightarrow}{\Rightarrow}{\Rightarrow}}
\def\implied{\mathchoice{\Longleftarrow}{\Leftarrow}{\Leftarrow}{\Leftarrow}}
\def\iff{\mathchoice{\Longleftrightarrow}{\Leftrightarrow}{\Leftrightarrow}{\Leftrightarrow}}

\def\To#1{\mathchoice{\xrightarrow{\labelstyle~~#1~}}{\stackrel{#1}{\longrightarrow}}{}{}}
\def\Mto#1{\mathrel{\mapstochar\To{#1}}}
\def\From#1{\mathchoice{\xleftarrow{\labelstyle~~#1~}}{\stackrel{#1}{\longleftarrow}}{}{}}

\newcommand{\isomto}{\overset{\sim}{\longrightarrow}}
\newcommand{\isomfrom}{\overset{\sim}{\longleftarrow}}

\newcommand{\IFF}{\unskip~$\iff$~\ignorespaces}
\newcommand{\THEN}{\unskip~$\implies$~\ignorespaces}

\let\leq\leqslant
\let\geq\geqslant
\let\le\leqslant
\let\ge\geqslant

\def\cf{{cf.}\kern.3em}
\def\Cf{{Cf.}\kern.3em}
\def\ie{{i.e.}\kern.3em}
\def\iev{{i.e.},\xspace}
%\def\loccito{loc.\kern2pt cit.\xspace}
%\def\loccitn{\textit{loc.\kern2pt cit.}\xspace}
%\def\loccit{\sisi{\loccito}{\loccitn}}
\def\EGA{\textup{EGA}\kern.2em}
\def\SGA{\textup{SGA}\kern.3em}

\let\oldS\S
\def\S{\oldS\kern .15em}
\def\ptbl{.\kern .2em}
\def\resp{resp.\kern.3em}
\def\Exp{Exp.\kern.3em}

\let\bar\overline
\let\hat\widehat
\let\tilde\widetilde
\let\dpl\displaystyle
\let\ldots\dots
\let\emptyset\varnothing
\let\skpt\skippointrait
\let\epsilon\varepsilon

\def\sfrac#1#2{{#1}/{#2}}

\DeclareMathOperator{\cd}{cd}
%\DeclareMathOperator{\Hhom}{\mathcal{H}\!\mathit{om}}
\DeclareMathOperator{\Hhom}{\mathbf{Hom}}
%\DeclareMathOperator{\Eext}{\mathcal{E}\!\mathit{xt}}
\DeclareMathOperator{\Eext}{\mathbf{Ext}}
\DeclareMathOperator{\Spec}{Spec}
\DeclareMathOperator{\Tr}{Tr}

\let\emptyset=\varnothing

%%%%%%%%%%%%%%%
%%% expose XVII
%%%%%%%%%%%%%%%

\makeatletter
\let\old@optexpose\optexpose
\newcommand{\empty@optexpose}[1]{}
\newcommand{\paren@optexpose}[1]{\old@optexpose{#1}(}
\newcommand{\refnoexpose}[1]{\let\optexpose\empty@optexpose\ref{#1}\let\optexpose\old@optexpose}%\refnoexpose{X.cor5.2} engendre "5.2" au lieu de "X 5.2"
\newcommand{\neweqref}[1]{\let\optexpose\paren@optexpose\textup{\ref{#1})}\let\optexpose\old@optexpose}%\neweqref{XII.eq4.1} engendre "XII (4.1)" au lieu de "(XII 4.1)"
\makeatother
\newcommand{\refvb}[2][VI${}_{\text{B}}$ ]{\sisi{#1\refnoexpose{#2}}{\ref{#2}}}

\newcommand{\leftexp}[2]{{\vphantom{#2}}^{#1}{#2}}
\newcommand{\AAA}{\mathbf{A}}
\newcommand{\Sf}{\mathfrak{S}}
\DeclareMathOperator*{\quotprojlim}{\text{``$\varprojlim$''}}
\DeclareMathOperator*{\quotinjlim}{\text{``$\varinjlim$''}}
\DeclareMathOperator{\Fl}{F\sisi{\ell}{l}}
\DeclareMathOperator{\Pl}{P\sisi{\ell}{l}}
%\DeclareMathOperator{\cl}{c\sisi{\ell}{l}}
\DeclareMathOperator{\tdf}{tdf}
\DeclareMathOperator{\rg}{rg}
\newcommand{\iInd}{\sisi{I}{i}nd}
\newcommand{\iind}{ind}
\newcommand{\ppro}{pro}
\newcommand{\tTor}{\sisi{T}{t}or}
\newcommand{\ttor}{tor}
\newcommand{\chap}{Chap.}
\newcommand{\Ffex}{\Ff\mathit{ex}}
\newcommand{\sheafk}{\underline{k}}
\newcommand{\sheafn}{\underline{n}}
\newcommand{\sheafGamma}{\underline{\Gamma}}

% Repeated references to an already typeset footnote must not use
% \footnotemark without a matching \footnotetext: Hyperref would then create
% a live annotation to a destination that cannot exist.  This helper keeps the
% current printed footnote mark and links it to an explicit stable note target.
\makeatletter
\newcommand{\sgarepeatedfootnotemark}[1]{%
  \begingroup
  \edef\@thefnmark{\thefootnote}%
  \hyperlink{#1}{\@makefnmark}%
  \endgroup}
\makeatother

\sisi{
\DeclareMathOperator{\simarrow}{\rule{.15mm}{1.8mm}\kern -3.4pt\rightsquigarrow}
\newcommand{\sheafell}{\underline{\ell}}
\newcommand{\sheafZZ}{\underline{\ZZ}}
}{
\newcommand{\nderef}[1]{la N.D.E.~(\ref{#1}) page~\pageref{#1}}
}

%%%%%%%%%%%%%%%%%
%%%%% expose XVII
%%%%%%%%%%%%%%%%%

% END native input: style/macros.tex
\providecommand{\SGAArchiveOnly}[1]{}
% BEGIN native input: style/localization.tex
% Chinese adaptation of this SGA 4 source's own AMS/SMF layout.
\usepackage{xeCJK}
\setCJKmainfont[BoldFont={Noto Serif CJK SC Bold},ItalicFont={Noto Serif CJK SC},SmallCapsFont={Noto Serif CJK SC}]{Noto Serif CJK SC}
\setCJKsansfont{Noto Sans CJK SC}
\setCJKmonofont{Noto Sans Mono CJK SC}
\def\chaptername{第}
\def\contentsname{目录}
\def\bibname{参考文献}
\def\appendixname{附录}
\def\proofname{证明}
\def\theoname{定理}\def\propname{命题}\def\conjname{猜想}
\def\coroname{推论}\def\lemmname{引理}\def\definame{定义}
\def\remaname{注}\def\exemname{例}
\def\cf{\ifmmode\text{参见}\else 参见\fi\kern.3em}\def\Cf{参见\kern.3em}
\def\ie{\ifmmode\text{即}\else 即\fi\kern.3em}\def\iev{即\xspace}
\def\resp{\ifmmode\text{分别}\else 分别\fi}\def\etcn{等等\xspace}\def\etco{等等\xspace}
\def\quaet{\quad\text{且}\quad}
\def\loccit{前引文献\xspace}
% Heading label follows original Roman exposé numbering.
\def\thechapter{\Roman{chapter}}
% SGA4 uses numbered paragraphs below subsubsections. The inherited AMS class
% stops at level3, suppressing those printed numbers despite explicit labels.
\setcounter{secnumdepth}{5}
\hypersetup{pdftitle={SGA 4：拓扑斯理论与概形的平展上同调，第 X--XIX 讲（中文首译稿）},pdfauthor={M. Artin; A. Grothendieck; J.-L. Verdier},pdfsubject={SGA 4 中文原生 LaTeX，覆盖状况见 PROGRESS.json},pdfkeywords={SGA 4, 拓扑斯, 平展上同调, 中文翻译, X--XIX}}

\emergencystretch=2em
\usepackage{etoolbox}
\makeatletter
\patchcmd{\@makechapterhead}{\enspace\thechapter}{\enspace\thechapter\enspace 讲}{}{}
\makeatother
\AtBeginDocument{\def\chaptername{第}\def\contentsname{目录}\def\bibname{参考文献}\def\proofname{证明}}
% BEGIN native input: style/external-references.tex
% Generated source-backed external reference registry. No synthetic content/pages.
\expandafter\def\csname SGAsource@5.3'\endcsname{3'}
\expandafter\def\csname SGAsource@5.4'\endcsname{4'}
\expandafter\def\csname SGAsource@I.4.3.1\endcsname{4.3.1}
\expandafter\def\csname SGAsource@I.8.12.1.1\endcsname{8.12.1.1}
\expandafter\def\csname SGAsource@I.8.9.12\endcsname{8.9.12}
\expandafter\def\csname SGAsource@I.9.12.1\endcsname{9.12.1}
\expandafter\def\csname SGAsource@I.condef1.3\endcsname{1.3}
\expandafter\def\csname SGAsource@I.cor2.10\endcsname{2.10}
\expandafter\def\csname SGAsource@I.cor2.3.2\endcsname{2.3.2}
\expandafter\def\csname SGAsource@I.cor2.9\endcsname{2.9}
\expandafter\def\csname SGAsource@I.cor3.2\endcsname{3.2}
\expandafter\def\csname SGAsource@I.cor3.3\endcsname{3.3}
\expandafter\def\csname SGAsource@I.cor3.5\endcsname{3.5}
\expandafter\def\csname SGAsource@I.cor5.5.1\endcsname{5.5.1}
\expandafter\def\csname SGAsource@I.cor6.3\endcsname{6.3}
\expandafter\def\csname SGAsource@I.cor7.11\endcsname{7.11}
\expandafter\def\csname SGAsource@I.cor7.12\endcsname{7.12}
\expandafter\def\csname SGAsource@I.cor7.3\endcsname{7.3}
\expandafter\def\csname SGAsource@I.cor7.5\endcsname{7.5}
\expandafter\def\csname SGAsource@I.cor7.7.1\endcsname{7.7.1}
\expandafter\def\csname SGAsource@I.cor8.1.7\endcsname{8.1.7}
\expandafter\def\csname SGAsource@I.cor8.11.3\endcsname{8.11.3}
\expandafter\def\csname SGAsource@I.cor8.12.7\endcsname{8.12.7}
\expandafter\def\csname SGAsource@I.cor8.12.8\endcsname{8.12.8}
\expandafter\def\csname SGAsource@I.cor8.13.3\endcsname{8.13.3}
\expandafter\def\csname SGAsource@I.cor8.13.4\endcsname{8.13.4}
\expandafter\def\csname SGAsource@I.cor8.5.2\endcsname{8.5.2}
\expandafter\def\csname SGAsource@I.cor8.6.5\endcsname{8.6.5}
\expandafter\def\csname SGAsource@I.cor8.7.6\endcsname{8.7.6}
\expandafter\def\csname SGAsource@I.cor8.7.7\endcsname{8.7.7}
\expandafter\def\csname SGAsource@I.cor8.9.2\endcsname{8.9.2}
\expandafter\def\csname SGAsource@I.cor8.9.3\endcsname{8.9.3}
\expandafter\def\csname SGAsource@I.cor8.9.8\endcsname{8.9.8}
\expandafter\def\csname SGAsource@I.cor9.15\endcsname{9.15}
\expandafter\def\csname SGAsource@I.cor9.17\endcsname{9.17}
\expandafter\def\csname SGAsource@I.cor9.18\endcsname{9.18}
\expandafter\def\csname SGAsource@I.cor9.19\endcsname{9.19}
\expandafter\def\csname SGAsource@I.cor9.20\endcsname{9.20}
\expandafter\def\csname SGAsource@I.cor9.21.10.1\endcsname{9.21.10.1}
\expandafter\def\csname SGAsource@I.cor9.23\endcsname{9.23}
\expandafter\def\csname SGAsource@I.cor9.24\endcsname{9.24}
\expandafter\def\csname SGAsource@I.cor9.25\endcsname{9.25}
\expandafter\def\csname SGAsource@I.cor9.26\endcsname{9.26}
\expandafter\def\csname SGAsource@I.cor9.7\endcsname{9.7}
\expandafter\def\csname SGAsource@I.cor9.8\endcsname{9.8}
\expandafter\def\csname SGAsource@I.cor9.9\endcsname{9.9}
\expandafter\def\csname SGAsource@I.def1.1\endcsname{1.1}
\expandafter\def\csname SGAsource@I.def1.2\endcsname{1.2}
\expandafter\def\csname SGAsource@I.def2.1\endcsname{2.1}
\expandafter\def\csname SGAsource@I.def2.2\endcsname{2.2}
\expandafter\def\csname SGAsource@I.def2.4.1\endcsname{2.4.1}
\expandafter\def\csname SGAsource@I.def2.5\endcsname{2.5}
\expandafter\def\csname SGAsource@I.def2.7\endcsname{2.7}
\expandafter\def\csname SGAsource@I.def4.1\endcsname{4.1}
\expandafter\def\csname SGAsource@I.def6.1\endcsname{6.1}
\expandafter\def\csname SGAsource@I.def7.1\endcsname{7.1}
\expandafter\def\csname SGAsource@I.def8.1.1\endcsname{8.1.1}
\expandafter\def\csname SGAsource@I.def8.1.8\endcsname{8.1.8}
\expandafter\def\csname SGAsource@I.def9.1\endcsname{9.1}
\expandafter\def\csname SGAsource@I.def9.12\endcsname{9.12}
\expandafter\def\csname SGAsource@I.def9.2\endcsname{9.2}
\expandafter\def\csname SGAsource@I.def9.3\endcsname{9.3}
\expandafter\def\csname SGAsource@I.def9.4\endcsname{9.4}
\expandafter\def\csname SGAsource@I.eq7.6.1\endcsname{1}
\expandafter\def\csname SGAsource@I.exem2.7.1\endcsname{2.7.1}
\expandafter\def\csname SGAsource@I.exem7.13\endcsname{7.13}
\expandafter\def\csname SGAsource@I.exem8.11.6\endcsname{8.11.6}
\expandafter\def\csname SGAsource@I.exem9.12.2\endcsname{9.12.2}
\expandafter\def\csname SGAsource@I.exer8.12.9\endcsname{8.12.9}
\expandafter\def\csname SGAsource@I.exer8.2.8\endcsname{8.2.8}
\expandafter\def\csname SGAsource@I.exer8.5.5\endcsname{8.5.5}
\expandafter\def\csname SGAsource@I.exer8.7.8\endcsname{8.7.8}
\expandafter\def\csname SGAsource@I.exer8.7.9\endcsname{8.7.9}
\expandafter\def\csname SGAsource@I.exer8.8.6\endcsname{8.8.6}
\expandafter\def\csname SGAsource@I.exer8.8.7\endcsname{8.8.7}
\expandafter\def\csname SGAsource@I.exer8.8.8\endcsname{8.8.8}
\expandafter\def\csname SGAsource@I.exer8.9.10\endcsname{8.9.10}
\expandafter\def\csname SGAsource@I.exer8.9.11\endcsname{8.9.11}
\expandafter\def\csname SGAsource@I.exer8.9.9\endcsname{8.9.9}
\expandafter\def\csname SGAsource@I.exer9.20.1\endcsname{9.20.1}
\expandafter\def\csname SGAsource@I.lem2.8.1\endcsname{2.8.1}
\expandafter\def\csname SGAsource@I.lem6.5.2\endcsname{6.5.2}
\expandafter\def\csname SGAsource@I.lem7.6.2\endcsname{7.6.2}
\expandafter\def\csname SGAsource@I.lem8.12.4\endcsname{8.12.4}
\expandafter\def\csname SGAsource@I.lem8.12.8.1\endcsname{8.12.8.1}
\expandafter\def\csname SGAsource@I.lem8.8.3.4\endcsname{8.8.3.4}
\expandafter\def\csname SGAsource@I.lem8.8.5.1\endcsname{8.8.5.1}
\expandafter\def\csname SGAsource@I.lem9.11.1\endcsname{9.11.1}
\expandafter\def\csname SGAsource@I.lem9.11.2\endcsname{9.11.2}
\expandafter\def\csname SGAsource@I.lem9.13.1\endcsname{9.13.1}
\expandafter\def\csname SGAsource@I.lem9.16.2\endcsname{9.16.2}
\expandafter\def\csname SGAsource@I.lem9.21.10\endcsname{9.21.10}
\expandafter\def\csname SGAsource@I.lem9.21.13\endcsname{9.21.13}
\expandafter\def\csname SGAsource@I.lem9.21.16\endcsname{9.21.16}
\expandafter\def\csname SGAsource@I.lem9.21.17\endcsname{9.21.17}
\expandafter\def\csname SGAsource@I.lem9.21.18\endcsname{9.21.18}
\expandafter\def\csname SGAsource@I.lem9.21.19\endcsname{9.21.19}
\expandafter\def\csname SGAsource@I.lem9.21.9\endcsname{9.21.9}
\expandafter\def\csname SGAsource@I.not5.9\endcsname{5.9}
\expandafter\def\csname SGAsource@I.prop1.4\endcsname{1.4}
\expandafter\def\csname SGAsource@I.prop2.11\endcsname{2.11}
\expandafter\def\csname SGAsource@I.prop2.12\endcsname{2.12}
\expandafter\def\csname SGAsource@I.prop2.3\endcsname{2.3}
\expandafter\def\csname SGAsource@I.prop2.3.3\endcsname{2.3.3}
\expandafter\def\csname SGAsource@I.prop2.6\endcsname{2.6}
\expandafter\def\csname SGAsource@I.prop2.8\endcsname{2.8}
\expandafter\def\csname SGAsource@I.prop3.1\endcsname{3.1}
\expandafter\def\csname SGAsource@I.prop3.4\endcsname{3.4}
\expandafter\def\csname SGAsource@I.prop3.6\endcsname{3.6}
\expandafter\def\csname SGAsource@I.prop4.2\endcsname{4.2}
\expandafter\def\csname SGAsource@I.prop5.1\endcsname{5.1}
\expandafter\def\csname SGAsource@I.prop5.11\endcsname{5.11}
\expandafter\def\csname SGAsource@I.prop5.14\endcsname{5.14}
\expandafter\def\csname SGAsource@I.prop5.2\endcsname{5.2}
\expandafter\def\csname SGAsource@I.prop5.4\endcsname{5.4}
\expandafter\def\csname SGAsource@I.prop5.5\endcsname{5.5}
\expandafter\def\csname SGAsource@I.prop5.6\endcsname{5.6}
\expandafter\def\csname SGAsource@I.prop5.8.3\endcsname{5.8.3}
\expandafter\def\csname SGAsource@I.prop6.2\endcsname{6.2}
\expandafter\def\csname SGAsource@I.prop6.4\endcsname{6.4}
\expandafter\def\csname SGAsource@I.prop6.5\endcsname{6.5}
\expandafter\def\csname SGAsource@I.prop7.10\endcsname{7.10}
\expandafter\def\csname SGAsource@I.prop7.2\endcsname{7.2}
\expandafter\def\csname SGAsource@I.prop7.4\endcsname{7.4}
\expandafter\def\csname SGAsource@I.prop7.6\endcsname{7.6}
\expandafter\def\csname SGAsource@I.prop7.7\endcsname{7.7}
\expandafter\def\csname SGAsource@I.prop7.8\endcsname{7.8}
\expandafter\def\csname SGAsource@I.prop8.1.3\endcsname{8.1.3}
\expandafter\def\csname SGAsource@I.prop8.1.6\endcsname{8.1.6}
\expandafter\def\csname SGAsource@I.prop8.11.2\endcsname{8.11.2}
\expandafter\def\csname SGAsource@I.prop8.11.4\endcsname{8.11.4}
\expandafter\def\csname SGAsource@I.prop8.12.2\endcsname{8.12.2}
\expandafter\def\csname SGAsource@I.prop8.12.5\endcsname{8.12.5}
\expandafter\def\csname SGAsource@I.prop8.13.2\endcsname{8.13.2}
\expandafter\def\csname SGAsource@I.prop8.3.1\endcsname{8.3.1}
\expandafter\def\csname SGAsource@I.prop8.5.1\endcsname{8.5.1}
\expandafter\def\csname SGAsource@I.prop8.6.3\endcsname{8.6.3}
\expandafter\def\csname SGAsource@I.prop8.6.4\endcsname{8.6.4}
\expandafter\def\csname SGAsource@I.prop8.7.3\endcsname{8.7.3}
\expandafter\def\csname SGAsource@I.prop8.7.5\endcsname{8.7.5}
\expandafter\def\csname SGAsource@I.prop8.8.2\endcsname{8.8.2}
\expandafter\def\csname SGAsource@I.prop8.8.3\endcsname{8.8.3}
\expandafter\def\csname SGAsource@I.prop8.8.5\endcsname{8.8.5}
\expandafter\def\csname SGAsource@I.prop8.9.1\endcsname{8.9.1}
\expandafter\def\csname SGAsource@I.prop8.9.5\endcsname{8.9.5}
\expandafter\def\csname SGAsource@I.prop8.9.7\endcsname{8.9.7}
\expandafter\def\csname SGAsource@I.prop9.11\endcsname{9.11}
\expandafter\def\csname SGAsource@I.prop9.13\endcsname{9.13}
\expandafter\def\csname SGAsource@I.prop9.14\endcsname{9.14}
\expandafter\def\csname SGAsource@I.prop9.16\endcsname{9.16}
\expandafter\def\csname SGAsource@I.prop9.5\endcsname{9.5}
\expandafter\def\csname SGAsource@I.prop9.6\endcsname{9.6}
\expandafter\def\csname SGAsource@I.rem1.1.2\endcsname{1.1.2}
\expandafter\def\csname SGAsource@I.rem4.2.1\endcsname{4.2.1}
\expandafter\def\csname SGAsource@I.rem5.5.2\endcsname{5.5.2}
\expandafter\def\csname SGAsource@I.rem5.7\endcsname{5.7}
\expandafter\def\csname SGAsource@I.rem8.12.6\endcsname{8.12.6}
\expandafter\def\csname SGAsource@I.rem8.3.4\endcsname{8.3.4}
\expandafter\def\csname SGAsource@I.rem8.7.5.2\endcsname{8.7.5.2}
\expandafter\def\csname SGAsource@I.rem8.8.4\endcsname{8.8.4}
\expandafter\def\csname SGAsource@I.rem8.9.4\endcsname{8.9.4}
\expandafter\def\csname SGAsource@I.rem8.9.6\endcsname{8.9.6}
\expandafter\def\csname SGAsource@I.rem9.10\endcsname{9.10}
\expandafter\def\csname SGAsource@I.rem9.11.3\endcsname{9.11.3}
\expandafter\def\csname SGAsource@I.rem9.13.2\endcsname{9.13.2}
\expandafter\def\csname SGAsource@I.sec0\endcsname{0}
\expandafter\def\csname SGAsource@I.sec1\endcsname{1}
\expandafter\def\csname SGAsource@I.sec1.0\endcsname{1.0}
\expandafter\def\csname SGAsource@I.sec1.1.1\endcsname{1.1.1}
\expandafter\def\csname SGAsource@I.sec1.3.1\endcsname{1.3.1}
\expandafter\def\csname SGAsource@I.sec1.3.2\endcsname{1.3.2}
\expandafter\def\csname SGAsource@I.sec1.3.3\endcsname{1.3.3}
\expandafter\def\csname SGAsource@I.sec1.4.1\endcsname{1.4.1}
\expandafter\def\csname SGAsource@I.sec1.4.2\endcsname{1.4.2}
\expandafter\def\csname SGAsource@I.sec10\endcsname{10}
\expandafter\def\csname SGAsource@I.sec10.1\endcsname{10.1}
\expandafter\def\csname SGAsource@I.sec10.10\endcsname{10.10}
\expandafter\def\csname SGAsource@I.sec10.11\endcsname{10.11}
\expandafter\def\csname SGAsource@I.sec10.2\endcsname{10.2}
\expandafter\def\csname SGAsource@I.sec10.3\endcsname{10.3}
\expandafter\def\csname SGAsource@I.sec10.4\endcsname{10.4}
\expandafter\def\csname SGAsource@I.sec10.5\endcsname{10.5}
\expandafter\def\csname SGAsource@I.sec10.6\endcsname{10.6}
\expandafter\def\csname SGAsource@I.sec10.7\endcsname{10.7}
\expandafter\def\csname SGAsource@I.sec10.8\endcsname{10.8}
\expandafter\def\csname SGAsource@I.sec10.9\endcsname{10.9}
\expandafter\def\csname SGAsource@I.sec2\endcsname{2}
\expandafter\def\csname SGAsource@I.sec2.1.1\endcsname{2.1.1}
\expandafter\def\csname SGAsource@I.sec2.2.1\endcsname{2.2.1}
\expandafter\def\csname SGAsource@I.sec2.3.1\endcsname{2.3.1}
\expandafter\def\csname SGAsource@I.sec2.4.2\endcsname{2.4.2}
\expandafter\def\csname SGAsource@I.sec2.5.0\endcsname{2.5.0}
\expandafter\def\csname SGAsource@I.sec3\endcsname{3}
\expandafter\def\csname SGAsource@I.sec3.4.0\endcsname{3.4.0}
\expandafter\def\csname SGAsource@I.sec4\endcsname{4}
\expandafter\def\csname SGAsource@I.sec4.3\endcsname{4.3}
\expandafter\def\csname SGAsource@I.sec4.3.2\endcsname{4.3.2}
\expandafter\def\csname SGAsource@I.sec4.3.3\endcsname{4.3.3}
\expandafter\def\csname SGAsource@I.sec5\endcsname{5}
\expandafter\def\csname SGAsource@I.sec5.0\endcsname{5.0}
\expandafter\def\csname SGAsource@I.sec5.10\endcsname{5.10}
\expandafter\def\csname SGAsource@I.sec5.12\endcsname{5.12}
\expandafter\def\csname SGAsource@I.sec5.13\endcsname{5.13}
\expandafter\def\csname SGAsource@I.sec5.3\endcsname{5.3}
\expandafter\def\csname SGAsource@I.sec5.8.0\endcsname{5.8.0}
\expandafter\def\csname SGAsource@I.sec5.8.2\endcsname{5.8.2}
\expandafter\def\csname SGAsource@I.sec6\endcsname{6}
\expandafter\def\csname SGAsource@I.sec6.1.1\endcsname{6.1.1}
\expandafter\def\csname SGAsource@I.sec6.5.1\endcsname{6.5.1}
\expandafter\def\csname SGAsource@I.sec7\endcsname{7}
\expandafter\def\csname SGAsource@I.sec7.1.1\endcsname{7.1.1}
\expandafter\def\csname SGAsource@I.sec7.1.2\endcsname{7.1.2}
\expandafter\def\csname SGAsource@I.sec7.1.3\endcsname{7.1.3}
\expandafter\def\csname SGAsource@I.sec7.1.3.1\endcsname{7.1.3.1}
\expandafter\def\csname SGAsource@I.sec7.5.1\endcsname{7.5.1}
\expandafter\def\csname SGAsource@I.sec7.5.3\endcsname{7.5.3}
\expandafter\def\csname SGAsource@I.sec7.9\endcsname{7.9}
\expandafter\def\csname SGAsource@I.sec8\endcsname{8}
\expandafter\def\csname SGAsource@I.sec8.1\endcsname{8.1}
\expandafter\def\csname SGAsource@I.sec8.1.2\endcsname{8.1.2}
\expandafter\def\csname SGAsource@I.sec8.1.2.2\endcsname{8.1.2.2}
\expandafter\def\csname SGAsource@I.sec8.1.4\endcsname{8.1.4}
\expandafter\def\csname SGAsource@I.sec8.1.5\endcsname{8.1.5}
\expandafter\def\csname SGAsource@I.sec8.10\endcsname{8.10}
\expandafter\def\csname SGAsource@I.sec8.10.10\endcsname{8.10.10}
\expandafter\def\csname SGAsource@I.sec8.10.14\endcsname{8.10.14}
\expandafter\def\csname SGAsource@I.sec8.10.15\endcsname{8.10.15}
\expandafter\def\csname SGAsource@I.sec8.11\endcsname{8.11}
\expandafter\def\csname SGAsource@I.sec8.11.1\endcsname{8.11.1}
\expandafter\def\csname SGAsource@I.sec8.11.5\endcsname{8.11.5}
\expandafter\def\csname SGAsource@I.sec8.12\endcsname{8.12}
\expandafter\def\csname SGAsource@I.sec8.12.1\endcsname{8.12.1}
\expandafter\def\csname SGAsource@I.sec8.12.2.2\endcsname{8.12.2.2}
\expandafter\def\csname SGAsource@I.sec8.12.2.3\endcsname{8.12.2.3}
\expandafter\def\csname SGAsource@I.sec8.12.3\endcsname{8.12.3}
\expandafter\def\csname SGAsource@I.sec8.13\endcsname{8.13}
\expandafter\def\csname SGAsource@I.sec8.13.1\endcsname{8.13.1}
\expandafter\def\csname SGAsource@I.sec8.2\endcsname{8.2}
\expandafter\def\csname SGAsource@I.sec8.2.1\endcsname{8.2.1}
\expandafter\def\csname SGAsource@I.sec8.2.2\endcsname{8.2.2}
\expandafter\def\csname SGAsource@I.sec8.2.3\endcsname{8.2.3}
\expandafter\def\csname SGAsource@I.sec8.2.4\endcsname{8.2.4}
\expandafter\def\csname SGAsource@I.sec8.2.5\endcsname{8.2.5}
\expandafter\def\csname SGAsource@I.sec8.2.6\endcsname{8.2.6}
\expandafter\def\csname SGAsource@I.sec8.2.7\endcsname{8.2.7}
\expandafter\def\csname SGAsource@I.sec8.3\endcsname{8.3}
\expandafter\def\csname SGAsource@I.sec8.3.2\endcsname{8.3.2}
\expandafter\def\csname SGAsource@I.sec8.3.2.3\endcsname{8.3.2.3}
\expandafter\def\csname SGAsource@I.sec8.4\endcsname{8.4}
\expandafter\def\csname SGAsource@I.sec8.4.3\endcsname{8.4.3}
\expandafter\def\csname SGAsource@I.sec8.4.4\endcsname{8.4.4}
\expandafter\def\csname SGAsource@I.sec8.5\endcsname{8.5}
\expandafter\def\csname SGAsource@I.sec8.5.3\endcsname{8.5.3}
\expandafter\def\csname SGAsource@I.sec8.5.4\endcsname{8.5.4}
\expandafter\def\csname SGAsource@I.sec8.6\endcsname{8.6}
\expandafter\def\csname SGAsource@I.sec8.6.1\endcsname{8.6.1}
\expandafter\def\csname SGAsource@I.sec8.6.2\endcsname{8.6.2}
\expandafter\def\csname SGAsource@I.sec8.7\endcsname{8.7}
\expandafter\def\csname SGAsource@I.sec8.7.1\endcsname{8.7.1}
\expandafter\def\csname SGAsource@I.sec8.7.1.6\endcsname{8.7.1.6}
\expandafter\def\csname SGAsource@I.sec8.7.1.7\endcsname{8.7.1.7}
\expandafter\def\csname SGAsource@I.sec8.7.2\endcsname{8.7.2}
\expandafter\def\csname SGAsource@I.sec8.7.4\endcsname{8.7.4}
\expandafter\def\csname SGAsource@I.sec8.8\endcsname{8.8}
\expandafter\def\csname SGAsource@I.sec8.8.1\endcsname{8.8.1}
\expandafter\def\csname SGAsource@I.sec8.9\endcsname{8.9}
\expandafter\def\csname SGAsource@I.sec9\endcsname{9}
\expandafter\def\csname SGAsource@I.sec9.0\endcsname{9.0}
\expandafter\def\csname SGAsource@I.sec9.1.1\endcsname{9.1.1}
\expandafter\def\csname SGAsource@I.sec9.1.2\endcsname{9.1.2}
\expandafter\def\csname SGAsource@I.sec9.15.1\endcsname{9.15.1}
\expandafter\def\csname SGAsource@I.sec9.2.1\endcsname{9.2.1}
\expandafter\def\csname SGAsource@I.sec9.21\endcsname{9.21}
\expandafter\def\csname SGAsource@I.sec9.21.14\endcsname{9.21.14}
\expandafter\def\csname SGAsource@I.sec9.21.15\endcsname{9.21.15}
\expandafter\def\csname SGAsource@I.sec9.21.20\endcsname{9.21.20}
\expandafter\def\csname SGAsource@I.sec9.21.21\endcsname{9.21.21}
\expandafter\def\csname SGAsource@I.sec9.3.1\endcsname{9.3.1}
\expandafter\def\csname SGAsource@I.them8.3.3\endcsname{8.3.3}
\expandafter\def\csname SGAsource@I.them9.22\endcsname{9.22}
\expandafter\def\csname SGAsource@I.var7.5.2\endcsname{7.5.2}
\expandafter\def\csname SGAsource@II.cor2.3\endcsname{2.3}
\expandafter\def\csname SGAsource@II.cor2.4\endcsname{2.4}
\expandafter\def\csname SGAsource@II.cor4.1.1\endcsname{4.1.1}
\expandafter\def\csname SGAsource@II.cor4.11\endcsname{4.11}
\expandafter\def\csname SGAsource@II.cor4.12\endcsname{4.12}
\expandafter\def\csname SGAsource@II.cor4.4.4\endcsname{4.4.4}
\expandafter\def\csname SGAsource@II.cor4.6.1\endcsname{4.6.1}
\expandafter\def\csname SGAsource@II.cor4.6.2\endcsname{4.6.2}
\expandafter\def\csname SGAsource@II.cor4.6.3\endcsname{4.6.3}
\expandafter\def\csname SGAsource@II.cor4.6.4\endcsname{4.6.4}
\expandafter\def\csname SGAsource@II.cor6.6\endcsname{6.6}
\expandafter\def\csname SGAsource@II.def1.1\endcsname{1.1}
\expandafter\def\csname SGAsource@II.def1.2\endcsname{1.2}
\expandafter\def\csname SGAsource@II.def1.3\endcsname{1.3}
\expandafter\def\csname SGAsource@II.def2.1\endcsname{2.1}
\expandafter\def\csname SGAsource@II.def2.5\endcsname{2.5}
\expandafter\def\csname SGAsource@II.def3.0.1\endcsname{3.0.1}
\expandafter\def\csname SGAsource@II.def3.0.2\endcsname{3.0.2}
\expandafter\def\csname SGAsource@II.def3.5\endcsname{3.5}
\expandafter\def\csname SGAsource@II.def4.5\endcsname{4.5}
\expandafter\def\csname SGAsource@II.def5.2\endcsname{5.2}
\expandafter\def\csname SGAsource@II.def6.1\endcsname{6.1}
\expandafter\def\csname SGAsource@II.eq4.10.1\endcsname{4.10.1}
\expandafter\def\csname SGAsource@II.eq4.10.2\endcsname{4.10.2}
\expandafter\def\csname SGAsource@II.eq4.6.1\endcsname{4.6.1}
\expandafter\def\csname SGAsource@II.exem4.5.1\endcsname{4.5.1}
\expandafter\def\csname SGAsource@II.lem3.1\endcsname{3.1}
\expandafter\def\csname SGAsource@II.not6.8\endcsname{6.8}
\expandafter\def\csname SGAsource@II.prop1.4\endcsname{1.4}
\expandafter\def\csname SGAsource@II.prop2.2\endcsname{2.2}
\expandafter\def\csname SGAsource@II.prop2.6\endcsname{2.6}
\expandafter\def\csname SGAsource@II.prop3.0.4\endcsname{3.0.4}
\expandafter\def\csname SGAsource@II.prop3.2\endcsname{3.2}
\expandafter\def\csname SGAsource@II.prop3.6\endcsname{3.6}
\expandafter\def\csname SGAsource@II.prop4.10\endcsname{4.10}
\expandafter\def\csname SGAsource@II.prop4.2\endcsname{4.2}
\expandafter\def\csname SGAsource@II.prop4.3\endcsname{4.3}
\expandafter\def\csname SGAsource@II.prop4.6\endcsname{4.6}
\expandafter\def\csname SGAsource@II.prop4.7\endcsname{4.7}
\expandafter\def\csname SGAsource@II.prop4.8\endcsname{4.8}
\expandafter\def\csname SGAsource@II.prop5.1\endcsname{5.1}
\expandafter\def\csname SGAsource@II.prop5.3\endcsname{5.3}
\expandafter\def\csname SGAsource@II.prop5.4\endcsname{5.4}
\expandafter\def\csname SGAsource@II.prop6.3.1\endcsname{6.3.1}
\expandafter\def\csname SGAsource@II.prop6.4\endcsname{6.4}
\expandafter\def\csname SGAsource@II.prop6.5\endcsname{6.5}
\expandafter\def\csname SGAsource@II.prop6.7\endcsname{6.7}
\expandafter\def\csname SGAsource@II.rem2.5.1\endcsname{2.5.1}
\expandafter\def\csname SGAsource@II.rem2.7\endcsname{2.7}
\expandafter\def\csname SGAsource@II.rem3.3\endcsname{3.3}
\expandafter\def\csname SGAsource@II.rem4.3.1\endcsname{4.3.1}
\expandafter\def\csname SGAsource@II.rem5.4.1\endcsname{5.4.1}
\expandafter\def\csname SGAsource@II.rem5.4.2\endcsname{5.4.2}
\expandafter\def\csname SGAsource@II.rem6.2\endcsname{6.2}
\expandafter\def\csname SGAsource@II.rem6.9\endcsname{6.9}
\expandafter\def\csname SGAsource@II.sec1\endcsname{1}
\expandafter\def\csname SGAsource@II.sec2\endcsname{2}
\expandafter\def\csname SGAsource@II.sec3\endcsname{3}
\expandafter\def\csname SGAsource@II.sec4\endcsname{4}
\expandafter\def\csname SGAsource@II.sec5\endcsname{5}
\expandafter\def\csname SGAsource@II.sec6\endcsname{6}
\expandafter\def\csname SGAsource@II.subsec4.9\endcsname{4.9}
\expandafter\def\csname SGAsource@II.subsec6.0\endcsname{6.0}
\expandafter\def\csname SGAsource@II.subsubsec1.1.1\endcsname{1.1.1}
\expandafter\def\csname SGAsource@II.subsubsec1.1.2\endcsname{1.1.2}
\expandafter\def\csname SGAsource@II.subsubsec1.1.3\endcsname{1.1.3}
\expandafter\def\csname SGAsource@II.subsubsec1.1.4\endcsname{1.1.4}
\expandafter\def\csname SGAsource@II.subsubsec1.1.5\endcsname{1.1.5}
\expandafter\def\csname SGAsource@II.subsubsec1.1.6\endcsname{1.1.6}
\expandafter\def\csname SGAsource@II.subsubsec1.3.1\endcsname{1.3.1}
\expandafter\def\csname SGAsource@II.subsubsec3.0.3\endcsname{3.0.3}
\expandafter\def\csname SGAsource@II.subsubsec3.0.5\endcsname{3.0.5}
\expandafter\def\csname SGAsource@II.subsubsec4.4.0\endcsname{4.4.0}
\expandafter\def\csname SGAsource@II.subsubsec5.3.1\endcsname{5.3.1}
\expandafter\def\csname SGAsource@II.subsubsec5.5.0\endcsname{5.5.0}
\expandafter\def\csname SGAsource@II.subsubsec6.3.0\endcsname{6.3.0}
\expandafter\def\csname SGAsource@II.subsubsec6.3.2\endcsname{6.3.2}
\expandafter\def\csname SGAsource@II.subsubsec6.3.3\endcsname{6.3.3}
\expandafter\def\csname SGAsource@II.theo3.4\endcsname{3.4}
\expandafter\def\csname SGAsource@II.theo4.1\endcsname{4.1}
\expandafter\def\csname SGAsource@II.theo4.4\endcsname{4.4}
\expandafter\def\csname SGAsource@II.theo5.5\endcsname{5.5}
\expandafter\def\csname SGAsource@III.cor3.3\endcsname{3.3}
\expandafter\def\csname SGAsource@III.cor3.4\endcsname{3.4}
\expandafter\def\csname SGAsource@III.def1.1\endcsname{1.1}
\expandafter\def\csname SGAsource@III.def2.1\endcsname{2.1}
\expandafter\def\csname SGAsource@III.eq1.5.1\endcsname{1.5.1}
\expandafter\def\csname SGAsource@III.eq2.3.1\endcsname{2.3.1}
\expandafter\def\csname SGAsource@III.eq2.3.2\endcsname{2.3.2}
\expandafter\def\csname SGAsource@III.exem1.9.1\endcsname{1.9.1}
\expandafter\def\csname SGAsource@III.exem1.9.2\endcsname{1.9.2}
\expandafter\def\csname SGAsource@III.exem1.9.3\endcsname{1.9.3}
\expandafter\def\csname SGAsource@III.exem2.7\endcsname{2.7}
\expandafter\def\csname SGAsource@III.lem3.6.1\endcsname{3.6.1}
\expandafter\def\csname SGAsource@III.not1.8\endcsname{1.8}
\expandafter\def\csname SGAsource@III.prop1.2\endcsname{1.2}
\expandafter\def\csname SGAsource@III.prop1.3\endcsname{1.3}
\expandafter\def\csname SGAsource@III.prop1.5\endcsname{1.5}
\expandafter\def\csname SGAsource@III.prop1.6\endcsname{1.6}
\expandafter\def\csname SGAsource@III.prop1.7\endcsname{1.7}
\expandafter\def\csname SGAsource@III.prop2.2\endcsname{2.2}
\expandafter\def\csname SGAsource@III.prop2.3\endcsname{2.3}
\expandafter\def\csname SGAsource@III.prop2.5\endcsname{2.5}
\expandafter\def\csname SGAsource@III.prop2.6\endcsname{2.6}
\expandafter\def\csname SGAsource@III.prop3.2\endcsname{3.2}
\expandafter\def\csname SGAsource@III.prop3.5\endcsname{3.5}
\expandafter\def\csname SGAsource@III.prop3.6\endcsname{3.6}
\expandafter\def\csname SGAsource@III.prop3.7\endcsname{3.7}
\expandafter\def\csname SGAsource@III.prop5.2\endcsname{5.2}
\expandafter\def\csname SGAsource@III.prop5.4\endcsname{5.4}
\expandafter\def\csname SGAsource@III.prop5.5\endcsname{5.5}
\expandafter\def\csname SGAsource@III.rem1.4\endcsname{1.4}
\expandafter\def\csname SGAsource@III.sec1\endcsname{1}
\expandafter\def\csname SGAsource@III.sec2\endcsname{2}
\expandafter\def\csname SGAsource@III.sec3\endcsname{3}
\expandafter\def\csname SGAsource@III.sec4\endcsname{4}
\expandafter\def\csname SGAsource@III.sec5\endcsname{5}
\expandafter\def\csname SGAsource@III.subsec1.11\endcsname{1.11}
\expandafter\def\csname SGAsource@III.subsec2.4\endcsname{2.4}
\expandafter\def\csname SGAsource@III.subsec3.1\endcsname{3.1}
\expandafter\def\csname SGAsource@III.subsec4.2\endcsname{4.2}
\expandafter\def\csname SGAsource@III.subsec4.3\endcsname{4.3}
\expandafter\def\csname SGAsource@III.subsec4.4\endcsname{4.4}
\expandafter\def\csname SGAsource@III.subsec5.1\endcsname{5.1}
\expandafter\def\csname SGAsource@III.subsec5.3\endcsname{3}
\expandafter\def\csname SGAsource@III.subsubsec1.2.1\endcsname{1.2.1}
\expandafter\def\csname SGAsource@III.subsubsec4.1.1\endcsname{4.1.1}
\expandafter\def\csname SGAsource@III.subsubsec4.1.2\endcsname{4.1.2}
\expandafter\def\csname SGAsource@III.theo4.1\endcsname{4.1}
\expandafter\def\csname SGAsource@IV.cor1.2.1\endcsname{1.2.1}
\expandafter\def\csname SGAsource@IV.cor1.5\endcsname{1.5}
\expandafter\def\csname SGAsource@IV.cor1.6\endcsname{1.6}
\expandafter\def\csname SGAsource@IV.cor1.7\endcsname{1.7}
\expandafter\def\csname SGAsource@IV.cor10.4\endcsname{10.4}
\expandafter\def\csname SGAsource@IV.cor10.7\endcsname{10.7}
\expandafter\def\csname SGAsource@IV.cor10.8\endcsname{10.8}
\expandafter\def\csname SGAsource@IV.cor12.12\endcsname{12.12}
\expandafter\def\csname SGAsource@IV.cor12.13\endcsname{12.13}
\expandafter\def\csname SGAsource@IV.cor12.4\endcsname{12.4}
\expandafter\def\csname SGAsource@IV.cor12.5\endcsname{12.5}
\expandafter\def\csname SGAsource@IV.cor12.6\endcsname{12.6}
\expandafter\def\csname SGAsource@IV.cor13.5\endcsname{13.5}
\expandafter\def\csname SGAsource@IV.cor13.6\endcsname{13.6}
\expandafter\def\csname SGAsource@IV.cor4.2.4\endcsname{4.2.4}
\expandafter\def\csname SGAsource@IV.cor6.5.1\endcsname{6.5.1}
\expandafter\def\csname SGAsource@IV.cor6.7.3\endcsname{6.7.3}
\expandafter\def\csname SGAsource@IV.cor9.1.5\endcsname{9.1.5}
\expandafter\def\csname SGAsource@IV.cor9.4.3\endcsname{9.4.3}
\expandafter\def\csname SGAsource@IV.cor9.4.4\endcsname{9.4.4}
\expandafter\def\csname SGAsource@IV.cor9.4.5\endcsname{9.4.5}
\expandafter\def\csname SGAsource@IV.cor9.4.6\endcsname{9.4.6}
\expandafter\def\csname SGAsource@IV.cor9.7.3\endcsname{9.7.3}
\expandafter\def\csname SGAsource@IV.def1.1\endcsname{1.1}
\expandafter\def\csname SGAsource@IV.def13.1\endcsname{13.1}
\expandafter\def\csname SGAsource@IV.def13.3\endcsname{13.3}
\expandafter\def\csname SGAsource@IV.def3.1\endcsname{3.1}
\expandafter\def\csname SGAsource@IV.def6.1\endcsname{6.1}
\expandafter\def\csname SGAsource@IV.def6.2\endcsname{6.2}
\expandafter\def\csname SGAsource@IV.def6.4.1\endcsname{6.4.1}
\expandafter\def\csname SGAsource@IV.def8.3\endcsname{8.3}
\expandafter\def\csname SGAsource@IV.def9.1.1\endcsname{9.1.1}
\expandafter\def\csname SGAsource@IV.def9.5.5\endcsname{9.5.5}
\expandafter\def\csname SGAsource@IV.eq10.2.1\endcsname{10.2.1}
\expandafter\def\csname SGAsource@IV.eq10.3.1\endcsname{10.3.1}
\expandafter\def\csname SGAsource@IV.eq11.2.2.1\endcsname{11.2.2.1}
\expandafter\def\csname SGAsource@IV.eq12.1.1\endcsname{12.1.1}
\expandafter\def\csname SGAsource@IV.eq12.1.2\endcsname{12.1.2}
\expandafter\def\csname SGAsource@IV.eq12.14.1\endcsname{12.14.1}
\expandafter\def\csname SGAsource@IV.eq12.8.1\endcsname{12.8.1}
\expandafter\def\csname SGAsource@IV.eq13.4.3\endcsname{13.4.3}
\expandafter\def\csname SGAsource@IV.eq14.10.1\endcsname{14.10.1}
\expandafter\def\csname SGAsource@IV.eq14.2.1\endcsname{14.2.1}
\expandafter\def\csname SGAsource@IV.eq14.4.1\endcsname{14.4.1}
\expandafter\def\csname SGAsource@IV.eq14.4.2\endcsname{14.4.2}
\expandafter\def\csname SGAsource@IV.eq4.1.1.1\endcsname{4.1.1.1}
\expandafter\def\csname SGAsource@IV.eq4.10.1\endcsname{4.10.1}
\expandafter\def\csname SGAsource@IV.eq4.10.1.1\endcsname{4.10.1.1}
\expandafter\def\csname SGAsource@IV.eq4.10.2.1\endcsname{4.10.2.1}
\expandafter\def\csname SGAsource@IV.eq4.10.3.1\endcsname{4.10.3.1}
\expandafter\def\csname SGAsource@IV.eq4.2.1.1\endcsname{4.2.1.1}
\expandafter\def\csname SGAsource@IV.eq4.3.2\endcsname{4.3.2}
\expandafter\def\csname SGAsource@IV.eq4.3.3\endcsname{4.3.3}
\expandafter\def\csname SGAsource@IV.eq4.4.1\endcsname{4.4.1}
\expandafter\def\csname SGAsource@IV.eq4.4.2\endcsname{4.4.2}
\expandafter\def\csname SGAsource@IV.eq4.6.2.1\endcsname{4.6.2.1}
\expandafter\def\csname SGAsource@IV.eq4.6.2.2\endcsname{4.6.2.2}
\expandafter\def\csname SGAsource@IV.eq4.6.3.1\endcsname{4.6.3.1}
\expandafter\def\csname SGAsource@IV.eq4.6.3.2\endcsname{4.6.3.2}
\expandafter\def\csname SGAsource@IV.eq4.7.1\endcsname{4.7.1}
\expandafter\def\csname SGAsource@IV.eq4.7.2\endcsname{4.7.2}
\expandafter\def\csname SGAsource@IV.eq4.7.3\endcsname{4.7.3}
\expandafter\def\csname SGAsource@IV.eq4.9.1.1\endcsname{4.9.1.1}
\expandafter\def\csname SGAsource@IV.eq5.10.1\endcsname{5.10.1}
\expandafter\def\csname SGAsource@IV.eq5.10.2\endcsname{5.10.2}
\expandafter\def\csname SGAsource@IV.eq5.10.3\endcsname{5.10.3}
\expandafter\def\csname SGAsource@IV.eq5.11.1.1\endcsname{5.11.1.1}
\expandafter\def\csname SGAsource@IV.eq5.11.1.2\endcsname{5.11.1.2}
\expandafter\def\csname SGAsource@IV.eq5.12.1\endcsname{5.12.1}
\expandafter\def\csname SGAsource@IV.eq5.2.1\endcsname{5.2.1}
\expandafter\def\csname SGAsource@IV.eq5.2.2\endcsname{5.2.2}
\expandafter\def\csname SGAsource@IV.eq5.4.1\endcsname{5.4.1}
\expandafter\def\csname SGAsource@IV.eq5.5.1\endcsname{5.5.1}
\expandafter\def\csname SGAsource@IV.eq5.5.2\endcsname{5.5.2}
\expandafter\def\csname SGAsource@IV.eq5.5.3\endcsname{5.5.3}
\expandafter\def\csname SGAsource@IV.eq5.8.1\endcsname{5.8.1}
\expandafter\def\csname SGAsource@IV.eq5.8.2\endcsname{5.8.2}
\expandafter\def\csname SGAsource@IV.eq5.8.3.1\endcsname{5.8.3.1}
\expandafter\def\csname SGAsource@IV.eq5.8.3.2\endcsname{5.8.3.2}
\expandafter\def\csname SGAsource@IV.eq6.10.1\endcsname{6.10.1}
\expandafter\def\csname SGAsource@IV.eq6.2.1.1\endcsname{6.2.1.1}
\expandafter\def\csname SGAsource@IV.eq6.3.1\endcsname{6.3.1}
\expandafter\def\csname SGAsource@IV.eq6.3.2\endcsname{6.3.2}
\expandafter\def\csname SGAsource@IV.eq6.3.3\endcsname{6.3.3}
\expandafter\def\csname SGAsource@IV.eq6.7.1\endcsname{6.7.1}
\expandafter\def\csname SGAsource@IV.eq6.7.2\endcsname{6.7.2}
\expandafter\def\csname SGAsource@IV.eq6.8.1\endcsname{6.8.1}
\expandafter\def\csname SGAsource@IV.eq6.8.3\endcsname{6.8.3}
\expandafter\def\csname SGAsource@IV.eq6.8.4\endcsname{6.8.4}
\expandafter\def\csname SGAsource@IV.eq6.8.5.1\endcsname{6.8.5.1}
\expandafter\def\csname SGAsource@IV.eq6.8.6.1\endcsname{6.8.6.1}
\expandafter\def\csname SGAsource@IV.eq6.8.7.1\endcsname{6.8.7.1}
\expandafter\def\csname SGAsource@IV.eq7.1.1\endcsname{7.1.1}
\expandafter\def\csname SGAsource@IV.eq7.1.2\endcsname{7.1.2}
\expandafter\def\csname SGAsource@IV.eq7.1.3\endcsname{7.1.3}
\expandafter\def\csname SGAsource@IV.eq7.1.4\endcsname{7.1.4}
\expandafter\def\csname SGAsource@IV.eq7.2.1\endcsname{7.2.1}
\expandafter\def\csname SGAsource@IV.eq7.2.2\endcsname{7.2.2}
\expandafter\def\csname SGAsource@IV.eq7.2.3\endcsname{7.2.3}
\expandafter\def\csname SGAsource@IV.eq7.2.4.1\endcsname{7.2.4.1}
\expandafter\def\csname SGAsource@IV.eq7.2.6.1\endcsname{7.2.6.1}
\expandafter\def\csname SGAsource@IV.eq7.2.6.2\endcsname{7.2.6.2}
\expandafter\def\csname SGAsource@IV.eq8.7.1\endcsname{8.7.1}
\expandafter\def\csname SGAsource@IV.eq8.7.2\endcsname{8.7.2}
\expandafter\def\csname SGAsource@IV.eq8.7.3\endcsname{8.7.3}
\expandafter\def\csname SGAsource@IV.eq8.7.4\endcsname{8.7.4}
\expandafter\def\csname SGAsource@IV.eq8.7.5\endcsname{8.7.5}
\expandafter\def\csname SGAsource@IV.eq8.7.6\endcsname{8.7.6}
\expandafter\def\csname SGAsource@IV.eq9.1.12.2\endcsname{9.1.12.2}
\expandafter\def\csname SGAsource@IV.eq9.1.14.1\endcsname{9.1.14.1}
\expandafter\def\csname SGAsource@IV.eq9.1.4.1\endcsname{9.1.4.1}
\expandafter\def\csname SGAsource@IV.eq9.1.5.1\endcsname{9.1.5.1}
\expandafter\def\csname SGAsource@IV.eq9.1.6.1\endcsname{9.1.6.1}
\expandafter\def\csname SGAsource@IV.eq9.1.7.1\endcsname{9.1.7.1}
\expandafter\def\csname SGAsource@IV.eq9.1.7.3\endcsname{9.1.7.3}
\expandafter\def\csname SGAsource@IV.eq9.1.8.1\endcsname{9.1.8.1}
\expandafter\def\csname SGAsource@IV.eq9.1.8.2\endcsname{9.1.8.2}
\expandafter\def\csname SGAsource@IV.eq9.2.2\endcsname{9.2.2}
\expandafter\def\csname SGAsource@IV.eq9.2.6.1\endcsname{9.2.6.1}
\expandafter\def\csname SGAsource@IV.eq9.3.2.1\endcsname{9.3.2.1}
\expandafter\def\csname SGAsource@IV.eq9.3.2.2\endcsname{9.3.2.2}
\expandafter\def\csname SGAsource@IV.eq9.3.2.3\endcsname{9.3.2.3}
\expandafter\def\csname SGAsource@IV.eq9.3.6.1\endcsname{9.3.6.1}
\expandafter\def\csname SGAsource@IV.eq9.3.6.2\endcsname{9.3.6.2}
\expandafter\def\csname SGAsource@IV.eq9.3.6.3\endcsname{9.3.6.3}
\expandafter\def\csname SGAsource@IV.eq9.4.3.1\endcsname{9.4.3.1}
\expandafter\def\csname SGAsource@IV.eq9.5.10.1\endcsname{9.5.10.1}
\expandafter\def\csname SGAsource@IV.eq9.5.10.2\endcsname{9.5.10.2}
\expandafter\def\csname SGAsource@IV.eq9.5.10.3\endcsname{9.5.10.3}
\expandafter\def\csname SGAsource@IV.eq9.5.3.1\endcsname{9.5.3.1}
\expandafter\def\csname SGAsource@IV.eq9.6.1\endcsname{9.6.1}
\expandafter\def\csname SGAsource@IV.eq9.6.2\endcsname{9.6.2}
\expandafter\def\csname SGAsource@IV.eq9.6.3\endcsname{9.6.3}
\expandafter\def\csname SGAsource@IV.eq9.6.4\endcsname{9.6.4}
\expandafter\def\csname SGAsource@IV.eq9.7.2.1\endcsname{9.7.2.1}
\expandafter\def\csname SGAsource@IV.eq9.7.2.2\endcsname{9.7.2.2}
\expandafter\def\csname SGAsource@IV.eq9.7.2.3\endcsname{9.7.2.3}
\expandafter\def\csname SGAsource@IV.eq9.7.2.4.1\endcsname{9.7.2.4.1}
\expandafter\def\csname SGAsource@IV.eq9.8.2.1\endcsname{9.8.2.1}
\expandafter\def\csname SGAsource@IV.exer11.3.5\endcsname{11.3.5}
\expandafter\def\csname SGAsource@IV.exer13.9\endcsname{13.9}
\expandafter\def\csname SGAsource@IV.exer2.7.5\endcsname{2.7.5}
\expandafter\def\csname SGAsource@IV.exer4.10.6\endcsname{4.10.6}
\expandafter\def\csname SGAsource@IV.exer5.14\endcsname{5.14}
\expandafter\def\csname SGAsource@IV.exer5.9\endcsname{5.9}
\expandafter\def\csname SGAsource@IV.exer6.10\endcsname{6.10}
\expandafter\def\csname SGAsource@IV.exer6.5.2\endcsname{6.5.2}
\expandafter\def\csname SGAsource@IV.exer6.9\endcsname{6.9}
\expandafter\def\csname SGAsource@IV.exer7.1.10\endcsname{7.1.10}
\expandafter\def\csname SGAsource@IV.exer7.1.11\endcsname{7.1.11}
\expandafter\def\csname SGAsource@IV.exer7.1.9\endcsname{7.1.9}
\expandafter\def\csname SGAsource@IV.exer7.2.6\endcsname{7.2.6}
\expandafter\def\csname SGAsource@IV.exer7.5\endcsname{7.5}
\expandafter\def\csname SGAsource@IV.exer7.6\endcsname{7.6}
\expandafter\def\csname SGAsource@IV.exer7.7\endcsname{7.7}
\expandafter\def\csname SGAsource@IV.exer7.8\endcsname{7.8}
\expandafter\def\csname SGAsource@IV.exer8.6\endcsname{8.6}
\expandafter\def\csname SGAsource@IV.exer8.7\endcsname{8.7}
\expandafter\def\csname SGAsource@IV.exer8.8\endcsname{8.8}
\expandafter\def\csname SGAsource@IV.exer9.1.10\endcsname{9.1.10}
\expandafter\def\csname SGAsource@IV.exer9.1.11\endcsname{9.1.11}
\expandafter\def\csname SGAsource@IV.exer9.1.12\endcsname{9.1.12}
\expandafter\def\csname SGAsource@IV.exer9.1.13\endcsname{9.1.13}
\expandafter\def\csname SGAsource@IV.exer9.1.6\endcsname{9.1.6}
\expandafter\def\csname SGAsource@IV.exer9.1.7\endcsname{9.1.7}
\expandafter\def\csname SGAsource@IV.exer9.1.7.2\endcsname{9.1.7.2}
\expandafter\def\csname SGAsource@IV.exer9.1.8\endcsname{9.1.8}
\expandafter\def\csname SGAsource@IV.exer9.1.9\endcsname{9.1.9}
\expandafter\def\csname SGAsource@IV.exer9.2.7\endcsname{9.2.7}
\expandafter\def\csname SGAsource@IV.exer9.4.10\endcsname{9.4.10}
\expandafter\def\csname SGAsource@IV.exer9.4.7\endcsname{9.4.7}
\expandafter\def\csname SGAsource@IV.exer9.4.8\endcsname{9.4.8}
\expandafter\def\csname SGAsource@IV.exer9.4.9\endcsname{9.4.9}
\expandafter\def\csname SGAsource@IV.exer9.5.10\endcsname{9.5.10}
\expandafter\def\csname SGAsource@IV.exer9.5.11\endcsname{9.5.11}
\expandafter\def\csname SGAsource@IV.exer9.5.12\endcsname{9.5.12}
\expandafter\def\csname SGAsource@IV.exer9.5.9\endcsname{9.5.9}
\expandafter\def\csname SGAsource@IV.exer9.6.5\endcsname{9.6.5}
\expandafter\def\csname SGAsource@IV.exer9.7.5\endcsname{9.7.5}
\expandafter\def\csname SGAsource@IV.exer9.8.1\endcsname{9.8.1}
\expandafter\def\csname SGAsource@IV.exer9.8.2\endcsname{9.8.2}
\expandafter\def\csname SGAsource@IV.exer9.8.3\endcsname{9.8.3}
\expandafter\def\csname SGAsource@IV.lem10.4.2\endcsname{10.4.2}
\expandafter\def\csname SGAsource@IV.lem5.11.2\endcsname{5.11.2}
\expandafter\def\csname SGAsource@IV.lem9.7.2.4\endcsname{9.7.2.4}
\expandafter\def\csname SGAsource@IV.par1.2.4.1\endcsname{1.2.4.1}
\expandafter\def\csname SGAsource@IV.par1.2.4.2\endcsname{1.2.4.2}
\expandafter\def\csname SGAsource@IV.par1.2.4.3\endcsname{1.2.4.3}
\expandafter\def\csname SGAsource@IV.par1.2.4.4\endcsname{1.2.4.4}
\expandafter\def\csname SGAsource@IV.par9.5.4.8\endcsname{9.5.4.8}
\expandafter\def\csname SGAsource@IV.par9.5.8.1\endcsname{9.5.8.1}
\expandafter\def\csname SGAsource@IV.par9.5.8.2\endcsname{9.5.8.2}
\expandafter\def\csname SGAsource@IV.par9.5.8.3\endcsname{9.5.8.3}
\expandafter\def\csname SGAsource@IV.par9.5.8.4\endcsname{9.5.8.4}
\expandafter\def\csname SGAsource@IV.par9.5.8.5\endcsname{9.5.8.5}
\expandafter\def\csname SGAsource@IV.par9.5.8.6\endcsname{9.5.8.6}
\expandafter\def\csname SGAsource@IV.par9.5.8.7\endcsname{9.5.8.7}
\expandafter\def\csname SGAsource@IV.par9.5.8.8\endcsname{9.5.8.8}
\expandafter\def\csname SGAsource@IV.prop1.4\endcsname{1.4}
\expandafter\def\csname SGAsource@IV.prop10.1\endcsname{10.1}
\expandafter\def\csname SGAsource@IV.prop10.3\endcsname{10.3}
\expandafter\def\csname SGAsource@IV.prop10.6\endcsname{10.6}
\expandafter\def\csname SGAsource@IV.prop11.1.6\endcsname{11.1.6}
\expandafter\def\csname SGAsource@IV.prop11.3.1\endcsname{11.3.1}
\expandafter\def\csname SGAsource@IV.prop11.3.3\endcsname{11.3.3}
\expandafter\def\csname SGAsource@IV.prop12.1\endcsname{12.1}
\expandafter\def\csname SGAsource@IV.prop12.10\endcsname{12.10}
\expandafter\def\csname SGAsource@IV.prop12.11\endcsname{12.11}
\expandafter\def\csname SGAsource@IV.prop12.14\endcsname{12.14}
\expandafter\def\csname SGAsource@IV.prop12.3\endcsname{12.3}
\expandafter\def\csname SGAsource@IV.prop12.7\endcsname{12.7}
\expandafter\def\csname SGAsource@IV.prop13.4\endcsname{13.4}
\expandafter\def\csname SGAsource@IV.prop14.10\endcsname{14.10}
\expandafter\def\csname SGAsource@IV.prop14.5\endcsname{14.5}
\expandafter\def\csname SGAsource@IV.prop14.6\endcsname{14.6}
\expandafter\def\csname SGAsource@IV.prop14.8\endcsname{14.8}
\expandafter\def\csname SGAsource@IV.prop4.9.4\endcsname{4.9.4}
\expandafter\def\csname SGAsource@IV.prop5.11\endcsname{5.11}
\expandafter\def\csname SGAsource@IV.prop5.12\endcsname{5.12}
\expandafter\def\csname SGAsource@IV.prop6.5\endcsname{6.5}
\expandafter\def\csname SGAsource@IV.prop6.5.3\endcsname{6.5.3}
\expandafter\def\csname SGAsource@IV.prop9.1.3\endcsname{9.1.3}
\expandafter\def\csname SGAsource@IV.prop9.1.4\endcsname{9.1.4}
\expandafter\def\csname SGAsource@IV.prop9.2.4\endcsname{9.2.4}
\expandafter\def\csname SGAsource@IV.prop9.3.3\endcsname{9.3.3}
\expandafter\def\csname SGAsource@IV.prop9.3.4\endcsname{9.3.4}
\expandafter\def\csname SGAsource@IV.prop9.4.1\endcsname{9.4.1}
\expandafter\def\csname SGAsource@IV.prop9.4.2\endcsname{9.4.2}
\expandafter\def\csname SGAsource@IV.prop9.7.2\endcsname{9.7.2}
\expandafter\def\csname SGAsource@IV.rem1.3\endcsname{1.3}
\expandafter\def\csname SGAsource@IV.rem1.8\endcsname{1.8}
\expandafter\def\csname SGAsource@IV.rem11.3.2\endcsname{11.3.2}
\expandafter\def\csname SGAsource@IV.rem11.3.4\endcsname{11.3.4}
\expandafter\def\csname SGAsource@IV.rem4.10.5\endcsname{4.10.5}
\expandafter\def\csname SGAsource@IV.rem4.6.4\endcsname{4.6.4}
\expandafter\def\csname SGAsource@IV.rem4.7.4\endcsname{4.7.4}
\expandafter\def\csname SGAsource@IV.rem5.13\endcsname{5.13}
\expandafter\def\csname SGAsource@IV.rem6.7.4\endcsname{6.7.4}
\expandafter\def\csname SGAsource@IV.rem7.2.4\endcsname{7.2.4}
\expandafter\def\csname SGAsource@IV.rem8.1.1\endcsname{8.1.1}
\expandafter\def\csname SGAsource@IV.rem8.2.1\endcsname{8.2.1}
\expandafter\def\csname SGAsource@IV.rem8.5.3\endcsname{8.5.3}
\expandafter\def\csname SGAsource@IV.rem9.1.2\endcsname{9.1.2}
\expandafter\def\csname SGAsource@IV.rem9.2.3\endcsname{9.2.3}
\expandafter\def\csname SGAsource@IV.rem9.2.4.1\endcsname{9.2.4.1}
\expandafter\def\csname SGAsource@IV.rem9.7.4\endcsname{9.7.4}
\expandafter\def\csname SGAsource@IV.sec0\endcsname{0}
\expandafter\def\csname SGAsource@IV.sec1\endcsname{1}
\expandafter\def\csname SGAsource@IV.sec10\endcsname{10}
\expandafter\def\csname SGAsource@IV.sec11\endcsname{11}
\expandafter\def\csname SGAsource@IV.sec12\endcsname{12}
\expandafter\def\csname SGAsource@IV.sec13\endcsname{13}
\expandafter\def\csname SGAsource@IV.sec14\endcsname{14}
\expandafter\def\csname SGAsource@IV.sec2\endcsname{2}
\expandafter\def\csname SGAsource@IV.sec3\endcsname{3}
\expandafter\def\csname SGAsource@IV.sec4\endcsname{4}
\expandafter\def\csname SGAsource@IV.sec5\endcsname{5}
\expandafter\def\csname SGAsource@IV.sec6\endcsname{6}
\expandafter\def\csname SGAsource@IV.sec7\endcsname{7}
\expandafter\def\csname SGAsource@IV.sec8\endcsname{8}
\expandafter\def\csname SGAsource@IV.sec9\endcsname{9}
\expandafter\def\csname SGAsource@IV.subsec0.1\endcsname{0.1}
\expandafter\def\csname SGAsource@IV.subsec0.2\endcsname{0.2}
\expandafter\def\csname SGAsource@IV.subsec0.3\endcsname{0.3}
\expandafter\def\csname SGAsource@IV.subsec0.4\endcsname{0.4}
\expandafter\def\csname SGAsource@IV.subsec0.5\endcsname{0.5}
\expandafter\def\csname SGAsource@IV.subsec10.2\endcsname{10.2}
\expandafter\def\csname SGAsource@IV.subsec10.5\endcsname{10.5}
\expandafter\def\csname SGAsource@IV.subsec12.2\endcsname{12.2}
\expandafter\def\csname SGAsource@IV.subsec12.8\endcsname{12.8}
\expandafter\def\csname SGAsource@IV.subsec12.9\endcsname{12.9}
\expandafter\def\csname SGAsource@IV.subsec13.2\endcsname{13.2}
\expandafter\def\csname SGAsource@IV.subsec13.7\endcsname{13.7}
\expandafter\def\csname SGAsource@IV.subsec13.8\endcsname{13.8}
\expandafter\def\csname SGAsource@IV.subsec14.1\endcsname{14.1}
\expandafter\def\csname SGAsource@IV.subsec14.2\endcsname{14.2}
\expandafter\def\csname SGAsource@IV.subsec14.3\endcsname{14.3}
\expandafter\def\csname SGAsource@IV.subsec14.4\endcsname{14.4}
\expandafter\def\csname SGAsource@IV.subsec14.7\endcsname{14.7}
\expandafter\def\csname SGAsource@IV.subsec14.9\endcsname{14.9}
\expandafter\def\csname SGAsource@IV.subsec2.0\endcsname{2.0}
\expandafter\def\csname SGAsource@IV.subsec2.1\endcsname{2.1}
\expandafter\def\csname SGAsource@IV.subsec2.2\endcsname{2.2}
\expandafter\def\csname SGAsource@IV.subsec2.3\endcsname{2.3}
\expandafter\def\csname SGAsource@IV.subsec2.4\endcsname{2.4}
\expandafter\def\csname SGAsource@IV.subsec2.5\endcsname{2.5}
\expandafter\def\csname SGAsource@IV.subsec2.6\endcsname{2.6}
\expandafter\def\csname SGAsource@IV.subsec2.7\endcsname{2.7}
\expandafter\def\csname SGAsource@IV.subsec2.8\endcsname{2.8}
\expandafter\def\csname SGAsource@IV.subsec3.2\endcsname{3.2}
\expandafter\def\csname SGAsource@IV.subsec3.4\endcsname{3.4}
\expandafter\def\csname SGAsource@IV.subsec4.1\endcsname{4.1}
\expandafter\def\csname SGAsource@IV.subsec4.10\endcsname{4.10}
\expandafter\def\csname SGAsource@IV.subsec4.2\endcsname{4.2}
\expandafter\def\csname SGAsource@IV.subsec4.3\endcsname{4.3}
\expandafter\def\csname SGAsource@IV.subsec4.4\endcsname{4.4}
\expandafter\def\csname SGAsource@IV.subsec4.5\endcsname{4.5}
\expandafter\def\csname SGAsource@IV.subsec4.6\endcsname{4.6}
\expandafter\def\csname SGAsource@IV.subsec4.7\endcsname{4.7}
\expandafter\def\csname SGAsource@IV.subsec4.8\endcsname{4.8}
\expandafter\def\csname SGAsource@IV.subsec4.9\endcsname{4.9}
\expandafter\def\csname SGAsource@IV.subsec5.1\endcsname{5.1}
\expandafter\def\csname SGAsource@IV.subsec5.10\endcsname{5.10}
\expandafter\def\csname SGAsource@IV.subsec5.2\endcsname{5.2}
\expandafter\def\csname SGAsource@IV.subsec5.3\endcsname{5.3}
\expandafter\def\csname SGAsource@IV.subsec5.4\endcsname{5.4}
\expandafter\def\csname SGAsource@IV.subsec5.5\endcsname{5.5}
\expandafter\def\csname SGAsource@IV.subsec5.6\endcsname{5.6}
\expandafter\def\csname SGAsource@IV.subsec5.7\endcsname{5.7}
\expandafter\def\csname SGAsource@IV.subsec5.8\endcsname{5.8}
\expandafter\def\csname SGAsource@IV.subsec6.0\endcsname{6.0}
\expandafter\def\csname SGAsource@IV.subsec6.3\endcsname{6.3}
\expandafter\def\csname SGAsource@IV.subsec6.6\endcsname{6.6}
\expandafter\def\csname SGAsource@IV.subsec6.7\endcsname{6.7}
\expandafter\def\csname SGAsource@IV.subsec6.8\endcsname{6.8}
\expandafter\def\csname SGAsource@IV.subsec7.1\endcsname{7.1}
\expandafter\def\csname SGAsource@IV.subsec7.2\endcsname{7.2}
\expandafter\def\csname SGAsource@IV.subsec7.3\endcsname{7.3}
\expandafter\def\csname SGAsource@IV.subsec7.4\endcsname{7.4}
\expandafter\def\csname SGAsource@IV.subsec8.1\endcsname{8.1}
\expandafter\def\csname SGAsource@IV.subsec8.2\endcsname{8.2}
\expandafter\def\csname SGAsource@IV.subsec8.4\endcsname{8.4}
\expandafter\def\csname SGAsource@IV.subsec8.5\endcsname{8.5}
\expandafter\def\csname SGAsource@IV.subsec9.1\endcsname{9.1}
\expandafter\def\csname SGAsource@IV.subsec9.2\endcsname{9.2}
\expandafter\def\csname SGAsource@IV.subsec9.3\endcsname{9.3}
\expandafter\def\csname SGAsource@IV.subsec9.4\endcsname{9.4}
\expandafter\def\csname SGAsource@IV.subsec9.5\endcsname{9.5}
\expandafter\def\csname SGAsource@IV.subsec9.6\endcsname{9.6}
\expandafter\def\csname SGAsource@IV.subsec9.7\endcsname{9.7}
\expandafter\def\csname SGAsource@IV.subsec9.7.1\endcsname{9.7.1}
\expandafter\def\csname SGAsource@IV.subsec9.8\endcsname{9.8}
\expandafter\def\csname SGAsource@IV.subsubsec1.1.1\endcsname{1.1.1}
\expandafter\def\csname SGAsource@IV.subsubsec1.1.2\endcsname{1.1.2}
\expandafter\def\csname SGAsource@IV.subsubsec1.2.3\endcsname{1.2.3}
\expandafter\def\csname SGAsource@IV.subsubsec1.2.4\endcsname{1.2.4}
\expandafter\def\csname SGAsource@IV.subsubsec10.4.1\endcsname{10.4.1}
\expandafter\def\csname SGAsource@IV.subsubsec10.4.3\endcsname{10.4.3}
\expandafter\def\csname SGAsource@IV.subsubsec10.5.1\endcsname{10.5.1}
\expandafter\def\csname SGAsource@IV.subsubsec10.5.2\endcsname{10.5.2}
\expandafter\def\csname SGAsource@IV.subsubsec10.5.3\endcsname{10.5.3}
\expandafter\def\csname SGAsource@IV.subsubsec10.5.4\endcsname{10.5.4}
\expandafter\def\csname SGAsource@IV.subsubsec10.6.1\endcsname{10.6.1}
\expandafter\def\csname SGAsource@IV.subsubsec11.1.1\endcsname{11.1.1}
\expandafter\def\csname SGAsource@IV.subsubsec11.1.2\endcsname{11.1.2}
\expandafter\def\csname SGAsource@IV.subsubsec11.1.3\endcsname{11.1.3}
\expandafter\def\csname SGAsource@IV.subsubsec11.1.4\endcsname{11.1.4}
\expandafter\def\csname SGAsource@IV.subsubsec11.1.5\endcsname{11.1.5}
\expandafter\def\csname SGAsource@IV.subsubsec11.2.1\endcsname{11.2.1}
\expandafter\def\csname SGAsource@IV.subsubsec11.2.2\endcsname{11.2.2}
\expandafter\def\csname SGAsource@IV.subsubsec11.2.3\endcsname{11.2.3}
\expandafter\def\csname SGAsource@IV.subsubsec13.1.1\endcsname{13.1.1}
\expandafter\def\csname SGAsource@IV.subsubsec13.2.1\endcsname{13.2.1}
\expandafter\def\csname SGAsource@IV.subsubsec13.2.2\endcsname{13.2.2}
\expandafter\def\csname SGAsource@IV.subsubsec13.2.3\endcsname{13.2.3}
\expandafter\def\csname SGAsource@IV.subsubsec13.3.1\endcsname{13.3.1}
\expandafter\def\csname SGAsource@IV.subsubsec2.7.1\endcsname{2.7.1}
\expandafter\def\csname SGAsource@IV.subsubsec2.7.2\endcsname{2.7.2}
\expandafter\def\csname SGAsource@IV.subsubsec2.7.3\endcsname{2.7.3}
\expandafter\def\csname SGAsource@IV.subsubsec2.7.4\endcsname{2.7.4}
\expandafter\def\csname SGAsource@IV.subsubsec3.1.1\endcsname{3.1.1}
\expandafter\def\csname SGAsource@IV.subsubsec3.1.2\endcsname{3.1.2}
\expandafter\def\csname SGAsource@IV.subsubsec3.1.3\endcsname{3.1.3}
\expandafter\def\csname SGAsource@IV.subsubsec3.2.2\endcsname{3.2.2}
\expandafter\def\csname SGAsource@IV.subsubsec3.3.1\endcsname{3.3.1}
\expandafter\def\csname SGAsource@IV.subsubsec3.3.2\endcsname{3.3.2}
\expandafter\def\csname SGAsource@IV.subsubsec3.4.1\endcsname{3.4.1}
\expandafter\def\csname SGAsource@IV.subsubsec4.1.1\endcsname{4.1.1}
\expandafter\def\csname SGAsource@IV.subsubsec4.1.2\endcsname{4.1.2}
\expandafter\def\csname SGAsource@IV.subsubsec4.1.3\endcsname{4.1.3}
\expandafter\def\csname SGAsource@IV.subsubsec4.10.2\endcsname{4.10.2}
\expandafter\def\csname SGAsource@IV.subsubsec4.10.3\endcsname{4.10.3}
\expandafter\def\csname SGAsource@IV.subsubsec4.10.4\endcsname{4.10.4}
\expandafter\def\csname SGAsource@IV.subsubsec4.2.1\endcsname{4.2.1}
\expandafter\def\csname SGAsource@IV.subsubsec4.2.2\endcsname{4.2.2}
\expandafter\def\csname SGAsource@IV.subsubsec4.2.3\endcsname{4.2.3}
\expandafter\def\csname SGAsource@IV.subsubsec4.3.4\endcsname{4.3.4}
\expandafter\def\csname SGAsource@IV.subsubsec4.3.5\endcsname{4.3.5}
\expandafter\def\csname SGAsource@IV.subsubsec4.3.6\endcsname{4.3.6}
\expandafter\def\csname SGAsource@IV.subsubsec4.5.1\endcsname{4.5.1}
\expandafter\def\csname SGAsource@IV.subsubsec4.5.2\endcsname{4.5.2}
\expandafter\def\csname SGAsource@IV.subsubsec4.6.1\endcsname{4.6.1}
\expandafter\def\csname SGAsource@IV.subsubsec4.6.2\endcsname{4.6.2}
\expandafter\def\csname SGAsource@IV.subsubsec4.6.3\endcsname{4.6.3}
\expandafter\def\csname SGAsource@IV.subsubsec4.7.5\endcsname{4.7.5}
\expandafter\def\csname SGAsource@IV.subsubsec4.9.1\endcsname{4.9.1}
\expandafter\def\csname SGAsource@IV.subsubsec4.9.2\endcsname{4.9.2}
\expandafter\def\csname SGAsource@IV.subsubsec4.9.3\endcsname{4.9.3}
\expandafter\def\csname SGAsource@IV.subsubsec4.9.5\endcsname{4.9.5}
\expandafter\def\csname SGAsource@IV.subsubsec4.9.6\endcsname{4.9.6}
\expandafter\def\csname SGAsource@IV.subsubsec5.10.4\endcsname{5.10.4}
\expandafter\def\csname SGAsource@IV.subsubsec5.11.1\endcsname{5.11.1}
\expandafter\def\csname SGAsource@IV.subsubsec5.12.2\endcsname{5.12.2}
\expandafter\def\csname SGAsource@IV.subsubsec5.8.3\endcsname{5.8.3}
\expandafter\def\csname SGAsource@IV.subsubsec5.8.4\endcsname{5.8.4}
\expandafter\def\csname SGAsource@IV.subsubsec6.1.1\endcsname{6.1.1}
\expandafter\def\csname SGAsource@IV.subsubsec6.1.2\endcsname{6.1.2}
\expandafter\def\csname SGAsource@IV.subsubsec6.2.1\endcsname{6.2.1}
\expandafter\def\csname SGAsource@IV.subsubsec6.4.0\endcsname{6.4.0}
\expandafter\def\csname SGAsource@IV.subsubsec6.4.2\endcsname{6.4.2}
\expandafter\def\csname SGAsource@IV.subsubsec6.4.3\endcsname{6.4.3}
\expandafter\def\csname SGAsource@IV.subsubsec6.8.2\endcsname{6.8.2}
\expandafter\def\csname SGAsource@IV.subsubsec6.8.5\endcsname{6.8.5}
\expandafter\def\csname SGAsource@IV.subsubsec6.8.6\endcsname{6.8.6}
\expandafter\def\csname SGAsource@IV.subsubsec6.8.7\endcsname{6.8.7}
\expandafter\def\csname SGAsource@IV.subsubsec7.1.5\endcsname{7.1.5}
\expandafter\def\csname SGAsource@IV.subsubsec7.1.6\endcsname{7.1.6}
\expandafter\def\csname SGAsource@IV.subsubsec7.1.7\endcsname{7.1.7}
\expandafter\def\csname SGAsource@IV.subsubsec7.1.8\endcsname{7.1.8}
\expandafter\def\csname SGAsource@IV.subsubsec7.2.5\endcsname{7.2.5}
\expandafter\def\csname SGAsource@IV.subsubsec8.4.1\endcsname{8.4.1}
\expandafter\def\csname SGAsource@IV.subsubsec8.4.2\endcsname{8.4.2}
\expandafter\def\csname SGAsource@IV.subsubsec8.4.3\endcsname{8.4.3}
\expandafter\def\csname SGAsource@IV.subsubsec8.4.4\endcsname{8.4.4}
\expandafter\def\csname SGAsource@IV.subsubsec8.4.5\endcsname{8.4.5}
\expandafter\def\csname SGAsource@IV.subsubsec8.5.1\endcsname{8.5.1}
\expandafter\def\csname SGAsource@IV.subsubsec8.5.2\endcsname{8.5.2}
\expandafter\def\csname SGAsource@IV.subsubsec8.5.4\endcsname{8.5.4}
\expandafter\def\csname SGAsource@IV.subsubsec8.5.5\endcsname{8.5.5}
\expandafter\def\csname SGAsource@IV.subsubsec9.1.14\endcsname{9.1.14}
\expandafter\def\csname SGAsource@IV.subsubsec9.2.1\endcsname{9.2.1}
\expandafter\def\csname SGAsource@IV.subsubsec9.2.5\endcsname{9.2.5}
\expandafter\def\csname SGAsource@IV.subsubsec9.2.6\endcsname{9.2.6}
\expandafter\def\csname SGAsource@IV.subsubsec9.3.1\endcsname{9.3.1}
\expandafter\def\csname SGAsource@IV.subsubsec9.3.2\endcsname{9.3.2}
\expandafter\def\csname SGAsource@IV.subsubsec9.3.5\endcsname{9.3.5}
\expandafter\def\csname SGAsource@IV.subsubsec9.3.6\endcsname{9.3.6}
\expandafter\def\csname SGAsource@IV.subsubsec9.5.1\endcsname{9.5.1}
\expandafter\def\csname SGAsource@IV.subsubsec9.5.2\endcsname{9.5.2}
\expandafter\def\csname SGAsource@IV.subsubsec9.5.3\endcsname{9.5.3}
\expandafter\def\csname SGAsource@IV.subsubsec9.5.6\endcsname{9.5.6}
\expandafter\def\csname SGAsource@IV.theo1.2\endcsname{1.2}
\expandafter\def\csname SGAsource@IV.theo9.5.4\endcsname{9.5.4}
\expandafter\def\csname SGAsource@IX.Thr3.2\endcsname{3.2}
\expandafter\def\csname SGAsource@IX.Thr3.5\endcsname{3.5}
\expandafter\def\csname SGAsource@IX.cor2.7.1\endcsname{2.7.1}
\expandafter\def\csname SGAsource@IX.cor2.7.2\endcsname{2.7.2}
\expandafter\def\csname SGAsource@IX.cor2.7.3\endcsname{2.7.3}
\expandafter\def\csname SGAsource@IX.cor2.7.4\endcsname{2.7.4}
\expandafter\def\csname SGAsource@IX.cor3.4\endcsname{3.4}
\expandafter\def\csname SGAsource@IX.cor4.2\endcsname{4.2}
\expandafter\def\csname SGAsource@IX.cor4.7\endcsname{4.7}
\expandafter\def\csname SGAsource@IX.cor5.2\endcsname{5.2}
\expandafter\def\csname SGAsource@IX.cor5.7\endcsname{5.7}
\expandafter\def\csname SGAsource@IX.def1.1\endcsname{1.1}
\expandafter\def\csname SGAsource@IX.def1.3\endcsname{1.3}
\expandafter\def\csname SGAsource@IX.def1.5\endcsname{1.5}
\expandafter\def\csname SGAsource@IX.def2.12\endcsname{2.12}
\expandafter\def\csname SGAsource@IX.def2.3\endcsname{2.3}
\expandafter\def\csname SGAsource@IX.eq3.1\endcsname{3.1}
\expandafter\def\csname SGAsource@IX.eq4.3.1\endcsname{4.3.1}
\expandafter\def\csname SGAsource@IX.eq4.4.1\endcsname{4.4.1}
\expandafter\def\csname SGAsource@IX.eq4.4.2\endcsname{4.4.2}
\expandafter\def\csname SGAsource@IX.eq4.5\endcsname{4.5}
\expandafter\def\csname SGAsource@IX.eq4.8\endcsname{4.8}
\expandafter\def\csname SGAsource@IX.eq5.1.1\endcsname{5.1.1}
\expandafter\def\csname SGAsource@IX.eq5.1.2\endcsname{5.1.2}
\expandafter\def\csname SGAsource@IX.eq5.1.3\endcsname{5.1.3}
\expandafter\def\csname SGAsource@IX.eq5.1.4\endcsname{5.1.4}
\expandafter\def\csname SGAsource@IX.eq5.4\endcsname{5.4}
\expandafter\def\csname SGAsource@IX.lem2.1\endcsname{2.1}
\expandafter\def\csname SGAsource@IX.lem2.10\endcsname{2.10}
\expandafter\def\csname SGAsource@IX.lem2.14.1\endcsname{2.14.1}
\expandafter\def\csname SGAsource@IX.lem2.2\endcsname{2.2}
\expandafter\def\csname SGAsource@IX.lem2.8.1\endcsname{2.8.1}
\expandafter\def\csname SGAsource@IX.lem4.1\endcsname{4.1}
\expandafter\def\csname SGAsource@IX.par2.7.2.1\endcsname{2.7.2.1}
\expandafter\def\csname SGAsource@IX.prop1.2\endcsname{1.2}
\expandafter\def\csname SGAsource@IX.prop1.4\endcsname{1.4}
\expandafter\def\csname SGAsource@IX.prop1.6\endcsname{1.6}
\expandafter\def\csname SGAsource@IX.prop2.11\endcsname{2.11}
\expandafter\def\csname SGAsource@IX.prop2.13\endcsname{2.13}
\expandafter\def\csname SGAsource@IX.prop2.14\endcsname{2.14}
\expandafter\def\csname SGAsource@IX.prop2.4\endcsname{2.4}
\expandafter\def\csname SGAsource@IX.prop2.5\endcsname{2.5}
\expandafter\def\csname SGAsource@IX.prop2.6\endcsname{2.6}
\expandafter\def\csname SGAsource@IX.prop2.7\endcsname{2.7}
\expandafter\def\csname SGAsource@IX.prop2.8\endcsname{2.8}
\expandafter\def\csname SGAsource@IX.prop2.9\endcsname{2.9}
\expandafter\def\csname SGAsource@IX.prop3.6\endcsname{3.6}
\expandafter\def\csname SGAsource@IX.prop5.5\endcsname{5.5}
\expandafter\def\csname SGAsource@IX.prop5.6\endcsname{5.6}
\expandafter\def\csname SGAsource@IX.rem5.8\endcsname{5.8}
\expandafter\def\csname SGAsource@IX.rems2.14.2\endcsname{2.14.2}
\expandafter\def\csname SGAsource@IX.rems3.7\endcsname{3.7}
\expandafter\def\csname SGAsource@IX.sec0\endcsname{0}
\expandafter\def\csname SGAsource@IX.sec1\endcsname{1}
\expandafter\def\csname SGAsource@IX.sec2\endcsname{2}
\expandafter\def\csname SGAsource@IX.sec3\endcsname{3}
\expandafter\def\csname SGAsource@IX.sec4\endcsname{4}
\expandafter\def\csname SGAsource@IX.sec5\endcsname{5}
\expandafter\def\csname SGAsource@IX.subsec2.0\endcsname{2.0}
\expandafter\def\csname SGAsource@IX.subsec3.0\endcsname{3.0}
\expandafter\def\csname SGAsource@IX.subsec4.0\endcsname{4.0}
\expandafter\def\csname SGAsource@IX.subsec4.3\endcsname{4.3}
\expandafter\def\csname SGAsource@IX.subsec4.4\endcsname{4.4}
\expandafter\def\csname SGAsource@IX.subsec5.1\endcsname{5.1}
\expandafter\def\csname SGAsource@IX.subsubsec2.3.1\endcsname{2.3.1}
\expandafter\def\csname SGAsource@IX.subsubsec3.3.1\endcsname{3.3.1}
\expandafter\def\csname SGAsource@IX.subsubsec4.1.1\endcsname{4.1.1}
\expandafter\def\csname SGAsource@IX.subsubsec5.2.1\endcsname{5.2.1}
\expandafter\def\csname SGAsource@IX.thm3.3\endcsname{3.3}
\expandafter\def\csname SGAsource@IX.thm4.6\endcsname{4.6}
\expandafter\def\csname SGAsource@Ii.prop2.5.3\endcsname{2.5.3}
\expandafter\def\csname SGAsource@V.1.10.1\endcsname{1.10.1}
\expandafter\def\csname SGAsource@V.1.10.2\endcsname{1.10.2}
\expandafter\def\csname SGAsource@V.1.10.3\endcsname{1.10.3}
\expandafter\def\csname SGAsource@V.1.10.4\endcsname{1.10.4}
\expandafter\def\csname SGAsource@V.1.10.5\endcsname{1.10.5}
\expandafter\def\csname SGAsource@V.1.4.2\endcsname{1.4.2}
\expandafter\def\csname SGAsource@V.1.4.3\endcsname{1.4.3}
\expandafter\def\csname SGAsource@V.1.5.1\endcsname{1.5.1}
\expandafter\def\csname SGAsource@V.1.5.2\endcsname{1.5.2}
\expandafter\def\csname SGAsource@V.1.5.3\endcsname{1.5.3}
\expandafter\def\csname SGAsource@V.1.5.4\endcsname{1.5.4}
\expandafter\def\csname SGAsource@V.1.5.5\endcsname{1.5.5}
\expandafter\def\csname SGAsource@V.2.2.1\endcsname{2.2.1}
\expandafter\def\csname SGAsource@V.2.2.2\endcsname{2.2.2}
\expandafter\def\csname SGAsource@V.2.3.3.1\endcsname{2.3.3.1}
\expandafter\def\csname SGAsource@V.2.4.2.1\endcsname{2.4.2.1}
\expandafter\def\csname SGAsource@V.2.4.2.2\endcsname{2.4.2.2}
\expandafter\def\csname SGAsource@V.2.4.5.1\endcsname{2.4.5.1}
\expandafter\def\csname SGAsource@V.2.4.5.2\endcsname{2.4.5.2}
\expandafter\def\csname SGAsource@V.2.4.5.3\endcsname{2.4.5.3}
\expandafter\def\csname SGAsource@V.2.4.5.4\endcsname{2.4.5.4}
\expandafter\def\csname SGAsource@V.2.4.6.1\endcsname{2.4.6.1}
\expandafter\def\csname SGAsource@V.2.5.1.1\endcsname{2.5.1.1}
\expandafter\def\csname SGAsource@V.2.5.1.2\endcsname{2.5.1.2}
\expandafter\def\csname SGAsource@V.2.5.1.3\endcsname{2.5.1.3}
\expandafter\def\csname SGAsource@V.3.2.1\endcsname{3.2.1}
\expandafter\def\csname SGAsource@V.3.3.1\endcsname{3.3.1}
\expandafter\def\csname SGAsource@V.3.4.1\endcsname{3.4.1}
\expandafter\def\csname SGAsource@V.3.4.2\endcsname{3.4.2}
\expandafter\def\csname SGAsource@V.3.4.3\endcsname{3.4.3}
\expandafter\def\csname SGAsource@V.3.6.1\endcsname{3.6.1}
\expandafter\def\csname SGAsource@V.4.11.1\endcsname{4.11.1}
\expandafter\def\csname SGAsource@V.5.1.1\endcsname{5.1.1}
\expandafter\def\csname SGAsource@V.5.3.1\endcsname{5.3.1}
\expandafter\def\csname SGAsource@V.5.3.2\endcsname{5.3.2}
\expandafter\def\csname SGAsource@V.6.0.1\endcsname{6.0.1}
\expandafter\def\csname SGAsource@V.6.0.2\endcsname{6.0.2}
\expandafter\def\csname SGAsource@V.6.0.3\endcsname{6.0.3}
\expandafter\def\csname SGAsource@V.6.0.4\endcsname{6.0.4}
\expandafter\def\csname SGAsource@V.6.1.1\endcsname{6.1.1}
\expandafter\def\csname SGAsource@V.6.1.2\endcsname{6.1.2}
\expandafter\def\csname SGAsource@V.6.1.3\endcsname{6.1.3}
\expandafter\def\csname SGAsource@V.6.10.1\endcsname{6.10.1}
\expandafter\def\csname SGAsource@V.6.13.5\endcsname{6.13.5}
\expandafter\def\csname SGAsource@V.6.13.6\endcsname{6.13.6}
\expandafter\def\csname SGAsource@V.6.13.7\endcsname{6.13.7}
\expandafter\def\csname SGAsource@V.6.13.8\endcsname{6.13.8}
\expandafter\def\csname SGAsource@V.6.14.1\endcsname{6.14.1}
\expandafter\def\csname SGAsource@V.6.14.2\endcsname{6.14.2}
\expandafter\def\csname SGAsource@V.6.15.1\endcsname{6.15.1}
\expandafter\def\csname SGAsource@V.6.2.1\endcsname{6.2.1}
\expandafter\def\csname SGAsource@V.6.2.2\endcsname{6.2.2}
\expandafter\def\csname SGAsource@V.6.3.3\endcsname{6.3.3}
\expandafter\def\csname SGAsource@V.6.3.5\endcsname{6.3.5}
\expandafter\def\csname SGAsource@V.6.3.6\endcsname{6.3.6}
\expandafter\def\csname SGAsource@V.6.3.7\endcsname{6.3.7}
\expandafter\def\csname SGAsource@V.6.4.1\endcsname{6.4.1}
\expandafter\def\csname SGAsource@V.6.4.2\endcsname{6.4.2}
\expandafter\def\csname SGAsource@V.6.4.3\endcsname{6.4.3}
\expandafter\def\csname SGAsource@V.6.5.1\endcsname{6.5.1}
\expandafter\def\csname SGAsource@V.6.5.2\endcsname{6.5.2}
\expandafter\def\csname SGAsource@V.6.5.3\endcsname{6.5.3}
\expandafter\def\csname SGAsource@V.6.5.4\endcsname{6.5.4}
\expandafter\def\csname SGAsource@V.6.5.5\endcsname{6.5.5}
\expandafter\def\csname SGAsource@V.6.8.1\endcsname{6.8.1}
\expandafter\def\csname SGAsource@V.6.8.2\endcsname{6.8.2}
\expandafter\def\csname SGAsource@V.6.8.3\endcsname{6.8.3}
\expandafter\def\csname SGAsource@V.6.8.4\endcsname{6.8.4}
\expandafter\def\csname SGAsource@V.6.8.5\endcsname{6.8.5}
\expandafter\def\csname SGAsource@V.6.8.6\endcsname{6.8.6}
\expandafter\def\csname SGAsource@V.6.8.7\endcsname{6.8.7}
\expandafter\def\csname SGAsource@V.6.9.1\endcsname{6.9.1}
\expandafter\def\csname SGAsource@V.6.9.2\endcsname{6.9.2}
\expandafter\def\csname SGAsource@V.7.2.2\endcsname{7.2.2}
\expandafter\def\csname SGAsource@V.7.4.0.1\endcsname{7.4.0.1}
\expandafter\def\csname SGAsource@V.7.4.0.2\endcsname{7.4.0.2}
\expandafter\def\csname SGAsource@V.7.4.0.3\endcsname{7.4.0.3}
\expandafter\def\csname SGAsource@V.7.4.0.4\endcsname{7.4.0.4}
\expandafter\def\csname SGAsource@V.8.2.7.1\endcsname{8.2.7.1}
\expandafter\def\csname SGAsource@V.cor1.3.3\endcsname{1.3.3}
\expandafter\def\csname SGAsource@V.cor1.7.1\endcsname{1.7.1}
\expandafter\def\csname SGAsource@V.cor1.7.2\endcsname{1.7.2}
\expandafter\def\csname SGAsource@V.cor2.3.5\endcsname{2.3.5}
\expandafter\def\csname SGAsource@V.cor3.3\endcsname{3.3}
\expandafter\def\csname SGAsource@V.cor3.4\endcsname{3.4}
\expandafter\def\csname SGAsource@V.cor3.5\endcsname{3.5}
\expandafter\def\csname SGAsource@V.cor4.5\endcsname{4.5}
\expandafter\def\csname SGAsource@V.cor8.2.9\endcsname{8.2.9}
\expandafter\def\csname SGAsource@V.def1.1\endcsname{1.1}
\expandafter\def\csname SGAsource@V.def1.8\endcsname{1.8}
\expandafter\def\csname SGAsource@V.def4.1\endcsname{4.1}
\expandafter\def\csname SGAsource@V.def6.12\endcsname{6.12}
\expandafter\def\csname SGAsource@V.def7.1.1\endcsname{7.1.1}
\expandafter\def\csname SGAsource@V.def7.3.3\endcsname{7.3.3}
\expandafter\def\csname SGAsource@V.def8.1.1\endcsname{8.1.1}
\expandafter\def\csname SGAsource@V.def8.1.2.1\endcsname{8.1.2.1}
\expandafter\def\csname SGAsource@V.def8.1.3\endcsname{8.1.3}
\expandafter\def\csname SGAsource@V.def8.1.4\endcsname{8.1.4}
\expandafter\def\csname SGAsource@V.def8.2.1\endcsname{8.2.1}
\expandafter\def\csname SGAsource@V.eq1.4.1\endcsname{1.4.1}
\expandafter\def\csname SGAsource@V.exe0.5.2\endcsname{0.5.2}
\expandafter\def\csname SGAsource@V.exem6.12.3\endcsname{6.12.3}
\expandafter\def\csname SGAsource@V.exer2.3.6\endcsname{2.3.6}
\expandafter\def\csname SGAsource@V.exer3.6\endcsname{3.6}
\expandafter\def\csname SGAsource@V.exer4.13\endcsname{4.13}
\expandafter\def\csname SGAsource@V.exer4.15\endcsname{4.15}
\expandafter\def\csname SGAsource@V.exer4.16\endcsname{4.16}
\expandafter\def\csname SGAsource@V.exer4.17\endcsname{4.17}
\expandafter\def\csname SGAsource@V.lem1.6.1\endcsname{1.6.1}
\expandafter\def\csname SGAsource@V.lem7.2.1\endcsname{7.2.1}
\expandafter\def\csname SGAsource@V.lem7.3.4\endcsname{7.3.4}
\expandafter\def\csname SGAsource@V.lem7.3.6\endcsname{7.3.6}
\expandafter\def\csname SGAsource@V.lem8.1.11\endcsname{8.1.11}
\expandafter\def\csname SGAsource@V.lem8.1.6\endcsname{8.1.6}
\expandafter\def\csname SGAsource@V.lem8.1.8\endcsname{8.1.8}
\expandafter\def\csname SGAsource@V.lem8.2.10\endcsname{8.2.10}
\expandafter\def\csname SGAsource@V.lem8.2.11\endcsname{8.2.11}
\expandafter\def\csname SGAsource@V.lem8.2.3\endcsname{8.2.3}
\expandafter\def\csname SGAsource@V.lem8.2.8\endcsname{8.2.8}
\expandafter\def\csname SGAsource@V.par7.3.1.1\endcsname{7.3.1.1}
\expandafter\def\csname SGAsource@V.par7.3.1.2\endcsname{7.3.1.2}
\expandafter\def\csname SGAsource@V.par7.3.1.3\endcsname{7.3.1.3}
\expandafter\def\csname SGAsource@V.par7.3.1.4\endcsname{7.3.1.4}
\expandafter\def\csname SGAsource@V.par7.3.1.6\endcsname{7.3.1.6}
\expandafter\def\csname SGAsource@V.par7.3.1.7\endcsname{7.3.1.7}
\expandafter\def\csname SGAsource@V.par7.3.6.0\endcsname{7.3.6.0}
\expandafter\def\csname SGAsource@V.par8.1.1.1\endcsname{8.1.1.1}
\expandafter\def\csname SGAsource@V.par8.1.13.0\endcsname{8.1.13.0}
\expandafter\def\csname SGAsource@V.par8.1.9.0\endcsname{8.1.9.0}
\expandafter\def\csname SGAsource@V.prob4.14\endcsname{4.14}
\expandafter\def\csname SGAsource@V.prop0.1\endcsname{0.1}
\expandafter\def\csname SGAsource@V.prop0.2\endcsname{0.2}
\expandafter\def\csname SGAsource@V.prop0.3\endcsname{0.3}
\expandafter\def\csname SGAsource@V.prop0.4\endcsname{0.4}
\expandafter\def\csname SGAsource@V.prop0.5\endcsname{0.5}
\expandafter\def\csname SGAsource@V.prop1.11\endcsname{1.11}
\expandafter\def\csname SGAsource@V.prop1.2\endcsname{1.2}
\expandafter\def\csname SGAsource@V.prop1.3.1\endcsname{1.3.1}
\expandafter\def\csname SGAsource@V.prop1.3.2\endcsname{1.3.2}
\expandafter\def\csname SGAsource@V.prop1.5\endcsname{1.5}
\expandafter\def\csname SGAsource@V.prop1.6\endcsname{1.6}
\expandafter\def\csname SGAsource@V.prop1.7\endcsname{1.7}
\expandafter\def\csname SGAsource@V.prop1.9\endcsname{1.9}
\expandafter\def\csname SGAsource@V.prop2.3.4\endcsname{2.3.4}
\expandafter\def\csname SGAsource@V.prop3.1\endcsname{3.1}
\expandafter\def\csname SGAsource@V.prop4.10\endcsname{4.10}
\expandafter\def\csname SGAsource@V.prop4.11\endcsname{4.11}
\expandafter\def\csname SGAsource@V.prop4.3\endcsname{4.3}
\expandafter\def\csname SGAsource@V.prop4.7\endcsname{4.7}
\expandafter\def\csname SGAsource@V.prop4.9\endcsname{4.9}
\expandafter\def\csname SGAsource@V.prop5.1\endcsname{5.1}
\expandafter\def\csname SGAsource@V.prop5.2\endcsname{5.2}
\expandafter\def\csname SGAsource@V.prop5.3\endcsname{5.3}
\expandafter\def\csname SGAsource@V.prop5.4\endcsname{5.4}
\expandafter\def\csname SGAsource@V.prop6.1\endcsname{6.1}
\expandafter\def\csname SGAsource@V.prop6.10\endcsname{6.10}
\expandafter\def\csname SGAsource@V.prop6.15\endcsname{6.15}
\expandafter\def\csname SGAsource@V.prop6.2\endcsname{6.2}
\expandafter\def\csname SGAsource@V.prop6.4\endcsname{6.4}
\expandafter\def\csname SGAsource@V.prop6.5\endcsname{6.5}
\expandafter\def\csname SGAsource@V.prop6.6\endcsname{6.6}
\expandafter\def\csname SGAsource@V.prop6.7\endcsname{6.7}
\expandafter\def\csname SGAsource@V.prop6.9\endcsname{6.9}
\expandafter\def\csname SGAsource@V.prop7.1.2\endcsname{7.1.2}
\expandafter\def\csname SGAsource@V.prop8.1.12\endcsname{8.1.12}
\expandafter\def\csname SGAsource@V.prop8.1.13.1\endcsname{8.1.13.1}
\expandafter\def\csname SGAsource@V.prop8.1.14\endcsname{8.1.14}
\expandafter\def\csname SGAsource@V.prop8.1.5\endcsname{8.1.5}
\expandafter\def\csname SGAsource@V.prop8.1.7\endcsname{8.1.7}
\expandafter\def\csname SGAsource@V.prop8.1.9\endcsname{8.1.9}
\expandafter\def\csname SGAsource@V.prop8.2.4\endcsname{8.2.4}
\expandafter\def\csname SGAsource@V.prop8.2.5\endcsname{8.2.5}
\expandafter\def\csname SGAsource@V.prop8.2.6\endcsname{8.2.6}
\expandafter\def\csname SGAsource@V.prop8.2.7\endcsname{8.2.7}
\expandafter\def\csname SGAsource@V.rem0.2.1\endcsname{0.2.1}
\expandafter\def\csname SGAsource@V.rem0.5.1\endcsname{0.5.1}
\expandafter\def\csname SGAsource@V.rem7.4.2\endcsname{7.4.2}
\expandafter\def\csname SGAsource@V.sec0\endcsname{0}
\expandafter\def\csname SGAsource@V.sec1\endcsname{1}
\expandafter\def\csname SGAsource@V.sec2\endcsname{2}
\expandafter\def\csname SGAsource@V.sec3\endcsname{3}
\expandafter\def\csname SGAsource@V.sec4\endcsname{4}
\expandafter\def\csname SGAsource@V.sec5\endcsname{5}
\expandafter\def\csname SGAsource@V.sec6\endcsname{6}
\expandafter\def\csname SGAsource@V.sec7\endcsname{7}
\expandafter\def\csname SGAsource@V.sec8\endcsname{8}
\expandafter\def\csname SGAsource@V.subsec1.10\endcsname{1.10}
\expandafter\def\csname SGAsource@V.subsec1.4\endcsname{1.4}
\expandafter\def\csname SGAsource@V.subsec2.1\endcsname{2.1}
\expandafter\def\csname SGAsource@V.subsec2.2\endcsname{2.2}
\expandafter\def\csname SGAsource@V.subsec2.3\endcsname{2.3}
\expandafter\def\csname SGAsource@V.subsec2.4\endcsname{2.4}
\expandafter\def\csname SGAsource@V.subsec2.5\endcsname{2.5}
\expandafter\def\csname SGAsource@V.subsec4.12\endcsname{4.12}
\expandafter\def\csname SGAsource@V.subsec4.2\endcsname{4.2}
\expandafter\def\csname SGAsource@V.subsec4.4\endcsname{4.4}
\expandafter\def\csname SGAsource@V.subsec4.6\endcsname{4.6}
\expandafter\def\csname SGAsource@V.subsec4.8\endcsname{4.8}
\expandafter\def\csname SGAsource@V.subsec5.0\endcsname{5.0}
\expandafter\def\csname SGAsource@V.subsec6.0\endcsname{6.0}
\expandafter\def\csname SGAsource@V.subsec6.11\endcsname{6.11}
\expandafter\def\csname SGAsource@V.subsec6.13\endcsname{6.13}
\expandafter\def\csname SGAsource@V.subsec6.14\endcsname{6.14}
\expandafter\def\csname SGAsource@V.subsec6.16\endcsname{6.16}
\expandafter\def\csname SGAsource@V.subsec6.3\endcsname{6.3}
\expandafter\def\csname SGAsource@V.subsec6.8\endcsname{6.8}
\expandafter\def\csname SGAsource@V.subsec7.1\endcsname{7.1}
\expandafter\def\csname SGAsource@V.subsec7.2\endcsname{7.2}
\expandafter\def\csname SGAsource@V.subsec7.3\endcsname{7.3}
\expandafter\def\csname SGAsource@V.subsec7.4\endcsname{7.4}
\expandafter\def\csname SGAsource@V.subsec8.1\endcsname{8.1}
\expandafter\def\csname SGAsource@V.subsec8.2\endcsname{8.2}
\expandafter\def\csname SGAsource@V.subsubsec0.1.1\endcsname{0.1.1}
\expandafter\def\csname SGAsource@V.subsubsec0.5.2\endcsname{0.5.2}
\expandafter\def\csname SGAsource@V.subsubsec2.1.1\endcsname{2.1.1}
\expandafter\def\csname SGAsource@V.subsubsec2.1.2\endcsname{2.1.2}
\expandafter\def\csname SGAsource@V.subsubsec2.3.1\endcsname{2.3.1}
\expandafter\def\csname SGAsource@V.subsubsec2.3.2\endcsname{2.3.2}
\expandafter\def\csname SGAsource@V.subsubsec2.3.3\endcsname{2.3.3}
\expandafter\def\csname SGAsource@V.subsubsec2.4.1\endcsname{2.4.1}
\expandafter\def\csname SGAsource@V.subsubsec2.4.2\endcsname{2.4.2}
\expandafter\def\csname SGAsource@V.subsubsec2.4.3\endcsname{2.4.3}
\expandafter\def\csname SGAsource@V.subsubsec2.4.4\endcsname{2.4.4}
\expandafter\def\csname SGAsource@V.subsubsec2.4.5\endcsname{2.4.5}
\expandafter\def\csname SGAsource@V.subsubsec2.4.6\endcsname{2.4.6}
\expandafter\def\csname SGAsource@V.subsubsec2.5.1\endcsname{2.5.1}
\expandafter\def\csname SGAsource@V.subsubsec2.5.2\endcsname{2.5.2}
\expandafter\def\csname SGAsource@V.subsubsec6.12.1\endcsname{6.12.1}
\expandafter\def\csname SGAsource@V.subsubsec6.12.2\endcsname{6.12.2}
\expandafter\def\csname SGAsource@V.subsubsec7.1.0\endcsname{7.1.0}
\expandafter\def\csname SGAsource@V.subsubsec7.2.0\endcsname{7.2.0}
\expandafter\def\csname SGAsource@V.subsubsec7.3.0\endcsname{7.3.0}
\expandafter\def\csname SGAsource@V.subsubsec7.3.1\endcsname{7.3.1}
\expandafter\def\csname SGAsource@V.subsubsec7.3.5\endcsname{7.3.5}
\expandafter\def\csname SGAsource@V.subsubsec7.3.7\endcsname{7.3.7}
\expandafter\def\csname SGAsource@V.subsubsec7.4.0\endcsname{7.4.0}
\expandafter\def\csname SGAsource@V.subsubsec8.1.0\endcsname{8.1.0}
\expandafter\def\csname SGAsource@V.subsubsec8.1.10\endcsname{8.1.10}
\expandafter\def\csname SGAsource@V.subsubsec8.1.2\endcsname{8.1.2}
\expandafter\def\csname SGAsource@V.subsubsec8.2.0\endcsname{8.2.0}
\expandafter\def\csname SGAsource@V.subusbsec1.6.2\endcsname{1.6.2}
\expandafter\def\csname SGAsource@V.thm3.2\endcsname{3.2}
\expandafter\def\csname SGAsource@V.thm7.3.2\endcsname{7.3.2}
\expandafter\def\csname SGAsource@V.thm7.4.1\endcsname{7.4.1}
\expandafter\def\csname SGAsource@V.thm8.2.12\endcsname{8.2.12}
\expandafter\def\csname SGAsource@V.thm8.2.2\endcsname{8.2.2}
\expandafter\def\csname SGAsource@VI.cor1.11\endcsname{1.11}
\expandafter\def\csname SGAsource@VI.cor1.12\endcsname{1.12}
\expandafter\def\csname SGAsource@VI.cor1.15\endcsname{1.15}
\expandafter\def\csname SGAsource@VI.cor1.16\endcsname{1.16}
\expandafter\def\csname SGAsource@VI.cor1.17\endcsname{1.17}
\expandafter\def\csname SGAsource@VI.cor1.17.1\endcsname{1.17.1}
\expandafter\def\csname SGAsource@VI.cor1.18\endcsname{1.18}
\expandafter\def\csname SGAsource@VI.cor1.19\endcsname{1.19}
\expandafter\def\csname SGAsource@VI.cor1.24\endcsname{1.24}
\expandafter\def\csname SGAsource@VI.cor1.24.2\endcsname{1.24.2}
\expandafter\def\csname SGAsource@VI.cor1.25\endcsname{1.25}
\expandafter\def\csname SGAsource@VI.cor1.26\endcsname{1.26}
\expandafter\def\csname SGAsource@VI.cor1.4\endcsname{1.4}
\expandafter\def\csname SGAsource@VI.cor1.8.1\endcsname{1.8.1}
\expandafter\def\csname SGAsource@VI.cor1.8.2\endcsname{1.8.2}
\expandafter\def\csname SGAsource@VI.cor2.1.1\endcsname{2.1.1}
\expandafter\def\csname SGAsource@VI.cor2.1.2\endcsname{2.1.2}
\expandafter\def\csname SGAsource@VI.cor2.6\endcsname{2.6}
\expandafter\def\csname SGAsource@VI.cor2.7\endcsname{2.7}
\expandafter\def\csname SGAsource@VI.cor2.8\endcsname{2.8}
\expandafter\def\csname SGAsource@VI.cor2.9.4\endcsname{2.9.4}
\expandafter\def\csname SGAsource@VI.cor3.3\endcsname{3.3}
\expandafter\def\csname SGAsource@VI.cor3.6\endcsname{3.6}
\expandafter\def\csname SGAsource@VI.cor3.7\endcsname{3.7}
\expandafter\def\csname SGAsource@VI.cor3.8\endcsname{3.8}
\expandafter\def\csname SGAsource@VI.cor4.11\endcsname{4.11}
\expandafter\def\csname SGAsource@VI.cor4.14\endcsname{4.14}
\expandafter\def\csname SGAsource@VI.cor4.3\endcsname{4.3}
\expandafter\def\csname SGAsource@VI.cor4.4\endcsname{4.4}
\expandafter\def\csname SGAsource@VI.cor4.5\endcsname{4.5}
\expandafter\def\csname SGAsource@VI.cor4.7\endcsname{4.7}
\expandafter\def\csname SGAsource@VI.cor5.10\endcsname{5.10}
\expandafter\def\csname SGAsource@VI.cor5.2\endcsname{5.2}
\expandafter\def\csname SGAsource@VI.cor5.3\endcsname{5.3}
\expandafter\def\csname SGAsource@VI.cor5.5\endcsname{5.5}
\expandafter\def\csname SGAsource@VI.cor5.8\endcsname{5.8}
\expandafter\def\csname SGAsource@VI.cor7.4.11\endcsname{7.4.11}
\expandafter\def\csname SGAsource@VI.cor8.3.14\endcsname{8.3.14}
\expandafter\def\csname SGAsource@VI.cor8.5.10\endcsname{8.5.10}
\expandafter\def\csname SGAsource@VI.cor8.5.4\endcsname{8.5.4}
\expandafter\def\csname SGAsource@VI.cor8.5.5\endcsname{8.5.5}
\expandafter\def\csname SGAsource@VI.cor8.5.6\endcsname{8.5.6}
\expandafter\def\csname SGAsource@VI.cor8.5.7\endcsname{8.5.7}
\expandafter\def\csname SGAsource@VI.cor8.7.4\endcsname{8.7.4}
\expandafter\def\csname SGAsource@VI.cor8.7.5\endcsname{8.7.5}
\expandafter\def\csname SGAsource@VI.cor8.7.6\endcsname{8.7.6}
\expandafter\def\csname SGAsource@VI.cor8.7.7\endcsname{8.7.7}
\expandafter\def\csname SGAsource@VI.cor8.7.9\endcsname{8.7.9}
\expandafter\def\csname SGAsource@VI.cor9.5\endcsname{9.5}
\expandafter\def\csname SGAsource@VI.def1.1\endcsname{1.1}
\expandafter\def\csname SGAsource@VI.def1.13\endcsname{1.13}
\expandafter\def\csname SGAsource@VI.def1.30\endcsname{1.30}
\expandafter\def\csname SGAsource@VI.def1.7\endcsname{1.7}
\expandafter\def\csname SGAsource@VI.def2.11\endcsname{2.11}
\expandafter\def\csname SGAsource@VI.def2.3\endcsname{2.3}
\expandafter\def\csname SGAsource@VI.def2.9.1\endcsname{2.9.1}
\expandafter\def\csname SGAsource@VI.def3.1\endcsname{3.1}
\expandafter\def\csname SGAsource@VI.def3.7\endcsname{3.7}
\expandafter\def\csname SGAsource@VI.def5.6\endcsname{5.6}
\expandafter\def\csname SGAsource@VI.def6.1\endcsname{6.1}
\expandafter\def\csname SGAsource@VI.def6.3\endcsname{6.3}
\expandafter\def\csname SGAsource@VI.def7.1.1\endcsname{7.1.1}
\expandafter\def\csname SGAsource@VI.def7.1.5\endcsname{7.1.5}
\expandafter\def\csname SGAsource@VI.def7.4.1\endcsname{7.4.1}
\expandafter\def\csname SGAsource@VI.def8.1.1\endcsname{8.1.1}
\expandafter\def\csname SGAsource@VI.def8.2.5\endcsname{8.2.5}
\expandafter\def\csname SGAsource@VI.eq1.23.1\endcsname{1.23.1}
\expandafter\def\csname SGAsource@VI.eq1.24.1\endcsname{1.24.1}
\expandafter\def\csname SGAsource@VI.eq4.10.1\endcsname{4.10.1}
\expandafter\def\csname SGAsource@VI.eq4.10.2\endcsname{4.10.2}
\expandafter\def\csname SGAsource@VI.eq4.10.3\endcsname{4.10.3}
\expandafter\def\csname SGAsource@VI.eq4.10.4\endcsname{4.10.4}
\expandafter\def\csname SGAsource@VI.eq4.10.5\endcsname{4.10.5}
\expandafter\def\csname SGAsource@VI.eq4.10.6\endcsname{4.10.6}
\expandafter\def\csname SGAsource@VI.eq4.13.1\endcsname{4.13.1}
\expandafter\def\csname SGAsource@VI.eq4.13.2\endcsname{4.13.2}
\expandafter\def\csname SGAsource@VI.eq4.13.3\endcsname{4.13.3}
\expandafter\def\csname SGAsource@VI.eq4.13.4\endcsname{4.13.4}
\expandafter\def\csname SGAsource@VI.eq4.9.1\endcsname{4.9.1}
\expandafter\def\csname SGAsource@VI.eq4.9.2\endcsname{4.9.2}
\expandafter\def\csname SGAsource@VI.eq4.9.4\endcsname{4.9.4}
\expandafter\def\csname SGAsource@VI.eq4.9.5\endcsname{4.9.5}
\expandafter\def\csname SGAsource@VI.eq4.9.6\endcsname{4.9.6}
\expandafter\def\csname SGAsource@VI.eq4.9.7\endcsname{4.9.7}
\expandafter\def\csname SGAsource@VI.eq5.5.1\endcsname{5.5.1}
\expandafter\def\csname SGAsource@VI.eq5.9.1\endcsname{5.9.1}
\expandafter\def\csname SGAsource@VI.eq5.9.2\endcsname{5.9.2}
\expandafter\def\csname SGAsource@VI.eq6.1.2.1\endcsname{6.1.2.1}
\expandafter\def\csname SGAsource@VI.eq6.1.2.2\endcsname{6.1.2.2}
\expandafter\def\csname SGAsource@VI.eq7.1.3.1\endcsname{7.1.3.1}
\expandafter\def\csname SGAsource@VI.eq7.1.3.2\endcsname{7.1.3.2}
\expandafter\def\csname SGAsource@VI.eq7.1.6.1\endcsname{7.1.6.1}
\expandafter\def\csname SGAsource@VI.eq7.1.6.2\endcsname{7.1.6.2}
\expandafter\def\csname SGAsource@VI.eq7.1.6.3\endcsname{7.1.6.3}
\expandafter\def\csname SGAsource@VI.eq7.1.6.4\endcsname{7.1.6.4}
\expandafter\def\csname SGAsource@VI.eq7.1.9.1\endcsname{7.1.9.1}
\expandafter\def\csname SGAsource@VI.eq7.1.9.2\endcsname{7.1.9.2}
\expandafter\def\csname SGAsource@VI.eq7.2.2.1\endcsname{7.2.2.1}
\expandafter\def\csname SGAsource@VI.eq7.2.2.2\endcsname{7.2.2.2}
\expandafter\def\csname SGAsource@VI.eq7.2.5.1\endcsname{7.2.5.1}
\expandafter\def\csname SGAsource@VI.eq7.2.6.1\endcsname{7.2.6.1}
\expandafter\def\csname SGAsource@VI.eq7.2.6.2\endcsname{7.2.6.2}
\expandafter\def\csname SGAsource@VI.eq7.2.6.3\endcsname{7.2.6.3}
\expandafter\def\csname SGAsource@VI.eq7.2.6.4\endcsname{7.2.6.4}
\expandafter\def\csname SGAsource@VI.eq7.2.6.5\endcsname{7.2.6.5}
\expandafter\def\csname SGAsource@VI.eq7.2.6.6\endcsname{7.2.6.6}
\expandafter\def\csname SGAsource@VI.eq7.2.6.7\endcsname{7.2.6.7}
\expandafter\def\csname SGAsource@VI.eq7.3.1.1\endcsname{7.3.1.1}
\expandafter\def\csname SGAsource@VI.eq7.3.1.2\endcsname{7.3.1.2}
\expandafter\def\csname SGAsource@VI.eq7.3.1.3\endcsname{7.3.1.3}
\expandafter\def\csname SGAsource@VI.eq7.3.1.4\endcsname{7.3.1.4}
\expandafter\def\csname SGAsource@VI.eq7.3.1.5\endcsname{7.3.1.5}
\expandafter\def\csname SGAsource@VI.eq7.3.2.1\endcsname{7.3.2.1}
\expandafter\def\csname SGAsource@VI.eq7.3.2.2\endcsname{7.3.2.2}
\expandafter\def\csname SGAsource@VI.eq7.3.2.3\endcsname{7.3.2.3}
\expandafter\def\csname SGAsource@VI.eq7.3.2.4\endcsname{7.3.2.4}
\expandafter\def\csname SGAsource@VI.eq7.3.2.5\endcsname{7.3.2.5}
\expandafter\def\csname SGAsource@VI.eq7.3.2.6\endcsname{7.3.2.6}
\expandafter\def\csname SGAsource@VI.eq7.3.2.7\endcsname{7.3.2.7}
\expandafter\def\csname SGAsource@VI.eq7.4.10.1\endcsname{7.4.10.1}
\expandafter\def\csname SGAsource@VI.eq7.4.10.2\endcsname{7.4.10.2}
\expandafter\def\csname SGAsource@VI.eq7.4.10.3\endcsname{7.4.10.3}
\expandafter\def\csname SGAsource@VI.eq7.4.10.4\endcsname{7.4.10.4}
\expandafter\def\csname SGAsource@VI.eq7.4.10.5\endcsname{7.4.10.5}
\expandafter\def\csname SGAsource@VI.eq7.4.10.6\endcsname{7.4.10.6}
\expandafter\def\csname SGAsource@VI.eq7.4.5.1\endcsname{7.4.5.1}
\expandafter\def\csname SGAsource@VI.eq7.4.5.2\endcsname{7.4.5.2}
\expandafter\def\csname SGAsource@VI.eq7.4.5.3\endcsname{7.4.5.3}
\expandafter\def\csname SGAsource@VI.eq7.4.5.4\endcsname{7.4.5.4}
\expandafter\def\csname SGAsource@VI.eq7.4.5.5\endcsname{7.4.5.5}
\expandafter\def\csname SGAsource@VI.eq7.4.5.6\endcsname{7.4.5.6}
\expandafter\def\csname SGAsource@VI.eq7.4.5.7\endcsname{7.4.5.7}
\expandafter\def\csname SGAsource@VI.eq7.4.5.8\endcsname{7.4.5.8}
\expandafter\def\csname SGAsource@VI.eq7.4.6.1\endcsname{7.4.6.1}
\expandafter\def\csname SGAsource@VI.eq7.4.9.1\endcsname{7.4.9.1}
\expandafter\def\csname SGAsource@VI.eq7.4.9.2\endcsname{7.4.9.2}
\expandafter\def\csname SGAsource@VI.eq8.1.1.1\endcsname{8.1.1.1}
\expandafter\def\csname SGAsource@VI.eq8.1.2.1\endcsname{8.1.2.1}
\expandafter\def\csname SGAsource@VI.eq8.1.2.2\endcsname{8.1.2.2}
\expandafter\def\csname SGAsource@VI.eq8.1.3.1\endcsname{8.1.3.1.}
\expandafter\def\csname SGAsource@VI.eq8.1.4.1\endcsname{8.1.4.1}
\expandafter\def\csname SGAsource@VI.eq8.1.4.2\endcsname{8.1.4.2}
\expandafter\def\csname SGAsource@VI.eq8.1.6.1\endcsname{8.1.6.1}
\expandafter\def\csname SGAsource@VI.eq8.1.6.2\endcsname{8.1.6.2}
\expandafter\def\csname SGAsource@VI.eq8.1.6.3\endcsname{8.1.6.3}
\expandafter\def\csname SGAsource@VI.eq8.1.6.4\endcsname{8.1.6.4}
\expandafter\def\csname SGAsource@VI.eq8.1.6.5\endcsname{8.1.6.5}
\expandafter\def\csname SGAsource@VI.eq8.2.11.1\endcsname{8.2.11.1}
\expandafter\def\csname SGAsource@VI.eq8.2.8.1\endcsname{8.2.8.1}
\expandafter\def\csname SGAsource@VI.eq8.3.11.1\endcsname{8.3.11.1}
\expandafter\def\csname SGAsource@VI.eq8.3.13.1\endcsname{8.3.13.1}
\expandafter\def\csname SGAsource@VI.eq8.4.1.1\endcsname{8.4.1.1}
\expandafter\def\csname SGAsource@VI.eq8.4.2.1\endcsname{8.4.2.1}
\expandafter\def\csname SGAsource@VI.eq8.4.3.1\endcsname{8.4.3.1}
\expandafter\def\csname SGAsource@VI.eq8.5.1.1\endcsname{8.5.1.1}
\expandafter\def\csname SGAsource@VI.eq8.5.1.2\endcsname{8.5.1.2}
\expandafter\def\csname SGAsource@VI.eq8.5.1.3\endcsname{8.5.1.3}
\expandafter\def\csname SGAsource@VI.eq8.5.1.4\endcsname{8.5.1.4}
\expandafter\def\csname SGAsource@VI.eq8.5.1.5\endcsname{8.5.1.5}
\expandafter\def\csname SGAsource@VI.eq8.5.1.6\endcsname{8.5.1.6}
\expandafter\def\csname SGAsource@VI.eq8.5.1.7\endcsname{8.5.1.7}
\expandafter\def\csname SGAsource@VI.eq8.5.10.1\endcsname{8.5.10.1}
\expandafter\def\csname SGAsource@VI.eq8.5.2.1\endcsname{8.5.2.1}
\expandafter\def\csname SGAsource@VI.eq8.5.2.2\endcsname{8.5.2.2}
\expandafter\def\csname SGAsource@VI.eq8.5.2.3\endcsname{8.5.2.3}
\expandafter\def\csname SGAsource@VI.eq8.5.2.4\endcsname{8.5.2.4}
\expandafter\def\csname SGAsource@VI.eq8.5.2.5\endcsname{8.5.2.5}
\expandafter\def\csname SGAsource@VI.eq8.5.2.6\endcsname{8.5.2.6}
\expandafter\def\csname SGAsource@VI.eq8.5.2.7\endcsname{8.5.2.7}
\expandafter\def\csname SGAsource@VI.eq8.5.2.8\endcsname{8.5.2.8}
\expandafter\def\csname SGAsource@VI.eq8.5.2.9\endcsname{8.5.2.9}
\expandafter\def\csname SGAsource@VI.eq8.5.3.1\endcsname{8.5.3.1}
\expandafter\def\csname SGAsource@VI.eq8.5.5.1\endcsname{8.5.5.1}
\expandafter\def\csname SGAsource@VI.eq8.5.7.1\endcsname{8.5.7.1}
\expandafter\def\csname SGAsource@VI.eq8.5.7.2\endcsname{8.5.7.2}
\expandafter\def\csname SGAsource@VI.eq8.5.8.1\endcsname{8.5.8.1}
\expandafter\def\csname SGAsource@VI.eq8.5.8.2\endcsname{8.5.8.2}
\expandafter\def\csname SGAsource@VI.eq8.5.9.1\endcsname{8.5.9.1}
\expandafter\def\csname SGAsource@VI.eq8.7.3.1\endcsname{8.7.3.1}
\expandafter\def\csname SGAsource@VI.eq8.7.4.1\endcsname{8.7.4.1}
\expandafter\def\csname SGAsource@VI.eq8.7.5.1\endcsname{8.7.5.1}
\expandafter\def\csname SGAsource@VI.eq8.7.6.1\endcsname{8.7.6.1}
\expandafter\def\csname SGAsource@VI.eq8.7.6.2\endcsname{8.7.6.2}
\expandafter\def\csname SGAsource@VI.eq8.7.6.3\endcsname{8.7.6.3}
\expandafter\def\csname SGAsource@VI.eq8.7.7.1\endcsname{8.7.7.1}
\expandafter\def\csname SGAsource@VI.eq8.7.7.2\endcsname{8.7.7.2}
\expandafter\def\csname SGAsource@VI.eq8.7.9.1\endcsname{8.7.9.1}
\expandafter\def\csname SGAsource@VI.exe6.6\endcsname{6.6}
\expandafter\def\csname SGAsource@VI.exe6.8\endcsname{6.8}
\expandafter\def\csname SGAsource@VI.exe7.3.3\endcsname{7.3.3}
\expandafter\def\csname SGAsource@VI.exe7.4.14\endcsname{7.4.14}
\expandafter\def\csname SGAsource@VI.exe7.4.15\endcsname{7.4.15}
\expandafter\def\csname SGAsource@VI.exe8.3.15\endcsname{8.3.15}
\expandafter\def\csname SGAsource@VI.exem1.33\endcsname{1.33}
\expandafter\def\csname SGAsource@VI.exem1.6.1\endcsname{1.6.1}
\expandafter\def\csname SGAsource@VI.exem2.15\endcsname{2.15}
\expandafter\def\csname SGAsource@VI.exem3.10\endcsname{3.10}
\expandafter\def\csname SGAsource@VI.exem7.3.1\endcsname{7.3.1}
\expandafter\def\csname SGAsource@VI.exem7.3.2\endcsname{7.3.2}
\expandafter\def\csname SGAsource@VI.exer1.28\endcsname{1.28}
\expandafter\def\csname SGAsource@VI.exer2.16\endcsname{2.16}
\expandafter\def\csname SGAsource@VI.exer2.17\endcsname{2.17}
\expandafter\def\csname SGAsource@VI.exer2.18\endcsname{2.18}
\expandafter\def\csname SGAsource@VI.exer3.11\endcsname{3.11}
\expandafter\def\csname SGAsource@VI.exer3.12\endcsname{3.12}
\expandafter\def\csname SGAsource@VI.exer4.8\endcsname{4.8}
\expandafter\def\csname SGAsource@VI.lem1.25.1\endcsname{1.25.1}
\expandafter\def\csname SGAsource@VI.lem2.1.3\endcsname{2.1.3}
\expandafter\def\csname SGAsource@VI.lem4.12.1\endcsname{4.12.1}
\expandafter\def\csname SGAsource@VI.lem4.2.1\endcsname{4.2.1}
\expandafter\def\csname SGAsource@VI.lem4.5.1\endcsname{4.5.1}
\expandafter\def\csname SGAsource@VI.lem5.4\endcsname{5.4}
\expandafter\def\csname SGAsource@VI.lem7.4.12\endcsname{7.4.12}
\expandafter\def\csname SGAsource@VI.lem7.4.2.2\endcsname{7.4.2.2}
\expandafter\def\csname SGAsource@VI.lem8.2.1\endcsname{8.2.1}
\expandafter\def\csname SGAsource@VI.lem8.2.2\endcsname{8.2.2}
\expandafter\def\csname SGAsource@VI.lem8.3.11\endcsname{8.3.11}
\expandafter\def\csname SGAsource@VI.lem8.3.4\endcsname{8.3.4}
\expandafter\def\csname SGAsource@VI.lem8.7.2\endcsname{8.7.2}
\expandafter\def\csname SGAsource@VI.lem9.1\endcsname{9.1}
\expandafter\def\csname SGAsource@VI.lem9.2\endcsname{9.2}
\expandafter\def\csname SGAsource@VI.lem9.3\endcsname{9.3}
\expandafter\def\csname SGAsource@VI.lem9.4\endcsname{9.4}
\expandafter\def\csname SGAsource@VI.par7.2.6.8\endcsname{7.2.6.8}
\expandafter\def\csname SGAsource@VI.par7.4.13.2\endcsname{7.4.13.2}
\expandafter\def\csname SGAsource@VI.par7.4.13.3\endcsname{7.4.13.3}
\expandafter\def\csname SGAsource@VI.par7.4.13.4\endcsname{7.4.13.4}
\expandafter\def\csname SGAsource@VI.par7.4.13.5\endcsname{7.4.13.5}
\expandafter\def\csname SGAsource@VI.par7.4.13.6\endcsname{7.4.13.6}
\expandafter\def\csname SGAsource@VI.par7.4.13.7\endcsname{7.4.13.7}
\expandafter\def\csname SGAsource@VI.par7.4.2.1\endcsname{7.4.2.1}
\expandafter\def\csname SGAsource@VI.par8.1.3.2\endcsname{8.1.3.2}
\expandafter\def\csname SGAsource@VI.par8.1.3.3\endcsname{8.1.3.3}
\expandafter\def\csname SGAsource@VI.prob2.9.5\endcsname{2.9.5}
\expandafter\def\csname SGAsource@VI.prop1.10\endcsname{1.10}
\expandafter\def\csname SGAsource@VI.prop1.14\endcsname{1.14}
\expandafter\def\csname SGAsource@VI.prop1.2\endcsname{1.2}
\expandafter\def\csname SGAsource@VI.prop1.3\endcsname{1.3}
\expandafter\def\csname SGAsource@VI.prop1.8\endcsname{1.8}
\expandafter\def\csname SGAsource@VI.prop2.1\endcsname{2.1}
\expandafter\def\csname SGAsource@VI.prop2.10\endcsname{2.10}
\expandafter\def\csname SGAsource@VI.prop2.2\endcsname{2.2}
\expandafter\def\csname SGAsource@VI.prop2.5\endcsname{2.5}
\expandafter\def\csname SGAsource@VI.prop2.9\endcsname{2.9}
\expandafter\def\csname SGAsource@VI.prop2.9.3\endcsname{2.9.3}
\expandafter\def\csname SGAsource@VI.prop3.2\endcsname{3.2}
\expandafter\def\csname SGAsource@VI.prop3.4\endcsname{3.4}
\expandafter\def\csname SGAsource@VI.prop3.5\endcsname{3.5}
\expandafter\def\csname SGAsource@VI.prop4.10\endcsname{4.10}
\expandafter\def\csname SGAsource@VI.prop4.2\endcsname{4.2}
\expandafter\def\csname SGAsource@VI.prop4.6\endcsname{4.6}
\expandafter\def\csname SGAsource@VI.prop5.7\endcsname{5.7}
\expandafter\def\csname SGAsource@VI.prop6.2\endcsname{6.2}
\expandafter\def\csname SGAsource@VI.prop6.4\endcsname{6.4}
\expandafter\def\csname SGAsource@VI.prop6.5\endcsname{6.5}
\expandafter\def\csname SGAsource@VI.prop6.7\endcsname{6.7}
\expandafter\def\csname SGAsource@VI.prop6.9\endcsname{6.9}
\expandafter\def\csname SGAsource@VI.prop7.2.7\endcsname{7.2.7}
\expandafter\def\csname SGAsource@VI.prop7.4.10\endcsname{7.4.10}
\expandafter\def\csname SGAsource@VI.prop7.4.2\endcsname{7.4.2}
\expandafter\def\csname SGAsource@VI.prop7.4.4\endcsname{7.4.4}
\expandafter\def\csname SGAsource@VI.prop7.4.7\endcsname{7.4.7}
\expandafter\def\csname SGAsource@VI.prop8.1.7\endcsname{8.1.7}
\expandafter\def\csname SGAsource@VI.prop8.3.3\endcsname{8.3.3}
\expandafter\def\csname SGAsource@VI.prop8.3.6\endcsname{8.3.6}
\expandafter\def\csname SGAsource@VI.prop8.5.2\endcsname{8.5.2}
\expandafter\def\csname SGAsource@VI.prop8.5.3\endcsname{8.5.3}
\expandafter\def\csname SGAsource@VI.prop8.5.8\endcsname{8.5.8}
\expandafter\def\csname SGAsource@VI.prop8.5.9\endcsname{8.5.9}
\expandafter\def\csname SGAsource@VI.prop9.0\endcsname{9.0}
\expandafter\def\csname SGAsource@VI.rem1.20.1\endcsname{1.20.1}
\expandafter\def\csname SGAsource@VI.rem1.23.2\endcsname{1.23.2}
\expandafter\def\csname SGAsource@VI.rem1.27\endcsname{1.27}
\expandafter\def\csname SGAsource@VI.rem1.5.1\endcsname{1.5.1}
\expandafter\def\csname SGAsource@VI.rem1.9.1\endcsname{1.9.1}
\expandafter\def\csname SGAsource@VI.rem2.2.1\endcsname{2.2.1}
\expandafter\def\csname SGAsource@VI.rem2.8.1\endcsname{2.8.1}
\expandafter\def\csname SGAsource@VI.rem3.9\endcsname{3.9}
\expandafter\def\csname SGAsource@VI.rem4.13\endcsname{4.13}
\expandafter\def\csname SGAsource@VI.rem4.15\endcsname{4.15}
\expandafter\def\csname SGAsource@VI.rem4.6.1\endcsname{4.6.1}
\expandafter\def\csname SGAsource@VI.rem7.4.11.1\endcsname{7.4.11.1}
\expandafter\def\csname SGAsource@VI.rem7.4.3\endcsname{7.4.3}
\expandafter\def\csname SGAsource@VI.rem7.4.8\endcsname{7.4.8}
\expandafter\def\csname SGAsource@VI.rem8.2.4\endcsname{8.2.4}
\expandafter\def\csname SGAsource@VI.rem8.5.11\endcsname{8.5.11}
\expandafter\def\csname SGAsource@VI.rem8.7.8\endcsname{8.7.8}
\expandafter\def\csname SGAsource@VI.sec0\endcsname{0}
\expandafter\def\csname SGAsource@VI.sec1\endcsname{1}
\expandafter\def\csname SGAsource@VI.sec2\endcsname{2}
\expandafter\def\csname SGAsource@VI.sec3\endcsname{3}
\expandafter\def\csname SGAsource@VI.sec4\endcsname{4}
\expandafter\def\csname SGAsource@VI.sec5\endcsname{5}
\expandafter\def\csname SGAsource@VI.sec6\endcsname{6}
\expandafter\def\csname SGAsource@VI.sec7\endcsname{7}
\expandafter\def\csname SGAsource@VI.sec8\endcsname{8}
\expandafter\def\csname SGAsource@VI.sec9\endcsname{9}
\expandafter\def\csname SGAsource@VI.subsec1.21\endcsname{1.21}
\expandafter\def\csname SGAsource@VI.subsec1.22\endcsname{1.22}
\expandafter\def\csname SGAsource@VI.subsec1.31\endcsname{1.31}
\expandafter\def\csname SGAsource@VI.subsec1.32\endcsname{1.32}
\expandafter\def\csname SGAsource@VI.subsec2.12\endcsname{2.12}
\expandafter\def\csname SGAsource@VI.subsec2.13\endcsname{2.13}
\expandafter\def\csname SGAsource@VI.subsec2.14\endcsname{2.14}
\expandafter\def\csname SGAsource@VI.subsec2.4\endcsname{2.4}
\expandafter\def\csname SGAsource@VI.subsec4.1\endcsname{4.1}
\expandafter\def\csname SGAsource@VI.subsec4.12\endcsname{4.12}
\expandafter\def\csname SGAsource@VI.subsec4.9\endcsname{4.9}
\expandafter\def\csname SGAsource@VI.subsec5.9\endcsname{5.9}
\expandafter\def\csname SGAsource@VI.subsec6.0\endcsname{6.0}
\expandafter\def\csname SGAsource@VI.subsec6.10\endcsname{6.10}
\expandafter\def\csname SGAsource@VI.subsec6.11\endcsname{6.11}
\expandafter\def\csname SGAsource@VI.subsec7.1\endcsname{7.1}
\expandafter\def\csname SGAsource@VI.subsec7.2\endcsname{7.2}
\expandafter\def\csname SGAsource@VI.subsec7.3\endcsname{7.3}
\expandafter\def\csname SGAsource@VI.subsec7.4\endcsname{7.4}
\expandafter\def\csname SGAsource@VI.subsec8.1\endcsname{8.1}
\expandafter\def\csname SGAsource@VI.subsec8.2\endcsname{8.2}
\expandafter\def\csname SGAsource@VI.subsec8.3\endcsname{8.3}
\expandafter\def\csname SGAsource@VI.subsec8.4\endcsname{8.4}
\expandafter\def\csname SGAsource@VI.subsec8.5\endcsname{8.5}
\expandafter\def\csname SGAsource@VI.subsec8.6\endcsname{8.6}
\expandafter\def\csname SGAsource@VI.subsec8.7\endcsname{8.7}
\expandafter\def\csname SGAsource@VI.subsubsec1.20.2\endcsname{1.20.2}
\expandafter\def\csname SGAsource@VI.subsubsec1.22.1\endcsname{1.22.1}
\expandafter\def\csname SGAsource@VI.subsubsec1.30.1\endcsname{1.30.1}
\expandafter\def\csname SGAsource@VI.subsubsec1.5.2\endcsname{1.5.2}
\expandafter\def\csname SGAsource@VI.subsubsec1.5.3\endcsname{1.5.3}
\expandafter\def\csname SGAsource@VI.subsubsec1.6.2\endcsname{1.6.2}
\expandafter\def\csname SGAsource@VI.subsubsec1.7.1\endcsname{1.7.1}
\expandafter\def\csname SGAsource@VI.subsubsec1.9.2\endcsname{1.9.2}
\expandafter\def\csname SGAsource@VI.subsubsec1.9.3\endcsname{1.9.3}
\expandafter\def\csname SGAsource@VI.subsubsec1.9.4\endcsname{1.9.4}
\expandafter\def\csname SGAsource@VI.subsubsec2.15.1\endcsname{2.15.1}
\expandafter\def\csname SGAsource@VI.subsubsec2.3.1\endcsname{2.3.1}
\expandafter\def\csname SGAsource@VI.subsubsec2.4.1\endcsname{2.4.1}
\expandafter\def\csname SGAsource@VI.subsubsec2.4.2\endcsname{2.4.2}
\expandafter\def\csname SGAsource@VI.subsubsec2.4.3\endcsname{2.4.3}
\expandafter\def\csname SGAsource@VI.subsubsec2.4.4\endcsname{2.4.4}
\expandafter\def\csname SGAsource@VI.subsubsec2.4.5\endcsname{2.4.5}
\expandafter\def\csname SGAsource@VI.subsubsec2.4.6\endcsname{2.4.6}
\expandafter\def\csname SGAsource@VI.subsubsec2.4.7\endcsname{2.4.7}
\expandafter\def\csname SGAsource@VI.subsubsec2.9.2\endcsname{2.9.2}
\expandafter\def\csname SGAsource@VI.subsubsec3.1.1\endcsname{3.1.1}
\expandafter\def\csname SGAsource@VI.subsubsec3.7.1\endcsname{3.7.1}
\expandafter\def\csname SGAsource@VI.subsubsec4.9.10\endcsname{4.9.10}
\expandafter\def\csname SGAsource@VI.subsubsec4.9.3\endcsname{4.9.3}
\expandafter\def\csname SGAsource@VI.subsubsec6.1.1\endcsname{6.1.1}
\expandafter\def\csname SGAsource@VI.subsubsec6.1.2\endcsname{6.1.2}
\expandafter\def\csname SGAsource@VI.subsubsec6.1.3\endcsname{6.1.3}
\expandafter\def\csname SGAsource@VI.subsubsec6.1.4\endcsname{6.1.4}
\expandafter\def\csname SGAsource@VI.subsubsec6.4.0\endcsname{6.4.0}
\expandafter\def\csname SGAsource@VI.subsubsec7.1.2\endcsname{7.1.2}
\expandafter\def\csname SGAsource@VI.subsubsec7.1.3\endcsname{7.1.3}
\expandafter\def\csname SGAsource@VI.subsubsec7.1.4\endcsname{7.1.4}
\expandafter\def\csname SGAsource@VI.subsubsec7.1.6\endcsname{7.1.6}
\expandafter\def\csname SGAsource@VI.subsubsec7.1.7\endcsname{7.1.7}
\expandafter\def\csname SGAsource@VI.subsubsec7.1.8\endcsname{7.1.8}
\expandafter\def\csname SGAsource@VI.subsubsec7.1.9\endcsname{7.1.9}
\expandafter\def\csname SGAsource@VI.subsubsec7.2.1\endcsname{7.2.1}
\expandafter\def\csname SGAsource@VI.subsubsec7.2.2\endcsname{7.2.2}
\expandafter\def\csname SGAsource@VI.subsubsec7.2.3\endcsname{7.2.3}
\expandafter\def\csname SGAsource@VI.subsubsec7.2.4\endcsname{7.2.4}
\expandafter\def\csname SGAsource@VI.subsubsec7.2.5\endcsname{7.2.5}
\expandafter\def\csname SGAsource@VI.subsubsec7.2.6\endcsname{7.2.6}
\expandafter\def\csname SGAsource@VI.subsubsec7.4.5\endcsname{7.4.5}
\expandafter\def\csname SGAsource@VI.subsubsec7.4.6\endcsname{7.4.6}
\expandafter\def\csname SGAsource@VI.subsubsec7.4.9\endcsname{7.4.9}
\expandafter\def\csname SGAsource@VI.subsubsec8.1.2\endcsname{8.1.2}
\expandafter\def\csname SGAsource@VI.subsubsec8.1.3\endcsname{8.1.3}
\expandafter\def\csname SGAsource@VI.subsubsec8.1.4\endcsname{8.1.4}
\expandafter\def\csname SGAsource@VI.subsubsec8.1.5\endcsname{8.1.5}
\expandafter\def\csname SGAsource@VI.subsubsec8.1.6\endcsname{8.1.6}
\expandafter\def\csname SGAsource@VI.subsubsec8.2.10\endcsname{8.2.10}
\expandafter\def\csname SGAsource@VI.subsubsec8.2.11\endcsname{8.2.11}
\expandafter\def\csname SGAsource@VI.subsubsec8.2.12\endcsname{8.2.12}
\expandafter\def\csname SGAsource@VI.subsubsec8.2.6\endcsname{8.2.6}
\expandafter\def\csname SGAsource@VI.subsubsec8.2.7\endcsname{8.2.7}
\expandafter\def\csname SGAsource@VI.subsubsec8.2.8\endcsname{8.2.8}
\expandafter\def\csname SGAsource@VI.subsubsec8.3.1\endcsname{8.3.1}
\expandafter\def\csname SGAsource@VI.subsubsec8.3.10\endcsname{8.3.10}
\expandafter\def\csname SGAsource@VI.subsubsec8.3.12\endcsname{8.3.12}
\expandafter\def\csname SGAsource@VI.subsubsec8.3.2\endcsname{8.3.2}
\expandafter\def\csname SGAsource@VI.subsubsec8.3.5\endcsname{8.3.5}
\expandafter\def\csname SGAsource@VI.subsubsec8.3.7\endcsname{8.3.7}
\expandafter\def\csname SGAsource@VI.subsubsec8.3.8\endcsname{8.3.8}
\expandafter\def\csname SGAsource@VI.subsubsec8.3.9\endcsname{8.3.9}
\expandafter\def\csname SGAsource@VI.subsubsec8.4.1\endcsname{8.4.1}
\expandafter\def\csname SGAsource@VI.subsubsec8.4.2\endcsname{8.4.2}
\expandafter\def\csname SGAsource@VI.subsubsec8.4.3\endcsname{8.4.3}
\expandafter\def\csname SGAsource@VI.subsubsec8.4.4\endcsname{8.4.4}
\expandafter\def\csname SGAsource@VI.subsubsec8.4.5\endcsname{8.4.5}
\expandafter\def\csname SGAsource@VI.subsubsec8.4.6\endcsname{8.4.6}
\expandafter\def\csname SGAsource@VI.subsubsec8.4.7\endcsname{8.4.7}
\expandafter\def\csname SGAsource@VI.subsubsec8.5.1\endcsname{8.5.1}
\expandafter\def\csname SGAsource@VI.subsubsec8.6.1\endcsname{8.6.1}
\expandafter\def\csname SGAsource@VI.subsubsec8.6.2\endcsname{8.6.2}
\expandafter\def\csname SGAsource@VI.subsubsec8.6.3\endcsname{8.6.3}
\expandafter\def\csname SGAsource@VI.subsubsec8.6.4\endcsname{8.6.4}
\expandafter\def\csname SGAsource@VI.subsubsec8.7.1\endcsname{8.7.1}
\expandafter\def\csname SGAsource@VI.thm1.23\endcsname{1.23}
\expandafter\def\csname SGAsource@VI.thm5.1\endcsname{5.1}
\expandafter\def\csname SGAsource@VI.thm7.4.13.1\endcsname{7.4.13.1}
\expandafter\def\csname SGAsource@VI.thm8.2.3\endcsname{8.2.3}
\expandafter\def\csname SGAsource@VI.thm8.2.9\endcsname{8.2.9}
\expandafter\def\csname SGAsource@VI.thm8.3.13\endcsname{8.3.13}
\expandafter\def\csname SGAsource@VI.thm8.7.3\endcsname{8.7.3}
\expandafter\def\csname SGAsource@VII.cor3.2\endcsname{3.2}
\expandafter\def\csname SGAsource@VII.cor4.4\endcsname{4.4}
\expandafter\def\csname SGAsource@VII.cor5.10\endcsname{5.10}
\expandafter\def\csname SGAsource@VII.cor5.11\endcsname{5.11}
\expandafter\def\csname SGAsource@VII.cor5.12\endcsname{5.12}
\expandafter\def\csname SGAsource@VII.cor5.13\endcsname{5.13}
\expandafter\def\csname SGAsource@VII.cor5.8\endcsname{5.8}
\expandafter\def\csname SGAsource@VII.cor5.9\endcsname{5.9}
\expandafter\def\csname SGAsource@VII.lem5.6\endcsname{5.6}
\expandafter\def\csname SGAsource@VII.prop3.1\endcsname{3.1}
\expandafter\def\csname SGAsource@VII.prop3.3\endcsname{3.3}
\expandafter\def\csname SGAsource@VII.prop4.1\endcsname{4.1}
\expandafter\def\csname SGAsource@VII.prop4.3\endcsname{4.3}
\expandafter\def\csname SGAsource@VII.rem2.1\endcsname{2.1}
\expandafter\def\csname SGAsource@VII.rem4.5\endcsname{4.5}
\expandafter\def\csname SGAsource@VII.rem5.14\endcsname{5.14}
\expandafter\def\csname SGAsource@VII.sec1\endcsname{1}
\expandafter\def\csname SGAsource@VII.sec2\endcsname{2}
\expandafter\def\csname SGAsource@VII.sec3\endcsname{3}
\expandafter\def\csname SGAsource@VII.sec4\endcsname{4}
\expandafter\def\csname SGAsource@VII.sec5\endcsname{5}
\expandafter\def\csname SGAsource@VII.subsec1.1\endcsname{1.1}
\expandafter\def\csname SGAsource@VII.subsec1.2\endcsname{1.2}
\expandafter\def\csname SGAsource@VII.subsec1.3\endcsname{1.3}
\expandafter\def\csname SGAsource@VII.subsec1.4\endcsname{1.4}
\expandafter\def\csname SGAsource@VII.subsec1.5\endcsname{1.5}
\expandafter\def\csname SGAsource@VII.subsec1.6\endcsname{1.6}
\expandafter\def\csname SGAsource@VII.subsec1.7\endcsname{1.7}
\expandafter\def\csname SGAsource@VII.subsec4.0\endcsname{4.0}
\expandafter\def\csname SGAsource@VII.subsec4.2\endcsname{4.2}
\expandafter\def\csname SGAsource@VII.subsec5.1\endcsname{5.1}
\expandafter\def\csname SGAsource@VII.subsec5.2\endcsname{5.2}
\expandafter\def\csname SGAsource@VII.subsec5.3\endcsname{5.3}
\expandafter\def\csname SGAsource@VII.subsec5.4\endcsname{5.4}
\expandafter\def\csname SGAsource@VII.subsec5.5\endcsname{5.5}
\expandafter\def\csname SGAsource@VII.subsubsec4.2.1\endcsname{4.2.1}
\expandafter\def\csname SGAsource@VII.subsubsec4.2.2\endcsname{4.2.2}
\expandafter\def\csname SGAsource@VII.thm5.7\endcsname{5.7}
\expandafter\def\csname SGAsource@VIII.cor1.2\endcsname{1.2}
\expandafter\def\csname SGAsource@VIII.cor2.2\endcsname{2.2}
\expandafter\def\csname SGAsource@VIII.cor2.3\endcsname{2.3}
\expandafter\def\csname SGAsource@VIII.cor2.4\endcsname{2.4}
\expandafter\def\csname SGAsource@VIII.cor2.5\endcsname{2.5}
\expandafter\def\csname SGAsource@VIII.cor3.6\endcsname{3.6}
\expandafter\def\csname SGAsource@VIII.cor3.7\endcsname{3.7}
\expandafter\def\csname SGAsource@VIII.cor3.8\endcsname{3.8}
\expandafter\def\csname SGAsource@VIII.cor4.7\endcsname{4.7}
\expandafter\def\csname SGAsource@VIII.cor4.8\endcsname{4.8}
\expandafter\def\csname SGAsource@VIII.cor5.6\endcsname{5.6}
\expandafter\def\csname SGAsource@VIII.cor5.7\endcsname{5.7}
\expandafter\def\csname SGAsource@VIII.cor6.2\endcsname{6.2}
\expandafter\def\csname SGAsource@VIII.cor6.4\endcsname{6.4}
\expandafter\def\csname SGAsource@VIII.cor7.5\endcsname{7.5}
\expandafter\def\csname SGAsource@VIII.cor7.6\endcsname{7.6}
\expandafter\def\csname SGAsource@VIII.cor8.5\endcsname{8.5}
\expandafter\def\csname SGAsource@VIII.cor8.6\endcsname{8.6}
\expandafter\def\csname SGAsource@VIII.cor9.2\endcsname{9.2}
\expandafter\def\csname SGAsource@VIII.def3.3\endcsname{3.3}
\expandafter\def\csname SGAsource@VIII.def4.2\endcsname{4.2}
\expandafter\def\csname SGAsource@VIII.def4.3\endcsname{4.3}
\expandafter\def\csname SGAsource@VIII.def7.2\endcsname{7.2}
\expandafter\def\csname SGAsource@VIII.eq5.1.1\endcsname{5.1.1}
\expandafter\def\csname SGAsource@VIII.eq5.1.2\endcsname{5.1.2}
\expandafter\def\csname SGAsource@VIII.eq5.4.1\endcsname{5.4.1}
\expandafter\def\csname SGAsource@VIII.eq5.4.2\endcsname{5.4.2}
\expandafter\def\csname SGAsource@VIII.exem1.3\endcsname{1.3}
\expandafter\def\csname SGAsource@VIII.lem7.3\endcsname{7.3}
\expandafter\def\csname SGAsource@VIII.lem8.2\endcsname{8.2}
\expandafter\def\csname SGAsource@VIII.lem9.3\endcsname{9.3}
\expandafter\def\csname SGAsource@VIII.prop2.1\endcsname{2.1}
\expandafter\def\csname SGAsource@VIII.prop3.9\endcsname{3.9}
\expandafter\def\csname SGAsource@VIII.prop4.6\endcsname{4.6}
\expandafter\def\csname SGAsource@VIII.prop5.5\endcsname{5.5}
\expandafter\def\csname SGAsource@VIII.prop6.1\endcsname{6.1}
\expandafter\def\csname SGAsource@VIII.prop7.4\endcsname{7.4}
\expandafter\def\csname SGAsource@VIII.prop8.1\endcsname{8.1}
\expandafter\def\csname SGAsource@VIII.prop8.4\endcsname{8.4}
\expandafter\def\csname SGAsource@VIII.prop9.1\endcsname{9.1}
\expandafter\def\csname SGAsource@VIII.rem1.4\endcsname{1.4}
\expandafter\def\csname SGAsource@VIII.rem3.13\endcsname{3.13}
\expandafter\def\csname SGAsource@VIII.rem3.2\endcsname{3.2}
\expandafter\def\csname SGAsource@VIII.rem5.3\endcsname{5.3}
\expandafter\def\csname SGAsource@VIII.rem5.8\endcsname{5.8}
\expandafter\def\csname SGAsource@VIII.rem8.3\endcsname{8.3}
\expandafter\def\csname SGAsource@VIII.sch3.10\endcsname{3.10}
\expandafter\def\csname SGAsource@VIII.sec1\endcsname{1}
\expandafter\def\csname SGAsource@VIII.sec2\endcsname{2}
\expandafter\def\csname SGAsource@VIII.sec3\endcsname{3}
\expandafter\def\csname SGAsource@VIII.sec4\endcsname{4}
\expandafter\def\csname SGAsource@VIII.sec5\endcsname{5}
\expandafter\def\csname SGAsource@VIII.sec6\endcsname{6}
\expandafter\def\csname SGAsource@VIII.sec7\endcsname{7}
\expandafter\def\csname SGAsource@VIII.sec8\endcsname{8}
\expandafter\def\csname SGAsource@VIII.sec9\endcsname{9}
\expandafter\def\csname SGAsource@VIII.subsec3.1\endcsname{3.1}
\expandafter\def\csname SGAsource@VIII.subsec3.11\endcsname{3.11}
\expandafter\def\csname SGAsource@VIII.subsec3.12\endcsname{3.12}
\expandafter\def\csname SGAsource@VIII.subsec3.4\endcsname{3.4}
\expandafter\def\csname SGAsource@VIII.subsec4.1\endcsname{4.1}
\expandafter\def\csname SGAsource@VIII.subsec4.4\endcsname{4.4}
\expandafter\def\csname SGAsource@VIII.subsec4.5\endcsname{4.5}
\expandafter\def\csname SGAsource@VIII.subsec5.1\endcsname{5.1}
\expandafter\def\csname SGAsource@VIII.subsec5.4\endcsname{5.4}
\expandafter\def\csname SGAsource@VIII.subsec6.5\endcsname{6.5}
\expandafter\def\csname SGAsource@VIII.subsec6.6\endcsname{6.6}
\expandafter\def\csname SGAsource@VIII.subsec7.1\endcsname{7.1}
\expandafter\def\csname SGAsource@VIII.subsec7.7\endcsname{7.7}
\expandafter\def\csname SGAsource@VIII.subsec7.8\endcsname{7.8}
\expandafter\def\csname SGAsource@VIII.subsubsec7.9.1\endcsname{7.9.1}
\expandafter\def\csname SGAsource@VIII.subsubsec7.9.2\endcsname{7.9.2}
\expandafter\def\csname SGAsource@VIII.subsubsec7.9.3\endcsname{7.9.3}
\expandafter\def\csname SGAsource@VIII.subsubsec9.4.1\endcsname{9.4.1}
\expandafter\def\csname SGAsource@VIII.subsubsec9.4.2\endcsname{9.4.2}
\expandafter\def\csname SGAsource@VIII.subsubsec9.4.3\endcsname{9.4.3}
\expandafter\def\csname SGAsource@VIII.subsubsec9.4.4\endcsname{9.4.4}
\expandafter\def\csname SGAsource@VIII.thm1.1\endcsname{1.1}
\expandafter\def\csname SGAsource@VIII.thm3.5\endcsname{3.5}
\expandafter\def\csname SGAsource@VIII.thm5.2\endcsname{5.2}
\expandafter\def\csname SGAsource@VIII.thm6.3\endcsname{6.3}
\expandafter\def\csname SGAsource@VIII.thm7.9\endcsname{7.9}
\expandafter\def\csname SGAsource@VIII.thm9.4\endcsname{9.4}
\expandafter\def\csname SGAsource@Vbis.cor1.2.11\endcsname{1.2.11}
\expandafter\def\csname SGAsource@Vbis.cor1.3.11\endcsname{1.3.11}
\expandafter\def\csname SGAsource@Vbis.cor1.3.12\endcsname{1.3.12}
\expandafter\def\csname SGAsource@Vbis.cor2.3.6\endcsname{2.3.6}
\expandafter\def\csname SGAsource@Vbis.cor2.3.7\endcsname{2.3.7}
\expandafter\def\csname SGAsource@Vbis.cor2.4.14\endcsname{2.4.14}
\expandafter\def\csname SGAsource@Vbis.cor3.1.7\endcsname{3.1.7}
\expandafter\def\csname SGAsource@Vbis.cor3.3.5\endcsname{3.3.5}
\expandafter\def\csname SGAsource@Vbis.cor4.1.2\endcsname{4.1.2}
\expandafter\def\csname SGAsource@Vbis.cor4.1.6\endcsname{4.1.6}
\expandafter\def\csname SGAsource@Vbis.cor4.1.7\endcsname{4.1.7}
\expandafter\def\csname SGAsource@Vbis.cor4.1.9\endcsname{4.1.9}
\expandafter\def\csname SGAsource@Vbis.cor4.3.5\endcsname{4.3.5}
\expandafter\def\csname SGAsource@Vbis.def1.2.1\endcsname{1.2.1}
\expandafter\def\csname SGAsource@Vbis.def1.2.14\endcsname{1.2.14}
\expandafter\def\csname SGAsource@Vbis.def1.2.2\endcsname{1.2.2}
\expandafter\def\csname SGAsource@Vbis.def1.3.1\endcsname{1.3.1}
\expandafter\def\csname SGAsource@Vbis.def1.3.2\endcsname{1.3.2}
\expandafter\def\csname SGAsource@Vbis.def2.1.2\endcsname{2.1.2}
\expandafter\def\csname SGAsource@Vbis.def2.2.2\endcsname{2.2.2}
\expandafter\def\csname SGAsource@Vbis.def2.2.4\endcsname{2.2.4}
\expandafter\def\csname SGAsource@Vbis.def2.2.6\endcsname{2.2.6}
\expandafter\def\csname SGAsource@Vbis.def2.4.0.1\endcsname{2.4.0.1}
\expandafter\def\csname SGAsource@Vbis.def2.4.10\endcsname{2.4.10}
\expandafter\def\csname SGAsource@Vbis.def2.4.11\endcsname{2.4.11}
\expandafter\def\csname SGAsource@Vbis.def2.4.15\endcsname{2.4.15}
\expandafter\def\csname SGAsource@Vbis.def2.4.16\endcsname{2.4.16}
\expandafter\def\csname SGAsource@Vbis.def2.4.3.1\endcsname{2.4.3.1}
\expandafter\def\csname SGAsource@Vbis.def2.4.8\endcsname{2.4.8}
\expandafter\def\csname SGAsource@Vbis.def3.1.2\endcsname{3.1.2}
\expandafter\def\csname SGAsource@Vbis.def3.2.2.4\endcsname{3.2.2.4}
\expandafter\def\csname SGAsource@Vbis.def5.1.1\endcsname{5.1.1}
\expandafter\def\csname SGAsource@Vbis.def5.1.5\endcsname{5.1.5}
\expandafter\def\csname SGAsource@Vbis.def5.2.6\endcsname{5.2.6}
\expandafter\def\csname SGAsource@Vbis.eq1.1\endcsname{1.1}
\expandafter\def\csname SGAsource@Vbis.eq1.2\endcsname{1.2}
\expandafter\def\csname SGAsource@Vbis.eq1.3\endcsname{1.3}
\expandafter\def\csname SGAsource@Vbis.eq1.3.10.2\endcsname{1.3.10.2}
\expandafter\def\csname SGAsource@Vbis.eq1.3.10.3\endcsname{1.3.10.3}
\expandafter\def\csname SGAsource@Vbis.eq1.3.10.4\endcsname{1.3.10.4}
\expandafter\def\csname SGAsource@Vbis.eq2.1\endcsname{2.1}
\expandafter\def\csname SGAsource@Vbis.eq2.2\endcsname{2.2}
\expandafter\def\csname SGAsource@Vbis.eq2.3\endcsname{2.3}
\expandafter\def\csname SGAsource@Vbis.eq2.3.10.2\endcsname{2.3.10.2}
\expandafter\def\csname SGAsource@Vbis.eq2.3.2.1\endcsname{2.3.2.1}
\expandafter\def\csname SGAsource@Vbis.eq2.3.3.1\endcsname{2.3.3.1}
\expandafter\def\csname SGAsource@Vbis.eq2.3.3.2\endcsname{2.3.3.2}
\expandafter\def\csname SGAsource@Vbis.eq2.3.4.1\endcsname{2.3.4.1}
\expandafter\def\csname SGAsource@Vbis.eq2.3.4.2\endcsname{2.3.4.2}
\expandafter\def\csname SGAsource@Vbis.eq2.3.4.3\endcsname{2.3.4.3}
\expandafter\def\csname SGAsource@Vbis.eq2.3.5.1\endcsname{2.3.5.1}
\expandafter\def\csname SGAsource@Vbis.eq2.3.5.2\endcsname{2.3.5.2}
\expandafter\def\csname SGAsource@Vbis.eq2.3.9.1\endcsname{2.3.9.1}
\expandafter\def\csname SGAsource@Vbis.eq2.4\endcsname{2.4}
\expandafter\def\csname SGAsource@Vbis.eq2.4.4.1\endcsname{2.4.4.1}
\expandafter\def\csname SGAsource@Vbis.eq2.5.1.1\endcsname{2.5.1.1}
\expandafter\def\csname SGAsource@Vbis.eq2.5.1.2\endcsname{2.5.1.2}
\expandafter\def\csname SGAsource@Vbis.eq2.5.2.1\endcsname{2.5.2.1}
\expandafter\def\csname SGAsource@Vbis.eq2.5.3.1\endcsname{2.5.3.1}
\expandafter\def\csname SGAsource@Vbis.eq2.5.3.2\endcsname{2.5.3.2}
\expandafter\def\csname SGAsource@Vbis.eq2.5.3.3\endcsname{2.5.3.3}
\expandafter\def\csname SGAsource@Vbis.eq2.5.5.1\endcsname{2.5.5.1}
\expandafter\def\csname SGAsource@Vbis.eq2.5.5.2\endcsname{2.5.5.2}
\expandafter\def\csname SGAsource@Vbis.eq2.5.5.3\endcsname{2.5.5.3}
\expandafter\def\csname SGAsource@Vbis.eq3.0.1.1\endcsname{3.0.1.1}
\expandafter\def\csname SGAsource@Vbis.eq3.0.1.2\endcsname{3.0.1.2}
\expandafter\def\csname SGAsource@Vbis.eq3.0.2.1\endcsname{3.0.2.1}
\expandafter\def\csname SGAsource@Vbis.eq3.1.1.1\endcsname{3.1.1.1}
\expandafter\def\csname SGAsource@Vbis.eq3.1.1.2\endcsname{3.1.1.2}
\expandafter\def\csname SGAsource@Vbis.eq3.1.1.3\endcsname{3.1.1.3}
\expandafter\def\csname SGAsource@Vbis.eq3.1.1.4\endcsname{3.1.1.4}
\expandafter\def\csname SGAsource@Vbis.eq3.1.1.5\endcsname{3.1.1.5}
\expandafter\def\csname SGAsource@Vbis.eq3.1.1.6\endcsname{3.1.1.6}
\expandafter\def\csname SGAsource@Vbis.eq3.1.1.7\endcsname{3.1.1.7}
\expandafter\def\csname SGAsource@Vbis.eq3.1.1.8\endcsname{3.1.1.8}
\expandafter\def\csname SGAsource@Vbis.eq3.2\endcsname{3.2}
\expandafter\def\csname SGAsource@Vbis.eq3.2.2.1\endcsname{3.2.2.1}
\expandafter\def\csname SGAsource@Vbis.eq3.2.2.2\endcsname{3.2.2.2}
\expandafter\def\csname SGAsource@Vbis.eq3.2.2.3\endcsname{3.2.2.3}
\expandafter\def\csname SGAsource@Vbis.eq3.2.3.1\endcsname{3.2.3.1}
\expandafter\def\csname SGAsource@Vbis.eq3.2.3.2\endcsname{3.2.3.2}
\expandafter\def\csname SGAsource@Vbis.eq3.3\endcsname{3.3}
\expandafter\def\csname SGAsource@Vbis.eq4.1\endcsname{4.1}
\expandafter\def\csname SGAsource@Vbis.eq4.1.3.1\endcsname{4.1.3.1}
\expandafter\def\csname SGAsource@Vbis.lem1.2.10\endcsname{1.2.10}
\expandafter\def\csname SGAsource@Vbis.lem1.2.10.1\endcsname{1.2.10.1}
\expandafter\def\csname SGAsource@Vbis.lem1.2.16\endcsname{1.2.16}
\expandafter\def\csname SGAsource@Vbis.lem1.3.10.1\endcsname{1.3.10.1}
\expandafter\def\csname SGAsource@Vbis.lem1.3.10.5\endcsname{1.3.10.5}
\expandafter\def\csname SGAsource@Vbis.lem2.2.3.1\endcsname{2.2.3.1}
\expandafter\def\csname SGAsource@Vbis.lem2.3.2.2\endcsname{2.3.2.2}
\expandafter\def\csname SGAsource@Vbis.lem2.4.3.2\endcsname{2.4.3.2}
\expandafter\def\csname SGAsource@Vbis.lem2.4.4.2\endcsname{2.4.4.2}
\expandafter\def\csname SGAsource@Vbis.lem2.4.5.1\endcsname{2.4.5.1}
\expandafter\def\csname SGAsource@Vbis.lem2.4.7.1\endcsname{2.4.7.1}
\expandafter\def\csname SGAsource@Vbis.lem3.0.2.3\endcsname{3.0.2.3}
\expandafter\def\csname SGAsource@Vbis.lem3.0.2.4\endcsname{3.0.2.4}
\expandafter\def\csname SGAsource@Vbis.lem3.1.1.9\endcsname{3.1.1.9}
\expandafter\def\csname SGAsource@Vbis.lem3.1.5.1\endcsname{3.1.5.1}
\expandafter\def\csname SGAsource@Vbis.lem3.2.3.3\endcsname{3.2.3.3}
\expandafter\def\csname SGAsource@Vbis.lem3.3.3.1\endcsname{3.3.3.1}
\expandafter\def\csname SGAsource@Vbis.lem3.3.3.2\endcsname{3.3.3.2}
\expandafter\def\csname SGAsource@Vbis.lem3.3.3.3\endcsname{3.3.3.3}
\expandafter\def\csname SGAsource@Vbis.lem3.3.4.1\endcsname{3.3.4.1}
\expandafter\def\csname SGAsource@Vbis.lem3.3.4.2\endcsname{3.3.4.2}
\expandafter\def\csname SGAsource@Vbis.lem4.1.3\endcsname{4.1.3}
\expandafter\def\csname SGAsource@Vbis.lem4.1.4\endcsname{4.1.4}
\expandafter\def\csname SGAsource@Vbis.lem4.1.5\endcsname{4.1.5}
\expandafter\def\csname SGAsource@Vbis.lem5.3.2\endcsname{5.3.2}
\expandafter\def\csname SGAsource@Vbis.par1.3.6.1\endcsname{1.3.6.1}
\expandafter\def\csname SGAsource@Vbis.par2.3.1.1\endcsname{2.3.1.1}
\expandafter\def\csname SGAsource@Vbis.par2.3.1.2\endcsname{2.3.1.2}
\expandafter\def\csname SGAsource@Vbis.par2.3.1.3\endcsname{2.3.1.3}
\expandafter\def\csname SGAsource@Vbis.par2.4.1.1\endcsname{2.4.1.1}
\expandafter\def\csname SGAsource@Vbis.par3.0.2.2\endcsname{3.0.2.2}
\expandafter\def\csname SGAsource@Vbis.prop1.2.12\endcsname{1.2.12}
\expandafter\def\csname SGAsource@Vbis.prop1.2.15\endcsname{1.2.15}
\expandafter\def\csname SGAsource@Vbis.prop1.2.9\endcsname{1.2.9}
\expandafter\def\csname SGAsource@Vbis.prop1.3.10\endcsname{1.3.10}
\expandafter\def\csname SGAsource@Vbis.prop1.3.7\endcsname{1.3.7}
\expandafter\def\csname SGAsource@Vbis.prop2.2.7\endcsname{2.2.7}
\expandafter\def\csname SGAsource@Vbis.prop2.3.5\endcsname{2.3.5}
\expandafter\def\csname SGAsource@Vbis.prop2.4.13\endcsname{2.4.13}
\expandafter\def\csname SGAsource@Vbis.prop2.4.6\endcsname{2.4.6}
\expandafter\def\csname SGAsource@Vbis.prop2.5.4\endcsname{2.5.4}
\expandafter\def\csname SGAsource@Vbis.prop2.5.5\endcsname{2.5.5}
\expandafter\def\csname SGAsource@Vbis.prop2.5.6\endcsname{2.5.6}
\expandafter\def\csname SGAsource@Vbis.prop2.5.7\endcsname{2.5.7}
\expandafter\def\csname SGAsource@Vbis.prop3.1.1\endcsname{3.1.1}
\expandafter\def\csname SGAsource@Vbis.prop3.1.4\endcsname{3.1.4}
\expandafter\def\csname SGAsource@Vbis.prop3.1.5\endcsname{3.1.5}
\expandafter\def\csname SGAsource@Vbis.prop3.1.6\endcsname{3.1.6}
\expandafter\def\csname SGAsource@Vbis.prop3.2.1\endcsname{3.2.1}
\expandafter\def\csname SGAsource@Vbis.prop3.2.4\endcsname{3.2.4}
\expandafter\def\csname SGAsource@Vbis.prop3.3.1\endcsname{3.3.1}
\expandafter\def\csname SGAsource@Vbis.prop3.3.4\endcsname{3.3.4}
\expandafter\def\csname SGAsource@Vbis.prop4.1.8\endcsname{4.1.8}
\expandafter\def\csname SGAsource@Vbis.prop4.2.1\endcsname{4.2.1}
\expandafter\def\csname SGAsource@Vbis.prop4.3.2\endcsname{4.3.2}
\expandafter\def\csname SGAsource@Vbis.prop4.3.3\endcsname{4.3.3}
\expandafter\def\csname SGAsource@Vbis.prop5.1.3\endcsname{5.1.3}
\expandafter\def\csname SGAsource@Vbis.prop5.1.7\endcsname{5.1.7}
\expandafter\def\csname SGAsource@Vbis.prop5.1.8\endcsname{5.1.8}
\expandafter\def\csname SGAsource@Vbis.prop5.2.3\endcsname{5.2.3}
\expandafter\def\csname SGAsource@Vbis.prop5.2.5\endcsname{5.2.5}
\expandafter\def\csname SGAsource@Vbis.prop5.2.8\endcsname{5.2.8}
\expandafter\def\csname SGAsource@Vbis.prop5.3.4\endcsname{5.3.4}
\expandafter\def\csname SGAsource@Vbis.rem1.3.13\endcsname{1.3.13}
\expandafter\def\csname SGAsource@Vbis.rem1.3.3\endcsname{1.3.3}
\expandafter\def\csname SGAsource@Vbis.rem2.2.5\endcsname{2.2.5}
\expandafter\def\csname SGAsource@Vbis.rem2.3.8\endcsname{2.3.8}
\expandafter\def\csname SGAsource@Vbis.rem2.4.9\endcsname{2.4.9}
\expandafter\def\csname SGAsource@Vbis.rem4.1.10\endcsname{4.1.10}
\expandafter\def\csname SGAsource@Vbis.rem4.3.4\endcsname{4.3.4}
\expandafter\def\csname SGAsource@Vbis.rem5.2.9\endcsname{5.2.9}
\expandafter\def\csname SGAsource@Vbis.sch1.2.17\endcsname{1.2.17}
\expandafter\def\csname SGAsource@Vbis.sec1\endcsname{1}
\expandafter\def\csname SGAsource@Vbis.sec2\endcsname{2}
\expandafter\def\csname SGAsource@Vbis.sec3\endcsname{3}
\expandafter\def\csname SGAsource@Vbis.sec4\endcsname{4}
\expandafter\def\csname SGAsource@Vbis.sec5\endcsname{5}
\expandafter\def\csname SGAsource@Vbis.subsec1.1\endcsname{1.1}
\expandafter\def\csname SGAsource@Vbis.subsec1.2\endcsname{1.2}
\expandafter\def\csname SGAsource@Vbis.subsec1.3\endcsname{1.3}
\expandafter\def\csname SGAsource@Vbis.subsec2.1\endcsname{2.1}
\expandafter\def\csname SGAsource@Vbis.subsec2.2\endcsname{2.2}
\expandafter\def\csname SGAsource@Vbis.subsec2.3\endcsname{2.3}
\expandafter\def\csname SGAsource@Vbis.subsec2.4\endcsname{2.4}
\expandafter\def\csname SGAsource@Vbis.subsec2.5\endcsname{2.5}
\expandafter\def\csname SGAsource@Vbis.subsec3.0\endcsname{3.0}
\expandafter\def\csname SGAsource@Vbis.subsec3.1\endcsname{3.1}
\expandafter\def\csname SGAsource@Vbis.subsec3.2\endcsname{3.2}
\expandafter\def\csname SGAsource@Vbis.subsec3.3\endcsname{3.3}
\expandafter\def\csname SGAsource@Vbis.subsec4.1\endcsname{4.1}
\expandafter\def\csname SGAsource@Vbis.subsec4.2\endcsname{4.2}
\expandafter\def\csname SGAsource@Vbis.subsec4.3\endcsname{4.3}
\expandafter\def\csname SGAsource@Vbis.subsec5.1\endcsname{5.1}
\expandafter\def\csname SGAsource@Vbis.subsec5.2\endcsname{5.2}
\expandafter\def\csname SGAsource@Vbis.subsec5.3\endcsname{5.3}
\expandafter\def\csname SGAsource@Vbis.subsubsec1.1.1\endcsname{1.1.1}
\expandafter\def\csname SGAsource@Vbis.subsubsec1.1.2\endcsname{1.1.2}
\expandafter\def\csname SGAsource@Vbis.subsubsec1.1.3\endcsname{1.1.3}
\expandafter\def\csname SGAsource@Vbis.subsubsec1.2.13\endcsname{1.2.13}
\expandafter\def\csname SGAsource@Vbis.subsubsec1.2.4\endcsname{1.2.4}
\expandafter\def\csname SGAsource@Vbis.subsubsec1.2.5\endcsname{1.2.5}
\expandafter\def\csname SGAsource@Vbis.subsubsec1.2.6\endcsname{1.2.6}
\expandafter\def\csname SGAsource@Vbis.subsubsec1.2.7\endcsname{1.2.7}
\expandafter\def\csname SGAsource@Vbis.subsubsec1.2.8\endcsname{1.2.8}
\expandafter\def\csname SGAsource@Vbis.subsubsec1.3.4\endcsname{1.3.4}
\expandafter\def\csname SGAsource@Vbis.subsubsec1.3.5\endcsname{1.3.5}
\expandafter\def\csname SGAsource@Vbis.subsubsec1.3.6\endcsname{1.3.6}
\expandafter\def\csname SGAsource@Vbis.subsubsec1.3.8\endcsname{1.3.8}
\expandafter\def\csname SGAsource@Vbis.subsubsec1.3.9\endcsname{1.3.9}
\expandafter\def\csname SGAsource@Vbis.subsubsec2.1.1\endcsname{2.1.1}
\expandafter\def\csname SGAsource@Vbis.subsubsec2.1.3\endcsname{2.1.3}
\expandafter\def\csname SGAsource@Vbis.subsubsec2.2.1\endcsname{2.2.1}
\expandafter\def\csname SGAsource@Vbis.subsubsec2.2.3\endcsname{2.2.3}
\expandafter\def\csname SGAsource@Vbis.subsubsec2.3.1\endcsname{2.3.1}
\expandafter\def\csname SGAsource@Vbis.subsubsec2.3.10\endcsname{2.3.10}
\expandafter\def\csname SGAsource@Vbis.subsubsec2.3.11\endcsname{2.3.11}
\expandafter\def\csname SGAsource@Vbis.subsubsec2.3.2\endcsname{2.3.2}
\expandafter\def\csname SGAsource@Vbis.subsubsec2.3.3\endcsname{2.3.3}
\expandafter\def\csname SGAsource@Vbis.subsubsec2.3.4\endcsname{2.3.4}
\expandafter\def\csname SGAsource@Vbis.subsubsec2.3.9\endcsname{2.3.9}
\expandafter\def\csname SGAsource@Vbis.subsubsec2.4.0\endcsname{2.4.0}
\expandafter\def\csname SGAsource@Vbis.subsubsec2.4.1\endcsname{2.4.1}
\expandafter\def\csname SGAsource@Vbis.subsubsec2.4.3\endcsname{2.4.3}
\expandafter\def\csname SGAsource@Vbis.subsubsec2.4.4\endcsname{2.4.4}
\expandafter\def\csname SGAsource@Vbis.subsubsec2.4.5\endcsname{2.4.5}
\expandafter\def\csname SGAsource@Vbis.subsubsec2.4.7\endcsname{2.4.7}
\expandafter\def\csname SGAsource@Vbis.subsubsec2.5.1\endcsname{2.5.1}
\expandafter\def\csname SGAsource@Vbis.subsubsec2.5.2\endcsname{2.5.2}
\expandafter\def\csname SGAsource@Vbis.subsubsec3.0.0\endcsname{3.0.0}
\expandafter\def\csname SGAsource@Vbis.subsubsec3.0.1\endcsname{3.0.1}
\expandafter\def\csname SGAsource@Vbis.subsubsec3.0.2\endcsname{3.0.2}
\expandafter\def\csname SGAsource@Vbis.subsubsec3.1.3\endcsname{3.1.3}
\expandafter\def\csname SGAsource@Vbis.subsubsec3.2.2\endcsname{3.2.2}
\expandafter\def\csname SGAsource@Vbis.subsubsec3.2.3\endcsname{3.2.3}
\expandafter\def\csname SGAsource@Vbis.subsubsec3.3.2\endcsname{3.3.2}
\expandafter\def\csname SGAsource@Vbis.subsubsec4.1.0\endcsname{4.1.0}
\expandafter\def\csname SGAsource@Vbis.subsubsec4.2.0\endcsname{4.2.0}
\expandafter\def\csname SGAsource@Vbis.subsubsec4.3.0\endcsname{4.3.0}
\expandafter\def\csname SGAsource@Vbis.subsubsec4.3.1\endcsname{4.3.1}
\expandafter\def\csname SGAsource@Vbis.subsubsec5.1.0\endcsname{5.1.0}
\expandafter\def\csname SGAsource@Vbis.subsubsec5.1.2\endcsname{5.1.2}
\expandafter\def\csname SGAsource@Vbis.subsubsec5.1.4\endcsname{5.1.4}
\expandafter\def\csname SGAsource@Vbis.subsubsec5.1.6\endcsname{5.1.6}
\expandafter\def\csname SGAsource@Vbis.subsubsec5.2.0\endcsname{5.2.0}
\expandafter\def\csname SGAsource@Vbis.subsubsec5.2.1\endcsname{5.2.1}
\expandafter\def\csname SGAsource@Vbis.subsubsec5.2.2\endcsname{5.2.2}
\expandafter\def\csname SGAsource@Vbis.subsubsec5.2.4\endcsname{5.2.4}
\expandafter\def\csname SGAsource@Vbis.subsubsec5.2.7\endcsname{5.2.7}
\expandafter\def\csname SGAsource@Vbis.subsubsec5.3.1\endcsname{5.3.1}
\expandafter\def\csname SGAsource@Vbis.subsubsec5.3.3\endcsname{5.3.3}
\expandafter\def\csname SGAsource@Vbis.susubsec2.4.2\endcsname{2.4.2}
\expandafter\def\csname SGAsource@Vbis.thm2.4.12\endcsname{2.4.12}
\expandafter\def\csname SGAsource@Vbis.thm3.3.3\endcsname{3.3.3}
\expandafter\def\csname SGAsource@Vbis.thm4.1.1\endcsname{4.1.1}
\expandafter\def\csname SGAsource@page112\endcsname{112}
\expandafter\def\csname SGAsource@page133\endcsname{133}
\expandafter\def\csname SGAsource@page6\endcsname{6}
\expandafter\def\csname SGAsource@page7\endcsname{7}
\expandafter\def\csname SGAsource@sga4.exp.I\endcsname{I}
\expandafter\def\csname SGAsource@sga4.exp.II\endcsname{II}
\expandafter\def\csname SGAsource@sga4.exp.III\endcsname{III}
\expandafter\def\csname SGAsource@sga4.exp.III.terminology.5.3\endcsname{3}
\expandafter\def\csname SGAsource@sga4.exp.IV\endcsname{IV}
\expandafter\def\csname SGAsource@sga4.exp.IX\endcsname{IX}
\expandafter\def\csname SGAsource@sga4.exp.V\endcsname{V}
\expandafter\def\csname SGAsource@sga4.exp.VI\endcsname{VI}
\expandafter\def\csname SGAsource@sga4.exp.VII\endcsname{VII}
\expandafter\def\csname SGAsource@sga4.exp.VIII\endcsname{VIII}
\expandafter\def\csname SGAsource@sga4.exp.Vbis\endcsname{Vbis}
\expandafter\def\csname SGAsource@sga4[1]\endcsname{sga4[1]}
\let\SGAoriginalhyperref\hyperref
\renewcommand{\hyperref}[2][]{\ifcsname SGAsource@#1\endcsname #2\else\SGAoriginalhyperref[#1]{#2}\fi}
\AtBeginDocument{\let\SGAoriginalref\ref\renewcommand{\ref}[1]{\ifcsname SGAsource@#1\endcsname\csname SGAsource@#1\endcsname\else\SGAoriginalref{#1}\fi}}

% END native input: style/external-references.tex
% Keep editorial note numbers local. The inherited macros let thefootnote
% escape, making later original notes [1], [2], [3] all show the last cnote.
\makeatletter
\renewcommand{\nde}[1]{\ifthenelse{\boolean{original}}{}{%
  \stepcounter{cnote}\begingroup
  \def\thefootnote{\arabic{cnote}}%
  \protect\footnote[\value{cnote}]{编者注: #1}\endgroup}}
\renewcommand{\ndemark}{\ifthenelse{\boolean{original}}{}{%
  \stepcounter{cnote}\begingroup
  \def\thefootnote{\arabic{cnote}}%
  \protect\footnotemark[\value{cnote}]\endgroup}}
\renewcommand{\ndetext}[1]{\ifthenelse{\boolean{original}}{}{%
  \begingroup\def\thefootnote{\arabic{cnote}}%
  \protect\footnotetext[\value{cnote}]{编者注: #1}\endgroup}}
% The two inherited repeated-author-note calls both refer to original note i.
% Explicit note numbers do not advance the ordinary footnote counter.
\renewcommand{\sgarepeatedfootnotemark}[2][1]{%
  \begingroup\c@footnote=#1\relax
  \edef\@thefnmark{\roman{footnote}}%
  \hyperlink{#2}{\@makefnmark}\endgroup}
\makeatother
\renewcommand{\nderef}[1]{第~\ref{#1}~条编者注（第~\pageref{#1}~页）}

% END native input: style/localization.tex
% END native input: style/preamble.tex
\begin{document}
\def\chaptername{第}
\def\contentsname{目录}
\def\bibname{参考文献}
\def\proofname{证明}
\frontmatter
\begin{center}
{\Large 拓扑斯理论与概形的平展上同调\\第一至三卷：第 X--XIX 讲}\\
\end{center}
\vskip1cm
\begin{center}
M.~Artin, A.~Grothendieck, J.-L. Verdier
\end{center}
\vskip1cm
\begin{center}
N.~Bourbaki, P.~Deligne, B.~Saint-Donat 参与
\end{center}
\vskip1cm
\begin{center}\small 中文首译稿；源文疑点与校订状态见随附说明.\end{center}
\clearpage
\tableofcontents
\mainmatter
\setcounter{chapter}{9}
% BEGIN native input: chapters/exp-10-zh.tex
\chapter{上同调维数：初步结果}
%\addtocontents{toc}{M. Artin}
\begin{center}
M. Artin
\end{center}
\phantomsection\MakeLinkTarget*{sga4.exp.X}\label{sga4.exp.X}

\setcounter{pageoriginal}{42}

\section{引言}\etiquette[1]{X.sec1}\pageoriginale

本讲中，除非另有明确说明，“层”（sheaf）均指“阿贝尔层”（abelian sheaf）。

设 $X$ 是概形（scheme），$\ell\in \PP$。“回顾”一下，$X$ 的
\emph{$\ell$-上同调维数}（cohomological dimension）是满足下述条件的最大整数
$\cd_{\ell}X=n$（若不存在这样的数则为 $=\infty$）：$X$ 上至少有一个
$\ell$-挠层（torsion sheaf）$F$ 使 $H^n(X,F)\neq 0$。
若 $X=\Spec A$，则记 $\cd_{\ell} A=\cd_{\ell}\Spec A$。考虑到术语对照（VIII 2），%% label
当 $A=k$ 是域时，这就给出 [CG] 中的“经典”概念。由于显然有公式
$$
\cd_{\ell}(\coprod_i X_i)=\displaystyle{\mathop{\sup}_i} \cd_{\ell}X_i
$$
经典理论显然可以推广到任意 Artin 环。\pageoriginale 这一理论见于 Serre 的
$(CG)=$ “\sisi{Galois cohomology}{Galois cohomology}”。%% Galois?
我们将在第 ${\rm n^{\circ}}$ \hyperref[X.sec2]{2} 节补充若干更精细的结果。

例如，考虑可分闭域（separably closed field）$k$ 上维数为 $n$ 的不可约代数簇
$X$，并固定 $\ell\neq \car k$。第 ${\rm n^{\circ}}$ \hyperref[X.sec4]{4} 节的一个定理
给出其上同调维数的上界 $2n$；这是 $X$ 的“真正”维数（当 $k=\CC$ 时的拓扑类比提示了这一点）。
一旦知道域的理论，这个结果便相当初等；事实上，除 $X$ 完备的情形外，上界 $2n$
还可以改进：另有上界 $2n-1$，只要 $X$“至少缺少一个点”，它就成立。
后者是正确的“稳定上界”，意思是每个维数为 $n$ 的 $X$ 都含有一个上同调维数为
$2n-1$ 的开子集。对仿射簇还可得到第三个上界，此时所得的上同调维数为 $n$。
这是关于超平面截面的 Lefschetz 定理之一的推广。后面这两个断言较深，将留待以后（XIV）证明。

注意，上同调维数不是一个局部概念（这与仿紧空间的情形不同）。例如，剩余域可分闭的
Hensel 环的上同调维数为零，但删去闭点之后，可以得到任意大的数；在最重要的情形下，
若 $n=\dim A$，这个数为 $2n-1$。这也符合由 $U-\{x\}$ 的上同调提供的拓扑类比，
其中 $U$ 是复解析空间中一个点 $x$ 的小开邻域。

\section{域上的辅助结果}\etiquette[2]{X.sec2}\pageoriginale

我们有如下结果。

\begin{enonce*}{定理 2.1}\etiquette[2.1]{X.thm2.1}
设 $k'/k$ 是域 $k$ 的有限型扩张，$\ell\in \PP$。则
$$
\cd_{\ell}k'\leqq \cd_{\ell}k+\deg.\tr.(k'/k){\quoi},
$$
若此不等式严格，且 $\ell\neq \car k$，则 $\ell=2$，$k$ 可序，而 $k'$ 不可序
（\ie $-1$ 在 $k'$ 中是平方和）。
\end{enonce*}

回顾，当 $\ell=\car k$ 时，有 $\cd_{\ell}k'\leqq 1$（CG II
\hyperref[X.thm2.2]{2.2}），因而 \hyperref[X.thm2.1]{2.1} 中的等号问题几乎不会出现。
此外，若 $k$ 可序，则 $\cd_2k=\infty$。事实上，$k$ 的代数闭包 $\bar{k}$ 含有一个子扩张
$\widetilde{k}$（$\bar{k}$ 的极大有序子域），使
$[\bar{k}:\widetilde{k}]=2$。显然，$\cd_2\widetilde{k}=\infty$，
而由（CG I 命题~14）有 $\cd_2\widetilde{k}\leqq \cd_2 k$。

证明。这里的内容几乎全是熟知的。上述不等式\sisi{}{，}以及当
$\cd_{\ell}k<\infty$ 时等号成立，都是 Tate 的结果（CG II 4.1.2）。
尚须证明，$\cd_{\ell}k'<\infty$ 和 $\cd_{\ell}k=\infty$
\sisi{蕴含}{蕴含} $k$ 可序且 $\ell=2$；只须处理以下两种情形：
\begin{itemize}
\item[a)] $[k':k]<\infty$，

\item[b)] $k'=k(x)$，其中 $x$ 在 $k$ 上超越。
\end{itemize}
在情形 a) 中，可以假设 $k'/k$ 可分。于是 Galois 群 $H=G(\bar{k}/k')$
是 $G=G(\bar{k}/k)$ 的开子群；由 Serre 的一个结果 [5]，除非 $G$ 含有阶为
$\ell$ 的元素，否则 $\cd_{\ell}H=\cd_{\ell}G$。代数闭域 $\bar{k}$
不可能含有有限指数的子域 $\widetilde{k}$，除非 $\ell=2$ 且
$\widetilde{k}$ 可序 [7]，因此得到所需结果。

考虑\pageoriginale 情形 b)。设 $R$ 是 $\Spec k[x]$ 在点 $x=0$ 处的 Hensel 化，
它是一个离散赋值环（discrete valuation ring），再设 $K$ 是 $R$ 的分式域。
于是 $\cd_{\ell}k(x)<\infty$ 蕴含 $\cd_{\ell}K<\infty$。
所以只须证明 $\cd_{\ell}K<\infty$ 蕴含 $\cd_{\ell}k<\infty$，
这可由下面的定理推出：

\begin{enonce*}{定理 2.2}\etiquette[2.2]{X.thm2.2}
\footnote{这个定理最近由 $Ax$ [1] 证明。}。
设 $R$ 是 Hensel 离散赋值环，分式域为 $K$，剩余域为 $k$。
设 $\ell\in \PP$。在下列两种情形中，有
$$
\cd_{\ell}K=\cd_{\ell}k+1
$$
成立：
\begin{enumerate}
\renewcommand{\theenumi}{\roman{enumi}}
\renewcommand{\labelenumi}{(\theenumi)}
\item $\ell\neq \car k$。

\item $\ell \neq \car K$，$k$ 完美，且 $\cd_{\ell}k<\infty$。
\end{enumerate}
\end{enonce*}

\begin{enonce*}[remark]{注}
自然，若 $\ell=\car K$，则在所有情形下都有 $\cd_{\ell}K\leq 1$
（CG II \hyperref[X.thm2.2]{2.2}）。因此，在 $\car k=\ell$ 且 $k$ 不完美的
不等特征情形中，尚须研究 $\cd_{\ell}K$ 与绝对微分模 $\Omega^1_k$
（\EGA IV 20.4.3）之间可能存在的关系。
\SGAArchiveOnly{\nde{冻结底本要求在此引用 Kato 的工作，但没有给出具体书目。}}
\end{enonce*}

定理的证明。与定理（\hyperref[X.thm2.1]{2.1}）一样，其大部分内容归功于 Tate。
当 $R$ 完备时，（CG II 命题~12）已明确处理情形 (ii)；利用下述引理便可归结到这一情形：

\begin{enonce*}{引理 2.2.1}\etiquette[2.2.1]{X.lem2.2.1}
设 $R$ 是分式域为 $K$ 的 Hensel 离散赋值环，$\widehat{K}$ 是 $R$ 的完备化
$\widehat{R}$ 的分式域。则
$$
G(\bar{\widehat{K}}/\widehat{K})\simeq G(\bar{K}/K){\quoi}.
$$
\end{enonce*}

证明。这就是说，从\pageoriginale 平展 $K$-代数范畴到平展 $\widehat{K}$-代数范畴的函子
$L\mapsto \widehat{L}=L\otimes_K \widehat{K}$ 是等价。
现在，$\widehat{K}$ 上每个平展代数 $\widehat{K}'$ 都由 $K$ 上某个适当的平展代数
$K'$ 诱导。事实上，若 $\widehat{K}'=\widehat{K}[\alpha]$，其中 $\alpha$
满足系数属于 $\widehat{R}$ 的可分不可约方程 $f(T)=0$，则只须取
$K'$ 为代数 $K[\alpha_0]$，其中 $\alpha_0$ 是系数属于 $R$ 的方程 $f_0(T)=0$ 的根，
且 $f_0\equiv f(\widehat{m}^N)$，这里 $N$ 足够大，$\widehat{m}$ 是
$\widehat{R}$ 的极大理想。这是通常的“Newton 方法”的一个应用（Bourbaki，
Alg.~Comm.~III 4.5）。尚须证明此函子全忠实（忠实性显然），这归结为证明：
若 $K'$ 是\emph{域}，则 $\widehat{K}'$ 也是域。设 $R'$ 是 $R$ 在 $K'$ 中的正规化；
由于 $K'/K$ 可分，$R'$ 在 $R$ 上有限。环 $R'$ 是离散赋值环，其完备化为
$\widehat{R}'=R'\otimes_R \widehat{R}$。因此 $\widehat{R}'$ 是离散赋值环，
% 源文校订：此处全分式域为完备化后的 $\widehat{K}'$。
而 $\widehat{K}'$ 是它的全分式环，所以确实是域。

考虑 \sisi{\eqhyperref[X.eq3.1]{3.1}}{\hyperref[X.thm2.2]{2.2}} 的断言 (i)。
若 $\cd_{\ell}k<\infty$，由 \hyperref[X.lem2.2.1]{2.2.1} 又可归结到
（CG II 命题~12）（该命题陈述中 $k$ 完美的假设，仅在剩余特征为 $\ell$ 的情形中
才被\sisi{用于}{用于}证明）。因此，尚须证明 $\cd_{\ell}K<\infty$ 蕴含
$\cd_{\ell}k<\infty$。

设 $\bar{G}=G(\bar{k}/k)=G(\widetilde{K}/K)$，其中 $\widetilde{K}$ 是 $K$ 的
极大非分歧扩张（\ie 剩余域可分闭的离散赋值环 $\widetilde{R}$ 的分式域，
该环是 $R$ 的严格 Hensel 化）；又设 $G=G(\bar{K}/K)$，
$H=G(\bar{K}/\widetilde{K})$，于是 $\bar{G}=G/H$。
设 $\bar{G}'\subset \bar{G}$ 是一个 $\ell$-Sylow 子群，$k'/k$、$K'/K$
是对应于 $\bar{G}'$ 的扩张。由（CG I 命题~14），$\cd_{\ell}k<\infty$
当且仅当 $\cd_{\ell}k'<\infty$；因而以它在 $K'$ 中的正规化代替 $R$，
便归结到 $\bar{G}$ 是 $\ell$-群的情形。此时只须证明：若 $\cd_{\ell}G<\infty$，
则对充分大的 $N$ 有 $H^N(\bar{G},\ZZ/\ell)=0$（CG I 命题~20）。

将已证明的 i) 的部分应用于 $\widetilde{R}$，\pageoriginale 得
$\cd_{\ell}H=1$。此外，有
\begin{equation*}
H^1(H, 𝛍_{\ell})=H^1(\Spec \widetilde{K},
𝛍_{\ell})=\widetilde{K}^{\ast}/\widetilde{K}^{\ast\ell}=\ZZ/\ell{\quoi}.\tag{2.2.2}
\end{equation*}\etiquette[2.2.2]{X.eq2.2.2}
事实上，由 Kummer 理论 \hyperrefext[IX][IX.Thr3.2]{3.2}，有
$H^1(\Spec \widetilde{K},𝛍_{\ell})=\widetilde{K}^{\ast}/\widetilde{K}^{\ast\ell}$。
我们有正合列
$$
0\longrightarrow \widetilde{R}^{\ast}\longrightarrow
\widetilde{K}^{\ast}\longrightarrow
\ZZ\longrightarrow 0{\quoi},
$$
而 $\widetilde{R}^{\ast}$ 是 $\ell$-可除的，因为 $\widetilde{R}$ 严格局部且
$\ell$ 在 $\widetilde{R}$ 中可逆，故得到
\sisi{\hyperref[X.lem2.2.1]{2.2.1}}{\eqhyperref[X.eq2.2.2]{2.2.2}}。

考虑 Hochschild--Serre 谱序列
$$
E^{pq}_2=H^p(\bar{G}, H^q(H,
𝛍_{\ell}))\Longrightarrow
H^{p+q}(G,𝛍_\ell){\quoi}.
$$
当 $q>1$ 时，有 $E^{pq}_2=0$。若还有 $\cd_{\ell}G<\infty$，则对于充分大的 $n$，
“边界”态射
$$
H^{n-1}(\bar{G}, H^1(H,𝛍_{\ell}))\longrightarrow
H^{n+1}(\sisi{\bar{G}}{G}, 𝛍_{\ell})
$$
必为双射，因而须证明，此时有
$$
H^{n-1}
(\bar{G},H^1(H,𝛍_{\ell}))=H^{n-1}(\bar{G},\ZZ/\ell)=0{\quoi}.
$$

令 $R'=R[x']$，其中 ${x'}^{\ell}=x$，$x$ 是 $R$ 的局部参数，
再令 $K'$ 是 $R'$ 的分式域，$G'=G(\bar{K}/\sisi{k}{K}')$，\etcn
环 $R'$ 的剩余域是 $k$，因为我们已经假设 $\bar{G}$ 是 $\ell$-群，
因而 $\ell$\sisi{次}{次}单位根属于 $K$。我们有图
$$
\xymatrix@M=.25cm@C=1.5cm{
H' \ar[r] \ar@{^{(}->}[d]& G' \ar[r] \ar@{^{(}->}[d]&\bar{G}\\
H  \ar[r]& G \ar[r] & \bar{G},
}
$$
从而\pageoriginale 得到 Hochschild--Serre 谱序列之间的态射，它给出交换图
$$
\xymatrix{
H^{n-1}(\bar{G},H^{1}(H',𝛍_l)) \ar[rd]& &\\
                                              &
H^{n+1}(\sisi{\bar{G}}{G},𝛍_l) &\\
H^{n-1}(\bar{G},H^{1}(H,𝛍_l))\ar[uu]^{\varepsilon}
\ar[ru] & &.
}
$$
由于 $\cd_{\ell}G<\infty$ 蕴含 $\cd_{\ell}G'<\infty$，当 $n$ 足够大时，
此图中所有态射都是同构。但诱导箭头 $\epsilon$ 的态射
$$
% 源文校订：由 (2.2.2)，分母应为 $\ell$。
\ZZ/\ell=H^1(H,𝛍_{\ell})\longrightarrow
H^1(H',𝛍_{\ell})=\ZZ/\ell
$$
显然为零，所以 $H^{n-1}(\bar{G},\ZZ/\ell)=0$，\cqfd

\begin{enonce*}{定理 2.3}\etiquette[2.3]{X.thm2.3}
设 $A$ 是维数为 $1$ 的 Noether Hensel 环，$\ell\in \PP$ 在 $A$ 中可逆。
设 $k$ 是 $A$ 的剩余域，$R(A)$ 是 $\Spec A$ 上的有理函数环。
%% say that this is the total ring of fractions? [with an EGA reference]?
则有
$$
\cd_{\ell}R(A)\leq \cd_{\ell}k+1{\quoi},
$$
且当 $\ell\neq 2$ 或 $k$ 不可序时，等号成立。
\end{enonce*}

证明。以 $A/P$ 代替 $A$，其中 $P$ 是 $A$ 的一个极小素理想，并应用 VIII 1，%% missing reference \XXX
可知能够假设 $A$ 整。设 $K$ 是它的分式域，$R$ 是它的正规化；后者是局部 Hensel 环，
所以是离散赋值环（由 Krull--Akizuki 定理，它是 Noether 环
（Bourbaki，Alg. Comm. 第~7 章，\S 2，命题~5））。又设 $k'$ 是 $R$ 的剩余域。
已知 $[k':k]<\infty$（\loccit）。现在 $K$ 是 $R$ 的分式域，因而由\pageoriginale
\hyperref[X.thm2.2]{2.2} 有
$$
\cd_{\ell}K=\cd_{\ell}k'+1\leq \cd_{\ell}k+1
$$
而根据 \hyperref[X.thm2.1]{2.1}，在 \hyperref[X.thm2.3]{2.3} 的最后一个假设下等号成立。

\begin{enonce*}[remark]{例}
当 $\ell=2$，且 $A$ 是 $\RR[x,y]/(x^2+y^2)$ 在原点处的 Hensel 化时，等号不成立。
\end{enonce*}

\begin{enonce*}{推论 2.4}\etiquette[2.4]{X.cor2.4}
设 $A$ 是剩余域为 $k$ 的 Noether 局部环，$\ell\in \PP$ 在 $A$ 中可逆。
假设 $\ell\neq 2$ \sisi{或}{或} $\Spec A$ 的任何剩余域都不可序。则
$$
\cd_{\ell}R(A)\geq \cd_{\ell}k+\dim A{\quoi},
$$
其中 $R(A)$ 是 $\Spec A$ 上的有理函数环。
\end{enonce*}

证明。可以假设 $A$ 整，\sisi{其}{其}分式域为 $R(A)=K$。
取 $x\in 𝔪_A$\sisi{ 为非零因子}{}。
%% harmonize maximal/prime ideals \XXX
则 $\dim A/(x)=\dim A-1$。对 $\dim A$ 归纳，有
$$
\cd_{\ell}R(A/(x))\geq \cd_{\ell}k+\dim A-1{\quoi}.
$$
设 $k'$ 是 $R(A/(x))$ 的一个剩余域，满足
$\cd_{\ell}k'=\cd_{\ell}R(A/(x))$（\cf VIII 1），并设 $P$ %% 𝔭
是 $A\rightarrow k'$ 的核这一素理想。环 $A_P$ 的维数为 $1$。
因而只须对局部环 $A_P$ 证明此推论，也就是说，归结到 $\dim A=1$ 的情形。
设 $\widetilde{A}$ 是 $A$ 的 Hensel 化。$R(\widetilde{A})$ 的每个剩余域
% 源文校订：这里扩张的基域是先前定义的 $R(A)=K$。
都可以等同于 $R(A)=K$ 的一个可分扩张，所以
$$
\cd_{\ell}R(A)=\cd_{\ell}K\geq \cd_{\ell}R(\widetilde{A}){\quoi}.
$$
这样就归结到 $A$ 为维数 1 的 Hensel 环的情形，此时结论由 \hyperref[X.thm2.3]{2.3} 得到。

\begin{enonce*}{推论 2.5}\etiquette[2.5]{X.cor2.5}
设\pageoriginale $X$ 是 Noether 概形，$\ell\in \PP$ 在 $X$ 上可逆。则
$$
\cd_{\ell}R(X)\geq \dim X{\quoi}.
$$
\end{enonce*}

事实上，取点 $x\in X$ 使 $\dim X=\dim \Oo_{X,x}$（若 $\dim X=\infty$，
则使 $\dim \Oo_{X,x}$ 很大）。设 $A$ 是 $\Oo_{X,x}$ 的严格 Hensel 化。
根据通常的论证，有 $\cd_{\ell}R(X)\geq \cd_{\ell}R(\Oo_{X,x})\geq \cd_{\ell}R(A)$。
此外，$\dim A=\dim X$（\resp 很大）。若 $\ell=2$，且 $\ell$ 在 $X$ 上可逆，
则方程 $T^2+1=0$ 在 $A$ 上可分；因为 $A$ 严格局部，$-1$ 在 $A$ 中是平方，
所以 $A$ 的剩余域都不可序。因此，在所有情形下，$A$ 都满足
\hyperref[X.cor2.4]{2.4} 的假设，且 $\cd_{\ell}R(A)\geq \dim A$，由此得出结果。


\section{严格局部环的分式域}\etiquette[3]{X.sec3}

\setcounter{subsection}{-1}
\subsection{}\etiquette[3.0]{X.subsec3.0}
在以下各节中，我们将迫切需要知道 $\cd_{\ell}R(A)$，这里 $A$ 是严格局部环
（strictly local ring），即剩余域可分闭的 Hensel 环。遗憾的是，一般还不知道
$\cd_{\ell}R(A)$，例如当 $A=\widehat{\ZZ}_p[[x]](p\neq \ell)$ 时。

不过，人们猜想\SGAArchiveOnly{\nde{冻结底本指出 O.~Gabber 已证明此事，但未提供具体书目。}}
若 $A$ 是 Noether 环，且 $\ell$ 在 $A$ 中可逆，则\sisi{将有}{有}
\begin{equation*}
\cd_{\ell}R(A)=\dim A\quad (\ell\mbox{ 在 } A\mbox{ 中可逆}){\quoi},\tag{3.1}
\end{equation*}\etiquette[3.1]{X.eq3.1}
至少在最重要的情形中如此，例如当 $A$ 是\emph{优环}（excellent ring）时
（\EGA IV 7.8.2）。

注意，若 $\dim A=1$，则由 \hyperref[X.thm2.3]{2.3}，
\eqhyperref[X.eq3.1]{3.1} 成立。此外，由 \hyperref[X.cor2.5]{2.5}，
不等式 $\geq$ 在所有情形中都成立。我们以后（XIX）将利用奇点解消
\cite{sga4.exp.X.bib.3}，证明若 $A$ 为优环且剩余\pageoriginale 特征为零，
则 \eqhyperref[X.eq3.1]{3.1} 也成立。这里，我们只在下面
\hyperref[X.prop3.2]{3.2} 的简单情形中证明 \eqhyperref[X.eq3.1]{3.1}。

自然，若 $A$ 的特征为 $\ell$，则有 $\cd_{\ell}R(A)\leq 1$，所以这种情形平凡。
此外，尚须分析剩余特征为 $\ell$ 的不等特征情形；我们并不预期
\eqhyperref[X.eq3.1]{3.1} 在这种\sisi{情形中}{情形中}成立。
当然必须考虑模 $\Omega^1_k$ 的秩（$k$ 为剩余域），\ie 次数 $[k:k^{\ell}]$。
\SGAArchiveOnly{\nde{冻结底本在此提及 Gabber 和 Orgogozo 的工作，但未提供具体书目。}}

\begin{enonce*}{命题 3.2}\etiquette[3.2]{X.prop3.2}
设 $X$ 是域 $k$ 的谱上或离散赋值环 $R$ 的谱上的有限型预概形（prescheme），
设 $\ell\in \PP$ 在 $k$（\resp~$R$）中可逆。则对 $X$ 的每个严格局部化环 $A$，
$\hyperref[X.eq3.1]{3.1}$ 都成立。
\end{enonce*}

\emph{证明}。取 $X$ 整（像通常一样，可归结到这种情形），并令 $n=\dim X$。

假设 $X$ 定义在\emph{可分闭}域 $k$ 上。那么，对 $X$ 在一个\emph{闭}点处的每个
严格局部化环 $A$，有 $\dim A=n$。由于 $R(A)$ 的每个剩余域 $K$ 都可以等同于
$R(X)$ 的一个代数扩张，且
$$
\deg.\tr.(R(X)/k)=\dim X=n{\quoi},
$$
%% \tr. etc. ugly: there is one space too many (tr .) \XXX
由 \hyperref[X.thm2.1]{2.1} 有 $\cd_{\ell}K\leq n$，故
$\cd_{\ell}R(A)\leq n$。反向不等式由 \hyperref[X.cor2.5]{2.5} 得到。

现在假设 $X$ 在 $\Spec R$ 上有限型，其中 $R$ 是离散赋值环，并设 $A$ 是
$X$ 在点 $x$ 处的严格局部化。若 $x$ 不位于 $\Spec R$ 的闭点上方，
则立即可归结到 $X$ 在域上有限型的情形（将 $R$ 局部化）。
出于同样的理由，还可以假设态射 $X\rightarrow \Spec R$ 为支配态射。

暂且\pageoriginale 再假设 $R$ 是\emph{严格局部}离散赋值环，且 $x$ 是
$X$ 的\emph{闭}点。于是 $\dim A=\dim X=n$（\EGA IV 5.6.5）。
因为 $R(A)$ 的每个剩余域 $K$ 都可以等同于 $R(X)$ 的一个可分扩张，由
\hyperref[X.thm2.1]{2.1} 和 \hyperref[X.thm2.2]{2.2} 有
$$
\cd_{\ell}R(A)\leq \cd_{\ell}R(X)\leq
\cd_{\ell}L+\deg.\tr.(R(X)/L)=1+(n-1)=n
$$
其中 $L$ 是 $R$ 的分式域；结合 \hyperref[X.cor2.5]{2.5}，便在此情形中得到结论。

利用下面的引理，就可以由这些特殊情形推出一般情形：

\begin{enonce*}{引理 3.3}\etiquette[3.3]{X.lem3.3}
\begin{enumerate}
\renewcommand{\theenumi}{\roman{enumi}}
\renewcommand{\labelenumi}{(\theenumi)}
\item 设 $X$ 在域 $k$ 的谱上有限型，$A$ 是 $X$ 在点 $x$ 处的严格局部化环。
  则存在可分闭域 $k'$、$\Spec k'$ 上的有限型概形 $X'$ 及 $X'$ 中的闭点 $x'$，
  使得 $A$ 同构于 $X'$ 在 $x'$ 处的一个严格局部化环。

\item 设 $X$ 在 $\Spec R$ 上有限型，其中 $R$ 是离散赋值环；设 $x$ 是
  $X$ 中位于 $\Spec R$ 的闭点上方的一个点，$A$ 是 $X$ 在 $x$ 处的严格局部化。
  则存在严格局部离散赋值环 $R'$、$\Spec R'$ 上的有限型概形 $X'$ 及 $X'$ 的闭点
  $x'$，它位于 $\Spec R'$ 的闭点上方，使 $A$ 同构于 $X'$ 在 $x'$ 处的一个严格局部化。
\end{enumerate}
\end{enonce*}

证明。
\begin{enumerate}
\renewcommand{\theenumi}{\roman{enumi}}
\renewcommand{\labelenumi}{(\theenumi)}
\item 可以取 $X$ 为仿射的。

设 $Y$ 是 $x$ 在 $X$ 中的闭包，赋予诱导的约化结构，并令 $d=\dim Y$。
由正规化引理（\cite{sga4.exp.X.bib.2} 第 V 章 \S 3 第 ${\rm n^{\circ}}$ 1 节），存在态射
$$
X\longrightarrow \Spec k[y_1,\ldots, y_d]=T
$$
诱导\pageoriginale 一个\emph{有限}态射 $Y\rightarrow T$。
因此，$\Oo_{X,x}$ 的严格局部化环同构于 $\Oo_{X',x'}$ 的严格局部化，
其中 $X'=\fprod{X}{\Spec\bar{R(T)}}{T}$，而 $x'$ 是 $X'$ 的一个闭点。

\item 这种情形的处理类似。可以假设 $X$ 仿射，其环为 $A$。
设 $Y$ 是 $x$ 在 $X$ 中的闭包，赋予诱导的约化结构；它是 $X$ 在 $\Spec R$ 上的
闭纤维 $X_0$ 的一个闭子概形，并令 $d=\dim Y$。
将 (i) 中考虑的 $A\otimes_R k$ 的元素 $y_1,\ldots,y_d$ 提升为 $A$ 的元素，便得到态射
$$
X\longrightarrow \Spec R[y_1,\ldots,y_d]=T
$$
诱导一个\emph{有限}态射 $Y\rightarrow T_0$。设 $R'$ 是
$R[y_1,\ldots,y_d]=B$ 在素理想 $𝔪B$ 处的严格局部化。
%% specify that 𝔪 is the maximal ideal of R? \XXX
则 $R'$ 是严格局部离散赋值环，可以取 $X'=\fprod{X}{\Spec R'}{T}$，
并取 $x'$ 为 $X'$ 中位于 $x$ 上方的一个闭点。
\end{enumerate}

\section{上同调维数：$ℓ$ 在 $𝒪_X$ 中可逆的情形}\etiquette[4]{X.sec4}

\begin{enonce*}{定理 4.1}\etiquette[4.1]{X.thm4.1}
设 $X$ 是 Noether 概形，$\ell\in \PP$，$\varphi$ 是 $X$ 上取值于 $\NN$ 的函数，
满足下述条件：
\begin{enumerate}
\renewcommand{\theenumi}{\roman{enumi}}
\renewcommand{\labelenumi}{(\theenumi)}
\item 对每个 $x\in X$，有 $\cd_{\ell}k(x)\leq \varphi (x)$。

\item 设 $x\in X$，$y\in Y=$ $x$ 的闭包，且 $y\neq x$。
设 $K$ 是 $Y$ 在 $y$ 上方某个严格局部化的有理函数环。
有 $\cd_{\ell}K<\varphi(x)-\varphi(y)$。
\end{enumerate}
则对 $X$ 上每个 $\ell$-挠层 $F$，有
\begin{gather*}
H^q(X,F)=0\\
q>\varphi(F)\mathop{=}_{\text{\rm def}}\Sup\{\varphi(x)\mid x\in X,
F_{\bar{x}}\neq 0\}{\quoi}.\tag*{若}
\end{gather*}
\end{enonce*}

证明。\pageoriginale
将 $\varphi$\sisi{，}{}限制在 $X$ 的任意子预概形上，假设仍然满足，
因而由 Noether 归纳法，可以假设定理对于 $X$ 的每个不等于 $X$ 的闭子预概形都成立。
给定 $F$，写成 $F=\varinjlim F_i$，其中 $F_i$ 是 $F$ 的可构造子层
（\hyperrefext[IX][IX.prop2.9]{2.9} (iii)）。

对每个 $F_i$，满足 $F_x\neq 0$ 的点 $x\in X$ 所成的集合是 $X$ 中的可构造集，
因此可对支集不稠密的每个 $F_i$ 应用归纳假设。由此可知，归纳假设实际上蕴含：
对至少在 $X$ 的一个极大点处为零的每个 $F$，定理成立。
以下便假设 $F$ 不满足这个条件。

现在设 $i:\Spec R(X)\rightarrow X$ 为包含映射，其中 $R(X)$ 是 $X$ 的有理函数环，
并考虑下列正合列，其中 $K$ 和 $C$ 分别定义为核与余核：
$$
0\longrightarrow K\longrightarrow F\longrightarrow i_{\ast}
i^{\ast}F\longrightarrow C\longrightarrow 0{\quoi}.
$$
这里，$K$ 和 $C$ 在 $X$ 的每个极大点处都为零，所以由 (ii)，
$\varphi(K)$（\resp $\varphi(C)$）严格小于 $n=\varphi(F)$，
故当 $q\geq n$ 时，$H^q(X,K)=H^q(X,C)=0$。因此，\sisi{}{，}
由上同调正合列，要证明对 $q>n$ 有 $H^q(X,F)=0$，只须证明
$H^q(X,i_{\ast}i^{\ast}F)=0$；所以可以假设 $F$ 形如 $i_{\ast}G$，
其中 $G$ 是 $\Spec R(X)$ 上适当的 $\ell$-层。

考虑 Leray 谱序列\sisi{。}{}
$$
E^{p,q}_2=H^p(X,R^qi_{\ast}G)\Longrightarrow H^{p+q}(\Spec
R(X),G){\quoi}.
$$
现在，$R(X)_{\text{\rm réd}}$ 是 $X$ 的各极大点 $x$ 的域 $k(x)$ 的乘积，
因而由 (i) 有 $\cd_{\ell}R(X)\leq n$；故此谱序列的极限在
\sisi{$\dim > n$}{严格大于~$n$ 的维数}处为零。
我们\pageoriginale 要证明，当 $p>n$ 时，$E^{p0}_2=0$。
考察谱序列可知，只须证明当 $q>0$ 且 $p>n-q-1$ 时 $E^{pq}_2=0$。
因而根据归纳假设，只须证明当 $q>0$ 时 $\varphi(R^qi_{\ast}G)<n-q$。
但由 VIII \hyperref[VIII.thm5.2]{5.2}，$R^qi_{\ast}G$ 在位于点 $y\in X$
上方的几何点 $\bar{y}$ 处的纤维恰为 $H^q(\Spec K, G_K)$，
其中 $K$ 是 $X$ 在 $\bar{y}$ 处的严格局部化的有理函数环，$G_K$ 是
$\Spec K$ 上的诱导层。由 (ii)，有 $\cd_{\ell}K<n-\varphi(y)$，
故当 $q\geq n-\varphi(y)$ 时 $(R^qi_{\ast}G)_{\bar{y}}=0$，这就得到所需结果。

\begin{enonce*}{推论 4.2}\etiquette[4.2]{X.cor4.2}
设 $X$ 是 Noether 概形，$\ell\in \PP$，$n\in \NN$，并假设满足以下条件：
\begin{enumerate}
\renewcommand{\theenumi}{\roman{enumi}}
\renewcommand{\labelenumi}{(\theenumi)}
\item 对 $X$ 的每个余维为 $c$ 的点 $y$，有 $\cd_{\ell}k(y)\leq n-2c$。

\item 对 $X$ 的不可约闭子概形 $Y$ 的每个严格局部化环 $A$，有
  $\cd_{\ell}R(A)\leq \dim A$。

则
$$
\cd_{\ell}X\leq n{\quoi}.
$$
更确切地，若 $F$ 是在余维 $<s$ 的每个点 $y$ 处为零的 $\ell$-层，则有
$$
H^q(X,F)=0\mbox{ 若 } q>n-2s{\quoi}.
$$
\end{enumerate}
\end{enonce*}

事实上，可以取 $\varphi(y)=n-2\codim y$ 并应用 \hyperref[X.thm4.1]{4.1}。
它的条件 (i)（\resp (ii)）由（\hyperref[X.cor4.2]{4.2}）中相应的条件 (i)
（\resp (ii)）推出；其中对 (ii) 要考虑不等式
$\dim \Oo_{Y,y}+\dim \Oo_{X,x}\leq \dim \Oo_{X,y}$。

\begin{enonce*}{推论 4.3}\etiquette[4.3]{X.cor4.3}
设 $X$ 是域 $k$ 上的有限型概形，维数为 $n$，且 $\ell\neq \car k$。
则 $\cd_{\ell}X\leq 2n+\cd_{\ell}k$。
\end{enonce*}

将前一推论\pageoriginale 中的 $n$ 换成 $2n+cd_{\ell}k$，并应用该推论。
利用 \hyperref[X.thm2.1]{2.1} 和 \hyperref[X.prop3.2]{3.2} 来验证
\hyperref[X.cor4.2]{4.2} 的条件 (i) 和 (ii)；其中对 (i) 要考虑不等式
$$
\deg.\tr.K(y)/k+\dim \Oo_{X,y}\leq \dim X{\quoi}.
$$

\begin{enonce*}{推论 4.4}\etiquette[4.4]{X.cor4.4}
设 $X$ 是 Noether 概形，$\ell\in \PP$ 在 $X$ 上可逆。
假设 $\hyperref[X.cor4.2]{4.2}$ 的条件 {\rm (ii)} 满足，并且
$\ell\neq 2$ \sisi{或}{或} $X$ 的任何剩余域都不可序。则有
$$
\cd_{\ell}X\leq \cd_{\ell}R(X)+\dim X\leq 2\cd_{\ell}R(X){\quoi}.
$$
\end{enonce*}

在 \hyperref[X.cor4.2]{4.2} 中令 $n=\cd_{\ell}R(X)+\dim X$，并应用
\hyperref[X.cor2.4]{2.4}。它蕴含 \hyperref[X.cor4.2]{4.2} 的条件 (i)，因为
\[\cd_{\ell}k(y)\leq \cd_{\ell}R(\Oo_{X,y})-\dim \Oo_{X,y}\leq
\cd_{\ell}R(X)-\dim\Oo_{X,y}\leq n-2\dim \Oo_{X,y}.
\]
\hyperref[X.cor4.4]{4.4} 中的第二个不等式由 \hyperref[X.cor2.5]{2.5} 得到。

\begin{enonce*}[remark]{例 4.5}\etiquette[4.5]{X.exem4.5}
\hyperref[X.cor4.4]{4.4} 中对剩余域所作的那个不大雅致的假设是必要的。
例如，取 $X=\Spec \RR[x,y,z]/(x^2+y^2+z^2)$，其余所有条件都满足，
且 $\dim X=\cd_{\ell}R(X)=\ell=2$；但集中在点 $(0,0,0)$ 上的层 $\ZZ/2$
在每个次数都有非零上同调，所以 $\cd_{\ell}X=+\infty$。
\end{enonce*}

\setcounter{subsection}{5}
\subsection{}\etiquette[4.6]{X.subsec4.6}
人们还会想尝试在 \hyperref[X.cor4.2]{4.2} 中仅对 $X$ 的闭点施加假设，但这是不可能的。
例如，存在剩余域代数闭、而函数域的上同调维数任意大的离散赋值环\footnote[1]{
  为看出这一点，只须注意：若 $X$ 是域 $k$ 上的光滑代数概形，且有一点 $x\in X(k)$，
  则可找到经过 $x$ 的“形式曲线弧”，其“代数包络”为 $X$。
  因而它的函数域含有 $X$ 的函数域 $K$，并在其上诱导剩余域为 $k$ 的离散赋值。}。
这样的环的谱同样具有很大的上同调维数。
%% to do (for the note): i -> (i) \XXX


\section{上同调维数：$ℓ=p$ 的情形}\etiquette[5]{X.sec5}\pageoriginale

\setcounter{subsection}{-1}
\subsection{}\etiquette[5.0]{X.subsec5.0}
利用 Artin\sisi{ - }{-}Schreier 理论（\hyperrefext[IX][IX.Thr3.5]{3.5}），
对于特征为~$p$ 的概形 $X$ 的 $p$-上同调维数，可以得到比 \hyperref[X.cor4.3]{4.3}
更好的上界。设 $\cdqc X$ 是最大的数 $n$，使 $X$ 上至少存在一个
\emph{拟凝聚}（quasi-coherent）模层 $F$（按 Zariski 意义），满足 $H^n(X,F)\neq 0$。
我们有如下结果。

\begin{enonce*}{定理 5.1}\etiquette[5.1]{X.thm5.1}
设 $X$ 是特征 $p>0$ 的 Noether 概形。则
$$
\cd_pX\leq \cdqc X+1{\quoi}.
$$
特别地，特征 $p>0$ 的 Noether 仿射概形 $X$ 的 $p$-上同调维数至多为 $1$。
\end{enonce*}

证明。应用 \hyperrefext[IX][IX.prop5.5]{5.5}：

考虑到当 $f$ 是有限态射且 $F$ 是阿贝尔层（VIII \hyperref[VIII.prop5.5]{5.5}）
%% reference \XXX
（\resp 拟凝聚层 $F$）时，$R^q f_{\ast}F$ 对 $q>0$ 为零，
% 源文校订：由下述正合列，所需范围为 $q>\cdqc X+1$。
只须证明：若 $i:U\rightarrow X$ 是开浸入且 $q>\cdqc X+1$，
则 $H^q(X,\sisi{i_1}{i_!}(\ZZ/p)_U)=0$。

设 $Y$ 是 $X$ 中支集为 $X-U$ 的闭子概形，$J$ 是定义 $Y$ 的凝聚理想层。
态射 $p:f\mapsto f^p-f$（\hyperrefext[IX][IX.Thr3.5]{3.5}）诱导一个态射
$J\rightarrow J$。考察正合列 \hyperrefext[IX][IX.Thr3.5]{3.5} 以及正合列
$$
0\longrightarrow i_{!}(\ZZ/p)_U\longrightarrow
(\ZZ/p)_X\longrightarrow
(\ZZ/p)_Y\longrightarrow 0{\quoi},
$$
可知实际上存在正合列
\begin{equation*}
0\longrightarrow i_{!}(\ZZ/p)_U\longrightarrow
J\longrightarrow J\longrightarrow 0{\quoi}.\tag{$\ast$}
\end{equation*}
因为 $J$ 凝聚，当 $q>\cdqc X$ 时，有 $H^q(X,J)=0$，
于是由上同调长正合列得到\pageoriginale 结论。

对于可分闭域 $k$ 上的拟射影概形，这个上界在某些情形中还可以\sisi{减去}{减去}~$1$。
事实上，有

\begin{enonce*}{推论 5.2}\etiquette[5.2]{X.cor5.2}
设 $X$ 是特征 $p>0$ 的域 $k$ 上的有限型概形。有
$$
\cd_pX\leq \dim X+1{\quoi},
$$
而若 $k$ 可分闭且 $X$ 拟射影\footnote[1]{利用 Chow 引理，容易将条件
  “$X$ 拟射影”替换为“$X$ 分离”。}，则有
$$
\cd_pX\leq \dim X{\quoi}.
$$
\end{enonce*}

事实上，第一个断言显然，因为 $\cdqc X\leq \dim X$
（\cite{sga4.exp.X.bib.6} 4.15.2）。下面处理第二个断言。
对这样的 $X$，对于凝聚层 $F$，上同调群 $H^n(X,F)$ $(n=\dim X)$
是 $k$ 上的有限维向量空间。令 $V=H^n(X,J)$，其中 $J$ 是正合列 $(\ast)$
中出现的理想层。由 $(\ast)$ 得到的态射 $\varphi:V\rightarrow V$
显然形如 $\varphi=\epsilon-\id$，\sisi{或}{其中} $\epsilon$ 对
$r\in k$ 和 $v\in V$ 满足 $\epsilon(rv)=r^p\epsilon(V)$。
因此，由 $\Phi$ 得到的仿射概形的映射 $\Phi$ 的 Jacobian 矩阵为 $-I$，
\sisi{$I$}{这里~$I$ 是}单位矩阵。所以 $\Phi$ 平展，且因其可加而为满射；
又因 $k$ 可分闭，便推出 $\Phi$ 为满射，结论立即得到。

\section{$\mathrm{Spec} ℤ$ 上有限型预概形的上同调维数}\etiquette[6]{X.sec6}\pageoriginale

\begin{enonce*}{命题 6.1}\etiquette[6.1]{X.prop6.1}
设 $X$ 在 $\Spec \ZZ$ 上拟有限。假设 $\ell\neq 2$ \sisi{或}{或}
$X$ 的每个特征零剩余域都完全虚。则 $\cd_{\ell}X\leq 3$。
\end{enonce*}

这是\sisi{由}{由}\sisi{定理}{定理} \hyperref[X.thm2.2]{2.2} (ii)、
\sisi{推论}{推论} \hyperref[X.cor4.2]{4.2} 以及下述事实推出的：
数域（若 $\ell =2$，则须完全虚）的 $\ell$-上同调维数 $\leq 2$（CG II 4.4）。

\sisi{事实上由类域论可知}{类域论的一个推论是}，若态射
$X\rightarrow \Spec \ZZ$ 固有且满射，则等号成立。

遗憾的是，若仅假设 $X$ 在 $\Spec \ZZ$ 上有限型，第 ${\rm n^{\circ}}$ 4 节的结果
便不能直接应用，因为（\hyperref[X.cor4.2]{4.2}）的假设 (ii) 对
\emph{剩余特征}为 $\ell$ 的局部环一般并不满足。
因此，必须利用结果（\hyperref[X.thm5.1]{5.1}）才能获得正确的上界；
现在就来利用定理（\hyperref[X.thm5.1]{5.1}）。

\begin{enonce*}{定理 6.2}\etiquette[6.2]{X.thm6.2}
设 $X$ 在 $\Spec \ZZ$ 上有限型，$\ell\in \PP$。若 $\ell=2$，
假设 $X$ 的任何剩余域都不可序。则
$$
\cd_{\ell}X\leq 2\dim X+1{\quoi}.
$$
更确切地，设 $F$ 是 $\ell$-挠层，对 $X$ 中满足 $F_{\bar{x}}\neq 0$ 的每个点 $x$，都有
$$
\dim\{\bar{x}\}\leq
\begin{cases}
(N-1)/2 &\mbox{若\quad } \ell\neq \car k(x)\\
N-1     &\mbox{若\quad } \ell=\car k(x)
\end{cases}
$$
（$\{\bar{x}\}$ 是 $\{x\}$ 的闭包）；则
$$
H^p(X,F)=0\mbox{ 若 } p>N{\quoi}.
$$
\end{enonce*}

证明。\pageoriginale 对 $X$ 归纳，可以假设定理对每个不等于 $X$ 的闭子概形都成立。
若 $F$ 如陈述所述，并写成 $F=\varinjlim F_i$，其中 $F_i$ 是 $F$ 的可构造子层
（\hyperrefext[IX][IX.prop2.9]{2.9} (iii)），则每个 $F_i$ 也满足相同的假设。
因此，只须在 $F$ 可构造时证明定理。

此时，满足 $F_{\bar{x}}\neq 0$ 的点 $x\in X$ 所成的集合 $E$ 是可构造子集；
归纳假设蕴含：对于 $E$ 不稠密的每个 $F$，定理成立。
此外，对于集中在 $X$ 位于 $\Spec \ZZ$ 的点 $(\ell)$ 上方的纤维上的层，
定理是 \hyperref[X.thm5.1]{5.1} 的推论；而对于集中在点 $(p)$ 上方纤维上的层，
其中 $p\neq \ell$，定理是 \hyperref[X.cor4.3]{4.3} 的推论，
这里要考虑 $\cd_{\ell}\ZZ/p \ZZ=1$（CG II 2.2）。

取 $X$ 约化，并令 $X=X_{0}\cup X_1$，其中 $X_1$ 是 $X$ 的特征为 $\ell$ 的
不可约分支之并这一闭子集，$X_{0}$ 是其余不可约分支之并；给 $X_{0}$ 和 $X_1$
赋予诱导的约化结构。设 $j_{\nu}:X_{\nu}\rightarrow X$ 为包含映射，
令 $d_{\nu}=\dim X_{\nu}$，$N=\sup \{2d_{0}+1, d_1+1\}$。
对于 $X$ 上每个可构造层 $F$，都有正合列\sisi{。}{}
\begin{equation*}
0\longrightarrow F\longrightarrow
{j_{0}}_{\ast} j_{0}{}^{\ast}F\oplus {j_1}_{\ast} j_1
{}^{\ast}F\longrightarrow C\longrightarrow 0{\quoi},\tag{$\ast$}
\end{equation*}
其中 $C$ 集中在 $X_{0}\cap X_1$ 上，后者的特征为 $\ell$，维数 $\leq d_1-1$。
于是
$$
H^q(X,C)=0\mbox{ 若 } q>N-1{\quoi},
$$
因为由 \hyperref[X.thm5.1]{5.1} 有
$\cd_{\ell}(X_{0}\cap X_1)\leq \dim
\sisi{}{(}X_{0}\cap X_1\sisi{}{)}+1\leq d_1\leq N-1$。
考察 $(\ast)$ 相应的上同调长正合列，便归结到 $X=X_{0}$ \sisi{或}{或} $X=X_1$ 的情形；
正如已经指出的，后一情形是 \hyperref[X.thm5.1]{5.1} 的推论。

因此，\pageoriginale 假设 $X$ 的每个极大点的特征均不同于 $\ell$，
并设 $i_{\nu}:x_{\nu}\rightarrow X$ 为各极大点的包含映射。
考察正合列
$$
0\longrightarrow K\longrightarrow F\longrightarrow \bigoplus_{\nu}
i_{\nu\ast}i_{\nu}{}^{\ast}F\longrightarrow C\longrightarrow 0
$$
由归纳法便归结到 $F=i_{\nu\ast}i_{\nu}{}^{\ast}F$ 的情形，也就是
$F$ 形如 $i_{\ast}G$，其中 $G$ 是 $\Spec k(x)$ 上的适当层，$x$ 是 $X$ 的一个点。
因此，可以假设 $X$ 不可约且约化，其一般点为 $x$，并由 \hyperref[X.cor4.3]{4.3}
假设 $k(x)$ 的特征为零。

考虑谱序列
\begin{equation*}
H^p(X,R^qi_{\ast}G)\Longrightarrow H^{p+q}(\Spec k(x),
G){\quoi}.\tag{${\ast}{\ast}$}
\end{equation*}
令 $s=\dim X$。我们须证明
\begin{equation*}
E^{p0}_2=0\mbox{ 对于 } p> 2s+1{\quoi}.\tag{$+$}
\end{equation*}
现在 $k(x)[i]$ 在 $\QQ[i]$ 上的超越次数为 $s-1$，
%% say that i = √(-1)? \XXX
而后者是上同调维数为~$2$ 的域。因为当 $\ell=2$ 时 $k(x)$ 不可序，
由 \hyperref[X.thm2.1]{2.1} 得 $\cd_{\ell}k(x)=s+1$。
这蕴含谱序列的目标当 $n>s+1$ 时为零。因此，为证明 $(+)$，只须证明
在 $({\ast}{\ast})$ 中有
\begin{equation*}
E^{pq}_2=0\mbox{ 若 } p>2s-q\mbox{ 且 } q>0{\quoi}.\tag{${+}{+}$}
\end{equation*}

设 $y$ 是 $X$ 的一个点，$K$ 是 $X$ 在位于 $y$ 上方的一个几何点 $\widetilde{y}$
处的严格局部化的分式环。则 $K$ 是若干域的直积，其中每个域在 $\ZZ$ 的一个
严格局部化上的超越次数为 $s-1$；故由 \hyperref[X.thm2.1]{2.1} 和
\hyperref[X.thm2.2]{2.2}，有 $\cd_{\ell}K\leq s$。
此外，若 $y$ 的剩余特征不为 $\ell$，可以应用 \hyperref[X.prop3.2]{3.2}。于是有
\begin{equation*}
(R^qi_{\ast}G)_{\widetilde{y}}=0
\begin{cases}
\mbox{ 若 $q>\codim y$，\ie $\dim \bar{y}>s-q$，且 $\ell\neq \car
k(y)$，}\\
\mbox{或在所有情形中，若 $q>s$。}
\end{cases}
\tag{$''$}
\end{equation*}
特别地，\pageoriginale 当 $q>s$ 时，有 $E^{pq}_2=0$。
为在 $0<q\leq s$ 的情形中验证 $({+}{+})$，对层 $R^qi_{\ast}(G)$ 应用归纳假设。

根据陈述和 $('')$，必须令 $N'=2s-q$，并验证以下两个不等式
\begin{align*}
& s-q\leq (N'-1)_{/2}\\
& s-1=\dim X_{\ell}\leq N'-1{\quoi},
\end{align*}
其中 $X_{\ell}$ 是 $X$ 在 $\Spec \ZZ$ 的\sisi{点}{点} $(\ell)$ 上方的
\sisi{纤维空间}{纤维}。当 $0<q\leq s$ 时，它们确实成立\sisi{，}{。} \cqfd

\begin{thebibliography}{99}
\bibitem{sga4.exp.X.bib.1}
Ax, J.  Proof of some conjectures of Serre on
cohomological dimension, Proc. Amer. Math. Society 16 (1965)
1214-1221.

\bibitem{sga4.exp.X.bib.2} Bourbaki, N. Algèbre \sisi{Commutative}{commutative}, Paris (1960)

\bibitem{sga4.exp.X.bib.3}
Hironaka, H.  Resolution of singularities of an algebraic surface
over a field of characteristic zero, Annals of Math. Vol. 79, No. 1, 2
(1964)


\bibitem{sga4.exp.X.bib.4} Nagata, M. Local Rings, New York (1962)

\bibitem{sga4.exp.X.bib.5}
Serre, J-P.  Sur la dimension cohomologique \sisi{de}{des} groupes profinis,
Topology Vol. 3, p. 413 (1965)


\bibitem{sga4.exp.X.bib.6} Godement, R. \sisi{Théorie des Faisceaux}{Topologie algébrique et théorie des faisceaux}, Paris (1958).

\bibitem{sga4.exp.X.bib.7}
E. Artin et A.  Schreier - Eine Kennzeichnung der \sisi{reel}{reell}
abgeschlossenen Körper, Hamb. Abh. 5 (1927) pp. 225-231.


\bibitem{sga4.exp.X.bib.8} Serre, J.P. Cohomologie \sisi{Galoisienne}{galoisienne}, Lecture Notes
${\rm n^{\circ}}$ 5 (Springer).
\end{thebibliography}

% END native input: chapters/exp-10-zh.tex
% BEGIN native input: chapters/exp-11-zh.tex
\chapter{与经典上同调的比较：光滑概形的情形}
%\addtocontents{toc}{M. Artin}
\begin{center}
M. Artin
\end{center}
\phantomsection\MakeLinkTarget*{sga4.exp.XI}\label{sga4.exp.XI}

\setcounter{pageoriginal}{63}

\section{引言}\etiquette[1]{XI.sec1}\pageoriginale

在第~\hyperref[XI.sec4]{4} 节中，我们将证明：对于 $\Spec\CC$ 上的光滑概形 $X$，以局部常值挠层（locally constant torsion sheaf）为系数的平展上同调（étale cohomology）与经典上同调（classical cohomology）一致。这个证明本质上就是 \cite{sga4.exp.XI.bib.1} 中给出的证明；只要有下面回顾的 Grauert 定理（\hyperref[XI.thm4.3]{4.3} (iii)），即 $X(\CC)$ 的每个有限覆盖都来自 $X$ 的一个平展覆盖，证明就是初等的。因此，尽管稍后我们将利用 Hironaka 的奇点消解定理 \cite{sga4.exp.XI.bib.3} 证明一个更完整的结果，这里仍给出这个证明\footnote[1]{还应指出，Hironaka 的奇点消解也给出了 Grauert 结果的一个非常简单的证明；从这个意义上说，无论如何都可以把它看作比较定理的基础。关于这一点，参见 Mme M. Raynaud 在 \SGA 1 XII（North Holland Pub. Cie）中的报告。}。此外，第~\hyperref[XI.sec3]{3} 节中构造的良好邻域，在其他\pageoriginale 情形中也可能有用。

注意，对于非挠层，平展上同调与经典上同调并不一致。例如，对每个正规概形 $X$ 都有 $H^1(X_{\et},\ZZ)=0$（\hyperrefext[IX][IX.prop3.6]{3.6} (ii)），但经典上同调未必为零。

\section{充分一般的超平面截面的存在性}\etiquette[2]{XI.sec2}

下面的结果是经典的，但尚缺少文献出处（有待 \EGA V）：

\begin{enonce*}{定理 2.1}\etiquette[2.1]{XI.thm2.1} %Why is the statement
                                %not formulated for a linear subspace
                                %varying around a fixed point?
                                %That is, (i) with $P\in L$ or
                                %$(ii)$ with $d=1$
设 $k$ 是一个域，$V\subset \PP^n_k$ 是一个纯 $r$ 维的射影代数概形。设 $V'$ 是由 $V$ 中在 $\Spec k$ 上光滑的点组成的开子概形，$P$ 是 $\PP^n$ 的一个有理点。则：
\begin{enumerate}
\renewcommand{\theenumi}{\roman{enumi}}
\renewcommand{\labelenumi}{(\theenumi)}
\item $\PP^n$ 中给定维数 $d$ 的一个“充分一般”的线性子空间 $L$，在 $L\cap V'$ 的每一点都与 $V$ 横截相交（meet transversely）。

\item 设 $\{H_1,\ldots,H_s\}$ $(s\leq r)$ 是一组 %a
                                %'hyperplane' is not linear here
  经过 $P$ 的 $d\geq 2$ 次超曲面，并设 $Y=H_1\cap\ldots \cap H_s$。若 $\{H_1,\ldots,H_s\}$ 充分一般，则 $Y$ 在 $Y\cap V'$ 的每一点都与 $V$ 横截相交。
\end{enumerate}
\end{enonce*}

“对于充分一般的线性子空间”是指：存在参数化 $d$ 维线性子空间的 Grassmann 概形的一个稠密开子集 $U$，只要 $L$ 的参数属于 $U$，所讨论的关于 $L$ 的断言就成立。因此，当 $k$ 不是无限域时，这并不蕴含这样的（在 $k$ 上定义的）$L$ 存在。同样，
% mention the Veronese embedding?
经过 $P$ 的 $d$ 次超曲面由一个射影空间 $\PP^N$ \pageoriginale 参数化，而“对于充分一般的 $\{H_1,\ldots,H_s\}$”是指：存在 $(\PP^N)^s$ 的一个稠密开子集 $U$，只要 $(H_1,\ldots,H_s)$ 的参数属于 $U$，所讨论的断言就成立。

(i) 的证明：由归纳法，立即归结为 $L$ 是超平面的情形。设 $(V'_{\nu})$ 是 $V'$ 的一个有限开覆盖。如果对每个 $\nu$，将符号 $V'$ 换成 $V'_{\nu}$ 后断言 (i) 都成立，那么它对 $V'$ 也成立，因为有限个稠密开子集的交仍是稠密开子集。因此，我们可以把问题换成仿射代数概形的情形，记作 $V\subset \EE^n_k=\Spec k[x_1,\ldots,x_n]$。超平面 $L$ 将是一个线性方程的零点集：
$$
\ell(x)=a_0+\sum^n_{i=1}a_ix_i{\quoi}.
$$

回顾光滑概形局部是完全交。因此，我们可以把 $V'$ 换成 $V$ 中点的一个开子集，使得存在在 $V$ 上为零的多项式 $f_1(x),\ldots,f_{n-r}(x)$，以及指标 $1\leq j_{\nu}\leq n$ $(\nu=1,\ldots,n-r)$，使行列式
$$
\left|\frac{\partial f_i}{\partial x_{j_{\nu}}}\right|\quad 
(i,\nu=1,\ldots,n-r)
$$
可逆。不妨设 $j_{\nu}=\nu$ $(\nu=1,\ldots,n-r)$。于是，我们可以由方程
$$
y_j=\sum^{n-r}_{i=1}\quad c_i\quad \frac{\partial f_i}{\partial x_j}
$$
解出\pageoriginale $c_i$，将其表示为变量 $x_i$、$y_j$ $(i=1,\ldots, n-r; j=1,\ldots,n)$ 的有理函数，而且这些函数 $c_i(x,y)$ 在 $V'$ 的每一点都有定义。

现在，$L$ 在 $V'$ 的一点 $Q$ 与 $V$ 横截相交，等价于说矩阵
$$
\begin{bmatrix}
\frac{\partial f_1}{\partial x_1} & \cdots & \frac{\partial
  f_1}{\partial x_n}\\
\vdots & &\vdots\\
\frac{\partial f_{n-r}}{\partial x_1} & & \frac{\partial
  f_{n-r}}{\partial x_n}\\
a_1 & \cdots & a_n
\end{bmatrix}
$$
在 $Q$ 的值的秩为 $n-r+1$。所以，$L$ 在 $V'$ 的每一点与 $V$ 横截相交的充要条件，是由以下 $r+1$ 个方程组成的方程组
\begin{align*}
a_j &= \sum^{n-r}_{i=1} c_i (x,a)\frac{\partial f_i}{\partial x_j}\qquad
(j=n-r+1,\ldots, n)\\
a_0 &+ \sum^n_{i=1} a_i x_i =0
\end{align*}
在 $V'$ 上无解。这样的条件显然关于 $(a)$ 是可构造的：设 $X=V'\times A$，其中 $A$ 是以 $a_0,\ldots,a_n$ 为坐标的仿射空间，并设 $Y\subset X$ 是这些方程的解组成的闭子集。所述条件就是 $Y/A$ 的纤维为空，这确实是一个可构造\pageoriginale 条件（\EGA IV 9.5.1）。因此，只需 % \ref
验证 $A$ 的一般点上的纤维为空，而这由方程的形式立即得到。(i) 得证。(ii)：考虑由经过 $P$ 的 $d$ 次（$d\geq 2$）超曲面线性系给出的有理映射 $\PP^n\rightarrow \PP^N$。众所周知（也容易验证），这是 $P$ 的一个“爆破”（blowing up），因此给出一个局部闭浸入
$$
f:\PP^n-\{P\}\longrightarrow \PP^N{\quoi}.
$$
设 $V''=V'-(\{P\}\cap V')$，$\bar{V}{}''=\sisi{f}{f}(V'')$，并设 $L_i$ 是 $\PP^N$ 中对应于超曲面 $H_i$ 的超平面 $(i=1,\ldots,s)$。则 $V''$ 与 $\bar{V}{}''$ 同构，而且根据 $f$ 的定义，$Y$ 与 $V''$ 横截相交，等价于 $L_1\cap\ldots \cap L_s$ 与 $\bar{V}{}''$ 横截相交。令 $\bar{V}=$ $\bar{V}{}''$ 的闭包，便归结为对 $\bar{V}$ 应用断言 (i)，由此得到将 $V'$ 换成 $V''$ 后的 (ii)。还需考虑点 $P$（如果它是 $V$ 的光滑点）。但显然，若 $\{H_1,\ldots,H_s\}$ 充分一般且 $P$ 是 $V$ 的光滑点，则 $Y$ 在 $P$ 与 $V$ 横截相交，结论得证。


\section{良好邻域的构造}\etiquette[3]{XI.sec3}

\begin{enonce*}{定义 3.1}\etiquette[3.1]{XI.def3.1}
初等纤维化（elementary fibration）是一个概形态射 ${f:X\rightarrow S}$，它可以嵌入一个交换图
$$
\xymatrix@=2cm{
X  \ar[r]^{j} \ar[rd]_{f}& \bar{X} \ar[d]^{\bar{f}}& Y \ar[l]_{i}  \ar[ld]^{g}\\
               &    S   &      
}
$$
并满足\pageoriginale 下列条件：
\begin{enumerate}
\renewcommand{\theenumi}{\roman{enumi}}
\renewcommand{\labelenumi}{(\theenumi)} 
% \item $i$ is a closed immersion.
\item $j$ 是开浸入，在每个纤维中稠密，且 $X=\bar{X}-Y$。

\item $\bar{f}$ 是光滑射影态射，其纤维为几何不可约的 $1$ 维概形。

\item $g$ 是平展覆盖，而且 $g$ 的每个纤维均非空。
\end{enumerate}
\end{enonce*}


容易验证，$X$ 的这样的嵌入如果存在，则在典范同构意义下唯一。

\begin{enonce*}{定义 3.2}\etiquette[3.2]{XI.def3.2}
相对于 $S$ 的良好邻域（good neighborhood）是一个 $S$-概形 $X$，使得存在 $S$-概形
$$
X=X_n,\ldots,X_0=S
$$
以及初等纤维化 $f_i:X_i\rightarrow X_{i-1}$，
% $f_i$ is an $S$-morphism
$i=1,\ldots, n$。
\end{enonce*}

\begin{enonce*}{命题 3.3}\etiquette[3.3]{XI.prop3.3}
\footnote{这个结果容易推广为半局部的断言。}
设 $k$ 是代数闭域，$X/\Spec k$ 是光滑概形，$x\in X$ 是一个有理点。则存在 $X$ 的一个包含 $x$ 的开子集，它是一个良好邻域（相对于 $\Spec k$）。
\end{enonce*}

证明。取 $X$ 不可约。对 $\dim X=n$ 作归纳，只需找到 $x$ 在 $X$ 中的一个邻域 $U$ 和一个初等纤维化 $f:U\rightarrow V$，使得 $V$ 光滑且维数为 $n-1$。事实上，将存在 $v=f(x)$ 的一个良好邻域 $V'$，于是可取 $U'=U\cap \sisi{f^{-1}(V')}{V'}$ 作为 $x$ 的良好邻域。

可以假定 $X\subset \EE^r$ 是仿射的。设 $X_0$ 是 $X$ 在 $\PP^r$ 中的闭包。设 $\bar{X}$ 是 $X_0$ 的正规化，并设 $Y=\bar{X}-X$，赋予它诱导的约化结构。此外，设\pageoriginale $S\subset \bar{X}$ 为奇异点组成的闭子集。有 $S\subset Y$，且
\begin{align*}
\dim \bar{X} &= \dim X=n{\quoi},\\
\dim Y &= n-1{\quoi},\\
\dim S &\leq n-2{\quoi}.
\end{align*}

利用形如 $M^{\otimes r}$ 的可逆层 $L$ 将 $\bar{X}$ 嵌入射影空间 $\PP^N$，其中 $M$ 很丰沛，$r\geq 2$。于是存在 $\PP^N$ 中的超平面 % linear
$H_1,\ldots, H_{n-1}$，其中 $H_i$ 是方程
$$
\sum^N_{\nu =0}a_{i\nu} x_{\nu} =0{\quoi},
$$
的零点集，使它们包含 $x$，而且交集 $L=H_1\cap\ldots \cap H_{n-1}$ 的维数为 $N-n+1$，并与 $\bar{X}$、$Y$ 横截相交（\hyperref[XI.thm2.1]{2.1}）。
% see the comment on Theorem 2.1
 %For the intersection with $\bar{X}$: $S\cap L=∅$ for sufficiently
 %general $L$, hence $\bar{X}\cap L=(\bar{X} - Y)\cap L$ and
 %$\bar{X}- Y$ is smooth.
 %For the intersection with $Y$: $Y$ is reduced and $k$ algebraically closed,
 %hence every irreducible component is generically smooth, and sufficiently
 %general $L$ does not meet the singular locus of $Y$.
交集 $\bar{X}\cap L$ 是一条光滑连通曲线（这由 Bertini 定理得到），而 $Y\cap L$ 的维数为 $0$。再选取另一个超平面
$$
H_0:\sum^N_{\nu=0}a_{0 \nu}x_{\nu}=0
$$
使得 $H_0$ 与 $\sisi{\bar{X}}{X}\cap L$ 横截相交，且 $H_0\cap Y\cap L=∅$。 % One must also choose $x\notin H_0$ so that
                    % the map $f'$ below is well defined at $x$

考虑通过“引入射影坐标”
$$
y_i=\sum^N_{\nu =0}a_{i_{\nu}}x_{\nu}{\quoi}.
$$
得到的投影 $\PP^N\rightarrow \PP^{n-1}$。这是在投影中心 $C=H_0\cap\ldots \cap H_{n-1}$ 之外有定义的有理映射。设 $\epsilon :P'\rightarrow \PP^N$ 为沿 $C$ 的爆破，从而有图
$$
\xymatrix@=1.5cm{
\PP^{N}  \ar[rd]& &P'\ar[ld]^{\pi}\ar[ll]_{\epsilon}\\
               &\PP^{n-1} &
}
$$
其中\pageoriginale $\pi$ 为一个态射。设 $\bar{X}{}'\subset P'$ 为 $\bar{X}$ 的“严格”逆像（proper inverse image），即 $\epsilon^{-1}(\bar{X}-(\bar{X}\cap C))$ 的闭包。由于按假设 $C$ 与 $\bar{X}$ 横截相交，态射 $\bar{X}{}'\rightarrow \bar{X}$ 是沿有限集 $\bar{X}\cap C$ 的爆破。设 $X'=X-\bar{X}\cap C$，也将它等同于 $\bar{X}{}'$ 的一个开子概形；并设 $Y'=\bar{X}{}'-X'$ 为赋予诱导约化结构的闭子概形。有如下态射图 $\DD$
$$
\xymatrix@=2cm{
  X'  \ar[r]^{j} \ar[rd]_{f'}& \bar{X}{}'  \ar[d]^{\bar{f}{}'} & Y'\ar[l]_{i} \ar[ld]^{g'}     &\\
                           & \PP^{n-1} &\qquad ,   
}
$$
而我断言，存在 $v=f'(x)$ 的一个邻域 $V$，使 $\DD$ 限制到 $V$ 后满足（\hyperref[XI.def3.1]{3.1}）中的条件 (i)、(ii)、(iii)，因而 $f'_{\mid V}$ 是一个初等纤维化。这将完成证明。

条件 (i) 是平凡的。对于 (ii)，注意按假设 $\bar{X}\cap L$ 是一条光滑曲线，而且立即可见，$\epsilon$ 诱导一个一一态射 ${\bar{f}{}'}^{-1}(v)\rightarrow \bar{X}\cap L$。
% one-to-one morphism?
因此，$\left({\bar{f}{}'}^{-1}(v)\right)_{\red}\rightarrow \bar{X}\cap L$ 是双射。为了验证 $\bar{f}{}'$ 在 $v$ 的一个邻域上光滑，
% remove the \left and \right?
由 Hironaka 引理（\SGA 1 II 2.6 %\ref
\footnote{以及 \EGA IV 5.12.10。}），只需验证它在 ${\bar{f}{}'}^{-1}(v)$ 的一般点光滑。在这一点，$\bar{X}{}'$ 与 $\bar{X}$ 同构，而该态射光滑，因为 $L$ 与 $\bar{X}$ 横截相交。

还需证明 $g'$ 在 $v$ 的一个邻域上是平展的（因为 $Y$ 的维数为 $n-1$，显然 $g'$ 的每个纤维均非空）。有 $Y'=\epsilon^{-1}(Y)\amalg D_1\amalg\ldots \amalg D_r$，其中 $\bar{X}\cap C=\{P_1,\ldots,P_r\}$，而 $D_i$ 是 $\bar{X}$ 中 $P_i$ 的爆破\pageoriginale。每个 $D_i$ 等同于 $\epsilon^{-1}(P_i)$，且同构地映到 $\PP^{n-1}$ 上\sisi{。此外，投影}{}在 $Y$ 的每一点都有定义%since $Y\cap C=∅$
，而且在 $Y$ 上诱导的态射就是 $g'_{\mid \epsilon^{-1}(Y)}$。由于 $L$ 与 $Y$ 横截相交，它在 $v$ 上是平展的，因此 $g'$ 在~$v$ 上平展。


\section{比较定理}\etiquette[4]{XI.sec4}

\setcounter{subsection}{-1}
\subsection{}\etiquette[4.0]{XI.subsec4.0}
设 $X$ 是 $\Spec \CC$ 上局部有限型的概形，并在 $X$ 的有理点空间 $X(\CC)$ 上考虑通常拓扑。以 $X_{\cl}$ 表示由\emph{局部同构}（local isomorphism）$f:U\rightarrow X(\CC)$ 组成的景，即由满足下述性质的拓扑空间态射 $f$ 组成的景：

\sisi{对}{对}每个 $x\in U$，存在开邻域 $U_x$，使得 $f_{\mid U_x}$ 是 $U_x$ 到 $f(x)$ 的一个开邻域的同胚。

$X_{\cl}$ 中的一族态射 $\{U_{\nu}\rightarrow U\}$ 称为\emph{覆盖}的，如果 $U$ 是各 $U_{\nu}$ 的像的并。

由于开浸入是局部同构，我们有一个景的态射（包含）（IV\quad ） % a reference is missing
\footnote{更确切地说，范畴的包含 $\Ouv(X(\CC))\rightarrow X_{\cl}$ 定义了反方向的景态射 \hyperref[XI.eq4.1]{(4.1)}。}
\begin{equation*}
\delta:X_{\cl}\longrightarrow X(\CC)\qquad
(\ast){\quoi}.\tag{4.1}
\end{equation*}\etiquette[4.1]{XI.eq4.1}
现在，对于每个局部同构 $U\rightarrow X(\CC)$，按定义，都存在由 $X(\CC)$ 的开子集给出的 $U$ 的一个“覆盖”，由（\quad ）可推出，这两个景所关联的拓扑斯 % a reference is missing
是等价的（通过 $\delta_{\ast}$）。特别地，计算通常上同调时，可以把 $X(\CC)$ 换成 $X_{\cl}$。

再设 $f:X'\rightarrow X$ 是一个平展概形态射。则 $f(\CC):X'(\CC)\rightarrow X(\CC)$ 是局部同构。这由\pageoriginale Jacobian 判别法（\SGA 1 II 4.10）和隐函数定理立即得到。因此，函子 $X'\mapsto X'(\CC)$ 给出一个景的态射
\begin{equation*}
\epsilon:X_{\cl}\longrightarrow X_{\et}{\quoi},\tag{4.2}
\end{equation*}\etiquette[4.2]{XI.eq4.2}
其中 $X_{\et}$ 是平展景。

回顾以下结果\footnote{这些结果见第~2 页脚注所引的那一讲。}： % identify the Exposés
                                % above to identify the Exposé
                                % mentioned in the note?

\begin{enonce*}{定理 4.3}\etiquette[4.3]{XI.thm4.3}
\begin{enumerate}
\renewcommand{\theenumi}{\roman{enumi}}
\renewcommand{\labelenumi}{(\theenumi)}
\item $X$ 连通且非空，当且仅当 $X(\CC)$ 连通且非空。

\item 函子 $\epsilon$ 是全忠实的。

\item （“Grauert--Remmert 定理”）。函子 $\epsilon$ 在下述两个范畴之间诱导等价：赋予局部紧空间的通常拓扑的 $X(\CC)$ 的有限覆盖范畴，以及 $X$ 的平展覆盖范畴。
\end{enumerate}
\end{enonce*}

我们把 (i) 作为已知结果接受。例如，它可由 (GAGA)
%\ref
\cite{sga4.exp.XI.bib.4} 推出。断言 (ii) 容易由 (i) 得到。事实上，为了验证 (ii) 中包含的满射性断言，设 $\varphi:X'(\CC)\rightarrow X''(\CC)$ 是一个 $X(\CC)$-态射。可以假定所有概形都是连通仿射的，因而是分离的。于是，$\varphi$ 的图像 $\Gamma$ 是 $\fprod{X'}{X''}{X}(\CC)$ 的一个连通分支，由 (i)，存在 $\fprod{X'}{X''}{X}$ 的某个连通分支 $Y$，使得 $\Gamma=Y(\CC)$。于是 $Y$ 是诱导 $\varphi$ 的一个态射 $f:X'\rightarrow X''$ 的图像，即投影 $Y\rightarrow X'$ 是同构：事实上，$Y(\CC)\rightarrow X'(\CC)$ 是同构，由此容易推出 $Y\rightarrow X'$ 也是同构（它是平展、根性（radicial）且满的！）。对于 (iii)，由 (ii)，问题完全归结为证明 $X(\CC)$ 的每个有限覆盖都形如 $X'(\CC)$，其中 $X'\rightarrow X$\pageoriginale 是平展的。

为此，利用下降（\SGA 1 IX 4.7）立即归结为
%\ref
$X$ 连通且正规的情形。该问题关于 $X$ 是局部的（对 Zariski 拓扑而言），因此可以假定 $X$ 是仿射的。设 $f:X\rightarrow \EE^n$ 是有限态射，$T'\rightarrow X(\CC)$ 是连通有限覆盖。立即可验证，复合态射 $\eta':T'\rightarrow \EE^n(\CC)$ 是（\cite{sga4.exp.XI.bib.2} \S 2，定义~3）意义下的“解析覆盖”（analytic covering）。由（\cite{sga4.exp.XI.bib.2} \S 2，Satz 8），\sisi{$\eta'$}{}可扩张为一个“解析覆盖” $\eta:T\rightarrow \PP^n(\CC)$，而只需证明 $\eta$ 是由适当的概形态射 $f:Y\rightarrow \PP^n$ \sisi{诱导}{诱导}的，其中 $Y$ 正规。现在，基本定理（\cite{sga4.exp.XI.bib.2} \S 13 Satz 42）断言，$T$ 上存在一个典范的正规解析结构，使得 $\eta$ 诱导一个解析空间态射。由于此时 $\eta$ 是有限态射，$T$ 上结构层的正像是 $\PP^n$ 上的一个凝聚解析代数层 $\mathfrak{a}$。由 \sisi{(GAGA)}{GAGA} \cite{sga4.exp.XI.bib.4} % \ref
可知，$\mathfrak{a}$ 是与 $\PP^n$ 上某个代数层 $A$ 关联的解析层，于是取 $Y=\sheaf{\Spec}A$。

现在可以证明以下结果：

\begin{enonce*}{定理 4.4}\etiquette[4.4]{XI.thm4.4}
设 $X$ 是 $\Spec \CC$ 上的光滑概形。
\begin{enumerate}
\renewcommand{\theenumi}{\roman{enumi}}
\renewcommand{\labelenumi}{(\theenumi)}
\item $X_{\et}$ 上的可构造局部常值挠层范畴，与 $X_{\cl}$ 上具有有限纤维的局部常值挠层范畴之间存在等价，此等价由互为拟逆的函子 $\epsilon_{\ast}$ 和 $\epsilon^{\ast}$ 给出。

\item 设 $F$ 是 $X_{\cl}$ 上具有有限纤维的局部常值挠层，并仍以 $F$ 表示 $X_{\et}$ 上诱导的层 $\epsilon_{\ast}F$。则
$$
R^q\epsilon_{\ast}F=0\quad \mbox{ 若 } q>0{\quoi}.
$$

\item 使用\pageoriginale {\rm(i)} 中的记号，典范态射
$$
H^q(X_{\et},F)\longrightarrow H^q(X_{\cl},F)
$$
对每个 $q$ 都是双射。特别地，对每个 $q$ 都有
$$
H^q(X_{\et}, \ZZ/\sisi{n\ZZ}{n})\xrightarrow{\sim} H^q(X_{\cl},\ZZ/\sisi{n\ZZ}{n})
$$
成立。
\end{enumerate}
\end{enonce*}

证明。对于 (i)，注意所讨论的层就是能由 $X$ 的平展覆盖（\resp $X(\CC)$ 的有限覆盖）表出的层。由于这些覆盖的范畴之间存在等价（\hyperref[XI.thm4.3]{4.3}），断言 (i) 立即得到。(iii) 中的同构由 (ii) 及 $\epsilon$ 的 Leray 谱序列 V \hyperref[V.prop5.2]{5.2} % \ref + 5.3 instead of 5.2?
得到。还需证明 (ii)。

现在，$R^q\epsilon_{\ast}F$ 是与预层 $\Rr^qF,$ 关联的层，其中，对平展态射 $X'\rightarrow X$，有 $(\Rr^qF)(X')=H^q(X'_{\cl},F)$。因此，为了证明当 $q>0$ 时 $R^q\epsilon_{\ast}F=0$，必须证明以下\sisi{引理}{}：

\begin{enonce*}{引理 4.5}\etiquette[4.5]{XI.lem4.5}
设 $\xi\in H^q(X_{\cl},F)$，$x\in X_{\cl}=X(\CC)$。存在像包含 $x$ 的平展态射 $X'\rightarrow X$，使得 $\xi$ 在 $H^q(X'_{\cl},F)$ 中的像为零。
\end{enonce*}

引理的证明。对 $n=\dim X$ 作归纳；问题关于 $X$ 对平展拓扑是局部的，\sisi{}{因而}可以假定 $F$ 为常值层，再利用分步化简（dévissage），甚至可以假定 $F=\ZZ/\sisi{n\ZZ}{n}$。此外，将 $X$ 换成 $x$ 的一个 Zariski 开邻域后，可以假定（\hyperref[XI.prop3.3]{3.3}）$X$ 有一个初等纤维化（\hyperref[XI.def3.1]{3.1}）$f:X\rightarrow S$。沿用（\hyperref[XI.def3.1]{3.1}）的记号：通过\pageoriginale 直接计算，可见当 $q>1$ 时有 $R^q{j}_{\ast}F=0$，$R^1j_{\cl\!\ast}F$ 是 $Y_{\cl}$ 上一个常值层的补零延拓（extension by $0$），最后，$j_{\cl\!\ast}F$ 是 $\bar{X}$ 上的常值层 $\ZZ/\sisi{n\ZZ}{n}$。此外，可以逐纤维计算 $R^q \bar{f}_{\cl\!\ast}(j_{\cl\!\ast}F)$。由 Leray 谱序列
$$
E^{pq}_2=R^p\bar{f}_{\cl\!\ast}(R^qj_{\cl\!\ast}F)\Longrightarrow
R^{p+q}f_{\cl\!\ast}F
$$
可推出 $R^{p+q}f_{\cl\!\ast}F$ 同样可以逐纤维计算，因此

$R^0f_{\cl\!\ast}\ZZ/\sisi{n\ZZ}{n}=\ZZ/\sisi{n\ZZ}{n}$，

$R^1f_{\cl\!\ast}\ZZ/\sisi{n\ZZ}{n}$ 是 $S_{\cl}$ 上的一个局部常值挠层，

当 $q>1$ 时，$R^qf_{\cl\!\ast}\ZZ/\sisi{n\ZZ}{n}=0$。

于是，Leray 谱序列

$E^{pq}_2=H^p(S_{\cl}, R^qf_{\cl\!\ast}F)\Longrightarrow
H^{p+q}(X_{\cl},F)$

\noindent
化为正合列
\begin{equation*}
\ldots\longrightarrow H^q(S_{\cl},f_{\cl\!\ast}F)\longrightarrow
H^q(X_{\cl},F)\longrightarrow H^{q-1}(S_{\cl},
R^1f_{\cl\!\ast}F)\longrightarrow\tag{$\ast$}
\end{equation*}
设 $s$ 为 $x$ 在 $S$ 中的像。由归纳假设及 (i)，对于每个类 $\eta\in H^q(S_{\cl},G)$（$G$ 为具有有限纤维的局部常值挠层），存在 $s$ 的一个“平展邻域” $S'\rightarrow S$，使 $\eta$ 在 $H^q(S'_{\cl},G)$ 中的像为零。现在，函子 $R^qf_{\cl\!\ast}$ 显然与平展基变换可交换，而且若 $S'\rightarrow S$ 是 $s$ 的平展邻域，则 $X'=\fprod{X}{S'}{S}\rightarrow X$ 是 $x$ 的平展邻域。因此，当 $q\geq 1$ 时，利用正合列 $(\ast)$，先取 $G=R^1f_{\cl\!\ast}F$，再\pageoriginale 取 $G=f_{\cl\!\ast}F$，便得到结论。 % The proof ends here; the
                                % following is a remark for the case $q=1$
当 $q=1$ 时，注意 $H^1(X_{\cl},\ZZ/n \ZZ)$（\resp $H^1(X,\ZZ/n\ZZ)$）分类 $X_{\cl}$（\resp $X$）的群为 $\ZZ/n\ZZ$ 的主平展覆盖；
% should one instead speak of finite coverings for $X_\cl$?
因而由 \hyperref[XI.thm4.3]{4.3} (iii)，映射 $H^1(X,\ZZ/n\ZZ)\rightarrow H^1(X_{\cl},\ZZ/n\ZZ)$ 是双射，由此立即得到 $q=1$ 时的消去结果 \hyperref[XI.lem4.5]{4.5}。

\begin{enonce*}[remark]{变体 4.6}\etiquette[4.6]{XI.var4.6}
还可以用下面这个归于 Serre 的巧妙方法证明（\hyperref[XI.lem4.5]{4.5}）：注意一个连通良好邻域（\hyperref[XI.def3.2]{3.2}）$X$ 满足下列两个条件：
\begin{enumerate}
\renewcommand{\theenumi}{\roman{enumi}}
\renewcommand{\labelenumi}{(\theenumi)}
\item 若 $n>1$，则 $\pi_n(X(\CC))=0$。

\item $\pi_1(X(\CC))$ 是有限型自由群的迭次扩张。
\end{enumerate}
\end{enonce*}

事实上，若 $X\rightarrow S$ 是一个初等纤维化，其中 $S$ 是良好邻域，则由归纳法，可以假定 (i)、(ii) 对 $S$ 成立。现在，容易看出 $X(\CC)\rightarrow S(\CC)$ 是局部平凡纤维空间，其纤维同构于 $X_0(\CC)$，其中 $X_0$ 为 $X/S$ 的任意一个纤维。此时 $X_0$ 是连通非奇异的\emph{非完备}曲线，而众所周知，这蕴含当 $n>1$ 时 $\pi_n(X_0(\CC))=0$，且 $\pi_1(X_0(\CC))$ 是自由群。 % \ref
于是 (i)、(ii) 可由同伦正合列得到：
$$
\longrightarrow \pi_n(X_0(\CC))\longrightarrow
\pi_n (X(\CC))\longrightarrow \pi_n
(S(\CC))\longrightarrow\ldots
$$

所以 $X(\CC)$ 是一个 $K(\pi,1)$ 空间，这蕴含每个上同调类 $\xi\in H^n(X(\CC),F)$ 在 $X(\CC)$ 的万有覆盖上变为零。只需证明 $\xi$ 在 $X$ 的某个有限覆盖 $X'$ 上已经变为零，这归结为证明 $\pi_1=\pi_1(X(\CC))$ 是\pageoriginale 一个\emph{好群}（good group）（\cf CG 1 16）， %\ref
也就是说，从 $\pi_1$ 到其完备化 $\sisi{\widehat}{}{\pi}_1=\varprojlim_N \pi_1/N$ 的态射（其中 $N$ 遍历 $\pi_1$ 的所有有限指数不变子群）
% 'invariant' = 'normal'?
对每个连续有限 $\pi_1$-模 $M$ 都诱导一个双射
$$
H^q(\widehat{\pi}_1,M)\xrightarrow{\sim} H^q(\pi_1,M)
$$
（其中左边的符号表示 $\widehat{\pi}_1$ 作为概有限群的上同调）。而有限型自由群的迭次扩张是好群（\cf CG I p.~15-16, exc.~1,2），所以结果得证。 % \ref

\begin{thebibliography}{99}
\bibitem{sga4.exp.XI.bib.1} GIRAUD J., ``Analysis Situs'', Séminaire
Bourbaki, Exposé nᵒ~256.

\bibitem{sga4.exp.XI.bib.2} GRAUERT H\sisi{.}{} and REMMERT R., ``Komplexe Räume'',
Math, Ann. Bd. 136, p. 345 (1958).

\bibitem{sga4.exp.XI.bib.3} HIRONAKA H., ``Resolution of singularities of an
algebraic variety over a field of characteristic zero'' Annals of
Math. Vol. 79, nᵒ~1, 2 (1964).

\bibitem{sga4.exp.XI.bib.4} SERRE J.P., ``Géométrie Algébrique et Géométrie
Analytique'', Annales de l'Institut Fourier, t. VI, p. 1-12 (1956).
\end{thebibliography}

% END native input: chapters/exp-11-zh.tex
% BEGIN native input: chapters/exp-12-zh.tex
\chapter{固有态射的基变换定理}
%\addtocontents{toc}{M. Artin}
\begin{center}
M. Artin
\end{center}
\phantomsection\MakeLinkTarget*{sga4.exp.XII}\label{sga4.exp.XII}

\setcounter{pageoriginal}{78}

\section{引言}\etiquette[1]{XII.sec1}\pageoriginale

本讲证明如下基本定理: 更精确地说, 对于概形间的\emph{固有}(proper) 态射 $f:X\to Y$ 及 $X$ 上的挠阿贝尔层 $F$, $R^qf_*F$ 在 $Y$ 的几何点 $\xi$ 处的茎同构于纤维 $X_\xi$ 的上同调 $H^q(X_\xi,F|X_\xi)$.

回顾一下, 仿紧空间间固有态射的类似结果是众所周知而且容易的 (Godement II 4.11). 对概形而言, 这一结果要微妙得多. 例如, 与仿紧空间的情形不同, $F$ 为挠层的假设是必不可少的.

证明分为几个阶段. 通过或多或少是形式的逐层剥离 (dévissage) 论证, 将问题归结为 $Y$ 是 Noether 概形、$f$ 的相对维数至多为 $1$ 且层 $F$ 为常层的情形. 为处理这一特殊情形, 要用到基本群的特殊化定理 (\hyperref[XII.thm5.9]{5.9}); 它是上述定理的一个非阿贝尔变体, 在这里所考虑的情形下可以直接证明 (\hyperrefext[XIII][XIII.sec2]{2}).

\section{一个例子}\etiquette[2]{XII.sec2}\pageoriginale

我们说明 $F$ 为挠层的假设是必要的. 设 $X$ 是代数闭域 $k$ 上的非奇异曲面, $Y$ 是 $k$ 上的非奇异曲线, $f:X\to Y$ 是固有态射, 其一般纤维光滑, 且在 $Y$ 的某个几何点 $\bar y$ 处的几何纤维不可约, 并有一个普通二重点 (ordinary double point). 取 $F$ 为常层 $\ZZ_X$.

回顾下述事实. 对于概形 $Z$ 的任意点 $z$, 设 $i:z\to Z$ 为嵌入, 则 $H^1(Z,i_*\ZZ_z)=0$ (\hyperrefext[IX][IX.prop3.6]{3.6} (i)). 由于 $X$ 正规, 因而是单枝的, 所以 $\ZZ_X\simeq i_*\ZZ_x$, 其中 $i:x\to X$ 为一般点的嵌入. 因此 $H^1(X,\ZZ_X)=0$, 同样也有 $R^1f_*\ZZ_X=0$. 但对于 $X/Y$ 的具有普通二重点的几何纤维 $X_\circ$, 我断言 $H^1(X_\circ,\ZZ_{X_\circ})\simeq\ZZ$\footnote[1]{从 $\Pi_1$ 的观点比较 \SGA 3 X (6, 例) 及 (1.6).}; 这就给出了反例, 因为 $\ZZ_{X_\circ}\simeq\ZZ_X|X_\circ$.

设 $i_\circ:x_\circ\to X_\circ$ 为 $X_\circ$ 的一般点的嵌入, $Q$ 为 $X_\circ$ 的奇点. 曲线 $X_\circ$ 在 $Q$ 处有两个“分枝”, 从而 $i_{\circ*}(\ZZ_{x_\circ})$ 在 $Q$ 处的茎同构于 $\ZZ\times\ZZ$. 在 $Q$ 之外, 曲线 $X_\circ$ 正规, 因此在 $X_\circ-Q$ 上有 $i_{\circ*}(\ZZ_{x_\circ})\simeq\ZZ_{X_\circ}$. 所以有正合列
$$
0\longrightarrow\ZZ_{X_\circ}\longrightarrow
i_{\circ*}(\ZZ_{x_\circ})\longrightarrow(\ZZ)_Q
\longrightarrow0{\quoi},
$$
其中末项是点 $Q$ 上的层 $\ZZ$ 的零延拓. 这三项的 $H^0$ 都同构于 $\ZZ$, 并且\linebreak $H^1(X_\circ,i_{\circ*}(\ZZ_{x_\circ}))=0$. 因此上同调长正合列给出 $H^1(X_\circ,\ZZ_{X_\circ})=\ZZ$\pageoriginale.

\section{非阿贝尔 $H^1$ 的回顾}\etiquette[3]{XII.sec3}

我们简要回顾下面将要用到的非阿贝尔上同调的一些结果. 大多数证明均略去; 它们在各种特殊情形下都是熟知的, 可见 Giraud 的学位论文 [1].

设 $X$ 为概形, $F$ 为 $X$ 上的群层. 对每个平展 $X'/X$, 可以定义带基点集合 (pointed set) $H^1(X',F)$, 它对 $F$ 具有协变函子性, 对 $X'$ 具有反变函子性. 函子 $H^1$ 与层的有限积可交换. 若 $f:X\to Y$ 为态射, 则可以定义带基点集合层 (即具有指定截面的层) $R^1f_*F$. 它是带基点集合预层 $\Rr^1F$ 的伴随层, 后者将每个平展 $Y'/Y$ 对应于集合 $H^1(\fprod{X}{Y'}{Y},F)$.

有通常的上同调正合列; 我们记录其中一个:

\begin{enonce*}{命题 3.1}\etiquette[3.1]{XII.prop3.1}
设 $0\to F\xrightarrow{u}G$ 为群层的单射, $C=G/F$ 为 $G$ 的齐性空间层. 有带基点集合层的正合列
$$
0\to f_*F\to f_*G\to f_*C\xrightarrow{\partial}
R^1f_*F\xrightarrow{u^1}R^1f_*G{\quoi}.
$$
\end{enonce*}

当然, 这一正合列由预层的正合列取伴随层而得; 后者在平展 $Y'/Y$ 上的值为
$$
0\to F(X')\to G(X')\to C(X')\xrightarrow{\partial}
H^1(X',F)\to H^1(X',G){\quoi},
$$
其中\pageoriginale $X'=\fprod{X}{Y'}{Y}$. $\partial$ 所诱导的等价关系就是 $f_*G$ 在 $f_*C$ 上的作用所给出的等价关系. 和通常一样, 可以利用“扭转”(twisting) 的概念 [1] 明确写出 $u^1$ 所诱导的等价关系.

\begin{enonce*}{命题 3.2}\etiquette[3.2]{XII.prop3.2}
设 $X\xrightarrow{f}Y\xrightarrow{g}Z$ 为概形态射. 对每个平展 $Z'/Z$, 置 $Y'=\fprod{Y}{Z'}{Z}$ 及 $X'=\fprod{X}{Z'}{Z}$, 则有正合列
$$
0\longrightarrow H^1(Y',f_*F)\longrightarrow H^1(X',F)
\longrightarrow H^0(Y',R^1f_*F){\quoi}.
$$
取伴随层, 得到 $Z$ 上带基点层的正合列:
$$
0\longrightarrow (R^1g_*)f_*F\longrightarrow R^1(gf)_*F
\longrightarrow g_*(R^1f_*)F{\quoi}.
$$
\end{enonce*}

\begin{enonce*}{命题 3.3}\etiquette[3.3]{XII.prop3.3}
形如 $H^1(X',F)$ (其中 $X'$ 在 $X$ 上平展) 以及形如 $R^1f_*F$ (其中 $f:X\to Y$ 为任意态射) 的函子族是可消去的 (effaceable): 存在单射 $F\xrightarrow{u}G$, 使所有诱导态射 $H^1(X',F)\to H^1(X',G)$ 及 $R^1f_*F\to R^1f_*G$ 均为零. 更精确地说, 存在这样的单射, 使每个 $H^1(X',G)$ (其中 $X'$ 在 $X$ 上平展) 以及每个带基点层 $R^1f_*G$ 都为零. 若 $X$ 是 Noether 概形, 且 $F$ 是可构造的 ind-$\LL$ 群层, 则存在这样的单射, 使 $G$ 也为 ind-$\LL$ 群层.
\end{enonce*}

\emph{证明.} 第一个断言大概在任意拓扑斯中都成立, 但对于平展拓扑, 利用“Godement 消解”来进行消去是方便的:
\begin{equation*}
F\longrightarrow G=\prod_{x\in X}i_{x*}i_x^*F{\quoi}.
\tag{3.4}
\end{equation*}\etiquette[3.4]{XII.eq3.4}
这里\pageoriginale $x$ 遍历 $X$ 的各点, $i_x:\bar x\to X$ 是 $x$ 上方的一个几何点. 对于平展 $X'/X$, 有
$$
G(X')=\prod_{x\in X}(i_x^*F)(\fprod{X'}{\bar x}{X})
=\prod_{x\in X}\left(\prod_{\bar x'/\bar x}F_{\bar x'}\right){\quoi},
$$
其中 $\bar x'$ 是 $X'$ 的几何点. 显然 $X'$ 的每个点上方至少有一个 $\bar x'$; 所以 $F(X')\subset G(X')$ (VIII 3), 因而 $F\subset G$. 此外, 每个平展 $\bar u/\bar x$ 都是可分闭域的谱的不交并, 所以对 $\bar x$ 上的每个层 $\bar G$ 都有 $H^1(\bar u,\bar G)=0$. 由 \hyperref[XII.prop3.2]{3.2} 可知, 对每个平展 $X'/X$ 都有 $H^1(X',i_{x*}i_x^*F)=0$. 由于 $H^1$ 与积可交换, 有 $H^1(X',G)=0$, 因而对每个 $f:X\to Y$ 都有 $R^1f_*G=0$; 这是由将 $R^1f_*G$ 描述为预层 $Y'\mapsto H^1(\fprod{X}{Y'}{Y},G)$ 的伴随层得到的.

再假设 $F$ 可构造, $X$ 是 Noether 概形. 立即可知, 存在 $X$ 的有限点集 $\{x_1,\ldots,x_n\}$, 使
$$
F\longrightarrow\prod_{1\leq i\leq n}i_{x_i*}i_{x_i}^*F
$$
已经为单射. 层 $i_{x_i*}i_{x_i}^*F$ 是 ind-有限群层 (IX \hyperref[IX.def1.5]{1.5}), 因而其有限积也是 (\hyperrefext[IX][IX.prop1.4]{1.4} 及 \hyperrefext[IX][IX.prop1.6]{1.6} (iii)). 以此层代替 $G$, 就得到最后一个断言.

\setcounter{subsection}{4}
\subsection{}\etiquette[3.5]{XII.subsec3.5}
与阿贝尔上同调的情形一样 (见 V \hyperref[V.subsec2.4]{2.4}), $H^1$ 可以用 \v{C}ech 方法计算 [1].

\section{基变换态射}\etiquette[4]{XII.sec4}\pageoriginale

设
$$
\xymatrix@C=2.75cm@R=1.5cm{
X \ar[d]_{f}&X'\ar[l]_{g'}\ar[d]^{f'}\\
S&S'\ar[l]_{g}
}
$$
是概形态射的笛卡尔方块. 由此按如下方式得到函子态射 $\varphi=\varphi_g$:
\begin{equation*}
g^{\ast}f_{\ast}\xrightarrow{\varphi}f'_{\ast}{g'}^{\ast}\tag{4.1}
\end{equation*}\etiquette[4.1]{XII.eq4.1}
由 $g^{\ast}$ 与 $g_{\ast}$ 之间的伴随性, 给定这样的态射等价于给定函子态射
\begin{equation*}
f_{\ast}\longrightarrow g_{\ast}{f'}_{\ast}{g'}^{\ast}{\quoi}.
\tag{4.1 adj}
\end{equation*}\MakeLinkTarget*{sga4.exp.XII.equation.4.1-adj}\etiquette[4.1 adj]{4.1 adj}
我们取为如下态射:
$$
f_{\ast}\xrightarrow{(\id\rightarrow {g'}_{\ast}{g'}^{\ast})}
f_{\ast}{g'}_{\ast}{g'}^{\ast}
=g_{\ast}{f'}_{\ast}{g'}^{\ast}{\quoi},
$$
这里利用了
$$
g_{\ast}{f'}_{\ast}=(gf')_{\ast}=(fg')_{\ast}
=f_{\ast}{g'}_{\ast}{\quoi}.
$$

\eqhyperref[XII.eq4.1]{4.1} 还有另一个候选. 事实上, 给定这样的态射也等价于给定
$$
{f'}^{\ast}g^{\ast}f_{\ast}\longrightarrow {g'}^{\ast},
$$
而我们可以取
$$
{f'}^{\ast}g^{\ast}f_{\ast}\simeq
{g'}^{\ast}f^{\ast}f_{\ast}
\xrightarrow{(f^{\ast}f_{\ast}\rightarrow\id)}{g'}^{\ast}{\quoi}.
$$

我不知道\pageoriginale 这两种选取是否总是一致\footnote[1]{\textsc{Deligne} 先生在广泛得多的背景下肯定地解决了这一疑问; \cf 第 XVII 讲.}; 无论如何, 我们只使用第一种. 称它为\emph{基变换态射}(base-change morphism).

将 \eqhyperref[XII.eq4.1]{4.1} 限制到阿贝尔层 (分别地, 群层), 就对每个 $q\geq0$ (分别地, $q=0,1$) 得到态射 $\varphi^q=\varphi^q_g$:
\begin{equation*}
g^{\ast}(R^qf_{\ast})\xrightarrow{\varphi^q}
(R^q{f'}_{\ast}){g'}^{\ast}{\quoi}.
\tag{4.2}
\end{equation*}\etiquette[4.2]{XII.eq4.2}
事实上, 这等价于给定
\begin{equation*}
(R^qf_{\ast})\longrightarrow
g_{\ast}(R^q{f'}_{\ast}){g'}^{\ast}{\quoi},
\tag{4.2 adj}
\end{equation*}\MakeLinkTarget*{sga4.exp.XII.equation.4.2-adj}\etiquette[4.2 adj]{4.2 adj}
而有如下“显然”的态射:
\begin{gather*}
(R^qf_{\ast})\xrightarrow{a}
(R^qf_{\ast}){g'}_{\ast}{g'}^{\ast}
\xrightarrow{b}R^q(fg')_{\ast}{g'}^{\ast}\\
=R^q(gf')_{\ast}{g'}^{\ast}
\xrightarrow{c}g_{\ast}(R^q{f'}_{\ast}){g'}^{\ast}{\quoi}.
\end{gather*}

\begin{enonce*}[remark]{注 4.3}\etiquette[4.3]{XII.rem4.3}
假设 $g$ (因而 $g'$) 为闭浸入. 则上述态射 $b$ 与 $c$ 为同构 (VIII \hyperref[VIII.cor5.7]{5.7}). 因此 \eqhyperref[XII.eq4.2]{4.2} 实质上正是上述态射 $a$, 它由对典范态射 $F\rightarrow {g'}_{\ast}{g'}^{\ast}F$ 应用 $R^qf_{\ast}$ 而得.
\end{enonce*}

\begin{enonce*}{命题 4.4}\etiquette[4.4]{XII.prop4.4}
(基变换态射的复合).
\begin{enumerate}
\renewcommand{\theenumi}{\roman{enumi}}
\renewcommand{\labelenumi}{(\theenumi)}
\item 设
$$
\xymatrix@C=2.5cm@R=1.5cm{
X\ar[d]_{f}&X'\ar[l]_{g'}\ar[d]_{f'}&
X''\ar[l]_{h'}\ar[d]_{f''}\\
S&S'\ar[l]_{g}&S''\ar[l]_{h}
}
$$
是概形态射的\pageoriginale 笛卡尔图表. 则下图交换:
$$
\xymatrix@C=2cm@R=1.5cm{
(gh)^{\ast}f_{\ast}\ar[rr]^{\varphi_{gh}}
\ar[d]^{\text{\rotatebox{90}{$\backsim$}}}&&
f''_{\ast}(g'h')^{\ast}
\ar[d]^{\text{\rotatebox{90}{$\backsim$}}}\\
h^{\ast}g^{\ast}f_{\ast}
\ar[r]^{h^{\ast}\varphi_g}&
h^{\ast}f'_{\ast}g'^{\ast}
\ar[r]^{\varphi_h g'^{\ast}}&
f''_{\ast}h'^{\ast}g'^{\ast}.}
$$
在有关函子的自然定义范围内, 将 $f_{\ast}$、$f'_{\ast}$ 及 $f''_{\ast}$ 分别换成 $R^qf_{\ast}$、$R^qf'_{\ast}$ 及 $R^qf''_{\ast}$ 所得的图表也同样交换 \textup{(\cf 第 $\textup{n}^{\circ}$ 3 节)}.

\item 设
$$
\xymatrix@C=3.25cm@R=1.25cm{
Y\ar[d]_{e}&Y'\ar[l]_{g''}\ar[d]^{e'}\\
X\ar[d]_{f}&X'\ar[l]_{g'}\ar[d]^{f'}\\
S&S'\ar[l]_{g}
}
$$
是概形态射的笛卡尔图表. 则下图交换:
$$
\xymatrix@C=2.75cm@R=1.5cm{
g^{\ast}(fe)_{\ast}\ar[rr]^{\varphi_g^{fe}}\ar@{=}[d]&&
(f'e')_{\ast}g''^{\ast}\ar@{=}[d]\\
g^{\ast}f_{\ast}e_{\ast}
\ar[r]^{\varphi_g^f e_{\ast}}&
f'_{\ast}g'{}^{\ast}e_{\ast}
\ar[r]^{f'_{\ast}\varphi_{g'}^e}&
f'_{\ast}e'_{\ast}g''^{\ast},
}
$$
其中三个水平箭头分别来自 $fe$、$f$ 及 $e$ 的基变换同态.

在阿贝尔情形下, $fe$ 的基变换态射由如下谱序列态射诱导:
$$
\xymatrix{
E^{p,q}_2=g^{\ast}(R^pf_{\ast})(R^qe_{\ast})F
\ar@{=>}[r]\ar[d]_{\varphi}&
g^{\ast}R^{p+q}(fe)_{\ast}F\ar[d]_{\varphi}\\
(R^pf'_{\ast})(R^qe'_{\ast})g''^{\ast}F
\ar@{=>}[r]&R^{p+q}(f'e')_{\ast}g''^{\ast}F,
}
$$
它在起始项上由 $e$ 与 $f$ 的基变换态射诱导. 对于群层, 有正合列态射
$$
\xymatrix{
0\ar[r]&g^{\ast}(R^1f_{\ast})e_{\ast}F
\ar[r]\ar[d]^{\varphi}&
g^{\ast}R^1(fe)_{\ast}F\ar[r]\ar[d]^{\varphi}&
g^{\ast}f_{\ast}(R^1e_{\ast})F\ar[d]^{\varphi}\\
0\ar[r]&(R^1f'_{\ast})e'_{\ast}g''^{\ast}F\ar[r]&
R^1(f'e')_{\ast}g''^{\ast}F\ar[r]&
f'_{\ast}(R^1e'_{\ast})g''^{\ast}F.
}
$$
\end{enumerate}
\end{enonce*}

证明留给读者.

\section{主要定理的陈述及若干变体}\etiquette[5]{XII.sec5}

定理如下:

\begin{enonce*}{定理 5.1}\etiquette[5.1]{XII.thm5.1}
(固有基变换定理). 设
$$
\xymatrix@C=3.25cm@R=1.5cm{
X \ar[d]_{f} & X' \ar[l]_{g'} \ar[d]^{f'}\\
S & S' \ar[l]_{g}
}
$$
是概形的笛卡尔图表, 其中 $f$ 固有.
\begin{enumerate}
\renewcommand{\theenumi}{\roman{enumi}}
\renewcommand{\labelenumi}{(\theenumi)}
\item 设 $F$ 是 $X$ 上的集合层. 则态射 \eqhyperref[XII.eq4.1]{4.1}
$$
\varphi:g^{\ast}f_{\ast}F\longrightarrow {f'}_{\ast}{g'}^{\ast}F
$$
为双射.

\item 设\pageoriginale $F$ 是 $X$ 上的群层 (分别地, ind-有限群层). 则态射 \eqhyperref[XII.eq4.2]{4.2}
$$
\varphi^1:g^{\ast}(R^1f_{\ast})F\longrightarrow
(R^1{f'}_{\ast}){g'}^{\ast}F
$$
为单射 (分别地, 双射).

\item 设 $F$ 是 $X$ 上的挠阿贝尔层. 则态射 \eqhyperref[XII.eq4.2]{4.2}
$$
\varphi^q:g^{\ast}(R^qf_{\ast})F\longrightarrow
(R^q{f'}_{\ast}){g'}^{\ast}F
$$
对每个 $q\geq0$ 均为双射.
\end{enumerate}
\end{enonce*}

先给出定理的一些直接推论. 若 $S'=\xi$ 是 $S$ 的几何点, 因而 $\xi$ 上的层由其整体截面决定, 则对 $X'$ 上的每个层 $G$ 都有 $H^0(\xi,R^q{f'}_{\ast}G)=H^q(X',G)$. 于是得到下述推论 (此外, 它实质上与 \hyperref[XII.thm5.1]{5.1} 等价):

\begin{enonce*}{推论 5.2}\etiquette[5.2]{XII.cor5.2}
设 $f:X\longrightarrow S$ 固有, $\xi\longrightarrow S$ 为几何点, $X_{\xi}$ 为 $f$ 在 $\xi$ 处的纤维.
\begin{enumerate}
\renewcommand{\theenumi}{\roman{enumi}}
\renewcommand{\labelenumi}{(\theenumi)}
\item 若 $F$ 为 $X$ 上的集合层, 则典范基变换态射
$$
(f_{\ast}F)_{\xi}\longrightarrow H^0(X_{\xi},F|X_{\xi})
$$
为双射.

\item 若 $F$ 为群层 (分别地, ind-有限群层), 则基变换态射
$$
(R^1f_{\ast}F)_{\xi}\longrightarrow H^1(X_{\xi},F|X_{\xi})
$$
为单射 (分别地, 双射).

\item 若\pageoriginale $F$ 为挠阿贝尔层, 则对每个 $q\geq0$, 基变换态射
$$
(R^qf_{\ast}F)_{\xi}\longrightarrow H^q(X_{\xi},F|X_{\xi})
$$
为双射.
\end{enumerate}
\end{enonce*}

考虑到 \hyperrefext[X][X.cor4.3]{4.3} 及 \hyperrefext[X][X.cor5.2]{5.2}, 得到:

\begin{enonce*}{推论 5.3}\etiquette[5.3]{XII.cor5.3}
在 \hyperref[XII.cor5.2]{5.2} 的假设下, 设 $n$ 为纤维 $X_{\xi}$ 的维数. 则对 $X$ 上的每个挠阿贝尔层 $F$, 茎 $(R^qf_{\ast}F)_{\xi}$ 在 $q>2n$ 时为零. 若 $F$ 为 $p$ 挠群层, 其中 $p$ 是 $\xi$ 的特征, 则 $R^qf_{\ast}F$ 在 $\xi$ 处的茎在 $q>n$ 时为零.
\end{enonce*}

回顾一下, 若每个纤维的维数均 $\leq n$, 就称 \emph{$X/S$ 的相对维数} $\leq n$. 因而有:

\begin{enonce*}{推论 5.3 bis}\etiquette[5.3 bis]{XII.corbis5.3}
设 $f:X\longrightarrow S$ 固有, $X/S$ 的相对维数 $\leq n$. 则对 $X$ 上的每个挠阿贝尔层 $F$, 当 $q>2n$ 时有 $R^qf_{\ast}F=0$. 若 $S$ 的特征为 $p$, 且 $F$ 为 $p$ 挠阿贝尔群层, 则当 $q>n$ 时有 $R^qf_{\ast}F=0$.
\end{enonce*}

将 \hyperref[XII.thm5.1]{5.1} 应用于可分闭域的谱之间的态射, 得到:

\begin{enonce*}{推论 5.4}\etiquette[5.4]{XII.cor5.4}
(固有情形下上同调在基域扩张下的不变性). 设 $k\subset K$ 为可分闭域, $X$ 为 $\Spec k$ 上的固有概形. 置 $X'=X_K$. 设 $F$ 为 $X$ 上的层, $F'$ 为它在 $X'$ 上的逆像层. 则:
\begin{enumerate}
\renewcommand{\theenumi}{\roman{enumi}}
\renewcommand{\labelenumi}{(\theenumi)}
\item $H^0(X,F)\xrightarrow{\sim}H^0(X',F')$.

\item 若\pageoriginale $F$ 为群层 (分别地, ind-有限群层), 则态射
$$
H^1(X,F)\longrightarrow H^1(X',F')
$$
为单射 (分别地, 双射).

\item 若 $F$ 为挠阿贝尔层, 则对每个 $q$ 都有
$$
H^q(X,F)\xrightarrow{\sim}H^q(X',F')
$$
\end{enumerate}
\end{enonce*}

\hyperref[XII.cor5.2]{5.2} 的一个略有不同的变体如下:

\begin{enonce*}{推论 5.5}\etiquette[5.5]{XII.cor5.5}
设 $S$ 为 Hensel 环的谱, $S'=s_{\circ}$ 为其闭点, $f:X\longrightarrow S$ 固有. 置 $X_{\circ}=X'$ 及 $F_{\circ}={g'}^{\ast}F$. 则:
\begin{enumerate}
\renewcommand{\theenumi}{\roman{enumi}}
\renewcommand{\labelenumi}{(\theenumi)}
\item $H^0(X,F)\xrightarrow{\sim}H^0(X_{\circ},F_{\circ})$

对每个集合层 $F$ 都成立.

\item 态射 $H^1(X,F)\longrightarrow H^1(X_{\circ},F_{\circ})$

对每个群层 (分别地, ind-有限群层) $F$ 为单射 (分别地, 双射).

\item $H^q(X,F)\xrightarrow{\sim}H^q(X_{\circ},F_{\circ})$

对每个挠阿贝尔层 $F$ 及每个 $q$ 都成立.
\end{enumerate}
\end{enonce*}

事实上, 当 $S$ 是严格局部概形时, 此陈述只是 \hyperref[XII.cor5.2]{5.2} 的特殊情形. 假设已知 \hyperref[XII.cor5.2]{5.2}. 设 $\bar S$ 为 $S$ 的严格局部化, $\bar s$ 为 $\bar S$ 的闭点, $G$ 为 $\bar S/S$ 的 Galois 群 (它也是 $\bar s/s_{\circ}$ 的 Galois 群). 则 $G$ 作用在 $R^q$ 的茎上, 并且采用显然的记号, 对集合层 $F$ 有
\begin{align*}
H^0(X,F)&\simeq H^0(\bar X,\bar F)^G,\\
H^0(X_{\circ},F_{\circ})&\simeq
H^0(\bar X_{\circ},\bar F_{\circ})^G
\end{align*}
\pageoriginale 由于 $H^0(\bar X,\bar F)\xrightarrow{\sim}
H^0(\bar X_{\circ},\bar F_{\circ})$, 而且该同构显然与 $G$ 的作用可交换, 得到断言 (i).

若 $F$ 为群层, 容易定义带基点集合正合列的态射
$$\vcenter{\etiquette[5.5.1]{XII.eq5.5.1}
\xymatrix{
0 \ar[r] & H^1(G,H^0(\bar X,\bar F)) \ar[r]\ar[d] &
H^1(X,F) \ar[r]\ar[d] & H^1(\bar X,\bar F)^G \ar[d]\\
0 \ar[r] & H^1(G,H^0(\bar X_{\circ},\bar F_{\circ})) \ar[r] &
H^1(X_{\circ},F_{\circ}) \ar[r] &
H^1(\bar X_{\circ},\bar F_{\circ})^G ,
}}\leqno{(5.5.1)}
$$
其中左边的竖直箭头为同构, 由 \hyperref[XII.cor5.2]{5.2}, 右边的竖直箭头为单射. 设 $\alpha,\beta\in H^1(X,F)$ 为在 $H^1(X_{\circ},F_{\circ})$ 中具有相同像的类, $F^{\alpha}$ 为将 $F$ 以 $\alpha$ 扭转所得的层 \cite{sga4.exp.XII.bib.1}. 设 $\beta'$ 为 $H^1(X,F^{\alpha})$ 中对应于 $\beta$ 的类. 由于对应于 $\alpha$ 的类 $\alpha'$ 是中性的, $\beta'$ 在 $H^1(X_{\circ},(F^{\alpha})_{\circ})$ 中的像是中性的.

考虑在 \eqhyperref[XII.eq5.5.1]{5.5.1} 中以 $F^{\alpha}$ 代替 $F$ 所得的图表. 它说明 $\beta'$ 在 $H^1(\bar X,\bar F^{\alpha})^G$ 中的像是中性的, 因而 $\beta'$ 来自 $H^1(G,H^0(\bar X,\bar F^{\alpha}))$ 的一个元素; 此元素本身必为中性的. 因此 $\beta'$ 是中性的, 即 $\beta'=\alpha'$, 从而 $\beta=\alpha$. 这就证明了 (ii) 中态射的单射性.

现在假设 $F$ 是 ind-有限群层, 并证明 (ii) 中箭头的满射性. 只须证明函子 $F\mapsto H^1(X_{\circ},F_{\circ})$ 可消去. 事实上, 设 $0\longrightarrow F\xrightarrow{u}G$ 是消去该函子的单射, 并置 $C=G/F$. 有正合列态射 (\hyperref[XII.prop4.4]{4.4})
$$
\xymatrix{
H^0(X,C) \ar[r]\ar[d] & H^1(X,F) \ar[d]_{\epsilon}\ar[r] &
H^1(X,G) \ar[d]\\
H^0(X_{\circ},C_{\circ}) \ar[r] &
H^1(X_{\circ},F_{\circ}) \ar[r]^{u_{\circ}^1} &
H^1(X_{\circ},G_{\circ})
}
$$
其中 $u_{\circ}^1$ 为零; 这确实蕴含 $\epsilon$ 的满射性.

来证明\pageoriginale 对于 ind-有限群层 $F$, 函子 $F\mapsto H^1(X_{\circ},F_{\circ})$ 可消去. 设 $F\xrightarrow{u}C$ 为“Godement 消解” \eqhyperref[XII.eq3.4]{3.4}. 有
$$
H^1(G,H^0(\bar X_{\circ},\bar C_{\circ}))\simeq
H^1(G,H^0(\bar X,\bar C))\subset H^1(X,C)=0{\quoi}.
$$

此外, 有交换方块
$$
\xymatrix{
H^1(\bar X,\bar F) \ar@{-}[r]^{\sim}\ar[d] &
H^1(\bar X_{\circ},\bar F_{\circ})\ar[d]\\
0=H^1(\bar X,\bar C) \ar[r] &
H^1(\bar X_{\circ},\bar C_{\circ})
}
$$
(事实上, $H^1(\bar X,\bar C)=\varinjlim H^1(X',C)=0$, 其中 $X'/X$ 遍历平展态射), 所以右边的竖直箭头为零. 于是有正合列态射
$$
\xymatrix{
0 \ar[r] & H^1(G,H^0(\bar X_{\circ},\bar F_{\circ}))
\ar[r]\ar[d] & H^1(X_{\circ},F_{\circ}) \ar[r]\ar[d]^{\beta} &
H^1(\bar X_{\circ},\bar F_{\circ})^G \ar[d]^{\alpha^G}\\
&0 \ar[r] & H^1(X_{\circ},C_{\circ}) \ar[r] &
H^1(\bar X_{\circ},\bar C_{\circ})^G
}
$$
其中 $\alpha^G$ 为零. 因而 $\beta$ 也为零, 所以该函子可消去.

最后证明 \hyperref[XII.cor5.5]{5.5} 的断言 (iii). 有谱序列态射
$$
\xymatrix{
**[l]E_2^{p,q}=H^p(G,H^q(\bar X,\bar F))
\ar@{=>}[r]\ar[d] & H^n(X,F)\ar[d]^{\varphi}\\
H^p(G,H^q(\bar X_{\circ},\bar F_{\circ}))
\ar@{=>}[r] & H^n(X_{\circ},F_{\circ}).
}
$$
由 \hyperref[XII.cor5.2]{5.2}, 态射 $H^q(\bar X,\bar F)\longrightarrow
H^q(\bar X_{\circ},\bar F_{\circ})$ 为同构. 因而 $E_2^{p,q}$ 项的态射为同构, 从而收敛项的态射也是同构. 结果得证, \cqfd

\setcounter{subsection}{5}
\subsection{}\etiquette[5.6]{XII.subsec5.6}
来考察\pageoriginale $F$ 为常层时的 \hyperref[XII.cor5.5]{5.5} (i) 与 (ii). 设 $F=T_X$ 为 $X$ 上以 $T$ 为值的常集合层. 则 $F(X)$ 就是 $X$ 到 $T$ 的局部常值映射的集合. 若 $X_{\circ}$ 为 $X$ 的子概形, 且 $T$ 至少有两个元素, 则 $F(X)\longrightarrow F(X_{\circ})$ 为双射当且仅当映射
$$
U\longmapsto U_{\circ}=U\cap X_{\circ}
$$
为双射, 这里它是从 $X$ 的既开又闭的子集所成集合 $\Theta(X)$ 到类似集合 $\Theta(X_{\circ})$ 的映射. 若 $X_{\circ}$ 的连通分支集合 $\Pi_{\circ}(X_{\circ})$ 有限, 则这等价于映射 $\Pi_{\circ}(X_{\circ})\longrightarrow\Pi_{\circ}(X)$ 为双射. 因而 \hyperref[XII.cor5.5]{5.5} (i) 也可重述如下 (注意在这种情形下, $X_{\circ}$ 在域上有限型, 因而是 Noether 概形, 所以 $\Pi_{\circ}(X_{\circ})$ 有限):

\begin{enonce*}{推论 5.7}\etiquette[5.7]{XII.cor5.7}
设 $S$ 为 Hensel 环的谱, $s_{\circ}$ 为其闭点, $f:X\longrightarrow S$ 为固有态射, $X_{\circ}$ 为 $X/S$ 的闭纤维. 则 $X_{\circ}\longrightarrow X$ 所诱导的映射
$$
\Pi_{\circ}(X_{\circ})\longrightarrow\Pi_{\circ}(X)
$$
为双射.
\end{enonce*}

(特别地, 这说明 $\Pi_{\circ}(X)$ 有限.) 我们将在 Noether 情形下直接证明 \hyperref[XII.cor5.7]{5.7}:

\begin{enonce*}{引理 5.8}\etiquette[5.8]{XII.lem5.8}
若 $S$ 为 Noether 概形, 则陈述 \hyperref[XII.cor5.7]{5.7} 成立\footnote[1]{\Cf \EGA IV 18.5.19.}.
\end{enonce*}

\emph{证明.} 在 Noether 情形下, $\Pi_{\circ}(X)$ 当然有限. 将 $X$ 换成一个连通分支, 可以假设 $\Pi_{\circ}(X)$ 只有一个元素, 即 $X$ 连通且非空. 必须证明 $X_{\circ}$ 也连通且非空. $X$ 在 $S$ 中的像因 $f$ 固有而闭, 又因 $X$ 非空而非空; 所以它包含 $s_{\circ}$, 这说明 $X_{\circ}$ 非空. 设
$$
\xymatrix@R=1cm@C=1.5cm{
X \ar[rd]^{f'} \ar[dd]_{f}& \\
  & S'\ar[ld]^{\pi}\\
S &
}
$$
为 $f$ 的\pageoriginale Stein 分解 (\EGA III 4.3.3), 于是 $S'/S$ 有限而 $S'$ 连通. 已知 $f'$ 的每个纤维都连通且非空, 而 $X_{\circ}$ 的底空间是 $S'$ 的各闭点上方的纤维的底空间之并. 只须证明 $S'$ 是某个\emph{局部}环的谱. 由于 $S$ 是 Hensel 环的谱, $S'$ 的环 $A'$ 是局部环之积 (VII \hyperref[VII.prop4.1]{4.1} (i)); 又因 $S'$ 连通, $A'$ 是局部环. 这就证明了 \hyperref[XII.lem5.8]{5.8}.

\setcounter{subsection}{9}
\setcounter{subsubsection}{-1}
\subsubsection{}\etiquette[5.9.0]{XII.subsubsec5.9.0}
现在假设 $F=G_X$ 为以有限群 $G$ 为值的常群层. 则 $H^1(X,G_X)$ 分类 $X$ 的群为 $G$ 的主 Galois 覆盖 (VII \hyperref[VII.rem2.1]{2.1}). 推论 \hyperref[XII.cor5.5]{5.5} (ii) 说态射 $X_{\circ}\longrightarrow X$ 诱导这些覆盖的同构类集合之间的双射. 根据 $\Pi_1$ 的定义, \hyperref[XII.cor5.5]{5.5} (ii) 对于常有限群层的结论 (连同 (i)) 等价于:

\begin{enonce*}{定理 5.9}\etiquette[5.9]{XII.thm5.9}
(基本群的特殊化定理). 采用 \hyperref[XII.cor5.5]{5.5} 的记号, 假设 $X$ 连通且非空. 则 $X_{\circ}$ 连通且非空; 并且, 若 $\bar x_{\circ}$ 为 $X_{\circ}$ 的几何点, 则基本群态射
$$
\Pi_1(X_{\circ},\bar x_{\circ})\longrightarrow
\Pi_1(X,\bar x_{\circ})
$$
为同构.
\end{enonce*}

另一方面\pageoriginale, 由 Galois 理论 (\SGA V), 陈述 \hyperref[XII.thm5.9]{5.9} (在 $X_{\circ}$ 连通的情形下; 由 \hyperref[XII.cor5.7]{5.7} 可以这样假设) 等价于:

\begin{enonce*}{定理 5.9 bis}\etiquette[5.9 bis]{XII.thmbis5.9}
设 $S$ 为 Hensel 环的谱, $s_{\circ}$ 为 $S$ 的闭点, $f:X\longrightarrow S$ 为固有态射, $X_{\circ}$ 为 $X/S$ 的闭纤维. 对概形 $Z$, 以 $\Et(Z)$ 表示 $Z$ 的平展覆盖范畴. 则限制函子
$$
\Et(X)\longrightarrow\Et(X_{\circ})
$$
为范畴等价.
\end{enonce*}

\begin{enonce*}[remark]{注 5.10}\etiquette[5.10]{XII.rem5.10}
回顾一下, 当 $S$ 为完备 Noether 局部环的谱时, 已经证明了 \hyperref[XII.thm5.9]{5.9} (\SGA 1 X 2.1). 遗憾的是, 这一结果似乎不能给出一般情形下 \hyperref[XII.thm5.9]{5.9} 的证明\footnote[1]{自写下这些文字以来, 情况已经改变; \cf M. Artin, Algebraic approximation of structures over complete local rings (待刊) 中的定理~3.1.}. 本讲及下一讲给出的 \hyperref[XII.thm5.9]{5.9} 的证明, 与 \SGA 1 X 中所给的证明很不相同.
\end{enonce*}

\setcounter{subsection}{10}
\subsection{}\etiquette[5.11]{XII.subsec5.11}
利用 J. \textsc{Giraud} 的 2-上同调理论 (\cf J. \textsc{Giraud} 正在撰写的学位论文), 可以给出 \hyperref[XII.cor5.5]{5.5} (ii) 的另一个“非交换”变体. 假设 $S$ 为 Noether 概形. 设 $L$ 为 $X$ 上的一个\emph{带}(band), 并假设 $L$ 在局部 (相对于平展拓扑) 由一个 ind-有限群层表示. 考虑映射
$$
\varphi^{(2)}:H^2(X,L)\longrightarrow H^2(X_{\circ},L_{\circ}){\quoi},
$$
其中 $L_{\circ}$ 为 $L$ 在 $X_{\circ}$ 上的限制. 则
\begin{equation*}
(\varphi^{(2)})^{-1}(H^2(X_{\circ},L_{\circ})')
=H^2(X,L)'{\quoi},\tag{5.11.1}
\end{equation*}\etiquette[5.11.1]{XII.eq5.11.1}
其中\pageoriginale “撇号”表示由中性元素组成的子集.

由 \hyperref[XII.cor5.5]{5.5} (iii), 这一陈述也等价于下述陈述:

\emph{对 $X$ 上带为上述 $L$ 的每个 gerbe（层群胚）$\mathcal G$, 以 $\mathcal G_{\circ}$ 表示它在 $X_{\circ}$ 上的限制. 则截面范畴之间的限制函子}
$$
\mathcal G(X)\longrightarrow\mathcal G_{\circ}(X_{\circ})
$$
\emph{为范畴等价.}


事实上, 此函子完全忠实由 \hyperref[XII.cor5.5]{5.5} (i) 得到; 当 $\mathcal G(X)\neq\varnothing$, 即 $\mathcal G$ 的类 $\xi\in H^2(X,\sisi{\mathcal G}{L})$ 为中性时, 此函子本质满由 \hyperref[XII.cor5.5]{5.5} (ii) 得到. 还须表达如下事实: $\mathcal G_{\circ}(X_{\circ})\neq\varnothing$ 蕴含 $\mathcal G(X)\neq\varnothing$, 而这恰为 \eqhyperref[XII.eq5.11.1]{5.11.1}.

这里不给出 \eqhyperref[XII.eq5.11.1]{5.11.1} 的证明. 只指出, 对任意 Noether 概形 $X$ 及 $X$ 的任意闭子概形, 事实上可以证明, 陈述 \hyperref[XII.cor5.5]{5.5} (i)、(ii) 的成立蕴含 \eqhyperref[XII.eq5.11.1]{5.11.1} 的成立. 证明相当初等, 利用下降结果 VIII \hyperref[VIII.thm9.4]{9.4} a).

\section{初步归约}\etiquette[6]{XII.sec6}

用缩写 $(f,F,g)$ 表示 \hyperref[XII.thm5.1]{5.1} 在该定理的 (i)、(ii)、(iii) 任一情形中的数据.

\begin{enonce*}{引理 6.1}\etiquette[6.1]{XII.lem6.1}
要使 \hyperref[XII.thm5.1]{5.1} 对数据 $(f,F,g)$ 成立, 其中 $f$ 与 $F$ 固定而 $g$ 任意, 只须定理对 $(\bar f,\bar F,\bar g)$ 成立, 这里 $\bar S$ 遍历如下 $S$ 概形: 它是某个在 $S$ 上局部\pageoriginale 有限型的概形 $S'$ 在其纤维中的一个闭点处的严格局部化. 这里 $\bar f$ 与 $\bar F$ 由 $f$ 与 $F$ 经基变换 $\bar S\to S$ 而得, $\bar g:\bar s\to\bar S$ 为 $\bar S$ 的闭点的嵌入.
\end{enonce*}

\noindent
\emph{证明.} 显然可以假设 $S$ 与 $S'$ 仿射, 设分别为环 $A$ 与 $A'$ 的谱. 则 $A'$ 为其有限型 $A$ 子代数 $A'_i$ 的滤过正向极限. 利用取极限的理论 (VII \hyperref[VII.cor5.11]{5.11} 及 \hyperref[VII.rem5.14]{5.14}), 只须对三元组 $(f,F,g_i)$ 证明 \hyperref[XII.thm5.1]{5.1}, 其中 $g_i:\Spec(A'_i)=S'_i\to S$; 这就归结为 $g$ 为有限型的情形. 为证明基变换同态为同构, 只须证明它们在有关层的茎上诱导同构, 并且只须考虑 $S'$ 的点 $s'$ 处的几何茎, 其中这些点在 $S$ 上的纤维中为闭点 (VIII \hyperref[VIII.rem3.13]{3.13}). 设 $\bar s'$ 为 $S'$ 的几何点, 对应于 $k(s')$ 的某个可分闭包. 由于 $s'$ 在 $S$ 的相应点 $s$ 上代数, $k(\bar s')$ 也是 $k(s)$ 的可分闭包. 以 $\bar s$ 表示视为 $S$ 上的预概形 $\bar s'$, 即视为 $S$ 的几何点, 并考虑 $S$ 与 $S'$ 分别在 $\bar s$ 与 $\bar s'$ 处的严格局部化 $\bar S$ 与 $\bar S'$. 于是有交换方块
$$
\vcenter{
\xymatrix@C=3cm@R=1.25cm{
\bar{s}'  \ar[r] \ar[d]& \bar{s}\ar[d]\\
\bar{S}'  \ar[r] &  \bar{S},
}}
$$
其中第一个水平箭头为同构. 由 VIII \hyperref[VIII.thm5.2]{5.2} 及 \hyperref[VIII.rem5.3]{5.3}, 与 $(f,F,g)$ 关联的基变换同态在 $\bar s'$ 处的茎上所诱导的同态, 等同于由 $\bar X'\to\bar X$ 诱导\pageoriginale 的同态
$$
H^{i}(\bar{X},\bar{F})\longrightarrow H^{i}(\bar{X}',\bar{F}')
$$
其中 $\bar X$ 与 $\bar X'$ 分别表示 $X$ 经基变换 $\bar S\to S$ 及 $\bar S'\to S$ 所得的概形. 上面的交换方块显然给出交换方块
$$
\xymatrix@C=3cm@R=1.5cm{
H^{i}(\bar{X}_{\bar{s}}, \bar{F}) \ar[r]&
H^{i}(\bar{X}'_{\bar{s}'},\bar{F}')\\
H^{i}(\bar{X},\bar{F})  \ar[r] \ar[u]&
H^{i}(\bar{X}',\bar{F}') \ar[u],
}
$$
其中第一个水平箭头仍为同构. 因此, 只要两个竖直箭头都是同构, 第二个水平箭头就是同构. 这恰是 \hyperref[XII.lem6.1]{6.1} 的假设所保证的, \cqfd

\setcounter{subsection}{1}
\subsubsection{}\etiquette[6.1.1]{XII.subsubsec6.1.1}
因此, 为在情形 $(F,f,-)$ 下证明 \hyperref[XII.thm5.1]{5.1}, 只须研究特殊情形 \hyperref[XII.cor5.5]{5.5}, 其中 $S$ 严格局部. 这引导我们一般地研究概形态射 $h:Y\to X$, 并寻求使相应同态
$$
H^{i}(X,F)\longrightarrow H^{i}(Y,h^{\ast}(F))
$$
为双射的一般条件. 本小节给出这一情形下的一些简单辅助结果. 为便于以后引用, 我们以比证明 \hyperref[XII.thm5.1]{5.1} 所需的更一般、更精确的形式陈述这些结果.

\begin{enonce*}{引理 6.2}\etiquette[6.2]{XII.lem6.2}
设 $C$ 与 $C'$ 为两个阿贝尔范畴, $T^{\bigdot}=(T^{i})$ 与 $T'^{\bigdot}=({T'}^{i})$ 为从 $C$ 到 $C'$ 的两个上同调函子, $\varphi=(\varphi^{i})$ 为从 $T^{\bigdot}$ 到 $T'^{\bigdot}$ 的上同调函子\pageoriginale 态射. 假设 $T^{\bigdot}$ 与 $T'^{\bigdot}$ 在次数 $<0$ 时为零, 且 $T^{i}$ 对 $i>0$ 可消去. 设 $n$ 为整数. 下述条件等价:
\begin{enumerate}
\renewcommand{\theenumi}{\alph{enumi}}
\renewcommand{\labelenumi}{\theenumi)}
\item $\varphi^{i}$ 对 $i\leq n$ 为同构, 对 $i=n+1$ 为单态射.

\item (若 $n\geq -1$.) 同态 $\varphi^{i}$ 对 $i=0$ 为单态射, 对 $i\leq n$ 为满态射.

\item (若 $n\geq 0$.) 同态 $\varphi^{\circ}$ 为同构, 且 ${T'}^{i}$ 对 $0<i\leq n$ 可消去.
\end{enumerate}

若这些条件成立且 $A\in\Ob(C)$, 则 $\varphi^{n+1}(A):T^{n+1}(A)\to {T'}^{n+1}(A)$ 为满态射 (从而为同构), 当且仅当 ${T'}^{n+1}(A)$ 可消去, 即存在单态射 $A\to B$, 使相应同态 ${T'}^{n+1}(A)\to {T'}^{n+1}(B)$ 为零. 若 $C'=(\Ab)$ 为阿贝尔群范畴, 且 $\xi\in {T'}^{n+1}(A)$, 则 $\xi$ 属于 $T^{n+1}(A)$ 的像, 当且仅当 $\xi$ 可消去, 即存在单态射 $A\to B$, 使 $\xi$ 在 ${T'}^{n+1}(B)$ 中的像为零.
\end{enonce*}

\noindent
\emph{证明.} 显然 (a) 蕴含 (b) 与 (c). 由于对 $T^{\bigdot}$ 的假设蕴含 $i>0$ 时的 $T^{i}$ 是 $T^{\circ}$ 的右卫星函子, 蕴含 (c) $\Rightarrow$ (a) 就由卫星函子的公理刻画得到. 它也可对 $n$ 作归纳, 从 \hyperref[XII.lem6.2]{6.2} 的最后一个断言得到; 后者可用熟知的论证证明, 这里不再重复 (比较下面 \hyperref[XII.cor6.6]{6.6} 的证明). 若 $n=-1$, 则蕴含 (b) $\Rightarrow$ (a) 是直接的. 若 $n\geq0$, 对 $n$ 作归纳: 归纳假设允许我们假设 $\varphi^{i}$ 对 $i\leq n-1$ 为同构, 对\pageoriginale $i=n$ 为单态射. 于是题设使它对 $i=n$ 也为同构, 只须利用 $\varphi^n$ 为同构, 证明它对 $i=n+1$ 为单态射. 设 $A\to B$ 为消去 $T^{n+1}$ 的单态射, 从而在 $C$ 中有正合列 $0\to A\to B\to C\to0$, 并在 $C'$ 中有正合列态射:
$$
\xymatrix@R=1.25cm{
T^{n}(B)  \ar[r] \ar[d]& T^{n}(C)  \ar[r]\ar[d] & T^{n+1}(A) \ar[d] \ar[r] & T^{n+1}(B)\ar[d]\\
T'^{n}(B)  \ar[r] & T'^{n}(C) \ar[r] & T'^{n+1}(A)  \ar[r] &
T'^{n+1}(B),
}
$$
其中由刚才的证明, 前两个竖直箭头为同构, 而 $T^{n+1}(A)\to T^{n+1}(B)$ 根据构造为零. 于是五引理蕴含 $T^{n+1}(A)\to {T'}^{n+1}(A)$ 为单态射, \cqfd

\begin{enonce*}{推论 6.3}\etiquette[6.3]{XII.cor6.3}
假设 $T^{\bigdot}$ 与 $T'^{\bigdot}$ 来自右导出范畴 $D^{+}(C)$ 上的上同调函子 (仍以相同符号表示), 取值于 $C'$ \cite{sga4.exp.XII.bib.3}; 不一定假设 $T^{i}$ 对 $i>0$ 可消去. 假设对 $K^{\bigdot}\in D^{+}(C)$, 条件“对 $i<0$ 有 $H^{i}(K^{\bigdot})=0$”蕴含“对 $i<0$ 有 $T^{i}(K^{\bigdot})={T'}^{i}(K^{\bigdot})=0$”. 设 $K^{\bigdot}\in\Ob D^{+}(C)$ 为 $C$ 中的下有界复形, 满足对 $i<0$ 有 $H^{i}(K^{\bigdot})=0$.
\begin{enumerate}
\renewcommand{\theenumi}{\roman{enumi}}
\renewcommand{\labelenumi}{(\theenumi)}
\item 若 \hyperref[XII.lem6.2]{6.2} 的条件 {\rm (a)} 满足, 则同态 $\varphi^{i}(K^{\bigdot})=T^{i}(K^{\bigdot})\to
  {T'}^{i}(K^{\bigdot})$ 对 $i\leq n$ 为同构, 对 $i=n+1$ 为单态射.

\item 再设 $L^{\bigdot}\in\Ob D^{+}(C)$ 满足对 $i<0$ 有 $H^{i}(L^{\bigdot})=0$, 且 $u:L^{\bigdot}\to K^{\bigdot}$ 为 $D^{+}(C)$ 中诱导单同态 $H^{\circ}(u):H^{\circ}(L^{\bigdot})\to
  H^{\circ}(K^{\bigdot})$ 的态射. 则同态 $T^{n+1}(L^{\bigdot})\to {T'}^{n+1}(L^{\bigdot})$ 为满射\pageoriginale, 当且仅当 ${T'}^{n+1}(L^{\bigdot})$ 在 ${T'}^{n+1}(K^{\bigdot})$ 中的像包含于 $T^{n+1}(K^{\bigdot})$ 的像. 若 $C'=(\Ab)$ 且 $\xi\in {T'}^{n+1}(L^{\bigdot})$, 则 $\xi$ 属于 $T^{n+1}(L^{\bigdot})$ 的像, 当且仅当其在 ${T'}^{n+1}(K^{\bigdot})$ 中的像属于 $T^{n+1}(K^{\bigdot})$ 的像.
\end{enumerate}
\end{enonce*}

回顾一下, $T^{\bigdot}$ 与 $T'^{\bigdot}$ 来自 $D^{+}(C)$ 上的上同调函子这一假设, 可以明确表述如下: 它们来自 $C$ 的下有界复形所成阿贝尔范畴上的上同调函子, 这些函子与次数平移可交换, 将零伦的复形态射映为 $C'$ 中的零态射, 并将每个为拟同构 (即诱导同构 $H^{i}(K^{\bigdot})\to H^{i}(L^{\bigdot})$) 的复形态射 $u:K^{\bigdot}\to L^{\bigdot}$ 映为 $C'$ 中的同构. 在实际中, 遇到的所有从 $C$ 到 $C'$ 的上同调函子都是这样来的.

\medskip
\noindent
\emph{\hyperref[XII.cor6.3]{6.3} 的证明.} 先证明 (i). 若对 $i\leq n+1$ 有 $H^{i}(K^{\bigdot})=0$, 断言显然成立, 因为加于 $T^{\bigdot}$ 与 $T'^{\bigdot}$ 的次数性质给出, 对 $i\leq n+1$ 有 $T^{i}(K^{\bigdot})={T'}^{i}(K^{\bigdot})=0$. 另一方面, 取整数 $m\geq0$ 使对 $i<m$ 有 $H^{i}(K^{\bigdot})=0$; 这样的整数存在, 例如取 $m=0$ 即可. 对 $m$ 作降归纳; 当 $m=n+2$ 时已经证明断言. 因此假设已对 $m'>m$ 证明断言, 并对 $m$ 证明之. 根据对 $K^{\bigdot}$ 次数的假设, 复形 $K^{\bigdot}$ 拟同构于一个集中在次数 $\geq m$ 的复形, 所以可以假设它本身就集中于这些次数. 于是有复形正合列
$$
0\longrightarrow H^m[-m]\longrightarrow K^{\bigdot}\longrightarrow
K'^{\bigdot}\longrightarrow 0{\quoi},
$$
其中 $H^m=H^m(K^{\bigdot})$, 记号 $[-m]$ 表示复形次数的 $-m$ 平移,\pageoriginale 而 $K'^{\bigdot}$ 定义为余核. 因此 $K'^{\bigdot}$ 的上同调对象, 在次数 $\geq m+1$ 为 $K^{\bigdot}$ 的上同调对象, 在次数 $\leq m$ 为零. 特别地, 归纳假设适用于 $K'^{\bigdot}$. 前述正合列给出正合列态射
$$
\xymatrix{
T^{i-1} (K'^{\;\bigdot}) \ar[d] \ar[r] & T^{i-m}(H^m)  \ar[r]\ar[d] &
T^{i}(K^{\bigdot}) \ar[r] \ar[d]& T^{i} (K'^{\;\bigdot})  \ar[r] \ar[d]&
T^{i+1-m}(H^m) \ar[d]\\
T'^{i-1}(K'^{\;\bigdot})  \ar[r] & T'^{i-m}(H^m)  \ar[r] &
T'^{i}(K^{\bigdot})  \ar[r] & T'^{i}(K'^{\;\bigdot})  \ar[r] & T'^{i+1-m}(H^m).
}
$$
在此图表中取 $i\leq n$, 利用归纳假设及 $m\geq0$, 得到紧邻中间竖直箭头两侧的竖直箭头为同构, 从而特别是满态射, 而最右边的竖直箭头为单态射. 因此五引理表明中间的竖直箭头为满态射. 取 $i\leq n+1$, 同样得到紧邻中间竖直箭头两侧的竖直箭头为单态射, 而最左边的竖直箭头为同构, 从而特别是满态射. 五引理于是表明中间的竖直箭头为单态射. 这就证明了 \hyperref[XII.cor6.3]{6.3} (i).

为证明第二个断言, 即 \hyperref[XII.cor6.3]{6.3} (ii), 同样使用正合列
$$
0\longrightarrow L^{\bigdot}\longrightarrow K^{\bigdot}\longrightarrow
K'^{\bigdot}\longrightarrow 0{\quoi},
$$
其中 $K'^{\bigdot}$ 定义为 $L^{\bigdot}\to K^{\bigdot}$ 的映射柱 (mapping cylinder), 并且仍仅在非负次数上有上同调. 在
$$
\xymatrix{
T^{n}(K^{\bigdot})  \ar[r] \ar[d]& T^{n}(K'^{\;\bigdot})  \ar[r] \ar[d]& T^{n+1}
(L^{\bigdot})  \ar[r]\ar[d] & {T^{n+1}(K^{\bigdot})}  \ar[r]
\ar[d]&{T^{n+1}(K'^{\;\bigdot})} \ar[d]\\
T'^{n}(K^{\bigdot}) \ar[r] & T'^{n}(K'^{\;\bigdot})  \ar[r] &
T'^{n+1}(L^{\bigdot}) \ar[r] & T'^{n+1}(K^{\bigdot}) \ar[r] & T'^{n+1}(K'^{\bigdot})
}
$$
中作通常的图追踪, 就得到\pageoriginale \hyperref[XII.cor6.3]{6.3} 的最后一个断言; 因为由已证的第一部分, 前两个竖直箭头为同构, 后两个为单态射. 至此完成了 \hyperref[XII.cor6.3]{6.3} 的证明.

\begin{enonce*}{引理 6.4}\etiquette[6.4]{XII.lem6.4}
若 $f$ 有限, 则定理 \hyperref[XII.thm5.1]{5.1} 成立.
\end{enonce*}

事实上, 这立即由 VIII \hyperref[VIII.prop5.5]{5.5} 及 \hyperref[VIII.rem5.8]{5.8} 得到.

\begin{enonce*}{命题 6.5}\etiquette[6.5]{XII.prop6.5}
设 $h:Y\to X$ 为拟紧且拟分离的概形态射. 在下面的 {\rm (i)} 中, $I$ 表示一个至少有两个元素的固定集合; 在 {\rm (ii)} 与 {\rm (iii)} 中, $\LL$ 表示素数集合 $\PP$ 的一个非空子集; 在 {\rm (iii)} 中, $n$ 表示整数.
\begin{enumerate}
\renewcommand{\theenumi}{\roman{enumi}}
\renewcommand{\labelenumi}{(\theenumi)}
\item 下述条件等价:
\begin{itemize}
\item[a)] 对 $X$ 上的每个集合层 $F$, 典范映射
$$
H^{\circ}(X,F)\longrightarrow H^{\circ}(Y,h^{\ast}(F))
$$
为单射 (分别地, 双射).

\item[b)] 对每个有限态射 $X'\to X$, 设 $h':Y'\to X'$ 为 $h$ 经基变换 $X'\to X$ 所得的态射. 则由 $h'$ 诱导的典范映射
$$
H^{\circ}(X', I_{X'})\longrightarrow H^{\circ}(Y',I_{Y'})
$$
为单射 (分别地, 双射).

\item[${\rm b')}$] 采用 {\rm (b)} 的记号, 从 $X'$ 的既开又闭的子集集合到 $Y'$ 的既开又闭的子集集合的映射 $U\mapsto {h'}^{-1}(U)$ 为单射 (分别地, 双射).

\item[${\rm b'')}$] (若\pageoriginale $X$ 局部 Noether.) 采用 {\rm (b)} 的记号, 所诱导的连通分支集合间的映射
$$
\Pi_{\circ}(h'):\Pi_{\circ}(Y')\longrightarrow \Pi_{\circ}(X')
$$
为满射 (分别地, 双射). 等价地, 若 $X'$ 非空 (分别地, 连通且非空), 则 $Y'$ 也如此.
\end{itemize}

此外, 若这些条件成立, 则对 $X$ 上的每个群层 $F$, 从 $F$ 下的挠子 (torsor, 齐性主丛) 范畴到 $h^{\ast}(F)$ 下的挠子范畴的函子 $P\mapsto h^{\ast}(P)$ 是忠实的 (分别地, 完全忠实的). 特别地, 在后一情形下, 典范映射
$$
H^1(X,F)\longrightarrow H^1(Y,h^{\ast}(F))
$$
为单射. 最后, 在同样的假设下, 从 $X$ 的平展覆盖范畴到 $Y$ 的平展覆盖范畴的函子 $X'\mapsto Y'=\fprod{X'}{Y}{X}$ 是忠实的 (分别地, 完全忠实的).

\item 下述条件等价:
\begin{itemize}
\item[a)] 对 $X$ 上的每个 ind-$\LL$ 有限群层 $F$, 典范映射
$$
H^i(X,F)\longrightarrow H^i(Y,h^{\ast}(F))
$$
对 $i=0,1$ 为双射.

\item[b)] 对每个有限态射 $X'\to X$, 设 $h':Y'\to X'$ 为 $h$ 经基变换所得的态射, $G$ 为任意通常的有限 $\LL$ 群. 则典范映射
$$
H^i(X',G_{X'})\longrightarrow H^i(Y',G_{Y'})
$$
对 $i=0,1$ 为双射.

\item[${\rm b')}$] (若\pageoriginale $\LL=\PP$.) 采用 {\rm (b)} 的记号, $h'$ 下的逆像函子诱导从 $X'$ 的平展覆盖范畴到 $Y'$ 的平展覆盖范畴的等价.

\item[${\rm b'')}$] (若 $X$ 是 Noether 概形且 $\LL=\PP$.) 采用 {\rm (b)} 的记号, 若 $X'$ 非空, 则 $Y'$ 非空. 若 $y'$ 为 $Y'$ 的几何点, $x'$ 为其在 $X'$ 中的像, 则典范映射
$$
\Pi_i(Y',y')\longrightarrow \Pi_i(X',x')\quav\mbox{ 对于 }i=0,1{\quoi},
$$
为双射.
\end{itemize}

\item 下述条件等价:
\begin{itemize}
\item[a)] 对 $X$ 上的每个 $\LL$ 挠群层 $F$, 同态
$$
H^{i}(X,F)\longrightarrow H^{i}(Y,h^{\ast}(F))
$$
对 $i\leq n$ 为同构, 对 $i=n+1$ 为单态射.

\item[${\rm a')}$] (若 $n\geq -1$.) 对每个上述 $F$, 有关同态对 $i=0$ 为单射, 对 $i\leq n$ 为满射.

\item[b)] (若 $n\geq -1$.) 对每个有限态射 $X'\to X$, 设 $h':Y'\to X'$ 为 $h$ 经基变换 $X'\to X$ 所得的态射. 则对每个 $\ell\in\LL$ 及 $\nu>0$, 典范同态
$$
H^{i}(X',(\ZZ/\ell^{\nu}\ZZ)_{X'})\longrightarrow
H^{i}(Y',(\ZZ/\ell^{\nu}\ZZ)_{Y'})
$$
对 $i=0$ 为单射, 对 $i\leq n$ 为满射.
\end{itemize}
\end{enumerate}
\end{enonce*}

此外, 对前述陈述还有如下补充:

\begin{enonce*}{推论 6.6}\etiquette[6.6]{XII.cor6.6}
\begin{enumerate}
\renewcommand{\theenumi}{\roman{enumi}}
\renewcommand{\labelenumi}{(\theenumi)}
\item 假设 \hyperref[XII.prop6.5]{6.5} {\rm (i) (a)} 中括号外的条件成立. 设 $F$ 为 $X$ 上的集合层, 且 $\xi\in H^{\circ}(Y,h^{\ast}(F))$. 假设\pageoriginale 存在集合层的单态射 $F\to G$, 使 $\xi$ 在 $H^{\circ}(Y,h^{\ast}(G))$ 中的像属于 $H^{\circ}(X,G)$ 的像. 则 $\xi$ 属于 $H^{\circ}(X,F)$ 的像.

\item 假设 \hyperref[XII.prop6.5]{6.5} {\rm (i) (a)} 中括号内的条件成立. 设 $F$ 为 $X$ 上的群层, 且 $\xi\in H^1(Y,h^{\ast}(F))$. 假设存在群层的单态射 $F\to G$, 使 $\xi$ 在 $H^1(Y,h^{\ast}(G))$ 中的像属于 $H^1(X,G)$ 的像. 则 $\xi$ 属于 $H^1(X,F)$ 的像.

\item 假设对某个固定的 $n$, 条件 \hyperref[XII.prop6.5]{6.5} {\rm (iii) (a)} 成立. 设 $F$ 为 $X$ 上的 $\LL$ 挠群层, 且 $\xi\in H^{n+1}(Y,h^{\ast}(F))$. 假设存在 $\LL$ 挠群层的单态射 $F\to G$, 使 $\xi$ 在 $H^{n+1}(Y,h^{\ast}(G))$ 中的像属于 $H^{n+1}(X,G)$ 的像. 则 $\xi$ 属于 $H^{n+1}(X,F)$ 的像.
\end{enumerate}
\end{enonce*}

\noindent
\emph{\hyperref[XII.prop6.5]{6.5} 及 \hyperref[XII.cor6.6]{6.6} 的证明.}
先证明 \hyperref[XII.cor6.6]{6.6}.

\hyperref[XII.cor6.6]{6.6} (iii) 的情形是 \hyperref[XII.cor6.3]{6.3} 最后一个断言的特殊情形: 取 $C$ 为 $X$ 上的阿贝尔 $\LL$ 挠层范畴, $C'$ 为范畴 $(\Ab)$, 并取 $T^{i}(F)=H^{i}(X,F)$、${T'}^{i}(F)=H^{i}(Y,h^{\ast}(F))$ (注意 $T'^{\bigdot}$ 确实是 $F$ 的上同调函子, 因为函子 $h^{\ast}$ 正合). 当然, 这里不必像 \hyperref[XII.cor6.3]{6.3} 中那样使用复形 (在后面的 \hyperref[XII.prop6.8]{6.8} 中这样做是方便的); 这里只用 \hyperref[XII.cor6.3]{6.3} 的最后一个断言.

为证明 (i), 考虑合并和 $H=G\amalg_F G$, 即图表
$$
\xymatrix@R=.5cm{
  & G\\
F \ar[ru] \ar[rd]&\\
  & G
}
$$
在 $X$ 上的集合层范畴中的余极限.\pageoriginale 由于 $F\rightarrow G$ 为单态射, 集合层拓扑斯通常的正合性性质表明, 有 $X$ 上的集合层的\emph{正合}图表
$$
F\to G\rightrightarrows H
$$
(这可表述为拓扑斯中每个单态射都是有效的, 并可照常归约到集合范畴予以证明). 由于函子 $H^{0}$ 与 $h^{\ast}$ 左正合, 得到集合的正合图表之间的态射
$$
\xymatrix@R=1.25cm{
H^{0}(X,F) \ar[r]\ar[d]& H^{0}(X,G)  \ar[d]\ar@<2pt>[r]\ar@<-2pt>[r]& H^{0}(X,H)\ar[d]\\
H^{0}(Y,h^{\ast}(F)) \ar[r]& H^{0}(Y,h^{\ast}(G)) \ar@<2pt>[r]\ar@<-2pt>[r]& H^{0}(Y,h^{\ast}(H)),
}
$$
其中竖直箭头均为单射. 直接作图追踪即证明断言 \hyperref[XII.cor6.6]{6.6} (i).

为证明 \hyperref[XII.cor6.6]{6.6} (ii), 先证明 \hyperref[XII.prop6.5]{6.5} (i) 的第二个断言, 即在 \hyperref[XII.prop6.5]{6.5} (i) 的条件 a) 下, 对 $X$ 上的每个群层 $F$, 从 $F$ 下的挠子范畴到 $h^{\ast}(F)$ 下的挠子范畴的函子 $P\mapsto h^{\ast}(P)$ 忠实 (\resp 完全忠实), 因而在相应情形下在 $H^1$ 上诱导单射. 为此, 若 $P$ 与 $P'$ 为 $F$ 下的两个挠子, 以 $\sheaf{\Isom}_F(P,P')$ 表示从 $P$ 到 $P'$ 的 $F$ 同构层. 立即可知, 该层的形成与任意基扩张可交换, 特别与 $h$ 下的基扩张可交换. (此断言对任意景态射都成立.) 将关于 $h^{\ast}$ 对 $H^{0}$ 的作用的假设应用于层 $\sheaf{\Isom}_F(P,P')$, 得到映射
$$
\Isom_F(P,P')\longrightarrow
\Isom_F(h^{\ast}(P),h^{\ast}(P'))
$$
为单射 (\resp 双射), 即函子 $P\mapsto h^{\ast}(P)$ \pageoriginale 忠实 (\resp 完全忠实). (注意群层下的挠子范畴中的每个态射都是同构.)

另一方面, 注意若 $F$ 是群层 $G$ 的子群层, 则给定 $F$ 下的挠子 $P$, 等价于给定 $G$ 下的挠子 $Q$ (即由 $P$ 沿 $F\rightarrow G$ 扩张结构群所得的挠子), \emph{连同} $Q/F$ 的一个截面 (这确实解释为将结构群从 $G$ 约化到 $F$). 采用 \hyperref[XII.cor6.6]{6.6} (ii) 的记号, $\xi$ 在 $H^1(Y,h^{\ast}(G))$ 中的像来自 $H^1(X,G)$ 的一个元素; 后者由 $G$ 下的挠子 $Q$ 表示, 所以 $\xi$ 的像是 $h^{\ast}(G)$ 下的挠子 $h^{\ast}(Q)$ 的类. 此类来自 $\xi\in H^1(Y,h^{\ast}(F))$, 明确地说就是 $\xi$ 是由 $h^{\ast}(Q)$ 及 $h^{\ast}(Q)/h^{\ast}(F)$ 的一个适当截面所定义的 $h^{\ast}(F)$ 挠子的类. 由于函子 $h^{\ast}$ 右正合, 最后这个层正是 $h^{\ast}(Q/F)$; 又因 $h^{\ast}$ 在集合层的 $H^{0}$ 上诱导双射, 所以 $h^{\ast}(Q)/h^{\ast}(F)$ 的所述截面来自 $Q/F$ 的一个截面. 这个截面通过结构群约化定义一个 $F$ 挠子, 该挠子在 $H^1(X,F)$ 中的类在 $H^1(Y,h^{\ast}(F))$ 中的像显然为 $\xi$. 这就证明了 \hyperref[XII.cor6.6]{6.6} (ii), 并完成了 \hyperref[XII.cor6.6]{6.6} 的证明.

现在证明 \hyperref[XII.prop6.5]{6.5}, 先从 \hyperref[XII.prop6.5]{6.5} (iii) 开始. 显然 a) 蕴含 ${\rm a'}$), 而反向蕴含是 \hyperref[XII.lem6.2]{6.2} 的特殊情形 (仍取 $C$ 为 $X$ 上的 $\LL$ 挠层范畴, $C'$ 为阿贝尔群范畴), 只须证明 $i>0$ 时函子 $H^{i}(X,F)$ 在 $C$ 中可消去. 这没有困难; 或者也可以\pageoriginale 不证明该结果, 而注意到 $F$ 是其子层 $F_i$ 的滤过余极限, 其中存在整数 $n>0$ 零化 $F_i$ (\hyperrefext[IX][IX.def1.1]{1.1}). 于是 $h^{\ast}(F)$ 是 $h^{\ast}(F_i)$ 的余极限; 考虑到上同调与滤过余极限可交换 (VII \hyperref[VII.prop3.3]{3.3}), 这一点来自 $X$ 与 $Y$ 的拟紧和拟分离假设 (这些假设在 \hyperref[XII.cor6.6]{6.6} 中并未使用), 只须对 $F_i$ 证明条件 a). 因而归结为存在整数 $m>0$ 使 $mF=0$ 的情形. 显然可以假设 $m$ 的每个素因子都属于 $\LL$, 于是可取 $C$ 为 $X$ 上的 $\ZZ/m\ZZ$ 模层范畴. 在此范畴中, 已知函子 $H^{i}(X,-)$ 对 $i\geq1$ 可消去. 这就证明了 \hyperref[XII.prop6.5]{6.5} (iii) 中条件 a) 与 ${\rm a'}$) 的等价性.

另一方面, 将 a) 应用于在 $X$ 上有限的概形 $X'$ 上的 $\LL$ 挠层 $F'$ 的正像; 根据 \hyperrefext[IX][IX.prop1.2]{1.2} (v) 及 \hyperrefext[IX][IX.prop1.6]{1.6} (iii), 此正像是 $\LL$ 挠层. 利用 \hyperref[XII.lem6.4]{6.4} 及 VIII \hyperref[VIII.prop5.5]{5.5}, 得到同态
$$
H^{i}(X',F')\longrightarrow H^{i}(Y',{h'}^{\ast}(F'))
$$
对 $i\leq n$ 为同构, 对 $i=n+1$ 为单态射, 因而更有条件 b). 只须证明此条件蕴含 ${\rm a'}$). 为此, 利用如下事实 (\hyperrefext[IX][IX.cor2.7.2]{2.7.2}): $X$ 上的每个 $\LL$ 挠层 $F$ 都是可构造 $\LL$ 挠层的滤过余极限. 上述余极限论证表明, 在 ${\rm a'}$) 中只须考虑可构造的 $F$. 根据 \hyperrefext[IX][IX.prop2.14]{2.14}, 存在有限多个有限态射 $p_i:X'_i\rightarrow X$、$X'_i$ 上的常 $\LL$ 挠层 $G_{X'_i}$, 以及单同态
$$
F\longrightarrow \prod_{i}p_{i\ast}(G_{X'_i}).
$$
现在\pageoriginale 若 $n=-1$, 即所求结论归结为
$$
H^{0}(X,F)\longrightarrow H^{0}(Y,h^{\ast}(F))
$$
为单射, 则在此问题中, 可以直接将 $F$ 换成一个包含 $F$ 的层 $G$, 因而可将 $F$ 换成各 $p_{i\ast}(G_{X'_i})$ 的积. 由 \hyperref[XII.lem6.4]{6.4}, 只须在将 $h:Y\rightarrow X$ 换成 $h':Y'\rightarrow X'$, 将 $F$ 换成 $X'$ 上以通常的有限 $\LL$ 挠群 $G$ 为值的常层之后, 证明单射性的断言. 但 $G$ 同构于形如 $\ZZ/\ell^{\nu}\ZZ$ 且 $\ell\in\LL$ 的群之直和. 因而确实归结为验证 b).

当 $n\geq0$ 时同样处理. 对 $n$ 作归纳, 可以假设已将 $n$ 换成 $n-1$ 证明了所求结论 (因而也证明了 a)). 应用 \hyperref[XII.cor6.6]{6.6} (iii), 可见仍可将 $F$ 换成一个包含 $F$ 的层 $G$, 且仍取 $G=p_{\ast}(G_{X'})$. 由 \hyperref[XII.lem6.4]{6.4}, 以及当 $p$ 有限时 $R^{i}p_{\ast}$ 对 $i\geq1$ 为零 (VIII \hyperref[VIII.prop5.5]{5.5}), 仍只须在将 $h:Y\rightarrow X$ 换成 $h':Y'\rightarrow X'$, 将 $F$ 换成 $X'$ 上以通常的有限 $\LL$ 挠群为值的常层之后, 证明 $H^n$ 的满射性. 根据归纳假设及蕴含 ${\rm a')}\Rightarrow{\rm a)}$, 已知单射性; 因而对 $G$ 作逐层剥离, 再次归结为 $G$ 形如 $\ZZ/\ell^{\nu}\ZZ$ 的情形, 即 \hyperref[XII.prop6.5]{6.5} (iii) 的条件 b) 所考虑的情形. 这就证明了 \hyperref[XII.prop6.5]{6.5} 的情形 (iii).

情形 (i) 与 (ii) 中, 关于 \hyperref[XII.prop6.5]{6.5} (i) \resp \hyperref[XII.prop6.5]{6.5} (ii) 中条件 a) 与 b) 的等价性, 证明完全相同: 利用\pageoriginale \hyperrefext[IX][IX.cor2.7.2]{2.7.2} 以及 $H^{0}$ 与 $H^1$ 与层的滤过余极限可交换 (VII \hyperref[VII.rem5.14]{5.14}), 归结为 $F$ 可构造的情形, 再利用 \hyperrefext[IX][IX.prop2.14]{2.14}、\hyperref[XII.cor6.6]{6.6} (i)、(ii) 及 VIII \hyperref[VIII.rem5.8]{5.8}, 归结为 b) 中所考虑的特殊情形. 此外, 在 (i) 中, 条件 b)、${\rm b'}$) 与 ${\rm b''}$) 的等价性是平凡的, 仅为便于引用而记录; 它们蕴含从 $X$ 的平展覆盖范畴到 $Y$ 的平展覆盖范畴的函子 $X'\mapsto Y'$ 忠实 (\resp 完全忠实), 这一事实也是如此. 在 \hyperref[XII.prop6.5]{6.5} (ii) 中补充给出的条件 b)、${\rm b'}$) 与 ${\rm b''}$) 的等价性也同样如此. 至此完成了 \hyperref[XII.prop6.5]{6.5} 的证明.

\begin{enonce*}[remark]{注 6.7}\etiquette[6.7]{XII.rem6.7}
\footnote{编者建议略去这些注以及 \hyperref[XII.rems6.13]{6.13}; 它们是由一位不敬的合作者偷偷加进来的 (他还对原文作了各种其他或好或坏的改动).}
\begin{enumerate}
\renewcommand{\labelenumi}{\theenumi)}
\item 证明 \hyperref[XII.prop6.5]{6.5} 与 \hyperref[XII.cor6.6]{6.6} 的论证显然十分一般, 实质上也很平凡; 若能将其抽出为与 \hyperref[XII.lem6.2]{6.2} 同类的抽象引理, 将是有益的 (其中各层 $p_{\ast}(G_{X'})$ 起所研究的范畴 $C$ 的余生成元的作用). 将这一愉快的练习留给读者, 只指出, 以 $R^{i}f_{\ast}$ 之类的函子代替函子 $H^{i}$, 仍可给出实质相同的陈述. 同样的评注也适用于下面的 \hyperref[XII.prop6.8]{6.8} 及 \hyperref[XII.prop6.11]{6.11}.


\item 当 $X$ 为 Noether 概形时, 上述论证表明, 在 \hyperref[XII.prop6.5]{6.5} (i)、(ii)、(iii) 的条件 b) 及其变体中, 可以只考虑\emph{整的} $X'$; 当 $X$ 泛日本时, 甚至可以取它们整且正规. 即使对 $X$ 不作日本性假设, 也可以如此,\pageoriginale 只须用\emph{整}态射 $p:X'\rightarrow X$ 代替有限态射. 本讨论班后面不会使用这些变体.

\item \hyperref[XII.cor6.6]{6.6} 的证明没有使用关于 $X$ 的拟紧和拟分离假设. 在 (i)、(ii)、(iii) 中, 条件 a) 蕴含其他条件, 以及条件 b)、${\rm b'}$)、${\rm b''}$) 彼此等价, 这两点也是如此. 还要注意, \hyperref[XII.prop6.5]{6.5} (i) 中条件的 b) 形式表明, 对 $\LL\neq\phi$, 它们是 \hyperref[XII.prop6.5]{6.5} (iii) 中 $n=-1$ \resp $n=0$ 时的特殊情形 (当 $X$ 拟紧且拟分离时). 此外, 即使没有所述关于 $X$ 的限制, 这一点仍然成立 (只考虑 a) 形式的条件). 事实上, 对每个集合层 $F$ 及每个素数 $\ell$, 可以找到两个阿贝尔 $\ell$ 挠层 $G$ 与 $H$, 以及底集合层的同态 $u:G\rightarrow H$, 使 $F$ 同构于零截面的逆像 (取 $G$ 为由 $F$ 生成的自由 $\ZZ/\ell\ZZ$ 模, 取 $H$ 为 $X$ 上的常层 $\ZZ/\ell\ZZ$). 同样的论证表明, \hyperref[XII.prop6.5]{6.5} (ii) a) (其中甚至只须取 $i=0$) 蕴含相应条件 \hyperref[XII.prop6.5]{6.5} (i) a), 而无需任何拟紧或拟分离限制. 由于后者已经蕴含对 $X$ 上的每个群层 $F$, $H^1(X,F)\rightarrow H^1(Y,h^{\ast}(F))$ 为单射, 因此在 \hyperref[XII.prop6.5]{6.5} (ii) a) 的陈述中, 对 $i=1$ 只须要求 $H^1(X,F)\rightarrow H^1(Y,h^{\ast}(F))$ 的\emph{满射性}. 遗憾的是, 不能将 \hyperref[XII.prop6.5]{6.5} (ii) 看成 \hyperref[XII.prop6.5]{6.5} (iii) 的特殊情形; 这使我们常常不得不在非交换情形中重复一个实质上已在交换情形中完成的论证. 还要注意, 除 \hyperref[XII.cor6.10]{6.10} 外, 后面的结果都表述为无需\pageoriginale 任何拟紧或拟分离假设仍然成立的形式.
\end{enumerate}
\end{enonce*}

\begin{enonce*}{命题 6.8}\etiquette[6.8]{XII.prop6.8}
(下降引理). 采用 \hyperref[XII.prop6.5]{6.5} 的一般记号, 再假设给定态射 $p:\bar{X}\rightarrow X$. 以 $\bar{h}:\bar{Y}\rightarrow\bar{X}$ 表示 $h$ 经基变换 $p$ 所得的态射, 以 $p_Y:\bar{Y}\rightarrow Y$ 表示典范态射. 假设 $p$ 为满射, 或者 $h$ 为闭浸入且 $X=p(\bar{X})\cup h(Y)$.
\begin{enumerate}
\renewcommand{\theenumi}{\roman{enumi}}
\renewcommand{\labelenumi}{(\theenumi)}
\item 假设 $\bar{h}:\bar{Y}\rightarrow\bar{X}$ 满足 \hyperref[XII.prop6.5]{6.5} {\rm(i) a)} 的单射性 (\resp 双射性) 条件. 在相应情形下, 再假设对 $\bar{X}$ 上的每个集合层 $\bar{F}$, 基变换态射
\begin{equation*}
h^{\ast}(p_{\ast}(\bar{F}))\longrightarrow
p_{Y\ast}(\bar{h}^{\ast}(\bar{F}))\tag{$\ast$}
\end{equation*}
为单射. 则态射 $h:Y\rightarrow X$ 也满足 \hyperref[XII.prop6.5]{6.5} {\rm(i) a)} 的单射性 (\resp 双射性) 条件.

\item 假设 $\bar{h}:\bar{Y}\rightarrow\bar{X}$ 满足条件 \hyperref[XII.prop6.5]{6.5} {\rm(ii) a)}, 且对 $\bar{X}$ 上的每个 ind-$\LL$ 群层 $\bar{F}$, 基变换态射 $(\ast)$ 为双射. 则 $h:Y\rightarrow X$ 也满足条件 \hyperref[XII.prop6.5]{6.5} {\rm(ii) a)}.

\item 假设 $\bar{h}:\bar{Y}\rightarrow\bar{X}$ 满足条件 \hyperref[XII.prop6.5]{6.5} {\rm(iii) a)}, 且对 $\bar{X}$ 上的每个阿贝尔 $\LL$ 挠层 $\bar{F}$, 基变换同态
$$
h^{\ast}(R^{i}p_{\ast}(\bar{F}))\longrightarrow
R^{i}p_{Y\ast}(\bar{h}^{\ast}(\bar{F}))
$$
对 $i\leq n-1$ 为双射, 对 $i=n$ 为单射, 其中 $n\geq1$ 为给定整数. 则 $h:Y\rightarrow X$ 也满足条件 \hyperref[XII.prop6.5]{6.5} {\rm(iii) a)}.
\end{enumerate}
\end{enonce*}

\noindent
\emph{\hyperref[XII.prop6.8]{6.8} 的证明.}\pageoriginale 在不假设 $p$ 为满射, 而假设 $h$ 为闭浸入且 $X=p(\bar{X})\cup h(Y)$ 的情形下, 考虑 $\bar{X}$ 与 $Y$ 的不交并概形 $\bar{X}'$, 以及由 $p$ 和 $h$ 诱导的态射 $p':\bar{X}'\rightarrow X$. 则 $p'$ 为满射; 此外, 在 (i)、(ii)、(iii) 三种情形中的每一种, 立即可知, 关于对 $(h,p)$ 的假设对于对 $(h,p')$ 仍然满足. 因而归结为 $p$ 为满射的情形, 以下均如此假设. 在这种情形下, 对 $X$ 上的每个层 $F$, 典范同态
$$
F\longrightarrow p_{\ast}(p^{\ast}(F))
$$
为单射. 在情形 (i) 中, 为验证 $H^{0}(X,F)\rightarrow H^{0}(Y,h^{\ast}(F))$ 的单射性, 立即归结为将 $F$ 换成 $p_{\ast}(\bar{F})$ 的情形, 其中 $\bar{F}=p^{\ast}(F)$. 考虑交换图表
$$
\xymatrix@C=2.75cm@R=.75cm{
**[r]H^{0}(X,p_{\ast}(\bar{F})) \ar[dd]^{\text{\rotatebox{90}{$\backsim$}}} \ar[r]& **[r]H^{0}(Y,h^{\ast}(p_{\ast}(\bar{F}))) \ar[d]\\
                           &**[r]H^{0}(Y,p_{Y\ast}(\bar{h}^{\ast}(\bar{F})))\ar[d]_{\text{\rotatebox{90}{$\backsim$}}}\\
**[r]H^{0}(\bar{X},\bar{F})  \ar[r] &
                           **[r]H^{0}(\bar{Y},\bar{h}^{\ast}(\bar{F})),
}
$$
其中左边的竖直箭头及右边的第二个竖直箭头为典范同构, 右边的第一个竖直箭头由 $h$ 的基变换同态诱导. 根据假设, 第二个水平箭头为单射, 立即可知第一个也是单射; 这证明了 (i) 中不带“分别地”的断言. 对于带“分别地”的断言, 假设第二个水平箭头为满射. 由于根据假设基变换同态为单态射, 右边的第一个竖直箭头为单态射, 从而右边两个\pageoriginale 竖直箭头的复合也为单态射. 立即可知第一个水平箭头也为满射, 至此完成了 (i) 的证明.

为证明 (ii), 由前面的论证及 \hyperref[XII.prop6.5]{6.5} (i), 只须证明 $H^1(X,F)\rightarrow H^1(Y,h^{\ast}(F))$ 的满射性. 应用 \hyperref[XII.cor6.6]{6.6} (ii), 又只须对形如 $p_{\ast}(\bar{F})$ 的 $F$ 证明它, 其中 $\bar{F}$ 为 $\bar{X}$ 上的 ind-$\LL$ 群层 (\hyperrefext[IX][IX.prop1.6]{1.6} (ii)). 对态射 $p$ 与 $p_Y$ 写出 $H^1$ 的正合列 (\hyperref[XII.prop3.2]{3.2} 的变体), 得到非交换正合列的态射:
$$
\xymatrix@R=1.5cm{
H^{1}(X,p_{\ast}(\bar{F})) \ar[r] \ar[d] & H^{1}(\bar{X},\bar{F}) \ar[r]\ar[d] &
H^{0}(X,R^{1}p_{\ast}(\bar{F})) \ar[d]&\\
H^{1}(Y,p_{Y\ast}(\bar{h}^{\ast}(\bar{F}))) \ar[r] &
H^{1}(\bar{Y},\bar{h}^{\ast}(\bar{F})) \ar[r] &
H^{0}(Y,R^{1}p_{Y\ast}(\bar{h}^{\ast}(\bar{F}))),
}
$$
其中每一行的第一个箭头为单射, 并将第一项等同于第三项指定点在第二个箭头下的逆像. 根据关于 $\bar{h}$ 的假设, 中间的竖直箭头为双射. 两侧的竖直箭头分别经由
$$
H^1(Y,h^{\ast}(p_{\ast}(\bar{F})))
\quad\hbox{及}\quad
H^{0}(Y,h^{\ast}(R^1p_{\ast}(\bar{F}))),
$$
分解, 方法是将同态 $H^{i}(X,-)\rightarrow H^{i}(Y,h^{\ast}(-))$ (分别对于 $i=1$ 及 $i=0$) 与对基变换同态应用 $H^{i}(Y,-)$ (对同样的 $i$ 值) 所得的同态复合; 这些基变换同态分别为
$$
h^{\ast}(p_{\ast}(\bar{F}))\longrightarrow
p_{Y\ast}(\bar{h}^{\ast}(\bar{F}))
$$
以及
$$
h^{\ast}(R^1p_{\ast}(\bar{F}))\longrightarrow
R^1p_{Y\ast}(\bar{h}^{\ast}(\bar{F})).
$$
由于根据假设第一个箭头为同构, 上图中第一个竖直箭头可等同于我们要研究的同态 $H^1(X,-)\rightarrow H^1(Y,h^{\ast}(-))$. 因而为证明它是双射, 只须证明图表中最后一个竖直箭头为单射. 根据刚才的说明, 它是两个映射的复合,\pageoriginale 第一个为典范同态 $H^{0}(X,-)\rightarrow H^{0}(Y,h^{\ast}(-))$, 从而为双射; 第二个为单射, 因为它由将左正合函子 $H^{0}(Y,-)$ 应用于单态射
$$
h^{\ast}(R^1p_{\ast}(\bar{F}))\longrightarrow
R^1p_{Y\ast}(\bar{h}^{\ast}(\bar{F})).
$$
而得. 最后这个事实容易从前一个次数 $0$ 的基变换同态为双射这一假设得到, 例如利用 \hyperref[XII.prop6.5]{6.5} (i) 的第二个断言及高阶正像的茎的通常计算 (VIII \hyperref[VIII.rem5.3]{5.3}). 这就证明了 \hyperref[XII.prop6.8]{6.8} (ii).

只须证明 \hyperref[XII.prop6.8]{6.8} (iii). 为此, 与其如上直接使用单射 $F\rightarrow p_{\ast}(p^{\ast}(F))$, 不如使用右导出范畴 $D^{+}(C)$ 中的同态
$$
F\longrightarrow Rp_{\ast}(\bar{F})
$$
这里 $C$ 表示 $X$ 上的阿贝尔层范畴. 回顾一下, 根据定义, $Rp_{\ast}(\bar{F})$ 为复形 $p_{\ast}(C(\bar{F}))$, 其中 $C(\bar{F})$ 为 $\bar{X}$ 上的阿贝尔层范畴中 $\bar{F}$ 的内射消解. 若 $K^{\bigdot}$ 为 $X$ 上的复形, 以 $H^{i}(X,K^{\bigdot})$ 表示 $X$ 以 $K^{\bigdot}$ 为系数的超上同调群, 并在 $Y$ 上使用类似记号. 由 \hyperref[XII.cor6.3]{6.3}, 只须证明如下同态的满射性:
\begin{equation*}
H^n(X,Rp_{\ast}(\bar{F}))\longrightarrow
H^n(Y,h^{\ast}(Rp_{\ast}(\bar{F}))).\tag{1}
\end{equation*}\etiquette[1]{XII.eq1}
根据假设, 基变换同态
\begin{equation*}
h^{\ast}(Rp_{\ast}(\bar{F}))\longrightarrow
Rp_{Y\ast}(\bar{h}^{\ast}(\bar{F}))\tag{2}
\end{equation*}\etiquette[2]{XII.eq2}
在次数 $\leq n-1$ 的上同调层上诱导同构, 在次数 $n$ 诱导单态射. 等价地, 上述同态的“映射柱”\pageoriginale $K^{\bigdot}$ 只在次数 $\geq n$ 有非零上同调层. 因而对 $i<n$ 有 $H^{i}(Y,K^{\bigdot})=0$, 而上同调正合列表明, 上述同态所诱导的同态
\begin{equation*}
H^{i}(Y,h^{\ast}(Rp_{\ast}(\bar{F})))\longrightarrow
H^{i}(Y,Rp_{Y\ast}(\bar{h}^{\ast}(\bar{F})))\tag{3}
\end{equation*}\etiquette[3]{XII.eq3}
对 $i\leq n-1$ 为同构, 对 $i=n$ 为单态射. 考虑 (1) 与 (3) 的复合:
\begin{equation*}
H^n(X,Rp_{\ast}(\bar{F}))\longrightarrow
H^n(Y,h^{\ast}(Rp_{\ast}(\bar{F})))\longrightarrow
H^n(Y,Rp_{Y\ast}(\bar{h}^{\ast}(\bar{F})))\tag{4}
\end{equation*}\etiquette[4]{XII.eq4}

根据定义, 两端的项分别同构于 $H^n(\bar{X},\bar{F})$ 与 $H^n(\bar{Y},\bar{h}^{\ast}(\bar{F}))$, 而 (4) 中的复合恰为 $\bar{h}$ 所诱导的同态, 由关于 $\bar{h}$ 的假设它为双射. 根据刚才的说明, (4) 中第二个箭头为单射, 因而第一个也为双射. 至此完成了 \hyperref[XII.prop6.8]{6.8} 的证明.

此外, 上面的证明立即给出下述略为精确和一般的结果:

\begin{enonce*}{推论 6.9}\etiquette[6.9]{XII.cor6.9}
记号同 \hyperref[XII.prop6.8]{6.8}. 再假设 $h$ 满足 \hyperref[XII.prop6.5]{6.5} \textup{(i)} 中括号外的条件 (分别地, \hyperref[XII.prop6.5]{6.5} \textup{(i)} 中括号内的条件; 分别地, \hyperref[XII.prop6.5]{6.5} \textup{(iii)} 中的条件, 其中 $n\geq -1$), 并且对 $\bar{X}$ 上的每个集合层 $\bar{F}$, 基变换同态
$$
h^{\ast}(p_{\ast}(\bar{F}))\longrightarrow
p_{Y\ast}(\bar{h}^{\ast}(\bar{F}))
$$
为单射 (分别地, 双射; 分别地, 对 $\bar{X}$ 上的每个 $\LL$ 挠群层 $\bar{F}$, 基变换同态
$$
h^{\ast}R^{i}p_{\ast}(\bar{F})\longrightarrow
R^{i}p_{Y\ast}(\bar{h}^{\ast}(\bar{F}))
$$
对 $i\leq n$ 为双射, 对 $i=n+1$ 为单射). 设 $F$ 为 $X$ 上的集合层 (分别地, ind-$\LL$ 群层; 分别地, $\LL$ 挠群层), 并设 $\xi\in H^{n+1}(Y,h^{\ast}(F))$ (其中在情形\pageoriginale \textup{(i)} 中取 $n=-1$, 在情形 \textup{(ii)} 中取 $n=0$). 置 $\bar{F}=p^{\ast}(F)$. 则 $\xi$ 属于 $H^{n+1}(X,F)$ 的像, 当且仅当它在
$$
H^{n+1}(\bar{Y},p_Y^{\ast}(h^{\ast}(F)))
=H^{n+1}(\bar{Y},\bar{h}^{\ast}(\bar{F}))
$$
中的逆像属于 $H^{n+1}(\bar{X},\bar{F})$ 的像.
\end{enonce*}

将利用 \hyperref[XII.prop6.8]{6.8} 的这一加强形式来证明下面的结果.

\begin{enonce*}{推论 6.10}\etiquette[6.10]{XII.cor6.10}
记号同 \hyperref[XII.prop6.5]{6.5}. 为简单起见, 假设 \hyperref[XII.prop6.5]{6.5} \textup{(i)} 中括号外的条件成立.
\begin{enumerate}
\renewcommand{\theenumi}{\roman{enumi}}
\renewcommand{\labelenumi}{(\theenumi)}
\item \hyperref[XII.prop6.5]{6.5} \textup{(i)} 中括号内的条件成立 (分别地, 条件 \hyperref[XII.prop6.5]{6.5} \textup{(ii) a)} 成立; 分别地, 条件 \hyperref[XII.prop6.5]{6.5} \textup{(iii) a)} 成立), 当且仅当下述条件满足. 对 $X$ 上的每个集合层 $F$ (分别地, ind-$\LL$ 群层; 分别地, 阿贝尔 $\LL$ 挠群层), 每个 $\xi\in H^{\circ}(Y,h^{\ast}(F))$ (分别地, 每个满足 $i=0,1$ 的 $\xi\in H^{i}(Y,h^{\ast}(F))$; 分别地, 每个满足 $i\leq n$ 的 $\xi\in H^{i}(Y,h^{\ast}(F))$), 以及 $X$ 的每个非空闭子集 $X'$, 若以 $Y'$ 表示它在 $Y$ 中的逆像, 则存在态射 $p:\bar{X}\rightarrow X$ 满足下述条件:
\begin{itemize}
\item[$1^{\circ}$)] 像 $p(\bar{X})$ 包含于 $X'$, 并包含 $X'$ 的一个非空开子集.

\item[$2^{\circ}$)] 对 $\bar{X}$ 上的每个集合层 $\bar{F}$, 基变换同态
$$
h^{\ast}(p_{\ast}(\bar{F}))\longrightarrow
p_{Y\ast}(\bar{h}^{\ast}(\bar{F}))
$$
为同构 (分别地, 对 $\bar{X}$ 上的每个 ind-$\LL$ 群层 $\bar{F}$, 基变换同态
$$
h^{\ast}(R^{i}p_{\ast}(\bar{F}))\longrightarrow
R^{i}p_{Y\ast}(\bar{h}^{\ast}(\bar{F}))
$$
对 $i=0,1$ 为双射; 分别地, 对 $\bar{X}$ 上的每个阿贝尔 $\LL$ 挠群层 $\bar{F}$, 基变换同态对 $i\leq n$ 为双射).

\item[$3^{\circ}$)] $\xi$ 在
$$
H^{i}(\bar{Y},p_Y^{\ast}(h^{\ast}(F)))
=H^{i}(\bar{Y},\bar{h}^{\ast}(p^{\ast}(F)))
$$
中的逆像属于 $H^{i}(\bar{X},p^{\ast}(F))$ 的像.
\end{itemize}

\item 设\pageoriginale $n\geq -1$, 并假设对 $X$ 上的每个集合层 $F$ (分别地, 每个 ind-$\LL$ 群层; 分别地, 每个阿贝尔 $\LL$ 挠群层), 同态
$$
H^{i}(X,F)\longrightarrow H^{i}(Y,h^{\ast}(F))
$$
对 $i\leq n$ 为双射, 对 $i=n+1$ 为单射 (在第一种情形下取 $n=-1$, 在第二种情形下取 $n=-1$ 或 $0$). 设 $\xi\in H^{n+1}(Y,h^{\ast}(F))$. 则 $\xi$ 属于 $H^{n+1}(X,F)$ 的像, 当且仅当它满足上述 \textup{(i)} 中的条件, 其中 $i=n+1$.
\end{enumerate}
\end{enonce*}

\noindent
\emph{\hyperref[XII.cor6.10]{6.10} 的证明.} 利用 \hyperref[XII.prop6.5]{6.5} 并对 $n$ 作直接的归纳, 可见 \hyperref[XII.cor6.10]{6.10} (i) 可由 \hyperref[XII.cor6.10]{6.10} (ii) 得到; 现在证明后者.

假设 $\xi$ 不属于 $H^{n+1}(X,F)$ 的像. 设 $\Phi$ 为 $X$ 的如下闭子集 $X'$ 所成的集合: $\xi$ 在 $H^{n+1}(Y',h^{\ast}(F)|Y')$ 中的像不属于 $H^{n+1}(X',F|X')$ 的像. 根据假设, $X\in\Phi$, 所以 $\Phi$ 非空. 用关系 $\supset$ 对 $\Phi$ 排序, 并证明 $\Phi$ 是归纳的. 只须证明, 若 $(X'_{\lambda})$ 为 $\Phi$ 中的一个全序元素族, $X'$ 为其交, 则 $X'\in\Phi$. 给各 $X'_{\lambda}$ 及 $X'$ 赋予其诱导的既约结构. 态射 $X'_{\lambda}\rightarrow X$ 是仿射的 (因为是闭浸入), 因而各 $X'_{\lambda}$ 所成逆向系满足 VII 5 中考虑的假设. 此外, 立即可知 $X'$ 是各 $X'_{\lambda}$ 的逆极限. 同样, $Y'$ 是各 $Y'_{\lambda}$ 所成逆向系的逆极限, 后者在 $Y$ 上仿射. 因而由 VII \hyperref[VII.cor5.8]{5.8},
$$
\begin{gathered}
\varinjlim H^{n+1}(X'_{\lambda},F|X'_{\lambda})\xrightarrow{\sim}
H^{n+1}(X',F|X')\quad,\\
\varinjlim H^{n+1}(Y'_{\lambda},h^{\ast}(F)|Y'_{\lambda})
\xrightarrow{\sim}H^{n+1}(Y',h^{\ast}(F)|Y'){\quoi}.
\end{gathered}
$$
立即可知, 若 $(\xi|Y')\in\Iim H^{n+1}(X',F|X')$, 则\pageoriginale 存在指标 $\lambda$, \linebreak 使 $(\xi|Y'_{\lambda})\in\allowbreak\Iim H^{n+1}(X'_{\lambda},F|X'_{\lambda})$, 矛盾. 因此确有 $X'\in\Phi$. 从而 $\Phi$ 是归纳的, 因而包含一个关于包含关系极小的元素, 记为 $X'$. 由于 $\varnothing\notin\Phi$, 集合 $X'$ 非空. 将 \hyperref[XII.cor6.10]{6.10} 的假设应用于 $\xi$ 与 $X'$, 得到满足 \hyperref[XII.cor6.10]{6.10} 中条件 $1^{\circ}$)--$3^{\circ}$) 的态射 $p:\bar{X}\rightarrow X$. 注意甚至可以假设 $p$ \emph{满射到 $X'$}. 事实上, 设 $U$ 为包含在 $p(\bar{X})$ 中的 $X'$ 的一个非空开子集, $X'_1$ 为其在 $X'$ 中的补集. 根据构造, $X'_1\notin\Phi$. 置 $\bar{X}_1=\bar{X}\amalg X'_1$, 并设 $p_1:\bar{X}_1\rightarrow X$ 为由 $p$ 及嵌入 $X'_1\hookrightarrow X$ 定义的态射. 利用关系 $X'_1\notin\Phi$, 立即可知 $p_1$ 仍满足与 $p$ 相同的假设; 而且它满射到 $X'$. 因而可以假设 $p$ 满射到 $X'$. 现在应用 \hyperref[XII.cor6.9]{6.9}, 将其中的 $h:Y\rightarrow X$ 换成 $h':Y'\rightarrow X'$ (例如给 $X'$ 赋予其诱导的既约结构), 并将 $\xi$ 换成 $\xi|Y'$. 由 \hyperref[XII.cor6.9]{6.9} 得到 $(\xi|Y')\in\Iim H^{n+1}(X',F|X')$, 与 $X'\in\Phi$ 矛盾. 至此完成了 \hyperref[XII.cor6.10]{6.10} 的证明.

将利用 \hyperref[XII.cor6.10]{6.10} 将 \hyperref[XII.cor5.5]{5.5} 归结为 $f$ 射影的情形. 但在非 Noether 情形下, 这需要 Chow 引理的非 Noether 变体, 将在下一节给出.

\begin{enonce*}{命题 6.11}\etiquette[6.11]{XII.prop6.11}
(传递性引理). 记号同 \hyperref[XII.prop6.8]{6.8}, 但对 $p$ 不加满射性假设. 假设 $h$ 满足 \hyperref[XII.prop6.5]{6.5} 的条件 \textup{(i) a)} (分别地, \textup{(ii) a)}; 分别地, \textup{(iii) a)}). 在情形 \textup{(i)} 中, 再假设对 $\bar{X}$ 上的每个集合层 $\bar{F}$, \hyperref[XII.prop6.8]{6.8} 中的基变换同态为单态射 (分别地, 同构); 在情形 \textup{(ii)} 中, 假设 $\bar{X}$ 是 Noether 概形, 并且对每个 ind-$\LL$ 群层 $\bar{F}$, 基变换同态
$$
h^{\ast}(R^ip_{\ast}(\bar{F}))\longrightarrow
R^ip_{Y\ast}(\bar{h}^{\ast}(\bar{F}))
$$
对 $i=0,1$ 为\pageoriginale 双射; 最后在情形 \textup{(iii)} 中, 假设基变换同态对 $i\leq n$ 为双射, 对 $i=n+1$ 为单射. 则 $\bar{h}:\bar{Y}\rightarrow\bar{X}$ 也满足 \hyperref[XII.prop6.5]{6.5} 的条件 \textup{(i)} (分别地, \textup{(ii)}; 分别地, \textup{(iii)}).
\end{enonce*}

\noindent
\emph{\hyperref[XII.prop6.11]{6.11} 的证明.} 在情形 (i) 中, 考虑 \hyperref[XII.prop6.8]{6.8} (i) 证明中的交换图表. 这里假设蕴含第一个水平箭头与右上方的竖直箭头为单射 (分别地, 双射), 因而第二个水平箭头也如此, 正如所求. 在情形 (ii) 中, 设 $\bar{F}$ 为 $\bar{X}$ 上的 ind-$\LL$ 群层. 须证明 $H^1(\bar{Y},\bar{h}^{\ast}(\bar{F}))$ 的每个元素都来自 $H^1(\bar{X},\bar{F})$ 的某个元素. 将 \hyperref[XII.cor6.6]{6.6} (ii) 应用于 $\bar{h}$, 只须找到 $\bar{X}$ 上的群层单态射 $\bar{F}\rightarrow\bar{G}$, 以\emph{消去} $H^1(\bar{Y},\bar{h}^{\ast}(\bar{F}))$. 考虑 \hyperref[XII.prop6.8]{6.8} (ii) 证明中所用图表的第二行. 只须依次找到单态射 $\bar{F}\rightarrow\bar{G}$, 其中 $\bar{G}$ 为 ind-$\LL$ 有限群层, 以消去最后一项 $H^{\circ}(Y,R^1p_{Y\ast}(\bar{h}^{\ast}(\bar{F})))$, 再找到群层单态射 $\bar{G}\rightarrow\bar{H}$, 以消去 $H^1(Y,p_{Y\ast}(\bar{h}^{\ast}(\bar{G})))$. 如 \hyperref[XII.prop6.8]{6.8} (i) 的证明中所指出, 根据这里对 $p$ 与 $h$ 的基变换同态的假设, 待消去的两项通过有关图表两端的竖直箭头分别同构于 $H^{\circ}(X,R^1p_{\ast}(\bar{F}))$ 与 $H^1(X,p_{\ast}(\bar{F}))$. 这些同构显然对 $\bar{F}$ 具有函子性, 所以只须消去后两项. 由 \hyperref[XII.prop3.3]{3.3}, 可以将 $\bar{F}$ 嵌入某个 ind-$\LL$ 群 $\GG$, 以消去 $R^1p_{\ast}(\bar{F})$ (从而更能消去 $H^{\circ}(X,R^1p_{\ast}(\bar{F}))$). 只须消去 $H^1(X,p_{\ast}(\bar{F}))$. 这是可能的, 因为\pageoriginale 根据有关图表, 后一项通过一个对 $\bar{F}$ 具有函子性的单态射映入 $H^1(\bar{X},\bar{F})$, 而 $H^1(\bar{X},\bar{F})$ 再次由 \hyperref[XII.prop3.3]{3.3} 可消去.

只须处理情形 (iii), 这里仍然倒读 \hyperref[XII.prop6.8]{6.8} 的证明\footnote[1]{也可不用 \hyperref[XII.cor6.3]{6.3}, 而使用 $p$ 与 $p_Y$ 的 Leray 谱序列.}. 须证明同态
$$
H^{i}(\bar{X},\bar{F})\longrightarrow
H^{i}(\bar{Y},\bar{h}^{\ast}(\bar{F}))
$$
对 $i\leq n$ 为同构, 对 $i=n+1$ 为单态射. 这些同态由复形同态
$$
\RR\Gamma_{\bar{X}}(\bar{F})\longrightarrow
\RR\Gamma_{\bar{Y}}(\bar{h}^{\ast}(\bar{F})){\quoi},
$$
诱导, 后者由一个内射消解 $C(\bar{F})$, 以及从 $\bar{h}^{\ast}(C(\bar{F}))$ (它是 $\bar{h}^{\ast}(\bar{F})$ 的消解) 到 $\bar{h}^{\ast}(\bar{F})$ 的一个内射消解的同态得到. 上述同态也等同于复合
$$
\RR\Gamma_X(Rp_{\ast}(\bar{F}))\longrightarrow
\RR\Gamma_Y(h^{\ast}(Rp_{\ast}(\bar{F})))\longrightarrow
\RR\Gamma_Y(Rp_{Y\ast}(\bar{h}^{\ast}(\bar{F}))){\quoi},
$$
其中第一个同态为 $h^{\ast}$ 对 $Rp_{\ast}(\bar{F})$ 所诱导的同态, 第二个为将 $\RR\Gamma_Y$ 应用于基变换同态
$$
h^{\ast}(Rp_{\ast}(\bar{F}))\longrightarrow
Rp_{Y\ast}(\bar{h}^{\ast}(\bar{F})).
$$
所得的同态. 因此, 根据关于 $h$ 的假设及 \hyperref[XII.cor6.3]{6.3} (i), 第一个同态在所需范围内于有关复形的上同调群上诱导同构. 根据关于基变换同态的假设, 使用 \hyperref[XII.prop6.8]{6.8} (iii) 证明中的论证, 第二个也如此. 因而它们的复合也有相同结论, 至此完成了 \hyperref[XII.prop6.11]{6.11} 的证明.

\begin{enonce*}[remark]{注 6.12}\etiquette[6.12]{XII.rem6.12}
对 \hyperref[XII.prop6.11]{6.11} 的情形 (i) 与 (iii) 给出的论证\pageoriginale 给出下述更精确的结果. 为对 $\bar{X}$ 上一个\emph{给定的}层 $\bar{F}$ 断言
$$
H^{i}(\bar{X},\bar{F})\longrightarrow
H^{i}(\bar{Y},\bar{h}^{\ast}(\bar{F}))
$$
对 $i\leq n$ 为双射, 对 $i=n+1$ 为单射, 只须假设
\begin{itemize}
\item[$1^{\circ}$)] 基变换同态
$$
h^{\ast}(R^{i}p_{\ast}(\bar{F}))\longrightarrow
R^{i}p_{Y\ast}(\bar{h}^{\ast}(\bar{F}))
$$
对 $i\leq n$ 为同构, 对 $i=n+1$ 为单态射, 且

\item[$2^{\circ}$)] 对每一对整数 $i,j\geq0$, 由 $h$ 诱导的同态
$$
H^{i}(X,R^{j}p_{\ast}(\bar{F}))\longrightarrow
H^{i}(Y,h^{\ast}(R^{j}p_{\ast}(\bar{F})))
$$
在 $i+j\leq n$ 时为同构, 在 $i+j=n+1$ 时为单态射.
\end{itemize}

情形 (ii) 中很可能也有类似结果, 但其陈述必然涉及非交换 $2$ 上同调. 为避免诉诸该理论, 在 \hyperref[XII.prop6.11]{6.11} (ii) 中采用了基于非交换 $H^1$ 的可消去性质的不同论证; 但它与上面的论证不同, 不能对固定的 $\bar{F}$ 给出结果.
\end{enonce*}


\begin{enonce*}[remark]{注 6.13}\etiquette[6.13]{XII.rems6.13}
\hyperref[XII.prop6.5]{6.5} (i)、(ii)、(iii) 的条件成立的最重要情形, 是 \hyperref[XII.cor5.5]{5.5} 中所考虑的情形; 下一讲将利用本讲所作的归约加以证明. 另一个有趣的情形 (其中 $h$ 不再像 \hyperref[XII.cor5.5]{5.5} 中那样为闭浸入) 将在 XV 中研究. 再提及与之密切相关的情形: $X$ 是某个 Noether 环 $A$ 的谱, 此环对于理想 $J$ 定义的拓扑分离且完备, $Y=\Spec(A/J)$ 是 $X$ 中由 $J$ 定义的闭子概形. 此时由 \EGA IV 18.3.2 可直接以\pageoriginale $\textup{ b'')}$ 的形式 (其中 $\LL=\PP$) 验证条件 \hyperref[XII.prop6.5]{6.5} (ii). 很有可能 \hyperref[XII.prop6.5]{6.5} (iii) 的条件也对每个 $n$ 成立; 即对 $X$ 上的每个挠层 $F$, 同态
$$
H^n(X,F)\longrightarrow H^n(Y,F|Y)
$$
为双射. 可以提出广泛得多的猜想, 它们与将 Hensel 化构造推广到非局部概形有关. 首先, 还不知道有任何例子, 其中 $h$ 为闭浸入, 且当 \hyperref[XII.prop6.5]{6.5} (i) 中括号内的条件成立, 即按 \EGA IV 18.5.5 的术语 $(X,Y)$ 为 \emph{Hensel 对}(henselian pair) 时, \hyperref[XII.prop6.5]{6.5} (ii) 与 (iii) (其中 $\LL=\PP$) 的条件不成立. 因而可以猜想, 每个这样的对都满足 \hyperref[XII.prop6.5]{6.5} (ii) 与 (iii) 的条件 (其中 $\LL=\PP$, 对每个 $n$).

另一方面, 可以问: 对一个由拟紧且拟分离概形 $X$ 及其闭子概形 $Y$ 构成的对 $(X,Y)$, 在什么条件下, 下述性质蕴含该对是 Hensel 对:

$(\ast)$ 对 $X$ 上的每个平展概形 $X'$, 置 $Y'=\fprod{X'}{Y}{X}$. 典范映射
$$
\Gamma(X'/X)\longrightarrow\Gamma(Y'/Y)
$$
为双射 (\cf \EGA IV 18.5.4 b)). 此条件\emph{先验地}显然弱于 Hensel 条件 \hyperref[XII.prop6.5]{6.5} (i), 而且事实上严格较弱: 例如, 当 $X$ 是域 $k$ 上维数为 $\geq2$ 的正规不可约射影概形, $Y$ 是 $X$ 的超平面截面时, 此条件成立 (这容易由 \SGA 2 XII 2.1 取 $S=\Spec(k)$ 得到), 但对 $(X,Y)$ 显然不是 Hensel 对. 不过, 当 $X$ 仿射时, 条件 $(\ast)$ 也许蕴含 $(X,Y)$ 是 Hensel 对. 如果如此, 则对每个仿射概形 $X$ 及 $X$ 的每个闭子集 $Y$, 应当能用一个\pageoriginale 类似局部环 Hensel 化的沿 $Y$ 的平展局部化过程, 与之关联一个 \emph{Hensel} 对 $(X',Y)$, 其中 $X'$ 在 $X$ 上“pro-平展”. 关于 $(X,Y)$ 的另一种假设 (由 Hironaka 与 Rossi 关于代数簇某些子簇可收缩性的结果所启发), 也许也能使 Hensel 性由较弱的性质 $(\ast)$ 推出, 即: 浸入 $Y\rightarrow X$ 是余维为 $1$ 的正则浸入, 且余法层 $\NN_{Y/X}=\JJ/\JJ^2$ (其中 $\JJ$ 是在 $X$ 中定义 $Y$ 的理想) 是 $Y$ 上的\emph{丰沛}可逆层. (相反, 在上面考虑的反例中, 这个余法层是反丰沛的, 即其逆是丰沛的.)
\end{enonce*}

\footnotetext{\hypertarget{sga4.exp.XII.note.section7}{}只对 Noether 情形下\hyperref[XII.sec5]{第 5 节}陈述感兴趣的读者可以略去本节.}

\section[Chow 引理的一个变体 ${}^{(*)}$.]{Chow 引理的一个变体
\texorpdfstring{\protect\hyperlink{sga4.exp.XII.note.section7}{${}^{(*)}$}}{(*)}.}\etiquette[7]{XII.sec7}

需要 (\EGA II 5.6.1) 的如下变体, 去掉所有 Noether 性及不可约性假设.

\begin{enonce*}{引理 7.1}\etiquette[7.1]{XII.lem7.1}
设 $S$ 为拟紧且拟分离概形, $f:X\rightarrow S$ 为有限型分离态射, 且 $X$ 非空. 存在交换图表
$$
\xymatrix@=1.25cm{
X \ar[rd]_{f}&    &\bar{X}\ar[ll]_{\pi}\ar[ld]^{\bar{f}}\\
  &  S &
}
$$
其中\pageoriginale $\pi$ 射影, $\bar{f}$ 拟射影, 并且存在 $X$ 的非空开子集 $U$, 使 $\pi$ 诱导从 $\bar{U}=\fprod{U}{\bar{X}}{X}=\pi^{-1}(U)$ 到 $U$ 的同构.
\end{enonce*}

\noindent
\emph{证明.} 先作下述观察. 设 $i:X'\rightarrow X$ 为 $X$ 的一个闭 $S$ 子概形, $V\subset X'$ 为在 $X$ 中开且在 $X'$ 中稠密的子概形. 若引理对 $X'$ 成立, 则对 $X$ 也成立. 事实上, 设
$$
\xymatrix@=1.25cm{
X' \ar[rd]_{f'}&    &\bar{X}'\ar[ll]_{\pi'}\ar[ld]^{\bar{f}'}\\
  &  S &
}
$$
为陈述中的那种图表, $U'\neq\varnothing$ 为 $X'$ 的开子集, 使 $\pi'$ 诱导同构 $\bar{U}'={\pi'}^{-1}(U')\rightarrow U'$. 由于 $V$ 在 $X'$ 中稠密, 有 $U'\cap V\neq\varnothing$, 因而可将 $U'$ 换成这个开子集. 则 $U'$ 在 $X$ 中开, 置 $\pi=i\pi'$、$\bar{f}=\bar{f}'$ 及 $U=U'$.

现在, $X$ 有由非空仿射开子集 $U_k$、$k=1,\ldots,n$ 构成的有限覆盖. 关于 $S$ 的假设蕴含每个 $U_k\rightarrow S$ 都拟仿射, 因而更是拟紧的. 对 $n$ 作归纳, 为每个 $X$ 证明引理. 若各 $U_k$ 的交 $U$ 在 $X$ 中稠密, 可以照搬 (\EGA II 5.6, B, C, D) 的证明, 并作必要的记号变更. 若此交非空, 则将 $X$ 换成 $U$ 的闭包 $X'$ (如上述观察中那样), 并将 $U_k$ 换成 $U'_k=U_k\cap X'$, 归结为交稠密的情形. 最后, 假设 $U$ 为空, 选取 $r$、$1\leq r<n$, 使
$$
V=\bigcap_{k=1}^{r}U_k\neq\varnothing
\qquad\hbox{但}\qquad
V\cap U_{r+1}=\varnothing;
$$
由于 $U_1\neq\varnothing$, 显然可以如此选取. 现在可以将 $X$ 换成 $V$ 的闭包 $X'$, 将每个 $U_k$ 换成 $U_k\cap X'$, 并舍去空的各项. 由于 $U_{r+1}\cap X'=\varnothing$ (它在 $X'$ 中开, 且与稠密开集 $V$ 不交), 结果由归纳假设得到.

\section{最后的归约}\etiquette[8]{XII.sec8}\pageoriginale

\begin{enonce*}{引理 8.1}\etiquette[8.1]{XII.lem8.1}
为证明断言 {\rm \hyperref[XII.thm5.1]{5.1} (i)、(ii)、(iii)} 中的任一个, 只须在 $f$ \emph{射影}且 $S$ 为 Noether 概形的假设下证明之.
\end{enonce*}

先假设已在 $f$ 射影时证明 \hyperref[XII.thm5.1]{5.1}; 说明如何对任意 $f$ 得到相同断言. 利用 \hyperref[XII.lem6.1]{6.1}, 归结为 \hyperref[XII.cor5.5]{5.5} 的情形, 其中 $S$ 严格局部. 现在取 $Y=X_{\circ}$ 应用 \hyperref[XII.cor6.10]{6.10}. 给定如 \hyperref[XII.cor6.10]{6.10} 中的 $(\xi,X')$, 应用 Chow 引理 \hyperref[XII.lem7.1]{7.1}, 找到 $S$ 态射 $p:\bar{X}\rightarrow X'$, 其中 $\bar{X}$ 在 $S$ 上射影, 使 $p(\bar{X})$ 包含 $X'$ 的一个非空开子集 (\hyperref[XII.cor6.10]{6.10} 的条件 $1^{\circ}$). 由于 $p$ 射影, 根据假设, $(p,h)$ 的基变换同态为同构 (\hyperref[XII.cor6.10]{6.10} 的条件 $2^{\circ}$). 最后, 由于 $\bar{X}$ 在 $S$ 上射影, 嵌入 $\bar{Y}=\bar{X}_{\circ}\rightarrow\bar{X}$ 根据假设在上同调群上诱导同构
$$
H^{i}(\bar{X},\bar{F})\longrightarrow
H^{i}(\bar{X}_{\circ},\bar{h}^{\ast}(\bar{F})),
$$
这正是 \hyperref[XII.cor6.10]{6.10} 的条件 $3^{\circ}$. 因而 \hyperref[XII.cor6.10]{6.10} 的所有条件均满足, 结论得证.

再说明, 为对 $f$ 射影的每个情形 $(f,F,g)$ 验证 \hyperref[XII.thm5.1]{5.1} 的三个断言中的任一个, 还可假设 $S$ 为 Noether 概形. 由 \hyperref[XII.lem6.1]{6.1}, 可假设 $S$ 严格局部, 且 $S'\rightarrow S$ 为闭点的嵌入; 因而只须验证 $S$ 严格局部时的 \hyperref[XII.cor5.5]{5.5}. 此外, 由于 $X$ 嵌入射影概形 $\PP^r_S$, 相容性 \hyperref[XII.prop4.4]{4.4} (ii) 允许将 $X$ 换成 $\PP^r_S$, 将 $F$ 换成其在 $\PP^r_S$ 上的正像. 因而归结为情形 $X=\PP^r_S$.

写成 $S=\Spec(A)$. 环 $A$ 是 Noether Hensel 局部子环 $A_i$ 的滤过正向极限, 其中嵌入 $A_i\rightarrow A$ 为局部同态: 可取 $A$ 的有限生成 $\ZZ$ 子代数在 $A$ 的极大理想所诱导的素理想处的严格 Hensel 化. 在这些条件下, $A$ 的剩余域 $k$ 同样是各 $A_i$ 的剩余域 $k_i$ 的滤过正向极限 (\EGA IV 5.13.1). 置 $S_i=\Spec(A_i)$ 及 $X_i=\PP^r_{S_i}$. 假设如 VII \hyperref[VII.thm5.7]{5.7} 中那样, 在各 $X_i$ 上给定一个层的正向系 $(F_i)$, 其层属于所考虑断言中的类型, 以 $F$ 表示它在 $X$ 上定义的层. 以 $F_{i\circ}$ (分别地, $F_{\circ}$) 表示 $F_i$ 在 $S_i$ 的闭点 $s_i$ 上的纤维 $X_{i\circ}$ 上所诱导的层 (分别地, $F$ 在 $S$ 的闭点 $s$ 上的纤维 $X_{\circ}$ 上所诱导的层). 于是有交换图表
$$
\xymatrix{
\displaystyle\mathop{\varinjlim}_{i} H^{n}(X_i,F_i) \ar[r] \ar[d]&
H^{n}(X,F)\ar[d]\\
\displaystyle\mathop{\varinjlim}_{i} H^{n}(X_{i\circ},F_{i\circ}) \ar[r] &
H^{n}(X_{\circ},F_{\circ}),
}
$$
由 VII \hyperref[VII.thm5.7]{5.7}, 两个水平箭头为同构. 因而为证明第二个竖直箭头为同构, 只须为第一个证明这一点; 故只须证明同态
$$
H^n(X_i,F_i)\longrightarrow H^n(X_{i\circ},F_{i\circ})
$$
为同构. 另一方面, $X$ 上所考虑类型的每个层 $F$ 都同构于上述类型的某个正向系的极限; 例如由 \hyperref[XII.cor5.2]{5.2} 及 \hyperrefext[IX][IX.prop1.6]{1.6} (iii), 可取 $F_i=u_{i*}(F)$, 其中 $u_i:X\rightarrow X_i$ 为典范态射. 至此完成了 \hyperref[XII.lem8.1]{8.1} 的证明.

\begin{enonce*}{推论 8.2}\etiquette[8.2]{XII.cor8.2}
定理 {\rm \hyperref[XII.thm5.1]{5.1} (i)} 成立.
\end{enonce*}

事实上, 只须结合 \hyperref[XII.lem8.1]{8.1}、判据 \hyperref[XII.prop6.5]{6.5} (i) ${\rm b'')}$ 及 \hyperref[XII.lem5.8]{5.8}.

\begin{enonce*}{引理 8.3}\etiquette[8.3]{XII.lem8.3}
为\pageoriginale 证明断言 {\rm \hyperref[XII.thm5.1]{5.1} (i)、(ii)、(iii)} 中的任一个, 只须在 $f$ 为相对维数 $\leq 1$ 的射影态射且 $S$ 为 Noether 概形时证明之.
\end{enonce*}

事实上, \hyperref[XII.lem8.1]{8.1} 已经表明, 可以只考虑 $f$ 射影且 $S$ 为 Noether 概形的情形, 前面用过的一个论证又允许假设 $X=\PP^r_S$. 对 $r$ 作归纳, 注意根据假设, 定理对 $r=1$ 成立. 因而可以假设 $r\geq2$, 且已经对 $\PP^{r-1}_S\rightarrow S$ 证明定理 (在每个基 $S$ 上). 按通常方式将 $Z=\PP^{r-2}_S$ 嵌入 $X=\PP^r_S$, 设 $p:\bar{X}\rightarrow X$ 为 $X$ 沿 $Z$ 的爆破. 直接计算 (\cf \EGA V) 给出一个自然态射 $p_1:\bar{X}\rightarrow X_1=\PP^1_S$, 使 $\bar{X}$ 成为一个局部同构于 $\PP^{r-1}_{X_1}$ 的 $X_1$ 概形 (它是 $X_1$ 上与 $X_1$ 上一个适当的秩为 $r$ 的局部自由层关联的射影丛). 另一方面, $p$ 的纤维维数为 $\leq1$, $f_1:X_1\rightarrow S$ 的纤维也是如此. 因而在交换图表
$$
\xymatrix{
   & \bar{X} \ar[ld]_{p} \ar[rd]^{p_1} &          \\
 X \ar[rd]_{f}&           & X_1\ar[ld]^{f_1}\\
   &    S      &
}$$
中, 可以对 $p$ 与 $f_1$ 应用相对维数为一的假设, 对 $p_1$ 应用归纳假设. 由于 $p$ 为满射, 根据 \hyperref[XII.prop6.8]{6.8}, 只须对复合 $\bar{f}=fp=f_1p_1$ 证明 \hyperref[XII.thm5.1]{5.1}. 但已经可以对 $f_1$ 与 $p_1$ 都应用 \hyperref[XII.thm5.1]{5.1}. 考虑到 \hyperref[XII.prop6.11]{6.11}, 得到 \hyperref[XII.thm5.1]{5.1} 对它们的复合 $f_1p_1$ 也成立. 至此完成了 \hyperref[XII.lem8.3]{8.3} 的证明.

\begin{enonce*}{引理 8.4}\etiquette[8.4]{XII.lem8.4}
为\pageoriginale 证明 {\rm \hyperref[XII.thm5.1]{5.1} (ii)}, 只须在 $f$ 为相对维数 $\leq1$ 的射影态射, 且 $S$ 为 Noether 严格局部概形的情形下, 证明“基本群的特殊化定理” {\rm \hyperref[XII.thmbis5.9]{5.9 bis}}.
\end{enonce*}

这立即由 \hyperref[XII.lem8.3]{8.3}、\hyperref[XII.lem6.1]{6.1} 及 \hyperref[XII.prop6.5]{6.5} (ii) (判据 b)) 得到.

\hyperref[XII.lem8.4]{8.4} 情形下的特殊化定理将在下一讲 (\hyperrefext[XIII][XIII.sec2]{2}) 中直接证明.

\begin{enonce*}{引理 8.5}\etiquette[8.5]{XII.lem8.5}
为证明 {\rm \hyperref[XII.thm5.1]{5.1} (iii)}, 只须证明 {\rm \hyperref[XII.thm5.1]{5.1} (ii)}, 并且再证明下述命题 (它将在下一讲 {\rm (\hyperrefext[XIII][XIII.sec3]{3})} 中证明, 或见 {\rm (\EGA IV 21.9.12)}):
\end{enonce*}

\begin{enonce*}{命题 8.6}\etiquette[8.6]{XII.prop8.6}
设 $S$ 为 Hensel Noether 概形, $f:X\rightarrow S$ 为相对维数 $\leq1$ 的射影态射, $X_{\circ}$ 为 $X/S$ 的闭纤维. 对每个概形 $Z$, 以 $\Pic Z=H^1(Z,(G_m)_Z)$ 表示 $Z$ 的 Picard 群. 则限制同态
$$
\Pic X\longrightarrow \Pic X_{\circ}
$$
为满射.
\end{enonce*}

\noindent
\emph{\hyperref[XII.lem8.5]{8.5} 的证明.} 由 \hyperref[XII.prop6.5]{6.5} (iii) 及 \hyperref[XII.cor8.2]{8.2}, 只须证明
$$
\varphi^q:H^q(X,\ZZ/\ell^\nu \ZZ)\longrightarrow
H^q(X_{\circ},\ZZ/\ell^{\nu}\ZZ)
$$
对每个 $q\geq1$ 及每个素数 $\ell$ 为满射. 对于 $q=1$, 满射性已由 \hyperref[XII.thm5.1]{5.1} (ii) 得知. 此外, 由于 $X_{\circ}$ 是可分闭域 $k$ 上维数 $\leq1$ 的射影概形, 有
$$
H^q(X_{\circ},\ZZ/\ell^{\nu}\ZZ)=0
$$
在 $\ell\neq\operatorname{char} k$ 时对 $q>2$ 成立 (\hyperrefext[X][X.cor4.3]{4.3}), 在 $\ell=\operatorname{char} k$ 时对 $q>1$ 成立 (\hyperrefext[X][X.cor5.2]{5.2}). 因而对于这些 $q$ 值, 满射性自动成立.

还\pageoriginale 须考虑素数 $\ell\neq\operatorname{char} k$ 时的 $q=2$. 置 $n=\ell^{\nu}$. $n$ 次单位根层 $(\mu_n)_S$ 局部同构于 $(\ZZ/n\ZZ)_S$, 由于 $S$ 严格局部, 存在非典范同构
$$
(\mu_n)_S\simeq(\ZZ/n\ZZ)_S,
$$
因而 $(\mu_n)_X\simeq(\ZZ/n\ZZ)_X$.

对 $X$ 与 $X_{\circ}$ 应用 Kummer 理论 \hyperrefext[IX][IX.Thr3.2]{3.2}, 得到正合列态射
$$
\xymatrix@R=1.5cm{
\Pic X \ar[r] \ar[d]^{\epsilon}&
H^{2}(X,(\mu_n)_X) \ar[d]^{\varphi^2}\ar[r]&
H^{2}(X,(G_m)_X)\ar[d]\\
\Pic X_{\circ} \ar[r]&
H^2(X_{\circ},(\mu_n)_{X_{\circ}}) \ar[r]&
H^2(X_{\circ},(G_m)_{X_{\circ}}).
}
$$
由 \hyperrefext[IX][IX.thm4.6]{4.6}, 有 $H^2(X_{\circ},(G_m)_{X_{\circ}})=0$. 因而 $\epsilon$ 的满射性蕴含 $\varphi^2$ 的满射性, 引理得证.

\begin{thebibliography}{99}
\bibitem{sga4.exp.XII.bib.1} \textsc{Giraud} J., \emph{Cohomologie non-abélienne}, 学位论文 (即将出版), 以及 C.R., vol. 260, pp. 2392 与 2656 (1965) 中的两篇短文.

\bibitem{sga4.exp.XII.bib.2} \textsc{Godement} R., \emph{Théorie des faisceaux}, Paris
(1958).

\bibitem{sga4.exp.XII.bib.3} \textsc{Verdier} J.L., \emph{Catégories dérivées},
Institut des Hautes Études Scientifiques (1964).
\end{thebibliography}



% END native input: chapters/exp-12-zh.tex
% BEGIN native input: chapters/exp-13-zh.tex
\chapter{固有态射的基变换定理：证明的完成}
%\addtocontents{toc}{M. Artin}
\begin{center}
M. Artin
\end{center}
\phantomsection\MakeLinkTarget*{sga4.exp.XIII}\label{sga4.exp.XIII}

\setcounter{pageoriginal}{131}

\section{射影且平坦的情形}\etiquette[1]{XIII.sec1}\pageoriginale

回顾 \hyperrefext[XII][XII.thm5.9]{5.9 bis} 在此处所考虑情形下的陈述：

\begin{enonce*}{命题 1.1}\etiquette[1.1]{XIII.prop1.1}
设 $S$ 为 Noether Hensel 环（henselian ring）的谱，
$f:X\rightarrow S$ 为射影且平坦的态射（projective and flat morphism）。设 $X_{\circ}$ 为
$X/S$ 的闭纤维。限制函子
$$
\Et(X)\longrightarrow \Et(X_{\circ})
$$
是范畴等价。
\end{enonce*}

如前所述，由 \hyperrefext[XII][XII.prop6.5]{6.5} (i) 可知该函子完全忠实。
还须证明，$X_{\circ}$ 的每个平展覆盖（étale covering）$Y_{\circ}$ 都由
$X$ 的某个平展覆盖 $Y$ 诱导。

\begin{enonce*}{引理 1.2}\etiquette[1.2]{XIII.lem1.2}
设 $S$ 局部 Noether，$f:X\rightarrow S$ 为射影且平坦的态射，
$Y\subset X$ 为子概形。终函子的子函子
$\epsilon=\prod_{X/S}Y/X$ 
“表示条件 $Y=X$”，即对 $S'\rightarrow S$ 有
$$
\epsilon(S')=
\begin{cases}
\phi & \mbox{若 }\quad Y'\neq X',\\
\{\phi\} &\mbox{若 }\quad Y'=X',
\end{cases}
$$
（其中\pageoriginale 撇号表示基变换（base change）
$S'\rightarrow S$ 的结果），可由 $S$ 的一个子概形 $Z$ 表示。若
$Y$ 是开的（\resp 闭的），则 $Z$ 亦然。
\end{enonce*}

\noindent
\emph{证明}。设 $Y$ 为 $X$ 的某个开子集 $U$ 中的闭子概形，
$C$ 为 $U$ 的闭补集。由于 $f$ 固有（proper），
$C$ 在 $S$ 中的像为闭子集 $D$，且显然条件
$U=X$ 由 $S$ 的开子集 $S-D$ 表示。因此可假设
已有 $U=X$，即 $Y$ 是闭子概形。
考虑 $\Oo_X$-模的满射态射
$\Oo_X\xrightarrow{\alpha}\Oo_Y$；条件 $Y=X$
等价于存在态射
$\Oo_Y\xrightarrow{\beta}\Oo_X$，使复合 $\beta\alpha$ 为
恒等态射，即存在 $\Hhom(\Oo_Y,\Oo_X)$ 的一个截面，
提升 $\Hhom(\Oo_X,\Oo_X)$ 的恒等截面。
由 \EGA III 7.7.8 可知，将 $S'/S$ 对应到
$\Hhom(\Oo_Y,\Oo_X)$（\resp $\Hhom(\Oo_X,\Oo_X)$）的
截面集的函子，可由 $S$ 上的向量丛（vector bundle）$V_1$（\resp $V_2$）表示。
设 $\varphi:V_1\rightarrow V_2$ 为由 $\Oo_X\rightarrow\Oo_Y$ 诱导的态射
$u\mapsto u\alpha$，$s:S\rightarrow V_2$ 为
对应于恒等态射 $\Oo_X\rightarrow\Oo_X$ 的截面。
截面 $s$ 在 $\varphi$ 下的逆像可与 $S$ 的一个闭子概形
$Z$ 等同，且显然 $Z$ 表示函子 $\epsilon$，
引理得证。

\begin{enonce*}{引理 1.3}\etiquette[1.3]{XIII.lem1.3}
设 $S$ 为局部 Noether 概形，$f:X\rightarrow S$ 为
射影且平坦的态射，$E$ 为 $X$ 上的局部自由层（locally free sheaf）。
设
$$
QE:(\Sch)/S\longrightarrow(\Ens)
$$
为如下函子：它将 $S'/S$ 对应到 $E'=E\otimes_{\Oo_S}\Oo_{S'}$ 的
赋有平展 $\Oo_{X'}$-代数结构的商层 $A'$ 所成的集合 $QE(S')$。
此函子可由 $S$ 上的局部有限型概形表示。
\end{enonce*}

\noindent
\emph{证明}。\pageoriginale 函子 $\sheaf{\Quot}(E)$ 可由一个局部有限型概形
（TDTE IV，第~221 讲）表示，记为 $Z/S$；其中
$\sheaf{\Quot}(E)(S')$ 是 $E'$ 的在 $S'$ 上平坦的商所成的集合。
设 $M_u$ 为 $X_Z$ 上 $E\otimes_{\Oo_S}\Oo_Z$ 的“泛”商。
$\sheaf{\Quot}(E)$ 中由在 $X'$ \emph{上}平坦的商所成的子函子，
由 $Z$ 的一个开子集 $U$ 表示。事实上，
设 $M'$ 为 $E'$ 的一个商，由某个 $S$-态射 $S'\rightarrow Z$ 给定。
由于 $M'$ 在 $S'$ 上平坦，它在 $X'$ 上平坦等价于：
对每个点 $s'\in S'$，模 $M'_{s'}$ 在纤维 $X'_{s'}$ 上平坦
（\EGA IV 5.9）。这又等价于
$M_u|_{X_z}$ 在纤维 $X_z$ 上平坦，其中 $z$ 为
$s'$ 的像。然而，由 \EGA IV 11.1.5，使
$M_u|_{X_z}$ 在 $X_z$ 上平坦的点 $z\in Z$ 所成的集合 $U$，
也就是使 $M_u$ 在 $X_z$ 的每一点处均在 $X_Z$ 上平坦的点集，
在 $Z$ 中是开的，故断言成立。

现在只须表示相对函子 $QE/U[f]$。于是将
$S$ 换为 $U$，$X$ 换为 $\fprod{X}{U}{S}$ 后，
问题就归结为以下引理：

\begin{enonce*}{引理 1.4}\etiquette[1.4]{XIII.lem1.4}
设 $S$ 为局部 Noether 概形，$f:X\rightarrow S$ 为
射影且平坦的态射，$M$ 为局部自由
$\Oo_X$-模。函子
$\Alget:(\Sch)/S\rightarrow(\Ens)$ 将 $S'/S$ 对应到
$M'$ 上的在 $\Oo_{X'}$ 上平展的代数结构所成的集合
$\Alget(S')$；此函子可由 $S$ 上的有限型概形表示。
\end{enonce*}

\noindent
\emph{证明}。为简便起见，以下说“可表示”时，
均理解为“可由一个\emph{有限型}相对概形表示”。
下文出现的所有向量丛均由\emph{凝聚}层定义，
因而都是有限型的。$M'$ 上的一个双线性合成律，
由局部自由层 $\Hhom(M'\otimes M',M')$ 的一个截面给定。
由于 $\Hhom(M'\otimes M',M')$ 在 $S$ 上平坦，
由\pageoriginale \EGA III 7.7.6 可知，“双线性合成律”函子
由 $S$ 上的一个向量丛 $Z$ 表示。还须验证，
合成律满足结合律、交换律及具有单位截面的附加条件，
可由 $Z$ 的一个子概形表示。考虑相对函子，
即可将 $S$ 换为 $Z$，因而可假设已给定
$M$ 上的一个合成律，只须证明表示上述条件的
$S$ 的子函子可由 $S$ 的子概形表示。

对结合律而言，条件是
$\Hhom(M\otimes M\otimes M,M)$ 的两个截面在基变换
$S'\rightarrow S$ 后相等。设 $Y$ 为这两个截面在对应于
$\Hhom(M\otimes M\otimes M,M)$ 的向量丛上的纤维积。
概形 $Y$ 可与 $X$ 的一个闭子概形等同，而在
$S'\rightarrow S$ 后成立的结合律表示为条件 $Y'=X'$。
因此，由 \hyperref[XIII.lem1.2]{1.2}，它定义了一个
可由 $S$ 的闭子概形表示的 $S$ 的子函子。
对交换律而言，条件是 $\Hhom(M\otimes M,M)$ 的两个截面相等，
由 \hyperref[XIII.lem1.2]{1.2}，这同样定义一个可由闭子概形表示的
$S$ 的子函子。最后，如今可假设合成律
已经满足结合律及交换律，并考虑单位截面存在的条件；
此时这样的截面必定唯一。将 $S'/S$ 对应到
$M'$ 的截面集的函子，由 \EGA III 7.7.6 可由一个向量丛表示，
于是只须表示此向量丛的对应于\emph{单位截面}的子函子
（这里利用刚才指出的唯一性）。如常转到相对情形，
即可假设已经给定 $M$ 的一个截面 $\epsilon$，
并证明表示它变为\pageoriginale 单位截面这一条件的
$S$ 的子函子可表示。然而，$\epsilon$ 为单位的条件，
正是显然的态射的复合
$M'\rightarrow M'\otimes M'\rightarrow M'$ 为恒等态射，
即 $\Hom(M,M)$ 的两个指定截面相等。
再次由 \hyperref[XIII.lem1.2]{1.2} 即得结论。

最后，假设局部自由模 $M$ 已赋有一个含幺交换代数结构。
它平展的点所成的集合在 $X$ 中是开的（\SGA I 4.5），
因此可在 $S$ 上由一个（开）子概形表示。
这就证明了 \hyperref[XIII.lem1.4]{1.4}，
从而也证明了 \hyperref[XIII.lem1.3]{1.3}。

\setcounter{subsection}{4}
\subsection{}\etiquette[1.5]{XIII.subsec1.5}
现设 $Z/S$ 为表示 \hyperref[XIII.lem1.3]{1.3} 中所考虑函子的对象，
并考虑 $Z$ 在某点 $z$ 处在 $S$ 上光滑的条件。
将 \SGA I III 3.1 译为函子 $QE$ 的语言，
即得到如下条件：

设 $B$ 为 Artin 局部环，$\Mm=\rad B$，$J$ 为
$B$ 的满足 $J\Mm=0$ 的理想，置 $S'=\Spec B$，$S''=\Spec B/J$，
并设给定态射 $S'\rightarrow S$。对于基变换
$S'\rightarrow S$ 及 $S''\rightarrow S$ 的结果，采用显然的记号。

\begin{itemize}
\item[$(\ast)$] 对 $E''$ 的每个赋有平展
$\Oo_{X''}$-代数结构、且对应于像为 $z$ 的
$S$-态射 $S''\rightarrow Z$ 的商 $A''$，都存在
$E'$ 的一个赋有平展 $\Oo_{X'}$-代数结构的商
$A'$，使其诱导出所给的商。
\end{itemize}

由于 $X'$ 和 $X''$ 的底空间相同，诱导
$A''$ 的代数 $A'$ 已在唯一同构意义下确定
（\SGA I 8.3）。\pageoriginale 因此，可将条件
$({\ast}{\ast})$ 表述如下：

\begin{itemize}
\item[$({\ast}{\ast})$] 设 $A'$ 为 $\Oo_{X'}$ 上的平展代数，
$A''$ 为 $A'$ 在 $X''$ 上诱导的代数。每个
满同态
$$
E''\xrightarrow{\varphi''}A''
$$
若对应于像为 $z$ 的态射 $S''\rightarrow Z$，则均可提升为同态
（由 Nakayama 引理，该同态必定为满同态）
$$
E'\xrightarrow{\varphi'}A'{\quoi}.
$$
\end{itemize}

由正合列
$$
0\longrightarrow J\longrightarrow B\longrightarrow B/J\longrightarrow
0
$$
得到 $X'$ 上的层的正合列：
$$
0\longrightarrow J\otimes\Hhom(E',A')\longrightarrow
\Hhom(E',A')\longrightarrow\Hhom(E'',A'')\longrightarrow
0{\quoi}.
$$
提升 $\varphi''$ 的障碍位于
$H^1(X',J\otimes\Hhom(E',A'))$ 中。设
$S_{\circ}=\Spec B/\Mm$。由于 $J\Mm=0$，有
$$
J\otimes\Hhom(E',A')\simeq
J\otimes\Hhom(E_{\circ},A_{\circ}){\quoi},
$$
且显然
$$
H^1(X',J\otimes\Hhom(E_{\circ},A_{\circ}))\simeq
J\otimes H^1(X_{\circ},\Hhom(E_{\circ},A_{\circ})){\quoi}.
$$
考虑到上同调与基域扩张可交换，得到：

\begin{enonce*}{推论 1.6}\etiquette[1.6]{XIII.cor1.6}
设 $Z/S$ 为表示 \hyperref[XIII.lem1.3]{1.3} 中函子的概形，
置 $X_Z=\fprod{X}{Z}{S}$，并设 $A_Z$ 为
$E_Z$ 的（“一般”）平展商代数。则在每个满足
$$
H^1(X_{\circ},\Hhom(E_{\circ},A_{\circ}))=0{\quoi},
$$
的点 $z_{\circ}$ 处，$Z/S$ 都是光滑的；其中\pageoriginale 下标零
表示限制到 $X_Z/Z$ 在点 $z_{\circ}$ 处的纤维
$X_{\circ}$。
\end{enonce*}

现在可以完成 \hyperref[XIII.prop1.1]{1.1} 的证明了。采用该命题的记号，
设 $Y_{\circ}/X_{\circ}$ 为一个平展覆盖，
由 $\Oo_{X_{\circ}}$ 上平展的代数层 $A_{\circ}$ 给定。
当 $n$ 充分大时，$A_{\circ}(n)$ 由其截面生成
（\EGA III 2.2.1）。于是有满射
$$
\Oo^N_{X_{\circ}}\longrightarrow A_{\circ}(n)\longrightarrow
0{\quoi},
$$
因而有满射
$$
\Oo^N_{X_{\circ}}(-n)=E_{\circ}\longrightarrow
A_{\circ}\longrightarrow 0{\quoi}.
$$
取 $n$ 充分大，使另有
\begin{equation*}
H^1(X_{\circ},\Hhom(E_{\circ},A_{\circ}))=0{\quoi},\tag{$+$}
\end{equation*}
由 \EGA III 2.2.1，这样的选取是可能的，因为
$$
\Hhom(E_{\circ},A_{\circ})
=[\Hhom(\Oo^N_{X_{\circ}},A_{\circ})](n)
=[A_{\circ}(n)]^N{\quoi}.
$$
置 $E=\Oo^N_X(-n)$，设 $Z/S$ 为表示
\hyperref[XIII.lem1.3]{1.3} 中所考虑函子的概形。
$E_{\circ}$ 的商 $A_{\circ}$ 对应于 $Z/S$ 的闭纤维中的
一个点 $z_{\circ}\in Z$，它是 $k(s)$ 上的有理点；
由 \hyperref[XIII.subsec1.5]{1.5} 及 $(+)$，
$Z/S$ 在 $z_{\circ}$ 处光滑。由于 $S$ 为\emph{Hensel} 环的谱，
$Z/S$ 有一个经过 $z_{\circ}$ 的截面；例如，
这可直接由光滑性的原始定义（\SGA II 1.1）得到。
因而，$A_{\circ}$ 由 $X$ 上平展的某个代数 $A$
（它是 $E$ 的商）诱导，这就完成了
\hyperref[XIII.prop1.1]{1.1} 的证明。


\section{相对维数至多为一的情形}\etiquette[2]{XIII.sec2}\pageoriginale

断言如下。

\begin{enonce*}{命题 2.1}\etiquette[2.1]{XIII.prop2.1}
采用 \textup{\hyperrefext[XII][XII.thm5.9]{5.9 bis}} 的记号，若
$S$ 为 Noether 严格局部概形，且 $f$ 为相对维数
$\leq1$ 的射影态射，则函子
$\Et(X)\rightarrow\Et(X_{\circ})$ 是范畴等价。

由 \textup{\hyperrefext[XII][XII.lem8.5]{8.5}}，这将完成
\textup{\hyperrefext[XII][XII.thm5.1]{5.1} (ii)} 的证明；
后者同时推广了 \hyperref[XIII.prop1.1]{1.1} 及
\hyperref[XIII.prop2.1]{2.1}。
\end{enonce*}

\begin{enonce*}{引理 2.2}\etiquette[2.2]{XIII.lem2.2}
设 $S$ 为 Noether 局部环 $A$ 的谱，$X/S$ 为
相对维数 $\leq n$ 的射影概形。则存在相对维数
$\leq n$ 的射影且平坦的 $S$-概形 $Z/S$，以及
$S$ 上的闭浸入 $X\hookrightarrow Z$。
\end{enonce*}

\noindent
\emph{证明。}设 $X\rightarrow\PP^N_S$ 为 $S$ 上的一个射影嵌入，
$J\subset A[x_{\circ},\ldots,x_N]$ 为定义 $X$ 的
齐次理想。置 $k=A/\mathfrak m_A$，$S_{\circ}=\Spec k$，
$X_{\circ}=X\otimes_Sk$。则 $X_{\circ}$ 由齐次理想
$J_{\circ}=\Iim(J\rightarrow k[x_{\circ},\ldots,x_N])$ 定义。
由于 $\dim X_{\circ}\leq n$，众所周知，存在
$J_{\circ}$ 中的元素 $f^{\circ}_1,\ldots,f^{\circ}_{N-n}$，
定义 $\PP^N_k$ 的一个维数为 $n$ 的子概形 $Z_{\circ}$。
设 $f_i$ $(1\leq i\leq N-n)$ 为 $J$ 中提升
$f^{\circ}_i$ 的元素，$Z$ 为由
$f_1,\ldots,f_{N-n}$ 定义的 $\PP^N_S$ 的子概形。
由 \EGA IV 11.3，$Z$ 在 $S$ 上平坦；由
\EGA IV 13.1.5，其相对维数为
$\dim Z_{\circ}=n$。\hyperref[XIII.lem2.2]{2.2} 得证。

\noindent
\emph{\hyperref[XIII.prop2.1]{2.1} 的证明。}设 $X/S$ 为相对维数
$\leq1$ 的射影概形，$S$ 为 Noether 严格局部概形。
可假设 $X$（从而 $X_{\circ}$）连通。设
$X\rightarrow Z$ 为 $S$ 上的一个闭浸入，其中 $Z$ 射影、
平坦，且相对维数 $\leq1$；见
\hyperref[XIII.lem2.2]{2.2}。设 $s_{\circ}$ 为
$S$ 的闭点。有相应的\pageoriginale 范畴图式
$$
\xymatrix@C=2.5cm@R=1.25cm{
\Et(Z)\ar[r]\ar[d]^{a}&\Et(X)\ar[d]^{b}\\
\Et(Z_{\circ})\ar[r]^{c}&\Et(X_{\circ}).
}
$$
要证明 $b$ 为等价。由 \hyperrefext[XII][XII.lem5.8]{5.8}，
$b$ 完全忠实，因此只须证明它本质满射。
由 \hyperref[XIII.prop1.1]{1.1}，函子 $a$ 是等价；
因而只须证明 $c$ 本质满射。略去下标
${}_{\circ}$，问题就归结为以下命题。

\begin{enonce*}{命题 2.3}\etiquette[2.3]{XIII.prop2.3}
设 $k$ 为可分闭域，$Z$ 为维数
$\leq1$ 的局部有限型 $k$-概形，$X\hookrightarrow Z$
为连通闭子概形。则 $X$ 的每个平展覆盖都由
$Z$ 的某个平展覆盖诱导。
\end{enonce*}

\noindent
\emph{证明。}可假设 $Z$ 和 $X$ 的维数均为 $1$，
因为当 $X$ 为零维时，结论由 $k$ 的假设立即得到。
此时 $Z=X\cup Y$，其中 $Y$ 是取 $Z-X$ 的闭包所得的
闭子概形。置 $V=X\cap Y$；这是一个离散的，
即局部 Artin 的概形。态射
$X\amalg Y\rightarrow Z$ 有限且满射，故为平展覆盖范畴的
严格下降态射（strict descent morphism）
（\SGA \hyperrefext[IX][IX.cor4.7]{4.7}）。显然
\begin{equation*}
\fprod{(X\amalg Y)}{(X\amalg Y)}{Z}
=X\amalg Y\amalg V^{(1)}\amalg V^{(2)}
\quad(V^{(i)}=V){\quoi},\tag{$\ast$}
\end{equation*}
其中 $\pr_i$ 将 $V^{(i)}$ 映到 $X\amalg Y$ 的分支 $X$，
而当 $j\neq i$ 时，将 $V^{(j)}$ 映到分支 $Y$。

设 $X'/X$ 为次数 $n$ 的平展覆盖，$Y''/Y$ 为
次数 $n$ 的一个覆盖（例如次数 $n$ 的平凡覆盖），
于是 $X'\amalg Y''$ 是\pageoriginale $X\amalg Y$ 的
次数 $n$ 的覆盖。只须找出此覆盖相对于
$X\amalg Y\rightarrow Z$ 的下降资料。考察
$(\ast)$ 可知，这样的下降资料恰好等价于一个同构
$$
\fprod{X'}{V}{X}=V'\xrightarrow{\sim}
V''=\fprod{Y''}{V}{Y}{\quoi}.
$$
由于 $k$ 可分闭，$V$ 各点的剩余域亦然；
所以 $V$ 的每个平展覆盖都完全分裂。
因此 $V'\simeq V''$，结论得证。

\begin{enonce*}[remark]{注 2.4}\etiquette[2.4]{XIII.rem2.4}
可以用一个直接论证，避开 \SGA IX 中的精细结果及下降理论：
只须证明，$Z$ 的平展覆盖范畴等价于三元组
$(X',Y',\theta)$ 所成的范畴，其中
$X'$（分别 $Y'$）为 $X$（分别 $Y$）的平展覆盖，
而 $\theta$ 为一个 $V$-同构
$\fprod{X'}{V}{X}\simeq\fprod{Y'}{V}{Y}$。
对任意可表为两个闭子概形 $X$ 与 $Y$ 之并的
概形 $Z$，这一结果都成立；它可直接由
$Z$ 上拟凝聚模的类似描述证明
（此时 $X=V(\mathfrak J)$，$Y=V(\mathfrak K)$，且
$\mathfrak J\cap\mathfrak K=0$）。
\end{enonce*}

\footnotetext[1]{这一结果也以稍一般的形式出现在 \EGA IV 21.9.11 及 21.9.12 中；
特别地，在 \hyperref[XIII.cor3.2]{3.2} 中，只须要求
$f$ \emph{分离}，不必要求它固有。}
\section{关于 Picard 群的一个辅助结果 ${}^{(*)}$}\etiquette[3]{XIII.sec3}

\begin{enonce*}{命题 3.1}\etiquette[3.1]{XIII.prop3.1}
设 $X$ 为 Noether 概形，$X_{\circ}$ 为 $X$ 的一个闭子概形，
满足 \textup{(a)} $\dim X_{\circ}\leq1$，以及
\textup{(b)} 对 $X$ 的每个非空连通闭子集 $Y$，
$Y_{\circ}=Y\cap X_{\circ}$ 非空且连通。
\begin{enumerate}
\renewcommand{\theenumi}{\roman{enumi}}
\renewcommand{\labelenumi}{(\theenumi)}
\item 设 $Z_{\circ}$ 为所有仅含一个点的集合
$\overline{\{x\}}\cap X_{\circ}$ 的并，其中 $x\in\Ass\Oo_X$。
对 $X_{\circ}$ 上的每个满足 $\supp(D_{\circ})$ 不与
$Z_{\circ}$ 相交的 Cartier 除子 $D_{\circ}$，
存在 $X$ 上的一个 Cartier 除子 $D$ 诱导出 $D_{\circ}$
（且当 $D_{\circ}\geq0$ 时，可取 $D\geq0$）。

\item 假设 $X_{\circ}$ 有一个包含有限集
$\Ass\Oo_{X_{\circ}}\cup Z_{\circ}$ 的仿射开子集 $U_{\circ}$
（若 $X_{\circ}$ 有丰沛可逆层，此条件自动满足）。
则 $\Pic(X)\rightarrow\Pic(X_{\circ})$ 为满射。
\end{enumerate}
\end{enonce*}

\noindent
\emph{证明。}
\begin{enumerate}
\renewcommand{\theenumi}{\roman{enumi}}
\renewcommand{\labelenumi}{(\theenumi)}
\item 由于 $\dim X_{\circ}\leq1$，$X_{\circ}$ 上的每个
Cartier 除子 $D_{\circ}$ 显然具有离散支集，
因而可写成两个有效 Cartier 除子之差，且它们的支集均包含在
$D_{\circ}$ 的支集中。因此，(i) 归结为
$D_{\circ}$ 有效的情形。对每个
$x\in\Supp D_{\circ}$，它在 $x$ 处由一个非零因子
$f_x\in\rad(\Oo_{X_{\circ},x})$ 定义。元素 $f_x$
由 $\Oo_X$ 在 $x$ 的某个开邻域 $U_x$ 上的一个截面
$g_x$ 诱导。断言：缩小 $U_x$ 后，可使截面
$g_x$ 成为 $\Oo_{U_x}$ 中的非零因子，即
$V(g_x)\cap\Ass\Oo_{U_x}=\varnothing$。

事实上，由条件 (a) 及 (b)，对每个
$z\in\Ass\Oo_X$，集合 $\overline{\{z\}}\cap X_{\circ}$
为 $X_{\circ}$ 的连通闭子集，故或者仅含一个闭点，
或者是 $X_{\circ}$ 的若干一维不可约分支的连通并。
由于 $x\notin Z_{\circ}$，在前一种情形下有
$x\notin\overline{\{z\}}$。在后一种情形下，
不可能同时有 $x\in\overline{\{z\}}$ 和 $z\in V(g_x)$，
否则 $g_x$，从而 $f_x$，将在
$\overline{\{z\}}\cap X_{\circ}$ 的某个不可约分支上恒为零，
而该分支也是 $X_{\circ}$ 的不可约分支；
这与 $f_x$ 为 $\Oo_{X_{\circ},x}$ 中的非零因子矛盾。
因此缩小 $U_x$ 后，可\pageoriginale 假设
$V(g_x)\cap\Ass\Oo_{U_x}=\varnothing$，故 $g_x$ 为非零因子。

设 $Y_x$ 为 $V(g_x)$ 在 $X$ 中的闭包。点 $x$ 在
$Y_x\cap X_{\circ}\cap U_x$ 中孤立，因而在
$Y_x\cap X_{\circ}$ 中也孤立。由条件 (b)，$Y_x$
有一个连通分支 $Y'_x$，满足 $Y'_x\cap X_{\circ}=\{x\}$。
而且 $Y'_x\subset U_x$：事实上，
$T_x=Y'_x-(Y'_x\cap U_x)$ 在 $X$ 中闭，且满足
$T_x\cap X_{\circ}=\varnothing$，于是由条件 (b) 得
$T_x=\varnothing$。设 $D'_x$ 为 $X$ 上的 Cartier 除子，
它在 $U_x$ 上由 $Y'_x$ 诱导，在 $X-Y'_x$ 上由零除子诱导。
则 $D=\sum_xD'_x$ 为 $X$ 上的有效 Cartier 除子，
在 $X_{\circ}$ 上诱导出 $D_{\circ}$。

\item 结论由 (i) 及如下事实得到：$X_{\circ}$ 上的每个
可逆模 $\LL_{\circ}$ 都同构于某个支集不与
$Z_{\circ}$ 相交的 Cartier 除子 $D_{\circ}$ 所定义的模。
后一断言意味着，$\LL_{\circ}$ 在某个包含
$\Ass\Oo_{X_{\circ}}\cup Z_{\circ}$ 的开集 $V_{\circ}$ 上
有一个截面，且该截面在此集合的每一点均不为零。
只须将 $X_{\circ}$ 换为仿射开集 $U_{\circ}$ 后证明这一点；
在此情形下，它由熟知的事实得到：半局部概形上的可逆模
与结构层同构。最后，若 $X_{\circ}$ 有丰沛可逆层，
则 $X_{\circ}$ 的每个有限子集都包含在某个仿射开子集中。
这就证明了 (ii)。
\end{enumerate}

有如下推论，它推广了 \hyperrefext[XII][XII.prop8.6]{8.6}，
并完成了 \hyperrefext[XII][XII.thm5.1]{5.1} (iii) 的证明。

\begin{enonce*}{推论 3.2}\etiquette[3.2]{XIII.cor3.2}
设 $S$ 为 Noether Hensel 环的谱，
$f:X\rightarrow S$ 为相对维数 $\leq1$ 的固有态射。
对 $X$ 的每个底空间与 $X/S$ 的闭纤维\pageoriginale
$X_{\circ}$ 相同的闭子概形 $X'_{\circ}$，典范映射
$\Pic(X)\rightarrow\Pic(X'_{\circ})$ 均为满射。
\end{enonce*}

事实上，将命题应用于 $(X,X'_{\circ})$。
条件 (a) 由假设成立，条件 (b) 由
\hyperrefext[XII][XII.cor5.7]{5.7} 得到。还须验证
(ii) 所要求的 $U_{\circ}$ 存在；这是因为
$X'_{\circ}$ 必为射影的，正如域上的每条固有曲线都是射影的一样。

\begin{enonce*}[remark]{注}
若读者不愿假定或重建域上的固有曲线均为射影这一事实，
可以只在 $f$ 为射影态射这一附加假设下使用此推论；
这对于此处所需的应用已经足够。反之，
若承认域上的固有代数曲线的射影性，则可得出：
在 \hyperref[XIII.cor3.2]{3.2} 的假设下，$X$
必在 $S$ 上射影。事实上，取一个丰沛元素
$\xi_{\circ}\in\Pic(X_{\circ})$，将它提升为某个
$\xi\in\Pic(X)$，再应用 \EGA III 4.7.1 即可。
\end{enonce*}

\begin{thebibliography}{99}
\bibitem{sga4.exp.XIII.bib.TDTE} Grothendieck A., Technique de descente et
théorèmes d'existence en Géométrie Algébrique, in Fondements de la
Géométrie Algébrique (Extraits du Sém. Bourbaki 1957--62), Secrétariat
mathématique, 11, rue Pierre Curie.
\end{thebibliography}


% END native input: chapters/exp-13-zh.tex
% BEGIN native input: chapters/exp-14-zh.tex
\chapter[固有态射的有限性与仿射上同调维数]{固有态射的有限性定理；仿射代数概形的上同调维数}
%\addtocontents{toc}{M. Artin}
\begin{center}
M. Artin
\end{center}
\phantomsection\MakeLinkTarget*{sga4.exp.XIV}\label{sga4.exp.XIV}

\setcounter{pageoriginal}{144}

\newcommand{\degtr}{{\rm deg.tr.\ }}

本\pageoriginale 讲包含两个重要定理，其证明用到固有态射的基变换定理。

\section{固有态射的有限性定理}\etiquette[1]{XIV.sec1}

\begin{enonce*}{定理 1.1}\etiquette[1.1]{XIV.thm1.1}
设 $f:X\rightarrow S$ 是有限表现的固有态射（proper morphism）。设 $A$ 是左 Noether 环，并且对于某个适当的 $n\in\NN$ 是一个 $\ZZ/n$-代数。设 $F$ 是 $X$ 上的\emph{可构造}（constructible）集合层（\resp 群层，\resp $A$-模层）。那么层 $f_*F$（\resp $R^1f_*F$，\resp 对每个 $q\geq0$ 的 $R^qf_*F$）也是可构造的。
\end{enonce*}

特别地，我们有：

\begin{enonce*}{推论 1.2}\etiquette[1.2]{XIV.cor1.2}
设 $X$ 是可分闭域 $k$ 的谱上的固有概形，$F$ 是 $X$ 上的可构造挠阿贝尔层。那么对每个 $q\geq0$，群 $H^q(X,F)$ 都是有限群。
\end{enonce*}

\begin{enonce*}[remark]{注 1.3}\etiquette[1.3]{XIV.rem1.3}
如果去掉 $f$ 固有这一条件，定理对于阿贝尔层便不成立：取 $F=\ZZ/p$，$p=\car k$，$X=\EE^1_k$ 为仿射直线，并取 $f$ 为 $X$ 到 $\PP^1_k$ 的包含态射或结构态射 $X\rightarrow\Spec k$，即可看出这一点。然而，有如下猜想\nde{Ofer Gabber 在 2005 年纪念 Luc Illusie 和 Pierre Deligne 的会议上报告了 \hyperref[XIV.thm1.1]{1.1} 的以下推广。设 $f:X\to Y$ 是拟优 Noether 概形之间的有限型态射（\EGA IV 7.8），$F$ 是 $X_{\text{ét}}$ 上的可构造 $\mathbf Z/n$-模层，其中 $n$ 在 $Y$ 上可逆。那么 $Y_{\text{ét}}$ 上的层 $R^if_*F$ 都是可构造的，并且除有限多个 $i$ 外均为零。对于可构造集合层以及阶在 $Y$ 上可逆的群层，也有类似的断言。\newline
在 Gabber 的结果之前，\hyperref[XIV.thm1.1]{定理 1.1} 已在以下情形中得到证明。设 $S$ 是维数 $\leq1$ 的正则概形，$f:X\to Y$ 是 $S$-概形之间的有限型态射，$F$ 是 $X_{\text{ét}}$ 上的左 $A$-模层。那么层 $R^if_*F$ 都是可构造的（\SGA 4 1/2, Finitude 1.1）。}：若 $F$ 的阶与 $S$ 的\pageoriginale 剩余特征互素，且 $S$ 是“优的”（excellent）（\EGA IV 7.8），则该定理成立。当奇点消解可用时（例如特征为零时），对于优的等特征概形 $S$ 可以证明这一点（\hyperrefext[XIX][XIX.thm5.1]{5.1}）。
\end{enonce*}

\noindent
\hyperref[XIV.thm1.1]{1.1} 的\emph{证明}。断言在 $S$ 上是局部的，所以可取 $S$ 为仿射概形。于是 $S$ 是 $\ZZ$ 上有限型环的谱的极限，而定理中的资料可以由一个固有态射 $f_0:X_0\rightarrow S_0$ 以及 $X_0$ 上的一个可构造层 $F_0$，经基变换 $S\rightarrow S_0$ 得到，其中 $S_0$ 在 $\ZZ$ 上有限型；见 \EGA IV 8 和 \hyperrefext[IX][IX.cor2.7.4]{2.7.4}。由于可构造层的逆像仍然可构造（\hyperrefext[IX][IX.prop2.4]{2.4} (iii)），将 \hyperrefext[XII][XII.thm5.1]{5.1} 应用于基变换态射 $S_0\leftarrow S$，便将问题归结为证明 $f_{0*}F_0$（\resp $R^qf_{0*}F_0$）的可构造性，由此得到下面的引理。

\begin{enonce*}{引理 1.4}\etiquette[1.4]{XIV.lem1.4}
只需在 $S$ 是 $\Spec\ZZ$ 上的有限型概形时验证 \hyperref[XIV.thm1.1]{1.1}。
\end{enonce*}

\begin{enonce*}{引理 1.5}\etiquette[1.5]{XIV.lem1.5}
\hyperref[XIV.thm1.1]{1.1} 中关于集合的断言成立：即若 $F$ 是可构造集合层，则 $f_*F$ 也是。
\end{enonce*}

\noindent
\emph{证明}。可设 $S$ 是 Noether 概形（\hyperref[XIV.lem1.4]{1.4}）。取一个单射 $F\rightarrow G=\prod_i\pi_{i*}(C_i)$，其中 $\pi_i:X_i\rightarrow X$ 是有限态射，且 $C_i$ 在 $X_i$ 上常值且可构造（\hyperrefext[IX][IX.prop2.14]{2.14}）。有 $f_*F\subset f_*G$；因此，由 \hyperrefext[IX][IX.prop2.9]{2.9} (ii)，只需处理 $F=G$ 的情形。现在 $f_*G=\prod_i(f\pi_i)_*C_i$，而 $S$ 上可构造层的有限积仍然可构造（\hyperrefext[IX][IX.prop2.6]{2.6} (ibis)）。用 $X_i$ 代替 $X$，便归结为 $F=D_X$ 是取值为 $D$ 的可构造常值层的情形。由 \hyperrefext[XII][XII.thm5.1]{5.1} (i)，纤维 $(f_*F)_{\bar s}$ 同构于 $D^{C(s)}$，其中 $C(s)=\pi_0(X_{\bar s})$；\pageoriginale 因而它是有限集。由 \hyperrefext[IX][IX.prop2.13]{2.13} (iii)，只需证明函数“$f_*F$ 的纤维中的元素个数”在 $S$ 上是可构造函数。这个函数为 $\card(D^C)$，其中 $C$ 是函数“几何纤维中的连通分支个数”，而 $C$ 是可构造函数（\EGA IV 9.7.9），这就证明了结论。

\begin{enonce*}{引理 1.6}\etiquette[1.6]{XIV.lem1.6}
设 $S$ 是 Noether 概形，$f:X\rightarrow S$ 固有，$F$ 是 $X$ 上的可构造群层（\resp $A$-模层），$q$ 为等于 $0$ 或 $1$ 的整数（\resp $\geq0$ 的整数）。设 $S_i$ $(1\leq i\leq n)$ 为 $S$ 的子概形，满足 $S=\bigcup_iS_i$；设 $F_i$ 和 $f_i:X_i\rightarrow S_i$ 是这些资料在 $S_i$ 上的“限制”。若对于每个 $i$，$R^qf_{i*}F_i$ 都可构造，则 $R^qf_*F$ 也可构造。
\end{enonce*}

\noindent
\emph{证明}。由 \hyperrefext[XII][XII.thm5.1]{5.1}，$R^qf_*F$ 在 $S_i$ 上的限制同构于 $R^qf_{i*}F_i$，而引理由 \hyperrefext[IX][IX.prop2.8]{2.8} 得出。

现在采用与 \hyperrefext[XII][XII.sec8]{8} 类似的论证：

\begin{enonce*}{引理 1.7}\etiquette[1.7]{XIV.lem1.7}
设 $S$ 是 Noether 概形，
$$
\xymatrix@=1.25cm{
X        \ar[rd]_{f} &           &\bar{X} \ar[ll]_{\pi}
 \ar[ld]^{\bar{f}}\\
                 &    S      &
}
$$
是固有态射的交换图，$j:U\rightarrow X$ 是一个开子概形，使得开子概形 $\bar U=\fprod{U}{\bar X}{X}$ 同构地映到 $U$。设 $Y=X-U$ 是约化闭子概形，$f_0:Y\rightarrow S$ 和 $i:Y\rightarrow X$ 是典范态射。若\hyperref[XIV.thm1.1]{定理 1.1} 对 $\bar f$、$f_0$ 和 $\pi$ 成立\nde{关于 $R^q\pi_*$ 的假设仅在 $q=0$ 时用于证明。在这一情形，可构造性已由 \hyperref[XIV.lem1.5]{1.5} 得知。}，则它对 $f$ 也成立。
\end{enonce*}

\noindent
\emph{证明}。\pageoriginale 先考虑 $X$ 上的 $A$-模层 $F$。有正合列
$$
0\longrightarrow j_!j^*F\longrightarrow F\longrightarrow
i_*i^*F\longrightarrow0{\quoi}.
$$
由相应的 $R^qf_*$ 的上同调正合列以及 \hyperrefext[IX][IX.prop2.6]{2.6}，只需验证两端项的可构造性。注意，由 \hyperrefext[IX][IX.prop2.4]{2.4} (iii) 和 \hyperrefext[IX][IX.prop2.14]{2.14} (i)，$i_*i^*F$ 可构造，因而 $j_!j^*F$ 也可构造（\hyperrefext[IX][IX.prop2.6]{2.6}）。由于 $i$ 是闭浸入，由 \hyperrefext[VIII][VIII.cor5.6]{5.6}，
$$
(R^qf_*)i_*(i^*F)\simeq(R^qf_{0*})(i^*F){\quoi},
$$
因而这一情形的可构造性由对 $f_0$ 的假设得出。此外，若 $\bar j:\bar U\rightarrow\bar X$ 表示包含态射，且 $\bar F=\pi^*F$，则
\begin{equation*}
j_!j^*F\simeq\pi_*\bar j_!\bar j^*\bar F{\quoi},\tag{$*$}
\end{equation*}
为此，先定义一个从左到右的态射，再利用 \hyperrefext[XII][XII.thm5.1]{5.1} 逐纤维证明其可逆即可。在 $Y$ 的每个几何点 $\bar y$ 处，$\bar j_!\bar j^*\bar F$ 在 $\bar X$ 关于 $X$ 的纤维 $\bar X_{\bar y}$ 上的限制为零；因此，由 \hyperrefext[XII][XII.thm5.1]{5.1} (iii)，对于每个 $q\geq0$，$(R^q\pi_*)(\bar j_!\bar j^*\bar F)$ 在 $\bar y$ 处的纤维为零。另一方面，$\pi$ 在 $X-Y$ 上诱导一个同构，所以对于 $q>0$，
$$
(R^q\pi_*)(\bar j_!\bar j^*\bar F)=0
$$
。因此，Leray 谱序列
$$
E^{p,q}_2=(R^pf_*)(R^q\pi_*)(\bar j_!\bar j^*\bar F)
\Longrightarrow R^n\bar f_*(\bar j_!\bar j^*\bar F)
$$
与 $(*)$ 一起给出同构
$$
(R^q\bar f_*)(\bar j_!\bar j^*\bar F)\simeq
(R^qf_*)(j_!j^*F)
$$
，对所有 $q\geq0$ 成立。由于 $\bar j_!\bar j^*\bar F$ 可构造（与 $j_!j^*F$ 的论证相同），由对 $\bar f$ 的假设，左端可构造。于是 $(R^qf_*)(j_!j^*F)$ 可构造，因而 $R^qf_*F$ 也可构造。\pageoriginale

现在设 $F$ 是 $X$ 上的可构造群层。用 $\pi\amalg i:(\bar X\amalg Y)\rightarrow X$ 代替 $\pi:\bar X\rightarrow X$。于是只需证明：若 $\pi:\bar X\rightarrow X$ 是\emph{满射}，且 \hyperref[XIV.thm1.1]{1.1} 对 $\pi$ 和 $\bar f$ 成立，则它对 $f$ 也成立。先设 $\bar F$ 是 $\bar X$ 上的可构造群层。那么 $\pi_*\bar F$ 在 $X$ 上可构造（\hyperref[XIV.lem1.5]{1.5}），且有单射
$$
(R^1f_*)\pi_*\bar F\hookrightarrow(R^1\bar f_*)\bar F.
$$
由假设右端可构造，因而左端也可构造（\hyperrefext[IX][IX.prop2.9]{2.9} (ii)）。由于 $\pi$ 满射，有单射 $F\rightarrow\pi_*\pi^*F=G$。层 $\bar F=\pi^*F$ 可构造，所以由刚才的证明，$G$ 和 $R^1f_*G$ 也可构造。这样就归结为下面的引理。

\begin{enonce*}{引理 1.8}\etiquette[1.8]{XIV.lem1.8}
设 $S$ 是 Noether 概形，$f:X\rightarrow S$ 是固有态射，而 $\displaystyle F\mathop{\hookrightarrow}^{u}G$ 是 $X$ 上可构造群层的单射。若 $R^1f_*G$ 可构造，则 $R^1f_*F$ 也可构造。
\end{enonce*}

\noindent
\emph{证明}。令 $C=G/F$，这是一个可构造的 $G$-齐性空间层（应用 \hyperrefext[IX][IX.prop2.6]{2.6}）。考虑正合列
\begin{equation*}
f_*C\xrightarrow{\delta}R^1f_*F\xrightarrow{u^1}R^1f_*G{\quoi}.
\tag{$*$}
\end{equation*}
为证明 $R^1f_*F$ 可构造，根据 Noether 归纳法和 \hyperref[XIV.lem1.6]{1.6}，只需证明：若 $S\neq\varnothing$，则 $S$ 有一个非空开子集，使 $R^1f_*F$ 在其上可构造。因而可设 $R^1f_*G$ 局部常值。由于可构造性关于平展（étale）拓扑是局部性质，还可进一步假设 $S$ 不可约，并且 $R^1f_*G$ 是取值为 $E=\{e_1,\ldots,e_n\}$ 的常值（且可构造）层。对于 $i=1,\ldots,n$，设 $D_i$ 是 $R^1f_*F$ 的子层\pageoriginale，它是截面 $e_i$ 在 $u^1$ 下的逆像。由于拓扑斯（topos）中“余积是泛的”，有 $R^1f_*F=\amalg_iD_i$；因而只需证明每个 $D_i$ 在某个适当的非空 $S$-平展概形 $S'$ 上可构造。

回忆 $R^1f_*F$ 是预层 $\Rr^1F$ 的伴随层，其中 $\Rr^1F(S')=H^1(X\times_SS',F)$。因此，可以假设（用非空的 $S$-平展概形 $S'$ 代替 $S$ 之后），$D_i$ 是“空层”（从而可构造），\emph{或}存在一个元素 $\alpha\in H^{1}(X,F)$，诱导 $D_i$ 的截面 $\sisi{\alpha}{\bar{\alpha}}$。在后一情形，设
$$
f_{\ast}C^{\alpha}\xrightarrow{\delta}R^{1}f_{\ast}F^{\alpha}\longrightarrow
R^{1}f_{\ast}G^{\alpha}
$$
是对 $(\ast)$ 用 $\alpha$ 作扭曲所得到的正合列。有
$$
R^{1}f_{\ast}F^{\alpha}\simeq R^{1}f_{\ast}F\quav
R^{1}f_{\ast}G^{\alpha}\simeq R^{1}f_{\ast}G{\quoi},
$$
并且 $R^{1}f_{\ast}F^{\alpha}$ 中与 $\bar{\alpha}$ 对应的截面是单位截面。因此，$R^{1}f_{\ast}F^{\alpha}$ 中与 $D_i$ 对应的子层 $D_i^{\alpha}$ 是 $R^{1}f_{\ast}G^{\alpha}$ 的单位截面的逆像；即 $D_i^{\alpha}$ 是 $f_{\ast}C^{\alpha}$ 在 $\delta$ 下的像。由于 $f_{\ast}C^{\alpha}$ 可构造（\hyperref[XIV.lem1.5]{1.5}），$D_i^{\alpha}$ 也可构造，因而 $D_i$ 也可构造（应用 \hyperrefext[IX][IX.prop2.6]{2.6}）。这就证明了引理。

\begin{enonce*}{引理 1.9}\etiquette[1.9]{XIV.lem1.9}
设 $S$ 是 Noether 概形，
$$
\xymatrix@=1.25cm{
X       \ar[rd]_{f}      &      &\bar{X} \ar[ll]_{\pi} \ar[ld]^{\bar{f}}\\
              & S    &
}
$$
是固有态射的交换图。若定理 \hyperref[XIV.thm1.1]{1.1} 对\pageoriginale $f$ 和 $\pi$ 成立，则它对 $\bar{f}$ 也成立。
\end{enonce*}

\noindent
\emph{证明}。对于可构造 $A$-模层 $\bar{F}$，使用 Leray 谱序列
$$
E^{p,q}_2=(R^{p}f_{\ast})(R^{q}\pi_{\ast})\bar{F}\Longrightarrow
(R^{n}\bar{f}_{\ast})\bar{F}{\quoi}.
$$
由假设，对每个 $p,q$，$E^{p,q}_2$ 都可构造；于是由 \hyperrefext[IX][IX.prop2.6]{2.6}，收敛项也可构造。

设 $\bar{F}$ 是可构造群层。有正合列（\hyperrefext[XII][XII.prop3.2]{3.2}）
\begin{equation*}
0\longrightarrow (R^{1}f_{\ast})\pi_{\ast}\bar{F}\longrightarrow
(R^{1}\bar{f}_{\ast})\bar{F}\longrightarrow
f_{\ast}(R^{1}\pi_{\ast})\bar{F}{\quoi},\tag{$\ast$}
\end{equation*}
由假设和 \hyperref[XIV.lem1.5]{1.5}，两端项可构造。仿照 \hyperref[XIV.lem1.8]{1.8} 的证明：由 Noether 归纳法和 \hyperref[XIV.lem1.6]{1.6}，只需假设 $S\neq ∅$，并证明在某个适当的非空开子集上的可构造性。因此可设 $f_{\ast}(R^{1}\pi_{\ast})\bar{F}$ 局部常值，并且在作平展局部化之后，甚至是取值为 $E=\{e_1,\ldots,e_n\}$ 的常值可构造层（\sisi{\hyperrefext[IX][IX.prop2.8]{2.8}}{\hyperrefext[IX][IX.prop2.4]{2.4}~(i)}）。设 $D_i$ 是 $(R^{1}\bar{f}_{\ast})\bar{F}$ 中 $e_i$ 的逆像子层，于是 $(R^{1}\bar{f}_{\ast})\bar{F}=\coprod_iD_i$。只需证明每个层 $D_i$ 在某个适当的非空连通平展概形 $S'/S$ 上诱导可构造层。若 $D_i$ 不是“空层”，则可归结为存在元素 $\alpha\in H^{1}(\bar{X},\bar{F})$ 诱导 $D_i$ 的截面 $\bar{\alpha}$ 的情形。设
$$
\sisi{
0\longrightarrow (R^{1}f_{\ast})\pi_{\ast}\bar{F}\longrightarrow
(R^{1}\bar{f}_{\ast})\bar{F}\longrightarrow
f_{\ast}(R^{1}\pi_{\ast})\bar{F}
}
{0\longrightarrow (R^{1}f_{\ast})\pi_{\ast}\bar{F}^{\alpha}\longrightarrow
(R^{1}\bar{f}_{\ast})\bar{F}^{\alpha}\longrightarrow
f_{\ast}(R^{1}\pi_{\ast})\bar{F}^{\alpha}}
$$
是对 $(\ast)$ 中的 $\bar{F}$ 用 $\alpha$ 作扭曲得到的正合列。$(R^{1}\bar{f}_{\ast})\bar{F}^{\alpha}$ 中与 $\bar{\alpha}$ 对应的截面是单位截面。因此，\hbox{$(R^{1}\bar{f}_{\ast})\bar{F}^{\alpha}$} 中与 $D_i$ 对应的子层 $D_i^{\alpha}$ 是 $f_{\ast}(R^{1}\pi_{\ast})\bar{F}^{\alpha}$ 的单位截面的逆像，\pageoriginale 从而 $D_i^{\alpha}\simeq (R^{1}f_{\ast})\pi_{\ast}\bar{F}^{\alpha}$，这确实是一个可构造层，\cqfd

\begin{enonce*}{引理 1.10}\etiquette[1.10]{XIV.lem1.10}
只需在 $S$ 是 $\Spec\ZZ$ 上的有限型概形，且 $f$ 是相对维数 $\leq 1$ 的射影态射时验证 \hyperref[XIV.thm1.1]{1.1}。
\end{enonce*}

\noindent
\emph{证明}。先设 $f$ 是射影态射，并对 $f$ 的相对维数作归纳；用 $S$ 的一个仿射开子集代替它后，可设这一相对维数有限。若 $f$ 的相对维数 $\leq n$，其中 $n\geq 2$，那么局部地在 $S$ 上可以找到一个有限态射 $\pi:X\rightarrow\PP^n_S$。由 \hyperref[XIV.lem1.9]{1.9} 和 \sisi{IX 2.15}{\hyperrefext[IX][IX.prop2.14]{2.14}} (i)，归结为 $f$ 是 $\PP^n_S$ 的结构态射的情形。考虑态射的交换图
$$
\xymatrix@C=2cm{
                   & \bar{X}\ar[ld]_{\sisi{}{\pi}} \ar[rd]^{a} \ar@{-}[dd]^{\bar{f}} &       \\
\PP^{n}_S \ar[rd]_{f}  &                  &\PP^{1}_S \ar[ld]^{b}\\
                   &  S                          &
}
$$
，它在 \hyperrefext[XII][XII.lem8.3]{8.3} 的证明中已经考虑过，其中 $\bar{X}$ 由 $\PP^n_S$ 沿闭子概形 $Y\simeq\PP^{n-2}_S$ 作爆破得到。$b$ 的相对维数 $\leq1$，$a$ 的相对维数 $\leq n-1$；所以由归纳假设，定理对 $a$ 和 $b$ 成立，因而由 \hyperref[XIV.lem1.9]{1.9}，也对 $\bar{f}$ 成立。此外，$\pi$ 的相对维数 $\leq1$，而 $f_0:Y\rightarrow S$ 的相对维数 $\leq n-1$。因此由 \hyperref[XIV.lem1.7]{1.7}，\hyperref[XIV.thm1.1]{1.1} 对 $f$ 成立，从而对每个射影态射都成立。现在令 $f$ 为任意固有态射，对 $X$ 作 Noether 归纳：可设定理对 $X$ 的每个真闭子概形都成立，并且 $X\neq∅$。设
$$
\xymatrix@=1.5cm{
X  \ar[rd]_{f}          &           &\bar{X} \ar[ll]_{\pi} \ar[ld]^{\bar{f}}\\
                        &    S       &
}
$$
是 Chow 引理 \hyperrefext[XII][XII.lem7.1]{7.1} 给出的那类\pageoriginale 图。对这个图应用 \hyperref[XIV.lem1.7]{1.7}，即得 $f$ 满足 \hyperref[XIV.thm1.1]{1.1}，引理得证。

\setcounter{subsection}{10}
\subsection{}\etiquette[1.11]{XIV.subsec1.11}
现在可以完成 \hyperref[XIV.thm1.1]{1.1} 的证明。若 $F$ 是可构造群层，可以找到一个单射 $F\hookrightarrow G=\prod_i\pi_{i\sisi{}{,}\ast}C_i$，其中每个 $\pi_i$ 有限，每个 $C_i$ 常值且可构造（\sisi{IX 2.15}{\hyperrefext[IX][IX.prop2.14]{2.14}}）；因此由 \hyperref[XIV.lem1.8]{1.8}，可设 $F$ 常值。同样，若 $F$ 是可构造 $A$-模层，则应用 \sisi{IX 2.15}{\hyperrefext[IX][IX.prop2.14]{2.14}} 可以找到 $F$ 的一个消解 $F\rightarrow G^{\bigdot}$，$G^{\bigdot}=\{0\rightarrow G^0\rightarrow G^1\rightarrow G^2\rightarrow\ldots\}$，使每个 $G^{\nu}$ 具有形式 $G^{\nu}=\prod_i\sisi{\eta}{\pi}_{i\sisi{}{,}\ast}C_i$。消解 $G^{\bigdot}$ 的谱序列
$$
E^{p,q}_2=H^{\sisi{q}{p}}((R^{q}f_{\ast})G^{\bigdot})\Longrightarrow
(R^{n}f_{\ast})F{\quoi},
$$
便将问题归结为 $F$ 是取值为 $M$ 的常值可构造层的情形。

回忆 $A$ 是一个 $\ZZ/n$-代数。设 $\ell$ 是 $n$ 的一个素因子。有正合列
$$
0\longrightarrow \ell M\longrightarrow M\longrightarrow M/\ell
M\longrightarrow 0{\quoi},
$$
而仍然只需对 $\ell M$ 和 $M/\ell M$ 证明 \hyperref[XIV.thm1.1]{1.1}。

一个简单的归纳便将问题归结为 $A$ 是 $\ZZ/\ell$-代数，$\ell\in\PP$，且 $M$ 是有限型 $A$-模的情形。于是 $M$ 是 $\ZZ/\ell$ 上的向量空间，而且对每个拟紧拟分离的 $Y$，典范态射
$$
H^q(Y,\ZZ/\ell)\otimes_{\ZZ/\ell}M\longrightarrow H^q(Y,M)
$$
都是双射（事实上，若 $M$ 的秩有限，这是显然的，而两端都与滤过正向极限交换\pageoriginale （\sisi{VI 1.6}{\hyperrefext[VI][VI.thm5.1]{5.1} 和 \hyperrefext[VII][VII.prop3.3]{3.3}}））。因此
$$
[R^qf_{\ast}(\ZZ/\ell)]\otimes_{\ZZ/\ell}M\xrightarrow{\sim}
R^qf_{\ast}(M){\quoi}.
$$
所以显然只需证明 $R^qf_{\ast}(\ZZ/\ell)$ 作为 $\ZZ/\ell$-模是可构造的；换言之，我们归结为 $A=\ZZ/\ell$ 的情形，以下均作此假设。

注意，对于任意 $F$，$f$ 的相对维数 $\leq0$（即 $f$ 有限）的情形，现在由 \hyperrefext[VIII][VIII.cor5.6]{5.6} 和 \hyperref[XIV.lem1.5]{1.5} 得出。

对 $S$ 作 Noether 归纳。由 \hyperref[XIV.lem1.6]{1.6}，可设 $S\neq∅$，并可用任意非空开子集代替 $S$。此外，可设 $S$ 整，泛点为 $s$。纤维 $X_s$ 是维数 $\leq1$ 的代数概形；因此存在 $k(s)=K$ 的纯不可分扩张 $K'$，使 $(X_s\otimes_KK')_{\text{\rm red}}$ 的正规化 $Y$ 在 $\Spec K'$ 上光滑\footnote[1]{\Cf \EGA IV 17.15.14.}。用非空开子集代替 $S$ 之后，可设存在有限、根性、满射的态射 $S'\rightarrow S$，使 $S'_s$ 与 $\Spec K'$ 为 $K$-同构；还存在光滑射影态射 $\bar{X}'\rightarrow S'$ 及有限满射的 $S'$-态射 $\pi:\bar{X}'\rightarrow X'=\fprod{X}{S'}{S}$，它在 $S$ 的泛点 $s$ 上的纤维上诱导 $Y\rightarrow X_s\otimes_KK'$（\EGA IV 9.6.1）。现在\sisi{}{，}由 \hyperrefext[VIII][VIII.thm1.1]{1.1}\sisi{}{，}作纯不可分扩张不影响问题。因此可用 $S'$ 代替 $S$，用 $X'$ 代替 $X$；也就是说，可设 $f$ 出现在图
$$
\xymatrix@=1.5cm{
X   \ar[rd]_{f}& &\bar{X} \ar[ll]_{\pi} \ar[ld]^{\bar{f}}\\
         &         S        &           ,
}
$$
中，其中\pageoriginale $\pi$ 有限且满射，$\bar f$ 光滑。此外，按构造，可设 $X$ 有一个闭子概形 $Y$，其泛纤维有限，且 $\pi$ 在 $X-Y$ 上诱导同构。再用一个非空开子集代替 $S$，便归结为 $f_0:Y\rightarrow S$ 也有限的情形。

由于 $F$ 常值，$\pi^{\ast}F$ 也常值，并且因 $\pi$ 满射，$F\rightarrow\pi_{\ast}\pi^{\ast}F$ 是单射。若 $F$ 是群层，只需证明 $(R^1f_{\ast})(\pi_{\ast}\pi^{\ast}F)\simeq (R^1\bar f_{\ast})(\pi^{\ast}F)$ 可构造（\hyperref[XIV.lem1.8]{1.8}）。若 $F$ 是 $\ZZ/\ell$-模层，有正合列
$$
0\longrightarrow F\longrightarrow\pi_{\ast}\pi^{\ast}F\longrightarrow
C\longrightarrow0{\quoi},
$$
其中 $C$ 的支集在 $Y$ 中，因为 $\pi$ 在 $X-Y$ 上诱导同构。由此，$R^qf_{\ast}C=R^qf_{0\sisi{}{,}\ast}C$ 可构造（且当 $q>0$ 时为零），因此仍然只需证明 $(R^qf_{\ast})(\pi_{\ast}\pi^{\ast}F)= (R^q\bar f_{\ast})(\pi^{\ast}F)$ 可构造。用 $\bar f$ 代替 $f$，用 $\pi^{\ast}F$ 代替 $F$，便归结为 $F$ 常值且可构造、$f$ 光滑的情形。

先处理常值群层 $\sisi{f}{F}=G_X$。$(R^1f_{\ast})G_X$ 的所有纤维都有限。事实上，为证明这一点，应用 \hyperrefext[XII][XII.cor5.2]{5.2} (ii)，可归结为 $S$ 是可分闭域的谱，且 $X$ 在 $S$ 上光滑、维数 $\leq1$ 的情形。在这一情形结论已知（\SGA 1 X 2.6）。因此只需证明特殊化态射是单射（\hyperrefext[IX][IX.prop2.13]{2.13} (ii)）。为此，归结为 $S$ 严格局部的情形；由基变换定理 \hyperrefext[XII][XII.thm5.1]{5.1} (ii)，甚至可用严格局部离散赋值环的谱代替 $S$。此外，可设 $X$ 连通。由于\pageoriginale $H^1(X,G)$ 分类以 $G$ 为群的 $X$ 的主 Galois 覆盖，只需验证：若 $X'\rightarrow X$ 是连通平展覆盖，则其几何纤维也连通。这很容易由 $f$ 光滑且固有、并具有几何连通纤维这一事实得出。

现在处理 $F=(\ZZ/\ell)_X$ 的情形。当 $q=0$ 和 $q=1$ 时，结论已知（\hyperref[XIV.lem1.5]{1.5} 及前面的论证）。\sisi{由 \hyperrefext[X][X.cor4.3]{4.3} 和}{由} \hyperrefext[XII][XII.corbis5.3]{5.3 bis}，只剩下 $q=2$ 的情形；而且，若 $S$ 的特征为 $\ell$，则 $R^2f_{\ast}(\ZZ/\ell)=0$。因此，用 $S$ 的一个适当非空开子集代替它，可归结为 $\ell$ 在 $S$ 上可逆的情形。于是，由 \hyperrefext[XII][XII.cor5.2]{5.2} \sisi{(ii)}{(iii)} 和 \sisi{X 4.7}{\hyperrefext[X][X.cor5.2]{5.2}}，$R^2f_{\ast}(\ZZ/\ell)$ 在 $s\in S$ 上方几何点 $\bar s$ 处的纤维同构于 $(\ZZ/\ell)^{c(s)}$，其中 $c(s)$ 为几何纤维 $X_{\bar s}$ 的一维不可约分支数目。函数 $c$ 在 $S$ 上可构造（\EGA IV 9.7.9），结论由此得出（\hyperrefext[IX][IX.prop2.13]{2.13} (iii)）。

\begin{enonce*}[remark]{注 1.11}\etiquette[1.11]{XIV.rem1.11}
若使用具有固有支集的上同调的形式体系，\hyperref[XIV.thm1.1]{1.1} 的证明会简单得多；这一形式体系将在 \hyperref[sga4.exp.XVII]{XVII} 中作为固有态射基变换定理 \hyperrefext[XII][XII.thm5.1]{5.1} 的直接推论加以建立。此外，\hyperref[XIV.thm1.1]{1.1} 的陈述可推广到如下情形：在 $\sisi{Y}{S}$ 上给定任意挠环层 $𝒜$，以及 $X$ 上的一个 $f^{\ast}(𝒜)$-模复形 $K^{\bigdot}$，满足某种“可构造性”条件；在 \hyperref[XIV.thm1.1]{1.1} 的特殊情形中，该条件仅表示 $K^{\bigdot}$ 的上同调层 $\Hh^i(K^{\bigdot})$ 是可构造 $A$-模，并在 $i$ 充分大时为零。
\end{enonce*}

\section{维数的一种变式}\etiquette[2]{XIV.sec2}\pageoriginale

\subsection{}\etiquette[2.1]{XIV.subsec2.1}
设 $X$ 是概形。对于 $x\in X$，记
\begin{equation*}
d(x)=\dim\bar{\{x\}}\tag{2.1.1}
\end{equation*}\etiquette[2.1.1]{XIV.eq2.1.1}
，其中 $\bar{\{x\}}$ 是 $\{x\}$ 的闭包。若 $F$ 是 $X$ 上的阿贝尔层，令
\begin{equation*}
d(F)=\sup \{d(x)\mid x\in X \quaet F_{\bar{x}}\neq 0\}{\quoi}.\tag{2.1.2}
\end{equation*}\etiquette[2.1.2]{XIV.eq2.1.2}
在局部情形中，需要对这一概念作如下修正：

\begin{enonce*}{定义 2.2}\etiquette[2.2]{XIV.def2.2}
设 $Y$ 是泛链状（universally catenary）Noether 局部环的谱（\EGA IV 5.6.2），$f:X\rightarrow Y$ 是有限型态射。设 $x\in X$，$y=f(x)$。令
$$
\delta(x)=\dim\bar{\{y\}}+\degtr k(x)/k(y){\quoi}.
$$
若 $F$ 是 $X$ 上的阿贝尔层，令
$$
\delta(F)=\sup \{\delta(x)\mid x\in X\quaet F_{\bar{x}}\neq 0\}{\quoi}.
$$
\end{enonce*}

\begin{enonce*}{命题 2.3}\etiquette[2.3]{XIV.prop2.3}
\begin{enumerate}
%%%
%%% The typesetting of this enumerate environment is terrible.
%%% The spacing between the dash and (i) is too big.
%%% Also, the items (ii), ... should better be without indent.
%%% If you use the package enum.sty, then you can build the necessary list environment.
%%%
\renewcommand{\theenumi}{\roman{enumi}}
\renewcommand{\labelenumi}{(\theenumi)}
\item 设 $i:X'\rightarrow X$ 是浸入，$x'\in X'$。那么 $\delta(i(x'))=\delta(x')$。
%%% new line break within formula in text: spend extra line to induce new line breaks?

\item 设 $X_{0}$ 是 $X/Y$ 的闭纤维。有
$$
\delta(x)=
\begin{cases}
d(x) & \mbox{若 } \bar{\{x\}}\cap X_{0}\neq ∅\\
d(x)+1 & \mbox{若 } \bar{\{x\}}\cap X_{0}=∅{\quoi};
\end{cases}
$$

\item 若 $f$ 固有，则 $\delta(x)=d(x)$。
\end{enumerate}
\end{enonce*}

\noindent
\emph{证明}。\pageoriginale 断言 (i) 直接由定义得出。由于当 $f$ 固有时 $\bar{\{x\}}\cap X_{0}$ 非空，(iii) 由 (ii) 得出。对于 (ii)，可用 $\bar{\{x\}}$ 代替 $X$，用 $\bar{\{y\}}$ 代替 $Y$（赋予它们诱导的约化结构），(i) 保证可以这样做。设 $x'\in X$，$y'=f(x')$。有公式（\EGA IV 5.6.1.1）
$$
\dim \Oo_{X,x'}=\dim
\Oo_{Y,y'}+\degtr k(x)/k(y) - \degtr k(x')/k(y'){\quoi}.
$$

设 $X_{0}\neq ∅$。显然，$\displaystyle{\mathop{\sup}_{x'}}$ $\{\dim \Oo_{X,x'}\}=\dim X$ 在取 $x'$ 为 $X_{0}$ 的一个闭点时达到，因此
$$
\dim X=\dim \Oo_{X,x'}=\dim Y + \degtr k(x)/k(y)=\delta(x){\quoi}.
$$
若 $X_{0}=∅$，显然上确界在取 $x'$ 为其纤维 $f^{-1}(y')$ 中的闭点、且 $\dim\bar{\{y'\}}=1$ 时达到。这样的点存在，因为满足 $\dim\bar{\{y'\}}=1$ 的 $y'\in Y$ 构成一个极稠密集合（\EGA IV 10.1.3，\sisi{10.5.9, 1.8.4}{10.3.1, 10.4.1, 10.5.9}）。于是
\begin{align*}
\dim X=\dim \Oo_{X,x'} &= \dim \Oo_{Y,y'}+\degtr k(x)/k(y)\\
&= \dim Y-1 +\degtr k(x)/k(y)=\delta(x) \sisi{=}{-} 1{\quoi},
\end{align*}
结论得证。

\begin{enonce*}{命题 2.4}\etiquette[2.4]{XIV.prop2.4}
设 $f:X\rightarrow Y$ 是域 $k$ 的谱上有限型概形之间的态射。设 $x\in X$，$y=f(x)$，$y_{0}$ 为 $y$ 的一个特殊化。设 $f':X'\rightarrow Y'$ 是 $f$ 在 $y_{0}$ 上方几何点 $\bar{y}_{0}$ 处的严格局部化，$x'\in X'$ 是 $x$ 上方的一个点。那么
$$
\delta(x')=d(x)-d(y_{0}){\quoi}.
$$
\end{enonce*}

\noindent
\emph{证明}。\pageoriginale 有
\begin{align*}
d(x) &= d(y)+\degtr k(x)/k(y)\\
\delta(x') &= d(y')+\degtr k(x')/k(y'){\quoi},
\end{align*}
其中 $y'=f'(x')$。此外，
$$
\degtr k(x)/k(y)=\degtr k(x')/k(y'){\quoi},
$$
因为扩张 $k(x')/k(x)$ 和 $k(y')/k(y)$ 都是代数扩张。因此只需证明
$$
d(y')=d(y)-d(y_{0}){\quoi}.
$$

可设 $y$ 是 $Y$ 中的极大点，因而 $y'$ 是 $Y'$ 中的极大点。于是
\begin{align*}
d(y') &= \dim \Oo_{Y',\bar{y}_{0}}\\
      &= \dim \Oo_{Y,y_{0}}\quad \text{\rm (\EGA IV 6.1.3)}\\
      &= \sisi{\dim Y}{\dim_{y_0} Y} -d(y_{0})\quad \text{\rm (\EGA IV 5.2.3.1)}\\
      &= d(y)-d(y_{0}){\quoi},
\end{align*}
结论得证。

\section{仿射代数概形的上同调维数}\etiquette[3]{XIV.sec3}

\begin{enonce*}{定理 3.1~(d)}\etiquette[3.1]{XIV.thm3.1}\nde{定理 \ref{XIV.thm3.1} 现在称为\textit{仿射 Lefschetz 定理}。}
设 $f:X\rightarrow Y$ 是域 $k$ 上有限型概形之间的仿射态射。设 $F$ 是 $X$ 上的挠层，满足 $d(F)\leq d$ \emph{（\cf \hyperref[XIV.subsec2.1]{2.1}）}。那么 $d(R^{q}f_{\ast}F)\leq d-q$。
\end{enonce*}

特别地，取 $Y=\Spec(k)$，得到：

\begin{enonce*}{推论 3.2}\etiquette[3.2]{XIV.cor3.2}
设\pageoriginale $X$ 是可分闭域 $k$ 上的有限型仿射概形。那么\nde{考虑复解析几何中的类似断言：设 $X$ 是 $\mathbf{C}$ 上的仿射簇，$F$ 是 $X$ 上的可构造层。那么对于每个 $m>\dim X$，有 $H^m(X^{\rm an},F)=0$，其中 $X^{\rm an}$ 是与 $X$ 相伴的解析空间。虽然这一断言由 \hyperref[XIV.cor3.2]{3.2} 及比较定理 \hyperrefext[XVI][XVI.thm4.1]{4.1} 得出，但还有两个完全解析的证明，分别由 Pierre Deligne 和 Hélène Esnault 给出；见 Esnault, H., Variation on Artin's vanishing theorem, \textit{Advances in Math.} \textbf{198} (2005), 435--438。}
$$
\cd X\leq \dim X{\quoi}.
$$
\end{enonce*}

\begin{enonce*}{推论 3.3}\etiquette[3.3]{XIV.cor3.3}
设 $X\subset \PP^N_k$ 是可分闭域 $k$ 上维数为 $n$ 的射影概形，$Y=H\cap X$ 是 $X$ 的一个“超平面截面”。那么，对于 $X$ 上的每个挠层 $F$，典范同态
$$
H^{q}_Y(X,F)\longrightarrow H^{q}(X,F)
$$
当 $q>n+1$ 时为双射，当 $q=n+1$ 时为满射。
\end{enonce*}

这一推论直接由 \sisi{\hyperref[XIV.def2.2]{2.2}}{\hyperref[XIV.cor3.2]{3.2}} 及闭子集的正合列（\sisi{IV 3.7.1}{\hyperrefext[V][V.6.5.4]{6.5.4}}）得出，
%%% This last reference should read \ref{V.eq6.5.4}, but the TeX-file of chapter V contains the
%%% label without "eq".
因为 $X-Y=\sisi{\mathcal{U}}{U}$ 是维数 $\leq n$ 的仿射概形。注意，\hyperref[XIV.cor3.3]{3.3} 推广了 Lefschetz 关于超平面截面的一个定理。事实上，若 $X$ 在 $k$ 上光滑，$F$ 局部常值且其阶与特征互素，则利用相对上同调纯性定理（\hyperrefext[XVI][XVI.sec3]{3}）及谱序列（\sisi{\quad}{\hyperrefext[V][V.6.4.3]{6.4.3}}），可以显式计算带支集上同调 $H^{q}_Y(X,F)$，
%%% This last reference should read \ref{V.eq6.4.3}, but the TeX-file of chapter V contains the
%%% label without "eq".
得到
$$
H^{q}_Y(X,\ZZ/n)\simeq H^{q-2}(Y,\mu_n^{-1}) \mbox{\rm\quad 对所有 $q$；}
$$
由 \hyperref[XIV.cor3.3]{3.3} 给出的同态 $H^{q-2}(Y,\mu_n^{-1})\rightarrow H^{q}(X,\ZZ/n)$ 正是“Gysin 同态”，我们将在后面研究它\footnote[1]{关于 Lefschetz 型定理的系统研究，\Cf \SGA 2 XIV。}。

定理 \hyperref[XIV.thm3.1]{3.1}~(d) 等价\nde{更确切地说，正文中的论证给出以下结果：对于 $d'\leq d$ 的断言 \hyperref[XIV.cor3.4]{3.4}~(d') 蕴含 \hyperref[XIV.thm3.1]{3.1}~(d)，而断言 \hyperref[XIV.thm3.1]{3.1}~(d) 蕴含 \hyperref[XIV.cor3.4]{3.4}~(d)。这足以得到下面的结论。}于下面的陈述：

\begin{enonce*}{推论 3.4~(d)}\etiquette[3.4]{XIV.cor3.4}
设 $Y'$ 是严格局部概形，且是某个域上代数概形的严格局部化。设 $f':X'\rightarrow Y'$ 是有限型仿射态射\pageoriginale，$F'$ 是 $X'$ 上满足 $\delta(F')\leq d$ 的挠层 \emph{（\cf （\hyperref[XIV.def2.2]{2.2}））}。那么当 $q>d$ 时，$H^{q}(X',F')=0$。
\end{enonce*}

事实上，设 \hyperref[XIV.cor3.4]{3.4}~(d) 成立。设 $f:X\rightarrow Y$ 是 $k$ 上代数概形之间的仿射态射，$F$ 是 $X$ 上满足 $d(F)\leq d$ 的挠层。设 $y$ 是 $Y$ 的一个点，$\xi$ 是 $y$ 上方的几何点。为证明 \sisi{\ref{XIV.subsec3.1}}{\hyperref[XIV.thm3.1]{3.1}}~(d)，只需证明当 $d(y)>d-q$ 时，茎 $(R^{q}f_{\ast}F)_{\xi}$ 为零。

设 $f':X'\rightarrow Y'$ 是 $f$ 在点 $\xi$ 处的严格局部化；即 $Y'$ 为 $Y$ 在 $\xi$ 处的严格局部化，$X'=\fprod{X}{Y'}{Y}$。

设 $F'$ 是 $F$ 在 $X'$ 上诱导的层，于是 $(R^{q}f_{\ast}F)_{\xi}=H^{q}(X',F')$（\hyperrefext[VIII][VIII.thm5.2]{5.2}）。
\sisi{有
$$
\delta(F')\leq d(F)-d(y)\leq d-d(y)\quad (\hyperref[XIV.prop2.4]{2.4}){\quoi},
$$}
{由（\hyperref[XIV.prop2.4]{2.4}），有
$$
\delta(F')\leq d(F)-d(y)\leq d-d(y){\quoi},
$$}
%%% I find it irritating that the reference to the proposition from which the above follows,
%%% is referred to in a style that might be confused with an equation number
%%% (although the correct equation number is 3.4.1)
并且应用 \hyperref[XIV.cor3.4]{3.4}~(d)，得到当 $q>d-d(y)$ 时 $H^{q}(X',F')=0$，这就证明了结论。

反过来，设 \sisi{\hyperref[XIV.subsec2.1]{2.1}}{\hyperref[XIV.thm3.1]{3.1}}~(d) 成立，$F'$ 和 $f':X'\rightarrow Y'$ 如 \hyperref[XIV.cor3.4]{3.4}~(d) 的陈述所述。由于上同调与滤过正向极限交换（\hyperrefext[VII][VII.prop3.3]{3.3}），且 $F'$ 是其可构造子层的正向极限（\hyperrefext[IX][IX.prop2.9]{2.9}），可设 $F'$ 可构造。现在 $Y'$ 是代数概形 $Y$ 在\emph{闭}点 $y_0$ 上方几何点 $\xi$ 处的严格局部化（\hyperrefext[X][X.lem3.3]{3.3}）。写成 $Y'=\varprojlim_{\alpha}Y_{\alpha}$，其中 $(Y_{\alpha})$ 是 $Y$ 上仿射平展概形的滤过逆向系统（\hyperrefext[VIII][VIII.subsec4.5]{4.5}）。由 \hyperrefext[IX][IX.cor2.7.4]{2.7.4}，当 $\alpha$ 充分大时，$F'$ 由 $\sisi{X=\fprod{X}{Y_K}{Y}}{X_\alpha=\fprod{X}{Y_\alpha}{Y}}$ 上的某个可构造挠层 \sisi{}{$F_\alpha$} 诱导。利用 \hyperref[XIV.prop2.4]{2.4} 和\pageoriginale \EGA IV 9.5.5，立即可见还可设 $d(F_{\alpha})=\delta(F')\leq d$。用 $Y_{\alpha}$ 代替 $Y$，归结为 $F'$ 由 $X$ 上满足 $d(F)\leq d$ 的层 $F$ 诱导的情形。于是，由 \sisi{\hyperref[XIV.subsec2.1]{2.1}}{\hyperref[XIV.thm3.1]{3.1}}~(d)，当 $q>\sisi{D}{d}$ 时 $R^{q}f_{\ast}F=0$，因而特别地
$$
(R^{q}f_{\ast}F)_{\xi}=H^{q}(X',F')=0\quad{\mbox{\rm 若}}\quad
q>d{\quoi},
$$
这就证明了 \hyperref[XIV.cor3.4]{3.4}~(d)。

作为 \hyperref[XIV.cor3.4]{3.4}~(d) 的特例，有：

\begin{enonce*}{推论 3.5}\etiquette[3.5]{XIV.cor3.5}
设 $Y'$ 是维数 $\leq d$ 的严格局部概形，且是域 $k$ 上某个代数概形的严格局部化。设 $U\subset Y'$ 为仿射开子集。那么 $\cd U\leq d$。
\end{enonce*}

\begin{enonce*}[remark]{注 3.6}\etiquette[3.6]{XIV.rem3.6}
很有可能\nde{Ofer Gabber 已成功证明 \ref{XIV.thm3.1} 的以下类似结果：设 $f:X\to Y$ 是有限型仿射态射，其中 $Y$ 拟优（\EGA IV 7.8），并赋有维数函数 $\delta_Y$。设 $\delta_X$ 为 $X$ 上相应的修正维数函数。若 $n$ 在 $Y$ 上可逆，那么对于 $X$ 上的每个 $\mathbf{Z}/n$-模层 $F$，以及每个 $q\geq0$，有 $\delta_Y(R^q f_\ast F)\leq\delta_X(F)-q$。这里与前面一样，对层 $G$ 定义 $\delta(G)=\sup\{\delta(x)\mid x\in X\quaet G_{\bar{x}}\neq0\}$。
}，\hyperref[XIV.cor3.4]{3.4}，从而 \hyperref[XIV.cor3.5]{3.5}，对于\emph{每个} Noether 局部概形 $Y'$（不必是代数概形的严格局部化）仍然成立，至少当 $Y'$ 是\emph{优的}（\EGA IV 7.8）时如此。在进一步假设 $Y$ 的特征为零时，将在（\sisi{XX 6}{\hyperrefext[XIX][XIX.sec6]{6}}）中证明这一点。
\end{enonce*}

\section[定理 3.1 的证明]{定理 \texorpdfstring{\hyperref[XIV.thm3.1]{3.1}}{3.1} 的证明}\etiquette[4]{XIV.sec4}

证明对 $d$ 作归纳。

\begin{enonce*}{引理 4.1}\etiquette[4.1]{XIV.lem4.1}
\emph{\hyperref[XIV.thm3.1]{3.1}~(0)} 和 \emph{\hyperref[XIV.thm3.1]{3.1}~(1)} 成立。
\end{enonce*}

\noindent
\emph{证明}。当 $d=0$ 时这是显然的，此时 $\supp F$ 在 $k$ 上有限，应用（\hyperrefext[VIII][VIII.prop5.5]{5.5}）即可。设 $F$ 和 $f:X\rightarrow Y$ 如 \hyperref[XIV.thm3.1]{3.1} 的陈述所述，并设 $d(F)\leq1$。用 $F$ 的一个可构造子层代替它（\hyperrefext[IX][IX.prop2.9]{2.9} 和 \hyperrefext[VII][VII.prop3.3]{3.3}），便归结为 $F$ 可构造的情形。于是 $F$ 的支集\pageoriginale 包含于 $X$ 的一个维数 $\leq1$ 的闭子概形中（\hyperrefext[IX][IX.def2.3]{2.3}\sisi{ (ii)}{}），可以用这个子概形代替 $X$；即可以假设 $\dim X\leq1$。由 \hyperrefext[IX][IX.cor5.7]{5.7}，可分闭域上维数为 $1$ 的仿射代数概形的上同调维数 $\leq1$。此外，只需处理 $k$ 可分闭的情形。事实上，设 $\bar f:\bar X\rightarrow\bar Y$ 为由 $f$ 经基变换 $\Spec\bar k\rightarrow\Spec k$ 得到的态射，其中 $\bar k$ 是 $k$ 的可分闭包。对于 $Y$ 的几何点 $\xi$，$f$ 在 $\xi$ 处的严格局部化 $f':X'\rightarrow Y'$ 可视为 $\bar f$ 的一个严格局部化，因而纤维 $(R^{q}f_{\ast}F)_\xi$ 可视为 $R^{q}\bar f_{\ast}\bar F$ 的一个纤维。由于 $f$ 仿射，得到当 $q>1$ 时 $R^{q}f_{\ast}F=0$。显然还有 $d(f_{\ast}F)\leq1$。对于 $q=1$，注意假设 $\dim X\leq1$ 蕴含 $Y$ 有一个有限闭子集 $Z$，使 $f$ 在 $Y-Z$ 上有限。由此，$R^{1}f_{\ast}F$ 在 $Y-Z$ 上为零（\hyperrefext[VIII][VIII.cor5.6]{5.6}），因而 $d(R^{1}f_{\ast}F)\leq0$。引理得证。

设 $d\geq2$，现在假设对于 $d'<d$，\hyperref[XIV.thm3.1]{3.1}~(d')（或等价地，\hyperref[XIV.cor3.4]{3.4}~(d')）成立。

\begin{enonce*}{引理 4.2}\etiquette[4.2]{XIV.lem4.2}
为在一般情形证明 \emph{\hyperref[XIV.cor3.4]{3.4}~(d)}，只需处理 $X'=\EE^1_{Y'}$ 是 $Y'$ 上的仿射直线，且 $f$ 是结构态射的情形。
\end{enonce*}

\noindent
\emph{证明}。显然可设对于某个适当的 $N$，有 $X'=\EE^N_{Y'}$，因为 $X'$ 可嵌入某个 $\EE^N_{Y'}$。设 $N>1$，且对于每个 $Y'$ 和 $r<N$，结论对于 $\EE^r_{Y'}$ 已知。设 $F'$ 是 $\EE^N_{Y'}$ 上的挠层，满足\pageoriginale $\delta(F')\leq d$。考虑图
$$
\xymatrix@=1.5cm{
\EE^{N}_{Y'} \ar[rd]_{f'}       \ar[rr]^{g'}&      &
\EE^{N-1}_{Y'} \ar[ld]^{h'}\\
                       & Y' & &,
}
$$
其中 $g'$ 和 $h'$ 为典范投影，所以 $\EE^N_{Y'}$ 同构于 $\EE^{N-1}_{Y'}$ 上的仿射直线。由于 $\EE^{N-1}_{Y'}$ 的每个严格局部化显然也是某个代数概形的严格局部化，可对 $R^{q}g'_{\ast}F'$ 的纤维应用归纳假设，得到
\begin{equation*}
\delta(R^{q}g'_{\ast}F')\leq d-q{\quoi}.\tag{$\ast$}
\end{equation*}
此外，由关于 $N$ 的归纳假设，\hyperref[XIV.cor3.4]{3.4}~(d') 对 $h'$ 和 $d'\leq d$ 成立。

考虑 Leray 谱序列
$$
E^{p,q}_2=H^{p}(\EE^{N-1}_{Y'},R^{q}g'_{\ast}(F'))
\Longrightarrow H^{n}(\EE^N_{Y'},F'){\quoi}.
$$
对态射 $h'$ 应用 $(\ast)$ 和 \hyperref[XIV.cor3.4]{3.4}~(d')，得到当 $p>d-q$ 时 $E^{p,q}_2=0$。因此收敛项在 $n>d$ 时为零，引理得证。

考虑下面的附加陈述，它是 \hyperref[XIV.cor3.5]{3.5} 的一个特例：

\begin{enonce*}{陈述 4.3~(d)}\etiquette[4.3]{XIV.enon4.3}
设 $X'$ 严格局部，维数 $\leq d$，并且是某个域上代数概形的严格局部化。设 $f\in\Gamma(X',\Oo_{X'})$，$U\subset X'$ 为开子集 $U=X'-V(f)$。那么 $\cd U\leq d$。
\end{enonce*}

\begin{enonce*}{引理 4.4}\etiquette[4.4]{XIV.lem4.4}
\emph{\hyperref[XIV.enon4.3]{4.3}~(d)} 蕴含 \emph{\hyperref[XIV.cor3.4]{3.4}~(d)}。
\end{enonce*}

\noindent
\emph{证明}。\pageoriginale 考虑图
$$
\xymatrix@=1.5cm{
\EE^1_{Y'} \ar[rr]^i \ar[rd]_{f'}&            &
\PP^1_{Y'} \ar[ld]^{g'}\\
                    &       Y'   &
}
$$
其中 $Y'$ 如 \sisi{2.6}{\hyperref[XIV.cor3.4]{3.4}} 的陈述所述，而 $\PP^1_{Y'}$ 是由 $\EE^1$“添上无穷远截面 $Y'_\infty$”得到的射影直线。设 $F'$ 是 $\EE^1$ 上满足 $d(\supp F')\leq d$ 的挠层。有谱序列
$$
H^{p}(\PP^1,R^{q}i_{\ast}F')\Longrightarrow
H^{n}(\EE^1,F'){\quoi}.
$$

当 $q>0$ 时，$R^{q}i_{\ast}F'$ 的支集在 $Y'_\infty$ 中，而它与 $Y'$ 为 $Y'$-同构。由于 $Y'$ 严格局部，得到当 $p,q>0$ 时 $H^{p}(\PP^1,R^{q}i_{\ast}F')=0$。此外，由\emph{固有基变换定理}（\hyperrefext[XII][XII.thm5.1]{5.1}），$\PP^1$ 上挠层的上同调可在闭纤维上计算，而闭纤维是上同调维数 $\leq2$ 的概形（\hyperrefext[X][X.cor4.3]{4.3}）。因此当 $p>2$ 时 $H^{p}(\PP^1,i_{\ast}F')=0$。由于我们要证明当 $n>d$ 时 $H^{n}(\EE^1,F')=0$，且 $d\geq2$，便归结为证明当 $q>d$ 时 $H^{0}(\PP^1,R^{q}i_{\ast}F')=0$。这一群还同构于 $H^{0}(Y'_\infty,R^{\sisi{1}{q}}i_{\ast}F')$，由（\hyperrefext[VIII][VIII.prop4.6]{4.6}），它同构于 $R^{q}i_{\ast}F'$ 在点 $Q$ 处的纤维，其中 $Q$ 是 $\PP^1$ 在闭纤维中的无穷远点。

只需处理 $F'$ 可构造的情形（\hyperrefext[VII][VII.prop3.3]{3.3} 和 \hyperrefext[IX][IX.prop2.9]{2.9}）。由于 $\delta(\supp F')\leq d$，于是存在 $\PP^1$ 的一个闭子概形 $X'$，满足 $\dim X'\leq d$，并包含 $F'$ 的支集。设 $\widetilde X'$ 为 $X'$ 在几何点 $Q$ 处的局部化，令
$$
\widetilde U=\widetilde X'-
\fprod{\widetilde X'}{Y'_\infty}{\PP^1},
$$
并设 $\widetilde F$ 为 $F'$ 经态射 $\widetilde U\rightarrow\PP^1$ 在 $\widetilde U$ 上诱导的层\pageoriginale。那么 $R^{q}i_{\ast}F'$ 在 $Q$ 处的纤维恰为 $H^{\sisi{p}{q}}(\widetilde U,\widetilde F)$（\hyperrefext[VIII][VIII.thm5.2]{5.2}），对严格局部概形 $\widetilde X'$ 及开子集 $\widetilde U$ 应用 \hyperref[XIV.enon4.3]{4.3}~(d) 即可。

以下引理与 \hyperref[XIV.lem4.4]{4.4} 一起将完成证明。

\begin{enonce*}{引理 4.5}\etiquette[4.5]{XIV.lem4.5}
对于 $d'<d$ 的 \emph{\hyperref[XIV.cor3.4]{3.4}~(d')} 蕴含对于 $d'\leq d$ 的 \emph{\hyperref[XIV.enon4.3]{4.3}~(d')}。
\end{enonce*}

\noindent
\emph{证明}。设 \hyperref[XIV.cor3.4]{3.4}~(d') 对 $d'<d$ 成立。由归纳法，可设 \hyperref[XIV.enon4.3]{4.3}~(d') 对 $d'<d$ 成立。设 $X'$、$f$ 和 $U$ 如陈述 \hyperref[XIV.enon4.3]{4.3}~(d) 所述。于是 $X'$ 是域 $k$ 上代数概形 $X_0$ 在一个\emph{闭}几何点处的严格局部化。用 $x_0$ 的一个邻域代替 $X_0$ 之后，可取 $X_0$ 仿射，维数 $\leq d$。于是 $X'$ 是仿射概形 $X_i$ 的逆向极限，这些概形在 $X_0$ 上平展（\hyperrefext[VIII][VIII.subsec4.5]{4.5}）。

设 $F$ 是 $U$ 上的挠层。须证当 $q>d$ 时 $H^q(U,F)=0$，而仍然只需处理 $F$ 可构造的情形。于是资料 $f$、$F$ 和 $U$ 来自某个 $X_i$ 上的类似资料（\sisi{\hyperrefext[IX][IX.prop2.7]{2.7}}{\hyperrefext[IX][IX.cor2.7.4]{2.7.4}}）。因此存在 $U_i=X_i-V(f_i)$ 上的层 $F_i$，它在 $U$ 上诱导 $F$。由于 $U=\varprojlim U_i$，有（\hyperrefext[VII][VII.cor5.8]{5.8}）
$$
H^q(U,F)=\varinjlim H^q(U_i,F_i){\quoi}.
$$
所以只需证明以下事实：对于每个 $i$，存在一个交换图
$$
\xymatrix@=1.5cm{
 U \ar[rr] \ar[rd]&        & U''\ar[ld]\\
   &  U_i   &
}
$$
其中 $U\rightarrow U_i$ 为给定的箭头，且 $\cd(U'')\sisi{<}{\leq}d$。事实上，这便蕴含\pageoriginale 当 $q>d$ 时，$H^q(U_i,F_i)$ 的每个元素在 $H^q(U,F)$ 中的像都为零。

设 $i=0$，考虑由 $f_0\in\Gamma(X_0,\Oo_{X_0})$ 定义的态射 $g:X_0\rightarrow\EE^1_k=Y_0$。设 $Y''$ 为 $\EE^1_k$ 在原点上方某个几何点处的严格局部化。于是 $Y''$ 为离散赋值环的谱。令 $X''=\fprod{X_0}{Y''}{Y_0}$。

由于 $X'$ 严格局部，态射 $X'\rightarrow X_0$ 通过 $X''$ 分解。此外，开子集 $U''=\fprod{X''}{U_0}{X_0}$ 恰为 $X''/Y''$ 的泛纤维。因此，它是 $Y''$ 的泛点上的仿射代数概形，而且显然 $\dim U''\leq d-1$。

验证 $\cd U''\leq d$，就完成了证明。设 $K$ 为 $Y''$ 在其泛点处的剩余域。这个域的上同调维数为 $1$；即 $\cd(G)=1$，其中 $G$ 为 $K$ 的可分闭包 $\bar K$ 的 Galois 群（\hyperrefext[X][X.thm2.2]{2.2}）。令 $\bar U''=\fprod{U''}{(\Spec\bar K)}{\Spec K}$。由关于 $d$ 的归纳假设，有 $\cd\bar U''\leq d-1$（这是\sisi{对维数 $\leq d$ 的概形应用 \hyperref[XIV.def2.2]{2.2}}{对维数 $<d$ 的概形应用 \hyperref[XIV.cor3.2]{3.2}}）。于是，由 Hochschild--Serre 谱序列（\hyperrefext[VIII][VIII.prop8.4]{8.4}），得 $\cd U''\leq d$：
$$
H^p(G,H^q(\bar U\sisi{'}{''},\bar F))\Longrightarrow H^n(U'',F){\quoi},
$$
这就证明了引理 \hyperref[XIV.lem4.5]{4.5}。

% END native input: chapters/exp-14-zh.tex
% BEGIN native input: chapters/exp-15-zh.tex
\chapter{无上同调态射}
%\addtocontents{toc}{M. Artin}
\begin{center}
M. Artin
\end{center}
\phantomsection\MakeLinkTarget*{sga4.exp.XV}\label{sga4.exp.XV}

\setcounter{pageoriginal}{167}

% 原底本的排版备忘:
% 统一限制记号, 例如 F|X -> F_{|X}.
% 添加 ~, 避免短数学表达式孤立在行首.
% 将 \ref{...} 改为 (\ref{}); 例如参见 1.15 的陈述.
% 注记不排成斜体 (参见注 1.18).

\section*{引言}\pageoriginale

设 $g:X\rightarrow Y$ 为概形态射. 在第一节中, 我们研究使 $g$ 无上同调（acyclic）的条件, 即对 $Y$ 上的每个挠层（torsion sheaf）$F$ 都有
$H^{q}(Y,F)\simeq H^{q}(X,g^{\ast}F)$, 并且将 $Y$ 换成其上的平展概形 $Y'$ 后这一性质仍然成立. 此处平展态射的底为 $Y$. 无上同调性的一个自然的必要条件是几何纤维无上同调, 即常值挠层的上同调平凡. 我们将看到, 这一条件加上一个称为 $g$ 的局部无上同调性（local acyclicity）的局部条件 (\hyperref[XV.def1.11]{1.11}), 就能推出 $g$ 无上同调.

基本定理是: 光滑态射对于阶与剩余特征互素的挠层局部无上同调, 见 \hyperref[XV.thm2.1]{2.1}. 由此将推出光滑态射的基变换定理 \hyperrefext[XVI][XVI.thm1.1]{1.1}; 下一讲将阐述这一定理及其最初的若干推论.

本讲独立于第 XI 至 XIV 讲, 特别是不依赖于“固有态射的基变换定理”. 从下一讲开始, 后一定理将与本讲的基本定理结合, 再次发挥关键作用.

\section{整体及局部无上同调态射的一般性质}\etiquette[1]{XV.sec1}\pageoriginale

% 此处未将 g' 定义为 g x_Y Y', 而 1.3 中作了这一定义.

\begin{enonce*}{命题 1.1}\etiquette[1.1]{XV.prop1.1}
设 $g:X\rightarrow Y$ 为概形态射. 下列断言 {\rm(i)} 至 {\rm(iii)} 等价:
\begin{itemize}
\item[(i)] 对每个平展态射 $Y'\rightarrow Y$ (其中可取 $Y'$ 在 $Y$ 的某个仿射开子集上为仿射的) 以及 $Y'$ 上的每个集合层 $F$, 典范态射 $F\rightarrow
  g'_{\ast}{g'}^{\ast}F$ 是单射 (\resp 双射).

\item[(ii)] 对每个平展态射 $Y'\rightarrow Y$ (其中可取 $Y'$ 在 $Y$ 的某个仿射开子集上为仿射的) 以及 $Y'$ 上的每个集合层 $F$, 典范态射 $H^{0}(Y',F)\rightarrow
  H^{0}(X',{g'}^{\ast}F)$ 是单射 (\resp 双射).

\item[(iii)] 对每个局部拟有限态射 $Y'\rightarrow Y$ 以及 $Y'$ 上的每个集合层 $F$, 典范态射 $H^{0}(Y',F)\rightarrow H^{0}(X',{g'}^{\ast}F)$ 是单射 (\resp 双射).

\item[(ii bis)] (设 $g$ 拟紧且拟分离.) 与 {\rm(ii)} 相同, 但只取形如 $I_{Y'}$ 的常值层 $F$, 其中 $I$ 是一个预先指定的集合, 且 $\card I\geq 2$.
\end{itemize}
\end{enonce*}

\noindent
\emph{证明}. (i) 与 (ii) 的等价性直接来自定义, 而 (iii) $\Rightarrow$ (ii) $\Rightarrow$ (ii bis) 是显然的. 有 (ii bis) $\Rightarrow$ (ii), 因为可设 $Y$ 仿射, 从而 $X$ 与 $Y$ 都拟紧且拟分离, 然后应用 \hyperrefext[XII][XII.prop6.5]{6.5} (i). 为证明 (i) $\Rightarrow$ (iii), 我们先给出下列引理:

\begin{enonce*}{引理 1.2}\etiquette[1.2]{XV.lem1.2}
设\pageoriginale $f:Y'\rightarrow Y$ 为局部拟有限态射. 则 $f$ “对于平展拓扑在 $Y'$ 上局部有限”, 即存在交换图
$$\vcenter{
\xymatrix@=1.75cm{
Y'_i        \ar[r]^{a_i}   \ar[d]^{b_i}& Y'\ar[d]^{f}\\
Y_i          \ar[r]^{c_i}              & Y
}}\leqno{(*)}
$$
其中 $a_i$ 与 $c_i$ 平展, $b_i$ 有限, 而 $a_i$ 构成 $Y'$ 的一个覆盖 \footnote[1]{\Cf \EGA IV 18.12.1.}.
\end{enonce*}

事实上, 设 $y'$ 为 $Y'$ 的一点, 并令 $y=f(y')$. 只须找到一个图 $({\ast})$, 使 $y'$ 属于 $Y'_i$ 的像. 设 $\widetilde{Y}$ 为环 $\Oo_{Y,y}$ 的 Hensel 化的谱 \hyperrefext[VIII][VIII.subsec4.1]{4.1}, 而 $Z$ 为 $\widetilde{Y}'=\fprod{Y'}{\widetilde{Y}}{Y}$ 在位于 $y'$ 及 $\widetilde{Y}$ 的闭点 $\widetilde{y}$ 之上的唯一一点处的 (通常意义的) 局部化. 由 \hyperrefext[VIII][VIII.subsec4.1]{4.1} (ii),
% (ii) 的标签?
$Z\rightarrow \widetilde{Y}$ 有限, 而 $Z$ 是 $\widetilde{Y}'$ 的诱导开子概形. 由于 $\widetilde{Y}$ 是 $Y$ 上平展概形 $Y_i$ 的极限, 由 \EGA IV 9 立即可知, 对适当的 $Y_i$, 存在 $\fprod{Y'}{Y_i}{Y}$ 的开子概形 $Z_i$, 使得 $Z=\fprod{Z_i}{\widetilde{Y}}{Y_i}$; 此时在 $({\ast})$ 中取 $Y'_i=Z_i$ 即可. \cqfd

现在假设 \hyperref[XV.prop1.1]{1.1} 的 (i) 成立. 显然, \hyperref[XV.prop1.1]{1.1} 的断言 (iii) 等价于将 (i) 中的平展态射 $Y'\rightarrow Y$ 换成局部拟有限态射后所得的断言. 在这一形式下, 对给定的 $Y$ 与 $Y'$, 该断言对于平展拓扑是 $Y$ 与 $Y'$ 上的局部性质; 因此该引理将问题化为 $f:Y'\rightarrow Y$ 为\emph{有限}态射的情形.\pageoriginale 设 $f':X'\rightarrow X$ 为 $f$ 沿 $g$ 作基变换所得的态射. 则基变换态射 $g^{\ast}f_{\ast}F\rightarrow {f'}_{\ast}{g'}^{\ast}F$ 是双射 \hyperrefext[VIII][VIII.cor5.6]{5.6}. 由 (i),
\sisi{伴随态射 $f_{\ast}F\rightarrow g_{\ast}g^{\ast}f_{\ast}F$ 为单射 (相应地, 为双射)}{伴随态射 $f_{\ast}F\rightarrow g_{\ast}g^{\ast}f_{\ast}F$ 为单射 (相应地, 为双射)}.
因而 (iii) 中的态射 $H^{0}(Y',F)\xrightarrow{\epsilon}
H^{0}(Y',g'_{\ast}{g'}^{\ast}F)$ 是复合态射 $H^{0}(Y',F)=H^{0}(Y,f_{\ast}F)\xrightarrow{\epsilon'}
H^{0}(Y,g_{\ast}g^{\ast}f_{\ast}F)=H^{0}(X,g^{\ast}f_{\ast}F)\simeq
H^{0}(X,{f'}_{\ast}{g'}^{\ast}F)=H^{0}(X',{g'}^{\ast}F)$.
由于 $\epsilon'$ 是单射 (\resp 双射), $\epsilon$ 也如此. 这就完成了 \hyperref[XV.prop1.1]{1.1} 的证明.

\begin{enonce*}{定义 1.3}\etiquette[1.3]{XV.def1.3}
具有 \hyperref[XV.prop1.1]{1.1} 中等价性质 {\rm(i)} 至 {\rm(iii)} 之一 (\resp 相应括号内性质) 的态射 $g:X\rightarrow Y$, 称为\emph{$(-1)$-无上同调态射} (\resp \emph{$0$-无上同调态射}). 若对每个基变换态射 $Y'\rightarrow Y$, 态射 $g'=\fprod{g}{Y'}{Y}:\fprod{X}{Y'}{Y}\rightarrow Y'$ 都是 $(-1)$-无上同调的 (\resp $0$-无上同调的), 则称原态射\emph{普遍} $(-1)$-无上同调 (\resp $0$-无上同调). 最后, 若概形 $X$ 非空 (\resp 连通且非空), 则称其为 $(-1)$-无上同调的 (\resp $0$-无上同调的).
\end{enonce*}

\begin{enonce*}{推论 1.4}\etiquette[1.4]{XV.cor1.4}
设 $g:X\rightarrow Y$ 为 $0$-无上同调态射. 则对每个局部拟有限态射 $Y'\rightarrow Y$ 以及 $Y'$ 上的每个群层 $F$, 典范态射
$$
H^{1}(Y',F)\longrightarrow H^{1}(X',{g'}^{\ast}F)
$$
是单射.
\end{enonce*}

这直接来自 \hyperrefext[XII][XII.prop6.5]{6.5} (i).

\begin{enonce*}{推论 1.5}\etiquette[1.5]{XV.cor1.5}
设\pageoriginale $g:X\rightarrow Y$ 为一个态射. 若 $g$ 满射 (\ie 其纤维非空, \ie 其纤维 $(-1)$-无上同调), 则 $g$ 为 $(-1)$-无上同调的. 若 $g$ 拟紧且拟分离, 则逆命题成立.
\end{enonce*}

通过考察有关层的茎, 第一项断言在 \hyperref[XV.prop1.1]{1.1} (ii) 的形式下是显然的. 第二项来自 \hyperref[XV.prop1.1]{1.1} (iii) 形式的定义及 \hyperrefext[VII][VII.sec5]{5} 的极限论证; 这里用到, 对每个 $y\in Y$, $y=\Spec k(y)$ 是 $Y$ 上局部拟有限仿射概形的投射极限, 即 $y$ 在 $\bar{\{y\}}$ (赋予诱导的既约结构) 中的仿射邻域的投射极限. \sisi{}{--} 注意, \hyperref[XV.cor1.5]{1.5} 表明, 当 $g$ 拟紧且拟分离时, $g$ 的 $(-1)$-无上同调性已经蕴含 $g$ \emph{普遍} $(-1)$-无上同调.

\begin{enonce*}{命题 1.6}\etiquette[1.6]{XV.prop1.6}
设 $g:X\rightarrow Y$ 为概形态射, $\LL\subset\PP$ 为素数集合, $n\in\NN$ 为自然数 (\resp 设 $n=1$). 下列断言 {\rm(i)} 至 {\rm(iii)} (\resp 相应括号内的断言) 等价. 若 $g$ 还拟紧且拟分离, 则 {\rm(i)} 至 {\rm(iii)} (\resp 相应括号内的断言) 也等价于 {\rm(iv)} (\resp 相应括号内的断言).
\begin{enumerate}
\renewcommand{\theenumi}{\roman{enumi}}
\renewcommand{\labelenumi}{(\theenumi)}
\item 对每个平展态射 $Y'\rightarrow Y$ 以及 $Y'$ 上的每个阿贝尔 $\LL$-挠层 (\resp $\text{\rm ind}$-$\LL$-群层) \emph{\hyperrefext[IX][IX.def1.5]{1.5}} $F$, 都有
$$
\begin{cases}
F\xrightarrow{\sim} {g'}_{\ast}{g'}^{\ast}F, &\\
(R^{q}{g'}_{\ast}){g'}^{\ast}F=0 & \mbox{若}\quad 1\leq q\leq n\quad
(\resp \mbox{若\; } q=1){\quoi}.
\end{cases}
$$

\item 对\pageoriginale 每个平展态射 $Y'\rightarrow Y$ 以及 $Y'$ 上的每个阿贝尔 $\LL$-挠层 (\resp ind-$\LL$-群层) $F$, 典范态射
$$
H^{q}(Y',F)\longrightarrow H^{q}(X',{g'}^{\ast}F)
$$
在 $q\leq n$ 时为双射, 在 $q=n+1$ 时为单射 (\resp 在 $q\leq1$ 时为双射).

\item 对每个局部拟有限态射 $Y'\rightarrow Y$ 以及 $Y'$ 上的每个阿贝尔 $\LL$-挠层 (\resp ind-$\LL$-群层) $F$, 态射
$$
H^{q}(Y',F)\longrightarrow H^{q}(X', {g'}^{\ast}F)
$$
在 $q\leq n$ 时为双射, 在 $q=n+1$ 时为单射 (\resp 在 $q\leq1$ 时为双射).

\item 态射 $g$ 满射, 并且对每个局部拟有限态射 $Y'\rightarrow Y$, 每个 $\nu>0$ 以及每个 $\ell\in\LL$, 典范态射
$$
H^{q}(Y',\ZZ/\ell^{\nu})\longrightarrow H^{q}(X',\ZZ/\ell^{\nu})
$$
在 $0\leq q\leq n$ 时为\emph{满射} (\resp 对每个局部拟有限态射 $Y'\rightarrow Y$ 以及每个通常的有限 $\LL$-群 $G$, 典范态射
$$
H^{i}(Y',G)\longrightarrow H^{i}(X',G)
$$
在 $i=0,1$ 时为\emph{满射}).
\end{enumerate}
\end{enonce*}

\noindent
\hyperref[XV.prop1.6]{1.6} 的\emph{证明}\pageoriginale. 蕴含 (ii) $\Leftarrow$ (iii) $\Rightarrow$ (iv) 是显然的.
% 满射性显然吗? (g 未必为凝聚态射.)
现在 $(R^{q}{g'}_{\ast}){g'}^{\ast}F$ 是预层 $R^{q}(Y'')=H^{q}(X'',{g''}^{\ast}F)$ ($Y''\rightarrow Y'$ 平展) 的伴随层. 若 (ii) 成立, 则有 $H^{q}(Y'',F)\xrightarrow{\sim}H^{q}(X'',{g''}^{\ast}F)$, 所以伴随层为零\sisi{ (当 $1\leq q\leq n$ 时)}{}, \ie (ii) $\Rightarrow$ (i). 在阿贝尔情形, 由 \hyperref[XV.prop1.1]{1.1} (i) (对应于 $n=0$ 的情形) 以及 Leray 谱序列
$$
H^{p}(Y',(R^{q}{g'}_{\ast}){g'}^{\ast}F)\Longrightarrow
H^{p+q}(X',{g'}^{\ast}F)
$$
可得 (i) $\Rightarrow$ (ii); 在群层情形则使用正合列
$$
0\longrightarrow H^{1}(Y',{g'}_{\ast}{g'}^{\ast}F)\longrightarrow
H^{1}(X',{g'}^{\ast}F)\longrightarrow
H^{0}(Y',(R^{1}{g'}_{\ast}){g'}^{\ast}F){\quoi}
$$

(i) $\Rightarrow$ (iii): 由于 (i) 与 (ii) 等价, 显然 (iii) 等价于将 (i) 中的“平展”换成“局部拟有限”后所得的陈述. 在这一形式下, 对给定的 $Y$ 与 $Y'$, 它对于平展拓扑是 $Y'$ 与 $Y$ 上的局部断言, 因而 \hyperref[XV.lem1.2]{1.2} 将问题化为 $f:Y'\rightarrow Y$ 为\emph{有限}态射的情形. 设 $f':X'\rightarrow X$ 为\sisi{$f$ 沿 $g$ 作基变换所得的态射}{$f$ 沿 $g$ 作基变换所得的态射}. 由 \hyperrefext[VIII][VIII.prop5.5]{5.5} (\resp \hyperrefext[VIII][VIII.rem5.8]{5.8}) (应用三次), 有
$$
f_{\ast}(R^{q}{g'}_{\ast}){g'}^{\ast}F\simeq
R^{q}(fg')_{\ast}{g'}^{\ast}F\simeq
R^{q}(gf')_{\ast}{g'}^{\ast}F\simeq
(R^{q}g_{\ast}){f'}_{\ast}{g'}^{\ast}F\simeq
(R^{q}g_{\ast})g^{\ast}(f_{\ast}F)
$$
若 (i) 成立, 则对于所考察的 $q$ 的取值, 上式最后一项为零. 因此 (i) $\Rightarrow$ (iii).

只剩下证明 (iv) $\Rightarrow$ (iii), 而这直接\pageoriginale 来自 \hyperrefext[XII][XII.prop6.5]{6.5}.

\begin{enonce*}{定义 1.7}\etiquette[1.7]{XV.def1.7}
设 $n\in\NN$ 且 $\LL\subset\PP$. 满足 \hyperref[XV.prop1.6]{1.6} 中等价条件 {\rm(i)} 至 {\rm(iii)} 之一 (\resp 相应括号内条件) 的态射 $g:X\rightarrow Y$, 称为关于 $\LL$ 的\emph{$n$-无上同调态射} (\resp \emph{$1$-非球面态射}（aspherical morphism）). 若 $g$ 对每个 $n$ 都关于 $\LL$ 为 $n$-无上同调的, 则称其关于 $\LL$ \emph{无上同调}. 若对每个 $Y'\rightarrow Y$, 态射 $g'=\fprod{g}{Y'}{Y}$ 关于 $\LL$ 为 $n$-无上同调的 (\resp $1$-非球面的), 则称 $g$ 关于 $\LL$ \emph{普遍 $n$-无上同调} (\resp \emph{普遍 $1$-非球面}). 最后, 若概形 $X$ 在 \hyperref[XV.def1.3]{1.3} 的意义下为 $0$-无上同调的, 并且对每个 $\ell\in\LL$ 和每个 $1\leq q\leq n$ 都有 $H^{q}(X,\ZZ/\ell)=0$ (\resp 并且对每个通常的有限 $\LL$-群 $G$ 都有 $H^{1}(X,G)=0$), 则称其关于 $\LL$ 为 $n$-\emph{无上同调}的 (\resp $1$-\emph{非球面}的).
\end{enonce*}

\begin{enonce*}[remark]{注 1.8}\etiquette[1.8]{XV.rem1.8}
\begin{enumerate}
\renewcommand{\theenumi}{\alph{enumi}}
\renewcommand{\labelenumi}{\theenumi)}
\item 若 $g$ 拟紧且拟分离, 则要使 $g$ 普遍 $n$-无上同调, 只须对每个局部有限表示态射 $Y'\rightarrow Y$, $g'$ 都为 $n$-无上同调的. 我们不需要这一事实, 将其证明留给读者, 作为一道关于过渡到极限的练习.

\item 我们不知道, 关于 $\LL$ 为 $n$-无上同调的态射, 是否必定关于 $\LL$ 普遍 $n$-无上同调. 下文 \hyperref[XV.def1.11]{1.11} 的局部变体也有同样的问题.

\item 当 $n=0$ 时, 就得到 \hyperref[XV.def1.3]{1.3} 的概念. 当然, \hyperref[XV.prop1.6]{1.6} 和 \hyperref[XV.def1.7]{1.7} 本可表述为同时包含 $n=-1$ 的情形. 以下命题在该情形也成立.\pageoriginale
\end{enumerate}
\end{enonce*}

\begin{enonce*}{命题 1.9}\etiquette[1.9]{XV.prop1.9}
\begin{enumerate}
\renewcommand{\theenumi}{\roman{enumi}}
\renewcommand{\labelenumi}{(\theenumi)}
\item 设 $X\xrightarrow{g}Y\xrightarrow{h}Z$ 为概形态射, $\LL\subset\PP$ 且 $n\in\NN$. 假设 $g$ 为 $0$-无上同调的 (\resp 关于 $\LL$ 为 $1$-非球面的, \resp 关于 $\LL$ 为 $n$-无上同调的). 令 $f=hg$. 则对 $Z$ 上的每个集合层 (\resp ind\sisi{-}{ }$\LL$-群层, \resp 阿贝尔 $\LL$-挠层) $F$, 有同构 $f_{\ast}f^{\ast}F\xrightarrow{\sim}h_{\ast}h^{\ast}F$ (\resp 当 $p\leq1$ 时有同构 $(R^{p}f_{\ast})f^{\ast}F\xrightarrow{\sim}
  (R^{p}h_{\ast})h^{\ast}F$, \resp 当 $p\leq n$ 时有这些同构). 特别地, $h$ 为 $0$-无上同调的 (\resp 关于 $\LL$ 为 $1$-非球面的, \resp 关于 $\LL$ 为 $n$-无上同调的), 当且仅当 $f$ 具有相应性质.

\item 设 $\{g_{\alpha}:X_{\alpha}\rightarrow Y_{\alpha}\}$ 为概形态射的投射系, 其中态射 $X_{\alpha}\rightarrow X_{\beta}$ 与 $Y_{\alpha}\rightarrow Y_{\beta}$ 仿射, $g_{\alpha}$ 拟紧且拟分离. 设 $g:X\rightarrow Y$ 为 $g_{\alpha}$ 的投射极限 \hyperrefext[VII][VII.subsec5.1]{5.1}. 若所有 $g_{\alpha}$ 都为 $0$-无上同调的 (\resp 关于 $\LL$ 为 $1$-非球面的, \resp 关于 $\LL$ 为 $n$-无上同调的), 则 $g$ 也如此.
\end{enumerate}
\end{enonce*}

\noindent
\emph{证明}. 对于 (i), 首先注意
$$
h_{\ast}h^{\ast}F\xrightarrow{\sim}
h_{\ast}(g_{\ast}g^{\ast})h^{\ast}F\simeq f_{\ast}f^{\ast}F{\quoi}.
$$
(i) 的第二项断言直接来自 Leray 谱序列
$$
(R^{p}h_{\ast})(R^{q}g_{\ast})f^{\ast}F\Longrightarrow
(R^{p+q}f_{\ast})f^{\ast}F{\quoi}.
$$
最后, 非阿贝尔断言来自正合列 \hyperrefext[XII][XII.prop3.2]{3.2}.

断言\pageoriginale (ii) 容易从过渡到极限的理论 \hyperrefext[VII][VII.sec5]{5} 及第 VI 讲推出. % 讲的标签?

\begin{enonce*}{推论 1.9.1}\etiquette[1.9.1]{XV.cor1.9.1}
设 $g:X\rightarrow Y$ 为拟紧且拟分离态射, 并设 $n\in \NN$, $\LL\subset \PP$. 若 $g$ 关于 $\LL$ 为 $n$-无上同调的 (\resp 关于 $\LL$ 为 $1$-非球面的), 则对 $Y$ 的每个几何点 $\bar{y}$, 只要它在 $Y$ 的某一点 $y$ 上代数, 几何纤维 $X_{\bar{y}}$ 就关于 $\LL$ 为 $n$-无上同调的 (\resp 关于 $\LL$ 为 $1$-非球面的).
\end{enonce*}

事实上, $\bar{y}$ 是 $Y$ 上局部拟有限仿射概形 $Y_i$ 的投射极限, 又因 $X_i=\fprod{X}{Y_i}{Y}\rightarrow Y_i$ 关于 $\LL$ 为 $n$-无上同调的 (\resp\ldots), 由 \hyperref[XV.prop1.9]{1.9} (ii), $X_{\bar{y}}\rightarrow \bar{y}$ 也具有这一性质, 从而得到所述结论.

\begin{enonce*}{命题 1.10}\etiquette[1.10]{XV.prop1.10}
设 $g:X\rightarrow Y$ 为概形态射, $\LL\subset
\PP$ 且 $n\in \NN$. 则下列条件 {\rm(i)} 至 {\rm(iv)} 等价.
\begin{enumerate}
\renewcommand{\theenumi}{\roman{enumi}}
\renewcommand{\labelenumi}{(\theenumi)}
\item 设 $\bar{y}$ 为 $Y$ 的一个几何点, $\bar{x}$ 为 $X$ 中位于 $\bar{y}$ 之上的几何点. 设 $\widetilde{X}$ 与 $\widetilde{Y}$ 分别为 $X$ 与 $Y$ 在 $\bar{x}$ 与 $\bar{y}$ 处的严格局部化（strict localization）, 并设 $\widetilde{g}:\widetilde{X}\rightarrow
\widetilde{Y}$ 为 $g$ 诱导的态射. 则 $\widetilde{g}$ 为 $0$-无上同调的 (\resp 关于 $\LL$ 为 $1$-非球面的, \resp 关于 $\LL$ 为 $n$-无上同调的).

\item 对每个方块均为笛卡儿方块的图
$$
\vcenter{
\xymatrix@C=2cm@R=1.25cm{
X   \ar[d]^{g}& X'\ar[l]_{j}  \ar[d]^{g'} & X'' \ar[l]_{j'} \ar[d]^{g''}\\
Y    & Y'\ar[l]_{i} &   Y'' \ar[l]_{i'} 
}}\leqno{(\ast)}
$$
其中\pageoriginale $i$ 平展, $i'$ 平展且有限表示, 以及 $Y''$ 上的每个集合层 (\resp ind-$\LL$-群层, \resp 阿贝尔 $\LL$-挠层) $F$, 基变换态射 \hyperrefext[XII][XII.eq4.1]{4.1}, \hyperrefext[XII][XII.eq4.2]{4.2}
$$
{g'}^{\ast}i'_{\ast}F\longrightarrow {j'}_{\ast}{g''}^{\ast}F
$$
为双射 (\resp 态射
$$
{g'}^{\ast}(R^{q}i'_{\ast})F\longrightarrow
(R^{q}j'_{\ast}){g''}^{\ast}F
$$
在 $q=0$, $1$ 时为双射, \resp 在 $0\leq q\leq
n$ 时为双射, 在 $q=n+1$ 时为单射).

\item 与 {\rm(ii)} 相同, 但只假设态射 $i'$ 拟有限且拟分离.

\item 与 {\rm(ii)} 相同, 但只假设 $i$ 局部拟有限, 而另一方面, 假设 $i'$ 为拟紧开浸入.
\end{enumerate}
\end{enonce*}

\noindent
\emph{证明}. (iii) $\Rightarrow$ (ii) 是显然的. 我们验证 (i) $\Rightarrow$ (iii). 只须逐茎验证 (iii) 中的态射. 设 $\bar{y}$ 为 $Y'$ 的几何点, $\bar{x}$ 为 $X'$ 中位于 $\bar{y}$ 之上的几何点. 设
$$
\xymatrix@=1.75cm{
\widetilde{X}'   \ar[d]^{\widetilde{g}'}& \widetilde{X}''
\ar[l]_{\widetilde{j}'} \ar[d]^{\widetilde{g}''}\\
 \widetilde{Y}'           & \widetilde{Y}'' \ar[l]_{\widetilde{i}'}
}
$$
为由 $({\ast})$ 导出的笛卡儿图, 其中 $\widetilde{X}'$ 与 $\widetilde{Y}'$ 分别为 $X'$ 与 $Y'$ 在 $\bar{x}$ 与 $\bar{y}$ 处的严格局部化, 且 $\widetilde{Y}''=\fprod{Y''}{\widetilde{Y}'}{Y'}$, $\widetilde{X}''=\fprod{X''}{\widetilde{X}'}{X'}$.\pageoriginale 由 \hyperrefext[VIII][VIII.thm5.2]{5.2} 和 \hyperrefext[VIII][VIII.rem5.3]{5.3}, 有
$$
H^{q}(\widetilde{Y}'',\widetilde{F})\simeq
(R^{q}i'{}_{\ast}F)_{\bar{y}}\simeq
({g'}^{\ast}(R^{q}i'_{\ast})F)_{\bar{x}}{\quoi},
$$
其中 $\widetilde{F}$ 为 $F$ 在 $\widetilde{Y}''$ 上的逆像, 并且
$$
H^{q}(\widetilde{X}'',{\widetilde{g}''}{}^{\ast}\widetilde{F})\simeq
((R^{q}j'_{\ast}){g''}{}^{\ast}F)_{\bar{x}}{\quoi}.
$$
由于 $\widetilde{i}'$ 拟有限, 可对态射
$$
H^{q}(\widetilde{Y}'',\widetilde{F})\longrightarrow
H^{q}(\widetilde{X}'',{\widetilde{g}''}{}^{\ast}\widetilde{F})
$$
以及所考察的 $q$ 的取值应用 \hyperref[XV.prop1.1]{1.1} (iii) (\resp \hyperref[XV.prop1.6]{1.6} (iii)), 这就证明了 (i) $\Rightarrow$ (iii).

(ii) $\Rightarrow$ (i). 假设 (ii), 我们对态射 $\widetilde{g}:\widetilde{X}\rightarrow \widetilde{Y}$ 验证无上同调性判据 \hyperref[XV.prop1.6]{1.6} (ii). 设 $\widetilde{Y}'\rightarrow \widetilde{Y}$ 平展, $\widetilde{Y}'$ 仿射, 且 $\widetilde{F}$ 为 $\widetilde{Y}'$ 上的集合层 (\resp 阿贝尔 $\LL$-挠层, \resp ind-$\LL$-群层). 必须证明
$$
H^{q}(\widetilde{Y}',\widetilde{F})\xrightarrow{\sim}
H^{q}(\widetilde{X}', \widetilde{g}'{}^{\ast}\widetilde{F}){\quoi},
$$
其中 $\widetilde{g}':\widetilde{X}'\rightarrow \widetilde{Y}'$ 是 $\widetilde{g}$ 作基变换所得的态射. 现写成 $\widetilde{Y}=\varprojlim Y_{\alpha}$, 其中 $Y_{\alpha}$ 仿射且在 $Y$ 上平展 (如前所述可假设 $Y$ 仿射). 可以将平展态射 $\widetilde{Y}'\rightarrow \widetilde{Y}$ “下降”到某个 $Y_{\alpha}$ 上, 例如下降为 $Y'_{0}\rightarrow Y_{0}$, 而且有 $\widetilde{Y}'=\varprojlim
Y'_{\alpha}=\fprod{Y}{Y'_{0}}{Y_{0}}$. $\widetilde{Y}'$ 上的每个层 $\widetilde{F}$ 是层 ${f'_{\alpha}}^{\ast}{f'_{\alpha}}_{\ast}\widetilde{F}$ 的归纳极限, 其中 $f'_{\alpha}:\widetilde{Y}'\rightarrow Y'_{\alpha}$ 为典范态射. 由 \sisi{XII 1.2 (v) 及 1.6 (iii)}{\hyperrefext[IX][IX.prop1.2]{1.2} 和 \hyperrefext[IX][IX.prop1.6]{1.6}}, 层 $f'_{\alpha}{}_{\ast}\widetilde{F}$ 为集合层 (\resp 阿贝尔 $\LL$-挠层, \resp ind-$\LL$-群层). 根据\pageoriginale \hyperrefext[VII][VII.prop3.3]{3.3}, 只须处理 $\widetilde{F}$ 可下降到某个 $Y'_{\alpha}$ 上的情形, 比如下降为 $Y'_{0}$ 上的层 $F_{0}$. 将 (ii) 应用于图
$$
\xymatrix@C=2.25cm@R=1.5cm{
X \ar[d]_{g}& X_0 \ar[l] \ar[d]_{g_0}& X'_0 \ar[l]_{j'_0} \ar[d]_{g'_0}\\
Y & Y_0 \ar[l]& Y'_0 \ar[l]_{i'_0},
}
$$
可得, 对相应的 $q$ 的取值, 有 ${g_{0}}^{\ast}(R^{q}{i'_0}_{\ast})F_{0}\xrightarrow{\sim}
(R^{q}{j'_0}_{\ast}){g'_0}^{\ast}F_{0}$. 几何点 $\bar{y}$ 与 $\bar{x}$ 典范地提升为 $Y_{0}$ 与 $X_{0}$ 的几何点 (仍用相同符号表示), 又由 \hyperrefext[VIII][VIII.thm5.2]{5.2} 和 \hyperrefext[VIII][VIII.rem5.3]{5.3}, 对所考察的 $q$ 的取值, 有
$$
H^{q}(\widetilde{Y}',\widetilde{F})=((R^{q}{i'_0}_{\ast})F_{0})_{\bar{y}}=({g_0}^{\ast}
(R^{q}{i'_0}_{\ast})F_0)_{\bar{x}}
$$
以及
$$
H^{q}(\widetilde{X}',{\widetilde{g}'}{}^{\ast}\widetilde{F})=((R^{q}j'_{0\ast}){g'_0}^{\ast}F_{0})_{\bar{x}}
$$
这就证明了结论.

只剩证明条件 (i) 至 (iii) 与 (iv) 等价. 利用形式 (i), 可知这些条件在局部拟有限基变换 $Y'\rightarrow Y$ 下稳定; 再利用形式 (ii), 便立即看出它们蕴含 (iv). 反之, (iv) 蕴含 (iii), 因为要验证 (iii), 例如利用 Leray 谱序列或下降结论 \hyperrefext[XII][XII.prop6.8]{6.8}, 立即可知可以假设 $Y'$ 与 $Y''$ 仿射. 此时, 由“主定理” \EGA IV 8.12.8, $Y''\rightarrow Y'$ 可分解为 $Y''\xrightarrow{i'_1}Y'_1\xrightarrow{i'_2}Y'$, 其中 $i'_1$ 是开浸入,\pageoriginale $i'_2$ 是\emph{有限}态射. 对 $i'_2$ 使用 \hyperrefext[VIII][VIII.prop5.5]{5.5}, \hyperrefext[VIII][VIII.rem5.8]{5.8}, 就将问题化为对 $i'_1$ 验证 (ii), 而这来自 (iv). 这就完成了 \hyperref[XV.prop1.10]{1.10} 的证明.

\begin{enonce*}{定义 1.11}\etiquette[1.11]{XV.def1.11}
满足 $\hyperref[XV.prop1.10]{1.10}$ 中等价条件 {\rm(i)} 至 {\rm(iv)} 之一的预概形态射 $g:X\rightarrow Y$, 称为\emph{局部 $0$-无上同调}态射 (\resp 关于 $\LL$ \emph{局部 $1$-非球面}, \resp 关于 $\LL$ \emph{局部 $n$-无上同调}). 以显然的方式 (比较 \hyperref[XV.def1.7]{1.7}), 可定义关于 $\LL$ 的\emph{局部无上同调}态射与\emph{普遍局部 $0$-无上同调}态射 (\resp\ldots) 的概念.
\end{enonce*}

当然, \hyperref[XV.prop1.10]{1.10} 的陈述 (i) 至 (iii) 可以作若干变化. 因而, 在 (iii) 和 (iv) 中, 只须限于形如 $G_{I}$ 的常值层 $F$,
% 这一记号的含义?
其中 $G$ 为有限集 (\resp 有限 $\LL$-群, \resp $G$ 形如 $\ZZ/\ell^{\nu}\ZZ$, 其中 $\ell\in\LL$, $\nu\in \NN$); 这容易从 \hyperrefext[VIII][VIII.thm5.2]{5.2}, \hyperrefext[VIII][VIII.rem5.3]{5.3} 以及 \hyperrefext[XII][XII.prop6.5]{6.5} 推出. 我们还给出以下变体:

\begin{enonce*}{推论 1.12}\etiquette[1.12]{XV.cor1.12}
假设 $g:X\rightarrow Y$ 局部有限型. 要验证局部 $0$-无上同调性 (\resp\ldots), 只须在几何点 $\bar{x}$ 为几何纤维 $X_{\bar{y}}$ 中的\emph{闭点}时验证 $\hyperref[XV.prop1.10]{1.10}$ {\rm(i)}.
\end{enonce*}

事实上, 在 \hyperref[XV.prop1.10]{1.10} 的 (i) $\Rightarrow$ (iii) 的证明中, 我们逐茎验证了 (iii), 只须对在其几何纤维中为闭点的几何点 $\bar{x}$ 处的茎验证即可,\pageoriginale 因为 $g$ 局部有限型 \hyperrefext[VIII][VIII.rem3.13]{3.13} b).

\begin{enonce*}{命题 1.13}\etiquette[1.13]{XV.prop1.13}
\begin{enumerate}
\renewcommand{\theenumi}{\roman{enumi}}
\renewcommand{\labelenumi}{(\theenumi)}
\item 设 $X\xrightarrow{g}Y\xrightarrow{h}Z$ 为概形态射, $\LL\subset \PP$ 且 $n\in \NN$. 假设 $g$ 满射且局部 $0$-无上同调 (\resp 关于 $\LL$ 局部 $1$-非球面, \resp 关于 $\LL$ 局部 $n$-无上同调). 则 $h$ 局部 $0$-无上同调 (\resp\ldots), 当且仅当 $f=hg$ 具有相应性质.

\item 设 $\{g_{\alpha}:X_{\alpha}\rightarrow Y_{\alpha}\}$ 为预概形态射的投射系, 其过渡态射 $X_{\alpha}\rightarrow X_{\beta}$ 与 $Y_{\alpha}\rightarrow Y_{\beta}$ 仿射. 设 $g$ 为 $\{g_{\alpha}\}$ 的投射极限 \hyperrefext[VII][VII.subsec5.1]{5.1}. 若每个 $g_{\alpha}$ 都局部 $0$-无上同调 (\resp\ldots), 则 $g$ 也如此.
\end{enumerate}
\end{enonce*}

这可由类似于 \hyperref[XV.prop1.9]{1.9} 的一连串推理得出.

\begin{enonce*}[remark]{注 1.14}\etiquette[1.14]{XV.rems1.14}
为简化工作, 作者在 \hyperref[XV.prop1.10]{1.10} 至 \hyperref[XV.prop1.13]{1.13} 中省略了对 $n=-1$ 情形的处理, 将验证有关陈述、定义和证明基本不变地适用于这一情形的工作留给读者. 因而, 用 \hyperref[XV.prop1.10]{1.10} 的记号, 条件 (i) 表示诱导的局部态射 $\widetilde{X}\rightarrow \widetilde{Y}$ 为 $(-1)$-无上同调的, 即由 \hyperref[XV.cor1.5]{1.5}, 它们\emph{满射}; 条件 (ii) 至 (iv) 表示基变换态射
${g'}^{\ast}({i'}_{\ast}(F))\rightarrow {j'}_{\ast}({g''}^{\ast}(F))$
是\emph{单态射}. 在这些条件下, 我们便称 $g$ \emph{局部 $(-1)$-无上同调}. 因而这一条件在局部拟有限基变换下稳定.\pageoriginale 例如, 平坦态射局部 $(-1)$-无上同调, 因为对这样的态射, 映射 $\widetilde{X}\rightarrow \widetilde{Y}$ 的满射性是其平坦性的一个众所周知的推论. 我们还注意到, 当 $f$ 局部有限表示时, $f$ 局部 $(-1)$-无上同调当且仅当它\emph{普遍开}, 至少在 $Y$ 局部 Noether 的情形如此; 这容易从 \EGA IV 14.4.9, 1.10.3 推出. 因此, 对这样的态射, 局部 $(-1)$-无上同调性蕴含\emph{普遍}局部 $(-1)$-无上同调性. (而且很可能\footnote[1]{事实上 M. Raynaud 已经验证了这一点. \Cf \EGA IV 14 的新版!} 在这些最后的陈述中, $Y$ 局部 Noether 的假设是不必要的; 至少当 $Y$ 的不可约分支集局部有限时的确如此, 见 \EGA IV 14.4.10.)
\end{enonce*}

\begin{enonce*}{定理 1.15}\etiquette[1.15]{XV.thm1.15}
设 $g:X\rightarrow Y$ 为拟紧且拟分离态射, $\LL\subset \PP$ 为素数集合, $n$ 为自然数 (\resp $n=1$). 假设 $g$ 关于 $\LL$ 为 $(n-1)$-无上同调的, 且关于 $\LL$ 局部 $(n-1)$-无上同调, 见 $\hyperref[XV.def1.3]{1.3}$, $\hyperref[XV.def1.7]{1.7}$, $\hyperref[XV.def1.11]{1.11}$, $\hyperref[XV.rems1.14]{1.14}$. 设 $F$ 为 $Y$ 上的一个层; 当 $n\geq 1$ 时, 假设它具有阿贝尔 $\LL$-挠层结构 (\resp 群层结构), 并设 $\xi$ 为 $H^{n}(X,g^{\ast}(F))$ 的一个元素. 要使 $\xi$ 来自 $H^{n}(Y,F)$ 的一个元素 (由于 $g$ 的 $(n-1)$-无上同调性假设, 该元素此时唯一确定), 充分必要条件是, 对 $Y$ 的每个在 $Y$ 的某一点 $y$ 上代数的几何点 $\bar{y}$, $\xi$ 在 $H^{n}(X_{\bar{y}},F|X_{\bar{y}})$ 中的像 $\xi_{\bar{y}}$ 属于 $H^{n}(\bar{y},F_{\bar{y}})$ 的像 (当 $n\geq 1$ 时, 这就\pageoriginale 表示 $\xi_{\bar{y}}=0$).
\end{enonce*}

立即可以注意到, \hyperref[XV.thm1.15]{1.15} 表明, 在对 $g$ 所加的前提下, 态射 $g$ 关于 $\LL$ 为 $n$-无上同调的 (\resp 关于 $\LL$ 为 $1$-非球面的), 当且仅当对上述每个 $\bar{y}$, $X_{\bar{y}}$ 都关于 $\LL$ 为 $n$-无上同调的 (\resp 关于 $\LL$ 为 $1$-非球面的): 充分性直接来自 \hyperref[XV.thm1.15]{1.15} 及定义, 必要性则总是直接来自 \hyperref[XV.cor1.9.1]{1.9.1}. 据此对 $n$ 作直接归纳, 即得下列结论.

\begin{enonce*}{推论 1.16}\etiquette[1.16]{XV.cor1.16}
设 $g$ 为拟紧且拟分离态射, 关于 $\LL$ 局部 $(n-1)$-无上同调, 其中 $n$ 为给定的自然数 (\resp $n=1$). 则 $g$ 关于 $\LL$ 为 $n$-无上同调的 (\resp 关于 $\LL$ 为 $1$-非球面的), 当且仅当对 $Y$ 的每个在 $Y$ 的某一点 $y$ 上代数的几何点 $\bar{y}$, 纤维 $X_{\bar{y}}$ 都关于 $\LL$ 为 $n$-无上同调的 (\resp 关于 $\LL$ 为 $1$-非球面的).
\end{enonce*}

\begin{enonce*}{推论 1.17}\etiquette[1.17]{XV.cor1.17}
设 $g:X\rightarrow Y$ 为预概形态射, $\LL\subset \PP$ 为素数集合, $n$ 为自然数 (\resp $n=1$). 要使 $g$ 局部

关于 $\LL$ 为 $n$-无上同调的 (\resp 关于 $\LL$ 局部 $1$-非球面的), 充分必要条件是: 对 $X$ 的每个几何点 $\bar{x}$, 若 $\widetilde{g}:\widetilde{X}\rightarrow
\widetilde{Y}$ 为 $g$ 诱导的 $X$ 与 $Y$ 的严格局部化之间的态射 (局部化分别在 $\bar{x}$ 及 $Y$ 中相应的几何点 $\bar{y}$ 处进行), 则对 $\widetilde{Y}$ 的每个在 $\widetilde{Y}$ 的某一点 $z$ 上代数的几何点 $\bar{z}$, $\widetilde{g}$ 的每个相应几何纤维 $\widetilde{X}_{\bar{z}}$ 都关于 $\LL$ 为 $n$-无上同调的 (\resp 关于 $\LL$ 为 $1$-非球面的).
\end{enonce*}

必要性来自判据 \hyperref[XV.prop1.10]{1.10} (i) 及 \hyperref[XV.cor1.9.1]{1.9.1}. 为\pageoriginale 证明充分性, 对 $n$ 归纳; 这使我们在 $n\geq 1$ 时可以假设 $g$ 已经关于 $\LL$ 局部 $(n-1)$-无上同调. 态射 $\widetilde{g}$ 于是也如此, 而 \hyperref[XV.cor1.16]{1.16} 蕴含 $\widetilde{g}$ 为 $n$-无上同调的 (\resp 关于 $\LL$ 为 $1$-非球面的); 由判据 \hyperref[XV.prop1.10]{1.10} (i), 这就蕴含 $g$ 关于 $\LL$ 局部 $n$-无上同调 (\resp 局部 $1$-非球面).

\begin{enonce*}{推论 1.18}\etiquette[1.18]{XV.cor1.18}
用 \hyperref[XV.cor1.17]{1.17} 的记号, 若 $g$ 局部有限型, 则在所述判据中, 只须考虑其底层点 $x$ 在自身纤维中为闭点的几何点 $\bar{x}$.
\end{enonce*}

这来自刚才的证明及 \hyperref[XV.cor1.12]{1.12}.

\noindent
\hyperref[XV.thm1.15]{1.15} 的\emph{证明}\footnote[1]{只关心 Noether 情形的读者, 请参阅\sisi{注 }{}\hyperref[XV.rem1.18]{1.18}.}.
% 本讲中有两个编号为 1.18 的陈述.
先证明可以假设 $Y$ 严格局部. 在 $n\geq 1$ 的阿贝尔情形, $g$ 与 $g^{\ast}(F)$ 的 Leray 谱序列, 连同在 $0<i\leq n-1$ 时成立的 $R^{i}g_{\ast}(g^{\ast}(F))=0$ 以及 $F\xrightarrow{\sim}g_{\ast}g^{\ast}(F)$ (这些关系表述了 $g$ 的 $(n-1)$-无上同调性假设), 给出典范正合列
$$
0\longrightarrow H^{n}(Y,F)\longrightarrow
H^{n}(X,g^{\ast}(F))\longrightarrow
H^{0}(Y,R^{n}g_{\ast}(g^{\ast}(F))){\quoi},
$$
这表明只须证明 $\xi$ 在第三项中的像为零. 这可以逐茎验证, 由 \hyperrefext[VIII][VIII.thm5.2]{5.2}, 就化为 $Y$ 严格局部的情形. 当 $n=0$ 且 $F$ 为集合层时同样处理: 对 $g$ 的假设蕴含 $F\rightarrow g_{\ast}g^{\ast}(F)$ 为单射, 而 $\xi$ 可视为第二个层的一个截面, 要证明它来自第一个层的截面. 这也可以逐\pageoriginale 茎验证, 然后应用 \hyperrefext[VIII][VIII.rem5.3]{5.3}. 非交换情形 $n=1$ 的论证基本相同, 只须使用正合列 \hyperrefext[XII][XII.prop3.2]{3.2}
$$
H^{1}(Y,F)\longrightarrow H^{1}(X,g^{\ast}(F))\longrightarrow
H^{0}(Y,R^{1}g_{\ast}(g^{\ast}(F))){\quoi},
$$
以及 \hyperrefext[VIII][VIII.rem5.3]{5.3}.

因此假设 $Y$ 严格局部, 从而 $X$ 拟紧且拟分离. 态射 $g:X\rightarrow Y$ 满足 \hyperrefext[XII][XII.cor6.10]{6.10} 的条件, 因而可将问题化为下述情形: 对 $Y$ 的每个不等于 $Y$ 的闭子集 $Y_1$, $\xi$ 在 $X_1=\fprod{X}{Y_1}{Y}$ 上的限制都属于 $H^{n}(Y_1,F|Y_1)\rightarrow H^{n}(X_1,g^{\ast}(F)|X_1)$ 的像; 并且只须在这一情形下验证, 可以找到有限态射 $f:Y'\rightarrow Y$, 使其在 $Y$ 中的像包含 $Y$ 的一个非空开子集, 且 $\xi$ 在 $X'=\fprod{X}{Y'}{Y}$ 上的逆像 $\xi'$ 属于 $H^{n}(Y',F')\rightarrow H^{n}(X',{g'}^{\ast}(F'))$ 的像. 注意, 根据 \hyperrefext[VIII][VIII.prop5.5]{5.5} 及 \hyperrefext[VIII][VIII.rem5.8]{5.8}, \loccit 的条件 $2^{\circ}$) 对这个态射 $f:Y'\rightarrow Y$ 自动满足.

为找到这样的 $f$, 可设 $Y$ 既约, 并考虑 $Y$ 的一个极大点 $y$. 设 $\bar{y}$ 为 $k(y)$ 的一个可分闭包所对应的几何点, 则 $\bar{y}$ 是有限整 $Y$-概形 $Y_i$ 的开子集 $U_i$ 的滤过投射极限, 每个这样的概形在 $Y$ 中的像都包含 $y$ 的一个邻域. 利用对 $\xi$ 的假设可知, 存在指标 $i$, 使 $\xi$ 在 $\fprod{X}{U_i}{Y}$ 上的逆像在 $n\geq 1$ 时为零, \resp 在 $n=0$ 时属于 $H^{0}(U_i,F_{U_i})$ 的像.

我们\pageoriginale 取 $Y'$ 为 $Y_i$, 并注意可设 $U'=U_i$ 是 $y$ 在 $Y$ 中的某个仿射开邻域 $U$ 的逆像. 于是 $\xi$ 在 $\fprod{X}{Z'}{Y}$ 上的逆像属于 $H^{n}(Z',F_{Z'})$ 的像 (其中 $Z'=Y'-U'$ 为 $Z=Y-U$ 在 $Y'$ 中的逆像). 用 $Y'$ 代替 $Y$ 后, 问题便化为以下情形: 存在 $Y$ 的仿射开子集 $U$, 其补集为 $Z$, 使 $\xi$ 在 $X_U$ 与 $X_Z$ 上的限制分别属于 $H^{n}(U,F_U)$ 的像与 $H^{n}(Z,F_Z)$ 的像; 此外, 当 $n\geq 1$ 时, 可以假设这些限制\emph{为零}.

当 $n\geq 1$ 时, 我们将证明在这些条件下 $\xi=0$. 为简化记号, 用 $Y'$ 代替 $X$, 以 $U'$ 与 $Z'$ 表示 $U$ 与 $Z$ 在 $Y'$ 中的逆像, 并以 $F'$ 表示逆像 $g^{\ast}(F)$. 设 $i:U\rightarrow Y$, $j:Z\rightarrow Y$, $i':U'\rightarrow Y'$ 及 $j':Z'\rightarrow Y'$ 为包含映射,
$$
\xymatrix@C=2cm@R=1.5cm{
U'  \ar[r]^{i'} \ar[d]^{g_U}& Y' = X  \ar[d]^{g} & Z'  \ar[l]_{j'} \ar[d]^{g_Z}\\
U   \ar[r]^{i} & Y & Z \ar[l]_{j},
}
$$
并考虑层的正合列
$$
0\longrightarrow F'_{U',Y'}\longrightarrow F'\longrightarrow
F'_{Z',Y'}\longrightarrow 0{\quoi}.
$$
% 这一记号的出处?
由于 $\xi$ 在 $H^{n}(Y',F'_{Z',Y'})=H^{n}(Z',F'_{Z'})$ 中的像为零, $\xi$ 来自 $H^{n}(Y',F'_{U',Y'})$ 的一个元素, 而该元素在 $U'$ 上的限制等于 $\xi$ 的限制, 因此为零. 由于 $F'_{U',Y'}$ 同构于 $F_{U,Y}$ 的逆像, 问题便\pageoriginale 化为形如 $F_{U,Y}$ 的层的情形; 因而可设 $F_Z=0$ (这计入了假设 $\xi|Z'=0$). 注意, 若 $F\rightarrow G$ 为阿贝尔 $\LL$-挠层 (\resp 群层) 的单态射, 则只须证明 $\xi$ 在 $H^{n}(Y',G')$ 中的像 $\eta$ 属于 $H^{n}(Y,G)$ 的像, \ie 为零 \hyperrefext[XII][XII.cor6.6]{6.6}. 这允许我们用 $i_{\ast}(F_U)$ 代替 $F$, 因为假设 $F_Z=0$ 蕴含典范同态 $F\rightarrow i_{\ast}(F_U)$ 为单射. 现在, 由于 $g$ 关于 $\LL$ 局部 $(n-1)$-无上同调, 从而关于 $\LL$ 为 $0$-无上同调的 (因为 $n\geq 1$), $i_{\ast}(F_U)$ 在 $Y'$ 上的逆像可等同于 $i'_{\ast}(F'_{U'})$. 因而只须证明 $\xi$ 在 $H^{n}(Y',i'_{\ast}(F'_{U'}))$ 中的像 $\eta$ 为零. $\eta$ 在 $U'$ 上的限制为零, 所以只须验证同态
\begin{equation*}
H^{n}(Y',i'_{\ast}(F'_{U'}))\longrightarrow H^{n}(U',F'_{U'})\tag{$\ast$}
\end{equation*}
为单态射. 在阿贝尔情形, 使用 $i'$ 与 $F'_{U'}$ 的 Leray 谱序列:
$$
E^{p,q}_2=H^p(Y',R^qi'_{\ast}(F'_{U'}))\Longrightarrow
H^{\ast}(U',F'_{U'}){\quoi},
$$
以及关系
\begin{equation*}
E^{p,q}_2=0\; \mbox{\rm 当 } 0<p\leq n-1\quav q\leq n-1{\quoi}.\tag{${\ast}{\ast}$}
\end{equation*}
为验证后一关系, 注意关于 $\LL$ 对 $g$ 的局部 $(n-1)$-无上同调性假设蕴含同构
$$
R^{q}i'_{\ast}(F'_{U'})\simeq g^{\ast}(R^{q}i_{\ast}(F_U))\;\mbox{\rm 当 } q\leq n-1{\quoi},
$$
又\pageoriginale 由于 $g$ 关于 $\LL$ 为 $(n-1)$-无上同调的, 因而对 $Y$ 上的每个阿贝尔 $\LL$-挠层 $G$, 当 $p\leq n-1$ 时有 $H^{p}(Y,G)\xrightarrow{\sim}H^{p}(Y',G')$, 结合在 $p>0$ 时成立的 $H^{p}(Y,G)=0$, 就得到 $({\ast}{\ast})$. 同理得到
\begin{equation*}
E^{0,q}_2\simeq H^{0}(Y,R{}^qi_{\ast}(F_U))\;\mbox{\rm
当}\quad q\leq n-1{\quoi}.\tag{${\ast}{\ast}{\ast}$}
\end{equation*}
所考察的 Leray 谱序列于是给出一个正合列, 构成下列交换图的第二行:
$$
\xymatrix{
**[r]H^{0}(Y,R^{n-1} i_{\ast}(F_U))  \ar[r] \ar[d]^{\text{\rotatebox{90}{$\sim$}}}&
**[r]H^n(Y,i_{\ast} (F_U)) = 0 \ar[d]\\
**[r]H^{0}(Y', R^{n-1} i'_{\ast}(F'_{U'}))  \ar[r] &
**[r]H^{n}(Y',i'_{\ast}(F'_{U'}))  \ar[r] & **[r]H^{n}(U',F'_{U'}),
}
$$
其中第一条竖箭头由 $({\ast}{\ast}{\ast})$ 为双射. 该图表明 $(\ast)$ 为单射, 从而完成这一情形的证明. 在非阿贝尔情形 $n=1$, 使用 $n=1$ 时相同的图:
$$
\xymatrix{
**[r]H^{0}(Y, i_{\ast}(F_U))  \ar[r] \ar[d]^{\text{\rotatebox{90}{$\sim$}}}&
**[r]H^1(Y,i_{\ast} (F_U)) = 0 \ar[d]\\
**[r]H^{0}(Y', i'_{\ast}(F'_{U'}))  \ar[r] &
**[r]H^{1}(Y',i'_{\ast}(F'_{U'}))  \ar[r] & **[r]H^{1}(U',F'_{U'}),&
}
$$
\hyperrefext[XII][XII.prop3.2]{3.2}, 其中由 $g$ 的局部及整体 $0$-无上同调性假设, 第一条竖箭头仍为双射; 因而再次得出 $(\ast)$ 在这一情形为单射.

最后考虑 $n=0$ 的情形. 根据 \sisi{\qquad}{\qquad}, 指定 $F'$ 的一个截面 $\xi$, 相当于指定 $F'_{U'}$ 的截面 $\xi_{U'}$ 与 $F'_{Z'}$ 的截面 $\xi_{Z'}$,
% 底本此处缺少引用.
使 $\xi_{Z'}$ 经自然同态
$$
F'_{Z'}\longrightarrow {j'}^{\ast}(i'_{\ast}(F'_{U'}))
$$
所得的像,\pageoriginale 等于右侧由 $\xi_{U'}$ 诱导的截面. 由假设, 分别存在 $F_Z$ 与 $F_U$ 的截面 $\eta_Z$ 与 $\eta_U$, 诱导 $\xi_{Z'}$ 与 $\xi_{U'}$. 只须验证它们满足类似的相容条件, 即它们是 $F$ 的同一个截面 $\eta$ 的限制. 必须验证 $j^{\ast}(i_{\ast}(F_U))$ 的两个截面相等. 由于 $Z'\rightarrow Z$ 满射 ($g$ 由假设为 $(-1)$-无上同调的, 因而满射), 只须验证
$$
g_Z{}^{\ast}(j^{\ast}(i_{\ast}(F_U)))
=j'{}^{\ast}(g^{\ast}(i_{\ast}(F_U)))
$$
的两个相应截面相等. 经基变换同态
\begin{equation*}
g^{\ast}(i_{\ast}(F_U))\longrightarrow
i'_{\ast}(g_U{}^{\ast}(F_U))=i'_{\ast}(F'_{U'}){\quoi},\tag{$\ast$}
\end{equation*}
所考察的层映到 ${j'}^{\ast}(i'_{\ast}(F'_{U'}))$, 而后一个同态为单射, 因为由 $g$ 局部 $(-1)$-无上同调的假设, $(\ast)$ 为单射. 因而只须验证 ${j'}^{\ast}(i'_{\ast}(F'_{U'}))$ 中作为上述两个截面之像的两个截面相等. 但立即可见, 它们正是上面所考察的两个截面, 而由假设它们相等. 证明完毕.

\begin{enonce*}{注 1.18}\etiquette[1.18]{XV.rem1.18}
\hyperref[XV.thm1.15]{1.15} 有一个加强形式, 当假设 $Y$ 局部 Noether 时, 它给出 \hyperref[XV.thm1.15]{1.15} 的一个简单得多的证明. 除了对 $g$ 与 $F$ 的原有假设之外, 假设 $Y$ 有有限个几何点 $f_i:\bar{y}_i\rightarrow Y$, 使态射
$$
F\longrightarrow \prod_if_{i\ast}(f_i{}^{\ast}(F))
$$
为\pageoriginale 单射 (当 $Y$ 为 Noether 概形且 $F$ 可构造时, 这一假设自动满足 \hyperrefext[IX][IX.prop2.14]{2.14}). 我断言, 在这些条件下, 即使只限于几何点 $\bar{y}_i$, 判据 \hyperref[XV.thm1.15]{1.15} 仍然成立. 事实上, 由 \hyperrefext[XII][XII.cor6.6]{6.6}, 可假设 $F$ 本身形如 $f_{\ast}(M_{\bar{y}})$, 其中 $M$ 为集合, \resp 群, \resp 阿贝尔 $\LL$-挠群, 且 $f:\bar{y}\rightarrow Y$ 为几何点. 设 $f':X'=X_{\bar{y}}\rightarrow X$ 为典范态射. 为明确起见, 假设 $n\geq 1$, 则对 $g$ 的局部假设蕴含
$$
g^{\ast}(f_{\ast}(M_{\bar{y}}))\simeq f'_{\ast}(M_{X'})\sisi{,\;}{\quav}
R^{i}f'_{\ast}(M_{X'})=0\quad\mbox{\rm 当}\quad 1\leq i\leq
n-1{\quoi},
$$
第二个关系来自其左侧同构于 $g^{\ast}(R^{i}f_{\ast}(M_{\bar{y}}))$, 而这一项为零, 因为显然当 $i\geq 1$ 时有 $R^{i}f_{\ast}(M_{\bar{y}})=0$. 由这些关系, 应用 $f'$ 的 Leray 谱序列可得
$$
H^{n}(X,g^{\ast}(F))\simeq H^{n}(X,f'_{\ast}(M_{X'}))\longrightarrow
H^{n}(X',M_{X'})
$$
为单射. 这证明, 若 $H^{n}(X,g^{\ast}(F))$ 的元素 $\xi$ 在 $H^{n}(X_{\bar{y}},M_{X_{\bar{y}}})$ 中诱导的元素为零, 则它本身为零. 要由此得到 $Y$ 局部 Noether 时 \hyperref[XV.thm1.15]{1.15} 的一个简化证明, 注意只须处理 $Y$ 为 Noether 的情形; 再利用 \hyperrefext[IX][IX.cor2.7.2]{2.7.2} 及 \hyperrefext[VII][VII.sec5]{5} 作一次简单的极限过渡, 就可化为 $F$ 也可构造的情形, 从而适用刚证明的陈述.
\end{enonce*}

\section{光滑态射的局部无上同调性}\etiquette[2]{XV.sec2}\pageoriginale

定理如下:

\begin{enonce*}{定理 2.1}\etiquette[2.1]{XV.thm2.1}
设 $g:X\rightarrow Y$ 为光滑态射, $\LL\subset
\PP$ 为 $X$ 的剩余特征集合的补集. 则 $g$ 关于 $\LL$ 局部无上同调且局部 $1$-非球面.
\end{enonce*}

先给出如下推论:

\begin{enonce*}{推论 2.2}\etiquette[2.2]{XV.cor2.2}
设 $g:\EE^m_Y\rightarrow Y$ 为仿射 $m$ 维空间的结构态射. 则 $g$ 关于 $Y$ 的剩余特征集合的补集 $\LL$ 无上同调且 $1$-非球面.
\end{enonce*}

事实上, 对 $m$ 归纳, 利用 \hyperref[XV.prop1.9]{1.9} (i), 就可化为 $m=1$ 的情形. 由 \hyperref[XV.thm2.1]{2.1}, $g$ 关于 $\LL$ 局部 $n$-无上同调且局部 $1$-非球面, 所以满足 \hyperref[XV.cor1.16]{1.16} 的\emph{局部}条件. 因而只须证明 $\EE^1_Y$ 的几何纤维关于 $\LL$ 无上同调且 $1$-非球面; 即当 $k$ 为特征为 $p\notin\LL$ 的可分闭域时, $\Spec k[t]$ 关于 $\LL$ 无上同调且 $1$-非球面. 无上同调性的断言来自 \hyperrefext[IX][IX.thm4.6]{4.6}, 而 $1$-非球面性的断言是众所周知的. 更一般地, 设 $X'$ 为 $X=\PP^1_k$ 的正规不可约 Galois 覆盖, 在点 $\infty$ 以外不分歧, 在该点温分歧. 我们证明该覆盖的次数 $n$ 为 $1$. Hurwitz 公式给出
$$
2(g'-1)=2n(g-1)+\deg \mathfrak{d}_{X'/X},
$$
其中 $g$ 与 $g'$ 为 $X$ 与 $X'$ 的亏格, $\mathfrak{d}_{X'/X}$ 为余不同除子. 有 $g=0$, 而温分歧蕴含 $\deg \mathfrak{d}_{X'/X}\leq n-1$; 因而 $2(g'-1)\leq-(n+1)$, 所以 $n\leq2(1-g')-1$. 由于 $g'\geq0$, 可得 $n\leq1$, 从而 $n=1$, \cqfd

\medskip
\noindent
\hyperref[XV.thm2.1]{2.1} 的\emph{证明}. 该断言对于平展拓扑是 $X$ 与 $Y$ 上的局部性质, 因而可假设 \sisi{(}{}\SGA 1 II 1.1\sisi{)}{} $g$ 为仿射 $m$ 维空间 $\EE^m_Y$ 的结构态射, 且 $Y$ 仿射. 由局部无上同调性的传递性 \hyperref[XV.prop1.13]{1.13} (i) 及恒等式 $\EE^m_Y=\EE^{m-1}_{\EE^1_Y}$, 立即可化为 $m=1$, 即 $X=\Spec\Oo_Y[t]$ 的情形.

应用 \hyperref[XV.cor1.18]{1.18}. 只须证明, 对应于一对几何点 $\bar x,\bar y$ 的严格局部概形之间的态射 $\widetilde g:\widetilde X\rightarrow\widetilde Y$ 的几何纤维, 关于 $\LL$ 无上同调且关于 $\LL$ 为 $1$-非球面的; 这里 $\bar x$ 在纤维 $X_{\bar y}$ 中为闭点. 显然可设 $Y=\widetilde Y$. 设 $Y'\rightarrow Y$ 为有限、局部且满射的态射, 从而 $Y'$ 严格局部; 设 $\bar y'$ 为 $Y'$ 中位于 $\bar y$ 之上的点, $\bar x'$ 为 $X'$ 中同时位于 $\bar x$ 与 $\bar y'$ 之上的点. 由于 $X'\rightarrow X$ 有限, $X'$ 在 $\bar x'$ 处的严格局部化恰为 $\widetilde X'=\fprod{\widetilde X}{Y'}{Y}$. 因而 $\widetilde X'/Y'$ 的每个几何纤维, 都由 $\widetilde X/Y$ 的某个几何纤维作底域的代数 (必为纯不可分) 扩张所得; 此外, $\widetilde X/Y$ 的每个几何纤维之上至少有 $\widetilde X'/Y'$ 的一个几何纤维. 所以, 由 \hyperrefext[VIII][VIII.thm1.1]{1.1}, 可用 $Y'$ 代替 $Y$. 适当选取 $Y'$, 就化为 $k(\bar x)=k(\bar y)$ 的情形. 此时 $\bar x$ 为纤维 $X_{\bar y}=\Spec k(\bar y)[t]$ 的一个有理点, 设为点 $t=a^0$. 将 $a^0$ 提升为 $\Oo_{\sisi{Y}{\widetilde Y}}$ 的截面 $a$, 并作“平移” $t_1=t-a$; 问题便化为 $\bar x$ 是 $\Spec k(\bar y)[t]$ 的点 $\{t=0\}$ 的情形. 令 $A=\Gamma(\sisi{Y}{\widetilde Y},
\Oo_{\sisi{Y}{\widetilde Y}})$, 则 $A$ 为严格局部环.

\begin{enonce*}[remark]{记号 2.3}\etiquette[2.3]{XV.not2.3}
设 $A$ 为 Hensel 局部环. 以 $A\{t\}$ 表示 $A[t]$ 在由 $\rad A$ 与 $t$ 生成的素理想处的 Hensel 化.
\end{enonce*}

\begin{enonce*}{推论 2.4}\etiquette[2.4]{XV.cor2.4}
设 $A$ 为 Hensel 环, $A'$ 为有限 $A$-代数. 则
$$
A'\{t\}\simeq A'\otimes_A A\{t\}{\quoi}.
$$
\end{enonce*}

这直接来自 \hyperrefext[VIII][VIII.subsec4.1]{4.1}. 注意, 若 $A$ 严格局部, 则 $A\{t\}$ 也严格局部.

定理的证明化为如下引理:

\begin{enonce*}{引理 2.5}\etiquette[2.5]{XV.lem2.5}
设 $A$ 为剩余特征为 $p$ 的 Hensel 环. 则 $\Spec A\{t\}$ 在 $\Spec A$ 上的每个几何纤维都关于 $\LL=\PP-\{p\}$ 无上同调且 $1$-非球面.
\end{enonce*}

由于 $A\{t\}$ 是 $A[t]$ 的 Hensel 化, 纤维 $\Spec(A\{t\}\otimes_A K)$ (其中 $K$ 为 $A$ 的一个剩余域) 是“直线” $\Spec K[t]$ 上有限型仿射概形 $X_i$ 的投射极限,\pageoriginale 即仿射代数曲线的极限. 因而 $\Spec A\{t\}$ 的每个几何纤维, 都是可分闭域 $\bar K$ 上的仿射曲线 $(X_i)_{\bar K}$ 的极限. 由 \hyperrefext[IX][IX.cor5.7]{5.7} 及 \hyperrefext[VII][VII.thm5.7]{5.7}, 几何纤维的上同调维数至多为 $1$. 因而只须验证几何纤维关于 $\LL$ 为 $1$-\emph{非球面}的; 即这样的几何纤维 $\bar Z$ 连通且非空, 并对每个有限 $\LL$-群 $G$, 都有 $H^1(\bar Z,G)=0$.

写成 $A=\varinjlim B_\alpha$, 其中 $B_\alpha$ 为 $\ZZ$ 上的有限型环. 态射 $B_\alpha\rightarrow A$ 通过 $B_\alpha$ 在由 $\rad A$ 诱导的素理想处的 Hensel 化 $A_\alpha$ 分解, 因而 $A=\varinjlim A_\alpha$. 也容易验证 $A\{t\}=\varinjlim A_\alpha\{t\}$. 设 $\bar y$ 为 $\Spec A$ 的一个几何点, $X_{\bar y}$ 为 $S=\Spec A\{t\}$ 在 $\Spec A$ 上于 $\bar y$ 处的几何纤维, $\bar y_\alpha$ 为 $\bar y$ 诱导的 $\Spec A_\alpha$ 的点, $X_{\alpha,\bar y_\alpha}$ 为 $\Spec A_\alpha\{t\}$ 在 $\Spec A_\alpha$ 上于 $\bar y_\alpha$ 处的几何纤维. 有 $X_{\bar y}=\varprojlim X_{\alpha,\bar y_\alpha}$. 将 \hyperref[XV.prop1.9]{1.9} (ii) 应用于态射 $X_{\alpha,\bar y_\alpha}\rightarrow\bar y_\alpha$, 便将 \hyperref[XV.lem2.5]{2.5} 的证明化为 $A=A_\alpha$ 的情形, 即 $A$ 为 $\ZZ$ 上某个有限型代数在给定素理想处的 Hensel 化的情形. 以下均作这一假设. 回忆这样的代数是\emph{优环}（excellent） \sisi{(}{}\EGA IV 7.8.3 (iii), 7.8.6 (i)\sisi{)}{}, 因而 $A_\alpha$ 也是优环 (这容易从定义推出; \cf \EGA IV 18.7.6).

我们验证 $0$-无上同调性. (更一般的结论见附录.) 由于 $A\{t\}/A$ 有截面“$t=0$”, 几何纤维\pageoriginale $\bar Z$ 非空. 为证明其连通, 将 $A$ 换成一个商环 \sisi{(}{}\hyperref[XV.cor2.4]{2.4}\sisi{)}{}, 就可化为 $A$ 整且 $\bar Z$ 为 $\Spec A$ 的泛点之上的几何纤维的情形. 只须证明, 对 $A$ 的分式域 $K$ 的每个有限可分扩张 $K'$, 概形 $Z_{K'}$ 连通, 其中 $Z=\Spec(A\{t\}\otimes_A K)$. 设 $A'$ 为 $A$ 在 $K'$ 中的正规化 (它是有限 $A$-代数, \sisi{(}{}\EGA IV 7.8.2\sisi{)}{}), $Z'$ 为 $\Spec A'\{t\}$ 在 $A'$ 上的泛纤维. 由 \hyperref[XV.cor2.4]{2.4}, 有 $A'\{t\}\simeq A'\otimes_A A\{t\}$, 因而 $Z_{K'}\simeq Z'$. 用 $A'$ 代替 $A$. 此时只须证明 $Z$ 整, 而为此只须证明 $A\{t\}$ 整. 由于 $A$ 正规, $A[t]$ 正规, $A\{t\}$ 也正规 \sisi{(}{}\SGA 1 I 9.5 (i)\sisi{)}{}. 断言得证.

因而 \hyperref[XV.lem2.5]{2.5}, 从而 \hyperref[XV.thm2.1]{2.1} 的证明, 化为以下引理, 将在下一节予以证明:

\begin{enonce*}{引理 2.6}\etiquette[2.6]{XV.lem2.6}
设 $A$ 为 Hensel 优环, $\bar Z$ 为 $\Spec A\{t\}/\Spec A$ 的一个几何纤维. 则 $\bar Z$ 的每个主 Galois 覆盖 $\bar Z'$, 只要其群 $G$ 的阶与 $A$ 的剩余特征 $p$ 互素, 就都是平凡的.
\end{enonce*}

\section{主要引理的证明}\etiquette[3]{XV.sec3}\pageoriginale

将 $A$ 换成一个有限 $A$-代数 $A'$ (\cf \hyperref[XV.cor2.4]{2.4}), 就将断言 \hyperref[XV.lem2.6]{2.6} 化为 $A$ 整且 $\bar{Z}$ 为 $\Spec A\{t\}$ 在 $\Spec A$ 上的几何泛纤维的情形. 设 $K$ 为 $A$ 的分式域, 并设 $Z=\Spec A\{t\}\otimes_A K$. 若 $\bar{Z}'/\bar{Z}$ 为主 Galois 覆盖, 可将其下降到某个 $Z_{K'}$ 上, 其中 $K'$ 为 $K$ 的适当的有限可分扩张. 将 $A$ 换成其在 $K'$ 中的正规化 (\hyperref[XV.cor2.4]{2.4}), 它是有限 $A$-代数 \sisi{(}{}\EGA IV 7.8.3 (vi)\sisi{)}{}, 便将问题化为证明下述形式的引理:

\begin{enonce*}{引理 3.1}\etiquette[3.1]{XV.lem3.1}
用 \hyperref[XV.lem2.6]{2.6} 的记号, 假设 $A$ 正规, 并设 $Z'/Z$ 为主平展 Galois 覆盖, 其群 $G$ 的阶与 $A$ 的剩余特征 $p$ 互素. 则 $Z'$ 由 $K$ 的一个 Galois 扩张 $K'$ 诱导; 即 \emph{(\SGA 1 IX 6.2)}, 由 $Z'$ 导出的 $\bar{Z}$ 的覆盖 $\bar{Z}'$ 平凡.
\end{enonce*}

\noindent
\emph{证明}. 设 $B'$ 为 $A\{t\}$ 在 $Z'$ 的有理函数环中的正规化, 并考虑 $A\{t\}$ 中使 $A\{t\}$-代数 $B'$ 分歧的高度为 $1$ 的素理想 $P$. 由于 $Z'/Z$ 平展, $P$ 不在 $\Spec A$ 的泛点之上; 即 $\mathfrak{p}=P\cap A\neq 0$. 因而 $A\{t\}/\mathfrak{p}A\{t\}=(A/\mathfrak{p})\{t\}$ 整, 且维数至多为 $\dim A$. 由于 $A\{t\}$ 为优环 \sisi{(}{}\EGA IV 7.8.6 (i) 及 18.7.6\sisi{)}{}, 从而为悬链环, 有 $\dim \Spec A\{t\}/P=(\dim A+1)-1=\dim A\geq
\dim A\{t\}/\mathfrak{p}A\{t\}$, 因此\pageoriginale $P=\mathfrak{p}A\{t\}$, 因为 $P\supset \mathfrak{p}A\{t\}$.

于是设 $\mathfrak{p}_1,\ldots,\mathfrak{p}_m$ 为 $A$ 中满足 $B'/A\{t\}$ 在理想 $\mathfrak{p}_iA\{t\}$ 处分歧的高度为 $1$ 的素理想; 设 $s$ 为 $B'$ 在位于 $\mathfrak{p}_iA\{t\}$ 之上的素理想处的惯性群 \sisi{(}{}\SGA IV 2\sisi{)}{} $G_i\subset G$ 的阶的最小公倍数. 这是整除 $G$ 的阶的整数, 因而与 $p$ 互素. 取一个在每个 $\mathfrak{p}_i$ 处都具有 $1$ 阶零点的元素 $f\in A$ (这样的 $f$ 存在, 例如可由 Bourbaki, Alg. Comm., Chap. II, \S 3, 命题 17 的推论得出). 将 $A$ \sisi{(}{}\hyperref[XV.cor2.4]{2.4}\sisi{)}{} 换成它在 $K$ 的有限扩张 $K_1=K[x]/(x^s-f)$ 中的正规化 $A_1$, 并将 $B'$ 换成 $A_1\otimes_A B'$ 的正规化; 由 Abhyankar 引理 \sisi{(}{}\SGA 1 X 3.6\sisi{)}{}, 就可化为 $B'$ 在 $A\{t\}$ 的任何高度为 $1$ 的素理想处都不分歧的情形.

此时代数 $B'/A\{t\}$ 在 $\Spec A$ 的每个余维为 $1$ 的点处的 $\Spec A\{t\}$ 的纤维上也平展. 事实上, 设 $\mathfrak{p}$ 为 $A$ 的高度为 $1$ 的素理想. 由于 $A$ 正规, $A_{\mathfrak{p}}$ 为离散赋值环, 从而正则, 所以 $A\{t\}\otimes_A A_{\mathfrak{p}}$ 也正则; 事实上它是 $A_{\mathfrak{p}}[t]$ 上平展概形的极限. 由假设, 代数 $B'/A\{t\}$ 在 $\Spec A\{t\}\otimes_A A_{\mathfrak{p}}$ 的每个余维 $\leq 1$ 的点处平展, 因而由 Zariski--Nagata 纯性定理 \sisi{(}{}\SGA 1 X 3.1\sisi{)}{}, 它处处平展; 证明可参看例如 \SGA 2, X 3.4.

将 $A$ \sisi{(}{}\hyperref[XV.cor2.4]{2.4}\sisi{)}{} 换成适当的有限 $A$-代数, 就可进一步假设 $A\{t\}/tA\{t\}\simeq A$ 上的代数 $B'/tB'$ 在 $\Spec A$ 的泛点之上完全分裂. 因而只须证明, 在这些条件下 $A\{t\}$ 上的代数 $B'$\pageoriginale “平凡”; 又由于 $A\{t\}$ 为 Hensel 环, 只须证明 $B'$ 在 $A\{t\}$ 上平展. 事实上, $B'/tB'$ 在 $\Spec A$ 的泛点之上分裂, 蕴含它在 $\Spec(A)$ 上分裂 \sisi{(}{}\SGA 1 I 10.1\sisi{)}{}; 因而 $B'$ 在 $A\{t\}$ 的极大理想之上的剩余扩张平凡, 然后应用 \EGA IV 18.5.14. 因此问题化为以下引理:

\begin{enonce*}{引理 3.2}\etiquette[3.2]{XV.lem3.2}
设 $B$ 为优正规局部环, $0\neq t\in \rad B$ 且使 $B/tB$ 正规. 设 $B'$ 为包含 $B$ 的有限整正规 $B$-代数, 满足:
\begin{enumerate}
\renewcommand{\theenumi}{\roman{enumi}}
\renewcommand{\labelenumi}{(\theenumi)}
\item $B/tB$ 上的代数 $B'/tB'$ 在 $\Spec(B/tB)$ 的极大点处完全分裂.

\item $B$ 上的代数 $B'$ 在 $\Spec(B/tB)$ 的余维为 $1$ 的点处平展.
\end{enumerate}
则 $B'$ 在 $B$ 上平展.
\end{enonce*}

\noindent
\emph{证明}. 由 (ii), 断言在 $\dim B\leq 2$ 时, 即在 $\dim B/tB\leq 1$ 时是显然的. 假设 $B'/B$ 不平展; 设 $x$ 为 $\Spec B/tB$ 中使 $B'/B$ 不平展的点所成集合的一个极大点, $C$ 为 $\Spec B$ 在 $x$ 处的局部环. 则 $C$ 仍为优正规环 (\EGA IV 7.8.2), 有 $0\neq t\in\rad C$, 并且 $C$-代数 $C'=C\otimes {}_BB'$ 满足 (i) 与 (ii). 因而可设 $B=C$; 即除了原点可能例外, $B'$ 在 $\Spec B/tB$ 的每个点之上都在 $B$ 上平展. 由于 $B/tB$ 正规\pageoriginale (事实上单分支就足够), 条件 (i) 蕴含代数 $B'/tB'$ 在原点以外完全分裂. 此外, 可设 $\dim B\geq 3$. 于是问题化为以下引理:

\begin{enonce*}{引理 3.3}\etiquette[3.3]{XV.lem3.3}
设 $B$ 为优正规局部环, $0\neq t\in \rad B$, 并假设 $\dim B\geq 3$. 设 $B'$ 为包含 $B$ 的有限整正规 $B$-代数, 使 $B/tB$-代数 $B'/tB'$ 在原点以外完全分裂. 则 $B'$ 在 $B$ 上平展.
\end{enonce*}

\noindent
\emph{证明}. 设 $X=\Spec B$, $X_1=X-\{x_0\}$, 其中 $x_0$ 为闭点; 设 $i:U\rightarrow X$ 为开浸入, 其像由使代数 $B'/B$ 平展的点 $x\in X$ 组成. 则 $(X-U)\cap V(t)=\{x_0\}$ 的维数为 $0$. 由于 $V(t)$ 由\emph{一个}方程定义, 可得 $\dim(X-U)\leq 1$ (\EGA $0_{\text{\rm IV}}$ 16.2.5); 因而 (由于 $X$ 为悬链概形), $x\in X-U\Rightarrow \dim \mathcal{O}_{X,x}\geq \dim B-1\geq 2$; 即 $\codim(X-U,X)\geq 2$.

设 $X'=\Spec(B')$, $U'$ 为 $U$ 在 $X'$ 中的逆像. 由于 $X$ 普遍悬链, 关系 $\codim(X-U,X)\geq 2$ 蕴含 $\codim(X'-U',X')\geq 2$ (\EGA IV 5.6.10). 由于 $X'$ 正规, 因而有
$$
B'=\Gamma(X',\Oo_{X'})\xrightarrow{\sim}
\Gamma(U',\Oo_{U'}){\quoi}.
$$
设 $\mathcal{B}'=\widetilde{B}'$ 为在 $X$ 上定义 $X'$ 的 $\Oo_X$-代数层, $\mathcal{C}$ 为其在 $U$ 上的限制. 上式于是等价于
$$
\mathcal{B}'\xrightarrow{\sim}i_{\ast}(\mathcal{C}){\quoi},
$$
因而\pageoriginale \hyperref[XV.lem3.3]{3.3} 化为下述陈述 (其中不再出现 $B'$): 设 $B$, $t$ 如 \hyperref[XV.lem3.3]{3.3}, 设 $X=\Spec(B)$, $Y=\Spec(B/tB)$, $U$ 为 $Y_1=Y-\{x_0\}$ 在 $X_1=X-\{x_0\}$ 中的一个开邻域, $U'$ 为 $U$ 的一个平展覆盖, 其在 $Y_1$ 上的限制完全分裂. 若 $\mathcal{C}$ 表示定义 $U'$ 的 $\Oo_U$-代数层, 则 (以 $i:U\rightarrow X$ 表示典范浸入) $i_{\ast}(\mathcal{C})$ 为 $X$ 上的凝聚\emph{平展}代数 (或者, 等价地, 考虑到上面指出的对 $x\in X-U$ 成立的关系 $\prof \Oo_{X,x}\geq 2$ 以及 \EGA IV 5.10.5, $U'$ 可延拓为 $X$ 的一个\emph{平展}覆盖 $X'$).

设 $\widehat{X}=\Spec\widehat{B}$, 其中 $\widehat{B}$ 为 $B$ 的完备化. 则 $\widehat{B}$ 正规 (\EGA IV 7.8.3 (v)), 且 $\widehat{B}/B$ 忠实平坦. 设
$$
\xymatrix@C=1.5cm@R=1cm{
\hat{U}      \ar[r]^{\sisi{\hat{i}}{\hat{i}}} \ar[d]^{f_1}& \hat{X} \ar[d]^{f}\\
U            \ar[r]^{i} &  X
}
$$
为由 $\widehat{X}\rightarrow X$ 导出的笛卡儿图. 由于 $f$ 忠实平坦, 函子 $i_{\ast}$ 与沿 $f$ 的基变换可交换 \sisi{(}{}\EGA IV 2.3.1\sisi{)}{}, 因而问题化为 \sisi{(}{}\EGA IV 17.7.1\sisi{)}{} 证明 $\hat{i}_{\ast}(\widehat{\mathcal{C}})$ 是 $\widehat{X}$ 上的平展代数. 所以可用 $\widehat{B}$ 代替 $B$; 即可假设 $B$ \emph{完备}. 在这一情形, 我们将证明 $U'$ 为 $U$ 的一个完全分裂覆盖 (这显然蕴含它可延拓为 $X$ 的平展覆盖, 从而完成证明). 这一断言现在直接来自“局部 Lefschetz 理论” \sisi{(}{}\SGA\pageoriginale 2 X 2.1 (i) 及 2.5\sisi{)}{}, 只须再次考虑在 $\dim \{x\}=1$ 时成立的关系 $\prof \Oo_{X,x}\geq 2$, 便立即可见.

\begin{enonce*}[remark]{注 3.4}\etiquette[3.4]{XV.rem3.4}
基本沿用相同的证明, 可以在更一般的条件下证明 \hyperref[XV.lem3.2]{3.2} 和 \hyperref[XV.lem3.3]{3.3}: 不必假设 $B$ 为优环且正规, 只须假设完备化 $\widehat{B}$ 正规 (由 \EGA IV 7 的结果, 这一条件在局部化下稳定), 或者只假设 $B$ 的 $(tB)$-进拓扑完备化正规且为某个正则环的商环.
\end{enonce*}

\section{附录: 局部 $0$-无上同调性的一个判据}\etiquette[4]{XV.sec4}\pageoriginale

本节结果在讨论班的后续部分中不会用到.

\begin{enonce*}{定理 4.1}\etiquette[4.1]{XV.thm4.1}
设 $g:X\rightarrow Y$ 为概形态射. 假设 $g$ 平坦且纤维可分, 并且 $X$ 与 $Y$ 局部 Noether \sisi{(}{}相应地, 假设 $g$ 局部有限表示\sisi{)}{}. 则 $g$ 普遍局部 $0$-无上同调.
\end{enonce*}

显然可设 $X$ 与 $Y$ 仿射; 在第二种情形, 利用 \EGA IV 9.7.7, \EGA IV 8 的标准方法便将问题化为 $Y$, 从而 $X$ 为 Noether 的情形. 因而只须证明第一项断言. 由于假设在有限型基变换下稳定, \hyperref[XV.prop1.13]{1.13}~(ii) 中的极限过渡将问题化为证明 $g$ 局部 $0$-无上同调. 为此, \hyperref[XV.cor1.17]{1.17} 将问题化为证明: 若还假设 $g$ 严格局部, 则 $g$ 的几何纤维连通. 这恰为 \EGA IV 18.9.8.

\begin{enonce*}{推论 4.2}\etiquette[4.2]{XV.cor4.2}
设 $g:X\rightarrow Y$ 为概形态射, $Y$ 离散. 则 $g$ 普遍局部 $0$-无上同调.
\end{enonce*}

由 \hyperrefext[VIII][VIII.thm1.1]{1.1}, 可假设 $Y$ 为完美域 $k$ 的谱. 还可假设 $X$ 仿射, 从而为 $k$ 上有限型仿射概形的投射极限; 根据 \hyperref[XV.prop1.13]{1.13}~(ii), 这将问题化为 $X$ 在 $k$ 上有限型的情形. 最后,\pageoriginale 显然可假设 $X$ 既约; 由于 $k$ 为完美域, 此时它在 $k$ 上可分. 结论即由 \hyperref[XV.thm4.1]{4.1} 得出.

\begin{enonce*}{推论 4.3}\etiquette[4.3]{XV.cor4.3}
假设 $g$ 满射且满足 \hyperref[XV.thm4.1]{4.1} 或 \hyperref[XV.cor4.2]{4.2} 的假设; 在 \hyperref[XV.thm4.1]{4.1} 的非括号情形, 再假设 $g$ 或者拟紧且拟分离, 或者普遍开. 则对于以变动预概形上的平展层组成的纤维范畴（fibered category）, $g$ 是普遍有效下降（universal effective descent）态射.
\end{enonce*}

在所考察的情形, $g$ 普遍为商映射, 由 \hyperrefext[VIII][VIII.prop9.1]{9.1}, $g$ 是有关纤维范畴的普遍下降态射. 因而就有效性问题而言, 只须处理 $Y$ 仿射的情形. 此外, 在所考察的每种情形, 都可选取 $X$ 的\emph{有限}个仿射开子集 $X_i$, 使其像覆盖 $Y$. 用 $X_i$ 的不交并代替 $X$, 因而可假设 $X$ 仿射, 从而更是 $Y$ 上拟紧且拟分离的. 此时可按 \hyperrefext[VIII][VIII.subsubsec9.4.1]{9.4.1} 的论证得出结论, 这里用到以下事实 \hyperrefext[XVI][XVI.thm1.1]{1.1}: 对拟紧且拟分离态射 $g$, $g_{\ast}(F)$ 的构造 (其中 $F$ 为 $X$ 上的平展层) 与每个局部 $0$-无上同调的基变换 $Y'\rightarrow Y$ 可交换.

\begin{enonce*}[remark]{注 4.4}\etiquette[4.4]{XV.rems4.4}
\begin{enumerate}
\renewcommand{\theenumi}{\alph{enumi}}
\renewcommand{\labelenumi}{\theenumi)}
\item 这将 \hyperrefext[VIII][VIII.thm9.4]{9.4}~d) 作为 \hyperref[XV.cor4.3]{4.3} 的特例加以证明, 补上了此前暂缺的证明. % d) 的标签?
还有一个更容易的证明, 不需要用到 \EGA IV 18.9.8 这个相当精细的\pageoriginale 结果: 注意, 由 \hyperrefext[VIII][VIII.prop9.1]{9.1}, 只须考虑域扩张 $K/k$ 所诱导的态射 $\Spec(K)\rightarrow \Spec(k)$ 的情形. 但 $\Spec(K)$ 在 $\Spec(k)$ 上普遍局部 $0$-无上同调, 因为它甚至关于 $\LL=\PP-\{p\}$ 普遍无上同调, 这里 $p=\car k$; 我们将在 \hyperrefext[XVI][XVI.cor1.5]{1.5} 中看到这一点.

\item 尚不知道, 每个拟紧且拟分离、满射、局部 $(-1)$-无上同调 (或普遍局部 $(-1)$-无上同调) 态射, 是否都是以变动预概形上的平展层组成的纤维范畴的有效下降态射. 这一陈述似乎相当可信, 并将加强 \hyperref[XV.cor4.3]{4.3} 以及 \hyperrefext[VIII][VIII.thm9.4]{9.4}~c), d). % c) 与 d) 的标签?
可比较 \textsc{Murre}, Sém. Bourbaki ${\rm n^{\circ}}$~293, p.~17 中的有效性陈述.

\item 在 \hyperref[XV.cor4.2]{4.2} 的假设下, $g$ 甚至关于 $\LL$ 局部无上同调且关于 $\LL$ 局部 $1$-非球面, 似乎也是可信的; 这里 $\LL$ 为不同于 $Y$ 的剩余特征的素数所成的集合. 对于这一问题, 显然可假设 $Y$ 为代数闭域 $k$ 的谱, 并可证明 (\SGA 1964/65), 只要 $k$ 上有限型概形的奇点消解成立, 答案就是肯定的. 因而利用 \textsc{Hironaka} 的结果, 可得当 $Y$ 的特征为零时答案是肯定的~\footnote[1]{\Cf \SGA 1 XIII.}.
\end{enumerate}
\end{enonce*}

% XII 4.1, 4.2 为公式引用.
% XII 1.2, 1.6 在第 501 行无引用.



% END native input: chapters/exp-15-zh.tex
% BEGIN native input: chapters/exp-16-zh.tex
\chapter{光滑基变换定理及其应用}
%\addtocontents{toc}{M. Artin}
\begin{center}
M. Artin
\end{center}
\phantomsection\MakeLinkTarget*{sga4.exp.XVI}\label{sga4.exp.XVI}

\setcounter{pageoriginal}{206}

\section{光滑基变换定理}\etiquette[1]{XVI.sec1}\pageoriginale

\begin{enonce*}{定理 1.1}\etiquette[1.1]{XVI.thm1.1}
设 $\LL\subset \PP$, $n\in \NN$, 且
$$\vcenter{
\xymatrix@C=1.5cm@R=1cm{
X   \ar[d]^{f}& X' \ar[l]_{g'} \ar[d]^{f'}\\
Y     & \ar[l]_{g} Y'
}}\leqno{(*)}
$$
是 Cartesian 图, 其中 $f$ 拟紧且拟分离\nde{即原文术语中的凝聚态射.}, $g$ 普遍局部 $0$-无上同调 (universally locally 0-acyclic) (\resp 对 $\LL$ 普遍局部 $1$-非球面, \resp 对 $\LL$ \sisi{普遍局部}{普遍局部} $n$-无上同调) \emph{(\hyperrefext[XV][XV.def1.11]{1.11}).} 则对于 $X$ 上任意集合层 $F$ (\resp \emph{ind}-$\LL$-群层, \resp $\LL$-挠阿贝尔层), 基变换态射
\emph{(XII \hyperref[XII.eq4.2]{4.2})}
$$
\varphi^{q}:g^{\ast}(R^{q}f_{\ast})F\longrightarrow
(R^{q}f'_{\ast}){g'}^{\ast}F
$$
在 $q=0$ (\resp $q\leq 1$, \resp $q\leq n$) 时是双射.
\end{enonce*}

特别地, 应用 \hyperrefext[XV][XV.thm2.1]{2.1}, 得到下述结果.

\begin{enonce*}{推论 1.2}\etiquette[1.2]{XVI.cor1.2}
(光滑基变换定理). 假设 $(\ast)$ 中的态射 $g$ 光滑, $f$ 拟紧且拟分离, $\LL\subset \PP$ 是 $Y$ 的剩余特征集合的补集. 则对于 $X$ 上任意集合层 $F$ (\resp \emph{ind}-$\LL$-群层, \resp $\LL$-挠阿贝尔层), 上述基变换\pageoriginale 态射 $\varphi^{q}$ 在 $q=0$ (\resp $q=0$, $1$, \resp 任意 $q$) 时是双射.
\end{enonce*}

\noindent
\emph{\hyperref[XVI.thm1.1]{1.1} 的证明}. 先处理 $f:X\rightarrow Y$ 拟射影的情形: 这只需结合固有基变换定理 \hyperrefext[XII][XII.thm5.1]{5.1} 和定义 \hyperrefext[XV][XV.prop1.9]{1.9}. 事实上, 有交换图
$$
\xymatrix@=1.25cm{
 X \ar[rr]^{i} \ar[rd]_{f}&        & \bar{X}\ar[ld]^{\bar{f}}\\
                         &   Y    &
}
$$
其中 $\bar{f}$ 射影, 从而固有, $i$ 是开浸入. 由 \hyperrefext[XII][XII.thm5.1]{5.1} 知, 无论态射 $g:Y'\rightarrow Y$ 如何, 基变换定理对 $\bar{f}$ 都成立. 由于 $\fprod{g}{\bar{X}}{Y}:\fprod{Y'}{\bar{X}}{Y}\rightarrow \bar{X}$ 仍是普遍 $0$-无上同调的 (\resp\ldots), 由 (\hyperrefext[XII][XII.prop4.4]{4.4} (ii)) 立即可知, 只须对态射 $i$ 证明本定理; 即归结为 $f$ 是开浸入的情形. 此时定理由定义 \hyperrefext[XV][XV.prop1.10]{1.10} (ii) 得出.

现在处理一般情形. 断言在 $Y$ 上是局部的, 因而可假设 $Y$ 仿射.

\begin{enonce*}{引理 1.3}\etiquette[1.3]{XVI.lem1.3}
还可假设 $X$ 仿射.
\end{enonce*}

\noindent
\emph{证明}. 仿照 \hyperrefext[XII][XII.lem6.1]{6.1} 的论证, 为了对给定的 $f$ 证明 \hyperref[XVI.thm1.1]{1.1}, 只须考虑由 $g$ 诱导的 $Y'$ 和 $Y$ 的严格局部化之间的局部态射 $\widetilde{g}:\widetilde{Y}'\rightarrow \widetilde{Y}$, 对给定情形作基变换 $\widetilde{Y}\rightarrow Y$, $\widetilde{Y}'\rightarrow Y'$ 所得\pageoriginale 的情形, 并证明同态
\begin{equation*}
H^{q}(\widetilde{X},F)\longrightarrow H^{q}(\widetilde{X}',F')\quad
(\mbox{\rm 其中 } F'={g'}^{\ast}(F))\tag{$\ast$}
\end{equation*}
在 $q\leq n$ 时是双射, 这里 $\widetilde{X}=\fprod{X}{\widetilde{Y}}{Y}$, $\widetilde{X}'=\fprod{X'}{\widetilde{Y}'}{Y'}$. 因而归结为 $Y$, $Y'$ 和 $g$ 均严格局部的情形, 且只须在已对仿射 $f$ 证明 \hyperref[XVI.thm1.1]{1.1} 的假设下证明上述映射的双射性. 取 $X$ 的有限开覆盖 $(X_i)$, 其中各 $X_i$ 为仿射开子概形, 令 $Z$ 为诸 $X_i$ 的不交并; 因而它仿射, 并带有满射态射 $h:Z\rightarrow X$, 从而有态射 $h':Z'\rightarrow X'$. 注意 $h$ 分离, 且当 $f$ 分离时它仿射. 应用下降引理 \hyperrefext[XII][XII.prop6.8]{6.8}, 可知为了证明对 $q\leq n$ 和任意 $F$, 映射 $(\ast)$ 都是双射, 只须证明相应映射 $H^{q}(Z,h^{\ast}(F))\rightarrow H^{q}(Z',{h'}^{\ast}(F'))$ 是双射, 条件是 $Z\rightarrow X$ 对基变换 $X'\rightarrow X$ 的基变换同态在维数 $\leq n$ 时为同构. 若 $f$ 分离, 则 $h$ 仿射, 根据假设确有此性质; 这就证明了 $f$ 分离时的 \hyperref[XVI.thm1.1]{1.1}. 一般情形下, 可将上述结果应用于总是分离的 $h$, 再次得到 \hyperref[XVI.thm1.1]{1.1} 对 $f$ 成立, 从而证明 \hyperref[XVI.lem1.3]{1.3}.

\begin{enonce*}[remark]{注}
也可使用 $X$ 的开覆盖 $(X_i)$ 的 Leray 谱序列及其非交换变形\pageoriginale, 但此法与使用 \hyperrefext[XII][XII.prop6.8]{6.8} 本质上并无区别, 缺点是迫使我们再次区分通常的三种情形, 而这项工作已在\loccit 完成.
\end{enonce*}

为在一般情形下完成定理的证明, 只须归结为 $f:X\rightarrow Y$ 仿射且有限型, 从而拟射影的情形. 由 \hyperref[XVI.lem1.3]{1.3}, 可假设 $X$ 仿射. 写成 $X=\varprojlim X_{\alpha}$, 其中态射 $f_{\alpha}:X_{\alpha}\rightarrow Y$ 仿射且有限型. 通常的极限论证 (\hyperrefext[VII][VII.sec5]{5}) 即完成证明.

\begin{enonce*}[remark]{补充 1.4}\etiquette[1.4]{XVI.compl1.4}
\begin{enumerate}
\renewcommand{\theenumi}{\alph{enumi}}
\renewcommand{\labelenumi}{\theenumi)}
\item 请读者核实, 当假设 $g$ 普遍局部 $(-1)$-无上同调时 (\hyperrefext[XV][XV.rems1.14]{1.14}), 同一证明仍可给出一个结论: 此时对 $X$ 上任意集合层 $F$, 基变换同态
$$
\phi^{0}:g^{\ast}(f_{\ast}(F))\longrightarrow f'_{\ast}({g'}^{\ast}(F))
$$
是\emph{单射}. 另一方面, 仍设 $n\geq 0$. 可对 \hyperref[XVI.thm1.1]{1.1} 补充下述单射性断言, 它也包含在上述证明中: 若 $g$ 普遍局部 $0$-无上同调, 则对于 $X$ 上任意群层 $F$, \hyperref[XVI.thm1.1]{1.1} 中的基变换同态 $\phi^{1}$ 是单态射; 若 $g$ 对 $\LL$ 普遍局部 $n$-无上同调, 则对于 $X$ 上任意 $\LL$-挠阿贝尔层 $F$, \hyperref[XVI.thm1.1]{1.1} 中的同态 $\phi^{n+1}$ 是单态射.

\item 很可能还可以通过引入 \textsc{Giraud} 学位论文中的非阿贝尔 $2$-上同调不变量, 从 $g$ 对 $\LL$ 局部 $1$-非球面的假设得到类似的单射性结论\pageoriginale; 比较 \hyperrefext[XII][XII.subsec5.11]{5.11}\footnote[1]{J. Giraud 的书中确实证明了这一点.}; 对下文的 (\hyperref[XVI.cor1.6]{1.6}) 和 (\hyperref[XVI.cor2.3]{2.3}) 也有同样的问题.
\end{enumerate}
\end{enonce*}

%raghu
\begin{enonce*}{推论 1.5}\etiquette[1.5]{XVI.cor1.5}
设 $K/k$ 为域扩张, $\LL=\PP-\{p\}$, 其中 $p= \car k$. 则 $\Spec(K)\rightarrow \Spec(k)$ 对 $\LL$ 普遍局部无上同调, 且对 $\LL$ 普遍局部 $1$-非球面; 若 $k$ 和 $K$ 可分闭, 则上述态射还对 $\LL$ 普遍无上同调, 且对 $\LL$ 普遍 $1$-非球面.
\end{enonce*}

立即记下 (\hyperref[XVI.cor1.5]{1.5}) 的如下推论.

\begin{enonce*}{推论 1.6}\etiquette[1.6]{XVI.cor1.6}
设 $K/k$ 为可分闭域的扩张, $X$ 为概形, $F$ 为 $X$ 上的\sisi{集合}{集合}层 (\resp ind-$\LL$-群层, \resp $\LL$-挠阿贝尔群层), 其中 $\LL=\PP-\{p\}$, $p=\car k$. 则映射
$$
H^{q}(X,F)\longrightarrow H^{q}(X_{K},F_{K})
$$
在 $q=0$ (\resp $q\leq 1$, \resp 任意 $q$) 时是双射.
\end{enonce*}

事实上, 这立即由 $\Spec(K)\rightarrow \Spec(k)$ 对 $\LL$ 普遍无上同调且对 $\LL$ 普遍 $1$-非球面这一事实及定义得出.\sisi{---~}{~}注意由 \hyperrefext[VIII][VIII.cor2.4]{2.4}, 上述 $H^{q}$ 之间的同态也可解释为与基变换同态
$$
g^{\ast}(R^{q}f_{\ast}(F))\longrightarrow
R^{q}{f'}_{\ast}({g'}^{\ast}(F)){\quoi},
$$
相同, 这里\pageoriginale $f:X\rightarrow \Spec(k)$, $f':X_K\rightarrow \Spec(K)$, $g:\Spec(K)\rightarrow \Spec(k)$ 和 $g':X_K\rightarrow X$ 是显然的态射. 因此, 当 $X$ 拟紧且拟分离时, 由 \hyperref[XVI.thm1.1]{1.1}, 这些同态的双射性也可看作 $\Spec(K)\rightarrow \Spec(k)$ 对 $\LL$ 普遍\emph{局部}无上同调且对 $\LL$ 普遍局部 $1$-非球面这一事实的推论. 因而后一事实已蕴涵拟紧且拟分离的 $X$ 所满足的 \hyperref[XVI.cor1.6]{1.6}, 而这显然足以推出 $\Spec(K)\rightarrow \Spec(k)$ 对 $\LL$ 普遍无上同调且对 $\LL$ 普遍 $1$-非球面 (\hyperrefext[XV][XV.def1.7]{1.7}, \hyperrefext[XV][XV.prop1.6]{1.6} (i)).

这说明, 为证明 \hyperref[XVI.cor1.5]{1.5}, 只须证明 \hyperref[XVI.cor1.5]{1.5} 的第一个断言. 由 \hyperrefext[VIII][VIII.thm1.1]{1.1}, 可以换成 $k$ 的完美闭包, 因而可假设 $k$ \emph{完美}. $k$ 的任意扩张 $K$ 都是其有限生成 $k$-子代数 $A$ 的滤过正向极限; 由于 $k$ 完美, 对这样的 $A$ 作局部化后, 可假设它光滑, 从而 $K$ 表示为光滑子代数的滤过正向极限. 由 \hyperrefext[XV][XV.thm2.1]{2.1}, 后者对 $\LL$ 普遍局部无上同调且对 $\LL$ 普遍局部 $1$-非球面, 故结论由 \hyperrefext[XV][XV.prop1.13]{1.13} (ii) 得出, \cqfd

\begin{enonce*}[remark]{注 1.7}\etiquette[1.7]{XVI.rem1.7}
应将不变性定理 \hyperref[XVI.cor1.5]{1.5} 与 \hyperrefext[XII][XII.cor5.4]{5.4} 比较: 后者不必假设所考虑的挠层与剩余特征互素, 但另一方面, 必须假设 $X$ 在 $k$ 上\emph{固有}. 即使只考虑 $H^{1}$ $(X,\ZZ/p\ZZ)$, 其中 $X$\pageoriginale 是特征 $p>0$ 的代数闭域 $k$ 上的仿射代数曲线 (例如仿射直线), 上述断言的类比也不成立, 原因在于无穷远处的“野分歧\nde{即“野的”.}”现象; 这些现象意味着在这种情形下, 群为 $\ZZ/p\ZZ$ 的主平展覆盖的分类本质上是“连续的”.
\end{enonce*}

\section{上同调群的特殊化定理}\etiquette[2]{XVI.sec2}

\begin{enonce*}{定理 2.1}\etiquette[2.1]{XVI.thm2.1}
设 $\LL\subset \PP$, $n\in \NN$, $f:X\rightarrow Y$ 为有限表示的固有态射\footnote[1]{关于比固有性更弱而能得到相同结论的假设, \cf \SGA 5 II 3 和 \SGA 1 XIII.}, 且它局部 $0$-无上同调 (\resp 对 $\LL$ 局部 $1$-非球面, \resp 对 $\LL$ 局部 $n$-无上同调). 设 $F$ 为 $X$ 上可构造且局部常值的集合层 (\resp \emph{ind}-$\LL$-群层, \resp $\LL$-挠阿贝尔层). 则在 $q=0$ (\resp $q=0$, $1$, \resp $q\leq n$) 时, $R^{q}f_{\ast}F$ 可构造且\emph{局部常值}, 并且对 $Y$ 的任意几何点 $\bar{y}$, 有
$$
(R^{q}f_{\ast}F)_{\bar{y}}\simeq
H^{q}(X_{\bar{y}},F_{\bar{y}}){\quoi},
$$
其中 $q$ 取上述相同的值.
\end{enonce*}

\begin{enonce*}[remark]{注}
注意最后一个断言正是固有基变换定理 \hyperrefext[XII][XII.cor5.2]{5.2}.
\end{enonce*}

由 \hyperrefext[XV][XV.thm2.1]{2.1}, 立即得到如下推论:\pageoriginale

\begin{enonce*}{推论 2.2}\etiquette[2.2]{XVI.cor2.2}
(上同调群的特殊化定理). 设 $f:X\rightarrow Y$ 为固有光滑态射, $\LL\subset \PP$ 为 $Y$ 的剩余特征集合的补集. 设 $F$ 为 $X$ 上可构造的集合层 (\resp ind-$\LL$-群层, \resp $\LL$-挠阿贝尔层). 则在 $q=0$ (\resp $q=0$, $1$, \resp 任意 $q$) 时, $R^{q}f_{\ast}(F)$ 可构造且\emph{局部常值}, 并且对 $Y$ 的任意几何点 $\bar{y}$, 有
$$
R^{q}f_{\ast}(F)_{\bar{y}}\simeq H^{q}(X_{\bar{y}},F_{\bar{y}})
$$
其中 $q$ 取上述相同的值.
\end{enonce*}

\begin{enonce*}[remark]{注}
取 $n=1$ 并翻译为基本群的语言, 就得到 (\SGA 1 X 3.8), 它比较固有光滑态射的各纤维的基本群.
\end{enonce*}

\noindent
\emph{\hyperref[XVI.thm2.1]{2.1} 的证明}. 已知 $R^{q}f_{\ast}(F)$ 可构造 \hyperrefext[XIV][XIV.thm1.1]{1.1}, 因而由 \hyperrefext[IX][IX.prop2.11]{2.11}, 只须证明对 $Y$ 的几何点的任意特殊化态射 $\bar{y}_1\rightarrow \bar{y}_{0}$, 相应的特殊化态射 (\hyperrefext[VIII][VIII.subsec7.7]{7.7})
\begin{equation*}
R^{q}f_{\ast}(F)_{\bar{y}_{0}}\longrightarrow
R^{q}f_{\ast}(F)_{\bar{y}_1}\tag{$\ast$}
\end{equation*}
对所考虑的 $q$ 值是双射. 为此, 我们有下述稍一般的结果 (其中不必假设\pageoriginale $f$ 有限表示或 $F$ 可构造):

\begin{enonce*}{推论 2.3}\etiquette[2.3]{XVI.cor2.3}
设 $\LL\subset \PP$, $n\in \NN$, $f:X\rightarrow Y$ 为局部 $0$-无上同调 (\resp 对 $\LL$ 局部 $1$-非球面, \resp 对 $\LL$ 局部 $n$-无上同调) 的固有态射, $F$ 为 $X$ 上局部常值的集合层 (\resp ind-$\LL$-群层, \resp $\LL$-挠阿贝尔层). 则对 $Y$ 的几何点的任意特殊化态射 $\bar{y}_1\rightarrow \bar{y}_0$, 上述相应特殊化同态 $(\ast)$ 在 $q=0$ (\resp $q=0$, $1$, \resp $q\leq n$) 时是双射. 此外, 当 $n=0$, 即 $f$ 局部 $0$-无上同调时, 对 $X$ 上任意局部常值群层 $F$, 同态 $R^{1}f_{\ast}(F)_{\bar{y}_0}\rightarrow R^{1}f_{\ast}(F)_{\bar{y}_1}$ 是单射; 对任意 $n$ 和 $X$ 上任意局部常值的 $\LL$-挠阿贝尔层 $F$, 特殊化同态 $(\ast)$ 在 $q=n+1$ 时是单射.
\end{enonce*}

存在整的严格局部概形 $Y'$ 及态射 $g:Y'\rightarrow Y$, 它将闭点 $y'_0$ 映到 $y_{0}$, 将一般点 $y'_1$ 映到 $y_1$: 例如先取 $Y$ 在 $\bar{y}_0$ 处的严格局部化, 再取它的一个位于 $y_1$ 之上的点所定义的整闭子概形即可. 此外, 将 $Y'$ 换成它在函数域的一个代数闭包中的正规化后, 可假设 $Y'$ 正规且 $k(y'_1)$ 代数闭, 于是 $y'_0$ 和 $y'_1$ 都是 $Y'$ 的几何点. 对 $(f,F,g)$ 使用基变换定理 (\hyperrefext[XII][XII.thm5.1]{5.1}), 并将\pageoriginale $Y$, $\bar{y}_0$, $\bar{y}_1$ 换成 $Y'$, $y'_0$, $y'_1$, 就将 \hyperref[XVI.cor2.3]{2.3} 的证明归结为证明 \hyperref[XVI.cor2.3]{2.3} 在 $Y$ 正规且严格局部, $\bar{y}_0$ 和 $\bar{y}_1$ 分别为其闭点 $y_0$ 和一般点 $y_1$ 的情形. 又由 \hyperrefext[XII][XII.sec4]{4}, 特殊化同态 $(\ast)$ 此时等同于同态
$$
H^{q}(X,F)\longrightarrow H^{q}(X_1,g^{\ast}_1(F)){\quoi},
$$
这里 $X_1=X_{y_1}$, 而 $g:y_1\rightarrow Y$ 和 $g_1:X_1\rightarrow X$ 是典范态射. 我们的断言现由下述引理得出 (其中不再假设 $f$ 固有):

\begin{enonce*}{引理 2.4}\etiquette[2.4]{XVI.lem2.4}
沿用 $\hyperref[XVI.cor2.3]{2.3}$ 开头的记号, 去掉 $f$ 固有的假设, 改为假设 $Y$ 正规且整, 其一般点 $y_1$ 满足 $k(y_1)$ 可分闭, 并令 $X_1=X_{y_1}$, $g_1:X_1\rightarrow X$ 为典范态射. 则同态
\begin{equation*}
H^{q}(X,F)\longrightarrow H^{q}(X_1,g^{\ast}_1(F))\tag{${\ast}{\ast}$}
\end{equation*}
在 $q=0$ 时为双射 (\resp 在 $q=0$, $1$ 时为双射, \resp 在 $q\leq n$ 时为双射且在 $q=n+1$ 时为单射). 此外, 当 $n=0$, 即 $f$ 局部 $0$-无上同调时, 对 $X$ 上任意群层 $F$, 同态 $({\ast}{\ast})$ 在 $q=1$ 时为单射.
\end{enonce*}

直接的论证 (\hyperrefext[XV][XV.prop1.6]{1.6} (i) $\Rightarrow$ (ii)) 表明, 只须证明关系
$$
F\xrightarrow{\sim} g_{1\ast}g_1{}^{\ast}(F){\quoi},
$$
以及\pageoriginale 在 $n\geq 1$ 时的附加关系
$$
R^{q}g_{1\ast}g_1{}^{\ast}(F)=0\quad\mbox{\rm 对于}\quad 1\leq
q\leq n{\quoi}.
$$
这些关系对 $X$ 的平展拓扑是局部的, $F$ 又对平展拓扑局部常值, 因而可假设 $F$ 常值, 从而具有形式 $G_X$, 其中 $G$ 为集合 (\resp ind-$\LL$-群, \resp $\LL$-挠阿贝尔群). 还应注意 $Y$ 正规的假设蕴涵 $g_{\ast}(G_{y_1})\xleftarrow{\sim}G_Y$ (\hyperrefext[IX][IX.lem2.14.1]{2.14.1}), $k(y_1)$ 可分闭的假设蕴涵 $R^{q}g_{\ast}(G_{y_1})$ 在 $q\geq 1$ 时为零. 另一方面, $f$ 局部 $0$-无上同调 (\resp\ldots) 的假设蕴涵: 在 $q\leq n$ 时, $R^{q}g_{\ast}(G_{y_1})$ 的形成与基变换 $f:X\rightarrow Y$ 交换. 这里用到如下事实: 可假设 $Y$ 仿射, 此时 $y_1$ 是其在 $Y$ 中拟紧的仿射开邻域的逆向极限, 因而可应用定义 \hyperrefext[XV][XV.def1.11]{1.11} 和过渡到极限的理论 \hyperrefext[VII][VII.sec5]{5}. 于是确有
$$
G_X\xrightarrow{\sim} g_{1\ast}(G_{X_1})\quav
R^{q}g_{1\ast}(G_{X_{1}})=0\quad\mbox{\rm 对于}\quad 1\leq q\leq
n{\quoi},
$$
这就完成了 \hyperref[XVI.lem2.4]{2.4} 的证明, 从而也完成了 \hyperref[XVI.thm2.1]{2.1} 的证明.

\begin{enonce*}[remark]{注}
上述通过 \hyperref[XVI.lem2.4]{2.4} 的证明比我们原来的证明更简单, 它归功于 M. Lubkin.
\end{enonce*}

\begin{enonce*}{推论 2.5}\etiquette[2.5]{XVI.cor2.5}
在 $\hyperref[XVI.thm2.1]{2.1}$ 的假设下, 设 $Y$ 连通, $q\leq n$. 则当 $\bar{y}$ 遍历 $Y$ 的几何点\pageoriginale 时, 所有群 $H^{q}(X_{\bar{y}},F)$ 都彼此同构.
\end{enonce*}

\section{相对上同调纯性定理}\etiquette[3]{XVI.sec3}

\begin{enonce*}{定义 3.1}\etiquette[3.1]{XVI.def3.1}
设 $S$ 为概形. 所谓\emph{光滑 $S$-对} (smooth S-pair) $(Y,X)$, 是指光滑 $S$-预概形之间的闭 $S$-浸入 $i:Y\rightarrow X$, 即如下预概形交换图
$$
\xymatrix@=1.25cm{
Y \ar[rr]^{i} \ar[rd]_{h}&     &X\ar[ld]^{f}\\
  & S   &
}
$$
其中 $f$ 和 $h$ 光滑, $i$ 是闭浸入. 记 $U=X-Y$, 用 $j:U\rightarrow X$ 表示开浸入, $g:U\rightarrow S$ 表示结构态射. $(Y,X)$ 在点 $y\in Y$ 处的\emph{余维数}是纤维 $Y_s$ 在 $X_s$ 中于 $y$ 处的余维数, 其中 $s=h(y)$.
\end{enonce*}

因而余维数是 $Y$ 上的局部常值函数. 当此函数实际上为常值\sisi{,}{,} 值为 $c$ 时, 就说“$(Y,X)$ 的余维数为 $c$”.

设 $(Y,X)$ 是光滑 $S$-对, $y\in Y$. 回顾 (\SGA II 4.10), 存在整数 $m$ 和 $n$, $y$ 在 $X$ 中的邻域 $X'$, 以及平展态射
$$
\phi:X'\longrightarrow \EE^n_S=\Spec \Oo_S[t_1,\ldots,t_n]
$$
使得\pageoriginale $Y'=Y\cap X'$ 是由方程
$$
t_{m+1}=\ldots=t_n=0{\quoi}.
$$
定义的闭子概形 $\EE^m_S$ 的逆像. 此结果可表述为:\emph{对平展拓扑局部地, 任意光滑 $S$-对 $(Y,X)$ 都同构于某个适当 $m$ 和 $n$ 对应的标准对 $(\EE^m_S, \EE^n_S)$,} 其中
$$
\EE^r_S=\Spec \Oo_S[t_1,\ldots,t_r]{\quoi}.
$$

本节拟计算光滑 $S$-对 $(Y,X)$ 的局部上同调层 $\HH^q_Y(X,F)=(R^{q}i^{!})F$ (\sisi{\hyperrefext[V][V.sec4]{4}}{\hyperrefext[V][V.sec6]{6}} 和 \hyperrefext[VIII][VIII.subsec6.6]{6.6}), 系数为与剩余特征互素的局部常值阿贝尔挠层 $F$. 先作若干预备讨论:

\begin{enonce*}{命题 3.2}\etiquette[3.2]{XVI.prop3.2}
设
$$
\xymatrix@=1.25cm{
U \ar[rr]^{j} \ar[rd]_{g}&     &X\ar[ld]^{f}\\
  & S   &
}
$$
是交换图, 其中 $f$ 光滑, $j$ 为开浸入, 且对任意 $s\in S$, 纤维 $U_s$ 在 $X_s$ 中稠密. 设 $F$ 为 $S$ 上的集合层. 则典范态射
$$
f^{\ast}F\longrightarrow j_{\ast}g^{\ast}F=j_{\ast}j^{\ast}f^{\ast}F
$$
是双射.
\end{enonce*}

\noindent
\emph{证明}.\pageoriginale 容易归结为 $U$ 在 $X$ 中拟紧的情形. 设 $\bar{s}$ 为 $S$ 的几何点, $\bar{x}$ 为位于 $\bar{s}$ 之上的 $X$ 的几何点, $\widetilde{f}:\widetilde{X}\rightarrow \widetilde{S}$ 为 $X$ 和 $S$ 在相应几何点处的严格局部化之间的态射, 并令 $\widetilde{U}=\fprod{U}{\widetilde{X}}{X}$. 采用显然的记号, 有
$$
(f^{\ast}F)_{\bar{x}}=H^{0}(\widetilde{X},\widetilde{f}^{\ast}\widetilde{F})\simeq
H^{0}(\widetilde{S},\widetilde{F}){\quoi}.
$$
因此, 只须证明 $\widetilde{g}:\widetilde{U}\rightarrow \widetilde{S}$ 是 $0$-无上同调的, 这将蕴涵
$$
(j_{\ast}g^{\ast}F)_{\bar{x}}\simeq
H^{0}(\widetilde{U},\widetilde{g}^{\ast}\widetilde{F})\simeq
H^{0}(\widetilde{S},\widetilde{F}){\quoi}.
$$
由于 $\widetilde{g}$ 是光滑态射的极限, 它局部 $0$-无上同调 (\hyperrefext[XV][XV.thm2.1]{2.1}, \hyperrefext[XV][XV.def1.11]{1.11}); 由 \hyperrefext[XV][XV.cor1.12]{1.12}, 只须证明 $\widetilde{U}/\widetilde{S}$ 的几何纤维 $0$-无上同调, 即连通且非空. 设 $\bar{s}'$ 为 $\widetilde{S}$ 的几何点. 则纤维 $\widetilde{X}_{\bar{s}'}$ 正则、连通且非空; 假设蕴涵 $\widetilde{U}_{\bar{s}'}$ 是 $\widetilde{X}_{\bar{s}'}$ 的稠密开子集, 从而连通且非空, 这就证明了结果.

\begin{enonce*}{推论 3.2.1}\etiquette[3.2.1]{XVI.cor3.2.1}
在 \hyperref[XVI.prop3.2]{3.2} 的假设下, 从 $X$ 的平展覆盖范畴到 $U$ 的平展覆盖范畴的函子 $X'\mapsto X'|U$ 是全忠实的. 特别地, 对任意有限群 $G$, 限制映射
$$
H^{1}(X,G)\longrightarrow H^{1}(U,G)
$$
是单射.
\end{enonce*}

这\pageoriginale 由 \hyperref[XVI.prop3.2]{3.2} 通过一个标准论证得出; 为便于引用, 将该论证列为下述引理.

\begin{enonce*}{引理 3.2.2}\etiquette[3.2.2]{XVI.lem3.2.2}
设 $f:X\rightarrow Y$ 为概形态射, $f^{\ast}:\Et(Y)\rightarrow \Et(X)$ 为从 $Y$ 的平展覆盖范畴 $\Et(Y)$ 到 $\Et(X)$ 的逆像函子 $Y'\mapsto \fprod{Y'}{X}{Y}$. 以下条件等价:
\begin{itemize}
\item[(i)] $f^{\ast}$ 忠实 (\resp 全忠实).

\item[(ii)] 对 $Y$ 的任意平展覆盖 $Y'$, 用 $f':X'\rightarrow Y'$ 表示由 $f$ 经基变换 $Y'\rightarrow Y$ 得到的态射, 则从 $Y'$ 的开闭子集集合到 $X'$ 的开闭子集集合的映射 $U\mapsto {f'}^{-1}(U)$ 是单射 (\resp 双射).

\item[(ii bis)] 采用 {\rm(ii)} 的记号, 对 $Y'$ 上任意常值集合层 $C$, 典范映射
$$
H^{0}(Y',C)\longrightarrow H^{0}(X',{f'}^{\ast}C)
$$
是单射 (\resp 双射).

\item[(iii)] 采用 {\rm(ii)} 的记号, 典范映射
$$
\Gamma(Y'/Y)\longrightarrow \Gamma(X'/X)
$$
是单射 (\resp 双射).
\end{itemize}
\end{enonce*}

证明留给读者 (\cf \SGA 2, IX 3.1 和 3.2).

\begin{enonce*}[remark]{注 3.2.3}\etiquette[3.2.3]{XVI.rems3.2.3}
命题 \hyperref[XVI.prop3.2]{3.2} 可以大幅推广\pageoriginale: 对集合层, 可削弱对 $f$ 的假设; 对群层, 则可得到关于 $R^{q}g_{\ast}(f^{\ast}F)$ 的结论 (\cf \SGA 2 XIV 1 和 \SGA 1 XIII). 类似的注也适用于 \hyperref[XVI.cor3.4]{3.4} 至 \hyperref[XVI.cor3.6]{3.6}.
\end{enonce*}

再“回顾” Zariski--Nagata 纯性定理的相对形式:

\begin{enonce*}{定理 3.3}\etiquette[3.3]{XVI.thm3.3}
(相对纯性定理). 设
$$
\xymatrix@=1.25cm{
U \ar[rr]^{j} \ar[rd]_{g}&     &X\ar[ld]^{f}\\
  & S   &
}
$$
是交换图, 其中 $f$ 光滑, $j$ 为开浸入, 且对任意 $s\in S$, $X_s-U_s$ 在 $X_s$ 中处处余维数 $\geq 2$. 则从 $X$ 的平展覆盖范畴到 $U$ 的平展覆盖范畴的函子 $X'\mapsto X'|U$ 是\emph{范畴等价.}
\end{enonce*}

\noindent
\emph{证明}. 已在 \hyperref[XVI.cor3.2.1]{3.2.1} 中看到此函子全忠实. 其\sisi{本质}{本质}满性将由经典定理得出. 给定 $U$ 的平展覆盖 $U'$, 必须证明它可延拓为 $X$ 的平展覆盖 $X'$. 可直接且容易地先归结为 $U$ 在 $X$ 中拟紧的情形, 再归结为 $S$ 在 $\Spec \ZZ$ 上仿射且有限型、$X$ 连通且 $f$ 有限型的情形.

对 $n=\dim S$ 作归纳\pageoriginale; 当 $n=0$ 时, 断言是经典定理 (\SGA 2 X 3.4) 的推论. 由于已有唯一性, 可将 \sisi{\SGA IX 4.7}{\SGA 1 IX 4.7} 的下降技巧应用于正规化 $S\leftarrow \bar{S}$, 从而归结为 $S$ (因而 $X$) 正规的情形 ($\bar{S}\rightarrow S$ 有限 (\EGA IV 7.8.3 (iii) (vi)). 又可扩大 $U$, 假设 $U$ 是 $U'$ 可在其上延拓的开子集中的极大者, 于是必须证明 $U=X$. 若不然, 取 $X-U$ 的极大点 $x$, 置 $s=f(x)$. 只须证明 $U'$ 可在 $x$ 的一个邻域上延拓. 可假设 $\dim \Oo_{S,s}=n\geq 1$, 且定理已在维数 $<n$ 时得证. 置 $B=\Oo_{X,x}$, 令 $B'$ 为 $B$ 在 $U$ 的覆盖 $U'$ 上的有理函数环中的正规化; 于是 $B'$ 是有限 $B$-代数 (\loccit). 显然只须证明 $B'$ 是平展 $B$-代数.

令 $A=\Oo_{S,s}$, $0\neq t\in \max A$, $A_{0}=A/tA$, $B_{0}=B/tB$. 由归纳假设, 此覆盖在 $U_{0}=U\otimes_A A_{0}$ 上的限制可延拓为 $\Spec B_{0}$ 的平展覆盖, 记为 $\Spec B^{\ast}_{0}$. 将 $X$ 换成 $x$ 的一个邻域的适当平展覆盖 (这是允许的) 后, 可假设 $\Spec B^{\ast}_{0}$ 在 $\Spec B_{0}$ 上完全分裂. 于是归结为 $B_{0}$-代数 $B'_{0}$ 在 $\Spec B_{0}$ 的闭点以外完全分裂的情形. 此时满足 \hyperrefext[XV][XV.lem3.3]{3.3} 的假设, 因而结果得证.

\begin{enonce*}{推论 3.4}\etiquette[3.4]{XVI.cor3.4}
在与 \hyperref[XVI.thm3.3]{3.3} 相同的假设下, 对 $S$ 上任意 ind-有限群层 $F$, 有 $(R^{1}j_{\ast})g^{\ast}F=0$.
\end{enonce*}

照常归结为\pageoriginale $S$ 为 Noether 概形、$f$ 有限型且 $F$ 可构造的情形. 设 $0\rightarrow F\rightarrow G$ 为 $S$ 上 ind-有限群层的单射, 置 $C=G/F$. 由 \hyperref[XVI.prop3.2]{3.2} 得函子 $j_{\ast}g^{\ast}$ 正合, 从而 $j_{\ast}g^{\ast}G\rightarrow j_{\ast}g^{\ast}C$ 为满射, 因而 (\hyperrefext[XII][XII.prop3.1]{3.1}) $(R^{1}j_{\ast})g^{\ast}F\rightarrow (R^{1}j_{\ast})g^{\ast}G$ 为单射. 所以可以将 $F$ 换成 $G$, 从而由 \hyperrefext[IX][IX.prop2.14]{2.14} 和 \hyperrefext[VIII][VIII.prop5.5]{5.5}, 归结为 $F=G$ 为常值层, 其值为通常的有限群 $G$ 的情形.
%% 源注: 非阿贝尔情形中此归约不完全清楚; 需给出 R^1 j_* p_*G = p_* R^1 j_*G 的参考文献.
设 $\bar{x}$ 为 $X$ 的几何点, $\widetilde{X}$ 为 $X$ 在 $\bar{x}$ 处的严格局部化, 并令 $\widetilde{U}=\fprod{\widetilde{X}}{U}{X}$. 由 \hyperref[XVI.thm3.3]{3.3} 得
$$
(R^{q}j_{\ast}g^{\ast}F)_{\bar{x}}=H^{q}(\widetilde{U},G_{\widetilde{U}})\simeq
H^{q}(\widetilde{X},G_{\widetilde{X}}){\quoi},
$$
而后者为零, 结果得证.

\begin{enonce*}{引理 3.5}\etiquette[3.5]{XVI.lem3.5}
(相对 Abhyankar 引理). 设 $(Y,X)$ 是余维数为 $1$ 的光滑 $S$-对, $Y$ 在 $X$ 中由方程 $t=0$ 定义. 设 $U=X-Y$, $V/U$ 是阶数 $n$ 与 $Y$ 的剩余特征互素的主 Galois 覆盖. 置 $X'=\Spec \Oo_X[z]/(z^n-t)$, $U'=\fprod{U}{X'}{X}$, $V'=\fprod{V}{X'}{X}$. 则覆盖 $V'$ 唯一地延拓为 $X'$ 的平展覆盖.
\end{enonce*}

\noindent
\emph{证明}. 立即归结为 $S$, $X$, $Y$ 均在 $\Spec \ZZ$ 上仿射且有限型的情形. 唯一性是 \hyperref[XVI.cor3.2.1]{3.2.1} 的特例. 将从 \hyperrefext[XV][XV.cor1.12]{1.12} 推出存在性: 设 $\bar{s}$ 为 $S$ 的几何点, $\bar{x}$ 为位于 $\bar{s}$ 之上的 $Y$ 的几何点. 设 $\widetilde{f}:\widetilde{X}\rightarrow \widetilde{S}$ 为 $f$ 诱导的相应严格局部化\pageoriginale 之间的态射, 并令 $\widetilde{U}=\fprod{U}{\widetilde{X}}{X}$, 等等\ldots 由唯一性和下降, 显然只须证明 $\widetilde{V}'$ 可延拓为 $\widetilde{X}'$ 的平展覆盖, 即 (因为 $\widetilde{X}'$ 没有非平凡覆盖) $\widetilde{V}'$ 是 $\widetilde{U}'$ 的平凡覆盖.

设

$\bar{U}=\varprojlim_m\widetilde{U}_{m}=\Spec \Oo_{\widetilde{U}}[t^{1/m}]$,
其中 $m$ 遍历与 $\widetilde{X}$ 的剩余特征互素的整数集合. 则 $\bar{U}/\widetilde{U}$ 是群为阿贝尔群 $G=\varprojlim_m \ZZ/m\ZZ$ 的“无限 Galois 覆盖”; 立即可知, 只须证明 $\bar{U}$ 对 $\widetilde{X}$ 的剩余特征的补集 $\LL\subset \PP$ 是 $1$-非球面的.

考虑态射 $\bar{f}:\bar{U}\rightarrow \widetilde{S}$. 它是光滑态射的极限, 因而对 $\LL$ 局部 $1$-非球面 (\hyperrefext[XV][XV.thm2.1]{2.1}, \hyperrefext[XV][XV.def1.11]{1.11} (ii)). 由 \hyperrefext[XV][XV.cor1.12]{1.12}, 只须证明 $\bar{U}/\widetilde{S}$ 的几何纤维对 $\LL$ 是 $1$-非球面的. 设 $\bar{s}$ 为 $\widetilde{S}$ 的几何点. 几何纤维 $\bar{U}_{\bar{s}}$ 是纤维 $(\widetilde{U}_m)_{\bar{s}}$ 的极限. 因而只须证明某个 $(\widetilde{U}_m)_{\bar{s}}$ 的任意阶数与\sisi{剩余特征}{剩余特征}互素的主 Galois 覆盖, 都在 $\bar{U}_{\bar{s}}$ 上诱导平凡覆盖. 将 $\widetilde{X}$ 换成 $\widetilde{X}_m=\Spec \Oo_{\widetilde{X}}[t^{1/m}]$ (它在 $\widetilde{S}$ 上仍光滑), 就归结为 $m=1$, 即 $\widetilde{U}_m=\widetilde{U}$ 的情形. 设 $V$ 为 $\widetilde{U}_{\bar{s}}$ 的这样一个覆盖. 这时 $\widetilde{X}_{\bar{s}}=Z$ 为正则严格局部概形, $W=\widetilde{U}_{\bar{s}}$ 是余维数为 $1$ 的正则子集的开补集. 因而可应用 Abhyankar 引理的通常形式 (\SGA 1 X 3.6), 得到: 对某个适当的 $m$, $V$ 诱导 $W_m$ 的覆盖\pageoriginale $V_m$, 它可延拓为 $Z_m$ 的平展覆盖. 但因 $Z$, 从而 $Z_m$, 严格局部, 此覆盖平凡; 所以 $V_m$ 是 $W_m$ 的平凡覆盖, \cqfd

\begin{enonce*}[remark]{注 3.5.1}\etiquette[3.5.1]{XVI.rem3.5.1}
若不假设 $G$ 的阶数 $n$ 与 $X$ 的剩余特征互素, 仍能证明如下结果: 设 $m$ 为整除 $G$ 的阶数且与 $Y$ 的剩余特征互素的最大整数, 并假设对任意点 $s\in S$, $U_s=X_s-Y_s$ 的覆盖 $V_s$ 在 $Y_s$ 的极大点处的惯性群的阶数与 $Y$ 的剩余特征互素. 则取 $X'=\Spec \Oo_X[z]/(z^m-t)$ 时, \hyperref[XVI.lem3.5]{3.5} 的结论仍成立. 更一般地, 也可以取 $m$ 为刚才考虑的惯性群的阶数的任意公倍数, 条件是这些阶数与相应的剩余特征互素.
\end{enonce*}

\begin{enonce*}{推论 3.6}\etiquette[3.6]{XVI.cor3.6}
(维数 $1$ 的上同调纯性). 设 $(Y,X)$ 为余维数 $1$ 的光滑 $S$-对, $F$ 为 $X$ 上可构造的局部常值群层, 其阶数与 $X$ 的剩余特征互素. 则采用 $\hyperref[XVI.def3.1]{3.1}$ 的记号, $(R^{1}j_{\ast})j^{\ast}F$ 作为带基点集合层局部同构于 $i_{\ast}(F_Y/\cint(F_Y))$, 其中 $F_Y=F|Y$, 而 $F_Y/\cint(F_Y)$ 表示 $F_Y$ 对自身的内自同构作用的商.
\end{enonce*}

\begin{enonce*}[remark]{注}
推论 \hyperref[XVI.cor3.6]{3.6} 的同构并非典范的.
\end{enonce*}

\noindent
\emph{证明}.\pageoriginale 立即归结为 $X$ 为 Noether 概形且连通的情形. 断言对 $X$ 的平展拓扑是局部的, 因而可假设 $F$ 为常值层, 其值为通常的有限群 $G$, 阶数 $m$ 与 $X$ 的剩余特征互素; 还可假设 $𝛍_{mX}\simeq (\ZZ/m\ZZ)_X$. 设 $U=X-Y$, $U'=\Spec \Oo_U[t^{1/m}]$; 因为 $𝛍_{mX}\simeq (\ZZ/m\ZZ)_X$, 这是 $U$ 的一个群为 $\ZZ/m\ZZ$ 的“本质”主 Galois 覆盖. 在 $U'$ 上平凡化的群为 $G$ 的 $U$ 的主覆盖, 在 $Y$ 上局部地由 $H=\Hom(\ZZ/m\ZZ,G)$ 模 $\cint(G)$ 分类, 后者是同构于 $G/\cint(G)$ 的底集合的带基点集合. 这给出 $Y$ 上常值层 $H_Y$ 到 $(R^{1}j_{\ast})G_U$ 的态射, 而只须证明此态射为双射, 这立即由 \hyperref[XVI.lem3.5]{3.5} 得出.

\begin{enonce*}{定理 3.7}\etiquette[3.7]{XVI.thm3.7}
\emph{(上同调纯性)}. 设 $(Y,X)$ 为余维数 $c>0$ 的光滑 $S$-对. 设 $F$ 为 $X$ 上的阿贝尔层, 对平展拓扑局部同构于形如 $f^{\ast}(G)$ 的层, 其中 $G$ 为 $S$ 上与剩余特征互素的挠层 (例如, $F$ 可以是 $X$ 上阶数与剩余特征互素的有限群的局部常值层). 则采用 $\hyperref[XVI.def3.1]{3.1}$ 的记号, 有
\begin{align*}
&\HH^{q}_Y(X,F)=(R^{q}i^{!})F=0\qquad\mbox{若}\quad q\neq 2c,\\
&\begin{cases}
 F\xrightarrow{\sim} j_{\ast}j^{\ast}F\quad\text{且} &\tag{\ie}\\
(R^{q}j_{\ast})j^{\ast}F=0 &\hspace{1.3cm} \mbox{若}\quad q\neq 0,\quad 2c-1,
\end{cases}
\end{align*}
且\pageoriginale
$$
\HH^{2c}_Y(X,F)=i^{\ast}[(R^{2c-1}j_{\ast})j^{\ast}F]
$$
是局部同构于 $i^{\ast}(F)$ 的层.
\end{enonce*}

\noindent
\emph{证明}. 断言对 $X$ 的平展拓扑是局部的; 因而可\footnote[1]{至少当假设 $F$ 局部常值时如此. 读者将确信, 下述证明本质上也适用于一般情形.}假设 $F$ 常值, 从而若愿意还可设 $F\simeq f^{\ast}(\ZZ/n\ZZ)_S$, 其中 $(\ZZ/n\ZZ)_S$ 为 $S$ 上的常值层.

先处理 $c=1$ 的情形. (至少在 $F$ 局部常值时, 关于 $(R^{1}j_{\ast})j^{\ast}F$ 的断言立即由 \hyperref[XVI.cor3.6]{3.6} 得出, 但我们将重新推导它.) 回顾 $(Y,X)$ 局部同构于对 $(\EE^{n-1}_S,\EE^n_S)$, 将 $S$ 换成 $\EE^{n-1}_S$, 就归结为 $X=\EE^{1}_S=\Spec \Oo_S[t]$ 且 $Y$ 为 $X/S$ 的截面 $t=0$ 的情形. 这种情形局部同构于射影直线 $\PP^{1}_S$ 中无穷远截面的包含. 因而可取 $X=\PP^{1}_S$, $Y$ 为无穷远截面, 并置 $U=\EE^{1}_S=X-Y$.

考虑 Leray 谱序列
$$
(R^{p}f_{\ast})(R^{q}j_{\ast})(\ZZ/n)\Longrightarrow
(R^{p+q}g_{\ast})(\ZZ/n){\quoi}.
$$
由 \hyperrefext[XV][XV.cor2.2]{2.2}, 其收敛项在 $p+q>0$ 时为零. 又当 $q>0$ 时, $R^{q}j_{\ast}$ 支撑于 $Y$ 上, 而 $f$ 将 $Y$ 同构地映到 $S$; 另一方面, 由 \hyperref[XVI.prop3.2]{3.2}, 有 $j_{\ast}(F|U)\xleftarrow{\sim}F$. 因而当 $p,q>0$ 时 $(R^{p}f_{\ast})(R^{q}j_{\ast})=0$, 当 $q>0$ 时 $f_{\ast}(R^{q}j_{\ast})$ 决定层 $(R^{q}j_{\ast})$, 最后还有 $\sisi{\RR}{R}^{p}f_{\ast}j_{\ast}(F|U)=\sisi{\RR}{R}^{p}f_{\ast}(F)$. 所以谱序列归结为同构\pageoriginale
$$
f_{\ast}(R^{q}j_{\ast})(F|U)\xrightarrow{\sim}
(R^{q+1}f_{\ast})(F)\quad\mbox{\rm 若}\quad q>0{\quoi}.
$$
但 $(R^{q}f_{\ast})(F)$ 可逐纤维计算 (\hyperrefext[XII][XII.cor5.2]{5.2}); 因而除 $q=0,2$ 外均为零, 并且当 $q=2$ 时是局部同构于 $G$ 的层, 这容易从 \hyperrefext[XII][XII.cor5.2]{5.2} 和 \hyperrefext[IX][IX.cor4.7]{4.7} 得出. 因此当 $q\neq 0,1$ 时 $(R^{q}j_{\ast})(F|U)=0$, 当 $q=1$ 时它是局部同构于 $i^{\ast}(F)$ 的层.

现在通过归纳很容易处理 $c>1$ 的情形: 设 $(Y,X)$ 为余维数 $c>1$ 的光滑 $S$-对. 显然在 $X$ 上局部地可以找到 $X$ 的余维数为 $1$、包含 $Y$ 且在 $S$ 上光滑的子概形 $Z$, 从而有交换图
$$
\xymatrix@=1.25cm{
Y       \ar[rr]^{u} \ar[rd]_{i}&    & Z\ar[ld]^v            \\
                  & X &,
}
$$
其中每个对都是光滑 $S$-对 (应称 $(Y,Z,X)$ 为光滑 $S$-三元组), $(Z,X)$ 的余维数为 $1$, $(Y,Z)$ 的余维数为 $c-1$.

有复合函子的谱序列
$$
(R^{p}u^{!})(R^{q}v^{!})F\Longrightarrow (R^{p+q}i^{!})F{\quoi}.
$$
由已经证明的情形, 当 $q\neq 2$ 时 $(R^{q}v^{!})F=0$, 当 $q=2$ 时它局部同构于 $v^{\ast}(F)$. 将归纳假设应用于 $(R^{q}v^{!})F$, 则当 $(p,q)\neq(2c-2,2)$ 时\pageoriginale $(R^{p}u^{!})(R^{q}v^{!})F$ 为零, 当 $(p,q)=(2c-2,2)$ 时它局部同构于 $u^{\ast}v^{\ast}(F)$. 关于 $(R^{p+q}i^{!})F$ 的结果立即得出, \cqfd

\begin{enonce*}{推论 3.8}\etiquette[3.8]{XVI.cor3.8}
设 $(Y,X)$ 为余维数 $c$ 的光滑 $S$-对, $n$ 为与 $X$ 的剩余特征互素的自然数, 并置
$$
T_{Y/X}=\HH^{2c}_Y(\ZZ/n\ZZ){\quoi},
$$
由 $\hyperref[XVI.thm3.7]{(3.7)}$, 这是 $Y$ 上对平展拓扑局部同构于层 $(\ZZ/n\ZZ)_Y$ 的层. 设 $F$ 为 $X$ 上被 $n$ 零化且满足 $\hyperref[XVI.thm3.7]{3.7}$ 中条件的阿贝尔层. 则 (为完整起见记下) 当 $i\neq 2c$ 时 $\HH^{i}_Y(F)=0$, 且有典范同构
$$
\HH^{2c}_Y(F)\simeq i^{\ast}(F)\otimes T_{Y/X}{\quoi}.
$$
\end{enonce*}

\noindent
\emph{证明}. 利用典范同态 $F\rightarrow\sheaf{\Hom}(\ZZ/n\ZZ,F)$, 容易定义典范同态
$$
\varphi:i^{\ast}(F)\otimes T_{Y/X}\longrightarrow \HH^{2c}_Y(F)
$$
即同态
$$
i^{\ast}(F)\longrightarrow
\sheaf{\Hom}(\HH^{2c}_Y(\ZZ/n\ZZ),\HH^{2c}_Y(F)){\quoi},
$$
. 只须证明 $\varphi$ 是同构, 沿用 \hyperref[XVI.thm3.7]{3.7} 的证明很容易核实这一点.

\begin{enonce*}{推论 3.8.1}\etiquette[3.8.1]{XVI.cor3.8.1}
在 \hyperref[XVI.cor3.8]{3.8} 的条件下, $\HH^{i}_Y(F)$ 的形成\pageoriginale 与任意基变换 $S'\rightarrow S$ 交换.
\end{enonce*}

这是显然的.

\begin{enonce*}{推论 3.9}\etiquette[3.9]{XVI.cor3.9}
(绝对上同调纯性定理). 设 $X$ 为正则概形, $Y$ 为处处余维数为 $c$ 的正则闭子预概形. 又假设 $X$ 在一个完美域 $k$ 上局部有限型. 则当 $n$ 与 $X$ 的剩余特征互素时, 有
$$
\HH^{i}_Y(\ZZ/n\ZZ)=0\mbox{ 若 } i\neq 2c{\quoi},
$$
且 $\HH^{2c}_Y(\ZZ/n\ZZ)$ 是 $Y$ 上对平展拓扑局部同构于常值层 $\ZZ/n\ZZ$ 的层.
\end{enonce*}

这为 \hyperref[XVI.thm3.7]{3.7} 的特例, 因为假设蕴涵 $X$ 和 $Y$ 在 $k$ 上光滑.

\begin{enonce*}[remark]{注 3.10}\etiquette[3.10]{XVI.rems3.10}
\begin{enumerate}
\renewcommand{\theenumi}{\alph{enumi}}
\renewcommand{\labelenumi}{\theenumi)}
\item 在 \hyperref[XVI.cor3.8]{3.8} 中可以明确描述层 $T_{Y/X}$ 的结构: 它典范同构于层 $(𝛍_n)_Y^{\otimes -c}$, 这里 $𝛍_n$ 表示 $n$ 次单位根层 (局部同构于常值层 $\ZZ/n\ZZ$). 此同构以 $Y$ 在 $X$ 中的局部\emph{基本类} (fundamental class) 的名义, 在讨论对偶性的第 XVIII 讲中, 于稍一般的相对完全交情形下研究.

\item 即使不假设存在基域 $k$, \hyperref[XVI.cor3.9]{3.9} 和 \hyperref[XVI.rems3.10]{3.10} a) 的结论也很可能\pageoriginale 仍成立, 至少在 $X$ 为“优”预概形时如此\SGAArchiveOnly{\nde{冻结底稿将这一后续进展归于 Gabber 的工作, 但未提供书目目标.}}. 当 $X$ 的特征为零时, 无论如何都会在第 XIX 讲中利用 Hironaka 的奇点解消证明这一点. 正如在 \SGA 5 中将看到的, 尤其需要“绝对纯性定理” (以及奇点解消) 来建立“局部对偶公式”, 后者本身又是 \textsc{Lefschetz--Verdier} 公式的主要组成部分之一.
\end{enumerate}
\end{enonce*}

\section{$\CC$ 上代数概形的上同调比较定理.}\etiquette[4]{XVI.sec4}

设 $X$ 为 $\Spec \CC$ 上的代数概形. 我们将比较 $X$ 的平展景和经典景, 并沿用 \hyperrefext[XI][XI.sec4]{4} 的记号. 比较定理的一般形式为 \hyperrefext[XI][XI.thm4.4]{4.4}, 其内容如下:

\begin{enonce*}{定理 4.1}\etiquette[4.1]{XVI.thm4.1}
设 $f:X\rightarrow S$ 为 $\Spec \CC$ 上局部有限型概形之间的有限型态射, 从而有如下景态射交换图
$$
\xymatrix@C=2cm@R=1.5cm{
X_{\et}  \ar[d]_{f_{\et}}                   &   X_{\cl}
\ar[l]_{\varepsilon} \ar[d]^{f_{\cl}} \\
S_{\et}                    &   S_{\cl}\ar[l]_{\varepsilon}
}
$$
设\pageoriginale $F$ 为集合层 (\resp ind-有限群层, \resp 阿贝尔挠层), 并假设满足下述情形之一:
\begin{enumerate}
\renewcommand{\theenumi}{\roman{enumi}}
\renewcommand{\labelenumi}{(\theenumi)}
\item $f$ 固有,

\item $F$ 可构造.
\end{enumerate}
则基变换态射 \sisi{\emph{XII 4.1.2}}{\hyperrefext[XII][XII.eq4.2]{4.2}}
\footnote[1]{实际上, 我们使用基变换同态的条件比第 XII 讲中更一般; 在那里只考虑了\emph{预概形}态射的 \emph{Cartesian} 方形. 读者容易核实, 此定义可逐字推广到\emph{景}态射的\emph{本质交换}方形. 关于这一点的进一步发展, 见下一讲.}
$$
\varphi:\epsilon^{\ast}(R^{q}f_{\et\ast})F\longrightarrow
(R^{q}f_{\cl\ast})\epsilon {}^{\ast}F
$$
在 $q=0$ (\resp $q=0$, $1$, \resp $q\geq 0$) 时是双射.
\end{enonce*}

注意, 在 \hyperrefext[XI][XI.thm4.4]{4.4} 的证明中所用的 Grauert--Remmert 定理 (\hyperrefext[XI][XI.thm4.3]{4.3} (iii)) 是 \hyperref[XVI.thm4.1]{4.1} 的特例.

证明\pageoriginale 时先处理 $f$ 固有的情形. 此时可通过一个本可放在第 XIV 讲中的初等论证, 从 GAGA 和固有基变换定理得到结果. 为处理一般情形, 还将使用奇点解消 [1] 和 \hyperref[XVI.thm3.7]{3.7}.

显然可假设 $S$, 从而 $X$, 在 $\CC$ 上有限型.

\emph{$f$ 固有的情形.} 回顾下述初等事实:
\begin{enumerate}
\renewcommand{\theenumi}{\alph{enumi}}
\renewcommand{\labelenumi}{(\theenumi)}
\item 拓扑空间 $X_{\cl}$ 局部紧.

\item 若 $f:X\rightarrow S$ 固有, 则 $f_{\cl}:X_{\cl}\rightarrow S_{\cl}$ 是拓扑空间的\emph{固有}映射.
\end{enumerate}

事实上, (a) 是平凡的, 因为局部地 $X_{\cl}$ 是某个 $\CC^n$ 的闭子集. 对于 (b), Chow 引理将问题归结为 $f$ 射影的情形, 而此时结论平凡.

设 $F$ 为 $X_{\et}$ 上的集合层 (\resp\ldots). 对两个态射 $f_{\et}$ 和 $f_{\cl}$, $R^{q}$ 的形成均与纤维的形成交换 (\hyperrefext[XII][XII.cor5.2]{5.2} 和 \cite{sga4.exp.XVI.bib.2}, 4.11.1). 因而只须在 $S_{\et}$ 的闭点处逐纤维核实定理.

在下述引理中, 不假设 $g$ 固有:

\begin{enonce*}{引理 4.2}\etiquette[4.2]{XVI.lem4.2}
设\pageoriginale
$$
\xymatrix@=1.5cm{
X  \ar[rd]_{g}& &\bar{X} \ar[ll]_{h}\ar[ld]^{\bar{g}}\\
& S &
}
$$
为交换图, $h$ 固有, $U$ 为 $X$ 的开子集, 且 $h$ 诱导 $\bar{X}$ 的开子集 $\bar{U}=\fprod{U}{\bar{X}}{X}$ 到 $U$ 的同构. 设 $j:U\rightarrow X$, $\bar{j}:\bar{U}\rightarrow \bar{X}$ 为包含映射, $F$ 为 $U$ 上的集合层 (\resp\ldots), $\bar{F}$ 为 $\bar{U}$ 上的诱导层. 若定理 $\hyperref[XVI.thm4.1]{4.1}$ 对 $(\bar{g},\bar{j}_{!}\bar{F})$ 成立, 则它对 $(g,j_{!}F)$ 也成立.
\end{enonce*}

\noindent
\emph{证明}. 定理对 $(h,\bar{j}_{!}\bar{F})$ 成立, 因为可以逐纤维核实, 此时分为两种情形: 层在纤维上的限制为零, 或态射 $h$ 为同构. 在两种情形下, 结果都平凡. 得到
$$
\epsilon^{\ast}j_{!}F\simeq
\epsilon^{\ast}h_{\text{\rm et\,}\ast}(\bar{j}_{!}\bar{F})\simeq
h_{\cl\ast}\epsilon^{\ast}(\bar{j}_{!}\bar{F})
$$
以及对所考虑的 $q$ 值, 有
$$
(R^{q}h_{\text{\rm et\,}\ast})\bar{j}_{!}\bar{F}
=(R^{q}h_{\cl\ast})\epsilon^{\ast}(\bar{j}_{!}\bar{F})=0
\text{ 若 }q>0
$$
从而
$$
(R^{q}g_{\text{\rm et\,}\ast})j_{!}F\simeq
(R^{q}\bar{g}_{\text{\rm et\,}\ast})\bar{j}_{!}\bar{F}
$$
以及
$$
(R^{q}g_{\cl \ast})j_{!}F\simeq (R^{q}\bar{g}_{\cl
\ast})\bar{j}_{!}\bar{F}{\quoi},
$$
这就证明了\hyperref[XVI.lem4.2]{引理 4.2}.\pageoriginale

\begin{enonce*}{引理 4.3}\etiquette[4.3]{XVI.lem4.3}
为对所有 $f$ 固有的资料 $(f,F)$ 核实 \hyperref[XVI.thm4.1]{4.1}, 只须在另外还满足 $f$ 的相对维数 $\leq 1$ 且 $S=\Spec \CC$ 为一点时核实.
\end{enonce*}

\noindent
\emph{证明}. 对相对维数作归纳. 设 $n>1$, 且已知结果在相对维数 $<n$ 时成立. 如前所述, 可逐纤维核实; 即可假设 $S=\Spec \CC$ 为一点, 从而 $X$ 的维数 $\leq n$. 又可假设 $X$ 约化. 设 $\varphi$ 为 $X$ 上的有理函数, 它在 $X$ 的\sisi{每个不可约分支}{每个不可约分支}上非常值; 取 $X$ 的稠密开子集 $U$, 使 $\varphi$ 在其上有定义; 令 $\bar{X}\subset X\times \PP^{1}$ 为态射 $\varphi:U\rightarrow \EE^{1}$ 的图的闭包. 有交换图
$$
\xymatrix@R=1.5cm@C=1.5cm{
&\bar{X}\ar[ld]_{h} \ar[rd]^{a} \ar[dd]^{\bar{g}}   &\\
X  \ar[rd]_{g}  &                                  &\PP^1 \ar[ld]_{b}\\
            & **[r]Y              =\Spec \CC&
}
$$
其中 $a$ 和 $b$ 的相对维数均 $<n$. 由归纳假设, 定理对 $(a,.)$ 和 $(b,.)$ 成立, 从而对 $(\bar{g},.)$ 也成立. 事实上, 当 $q=0$ 时这是平凡的; 若 $F$ 为 $\bar{X}$ 上的阿贝尔挠层, 则由下述 Leray 谱序列的态射\pageoriginale 得出:
$$
\xymatrix{
\varepsilon^{\ast}(R^{p}b_{\text{et}\ast})(R^{q}a_{\text{et}\ast})F \ar@{=>}[r] \ar[d]&
\varepsilon^{\ast}(R^{p+q}\bar{g}_{\text{et}\ast})F  \ar[d]\\
(R^{p}b_{\text{cl}\ast}(R^{q}{a_{\text{cl}\ast}}))
\varepsilon^{\ast}F  \ar@{=>}[r] & (R^{p+q}\bar{g}_{\text{cl}\ast})\varepsilon^{\ast}F;
}
$$
最后, 若 $F$ 为 $\bar{X}$ 上的 ind-有限群层, 注意 $\epsilon^{\ast}(R^{1}\bar{g}_{\text{\rm et\,}\ast})$ 是可抹去函子, 故鉴于 $q=0$ 的结果, 只须证明 $(R^{1}\bar{g}_{\cl\ast})\epsilon^{\ast}$ 也可抹去 (\cf \hyperrefext[XII][XII.cor8.2]{8.2}). 这由正合列
$$
0\longrightarrow (R^{1}b_{\cl\ast})a_{\cl
\ast}\epsilon^{\ast}F\longrightarrow (R^{1}\bar{g}_{\cl
\ast})\epsilon^{\ast}F\longrightarrow b_{\cl \ast}(R^{1}a_{\cl
\ast})\epsilon^{\ast}F{\quoi},
$$
以及根据归纳假设 $(R^{1}b_{\cl\ast})\epsilon^{\ast}$ 和 $(R^{1}a_{\cl\ast})\epsilon^{\ast}$ 可抹去这一事实得出.

现在设 $F$ 为 $X$ 上的集合层. 有 $F\hookrightarrow h_{\text{\rm et\,}\ast}h_{\et}{}^{\ast}F=G$. 正合列
\begin{equation*}
F\rightarrow G\rightrightarrows G\amalg_F G{\quoi},\tag{4.3.1}
\end{equation*}\etiquette[4.3.1]{XVI.eq4.3.1}
将问题归结为 (对于 $q=0$) 对 $G$ 证明 \hyperref[XVI.thm4.1]{4.1}, 从而归结为形如 $h_{\text{\rm et\,}\ast}\bar{F}$ 的层. 但对这种层, 结果从对 $(\bar{g},.)$ 和 $(h,.)$ 的定理得出; 而 $h$ 的相对维数 $<n$, 这是直接可见的.

设 $F$ 为 $X$ 上的阿贝尔挠层. 有正合列
\begin{equation*}
0\longrightarrow j_{!}j^{\ast}F\longrightarrow F\longrightarrow
i_{\ast}i^{\ast}F\longrightarrow 0{\quoi}.\tag{$\ast$}
\end{equation*}
由于 $U$ 在 $X$ 中稠密, 有 $\dim Y<n$, 故 \hyperref[XVI.thm4.1]{4.1} 对 $(g|Y,i^{\ast}F)$ 成立\pageoriginale, 从而对 $(g,i_{\ast}i^{\ast}F)$ 成立. 因而由五引理, 只须对 $(g,j_{!}j^{\ast}F)$ 核实 \hyperref[XVI.thm4.1]{4.1}, 而这由对 $\bar{g}$ 的定理 (\hyperref[XVI.lem4.2]{4.2}) 得出.

设 $F$ 为 $X$ 上的 ind-有限群层. 只须证明函子 $(R^{1}g_{\cl\ast})\epsilon^{\ast}$ 对层 $F$ 可抹去. 这由 ($\ast$) 及对 $j_{!}j^{\ast}F$ 和 $i_{\ast}i^{\ast}F$ 的可抹去性得出, 即由对 $(g,j_{!}j^{\ast}F)$ 的定理 (\hyperref[XVI.lem4.2]{4.2}) 和对 $(g,i_{\ast}i^{\ast}F)$ 的定理 (归纳假设) 得出, \cqfd

\begin{enonce*}{引理 4.4}\etiquette[4.4]{XVI.lem4.4}
为在 $f$ 固有时证明 $\hyperref[XVI.thm4.1]{(4.1)}$, 只须证明如下断言: 设 $X$ 为 $\Spec \CC$ 上的完备正则曲线, $F$ 为 $X$ 上的常值可构造集合层 (\resp\ldots). 则态射
$$
H^{q}(X_{\et},F)\longrightarrow H^{q}(X_{\cl},F)
$$
对所考虑的 $q$ 值是双射.
\end{enonce*}

\noindent
\emph{证明}. 由于只须逐纤维核实定理, \hyperref[XVI.lem4.3]{4.3} 将问题归结为 $X$ 固有且维数 $\leq 1$ 的情形. 维数 $0$ 的情形平凡, 从而定理对有限态射成立. 由于 $f$ 固有, $X_{\et}$ 和 $X_{\cl}$ 的上同调与滤过正向极限交换 (\hyperrefext[VII][VII.prop3.3]{3.3} 和 [2], 4.12.1), 因而可归结为 $F$ 可构造的情形 (\hyperrefext[IX][IX.prop2.9]{2.9} (iii)).

应用 \hyperrefext[IX][IX.prop2.14]{2.14}. 对集合层, 正合列 \eqhyperref[XVI.eq4.3.1]{4.3.1} 和 \hyperrefext[IX][IX.prop2.14]{2.14} 表明可取 $F=\pi{}_{\ast}C$, 其中\pageoriginale $\pi:X'\rightarrow X$ 有限, $X'$ 正规从而正则, $C$ 为 $X'$ 上的常值层. 对阿贝尔挠层, 取 $F$ 的一个形如 $\pi_{\ast}C$ 的层的有限乘积组成的消解, 并应用消解的谱序列, 再次归结为 $F=\pi_{\ast}C$ 的情形. 对 ind-有限群层, 回顾只须证明 $(R^{1}g_{\cl\ast})\epsilon^{\ast}$ 对层 $F$ 可抹去, 从而对某个 $F$ 可嵌入其中的层 $G$ 可抹去; 因此又归结为 $F=\pi_{\ast}C$ 的情形. 由 (\hyperrefext[VIII][VIII.prop5.5]{5.5} 和 \hyperrefext[VIII][VIII.rem5.8]{5.8}) 现已明显可以将 $X$ 换成 $X'$, 引理得证.

为处理最后这种情形, 鉴于上同调维数的结果 (\sisi{IX 7.7}{\hyperrefext[IX][IX.cor5.7]{5.7}}) 和 \cite{sga4.exp.XVI.bib.2}, 4.14.1, 显然只须对完备正则曲线证明:
\begin{itemize}
\item[$(q=0)$:] $X$ 连通且非空当且仅当 $X_{\cl}$ 连通且非空.

\item[$(q=1)$:] 基变换态射诱导 $X$ 的平展覆盖范畴与 $X_{\cl}$ 的有限解析平展覆盖范畴之间的等价.

\item[$(q=2)$:] 有同构 $\epsilon:H^2(X_{\et},\mu_n)\xrightarrow{\sim}H^2(X_{\cl},\mu_n)$.
\end{itemize}

前两个断言立即由 GAGA 得出. 对最后一个, 注意 $H^2(X_{\et},\GG_m)=0$ (\hyperrefext[IX][IX.eq4.5]{4.5}). 同样, 若 $\Oo^*$ 表示 $X_{\cl}$ 上的可逆全纯函数层, 则 $H^2(X_{\cl},\Oo^*)=0$. 这由正合列
$$
0\longrightarrow\ZZ\longrightarrow\Oo\xrightarrow{\exp}
\Oo^*\longrightarrow0
$$
及\pageoriginale 根据 GAGA 有 $H^2(X_{\cl},\Oo)=H^2(X_{\et},\Oo_X)=0$ 这一事实得出. 因此, Kummer 正合列 \hyperrefext[IX][IX.Thr3.2]{3.2} (相应地, $X_{\cl}$ 上的正合列
$$
0\rightarrow\mu_n\rightarrow
\Oo^*\xrightarrow{n}\Oo^*\rightarrow0)
$$
给出
\begin{align*}
&H^1(X_{\et},\GG_m)/n\xrightarrow{\sim}
H^2(X_{\et},\mu_n)\\
(\resp\quad &H^1(X_{\cl},\Oo^*)/n\xrightarrow{\sim}
H^2(X_{\cl},\mu_n)){\quoi}.
\end{align*}
由 GAGA, 典范态射
$$
H^1(X_{\et},\GG_m)=\Pic X\longrightarrow
\Pic X_{\cl}=H^1(X_{\cl},\Oo^*)
$$
是双射. 结果得证, 从而证明了 $f$ 固有时的 \hyperref[XVI.thm4.1]{4.1}.

\medskip
\noindent
\emph{$F$ 可构造的情形}. 对 $\dim X$ 作归纳. 设已知结果在维数 $<n$ 时成立, 来证明它在 $\dim X\leq n$ 时成立.

\begin{enonce*}{引理 4.5}\etiquette[4.5]{XVI.lem4.5}
可假设 $f:X\rightarrow S$ 是稠密开浸入, 且 $F$ 常值.
\end{enonce*}

\noindent
\emph{证明}. 仿照 \hyperref[XVI.lem1.3]{1.3} 的证明, 可假设 $X$ 和 $S$ 仿射, 从而 $f$ 拟射影. 再用 \hyperref[XVI.lem4.4]{4.4} 的论证归结为 $F=\pi_*C$, 其中 $\pi:X'\rightarrow X$ 有限且 $C$ 常值. 将 $X$ 换成 $X'$, 因而归结为\pageoriginale $F$ 常值的情形. 可将 $X$ 换成 $X_{\red}$. 取交换图
$$
\xymatrix@=1.25cm{
  X  \ar[rr]^{h}\ar[rd]_{f} &      &\bar{X} \ar[ld]^{\bar{f}}\\
             &     S  &
}
$$
其中 $\bar X$ 约化, $\bar f$ 射影, $h$ 为稠密开浸入. 只须对两个态射 $h$ 和 $\bar f$ 核实定理 (比较 \hyperref[XVI.lem4.3]{4.3} 的证明); 因为 $\bar f$ 固有, 只须对 $h$ 核实即可. 因而归结为稠密开浸入.

\begin{enonce*}{引理 4.6}\etiquette[4.6]{XVI.lem4.6}
还可假设 $S$ 正则.
\end{enonce*}

\noindent
\emph{证明}. 设 $S'\xrightarrow{h}S$ 为“$S$ 的奇点解消”, 即满射、固有、双有理态射, 且 $S'$ 正则, 并设
$$
\xymatrix@C=2cm@R=1.25cm{
X'     \ar[r]^{f'} \ar[d]^{h'}&     S'\ar[d]^{h}\\
X   \ar[r]^{f}  & S
}
$$
为由 $h$ 得到的 Cartesian 图.

设 $F$ 为 $X$ 上的常值集合层. 有 $F\hookrightarrow G=h'_*{h'}^*F$, 从而当 $q=0$ 时, 只须对 $G$ 核实 \hyperref[XVI.thm4.1]{4.1} (比较 \eqhyperref[XVI.eq4.3.1]{4.3.1}), 即对形如 ${h'}_*F'$ 的层核实, 其中 $F'$ 在 $X'$ 上常值. 有交换\pageoriginale 图 (其中省去符号“ét”)
$$
\xymatrix@C=2cm@R=1.5cm{
\varepsilon^*f_* (h'_*F') = \varepsilon^*h_*f'_*F'  \ar[d]^{d}\ar[r]^(.6){a} & h_{\cl *}
	    \varepsilon^*f'_*F' \ar[d]^{b}\\
	    f_{\cl *}\varepsilon^*(h'_*F')  \ar[r]^(.4){c} &
	    f_{\cl *} h'_{\cl *}\varepsilon^*F' = h_{\cl *}f'_{\cl *}
	    \varepsilon^*F',
}
$$
根据固有态射的结果, $a$ 和 $c$ 为双射. 为证明 $d$ 为双射, 只须证明 $b$ 为双射, 即
$$
\epsilon^*f'_*F'\xrightarrow{\sim}
f'_{\cl *}\epsilon^*F'.
$$
因而归结为 $S=S'$ 且 $S$ 正则.

\noindent
(注: \sisi{我们}{我们}用到了固有态射的有限性定理 \hyperrefext[XIV][XIV.thm1.1]{1.1}, 而未明确提及.)

设 $F$ 为 $X$ 上的常值阿贝尔挠层. 只须处理 $F=(\ZZ/n)_X$. 态射 $h'$ 在一个稠密开子集 $j:U\rightarrow X$ 上为同构; 设 $j':U'\rightarrow X'$ 为开子集 ${h'}^{-1}(U)\xrightarrow{\sim}U$, 置 $Y=X-U$. 因为 $U$ 稠密, 有 $\dim Y<n$, 故定理对 $(f,(\ZZ/n)_Y)$ 成立. 正合列
$$
0\longrightarrow j_!(\ZZ/n)_U\longrightarrow
(\ZZ/n)_X\longrightarrow(\ZZ/n)_Y\longrightarrow0
$$
和五引理将问题归结为对 $j_!(\ZZ/n)_U$ 和 $f$ 证明定理, 再由 \hyperref[XVI.lem4.2]{4.2} 归结为对 $(hf',{j'_!}(\ZZ/n)_{U'})$ 证明定理. 有正合列
$$
0\longrightarrow {j'_!}(\ZZ/n)_{U'}\longrightarrow
(\ZZ/n)_{X'}\longrightarrow(\ZZ/n)_{Y'}\longrightarrow0{\quoi},
$$
且\pageoriginale 由归纳假设, 定理对 $(hf',(\ZZ/n)_{Y'})$ 成立. 因而只须对 $(hf',(\ZZ/n)_{X'})$ 证明. 但固有情形给出对 $(h,\cdot)$ 的定理, 而态射对 $h,f'$ 的 Leray 谱序列表明只须对 $(f',\ZZ/n)$ 证明即可. 这证明了阿贝尔情形的引理.

最后, 假设 $F$ 为常值有限群层. 回顾只须证明函子 $(R^1f_{\cl *})\epsilon^*(\mathord\cdot)$ 对 $F$ 可抹去, 从而对某个 $F$ 可嵌入其中的层 $G$ 可抹去; 因此对 $G={h'}_*{h'}^*F$, 即形如 ${h'}_*F'$ 的层可抹去即可, 其中 $F'$ 在 $X'$ 上常值. 有
$$
\epsilon^*({h'}_*F')\xrightarrow{\sim}
h'_{\cl *}(\epsilon^*F')
$$
及
$$
0\longrightarrow
(R^1f_{\cl *})h'_{\cl *}\epsilon^*(\mathord\cdot)\longrightarrow
(R^1(f_{\cl}h'_{\cl})_*)\epsilon^*(\mathord\cdot)
=(R^1(h_{\cl}f'_{\cl})_*)\epsilon^*(\mathord\cdot){\quoi}.
$$
因而只须证明最后的函子对常值层可抹去. 但正合列
$$
0\longrightarrow (R^1h_{\cl *})f'_{\cl *}
\epsilon^*(\mathord\cdot)\longrightarrow
(R^1(h_{\cl}f'_{\cl})_*)\epsilon^*(\mathord\cdot)\longrightarrow
h_{\cl *}(R^1f'_{\cl *})\epsilon^*(\mathord\cdot){\quoi}
$$
表明第一项可抹去, 例如可由 Godement 消解抹去 (\hyperrefext[XII][XII.prop3.3]{3.3}), 因为 $h$ 固有. 因而只须证明右端项可由 Godement 消解抹去, 从而只须 $(R^1f'_{\cl *})\epsilon^*(\mathord\cdot)$ 可如此抹去, 即定理对 $(f',F')$ 成立, 其中 $F'$ 常值. 引理得证.

\noindent
\emph{证明的结束}.\pageoriginale 现在设 $F$ 为常值层, $f:X\rightarrow S$ 为稠密开浸入, 且 $S$ 正则. 集合层的断言立即由 \hyperref[XVI.prop3.2]{3.2} 及经典拓扑的类似的平凡引理得出\sisi{---事实上}{. 事实上}, 它们表明 $\epsilon^*f_*F$ 和 $f_{\cl *}\epsilon^*F$ 为取同一值的常值层, 从而在 $\varphi$ 下同构.

对 $q>0$, 先进一步归结为 $X=S-Y$ 且 $Y$ 正则的情形, 即 $(Y,S)$ 为光滑 $\Spec\CC$-对的情形. 事实上, 写成 $X=S-C_1$, 其中 $C_1$ 赋予诱导的约化结构, 并递归定义

$Y_\nu=$ $C_\nu$ 的正则点组成的稠密开子集,

$C_{\nu+1}=C_\nu-Y_\nu$, 赋予诱导的约化结构.

\noindent
因 $Y_\nu$ 在 $C_\nu$ 中稠密, 有 $\dim C_{\nu+1}<\dim C_\nu$, 从而得到序列
$$
X=X_1\subset X_2=S-C_2\subset X_3=S-C_3\subset\cdots\subset X_r=S{\quoi},
$$
其中包含映射 $i_\nu:X_\nu\rightarrow X_{\nu+1}$ 是开浸入, 且
$$
X_\nu=X_{\nu+1}-Y_\nu,\qquad
Y_\nu\ \mbox{在 $X_{\nu+1}$ 中闭且正则}{\quoi}.
$$
$X$ 上常值层的正像 $i_{\nu *}F$ 常值 (\hyperref[XVI.prop3.2]{3.2}), 且对所考虑的 $q$ 值, $R^qi_{\nu *}F$ 可构造, 集中于 $Y$ 的维数至多 $n-1$ 的子簇上 (\hyperref[XVI.cor3.4]{3.4}, \hyperref[XVI.cor3.6]{3.6}, \hyperref[XVI.thm3.7]{3.7}). 因而归纳假设和 Leray 谱序列 (相应地, 正合列 \hyperrefext[XII][XII.prop3.2]{3.2}) 将问题归结为\pageoriginale $i_\nu$ 的情形, 这就证明了所述归约.

但对光滑对 $(Y,S)$, $(R^qf_*)F$ 的值已就平展拓扑明确算出 (\hyperref[XVI.cor3.4]{3.4}, \hyperref[XVI.cor3.6]{3.6}, \hyperref[XVI.thm3.7]{3.7}); 对经典拓扑作类似的简单计算可得相应结果, 留给读者. 直接比较就证明了对 $f$ 的定理.

\section{特征零代数概形的有限性定理}\etiquette[5]{XVI.sec5}

\begin{enonce*}{定理 5.1}\etiquette[5.1]{XVI.thm5.1}
设 $k$ 为特征零域, $f:X\rightarrow S$ 为 $k$ 上局部有限型概形之间的有限型态射. 设 $F$ 为 $X$ 上可构造的集合层 (\resp 有限群层, \resp 阿贝尔挠层). 则在 $q=0$ (\resp $q=0,1$, \resp 任意 $q$) 时, $R^qf_*F$ 可构造.
\end{enonce*}

要求 $k$ 的特征为零, 是因为将用到奇点解消定理 \cite{sga4.exp.XVI.bib.1}. 在特征 $p>0$ 时, 若假定奇点解消成立, 该证明也有效, \emph{条件是 $F$ 的阶数与 $p$ 互素}. 因而应用 Abhyankar 的结果 \cite{sga4.exp.XVI.bib.3}, 可以在非零特征、$\dim X\leq2$ 时得到一个结果. 除了始终用到的\pageoriginale 固有基变换定理外, 证明的另一个要素是“绝对上同调纯性定理” \hyperref[XVI.cor3.9]{3.9}\SGAArchiveOnly{\nde{冻结底稿将这一后续进展归于 Deligne 和 Gabber 的工作, 但未提供书目目标.}}.

回顾, 对固有态射和任意 $S$, 定理已经证明 (\hyperrefext[XIV][XIV.thm1.1]{1.1}). 此外, 第 XIX 讲中还将 (仍在特征零时) 在弱得多的假设, 即 $S$ 为\emph{优}概形且 $f$ 有限型下得到有限性定理; 方法是对这类概形证明 \hyperref[XVI.rems3.10]{3.10} 所述形式的绝对纯性定理.

在非零特征、维数至少为 $3$ 时, 目前只有下面这个 \hyperrefext[XI][XI.prop3.3]{3.3} 的简单推论:

\begin{enonce*}{定理 5.2}\etiquette[5.2]{XVI.thm5.2}
设 $X$ 为 $\Spec k$ 上的光滑代数概形, 其中 $k$ 为可分闭域, $F$ 为局部常值有限阿贝尔群层 (\resp 有限群层), 阶数与 $k$ 的特征互素. 则对任意 $q$ (\resp $q=0,1$), $H^q(X,F)$ 是有限群 (\resp 有限带基点集合).
\end{enonce*}

\noindent
\emph{\hyperref[XVI.thm5.2]{5.2} 的证明}. 可假设 $k$ 代数闭 (\hyperrefext[VIII][VIII.thm1.1]{1.1}). 对 $n=\dim X$ 作归纳. 由 \hyperrefext[XI][XI.prop3.3]{3.3}, $X$ 有一个由“良好邻域”开子集组成的超覆盖. 超覆盖的谱序列立即将问题归结为 $X$ 为良好邻域的情形, 因而它具有初等纤维化 (\hyperrefext[XI][XI.def3.1]{3.1})
$$
\xymatrix@=1.5cm{
X \ar[rr]^{j} \ar[rd]_{f}&         &
\bar{X}\ar[ld]^{\bar{f}} \\
             &      S  & .
}
$$
由\pageoriginale \hyperref[XVI.prop3.2]{3.2}, \hyperref[XVI.cor3.6]{3.6} 和 \hyperref[XVI.thm3.7]{3.7}, 层 $j_*F$ 在 $\bar X$ 上局部常值, $(R^1j_*)F$ 在 $Y=\bar X-X$ 上局部常值, 且当 $q>1$ 时 $(R^qj_*)F=0$. 由 \hyperref[XVI.thm2.1]{2.1}, 对所考虑的 $q$ 值, $(R^p\bar f_*)(R^qj_*)F$ 是 $S$ 上的局部常值层. 结果由归纳假设和 Leray 谱序列 (\resp 正合列 \hyperrefext[XII][XII.prop3.2]{3.2}) 得出.

\hyperref[XVI.thm5.1]{5.1} 的证明与可构造情形下 \hyperref[XVI.thm4.1]{4.1} 的证明非常接近. 只指出主要步骤. 对 $\dim X=n$ 作归纳. 先应用 \hyperref[XVI.lem1.3]{1.3} 证明中所用的方法, 归结为 $X$ 和 $S$ 仿射, 从而 $f$ 拟射影的情形. 再应用 \sisi{IX 2.15}{\hyperrefext[IX][IX.prop2.14]{2.14}} 和 \hyperrefext[VIII][VIII.prop5.5]{5.5}, \hyperrefext[VIII][VIII.rem5.8]{5.8}, 归结为 $F$ 常值的情形. 有交换图
$$
\xymatrix@=1.5cm{
X \ar[rr]^{h} \ar[rd]_{f}&         & \bar{X}\ar[ld]^{\bar{f}}\\
             &      S  & ,
}
$$
其中 $h$ 为稠密开浸入, $\bar f$ 固有, 故定理对 $\bar f$ 成立 (\hyperrefext[XIV][XIV.thm1.1]{1.1}). 因而只须对 $(h,F)$ 证明定理, 即考虑稠密开浸入\pageoriginale 和常值层. 设 $h:S'\rightarrow S$ 为 $S$ 的奇点解消, 并设
$$
\xymatrix@C=1.75cm@R=1cm{
X  \ar[d]^{f} & X' \ar[l]_{h'} \ar[d]^{f'}\\
S          & S' \ar[l]_{h}
}
$$
为由 $h$ 得到的 Cartesian 图. 可归结为对 $f'$ 证明定理, 从而归结为 $S=S'$ 正则的情形. 在 \hyperref[XVI.thm4.1]{4.1} 证明末尾使用的归纳法, 将问题归结为 $Y=S-X$ 非奇异, 即 $(Y,S)$ 为光滑对的情形. 应用 \hyperref[XVI.prop3.2]{3.2}, \hyperref[XVI.cor3.4]{3.4}, \hyperref[XVI.cor3.6]{3.6} 和 \hyperref[XVI.thm3.7]{3.7} 即完成证明.

\begin{enonce*}[remark]{注 5.3}\etiquette[5.3]{XVI.rem5.3}
刚才概述的 \hyperref[XVI.thm5.1]{5.1} 的证明也适用于如下情形: 给定被整数 $n>0$ 零化的左 Noether 环 $A$ (若不假设特征为零, 必须假设该整数与特征互素; 比较 \hyperref[XVI.thm5.1]{5.1} 后的注), 考虑可构造的左 $A$-模层.
\end{enonce*}

\begin{thebibliography}{99}
\bibitem{sga4.exp.XVI.bib.1} Hironaka, H. Resolution of singularities of an
algebraic variety over a field of characteristic zero, Annals of
Math., vol. 79 (1964), p. 109.

\bibitem{sga4.exp.XVI.bib.2} Godement,\pageoriginale R. Théorie des
\sisi{Faisceaux}{faisceaux}, Paris, 1958\sisi{}{.}

\bibitem{sga4.exp.XVI.bib.3} Abhyankar, S. Local uniformization on algebraic
surfaces over ground fields of characteristic $p\neq0$, Ann. of
Math., vol. 63 (1956), p. 491.

\bibitem{sga4.exp.XVI.bib.4} J.-P. Serre, \sisi{Géométrie Algébrique et Géométrie
Analytique}{Géométrie algébrique et géométrie analytique},
Annales Inst. Fourier 1956, pp. 1--42 (\sisi{引作}{引作} GAGA).
\end{thebibliography}

% 源注: XII 4.1 无引用编号, 原行78; XII 4.1.2 无引用编号, 原行992.
% 源注: 4.11.1 无引用编号, 原行1040; XI 7.7 无引用编号, 原行1210.
% 源注: IX 4.5 为公式引用; IX 2.15 无引用编号, 原行1504.

% END native input: chapters/exp-16-zh.tex
\def\danslelivre{1}
% BEGIN native input: chapters/exp-17-part1-zh.tex
\chapter{固有支集上同调}\MakeLinkTarget*{sga4.exp.XVII}\label{sga4.exp.XVII}
%\addtocontents{toc}{P. Deligne}
\begin{center}
P. Deligne
\end{center}

% The edition is fixed to the corrected (second) reading. Keeping the
% generic ifthen-based selector inside Xy-pic matrices breaks current
% Xy-pic; this equivalent expandable selector avoids that toolchain bug.
\renewcommand{\sisi}[2]{#2}

\setcounter{pageoriginal}{251}

\footnotetext[1]{本讲及下一讲写于 1968 年和 1969 年，重现并补充了
  A.~\textsc{Grothendieck} 的口头讲演（1964 年春）。这份书面稿的作者没有参加
  当时的讨论班，撰写时部分依据了 A.~\textsc{Grothendieck} 的笔记。}


\section*{引言}\pageoriginale

本讲建立固有支集上同调（cohomology with proper supports）的形式理论。
口头讨论班中仅作了合理性说明的协变、反变性质问题，在这里得到了细致的处理。
这也解释了本讲篇幅之长；其中第 \hyperref[XVII.sec1]{1} 至
\hyperref[XVII.sec4]{4} 节讨论范畴和一般拓扑。
强烈建议读者只阅读第 \hyperref[XVII.sec5]{5} 和 \hyperref[XVII.sec6]{6} 节，
本讲的几何实质集中在那里（\S \hyperref[XVII.sec5]{5}：固有支集上同调的构造与
协变、反变性质，基变换，有限性，有限 \tTor-维数及比较定理；
\S \hyperref[XVII.sec6]{6}：拟有限平坦态射的\sisi{迹}{迹}态射理论）。

第 \hyperref[XVII.sec1]{1} 和 \hyperref[XVII.sec2]{2} 节“回顾”导出范畴的若干结果。
\S \hyperref[XVII.sec1]{1} 第 ${\rm n^{\circ}}$ 1 节专门讨论符号规则。
在 \S \hyperref[XVII.sec3]{3}，我们处理两套协变、反变形式理论的黏合问题。
在 \S \hyperref[XVII.sec4]{4}，我们引入平坦解消，并研究 \textsc{Godement}
松层解消的特殊性质。

本讲\pageoriginale 仅在 $f$ 为可紧化态射（compactifiable morphism）
（\hyperref[XVII.def3.2.1]{3.2.1}）时定义固有支集高阶直像函子 $R^{q}f_{!}$
以及产生这些函子的函子 $Rf_{!}$。由 \textsc{B. Saint-Donat} 撰写的附录说明，
如何将定义推广到目标拟紧且拟分离的有限型分离态射。
\nde{\MakeLinkTarget*{sga4.exp.XVII.editor-note.introduction.nagata-compactification}\label{XVII.ndeConrad}
Nagata 紧化定理保证，目标拟紧且拟分离的每个有限型分离态射都可紧化。
见 B.~Conrad，“Deligne's notes on Nagata compactifications”，
\emph{J. Ramanujan Math.\ Soc.} \textbf{22} (2007)，205--257，
下文 \hyperref[XVII.subsec4.1]{4.1} 引为 [C]；另见
“Erratum for `Deligne's notes on Nagata compactifications'”，
\emph{J. Ramanujan Math.\ Soc.} \textbf{24} (2009)，427--428。}

\S \hyperref[XVII.sec5]{5} 第 ${\rm n^{\circ}}$ \hyperref[XVII.subsec5.5]{5} 节
（对称积的上同调）在本讨论班的后文中不再使用。它将在 \SGA 5 中用于加强
$L$-函数的有理性定理。\nde{见 \SGA 4$\frac{1}{2}$，
模 $\ell^n$ 和模 $p$ 的 $L$-函数。}

\S \hyperref[XVII.sec6]{6} 第 ${\rm n^{\circ}}$ \hyperref[XVII.subsec6.3]{3} 节
（\emph{连续}系数的迹理论）在第 XVIII 讲中仅用于光滑群这一相对简单的情形。
若读者关心 Poincaré 对偶定理（整体对偶），并愿意接受一个超越论证，
甚至可以完全不读此 \S \hyperref[XVII.sec6]{6} 第 ${\rm n^{\circ}}$
\hyperref[XVII.subsec6.3]{3} 节，以及 XVIII 中 \S \refnoexpose{XVIII.sec1} 的大部分内容。

\setcounter{section}{-1}
\section{术语预备}\etiquette[0]{XVII.sec0}

\subsection{}\etiquette[0.1]{XVII.subsec0.1}
两个表示范畴中对象的符号组之间的符号 $=$，有时将滥用记号，表示这些对象典范同构。
原则上，范畴与该典范同构都应事先定义。在图中，符号 $=$ 此时表示该同构本身。

\hypertarget{sga4.exp.XVII.subsection.0.2}{}\setcounter{subsection}{2}
\setcounter{subsubsection}{-1}
\subsubsection{}\etiquette[0.2.0]{XVII.subsubsec0.2.0}
设 $f:S\rightarrow S'$ 是景（site）的态射
（\hyperrefext[IV][IV.subsubsec4.9.3]{4.9.3}）。若 $U$ 与 $U'$ 分别是 $S$ 和 $S'$
的对象，则定义从 $U$ 到 $U'$ 的 $f$-态射为从 $U$ 到 $f^{\ast}U'$ 的态射。
对概形的平展景，这就回到通常的概念。

\subsubsection{}\etiquette[0.2.1]{XVII.subsubsec0.2.1}
若\pageoriginale $\mathcal{F}$ 和 $\mathcal{F}'$ 分别是 $S$ 和 $S'$ 上的层，
则从 $\mathcal{F}$ 到 $\mathcal{F}'$ 的 $f$-态射等价地指：
\begin{enumerate}
\renewcommand{\theenumi}{\roman{enumi}}
\renewcommand{\labelenumi}{(\theenumi)}
\item 从 $f^{\ast}\mathcal{F}'$ 到 $\mathcal{F}$ 的态射；

\item 从 $\mathcal{F}'$ 到 $f_{\ast}\mathcal{F}$ 的态射；

\item 一个函数，它给从 $S$ 的对象 $U$ 到 $S'$ 的对象 $U'$ 的每个 $f$-态射
  $\varphi$ 赋予一个从 $\mathcal{F}'(U')$ 到 $\mathcal{F}(U)$ 的函数，
  且这一赋予与 $\varphi$ 同 $S$ 或 $S'$ 中箭头的复合相容
  \footnote{这些术语不与将 $S$ 的对象等同于其伴随层的做法相容。}。
\end{enumerate}

由 (i)、(ii) 和 (iii)，集合层在景的范畴上组成一个纤维化且余纤维化范畴。

赋环景上的左模层同样如此。概形上的平展层也一样；这并非立即就是上述陈述的特例，
因为景实际上组成一个 2-范畴，而“$X$ 的平展景”仅是关于 $X$ 的一个伪函子
（\hyperrefext[VII][VII.subsec1.4]{1.4}）。

上述定义使变化的景上的集合层（\resp\ldots）成为景的范畴上的一个纤维化范畴，
其纤维范畴是通常集合层（\resp\ldots）范畴的\emph{反范畴}。

\subsection{}\etiquette[0.3]{XVII.subsec0.3}
设 $I$ 是有限集，$\epsilon$ 是 $I$ 上取值于 $\{+1,-1\}$ 的函数，
$(\Aa_i)_{i\in I}$ 是一族由平移函子 $T_i$ 分次的加性范畴
（\cite{sga4.exp.XVII.bib.CD} I \sisi{1.10}{1.1.0}），$F$ 是从诸范畴 $\Aa_i$
到加性范畴 $\Aa$ 的多重加性多函子，后者由平移函子 $T$ 分次。
假设 $F$ 对于满足 $\epsilon_i=1$（\resp $\epsilon_i=-1$）的变量 $i$
是协变（\resp 反变）的。修改并补充\pageoriginale
（\cite{sga4.exp.XVII.bib.CD} \sisi{I 1.1}{II 2.1.1}），若给定一族
$(\varphi_i)_{i\in I}$，其中每个都是函子 $T\circ F$ 与
$F\circ T_i^{\varepsilon_{i}}$ 之间的同构，且对于 $I$ 中的 $i\neq j$，
下面的图反交换，则称 $F$ 是\emph{分次函子}（graded functor）：
\[
\xymatrix{
F \circ (T_i^{\varepsilon_i}, T_j^{\varepsilon_j})  \ar[rr]^{\varphi_j
  \ast T_i^{\varepsilon_i}} \ar[dd]^{\varphi_i\ast T_j^{\varepsilon_j}}& &T\circ F \circ T_i^{\varepsilon_i} \ar[dd]^{T\ast
  \varphi_i}\\
 & - & \\
 T\circ F \circ T_j^{\varepsilon_j}  \ar[rr]^{T\ast \varphi_j} & &T\circ T
  \circ F.
}
\]
两个分次函子的复合的定义及其仍为分次函子的验证，留给读者。

\subsection{}\etiquette[0.4]{XVII.subsec0.4}
设 $\Aa$ 是加性范畴，$I$ 是有限集。对 $i\in I$，以 $1_i$ 表示
$\ZZ^{I}$ 的第 $i$ 个基向量。$\Aa$ 的一个\emph{$I$-重复形}
（\resp 一个\sisi{朴素}{\emph{朴素}}\emph{$I$-重复形}）由下述数据组成：
\begin{enumerate}
\renewcommand{\theenumi}{\alph{enumi}}
\renewcommand{\labelenumi}{(\theenumi)}
\item $\Aa$ 的一族对象 $(K^{k})_{k\in \ZZ^{I}}$；

\item 对每个 $i\in I$ 和每个 $k\in \ZZ^{I}$，一个箭头
  $d^{k}_i:K^{k}\rightarrow K^{k+1_i}$，满足 $d^{k+1_i}_i d^{k}_i=0$，
  且当 $i\neq j$ 时，$d^{k+1_j}_i d^{k}_j+d^{k+1_i}_j d^{k}_i=0$
  （\resp $d^{k+1_j}_i d^{k}_j=d^{k+1_i}_j d^{k}_i$）。
\end{enumerate}

\subsection{}\etiquette[0.5]{XVII.subsec0.5}
范畴 $\sisi{C}{\Cc}$ 的一个\pageoriginale \emph{\ppro-对象}，是从
$\sisi{C}{\Cc}$ 到 $(\Ens)$ 的一个函子，它是可表函子的滤过正向极限
（指标范畴为小滤过范畴）（\cf \hyperrefext[I][I.sec8.10]{8.10}）。
每个 \ppro-对象都是以小滤过偏序集为指标的可表函子的正向极限。
若 $X_{i}$ 是 $\sisi{C}{\Cc}$ 中以小滤过范畴为指标的对象逆系统，
则以 $\quotprojlim X_{i}$ 表示 \ppro-对象 $\varinjlim h^{X_i}$。

\subsection{}\etiquette[0.6]{XVII.subsec0.6}
读者可将 \hyperref[XVII.subsec0.5]{0.5} 对偶化，得到 \iind-对象的情形
（\Cf \hyperrefext[I][I.sec8.2]{8.2}）。

\subsection{}\etiquette[0.7]{XVII.subsec0.7}
按照新术语，过去称为预概形的对象现在称为\emph{概形}，
过去称为概形的对象现在称为\emph{分离概形}。

\subsection{}\etiquette[0.8]{XVII.subsec0.8}
\emph{赋环概形}范畴的对象是平展景上赋有环层的概形；从 $(S,\Aa)$ 到 $(T,\Bb)$
的箭头是一个二元组，由从 $S$ 到 $T$ 的概形态射 $f$ 以及从 $f^{\ast}\Bb$
到 $\Aa$ 的环层同态 $\varphi$ 组成。称概形态射 $f$ \emph{由} $(f,\varphi)$ \emph{诱导}。

\subsection{}\etiquette[0.9]{XVII.subsec0.9}
若 $X$ 是 $Y$ 上的概形，则以 $(X/Y)^n$ 表示 $X$ 在 $Y$ 上的 $n$ 重纤维积。

\subsection{}\etiquette[0.10]{XVII.subsec0.10}
设 $S$ 是概形。$S$ 的\emph{大平展景}（\resp \emph{大 fppf 景}，
\resp \emph{大 fpqc 景}）是景 $\Sch/{S}$，赋予平展拓扑
（\resp fppf 拓扑，\resp fpqc 拓扑）（\SGA 3 IV 6.3）。
务必注意，它不是 $\sisi{U}{\Uu}$-景（$\sisi{U}{\Uu}$ 是
\sisi{有限}{固定的}宇宙，而 $\Sch/{S}$ 由 $\in \sisi{U}{\Uu}$ 的概形组成）。

$S$ 的\pageoriginale \emph{小 fpqc 景}（\resp \emph{小 fppf 景}）
是 $S$ 上平坦（\resp 平坦且有限表现）概形所成的景，赋予 fpqc 拓扑
（\resp fppf 拓扑）。需要避免歧义时，$S$ 的\emph{小平展景}指景
$S_{\et}$（\hyperrefext[VII][VII.subsec1.2]{1.2}）。

\subsection{}\etiquette[0.11]{XVII.subsec0.11}
在许多地方，本应说本质交换的图，即相差同构意义下交换的图时，
作者径直说成了景态射的交换图（\cf IV \hyperref[IV.subsubsec3.2.2]{3.2.2}）。
读者可验证，我们给出的论证也适用于一般情形。

\subsection{}\etiquette[0.12]{XVII.subsec0.12}
A.~Grothendieck 在若干地方不动声色地引入了“\emph{凝聚概形}”这一术语，
用来表示“拟紧拟分离概形”。

\subsection{}\etiquette[0.13]{XVII.subsec0.13}
概形 $S$ 上的挠层 $F$ 称为\emph{与 $S$ 的剩余特征互素}的，
如果它是其被 $S$ 上可逆整数 $n$ 零化的子层的正向极限。

\section{导出范畴}\etiquette[1]{XVII.sec1}

\subsection{正合函子（符号规则）}\etiquette[1.1]{XVII.subsec1.1}

\subsubsection{}\etiquette[1.1.1]{XVII.subsubsec1.1.1}
关于导出范畴的基本定理，请参阅 \textsc{Verdier} \cite{sga4.exp.XVII.bib.CD}
及 \cite{sga4.exp.XVII.bib.V2}。本小节特别仔细地处理符号问题。

回顾，若 $f:X\rightarrow Y$ 是复形的态射（所在范畴为加性范畴，以下对此不再说明），
则它的\emph{锥}（cone）$C(f)$\pageoriginale 定义为
$$\etiquette[1.1.1.1]{XVII.eq1.1.1.1}
C(f)^{n}=X^{n+1}\oplus Y^{n}\quad d^{n}=-d^{n+1}_X+f^{n+1}+d^{n}_Y{\quoi}.\leqno(1.1.1.1)
$$
当 $X=0$（\resp $Y=0$）时，有 $C(f)=Y$（\resp $C(f)=X[1]$），
从而交换图
\[
\xymatrix{
0        \ar@{=}[r] \ar[d]&   0   \sisi{\ar@{=}[r]}{\ar[r]} \ar[d]&  X
\ar@{=}[r]\ar[d]^{f}  & X \ar[d]\\
X        \ar[r]^{f}   &  Y    \ar@{=}[r] & Y     \ar[r]        & 0
}
\]
定义一个“三角”
\[
X\xrightarrow{f}Y\rightarrow C(f)\xrightarrow{j}X[1]{\quoi}.
\]

在复形模同伦的范畴中，若一个三角同构于上述类型的三角，则称它为\emph{特殊三角}
（distinguished triangle）；若将各箭头的符号改变后成为特殊三角，则称为
\emph{反特殊三角}（antidistinguished triangle）。容易验证，由逐次分裂的复形正合列
所定义的三角（\cite{sga4.exp.XVII.bib.CD} 第~10 页，\cf
\hyperref[XVII.lem1.1.5.4]{1.1.5.4} 的证明）是反特殊三角；
而在同构意义下，每个反特殊三角都可如此得到。

\subsubsection*{${}^{\ast}$1.1.2.}\etiquette[*1.1.2]{XVII.subsubsec*1.1.2}
使用下述“形式规则”进行符号计算最为方便。
将 $C(f)^{n}$ 的元素写成 $1\otimes x+y$ $(x\in X^{n+1}, y\in Y^{n})$，
其中 $1$ 是一个维数为 $1$ 的“胞腔”，并令 $d(1)\otimes x=f(x)$。于是
\[
d(1\otimes x+y)=d1\otimes x-1\otimes dx+dy=-1\otimes
dx+f(x)+dy{\quoi}._{\ast}
\]

\begin{enonce*}{定义 1.1.3}\etiquette[1.1.3]{XVII.def1.1.3}
\begin{enumerate}
\renewcommand{\theenumi}{\roman{enumi}}
\renewcommand{\labelenumi}{(\theenumi)}
\item 从三角范畴 $\Aa$ 到三角范畴 $\Bb$ 的\pageoriginale 分次函子
  （\hyperref[XVII.subsec0.3]{0.3}），若将特殊三角映为特殊三角，则称为正合的。

\item 从三角范畴 $\Aa$ 到三角范畴 $\Bb$ 的反变分次函子，若对 $\Aa$ 的每个特殊三角
  $(X,Y,Z,u,v,w)$，三角 $(F(Z),F(Y), F(X), F(v),F(u),TF(w))$ 都是特殊三角，
  则称为\emph{正合}的（此定义中借助 $F$ 的分次，将 $FX$ 等同于 $TFTX$）。

\item 设 $(\Aa_i)_{i\in I}$ 是三角范畴的有限族，$\epsilon$ 是从 $I$ 到
  $\{+1,-1\}$ 的函数，$\Aa$ 是三角范畴。设 $F$ 是从诸 $\Aa_i$ 的乘积到
  $\Aa$ 的分次函子，对满足 $\epsilon(i)=1$ 的 $i$ 协变，对满足 $\epsilon(i)=-1$
  的 $i$ 反变。若对每个 $i\in I$，将其余变量固定后，
  由 $F$ 得到的从 $\Aa_i$ 到 $\Aa$ 的函子都是正合函子，则称该函子\emph{正合}。
\end{enumerate}
\end{enonce*}

正合函子的复合仍为正合函子。

\setcounter{subsubsection}{3}
\subsubsection{}\etiquette[1.1.4]{XVII.subsubsec1.1.4}
若 $(K^{n_1\ldots n_p}, d_1,\ldots,d_p)$ 是多重复形
（\hyperref[XVII.subsec0.4]{0.4}），则其伴随单复形定义为
$$\etiquette[1.1.4.0]{XVII.eq1.1.4.0}
K^{n}=\prod_{\sum n_i=n}K^{(n_i)}\quav d=\sum d_i{\quoi}.\leqno(1.1.4.0)
$$

在个别情况下若采用直和而非乘积来定义，将明确说明。

设 $I$ 是有限集。按照 \textsc{Cartan--Eilenberg} 和 \textsc{J.~P. Serre}，
我们说明如何给每个朴素 $I$-重复形 $K$ 典范地\pageoriginale 伴随一个
$I$-重复形（\hyperref[XVII.subsec0.4]{0.4}）。

对 $I$ 上的每个全序 $<$，设 $K(<)$ 为下述 $I$-重复形：
$$\etiquette[1.1.4.1]{XVII.eq1.1.4.1}
\begin{cases}
K(<)^{\sheafk }=K^{\sheafk }\\
d(<)_i^{\sheafk }=(-1)^{\sum\limits_{j<i} k_j} d_i{\quoi}.
\end{cases}\leqno{(1.1.4.1)}
$$

若 $<_1$ 和 $<_2$ 是 $I$ 上的两个\sisi{全}{全}序，则定义
\emph{$K(<_1)$ 与 $K(<_2)$ 之间的典范同构 $\tau$}为下式给出的同构：
$$\etiquette[1.1.4.2]{XVII.eq1.1.4.2}
\tau^{\sheafk }=(-1)^{\epsilon(<_1,<_2)},\quad\mbox{\rm 其中}\quad
\epsilon(<_1,<_2)=\sum_{\substack{i<_1j\\ \sisi{i<_2 j}{j<_2 i}}}k_ik_j{\quoi}.\leqno{(1.1.4.2)}
$$

\subsubsection*{1.1.4.2.1}\etiquette[1.1.4.2.1]{XVII.par1.1.4.2.1}
容易验证，当 $<$ 遍历 $I$ 上的全序时，这些典范同构在各 $K(<)$ 之间组成一个传递同构系。
因此，这些 $K(<)$ 决定一个在唯一同构意义下唯一的 $I$-重复形，本小节将其记为
$K(\ast)$，称为\emph{朴素 $I$-重复形 $K$ 的伴随 $I$-重复形}。
注意，$K(\ast)^{\sheafk}$ \emph{并不典范同构于 $K^{\sheafk}$：
只有在 $I$ 上选定全序，才能将这两个对象等同。}

若 $K$ 是朴素 $I$-重复形，$i\in I$，则以 $K[1_i]$ 表示下述朴素复形，
其微分记作 $d_j[1_i]$：
$$\etiquette[1.1.4.3]{XVII.eq1.1.4.3}
\begin{cases}
K[1_i]^{\sheafk } = K^{\sheafk +1_i}\\
d_j[1_i]^{\sheafk } = d_j^{\sheafk +1_{i}}\quad{\mbox{\rm 若 }} i\neq
j\\
d_i[1_i]^{\sheafk }= -d_i^{\sheafk +1_i}{\quoi}.
\end{cases}\leqno{(1.1.4.3)}
$$

若\pageoriginale $K$ 是 $I$-重复形，则令
$$\etiquette[1.1.4.4]{XVII.eq1.1.4.4}
\begin{cases}
K[1_i]^{\sheafk } = K^{\sheafk +1_i}\\
d_j[1_i]^{\sheafk } = -d_j^{\sheafk +1_i}\sisi{}{.}
\end{cases}\leqno{(1.1.4.4)}
$$

若 $K$ 是朴素 $I$-重复形，$<$ 是 $I$ 上的全序，从而 $K(\ast)$ 由
\eqhyperref[XVII.eq1.1.4.1]{1.1.4.1} 给出，则用公式
$$\etiquette[1.1.4.5]{XVII.eq1.1.4.5}
\sigma^{\sheafk }=(-1)^{\sum\limits_{j<i}k_j} \id_{K^{\sheafk +1_i}}{\quoi}.\leqno{(1.1.4.5)}
$$
定义 $K[1_i](\ast)$ 与 $K(\ast)[1_i]$ 之间的同构 $\sigma$。

\setcounter{subsubsection}{4}
\setcounter{paragraph}{5}
\paragraph{}\etiquette[1.1.4.6]{XVII.par1.1.4.6}
容易验证，这个同构与全序 $<$ 的选取无关，因而定义一个
$K[1_i](\ast)$ 与 $K(\ast)[1_i]$ \emph{之间的同构，称为典范同构}。

\paragraph{}\etiquette[1.1.4.7]{XVII.par1.1.4.7}
对任意 $i,j\in I$，若 $K$ 是 $I$-重复形或朴素 $I$-重复形，
则以显然的方式将 $K[1_i][1_j]$ 与 $K[1_j][1_i]$ 等同。

此时，对于 $I$ 中的 $i\neq j$，下列由 \hyperref[XVII.par1.1.4.6]{1.1.4.6}
或 \hyperref[XVII.par1.1.4.7]{1.1.4.7} 型典范同构组成的图\emph{反交换}：
$$\etiquette[1.1.4.7.1]{XVII.eq1.1.4.7.1}
\vcenter{
\xymatrix{
K[1_i][1_j](\ast) \ar@{=}[r] \ar@{<->}[d]& K[1_j][1_i](\ast)
\ar@{<->}[r] & K[1_j](\ast)[1_i]\ar@{<->}[d] \\
K[1_i](\ast)[1_j] \ar@{<->}[r] & K(\ast)[1_i][1_j]  \ar@{=}[r]& K(\ast)[1_j][1_i]\rlap{.}
}}\leqno{(1.1.4.7.1)}
$$


\subsubsection{}\etiquette[1.1.5]{XVII.subsubsec1.1.5}
设\pageoriginale $F:\Aa\rightarrow \Bb$ 是加性范畴之间的反变加性函子。
它的\emph{对复形的延拓} $F:C(\Aa)\rightarrow C(\Bb)$ 定义如下。
若 $K$ 是 $\Aa$ 的复形，则令
$$\etiquette[1.1.5.1]{XVII.eq1.1.5.1}
\begin{cases}
F(K)^{k}=F(K^{-k})\\
d^{k}_{F(K)}=(-1)^{k+1}F(d^{-k-1})\sisi{}{.}
\end{cases}\leqno(1.1.5.1)
$$

\emph{$F(K)[1]$ 与 $F(K[-1])$ 之间的典范同构 $\sigma$}定义为
$$\etiquette[1.1.5.2]{XVII.eq1.1.5.2}
\sigma^{\sisi{\sheafk }{k}}=(-1)^{k+1}\id_{F(K^{-k-1})}{\quoi}.\leqno(1.1.5.2)
$$

若 $F:\Aa\rightarrow \Bb$ 和 $G:\Bb\rightarrow \sisi{C}{\Cc}$
是两个反变加性函子，则用下式定义 $G{\circ}F(K)$ \emph{与} $G(F(K))$
之间的典范同构 $\rho$：
$$\etiquette[1.1.5.3]{XVII.eq1.1.5.3}
\rho^{k}=(-1)^{k}\sisi{}{\id_{G\circ F(K^k)}}{\quoi}.\leqno(1.1.5.3)
$$

\begin{enonce*}{引理 1.1.5.4}\etiquette[1.1.5.4]{XVII.lem1.1.5.4}
\begin{enumerate}
\renewcommand{\theenumi}{\roman{enumi}}
\renewcommand{\labelenumi}{(\theenumi)}
\item 对每个反变加性函子 $F:\Aa\rightarrow \Bb$，公式
  \eqhyperref[XVII.eq1.1.5.1]{1.1.5.1} 和 \eqhyperref[XVII.eq1.1.5.2]{1.1.5.2}
  定义一个从 $K(\Aa)$ 到 $K(\Bb)$ 的正合反变函子。

\item 对每对可复合的反变加性函子 $(G,F)$，公式
  \eqhyperref[XVII.eq1.1.5.3]{1.1.5.3} 定义一个分次函子同构，
  介于 $GF$ 对复形的延拓与 $G$、$F$ 各自对复形的延拓的复合之间。
\end{enumerate}
\end{enonce*}

\begin{enonce*}[remark]{证明}
一个\pageoriginale 次数为 $i$ 的态射 $f:K\rightarrow L$，可以看成下列任一种数据：
\begin{enumerate}
\renewcommand{\theenumi}{\alph{enumi}}
\renewcommand{\labelenumi}{(\theenumi)}
\item 从 $K$ 到 $L[i]$ 的态射；

\item 从 $K[-i]$ 到 $L$ 的态射；

\item 一族 $f^{n}:K^{n}\rightarrow L^{n+i}$，按照 $i$ 的奇偶性与微分交换或反交换。
\end{enumerate}
\end{enonce*}

若 $f:K\rightarrow L[i]$ 是次数为 $i$ 的态射，则
$F(f):F(L)[-i]\rightarrow F(K)$ 仍是次数为 $i$ 的态射；对于 $i=1$，有
$$\etiquette[1.1.5.5]{XVII.eq1.1.5.5}
F(f)^{n}=(-1)^{n+1}F(f^{-n-1}){\quoi}.\leqno(1.1.5.5)
$$

若 $0\rightarrow K\xrightarrow{u} L\xrightarrow{v} M\rightarrow 0$
是逐次分裂的复形正合列，则选定一个分裂，便可将 $L$ 的微分写成
\[
d_L=d_K+d_M+f{\quoi},
\]
其中 $f$ 是从 $M$ 到 $K$ 的一次态射。类似地，对偶正合列
\[
0\longrightarrow F(M)\longrightarrow F(L)\longrightarrow
F(K)\longrightarrow 0
\]
定义一个从 $F(K)$ 到 $F(\sisi{L}{M})$ 的一次态射 $f^{\ast}$。
有 $f^{\ast}=F(f)$。由此可知，短正合复形列定义的三角 $(K,L,M,u,v,f)$
在 $F$ 下的像仍是同一类型的三角，而由 \hyperref[XVII.subsubsec1.1.1]{1.1.1}，
这就证明了 (i)。

在 (ii) 的假设下，对一次的 $f$，有
$G(F(f))^{n}=(-1)^{n+1}(-1)^{-n}G\circ F(f)^{n}=-G\circ F(f)^{n}$；
特别地，分次同构 $G(F(K[1]))\simeq G(F(K))[1]$ 是各分量均为 $-1$ 的同构。

因而\pageoriginale 图
\[
\xymatrix{
G\circ F(K)[1]  \ar@{<->}[r]^{\rho[1]} \ar@{=}[d] & G(F(K))[1] \ar@{<->}[d]\\
G\circ F(K[1])  \ar@{<->}[r]^{\rho} & G(F(K[1]))
}
\]
交换，这就证明了 (ii)。

\subsubsection{}\etiquette[1.1.6]{XVII.subsubsec1.1.6}
设 $s=+$ 或 $s=-$；设 $I$ 是有限集，$\epsilon:I\rightarrow \{+,-\}$ 是函数，
$(\Aa_i)_{i\in I}$ 和 $\Aa$ 是加性范畴，$F$ 是从诸范畴 $\Aa_i$ 到
$\Aa$ 的多重加性多函子，对于满足 $\epsilon(i)=+$ 的 $i$ 所标记的变量协变，
对于其余变量反变。我们将定义 $F$ 对复形的延拓 $K^{s}(F)$（或简称 $F$），
它是从诸范畴 $K^{s\epsilon(i)}(\Aa_i)$（$s\epsilon(i)=\pm$，遵循符号规则）
到 $K^{s}(\Aa)$ 的正合多函子\footnote[1]{这个构造也给出 $F$ 对复形的延拓，
  它是从诸范畴 $C^{s\epsilon(i)}(\Aa_i)$ 到 $C^{s}(\Aa)$ 的多函子。}，
其协变、反变性质与 $F$ 相同。然后我们将说明，这一构造在何种意义下与函子的复合相容。
对于\emph{次数}还可以采用别的约定，后文也会使用；这个问题与此处所考虑的\emph{符号}问题无关。

\paragraph{}\etiquette[1.1.6.1]{XVII.par1.1.6.1}
若 $F$ 是单变量反变函子，则 $K^{s}(F)$ 由 \hyperref[XVII.subsubsec1.1.5]{1.1.5}
的公式定义。

\paragraph{}\etiquette[1.1.6.2]{XVII.par1.1.6.2}
若 $F$ 协变，则函子 $K^{s}(F)$ 由以下函子的复合取商而得：
从诸范畴 $C^{s}(\Aa_i)$ 的乘积到朴素 $I$-重复形范畴的显然函子；
从朴素 $I$-重复形范畴到\pageoriginale $I$-重复形范畴的函子
（\hyperref[XVII.subsubsec1.1.4]{1.1.4}，\hyperref[XVII.par1.1.4.2.1]{1.1.4.2.1}）；
以及伴随单复形函子 \eqhyperref[XVII.eq1.1.4.0]{1.1.4.0}：
\[
K^{s}(F)=C^{s}(F)(\ast)_{s.a}{\quoi}.
\]

设 $K_j\in K^{s}(\Aa_j)$ 是复形，$i\in I$；当 $j\neq i$ 时令
$K_j[1_i]=K_j$，并令 $K_i[1_i]=K_{i}[1]$。此时有显然的同构
\[
C^{s}(F)(K_{j}[1_i])\simeq C^{s}(F)(K_j)[1_i]{\quoi},
\]
以及
\[
C^{s}(F)(K_j)(\ast)[1_i]_{s.a}\simeq
C^{s}(F)(K_j)(\ast)_{s.a}[1]{\quoi},
\]
从而由 \eqhyperref[XVII.eq1.1.4.5]{1.1.4.5} 得到
$K^{s}(F)(K_{j}[1_i])$ 与 $K^{s}(F)(K_{j})[1]$ 之间的\emph{典范同构} $\sigma$。
由 \hyperref[XVII.par1.1.4.7]{1.1.4.7} 可知，这些同构定义 $K^{s}(F)$ 的一个分次。

\paragraph{}\etiquette[1.1.6.3]{XVII.par1.1.6.3}
在一般情形中，$F$ 是典范反变函子 $\Aa_i\rightarrow \Aa^{\circ}_i$
（对 $\epsilon(i)=-$）与一个协变函子 $F^{+}$ 的复合；利用
\hyperref[XVII.par1.1.6.1]{1.1.6.1} 和 \hyperref[XVII.par1.1.6.2]{1.1.6.2}，
通过复合定义函子 $K^{s}(F)$。

\begin{enonce*}{命题 1.1.7}\etiquette[1.1.7]{XVII.prop1.1.7}
在上述假设下，函子 $K^{s}(F)$ 是正合多重加性函子。
\end{enonce*}

只须在 \hyperref[XVII.par1.1.6.1]{1.1.6.1} 的假设下
（此时命题就是 \hyperref[XVII.lem1.1.5.4]{1.1.5.4}~(i)），
或在 \hyperref[XVII.par1.1.6.2]{1.1.6.2} 的假设下验证 \hyperref[XVII.prop1.1.7]{1.1.7}。
要验证将除第 $i$ 个变量以外的所有变量固定后，由 $K^{s}(F)$ 得到的函子正合，
可以选择 $I$ 上以 $i$ 为最小元素的全序，显式描述 $K^{s}(F)$；
此时，标准\pageoriginale 形式 \eqhyperref[XVII.eq1.1.1.1]{1.1.1.1} 的特殊三角
便映为同一形式的特殊三角。

\setcounter{subsubsection}{7}
\subsubsection{}\etiquette[1.1.8]{XVII.subsubsec1.1.8}
设 $I$ 和 $J$ 是两个有限集，$\Psi:I\rightarrow J$ 是满射，
$\epsilon_{I}:I\rightarrow \{+,-\}$ 和

$\epsilon_{J}:J\rightarrow \{+,-\}$ 是函数，并设给定多重加性多函子，
它们对满足 $\epsilon=+$ 的变量协变，对其余变量反变：
\begin{align*}
G_j&: \prod_{\Psi(i)=j}\Aa_i\longrightarrow
\Bb_j\qquad (j\in J)\\
F&:  \prod_{j\in J} \Bb_j\longrightarrow \sisi{C}{\Cc.}
\end{align*}

在 $I$ 和 $J$ 上选取全序，使 $\Psi$ 是递增映射。
我们用下式定义\emph{$F\circ (G_j)$ 对复形的延拓，与 $F$ 及各 $G_j$
对复形的延拓的复合之间的典范同构 $\rho$}：
$$\etiquette[1.1.8.1]{XVII.eq1.1.8.1}
\rho^{\sheafk }=\sisi{\prod_i \eta_i^{k_i}}{(-1)^{A(\sheafk )}}:\; \mbox{\rm 下式的自同构：}
F\circ (G_j)(K_{i}^{k_i\epsilon\sisi{}{_I}(i)\epsilon\sisi{}{_J}(\Psi(i))}){\quoi}.\leqno(1.1.8.1)
$$%
其中 \sisi{若 $\epsilon_I(i)=\epsilon_{J}(\Psi(i))=-$，则 $\eta_i=-1$，
否则 $\eta_i=1$。}{\[
A(\sheafk )=\sum_{\epsilon_I(i)=\epsilon_J(\Psi(i))=-} k_i + \sum_{\epsilon_J(j)=-} \sum_{\substack{\Psi(a)=\Psi(b)=j\\a<b}} k_a k_b.
\]}%See the Erratum in SGA 4 1/2, p. 312

此公式推广了 \eqhyperref[XVII.eq1.1.5.3]{1.1.5.3}；容易验证，$\rho$ 不依赖于
$I$ 与 $J$ 上所选的全序，并且是分次函子的同构
（\cf \hyperref[XVII.lem1.1.5.4]{1.1.5.4} (ii)）。

这些同构 $\rho$ 对三重复合满足显然的上循环条件。

\subsubsection{}\etiquette[1.1.9]{XVII.subsubsec1.1.9}
设\pageoriginale $I$ 是有限集，$i$ 和 $j$ 是 $I$ 的两个不同元素，
$I'=I\backslash\{i,j\}$，$(\Aa_i)_{i\in I}$ 与 $\Aa$ 是加性范畴，
$F:(\Aa_i)_{i\in I}\rightarrow \Aa$ 和 $G:(\Aa_i)_{i\in I'}\rightarrow \Aa$
是多重加性多函子。假设 $\Aa_i=\Aa_j$，且 $F$ 对第 $i$ 个变量反变，
对第 $j$ 个变量协变。对 $X$、$Y\in \Ob \Aa_i$，以 $F_{X,Y}$ 表示
将第 $i$ 和第 $j$ 个变量固定后，得到的从 $(\Aa_i)_{i\in I'}$ 到 $\Aa$ 的多函子：
\[
F_{X,Y}(\ldots)=F(\ldots\sisi{}{,} X,\ldots,Y,\ldots){\quoi}.
\]
一个\emph{缩并态射}（contraction morphism）$c:F\rightarrow G$ 由以下数据组成：
对每个 $X\in \Aa_i$，给定函子态射
\[
c_{X}:F_{X,X}\longrightarrow G{\quoi},
\]
使得对 $\Aa_i$ 的每个箭头 $f:X\rightarrow Y$，图
\[
\xymatrix{
F_{Y,X} \ar[d] \ar[r] & F_{X,X} \ar[d]^{c_X}\\
F_{Y,Y}  \ar[r]_{c_Y}     & G
}
\]
交换。

\hyperref[XVII.subsubsec1.1.6]{1.1.6} 的约定来源于下面的引理。

\begin{enonce*}{引理 1.1.9.1}\etiquette[1.1.9.1]{XVII.lem1.1.9.1}
设 $c:F\rightarrow G$ 是缩并态射，$(K_{\ell})_{\ell \in I}$ 是诸 $\Aa_{\ell}$
中的复形，且 $K_{i}=K_{j}$。设 $<$ 是 $I$ 上的一个全序，其中 $i$ 紧邻在 $j$ 之前。
沿用 \hyperref[XVII.subsubsec1.1.6]{1.1.6} 的记号，对 $k\in \ZZ$，设
\[
c^{k}:K^{s}(F)(K_{\ell})^{k}\longrightarrow K^{s}(G)(K_{\ell})^{k}
\]
为\pageoriginale 如下态射：它在 $F(\ldots K_{i}^{k_i}K_{j}^{k_j}\ldots)$ 上的限制，
当 $k_i\neq k_j$ 时为 $0$，当 $k_{i}=k_{j}$ 时为 $c_{K_i^{k_i}}$。
则 $(c^{k})_{k\in \ZZ}$ 是复形态射。
\end{enonce*}

此引理的验证，以及\emph{缩并态射} $(c^{k})$ 不依赖于 $I$ 上所选全序这一事实的验证，
都留给读者。

缩并态射的标准例子是复合律
\[
\Hom(Y,Z)\otimes \Hom(X,Y)\longrightarrow \Hom(X,Z)\sisi{}{.}
\]

\begin{enonce*}[remark]{例 1.1.10}\etiquette[1.1.10]{XVII.exem1.1.10}
\begin{enumerate}
\renewcommand{\theenumi}{\roman{enumi}}
\renewcommand{\labelenumi}{(\theenumi)}
\item 将 \hyperref[XVII.subsubsec1.1.6]{1.1.6} 的定义应用于\emph{张量积}函子，
  便得到两个复形的张量积的通常概念，其微分为
\[
d(x\otimes y)=dx\otimes y+(-1)^{\deg x}x\otimes dy{\quoi}.
\]

\item 将 \hyperref[XVII.subsubsec1.1.6]{1.1.6} 的定义应用于函子 $\Hom(X,Y)$，
  \emph{将 $Y$ 看作第一个变量、$X$ 看作第二个变量}；我们将始终遵循这个约定。
  按此约定，对于两个复形 $X$、$Y$，\hyperref[XVII.subsubsec1.1.6]{1.1.6}
  定义的复形 $\Hom(X,Y)$ 的第 $k$\sisi{个}{个}分量等同于
  $\prod_n\Hom(X^{n},Y^{n+k})$，并得到通常的微分：对
  $f\in \Hom(X^{n},Y^{n+k})$，有
\[
df=d\circ f-(-1)^{k}f\circ d{\quoi}.
\]

\item 上述约定的优点是，下列箭头系均为复形态射，无须插入任何符号
  \pageoriginale（\hyperref[XVII.subsubsec1.1.9]{1.1.9}）：
\begin{align*}
& \Hom^{\bigdot}(L,M)\otimes \Hom^{\bigdot}(K,L)\longrightarrow
  \Hom^{\bigdot}(K,M):f\otimes g\mapsto f\circ g\\
& \Hom^{\bigdot}(L,M)\otimes L\longrightarrow M: f\otimes x\mapsto f(x){\quoi}.
\end{align*}
\end{enumerate}
\end{enonce*}

\subsubsection*{${}^{\ast}$1.1.10.1}\etiquette[*1.1.10.1]{XVII.par1.1.10.1}
对于 $F=\otimes$，约定 \eqhyperref[XVII.eq1.1.4.5]{1.1.4.5}--
（\hyperref[XVII.par1.1.6.2]{1.1.6.2}）的理由如下：
$((K\otimes L)[1])^{n}$ 的元素写成（\hyperref[XVII.subsubsec*1.1.2]{*1.1.2}）
（回顾 $(K\otimes L)[1]=C(K\otimes L\rightarrow 0)$
\eqhyperref[XVII.eq1.1.1.1]{1.1.1.1}）$\sum 1\otimes(x^{p}\otimes y^{q})$；
$(K[1]\otimes L)^{n}$ 的元素写成 $\sum (1\otimes x^{p})\otimes y^{q}$；
$(K\otimes L[1])^{n}$ 的元素写成 $\sum x^{p}\otimes (1\otimes y^{q})$。
符号在交换 $x^{p}$ 与 $1$ 时出现。
对于 $F=$ 对偶，约定 \eqhyperref[XVII.eq1.1.5.2]{1.1.5.2} 的理由如下：
$K^{\ast}[1]$ 的一个元素写成 $1\otimes \omega$；它在 $K[-1]$ 的元素
$1'\otimes x$（$1'$ 是维数为 $-1$ 的胞腔）上的值为
$<1\otimes \omega, 1'\otimes x>=(-1)^{\deg\omega}<1, 1'><\omega,x>=(-1)^{\deg\omega}<\omega,x>$；
这正是所采用的约定。${}_{\ast}$

\begin{enonce*}[remark]{注 1.1.11}\etiquette[1.1.11]{XVII.rem1.1.11}
“回顾”，\emph{预加性范畴}（preadditive category）是其箭头集 $\Hom(X,Y)$
赋有阿贝尔群结构，且态射复合为双加性的范畴。
预加性范畴之间的\emph{加性函子} $F$ 是满足下式的函子：
\[
F(f+g)=F(f)+F(g)
\]
其中 $f$、$g\in \Hom(X,Y)$。
从预加性范畴 $\Aa$ 到另一预加性范畴 $\Bb$ 的加性函子所成的范畴，记为
$\Hhom_{\text{\rm add}}(\Aa,\Bb)$。
\end{enonce*}

设\pageoriginale $(D)$ 是唯一的预加性范畴，其对象集以 $\ZZ$ 为指标，
记为 $(X\sisi{_n}{^n})_{n\in \ZZ}$，且满足
$$\etiquette[1.1.11.1]{XVII.eq1.1.11.1}
\begin{cases}
\Hom(X^{n},X^{n})=\ZZ,\;\mbox{\rm 由恒等态射生成}\\
\Hom(X^{n},X^{n+1})=\ZZ,\;\mbox{\rm 由记为 $d_n$ 的箭头生成}\\
\Hom(X^{i},X^{j})=0\;\mbox{若 } j-i\neq 0\;\mbox{\rm 或 } 1{\quoi}.
\end{cases}\leqno(1.1.11.1)
$$

设 $D$ 是 $(D)$ 中由 $D^{n}=X^{n}$ 及 $d^{n}=d_n$ 定义的复形。

对每个预加性范畴 $\Aa$，从 $\Hhom_{\text{\rm add}}((D),\Aa)$ 到 $C(\Aa)$
的函子 $G\mapsto G(D)$ 是同构。

\hyperref[XVII.subsubsec1.1.6]{1.1.6} 的构造可以推广到从预加性范畴 $\Aa_i$
到加性范畴 $\Aa$ 的多重加性函子的情形。

设
\[
F^{\bigdot}=(F^{n},d^{n})_{n\in \ZZ}
\]
是从加性范畴 $\Aa_i$ 到加性范畴 $\Aa$ 的多重加性函子复形。
可以将 $F^{\bigdot}$ 看作从 $(D)$ 到多函子范畴的函子，或者等价地，
看作多重加性函子
\[
F_1:\prod\Aa_i\times (D)\longrightarrow \Aa{\quoi}.
\]
\emph{$F\sisi{}{^{\bigdot}}$ 对复形的延拓}由下式定义：
$$\etiquette[1.1.11.2]{XVII.eq1.1.11.2}
F^{\bigdot}(K_{1}\ldots K_{n})=F_{1}(K_{1}\ldots K_{n},D){\quoi}.\leqno(1.1.11.2)
$$

\setcounter{subsubsection}{11}
\subsubsection{}\etiquette[1.1.12]{XVII.subsubsec1.1.12}
设 $\Aa$ 和 $\Bb$ 是两个加性范畴，并设
\[
F^{\bigdot}=(F^{n},d^{n})_{n\in \ZZ}
\]
是从 $\Aa$ 到 $\Bb$ 的加性函子复形。假设函子\pageoriginale $F^{i}$ 有右伴随
$F_i$；设 $\leftexp{t}{d}^{i}:F_{i+1}\rightarrow F_i$ 是 $d^{i}$ 的转置，
$F_{\bigdot}$ 是从 $\Bb$ 到 $\Aa$ 的加性函子复形，其分量为 $F_{i}$，微分算子为
\[
(d_\bigdot)^{i}=(-1)^{i+1}\,\leftexp{t}{d}^{-i-1}:F_{-i}\longrightarrow
F_{-i-1}{\quoi}.
\]
设 ${}^{t}$ 是从 $(D)$ 到 $(D)$ 的\emph{反变}函子，给定为
\[
\leftexp{t}{X}^{n}=X^{-n}\quav \leftexp{t}{d}^{n}=d^{-n-1}{\quoi}.
\]
下图在典范同构意义下交换：
$$\etiquette[1.1.12.1]{XVII.eq1.1.12.1}
\vcenter{
\xymatrix@C=6em{
\mathcal{A} \times (D) \times \mathcal{B}  \ar[d]^{(\Id,{}^t,\Id)}\ar[r]^(.55){(F^{n}(A), \Id)} &
\mathcal{B}\times \mathcal{B}  \ar[r]^{\Hom_{\mathcal{B}}}&
(\Ab)\ar@{=}[d] \\
\mathcal{A} \times (D) \times \mathcal{B}
\ar[r]^(.6){(\Id,(F_{\bigdot})^n(B))} & \mathcal{A} \times \mathcal{A}
\ar[r]^{\Hom_{\mathcal{A}}} & (\Ab) \rlap{.}
}}\leqno{(1.1.12.1)}
$$
它表达伴随公式
\begin{equation*}
\Hom(F^{n}(A),B)\overset{\sim}{\longleftrightarrow}\Hom(A,F_{n}(B)){\quoi}.
\end{equation*}
因此，\hyperref[XVII.def1.1.3]{1.1.3} 的一般约定给出同构
\begin{equation*}
\Hom^{\bigdot}(F^{\bigdot}(K),L)\overset{\sim}{\longleftrightarrow}\Hom\sisi{}{^\bigdot}(K,F_{\bigdot}(L)){\quoi}.
\end{equation*}

\begin{enonce*}{定义 1.1.13}\etiquette[1.1.13]{XVII.def1.1.13}
若 $K$ 是阿贝尔范畴 $\Aa$ 中的复形，则 $K$ 在维数 $n$ 处的
\emph{右截断（\resp 左截断）}，记为 $\tau_{\leq n}(K)$
\emph{（\resp $\tau_{\geq n}(K)$）}，指下列复形：
\begin{align*}
\tau_{\leq n}(K) &: \ldots K^{n-2}\longrightarrow
K^{n-1}\longrightarrow \Ker(d^n)\longrightarrow 0\\
\tau_{\geq n}(K) &: 0\longrightarrow \coker (d^{n-1})\longrightarrow
K^{n+1}\longrightarrow K^{n+2}\ldots.
\end{align*}
\end{enonce*}

\begin{enonce*}{命题 1.1.14}\etiquette[1.1.14]{XVII.prop1.1.14}
\begin{enumerate}
\renewcommand{\theenumi}{\roman{enumi}}
\renewcommand{\labelenumi}{(\theenumi)}
\item 当 $i>n$ 时，$H^{i}(\tau_{\leq n}(K))=0$\pageoriginale；
当 $i<n$ 时，$H^{i}(\tau_{\geq n}(K))=0$。

\item 典范态射 $\tau_{\leq n}(K)\rightarrow K$（\resp $K\rightarrow \tau_{\geq n}K$）
在 $i\leq n$（\resp $i\geq n$）时，诱导 $H^{i}$ 上的同构。
\end{enumerate}
\end{enonce*}

函子 $\tau_{\leq n}$（\resp $\tau_{\geq n}$）将同伦的态射映为同伦的态射。
因此由 \sisi{\hyperref[XVII.subsubsec1.1.4]{1.1.4}}{\hyperref[XVII.prop1.1.14]{1.1.14}}，
\hyperref[XVII.def1.1.13]{1.1.13} 的函子在导出范畴中仍有意义。

\setcounter{subsubsection}{14}
\subsubsection{}\etiquette[1.1.15]{XVII.subsubsec1.1.15}
若 $K$ 是双复形，第一（\resp 第二）微分为 $d'$（\resp $d''$），
则以 $\tau''_{\leq n}(K)$ 表示 $K$ 的如下双子复形：
\begin{equation*}
\tau''_{\leq n}(K)^{pq}=
\begin{cases}
K^{pq}&\mbox{\rm 若 } q<n\\
\Ker(d'') &\mbox{\rm 若 } q=n\\
0 &\mbox{\rm 若 } q>n
\end{cases}
\end{equation*}
并类似定义 $\tau''_{\geq n}(K)$、$\tau'_{\leq n}(K)$ 以及 $\tau'_{\geq n}(K)$。

\subsubsection{}\etiquette[1.1.16]{XVII.subsubsec1.1.16}
我们也会考虑复形 $K$ 的“愚笨截断”，定义为
\begin{xalignat*}{2}
\sigma_{\leq n}(K)^{i}=
\begin{cases}
K^{i} &\mbox{\rm 若 } i\leq n\\
0 &\mbox{\rm 若 } i>n
\end{cases}
&&
\sigma_{\geq n}(K)^{i}=
\begin{cases}
K^{i} &\mbox{\rm 若 } i\geq n\\
0 &\mbox{\rm 若 } i<n{\quoi}.
\end{cases}
\end{xalignat*}

\subsection{导出函子}\etiquette[1.2]{XVII.subsec1.2}\pageoriginale

设 $\Aa$ 是三角范畴，$S$ 是饱和乘法系
（\cite{sga4.exp.XVII.bib.CD} \sisi{}{I }2.1.2）。若 $A\in \Ob \Aa$，
则以 $A^{+}$ 和 $A^{-}$ 分别表示 \iind-对象和 \ppro-对象
（（\hyperref[XVII.subsec0.5]{0.5}）和（\hyperref[XVII.subsec0.6]{0.6}））
\begin{equation*}
A^{+}=\quotinjlim_{A\xrightarrow{s}A'} A' \qquad A^{-}=\quotprojlim_{A'\xrightarrow{s}A}A'{\quoi},
\end{equation*}
其中 $s$ 分别遍历 $S$ 中以 $A$ 为源或为目标的箭头组成的滤过范畴。

函子 $A\mapsto A^{+}$ 和 $A\mapsto A^{-}$ 使 $S$ 的元素可逆，
并分别定义从 $\Aa(S^{-1})$ 到 $\Aa$ 的 Ind-对象范畴和 \ppro-对象范畴的全忠实函子。

\begin{enonce*}{定义 1.2.1}\etiquette[1.2.1]{XVII.def1.2.1}
设 $F$ 是从三角范畴 $(\Aa_i)_{0<i\leq n}$ 到三角范畴 $\Aa$ 的正合协变多函子，
$S_i$ 和 $S$ 分别是诸范畴 $\Aa_i$ 和 $\Aa$ 中的饱和乘法系。
\begin{enumerate}
\renewcommand{\theenumi}{\roman{enumi}}
\renewcommand{\labelenumi}{(\theenumi)}
\item $F$ 的右导出函子 $RF$（相对于诸 $S_i$ 和 $S$），是从诸范畴
$\Aa_i(S^{-1}_i)$ 到 $\Aa(S^{-1})$ 的 \iInd-对象范畴的函子，使下图交换。
\[
\xymatrix@C=5em{%csy
  (\Aa_i)  \ar[d]\ar[r]^{(A_i)\sisi{\simarrow}{\mapsto}  F(A^{+}_i)} & \Ind
  \Aa\ar[d] \\
\Aa_i(S^{-1}_i)\ar[r]^{RF} & \Ind\Aa(S^{-1}).
}
\]

\item 若 $\Aa(S^{-1})$ 的 \iInd-对象 $RF(A_i)$ “本质常值”，
\ie 来自 $\Aa(S^{-1})$ 的一个对象，则\pageoriginale 称 $RF$ 在 $(A_i)$ 处有定义；
该对象称为 $RF$ 在 $(A_i)$ 处\emph{的值}。

\item 若 $RF$ 处处有定义，则称 $F$ \emph{可右导出}；
此时仍以 $RF$ 表示所得到的从诸 $\Aa_i(S^{-1}_i)$ 到 $\Aa(S^{-1})$ 的函子。

\item 诸 $\Aa_i$ 的一族对象 $(A_i)$，若从 $F(A_i)$ 到 $RF(A_i)$ 的典范态射是同构，
则称它\emph{适合于} $RF$。
\end{enumerate}
\end{enonce*}

通过对偶性定义 $F$ 的左导出函子 $LF$（取值于 $\Pro \Aa(S^{-1})$），
并将上述定义推广到 $F$ 对某些变量协变、对另一些变量反变的情形，这些工作留给读者。

当 $RF$ 处处有定义时，这里的定义与 Verdier 的定义一致
（\cite{sga4.exp.XVII.bib.CD} 第~39 页），容易验证。

\begin{enonce*}{命题 1.2.2}\etiquette[1.2.2]{XVII.prop1.2.2}
设 $(A'_1, A_1,A''_1)$ 是 $\Aa_1$ 中的特殊三角，并设\break
$A_i\in \Ob \Aa_i$ $(i\neq 1)$。
\begin{enumerate}
\renewcommand{\theenumi}{\roman{enumi}}
\renewcommand{\labelenumi}{(\theenumi)}
\item 若 $(A'_{1},(A_{i}))$ 和 $(A''_{1},(A_{i}))$ 适合于 $RF$，
  则族 $(A_{1},(A_{i}))$ 也适合于 $RF$。

\item 若 $RF$ 在 $(A'_1, (A_i))$ 和 $(A''_1, (A_i))$ 处有定义，
则它在 $(A_1, (A_i))$ 处也有定义，且三角
$(RF(A'_1,(A_i)), RF(A_1,(A_i)),RF(A''_1,(A_i)))$
是 $\Aa(S^{-1})$ 中的特殊三角。
\end{enumerate}
\end{enonce*}

先证明\pageoriginale (ii) $\Rightarrow$ (i)。因为 $F$ 正合，箭头
$$(F(A'_1,(A_i)), F(A_1,(A_i)), F(A''_1,(A_i)))\rightarrow
(RF(A'_1,(A_i)), RF(A_1,(A_i)), RF(A''_1,(A_i)))$$
是特殊三角的态射。由假设，此箭头在两个顶点处是同构，所以在第三个顶点处也是同构，
这就是所需结论。

\begin{enonce*}{引理 1.2.2.1}\etiquette[1.2.2.1]{XVII.lem1.2.2.1}
若 $\Delta=(X,Y,Z)$ 是三角范畴 $\Aa$ 中的特殊三角，$S$ 是 $\Aa$ 中的饱和乘法系，
则 $\Ind\Aa$ 中的三角 $\Delta^{+}=(X^{+},Y^{+},Z^{+})$ 是
$\Aa$ 中特殊三角的极限。
\end{enonce*}

以 $\mathcal{L}$ 表示源为 $\Delta$ 且其各箭头均属于 $S$ 的特殊三角态射所成的范畴。
\begin{enumerate}
\renewcommand{\theenumi}{\alph{enumi}}
\renewcommand{\labelenumi}{(\theenumi)}
\item 若 $\Delta=(X,Y,Z)$ 与 $\Delta'=(X',Y',Z')$ 是 $\Aa$ 中的两个特殊三角，
且 $f:\Delta\rightarrow \Delta'$ 是三角的 $\Aa(S^{-1})$-态射，
则在 $\Aa$ 中存在三角态射 $s:\Delta'\rightarrow\Delta''$ 和
$f_1:\Delta\rightarrow \Delta''$，使 $s$ 的各分量属于 $S$，且图
\[
\vcenter{
\xymatrix{
\Delta                \ar[rd]^{f_1}  \ar[d]_{f}  &\\
\Delta'          \ar[r]^{s}    & \Delta''
}}\tag{1}
\]
在 $\Aa(S^{-1})$ 中交换。


\sisi{\noindent}{}%
存在交换图
\[
\vcenter{
\xymatrix{
X \ar[rd]^{f_1} \ar[d]_{f}&  &\\
X' \ar[r]^{s\sisi{}{_X}} & X''\sisi{}{_0}&
}}(s\sisi{}{_X} \in S, f_1 \in \Fl \Aa)\tag{2}
\]
以及\pageoriginale $S$ 中的态射 $\Delta'\xrightarrow{s\sisi{_1}{}}\Delta''\sisi{_1}{_0}$，
它以 $s\sisi{}{_X}$ 为一个分量。在其他顶点处类似论证，便得到类型 (1) 的图。
不过，$f_1$ 此时还不是\emph{在} $\Aa$ \emph{中}的三角态射，因为例如图
\[
\vcenter{
\xymatrix{
Z  \ar[r]^{f\sisi{}{_1}} \ar[d]_{w}& Z''_{0} \ar[d]_{w\sisi{'}{''}}\\
X         \ar[r]^{f\sisi{}{_1}} & X''_{0}
}}\tag{3}
\]
可能不在 $\Aa$ \emph{中}交换，而仅在 $\Aa(S^{-1})$ 中交换。
存在 $S$ 中的 $t\sisi{_0}{_X}:X''_0\rightarrow X''_1$，使
$t\sisi{}{_X}w''f\sisi{}{_1}=t\sisi{}{_X}f\sisi{}{_1}w$，
以及以 $t\sisi{_0}{_X}$ 为一个分量的 $t:\Delta''_0\rightarrow \Delta''_1$。
用 $\Delta''_1$ 代替 $\Delta''_0$，用 $ts$ 代替 $s$，用 $tf_1$ 代替 $f_1$；
在所得的类型 (1) 的图中，(3) 现在交换。在其他顶点处类似处理，便得到 (1)。
以下事实的类似验证留给读者：

\item 若 $f$、$g:\Delta\rightrightarrows \Delta'$ 是 $\Aa$ 中的两个特殊三角态射，
它们在 $\Aa(S^{-1})$ 中相等，则存在各箭头均在 $S$ 中的特殊三角态射
$s:\Delta'\rightarrow \Delta''$，使 $\sisi{S}{s}f=\sisi{S}{s}g$。

\item 范畴 $\Ll$ 是向右滤过的。

\item 注意到
\[
\Delta^{+}=\quotinjlim_{s\in \Ll}{\mbox{\rm （$s$ 的目标）。}}
\]
便完成引理的证明。
\end{enumerate}

现在证明\pageoriginale \hyperref[XVII.prop1.2.2]{1.2.2} (ii)。
设 $X'$ 和 $X''$ 是 $\Aa(S^{-1})$ 的对象，分别代表 $RF(A'_1,(A_i))$
与 $RF(A''_1,(A_i))$，并将从 $X''$ 到 $X'$ 的一次态射补全为特殊三角 $(X',X,X'')$。
将 $A_i$、$A'_1$、$A_1$ 和 $A''_1$ 分别换成
\sisi{$S$}{$S_i$ 或 $S_1$} 中以 $A_i$、$A'_1$、$A_1$ 或 $A''_1$ 为源的箭头的目标，
便归结到存在以下交换图的情形：
\[
\xymatrix{
F(A''_1, (A_i)) \ar[d] & X''  \ar[l]_-{a''} \ar[d]\\
F(A'_1, (A_i))    & X'  \ar[l]_-{a'}\rlap{,}
}
\]
其中 $a''$ 和 $a'$ 分别是从 $F(A''_1,(A_i))$ 到 $X''$ 以及
从 $F(A'_1,(A_i))$ 到 $X'$ 的箭头的截面。
由 \cite{sga4.exp.XVII.bib.CD} 的公理 $TR$ 3，此图可补全为三角态射
\begin{align*}
(X'XX'')& \longrightarrow
(F(A'_1,(A_i)),F(A_1,(A_i)),F(A''_1,(A_i)))\\
& \longrightarrow
(RF(A'_1,(A_i)),RF(A_1,(A_i)),RF(A''_1,(A_i))){\quoi}.
\end{align*}

复合箭头是三角态射，其源为特殊三角，目标为特殊三角的极限
（在 $\Ind\Aa(S^{-1})$ 中）（\hyperref[XVII.lem1.2.2.1]{1.2.2.1}）。
对 $\Aa(S^{-1})$ 中的任意 $Y$，下图交换且各行正合：
\[
\text{\tiny
$\xymatrix@C=0.5em{
\Hom\sisi{'}{(}Y, X''[\sisi{}{-}1]) \ar[r]\ar[d]& \Hom(Y,X') \ar[r] \ar[d]& \Hom(Y,X)
\ar[r]\ar[d]& \Hom(Y,X'') \ar[r] \ar[d]&
\Hom(Y,X'[1]) \ar[d]\\
\Hom(Y,RF(A''\sisi{}{_1},A_i)[\sisi{}{-}1])\ar[r] & \Hom(Y,RF(A'_1,A_i)\sisi{}{)}\ar[r] &  \Hom(Y,RF(A_1,A_i))\ar[r] &
\Hom(Y,RF(A''_1,A_i)\sisi{}{)}\ar[r] & \Hom(Y,RF(A'_1,A_i)[1])
}$}
\]
由\pageoriginale 假设，除了中间的箭头，其余箭头都是同构。
\hyperref[XVII.sec5]{5}-引理表明，对每个 $Y$，中间箭头也为同构，
从而特殊三角 $(X',X,X'')$ 同构于三角（$RF(A'_1,(A_i))$、
$RF(A_1,(A_i))$、$RF(A''_1,(A_i))$）。证明至此完成。

若 $\Aa$ 和 $\Bb$ 是两个阿贝尔范畴，$F$ 是从 $\Aa$ 到 $\Bb$ 的加性协变函子，
则 $F$ 定义从 $K(\Aa)$ 到 $K(\Bb)$、从 $K^{+}(\Aa)$ 到 $K^{+}(\Bb)$
以及从 $K^{-}(\Aa)$ 到 $K^{-}(\Bb)$ 的正合函子
（\hyperref[XVII.prop1.1.7]{1.1.7}）。

\begin{enonce*}{定义 1.2.3}\etiquette[1.2.3]{XVII.def1.2.3}
\begin{enumerate}
\renewcommand{\theenumi}{\roman{enumi}}
\renewcommand{\labelenumi}{(\theenumi)}
\item \emph{$F$ 的右导出函子}，是由 $F$ 对复形的延拓
$K^{+}(\Aa)\rightarrow K^{+}(\Bb)$ 和 $K(\Aa)\rightarrow K(\Bb)$
右导出得到的、从 $D^{+}(\Aa)$ 到 $\Ind D^{+}\sisi{\Bb}{(\Bb)}$
以及从 $D(\Aa)$ 到 $\Ind D(\Bb)$ 的函子。

\item $\Aa$ 的一个对象 $A$，若在次数 $0$ 集中于 $A$ 的复形适合于 $RF$，
则称为关于 $F$ \emph{右无环}的，或滥用术语，称为\emph{关于} $RF$ \emph{无环}的。
\end{enumerate}
\end{enonce*}

我们主要关心 $RF$ 处处有定义的情形；此时，按照所考虑的情形，可以将它看作
从 $D^{+}(\Aa)$ 到 $D^{+}(\Bb)$，或从 $D(\Aa)$ 到 $D(\Bb)$ 的函子。

\begin{enonce*}{命题 1.2.4}\etiquette[1.2.4]{XVII.prop1.2.4}
下图交换。
\[
\xymatrix{
D^{+}(\mathcal{A})  \ar[r] \ar[d]^{RF}& D(\mathcal{A})\ar[d]^{RF}        \\
 \Ind \;D^{+}(\mathcal{B}) \ar[r] & \Ind \; D(\mathcal{B}).
}
\]
\end{enonce*}

若\pageoriginale $K$ 是在次数 $<n$ 处为零的复形，且 $s:K\rightarrow L$
是拟同构，则复合态射 $K\xrightarrow{s}L\rightarrow \sigma_{\geq n}(L)$
（\hyperref[XVII.subsubsec1.1.6]{1.1.6}）仍为拟同构。因此，从 $K$ 到下有界复形的
拟同构组成以 $K$ 为源的所有拟同构范畴的一个共尾（满）子范畴。
命题形式地随之得到。

特别地，若函子 $RF$\sisi{，}{}在 $D(\Aa)$ 上处处有定义，
则它在 $D^{+}(\sisi{\Bb}{\Aa})$ 上处处有定义，且导出函子
$RF:D^{+}(\Aa)\rightarrow D^{+}(\Bb)$ 是导出函子
$RF:D(\Aa)\rightarrow D(\Bb)$ 的限制。

\setcounter{subsubsection}{4}
\subsubsection{}\etiquette[1.2.5]{XVII.subsubsec1.2.5}
以下工作留给读者：
\begin{enumerate}
\renewcommand{\theenumi}{\roman{enumi}}
\renewcommand{\labelenumi}{(\theenumi)}
\item 将定义 \hyperref[XVII.def1.2.3]{1.2.3} 推广到多函子复形的情形，
它们对某些变量协变，对另一些变量反变；

\item 将定义 \hyperref[XVII.subsubsec1.2.5]{1.2.5} (i) 对偶化，定义左导出函子；

\item 将命题 \hyperref[XVII.prop1.2.4]{1.2.4} 推广到有界多函子复形的情形；

\item 考虑左导出函子以及 $D^{-}(\Bb)$ 到 $D(\Bb)$ 的包含，
将命题 \hyperref[XVII.subsubsec1.2.5]{1.2.5} (iii) 对偶化。
\end{enumerate}

\begin{enonce*}{命题 1.2.6}\etiquette[1.2.6]{XVII.prop1.2.6}
设 $\Aa$ 和 $\Bb$ 是两个阿贝尔范畴，$𝒫$ 是 $\Ob(\Aa)$ 的一个子集，
$F$ 是从 $\Aa$ 到 $\Bb$ 的加性函子。假设：
\begin{enumerate}
\renewcommand{\theenumi}{\alph{enumi}}
\renewcommand{\labelenumi}{(\theenumi)}
\item $\Aa$ 的每个对象都是 $𝒫$ 的一个元素的商；

\item 对每个短正合列
\[
0\longrightarrow P\longrightarrow Q\longrightarrow R\longrightarrow
0{\quoi},
\]
其中 $Q$ 和 $R$ 属于 $𝒫$，则 $P$ 也属于 $𝒫$，且 $0$-列
$0\rightarrow FP\rightarrow FQ\rightarrow FR\rightarrow 0$\pageoriginale 正合。
则 $𝒫$ 的每个对象都关于 $LF$ 无环。
\end{enumerate}
\end{enonce*}

若 $A\in \Ob \Aa$，则 $A$ 的由 $𝒫$ 中对象组成的左解消，在 $A\sisi{_0}{}$
的左解消的滤过范畴中组成一个共尾系。因此只须证明：若复形
\[
\ldots\longrightarrow A_2\longrightarrow A_1\longrightarrow A_0
\longrightarrow A\longrightarrow 0
\]
无环，且其对象属于 $𝒫$，则它的像也无环。将它拆成短正合列
$$\etiquette[1.2.6.i]{XVII.eq1.2.6.i}
0\longrightarrow \Iim(d_{i-1})\longrightarrow A_i\longrightarrow
\Iim(d_i)\longrightarrow 0\qquad\sisi{(1.2.6. i)\quad}{}(i\geq 0)\sisi{}{\leqno(1.2.6.i)}%\tag{1.2.6.$i$}
$$
对 $i$ 归纳可知，$\Iim(d_i)\in \PP$，且
\sisi{(\hyperref[XVII.prop1.2.6]{1.2.6} $i$)}{\eqhyperref[XVII.eq1.2.6.i]{1.2.6.i}}
在 $F$ 下的像为短正合列。因此，该复形的像无环。

\begin{enonce*}{命题 1.2.7}\etiquette[1.2.7]{XVII.prop1.2.7}
设 $F$ 是从一个阿贝尔范畴到阿贝尔范畴 $\Bb$ 的加性函子。
若 $\Aa$ 的每个对象都是关于 $LF$（\resp $RF$）无环的一个对象的商
（\resp 子对象）（\hyperref[XVII.def1.2.3]{1.2.3}），则导出函子
$LF:D^{-}(\Aa)\rightarrow D^{-}(\Bb)$（\resp $RF:D^{+}(\Aa)\rightarrow D^{+}(\Bb)$）
处处有定义，且关于 $LF$（\resp $RF$）无环的对象组成的每个上有界
（\resp 下有界）复形都是适合的。
\end{enonce*}

只须考虑左导出函子 $LF$。若 $K$ 是上有界复形，则 $L$ 上有界且由无环对象组成的
拟同构 $s:L\rightarrow K$，在目标为 $K$ 的拟同构范畴中组成共尾系。
只须证明，当 $K$ 由无环对象组成时，$F(s):F(L)\rightarrow F(K)$ 是拟同构；
\pageoriginale 此时，在一般情形中将有 $F(L)=LF(K)$。
设 $M$ 是 $s$ 的锥。此复形由无环对象\sisi{}{们}组成，且 $H^{\ast}(M)=0$；
须证明 $H^{\ast}(F(M))=0$。

设 $i$ 是整数。适合的复形组成一个三角子范畴，所以无环对象组成的每个有界复形，
特别是 $\tau_{\geq i}(M)$，都是适合的。
因为对 $j>i$ 有 $H^{j}(\tau_{\geq i}(\sisi{L}{M}))=0$，所以对 $j>i$，有
$H^{j}(F(M))=H^{j}(F(\tau_{\geq i}(M)))=H^{j}LF(\tau_{\geq i}(M))=0$。

\begin{enonce*}{命题 1.2.8}\etiquette[1.2.8]{XVII.prop1.2.8}
设 $F^{\bigdot}$ 是从 $\Aa$ 到 $\Bb$ 的下有界函子复形。
若对每个 $K\in K^{+}(\Aa)$ 都存在拟同构 $L\xrightarrow{s}K$，其中
$L\in K^{+}(\Aa)$ 适合于所有 $RF^{i}$，则导出函子
$RF^{\bigdot}:D^{+}(\Aa)\rightarrow D^{+}(\Bb)$ 存在，
且适合于所有 $RF^{i}$ 的每个\emph{下有界}复形都适合于 $RF^{\bigdot}$。
\end{enonce*}

\sisi{\noindent}{}%
证明留给读者，对偶陈述的证明也一样；取 $F$ 为反变函子，
或取有界函子复形并在 $D(\Aa)$ 与 $D(\Bb)$ 中考虑所得到的变体，其证明也留给读者。

\begin{enonce*}{定义 1.2.9}\etiquette[1.2.9]{XVII.def1.2.9}
设 $F$ 是从 $D^{b}(\Aa)$ 到 $D(\Bb)$ 的函子，$d$ 是 $\ZZ$ 中的一个区间。
若对每个 $A\in \Ob \Aa$ 以及 $i\not\in d$，都有 $H^{i}F(A)=0$，
则称 $F$ 的\emph{上同调幅度 $\subset d$}。
\end{enonce*}

很多时候，$d$ 是形如 $[0,x]$ 或 $[-x,0]$ 的区间；
此时称 $F$ 的上同调维数（\resp 同调维数）$\leq x$。
若存在 $\ZZ$ 中的一个有限区间 $d$，使 $F$ 的幅度 $\subset d$，则称 $F$ 的维数有限。

\begin{enonce*}{命题 1.2.10}\etiquette[1.2.10]{XVII.prop1.2.10}
假设\pageoriginale \hyperref[XVII.prop1.2.7]{1.2.7} 的假设成立，且函子
$RF:D^{+}(\Aa)\rightarrow D^{+}(\Bb)$（\resp $LF:D^{-}(\Aa)\rightarrow D^{-}(\Bb)$）
的维数有限。则关于 $RF$（\resp $LF$）无环的对象组成的每个复形都是适合的，
且函子 $RF:D(\Aa)\rightarrow D(\Bb)$（\resp $LF:D(\Aa)\rightarrow D(\Bb)$）处处有定义。
\end{enonce*}

证明参见 Verdier \cite{sga4.exp.XVII.bib.CD}。
对于有界函子复形有类似陈述。上标 $+$（\resp $-$）也可换成 $b$。

% END native input: chapters/exp-17-part1-zh.tex
% BEGIN native input: chapters/exp-17-part2-zh.tex
\section{纤维为导出范畴的范畴}\etiquette[2]{XVII.sec2}

\subsection{引言}\etiquette[2.1]{XVII.sec2.1}

撰稿人强烈建议读者不要阅读本\S。这里给出平凡对偶定理的形式体系，并解决 Artin 在 \hyperrefext[XII][XII.sec4]{4} 中提出的疑惑。

关于一个范畴之上的范畴、纤维化范畴（fibered category）、余纤维化范畴（cofibered category）以及劈分（cleavage）的概念，参见 \SGA 1 VI。该讲特别建立了范畴 $\sisi{C}{\Cc}$ 上的正规化分裂纤维化范畴与从 $\sisi{C}{\Cc}^{\circ}$ 到 $(\Cat)$ 的伪函子（pseudofunctor）之间的等价。

\begin{enonce*}{定义 2.1.1}\etiquette[2.1.1]{XVII.def2.1.1}
若 $X$ 是环化景（ringed site），则 $D(X)$ 表示 $X$ 上左模层的阿贝尔范畴的导出范畴（derived category）。
\end{enonce*}

需要避免混淆时，改写为 $D(X,\Aa)$，其中 $\Aa$ 表示环层。

若\pageoriginale $f:X\rightarrow Y$ 是环化景的态射，则 $f$ 定义导出函子
\begin{itemize}
\item[(1)] $Rf_{\ast}:D^{+}(X)\rightarrow D^{+}(Y)$，以及

\item[(2)] $Lf^{\ast}:D^{-}(Y)\rightarrow D^{-}(X)$，
\end{itemize}
转到反范畴\sisi{}{}，则得到函子
\begin{itemize}
\item[$\rm(1^{\circ})$] $Rf_{\ast}:D^{+}(X)^{\circ}\rightarrow
D^{+}(Y)^{\circ}$

\item[$\rm(2^{\circ})$] $Lf^{\ast}:D^{-}(Y)^{\circ}\rightarrow
D^{-}(X)^{\circ}$，
\end{itemize}
它们关于 $f$ 是“函子性的”。更确切地说，$(1^{\circ})$ 和 $(2^{\circ})$ 定义了从景的范畴到三角范畴（triangulated category）的范畴的两个伪函子，一个协变，另一个反变。由此得到景的范畴之上的两个范畴：一个是余纤维化的，纤维为 $D^{+}(X)^{\circ}$，正像函子为 $Rf_{\ast}$；另一个是纤维化的，纤维为 $D^{-}(X)^{\circ}$，逆像函子为 $\sisi{\LL}{L} f^{\ast}$。若函子 $f^{\ast}$ 具有有限同调维数（\ie 有限的“\tTor-维数”），则 (2) 中的上标可换成 $+$。这样，在由环化景及有限 \tTor-维数的环化景态射组成的范畴 $\Ss_{\tdf}$ 之上，便得到两个范畴\sisi{}{}；它们的纤维都是范畴 $D^{+}(X)^{\circ}$。\textsc{Verdier} 的“平凡对偶定理”（\cite{sga4.exp.XVII.bib.CD} 第~48 页）断言这两个范畴典范同构；所以它们既是纤维化的也是余纤维化的；其“正像”函子为 $Rf_{\ast}$\sisi{}{ }（原文如此），而“逆像”函子为 $Lf^{\ast}$\sisi{}{ }（原文如此）。

\setcounter{subsection}{1}
\setcounter{subsubsection}{1}
\subsubsection{}\etiquette[2.1.2]{XVII.subsubsec2.1.2}

纤维范畴的语言很容易解决 \textsc{Artin} 在 \hyperrefext[XII][XII.sec4]{4} 中提出的疑惑：只需把下面的命题应用于景的范畴之上的一个既纤维化又余纤维化的范畴\pageoriginale，其纤维为相应景上层范畴的反范畴（\cf \hyperlink{sga4.exp.XVII.subsection.0.2}{0.2}）。

设 $\Ee$ 是 $\Bb$ 上的范畴，$f:y\rightarrow x$ 是 $\Bb$ 的箭头，$F\in \Ob(\Ee_x)$。回忆，若对所有 $g:\sisi{Z}{z}\rightarrow y$ 和 $H\in \Ob
\Ee\sisi{_Z}{_z}$ 都有
\[
\Hom_g(H,G)\xrightarrow[\sim]{\varphi}\Hom_{fg}(H,F){\quoi}.
\]
则称二元组 $(G,\varphi)$ $(G\in \Ob \Ee_y, \varphi\in
\Hom_f(G,F))$ 为 $F$ 关于 $f$ 的\emph{严格逆像}（strict inverse image）。严格正像的定义类似。若 $\Ee$ 在 $\Bb$ 上是纤维化或余纤维化的，则这一概念归结为逆像的概念。

\begin{enonce*}{命题 2.1.3}\etiquette[2.1.3]{XVII.prop2.1.3}
设 $\epsilon$ 是 $\Bb$ 上的范畴，$y\in \Ob
\Bb$，$F\in \Ob \Ee_y$，并考虑 $\Bb$ 中的交换图：
$$\etiquette[2.1.3.1]{XVII.eq2.1.3.1}
\vcenter{
\xymatrix{
y'  \ar[r]^{g'} \ar[d]_{f'}& y\ar[d]^{f}         \\
x'          \ar[r]    & x.
}}\leqno{(2.1.3.1)}
$$
假设正像和逆像 $f_{\ast}F$、$g^{\ast}f_{\ast}F$、${g'}^{\ast}F$ 及 $f'_{\ast}{g'}^{\ast}F$ 在严格意义下存在。那么存在唯一的箭头 $\varphi \in \Fl \Ee_{x'}$，称为\emph{基变换箭头}，使图
$$\etiquette[2.1.3.2]{XVII.eq2.1.3.2}
\vcenter{
\xymatrix{
                             &  g'^{\ast}F \ar[ld]  \ar[r]    & F \ar[d]\\
f'_{\ast}g'^{\ast}F \ar[rd]_{\varphi}  &                                & f_{\ast}F\\
  & g^{\ast}f_{\ast}F \ar[ru]&.
}}\leqno{(2.1.3.2)}
$$
交换。若\pageoriginale 此外正像 $g'_{\ast}{g'}^{\ast
}F$ 和 $g'_{\ast}f'_{\ast}{g'}^{\ast}F$ 在严格意义下存在，则箭头 $\varphi$ 与\emph{\hyperrefext[XII][XII.sec4]{4}} 中定义的箭头一致。若 ${f'}^{\ast}f'_{\ast}{g'}^{\ast}F$ 和 ${f'}^{\ast}f'_{\ast}F$ 在严格意义下存在，则它还与\emph{\hyperrefext[XII][XII.sec4]{4}} 中定义的第 $2$\sisi{个}{个}箭头一致。

特别地，\emph{\hyperrefext[XII][XII.sec4]{4}} 中的两个基变换箭头相同。
\end{enonce*}

由假设，图 \eqhyperref[XVII.eq2.1.3.2]{2.1.3.2} 中的箭头诱导双射
\sisi{\begin{align*}
& \Hom_y(F,F)=\Hom_{g'}({g'}^{\ast}F,F)=\Hom_{fg'}({g'}^{\ast}F,f_{\ast}F)
\quad\mbox{\rm 且}\\
&
\Hom_{x'}(f'_{\ast}{g'}^{\ast}F,g^{\ast}f_{\ast}F)=\Hom_{f'}({g'}^{\ast}F,
g^{\ast}f_{\ast}F)=\Hom_{gf'}({g'}^{\ast}F, f_{\ast}F){\quoi}.
\end{align*}}{\[\Hom_{x'}(f'_{\ast}{g'}^{\ast}F,g^{\ast}f_{\ast}F)=\Hom_{f'}({g'}^{\ast}F,
g^{\ast}f_{\ast}F)=\Hom_{gf'}({g'}^{\ast}F, f_{\ast}F).\]}
由于 $fg'=gf'$，$\Hom_{fg'}({g'}^{\ast}F,f_{\ast}F)$ 与 $\Hom_{gf'}({g'}^{\ast}F,f_{\ast}F)$ 相同，而图 \sisi{(2.1.2.2)}{\eqhyperref[XVII.eq2.1.3.2]{2.1.3.2}} 中的箭头诱导\sisi{双射}{映射}
\[
\Hom_y(F,F)\sisi{=}{\to}\Hom_{fg'}({g'}^{\ast}F,f_{\ast}F)=\Hom_{x'}(f'_{\ast}{g'}^{\ast}F,g^{\ast}f_{\ast}F){\quoi},
\]
并且 $1_F$ 的像是 $\Ee_{x'}$ 中使 \eqhyperref[XVII.eq2.1.3.2]{2.1.3.2} 交换的唯一箭头。%See the Erratum in SGA 4 1/2, p. 312.

按 \hyperrefext[XII][XII.sec4]{4} 中的定义，基变换箭头 $\psi$ 是使下图交换的箭头：
$$\etiquette[2.1.3.3]{XVII.eq2.1.3.3}
\vcenter{
\xymatrix{
f'_{\ast}g'{}^{\ast}F \ar[d]^{\psi}  & g'{}^{\ast}F \ar[d]^{u_4} \ar[l] \ar@{=}[r] &
g'{}^{\ast}F \ar[r] \ar[d]^{u_3}   & F \ar@{=}[r] \ar[d]^{u_2}&  F \ar[d]^{u_1}\\
g'_{\ast}f_{\ast}F   & {f'}^{\ast}g^{\ast}f_{\ast}F  \ar[l] \ar@{=}[r] &
g'f^{\ast}f_{\ast}F   \ar[r] & f^{\ast}f_{\ast}F \ar[r] &  f_{\ast}F.
}}\leqno{(2.1.3.3)}
$$
水平箭头和 $u_1$ 是典范箭头；因此 $u_2$ 是伴随箭头，$u_3={g\sisi{^{1}}{'}}^{\ast}(u_2)$，$u_4$ 由 ${f'}^{\ast}g^{\ast}$ 与 ${g'}^{\ast}f^{\ast}$ 之间的同构导出，而箭头 $\psi$ 由 $u_4$ 通过伴随导出。

源为 ${g'}^{\ast}F$、靶为 $f_{\ast}F$ 的复合\pageoriginale 映射就是 \eqhyperref[XVII.eq2.1.3.2]{2.1.3.2} 中出现的那个映射，因而在下面的交换图中也是如此；其左侧方块取自 \sisi{(2.1.2.3)}{\eqhyperref[XVII.eq2.1.3.3]{2.1.3.3}}：
\[
\xymatrix{
g'{}^{\ast}F \ar[r]^{u_4} \ar[d]& f'{}^{\ast}g{}^{\ast}f_{\ast}F
\ar[r] \ar[d]&
f_{\ast}F \ar@{=}[d]\\
f'_{\ast}g'{}^{\ast}F  \ar[r]^{\psi} & g^{\ast}f_{\ast}F \ar[r] &
f_{\ast}F\rlap{.}
}
\]
这表明 $\psi$ 满足 $\varphi$ 的定义性质。

$\varphi$ 的定义是自对偶的，因而不必证明 \hyperref[XVII.prop2.1.3]{2.1.3} 的第二部分。

\subsection{纤维为三角范畴的范畴}\etiquette[2.2]{XVII.subsec2.2}

\begin{enonce*}{定义 2.2.1}\etiquette[2.2.1]{XVII.def2.2.1}
“范畴 $\Bb$ 上的加性范畴”是指 $\Bb$ 上的范畴 $\Aa$，其箭头集 $\Hom_f(X,Y)$（其中 $f\in \Fl (\Bb)$，$X$（\resp $Y$）位于 $f$ 的源（\resp 靶）之上）赋有阿贝尔群结构，使得：
\begin{enumerate}
\renewcommand{\theenumi}{\roman{enumi}}
\renewcommand{\labelenumi}{(\theenumi)}
\item 对于 $\Bb$ 中任意的 $x\xrightarrow{f}y\xrightarrow{g}z$，以及 $\Aa$ 中分别位于 $\sisi{s}{x}$、$y$、$z$ 之上的 $X$、$Y$、$Z$，复合法则
\[
\Hom_g(Y,Z)\times \Hom_f
(X,Y)\xrightarrow{\circ}\Hom_{gf}(X,Z)
\]
是双加性的。

\item 各纤维范畴是加性的。
\end{enumerate}
\end{enonce*}

这些阿贝尔群结构若存在，便由 $\Aa$、$\Bb$ 和 $\rho$ 唯一确定；这里不需要这一事实。\nde{这里 $\rho$ 表示结构函子 $\Aa\to \Bb$。} 若 $\Bb$ 上加性范畴之间的 $\Bb$-函子 $F:\Aa_1\rightarrow \Aa_2$ 诱导同态
\[
F:\Hom_{\Aa_1,f}(X,Y)\longrightarrow
\Hom_{\Aa_{2},f}(X,Y){\quoi}.
\]
，则称其为\emph{加性}函子。

\begin{enonce*}{定义 2.2.2}\etiquette[2.2.2]{XVII.def2.2.2}
“范畴 $\Bb$ 上的阿贝尔范畴”\pageoriginale 是指 $\Bb$ 上的加性范畴 $\Aa$，其纤维是阿贝尔范畴，并且对 $\Bb$ 的每个箭头 $f:x\rightarrow y$，双函子
\[
\Hom_f(X,Y):\Aa^{\circ}_x\times \Aa_y\longrightarrow
(\Ab)
\]
关于 $X$ 和 $Y$ 均左正合。
\end{enonce*}

若 $\Aa$ 是 $\Bb$ 上的纤维且余纤维范畴（\SGA 1 VI 6），则只要其纤维是阿贝尔范畴，它就是 $\Bb$ 上的阿贝尔范畴；这里也不需要这一事实。

\begin{enonce*}{定义 2.2.3}\etiquette[2.2.3]{XVII.def2.2.3}
“范畴 $\Bb$ 上的三角范畴”是指 $\Bb$ 上的加性范畴 $\Aa$，赋有一个函子 $T:\Aa\rightarrow \Aa$，并且对每个 $x\in \Ob \Bb$，纤维 $\Aa_x$ 赋有三角形集合 $\Delta_x$，满足：
\begin{enumerate}
\renewcommand{\theenumi}{\roman{enumi}}
\renewcommand{\labelenumi}{(\theenumi)}
\item $T$ 是 $\Aa$ 的加性 $\Bb$-自同构。

\item 每个纤维范畴 $\Aa_x$，赋予由 $T$ 诱导的函子 $T_x$ 以及集合 $\Delta_x$ 后，是一个三角范畴 \emph{（\cite{sga4.exp.XVII.bib.CD} \sisi{}{I 1.}1.1）}。

\item 若三角形 $(X,Y,Z,u,v,w)$ 和 $(X',Y',Z',u',v',w')$ 分别属于 $\Delta_x$ 和 $\Delta_{x'}$，则每个满足 $u'f=gu$ 的图
\[
\xymatrix{
            & Z   \ar[ldd]_{w}&            &  &           & Z'   \ar[ldd]|\hole_(.3){w'} & \\
            &   & Y\ar[lu]_{v} \ar[rrrr]_(.4){\sisi{j}{g}}&  &     &       &Y'\ar[lu]_{v'} \\
X   \ar[rru]^{u} \ar[rrrr]_{f}&             &          &  &   X'   \ar[rru]^{u'} &              &
}
\]
都能补成三角形的态射 $(f,g,h):(X,Y,Z)\rightarrow (X',Y',Z')$。
\end{enumerate}
\end{enonce*}

称 $T$\pageoriginale 为平移函子（translation functor）；常用 $A[1]$ 代替 $T(A)$，此时以 $A[n]$ 表示 $T^{n}(A)$ $(n\in\ZZ)$。

属于集合 $\Delta_x$ 的三角形称为\emph{特殊三角}（distinguished triangle）或\emph{正合三角形}。

\begin{enonce*}{命题 2.2.4}\etiquette[2.2.4]{XVII.prop2.2.4}
设 $\Aa$ 是 $\Bb$ 上的三角范畴，$x$ 是 $\Bb$ 的一个对象，$(X,Y,Z,u,v,w)$ 是 $\Aa_x$ 的一个特殊三角。对于每个 $f:y\rightarrow x$ 以及每个对象 $A\in
\Ob \Aa_y$，下面的无限序列正合：
\begin{align*}
\ldots  \longrightarrow \Hom_f(A,X) \longrightarrow \Hom_f(A,Y)&\longrightarrow
\Hom_f(A,Z)\\
 & \longrightarrow \Hom_f(A,X[1])\longrightarrow
\Hom_f(A,Y[1])\longrightarrow\ldots
\end{align*}
\end{enonce*}

对于 $f:x\rightarrow y$ 和 $A\in \Ob \Aa_y$，有类似的陈述。

读者可仿照 \textsc{Verdier} \cite{sga4.exp.XVII.bib.CD} 第~4 页证明这一命题。下面的命题的含义及其验证也留给读者：

\begin{enonce*}{命题 2.2.5}\etiquette[2.2.5]{XVII.prop2.2.5}
设 $\Aa$ 是 $\Bb$ 上的加性（\resp 阿贝尔，\resp 三角）范畴。对于每个靶为 $\Bb$ 的函子 $\Bb'\rightarrow \Bb$，范畴 $\Aa'=\fprod{\Aa}{\Bb'}{\Bb}$ 是 $\Bb'$ 上的加性（\resp 阿贝尔，\resp 三角）范畴。
\end{enonce*}

\begin{enonce*}[remark]{例 2.2.6}\etiquette[2.2.6]{XVII.exem2.2.6}
变动环化景上的模层范畴是环化景范畴之上的阿贝尔范畴。
\end{enonce*}

\begin{enonce*}[remark]{例 2.2.7}\etiquette[2.2.7]{XVII.exem2.2.7}
设 $\Aa$ 是 $\Bb$ 上的加性范畴。$\Bb$ 上 $\Aa$ 的复形范畴 $C(\Aa/\Bb)$ 的对象为各纤维范畴的复形。从 $(X^n,d)$ 到 $(Y^n,d)$ 的箭头 $f$ 是一族\pageoriginale 箭头 $f^{n}:X^{n}\rightarrow Y^{n}$，满足 $\sisi{fd}{df}^{n}=f^{n+1} d$。令 $p(f)=p(f^{n})$（这个箭头与 $n$ 无关，因为 $p(d)$ 都是恒等箭头）。$C(\Aa/\Bb)$ 是 $\Bb$ 上的加性范畴。令
\[
(K[1])^{n}=K^{n+1}\quav (d[1])\sisi{_n}{^n}=-d\sisi{_{n+1}}{^{n+1}}{\quoi}.
\]
即定义了“平移”函子。
\end{enonce*}

\begin{enonce*}[remark]{例 2.2.8}\etiquette[2.2.8]{XVII.exem2.2.8}
设 $\Aa$ 是 $\Bb$ 上的加性范畴。\emph{$\Bb$ 上 $\Aa$ 的复形模去同伦所得的范畴 $K(\Aa/\Bb)$} 与 $C(\Aa/\Bb)$ 有相同的对象。它由 $C(\Aa/\Bb)$ 中与零同伦的箭头置零得到；换言之，对于 $\Bb$ 中的 $f:x\rightarrow y$，以及 $X^{\bigdot}\in \Ob C(\Aa/\Bb)_x$ 和 $Y^{\bigdot}\in
\Ob C(\Aa/\Bb)_y$，令
\[
\Hom_f(X^{\bigdot},
Y^{\bigdot})=H^{0}\Hom^{\bigdot}_f(X^{\bigdot},Y^{\bigdot})
\]
\end{enonce*}

平移函子降为商上的函子。

若 $F:\Bb'\rightarrow \Bb$ 是以 $\Bb$ 为靶的函子，且 $\Aa'=\fprod{\Aa}{\Bb'}{\Bb}$，则将 $C(\Aa'/\Bb')$ 与 $\fprod{C(\Aa/\Bb)}{\Bb'}{\Bb}$、$K(\Aa'/\Bb')$ 与 $\fprod{K(\Aa/\Bb)}{\Bb'}{\Bb}$ 等同的工作留给读者。特别地，范畴 $K(\Aa/\Bb)$ 的纤维为范畴 $K(\Aa_x)$，因此是三角范畴。

\begin{enonce*}{命题 2.2.9}\etiquette[2.2.9]{XVII.prop2.2.9}
\begin{enumerate}
\renewcommand{\theenumi}{\roman{enumi}}
\renewcommand{\labelenumi}{(\theenumi)}
\item 赋予 \hyperref[XVII.exem2.2.8]{2.2.8} 中定义的附加结构后，范畴 $K(\Aa/\Bb)$ 是 $\Bb$ 上的三角范畴。

\item 加性（\resp 三角）范畴 $C(\Aa/\Bb)$（\resp $K(\Aa/\Bb)$）的构造与任意底范畴变换 $\Bb'\rightarrow \Bb$ 相容。
\end{enumerate}
\end{enonce*}

定义 \hyperref[XVII.def2.2.3]{2.2.3} 的条件\pageoriginale (i) 和 (ii) 逐纤维验证。仿照 \textsc{Verdier} \cite{sga4.exp.XVII.bib.CD} 第~4 页即可验证条件 (iii)。

\begin{enonce*}[remark]{注 2.2.10}\etiquette[2.2.10]{XVII.rem2.2.10}
前面的讨论也适用于范畴\sisi{}{} $K^{+}(\Aa/\Bb)$、$K^{-}(\Aa/\Bb)$ 及 $K^{b}(\Aa/\Bb)$，它们分别由限制到下有界、上有界以及有界复形得到。
\end{enonce*}

\begin{enonce*}[remark]{注 2.2.11}\etiquette[2.2.11]{XVII.rem2.2.11}
取 $\Bb$ 为终范畴，就得到通常的复形范畴。
\end{enonce*}

\setcounter{subsection}{2}
\setcounter{subsubsection}{11}
\subsubsection{}\etiquette[2.2.12]{XVII.subsubsec2.2.12}
设 $F$ 是从范畴 $\Bb^{\circ}$ 到 $(\Cat)$ 的伪函子。已知（\SGA 1 VI），$F$ 定义一个 $\Bb$ 上的纤维化范畴，其纤维为范畴 $F(x)$。对于 $\Bb$ 中任意的 $f:x\rightarrow y$ 及 $X\in \Ob F(\sisi{n}{x})$、$Y\in \Ob
F(y)$，按定义有
$$\etiquette[2.2.12.1]{XVII.eq2.2.12.1}
\Hom_f(X,Y)=\Hom_{F(x)}(X,F(f)(Y)){\quoi}.\leqno(2.2.12.1)
$$
已知（同上），前面的构造定义了从 $\Bb^{\circ}$ 到 $(\Cat)$ 的伪函子的 2-范畴与 $\Bb$ 上纤维化范畴的 2-范畴之间的等价。

从一个由 $\Bb^{\circ}$ 到如下 2-范畴的伪函子出发：其对象为加性（\resp 阿贝尔，\resp 三角）范畴，其 1-箭头为加性（\resp 左正合，\resp 三角）函子。立即可验证，由 \sisi{(2.1.12.1)}{\eqhyperref[XVII.eq2.2.12.1]{2.2.12.1}} 定义的范畴是加性（\resp 阿贝尔，\resp 三角）的。反过来：

\begin{enonce*}{命题 2.2.13}\etiquette[2.2.13]{XVII.prop2.2.13}
上述\pageoriginale 构造建立了以下两者之间的等价：
\begin{enumerate}
\renewcommand{\theenumi}{\alph{enumi}}
\renewcommand{\labelenumi}{(\theenumi)}
\item 从 $\Bb^{\circ}$ 到如下 $2$-范畴的伪函子的 $2$-范畴：其对象为加性（\resp 阿贝尔，\resp 三角）范畴，其 $1$-箭头为加性（\resp 左正合，\resp 三角）函子。

\item $\Bb$ 上加性（\resp 阿贝尔，\resp 三角）范畴中在 $\Bb$ 上纤维化的那些范畴的 $2$-范畴。
\end{enumerate}
\end{enonce*}

\SGA 1 VI 中建立的理论表明，只需验证下面这一点：

\setcounter{subsubsection}{13}
\setcounter{paragraph}{0}
\paragraph{}\etiquette[2.2.13.1]{XVII.par2.2.13.1}
设 $\Aa$ 是 $\Bb$ 上的加性（\resp 阿贝尔，\resp 三角）范畴，$f:x\rightarrow y$ 是 $\Bb$ 的箭头，并假设逆像函子 $f^{\ast}:\Aa_y\rightarrow \Aa\sisi{_n}{_x}$ 存在；那么 $f^{\ast}$ 是加性的（\resp 左正合的，\resp 三角的）。

函子 $f^{\ast}$ 给出同构
\[
\Hom_f(X,Y)\simeq \Hom_x (X,f^{\ast}(Y)){\quoi}.
\]
在三角情形\sisi{}{}，还须证明 $f^{\ast}$ 关于使下图交换的同构 $Tf^{\ast}=f^{\ast}T$（其中 $T$ 是平移函子）为三角函子：
$$\etiquette[2.2.13.2]{XVII.eq2.2.13.2}
\vcenter{
\xymatrix{
\Hom_f(X,Y)  \ar@{=}[r] \ar@{=}[d]& \Hom_x(X,f^{\ast}Y)\ar@{=}[rd]\\
\Hom_f(TX,TY)\ar@{=}[rd]           &      & \Hom_x(TX,Tf^{\ast}Y) \\
             &      \Hom_x(TX,f^{\ast}TY)\ar@{=}[ru] &
}}\leqno{(2.2.13.2)}
$$

\begin{itemize}
\item[\emph{加性情形}：] 考虑\pageoriginale 图
$$\etiquette[2.2.13.3]{XVII.eq2.2.13.3}
\vcenter{
\xymatrix@C=2.5cm@R=1.75cm{
f^{\ast} X \ar@<-3pt>[d]_{f^{*}(u)}\ar@<3pt>^{f^{\ast}(v)}[d]        \ar[r]^{p} & X
\ar@<-3pt>_{u}[d] \ar@<3pt>^{v}[d]\\
f^{\ast} X' \ar[r]^{p'}& X'\rlap{.}
}}\leqno{(2.2.13.3)}
$$
有 $p'f^{\ast}(u+v)=(u+v)p=up+vp=p'f^{\ast}(u)+p'f^{\ast}(v)=p'(f^{\ast}(u)+f^{\ast}(v))$，因而 $f^{\ast}(u+v)=f^{\ast}(u)+f^{\ast}(v)$。

\item[\emph{阿贝尔情形}：] 由定义（\hyperref[XVII.subsec2.2]{2.2}），对每个 $X$，函子 $\Hom(X,f^{\ast}Y)=\Hom_f(X,Y)$ 关于 $Y$ 左正合，故 $f^{\ast}$ 左正合。

\item[\emph{三角情形}：] 设 $(Y',Y,Y'')$ 是 $\Aa_y$ 中的特殊三角，并设 \sisi{$(f^{\ast}Y',f^{\ast}Y)$}{$(f^{\ast}Y',f^{\ast}Y, X'')$ 是 $\Aa_x$ 中的特殊三角}。由 \hyperref[XVII.def2.2.3]{2.2.3}，这些三角形的底之间的典范态射可以补成三角形的态射，给出下图：
\[
\xymatrix@C=1.8em{
&  X''       \ar[rrr]^{\sisi{}{s}}\ar[ldd]   &     & &**[r]f^{\ast}Y''\ar[rrr]\ar[ldd]   &  & & Y''\ar[ldd] \\                      &     & f^{\ast}Y   \sisi{\ar[rrr]|\hole}{\ar@{=}[rrr]|\hole}\ar[lu]& &      &f^{\ast}Y \ar[rrr]|\hole\sisi{}{\ar[ul]}&&& Y\ar[lu]\\
  f^{\ast}Y' \sisi{\ar[rrr]}{\ar@{=}[rrr]}\ar[rru]         &&& f^{\ast}Y'\sisi{}{\ar[rru]}   \ar[rrr] && & Y'\ar[rru]&&.
}
\]

证明 $s$ 是同构，便完成了证明。对于 $\Ob \Aa_x$ 中的每个 $X$，有交换图：
\[
\text{\scriptsize
$\xymatrix@C=1em{
\Hom_x(X,f^{\ast}Y') \ar[r] \ar@{=}[d]& \Hom_x(X,f^{\ast}Y) \ar@{=}[d]\ar[r]& \Hom_x(X,X'')\ar[d]
\ar[r]& \Hom_x(X,f^{\ast}Y'[1])\ar[r]\ar@{=}[d]
& \Hom_x(X,f^{\ast}Y[1]) \ar@{=}[d]\\
\Hom_x(X,f^{\ast}Y') \ar[r]\ar@{=}[d]&  \Hom_x(X,f^{\ast}Y)\ar[r] \ar@{=}[d]&
\Hom_x(X,f^{\ast}Y'') \ar[r]\ar@{=}[d]&  \Hom_x(X,f^{\ast}Y'[1])\ar@{=}[d] \ar[r]&
\Hom_x(X,f^{\ast}Y[1]) \ar@{=}[d]\\
\Hom_f(X,Y)  \ar[r]& \Hom_f(X,Y)\ar[r]  &
\Hom_f(X,Y'') \ar[r]& \Hom_f(X,Y'[1])\ar[r] &
\Hom_x(X,Y[1]).
}$}
\]

由 \hyperref[XVII.prop2.2.4]{2.2.4}，\pageoriginale 第一行和最后一行正合。由五引理，$s$ 诱导的从 $\Hom(X,X'')$ 到 $\Hom(X,f^{\ast}Y'')$ 的映射总是双射，因此 $s$ 是同构。
\end{itemize}

当然，对于余纤维化范畴也有对偶的结果。

\subsection{平凡对偶公式}\etiquette[2.3]{XVII.subsec2.3}

设 $\Aa$ 是底范畴~$\Bb$ 上的范畴，$S$ 是 $\Aa$ 的一组箭头，每个都投到 $\Bb$ 的恒等箭头上。由 $\Aa$ 通过分式演算得到的范畴 $\Aa(S^{-1})$ 是 $\Bb$ 上的范畴（由泛性质）。

设 $\Bb'\rightarrow \Bb$ 为函子；以 $\Aa'$ 表示纤维积 $\fprod{\Aa}{\Bb'}{\Bb}$，以 $S'$ 表示 $S$ 在 $\Aa'$ 中的逆像。复合函子 $\Aa'\rightarrow \Aa\rightarrow \Aa(S^{-1})$ 使 $S'$ 的箭头可逆，因而通过 $\Aa'({S'}^{-1})$ 分解。于是得到如下的典范 $\Bb'$-函子：
$$\etiquette[2.3.1]{XVII.eq2.3.1}
\Aa'(S^{-1})\longrightarrow
\fprod{\Bb'}{\Aa(S^{-1})}{\Bb}{\quoi}.\leqno(2.3.1)
$$

\begin{enonce*}{引理 2.3.2}\etiquette[2.3.2]{XVII.lem2.3.2}
若 $\Bb'$ 是 $\Bb$ 的子范畴，并且每当 $\Bb$ 的复合箭头 $h=f\circ g$ 属于 $\Bb'$ 时，$f$ 和 $g$ 都属于 $\Bb'$，那么函子 \eqhyperref[XVII.eq2.3.1]{2.3.1} 为同构。
\end{enonce*}

注意\pageoriginale，\eqhyperref[XVII.eq2.3.1]{2.3.1} 总在对象上诱导双射，而且在当前情形，$\fprod{\Bb'}{\Aa(S^{-1})}{\Bb}$ 是 $\Aa(S^{-1})$ 的子范畴。

回忆 $\Aa(S^{-1})$ 的“显式”构造 \cite[\sisi{第}{\chap}~I 章]{sga4.exp.XVII.bib.GZ}。给定 $\Aa$ 的两个对象 $X$ 和 $Y$，考虑“形式复合”$f_1\ldots f_n$，其中每个 $f_i$ 或为 $\Aa$ 的箭头，或为 $S$ 的箭头的逆。集合 $\Hom_{\Aa(S^{-1})}(X,Y)$ 是这些“形式复合”的集合模去由以下关系生成的等价关系所得的商：
\begin{alignat*}{2}
& f_1\ldots uv\ldots f_n\equiv f_1\ldots (u\circ v)\ldots f_n &\quad &
\mbox{\rm 其中 } u,v\;\mbox{\rm 属于 } \Fl(\Aa)\\
& f_1\ldots ss^{-1}\ldots f_n\equiv f_1\ldots f_n && \mbox{\rm 对 }
s\in S\\
& f_1\ldots s^{-1}s\ldots f_n\equiv f_1\ldots f_n && \mbox{\rm 对 }
s\in S{\quoi}.
\end{alignat*}

若对 $\Aa$ 和子范畴 $\Aa'$ 应用这些构造，可验证 $\Aa'(S^{-1})$ 与 $\Aa(S^{-1})$ 中 $\Bb'$ 的逆像子范畴相同，因为若一个“形式复合”$f_1\ldots f_n$ 在 $\Bb$ 中的复合为 $\Bb'$ 的箭头，则每个 $f_i$ 的像也都是 $\Bb'$ 的箭头。

现在设 $\Aa$ 是 $\Bb$ 上的三角范畴，并且对每个 $x\in \Ob(\Bb)$，集合 $S_x=S\cap \Fl(\Aa_x)$ 是饱和乘法系（\cite{sga4.exp.XVII.bib.CD} \sisi{第}{\chap}~I 章 \S~2\sisi{. 1}{ 第 1 号}）。

\begin{enonce*}{定义 2.3.3}\etiquette[2.3.3]{XVII.def2.3.3}
在前面的假设下，若对每个函子 $F:\Bb'\rightarrow\Bb$，函子 \eqhyperref[XVII.eq2.3.1]{2.3.1} 为同构（\resp 并且 $\Aa(S^{-1})$ 在 $\Bb$ 上纤维化，\resp 余纤维化），则称范畴 $\Aa/\Bb$ 关于 $S$ \emph{几乎可导出}（almost derivable）（\resp 可导出，\resp 可余导出）。
\end{enonce*}

以\pageoriginale $I$ 表示下图所给的范畴：
$$\etiquette[2.3.4]{XVII.eq2.3.4}
I:\underset{0}{\times} \xrightarrow{F} \underset{1}{\times}{\quoi}.\leqno{(2.3.4)}
$$

由 \hyperref[XVII.prop2.2.13]{2.2.13}，给定 $I$ 上的一个三角范畴 $\Aa$，在 $\sisi{\Aa}{I}$ 上纤维化（\resp 余纤维化），等价于给定其纤维 $\Aa_0$ 和 $\Aa_1$ 以及三角逆像函子 $F^{\ast}:\Aa_1\rightarrow\Aa_0$（\resp 以及三角正像函子 $F_{\ast}:\Aa_0\rightarrow\Aa_1$）。若 $A_0\in\Ob\Aa_0$ 和 $A_1\in\Ob\Aa_1$，则由 \eqhyperref[XVII.eq2.2.12.1]{2.2.12.1}，
\[
\Hom_F(A_0,A_1)=\Hom(A_0,F^{\ast}A_1)
\]
（\resp $\Hom_F(A_0,A_1)=\Hom(F_{\ast}A_0,A_1)$）。

\begin{enonce*}{命题 2.3.5}\etiquette[2.3.5]{XVII.prop2.3.5}
设 $\Aa$ 是 $I$ 上的三角范畴，$S$ 如（\hyperref[XVII.def2.3.3]{2.3.3}）所述。
\begin{enumerate}
\renewcommand{\theenumi}{\roman{enumi}}
\renewcommand{\labelenumi}{(\theenumi)}
\item $\Aa$ 关于 $S$ 几乎可导出，并且 $\Aa(S^{-1})$ 在 $I$ 上为三角范畴。

\item 若 $\Aa$ 是纤维化的（\resp 余纤维化的），则函子 $F^{\ast}$（\resp $F_{\ast}$）关于 $S_0$ 和 $S_1$ 右可导出（\resp 左可导出）（\hyperref[XVII.def1.2.1]{1.2.1}），当且仅当 $\Aa$ 在 $I$ 上可导出（\resp 可余导出）。
\end{enumerate}
\end{enonce*}

由引理 \hyperref[XVII.lem2.3.2]{2.3.2}，$\Aa(S^{-1})$ 的纤维范畴为 $\Aa_0(S^{-1}_0)$ 和 $\Aa_1(S^{-1}_1)$，其中 $S_i=S\cap\Fl(\Aa_i)$ $(i=0,1)$。沿用 \S \hyperref[XVII.sec1]{1} ${\rm n^{\circ}}$ \hyperref[XVII.sec2]{2} 的记号，函子 $A_0\rightarrow A^{-}_0$ 和 $A_1\rightarrow A^{+}_1$ 将 $\Aa_0(S^{-1}_0)$ 和 $\Aa_1(S^{-1}_1)$ 分别等同于 $\Aa_0$ 和 $\Aa_1$ 的 \ppro-对象与 \iind-对象的范畴。这样，可以定义 $I$ 上的一个范畴，其纤维为 $\Aa_0(S^{-1}_0)$ 和 $\Aa_1(S^{-1}_1)$：对于 $A_0\in\Aa_0$ 和 $A_1\in\Aa_1$，令
\[
\Hom_F(A_0,A_1)=\Hom_F(A^{-}_0,A^{+}_1)=\varinjlim_{A'_0{\xrightarrow[s]{}}A_{\sisi{p}{0}}}
\varinjlim_{A_1{\xrightarrow[t]{}}{A'_1}}\Hom(A'_0,A'_1){\quoi},
\]
其中\pageoriginale $s$ 和 $t$ 遍历 $S_0$ 和 $S_1$ 中以 $A_0$ 为靶（分别以 $A_1$ 为源）的元素。范畴 $\Aa$ 映到这个范畴，而且立即可验证，它在 $I$ 上以及在每次基变换之后都满足 $\Aa(S^{-1})$ 的泛性质。

为验证 (i)，只需证明这个范畴，即 $\Aa(S^{-1})$，是三角范畴，而为此只需验证条件 \hyperref[XVII.def2.2.3]{2.2.3}~(iii)。于是设 $(X_0,Y_0,Z_0)$ 和 $(X_1,Y_1,Z_1)$ 分别为 $\Aa_0$ 和 $\Aa_1$ 的特殊三角，它们给出 $\Aa_0(S^{-1}_0)$ 和 $\Aa_1(S^{-1}_1)$ 中的特殊三角 $(X^{-}_0,Y^{-}_0,Z^{-}_0)$ 和 $(X^{+}_1,Y^{+}_1,Z^{+}_1)$。
$$\etiquette[2.3.5.1]{XVII.eq2.3.5.1}
\vcenter{
\xymatrix{
X^{-}_{0}  \ar[r]^{u} \ar[d]& X^{+}_1  \ar[d]\\
Y^{-}_{0}  \ar[r]^{v} & Y^{+}_1\rlap{.}\\
}}\leqno{(2.3.5.1)}
$$

由引理 \hyperref[XVII.lem1.2.2.1]{1.2.2.1}，三角形 $(X^{-}_0,Y^{-}_0,Z^{-}_0)$ 和 $(X^{+}_1,Y^{+}_1,Z^{+}_1)$ 是特殊三角的极限；因此存在特殊三角的态射
\begin{equation*}
\begin{split}
(X^{-}_0,Y^{-}_0,Z^{-}_0) &\longrightarrow (X'_0, Y'_0,Z'_0)\\
(X'_1, Y'_1, Z'_1) &\longrightarrow (X^{+}_1, Y^{+}_1, Z^{+}_1){\quoi},
\end{split}\tag{2.3.5.2}
\end{equation*}\etiquette[2.3.5.2]{XVII.eq2.3.5.2}%Removing this comment would introduce unwanted white space.
以及\pageoriginale 一个与 \eqhyperref[XVII.eq2.3.5.1]{2.3.5.1} 相容的图
\[
\xymatrix{
X'_{0}  \ar[r]^{u'} \ar[d]& X'_1 \ar[d] \\
Y'_{0}  \ar[r]^{v'} & Y'_1\rlap{,}
}
\]
。由于 $\Aa$ 在 $I$ 上为三角范畴，此图可以扩充为三角形的态射
\[
(X'_0,Y'_0,Z'_0)\longrightarrow (X'_1, Y'_1, Z'_1){\quoi}.
\]
将后者与态射 \eqhyperref[XVII.eq2.3.5.2]{2.3.5.2} 复合，可见图 \eqhyperref[XVII.eq2.3.5.1]{2.3.5.1} 能扩充为三角形的态射，这就证明了 \hyperref[XVII.def2.2.3]{2.2.3}~(iii)。

验证 (ii)，例如假设 $\Aa$ 在 $I$ 上纤维化。若 $A_1\in\Ob(\Aa_1(S^{-1}_1))$，则 $A_1$ 在 $\Aa(S^{-1})$ 中的逆像是 $\Aa_0(S^{-1}_0)$ 的一个对象 $B$，满足在 $\Aa_0(S^{-1}_0)$ 中，
\begin{align*}
\Hom_{\Aa_0(S^{-1}_0)}(A_0,B) & =\Hom_F(A^{-}_0,
A^{+}_1)=\Hom(A^{-}_0,
F^{\ast}(A^{+}_1))\\
&=\Hom_{\Aa_0(S^{-1}_0)}
(A_0,F^{\ast}(A^{+}_1)){\quoi}.
\end{align*}
这样的对象 $B$ 存在，当且仅当 $\Aa_0(S^{-1}_0)$ 的 \iInd-对象 $F^{\ast}(A^{+}_1)$ 本质常值（\ie 可表）；这时 $B$ 表示这个 \iind-对象，它恰为 $RF(A_1)$（\hyperref[XVII.def1.2.1]{1.2.1}）；这就证明了断言。

\begin{enonce*}{推论 2.3.6}\etiquette[2.3.6]{XVII.cor2.3.6}
沿用 \hyperref[XVII.def2.3.3]{2.3.3} 的记号，若 $\Aa$ 在 $\Bb$ 上几乎可导出，则 $\Aa(S^{-1})$ 是 $\Bb$ 上的三角范畴。
\end{enonce*}

首先，纤维范畴 $\Aa(S^{-1})_x$ 可与 $\Aa_x(S^{-1}_x)$ 等同：将假设应用于基变换 $\{x\}\rightarrow\Bb$ 即可。这个\pageoriginale 范畴是三角范畴，赋有三角形集合 $\Delta_x$，因而前面的陈述有了意义。为验证条件 \hyperref[XVII.def2.2.3]{2.2.3}~(iii)，只需在每次基变换 $I\rightarrow\Bb$ 后验证，而断言由 \hyperref[XVII.prop2.3.5]{2.3.5}~(i) 得出。

由 \hyperref[XVII.prop2.3.5]{2.3.5}~(ii) 的证明立即得到：

\begin{enonce*}{定理 2.3.7}\etiquette[2.3.7]{XVII.thm2.3.7}
\emph{（平凡对偶公式）。}

设 $\Aa$ 是 $I$ 上的三角纤维化且余纤维化范畴，并以 $F_{\ast}$ 和 $F^{\ast}$ 表示正像与逆像函子。设 $S=S_0\cup S_1$ 是如 \hyperref[XVII.def2.3.3]{2.3.3} 中的乘法系。假设 $LF_{\ast}$ 在 $A_0\in\Ob\Aa(S^{-1}_0)$ 处有定义，$RF^{\ast}$ 在 $A_1\in\Ob\Aa_1(S^{-1}_1)$ 处有定义。那么
\[
\Hom_{\Aa_0(S^{-1}_0)}(A_0,RF^{\ast}A_1)=\Hom_{\Aa\sisi{_{(S^{-1})}}{(S^{-1})}}(A_0,A_1)=\Hom_{\Aa_1(S^{-1}_1)}(LF_{\ast}A_0,A_1){\quoi}.
\]
\end{enonce*}

\hyperref[XVII.thm2.3.7]{2.3.7} 中符号 $R$、$L$、${}^{\ast}$ 和 ${}_{\ast}$ 的位置只是表面上反常：事实上，在变动景上的层构成的纤维化且余纤维化范畴 \hyperlink{sga4.exp.XVII.subsection.0.2}{0.2} 中，纤维范畴是通常层范畴的反范畴。

\subsection{纤维为导出范畴的范畴}\etiquette[2.4]{XVII.subsec2.4}

\setcounter{subsubsection}{-1}
\subsubsection{}\etiquette[2.4.0]{XVII.subsubsec2.4.0}
在本 ${\rm n^{\circ}}$ 中，$\Aa$ 表示 $\Bb$ 上余纤维化的三角范畴，$S$ 表示 $\Aa$ 的一个态射乘法系，其中的态射都位于恒等箭头之上，并且对于每个 $x\in\Ob(\Bb)$，$S_x=S\cap\Fl(\Aa_x)$ 都是 $\Aa_x$ 的饱和乘法系 \cite[\sisi{}{I }1.2.1]{sga4.exp.XVII.bib.CD}。

我们\pageoriginale 要为 $\Aa$ 关于 $S$ 可余导出（\hyperref[XVII.def2.3.3]{2.3.3}）寻找一个判据；这蕴含，将所有与 $\Aa$ 的特殊三角的像同构的三角组成的集合赋给 $\Aa(S^{-1})$，它就是 $\Bb$ 上的三角范畴（\hyperref[XVII.prop2.3.5]{2.3.5}）。

\begin{enonce*}{命题 2.4.1}\etiquette[2.4.1]{XVII.prop2.4.1}
设 $\Aa$、$\Bb$ 和 $S$ 如上。要使 $\Aa$ 关于 $S$ 可余导出，只需满足以下条件：
\begin{itemize}
\item[$(\ast)$] $\forall x\in\Ob\Bb$，$\forall A\in\Ob(\Aa_x)$，$\exists s:A'\rightarrow A$ 属于 $S$\sisi{}{，} $\forall g$，$f:x\xrightarrow{f}y\xrightarrow{g}\sisi{Z}{z}$，$f_{\ast}A'$ 对函子 $g_{\ast}$ 左适配。
\end{itemize}
\end{enonce*}

以 $\Aa^{1}$ 表示 $\Aa$ 的满子范畴，其对象为所有正像都对所有正像函子适配的对象。于是 \hyperref[XVII.prop1.2.2]{1.2.2} 表明 $\Aa^{1}$ 是 $\Aa$ 的三角子范畴，显然是余纤维化的。令 $S^{1}=S\cap\text{Fl}(\Aa^{1})$。

\begin{enonce*}{引理 2.4.1.1}\etiquette[2.4.1.1]{XVII.lem2.4.1.1}
假设 $S$ 在正像下稳定。那么 $\Aa/\Bb$ 可余导出，而 $\Aa(S^{-1})$ 可用右分式演算计算。
\end{enonce*}

验证 \cite[I 2.2]{sga4.exp.XVII.bib.GZ} 的条件 c)，它或许是各条件中最不显然的一个。给定 $f$ 和 $s$ $(s\in S)$，须找出一个满足 $t\in S$ 的交换\pageoriginale 图 \eqhyperref[XVII.eq2.4.1.2]{2.4.1.2}：
$$\etiquette[2.4.1.2]{XVII.eq2.4.1.2}
\vcenter{
\xymatrix{
A'        \ar@{-->}[r]      & B'          \\
A  \ar[u]^{s}   \sisi{\ar[r]}{\ar[r]_f}          & B\rlap{,}  \ar@{-->}[u]^{t}
}}\leqno{(2.4.1.2)}
$$
设 $F$ 是 $f$ 在 $\Bb$ 中的像。由假设，$F_{\ast}(s)$ 仍属于 $S$，所以由 \textsc{Verdier}，存在交换图
\[
\xymatrix{
F_{\ast}(A')  \ar[r] & B'\\
F_{\ast}(A)    \ar[r]_{f}\sisi{\ar[u]_{s}}{\ar[u]^{F_\ast(s)}}   & B\ar[u]_{t}
}
\]
，因而也存在图 \eqhyperref[XVII.eq2.4.1.2]{2.4.1.2}。

由于 $\Aa(S^{-1})$ 用右分式演算计算，其构造与基变换 $c:\Bb'\rightarrow\Bb$ 交换是立即可见的。

此引理特别适用于 $\Aa^{1}$ 和 $S^{1}$；因此，下面的引理蕴含 \hyperref[XVII.prop2.4.1]{2.4.1}：

\begin{enonce*}{引理 2.4.1.3}\etiquette[2.4.1.3]{XVII.lem2.4.1.3}
在 $\Bb$ 上，以及在每次基变换 $c:\Bb'\rightarrow\Bb$ 后，典范函子
$\Phi:\Aa\sisi{_1}{^1}(\sisi{{S^1}}{(S^1)}^{-1})\rightarrow\Aa(S^{-1})$
都是范畴等价。
\end{enonce*}

构造 $\Phi$ 的一个拟逆。对于每个 $A\in\Ob\Aa$，选择以 $A$ 为靶、以 $A^1\in\Ob\Aa^{1}$ 为源的 $s_A\in S$，并且当 $A\in\Ob\Aa^{1}$ 时取 $s_A=\Id_A$。设 $f$ 是 $\Aa$ 中以 $B$ 为靶、以 $A$ 为源、在 $\Bb$ 中的像为 $F$ 的箭头。由于 $s_B\in S$ 以及 $F_{\ast}A^{1}\in\Ob\Aa^{1}$，可以找到 $B\sisi{^{1}}{'}$、\sisi{$v$ 和 $s\in S^{1}$}{$s\in S^1$ 和 $u$} 使下图交换：
$$\etiquette[2.4.1.4]{XVII.eq2.4.1.4}
\vcenter{
\xymatrix{
A      \ar[r]^{f} &  B      \ar@{<-}[r]_{s_B} & B\sisi{_1}{^1}       \\
A\sisi{_1}{^1} \ar[u]_{s_A}\ar[r]^{a} & F_{\ast}A\sisi{_1}{^1}  \ar[u]& B'\rlap{.} \ar[l]_{s} \ar[u]_u
}}\leqno{(2.4.1.4)}
$$

于是 $\Aa^{1}(\sisi{{S^1}}{(S^1)}^{-1})$ 的箭头\pageoriginale $f^{1}=us^{-1}a$ 只依赖于 $f$。

设 $f$ 和 $g$ 为两个可复合的箭头，其在 $\Bb$ 中的像为 $F$ 和 $G$。为证明 $(g\circ f)^{1}=g^{1}\circ f^{1}$，考虑下图，它包含了 $f$ 和 $g$ 的图 \eqhyperref[XVII.eq2.4.1.4]{2.4.1.4}：
\[
\xymatrix@C=2em{
                      & A \ar[rr]^{f}&      & B\ar[rrr]^{g}     & &
                      & C  &       \\
 A^{1} \ar[ru] \ar[rr]^{a}&          &F_{\ast}A^{1}\ar[ru] \ar[rrr]^(.6){a'}& & & (GF)_{\ast}A^1\ar[ru] &
              & &C^{1}\ar[llu]\\
& & & B^1 \ar[rrr]^(.4){b}|(.58)\hole\ar[uu]|\hole & & & G_{\ast}B^1\ar[uu] &     \\
&  & B'\ar[uu]^{s}\ar[rrr]^{b'}\ar[ru]_{u}       &   & &G_{\ast}B'\ar[ru]\ar[uu]^(.3){G_{\ast}s}    &  &C''\ar[ll]^{r}\ar[r]_{w} & C'\ar[llu]_t\ar[uu]_{v}\rlap{.}
}
\]

由于 $t\in S$，存在 $C''$、$r\in S^{1}$ 和 $w$ 使图交换。由定义，\break $g^{1}\circ f^{1}=vt^{-1}b\,u\,s^{-1}a$ 且 $(g\circ f)^{1}=vw(G_{\ast}(s)r)^{-1}a'a$，结论由图的交换性得出。

这样便定义了从 $\Aa$ 到 $\Aa^{1}(\sisi{{S^1}}{(S^1)}^{-1})$ 的函子；由 $\Aa(S^{-1})$ 的泛性质，它通过 $\Psi:\Aa(S^{-1})\rightarrow\Aa^{1}(\sisi{{S^1}}{(S^1)}^{-1})$ 分解。复合 $\Phi\Psi$ 同构于恒等函子，因为它与函子 $J:\Aa\rightarrow\Aa(S^{-1})$ 的复合同构于 $J$。类似地可验证 $\Psi\Phi$ 同构于恒等函子。

\begin{enonce*}{推论 2.4.2}\etiquette[2.4.2]{XVII.cor2.4.2}
沿用 \hyperref[XVII.prop2.4.1]{2.4.1} 的记号，只要对于 $\Bb$ 的任意有限箭头集 $B$，当只考虑 $B$ 中箭头下的正像时条件 $(\ast)$ 成立，$\Aa$ 就关于 $S$ 可余导出。
\end{enonce*}

对于\pageoriginale 每个箭头集有限的范畴 $\Bb'$，在每次基变换 $\Bb'\rightarrow\Bb$ 后都可应用命题 \hyperref[XVII.prop2.4.1]{2.4.1}。我们将取 $\Bb'$ 为以下各图给出的范畴：
\[
  \begin{gathered}
    *\\
    (0)
  \end{gathered}\qquad
  \begin{gathered}
    *\to*\\
    (1)
  \end{gathered}\qquad
  \begin{gathered}
    *\to*\to*\\
    (2)
  \end{gathered}\qquad
  \begin{gathered}
    *\to*\to*\to*\\
    (3)
  \end{gathered}
\]
考虑范畴 $\Aa_x(S^{-1}_x)$ $(x\in\Ob\Bb)$。给定 $f\in\Fl(\Bb)$ 以及这些范畴中位于 $f$ 的源与靶之上的 $A_1,A_2$，将 $\Hom_f(A_1,A_2)$ 定义为由 $f$ 确定的 (1) 型基变换后所定义的类似集合。利用 (3) 型基变换定义态射的复合。这样构造出 $\Bb$ 上的范畴 $\Aa(S^{-1})^{\ast}$，其构造与任意基变换相容；利用 (0)、(1) 和 (2) 型基变换，可验证它解决了定义 $\Aa(S^{-1})$ 的泛问题。

以下结果与 \hyperref[XVII.cor2.4.2]{2.4.2} 对偶。

\begin{enonce*}{命题 2.4.3}\etiquette[2.4.3]{XVII.prop2.4.3}
设 $\Aa$ 是 $\Bb$ 上纤维化的三角范畴，$S$ 如 \hyperref[XVII.prop2.4.1]{2.4.1} 所述。要使 $\Aa$ 关于 $S$ 可导出，只需对于 $\Bb$ 的每个有限箭头集 $B$，有
\begin{itemize}
\item[$(\ast')$] $\forall x\in\Ob(\Bb)$，$\forall A\in\Ob(\Aa_x)$，$\exists s:A\rightarrow A'$ 属于 $S$，对于 $B$ 中的 $\forall g,f:Z\xrightarrow{g}Y\xrightarrow{f}X$，$f^{\ast}A'$ 对函子 $g^{\ast}$ 右适配。
\end{itemize}
\end{enonce*}

\section{纤维化或余纤维化范畴的粘合}\etiquette[3]{XVII.sec3}\pageoriginale

\subsection{引言}\etiquette[3.1]{XVII.subsec3.1}

设 $f:Y\rightarrow X$ 是从概形 $Y$ 到概形 $X$ 的有限型分离态射，或是从局部紧空间到局部紧空间的连续映射。这时可定义一个从 $Y$ 上阿贝尔层的范畴到 $X$ 上阿贝尔层的范畴的“具有固有支集的正像”函子，记为 $f_{!}$（概形的情形，\cf \sisi{\hyperref[XVII.subsec5.3]{5.3}}{\hyperref[XVII.subsec6.1]{6.1}}；分离局部紧空间的情形，见 \textsc{Verdier} \cite{sga4.exp.XVII.bib.V3}）。当 $Y$ 和 $X$ 是 $\CC$ 上的有限型概形时，这两个定义在显然的意义下相容。遗憾的是，概形的情况与局部紧空间不同，$f_{!}$ 的导出函子具有病态行为；尽管如此，仍可直接定义从 $D^{+}(Y)$ 到 $D^{+}(X)$ 的合理的“具有固有支集的正像”函子。为此，把 $f$ 分解为开浸入后接固有态射：$f=\bar f j$。所求函子 $Rf_{!}$ 就定义为“零延拓”$j_{!}$ 与函子 $R\bar f_{\ast}$ 的复合，后者由 $\bar f$ 下的正像导出。

这些“具有固有支集的正像”函子 $Rf_{!}$，只有在验证了一套“协变反变形式体系”之后才便于使用；对于复合 $f=gh$，这套形式体系包含公式
\[
Rf_{!}=Rg_{!}Rh_{!}
\]
，而这最宜用余纤维化范畴的语言表述。对于“零延拓”和固有态射下的\pageoriginale 正像，分别已有这样的形式体系，问题在于将它们粘合。为避免以后陷入无谓的计算，本\S 拟列出完成这一粘合所需的初等验证。我们将在比实际需要稍抽象一些的框架下进行，希望结果以后还能用到。

上面概述的 $Rf_{!}$ 的构造要求态射 $f$ 能分解为开浸入后接固有态射。依照 \textsc{Hartshorne} \cite[III 8 第~189 页]{sga4.exp.XVII.bib.RD}，在第 \hyperref[XVII.subsec3.2]{2} 号中引入适当的可紧化态射概念。

\subsection{可紧化态射}\etiquette[3.2]{XVII.subsec3.2}

\begin{enonce*}{定义 3.2.1}\etiquette[3.2.1]{XVII.def3.2.1}
若从 $S$-概形 $X$ 到拟紧拟分离的 $S$-概形 $Y$ 的态射 $f$，存在一个在 $S$ 上固有的 $S$-概形 $P$，以及 $f$ 的如下分解：先由分离拟有限态射 $j$ 从 $X$ 映到 $\fprod{P}{Y}{S}$，再由 $\fprod{P}{Y}{S}$ 投影到 $Y$，则称该态射是\emph{$S$-可紧化的}（compactifiable）：
\[
\xymatrix{
                                & &  &\fprod{P}{Y}{S} \ar[rd]\ar[ld]& \\
X \ar[rr]^(.7){f} \ar[rrru]^{j}\ar[rrrd]& &Y \ar[rd]&                & P\ar[ld]\\
                                & &  &       S        &.
}
\]
\end{enonce*}

当 $S=Y$ 时，简称为可紧化态射。

\begin{enonce*}{命题 3.2.2}\etiquette[3.2.2]{XVII.prop3.2.2}
\begin{enumerate}
\renewcommand{\theenumi}{\roman{enumi}}
\renewcommand{\labelenumi}{(\theenumi)}
\item $S$-可紧化态射\pageoriginale 是分离且有限型的。

\item 若 $Y$ 拟紧且拟分离，则靶为 $Y$ 的分离拟有限态射是\sisi{}{$S$-}可紧化的。

\item 两个 $S$-可紧化态射的复合仍为 $S$-可紧化的。

\item 设有拟紧拟分离 $S$-概形的图
\[
\xymatrix{
\fprod{X}{Y'}{Y} \ar[d]_{f'} \ar[r] & X \ar[d]^{f}      \\
Y' \ar[r]& Y\rlap{.}
}
\]
若 $f$ 为 $S$-可紧化的，则 $f'$ 为 $S$-可紧化的。

\item 设 $g$ 是 $S$-可紧化态射；\nde{只需假设 $g$ 是 $S$-概形之间的拟紧拟分离态射。}要使复合态射 $g\circ f$ 为 $S$-可紧化的，必须使 $\sisi{h}{f}$ 也如此。
\end{enumerate}
\end{enonce*}

断言 (i)、(ii) 和 (iv) 是立即可见的。为证明 (iii) 和 (v)，考虑复合 $f=gh$，其中 $g$ 是\sisi{}{$S$-}可紧化的。

设 $gh$ 是\sisi{}{$S$-}可紧化的，于是得到图
\[
\xymatrix{
                     & \fprod{P}{Y}{S} \ar[d] \ar[r]^{v} & \fprod{P}{Z}{S}\ar[d] \\  X  \ar[ru]^{u} \ar[r]^{h}   & Y                \ar[r]^{g} &    Z\rlap{.}
}
\]
复合 $vu$ 拟有限；因此 $u$ 拟有限，(v) 得证。

若 $h$ 可紧化，则有图
\[
\xymatrix{
  &&                & &\sisi{Q\times P \times Z}{\fprod{\fprod{Q}{P}{S}}{Z}{S}} \ar[d]&\\
  &&\fprod{Q}{Y}{S} \ar[d]\ar[rru]\ar[rdd]& & \fprod{P}{Z}{S} \ar[d]\ar[rdd] & \\
X \ar[rrrdd]\ar[rru] \ar[rr]^{h}&& Y \sisi{\ar[rru]}{\ar[rru]|(.2)\hole} \ar[rdd]\sisi{\ar[rr]^{g}}{\ar[rr]^g|(.25)\hole}            & & Z\ar[ldd]                & \\
  &&                &Q\ar[d]&                  &P\ar[lld] \\
  &&                &S&                   &,
}
\]
并且\pageoriginale $f=gh$ 通过一个映到 $\fprod{(\fprod{Q}{P}{S})}{Z}{S}$ 的分离拟有限态射分解。

Nagata 在 \sisi{[5]}{\cite{sga4.exp.XVII.bib.Na}} 中断言，分离整 Noether 概形之间的每个有限型态射都可紧化。\nde{更一般地，靶拟紧且拟分离的每个有限型分离态射都可紧化。见 \nderef{XVII.ndeConrad}。}撰稿人承认自己未能理解这一证明。

现在固定概形 $S$，以 $(S)$ 表示 $S$-概形范畴的子范畴，其对象为拟紧拟分离的 $S$-概形，箭头为 $S$-可紧化的 $S$-态射。\nde{由紧化定理，也可以在（绝对）概形范畴的如下子范畴中工作：其对象为拟紧拟分离概形，箭头为有限型分离态射（\ie 可紧化态射）。注意，可紧化态射在复合下的稳定性等价于紧化定理的一般形式。}

\begin{enonce*}{命题 3.2.3}\etiquette[3.2.3]{XVII.prop3.2.3}
在范畴 $(S)$ 中：
\begin{enumerate}
\renewcommand{\theenumi}{\roman{enumi}}
\renewcommand{\labelenumi}{(\theenumi)}
\item 开浸入（\resp 固有态射）定义 $(S)$ 的子范畴 $(S,i)$（\resp $(S,p)$），其对象与 $(S)$ 相同。

\item $(S,p)$ 中存在纤维积，且它们也是 $(S)$ 中的纤维积。

\item 每个态射 $f$ 都是固有态射 $p$ 与开浸入的复合 $f=pj$。
\end{enumerate}
\end{enonce*}

由\pageoriginale \hyperref[XVII.prop3.2.2]{3.2.2} (iii)、(iv) 和 (v)，$(S)$ 中存在纤维积，且它们也是所有概形组成的范畴中的纤维积。由于两个在 $Y$ 上固有的概形的纤维积 $X=\fprod{X_1}{X_2}{Y}$ 仍在 $Y$ 上固有，并且一个概形在 $X$ 上固有，当且仅当它在 $X_1$ 和 $X_2$ 上都固有，条件 (ii) 得到满足。

若 $f$ 是 $(S)$ 的态射，则由假设，它在 $(S)$ 中有分解 $f=pj'$，其中 $p$ 固有，$j\sisi{}{'}$ 分离且拟有限。由 \EGA IV 18.12.13，$j'$ 有分解 $j'=qj$，其中 $q$ 有限，$j$ 为开浸入；此外，由 \hyperref[XVII.prop3.2.2]{3.2.2} (ii)，$q$ 和 $j$ 属于 $(S)$。于是有 $f=(pq)j$，这就证明了 (iii)。断言 (i) 由 \hyperref[XVII.prop3.2.2]{3.2.2} (iii) 得出。

\setcounter{subsubsection}{3}
\subsubsection{}\etiquette[3.2.4]{XVII.subsubsec3.2.4}
本\S 余下部分中，设给定一个范畴 $(S)$，赋有两个子范畴 $(S,i)$ 和 $(S,p)$，使得
\begin{enumerate}
\renewcommand{\theenumi}{\roman{enumi}}
\renewcommand{\labelenumi}{(\theenumi)}
\item $\Ob(S)=\Ob(S,i)=\Ob(S,p)$；

\item $(S,p)$ 中存在纤维积，且它们也是 $(S)$ 中的纤维积；

\item $\forall f\in \Fl(S)$，$\exists p\in \Fl(S,p)$，$\exists
j\in \Fl(S,i)$，$f=pj$。
\end{enumerate}

$(S,i)$（\resp $(S,p)$）的箭头将称为开浸入（\resp 固有态射）。

\begin{enonce*}{定义 3.2.5}\etiquette[3.2.5]{XVII.def3.2.5}
\begin{enumerate}
\renewcommand{\theenumi}{\roman{enumi}}
\renewcommand{\labelenumi}{(\theenumi)}
\item $(S)$ 的态射 $f$（\resp 一对可复合态射 $(f,g)$，\resp 三个可复合态射 $(f,g,h)$）的\emph{紧化}，是指水平箭头为开浸入、斜箭头为固有态射的交换图 \eqhyperref[XVII.eq3.2.5.1]{3.2.5.1}（\resp \eqhyperref[XVII.eq3.2.5.2]{3.2.5.2}，\resp \eqhyperref[XVII.eq3.2.5.3]{3.2.5.3}）。

\item \emph{紧化之间的态射}\pageoriginale 是图 \eqhyperref[XVII.eq3.2.5.1]{3.2.5.1}（\resp\ldots）之间的态射，它在第 1\sisi{个}{个}竖列上为恒等态射，并且在每个顶点处都是固有态射：
$$
\etiquette[3.2.5.1]{XVII.eq3.2.5.1}
\etiquette[3.2.5.2]{XVII.eq3.2.5.2}
\etiquette[3.2.5.3]{XVII.eq3.2.5.3}
\begin{gathered}[b]
\xymatrix{
Y        \ar[d]_{f} \ar@{^{(}->}[r]^{i}                & \bar{Y} \ar[ld]^{\bar{f}}\\
X      & &}\\
\hyperref[XVII.eq3.2.5.1]{(3.2.5.1)}
\end{gathered}\qquad
\begin{gathered}[b]
\xymatrix{
Z        \ar@{^{(}->}[r]^{j} \sisi{\ar[d]}{\ar[d]_g}& \bar{Z}\ar@{^{(}->}[r]^{j'}\sisi{\ar@{-}[ld]^{\bar{g}}}{\ar[ld]^{\bar{g}}}
&\bar{\bar{Z}} \ar[ld]^{\bar{\bar{g}}}\\
Y \ar[d]_{f}\ar@{^{(}->}[r]^{i} & {\bar{Y}} \ar[ld]^{\bar{f}}&\\
X      & &}\\
\hyperref[XVII.eq3.2.5.2]{(3.2.5.2)}
\end{gathered}\qquad
\begin{gathered}[b]
\xymatrix@=3em{
\bigdot \ar@{^{(}->}[r]_(.4){\sisi{h}{k}} \ar[d]_{h}& \bigdot \ar@{^{(}->}[r]_(.4){\sisi{h}{k}'}\ar[ld]^(.4){\bar{h}} & \bigdot
\ar@{^{(}->}[r]_(.4){\sisi{h}{k}''} \ar[ld]^(.4){\bar{\bar{h}}}&\bigdot \ar[ld]^(.4){\bar{\bar{\bar{h}}}}\\
\bigdot \ar@{^{(}->}[r]_(.4){j}\ar[d]_{g} & \bigdot \ar@{^{(}->}[r]_(.4){j'}
\ar[ld]^(.4){\bar{g}}& \bigdot \ar[ld]^(.4){\bar{\bar{g}}}\\
\bigdot
\ar@{^{(}->}[r]_(.4){i} \ar[d]_{f}&\bigdot \ar[ld]^(.4){\bar{f}}\\
\bigdot}\\
\hyperref[XVII.eq3.2.5.3]{(3.2.5.3)}
\end{gathered}
$$
\end{enumerate}
\end{enonce*}

\begin{enonce*}{命题 3.2.6}\etiquette[3.2.6]{XVII.prop3.2.6}
\begin{enumerate}
\renewcommand{\theenumi}{\roman{enumi}}
\renewcommand{\labelenumi}{(\theenumi)}
\item 每个态射（\resp 每对可复合态射 $(f,g)$，\resp 每组三个可复合态射 $(f,g,h)$）都存在紧化。

\item $f$ 的紧化范畴是左滤过的。\nde{“左滤过”意为“余滤过”（\hyperrefext[I][I.def2.7]{2.7}）。关于 \hyperref[XVII.prop3.2.6]{3.2.6} 的证明中的公理 $L_1$ 和 $L_2$，\cf \hyperrefext[V][V.par8.1.1.1]{8.1.1.1}。}

\item 若图 $\sisi{(3.2.4.2)}{\eqhyperref[XVII.eq3.2.5.2]{3.2.5.2}}_i$ $(i=1,2)$ 是 $(f,g)$ 的紧化，则存在 $(f,g)$ 的一个紧化 $\sisi{(3.2.4.2)}{\eqhyperref[XVII.eq3.2.5.2]{3.2.5.2}}_3$，以及从 $\sisi{(3.2.4.2)}{\eqhyperref[XVII.eq3.2.5.2]{3.2.5.2}}_3$ 到各紧化 $\sisi{(3.2.4.2)}{\eqhyperref[XVII.eq3.2.5.2]{3.2.5.2}}_i$ 的紧化态射。
\end{enumerate}
\end{enonce*}

\begin{enumerate}
\renewcommand{\theenumi}{\roman{enumi}}
\renewcommand{\labelenumi}{(\theenumi)}
\item 一个态射的情形正是 \hyperref[XVII.subsubsec3.2.4]{3.2.4} (iii)。处理三个态射的情形；沿用 \eqhyperref[XVII.eq3.2.5.3]{3.2.5.3} 的记号，只需依次构造 $f$、$g$ 和 $h$ 的紧化，$i\bar{g}$ 和 $j\bar{h}$ 的紧化，最后构造 $j'\bar{\bar{h}}$ 的紧化。

\item 设\pageoriginale $\eqhyperref[XVII.eq3.2.5.1]{3.2.5.1}_1$ 和 $\eqhyperref[XVII.eq3.2.5.1]{3.2.5.1}_2$ 是 $f$ 的两个紧化。从 $\fprod{\bar{Y}_1}{\bar{Y}_2}{S}$ 到 $X$ 的投影 $\sisi{\rho}{p}$ 固有，而 $Y$ 通过 $(i_1,i_2)$ 映到它上面；根据（\hyperref[XVII.subsubsec3.2.4]{3.2.4} (iii)），后者可紧化为 $(i_1,i_2)=qj$。紧化 $f=(p\sisi{\cdot}{\circ} q)j$ 控制 $\sisi{(3.2.4.1)}{\eqhyperref[XVII.eq3.2.5.1]{3.2.5.1}}_1$ 和 $\sisi{(3.2.4.1)}{\eqhyperref[XVII.eq3.2.5.1]{3.2.5.1}}_2$。

若 $r$ 和 $s$ 是从 $\sisi{(3.2.4.1)}{\eqhyperref[XVII.eq3.2.5.1]{3.2.5.1}}_1$ 到 $\sisi{(3.2.4.1)}{\eqhyperref[XVII.eq3.2.5.1]{3.2.5.1}}_2$ 的两个态射，则箭头对 $\bar{Y}_{1} \overset{r}{\underset{s}\rightrightarrows} \bar{Y}_{2}$ 的核 $K$ 通过固有态射映到 $\bar{Y}_1$，而 $i_1$ 通过 $K$ 分解：写成 $i_1=ki$。若 $i=pi_3$ 是 $i$ 的紧化，则紧化 $f=(\bar{f}_1kp)i_3$ 的存在性表明公理 $L_2$ 得到满足。

\item 对从 $\sisi{(3.2.4.2)}{\eqhyperref[XVII.eq3.2.5.2]{3.2.5.2}}_i$ $(i=1,2)$ 导出的 $f$ 和 $g$ 的紧化，应用滤过范畴的公理 $L_1$。由此得到 $\bar{Y}_3$ 和 $\bar{Z}_3$，给出交换图 $\eqhyperref[XVII.eq3.2.5.1]{3.2.5.1}_i$ $(i=1,2)$：
\[
\xymatrix{
&         &\bar{Z}_i \ar[d] \ar@{^{(}->}[r]& \bar{\bar{Z}}_i\ar[d]\\
\bar{Z}_3 \ar[rru] \ar[d]^{\bar{g}_3} &&Y   \ar@{^{(}->}[r]      & \bar{Y}_i\\
Y \ar@{=}[rru]\ar[r]_{i_3} & \bar{Y}_3 \ar[rru]&&\sisi{(3.2.5.1)_i}{},
}
\]
只需再找到通过 $\bar{\bar{Z}}_3$ 的 $i_3\bar{g}_3$ 的紧化，使得 $\bar{\bar{Z}}_3$ 能通过固有态射映到 $\bar{\bar{Z}}_i$，从而给出交换立方体 $(i=1,2)$。若 $i_3\bar{g}_3$ 的紧化之间的某个态射的靶具有这一性质，则其源也具有这一性质；由 (ii)，只需分别考虑 $i=1$ 和 $i=2$ 的情形。设\pageoriginale $T$ 为 $\bar{Y}_3$ 和 $\bar{\bar{Z}}_i$ 在 $\bar{Y}_i$ 上的纤维积：于是 $T$ 出现在一个与所求立方体类似的交换立方体中，只是从 $\bar{Z}_3$ 到 $T$ 的箭头 $k$ 可能不是开浸入。写成 $k=pj$（$p$ 固有，$j$ 为开浸入），并用 $p$ 的源 $\bar{\bar{Z}}_3$ 代替 $T$，便得到所求立方体。
\end{enumerate}

\begin{enonce*}{注 3.2.7}\etiquette[3.2.7]{XVII.rem3.2.7}
在可紧化态射的定义 \hyperref[XVII.def3.2.1]{3.2.1} 中，即使去掉 $Y$ 拟紧且拟分离的要求，改为要求 $j$ 是拟紧开浸入而不只是拟有限态射，命题 \hyperref[XVII.prop3.2.3]{3.2.3} 仍然成立。
\end{enonce*}

\subsection{粘合}\etiquette[3.3]{XVII.subsec3.3}

回忆，我们假设 \hyperref[XVII.subsubsec3.2.4]{3.2.4} 中的条件得到满足。

\subsubsection{}\etiquette[3.3.1]{XVII.subsubsec3.3.1}
设给定：
\begin{itemize}
\item[$\alpha$)] 对每个 $X\in \Ob (S)$，给定一个范畴 $F(X)$，

\item[$\beta$)] 给定一个 $(S,i)$-余纤维化范畴 $F_i$ 及一个 $(S,p)$-余纤维化范畴 $F_{p}$，

\item[$\gamma$)] 对每个 $X\in \Ob(S)$，给定 $F(X)$、$(F_{i})_X$ 和 $(F_{p})_X$ 之间的同构，

\item[$\delta$)] 给定 $F_{i}$ 的一个正规化劈分以及 $F_{p}$ 的一个正规化劈分。
\end{itemize}

由（\SGA 1 VI），这与改为给定 $\alpha$) 及以下资料“是一回事”：
\begin{itemize}
\item[(i)] 对每个固有态射 $p:X\rightarrow Y$，给定函子 $p_{\ast}:F(X)\rightarrow F(Y)$，

\item[(ii)] 对每对可复合的固有态射 $(p,q)$，给定函子同构\pageoriginale $c_{p,q}:p_{\ast}q_{\ast}\leftrightarrow (pq)_{\ast}$，

\item[(i\sisi{}{'}) 及 (ii\sisi{)'}{')},] 对 $(S,i)$ 给定类似的资料，
\end{itemize}
这些资料满足以下条件：
\begin{itemize}
\item[(a)] 若 $p$ 为恒等态射，则 $p_{\ast}$、$c_{p,q}$ 和 $c_{q,p}$ 都为恒等态射，

\item[(b)] 对任意三个可复合的固有态射，有
\[
c_{p,qr}\circ (p_{\ast}\ast c_{q,r})=c_{pq,r}\circ (c_{p,q}\ast
r_{\ast}){\quoi},
\]

\item[(a\sisi{)'}{')} 及 (b\sisi{)'}{')},] 对 $(S,p)$ 有类似的条件。
\end{itemize}

此外，假设\emph{给定}：
\begin{itemize}
\item[(iii)] 对每个交换图 $D$
\[
\vcenter{
\xymatrix{
X \ar[r]^{j'} \ar[d]^{p'}& Y\ar[d]^{p} \\
Z \ar[r]_{j}             &           T
}}\quad j,j' \in \Fl(S,i)\text{; }
  p,p' \in \Fl(S,p),
\]
给定函子同构 $d(D):p_{\ast}j'_{\ast}\leftrightarrow
j_{\ast}p'_{\ast}$。
\end{itemize}

\begin{enonce*}{命题 3.3.2}\etiquette[3.3.2]{XVII.prop3.3.2}
若资料 {\rm (i) (ii) ${\rm (i')}$ ${\rm (ii')}$ (iii)} 满足条件 {\rm (a) (b) ${\rm (a')}$ ${\rm (b')}$} 以及下面陈述的 {\rm(c)} 和 ${\rm (c')}$，则存在一个且“本质上”唯一的 $S$ 上余纤维化的范畴 $F$，赋有同构 $\fprod{F}{(S,i)}{S}=F_i$ 和 $\fprod{F}{(S,p)}{S}=F_p$，与资料 \hyperref[XVII.subsubsec3.3.1]{3.3.1} $\gamma$) 相容，并使 {\rm(iii)} 中的同构 $d(D)$ 成为复合同构 $j_{\ast}p'_{\ast}=(jp')_{\ast}=(pj')_{\ast}=p_{\ast}j'_{\ast}$。
\end{enonce*}

\begin{itemize}
\item[(c)] 对于以下类型的每个交换图
\[
\xymatrix{
X_1 \ar@{^{(}->}[r]^{i} \ar[d]^{p_1}& X_2 \ar@{^{(}->}[r]^{i'}
\ar[d]^{p_2}& X_3 \ar[d]^{p_3}\\
Y_1 \ar@{^{(}->}[r]^{j}& Y_2 \ar@{^{(}->}[r]^{j'}& Y_3,\\
}
\]
，其中\pageoriginale 水平箭头是开浸入，竖直箭头是固有态射，由 (iii) 中的同构 $d$ 构成的以下三角形交换：
\[
\xymatrix@R=1cm{
p_{3{\ast}} i'_{\ast} i_{\ast} \ar[rr] \ar@{-}[rd]&  & j'_{\ast}j_{\ast}p_{1{\ast}}\ar@{-}[ld]\\
                               &j'_{\ast}p_{2{\ast}}
                                i_{\ast}&.
}
\]

\item[(c')] 对于以下类型的每个交换图
\[
\xymatrix{
X_1  \ar[r]^{p} \ar[d]^{i_1}& X_2  \ar[r]^{p'} \ar[d]^{i_2}& X_3 \ar[d]^{i_3}          \\
Y_1         \ar[r]^{q} & Y_2  \ar[r]^{q'} & Y_3\rlap{,}
}
\]
，其中水平箭头是固有态射，竖直箭头是开浸入，由 (iii) 中的同构 $d$ 构成的以下三角形交换：
\[
\xymatrix{
i_{3{\ast}} p'_{\ast} p_{\ast} \ar[rr] \ar@{-}[rd]&    &q'_{\ast} q_{\ast}
i_{1{\ast}} \\
&q'_{\ast} i_{2{\ast}} p_{\ast}  \ar@{-}[ru]&.
}
\]
\end{itemize}

\setcounter{subsubsection}{2}
\paragraph{$f_{\ast}$ 的定义}\etiquette[3.3.2.1]{XVII.par3.3.2.1}\pageoriginale
设 $f$ 是从 $Y$ 到 $X$ 的态射。$f$ 的每个紧化，其图为 \sisi{(\hyperref[XVII.def3.2.5]{3.2.5} (i))}{\eqhyperref[XVII.eq3.2.5.1]{3.2.5.1}}，都定义一个从 $F(Y)$ 到 $F(X)$ 的复合函子 $\bar{f}_{\ast}i_{\ast}$。紧化之间的一个态射定义交换图
$$\etiquette[3.3.2.2]{XVII.eq3.3.2.2}
\vcenter{
\xymatrix{
Y  \ar@{^{(}->}[r]^{i_1} \ar@{=}[d]& \bar{Y}_1  \sisi{\ar[d]}{\ar[d]^p}\\
Y  \ar@{^{(}->}[r]^{i_2} \ar[d]        & \bar{Y}_2 \ar[ld]^{\bar{f}_2}\\
X                         &.
}}\leqno{(3.3.2.2)}
$$
因此资料 (iii) 给出一个同构
\[
\bar{f}_{2{\ast}} i_{2{\ast}}=\bar{f}_{2\ast} \circ
(i_{2\ast}\circ 1_{Y{\ast}})=\bar{f}_{2{\ast}}\circ
(p_{\ast}\circ i_{1{\ast}})=\bar{f}_{1{\ast}}\circ
i_{1{\ast}}{\quoi}.
\]
公理 ${\rm(c')}$ 还保证：与紧化态射的复合相伴的同构，就是与每个紧化态射相伴的同构的复合。由于紧化态射的范畴是余滤过的（\hyperref[XVII.prop3.2.6]{3.2.6} (ii)），由此得到与 $f$ 的各个紧化相伴的函子之间的传递同构系。选定其中一个，就定义了函子 $f_{\ast}$。

\setcounter{paragraph}{2}
\paragraph{传递同态}\etiquette[3.3.2.3]{XVII.par3.3.2.3}

每对可复合态射 $(f,g)$ 的紧化（图 \eqhyperref[XVII.eq3.2.5.2]{3.2.5.2}）都定义 $f$ 的一个紧化、$g$ 的一个紧化、$f\circ g=(\bar{f}\bar{\bar{g}})(j'j)$ 的一个紧化，以及同构 $c_{f,g}:(fg)_{\ast}\isomto f_{\ast}g_{\ast}$，它是以下同构的复合：
\[
(fg)_{\ast}=(\bar{g}\,\bar{\bar{g}})_{\ast}(j'j)_{\ast}=\bar{f}_{\ast}(\bar{\bar{g}}_{\ast}j'_{\ast})j_{\ast}=\bar{f}_{\ast}(i_{\ast}\bar{g}_{\ast})j_{\ast}=(\bar{f}_{\ast}i_{\ast})(\bar{g}_{\ast}j_{\ast})=f_{\ast}g_{\ast}
\]
（其中\pageoriginale 中间的同构来自 (iii)）。

证明这个同构不依赖为 $(f,g)$ 选择的紧化。由 \hyperref[XVII.prop3.2.6]{3.2.6} (iii)，只需比较从两个紧化 $\sisi{(3.2.4.2)}{\eqhyperref[XVII.eq3.2.5.2]{3.2.5.2}}_i$ $(i=1,2)$ 得到的同构 $c_{f,g}$，并假设存在从第一个紧化到第二个紧化的紧化态射。这些态射中的箭头记为 $p$。

考虑以下函子同构的图。为便于阅读，省略了 ${\ast}$。
\[
\xymatrix{
\bar{f}_2 \bar{\bar{g}}_2j'_2j_2\sisi{}{\ar@{}[rrd]|{(1)}} \ar@{-}[rr] \ar@{-}[ddd] \ar@{-}[rd]& &
\bar{f}_2\bar{\bar{g}}_2pj'_1j_1
\ar@{-}[rr] \ar@{-}[ld] \ar@{-}[rd]& &\bar{f}_1\bar{\bar{g}}_1j'_1j_1 \ar@{-}[ddd]\sisi{}{\ar@{}[lld]|{(2)}}\\
 \ar@{}[ddr]|{(3)} & \bar{f}_2\bar{\bar{g}}_2j'_2pj_1\ar@{-}[ru] \ar@{-}[dd]& &
\bar{f}_2p\bar{\bar{g}}_1j'_1j_1 \ar@{-}[dd]\ar@{-}[ru]& \ar@{}[ldd]|{(5)}\\
 &  &  (4)  & &\\
\bar{f}_2i_2\bar{g}_2j_2 \ar@{-}[r] & \bar{f}_2i_2\bar{g}_2pj_1 & &
\bar{f}_2pi_1\bar{g}_1j_1 \ar@{-}[r]&  \bar{f}_1i_1\bar{g}_1j_1\\
& &  \bar{f}_2i_21_Y\bar{g}_1j_1 \ar@{-}[ru]\ar@{-}[lu]& & .\\
}
\]

\begin{enumerate}
\renewcommand{\theenumi}{\alph{enumi}}
\renewcommand{\labelenumi}{\theenumi)}
\item 沿第 1\sisi{行}{行}的复合同构\pageoriginale，就是 $(fg)_{\ast}$ 的两个定义之间的同构 \hyperref[XVII.par3.3.2.1]{3.3.2.1}。

\item 沿最后一行的复合同构，就是 $f_{\ast}g_{\ast}$ 的两个定义之间的同构 \hyperref[XVII.par3.3.2.1]{3.3.2.1}。

\item 两侧最外面的竖直同构，就是从两个紧化之一导出的 $(fg)_{\ast}$ 与 $f_{\ast}g_{\ast}$ 之间的同构 $c_{f,g}$。
\end{enumerate}
因此，问题在于证明此图的外边界交换。一个交换的函子图，在各函子都与同一函子复合后仍然交换。考虑到这一点，(\sisi{i}{1}) 的交换性可由对图
\[
\xymatrix{
Z  \ar@{^{(}->}[r]^{j_1} \ar@{=}[d]& \bar{Z}_1  \ar@{^{(}->}[r]^{j'_1} \ar[d]& \bar{\bar{Z}}_1 \ar[d]\\
Z  \ar@{^{(}->}[r]^{j_2} & \bar{Z}_2  \ar@{^{(}->}[r]^{j'_2} & \bar{\bar{Z}}_2\rlap{.}
}
\]
应用 (c) 得出。(3) 和 (5) 的交换性是显然的；(2) 本质上涉及固有态射，其交换性也是显然的。只需考虑六边形 (4)。

考虑以下交换立方体
\[
\xymatrix{
                              &**[r]\bar{Z}_2  \ar@{^{(}->}[rrr]^{ j'_2}
                              \ar[dd]^(.35){\bar{g}_2}|\hole &    &   &
                              **[r]\bar{\bar{Z}}_2 \ar[dd]^{\bar{\bar{g}}_2}\\
  \bar{Z}_1\sisi{}{\ar@{-->}[rd]}  \ar@{^{(}->}[rrr]^(.65){ j'_1} \ar[ru] \ar[dd]^{\bar{g}_1}   &           &               &
                              \bar{\bar{Z}}_1 \ar[ru]\ar@{-->}[rd] \ar[dd]^(.35){\bar{\bar{g}}_1}|\hole&\\
& **[r] Y\ar[rrr]_{i_2}  \ar@{^{(}->}[rrr]  & && **[r]\bar{Y}_2 \\
Y  \ar@{^{(}->}[rrr]_{i_1}   \ar@{=}[ru]  &           &               &
                              \bar{Y}_1\ar[ru]  &.\\
}
\]

从\pageoriginale $\bar{Z}_1$ 沿立方体的棱走到 $\bar{Y}_2$，共有六种走法，每种都与另外两种“相邻”；与两条相邻路径（\ie 只隔着一个面的路径）相伴的复合函子之间有一个同构，而 (4) 的交换性将由所得六边形的交换性得出。

这类图的意义在于：一条棱表示一个函子，一个面表示复合函子之间的同构，一个体表示相容性条件（即其边界上各面之间的相容性）。我们沿“平面”$\bar{Z}_1 \bar{\bar{Z}}_1 Y \bar{Y}_2$ 把立方体分解成两个三棱柱，以归结到假设 ${\rm(c')}$。

更确切地，把六边形补成以下同构图（为便于阅读，省略了 ${\ast}$）：
\[
\xymatrix{
                   & & \bar{\bar{g}}_2 j'_2p      & &   \\
i_2\bar{g}_2 p  \ar@{-}[rru] \ar@{-}[rd]\ar@{-}[dd] &   &   &&   \bar{\bar{g}}_2 p j'_1 \ar@{-}[ld]\ar[dd]\ar@{-}[llu]\\
 & i_2(g_2p) \ar@{-}[rr]&      & (\bar{\bar{g}}_2p)j'_1 & \\
i_21_{Y}\bar{g}_1 \ar@{-}[ru] \ar@{-}[rrd]&      &&   &
  p\bar{\bar{g}}_1 j'_1 \ar@{-}[lu] \ar@{-}[lld]\\
&  &  pi_1 \bar{g}_1  & &
}
\]
这些\pageoriginale 三角形显然交换，而五边形由 ${\rm(c')}$ 交换。

\paragraph{上闭链条件}\etiquette[3.3.2.4]{XVII.par3.3.2.4}

验证在 \hyperref[XVII.par3.3.2.3]{3.3.2.3} 中构造的同构 $c_{f,g}$ 对于任意三个可复合态射 $(f,g,h)$ 满足上闭链条件（cocycle condition）。引入 $(f,g,h)$ 的一个紧化，其图为 \sisi{(3.2.4.3)}{\eqhyperref[XVII.eq3.2.5.3]{3.2.5.3}}，并考虑以下函子同构的图：
\[
\sisi{\xymatrix{
\bar{f} i \bar{g} j \bar{h} k \ar@{-}[rr] \ar@{-}[dd]&
&\bar{f}\bar{\bar{g}} j' j\bar{h} k \ar@{-}[dd]\\
 & \bar{f}\bar{\bar{g}}j' \bar{\bar{h}}
k' k \ar@{-}[rd]&\\
\bar{f} i \bar{g}\bar{\bar{h}}k' k \ar@{-}[ru] \ar@{-}[rr]&    &
\bar{f}\bar{\bar{g}}\bar{\bar{h}}k'' k' k.
}}{\xymatrix{
\bar{f} i \bar{g} j \bar{h} k \ar@{-}[rr]^{c_{f,g}*h_*} \ar@{-}[dd]_{f_**c_{g,h}}&&
\bar{f}\bar{\bar{g}} j' j\bar{h} k \ar@{-}[dd]\ar@{-}[rddd]^{c_{fg,h}}\\
&\\
\bar{f} i \bar{g}\bar{\bar{h}}k' k \ar@{-}[rrrd]_{c_{f,gh}} \ar@{-}[rr] && \bar{f}\bar{\bar{g}}j' \bar{\bar{h}}
k' k \ar@{-}[rd]\\
& &   & \bar{f}\bar{\bar{g}}\bar{\bar{\bar{h}}}k'' k' k.
}}
\]
方块的交换性显然，而三角形的交换性由 (c) 和 ${\rm(c')}$ 得出。所要证明的正是外边界的交换性，因而这就完成了 \hyperref[XVII.prop3.3.2]{3.3.2} 的证明。

\subsubsection{}\etiquette[3.3.3]{XVII.subsubsec3.3.3}

设在 $(S,i)$ 和 $(S,p)$ 上分别给定赋有正规化劈分的纤维化范畴 $F_i$ 和 $F_p$，并假设对于每个 $X\in \Ob S$，$F_i$ 和 $F_p$ 有相同的纤维 $F(X)$。

根据 \SGA 1 VI，这与改为给定以下资料“是一回事”：
\begin{itemize}
\item[(i bis)] 对每个固有态射 $p:X\rightarrow Y$，给定函子 $p^{\ast}:F(Y)\rightarrow F(X)$，

\item[(ii bis)] 对每对可复合的固有态射 $(p,q)$，给定函子同构 $c_{p,q}:q^{\ast}p^{\ast}\leftrightarrow (pq)^{\ast}$，

\item[(${\rm i'}$ bis)] 及\pageoriginale (${\rm ii'}$ bis)，对 $(S,i)$ 给定类似的资料；这些资料满足与 \hyperref[XVII.subsubsec3.3.1]{3.3.1} (a) (b) ${\rm (a')}$ ${\rm (b')}$ 对偶的条件 (a bis) (b bis) (${\rm a'}$ bis) (${\rm b'}$ bis)。

此外，假设\emph{给定}：

\item[(iii bis)] 对每个交换图 $D$
\[
\vcenter{
\xymatrix{
X \ar[r]^{j'} \ar[d]^{p'}&  Y\ar[d]^{p} \\
Z \ar[r]^{j}& T
}}\quad j,j' \in \Fl(S,i)\text{; }p\sisi{}{,} p' \in \Fl(S,p)
\]
给定函子同构 $d(D):{j'}^{\ast}p^{\ast}\leftrightarrow
{p'}^{\ast}j^{\ast}$。
\end{itemize}

与 \hyperref[XVII.prop3.3.2]{3.3.2} (c) 和 ${\rm (c')}$ 对偶的条件 (c bis) 和 (${\rm c'}$ bis) 的陈述留给读者。

与前面平行的论证即可证明下面的命题。

\begin{enonce*}{命题 3.3.4}\etiquette[3.3.4]{XVII.prop3.3.4}
若资料 {\rm (i bis) (ii bis) (${\rm i'}$ bis) (${\rm ii'}$ bis) (iii bis)} 满足条件 {\rm (a bis) (b bis) (${\rm
a'}$ bis) (${\rm b'}$ bis) (c bis) (${\rm c'}$ bis)}，则存在一个且本质上唯一的 $S$ 上纤维化的范畴 $F$，赋有同构 $\fprod{F}{(S,i)}{S}=F_i$、$\fprod{F}{(S,p)}{S}=F_p$；这些同构与等同 $F_{iX}=F_{pX}$ $(X\in \Ob S)$ 相容，并使 {\rm (iii bis)} 中的同构 $d(D)$ 成为复合同构 ${p'}^{\ast}j^{\ast}=(jp')^{\ast}=(pj')^{\ast}={j'}^{\ast}p^{\ast}$。
\end{enonce*}

然而应注意，\hyperref[XVII.prop3.3.4]{3.3.4} 并不与 \hyperref[XVII.prop3.3.2]{3.3.2} 对偶，因为对 $(S,i)$ 和 $(S,p)$ 所加的假设并不自对偶。


% END native input: chapters/exp-17-part2-zh.tex
% BEGIN native input: chapters/exp-17-part3-zh.tex
\section{消解及其对基变换态射的应用}\etiquette[4]{XVII.sec4}\pageoriginale

\subsection{平坦消解}\etiquette[4.1]{XVII.subsec4.1}

回顾一下, 若 $f:(X,\Aa)\rightarrow (Y,\Bb)$ 为赋环景 (或赋环拓扑斯) 的态射, 则左 $\Bb$ 模层 $\Mm$ 在 $f$ 下的逆像为左 $\Aa$ 模层 $\Aa\otimes_{f^{\ast}\Bb} f^{\ast}\Mm$.

\begin{enonce*}{命题 4.1.1}\etiquette[4.1.1]{XVII.prop4.1.1}
设 $f:(X,\Aa)\rightarrow (Y,\Bb)$ 为赋环景态射, $K$ 为右 $\Aa$ 模层, $0\rightarrow L\rightarrow M\rightarrow N\rightarrow 0$ 为左 $\Bb$ 模层的正合列. 若 $M$ 与 $N$ 平坦, 则
\begin{enumerate}
\renewcommand{\theenumi}{\roman{enumi}}
\renewcommand{\labelenumi}{(\theenumi)}
\item $L$ 平坦,

\item 下列序列正合:
\[
0\longrightarrow K\otimes_{\Aa}f^{\ast}L\longrightarrow
K\otimes_{\Aa}f^{\ast}M\longrightarrow
K\otimes_{\Aa}f^{\ast}N\longrightarrow 0{\quoi}.
\]
\end{enumerate}
\end{enonce*}

将 $f$ 分解为箭头 $(X,\Aa)\xrightarrow{u}(X,f^{\sisi{\ast}{-1}}\Bb)\xrightarrow{v}(Y,\Bb)$. 函子 $v^{\ast}$ 正合且将平坦模映为平坦模 (\sisi{IV 15}{\hyperrefext[V][V.sec1]{1}}), 所以只须对 $u$ 证明 \hyperref[XVII.prop4.1.1]{4.1.1}. 同构
$$\etiquette[4.1.1.1]{XVII.eq4.1.1.1}
K\otimes_{\Aa}(\Aa\sisi{_{f^{\ast}\Bb}\otimes}{\otimes_{f^{-1}\Bb}} P)=K\otimes_{f^{\sisi{\ast}{-1}}\Bb}P\leqno(4.1.1.1)
$$
又将问题归结为 $f$ 为恒等态射的情形.

取正合列 $0\rightarrow R\rightarrow P\rightarrow
K\rightarrow 0$, 其中 $P$ 平坦\sisi{ (\hyperref[XVII.subsubsec4.1.2]{4.1.2})}{}. 下图的各行及各列均正合:
\[
\xymatrix{
 &      &                          &  0\ar[d]                        &  &\\
 &      & R\otimes L  \ar[r]\ar[d] & P\otimes L  \ar[r]\ar[d] & K\otimes L
 \ar[r]\ar[d]& 0\\
 & 0\ar[r]     & R\otimes M  \ar[r]\ar[d] & P\otimes M  \ar[r]\ar[d] & K\otimes M
 \ar[r]\ar[d]& 0\\
 & 0\ar[r]     & R\otimes N  \ar[r] \ar[d]& P\otimes N  \ar[r]\ar[d] & K\otimes M
 \ar[r] \ar[d]& 0\\
 &      & 0         &   0   &   0   &.
}
\]

对前两列应用\pageoriginale 蛇引理, 得到
\[
0\xrightarrow{j}K\otimes L\rightarrow K\otimes M{\quoi},
\]
这就证明了 (ii). 现在去掉 $P$ 平坦的假设. 在前图中, 三列及两行为短正合列; 因而第一行正合, 即 $L$ 平坦.

\setcounter{subsubsection}{1}
\subsubsection{}\etiquette[4.1.2]{XVII.subsubsec4.1.2}
设 $\Ss$ 为赋环景范畴的一个子范畴.\nde{\MakeLinkTarget*{sga4.exp.XVII.editor-note.4.1.2.ringed-site-category}\label{XVII.ndesous}更一般地, 该理论对由赋环景组成的范畴 $\Ss$ 也成立.} $\Ss$ 中各赋环景上的模层构成范畴 $\Mm$, 它在 $\Ss$ 上阿贝尔 (\hyperref[XVII.def2.2.2]{2.2.2}), 且为纤维化及余纤维化, 其纤维是景上通常的模层范畴的对偶范畴. 以 $K(\Ss)$ 表示范畴 $K(\Mm)$ (\hyperref[XVII.exem2.2.8]{2.2.8}), 以 $K^{+}(\Ss)$ (\resp $K^{-}(\Ss)$) 表示 $K(\Ss)$ 的全子范畴, 其纤维范畴是 $S\in \Ob \Ss$ 时各范畴 $K^{+}(S)$ (\resp $K^{-}(S)$) 的对偶范畴. 范畴 $K(\Ss)$、$K^{+}(\Ss)$ 及 $K^{-}(\Ss)$ 在 $\Ss$ 上三角, 且为纤维化及余纤维化.

设\pageoriginale $f:S\rightarrow S'$ 为 $\Ss$ 中的箭头. 每个内射对象对于函子 $Rf_{\ast}$ 无上同调 (\hyperref[XVII.def1.2.3]{1.2.3}), 每个平坦模层对于函子 $Lf^{\ast}$ 无上同调 (\hyperref[XVII.prop4.1.1]{4.1.1} 及 \hyperref[XVII.prop1.2.6]{1.2.6}), 从而得到导出函子 (\hyperref[XVII.prop1.2.7]{1.2.7})
\begin{align*}
Rf_{\ast}:D^{+}(S) &\longrightarrow D^{+}(S')\tag{4.1.2.1}\\
Lf^{\ast}:D^{-}(S\sisi{}{'}) &\longrightarrow D^{-}(S\sisi{'}{}){\quoi}.\tag{4.1.2.2}
\end{align*}\etiquette[4.1.2.1]{XVII.eq4.1.2.1}
\etiquette[4.1.2.2]{XVII.eq4.1.2.2}%原文提示: 此段后的竖直空白过大.

若 $Rf_{\ast}$ (\resp $Lf^{\ast}$) 具有有限维数, 则在 \eqhyperref[XVII.eq4.1.2.1]{4.1.2.1} (\resp \eqhyperref[XVII.eq4.1.2.2]{4.1.2.2}) 中, 可以将上标 $+$ (\resp $-$) 换成 $b$、换成 $-$ (\resp $+$), 或将其省略 (\hyperref[XVII.prop1.2.10]{1.2.10}).\nde{\MakeLinkTarget*{sga4.exp.XVII.editor-note.4.1.2.spaltenstein-resolutions}\label{XVII.ndeKS}使用 Spaltenstein 消解, 可以无条件地省略这些上标. 交换情形可见 M.~Kashiwara 与 P.~Schapira, \emph{Categories and Sheaves}, Springer, 2010 (下文引用为 [KS]), 18.6.}

以 $D(\Ss)$ (\resp $D^{+}(\Ss)$, \resp $D^{-}(\Ss)$) 表示在 $K(\Ss)$ (\resp $K^{+}(\Ss)$, \resp $K^{-}(\Ss)$) 中将拟同构取逆所得的范畴.

\begin{enonce*}{附论 4.1.3}\etiquette[4.1.3]{XVII.sch4.1.3}
(“平凡对偶定理”).
\begin{enumerate}
\renewcommand{\theenumi}{\alph{enumi}}
\renewcommand{\labelenumi}{\theenumi)}
\item 假设所有函子 $Rf_{\ast}$ (\resp $Lf^{\ast}$) 都具有有限维数. 则:
\begin{itemize}
\item[(i)] $D(\Ss)$ 的形成与基范畴 $\Ss'\rightarrow \Ss$ 的任意变换可交换. 范畴 $D(\Ss)$ 在 $\Ss$ 上三角 \emph{(\hyperref[XVII.def2.2.3]{2.2.3})}, 且其纤维等同于 $S\in \Ob
\Ss$ 时各范畴 $D(S)$ 的对偶范畴.

\item[(ii)] 范畴 $D(\Ss)$ 在 $\Ss$ 上为余纤维化 (\resp 纤维化). 设 $f:S\rightarrow S'$; 若 $Lf^{\ast}$ (\resp $Rf_{\ast}$) 具有有限维数, 则 $D(S')$ (\resp $D(S)$) 的每个对象在 $\Ss$ 上的范畴 $D(\Ss)$ 意义下都有 $f$ 下的逆像 (\resp 正像). 逆像函子及正像函子等同于 $Lf^{\ast}$ 及 $Rf_{\ast}$.

\item[(iii)] 每个\pageoriginale $D^{-}(S')$ (\resp $D^{+}(S)$) 的对象在 $D(\Ss)$ 意义下都有逆像 (\resp 正像), 且逆像 (\resp 正像) 函子等同于 $Lf^{\ast}$ (\resp $Rf_{\ast}$).\nde{(a) 中的有限维数假设是多余的. \Cf 前面的编者注.}
\end{itemize}

\item \begin{itemize}
\item[(i)] $D^{-}(\Ss)$ (\resp $D^{+}(\Ss)$) 的形成与基范畴 $\Ss'\rightarrow \Ss$ 的任意变换可交换. 范畴 $D^{-}(\Ss)$ (\resp $D^{+}(\Ss)$) 在 $\Ss$ 上三角, 且其纤维等同于各范畴 $D^{-}(S)$ (\resp $D^{+}(S)$) 的对偶范畴.

\item[(ii)] 范畴 $D^{-}(\Ss)$ (\resp $D^{+}(\Ss)$) 在 $\Ss$ 上为纤维化 (\resp 余纤维化); “逆像”(\resp “正像”) 函子等同于导出函子 $Lf^{\ast}$ (\resp $Rf_{\ast}$).

\item[(iii)] 若 $f:S\rightarrow S'$ 使 $Rf_{\ast}$ (\resp $Lf^{\ast}$) 具有有限维数, 则 $D^{-}(S)$ (\resp $D^{+}(S')$) 的每个对象在 $D^{-}(\Ss)$ (\resp $D^{+}(\Ss)$) 意义下都有正像 (\resp 逆像), 且“正像”(\resp “逆像”) 函子等同于 $Rf_{\ast}$ (\resp $Lf^{\ast}$).
\end{itemize}
\end{enumerate}
\end{enonce*}

只须验证每种情形下均满足 \hyperref[XVII.prop2.4.1]{2.4.1} 或 \hyperref[XVII.prop2.4.3]{2.4.3} 的假设; 其他断言在作基变换 $\{S\}\rightarrow \Ss$ 或 $I\rightarrow
\Ss$ 后, 由 \hyperref[XVII.prop2.3.5]{2.3.5} 及 \hyperref[XVII.prop1.2.10]{1.2.10} 得到. 在情形 (b) 中, 平坦 (\resp 松弛) 层的复形满足 \sisi{\hyperref[XVII.prop2.4.1]{2.4.1}}{\hyperref[XVII.prop2.4.3]{2.4.3} (\resp \hyperref[XVII.prop2.4.1]{2.4.1})} 的假设. 在情形 (a) 中, 使用 \sisi{\hyperref[XVII.prop2.4.3]{2.4.3}}{\hyperref[XVII.prop2.4.1]{2.4.1} (\resp \hyperref[XVII.prop2.4.3]{2.4.3})}, 所以只须考虑那些使 $Rf_{\ast}$ (\resp $Lf^{\ast}$) 的维数被固定的 $N$ 所界的 $f$. 来证明存在“足够多”的复形 $K$, 其各分量的正像 (\resp 逆像) 对于所考虑的正像 (分别地, 逆像) 函子全部无上同调. 设 $K$ 为复形,\pageoriginale $K^{\ast}$ (\resp $K_{\ast}$) 为 $K$ 的一个 Cartan--Eilenberg 消解, 由分量为松弛 (\resp 平坦) 层的复形组成.\nde{只要求对每个 $p$, 增广 $K^p\to (K^*)^{p\bigdot}$ (\resp $(K_*)^{p\bigdot}\to K^p$) 为松弛 (\resp 平坦) 消解.} 若 $n>N$, 可知 $f_{\ast}\tau''_{<n}(K^{\ast})=\tau''_{<\sisi{}{n}}f_{\ast}(K^{\ast})$ (\resp $f^{\ast}\tau''_{>-n}(K_{\ast})=\tau''_{>-n}f^{\ast}(K_{\ast})$), 其中 $\tau\sisi{^{\ast}}{''}$ 表示对于第二个微分的截断. 对 $n>2N$, 由此可知 $\tau''_{<n}(K^{\ast})$ (\resp $\tau''_{>-n}(K_{\ast})$) 满足 \sisi{\hyperref[XVII.prop2.4.3]{2.4.3}}{\hyperref[XVII.prop2.4.1]{2.4.1} (\resp \hyperref[XVII.prop2.4.3]{2.4.3})} 的假设, 并且每个复形 $K$ 都是某个拟同构的源 (\resp 靶), 而该拟同构的靶 (\resp 源) 是这样的复形.

\setcounter{subsubsection}{3}
\subsubsection{}\etiquette[4.1.4]{XVII.subsubsec4.1.4}
考虑赋环景的交换方块 \eqhyperref[XVII.eq4.1.4.1]{4.1.4.1} (\cf~\hyperref[XVII.subsec0.11]{0.11})
$$\etiquette[4.1.4.1]{XVII.eq4.1.4.1}
\vcenter{
\xymatrix{
  X' \ar[r]^{g'}\ar[d]_{f'}& X\ar[d]^{f}\\
S' \ar[r]^{g} & S\rlap{.}
}}\leqno{(4.1.4.1)}
$$
可以考虑由此图表组成的景范畴的子范畴;\nde{应当考虑此图表所定义的赋环景范畴, 一般说来严格地讲它并不是赋环景范畴的子范畴. \Cf \nderef{XVII.ndesous}.} 然后利用 \sisi{\hyperref[XVII.subsubsec2.1.2]{2.1.2}}{\hyperref[XVII.prop2.1.3]{2.1.3}}, 在 \hyperref[XVII.sch4.1.3]{4.1.3} 适用时定义从 $Lg^{\ast}Rf_{\ast}$ 到 $Rf'_{\ast}L{g'}^{\ast}$ 的“基变换态射”:

\begin{enonce*}{命题 4.1.5}\etiquette[4.1.5]{XVII.prop4.1.5}
\sisi{\hyperref[XVII.subsubsec2.1.2]{2.1.2}}{\hyperref[XVII.prop2.1.3]{2.1.3}} 中从 $Lg^{\ast}Rf_{\ast}$ 到 $Rf'_{\ast}L{g'}^{\ast}$ 的基变换态射在以下每种情形下都有定义:
\begin{enumerate}
\renewcommand{\theenumi}{\roman{enumi}}
\renewcommand{\labelenumi}{(\theenumi)}
\item 在范畴 $D^{+}$ 中工作, 且 $Lg^{\ast}$ 与 $L{g'}^{\ast}$ 具有有限维数.

\item 在范畴 $D^{-}$ 中工作, 且 $Rf_{\ast}$ 与 $Rf'_{\ast}$ 具有有限维数.

\item 在无界导出范畴中工作, 且 $Rf_{\ast}$、$Rf'_{\ast}$、$Lg^{\ast}$ 及 $L{g'}^{\ast}$ 具有有限维数, $Lf^{\ast}$ 与 $L{f'}^{\ast}$ (或者: $Rg_{\ast}$ 与 $Rg'_{\ast}$) 也如此.\nde{(iii) 中的有限维数假设是多余的. \Cf \nderef{XVII.ndeKS}.}
\end{enumerate}
\end{enonce*}

应当\pageoriginale 注意, (iii) 中的最后一个假设颇像一个恶作剧. 除了所考虑的景具有足够多点的情形 (\hyperref[XVII.subsubsec4.2.12]{4.2.12}), 编者不知道怎样去掉它. 这一基变换态射的函子性质立即由 \sisi{\hyperref[XVII.subsubsec2.1.2]{2.1.2}}{\hyperref[XVII.prop2.1.3]{2.1.3}} 中的描述得到.

\setcounter{subsubsection}{5}
\subsubsection{}\etiquette[4.1.6]{XVII.subsubsec4.1.6}
设 $(S,\Aa)$ 为赋环景, 考虑“张量积”函子
\[
\otimes:\Mod\sisi{_d}{}(S,\Aa^{\circ})\times
\Mod\sisi{_g}{}(S,\Aa)\longrightarrow \Mod(S,C)
\]%源的第一个分量取右 Aa 模, 或左 Aa-circ 模.
其中 $C$ 为 $\Aa$ 的中心.

\hyperref[XVII.subsubsec1.1.6]{1.1.6} 的方法 (应用此方法时, 与双复形关联的单复形用直和定义) 将此函子延拓为仍记作 $\otimes$ 的函子:
\[
\otimes:K(S,\Aa^{\circ})\times K(S,\Aa)\longrightarrow
K(S,C){\quoi}.
\]

\begin{enonce*}{命题 4.1.7}\etiquette[4.1.7]{XVII.prop4.1.7}
采用上述记号:
\begin{enumerate}
\renewcommand{\theenumi}{\roman{enumi}}
\renewcommand{\labelenumi}{(\theenumi)}
\item 张量积函子可左导出 (\hyperref[XVII.def1.2.1]{1.2.1}) 为“全张量积”函子
\[
{}\overset{L}{\otimes}{}:D^{-}(S,\Aa^{\circ})\times
D^{-}(S,\Aa)\longrightarrow D^{-}(S,C){\quoi}\text{;}
\]

\item 一对上有界复形 $(K,L)$ 只要 $K$ 或 $L$ 的分量平坦, 就适合于 $\overset{L}{\otimes}$.
\end{enumerate}
\end{enonce*}

证明是标准的.

\begin{enonce*}{记号 4.1.8}\etiquette[4.1.8]{XVII.not4.1.8}
若\pageoriginale $K\in D^{-}(S,\Aa^{\circ})$ 且 $L\in D^{-}(S,\Aa)$, 以 $K\overset{L}{\otimes}L$ 表示 $K$ 与 $L$ 的全张量积. 层 $H_{k}(K\overset{L}{\otimes}L)=H^{-k}(K\overset{L}{\otimes}L)$ 记作 $\Tor_k(K,L)$, 称为 $K$ 与 $L$ 的第 $k$ 个局部超 Tor (local hypertor).
\end{enonce*}

设 $K\in D^{-}(S,\Aa^{\circ})$ 及 $L\in D^{-}(S,\Aa)$. 若对 $i>k$ 有 $H^{i}(K)=0$, 对 $i>\ell$ 有 $H^{i}(L)=0$, 可以验证, 对 $i>k+\ell$ 有 $H^{i}(K\overset{L}{\otimes}L)=0$, 并且 $H^{k+\ell}(K\overset{L}{\otimes}L)=H^{k}(K)\otimes H^{\ell}(L)$. 固定 $K$, 将 $K\overset{L}{\otimes}L$ 看作 $L$ 的函子. 若将上同调幅的定义 \hyperref[XVII.def1.2.9]{1.2.9} 应用为 $\ZZ$ 中的区间, 则此区间的上端点具有平凡性质: 它是使 $H^{i}(K)\neq0$ 的最大整数 $i$. 下端点更为有趣:

\begin{enonce*}{定义 4.1.9}\etiquette[4.1.9]{XVII.def4.1.9}
若下述等价条件成立, 就称复形 $K\in\Ob D^{b}(S,\Aa^{\circ})$ 的 \emph{\tTor 维数 $\leq n$}:
\begin{enumerate}
\renewcommand{\theenumi}{\roman{enumi}}
\renewcommand{\labelenumi}{(\theenumi)}
\item 对每个 $\Aa$ 模 $L$, 当 $k>n$ 时有 $\Tor_k(K,L)=0$.

\item 对每个 $\Aa$ 模复形 $L$, 若对 $i<\ell$ 有 $H^i(L)=0$, 则对 $i<\ell-n$ 有 $H^i(K\overset{L}{\otimes}L)=0$.

\item 存在以 $K$ 为靶、以复形 $K'$ 为源的拟同构, 其中该复形\sisi{}{有界}, 各分量平坦, 且对 $i<-n$ 为零.
\end{enumerate}
\end{enonce*}

\setcounter{subsubsection}{9}
\subsubsection{}\etiquette[4.1.10]{XVII.subsubsec4.1.10}
有如下可能的变体:
\begin{enumerate}
\renewcommand{\theenumi}{\roman{enumi}}
\renewcommand{\labelenumi}{(\theenumi)}
\item 若 $K\in\Ob D^{b}(S,\Aa^{\circ})$ 有限 \ttor 维, 且 $L\in\Ob D(S,\Aa)$, 则函子 $\overset{L}{\otimes}$ 在 $(K,L)$ 处有定义 (\hyperref[XVII.def1.2.1]{1.2.1}).

\item 若 $K\in\Ob K^{b}(S,\Aa^{\circ})$ 的各分量具有有限 \tTor 维数, 则函子 $L\mapsto K\otimes L$ $(L\in K(S,\Aa))$ 可左导出.

\item 若每个模都具有有限 \tTor 维数,\nde{此条件是多余的. 在 $\Aa$ 交换的情形下, 见 \nderef{XVII.ndeKS} 中所引 [KS], 18.6.} 则函子 $\otimes$ 导出为函子
\[
{}\overset{L}{\otimes}{}:D(S,\Aa^{\circ})\times
D(S,\Aa)\longrightarrow D(S,\Cc){\quoi}.
\]
\end{enumerate}

\subsection{Godement 松弛消解}\etiquette[4.2]{XVII.subsec4.2}\pageoriginale

\subsubsection{}\etiquette[4.2.1]{XVII.subsubsec4.2.1}
设 $S$ 为拓扑斯, $X$ 为集合. 给定 $S$ 的一个以 $X$ 为指标的点族\sisi{}{}, 等价于给定从离散拓扑空间 $X$ 所定义的拓扑斯到 $S$ 的态射 $k:\Top(X)\rightarrow S$. 此族保守 (IV \hyperref[IV.def6.4.1]{6.4.1}) 的充要条件为: 对 $S$ 上的每个层 $F$, 从 $F$ 到 $k_{\ast}k^{\ast}F$ 的伴随箭头为单态射 (\cf IV \hyperref[IV.subsubsec6.4.0]{6.4.0}).

\begin{enonce*}{定义 4.2.2}\etiquette[4.2.2]{XVII.def4.2.2}
设 $k:\Top(X)\rightarrow S$ 为拓扑斯 $S$ 的保守点族. $S$ 上阿贝尔层 $F$ 的 \emph{Godement 消解}或\emph{典范松弛消解} (相对于 $X$), 记作 $\Cc_X^{\ast}(F)$ (或简记 $\Cc^{\ast}(F)$), 是 $F$ 的如下右消解:
\begin{enumerate}
\renewcommand{\labelenumi}{\theenumi)}
\item $\Cc^{0}(F)=k_{\ast}k^{\ast}F$, 且 $\epsilon:F\rightarrow\Cc^{0}(F)$ 为伴随箭头.

\item 置 $d^{-1}=\epsilon$; 对 $n\geq0$, 递归定义

$\Cc^{n+1}(F)=\Cc^{0}(\coker d^{n-1})$, 并将 $d^n$ 定义为复合箭头
\[
d^{n}:\Cc^n(F)\rightarrow
\coker(d^{n-1})\rightarrow\Cc^{0}(\coker(d^{n-1})){\quoi}.
\]
\end{enumerate}
\end{enonce*}

\begin{enonce*}{命题 4.2.3}\etiquette[4.2.3]{XVII.prop4.2.3}
在 \emph{\hyperref[XVII.def4.2.2]{4.2.2}} 的假设下,
\begin{enumerate}
\renewcommand{\theenumi}{\roman{enumi}}
\renewcommand{\labelenumi}{(\theenumi)}
\item $\Cc^n(F)$ 为松弛层 \emph{\sisi{(V 3.7)}{\hyperrefext[V][V.def4.1]{4.1}}},

\item $\Cc^n(F)$ 是 $F$ 的正合函子,

\item 复形 $\Cc^{\ast}(F)$ 在点 $x\in X$ 处的茎是 $F_x$ 的一个典范分裂的消解.
\end{enumerate}
\end{enonce*}

层\pageoriginale $\Cc^n(F)$ 松弛, 因为它们是离散拓扑空间 $X$ 上的层 (自动松弛) 的正像.

函子 $k^{\ast}$ 与 $k_{\ast}$ 正合. 因而函子 $\Cc^0$ 正合. 对 $n\geq0$ 作归纳, 证明函子 $\mathcal{Z}^{n}(F)=\coker(d^{n-1})$ 正合, 从而 $\Cc^{n+1}=\Cc^0\mathcal{Z}^{n}$ 也正合. 置 $\Zz^{-1}(F)=F$. 层的正合列
\[
0\longrightarrow F'\longrightarrow F\longrightarrow F''\longrightarrow0
\]
对 $n\geq0$ 定义图表
\[
\xymatrix{
        &   0   \ar[d]             &  0   \ar[d]            &  0 \ar[d]\\
0  \ar[r] & \Zz^{n-1}(F\sisi{^1}{'})  \ar[r] \ar[d]& \Zz^{n-1}(F)  \ar[r] \ar[d]& \Zz^{n-1}(F'')
\ar[r] \ar[d]& 0 \\
0  \ar[r] & \Cc^n(F')  \ar[r] \ar[d]& \Cc^n(F)  \ar[r] \ar[d]&
\Cc^n (F'')  \ar[r] \ar[d]& 0 \\
 0  \ar[r] & \Zz^{n}(F')  \ar[r]\ar[d] & \Zz^{n}(F)  \ar[r] \ar[d]& \Zz^{n}(F'')  \ar[r] \ar[d]&
0\\
        &  0               &  0             &  0      &,
}
\]
其各列正合, 并且由归纳假设, 前两行也正合. 因而第三行正合, 这就证明了 (ii).

短正合列
\[
0\longrightarrow\mathcal{Z}^{n+1}(F)\longrightarrow
\Cc^{n}(F)\longrightarrow\mathcal{Z}^{n}(F)\longrightarrow0
\]
是短正合列
$$\etiquette[4.2.3.1]{XVII.eq4.2.3.1}
0\rightarrow F\xrightarrow{\epsilon}\Cc^{0}(F)\rightarrow
\mathcal{Z}^{0}(F)\rightarrow0\leqno(4.2.3.1)
$$
的特殊情形 (作代换\pageoriginale $F\mapsto\mathcal{Z}^{n-1}(F)$). 为证明 (iii), 只须证明 \eqhyperref[XVII.eq4.2.3.1]{4.2.3.1} 在 $k$ 下的逆像序列典范地分裂. 事实上, 复合箭头
\[
k^{\ast}F\xrightarrow{k^{\ast}(\epsilon)}k^{\ast}\Cc^{0}F=k^{\ast}(k_{\ast}k^{\ast})F=(k^{\ast}k_{\ast})k^{\ast}F\longrightarrow
k^{\ast}F
\]
为恒等箭头.

\begin{enonce*}[remark]{注 4.2.4}\etiquette[4.2.4]{XVII.rem4.2.4}
若 $(x_i)_{i\in X}$ 为某个景的保守点族, 则此族对应于伴随拓扑斯的一个保守点族, 因而仍可定义典范松弛消解.
\end{enonce*}

\begin{enonce*}[remark]{注 4.2.5}\etiquette[4.2.5]{XVII.rem4.2.5}
若 $S$ 为赋环拓扑斯, 则函子 $\Cc^n$ 将模映为模.
\end{enonce*}

\begin{enonce*}[remark]{注 4.2.6}\etiquette[4.2.6]{XVII.rem4.2.6}
设 $\Ss$ 为一个由景组成的范畴. 假设对每个 $S\in\Ob\Ss$, 给定一个集合 $S^d$ 及 $S$ 的一个以 $S^d$ 为指标的保守点族 $k_S:\Top(S^d)\rightarrow S$. 假设 $S^d$ 对 $S$ 具有函子性, 且图表
\[
\xymatrix{
\Top(S^d)   \ar[r]^{f^d} \ar[d]^{k_S}& \Top(T^d)\ar[d]^{k_T}\\
S        \ar[r]^{f}  & T
}
\]
交换. 则相应的典范松弛消解“对 $S$ 具有函子性”.
\end{enonce*}

若\pageoriginale $\Ss$ 为拓扑空间范畴, 可以取 $S^d$ 为底集合.

若 $\Ss$ 为在概形 $S_0$ 上局部有限型的各概形 $S$ 的平展景所成的范畴, 且对每个 $s\in S_0$, $\bar{k(s)}$ 为 $k(s)$ 的代数闭包, 则可取 $S^d$ 为 $S$ 在某个 $\bar{k(s)}$ 中取值的点的集合.

若将代数闭包 $\bar{k(s)}$ 换成足够大的代数闭扩张, 则可将“有限型”条件换成一个基数条件.

\setcounter{subsubsection}{6}
\subsubsection{}\etiquette[4.2.7]{XVII.subsubsec4.2.7}
典范松弛消解的一个优点是, 它们使我们可以构造同时适合“张量积”型函子及“正像”型函子的复形. \sisi{这}{它们}因而给出了基变换态射的一种也许更易理解的构造 (\hyperref[XVII.prop4.1.5]{4.1.5}).


\begin{enonce*}{引理 4.2.8}\etiquette[4.2.8]{XVII.lem4.2.8}
设 $f:(S,\Aa)\rightarrow(S',\Bb)$ 为赋环景态射, $K\in\Ob K(S,\Aa^{\circ})$ 与 $L\in\Ob K(S',\Bb)$ 为复形, $k:\Top(X)\rightarrow S$ 为 $S$ 的一个保守点集. 假设下述条件之一成立:
\begin{enumerate}
\renewcommand{\theenumi}{\alph{enumi}}
\renewcommand{\labelenumi}{\theenumi)}
\item $K$ 与 $L$ 的上同调上有界;

\item $\Aa$ 在 $f^{\sisi{\ast}{-1}}\Bb$ 上具有有限 \tTor 维数, 且 $K\in\Ob D^b(S,\Aa^{\circ})$ 具有有限 \tTor 维数;

\item $L\in\Ob D^b(S',\Bb)$ 具有有限 \tTor 维数.
\end{enumerate}

则\pageoriginale 函子 $(K,L)\mapsto K\otimes f^{\ast}L$ 在 $(K,L)$ 处可\sisi{}{左}导出 \emph{(\hyperref[XVII.def1.2.1]{1.2.1})}; 以 ${}\overset{\LL}{\otimes}\LL f^{\ast}$ 表示其导出函子.
\begin{enumerate}
\renewcommand{\theenumi}{\roman{enumi}}
\renewcommand{\labelenumi}{(\theenumi)}
\item 若 $x$ 为 $S$ 的一个点, 则 $(K\overset{\LL}{\otimes}\LL f^{\ast}L)_x\simeq
K_x\overset{\LL}{\otimes}_{\Bb_{f(x)}}L_{f(x)}$.

\item 对 ${}\overset{\LL}{\otimes}f^{\ast}$ \eqhyperref[XVII.def1.2.1]{1.2.1}, 对 $(K,L)$ 左适合, 当且仅当对每个 $x\in X$, 对 $(K_x,L_{f(x)})$ 左适合于函子 $\otimes_{\Bb_{f(x)}}$.

\item 在假设 $(a)$ 或 $(b)$ 下, $L$ 中的函子 $K\overset{\LL}{\otimes}f^{\ast}L$ 的上同调幅为 $\subset d$ \emph{(\hyperref[XVII.def1.2.9]{1.2.9})}, 当且仅当对每个 $x\in X$, $\Bb_{f(x)}$ 模复形 $K_x$ 的 \tTor 幅为 $\subset d$.\nde{假设 (a)、(b)、(c) 全部是多余的. \Cf \nderef{XVII.ndeKS}.}
\end{enumerate}
\end{enonce*}

此引理的证明是标准的; 它表达了所考虑函子的“逐点性质”.

\setcounter{subsubsection}{8}
\subsubsection{}\etiquette[4.2.9]{XVII.subsubsec4.2.9}

设 $L$ 为复形. $L$ 的 \emph{Godement 消解} (相对于 $X$) 是与双复形 $\sisi{C}{\Cc}^{\ast}L$ 关联的单复形 (\emph{用直和定义}) (\hyperref[XVII.subsubsec1.1.4]{1.1.4}). \emph{截断至 $n$ 阶的 Godement 消解} $\tau''_{\leq n}\sisi{C}{\Cc}^{\ast}L$, 是与由 $\sisi{C}{\Cc}^{\ast}L$ 沿 Godement 微分方向截断 (\hyperref[XVII.subsubsec1.1.15]{1.1.15}) 所得双复形关联的单复形. 根据 \hyperref[XVII.prop4.2.3]{4.2.3}~(iii), $L$ 在 $X$ 上逐点同伦等价于复形 $\tau''_{\leq n}\sisi{C}{\Cc}^{\ast}L$ 及 $\sisi{C}{\Cc}^{\ast}L$; 因而对于“张量积”型函子, 这些复形“与 $L$ 一样好” (\hyperref[XVII.lem4.2.8]{4.2.8}).

\begin{enonce*}{命题 4.2.10}\etiquette[4.2.10]{XVII.prop4.2.10}
设 $(S,\Aa)$ 为赋环景, $X$ 为 $S$ 的一个保守点集, $N,M\in\NN\cup\{\infty\}$, 且 $N$ 或 $M$ 有限. 设 $A$ 为 $K(S)$ 中由满足下列条件的复形 $L$ 组成的子范畴:
\begin{enumerate}
\renewcommand{\theenumi}{\alph{enumi}}
\renewcommand{\labelenumi}{\theenumi)}
\item 若 $N=\infty$, 则对充分大的 $i$ 有 $H^i(L)=0$; 若 $M=\infty$, 则对充分小的 $i$ 有 $H^i(L)=0$.

\item 对每个\pageoriginale $x\in X$ 及每个 \tTor 维数 $\leq N$ 的 $\Aa_x$ 模 $Q$, 复形 $L_x$ 左适合于函子 $\sisi{M_x}{Q}\otimes_{\Aa_x}\ast$.

\item 对每个上同调维数 $\leq M$ 的拓扑斯态射 $f:\sisi{X^{\sim}}{S}\rightarrow E$, 复形 $L$ 右适合于函子 $f_{\ast}$.
\end{enumerate}

则 $A$ 到 $K(S)$ 的嵌入, 诱导在 $A$ 中将拟同构取逆所得的范畴与范畴 $D^{\ast}(S)$ 之间的等价, 其中若 $N<M=\infty$ 则 $\ast=+$, 若 $M<N=\infty$ 则 $\ast=-$, 若 $N,M<\infty$ 则 $\ast$ 留空.
\end{enonce*}

设 $K$ 为满足 (a) 的复形. 若 $N=\infty$, 则对充分大的 $i$, 典范箭头 $a:K_1=\sisi{\sigma}{\tau}_{\leq i}K\rightarrow K$ 为拟同构. 若 $N<\infty$, 置 $K_1=K$. 设 $K_1^{\ast}$ 为一个平坦 Cartan--Eilenberg 消解, 并置 $K_2=\tau''_{\geq\sisi{}{-}N}K_1$ (沿新微分方向的截断). 从 $K_2$ 到 $K$ 的复合箭头为拟同构, 且 $K_2$ 满足 (b).

设 $K$ 为满足 (a) 及 (b) 的复形. 若 $M=\infty$, 则 $N<\infty$; 对充分大的 $i$, $b:K\rightarrow\sisi{\sigma}{\tau}_{\geq i}K_1=K_3$ 为拟同构, \sisi{从}{且} $K_3$ 仍满足 (a) 及 (b). 若 $M<\infty$, 置 $K=K_3$. 设 $K_4=\sisi{\sigma}{\tau}''_{\leq M}\sisi{C}{\Cc}^{\ast}K_3$. 则由 \hyperref[XVII.prop4.2.3]{4.2.3}~(iii), $K_4$ 满足 \sisi{(a)、(c) 及 (b)}{(a)、(b) 及 (c)}.

因此, 连续两次应用 \cite[I 2.4.2]{sga4.exp.XVII.bib.CD} 即得到此命题.

\begin{enonce*}[remark]{变体 4.2.11}\etiquette[4.2.11]{XVII.var4.2.11}
若 $N<\infty$ (\resp 若 $M<\infty$, \resp 若 $N,M<\infty$), 同样也可以用下有界 (\resp 上有界, \resp 有界) 复形来工作.
\end{enonce*}

\setcounter{subsubsection}{11}
\subsubsection{}\etiquette[4.2.12]{XVII.subsubsec4.2.12}
当所考虑的景有足够多点时, 可以如下\pageoriginale 利用这些构造, 在比 \loccit{} 更一般的假设下定义 \hyperref[XVII.subsubsec4.1.4]{4.1.4} 的\emph{基变换态射}.

设给定景的交换图表 (\cf (\hyperref[XVII.subsec0.11]{0.11}))
$$\etiquette[4.2.12.1]{XVII.eq4.2.12.1}
\vcenter{
\xymatrix{
X'          \ar[r]^{g'} \ar[d]^{f'}& X\ar[d]^{f}\\
S'          \ar[r]^{g}  & S\rlap{,}
}}\leqno{(4.2.12.1)}
$$
设 $\Aa$ 为 $S$ 上的环层, $\Aa'$ 为 $S'$ 上的环层, $L$ 为 $S'$ 上的 $(\Aa',g^{\ast}\Aa)$ 双模复形. 假设:
\begin{enumerate}
\renewcommand{\theenumi}{\roman{enumi}}
\renewcommand{\labelenumi}{(\theenumi)}
\item 景 $X$ 与 $X'$ 有保守点族 $k:Q\rightarrow X$ 与 $k':Q'\rightarrow X'$, 组成交换图表
$$\etiquette[4.2.12.2]{XVII.eq4.2.12.2}
\vcenter{
\sisi{\xymatrix{
X'         \ar[r]^{g'}  \ar[d]^{k'} & X     \ar[d]^{k}  & \\
Q'         \ar[r]^{g'_0} & Q &.
}}{\xymatrix{
Q'         \ar[r]^{g'_0}  \ar[d]^{k'} & Q     \ar[d]^{k}  & \\
X'         \ar[r]^{g'} & X \rlap{.}
}}}\leqno{(4.2.12.2)}
$$

\item 下述假设之一成立:
\begin{enumerate}
\renewcommand{\theenumi}{\alph{enumi}}
\renewcommand{\labelenumi}{\theenumi)}
\item 函子 $Rf_{\ast}$ 与 $Rf'_{\ast}$ 具有有限上同调维数, 且 $L$ 上有界. 置 $s=-$.

\item $L$ 下有界, 各分量的 \tTor 维数一致有界 (作为 $g^{\ast}\Aa$ 模). 置 $s=+$.

\item $L$ 有界, 且 (a) 与 (b) 的假设成立. 置 $s=$ 为空.
\end{enumerate}
\end{enumerate}

以\pageoriginale $L\overset{\LL}{\otimes}\LL g^{\ast}:D^{s}(S,\Aa)\rightarrow
D^{s}(S', \Aa')$ 表示从 $K^{s}(S,\Aa)$ 到 $K^{s}(S',\Aa')$ 的函子 $K\mapsto L\otimes_{g^{\ast}\Aa}g^{\ast}K$ 的左导出函子 (\hyperref[XVII.def1.2.1]{1.2.1}).

将对 $D^{s}(X,\Aa)$ 中的 $K$ 定义基变换态射
$$\etiquette[4.2.12.3]{XVII.eq4.2.12.3}
\varphi:L\overset{\LL}{\otimes}\LL g^{\ast}Rf_{\ast}K\longrightarrow
Rf'_{\ast}({f'}^{\ast}L\overset{\LL}{\otimes}\LL{g'}^{\ast}K)\leqno(4.2.12.3)
$$
\nde{\MakeLinkTarget*{sga4.exp.XVII.editor-note.4.2.12.3.unbounded-adjunction}\label{XVII.ndecb}无界导出范畴中的伴随使我们能够不加假设 (i) 与 (ii) 而定义这样的态射. \Cf \nderef{XVII.ndeKS}.}

由 \hyperref[XVII.prop4.2.10]{4.2.10}, 只须对属于 \hyperref[XVII.prop4.2.10]{4.2.10} 中范畴 $A$ 的复形, 对 $K$ 具有函子性地定义它, 其中
\begin{align*}
N & \geq\sup_i \text{ ($L^{i}$ 在 $g^{\ast}$ 上的 tor\sisi{}{-}维数)}\\
M & \geq\text{$Rf_{\ast}$ 的上同调维数}.
\end{align*}
这样的复形同时适合函子 $Rf_{\ast}$ 及函子 ${f'}^{\ast}L\overset{\LL}{\otimes}_{\Aa'}\sisi{}{\LL}{g'}^{\ast}$ (\hyperref[XVII.lem4.2.8]{4.2.8}). 因而典范态射
\begin{align*}
\alpha &: f_{\ast}K\longrightarrow Rf_{\ast}K\\
\beta &:
{f'}^{\ast}L\overset{\LL}{\otimes}\LL \sisi{g}{{g'}}^{\ast}K\longrightarrow
\sisi{}{{f'}^*}L\otimes \sisi{g}{{g'}}^{\ast}K
\end{align*}
为拟同构. 所以为计算 \eqhyperref[XVII.eq4.2.12.3]{4.2.12.3} 的左端 $I$, 只须取左消解 (即拟同构)
\[
\gamma:K_{1}\longrightarrow f_{\ast}K{\quoi},
\]
其中 $K_{1}$ 适合于 $L\overset{\LL}{\otimes}\LL g^{\ast}$; 有
\[
I\simeq L\otimes \sisi{\LL}{} g^{\ast}K_{1}{\quoi}.
\]

同样,\pageoriginale 为计算 \eqhyperref[XVII.eq4.2.12.3]{4.2.12.3} 的右端 $II$, 只须取右消解
\[
\delta:{f'}^{\ast}L\otimes {g'}^{\ast}K\longrightarrow K_2{\quoi},
\]
其中 $K_{2}$ 适合于 $Rf'_{\ast}$; 则
\[
II\simeq f'_{\ast}K_{2}{\quoi}.
\]
于是将基变换态射定义为复合
\[
I\simeq L\otimes g^{\ast}K_{1}\xrightarrow{\gamma}L\otimes
g^{\ast}f_{\ast}K\xrightarrow{c}f'_{\ast}({f'}^{\ast}L\otimes
{g'}^{\ast}K)\xrightarrow{\sisi{\alpha}{\delta}}f'_{\ast}K_{2}\simeq II{\quoi},
\]
其中 $c$ 为实际复形层次上的通常基变换态射.

此基变换态射有下述刻画性质, 原则上使其相容性性质易于验证.

\begin{enonce*}{命题 4.2.13}\etiquette[4.2.13]{XVII.prop4.2.13}
在上述假设下, 设 $K\in \Ob
C^{s}(X\sisi{'}{},\Aa)$、$K_{1}\in \Ob C^{s}(\sisi{X}{S},\Aa)$ 及 $K_{2}\in \Ob C^{s}(\sisi{S}{X'},\Aa)$. 此外, 设 $u$ 为从 $K$ 到 $K_{1}$ 的 $f$ 态射, 即 $u\in
\Hom(K_{1},f_{\ast}K)\simeq \Hom(f^{\ast}K_{1},K)$, $v$ 为从 ${f'}^{\ast}L\otimes {g'}^{\ast}K$ 到 $K_{2}$ 的态射. 由 $u$、$v$ 及通常的基变换态射, 得到复形态射 $c_{u,v}:L\otimes g^{\ast}K_{1}\rightarrow
f'_{\ast}\sisi{k}{K}_2$. 以 $u'$、$v'$ 及 $c'_{uv}$ 表示复合态射
\begin{align*}
u' &: K_{1}\xrightarrow{u} f_{\ast}K\rightarrow Rf_{\ast}K\\
v' &: L\overset{\LL}{\otimes}\LL{g'}^{\ast}K\rightarrow
L\otimes {g'}^{\ast}K\xrightarrow{v} K_{2}\\
c'_{u,v} &:
L\overset{\LL}{\otimes}\LL g^{\ast}K_{1}\rightarrow
L\otimes g^{\ast}K_{1}
\xrightarrow{c_{uv}}f'_{\ast}K_{2}\rightarrow Rf'_{\ast}K_{2}{\quoi}.
\end{align*}
下图\pageoriginale 交换:
$$\etiquette[4.2.13.1]{XVII.eq4.2.13.1}
\vcenter{
\xymatrix{
L \overset{\LL}{\otimes} \sisi{L}{\LL}g^{\ast}Rf_{\ast}K
\ar[r]^{\varphi} & Rf'_{\ast}(f'^{\ast}L
\overset{\LL}{\otimes}\sisi{L}{\LL} g'^{\ast}K) \ar[d]^{Rf'_{\ast}(v')}&\\
L \overset{\LL}{\otimes} \sisi{L}{\LL}g^{\ast}K_1 \ar[u]_{L \overset{\LL}{\otimes} \sisi{L}{\LL}g^{\ast}(u')}
\ar[r]^{c'_{uv}} & Rf'_{\ast}K_2.
}}\leqno{(4.2.13.1)}
$$
\end{enonce*}

先假设 $K_{1}=f_{\ast}K$ 且 $K_{2}={f'}^{\ast}L\otimes {g'}^{\ast}K$, 考虑下述复形交换图表:
\[
\xymatrix{
K^{a}  \ar[r] \ar[d]& K\ar[d]\\
K^{c}  \ar[r] & K^{b}
}
\]
其中 $K^{a}$ 适合于 $\sisi{L}{}{f'}^{\ast}L\overset{\LL}{\otimes}\sisi{L}{\LL}{g'}^{\ast}$, $K^{b}$ 适合于 $Rf_{\ast}$, $K^{c}$ 同时适合于这两个函子. 为构造此图表, 先构造 $K^{a}$, 再对 $K^{a}$ 与 $K$ 应用典范松弛消解 (必要时作截断).

对此图表应用函子 $f_{\ast}$ (\resp~$L\otimes {g'}^{\ast}$), 并对图表 $(K_{1},K^{a}_{1},K^{b}_{1}, K^{c}_{1})$ (\resp $(K_{2}, K^{a}_{2},
K^{b}_{2}, K^{c}_{2})$) 中的复形取消解, 得到一个交换立方体, 其中 $\approx$ 表示拟同构:
\[
\xymatrix@C=2em{
          &         &K'_1 \ar[rdd] \ar[rrrrrr]^{\approx}&        && &         &
          & K_1\ar[rdd]\\  K_1^{a'}\ar[rru]\ar[rdd]\ar[rrrrrr]^(.75){\approx} &         &      &        && & K_1^{a}\ar[rru] \ar[rdd]&       &\\
  &         &      &K^{b'}_1 \ar[rrrrrr] ^(.25){\approx}&& &         &         &   & K^{b}_1\\
          &K_1^{c'} \ar[rru]^{\approx}\ar[rrrrrr] ^(.65){\approx}&      &        && &         &K_1^{c}\ar[rru]^{\approx} \\
}
\]
其中\pageoriginale $K'_{1}$、$K^{a'}_{1}$、$K^{b'}_{1}$ 及 $K^{c'}_{1}$ 适合于函子 $L\overset{\LL}{\otimes}\sisi{L}{\LL}g^{\ast}$ (\resp \[
\xymatrix@C=2em{
        &         &K_2 \ar[rdd] \ar[rrrrrr]^{\approx}&        && &         &
          & K'_2\ar[rdd]\\
K_2^{a}\ar[rru]\ar[rdd]^{\approx}\ar[rrrrrr]^(.75){\approx} &         &      &        && & K_2^{a'}\ar[rru] \ar[rdd]^(.65){\approx}&       &\\
  &         &      &K^{b}_2 \ar[rrrrrr] ^(.25){\approx}&& &         &         &   & K^{b'}_2\\
          &K_2^{c} \ar[rru]\ar[rrrrrr] ^(.65){\approx}&      &
        && &         &K_2^{c'}\ar[rru] & &, \\
}
\], 其中 $K'_{2}$、$K^{a'}_{2}$、$K^{b'}_{2}$ 及 $K^{c'}_{2}$ 适合于函子 $Rf'_{\ast}$). 利用通常的基变换态射, 得到 $S$ 上复形的一个交换棱柱图.
%%%%%%
\newcommand{\circlenum}[1]{\bigcirc \kern -.255cm \lower
  .05em\hbox{${#1}$}}

\[
\def\objectstyle{\scriptstyle}
\xymatrix@C=-.34cm{
        &         & \scriptstyle{L\otimes g^{\ast}K'_1}\ar[rrrrrr]_*+[o][F-]{\scriptscriptstyle{4}} \ar[rdd]^(.35)*+[o][F-]{\scriptscriptstyle{2}}&        && &         &
          & \scriptstyle{L\otimes g^{\ast}K_1} \ar[rdd] \ar[rrrrr]_*+[o][F-]{\scriptscriptstyle{4}}& & & & & \scriptstyle{f'_{\ast}K_2} \ar[rrrrr]_*+[o][F-]{\scriptscriptstyle{4}}\ar[rdd] & & & && \scriptscriptstyle{f'_{\ast}K'_2} \ar[rdd]\\
\scriptstyle{L\otimes g^{\ast}K_1^{a'}} \ar[rru]\ar[rdd]\ar[rrrrrr] &         &
        &        && & \scriptstyle{L\otimes g^{\ast}K_1\sisi{}{^a}} \ar[rru] \ar[rdd]\ar[rrrrr]&       &
        &&& \scriptstyle{f'_{\ast}K^{a}_2}\ar[rdd]\ar[rru] \ar[rrrrr]&
        & & & &
        \scriptstyle{f'_{\ast}K^{a'}_2}\ar[rru]_*+[o][F-]{\scriptscriptstyle{3}}\ar[rdd]^(.7)*+[o][F-]{\scriptscriptstyle{3}}^{\quad \raisebox{1cm}{}{\scriptscriptstyle{-1}}}_(.7){\text{\rotatebox{110}{$\approx$}}}\\
  &         &      & \scriptstyle{L\otimes g^{\ast}K^{\sisi{c}{b}'}_1}  \ar[rrrrrr]
        && &         &         &   & \scriptstyle{L\otimes g^{\ast}K^{b}_2}
        \ar[rrrrr]& & & & & \scriptstyle{f'_{\ast}K^{b}_2} \ar[rrrrr]& & & & & \scriptstyle{f'_{\ast}K^{b'}_2}\\
          &\scriptstyle{L\otimes g^{\ast}K^{c'}_1}
 \ar[rrrrrr]_*+[o][F-]{\scriptscriptstyle{1}}\ar[rru]_{\text{\rotatebox{50}{$\approx$}}}^*+[o][F-]{\scriptscriptstyle{2}}^{\qquad\qquad \raisebox{-.3cm}{}{\scriptscriptstyle{-1}}}
&      &        && &         &\scriptstyle{L\otimes g^{\ast}K^{c}_2} \ar[rru] \ar[rrrrr]_*+[o][F-]{\scriptscriptstyle{1}}& & & & &\scriptstyle{f'_{\ast}K^{c}_2}\ar[rru]\ar[rrrrr]_*+[o][F-]{\scriptscriptstyle{1}}&
        & & & & \scriptstyle{f'_{\ast}K^{c'}_{2'}}\ar[rru]\\
}
\]
注意 \eqhyperref[XVII.eq4.2.13.1]{4.2.13.1} 图表中的态射等同于下列复合,\pageoriginale 即可得出结论:

对于 $\varphi$: 各态射 \ding{192} 的复合;

对于 $L\overset{\LL}{\otimes}\sisi{L}{\LL}g^{\ast}(u')$: 各态射 \ding{193} 的复合, 其中一个态射是某拟同构的逆;

对于 $Rf'_{\ast}(v')$: 各态射 \ding{194} 的复合, 其中一个态射是某拟同构的逆;

对于 $\sisi{C}{c}'_{uv}$: 各态射 \ding{195} 的复合.

在一般情形下, 置 $K'_{1}=f_{\ast}K$ 及 $K'_{2}={f'}^{\ast}L\otimes {g'}^{\ast}K$. 于是有图表
\[
\xymatrix{
&L \overset{\LL}{\otimes} \sisi{L}{\LL}g^{\ast}Rf_{\ast}K
\ar[r]^\varphi&
Rf'_{\ast}(f'^{\ast}L \overset{\LL}{\otimes}\sisi{L}{\LL}
g'^{\ast}K) \ar[d] \ar[rd]& \\
L \overset{\LL}{\otimes} \sisi{L}{\LL}g^{\ast}K_1 \ar[r]\ar[ru]& L
\overset{\LL}{\otimes} \sisi{L}{\LL}g^{\ast}K'_1 \ar[r] \ar[u]&
Rf'_{\ast}K'_2 \ar[r]&  Rf'_{\ast}K_2,
}
\]
其中各三角形交换, 内部的矩形根据上述论证也交换. 因而外边界交换, 这就证明了 \hyperref[XVII.prop4.2.13]{4.2.13}.

\subsection{基变换定理}\etiquette[4.3]{XVII.subsec4.3}\pageoriginale

将用导出范畴的语言重新表述固有态射的基变换定理 \hyperrefext[XII][XII.thm5.1]{5.1} (iii).

\begin{enonce*}{定理 4.3.1}\etiquette[4.3.1]{XVII.thm4.3.1}
设
\[
\xymatrix{
X' \ar[rr]^{g'} \ar[d]^{f'}&& X\ar[d]^{f}\\
S' \ar[rr]^{g}& & S
}
\]
为概形的笛卡尔图表, 其中 $f$ 固有. 设 $\Aa$ 与 $\Aa'$ 为 $S$ 及 $S'$ 上的挠环层, $L$ 为上有界的 $(\Aa',g^{\ast}\Aa)$ 双模复形. 假设 $f$ 的各纤维维数有界. 则 \hyperref[XVII.subsubsec4.2.12]{4.2.12} 的假设 {\rm(a)} 满足, 从 $D^{-}(X,f^{\ast}\Aa)$ 到 $D^{-}(S',\Aa')$ 的函子之间的基变换态射 \eqhyperref[XVII.eq4.2.12.3]{4.2.12.3}
\[
\varphi_K:L\overset{\LL}{\otimes}\sisi{L}{\LL}g^{\ast}\sisi{}{Rf_*}K\longrightarrow
Rf'_{\ast}({f'}^{\ast}L\overset{\LL}{\otimes}\sisi{L}{\LL}{g'}^{\ast}K){\quoi},
\]
为同构.
\end{enonce*}

所考虑的函子是三角的且远端有界 (way-out) \cite[I 7]{sga4.exp.XVII.bib.RD}, 所以为验证 $\varphi_{K}$ 为同构, 只须验证各 $\varphi_{ℋ^{i}(K)}$ 为同构. 因而可假设 $K$ 为集中在次数 $0$ 的挠 $f^{\ast}\Aa$ 模层.

为\pageoriginale 验证 $\varphi_{K}$ 为同构, 只须在应用从 $\Aa'$ 模范畴到 $\ZZ/n\ZZ$ 模范畴的遗忘函子后验证之, 其中 $n\Aa=n\Aa'=0$.

于是归结为 $\Aa'=\ZZ/n\ZZ$ 且 $n\Aa=0$ 的情形. 此时可以将 $L$ 换成导出范畴中与之同构的一个上有界复形, 其各分量在 $g^{\ast}\Aa$ 上平坦. 由于所考虑的函子对 $L$ 是三角的且远端有界 (way-out), 只须在将 $L$ 换成其某个分量后验证 $\varphi_{K}$ 为同构. 因而可以假设 $L$ 为置于次数 $0$ 的平坦 $g^{\ast}\Aa$ 模层. 须证明态射
\[
\varphi^{q}:L\otimes g^{\ast}R^{q}f_{\ast}K\longrightarrow
R^{q}f'_{\ast}({f'}^{\ast}L\otimes g^{\ast}K)
\]
为同构. 设 $s$ 为 $S'$ 的几何点, $X_s$ 为 $X'$ 在 $s$ 处的纤维, $K_s$ 为 $K$ 在 $X_s$ 上的逆像. 根据 \hyperrefext[XII][XII.thm5.1]{5.1} (iii), $\varphi^{q}$ 在 $s$ 处的茎为态射
\[
\varphi^{q}_s:L_{s}\otimes H^{q}(X_s,K)\longrightarrow H^{q}(X_s,
L_s\otimes K){\quoi}.
\]
根据 $D$. Lazard 的定理 \cite{sga4.exp.XVII.bib.La}, 平坦 $\Aa_s$ 模 $L_s$ 是自由 $\Aa_s$ 模 $L_{s,i}$ 的滤过正向极限; 函子 $H^{q}(X_s,\ast)$ 与滤过正向极限可交换\sisi{}{ (\hyperrefext[VII][VII.prop3.3]{3.3})}. 因而态射 $\varphi^{q}$ 为态射
\[
\varphi^{q}_i:L_{s,i}\otimes H^{q}(X_s,K)\longrightarrow H^{q}(X_s,
L_{s,i}\otimes K){\quoi}.
\]
的正向极限. 后者显然是同构; 因而 $\varphi^{q}_s$ 也为同构, $\varphi\sisi{}{^q}$ 为同构.

\begin{enonce*}[remark]{变体 4.3.2}\etiquette[4.3.2]{XVII.var4.3.2}
\begin{enumerate}
\renewcommand{\theenumi}{\roman{enumi}}
\renewcommand{\labelenumi}{(\theenumi)}
\item 当\pageoriginale $L$ 为有界复形, 且其各分量在 $g^{\ast}\Aa$ 上具有有限 \tTor 维数时, \hyperref[XVII.subsubsec4.2.12]{4.2.12} 的假设 (c) 满足, 可以在无界导出范畴中工作.\nde{不加这些假设也可以在无界导出范畴中工作. \Cf \nderef{XVII.ndeKS}.}

\item 若 $L$ 下有界, 且其各分量在 $g^{\ast}\Aa$ 上的 \tTor 维数\sisi{有限}{一致有界}, 则即使 $\Aa$ 与 $\Aa'$ 不是挠的, \hyperref[XVII.subsubsec4.2.12]{4.2.12} 的假设 (b) 仍满足; 若层 $ℋ^{i}(K)$ 为挠层, 基变换态射 $\varphi_{K}$ $(K\in
D^{+}(X,f^{\ast}\Aa)$) 将为同构.
\end{enumerate}
\end{enonce*}

\section{固有支集正像函子}\etiquette[5]{XVII.sec5}

\subsection{基本构造}\etiquette[5.1]{XVII.subsec5.1}

将遵循 \hyperref[XVII.subsec3.1]{3.1} 中所列的纲领, 对可紧化态射 $f$ 定义函子 $Rf_{!}$ (\hyperref[XVII.subsec3.2]{3.2}). 附录中将说明如何处理任意有限型分离态射的情形.\nde{\Cf \nderef{XVII.ndeConrad}.}

\subsubsection{}\etiquette[5.1.1]{XVII.subsubsec5.1.1}

设 $X$ 为拓扑斯, $U$ 为 $X$ 的开子拓扑斯 (\hyperrefext[IV][IV.def8.3]{8.3}), $j$ 为从 $U$ 到 $X$ 的嵌入态射. 在 \hyperrefext[IV][IV.subsec5.2]{5.2} 中, 定义了从 $U$ 上的集合层范畴到 $X$ 上的集合层范畴的一个“\emph{空集延拓}”函子. 此函子 $j_{!}$ 左伴随于限制函子 $j^{\ast}$.

仍以 $j_{!}$ 表示从 $U$ 上的带基点集合层范畴 (\resp 群层范畴, \resp $j^{\ast}\Aa$ 模层范畴, 其中 $\Aa$ 为 $X$ 上的环层) 到\pageoriginale $X$ 上的类似范畴的函子, 它左伴随于限制函子 $j^{\ast}$ (\sisi{IV\quad}{\hyperrefext[IV][IV.prop11.3.1]{11.3.1}}). 这三个函子 $j_{!}$ 与从模范畴到群范畴、从群范畴到带基点集合范畴的遗忘函子相容\footnote[1]{但与到集合范畴的遗忘函子不相容!
}. 若 $i:F\rightarrow X$ 为 $U$ 的闭补的嵌入, 则 $j^{\ast}j_{!}$ (典范同构于) 恒等函子, $i^{\ast}j_{!}$ 为以始对象 (或终对象, 这里两者相同) 为值的常函子. 这也刻画了 $j_{!}$. 因而所考虑的三个函子 $j_{!}$ 正合, 并嵌入相应的函子 $j_{\ast}$. 称其为\emph{零延拓}函子. 公式 $(j_1j_2)^{\ast}=j^{\ast}_2 j^{\ast}_1$ 转置为传递性同构
\[
(j_1 j_2)_{!}=j_{1!}j_{2!}{\quoi}.
\]
函子 $j_{!}$ 与基变换相容:

\begin{enonce*}{引理 5.1.2}\etiquette[5.1.2]{XVII.lem5.1.2}
设 $f:X\rightarrow Y$ 为拓扑斯态射, $V$ 为 $Y$ 的开子拓扑斯, $U$ 为它在 $X$ 中的逆像:
\[
\xymatrix{
U \ar[rr]^{f'}  \ar[d]^{j'}& & V\ar[d]^{j}\\
X \ar[rr]^{f}   && Y
}
\]
则存在唯一的基变换同构 $f^{\ast}j_{!}\isomto j'_{!}{f'}^{\ast}$, 使图表
$$\etiquette[5.1.2.1]{XVII.eq5.1.2.1}
\vcenter{
\xymatrix{
f^{\ast}j_{!}    \ar[r]^{\sim} \ar@{^{(}->}[d] & j'_{!}f'^{\ast} \ar@{^{(}->}[d]     \\
f^{\ast}j_{\ast}  \ar[r]^{\alpha}&j'_{\ast}f'^{\ast},
}}\leqno{(5.1.2.1)}
$$
交换, 其中 $\alpha$ 为基变换态射.
\end{enonce*}

\begin{enonce*}[remark]{证明}
基变换态射\pageoriginale $\varphi:j^{\ast}f_{\ast}\rightarrow f'_{\ast}{j'}^{\ast}$ 显然为同构. 它转置为基变换同构 $\varphi':j'_{!}{f'}^{\ast}\isomto
f^{\ast}j_{!}$. 为验证图表 \eqhyperref[XVII.eq5.1.2.1]{5.1.2.1} 交换, 只须在限制到 $U$ 后验证之, 而这显然成立.
\end{enonce*}

\setcounter{subsubsection}{2}
\subsubsection{}\etiquette[5.1.3]{XVII.subsubsec5.1.3}
在 \hyperref[XVII.subsubsec5.1.1]{5.1.1} 的假设下, 以环层 $\Aa$ 赋环于 $X$, 模上的函子 $j_{!}$ 正合, 因而平凡地下降到导出范畴, 定义一个仍记作 $j_{!}$ 的从 $D(U,j^{\ast}\Aa)$ 到 $D(X,\Aa)$ 的函子. 仍有传递性同构
$$\etiquette[5.1.3.1]{XVII.eq5.1.3.1}
(j_1 j_2)\sisi{}{!}=j_{1!}j_{2!}{\quoi}.\leqno(5.1.3.1)
$$

在导出范畴层次上, 仍有伴随公式
$$\etiquette[5.1.3.2]{XVII.eq5.1.3.2}
\Hom_{D(X)}(j_{!}K,L)\isomto \Hom_{D(U)}
(K,j^{\ast}L){\quoi}.\leqno(5.1.3.2)
$$

\subsubsection{}\etiquette[5.1.4]{XVII.subsubsec5.1.4}
设 $f:X\rightarrow Y$ 为固有态射, $\Aa$ 为 $Y$ 上的\emph{挠}环层, 以 $\Aa$ 的逆像环层赋环于 $X$. 假设 $f$ 的各纤维维数以上界整数 $n$ 为界 (当 $Y$ 拟紧时即如此). 则已知对 $X$ 上的每个 $f^{\ast}\Aa$ 模层 $F$, 当 $q>2n$ 时层 $R^{q}f_{\ast}F$ 为零 (\hyperrefext[XII][XII.corbis5.3]{5.3 bis}). 因而函子 $f_{\ast}$ 有导出函子 (\hyperref[XVII.prop1.2.10]{1.2.10})
\[
Rf_{\ast}:D(X,f^{\ast}\Aa)\longrightarrow
D(Y,\Aa){\quoi}.
\]

若\pageoriginale $f$ 为两个固有态射的复合 $gh$, 则有传递性公式 $Rf_{\ast}=Rg_{\ast}Rh_{\ast}$, 和通常一样, 它可用余纤维化范畴的语言表述 (附论 \sisi{\hyperref[XVII.subsubsec5.1.3]{5.1.3}}{\hyperref[XVII.sch4.1.3]{4.1.3}}).


\subsubsection{}\etiquette[5.1.5]{XVII.subsubsec5.1.5}
考虑概形的交换图表
$$\etiquette[5.1.5.1]{XVII.eq5.1.5.1}
\vcenter{
\xymatrix{
U        \ar@{^{(}->}[r]_{j_1} \ar[d]^{g}      & X  \ar[d]^{f}\\
V        \ar@{^{(}->}[r]_{j_2}       & Y,
}}\leqno{(5.1.5.1)}
$$
其中 $j_1$ 与 $j_2$ 为开浸入, $f$、$g$ 为固有态射, 其纤维维数以上界整数 $n$ 为界. 设 $\Aa$ 为 $Y$ 上的挠环层, 给 $U$、$V$ 及 $X$ 赋予 $\Aa$ 的适当逆像. 由 \eqhyperref[XVII.eq5.1.3.2]{5.1.3.2}, 若 $K\in D(U)$, 则定义态射
$$\etiquette[5.1.5.2]{XVII.eq5.1.5.2}
d:j_{2!}Rg_{\ast}K\longrightarrow Rf_{\ast}j_{1!}K\leqno{(5.1.5.2)}
$$
等价于定义态射
\[
d':Rg_{\ast}K\longrightarrow j^{\ast}_2 Rf_{\ast}j_{1!}K{\quoi}.
\]

纤维积 $X'=V\times {}_YX$ 组成交换图表
$$\etiquette[5.1.5.3]{XVII.eq5.1.5.3}
\vcenter{
\xymatrix{
U \ar@{^{(}->}[r]^{k} \ar[d]^{g} & X'  \ar@{^{(}->}[r]_{j'_2} \ar[d]^{f'}& X \ar[d]^{f}          &\\
V         \ar@{=}[r] & V         \ar@{^{(}->}[r]_{j_2}  & Y   \rlap{,}
}}\leqno{(5.1.5.3)}
$$
并且,\pageoriginale 显然有 $j^{\ast}_2Rf_{\ast}j_{1!}K\simeq
Rf'_{\ast}(j'_2{}^\ast j_{1!}K)\simeq Rf'_{\ast}k_{!}K$.

态射 $k$ 为固有的开浸入; 因此 $k_{!}=k_{\ast}$, 得到同构 $Rf'_{\ast}k_{!}K\simeq
Rf'_{\ast}k_{\ast}K\simeq Rg_{\ast}K$, 从而给出所求同构 $d'$ 及态射 $d$. 此外, $d$ 在 $V$ 上的限制为同构.

\begin{enonce*}{引理 5.1.6}\etiquette[5.1.6]{XVII.lem5.1.6}
态射 \eqhyperref[XVII.eq5.1.5.2]{5.1.5.2} 为同构.
\end{enonce*}

只须验证 $d$ 在 $Y-V=F$ 上的限制为同构. 设 $X''=\fprod{X}{F}{Y}$, 将固有基变换定理 (\hyperref[XVII.thm4.3.1]{4.3.1}) 应用于笛卡尔方块
\[
\xymatrix{
X'' \ar[r]_{i'} \ar[d]^{f''}& X\ar[d]^{f\sisi{'}{}}\\
F   \ar[r]_{i}& Y\rlap{,}
}
\]
其中取 $\Aa'=i^{\ast}\Aa$ 及 $L=\Aa'[0]$. 得到 $i^{\ast}Rf_{\ast}j_{1!}K\simeq Rf''_{\ast}{i'}^{\ast}j_{1!}K\simeq
Rf''_{\ast}0\simeq 0$, 因而 \eqhyperref[XVII.eq5.1.5.2]{5.1.5.2} 的两端在 $V$ 之外都为零, 且 $d$ 为同构.

\setcounter{subsubsection}{6}
\subsubsection{}\etiquette[5.1.7]{XVII.subsubsec5.1.7}
设 $S$ 为以挠环层 $\Aa$ 赋环的概形. 对每个 $S$ 概形 $\varphi:X\rightarrow S$, 置 $\Aa_X=\varphi^{\ast}\Aa$. 在由拟紧、拟分离的 $S$ 概形及 $S$ 可紧化态射组成的范畴 $(S)$ 中工作 (\hyperref[XVII.subsec3.2]{3.2}). 由 \hyperref[XVII.prop3.2.3]{3.2.3}, 此范畴、开浸入子范畴 $(S,i)$ 及固有态射子范畴 $(S,p)$ 满足 \hyperref[XVII.subsubsec3.2.4]{3.2.4} 的假设.

拟对下述数据\pageoriginale 应用 \hyperref[XVII.prop3.3.2]{3.3.2} 的理论 (采用 \hyperref[XVII.subsubsec3.3.1]{3.3.1} 的记号)

\begin{itemize}
\item[$\alpha$)] 对 $S$ 中的 $\varphi:X\rightarrow
S$ 取 $F(X)=D(X,\Aa_X)$:
\begin{itemize}
\item[(i)] 对固有的 $p:X\rightarrow Y$, 取函子 $p_{\ast}$ 为函子 $Rp_{\ast}$.

\item[(ii)] 使用通常的传递性同构 $R(pq)_{\ast}=Rp_{\ast}Rq_{\ast}$.

\item[${\rm(i')}$] 对开浸入 $i:X\rightarrow Y$, 取函子 $i_{\ast}$ 为函子 $i_{!}$.

\item[${\rm(ii')}$] 使用传递性同构 \eqhyperref[XVII.eq5.1.3.1]{5.1.3.1}.

\item[${\rm(iii)}$] 对每个交换图表 \eqhyperref[XVII.eq5.1.5.1]{5.1.5.1}, 取同构 \eqhyperref[XVII.eq5.1.5.2]{5.1.5.2} 为 $d$.
\end{itemize}
\end{itemize}

\hyperref[XVII.prop3.3.2]{3.3.2} 的条件 (a)、(b)、(c)、${\rm(a')}$、${\rm(b')}$、${\rm(c')}$ 均经验证成立, 于是证明了下述定理的 (a) 部分:

\begin{enonce*}{定理 5.1.8}\etiquette[5.1.8]{XVII.thm5.1.8}
\begin{enumerate}
\renewcommand{\theenumi}{\alph{enumi}}
\renewcommand{\labelenumi}{\theenumi)}
\item 在 \hyperref[XVII.subsubsec5.1.7]{5.1.7} 定义的范畴 $(S)$ 中, 存在且本质上唯一一种方式, 对每个箭头 $f:X\rightarrow Y$ 定义函子
\[
Rf_{!}:D(X,\Aa_X)\longrightarrow D(Y,\Aa_Y){\quoi},
\]
并对每个复合态射 $f=gh:X\xrightarrow{h}Y\xrightarrow{g}Z$ 定义 $Rf_{!}$ 与 $Rg_{!}$ $Rh_{!}$ 之间的传递性同构 $c_{\sisi{f,g}{g,h}}$, 使得
\begin{enumerate}
\renewcommand{\theenumii}{\roman{enumii}}
\renewcommand{\labelenumii}{(\theenumii)}
\item 各 $c_{f,g}$ 满足“上闭链条件”
\[
(c_{f,g}\ast Rh_{!})\circ c_{fg,h}=(Rf_{!}\ast c_{g,h})\circ
c_{f,gh}{\quoi}.
\]

\item 限制于固有态射 (\resp 开浸入) 时, 得到 \hyperref[XVII.subsubsec5.1.4]{5.1.4} (\resp~\hyperref[XVII.subsubsec5.1.3]{5.1.3}) 的变性理论.

\item 对\pageoriginale 每个图表 \eqhyperref[XVII.eq5.1.5.1]{5.1.5.1}, \eqhyperref[XVII.eq5.1.5.2]{5.1.5.2} 中的同构 $d$ 为复合同构 $c_{f,j\sisi{'}{_1}} \circ c^{-1}_{j\sisi{,f'}{_2,g}}$.
\end{enumerate}

\item 函子 $Rf_{!}$ 远端有界 \emph{\cite[I 7]{sga4.exp.XVII.bib.RD}}, 且为三角函子.
\end{enumerate}
\end{enonce*}

每个函子 $Rf_{!}$ 为函子 $Rg_{\ast}$ 与函子 $j_{!}$ 的复合, 其中 $g$ 固有且其靶拟紧, 因而纤维维数有界, $j$ 为开浸入. 这些函子远端有界且为三角函子, 因而 $Rf_{!}$ 也是如此.

\setcounter{subsubsection}{8}
\paragraph{}\etiquette[5.1.8.1]{XVII.par5.1.8.1}

\hyperref[XVII.thm5.1.8]{5.1.8} (a) 的传递性同构, 常通过由它们推导的“复合函子的谱序列”使用; 对 $K\in \Ob D(X,\Aa_X)$, 此谱序列为:
$$\etiquette[5.1.8.2]{XVII.eq5.1.8.2}
E^{pq}_2=R^{p}g_! R^{q}h_{!}(K)\Longrightarrow
R^{p+q}f_{!}(K){\quoi}.\leqno{(5.1.8.2)}
$$

为得到它, 注意根据 \hyperref[XVII.thm5.1.8]{5.1.8} b) 及 \textsc{Verdier} \cite{sga4.exp.XVII.bib.V2}, 对每个 $L\in \Ob D(Y,\Aa_Y)$, 有谱序列
$$\etiquette[5.1.8.3]{XVII.eq5.1.8.3}
E^{pq}_2=R^{p} g_!(ℋ^{q}(L))\Longrightarrow
R^{p+q}g_!(L){\quoi},\leqno{(5.1.8.3)}
$$
而 \eqhyperref[XVII.eq5.1.8.2]{5.1.8.2} 是 \eqhyperref[XVII.eq5.1.8.3]{5.1.8.3} 在 $L=Rh_{!}K$ 时的特殊情形.

\begin{enonce*}{定义 5.1.9}\etiquette[5.1.9]{XVII.def5.1.9}
\begin{enumerate}
\renewcommand{\theenumi}{\roman{enumi}}
\renewcommand{\labelenumi}{(\theenumi)}
\item 称 \hyperref[XVII.thm5.1.8]{5.1.8} 中构造的函子为“\emph{\sisi{全}{全}固有支集正像函子}”. 在不致混淆时, 将省略形容词“全”.

\item 第 $q$ 个\emph{固有支集正像函子}, 记作 $R^{q}f_{!}$, 是复合函子 $\Hh^{q}\circ Rf_{!}$.

\item 若 $X$ 为可分闭域 $k$ 上的概形, 可嵌入 $k$ 上的一个固有概形, 将\pageoriginale $\Spec (k)$ 上的模层等同于通常的模, 并以记号 $H^{q}_c(X,\quad)$、$H^{q}_{!}(X,\quad)$ 之一, 或在不致混淆时简记 $H^{q}_c(\quad)$, 表示函子 $R^{q}f_{!}$. 此时称它为“第 $q$ 个\emph{固有支集超上同调群}”, 或在它应用于集中于次数 $0$ 的复形时, 称为“第 $q$ 个\emph{固有支集上同调群}”.
\end{enumerate}
\end{enonce*}

\setcounter{subsubsection}{9}
\setcounter{paragraph}{0}
\paragraph{}\etiquette[5.1.9.1]{XVII.par5.1.9.1}
记号 $Rf_{!}$ 有滥用之处, 因为\emph{函子 $Rf_{!}$ 不是将于 \hyperref[XVII.def6.1.2]{6.1.2} 中定义的函子 $f_{!}$ 的右导出函子}. 有人更喜欢记号 $R_{!}f$.

\subsubsection{}\etiquette[5.1.10]{XVII.subsubsec5.1.10}
\hyperref[XVII.thm5.1.8]{5.1.8} 证明所需的要素, 是固有基变换定理及 $j$ 为固有开浸入时 $j_{\ast}$ 与 $j_{!}$ 的恒等. 由此可陈述 \hyperref[XVII.thm5.1.8]{5.1.8} 的各种变体; 其精确含义及验证留给读者.

\begin{enonce*}{变体 5.1.11}\etiquette[5.1.11]{XVII.var5.1.11}
在范畴 $(S)$ 中, 存在且本质上唯一一种方式, 对每个箭头 $f:X\rightarrow Y$ 定义从 $X$ 上的带基点集合层范畴到 $Y$ 上的带基点集合层范畴的函子 $f_{!}$, 并对每个复合态射 $f=gh$ 定义传递性同构, 满足与 \sisi{\hyperref[XVII.subsubsec5.1.4]{5.1.4}}{\hyperref[XVII.thm5.1.8]{5.1.8}} 的条件类似的条件 {\rm (i)\sisi{,}{} (ii) (iii)}.
\end{enonce*}

此函子的直接描述见 \hyperref[XVII.def6.1.2]{6.1.2}.

\begin{enonce*}{变体 5.1.12}\etiquette[5.1.12]{XVII.var5.1.12}
对\pageoriginale 每个可紧化态射 $f:X\rightarrow Y$, 可以定义从 $X$ 上的 \iInd 有限群层范畴到 $Y$ 上的带基点集合层范畴的函子 $R^{1}f_{!}$, 并且对 $f$ 的每个紧化 $f=\bar{f}j$, 有同构 $R^{1}f_{!}\simeq R^{1}\bar{f}_{\ast}\circ j_{!}$.
\end{enonce*}

这里不要求传递性公式; 唯一的问题是证明函子 $R^{1}f_{!}$ 不依赖于所选取的 $f$ 的紧化. 设 $p:(\bar{f}_1, i_{1})\rightarrow
(\bar{f}_{2}, i_{2})$ 为紧化之间的态射, 即图表
$$\etiquette[5.1.12.1]{XVII.eq5.1.12.1}
\vcenter{
\xymatrix{
X   \ar[rr]^{i_1} \ar[rrd]^{i_2} \ar[dd]_{f}& & \bar{X}_1 \ar[d]^{p}\\
    & &  \bar{X}_2 \ar[lld]^{\bar{f}_2}\\
Y   & & .
}}\leqno{(5.1.12.1)}
$$
其中 $\bar{f}_{1}=\bar{f}_2\circ p$, 且 $p$ 固有. 若 $G$ 为 $X$ 上的群\sisi{层}{层}, 则有带基点集合的正合列 (其中 $\alpha$ 为单射)
\[
0\rightarrow
R^{1}\sisi{f}{\bar{f}}_{2\ast}(p_{\ast}i_{1!}G)\xrightarrow{\alpha}R^{1}\sisi{f}{\bar{f}}_{1\ast}(i_{1!}G)\rightarrow
\bar{f}_{2\ast}R^{1}p_{\ast}(i_{1!}G){\quoi},
\]
以及从 $i_{2!}G$ 到 $p_{\ast}i_{1!}G$ 的态射, 从而经复合得到从 $R^{1}\sisi{f}{\bar{f}}_{2\ast}i_{2!}G$ 到 $R^{1}\sisi{f}{\bar{f}}_{1\ast}i_{1!}G$ 的态射 $\beta$. 与 \hyperref[XVII.par3.3.2.1]{3.3.2.1} 中一样推理, 可见问题在于证明当 $G$ 为 \iInd 有限时, 此态射总是同构.

将固有基变换定理\pageoriginale 应用于 $p$ 及集合层 $i_{1!}G$, 验证得到 $p_{\ast}i_{1!}G=i_{2!}G$, 从而得到正合列
\[
0\rightarrow R^{1}
\sisi{f}{\bar{f}}_{2\ast}(i_{2!}G)\xrightarrow{\beta}R^{1}\sisi{f}{\bar{f}}_{1\ast}(i_{1!}G)\rightarrow
\bar{f}_{2\ast}R^{1}p_{\ast}i_{1!}G{\quoi}.
\]
态射 $\beta$ 为单射 (因为 $\alpha$ 如此). 将固有基变换定理应用于 $p$ 及假设为 \iInd 有限的群层 $G$, 得到 $R^{1}p_{\ast}i_{1!}G=0$. 因此 $\beta$ 为满态射, 至此完成了 \hyperref[XVII.var5.1.12]{5.1.12} 的证明.

\begin{enonce*}{定义 5.1.13}\etiquette[5.1.13]{XVII.def5.1.13}
对每个景 $X$, 以 $D_{\tors}(X)$ 表示 $X$ 上的阿贝尔层范畴的导出范畴中, 由其上同调层为挠层的复形组成的全子范畴.
\end{enonce*}

\begin{enonce*}{变体 5.1.14}\etiquette[5.1.14]{XVII.var5.1.14}
在范畴 $(S)$ 中, 存在且本质上唯一一种方式, 对每个箭头 $f:X\rightarrow Y$ 定义一个三角且远端有界的函子 \emph{\cite[I 7]{sga4.exp.XVII.bib.RD}}
\[
Rf_{!}:D^{+}_{\tors}(X)\longrightarrow D^{+}_{\tors}(Y)
\]
并对每个复合态射 $f=gh$ 定义传递性同构, 满足与 \hyperref[XVII.subsubsec5.1.4]{5.1.4} 的条件类似的条件 \mbox{\rm (i) (ii) (iii)}.
\end{enonce*}

\begin{enonce*}{变体 5.1.\sisi{}{1}5}\etiquette[5.1.15]{XVII.var5.1.15}
在以挠环层赋环的 $S$ 概形范畴中工作, 其箭头为赋环概形的态射, 诱导 (\hyperref[XVII.subsec0.8]{0.8}) 一个 $S$ 可紧化的概形态射. 对每个箭头
\[
f:(X,\Aa)\longrightarrow (Y,\Bb),
\]
以\pageoriginale $Rf_{!}$ 表示从 $\Aa$ 模范畴到 $f^{\sisi{\ast}{-1}}\Bb$ 模范畴的遗忘函子与其后的通常的 $Rf_{!}$ 的复合. 这些函子有通常函子 $Rf_{!}$ 的变性性质.
\end{enonce*}

但要注意, 这里的记号 $Rf_{!}$ 仅当 $\Aa$ 在 $f^{\sisi{\ast}{-1}}\Bb$ 上有限时才合理, 取 $X=Y=\Spec(k)$ 且 $k$ 代数闭即可看出这一点.

\setcounter{subsubsection}{15}
\subsubsection{}\etiquette[5.1.16]{XVII.subsubsec5.1.16}
设 $f:X\rightarrow S$ 为可紧化态射, $j:U\rightarrow
X$ 为到 $X$ 的开浸入, $i:F\rightarrow X$ 为其闭补, $\Aa_S$ 为 $S$ 上的挠环层, $\Aa_X$ 为它在 $X$ 上的逆像, 并设 $K\in \Ob
D(X,\Aa_X)$.

复形正合列
\[
0\longrightarrow j_{!}j^{\ast}K\longrightarrow K\longrightarrow
i_{!}i^{\ast}K\longrightarrow 0
\]
由 \textsc{Verdier} \cite[\sisi{I 2.4}{II 1.1.5}]{sga4.exp.XVII.bib.CD} 在 $D(X,\Aa_X)$ 中给出一个特殊三角 (distinguished triangle). 它在 $Rf_{!}$ 下的像是特殊三角 (\hyperref[XVII.thm5.1.8]{5.1.8} (b))
$$\etiquette[5.1.16.1]{XVII.eq5.1.16.1}
R(fj)_{!}j^{\ast}K\longrightarrow Rf_{!}K\longrightarrow
R(fi)_{!}i^{\ast}K\longrightarrow R(fj)_{!}j^{\ast}K[1]\ldots\leqno{(5.1.16.1)}
$$
从而定义上同调长正合列
$$\etiquette[5.1.16.2]{XVII.eq5.1.16.2}
\ldots \sisi{}{\to}R^{q}(fj)_{!}j^{\ast}K\rightarrow
R^{q}f_{!}K\rightarrow R^{q}(fi)_{!}i^{\ast}K \xrightarrow{j} R^{q+1}(fj)_{!}j^{\ast}K\rightarrow\ldots.\leqno{(5.1.16.2)}
$$

在 $S$ 为可分闭域的谱, $K$ 为集中在次数 $0$ 的层 $𝒢$ 的特殊情形下, 序列 \eqhyperref[XVII.eq5.1.16.2]{5.1.16.2} 成为
$$\etiquette[5.1.16.3]{XVII.eq5.1.16.3}
\ldots H^{q}_c(U,𝒢)\rightarrow H^{q}_c (X,𝒢)\rightarrow
H^{q}_c(F,𝒢)\xrightarrow{j} H^{q+1}_c(U,𝒢)\ldots\leqno{(5.1.16.3)}
$$

\subsubsection{}\etiquette[5.1.17]{XVII.subsubsec5.1.17}

考虑\pageoriginale $S$ 可紧化概形的图表
\[
\xymatrix{
X \ar[rr]^{u} \ar[rd]_{f}& & Y\ar[ld]^{g}\\
  & S &,
}
\]
其中 $u$ \emph{固有}. 以 $S$ 上的一个挠环层 $\Aa_S$ 的逆像赋环于 $X$ 与 $Y$. 对 $K\in
D(Y,\Aa_Y)$, 显然的 (伴随) 态射 $K\rightarrow
Ru_{\ast}u^{\ast}K$ 在应用函子 $Rg_{!}$ 后定义态射
$$\etiquette[5.1.17.1]{XVII.eq5.1.17.1}
u^{\ast}:Rg_{!}K\longrightarrow
Rg_{!}Ru_{\ast}u^{\ast}K=Rg_{!}Ru_{!}u^{\ast}K\sisi{\eqsim}{\simeq} Rf_{!}u^{\ast}K\leqno(5.1.17.1)
$$
(\emph{固有支集上同调对于固有态射的反变性}).

当 $S$ 为代数闭域的谱, $K$ 为置于次数 $0$ 的层 $𝒢$ 时, 此态射成为
$$\etiquette[5.1.17.2]{XVII.eq5.1.17.2}
u^{\ast}:H^{q}_c(Y,𝒢)\longrightarrow H^{q}_c(X,𝒢)\leqno(5.1.17.2)
$$

\subsection{基变换定理}\etiquette[5.2]{XVII.subsec5.2}

\subsubsection{}\etiquette[5.2.1]{XVII.subsubsec5.2.1}

设 $f:X\rightarrow Y$ 为拓扑斯态射, $V$ 为 $Y$ 的开集, $U=f^{-1}(V)$ 为它在 $X$ 中的逆像, $\Aa$ 为 $Y$ 上的环层, $\Aa'$ 为 $X$ 上的环层, $L$ 为上有界的 $(\Aa',
f^{\ast}\Aa)$ 双模复形:
$$\etiquette[5.2.1.1]{XVII.eq5.2.1.1}
\vcenter{
\xymatrix{
U \ar[rr]^{f'} \ar@{^{(}->}[d]^{j'}& & V \ar@{^{(}->}[d]^{j}\\
X \ar[rr]_{f} & & Y.
}}\leqno{(5.2.1.1)}
$$

于是\pageoriginale 有从 $D(V,\Aa)$ 到 $D(X,f^{\ast}\Aa):f^{\ast}j_{!}\isomto
j'_{!}{f'}^{\ast}$ 的函子之间的基变换同构 \eqhyperref[XVII.eq5.1.2.1]{5.1.2.1}. 此外, 从 $D^{-}(U,{f'}^{\ast}\Aa)$ 到 $D^{-}(X,\Aa'):L\overset{\LL}{\otimes}_{f^{\ast}\Aa}j'_{!}K\simeq
j'_{!}({j'}^{\ast}L\overset{\LL}{\otimes}_{f^{\ast}\Aa}K)$ 的函子之间有一个显然的同构. 为构造它, 只须注意: 若复形 $K$ 适合于函子 ${j'}^{\ast}L\overset{\LL}{\otimes}_{f^{\ast}\Aa}$, 则复形 $j'_{!}K$ 适合于函子 $L\overset{\LL}{\otimes}_{f^{\ast}\Aa}$, 且在复形层次上存在上述类型的同构. 将这两个同构复合, 得到从 $D^{-}(V,\Aa)$ 到 $D^{-}(X,\Aa')$ 的函子之间的基变换同构:
$$\etiquette[5.2.1.2]{XVII.eq5.2.1.2}
L\overset{\LL}{\otimes}_{f^{\ast}\Aa}f^{\ast}j_{!}K\isomto j'_{!}({j'}^{\ast}L\overset{\LL}{\otimes}_{{f\sisi{'}{}}^{\ast}\Aa}{f'}^{\ast}K){\quoi}.\leqno(5.2.1.2)
$$

若 $L$ 有界, 其各分量在

$f^{\ast}\Aa$ 上具有有限 \ttor 维数,\nde{此条件是多余的. \Cf \nderef{XVII.ndeKS}.} 则此同构对 $K\in\Ob D(V,\Aa)$ 仍有定义.

\subsubsection{}\etiquette[5.2.2]{XVII.subsubsec5.2.2}

设 $f:X\rightarrow S$ 与 $g:S'\rightarrow S$ 为概形态射, 并假设 $f$ 为开浸入 $j$ 及固有态射 $p$ 的复合, 后者各纤维的维数以整数 $n$ 为界 (当 $S$ 拟紧时自动成立). 这些态射组成各方块均为笛卡尔方块的图表
\[
\xymatrix{
X' \ar[rrrr]_{g'} \ar@{^{(}->}[rd]^{j'} \ar[dd]_{f'}& &         &   &
X\ar@{^{(}->}[rd]^{j} \ar[dd]_(.35){f} & \\
             &\bar{X}'\ar[rrrr]|(.75)\hole_{\bar{g}}\ar[ld]^{p'}  &&   &   & \bar{X}\ar[ld]^{p}\\
S' \ar[rrrr]_{g}  & & &&  S &.
}
\]
设\pageoriginale $\Aa$ 为 $S$ 上的挠环层, $\Aa'$ 为 $S'$ 上的挠环层, $L$ 为上有界的 $(\Aa',g^{\ast}\Aa)$ 双模复形. 以 $\Aa$ (\resp $\Aa'$) 在 $X$ 与 $\bar X$ (\resp $X'$ 与 $\bar X'$) 上的逆像 $\Aa_X$ 与 $\Aa_{\bar X}$ (\resp $\Aa'_{X'}$ 与 $\Aa'_{\bar X'}$), 给 $X$ 与 $\bar X$ (\resp $X'$ 与 $\bar X'$) 赋环. 于是, 从 $D^{-}(X,\Aa_X)$ 到 $D^{-}(\bar X',\Aa'_{\bar X'})$ 的函子 ${p'}^{\ast}L\overset{\LL}{\otimes}\bar g^{\ast}j_{!}K$ 与 $j'_{!}({f'}^{\ast}L\sisi{\otimes}{\overset{\LL}{\otimes}}
{g'}^{\ast}K)$ 之间有基变换同构 \eqhyperref[XVII.eq5.2.1.2]{5.2.1.2}, 从 $D^{-}(\bar X,\Aa_{\bar X})$ 到 $D^{-}(S',\Aa')$ 的函子 $L\overset{\LL}{\otimes}g^{\ast}Rp_{\ast}K$ 与 $Rp'_{\ast}({p'}^{\ast}L\overset{\LL}{\otimes}\bar g^{\ast}K)$ 之间有基变换同构 (\hyperref[XVII.thm4.3.1]{4.3.1}). 经复合, 得到从 $D^{-}(X,\Aa_X)$ 到 $D^{-}(S',\Aa')$ 的函子之间的基变换同构:
$$\etiquette[5.2.2.1]{XVII.eq5.2.2.1}
L\overset{\LL}{\otimes}g^{\ast}Rp_{\ast}j_{!}K\isomto
Rp'_{\ast}j'_{!}({f'}^{\ast}L\overset{\LL}{\otimes}{g'}^{\ast}K){\quoi}.\leqno(5.2.2.1)
$$

当 $L$ 有界且其各分量具有有限 \ttor 维数时, 上述论证及 \hyperref[XVII.var4.3.2]{4.3.2}~(i) 对 $K\in\Ob D(X,\Aa_X)$ 定义了基变换同构.\nde{结合前面的编者注及 \nderef{XVII.ndecb}, 可以不对 $L$ 加任何条件, 也不假设 $p$ 的各纤维维数有界, 而定义此同构.}

为引用 \hyperref[XVII.thm5.1.8]{5.1.8}, 假设 $S$ 拟紧且拟分离. 则函子 $Rp_{\ast}j_{!}=Rf_{!}$ 与 $Rp'_{\ast}j'_{!}=Rf'_{!}$ 只依赖于 $f$ 及 $f'$, 因而下述陈述有意义.

\begin{enonce*}{引理 5.2.3}\etiquette[5.2.3]{XVII.lem5.2.3}
基变换同构 \eqhyperref[XVII.eq5.2.2.1]{5.2.2.1} 不依赖于所选取的 $X$ 的紧化 $f=pj$.
\end{enonce*}

设 $q:\bar X_1\rightarrow\bar X_2$ 为紧化之间的态射
$$\etiquette[5.2.3.1]{XVII.eq5.2.3.1}
\vcenter{
\xymatrix{
X \ar[rr]^{j_1} \ar[rrd]|(.65)\hole^{j_2} \ar[dd]_{f}& & \bar{X}_1 \ar[d]_{q} \ar[lldd]^{p_1}\\
  & &  \bar{X}_2 \ar[lld]^{p_2}\\
S & & .
}}\leqno{(5.2.3.1)}
$$
相对于\pageoriginale $\bar X_1$ 的基变换同构 $c_1$, 是将相对于 $j_1$ 与 $p_1$ 的基变换同构“复合”而得, 或等价地, 将相对于 $j_1$、$q$ 及 $p_2$ 的基变换同构复合而得. 相对于 $\bar X_2$ 的同构 $c_2$, 由将相对于 $j_2$ 与 $p_2$ 的同构复合而得. 为验证 $c_1=c_2$, 只须验证相对于 $j_2$ 的同构 $c'_1$ 是相对于 $j_1$ 与 $q$ 的同构的复合 $c'_2$. 用撇号表示由各箭头 \eqhyperref[XVII.eq5.2.3.1]{5.2.3.1} 经基变换 $g$ 所得的箭头; $c'_1$ 与 $c'_2$ 为同构
\[
c'_{1}, c'_{2}:L\overset{\LL}{\otimes}\bar{g}_{2}^*\ast
j_{2!}K\isomto  j'_{2!}({j'}_{2}^*\ast
L\overset{\LL}{\otimes} {g'}^{\ast}K){\quoi}.
\]
由 \eqhyperref[XVII.eq5.1.3.2]{5.1.3.2}, 为验证 $c'_1=c'_2$, 只须在限制到 $\sisi{X}{\bar X}'_2$ 后验证之, 而这立即成立.

因而, 两个紧化之间只要存在态射, 它们对应的基变换同构就相等; 结论由 \hyperref[XVII.prop3.2.6]{3.2.6}~(ii) 得到.

\begin{enonce*}{引理 5.2.4}\etiquette[5.2.4]{XVII.lem5.2.4}
设 $f=uv:X\rightarrow Y\rightarrow S$ 为 $S$ 可紧化态射的复合. 相对于 $u$、$v$ 及 $f$ 的基变换同构 \eqhyperref[XVII.eq5.2.2.1]{5.2.2.1} 给出交换图表 (记号同下面的 \eqhyperref[XVII.eq5.2.4.1]{5.2.4.1}):
\[
\xymatrix{
L\overset{\LL}{\otimes} g^{\ast} R\sisi{uv}{(uv)}_{!} K
\ar@<-4pt>[rr]^{\sim} \ar@{=}[d]& & R\sisi{u'v'}{(u'v')}_{!}((u'v')^{\ast}L
\overset{\LL}{\otimes}g''^{\ast}K)\ar@{=}[d]\\
L\overset{\LL}{\otimes} g^{\ast}Ru_{!}Rv_{!}K \ar@<-4pt>[r]^(.45){\sim}& Ru'_{!}(u'^{\ast}L
\overset{\LL}{\otimes}g'^{\ast}Rv_{!}\sisi{}{K}) \ar@<-4pt>[r]^(.45){\sim}& Ru'_{!}Rv'_{!}(v'^{\ast}u'^{\ast}L\overset{\LL}{\otimes}g''^{\ast}K).
}
\]
\end{enonce*}

选取\pageoriginale $(u,v)$ 的一个紧化 (\hyperref[XVII.def3.2.5]{3.2.5}):
$$\etiquette[5.2.4.1]{XVII.eq5.2.4.1}
\vcenter{
\xymatrix{
X' \ar[rr]_{g''} \ar[d]_{\sisi{u}{v}'}& & X \ar@{^{(}->}[r]^{j_1} \ar[d]_{\sisi{u}{v}}& \bar{X}
\ar[ld]_{p_1}\ar@{^{(}->}[r]^{j_2} \ar[d]& \bar{\bar{X}} \ar[ld]^{p_3}& \\
Y' \ar[rr]_{g'} \ar[d]_{\sisi{v}{u}'}& & Y \ar@{^{(}->}[r]^{j_3}\ar[d]_{\sisi{v}{u}}& \bar{Y}\ar[ld]^{p_2} &  &\\
S' \ar[rr]_{g} &  & S &      &  .
}}\leqno{(5.2.4.1)}
$$

引理 \hyperref[XVII.lem5.2.4]{5.2.4} 说相对于 $f$ 的基变换同构 $c_1$ 是相对于 $u$ 与 $v$ 的基变换同构的复合 $c_2$. 同构 $c_1$ 由相对于 $j_1$、$j_2$、$p_3$ 及 $p_2$ 的基变换同构复合而得. 对 $c_2$, 则复合相对于 $j_1$、$p_1$、$j_3$ 及 $p_2$ 的基变换同构. 因此只须比较两个同构: 相对于 $p_3$ 与 $j_2$ 的同构的复合 $c'_1$, 以及相对于 $j_3$ 与 $p_1$ 的同构的复合 $c'_2$.

考虑笛卡尔图表
\[
\xymatrix{
  \bar{X}' \ar[ddd]^{p'_1} \ar[rrrr]^{\bar{g}''} \ar@{^{(}->}[rrd]^{j'_2} &    &       &    & \bar{X}\ar@{^{(}->}[rrd]^{j_2} \ar[ddd]^{\sisi{}{p_1}}|(.38)\hole\\
                          & &\bar{\bar{X}}' \ar[rrrr]_(.3){\bar{\bar{g}}} \ar[ddd]^{p'_3}&   &&
& \bar{\bar{X}} \ar[ddd]^{\sisi{}{p_3}}\\
 & & &\\
  Y'  \ar@{^{(}->}[rrd]_{j'_3}\ar[rrrr]^(.7){g'}|(.52)\hole \ar[rdd]_{\sisi{v}{u}'} &    &&
& Y\ar[rdd]_(.65){\sisi{v}{u}} \ar@{^{(}->}[rrd]^{j_3}\\
  & &\bar{Y}' \ar[rrrr]_{\bar{g}'} \ar[ld]^{p'_2}&        &  & &\bar{Y}\ar[ld]^{p_2}\\
 &S'  \ar[rrrr]_{g} &&& & S &,
}
\]
$L'={p'_2}^{\ast}L$\sisi{}{, $K'=j_{1!}K$}\pageoriginale 以及
\[
c'_{1}, c'_{2}: L'\overset{\LL}{\otimes} {\bar{g}'}^{\ast}
R(j_{3}p_{1})_{!} K\sisi{}{'}\isomto  R(j'_{3} p'_{1})_{!} ((j'_{3}
p'_{1})^{\ast} L'\otimes {\bar{g}''}^{\ast}K\sisi{}{'}){\quoi}.
\]
由 \eqhyperref[XVII.eq5.1.3.2]{5.1.3.2}, 为验证 $c'_1=c'_2$, 只须验证 ${j'_3}^{\ast}(c'_1)={j'_3}^{\ast}(c'_2)$, 而这立即成立.

\begin{enonce*}{引理 5.2.5}\etiquette[5.2.5]{XVII.lem5.2.5}
设给定以挠环层赋环的概形的态射
\[
(S'', \Aa'')\xrightarrow{g_{2}} (S',
\Aa')\xrightarrow{g_{1}} (S,\Aa){\quoi},
\]
并设 $f:X\rightarrow S$ 为可紧化态射, 从而给出图表:
\[
\xymatrix{
X'' \ar[r]_{g'_2} \ar[d]^{f''}& X' \ar[r]_{g'_1}\ar[d]^{f'}& X \ar[d]^{f}\\
S'' \ar[r]_{g_2}& S' \ar[r]_{g_1}& S \rlap{.}
}
\]
分别以\pageoriginale $f^{\ast}\Aa$、${f'}^{\ast}\Aa'$ 及 ${f''}^{\ast}\Aa''$ 赋环于 $X$、$X'$ 及 $X''$. 则图表
\[
\xymatrix{
L(g_1g_2)^{\ast}Rf_{!}  \ar[rr]^{\sim} \ar@{=}[d]&& Rf''_{!}L(g'_1g'_2)^{\ast}\ar@{=}[d]\\
Lg_{\sisi{1}{2}}^{\ast}Lg^{\ast}_{\sisi{2}{1}} R f_{!} \ar[r]^{\sim}& Lg^{\ast}_{\sisi{1}{2}} Rf'_{!}
L\sisi{g^\ast_2}{g'^\ast_1}  \ar[r]^{\sim} & Rf''_{!} Lg'^{\ast}_{\sisi{1}{2}}  L g'^{\ast}_{\sisi{2}{1}}
}
\]
交换.
\end{enonce*}

此引理立即由函子 $j_{!}$ 以及 $p$ 固有时函子 $Rp_{\ast}$ 的类似结果得到. 含模的精确版本留给读者表述.

上述陈述归纳为下述定理.

\begin{enonce*}{定理 5.2.6}\etiquette[5.2.6]{XVII.thm5.2.6}
函子 $Rf_{!}$ 与基变换可交换.
\end{enonce*}

\begin{enonce*}[remark]{变体 5.2.7}\etiquette[5.2.7]{XVII.var5.2.7}
\begin{enumerate}
\renewcommand{\theenumi}{\roman{enumi}}
\renewcommand{\labelenumi}{(\theenumi)}
\item 函子 $f_{!}$ (对于带基点集合层: \hyperref[XVII.var5.1.11]{5.1.11}) 及函子 $R^1f_{!}$ (对于群层: \hyperref[XVII.var5.1.12]{5.1.12}) 与基变换可交换.

\item 函子 $Rf_{!}$ (\hyperref[XVII.var5.1.14]{5.1.14}) 与基变换可交换.


\item 在 (\hyperref[XVII.subsubsec5.2.2]{5.2.2}) 的假设下, 若 $L$ 有界且各分量具有有限 \ttor 维数,\nde{如 \hyperref[XVII.subsubsec5.2.2]{5.2.2} 中所指出, 这些假设是多余的.} 则函子 $Rf_{!}:D(X)\rightarrow D(S)$ 与基变换可交换.
\end{enumerate}
\end{enonce*}

下述\pageoriginale 两个命题 \hyperref[XVII.prop5.2.8]{5.2.8} 及 \hyperref[XVII.prop5.2.9]{5.2.9} 合在一起, 实质上等价于定理 \hyperref[XVII.thm5.2.6]{5.2.6}.

\begin{enonce*}{命题 5.2.8}\etiquette[5.2.8]{XVII.prop5.2.8}
设 $f:X\rightarrow S$ 为可紧化态射, $\mathcal F$ 为 $X$ 上的挠层. 对 $S$ 的每个几何点 $s$, 以 $X_s$ 表示 $X$ 在 $s$ 处的几何纤维, 以 $\mathcal F_s$ 表示所诱导的层. 则典范地有
\[
(R^{q}f_{!}(\mathcal{F}))_s=H^{q}_c(X_s,\mathcal{F}_s){\quoi}.
\]
\end{enonce*}

\begin{enonce*}{推论 5.2.8.1}\etiquette[5.2.8.1]{XVII.cor5.2.8.1}
若 $f$ 的纤维维数 $\leq d$, 则函子 $Rf_{!}$ 的上同调维数 $\leq2d$ (\hyperref[XVII.def1.2.9]{1.2.9}): 对每个挠层 $\sisi{F}{ℱ}$, 当 $i>2d$ 时有 $R^if_{!}\sisi{F}{ℱ}=0$. 因而函子 $R^{2d}f_{!}$ 在挠层范畴上右正合.
\end{enonce*}

通过基变换 (\hyperref[XVII.prop5.2.8]{5.2.8}), 归结为 $S$ 为代数闭域的谱的情形. 此时若 $\bar X$ 为 $S$ 上的固有概形, 以 $X$ 为稠密 Zariski 开子概形, 且 $j:X\rightarrow\bar X$, 则 $\dim(\bar X)=\dim(X)\leq\sisi{D}{d}$. 因而对 $X$ 上的挠层 $ℱ$ 及 $i>2d$, 由 \hyperrefext[X][X.cor4.3]{4.3}、\refnoexpose{X.cor5.2},
\[
H^{i}_c(X,ℱ)=H^{i}(\bar{X},j_{!}\sisi{F}{ℱ})=0{\quoi}.
\]

\begin{enonce*}{命题 5.2.9}\etiquette[5.2.9]{XVII.prop5.2.9}
设 $f:X\rightarrow S$ 为可紧化态射, $\Aa$ 为 $S$ 上的挠环层, $K$ 为左 $f^{\ast}\Aa$ 模的上有界复形, $L$ 为右 $\Aa$ 模的上有界复形. 则典范地有
\[
L\overset{\LL}{\underset{\Aa}{\otimes}}
Rf_{!}(K)=Rf_{!}(f^{\ast}L\overset{\LL}{\underset{f^{\ast}\Aa}{\otimes}}K){\quoi}.
\]
\end{enonce*}

此\pageoriginale 命题形式地蕴含下述结果:

\begin{enonce*}{定理 5.2.10}\etiquette[5.2.10]{XVII.thm5.2.10}
设 $f:X\rightarrow S$ 为可紧化态射, $\Aa$ 为 $S$ 上的挠环层, $K\in D^{-}(X,f^{\ast}\Aa)$, 且 $d\in\ZZ$. 若 $K$ 有有限 \tTor 维数 $\leq d$ (\hyperref[XVII.def4.1.9]{4.1.9}), 则 $Rf_{!}K$ 有有限 \tTor 维数 $\leq d$.
\end{enonce*}

只须验证, 对 $S$ 上的每个右 $\Aa$ 模层 $L$,
\[
\sisi{H}{\Hh}^{-i}(L\overset{\LL}{\otimes}Rf_{!}K)=0\quad\mbox{\rm
对于}\quad i>d{\quoi}.
\]
由 \hyperref[XVII.prop5.2.9]{5.2.9}, 已知
\[
\sisi{H}{\Hh}^{-i}(L\overset{\sisi{L}{\LL}}{\otimes}Rf_{!}K)=R^{-i}f_{!}(f^{\ast}L\overset{\sisi{L}{\LL}}{\otimes}K){\quoi},
\]
且根据假设,
\[
\sisi{H}{\Hh}^{-i}(f^{\ast}L\overset{\sisi{L}{\LL}}{\otimes}K)=0\quad\mbox{\rm 对于}\quad
i>d{\quoi}.
\]
注意若 $X$ 上的模层复形 $M$ 满足对 $i>d$ 有 $\sisi{H}{\Hh}^{-i}(M)=0$, 则仍有对 $i>d$ 成立的 $\sisi{H}{\Hh}^{-i}(Rf_{!}M)=0$, 即得结论.

对于常 Noether 环层, 有下述更一般的结果:

\begin{enonce*}{定理 5.2.11}\etiquette[5.2.11]{XVII.thm5.2.11}
设 $f:X\rightarrow S$ 为景态射, $A$ 为环, $C$ 为其中心, 并设 $d\in\ZZ$. 假设
\begin{enumerate}
\renewcommand{\theenumi}{\roman{enumi}}
\renewcommand{\labelenumi}{(\theenumi)}
\item 函子 $Rf_{\ast}:D^{+}(X,C)\rightarrow D^{+}(S,C)$ 具有有限上同调维数;

\item $A$ 为右 Noether 环;

\item $S$ 有足够多点.
\end{enumerate}
对\pageoriginale 每个左 $A$ 模层的复形 $K\in\Ob D^{-}(X,A)$, 若 $K$ 有有限 \tTor 维数 $\leq d$, 则复形 $Rf_{\ast}K\in\Ob D^{-}(S,A)$ 的 \tTor 维数 $\leq d$.
\end{enonce*}

\begin{enonce*}[remark]{证明}
由 (i), 函子 $f_{\ast}$ 可导出为函子
\[
Rf_{\ast}:D(X,A)\longrightarrow D(S,A)
\]
及
\[
Rf_{\ast}:D(X,C)\longrightarrow D(S,C){\quoi}.
\]
\end{enonce*}

由 \hyperref[XVII.lem4.2.8]{4.2.8}, 只须验证对 $S$ 的每个点 $s$, $A$ 模复形 $(Rf_{\ast}K)_s$ 的 \tTor 维数 $\leq d$; 为此只须对每个有限生成右 $A$ 模 $L$ 有
\[
H^{-i}(L\overset{L}{\otimes}Rf_{\ast}K)=0\quad\mbox{\rm 对于}\quad
i>d{\quoi}.
\]

与 \hyperref[XVII.thm5.2.10]{5.2.10} 中一样推理, 可知只须“基变换箭头”
$$\etiquette[5.2.11.1]{XVII.eq5.2.11.1}
L\overset{L}{\otimes}Rf_{\ast}K\longrightarrow Rf_{\ast}(L\overset{L}{\otimes}K)\leqno(5.2.11.1)
$$
为同构.

设 $L^{\ast}$ 为 $L$ 的一个由有限生成自由 $A$ 模组成的消解. 由于所考虑的函子向右远端有界, 只须在将 $L$ 换成 $L\sisi{_{\ast}}{^\ast}$ 的任意分量后, 证明 \eqhyperref[XVII.eq5.2.11.1]{5.2.11.1} 为同构; 于是归结为 $L$ 有限生成且自由的显然情形.

\begin{enonce*}[remark]{注 5.2.12}\etiquette[5.2.12]{XVII.rems5.2.12}
\begin{enumerate}
\renewcommand{\theenumi}{\roman{enumi}}
\renewcommand{\labelenumi}{(\theenumi)}
\item \hyperref[XVII.thm5.2.11]{5.2.11} 的假设\pageoriginale (iii) 很可能不必要.

\item 若省略 $A$ 的 Noether 假设, 当函子 $R^if_{\ast}$ 与滤过正向极限可交换时, \hyperref[XVII.thm5.2.11]{5.2.11} 仍成立.

\item 以\emph{常}环层赋环于 $S$ 的假设是必不可少的. 可取 $S$ 为严格局部离散赋值环的谱, $f$ 为一般点 $\eta$ 的嵌入, 环层 $\Aa$ 为在一般点等于 $\ZZ/n\ZZ$、在闭点 $s$ 等于 $\ZZ/n^2\ZZ$ 的层 ($n>1$ 与剩余特征互素), 并取所考虑的层为模 $f^{\ast}\Aa$, 即得到反例.
\end{enumerate}
\end{enonce*}

\subsection{比较定理与有限性定理}\etiquette[5.3]{XVII.subsec5.3}

\subsubsection{}\etiquette[5.3.1]{XVII.subsubsec5.3.1}
对于 $\CC$ 上的有限型概形 $X$, 可以关联其平展景 $X_{\text{\rm et}}$、超越景 $X_{\cl}$ 以及景态射 (\hyperrefext[XI][XI.sec4]{4})
\[
\epsilon:X_{\cl}\longrightarrow X_{\et}{\quoi}.
\]

若 $j:U\rightarrow X$ 为开浸入, 则交换图表
$$
\xymatrix{
U_{\cl}  \ar[r]\ar[d]   & U_{\et} \ar[d]\\
X_{\cl}  \ar[r] & X_{\et}
}
$$
将 $U^{\sim}_{\cl}$ 等同于 $X^{\sim}_{\cl}$ 的一个开子拓扑斯, 它是 $X_{\widetilde{\text{\rm et}}}$ 的开子拓扑斯 $U_{\widetilde{\text{\rm et}}}$ 的逆像.

若 $p:Y\rightarrow X$ 为固有态射, 则连续映射 $p^{\an}:Y^{\an}\rightarrow X^{\an}$ 固有且分离. 为说明这一点, 用 Chow 引理归结为 $p$ 射影的情形.

设\pageoriginale $f:X\rightarrow S$ 为 $\CC$ 上有限型概形之间的可紧化态射, $f=pj$ 为 $f$ 的一个紧化.
$$\etiquette[5.3.1.1]{XVII.eq5.3.1.1}
\vcenter{
\xymatrix{
X_{\cl} \ar[rrr]_{\epsilon} \ar[rd]^{j_{\cl}} \ar[dd]_{f_{\cl}}&  & & X \ar[rd]^{j}\ar[dd]_(.6){f}|\hole & \\
            &\bar{X}_{\cl}\ar[ld]^{p_{\cl}}  \ar[rrr]_{\epsilon} &  & &\bar{X}\ar[ld]^{p} \\
S_{\cl}  \ar[rrr]_{\epsilon} &                 &   &  S & .
}}\leqno{(5.3.1.1)}
$$
函子 $p_{\cl\ast}j_{\cl !}$ 恰为固有支集正像函子 $f_{\cl !}$ (见 \textsc{Verdier} \cite{sga4.exp.XVII.bib.V3}). 若 $S_{\cl}$, 因而 $X_{\cl}$ 分离, 则函子 $j_{\cl !}$ 将 $X_{\cl}$ 上的软层映为 $\bar X_{\cl}$ 上的软层 \cite[II 3.5.5]{sga4.exp.XVII.bib.TF}. 由于此问题在 $S$ 上是局部的, 无论如何它将 $X_{\cl}$ 上的松弛层映为 $\sisi{X}{\bar X}_{\cl}$ 上对函子 $p_{\cl\ast}$ 无上同调的层. 因而 (\cf \textsc{Verdier} \cite{sga4.exp.XVII.bib.V3})
$$
Rf_{\cl !}=Rp_{\cl\ast}j_{\cl !}{\quoi}.
$$

\subsubsection{}\etiquette[5.3.2]{XVII.subsubsec5.3.2}

由 \hyperref[XVII.lem5.1.2]{5.1.2}, 在从 $D(X)$ 到 $D(\bar X_{\cl})$ 的函子之间有基变换同构:
$$
\epsilon^{\ast}j_{!}\isomto j_{\cl !}\epsilon^{\ast}{\quoi}.
$$

由\pageoriginale \hyperref[XVII.subsubsec4.2.12]{4.2.12} 或 \hyperref[XVII.prop4.1.5]{4.1.5}, 在从 $D^{+}(\bar X)$ 到 $D^{+}(S_{\cl})$ 的函子之间有基变换态射:
$$
\epsilon^{\ast}Rp_{\ast}\longrightarrow Rp_{\cl
\ast}\epsilon^{\ast}{\quoi}.
$$
经复合得到\emph{从} $D^{+}_{\tors}(X)$ \emph{到} $D^{+}_{\tors}(S_{\cl})$ \emph{的函子之间的基变换态射} (\hyperref[XVII.def5.1.13]{5.1.13}):
$$\etiquette[5.3.2.1]{XVII.eq5.3.2.1}
\epsilon^{\ast}Rf_{!}\longrightarrow Rf_{\cl
!}\epsilon^{\ast}{\quoi}.\leqno(5.3.2.1)
$$

与 \hyperref[XVII.lem5.2.3]{5.2.3} 中一样, 可以验证它不依赖于所选取的 $f$ 的紧化; 它满足 \hyperref[XVII.lem5.2.4]{5.2.4} 类型的传递性公式.

\begin{enonce*}{定理 5.3.3}\etiquette[5.3.3]{XVII.thm5.3.3}
\emph{(比较定理).} 态射 \eqhyperref[XVII.eq5.3.2.1]{5.3.2.1} 为同构.
\end{enonce*}

根据构造, 只须处理 $f$ 固有的情形. 所考虑的函子为三角且远端有界 \cite[I 7]{sga4.exp.XVII.bib.RD}. 因而只须处理 $K$ 为置于次数 $0$ 的挠层的情形, 结果由 \hyperrefext[XVI][XVI.thm4.1]{4.1}~(i) 得到.

\begin{enonce*}[remark]{变体 5.3.4}\etiquette[5.3.4]{XVII.var5.3.4}
\begin{enumerate}
\renewcommand{\theenumi}{\roman{enumi}}
\renewcommand{\labelenumi}{(\theenumi)}
\item 若 $\Aa$ 为 $S_{\text{\rm et}}$ 上的挠环层, 同样得到从 $D(X,f^{\ast}\Aa)$ 到 $D(S,\epsilon^{\ast}\Aa)$ 的函子之间的同构 \eqhyperref[XVII.eq5.3.2.1]{5.3.2.1}.

\item 在低次时, 对带基点集合层及 \iInd 有限群层得到通常的变体.
\end{enumerate}
\end{enonce*}

\begin{enonce*}{推论 5.3.5}\etiquette[5.3.5]{XVII.cor5.3.5}
设\pageoriginale $X$ 为 $\CC$ 上的有限型分离概形, 可紧化于 $\CC$ 上的固有概形, $F$ 为 $X$ 上的挠层. 则
$$
H^{q}_{!}(X,F)\isomto H^{q}_c(X_{\cl},\epsilon^{\ast}F){\quoi}.
$$
\end{enonce*}

\begin{enonce*}{定理 5.3.6}\etiquette[5.3.6]{XVII.thm5.3.6}
设 $f:X\rightarrow S$ 为可紧化的有限表示态射, $A$ 为左 Noether 挠环. 函子 $Rf_{!}:D(X,A)\rightarrow D(S,A)$ 将具有可构造上同调的复形映为具有可构造上同调的复形.
\end{enonce*}

此断言在 $S$ 上是局部的, 可以假设其仿射. 通过取极限, 归结为 $S$ 为 Noether 概形的情形. 具有可构造上同调的复形组成 $D(S,A)$ 的三角子范畴. 由于函子 $Rf_{!}$ 远端有界, 只须验证若 $\mathcal F$ 为可构造的 $A$ 模层, 则层 $R^qf_{!}\mathcal F$ 可构造. 写成 $f=pj$, 其中 $p$ 固有, $j$ 为开浸入. 则 $R^qf_{!}\mathcal F=R^qp_{\ast}(j_{!}\mathcal F)$. 层 $j_{!}\mathcal F$ 可构造, 结论由 \hyperrefext[XIV][XIV.thm1.1]{1.1} 得到.

\setcounter{subsubsection}{6}
\subsubsection{}\etiquette[5.3.7]{XVII.subsubsec5.3.7}
在低次时, 对一个可构造的 \iInd 有限集合层或群层的固有支集正像, 有类似陈述.

\begin{enonce*}{推论 5.3.8}\etiquette[5.3.8]{XVII.cor5.3.8}
设 $X$ 为可分闭域 $k$ 上的有限型分离概形, $\Ff$ 为 $X$ 上的可构造挠层. 假设 $X$ 可紧化于 $k$ 上的一个固有概形.\nde{此假设自动满足. 见 \nderef{XVII.ndeConrad}.} 则群 $H^q_c(X,\mathcal F)$ 有限.
\end{enonce*}

\subsection{Künneth 公式}\etiquette[5.4]{XVII.subsec5.4}\pageoriginale

\subsubsection{}\etiquette[5.4.1]{XVII.subsubsec5.4.1}
设有景的交换图表 (\hyperref[XVII.subsec0.11]{0.11})
$$\etiquette[5.4.1.1]{XVII.eq5.4.1.1}
\vcenter{
\xymatrix@R=.5cm@C=1cm{
             &  Z'\ar[ld]_{p'}\ar[rd]^{q'} \ar[dd]^{h}&    \\
X' \ar[dd]_{f}&                   &  Y' \ar[dd]^{g}\\
             &       Z\ar[ld]_{p} \ar[rd]^{q}   &    \\
X \ar[rd]           &                    & Y \ar[ld]\\
             &      S             & &.
}}\leqno{(5.4.1.1)}
$$
假设 $S$ 以环层 $\Aa$ 赋环, 其中心为 $\Cc$; 仍以 $\Aa$ 与 $\Cc$ 表示这些层在 \eqhyperref[XVII.eq5.4.1.1]{5.4.1.1} 的各景上的逆像.

若 $K$ 为 $X$ 上的右 $\Aa$ 模层, $L$ 为 $Y$ 上的左 $\Aa$ 模层, 置
\[
K\mathop{\boxtimes}_{\Aa}L=p^{\ast}\sisi{}{K}
\underset{\Aa}{\otimes}q^{\ast}L
\]
(外张量积). $X'$ (\resp $Y'$) 上的 $K$ (\resp $L$) 也类似定义. 现设 $K$ 为 $X'$ 上的右 $\Aa$ 模层, $L$ 为 $Y'$ 上的左 $\Aa$ 模层. 基变换箭头 (\sisi{}{\hyperref[XVII.prop2.1.3]{2.1.3}, }\sisi{\hyperrefext[XII][XII.eq4.1]{4.1}}{\neweqhyperrefext[XII][XII.eq4.1]{4.1}}) 定义
$$\etiquette[5.4.1.2]{XVII.eq5.4.1.2}
\begin{cases}
p^{\ast}f_{\ast}K\longrightarrow h_{\ast}{p'}^{\ast}K\\
q^{\ast}g_{\ast}L\longrightarrow h_{\ast}{q'}^{\ast}L{\quoi},
\end{cases}\leqno(5.4.1.2)
$$
从而\pageoriginale 取其张量积, 得到箭头
$$\etiquette[5.4.1.3]{XVII.eq5.4.1.3}
f_{\ast}K\boxtimes_{\Aa} g_{\ast}L\longrightarrow
h_{\ast}{p'}^{\ast}K\otimes_{\Aa} h_{\ast}
{q'}^{\ast}L\longrightarrow h_{\ast}(K\boxtimes_{\Aa}L){\quoi}.\leqno(5.4.1.3)
$$

假设函子
\begin{align*}
Rf_{\ast} &:D^{+}(X',\Aa^{\circ})\longrightarrow D^{+}(X,\Aa^{\circ})\\
Rg_{\ast} &:D^{+}(Y',\Aa)\longrightarrow D^{+}(Y,\Aa)\\
Rh_{\ast} &:D^{+}(Z',\Cc)\longrightarrow D^{+}(Z,\Cc)
\end{align*}
具有有限上同调维数, 从而延拓到无界导出范畴. 因而对 $K\in\Ob D^{-}(X',\Aa^{\circ})$ 与 $L\in\Ob D^{-}(Y',\Aa)$, 自然要定义 \emph{Künneth 箭头}
$$\etiquette[5.4.1.4]{XVII.eq5.4.1.4}
Rf_{\ast}K\overset{\LL}{\boxtimes}_{\Aa}
Rg_{\ast}L\longrightarrow Rh_{\ast}(K
\overset{\LL}{\boxtimes}_{\Aa}L){\quoi}.\leqno(5.4.1.4)
$$

只在所考虑的景有足够多点时定义它, 以便使用 \hyperref[XVII.prop4.2.10]{4.2.10} 及 \hyperref[XVII.lem4.2.8]{4.2.8}. 由 \hyperref[XVII.prop4.2.10]{4.2.10}, 只须对左适合于函子 $(K,L)\mapsto K\boxtimes_{\Aa}L$ 的一对 $(K,L)$ 定义它, 其中 $K$ 与 $L$ 上有界, 并右适合于函子 $f_{\ast}$ 与 $g_{\ast}$.

设 $K_1$ 为 $f_{\ast}K$ 的平坦消解, $L_1$ 为 $g_{\ast}L$ 的平坦消解, $u:K\overset{\LL}{\boxtimes}\sisi{}{_\Aa}L\rightarrow M$ 为 $K\overset{\LL}{\boxtimes}\sisi{}{_\Aa}L$ 的适合于函子 $Rh_{\ast}$ 的消解. 有态射
\begin{gather*}
Rf_{\ast}K\overset{\LL}{\boxtimes}_{\Aa}Rg_{\ast}L\simeq
K_1\boxtimes_{\Aa}L_1\longrightarrow
f_{\ast}K\boxtimes_{\Aa}g_{\ast}L\longrightarrow\\
\longrightarrow h_{\ast}(K\boxtimes_{\Aa}L)\longrightarrow
h_{\ast}M\simeq Rh_{\ast}(K\overset{\LL}{\boxtimes}_{\Aa}L)
\end{gather*}
将 \eqhyperref[XVII.eq5.4.1.4]{5.4.1.4} 定义为其复合.

此\pageoriginale 箭头有下述刻画性质, 其证明与 \hyperref[XVII.prop4.2.13]{4.2.13} 的证明完全类似:

\setcounter{paragraph}{4}
\paragraph{}\etiquette[5.4.1.5]{XVII.par5.4.1.5}
设 $K\in\Ob C^{-}(X',\Aa^{\circ})$ 及 $L\in\Ob C^{-}(Y',\Aa)$. 置 $K_1=f_{\ast}K$、$L_1=g_{\ast}L$ 及 $M_2=K\boxtimes\sisi{}{_\Aa}L$. 则下图交换:
$$
\xymatrix@R=1.25cm{
Rf_{\ast}K
\overset{\LL}{\boxtimes}_{\Aa}Rg_{\ast}
L \ar[rrr]& & & Rh_{\ast}(K
\overset{\LL}{\boxtimes}_{\Aa}L) \ar[d]\\
K_1\overset{\LL}{\boxtimes}_{\Aa}L_{\sisi{k}{1}}
\ar[r] \ar[u]&
K_1\boxtimes\sisi{}{_\Aa}L_1 \ar[r]& h_{\ast}M_2 \ar[r]& Rh_{\ast}M_2.
}
$$

若改从态射 $K_1\rightarrow f_{\ast}K$、$L_1\rightarrow g_{\ast}L$ 及 $K\boxtimes\sisi{}{_\Aa}L\rightarrow M_2$ 出发, 也有类似的交换性 (\cf \hyperref[XVII.prop4.2.13]{4.2.13}).

\hyperref[XVII.par5.4.1.5]{5.4.1.5} 中的交换性原则上使 Künneth 态射的相容性性质易于验证.

\subsubsection{}\etiquette[5.4.2]{XVII.subsubsec5.4.2}
设 \eqhyperref[XVII.eq5.4.1.1]{5.4.1.1} 为概形的交换图表, $\Aa$ 为 $S$ 上的挠层. 假设 $f$ 与 $g$ 可紧化 (\hyperref[XVII.def3.2.1]{3.2.1}), $Z=\fprod{X}{Y}{S}$, 且 $Z'=\fprod{X'}{Y'}{S}$. 将 $f$ 与 $g$ 分解为固有态射与开浸入, 得到图表
$$
\xymatrix@C=1.75cm@R=1cm{
  X'      \sisi{\ar[r]^{p'}}{} \ar[d]^{j_1} \ar@/_1.75pc/[dd]_{f}  & Z'   \ar[d]^{j_3}\sisi{}{\ar[l]_{p'}\ar[r]^{q'}}
& \sisi{Y'\ar[l]_{q'}}{} \ar[d]^{j_2} \ar@/^1.75pc/[dd]^{g}\\
\bar{X}' \sisi{\ar[r]^{\bar{p}}}{} \ar[d]^{\bar{f}}  & \bar{Z}' \ar[d]^{\bar{h}}\sisi{}{\ar[l]_{\bar{p}}\ar[r]^{\bar{q}}} & \bar{Y}' \sisi{\ar[l]_{\bar{q}}}{}\ar[d]^{\bar{g}}\\
X            &    Z \ar[l]_{p} \ar[r]^{q}  & Y
}
$$
其中\pageoriginale $\bar Z'=\fprod{\bar X'}{\bar Y'}{S}$. 对 $X'$ 与 $Y'$ 上的复形 $K$ 与 $L$, 有一个显然的同构
$$\etiquette[5.4.2.1]{XVII.eq5.4.2.1}
j_{3!}(K\boxtimes_{\Aa}L)\isomto
j_{1!}K\boxtimes_{\Aa}j_{2!}L{\quoi}.\leqno(5.4.2.1)
$$

将此同构的逆与相对于 $(\bar f,\bar g,\bar h,j_{1!}K,j_{2!}L)$ 的 Künneth 箭头 \eqhyperref[XVII.eq5.4.1.4]{5.4.1.4}“复合”, 对\emph{ $K\in\Ob D^{-}(X',\Aa)$ 及 $L\in\Ob D^{-}(Y',\Aa)$, 得到 Künneth 箭头}
$$\etiquette[5.4.2.2]{XVII.eq5.4.2.2}
Rf_{!}K\overset{\LL}{\boxtimes}\sisi{}{_\Aa}Rg_{!}L\longrightarrow
Rh_{!}(K\overset{\LL}{\boxtimes}\sisi{}{_\Aa}L){\quoi},\leqno(5.4.2.2)
$$
可以验证它与所选取的紧化无关.

\begin{enonce*}{定理 5.4.3}\etiquette[5.4.3]{XVII.thm5.4.3}
在 \hyperref[XVII.subsubsec5.4.2]{5.4.2} 的假设下, 固有支集上同调的 Künneth 箭头 \eqhyperref[XVII.eq5.4.2.2]{5.4.2.2} 为同构.
\end{enonce*}

\begin{enonce*}[remark]{证明}
\begin{enumerate}
\renewcommand{\theenumi}{\alph{enumi}}
\renewcommand{\labelenumi}{\theenumi)}
\item 根据构造, 只须在 $f$ 与 $g$ (因而 $h$) 固有时验证结果.

\item 设给定交换图表
$$\etiquette[5.4.3.1]{XVII.eq5.4.3.1}
\vcenter{
\xymatrix@C=1.75cm@R=1.25cm{
X_1 \ar[r]^{p_1} \ar[d]& S_1 \sisi{\ar[r]^{q_1}}{}\ar[d]& Y_1\ar[d]\sisi{}{\ar[l]_{q_1}}\\
X   \ar[r]^{p}& S  \sisi{\ar[r]^{q}}{} &  Y\sisi{}{\ar[l]_q}.
}}\leqno{(5.4.3.1)}
$$
并设 $Z_1=\fprod{X_1}{Y_1}{S_1}$、$X'_1=\fprod{X'}{X_1}{X}$、$Y'_1=\fprod{Y'}{Y_1}{Y}$ 及 $Z'_1=\fprod{X'_1}{Y'_1}{S_1}=\fprod{Z'}{Z\sisi{_{1}}{}}{Z_1}$, 从而有图表 \sisi{(\hyperref[XVII.subsubsec5.4.1]{5.4.1})}{\eqhyperref[XVII.eq5.4.1.1]{5.4.1.1}} 之间的态射
\[
(S_1,X_1,Y_1,Z_1,X'_1,Y'_1,Z'_1)\longrightarrow
(S,X,Y,Z,X',Y',Z')
\]
其所有箭头统记作 $C$.

将承认\pageoriginale 下述引理:

\begin{enonce*}{引理 5.4.3.2}\etiquette[5.4.3.2]{XVII.lem5.4.3.2}
下图交换:
$$
\xymatrix{
C^{\ast}(Rf_{!} K \overset{L}{\boxtimes}_{\Aa}
Rg_{!} L) \ar@{=}[r] \ar[d]& C^{\ast}Rf_{!}K \overset{L}{\boxtimes}_{\Aa}
C^{\ast} Rg_{!} L \ar[r]^{\sim}&  Rf_{1!} C^{\ast}K
\overset{L}{\boxtimes}_{\Aa} Rg_{1!} C^{\ast}
L \ar[d]\\
C^{\ast}(Rh_{!}(K \overset{L}{\boxtimes}_{\Aa}
L)) \ar[r]^{\sim}& Rh_{\sisi{}{1}!}(C^{\ast}(K
 \overset{L}{\boxtimes}_{\Aa}  L))\ar@{=}[r] &Rh_{\sisi{}{1}!}(C^{\ast}K
 \overset{L}{\boxtimes}_{\Aa}  C^*L ).
}
$$
\end{enonce*}

在此图表中, 水平箭头为基变换箭头, 竖直箭头为所考虑的两个图表 \eqhyperref[XVII.eq5.4.1.1]{5.4.1.1} 的 Künneth 箭头.

为证明 \hyperref[XVII.thm5.4.3]{5.4.3}, 只须证明 Künneth 箭头在 $Z$ 的每个几何点处为同构. 由 \hyperref[XVII.lem5.4.3.2]{5.4.3.2}, 只须在任意使 $S_1$ 为代数闭域的谱且 $p_1,q_1$ 为同构的基变换 \eqhyperref[XVII.eq5.4.3.1]{5.4.3.1} 后验证之.

作这样的基变换后, 图表 \eqhyperref[XVII.eq5.4.1.1]{5.4.1.1} 归结为图表
$$\vcenter{\etiquette[5.4.3.3]{XVII.eq5.4.3.3}
\xymatrix@R=1.5cm@C=2cm{
    & \fprod{X}{Y}{S} \ar[ld]_{p} \ar[rd]^{q} \ar[dd]^{f\times g}&     \\
X \ar[rd]_{f}  &                   &    Y \ar[ld]^{g}\\
    &      S             &\sisi{.}{}
}}\leqno{(5.4.3.3)}
$$
而 Künneth 箭头归结为箭头 \eqhyperref[XVII.eq5.4.4.1]{5.4.4.1}
$$\etiquette[5.4.3.4]{XVII.eq5.4.3.4}
k:Rf_{!}K\overset{\LL}{\otimes}\sisi{}{_\Aa} Rg_{!} L\longrightarrow
R(f\times g)_{!}(K\overset{\LL}{\boxtimes}\sisi{}{_\Aa}L){\quoi}.\leqno(5.4.3.4)
$$

\item 设\pageoriginale \eqhyperref[XVII.eq5.4.3.3]{5.4.3.3} 为一个其中 $f$ 与 $g$ 固有的图表. 将承认下述引理:

\begin{enonce*}{引理 5.4.3.5}\etiquette[5.4.3.5]{XVII.lem5.4.3.5}
下图交换:
$$
\xymatrix@C=1.75cm{
Rf_{!}K \overset{\sisi{L}{\LL}}{\underset{\Aa}{\otimes}} Rg_{!}L
\ar[r] \ar[d]& R(f\times g)_{!}(K
\overset{\sisi{L}{\LL}}{\underset{\Aa}{\otimes}} L) \ar@{=}[d]\\
Rf_{!}(K \overset{\sisi{L}{\LL}}{\sisi{\boxtimes}{\otimes}}_{\Aa}
f^{\ast}Rg_{!}L) \ar[r]& Rf_{!} Rp_{!}(p^{\ast}K \overset{\sisi{}{\LL}}{\otimes}\sisi{}{_\Aa} q^{\ast}L).
}
$$
\end{enonce*}

在此图表中, 水平箭头为 Künneth \sisi{箭头}{箭头}及基变换箭头. 竖直箭头为 $f$ 下的基变换箭头及一个传递性同构. 定理 \sisi{\hyperref[XVII.subsubsec5.5.9]{5.5.9}}{\hyperref[XVII.thm5.4.3]{5.4.3}} 由此引理及 \hyperref[XVII.thm4.3.1]{4.3.1} 得到.

勤勉的读者可以验证, 即使 $f$ 与 $g$ 不一定固有, 此引理仍然成立.
\end{enumerate}
\end{enonce*}

\setcounter{subsubsection}{3}
\subsubsection{}\etiquette[5.4.4]{XVII.subsubsec5.4.4}
设 $S$ 为概形, $f_i:X'_i\rightarrow X_i$ $(i\in I)$ 为 $S$ 概形之间的可紧化态射的有限族, $X$ 为各 $X_i$ 在 $S$ 上的积, $X'$ 为各 $X'_i$ 在 $S$ 上的积, $f:X'\rightarrow X$ 为各 $f_i$ 的积. 若如 \hyperref[XVII.subsubsec5.4.1]{5.4.1} 那样以 $\Aa$ 赋环于 $S$, 并以 $\Aa$ 的逆像赋环于各 $X_i$ 与 $X'_i$, 则更一般地, 对 $K_i\in\Ob D^{-}(X'_i,\Aa)$ 有同构
$$\etiquette[5.4.4.1]{XVII.eq5.4.4.1}
\mathop{\boxtimes}^{\LL}_{i\in I} Rf_{i!}(K_i)
\isomto Rf_{!}(\mathop{\boxtimes}^{\LL}_{i\in I}K_i){\quoi},\leqno(5.4.4.1)
$$
此同构及 \eqhyperref[XVII.eq5.4.4.1]{5.4.4.1} 两端的定义都\emph{不要求}在 $I$ 上选取全序.



% END native input: chapters/exp-17-part3-zh.tex
% BEGIN native input: chapters/exp-17-part4-zh.tex
\subsection{对称幂的上同调}\etiquette[5.5]{XVII.subsec5.5}
\pageoriginale

本 ${\rm n^{\circ}}$ 的主要结果是定理 \hyperref[XVII.thm5.5.21]{5.5.21}; 特别地, 对适当系数 (例如 $\ZZ/n\ZZ$), 它用代数闭域 $k$ 上概形 $X$ 的固有支集上同调, 给出其对称幂的固有支集上同调, 也就是以 $X$ 的固有支集上同调来表示. 这一定理在半单纯拓扑中的类似物是众所周知的, 至少对以交换环 $A$ 为系数的上同调而言如此; 它可在半单纯上链的层面验证. 事实上, 若 $X_{\ast}$ 为半单纯集, 则从 $X_{\ast}$ 的第 $n\sisi{}{}$ 对称幂 $(X_{\ast})^n/\Sf_n$ 到 $X_{\ast}$ 的几何实现的第 $n$ 对称幂, 有一个典范连续双射, 并且它在上同调上诱导同构.

\subsubsection{}\etiquette[5.5.1]{XVII.subsubsec5.5.1}
回顾平坦模的一些性质 (归功于 \textsc{D. Lazard}), 函子 $TS^n$ 与 $T^n$ 的一些性质 (归功于 \textsc{Roby}), 以及非加性函子的导出函子的一些性质 (归功于 \textsc{Dold} 和 \textsc{Puppe}).

\paragraph{}\etiquette[5.5.1.1]{XVII.par5.5.1.1}
\emph{环 $A$ 上的每个平坦模都是有限生成自由模的滤过归纳极限} (\textsc{Lazard} \cite[I 1.2]{sga4.exp.XVII.bib.La}). 这一结果蕴含下述陈述:

\paragraph{}\etiquette[5.5.1.2]{XVII.par5.5.1.2}
函子 $\varinjlim$ 给出一个等价, 从有限生成自由 (\resp 有限生成投射) $A$-模范畴的 \iInd-对象 (\hyperrefext[I][I.sec8.1]{8.1}) 范畴到平坦 $A$-模范畴.

\paragraph{}\etiquette[5.5.1.3]{XVII.par5.5.1.3}
若\pageoriginale $A$ 与 $B$ 为两个环, 则每个将有限生成自由模映为平坦模并与滤过归纳极限可交换的函子 $T:\Mod(A)\rightarrow\Mod(B)$, 都将平坦模映为平坦模 (\textsc{Lazard} \cite[I 1.5]{sga4.exp.XVII.bib.La}).

\paragraph{}\etiquette[5.5.1.4]{XVII.par5.5.1.4}
每个 $A$-模的正合列, 若 $L''$ 平坦,

$$
0\longrightarrow L'\longrightarrow L\longrightarrow L''\longrightarrow
0
$$
就是分裂正合列的滤过归纳极限 (\textsc{Lazard} \cite[I 2.3]{sga4.exp.XVII.bib.La}).

\subsubsection{}\etiquette[5.5.2]{XVII.subsubsec5.5.2}
设 $A$ 为交换环, $M$ 为 $A$-模.

\paragraph{}\etiquette[5.5.2.1]{XVII.par5.5.2.1}
对每个整数 $n$, 以 $TS^{n}(M)$ (或 $TS^{n}_A(M)$) 表示 $M$ 的 $n$ 次对称张量（symmetric tensor）模. 以 $\Sf_n$ 表示 $n$ 次对称群, 则按定义有
$$
TS^{n}(M)=(\bigotimes^{n}M)^{\Sf_n}{\quoi}.
$$

若 $B$ 为 $A$-代数, 则 $TS^{n}(B)$ 是 $\bigotimes\limits^{n}B$ 的子代数. 若 $B$ 还交换, 则按定义 $\Spec (TS^{n}(B))$ 为 $\Spec(B)$ 在 $\Spec(A)$ 上的第 $n$ 对称幂:
$$
\Spec(TS^{n}_{A}(B))=\Sym^{n}_{\Spec(A)}(\Spec(B)\underset{\dfn}{=}(\Spec(B)/\Spec(A))^{n}/\Sf_n{\quoi}.
$$

\paragraph{}\etiquette[5.5.2.2]{XVII.par5.5.2.2}
回顾 (\textsc{Roby} \cite[I \sisi{1 pg.226}{2 p.~219}]{sga4.exp.XVII.bib.Ro}), 从 $A$-模 $M$ 到 $A$-模 $N$ 的一个\emph{多项式律}（polynomial law）$f$, 是对每个交换 $A$-代数 $A'$ 给定映射 $f_{A'}:M\otimes_{A}A'\rightarrow N\otimes_{A}A'$, 并且这些映射满足: 对每个 $A$-代数同态 $u:A'\rightarrow A''$, 图
\[
\xymatrix{
M \otimes_A A'  \ar[r]^{f_{A'}} \ar[d]^{u}& N \otimes_A A' \ar[d]^{u}\\
M \otimes_A A'' \ar[r]^{f_{A''}}& N \otimes_A A''
}
\]
均\pageoriginale 交换.

若对每个 $a'\in A'$ 恒有
\[
f_{A'}(a'x)={a'}^{n}f_{A'}(x){\quoi}.
\]
则称多项式律 $f$ \emph{为 $n$ 次齐次的} (\textsc{Roby} \cite[I \sisi{1}{8} p\sisi{g}{}.~226]{sga4.exp.XVII.bib.Ro}). 这样的多项式律也称为\emph{$n$ 次映射} (二次、三次、\ldots); 它们所成的集合记为
\[
\Hom n^{\text{\rm ique}} (M,N){\quoi}.
\]

更一般地, 若 $(M_{i})_{i\in I}$ 为一族有限个 $A$-模, $\sheafn =(n_{i})_{i\in I}$ 为一族整数, 则从模 $M_{i}$ 到 $N$ 的多重次数为 $\sheafn $ 的多重齐次多项式律, 是从 $\sisi{\sum}{\oplus} M_{i}$ 到 $N$ 的多项式律, 使对 $a_{i}\in A'$ 恒有
\[
f_{A'}((a_{i}x_{i})_{i\in I})=\prod_{i\in I} a^{n_{i}}_i\cdot
f_{A'}((x_{i})_{i\in I}){\quoi}.
\]
以 $\Hom \sheafn ^{\text{\rm ique}}((M_{i}),N)$ 表示这些多项式律所成的集合.

\paragraph{}\etiquette[5.5.2.3]{XVII.par5.5.2.3}
对每个模 $M$, 从 $A$-模到集合的协变函子
$$
N\longmapsto \Hom n^{\text{\rm ique}} (M,N)
$$
由一个 $A$-模 $\Gamma^{n}(M)$ (或 $\Gamma^{n}_{A}(M)$) 表示, 该模配有\emph{泛\pageoriginale $n$ 次映射} $\gamma^{n}:M\rightarrow
\Gamma^{n}(M)$ (\textsc{Roby} \cite[IV 1 \sisi{pg.226}{p.~265}]{sga4.exp.XVII.bib.Ro}). 给定如 \hyperref[XVII.par5.5.2.2]{5.5.2.2} 中的 $(M_{i})_{i\in I}$ 与 $\sheafn $, 协变函子
$$
N\longmapsto \Hom \sheafn ^{\text{\rm ique}} ((M_{i}),N)
$$
由模 $\bigotimes\limits_{i\in
I}\Gamma^{n_i}(M_{i})$ 表示, 该模配有\emph{泛多重齐次多项式律 (多重次数为 $\sheafn $)} $(m_{i})\mapsto
\bigotimes\limits_{i\in I}\gamma^{n_{i}}(m_{i})$.

\paragraph{}\etiquette[5.5.2.4]{XVII.par5.5.2.4}
函子 $TS^{n}$ 与 $\Gamma^{n}$ 都与滤过归纳极限可交换 (后者见 \textsc{Roby} \cite[IV \sisi{6 pg.275}{7 p.~277}]{sga4.exp.XVII.bib.Ro}). 它们将自由模映为自由模, 因而也将平坦模映为平坦模 (\hyperref[XVII.par5.5.1.2]{5.5.1.2}).

\paragraph{}\etiquette[5.5.2.5]{XVII.par5.5.2.5}
由 $m\mapsto \bigotimes\limits^{n}m$ 给出的从 $M$ 到 $TS^{n}(M)$ 的 $n^{\text{\rm ique}}$ 次映射, 确定了一个从 $\Gamma^{n}(M)$ 到 $TS^{n}(M)$ 的映射 $\varphi$, 称为典范映射, 满足
$$
\varphi(\gamma^{n}(m))=\bigotimes^{n}m{\quoi}.
$$
当 $M$ 有限自由时, 这一映射为同构 (\textsc{Roby} \cite[IV \sisi{pg.272}{5 p.~272}]{sga4.exp.XVII.bib.Ro}); 因而当 $M$ 平坦时也为同构 (\hyperref[XVII.par5.5.2.4]{5.5.2.4} 及 \hyperref[XVII.par5.5.1.1]{5.5.1.1}).

设 $M$ 为 $A$-模, $P$ 为正合列
$$
P:P_{1}\xrightarrow{d} P_{0} \xrightarrow{\epsilon} M\longrightarrow
0{\quoi}.
$$
令 $P'_{1}=P_{1}\times P_{0}$, 并设 $j_{0}$ 与 $j_{1}$ 为从 $P'_{1}$ 到 $P_{0}$ 的态射, 坐标分别为 $(d,\Id)$ 与 $(0,\Id)$. 则序列
$$
\xymatrix@C=1.75cm{
\Gamma^{n}(P'_1) \ar@<2pt>[r]^{\Gamma^{n}(j_0)} \ar@<-2pt>[r]_{\Gamma^{n}(j_1)}& \Gamma^{n}(P_0) \ar[r]& \Gamma^{n}(M) \ar[r]& 0
}
$$
正合 (\textsc{Roby} \cite[\sisi{pg.284}{IV 10 p.~284}]{sga4.exp.XVII.bib.Ro}). \emph{当 $P$ 为 $M$ 的平坦表示时},\pageoriginale 这就用函子 $TS^{n}$ 给出了函子 $\Gamma^{n}$ 的描述:
$$\etiquette[5.5.2.5.1]{XVII.eq5.5.2.5.1}
TS^{n}(P'_1)\rightrightarrows TS^{n}(P_{0})\rightarrow
\Gamma^{n}(M)\rightarrow 0{\quoi}.\leqno(5.5.2.5.1)
$$

\paragraph{}\etiquette[5.5.2.6]{XVII.par5.5.2.6}
若 $M=M'\oplus M''$, 则 $n^{\text{\rm ique}}$ 次映射
\[
m\longmapsto (\gamma^{n'}m\otimes \gamma^{n''}m)_{n'+n''=n}
\]
定义一个称为典范同构的同构 (\textsc{Roby} \cite[III \sisi{4}{9 p.~262}]{sga4.exp.XVII.bib.Ro})
$$\etiquette[5.5.2.6.1]{XVII.eq5.5.2.6.1}
\Gamma^{n}(M)\isomto \sisi{\sum}{\bigoplus}_{n'+n''=n} \Gamma^{n'}(M')\otimes \Gamma^{n''}(M''){\quoi}.\leqno(5.5.2.6.1)
$$

若 $\Sigma$ 为正合列
\[
\Sigma:0\longrightarrow M'\longrightarrow M\longrightarrow
M''\longrightarrow 0{\quoi},
\]
则模 $\Gamma^{n}(M)$ 配有典范的 (递减) 滤过及态射
$$\etiquette[5.5.2.6.2]{XVII.eq5.5.2.6.2}
\Gamma^{i}(M'')\otimes \Gamma^{n-i}(M')\longrightarrow \Gr^{i}(\Gamma^{n}(M)){\quoi}.\leqno(5.5.2.6.2)
$$
这些数据的特征为: 它们关于 $\Sigma$ 具有函子性, 且当 $\Sigma$ 分裂时满足与 \eqhyperref[XVII.eq5.5.2.6.1]{5.5.2.6.1} 的显然相容性. 当 $\Sigma$ 分裂时, 箭头 \eqhyperref[XVII.eq5.5.2.6.2]{5.5.2.6.2} 为同构, 因而当 $M''$ 平坦时也为同构 (\hyperref[XVII.par5.5.1.4]{5.5.1.4}).

\paragraph{}\etiquette[5.5.2.7]{XVII.par5.5.2.7}
函子 $\Gamma^{n}$ 与标量扩张可交换 (\textsc{Roby} \cite[III \sisi{3 pg.262}{8 p.~262}]{sga4.exp.XVII.bib.Ro}) (对于 $TS^{n}$ 的类似性质一般不成立).

\paragraph{}\etiquette[5.5.2.8]{XVII.par5.5.2.8}
若 $u:M_{1}\times M_{2}\rightarrow N$ 为双线性映射, 则 $\gamma^{n}\circ u$ 为双次数 $(n,n)$ 的双齐次多项式律, 从而确定 (\hyperref[XVII.par5.5.2.3]{5.5.2.3}) 一个满足 $u^{n}(\gamma^{n}x\otimes
\gamma^{n}y)=\gamma^{n}u(x,y)$ 的 $u^{n}:\Gamma^{n}(M_{1})\otimes
\Gamma^{n}(M_{2})\rightarrow \Gamma^{n}(N)$. 特别地,\pageoriginale 若 $B$ 为 $A$-代数, 乘法 $B\times B\rightarrow B$ 在 $\Gamma^{n}(B)$ 上诱导一个乘积. 这一乘积使 $\Gamma^{n}(B)$ 成为 $A$-代数, 并且从 $\Gamma^{n}(B)$ 到 $TS^{n}(B)$ 的典范映射是代数同态.

\subsubsection{}\etiquette[5.5.3]{XVII.subsubsec5.5.3}
设 $\Aa$ 为阿贝尔范畴. 以 $\text{\rm ssimpl}^{+}(\Aa)$ 表示 $\Aa$ 中的余半单纯对象\nde{现代术语称为\emph{余单纯}对象（cosimplicial objects）.}所成的范畴, 以 $C^{\geq 0}(\Aa)$ 表示满足在 $i<0$ 时 $K^{i}=0$ 的微分复形 $K$ 所成的范畴, 以 $K^{\geq
  0}(\Aa)$ 表示 $C^{\geq 0}(\Aa)$ 在 $K(\Aa)$ 中的像. 若 $K\in \Ob
\text{ssimpl}^{+}(\Aa)$, 则伴随的正规化复形（normalized complex）$N(K)\in \Ob C^{\geq 0}(\Aa)$ 定义为
$$\etiquette[5.5.3.1]{XVII.eq5.5.3.1}
\begin{cases}
N(K)^{n} =\text{退化算子的核的交}\\
\qquad\qquad\text{(从 $K^{n}$ 到 $K^{n-1}$)}\\
d=\sum(-1)^{i}\partial_{i}{\quoi}.
\end{cases}\leqno(5.5.3.1)
$$

若 $K$ 也表示各分量为 $K^{n}$, 微分为 $d=\sum(-1)^{i}\partial_{i}$ 的复形, 则 $N(K)$ 为 $K$ 的直和因子子复形, 且 $K/N(K)$ 零伦.

由 \textsc{Dold-Puppe} \cite{sga4.exp.XVII.bib.DP}, 函子 $N$ 为范畴等价
$$
N:\;\text{\rm ssimpl}^{+}(\Aa)\longrightarrow C^{\geq
0}(\Aa){\quoi}.
$$
这一等价及其逆都将同伦态射映为同伦态射. 此外, 对每个 $K\in \Ob\text{\rm
ssimpl}^{+}(\Aa)$, $K^{n}$ 为若干个 $N(K)^{i}$ 的直和 $(i\leq n)$; 若 $\PP$ 为 $\Aa$ 中对象的一个性质, 在有限直和与直和因子下稳定, 则 $K^{i}$ 在 $i\leq n$ 时具有性质 $\PP$, 当且仅当 $N(K)^{i}$ 在 $i\leq n$ 时具有这一性质.

\textsc{Dold} 和 \textsc{Puppe} 利用等价\pageoriginale $N$, 将每个从 $\Aa$ 到任意范畴 $\Bb$ 的函子 $T:\Aa\rightarrow \Bb$ (未必加性), 延拓为从 $C^{\geq 0}(\Aa)$ 到 $C^{\geq
0}(\Bb)$ 的函子 $C^{\geq 0}
(T)$ (或简记 $T$):
$$\etiquette[5.5.3.2]{XVII.eq5.5.3.2}
C^{\geq 0}(T)=NTN^{-1}{\quoi}.\leqno(5.5.3.2)
$$

函子 $C^{\geq 0}(T)$ 下降到商范畴后定义
$$\etiquette[5.5.3.3]{XVII.eq5.5.3.3}
T\quad\mbox{\rm 或}\quad K^{\geq 0}(T):K^{\geq
0}(\Aa)\longrightarrow K^{\geq 0}(\Bb){\quoi}.\leqno(5.5.3.3)
$$

\setcounter{paragraph}{3}
\paragraph{}\etiquette[5.5.3.4]{XVII.par5.5.3.4}
若 $T$ \emph{为 $\leq n$ 次的}, \ie 第 $(n+1)$ 个“交叉效应”（cross-effect）为零 (\textsc{Dold-Puppe} \cite[4.18 \sisi{pg.232}{p.~232}]{sga4.exp.XVII.bib.DP})\sisi{}{,} 则对每个在 $i>m$ 时满足 $K^{i}=0$ 的 $K\in \Ob C^{\geq 0}(\Aa)$, 当 $i>nm$ 时有 $[T(K)]^{i}=0$ (\textsc{Dold-Puppe} \cite[4.2\sisi{.}{}3 \sisi{pg.233}{p.~233}]{sga4.exp.XVII.bib.DP}).

\paragraph{}\etiquette[5.5.3.5]{XVII.par5.5.3.5}
特别地, 若 $T$ 为从 $A$-模范畴到 $B$-模范畴的有限次函子, 并将平坦对象映为平坦对象, 则有界平坦复形的像仍是有界平坦复形.

\begin{enonce*}{引理 5.5.4}\etiquette[5.5.4]{XVII.lem5.5.4}
设 $\Aa$ 为阿贝尔范畴.
\begin{enumerate}
\renewcommand{\theenumi}{\roman{enumi}}
\renewcommand{\labelenumi}{(\theenumi)}
\item $K^{\geq 0}(\Aa)$ 中的拟同构容许左分式演算与右分式演算.

\item 显然的函子给出一个等价, 从 $K^{\geq 0}(\Aa)$ 中将拟同构变为可逆态射所得的范畴到范畴 $D^{\geq 0}(\Aa)$.
\end{enumerate}
\end{enonce*}

函子\pageoriginale $\tau_{\geq 0}:C(\Aa)\rightarrow
C^{\geq 0}(\Aa)$ (\hyperref[XVII.def1.1.13]{1.1.13}) 下降到商范畴后定义函子 $\tau_{\geq 0}:K(\Aa)\rightarrow K^{\geq
0}(\Aa)$, 它是 $K^{\geq
0}(\Aa)$ 到 $\sisi{\Aa}{K(\Aa)}$ 的包含函子的左伴随. 函子 $\tau_{\geq 0}$ 将拟同构映为拟同构, 并且伴随箭头 $X\rightarrow \tau_{\geq 0}(X)$ 为拟同构, 当且仅当在 $i<0$ 时 $H^{i}(X)=0$. 因而引理 \hyperref[XVII.lem5.5.4]{5.5.4} 来自 \textsc{Verdier} \cite{sga4.exp.XVII.bib.CD} 以及下述引理的 (i) 和 (ii); 其证明留给读者:

\begin{enonce*}{引理 5.5.4.1}\etiquette[5.5.4.1]{XVII.lem5.5.4.1}
设 $K$ 为范畴, $K'$ 为满子范畴, $S$ 为 $K$ 中的一组态射, 它等于 $K(S^{-1})$ 的同构集合的逆像 (\cf \cite[\sisi{chap 1}{\chap~1}]{sga4.exp.XVII.bib.GZ}); $S'$ 为它在 $K'$ 中的迹. 假设 $K'$ 到 $K$ 的包含函子 $i$ 有左伴随 $\tau$, 且 $\tau(S)\subset S'$.
\begin{enumerate}
\renewcommand{\theenumi}{\roman{enumi}}
\renewcommand{\labelenumi}{(\theenumi)}
\item 从 $K'({S'}^{-1})$ 到 $K(S^{-1})$ 的显然函子为全忠实函子, 并以 $\tau$ 为左伴随.

\item 若 $K(S^{-1})$ 容许左分式演算 (\resp 右分式演算), 则 $K'({S'}^{-1})$ 也具有相同性质.

若包含函子 $i$ 未必有左伴随, 且以 $\tau$ 表示与 $i$ 左伴随的\sisi{部分}{不处处}定义的函子, 则有:

\item 若 $K(S^{-1})$ 容许左分式演算 (\resp 右分式演算), 并且对每个终点 (\resp 起点) 在 $K'$ 中的 $u\in S$, $\tau(u)$ 都有定义且属于 $S'$, 则 $K'({S'}^{-1})$ 容许左分式演算 (\resp 右分式演算), 且函子 $K'({S'}^{-1})\rightarrow K(S^{-1})$ 全忠实.
\end{enumerate}
\end{enonce*}

\setcounter{subsubsection}{4}
\subsubsection{}\etiquette[5.5.5]{XVII.subsubsec5.5.5}
设\pageoriginale $T:\Aa\rightarrow \Bb$ 为阿贝尔范畴之间的函子. 如 \hyperref[XVII.def1.2.1]{1.2.1} 中那样, 由公式
\begin{align*}
LT(K) &= \quotprojlim_{s:K'\rightarrow K}T(K')\\
RT(K) &= \quotinjlim_{s:K\rightarrow K'}T(K'){\quoi},
\end{align*}
定义 $T$ 的导出函子
\begin{align*}
LT &: D^{\geq 0}(\Aa)\longrightarrow \Pro(D^{\geq
0}(\Bb))\\
RT &: D^{\geq 0}(\Aa)\longrightarrow \Ind(D^{\geq 0}(\Bb))
\end{align*}
其中 $s$ 遍历终点 (\resp 起点) 为 $K$ 的拟同构所成的滤过范畴.

若 $LT(K)$ $(RT(K))$ 可表, 则称 $T$ \emph{在 $K$ 处左 (\resp 右) 可导}. 若下述典范箭头为同构, 则称 $K$ \emph{左 (\resp 右) 适合于} $T$ (或滥用术语, 称其适合于 $LT$ (\resp $RT$)):
\[
LT(K)\isomto  TK\text{ (\resp $TK \isomto  RT(K)$)}{\quoi}.
\]

\begin{enonce*}{命题 5.5.6}\etiquette[5.5.6]{XVII.prop5.5.6}
设 $F$ 为从环 $A$ 上的模范畴到阿贝尔范畴 $\Bb$ 的函子. 假设 $\Bb$ 中的滤过归纳极限可表且正合, 而 $F$ 与滤过归纳极限可交换. 则每个平坦模的有界复形 $M^{\bigdot}\in \Ob K^{\geq
0}(A)$ 都适合于 $LF$.
\end{enonce*}

命题\pageoriginale 直接来自定义、下述将于 \hyperref[XVII.par5.5.6.4]{5.5.6.4} 证明的引理, 以及下述事实: 分量平坦的 $N^{\bigdot}\in \sisi{B}{C}^{b,\geq 0}(A)$ 所给出的拟同构 $u:N^{\bigdot}\rightarrow M^{\bigdot}$, 在终点为 $M^{\bigdot}$ 的拟同构范畴中构成一个共尾集合.

\begin{enonce*}{引理 5.5.6.1}\etiquette[5.5.6.1]{XVII.lem5.5.6.1}
设 $u:N^{\bigdot}\rightarrow M^{\bigdot}$ 为平坦复形之间的拟同构, 这些复形是 $C^{b,\geq 0}(A)$ 的对象. 则 $F(u)$ 为拟同构.
\end{enonce*}

\begin{enonce*}{引理 5.5.6.2}\etiquette[5.5.6.2]{XVII.lem5.5.6.2}
设 $\Aa$ 为加性范畴, $d$ 为 $\ZZ$ 中的上有界区间. 以 $C^{d}$ (\resp $C^{d,b}$, \resp $K^{d}$, \resp $K^{d,b}$) 表示在次数 $i\not\in d$ 处为零的复形 (\resp 且有界, \resp 模去同伦) 所成的范畴. 有一个等价
$$
\sisi{\sheafell }{\ell}:\Ind(C^{d,b}(\Aa)\sisi{}{)}\isomto
C^{d}(\Ind\Aa){\quoi};
$$
这一等价的逆诱导函子
$$
\sisi{\sheafell }{\ell}':K^{d}(\Ind\Aa)\longrightarrow
\Ind(K^{d,b}(\Aa)){\quoi}.
$$
\end{enonce*}

\begin{enumerate}
\renewcommand{\theenumi}{\Alph{enumi}}
\renewcommand{\labelenumi}{\theenumi)}
\item 证明 $\sisi{\sheafell }{\ell}$ 全忠实. 设 $K$ 为 $\Aa$ 中对象的有界复形, $L=\quotinjlim
L_{\alpha}$ 为 \iInd-复形. 有
\begin{align*}
\Hom(K,L)\overset{\dfn}{=} \varinjlim \Hom(K,L_{\alpha})  = & \varinjlim
Z^{0}(\Hom^{\bigdot} (K,L_{\alpha}))\simeq Z^{0}
\varinjlim_{\alpha} \Hom^{\bigdot} (K,L_{\alpha})\\
& \isomto
Z^{0} \Hom^{\bigdot} (K,\ell L)=\Hom(K,\ell L){\quoi},
\end{align*}
倒数第二个同构来自 $K$ 有限这一事实, 因而 $\varinjlim$ 与从双复形过渡到全复形可交换. 断言得证.

\item 我们\pageoriginale 将验证 $\sisi{\sheafell }{\ell}$ 本质满射的 (颇繁琐的) 工作留给读者: 从任意选定的次数开始, 从左向右构造映到一个上有界 \iInd-对象复形 $K$ 的“足够多”的有界复形; 同样, 留给读者展开如下论证:

\item 若 $\ell f:\ell \varinjlim K_{\alpha}\rightarrow \ell \varinjlim
L_{\beta}$ 零伦, $f=dH+Hd$, 则对每个 $\alpha$, $H$ 都提升为 $H^{i}_{\alpha}:K^{i}_{\alpha}\rightarrow
L^{i\sisi{}{-1}}_{\beta(\alpha)}$, 而当 $\beta$ 足够大时, 它给出 $f_{\alpha}:K_{\alpha}\rightarrow L_{\beta(\alpha)}$ 的零伦.
\end{enumerate}

\setcounter{subsubsection}{6}
\setcounter{paragraph}{2}
\paragraph{}\etiquette[5.5.6.3]{XVII.par5.5.6.3}
设 $A$ 为环, $d$ 为 $\ZZ$ 中的上有界区间, $P(A)$ 为有限投射 $A$-模的范畴, $\Pl(A)$ 为平坦 $A$-模的范畴. 以 $D^{\Tor\subset d}(A)$ 表示 $C^{d}(\Pl(A))$ 在 $D(A)$ 中的像, 以 $D^{\parf\subset d}(A)$ 表示 $C^{d,b}(P(A))$ 在 $D(A)$ 中的像. 由 \hyperref[XVII.par5.5.1.2]{5.5.1.2}, 有 $\Pl(A)\simeq \Ind
P(A)$. 另一方面, $K^{d,b}(P(A))\simeq D^{\parf\subset d}(A)$. 于是由 \hyperref[XVII.lem5.5.6.2]{5.5.6.2} 以及 $H^{i}$ 与滤过归纳极限可交换这一事实, 存在函子的交换图
$$
\xymatrix@C=2.25cm{
C^{d}(\text{Pl}(A)) \ar[d]& \Ind C^{d,b}(P(A)) \ar[d]\ar[l]_{\substack{\varinjlim\\\sim}}\\
K^{d}(\text{Pl}(A)) \ar[d]\ar[r]^{\ell'}& \Ind K^{d,b}(P(A)) \ar[d]^{\sisi{}{\simeq}}\\
D^{\Tor\subset d}(A) \ar[r]^{\ell''}& \Ind D^{\parf \subset d}(A).
}
$$

\setcounter{subsubsection}{6}
\paragraph{}\etiquette[5.5.6.4]{XVII.par5.5.6.4}
我们证明\pageoriginale \hyperref[XVII.lem5.5.6.1]{5.5.6.1}. 用 \hyperref[XVII.par5.5.6.3]{5.5.6.3} 的记号, 若 $d$ 为区间 $[0,k]$, 则 (由于 $F$ 与 $\varinjlim$ 可交换) 有交换图
$$
\xymatrix@C=2.5cm{
K^{d}(\text{Pl}(A))          \ar[r]^{\ell'} \ar[d]_{K^{\geq 0}(F)} & \Ind K^d(P(A))  \ar[d]_{\Ind(K^{\geq 0}(F))}             \\
K^{+}(\Bb) \ar[d]_{H^i}                        & \Ind K^{+}(\Bb)     \ar[d]_{\Ind(H^i)}       \\
\Bb                  & \Ind \Bb  \ar[l]_{\varinjlim}.
}
$$
若 $u$ 如 \hyperref[XVII.lem5.5.6.1]{5.5.6.1}, 则由 \hyperref[XVII.par5.5.6.3]{5.5.6.3}, $\ell'(u)$ 为同构, 因而 $H^iK^{\geq0}(F)(u)$ 也为同构, 如所欲证.

\subsubsection{}\etiquette[5.5.7]{XVII.subsubsec5.5.7}
设 $X$ 为概形 $S$ 上的概形, $k$ 为整数. 假设对称幂 $\Sym^k_S(X)$ 有定义. 若读者不愿作这一假设, 可在代数空间范畴中进行 (\textsc{Artin}\sisi{ [1]}{\nde{原文此处引用了一篇未指明的著作, 可能是第 \pageref{XVII.noteArtin1} 页和第 \pageref{XVII.noteArtin2} 页注释中所引的 Artin 的两篇著作之一.}}). 以\emph{交换}环层 $\Aa$ 给 $S$ 赋环. 从 \hyperref[XVII.subsubsec5.5.19]{5.5.19} 开始, 将假设 $\Aa$ 为\emph{挠的}. 滥用记号, 我们\emph{仍以} $\Aa$ 表示 $\Aa$ 在任一 $S$-概形上的逆像.

设 $ℱ$ 为 $X$ 上的模层; $ℱ$ 的 $k$ 份外张量积所成的层 $\mathop{\boxtimes}\limits^kℱ$ 是 $X^k/S$ 上的层, 对称群 $\Sf_k$ 在 $X^k/S$ 上作用时它具有等变性. 取不变量后,\pageoriginale 它定义 $\Sym^k_S(X)$ 上的一个层 $TS^k_{\ext}(ℱ)$, 称为 \emph{$ℱ$ 的外对称张量幂}（external symmetric tensor power）\sisi{}{.}

为了“逐点”论证, 采用如下描述的 $\Sym^k_S(X)$ 的点是有用的.

\paragraph{}\etiquette[5.5.7.1]{XVII.par5.5.7.1}
设 $s$ 为 $S$ 的一个几何点. $\Sym^k_S(X)$ 的一个取值于 $k(s)$ 的点 $x$ 定义 $\Sym^k_S(X)$ 的一个几何点. 以 $(x_1\ldots x_k)$ 表示 $(X/S)^k$ 的一个取值于 $k(s)$ 且像为 $x$ 的点, 设 $y_1\ldots y_\ell$ 为序列 $x_1\ldots x_k$ 中互异的元素, $n_i$ $(1\leq i\leq\ell)$ 为 $y_i$ 在该序列中的重数 $(\sum n_i=k)$.

用 \hyperref[XVII.par5.5.7.1]{5.5.7.1} 的记号, 并以 $\Sf$ 表示对称群, 有
$$\etiquette[5.5.7.2]{XVII.eq5.5.7.2}
TS^{k}_{\ext}(F)_x\simeq (\bigotimes^{\ell}_{1}
(\bigotimes^{n_i} F_{y_i}))^{\prod^{\ell}_{1}\Sf_{n_i}}{\quoi},\leqno(5.5.7.2)
$$
其中张量积在 $\Aa_S$ 上取. 特别地, \emph{当 $F$ 平坦时}, 由 \hyperref[XVII.par5.5.2.5]{5.5.2.5}, 有
$$\etiquette[5.5.7.3]{XVII.eq5.5.7.3}
TS^{k}_{\ext}(F)_x\isomfrom \bigotimes^{\ell}_{1}
\Gamma^{n_i}F_{y_i}{\quoi},\leqno(5.5.7.3)
$$
这里对 $\ell$ 归纳, 利用 $\Gamma^{n_i}F_{y_i}$ 的平坦性 (\hyperref[XVII.par5.5.2.4]{5.5.2.4}) 以及下述事实: 若 $G$ (\resp $H$) 为作用于 $A$-模 $M$\sisi{}{ }(\resp $N$) 的有限群, 且 $M$ 与 $N^H$ 平坦, 则 $M^G\otimes_A N^H\isomto(M\otimes_A N)^{G\times H}$.

\subsubsection{}\etiquette[5.5.8]{XVII.subsubsec5.5.8}
设 $F$ 为 $X$ 上的模, $P$ 为 $F$ 的一个平坦表示:
$$
P:P_1\longrightarrow P_0\longrightarrow F\longrightarrow0{\quoi}.
$$
如 \hyperref[XVII.par5.5.2.5]{5.5.2.5} 中那样定义\pageoriginale $P'_1$, $j_0$ 及 $j_1$; 将 $\Gamma^k_{\ext}(F)$ (即\emph{外 $\Gamma$ 函子}) 定义为一对箭头 $(TS^k_{\ext}(j_0),TS^k_{\ext}(j_1))$ 的余核. 用 \hyperref[XVII.par5.5.7.1]{5.5.7.1} 的记号, 由 \eqhyperref[XVII.eq5.5.7.3]{5.5.7.3} 及 \hyperref[XVII.par5.5.2.5]{5.5.2.5} 得到同构
$$\etiquette[5.5.8.1]{XVII.eq5.5.8.1}
\Gamma^k_{\ext}(F)_x\simeq\bigotimes^{\ell}_{1}
\Gamma^{n_i}(F_{y_i}){\quoi}.\leqno(5.5.8.1)
$$

这一公式表明, $\Gamma^k_{\ext}(F)$ 先验地函子依赖于 $P$, 但在相差唯一同构的意义下只依赖于 $F$.

\subsubsection{}\etiquette[5.5.9]{XVII.subsubsec5.5.9}
存在函子性态射
$$
\alpha:\Gamma^k_{\ext}(F)\longrightarrow TS^k_{\ext}(F){\quoi}.
$$
对 $(\Gamma,\alpha)$ 由下列两项性质刻画:
\begin{enumerate}
\renewcommand{\theenumi}{\roman{enumi}}
\renewcommand{\labelenumi}{(\theenumi)}
\item 当 $F$ 平坦时, $\alpha$ 为同构.

\item 用 \hyperref[XVII.par5.5.2.5]{5.5.2.5} 的记号, 对 $F$ 的每个表示
$$
P_1\xrightarrow{d}P_0\xrightarrow{\epsilon}F\longrightarrow0
$$
序列
$$
\Gamma^k_{\ext}(P'_1)\rightrightarrows
\Gamma^k_{\ext}(P_0)\rightarrow\Gamma^k_{\ext}(F)\rightarrow0
$$
都正合.

此外:

\item 函子 $TS^k$ 与 $\Gamma^k$ 都与滤过归纳极限可交换 (由 \sisi{\hyperref[XVII.eq5.5.7.2]{5.5.7.2}}{\eqhyperref[XVII.eq5.5.7.2]{5.5.7.2}}, \sisi{\hyperref[XVII.eq5.5.8.1]{5.5.8.1}}{\eqhyperref[XVII.eq5.5.8.1]{5.5.8.1}}, \hyperref[XVII.par5.5.2.4]{5.5.2.4}).

\item 函子 $\Gamma^k$ 与底环的每个扩张 $\Aa\rightarrow\Aa'$ 可交换.

\item 函子\pageoriginale $\Gamma^k$ 与每个基变换 $S'\rightarrow S$ 可交换: 对每个笛卡儿方块
$$
\xymatrix@C=2.5cm@R=1.5cm{
X' \ar[r]^{g'}\ar[d]^{f'}& X \ar[d]^{f}\\
S' \ar[r]^{g}& S,
}
$$
有
$$
\Sym^k_g(g')^{\ast}\Gamma^k_{\ext}\simeq
\Gamma^k_{\ext}{g'}^{\ast}{\quoi}.
$$
(\sisi{\hyperref[XVII.eq5.5.8.1]{5.5.8.1}}{\eqhyperref[XVII.eq5.5.8.1]{5.5.8.1}} 及 \hyperref[XVII.par5.5.2.7]{5.5.2.7}).

\item 函子 $\Gamma^k$ 与沿每个\emph{根}态射（radicial morphism）$i:X'\rightarrow X$ (例如闭浸入) 的逆像可交换: 有
$$
\Sym^k_S(i)^{\ast}\Gamma^k_{\ext}\simeq
\Gamma^k_{\ext}i^{\ast}
$$
\eqhyperref[XVII.eq5.5.8.1]{5.5.8.1}.

\item 对\emph{浸入} $j:X'\rightarrow X$, 有
$$
\Sym^k_S(j)_{!}\Gamma^k_{\ext}\simeq\Gamma^k_{\ext}j_{!}
$$
\eqhyperref[XVII.eq5.5.8.1]{5.5.8.1}.
\end{enumerate}

\subsubsection{}\etiquette[5.5.10]{XVII.subsubsec5.5.10}

设 $k_1,k_2$ 及 $k=k_1+k_2$ 为三个整数 $\geq0$. 则有有限态射 (\quad):\nde{我未找到一般情形下 $\vee_{k_1k_2}$ 有限性的参考文献.}
$$
\vee_{k_1k_2}:\Sym^{k_1}_S(X)\times_S\Sym^{k_2}_S(X)
\longrightarrow\Sym^k_S(X){\quoi}.
$$
若 $F_i$ $(i=1,2)$ 为 $\Sym_S^{k_i}(X)$ 上的模, 以 $F_1\vee F_2$ 表示 $\Sym_S^k(X)$ 上的层
$$
F_1\vee F_2=\vee_{k_1k_2\ast}(F_1\boxtimes F_2){\quoi}.
$$

\subsubsection{}\etiquette[5.5.11]{XVII.subsubsec5.5.11}
设\pageoriginale $X$ 上有模的正合列
\[
\Sigma:0\longrightarrow F'\longrightarrow F\longrightarrow
F''\longrightarrow0{\quoi}.
\]
用 \hyperref[XVII.par5.5.7.1]{5.5.7.1} 的记号, 由 \eqhyperref[XVII.eq5.5.8.1]{5.5.8.1} 及 \hyperref[XVII.par5.5.2.6]{5.5.2.6}, $\Gamma^k_{\ext}(F)_x$ 上有一个滤过, 并有态射
\[
\sum_{\sum^{\ell}_{1}i_\alpha=i}
\bigotimes^{\ell}_{1}\Gamma^{i_\alpha}(F'_{y_\alpha})\otimes
\bigotimes^{\ell}_{1}\Gamma^{n_i-i_\alpha}(F''_{y_\alpha})
\longrightarrow\Gr^i(\Gamma^k_{\ext}(F)_x){\quoi}.
\]
第一项恰为 $(\Gamma^i_{\ext}(F')\vee\Gamma^{k-i}_{\ext}(F''))_x$. 容易验证, 这些逐点构造的数据来自 $\Gamma^k_{\ext}(F)$ 上的 (唯一) $k+1$ 步滤过, 以及 (唯一的) 态射
$$\etiquette[5.5.11.1]{XVII.eq5.5.11.1}
\Gamma^i_{\ext}(F')\vee\Gamma^{k-i}_{\ext}(F'')
\longrightarrow\Gr^i(\Gamma^k_{\ext}(F)){\quoi}.\leqno(5.5.11.1)
$$

若正合列 $\Sigma$ 分裂, 更准确地有 (由 \sisi{\hyperref[XVII.eq5.5.2.6.2]{5.5.2.6.2}}{\eqhyperref[XVII.eq5.5.2.6.1]{5.5.2.6.1}} 逐点验证):
$$\etiquette[5.5.11.2]{XVII.eq5.5.11.2}
\Gamma^k_{\ext}(F'\oplus F'')\simeq
\sisi{\sum}{\bigoplus}_{k_1+k_2=k}\Gamma^{k_1}_{\ext}(F')\vee
\Gamma^{k_2}_{\ext}(F''){\quoi}.\leqno(5.5.11.2)
$$

若正合列 $\Sigma$ 逐点分裂, 或者 $F''$ 平坦, 则态射 \eqhyperref[XVII.eq5.5.11.1]{5.5.11.1} 为同构 (\cf \hyperref[XVII.par5.5.1.4]{5.5.1.4}).

\subsubsection{}\etiquette[5.5.12]{XVII.subsubsec5.5.12}
设 $u:X\rightarrow Y$ 为 $S$-概形态射. 假设对称幂 $\Sym^k_S(X)$ 与 $\Sym^k_S(Y)$ 存在. 设 $F$ 为 $X$ 上的模, $G$ 为 $Y$ 上的模, 并设
$v\in\Hom(u^{\ast}G,F)\underset{\sisi{n}{}}{\simeq}\Hom(G,u_{\ast}F)$
为从 $F$ 到 $G$ 的 $u$-同态 \hyperlink{sga4.exp.XVII.subsection.0.2}{(0.2)}. 若 $\prod^n u$ 为 $u$ 的第 $n$ 笛卡儿幂, 则同态 $v$ 诱导一个从 $\mathop{\boxtimes}\limits^nF$ 到 $\mathop{\boxtimes}\limits^nG$ 的 $\prod^n u$-同态 $\mathop{\boxtimes}\limits^n v$. 取不变量后,\pageoriginale 后者定义
\[
TS^n(v):TS^n_{\ext}(F)\xrightarrow[\Sym_S^n(u)]{}TS^n_{\ext}(G){\quoi}.
\]

由定义 \hyperref[XVII.subsubsec5.5.8]{5.5.8}, 存在唯一的 $\Sym^n(u)$-同态 $\Gamma^n_{\ext}(v)$, 关于 $(F,G,v)$ 具有函子性,
$$\etiquette[5.5.12.0]{XVII.eq5.5.12.0}
\Gamma^n_{\ext}(v):\Gamma^n_{\ext}(F)
\xrightarrow[\Sym^n_S(u)]{}\Gamma^n_{\ext}(G)\leqno(5.5.12.0)
$$
使下图交换 (关于箭头的不寻常方向, \cf \hyperlink{sga4.exp.XVII.subsection.0.2}{(0.2)}):
$$
\xymatrix@C=3.5cm@R=2cm{
\Gamma^n_{\ext}(F) \ar[r]^{\Gamma^n_{\ext}(v)}_{\text{在 $\Sym^n_S(u)$ 上}}& \Gamma^n_{\ext}(G)\\
TS^n_{\ext}(F) \ar[r]^{TS^n_{\ext}(v)}_{\text{在 $\Sym^n_S(u)$ 上}}\ar[u]^{\alpha}_{\text{在 $\Id$ 上}} & TS^n_{\ext}(G)\ar[u]^{\alpha}_{\text{在 $\Id$ 上}}.
}
$$

我们将主要把这条箭头用作“对称 Künneth 箭头”: 对 $G=u_{\ast}F$ 以及典范箭头 $v$, 箭头 $\Gamma^n_{\ext}(v)$ 等同于模的同态\sisi{s}{}
$$\etiquette[5.5.12.1]{XVII.eq5.5.12.1}
k^n:\Gamma^n_{\ext}(u_{\ast}F)\longrightarrow
\Sym^n(u)_{\ast}\Gamma^n_{\ext}(F){\quoi}.\leqno(5.5.12.1)
$$

\subsubsection{}\etiquette[5.5.13]{XVII.subsubsec5.5.13}
考虑条件:

\paragraph{}\etiquette[5.5.13.1]{XVII.par5.5.13.1}
$K\in\Ob K^{b,\geq0}(X,\Aa)$, 并且对 $X$ 的每个几何点 $x$, 复形 $K_x$ 同伦等价于一个平坦模复形 $L_x\in\Ob C^{b,\geq0}(\Aa_x)$.

\begin{enonce*}{引理 5.5.13.2}\etiquette[5.5.13.2]{XVII.lem5.5.13.2}
$D^b(X,\Aa)$ 中由 \tTor-维数 $\leq0$ \emph{\hyperref[XVII.def4.1.9]{4.1.9}} 的复形组成的子范畴\pageoriginale $D^{b,\tor\leq0}(X,\Aa)$, 等价于将满足\emph{(\hyperref[XVII.par5.5.13.1]{5.5.13.1})} 的复形 $K\in\Ob K^{b,\geq0}(X,\Aa)$ 的范畴中的拟同构变为可逆态射所得的范畴. 这一构造可以同样用左分式演算或右分式演算进行.
\end{enonce*}

设 $P$ 为 $K^{-}(X,\Aa)$ 中由逐点同伦等价于平坦模复形的复形组成的子范畴. 已知 (\textsc{Verdier} \sisi{\cite{sga4.exp.XVII.bib.CD}}{\cite[I 2.4.2]{sga4.exp.XVII.bib.CD}}), $D^{-}(X,\Aa)$ 是从 $P$ 通过左分式演算或右分式演算将拟同构变为可逆态射而得到的. 将 \hyperref[XVII.lem5.5.4.1]{5.5.4.1} 应用于 $K^{b,\geq0}(X,\Aa)\cap P$ 到 $P$ 的包含函子. 事实上, 若 $K$ 的 \tTor-维数为 $\leq0$, 并逐点同伦等价于平坦模复形 $L_x$, 则 $\tau_{\geq0}(K)$ 逐点同伦等价于 $\tau_{\geq0}(L_x)$, 而后者仍为平坦模复形, 因而 $\tau_{\geq0}(K)$ 满足 \hyperref[XVII.par5.5.13.1]{5.5.13.1}.

\subsubsection{}\etiquette[5.5.14]{XVII.subsubsec5.5.14}
由 \hyperref[XVII.lem5.5.6.1]{5.5.6.1}, \sisi{\hyperref[XVII.eq5.5.7.2]{5.5.7.2}}{\eqhyperref[XVII.eq5.5.7.2]{5.5.7.2}} 或 \sisi{\hyperref[XVII.eq5.5.8.1]{5.5.8.1}}{\eqhyperref[XVII.eq5.5.8.1]{5.5.8.1}} 以及 \hyperref[XVII.par5.5.2.4]{5.5.2.4}, 满足 (\hyperref[XVII.par5.5.13.1]{5.5.13.1}) 的复形之间的每个拟同构 $v:K'\rightarrow K$ 都诱导拟同构
\begin{align*}
TS^k(v)&:TS^k_{\ext}(K')\xrightarrow{\approx}TS^k_{\ext}(K)\\
\Gamma^k(v)&:\Gamma^k_{\ext}(K')\xrightarrow{\approx}\Gamma^k_{\ext}(K){\quoi}.
\end{align*}
因此函子 $\Gamma^k$ 与 $TS^k$ 在 $D^b(X,\Aa)$ 的子范畴 $D^{b,\tor\leq0}(X,\Aa)$ 上可导 (\hyperref[XVII.subsubsec5.5.5]{5.5.5}), 而满足 \hyperref[XVII.par5.5.13.1]{5.5.13.1} 的复形左适合于这些函子. 由 \hyperref[XVII.par5.5.2.4]{5.5.2.4},\pageoriginale $\Gamma^k$ 与 $TS^k$ 诱导函子
\begin{align*}
LTS^k &:D^{b,\tor\leq0}(X,\Aa)\longrightarrow
D^{b,\tor\leq0}(\Sym^k_S(X),\Aa)\tag{5.5.14.1}\\
L\Gamma^k &:D^{b,\tor\leq0}(X,\Aa)\longrightarrow
D^{b,\tor\leq0}(\Sym^k_S(X),\Aa){\quoi}.\tag{5.5.14.2}
\end{align*}\etiquette[5.5.14.1]{XVII.eq5.5.14.1}
\etiquette[5.5.14.2]{XVII.eq5.5.14.2}% 删除此注释会引入不需要的空白.
此外, 由 \hyperref[XVII.subsubsec5.5.9]{5.5.9}~(i), 典范态射 $\alpha$ (\hyperref[XVII.subsubsec5.5.9]{5.5.9}) 在函子 \eqhyperref[XVII.eq5.5.14.1]{5.5.14.1} 与 \eqhyperref[XVII.eq5.5.14.2]{5.5.14.2} 之间诱导同构:
$$\etiquette[5.5.14.3]{XVII.eq5.5.14.3}
\alpha:L\Gamma^k\isomto LTS^k{\quoi}.\leqno(5.5.14.3)
$$

\setcounter{paragraph}{3}
\paragraph{}\etiquette[5.5.14.4]{XVII.par5.5.14.4}
\hyperref[XVII.subsubsec5.5.9]{5.5.9}~(iv), (v), (vi) 及 (vii) 中的同构过渡到导出范畴后定义同构:
\begin{itemize}
\item[${\rm(iv')}$] 对 $\Aa$-代数 $\Aa'$, 有
$$
L\Gamma^k_{\ext_{\Aa}}(K)\overset{\LL}{\otimes}\Aa'\simeq
L\Gamma^k_{\ext_{\Aa'}}(K\overset{\LL}{\underset{\Aa}{\otimes}}\Aa'){\quoi}.
$$

\item[${\rm(v')}$] 对每个基变换 $g:S'\rightarrow S$, 有
$$
\Sym^k_g(g')^{\ast}L\Gamma^k_{\ext}K\simeq
L\Gamma^k_{\ext}({g'}^{\ast}K){\quoi}.
$$

\item[${\rm(vi')}$] 对每个\emph{根}态射 $i:X'\rightarrow X$, 有
$$
\Sym^k_S(i)^{\ast}L\Gamma^k_{\ext}K\simeq
L\Gamma^k_{\ext}i^{\ast}K{\quoi}.
$$

\item[${\rm(vii')}$] 对\emph{浸入} $j:X'\rightarrow X$, 有
$$
\Sym^k_S(j)_{!}L\Gamma^k_{\ext}\simeq
L\Gamma^k_{\ext}j_{!}{\quoi}.
$$
\end{itemize}

\subsubsection{}\etiquette[5.5.15]{XVII.subsubsec5.5.15}
设 $C^{\geq0}(X,\Aa)$ 中有平坦复形的正合列,
$$
0\longrightarrow K'\longrightarrow K\longrightarrow K''\longrightarrow0{\quoi},
$$
并应用函子 $\Gamma^k_{\ext}$ (\hyperref[XVII.subsubsec5.5.3]{5.5.3}). 可见 (如 \hyperref[XVII.subsubsec5.5.11]{5.5.11} 中一样), $\Sym^k_S(X)$ 上的半单纯模 $\Gamma^k_{\ext}(N^{-1}K)$ 有一个滤过, 其\pageoriginale 逐次商为半单纯模 $\Gamma^{k'}_{\ext}(N^{-1}K')\vee
\Gamma^{k''}_{\ext}(N^{-1}K'')$ $(k'+k''=k)$.

过渡到导出范畴可见, $L\Gamma^k_{\ext}(K)$ 由一个滤过复形表示, 其逐次商表示导出范畴的对象 $L\Gamma^i_{\ext}(K'\sisi{}{)}\overset{\LL}{\vee}
L\Gamma^{k-i}_{\ext}(K'')$ (应用 \hyperref[XVII.par5.5.2.4]{5.5.2.4} 及 \textsc{Eilenberg--Zilber--Cartier} 定理, 即双半单纯对象与其对角对象在 $\mathcal K$ 中同构: \textsc{Dold\sisi{ }{-}Puppe} \cite[2.9 \sisi{pg.213}{p.~213}]{sga4.exp.XVII.bib.DP}).

在导出范畴中, 从满足 \hyperref[XVII.par5.5.13.1]{5.5.13.1} 的复形的一个正合列出发, 已经可以进行这一构造.

${}^{\ast}$ 用更内在的术语说, 设 $\Delta$ 为 $D^{b,\tor\leq0}(X,\Aa)$ 的对象的一个“\emph{真三角}”, \ie 为“滤过导出范畴”的一个对象, 在 $i\neq0,1$ 时满足 $\Gr^i(\Delta)=0$, 且 $\Gr^i(\Delta)\in\Ob D^{b,\tor\leq0}(X,\Aa)$ (\textsc{Deligne}, 见 \textsc{Verdier} \cite{sga4.exp.XVII.bib.V2} 的附录, 或 \refvb[VI B nᵒ]{Vbis.subsubsec2.5.1}, \refnoexpose{Vbis.subsubsec2.5.2}).

前述构造于是给出 $\Sym^k_S(X)$ 的滤过导出范畴的对象 $L\Gamma^k_{\ext}(\Delta)$、同构
\[
\Gr^i(L\Gamma^k_{\ext}(\Delta))\simeq
L\Gamma^i_{\ext}(\Gr^1(\Delta))\overset{\LL}{\vee}
L\Gamma^{k-i}_{\ext}(\Gr^0(\Delta)\sisi{}{)}
\]
以及一个同构 (以 $\lvert\quad\rvert$ 表示滤过复形的底层复形)
\[
\lvert L\Gamma^k_{\ext}(\Delta)\rvert\simeq
L\Gamma^k_{\ext}(\lvert\Delta\rvert){\quoi}.{}_{\ast}
\]

\begin{enonce*}{引理 5.5.16}\etiquette[5.5.16]{XVII.lem5.5.16}
设 $N$ 为整数. 范畴 $D^{b,\tor\leq0}(X,\Aa)$ 可由满足下列条件 {\rm(i)} 和 {\rm(ii)} 的复形 $K\in\Ob K^{b,\geq0}(X,\Aa)$ 的范畴, 通过将拟同构变为可逆态射而得到:
\begin{enumerate}
\renewcommand{\theenumi}{\roman{enumi}}
\renewcommand{\labelenumi}{(\theenumi)}
\item 条件\pageoriginale (\hyperref[XVII.par5.5.13.1]{5.5.13.1}).

\item 复形 $K$ 适合于每个关于拓扑斯态射 $f:\widetilde X_{\text{\rm et}}\rightarrow T$ 的上同调维数 $\leq N$ 的正像函子.
\end{enumerate}
\end{enonce*}

右分式演算的存在性来自 \hyperref[XVII.subsubsec5.5.13]{5.5.13} 以及以下事实: 利用截断的典范松软消解（canonical flasque resolution）(\hyperref[XVII.subsubsec4.2.9]{4.2.9}), 每个满足 (i) 的复形 $K$ 都映到一个满足 (i) 和 (ii) 的复形.

\setcounter{subsubsection}{16}

\subsubsection{}\etiquette[5.5.17]{XVII.subsubsec5.5.17}
设 $u:X\rightarrow Y$ 为 $S$-概形态射. 假设 $\Sym^n_S(X)$ 与 $\Sym^n_S(Y)$ 存在, 存在整数 $N$ 使函子 $u_{\ast}$ 在 $\Aa$-模范畴上的维数为 $\leq N$, 且函子 $Ru_{\ast}$ 将 \tTor-维数 $\leq0$ 的有界复形映为 \tTor-维数 $\leq0$ 的复形. 若 $u$ 固有, $\Aa$ 为挠的, 且 $Y$ 凝聚 (\ie 拟紧且拟分离), 则后两项条件自动满足 (\hyperref[XVII.cor5.2.8.1]{5.2.8.1}, \hyperref[XVII.thm5.2.10]{5.2.10}).\nde{$Y$ 的凝聚性可以换成较弱的条件, 即 $u$ 的纤维维数有界.}

设 $K\in\Ob K^{b,\geq0}(X,\Aa)$ 满足 \hyperref[XVII.lem5.5.16]{5.5.16} 的条件 (i) 与 (ii), 其中 $N$ 如上. 由 \hyperref[XVII.subsubsec5.5.14]{5.5.14}, 箭头
$$
\beta:L\Gamma^n_{\ext}(K)\isomto\Gamma^n_{\ext}(K)
$$
为拟同构. 同样, 箭头
$$
\gamma:u_{\ast}K\isomto Ru_{\ast}K
$$
为拟同构.

将\pageoriginale 从 $L\Gamma^n_{\ext}Ru_{\ast}(K)$ 到 $R\Sym^n(u)_{\ast}L\Gamma^n_{\ext}(K)$ 的\emph{对称 Künneth 箭头}定义为复合态射
\begin{multline*}
k^{n}: L\Gamma^{n}_{\ext}Ru_{\ast}K\xleftarrow[\gamma]{\sim}
L\Gamma^{n}_{\ext}u_{\ast}K\rightarrow
\Gamma^{n}_{\ext}u_{\ast}K\xrightarrow[\protect\hyperlink{sga4.semantic.XVII.eq5.5.12.1}{5.5.12.1}]{}\Sym^{n}_{S}(u)_{\ast}\Gamma^{n}_{\ext}(K)\\
\longrightarrow R
\Sym^{n}_{S}(u)_{\ast}\Gamma^{n}_{\ext}(K)\xleftarrow[\beta]{\sim} R \Sym^{n}_{S}(u)_{\ast}L\Gamma^{n}_{\ext}(K){\quoi}.\tag{5.5.17.1}
\end{multline*}\etiquette[5.5.17.1]{XVII.eq5.5.17.1}% 本段之后的竖直间距过大.

由 \hyperref[XVII.lem5.5.16]{5.5.16}, 态射 \eqhyperref[XVII.eq5.5.17.1]{5.5.17.1} 定义从 $D^{b,\tor\leq0}(X,\Aa)$ 到 $D^{+}(\Sym^n_S(Y),\Aa)$ 的函子之间的态射:
$$\etiquette[5.5.17.2]{XVII.eq5.5.17.2}
k^n:L\Gamma^n_{\ext}Ru_{\ast}\longrightarrow
R\Sym^n_S(u)_{\ast}L\Gamma^n_{\ext}{\quoi},\leqno(5.5.17.2)
$$
也称为\emph{对称 Künneth 态射}.

这一态射由下述命题 \hyperref[XVII.prop5.5.18]{5.5.18} 刻画.

\setcounter{paragraph}{2}
\paragraph{}\etiquette[5.5.17.3]{XVII.par5.5.17.3}
设 $S$, $X$, $Y$ 及 $u$ 如上, $K\in\Ob C^{b,\geq0}(X,\Aa)$ 为 \tTor-维数 $\leq0$ 的复形, $K_1\in\Ob C^{b,\geq0}(Y,\Aa)$ 且 $K_2\in\Ob C^b(\Sym^n_S(X),\Aa)$. 设 $v$ 为从 $K$ 到 $K_1$ 的 $u$-同态 \hyperlink{sga4.exp.XVII.subsection.0.2}{(0.2)}, $t$ 为从 $\Gamma^n_{\ext}(K)$ 到 $K_2$ 的同态 (是通常意义的同态, 以下所有同态也如此). 有态射
\begin{enumerate}
\renewcommand{\theenumi}{\alph{enumi}}
\renewcommand{\labelenumi}{\theenumi)}
\itemsep=8pt
\item $\Gamma^n_{\ext}(v):\Gamma^n_{\ext}(K_1)\rightarrow
\Sym^n_S(u)_{\ast}\Gamma^n_{\ext}(K)$\qquad
\eqhyperref[XVII.eq5.5.12.0]{5.5.12.0}

\item
$\Sym^n_S(u)_{\ast}(t):\Sym^n_S(u)_{\ast}\Gamma^n_{\ext}(K)
\rightarrow\Sym^n_S(u)_{\ast}K_2$,
\end{enumerate}
它们的复合
$$
\varphi':\Gamma^n_{\ext}(K_1)\longrightarrow
\Sym^n_S(u)_{\ast}K_2{\quoi},
$$
通过复合定义
$$
\varphi:L\Gamma^n_{\ext}(K_1)\longrightarrow
R\Sym^n(u)_{\ast}K_2{\quoi}.
$$
此外,\pageoriginale $v$ 与 $t$ 诱导
$$
\begin{cases}
v':K_1\longrightarrow Ru_{\ast}K\\
t':L\Gamma^n_{\ext}K\longrightarrow K_2{\quoi}.
\end{cases}
$$

于是:

\begin{enonce*}{命题 5.5.18}\etiquette[5.5.18]{XVII.prop5.5.18}
在 \hyperref[XVII.par5.5.17.3]{5.5.17.3} 的假设下, 下图在 $D^{+}(\Sym^n_S(Y),\Aa)$ 中交换:
\[
\xymatrix@C=4cm@R=1.75cm{
  L\Gamma^n_{\ext}Ru_{\ast}K \ar[r]^{k^n}_{\text{\textsuperscript{(对称 Künneth)}}}& R\Sym^n_S(u)_{\ast}L
\Gamma^n_{\ext}K \ar[d]_{R
  \Sym^n_S(u)_{\ast}(t')}\\
L\Gamma^n_{\ext}K_1 \ar[r]^{\varphi}\ar[u]^{L\Gamma^n_{\ext}(v')} & R\Sym^n_S(u)_{\ast}K_2.
}
\]
\end{enonce*}

证明与 \hyperref[XVII.par5.4.1.5]{5.4.1.5} 的证明完全类似, 留给读者.

\setcounter{subsubsection}{18}
\subsubsection{从现在起假设 $\Aa$ 为挠的}\etiquette[5.5.19]{XVII.subsubsec5.5.19}

设 $u:X\rightarrow Y$ 为可紧化的 $S$-概形态射:
\[
\xymatrix@C=3cm@R=1.75cm{
   & \bar{X} \ar[rd]^{\bar{u}}&\\
X \ar@{^{(}->}[ru]^{j} \ar[rr]^{u} \ar[rd]&         & Y \ar[ld]\\
   &   S     & .
}
\]
假设\pageoriginale $\Sym^n_S(X)$, $\Sym^n_S(\bar X)$ 及 $\Sym^n_S(Y)$ 存在. 对 $X$ 上的每个模 $F$, 都有
$$\etiquette[5.5.19.1]{XVII.eq5.5.19.1}
\Gamma^n_{\ext}(j_{!}F)\simeq
\Sym^n_S(j)_{!}\Gamma^n_{\ext}(F){\quoi},\leqno(5.5.19.1)
$$
这一同构立即过渡到导出范畴 \sisi{}{(\hyperref[XVII.par5.5.14.4]{5.5.14.4}~(vii'))}. 将从 $D^{b,\tor\leq0}(X,\Aa)$ \linebreak 到 $D^{b,\tor\leq0}(\Sym^n_S(Y),\Aa)$\nde{一般情况下如何证明 $\Sym^n_S(u)$ 为有限型? 这一条件是 $R\Sym^n_S(u)_!$ 有定义的必要条件.} 的函子之间的\emph{对称 Künneth 态射}
$$\etiquette[5.5.19.2]{XVII.eq5.5.19.2}
k^n_u\quad\mbox{\rm 或}\quad
k^n:L\Gamma^n_{\ext}Ru_{!}\longrightarrow
R\Sym^n_S(u)_{!}L\Gamma^n_{\ext}\leqno(5.5.19.2)
$$
定义为复合态射
\begin{align*}
L\Gamma^n_{\ext}Ru_{!}\stackrel{\dfn}{\simeq}L\Gamma^n_{\ext}R\bar u_{\ast}j_{!}
&\xrightarrow{\eqhyperref[XVII.eq5.5.17.2]{5.5.17.2}}
R\Sym^n_S(\bar u)_{\ast}L\Gamma^n_{\ext}j_!\\
&\xrightarrow[\sisi{}{\text{(\hyperref[XVII.par5.5.14.4]{5.5.14.4} (vii'))}}]{\sisi{\substack{\eqhyperref[XVII.eq5.5.19.1]{5.5.19.1}\\\sim}}{\sim}}
R\Sym^n_S(\bar u)_*\Sym^n_S(j)_{!}L\Gamma^n_{\ext}\\
&\xrightarrow{\sisi{\text{(5.5.13.5 (vii'))}}{}}
R\Sym^n_S(u)_{!}L\Gamma^n_{\ext}.
\end{align*}

\setcounter{subsubsection}{19}
\subsubsection{}\etiquette[5.5.20]{XVII.subsubsec5.5.20}
以下验证留给读者, 利用 \hyperref[XVII.prop5.5.18]{5.5.18}, 原则上它们都是容易的.

\paragraph{}\etiquette[5.5.20.1]{XVII.par5.5.20.1}
(内在性). 态射 \eqhyperref[XVII.eq5.5.19.2]{5.5.19.2} 与 $u$ 的紧化 $\bar u$ 的选取无关.

\paragraph{}\etiquette[5.5.20.2]{XVII.par5.5.20.2}
(传递性). 设有 $S$-概形的图
$$
X\xrightarrow{u}Y\xrightarrow{v}Z{\quoi},
$$
其中\pageoriginale $u$ 与 $v$ 拟射影; 假设 $\Sym^n_S(Z)$ 存在 (从而 $\Sym^n_S(Y)$ 与 $\Sym^n_S(X)$ 也存在 (\SGA 1 V 1.7, \EGA II 4.5.4)). 则下列函子图交换:
$$
\xymatrix@C=1.5cm@R=1.75cm{
L\Gamma^n_{\ext}R(vu)_{!} \ar@{=}[d] \ar[rr]^{k^n_{vu}}& & R\Sym^n_S(vu)_{!}L\Gamma^n_{\ext}\ar@{=}[d]\\
L\Gamma^n_{\ext}Rv_{!}Ru_{!} \ar[r]^{k^n_v}& R\Sym^n_S(v)_!
L\Gamma^n_{\ext}Ru_{!} \ar[r]^{k^n_u}& R\Sym^n_S(v)_{!}R\Sym^n_S(u)_{!}L\Gamma^n_{\ext}.
}
$$

\paragraph{}\etiquette[5.5.20.3]{XVII.par5.5.20.3}
对称 Künneth 态射与 \hyperref[XVII.subsubsec5.5.9]{5.5.9}~(iv), (v) 及 (vi) 类型的底或空间变换相容. 更准确地, 与态射 \sisi{5.5.13.5}{\hyperref[XVII.par5.5.14.4]{5.5.14.4}} ${\rm(iv')}$, ${\rm(v')}$ 及 ${\rm(vi')}$ 有如下相容性.
\begin{itemize}
\item[${\rm(iv'')}$] 对每个 $\Aa$-代数 $\Aa'$ 以及每个拟射影 $S$-态射 $u:X\rightarrow Y$, 图
\[\def\objectstyle{\scriptstyle}
\xymatrix@C=.8cm@R=1.75cm{
L\Gamma^n_{\ext}Ru_{!}(K)\overset{\LL}{\otimes}_{\Aa}\Aa'
\ar@{-}[r]^{\sim} \ar[d]^{k^n_u}&
L\Gamma^n_{\ext}(Ru_{!}(K)\overset{\LL}{\otimes}_{\Aa}\Aa')
\ar@{->}[r]^{\sim}&
L\Gamma^n_{\ext}Ru_{!}(K\overset{\LL}{\otimes}_{\Aa}\Aa') \ar[d]^{k^n_u}\\
R\Sym^n_S(u)_{!}L\Gamma^n_{\ext}(K)\overset{\LL}{\otimes}_{\Aa}\Aa'
\ar@{->}[r]^{\sim}&
R\Sym^n_S(u)_{!}(L\Gamma^n_{\ext}(K)\overset{\LL}{\otimes}_{\Aa}\Aa')
\ar@{-}[r]^{\sim}&
R\Sym^n_S(u)_{!}L\Gamma^n_{\ext}(K\overset{\LL}{\otimes}_{\Aa}\Aa')
}
\]
交换.

\item[${\rm(v'')}$] 对\pageoriginale 每个基变换 $g:S'\rightarrow S$, 它给出一个\sisi{笛卡儿}{方块均为笛卡儿方块的}图
$$
\xymatrix{
X' \ar[rrr]^{g'}\ar[rd]^{u'}\ar[dd]&   & & X \ar[rd]^{u}\ar[dd]|\hole &   \\
   & Y' \ar[ld]\ar[rrr]^{g''}&  & & Y \ar[ld]\\
S' \ar[rrr]_{g}&  &  & S &   ,
}
$$
完全类似于 ${\rm(iv'')}$ 的函子图交换:
$$\def\objectstyle{\scriptstyle}
\xymatrix@R=2.25cm{
\Sym^n_g(g'')^{\ast}L\Gamma^n_{\ext}Ru_{!} \ar@{-}[r]^{\sim} \ar[d]^{k^n_u}&
L\Gamma^n_{\ext}g''^{\ast}Ru_{!} \ar[r]^{\sim}& L\Gamma^n_{\ext}Ru'_{!}
g'^{\ast}\ar[d]^{k^n_{u'}} \\
\Sym^n_g(g'')^{\ast}R\Sym^n_S(u)_{!}L\Gamma^n_{\ext} \ar[r]^{\sim}& R
\Sym^n_S(u')_{!}\Sym^n_g(g')^{\ast}L\Gamma^n_{\ext} \ar@{-}[r]^\sim& R\Sym^n_S(u')_{!}L\Gamma^n_{\ext}g'^{\ast}.
}
$$

\item[${\rm(vi'')}$] 对每个\emph{根}基变换 $g:Y'\rightarrow Y$, 它给出 $S$-概形的笛卡儿图
$$
\xymatrix@C=2.5cm@R=1.5cm{
X' \ar[r]^{g'}\ar[d]^{u'}& X \ar[d]^{u}\\
Y' \ar[r]^{g}& Y,
}
$$
下列函子图交换:
\[
\xymatrix{
\Sym^n_S(g)^{\ast}L\Gamma^n_{\ext}Ru_{!} \ar@{-}[r]^{\sim} \ar[d]^{k^n_u}&
L\Gamma^n_{\ext}g^{\ast}Ru_{!} \ar[r]^{\sim}& L\Gamma^n_{\ext}
Ru'_{!}g'^{\ast} \ar[d]^{k^n_{u'}}\\
\Sym^n_S(g)^{\ast}R\Sym^n_S(u\sisi{_!)}{)_!L\Gamma^n_{\ext}} \ar[r]^{\sim}& R
\Sym^n_S(u')_{!}{g'}^{\ast}L\Gamma^n_{\ext} \ar@{-}[r]^{\sim}& R\Sym^n_S(u')_{!}L\Gamma^n_{\ext}g'^{\ast}
}
\]
\end{itemize}

\begin{enonce*}{定理 5.5.21}\etiquette[5.5.21]{XVII.thm5.5.21}
{\rm(对称\pageoriginale Künneth 公式).} 设 $S$ 为凝聚概形 (= 拟紧且拟分离), 以交换挠环层 $\Aa$ 赋环, 且其茎均为 Noether 环. 以 $\Aa$ 的逆像给每个 $S$-概形赋环. 设 $u:X\rightarrow Y$ 为拟射影 $S$-概形之间的态射. 对每个\sisi{}{ }整数 $n\geq0$ 以及每个 $K\in\Ob D^{b,\tor\leq0}(X,\Aa)$, 对称 Künneth 箭头 \eqhyperref[XVII.eq5.5.17.2]{5.5.17.2}
$$
k^n:L\Gamma^n_{\ext}Ru_{!}K\longrightarrow
R\Sym^n_S(u)_{!}L\Gamma^n_{\ext}K
$$
是 $D^{b,\tor\leq0}(\Sym^n_S(Y),\Aa)$ 中的同构.
\end{enonce*}

这一定理的证明将占据 \hyperref[XVII.subsec5.5]{5.5} 的其余部分.

\begin{enonce*}[remark]{注 5.5.21.1}\etiquette[5.5.21.1]{XVII.rem5.5.21.1}
拟射影性假设的作用仅在于保证 $u$ 容许紧化 $\bar u$
$$
X\xrightarrow{j}\bar X\xrightarrow{\bar u}Y
$$
使概形 $\Sym^n_S(X)$, $\Sym^n_S(\bar X)$ 及 $\Sym^n_S(Y)$ 对每个 $n$ 都存在. 代数空间的语言应能使我们将这一假设换成仅假设 $u$ 可紧化.
\end{enonce*}

证明的\pageoriginale 第一部分是通过分解归约（dévissage）, 将 \hyperref[XVII.thm5.5.21]{5.5.21} 化为一个“不可约”的情形:

\begin{enonce*}{引理 5.5.22}\etiquette[5.5.22]{XVII.lem5.5.22}
定理 \hyperref[XVII.thm5.5.21]{5.5.21} 来自下述特例:
\begin{itemize}
\item[${(\ast)}$] 对每个代数闭域 $k$, $k$ 上的每条光滑射影曲线 $X$, 以及每个素数 $\ell$, 对称 Künneth 箭头 (取 $Y=S=\Spec(k)$, $\Aa=\ZZ/\ell$, $K=\ZZ/\ell$ 所得; 见下文说明)
$$\etiquette[5.5.22.1]{XVII.eq5.5.22.1}
L\Gamma^n(R\Gamma(X,\ZZ/\ell))\longrightarrow
R\Gamma(\Sym^n_k(X),\ZZ/\ell)\leqno(5.5.22.1)
$$
都为同构 $(n\geq0)$.
\end{itemize}
\end{enonce*}

在 \sisi{\hyperref[XVII.eq5.5.22.1]{5.5.22.1}}{\eqhyperref[XVII.eq5.5.22.1]{5.5.22.1}} 中, 符号 $\Gamma$ 不幸有双重用途: 在 $R\Gamma$ 中它表示“整体截面”函子, 而在 $L\Gamma^n$ 中则表示 \hyperref[XVII.par5.5.2.3]{5.5.2.3} 的函子. 这一公式使用了同构 $\Gamma^n_{\ext}(\ZZ/\ell)\simeq\ZZ/\ell$.

设 $S$, $X$, $Y$, $u$ 及 $\Aa$ 如 \hyperref[XVII.thm5.5.21]{5.5.21} 中所述.

\begin{enonce*}{引理 5.5.22.2}\etiquette[5.5.22.2]{XVII.lem5.5.22.2}
设 $C^{b,\geq0}(X,\Aa)$ 中有 \ttor-维数 $\leq0$ 的复形的正合列:
$$
0\longrightarrow K'\longrightarrow K\longrightarrow K''\longrightarrow0{\quoi}.
$$
若关于复形 $K'$, $K$ 及 $K''$ 中任意两个的所有 Künneth 箭头 $k^n$ $(n\geq0)$ 均为同构, 则关于第三个的 Künneth 箭头也都为同构.
\end{enonce*}

我们\pageoriginale 化为 $u$ 射影、且复形 $K'$, $K$ 及 $K''$ 满足 \hyperref[XVII.lem5.5.16]{5.5.16} 的条件 (i) 与 (ii) 的情形, 其中 $N$ 至少为 $u$ 的纤维维数的两倍. 对称 Künneth 箭头此时由 \eqhyperref[XVII.eq5.5.17.1]{5.5.17.1} 给出.

对 $n$ 归纳证明: 若对于 $i\leq n$ 以及复形 $K'$, $K$ 及 $K''$ 中的任意两个, 箭头 $k^i$ 都为同构, 则第三个复形的箭头 $k^i$ $(i\leq n)$ 也都为同构.

为此注意, 由 \sisi{\hyperref[XVII.subsubsec5.5.14]{5.5.14}}{\hyperref[XVII.subsubsec5.5.15]{5.5.15}}, 态射
\begin{equation*}
k^n:L\Gamma^n_{\ext}Ru_{\ast}K\longrightarrow
R\Sym^n(u)_{\ast}L\Gamma^n_{\ext}K\tag{a}
\end{equation*}
可以表示为有限滤过的滤过复形之间的态射
$$
{k'}^n:K_0\longrightarrow K_1
$$
在逐次商上, 这一态射诱导的箭头同构于复合箭头
\begin{equation*}
L\Gamma^i_{\ext}Ru_{\ast}K'\overset{\LL}{\vee}
L\Gamma^{n-i}_{\ext}Ru_{\ast}K''\longrightarrow
R\Sym^n(u)_{\ast}(L\Gamma^i_{\ext}K'\overset{\sisi{}{\LL}}{\vee}
L\Gamma^{n-i}_{\ext}K'')\tag{$\text{b}_i$}
\end{equation*}
具体如下:
\begin{gather*}
L\Gamma^i_{\ext}Ru_{\ast}K'\overset{\LL}{\vee}
L\Gamma^{n-i}_{\ext}Ru_{\ast}K''\xrightarrow{k^i\overset{\LL}{\vee}k^{n-i}}
R\Sym^i(u)_{\ast}L\Gamma^i_{\ext}K'\overset{\LL}{\vee}
R\Sym^{n-i}(u)_{\ast}L\Gamma^{n-i}_{\ext}(K'')\\
\underset{\defn}{=}
\vee_{i,n-i{\ast}}(R\Sym^i(u)_{\ast}L\Gamma^i_{\ext}K'\overset{\LL}{\boxtimes}
R\Sym^{n-i}(u)_{\ast}L\Gamma^{n-i}_{\ext}(K''))\\
\xrightarrow{\substack{\hyperref[XVII.thm5.4.3]{5.4.3}\\\sim}}
\vee_{i,n-i{\ast}}\sisi{(R\Sym^i(u)_{\ast}\times R\Sym^{n-i}(u)_{\ast})}{(R\Sym^i(u)\times R\Sym^{n-i}(u))_\ast}
(L\Gamma^i_{\ext}K'\overset{\LL}{\boxtimes}
L\Gamma^{n-i}_{\ext}K'')\\
=R\Sym^n(u)_{\ast}\vee_{i,n-i{\ast}}(L\Gamma^i_{\ext}K'\overset{\LL}{\boxtimes}L\Gamma^{n-i}_{\ext}K'')
\overset{\dfn}{=}R\Sym^n(u)_{\ast}(L\Gamma^i_{\ext}K'\overset{\LL}{\vee}L\Gamma^{n-i}_{\ext}K''){\quoi}.
\end{gather*}

由\pageoriginale 归纳假设 (及 Künneth 公式), 这些箭头在 $i\neq0,n$ 时为拟同构. 因而下列三个断言中的任意两个蕴含第三个: (a) 为拟同构; $(\text{b}_0)$ 为拟同构; $(\text{b}_n)$ 为拟同构. 这些恰是所考察的三条箭头, \cqfd

\setcounter{subsubsection}{22}
\setcounter{paragraph}{2}
\paragraph{}\etiquette[5.5.22.3]{XVII.par5.5.22.3}
假设 \hyperref[XVII.lem5.5.22]{5.5.22} 的断言 $(\ast)$, 来证明 \hyperref[XVII.thm5.5.21]{5.5.21}. 因而设 $S$, $X$, $Y$, $\Aa$, $K$ 及 $u$ 如 \hyperref[XVII.thm5.5.21]{5.5.21}. 问题在 $S$ 上是局部的, 可设它仿射. 每个复形 $K$ 都是形如 $i_{!}i^{\ast}K$ 的复形的逐次扩张, 其中 $i$ 为局部闭仿射子集到 $X$ 的包含映射. 通过分解归约 (\hyperref[XVII.lem5.5.22.2]{5.5.22.2}), 可假设 $K$ 形如 $i_{!}K_1$, 其中 $i:X_1\rightarrow X$ 为起点仿射的浸入. 相容性 \hyperref[XVII.par5.5.20.2]{5.5.20.2} 允许用 $X_1$ 代替 $X$, 用 $K_1$ 代替 $K$, 因而可设 $u$ 仿射. 将 $u$ 分解为纤维维数 $\leq1$ 的态射的复合 $u=u_1\ldots u_k$. 相容性 \hyperref[XVII.par5.5.20.2]{5.5.20.2} 表明, 只须证明关于 $u_i$ 及任意复形 $K$ 的箭头 $k^n$ 为同构. 因而问题化为 $S$ \emph{仿射且 $u$ 的相对维数为 $\leq1$}的情形.

\paragraph{}\etiquette[5.5.22.4]{XVII.par5.5.22.4}
要使 $k^n$ 为同构, 只须它在 $\Sym^n_S(Y)$ 的几何纤维上的限制为同构; 因而相容性 \hyperref[XVII.par5.5.20.3]{5.5.20.3} ${\rm(v'')}$ 将问题化为 $S$ 为\emph{代数闭域 $k$ 的谱}的情形.

\paragraph{}\etiquette[5.5.22.5]{XVII.par5.5.22.5}
对\pageoriginale $\Sym^n_k(Y)$ 的每个有理点 $y$, 存在 $Y$ 的一个有限既约子概形 $Y'$, 使 $\sisi{x}{y}$ 属于 $\Sym^n_k(Y')$ 的像. 因而相容性 \hyperref[XVII.par5.5.20.3]{5.5.20.3} ${\rm(vi'')}$ 允许假设 $Y$ \emph{为有限个} $\Spec(k)$ \emph{的副本的不交并}.

\paragraph{}\etiquette[5.5.22.6]{XVII.par5.5.22.6}
若 $Y=\amalg Y_i$, 则 $X$ 上的每个复形 $K$ 都是支集在某个概形 $X_i=u^{-1}(Y_i)$ 上的复形之和. 由分解归约 (\hyperref[XVII.lem5.5.22.2]{5.5.22.2}) 以及相容性 \hyperref[XVII.par5.5.20.3]{5.5.20.3} ${\rm(vi'')}$, 只须对每个态射 $u_i:X_i\rightarrow Y_i$ 证明 \hyperref[XVII.thm5.5.21]{5.5.21}. 根据 \hyperref[XVII.par5.5.22.5]{5.5.22.5}, 这将问题化为
$$
\text{\rm $Y=S=\Spec(k)$, $k$ 代数闭, $\dim X\leq1$.}
$$

层 $\Aa$ 此时化为一个 Noether 环 $A$. 下述命题使我们能化为 $A$ 为域的情形.

\begin{enonce*}{命题 5.5.23}\etiquette[5.5.23]{XVII.prop5.5.23}
设 $A$ 为 (交换) Noether 环, $K\in\Ob D^{-}(A)$ 为 $A$-模的上有界复形, $n$ 为整数. 下列条件等价:

\begin{enumerate}
\renewcommand{\theenumi}{\roman{enumi}}
\renewcommand{\labelenumi}{(\theenumi)}
\item $K$ 的 \tTor-维数为 $\leq n$.

\item 对每个 $x\in \Spec (A)$ 以及每个 $i<-n$, 都有 $H^{i}(K\overset{\LL}{\underset{A}{\otimes}} k(x))=0$.
\end{enumerate}
\end{enonce*}

显然有 (i) $\Rightarrow$ (ii)\sisi{,}{.} 假设 (ii). 为证明 (i), 对 $\Spec(A)$ 的闭子集 $F$ 作 Noether 归纳, 证明若模 $N$ 在 $x\not\in F$ 时满足 $N_{x}=0$, 则在 $i<-n$ 时有 $H^{i}(K\overset{\LL}{\otimes}_{A}N)=0$.

设\pageoriginale $I$ 为定义 $F$ 且使 $A/I$ 既约的唯一理想, $N_{k}$ 为 $N$ 中被 $I^{k}$ 零化的子模. 有 $N=\varinjlim N_{k}$. 由于 Tor 函子与滤过归纳极限可交换, 可化为存在 $k$ 使 $N=N_{k}$ 的情形, 从而 $N$ 是被 $I$ 零化的模的逐次扩张; 上同调长正合列将问题化为 $N$ 本身已是 $A/I$-模的情形.

设 $\eta$ 为 $F$ 的一个极大点, 并由正合列
$$\etiquette[5.5.23.1]{XVII.eq5.5.23.1}
0\longrightarrow M_{1}\longrightarrow N \xrightarrow{i} N\otimes
k(\eta)\longrightarrow M_{2} \longrightarrow 0{\quoi}.\leqno(5.5.23.1)
$$
定义 $A/I$-模 $M_{1}$ 与 $M_{2}$.

有 $(M_{j})_{\eta}=0$ $(j=1,2)$. 对每个满足 $\eta\not\in G$ 的闭子集 $G\subset F$, 若以 $M_{j,G}$ 表示 $M_{j}$ 中由 $M_{j}$ 的支集在 $G$ 中的截面组成的子模, 则
$$
M_{j}=\varinjlim_{\eta\not\in G} M_{j,G}{\quoi}.
$$
由归纳假设及极限过渡, 因而有
$$
H^{i}(K\overset{\LL}{\otimes}_{A} M_j)=0\quad\mbox{\rm
当}\quad i< -n{\quoi}.
$$

模 $N\otimes k(\eta)$ 为 $k(\eta)$ 上的向量空间, 因而为若干个 $k(\eta)$ 的直和. 由假设, 有
$$
H^{i}(K\overset{\LL}{\otimes} (N\otimes
k(\eta))\sisi{}{)}=0\quad\mbox{\rm 当}\quad i<-n{\quoi}.
$$

若 $I$ 为 \eqhyperref[XVII.eq5.5.23.1]{5.5.23.1} 中箭头 $i$ 的像, 则有正合列
\begin{align*}
& H^{i-1}(K\overset{\LL}{\otimes}M_{2})\longrightarrow
H^{i}(K\overset{\LL}{\otimes}I)\longrightarrow
H^{i}(K\overset{\LL}{\otimes}(N\otimes k(\eta)))\\
& H^{i}(K\overset{\LL}{\otimes}M_{1})\longrightarrow
H^{i}(K\overset{\LL}{\otimes}N)\longrightarrow
H^{i}(K\overset{\LL}{\otimes}I){\quoi},
\end{align*}
从而\pageoriginale 在 $i<-n$ 时有 $H^{i}(K\overset{\LL}{\otimes}N)=0$,
\cqfd

\begin{enonce*}{推论 5.5.24}\etiquette[5.5.24]{XVII.cor5.5.24}
要使 $D^{-}(A)$ 中的态射 $u:K\rightarrow L$ 为同构, 只须对每个 $x\in \Spec(A)$, 态射
$$
u_{x}:K\overset{\LL}{\otimes}_{A}k(x)\longrightarrow
L\overset{\LL}{\otimes}_{A}k(x)
$$
都是 $D^{-}(k(x))$ 中的同构.
\end{enonce*}

\sisi{\noindent}{}设 $M$ 为 $u$ 的锥 (映射柱). 由假设, 对每个 $x\in\Spec(A)$, 都有 $M\overset{\LL}{\otimes}_{A}k(x)=0$. 因而由 \hyperref[XVII.prop5.5.23]{5.5.23}, 有 $M=0$, 所以 $u$ 为同构.

\setcounter{subsubsection}{25}
\setcounter{paragraph}{0}
\paragraph{}\etiquette[5.5.25.1]{XVII.par5.5.25.1}
继续在 \hyperref[XVII.prop5.5.23]{5.5.23} 处中断的 \hyperref[XVII.lem5.5.22]{5.5.22} 的证明. 由相容性 \hyperref[XVII.par5.5.20.3]{5.5.20.3} ${\rm(iv'')}$ 以及 \hyperref[XVII.cor5.5.24]{5.5.24}, 与 $K\in \Ob D^{b,\tor\leq 0}(X,A)$ 相伴的态射 $k^{n}$ 为同构, 当且仅当对 $x\in
\Spec(A)$, 与 $K\overset{\LL}{\otimes}_{A}k(x)\in \Ob
D^{b,\tor\leq 0}(X,k(x))$ 相伴的态射 $k^{n}$ 为同构. 这将问题化为 $A$ \emph{为特征} $\ell > 0$ \emph{的域}的情形.

\paragraph{}\etiquette[5.5.25.2]{XVII.par5.5.25.2}
正合列 (其中 $C$ 表示一个锥)
$$
0\longrightarrow K[-1]\longrightarrow CK[-1]\longrightarrow
K\longrightarrow 0
$$
表明, 对 $K$ 的态射 $k^{n}$ 为同构, 当且仅当对 $K[-1]$ 的相应态射为同构 (\hyperref[XVII.lem5.5.22.2]{5.5.22.2}).

\paragraph{}\etiquette[5.5.25.3]{XVII.par5.5.25.3}
正合列\pageoriginale
$$
0\longrightarrow \sigma_{\geq n} K\longrightarrow K\longrightarrow
\sigma_{\leq n-1} K\longrightarrow 0
$$
表明, 要使 $K$ 为平坦复形时的 Künneth 态射为同构, 只须它们对 $K$ 的各分量放在次数 $0$ 处时为同构 (\hyperref[XVII.lem5.5.22.2]{5.5.22.2}\sisi{)}{} 及 \sisi{(}{}\hyperref[XVII.par5.5.25.2]{5.5.25.2}).

\paragraph{}\etiquette[5.5.25.4]{XVII.par5.5.25.4}
于是设 $\sisi{F}{K}$ 为 $A$-向量空间层. 这一层为可构造层的滤过归纳极限 (\hyperrefext[IX][IX.cor2.7.2]{2.7.2}); 由于 $A$ 为域, 这些层自动平坦; 通过极限过渡, 问题化为 $K$ \emph{可构造}且集中于零次的情形.

\paragraph{}\etiquette[5.5.25.5]{XVII.par5.5.25.5}
若复形 $K$ 为复形 $L$ 的直和因子, 则与 $K$ 相伴的态射 $k^{n}$ 为与 $L$ 相伴的态射 $k^{n}$ 的直和因子, 因而后者为同构时前者也为同构. 所以迹方法 (\hyperrefext[IX][IX.rem5.8]{5.8}) 适用, 只须处理 $F$ 形如 $j_{!}A_{X'}$ 且 $j:X'\rightarrow X$ 拟有限的情形. 相容性 \hyperref[XVII.par5.5.20.2]{5.5.20.2} 允许用 $X'$ 代替 $X$, 从而得到归约
$$
\mbox{$X$ \emph{的维数为 $\leq 1$, 且} $K=A_{X}${\quoi}.}
$$

\paragraph{}\etiquette[5.5.25.6]{XVII.par5.5.25.6}
可假设 $X$ 既约. 若 $j:U\rightarrow X$ 为 $X$ 的光滑部分, 补集为 $i:F\rightarrow X$, 则正合列
$$
0\longrightarrow j_{!}j^{\ast}A_{X}\longrightarrow
A_{X}\longrightarrow i_{!}i^{\ast}A_{X}\longrightarrow 0
$$
通过分解归约, 将问题化为 $X$ \emph{或为光滑曲线, 或为有限集}的情形. 此外, 由公式
$$
A_{X}=A\overset{\LL}{\otimes}_{\ZZ/\ell}(\ZZ/\ell)_{X}
$$
及 \hyperref[XVII.par5.5.20.3]{5.5.20.3} ${\rm(iv'')}$, 问题化为 $A_{X}=\ZZ/\ell$ 的情形.

若\pageoriginale $X$ 为曲线, 则将 $X$ 嵌入相应的光滑射影曲线, 再作一次分解归约, 即化为下列两种情形之一:
\begin{enumerate}
\renewcommand{\theenumi}{\roman{enumi}}
\renewcommand{\labelenumi}{(\theenumi)}
\item $X$ 为光滑射影曲线, $Y=S=\Spec(k)$ 且 $k$ 代数闭, $\Aa=\ZZ/\ell$, $K=\ZZ/\ell$.

\item $X$ 为有限个 $\Spec(k)$ 的副本的不交并, $Y=S=\Spec(k)$ ($k$ 代数闭), $\Aa=\ZZ/\ell$, $K=\ZZ/\ell$.
\end{enumerate}

最后将 $\sisi{}{(}\ZZ/\ell\sisi{}{)}_{X}$ 分解为层 $\sisi{}{(}\ZZ/\ell\sisi{}{)}_{X\sisi{}{'}}$, 其中 $X'$ 为 $X$ 的连通分支, 就可在 (i) 及 (ii) 中假设 $X$ 连通. 此时情形 (ii) 是平凡的. 这就证明了 \hyperref[XVII.lem5.5.22]{5.5.22}.

\begin{enonce*}{引理 5.5.26}\etiquette[5.5.26]{XVII.lem5.5.26}
定理\emph{\hyperref[XVII.thm5.5.21]{5.5.21}} 来自\emph{\hyperref[XVII.lem5.5.22]{5.5.22}} 的断言 $(\ast)$ 的下列两个特例:
\begin{itemize}
\item[\emph{情形} 1.] 假设 $k=\CC$.

\item[\emph{情形} 2.] 假设 $k$ 的特征为 $p>0$ 且 $\ell=p$.
\end{itemize}
\end{enonce*}

用 \hyperref[XVII.lem5.5.22]{5.5.22} 的记号, 以 $p$ 表示 $k$ 的特征. 若 $p=0$, 则 Lefschetz 原理 (这里依赖于相容性 \hyperref[XVII.par5.5.20.3]{5.5.20.3} ${\rm(v'')}$ 以及 XII \hyperref[XII.cor5.4]{5.4}) 允许化为 $k=\CC$ 的情形.

只剩处理 $p>0$ 且 $\ell\neq p$ 的情形. 曲线 $X$ 可提升为 $k$ 上的 Witt 向量环 $W(k)$ 上的一条曲线 $X'$ (\SGA 1 III 7.4). 令 $S=\Spec(W(k))$, 并考虑图
$$
\xymatrix@C=2.25cm@R=1.5cm{
X'_{\eta} \ar[d]^{f_\eta}          \ar[r]              & X' \ar[d]^{f'}& X\ar[l]\ar[d]^{f}\\
\eta               \ar[r]         & S   & \Spec(k), \ar[l],
}
$$
并\pageoriginale 设 $f_{n}$, $f'_{n}$ 及 $f_{\eta n}$ 分别为 $\Sym^{n}_{k}(X)$, $\Sym^{n}_{S}(X')$ 及 $\Sym^{n}_{\eta}(X'_{\eta})$ 到 $\Spec(k)$, $S$ 及 $\eta$ 的投影. 已知 $f'_{n}$ 为光滑射影态射 (射影性显然; 对光滑性, 最简单的方法是化为 $X'=\PP^{1}_{S}$ 的情形, 此时由初等对称函数定理有 $\Sym^{n}_{S}(X')\simeq
\PP^{n}$). 由特殊化定理 \hyperrefext[XVI][XVI.cor2.2]{2.2}, 层 $R^{i}f'_{n{\ast}}$ $(\ZZ/\ell)$ 在 $S$ 上局部常值. 层 $L^{i}\Gamma^{n}(Rf'_{\ast}\ZZ/\ell)$ 也如此. 因而, 要验证与 $k$ 上的 $X$ 相伴的 \eqhyperref[XVII.eq5.5.22.1]{5.5.22.1} 中的态射为同构, 由 \hyperref[XVII.par5.5.20.3]{5.5.20.3} ${\rm(v'')}$, 只须证明 $X$ 在 $S$ 上的几何泛纤维 $X_{\bar{\eta}}$ 的类似陈述. 域 $k(\bar{\eta})$ 的特征为 $0$, 因而根据第一次归约, 这就证明了 \hyperref[XVII.lem5.5.26]{5.5.26}.

为在 \hyperref[XVII.lem5.5.26]{5.5.26} 的两个特例中证明 \hyperref[XVII.thm5.5.21]{5.5.21}, 我们将对称 Künneth 态射一方面与其超越类似物比较 (\hyperref[XVII.subsubsec5.5.27]{5.5.27} 及 \hyperref[XVII.subsubsec5.5.28]{5.5.28}), 另一方面与其凝聚类似物比较 (\hyperref[XVII.subsubsec5.5.29]{5.5.29} 至 \hyperref[XVII.lem5.5.37]{5.5.37}).

\setcounter{subsubsection}{26}
\subsubsection{}\etiquette[5.5.27]{XVII.subsubsec5.5.27}
若 $p:X\rightarrow S$ 为拓扑空间之间的连续映射, $S$ 以交换环层 $\Aa$ 赋环, 则 \hyperref[XVII.subsubsec5.5.7]{5.5.7} 及 \hyperref[XVII.subsubsec5.5.8]{5.5.8} 的构造原样适用, 定义从 $X$ 上的模到 $\Sym^{n}_{S}(X)=X^{n}/\Sf_{n}$ 上的模的函子 $\Gamma^{n}_{\ext}$. \hyperref[XVII.subsubsec5.5.7]{5.5.7} 至 \hyperref[XVII.prop5.5.18]{5.5.18} 的形式论证均可照搬. 同样, 若 $S$, $X$ 及 $Y$ 为局部紧 Hausdorff 拓扑空间,
$$
\xymatrix@C=1.5cm@R=1.25cm{
X \ar[rr]^{u}\ar[rd]&  &  Y\ar[ld]\\
            & S &,
}
$$
并且\pageoriginale $u$ 的纤维的上同调维数有界, 则 \hyperref[XVII.subsubsec5.5.19]{5.5.19} 及 \hyperref[XVII.subsubsec5.5.20]{5.5.20} 中关于固有支集上同调中的对称 Künneth 态射的形式体系也均可照搬.

若 $p:X\rightarrow S$ 为 $\CC$ 上有限型概形之间的态射, 则函子 $\Gamma^{n}$ 与 $L\Gamma^{n}$ 同“过渡到超越情形”的函子可交换. 最后, 若 $X$, $Y$, $S$, $\Aa$ 及 $K$ 如 \hyperref[XVII.thm5.5.21]{5.5.21}, $S$ 在 $\CC$ 上分离且有限型, 并且以 $\epsilon$ 统称 $\epsilon:X_{\cl}\rightarrow
X_{\et}$ 类型的景态射, 则下图交换:
$$
\vcenter{\etiquette[5.5.27.1]{XVII.eq5.5.27.1}
\xymatrix@C=1.75cm@R=1.5cm{
\epsilon^{\ast} L\Gamma^{n}_{\ext} Ru_{!}K  \ar[r]^{k^{\sisi{k}{n}}} \ar[d]&
\epsilon^{\ast} R \Sym^{n}_{S}(u)_{!} L \Gamma^{n}_{\ext} K \ar[d]\\
L\Gamma^{n}_{\ext} R u^{\an}_{!} \epsilon^{\ast}K  \ar[r]^{k^{n \an}}
& R \Sym^{n}_{S}(u^{\an})_{!} L\Gamma^{n}_{\ext} \epsilon^{\ast}K.
}}\leqno{(5.5.27.1)}
$$

\subsubsection{}\etiquette[5.5.28]{XVII.subsubsec5.5.28}
由相容性 \eqhyperref[XVII.eq5.5.27.1]{5.5.27.1} 以及比较定理的一个初等情形 (\hyperrefext[XI][XI.thm4.4]{4.4} 或 \hyperrefext[XVI][XVI.thm4.1]{4.1}), 要处理 \hyperref[XVII.lem5.5.26]{5.5.26} 的情形 1, 只须验证: 若 $X$ 为光滑复射影曲线, 则 \eqhyperref[XVII.eq5.5.22.1]{5.5.22.1} 的拓扑类似物为同构.

拓扑空间 $X^{\an}$ 可三角剖分 (这一事实在此是平凡的, 因为 $X$ 为 Riemann 球面的分歧覆盖). 分解归约引理 \hyperref[XVII.lem5.5.22.2]{5.5.22.2} 的拓扑类似物, 于是将问题化为证明
\begin{itemize}
\item[$(\ast)$] 若\pageoriginale $D$ 为闭单形 (维数为 0、1 或 2), 则拓扑的对称 Künneth 态射
$$\etiquette[5.5.28.1]{XVII.eq5.5.28.1}
k^{n}:L\Gamma^{n} R\Gamma(D,\ZZ/\ell)\longrightarrow R\Gamma(\Sym^{n}(D\sisi{}{)},\ZZ/\ell)\leqno(5.5.28.1)
$$
为同构.
\end{itemize}

有
$$
R\Gamma(D,\ZZ/\ell)=\ZZ/\ell\;\mbox{\rm 置于次数 }
0{\quoi},
$$
因而
$$
L\Gamma^{n}R\Gamma(D,\ZZ/\ell)=\ZZ/\ell\; \mbox{\rm
置于次数 $0$}{\quoi}.
$$

另一方面, $\Sym^{n}(D)$ 可缩, 于是仍有
$$
R\Gamma(\Sym^{n}(D), \ZZ/\ell)=\ZZ/\ell\;\mbox{\rm 置于
次数 $0$}{\quoi}.
$$

因而 \eqhyperref[XVII.eq5.5.28.1]{5.5.28.1} 两边的 $H^{i}$ 在非零次数都为零, 而在 $H^0$ 上诱导的态射等同于 $\ZZ/\ell$ 的恒等态射.

\subsubsection{}\etiquette[5.5.29]{XVII.subsubsec5.5.29}
对于我们打算作的应用, 将概形 $X$ 上的拟凝聚层视为 $X$ 的平展景上的拟凝聚 $\Oo$-模层是有用的.

设 $p:X\rightarrow \sisi{X}{S}$ 为概形的拟射影态射 (参见 \hyperref[XVII.rem5.5.21.1]{5.5.21.1} \footnote[1]{以下直至包括 \hyperref[XVII.prop5.5.34]{5.5.34} 的全部内容, 在仅假设 $X$ 与 $Y$ 容许 $S$ 上的相对丰沛可逆层的条件下仍然成立 (不假设 $X$ 与 $Y$ 在 $S$ 上有限型).}). 若 $ℱ$ 为 $X_{\text{\rm et}}$ 上的 $\Oo_{X}$-模层, 则
\begin{enumerate}
\renewcommand{\theenumi}{\alph{enumi}}
\renewcommand{\labelenumi}{\theenumi)}
\item 以 $\bigotimes\limits^{n}ℱ$ 表示 $(X/S)^{n}_{\text{\rm et}}$ 上的模层, 它由 $ℱ$ 在 $(X/S)^{n}_{\et}$ 上的逆像 $\Oo_{(X/S)^n}$-模的张量积给出;

\item 层 $\bigotimes\limits^{n}ℱ$ 为 $\mathcal{G}_n$-等变层. 取不变量后, 它定义 $\Sym^{n}_{S}(X)$ 上的一个层 $TS^{n}_{\text{\rm ext}}(ℱ)$: 若 $\sigma$ 为 $(X/S)^{n}$ 到 $\Sym^{n}_{S}(X)$ 的投影, 则有
$$\etiquette[5.5.29.1]{XVII.eq5.5.29.1}
TS^{n}_{\ext}(ℱ)=(\sigma_{\ast}\bigotimes^{n}ℱ)^{\Sf^{n}}{\quoi}.\leqno(5.5.29.1)
$$
\end{enumerate}

\emph{假设\pageoriginale $p$ 平坦}; 如 \hyperref[XVII.subsubsec5.5.8]{5.5.8} 中那样, 将函子 $\Gamma^{n}_{\text{\rm ext}}$ 定义为 $TS^{n}_{\ext}$ 的第 $0$\sisi{阶}{阶}左导出函子. 用 \hyperref[XVII.par5.5.2.5]{5.5.2.5} 及 \hyperref[XVII.subsubsec5.5.8]{5.5.8} 的记号, 若
$$
P_{1}\longrightarrow P_{0}\longrightarrow ℱ\longrightarrow 0
$$
为 $ℱ$ 的平坦表示, 则序列
$$
\xymatrix@C=1.75cm@M=.25cm{
TS^{n}_{\ext} (P'_1) \ar@<3pt>[r]^{j_0} \ar@<-3pt>[r]_{j_1}& TS^{n}_{\ext}(P_0)\ar[r] & \Gamma^{n}_{\ext}(\sisi{F}{ℱ})
}
$$
正合. 若 $x$ 为如 \hyperref[XVII.par5.5.7.1]{5.5.7.1} 中所述的 $\Sym^{n}_{S}(X)$ 的几何点, 则有 (\cf \eqhyperref[XVII.eq5.5.8.1]{5.5.8.1})
\begin{align*}
& \Gamma^{n}_{\text{\rm
ext}}(ℱ)_{x}=\Oo_{x}\otimes_{C}(\bigotimes^{\ell}_{1}
\Gamma_{\Oo_{S}}^{n_{i}}(\sisi{F}{ℱ}_{y_{i}})){\quoi},\tag{5.5.29.2}\\
& C=\bigotimes^{\ell}_{1} \Sym^{n_{i}}_{\Oo_{S}}(\Oo_{y_{i}}){\quoi}.\tag*{其中}
\end{align*}\etiquette[5.5.29.2]{XVII.eq5.5.29.2}% 删除此注释会引入不需要的空白.
事实上, 只须在 $\sisi{F}{ℱ}$ 的茎为自由模的平凡情形下验证这一公式. 注意, $\Oo_{x}$ 在 $C$ 上平坦, 因为它是 $C$ 在 $x$ 的像这一闭点处的严格 Hensel 化.

\subsubsection{}\etiquette[5.5.30]{XVII.subsubsec5.5.30}
假设 $X$ 与 $S$ 仿射, $p$ 平坦且 $ℱ$ 拟凝聚. 设 $S=\Spec(A)$, $X=\Spec(B)$, 并令 $F=H^0(X,ℱ)$, 则 $ℱ$ 为 $B$-模 $F$ 定义的层 $F^{\sim}$. 由 \eqhyperref[XVII.eq5.5.29.1]{5.5.29.1}, 有
$$
TS^{n}_{\ext}(ℱ)=TS^{n}_{A}(F)^{\sim}{\quoi},
$$
因而导出后有
$$\etiquette[5.5.30.1]{XVII.eq5.5.30.1}
\Gamma^{n}_{\text{\rm
ext}}(ℱ)=\Gamma^{n}_{A}(F)^{\sim}{\quoi}.\leqno(5.5.30.1)
$$
所以\pageoriginale 函子 $\Gamma^{n}_{\text{\rm ext}}$ 表现为一个“相对于 $A$”的函子 $\Gamma^{n}$, 正如 \eqhyperref[XVII.eq5.5.29.2]{5.5.29.2} 已在局部所表述的.

\subsubsection{}\etiquette[5.5.31]{XVII.subsubsec5.5.31}
以 $D^{b,\tor_{S}\leq 0}(X,\Oo_{X})$ 表示 $D^{b}(X,\Oo_{X})$ 中由 $\Oo_{X}$-模复形组成的子范畴, 这些复形作为 $p^{-1}(\Oo_{S})$-模复形的 \tTor-维数为 $\leq 0$.

考虑条件

\paragraph{}\etiquette[5.5.31.1]{XVII.par5.5.31.1}
$K\in \Ob C^{b,\geq 0}(X,\sisi{O}{\Oo}_{X})$, 并且在 $X$ 的每个几何点 $x$ 处, $K_{x}$ 同伦于一个非负次数的 $\Oo_{X,x}$-模复形, 其分量在 $\Oo_{S,p(x)}$ 上平坦.

由于态射 $p$ 平坦, 从而有“足够多”的在 $S$ 上平坦的 $\Oo_{X}$-模, 如 \hyperref[XVII.lem5.5.13.2]{5.5.13.2} 中一样可验证, $D^{b,\tor_{S}\leq 0}(X,\Oo_{X})$ 是从满足 (\hyperref[XVII.par5.5.31.1]{5.5.31.1}) 的复形 $K\in \Ob K^{b,\geq 0}(X,\Oo_{X})$ 的范畴, 通过将拟同构变为可逆态射而得到的.

由 \eqhyperref[XVII.eq5.5.29.2]{5.5.29.2} 及 \hyperref[XVII.lem5.5.6.1]{5.5.6.1}, 若 $u:K\rightarrow L$ 为满足 (\hyperref[XVII.par5.5.31.1]{5.5.31.1}) 的复形之间的拟同构, 则 $\Gamma^{n}_{\text{\rm ext}}(u)$ (在 \textsc{Dold-Puppe} 的意义下) 仍为拟同构. 因而可以导出函子 $\Gamma^{n}_{\text{\rm ext}}$, 得到
$$\etiquette[5.5.31.2]{XVII.eq5.5.31.2}
L\Gamma^{n}_{\text{\rm ext}}:D^{b,\tor_{S}\leq
0}(X,\Oo_{X})\longrightarrow D^{b,\tor_{S}\leq 0}(\Sym^{n}_{S}(X), \Oo_{\Sym^{n}_{S}(X)})\leqno(5.5.31.2)
$$

\subsubsection{}\etiquette[5.5.32]{XVII.subsubsec5.5.32}
以 $D^{b,\tor_{S}\leq 0}_{\text{\rm qcoh}}(X,\Oo_{X})$ 表示 $D^{b,\tor_{S}\leq 0}(X,\Oo_{X})$ 中由上同调拟凝聚的复形组成的子范畴.

标准的\pageoriginale 论证使我们可以如 \hyperref[XVII.subsubsec5.5.17]{5.5.17} 中那样定义如下类型的对称 Künneth 箭头:

设 $u:X\rightarrow Y$ 为在 $S$ 上平坦的拟射影 $S$-概形之间的态射; 对 $K\in \Ob D^{b,\tor_{S}\leq 0}_{\text{\rm
qcoh}}(X,\Oo_{X})$, $D^{b,\tor_{S}\leq 0}_{\text{\rm qcoh}}(\Sym^{n}_{S}(Y),\Oo)$ 中的对称 Künneth 箭头 $k^{n}$ 为箭头
$$\etiquette[5.5.32.1]{XVII.eq5.5.32.1}
k^{n}:L\Gamma^{n}_{\text{\rm ext}}Ru_{\ast}K\longrightarrow
R\Sym^{n}_{S}u_{\ast}L\Gamma^{n}_{\text{\rm ext}} K{\quoi}.\leqno(5.5.32.1)
$$

\eqhyperref[XVII.eq5.5.32.1]{5.5.32.1} 将用到的形式性质的验证留给读者, 特别是如下性质:

\begin{enonce*}{引理 5.5.33}\etiquette[5.5.33]{XVII.lem5.5.33}
若 $S$ 的特征为 $p>0$, 且 $u:X\rightarrow S$ 固有且平坦, 则对称 Künneth 箭头 \eqhyperref[XVII.eq5.5.17.2]{5.5.17.2} 与 \eqhyperref[XVII.eq5.5.32.1]{5.5.32.1} 的图交换 ($Y=S, K=\sisi{\sheafZZ}{\ZZ}/p$ 或 $\Oo_{X}$):
\[
\xymatrix{
L\Gamma^{n}_{\sisi{\sheafZZ}{\ZZ}/p}Ru_{\ast}(\sisi{\sheafZZ}{\ZZ}/p)
\ar[r]^{k^n} \ar[d]&R
\Sym^{n}_S(u)_{\ast} (\sisi{\sheafZZ}{\ZZ}/p) \ar[d]&\\
L\Gamma^{n}_{\Oo} Ru_{\ast}(\Oo) \ar[r]^{k^n}& R
\Sym^{n}_{S}(u)_{\ast}\Oo &.
}
\]
\end{enonce*}

\begin{enonce*}{命题 5.5.34}\etiquette[5.5.34]{XVII.prop5.5.34}
在 \hyperref[XVII.subsubsec5.5.32]{5.5.32} 的假设下, 对称 Künneth 箭头 \eqhyperref[XVII.eq5.5.32.1]{5.5.32.1} 为同构.
\end{enonce*}

\begin{enumerate}
\renewcommand{\theenumi}{\alph{enumi}}
\renewcommand{\labelenumi}{\theenumi)}
\item 问题在 $S$ 上是局部的, 可以假设它仿射.

\item 对 $\Sym^{n}_{S}(Y)$ 的每个点 $x$, 存在 $Y$ 的一个仿射开集 $Y_{1}$, 使 $x$ 属于 $\Sym^{n}_{S}(Y_{1})$. 因而通过局部化可以假设 $Y$ 仿射.

\item 所考察类型的每个\pageoriginale 复形 $K$ 都可表示为 $X$ 上的一个拟凝聚模复形 $K_{0}$, 在次数 $<0$ 处为零. 若 $\mathcal{U}$ 为 $X$ 的一个仿射覆盖, 则 $K$ 容许由形如 $j_{\ast}j^{\ast}K$ 的复形的和组成的有限消解, 其中 $j$ 为仿射开集的包含映射 (交错 Čech 复形). 利用分解归约引理 \hyperref[XVII.lem5.5.22.2]{5.5.22.2} 的拟凝聚类似物, 可化为存在 $X$ 的仿射开集 $j:X_{1}\hookrightarrow X$ 及 $K_{1}\in D^{b,\tor_{S}\leq 0}_{\text{\rm
qcoh}}(X_{1},\Oo)$ 使 $K=Rj_{\ast}K_{1}$ 的情形.

\item 根据传递性 \hyperref[XVII.par5.5.20.2]{5.5.20.2} 的拟凝聚类似物, 只须对 $uj$ 与 $j$ 证明 \hyperref[XVII.prop5.5.34]{5.5.34}. 将归约 b) 应用于 $j$, 最终可化为 $X$, $Y$ 及 $S$ 仿射的情形. 此时可将 $K$ 表示为分量为 \sisi{}{$S$-}平坦拟凝聚模的复形 $K_{2}\in C^{b,\geq
0}(X,\Oo_{X})$ (\sisi{\quad}{\hyperref[XVII.def4.1.9]{4.1.9}}). 于是计算函子 $Ru_{\ast}$ 与 $L\Gamma^{n}$ 无须再消解 $K_{2}$, 只须注意, 由 \eqhyperref[XVII.eq5.5.30.1]{5.5.30.1}, 对 $X$ 上的拟凝聚模 $ℱ$, 有
\[
\Gamma^{n}_{\text{\rm ext}} u_{\ast}ℱ\isomto
\Sym^{n}_{S}(u)_{\ast} \Gamma^{n}_{\text{\rm ext}}ℱ{\quoi},
\]
两边均等同于
\[
\Gamma^{n}_{A}(F)^{\sim} \quad\mbox{\rm 当}\quad
ℱ=F^{\sim}{\quoi},
\]
\cqfd
\end{enumerate}

\setcounter{subsubsection}{34}
\subsubsection{}\etiquette[5.5.35]{XVII.subsubsec5.5.35}
处理 \hyperref[XVII.lem5.5.26]{5.5.26} 的情形 $2$. 用 \hyperref[XVII.lem5.5.22]{5.5.22} 的记号, 下图交换 (\hyperref[XVII.lem5.5.33]{5.5.33}):
$$\vcenter{\etiquette[5.5.35.1]{XVII.eq5.5.35.1}
\xymatrix@C=2cm{
  L\Gamma^{n} (R\Gamma(X,\ZZ/p)) \ar[r]^{k^n}_*+[o][F-]{\scriptstyle{1}}
  \ar[d]& R\Gamma(\Sym^{n}_{k}(X),\ZZ/p)\ar[d] \\
L\Gamma^{n}(R\Gamma(X,\Oo)) \ar[r]^{k^n}_*+[o][F-]{\scriptstyle{2}} & R\Gamma(\Sym^{n}_{k}(X),\Oo),
}}\leqno{(5.5.35.1)}
$$
并且\pageoriginale 由 \hyperref[XVII.prop5.5.34]{5.5.34}, 箭头 \ding{193} 为同构. $\Oo$ 的 Frobenius 自同态 $F:x\rightarrow x^{p}$ 在 $R\Gamma(X,\Oo)$ 及 $R\Gamma(\Sym^{n}_{k}(X),\Oo)$ 上 $p$-线性地作用. 因而它也在 $L\Gamma^{n}(R\Gamma(X,\Oo))$ 上 $p$-线性地作用. 可验证 $k^{n}$ 与 $F$ 可交换.

Artin--Schreier 理论给出正合列
$$
0\longrightarrow H^{i}(X,\ZZ/p)\longrightarrow
H^{i}(X,\Oo)\xrightarrow{F-I} H^{i}(X,\Oo)\longrightarrow
\sisi{\Oo}{0}{\quoi},
$$
对 $\Sym^{n}_{S}(X)$ 也类似. 要由 \ding{193} 为同构推出 \ding{192} 为同构, 只须验证以下引理.

\begin{enonce*}{引理 5.5.36}\etiquette[5.5.36]{XVII.lem5.5.36}
设 $k$ 为特征 $p>0$ 的完美域, $K\in C^{\geq
0}(\ZZ/p)$ 为有限维 $\ZZ/p$-向量空间复形, $L\in C^{\geq 0}(k)$ 为有限维 $k$-向量空间复形, $F$ 为 $L$ 的 $p$-线性自同态, 并设 $u:K\rightarrow L$ 为一个态射. 假设 $H^{\ast}(Fu)=H^{\ast}(u)$, 且序列
$$
0\longrightarrow H^{i}(K)\longrightarrow
H^{i}(L)\xrightarrow{F-I}H^{i}(L)\longrightarrow 0
$$
正合. 则\pageoriginale 序列
$$
0\longrightarrow H^{i}(\Gamma^{n}K)\longrightarrow
H^{i}(\Gamma^{n}L)\xrightarrow{F-I}H^{i}(\Gamma^{n}L)\longrightarrow 0
$$
正合.
\end{enonce*}

由于 $H^{\ast}(Fu)=H^{\ast}(u)$ 且 $\ZZ/p$ 为域, 存在同伦
$$
Fu-u=dH+Hd{\quoi}.
$$
由于域 $k$ 完美, 可以写成
$$
FH_{0}-H_{0}=H{\quoi}.
$$
用 $u-(dH_{0}+H_{0}d)$ 代替 $u$ 后, 可设 $Fu=u$.

函子 $S$ 将每个配有 $p$-线性自同态 $F$ 的 $k$-向量空间 $V$, 映为由满足 $Fv=v$ 的 $v\in
V$ 所成的 $\ZZ/p$-向量空间 $V^{S}$; 它是一个正合函子. 因而假设意味着态射
$$
K\longrightarrow L^{S}
$$
为拟同构. 只须证明函子 $S$ 与函子 $\Gamma^{n}$ 可交换. 这容易由下述众所周知的引理推出, 其证明留给读者:

\begin{enonce*}{引理 5.5.37}\etiquette[5.5.37]{XVII.lem5.5.37}
完美域 $k$ 上配有 $p$-线性自同态 $F$ 的有限维向量空间 $V$, 容许唯一的分解
$$
V=V_{0}+V_{n}
$$
满足 $V^{S}\otimes k\isomto  V_{0}$, $F(V_{n})\subset V_{n}$, 且 $F|V_{n}$ 幂零.
\end{enonce*}

这就完成了 \hyperref[XVII.thm5.5.21]{5.5.21} 的证明.



% END native input: chapters/exp-17-part4-zh.tex
% BEGIN native input: chapters/exp-17-part5-zh.tex
\section[函子 $f_{!}$]{函子 $f_{!}$}\etiquette[6]{XVII.sec6}\pageoriginale

我们将直接构造 \hyperref[XVII.var5.1.11]{5.1.11} 的函子 $f_{!}$。
然后利用这一构造，为有限表现的拟有限平坦态射定义“迹态射”
（\hyperref[XVII.thm6.2.3]{6.2.3}）。

\subsection[函子 $f_{!}$ 的新构造]{函子 $f_{!}$ 的新构造}\etiquette[6.1]{XVII.subsec6.1}

设 $f:X\rightarrow S$ 是局部有限型态射。$X$ 的一个子集 $P$，
若是 $X$ 的某个在 $S$ 上固有的子概形的像，则称它在 $S$ 上\emph{固有}。

\begin{enonce*}{命题 6.1.1}\etiquette[6.1.1]{XVII.prop6.1.1}
设 $f:X\rightarrow S$ 是局部有限型且拟分离的态射。
\begin{enumerate}
\renewcommand{\theenumi}{\roman{enumi}}
\renewcommand{\labelenumi}{(\theenumi)}
\item $X$ 的一个在 $S$ 上固有的子集是局部闭的；若 $f$ 分离，则该子集闭。

\item 若 $f$ 分离，则 $X$ 的两个在 $S$ 上固有的子集之并仍在 $S$ 上固有。

\item 对每个 $S$-概形 $U$，以 $\Phi_{X}(U)$ 表示 $\fprod{U}{X}{S}$ 的
在 $U$ 上固有的子集所成的集合。函子 $\Phi_{X}$ 是 \emph{fpqc} 拓扑下的层。
若 $f$ 分离，则就包含关系而言，$\Phi\sisi{}{_X}$ 是滤过偏序集层。
\end{enumerate}
\end{enonce*}

断言\pageoriginale (i) 和 (ii) 平凡，(iii) 的第二个断言由 (ii) 得到。
对于覆盖均为满射的每种拓扑，函子 $U\mapsto \mathcal{P}(\fprod{U}{X}{S})$ 都是层。
由 \EGA IV 2.3.14 及 \EGA IV 2.7.1 (vii)，$\Phi_{X}$ 是它的子层。

若 $\Ff $ 是概形 $X$ 上的带点集平展层，回顾（IV \hyperref[IV.subsubsec8.5.2]{8.5.2}），
$\Ff $ 的一个截面 $s$ 的支集，是使 $s=0$ 的最大开集的补集，
其中 $0$ 表示指定截面。

\begin{enonce*}{定义 6.1.2}\etiquette[6.1.2]{XVII.def6.1.2}
设 $f:X\rightarrow S$ 是局部有限型分离态射，$\Ff $ 是 $X$ 上的带点集层。
以 $f_{!}\Ff $ 表示 $f_{\ast}\Ff $ 的一个子层，它在 $U$ 上的截面
（它在 $S$ 上平展），是 $\Ff $ 在 $\fprod{U}{X}{S}$ 上的支集在 $U$ 上固有的截面。
\end{enonce*}

由 \hyperref[XVII.prop6.1.1]{6.1.1}，$f_{!}\Ff $ 确实是层。

容易验证如下引理：

\begin{enonce*}{引理 6.1.3}\etiquette[6.1.3]{XVII.lem6.1.3}
若 $f$ 是两个局部有限型分离态射的复合 $gh$，则存在唯一的函子同构
$f_{!}\simeq g_{!}h_{!}$，使下图交换
$$
\xymatrix@C=1cm@R=1.25cm@M=.2cm{
f_{!}  \ar@{-}[r]^{\sim} \ar@{^{(}->}[d]& g_{!}h_{!} \ar@{^{(}->}[d]\\
f_{\ast} \ar@{=}[r]      & g_{\ast} h_{\ast}.
}
$$
\end{enonce*}

若\pageoriginale $f$ 是开浸入（\resp 固有态射），则函子 $f_{!}$ 是补零延拓函子
（\resp 函子 $f_{\ast}$）。因此，由 \hyperref[XVII.lem6.1.3]{6.1.3}，
定义 \hyperref[XVII.def6.1.2]{6.1.2} 和 \hyperref[XVII.var5.1.11]{5.1.11}
在共同适用的范围内一致。

在 \hyperref[XVII.sec5]{\S 5} 验证的函子 $f_{!}$ 的性质，
在更一般的 \hyperref[XVII.def6.1.2]{6.1.2} 的情形中仍然有效。

\begin{enonce*}{命题 6.1.4}\etiquette[6.1.4]{XVII.prop6.1.4}
设 $f:X\rightarrow S$ 是局部有限型分离态射。
\begin{enumerate}
\renewcommand{\theenumi}{\roman{enumi}}
\renewcommand{\labelenumi}{(\theenumi)}
\item 函子 $f_{!}$ 左正合（\ie 与有限逆向极限交换），且与滤过正向极限交换。

\item 若 $f$ 拟有限，则函子 $f_{!}$ 是从 $X$ 上的阿贝尔层范畴到
$S$ 上的阿贝尔层范畴的正合函子。

\item 对每个笛卡尔图
$$
\xymatrix@C=1.5cm@R=1.25cm{
X'          \ar[r]^{g'} \ar[d]^{f'}& X \ar[d]^{f}\\
S'          \ar[r]^{g}  & S
}
$$
存在基变换同构，使图
$$
\xymatrix@C=1.5cm@R=1.25cm{
g^{\ast}f_{!}  \sisi{\ar@{=}[r]}{\ar@{-}[r]^\sim} \ar@{^{(}->}[d]& f'_{!}g'^{\ast}\ar@{^{(}->}[d]\\
g^{\ast}f_{\ast}  \ar[r]    & f'_{\ast}g'^{\ast},
}
$$
交换，其中第二个水平箭头是基变换箭头
\emph{\neweqhyperrefext[XII][XII.eq4.1]{4.1}}。
\end{enumerate}
\end{enonce*}

\begin{enonce*}[remark]{证明}
对\pageoriginale 每个 $U\in \Ob S_{\et}$，有
\[
f_{!}\Ff (U)=\varinjlim_{P\in
\Phi(U)}\Gamma_{\sisi{p}{P}}(\Ff ){\quoi},
\]
而当 $U$ 拟紧且拟分离时，函子 $\Ff \mapsto f_{!}\Ff (U)$
与有限逆向极限和滤过正向极限交换，因为它是具有这些性质的函子的滤过正向极限
（\hyperref[XVII.prop6.1.1]{6.1.1}\sisi{)}{} (iii)\sisi{}{)}。

这证明了 (i)。断言 (ii) 立即由 (iii) 得到（取 $S$ 的各几何点作为基变换）。

断言 (iii) 在 $S$ 上是局部的，因此可假设它仿射。
设 $I$ 是 $X$ 的在 $S$ 上有限型的开子集所成的滤过偏序集。
若 $U\overset{j}{\hookrightarrow}X$ 是 $X$ 的开子集，则
$(fj)_{!}(\Ff |U)=f_{!}(j_{!}(\Ff |U))$ 映到 $f_{!}(\Ff\sisi{|U}{})$。
此外，
$$
f_{!}(\Ff )=\varinjlim_{U\in I}
(fj\sisi{}{)}_{!}(\Ff |U){\quoi}.
$$
这样便可假设 $X$ 在 $S$ 上有限型，并应用 Chow 引理
（\EGA II 第 2\sisi{}{} 版）\nde{见 \nderef{XVII.ndeConrad} 中引用的 [C, 2.6]。}
找到一个射影满态射 $p^{0}:X^{0}\rightarrow X$，使 $fp\sisi{_{0}}{^0}$ 可紧化。
设 $X^{1}=\fprod{X^{0}}{X^{0}}{X}$，$p^{1}$ 是从 $X^{1}$ 到 $X$ 的自然箭头。
序列
$$
0\rightarrow \Ff \rightarrow
p^{0}_{\ast}p^{0{\ast}}\Ff \rightrightarrows
p^{1}_{\ast}p^{1{\ast}}\Ff
$$
正合（\sisi{VII 9.2}{\hyperrefext[VIII][VIII.cor9.2]{9.2}}）。

考虑到基变换箭头的函子性和 (i)，只须对层
$p^{0}_{\ast}p^{0^{\ast}}\Ff $ 与 $p^{1}_{\ast}p^{1^{\ast}}\Ff $ 验证 (iii)。
通过记号变换 $(p^{1}\mapsto p^{0})$，我们\pageoriginale 便归结到假设
$\Ff $ 形如 $p^{0}_{\ast}\mathcal{G}$。对二元组 $(f,p^{0}_{\ast}\mathcal{G})$
的断言 (iii)，等价于对二元组 $(p^{0}_{\ast}\sisi{}{,}\mathcal{G})$ 和
$(fp^{0}_{\ast},\mathcal{G})$ 的断言 (iii)。由 \hyperref[XVII.var5.2.7]{5.2.7} (i)
即可得出结论；它之所以适用，是因为 $p^{0}$ 和 $fp^{0}$ 可紧化。
\end{enonce*}

\begin{enonce*}[remark]{反例 6.1.6}\etiquette[6.1.6]{XVII.contrexem6.1.6}
函子 $Rf_{!}$ 一般不与函子 $f_{!}$ 的导出函子一致。
事实上，若 $X$ 是代数闭域 $k$ 上的仿射概形，则有
$$
f_{!}\Ff =\sisi{\sum}{\bigoplus}_{x\in X(k)}H^{0}_{x}(\Ff ){\quoi},
$$
因而 $f_{!}$ 的第 $i$ 个导出函子在 $\Ff $ 处的值为
$$
\sisi{\sum}{\bigoplus}_{x\in X(k)}H^{i}_{x}(\Ff ){\quoi}.
$$
另一方面，若 $X$ 光滑、不可约且维数为 1，则对于与 $k$ 的特征互素的 $\ell$，有
$$
H^{2}_{c}(X,\mu_{\ell^{n}})\simeq \ZZ/\ell^{n}
$$
成立。
\end{enonce*}

\subsection{相对维数 0 的迹态射}\etiquette[6.2]{XVII.subsec6.2}

本 ${\rm n^{\circ}}$ 的主要结果是定理 \hyperref[XVII.thm6.2.3]{6.2.3}。
建议读者不读技术预备部分 \hyperref[XVII.subsubsec6.2.1]{6.2.1}
和 \hyperref[XVII.prop6.2.2]{6.2.2}。
它们使 \hyperref[XVII.thm6.2.3]{6.2.3} 的证明只须考虑有限态射 $f$；
借助 \EGA IV 18.12\sisi{.}{，}，这一归约无论如何都相当显然。

\subsubsection{}\etiquette[6.2.1]{XVII.subsubsec6.2.1}
设\pageoriginale $f:X\rightarrow S$ 是\emph{拟有限分离}态射。
对每个 $U\in \Ob S^{\et}$，以 $\Psi_{f}(U)$ 表示下述二元组所成的范畴：
一个有限带点集 $I$，其指定点记为 $0$；以及将 $f^{\ast}U$ 分成开闭子集的一个分解
$f^{\ast}U=\coprod_{i\in I}V_{i}$，其中当 $i\neq 0$ 时 $V_{i}$ 在 $U$ 上有限。
从 $\varphi=(V_{i})_{i\in I}$ 到 $\varphi'=(V'_{i})_{i\in I'}$ 的态射，
是从 $I'$ 到 $I$ 的一个映射 $\sigma$，使对 $i\in I$ 有
$V_{i}=\bigcup_{j\in \sigma^{-1}(i)}V'_{j}$。
若存在从 $\varphi$ 到 $\varphi'$ 的态射，则称 $\varphi'$ \emph{加细} $\varphi$。

设 $\Ff $ 是 $S$ 上的阿贝尔群层，$U\in \Ob S_{\text{\rm et}}$，
$\varphi=(V_{i})_{i\in I}\in \Ob\Psi_{f}(U)$，$I^{\ast}=I\backslash \{0\}$，
$f_{i}$ 是 $V_{i}$ 到 $U$ 的投影。有
$$
f_{!}f^{\ast}\Ff |U=f_{0!}f_{0}^{\ast}(\Ff |U)\oplus
\bigoplus_{i\neq 0} f_{i\ast}f_{i}^{\ast}(\Ff |U){\quoi},
$$
从而有复合箭头
$$\etiquette[6.2.1.1]{XVII.eq6.2.1.1}
\tau_{\varphi}:\Ff ^{I^{\ast}}|U\longrightarrow
\bigoplus_{i\neq 0}
f_{i\ast}f_{i}^{\ast}(\Ff |U)\longrightarrow f_{!}f^{\ast}\Ff |U{\quoi}.\leqno(6.2.1.1)
$$

若 $\sigma:I\rightarrow J$ 是带点集态射，则以
$\leftexp{t}{\sigma}:\Ff ^{J^{\ast}}\rightarrow \Ff ^{I^{\ast}}$ 表示如下箭头：
若 $\sigma(i)=j$，则其坐标 $\leftexp{t}{\sigma}^i_{\phantom{i}j}=\Id$；
若 $\sigma(i)\neq j$，则 $\leftexp{t}{\sigma}^i_{\phantom{i}j}=0$。
若 $\sigma:I'\rightarrow I$ 是从 $\varphi$ 到 $\varphi'$ 的态射，则图
$$
\vcenter{\etiquette[6.2.1.2]{XVII.eq6.2.1.2}
\xymatrix@R=1.25cm{
\Ff^{I^{\ast}}|U \ar[rr]^{\leftexp{t}{\sigma}}   \ar[rd] \ar[rd]_{\tau_{\varphi}}&&
\Ff^{I'^{\ast}}|U     \ar[ld]^{\tau_{\varphi'}}\\
      &          f_{!}f^{\ast}\Ff|U  &
}}\leqno{(6.2.1.2)}
$$
\pageoriginale 交换。此外，${}^{t}(\sigma\rho)={}^{t}\rho{}^{t}\sigma$。
各 $\Psi_{f}(U)$ 组成 $S_{\text{\rm et}}$ 上的一个纤维化范畴 $\Psi_{f}$；
上述构造与局部化交换，因而定义一个 $S_{\text{\rm et}}$-函子
$\varphi=(V_{i})_{i\in I} \in \Psi_{f}(U)\mapsto \Ff ^{I^{\ast}}|U$，
从 $\Psi$ 到 $S$ 的平展开子概形上的层所成的 $S_{\text{\rm et}}$-范畴，
并定义一个函子态射 $\tau_{\varphi}:\Ff ^{I^{\ast}}|U\rightarrow f_{!}f^{\ast}\Ff |U$。

\begin{enonce*}{命题 6.2.2}\etiquette[6.2.2]{XVII.prop6.2.2}
\begin{enumerate}
\renewcommand{\theenumi}{\roman{enumi}}
\renewcommand{\labelenumi}{(\theenumi)}
\item 范畴 $\Psi_{f}$ 在 $S_{\text{\rm et}}$ 上局部向右滤过
\emph{（\sisi{IV 14}{\hyperrefext[V][V.def8.1.1]{8.1.1}}）}。

\item 箭头 \eqhyperref[XVII.eq6.2.1.1]{6.2.1.1} 定义一个同构\sisi{们}{}
\[
\tau:\varinjlim_{\varphi=(V_{i})_{i\in I}\in
\Psi_{f}(U)}\Ff ^{I^{\ast}}|U\isomto  f_{!}f^{\ast}\Ff\sisi{}{|U} {\quoi}.
\]
\end{enumerate}
\end{enonce*}

若 $S$ 严格局部，闭点为 $s$，则 $S$-概形 $X$ 有开闭部分的一个分解
$\varphi\in\Psi_{f}(S)$，$X=\coprod_{i\in I}X_{i}$，以带点集 $I$ 为指标，
使 $X_{0}$ 在 $s$ 处的纤维为空，各 $X_{i}$ $(i\neq 0)$ 在 $S$ 上有限，
且 $X_{i}$ $(i\neq 0)$ 在 $s$ 处的纤维仅含一点（\EGA IV 18.12.1）。
这样的分解是 $\Psi\sisi{}{_f}(S)$ 的终对象。
此外，由 \hyperref[XVII.prop6.1.4]{6.1.4} (iii)（或直接由显然性），
箭头 $\tau_{\varphi}$ \eqhyperref[XVII.eq6.2.1.1]{6.2.1.1} 在 $s$ 处的纤维是同构。

在一般情形中，设 $x$ 是 $S$ 的一个几何点，$\Psi_{x}$ 是 $\Psi_{f}$ 在 $x$ 处的纤维，
\[
\Psi_{x}=\varinjlim_{U(x)} \Psi_{f}(U){\quoi},
\]
其中 $U$ 遍历 $x$ 的平展邻域。设 $S_{x}$ 是 $S$ 在 $x$ 处的严格局部化，
$X_{x}=\fprod{X}{S_{x}}{S}$，$f_{x}:X_{x}\rightarrow S_{x}$。
标准的极限论证表明\pageoriginale，范畴 $\Psi_{x}$ 等价于范畴
$\Psi_{f_{x}}(S_{x})$，因而向右滤过。由此，$\Psi\sisi{}{_f}$ 局部向右滤过。
此外，$\tau$ 在 $x$ 处的纤维等同于箭头
$$
\tau_{x}:\varinjlim_{\Psi_{x}}\Ff _{x}^{I_{x}}\longrightarrow
(f_{!}f^{\ast}\Ff )_{x}
$$
而根据上述讨论，这个箭头是同构。

\begin{enonce*}{定理 6.2.3}\etiquette[6.2.3]{XVII.thm6.2.3}
存在唯一的方式，对每个分离、平坦、有限表现且拟有限的态射 $f:X\rightarrow S$
以及 $S$ 上的每个阿贝尔层 $\Ff $，定义迹态射
$$
\Tr_{f}:f_{!}f^{\ast}\Ff \longrightarrow \Ff
$$
使以下条件满足：
\begin{itemize}
\item[\textsuperscript{(Var 1)}] $\Tr_{f}$ 对 $\Ff $ 具有函子性。

\item[\textsuperscript{(Var 2)}] $\Tr_{f}$ 与任意基变换相容，\ie 对每个笛卡尔图
$$
\xymatrix{
X' \ar[rr]^{g'} \ar[d]^{f'}& & X\ar[d]^{f}\\
S' \ar[rr]^{g} & &  S
}
$$
下列函子图交换：
$$
\xymatrix{
f'_{!} f'^{\ast} g^{\ast} \ar@{=}[r] \ar[d]^{\Tr_{f'}\ast g^{\ast}}& f'_{!}g'^{\ast} f^{\ast}\ar@{=}[r] &
g^{\ast}f_{!}f^{\ast} \ar[d]^{g^{\ast}\ast
  \Tr_f}\\
g^{\ast} \ar@{=}[rr]&  & g^{\ast} &.
}
$$

\item[\textsuperscript{(Var 3)}] $\Tr_{f}$\pageoriginale 与复合相容，
\ie 若 $f$ 是两个分离、平坦、有限表现且拟有限的态射的复合 $f=gh$，则下图交换：
$$
\xymatrix{
f_{!}f^{\ast}\ar[rrdd]_{\Tr_f} \ar@{=}[rr]& & g_{!}h_{!} h^{\ast}g^{\ast}\ar[d]^{g_{!}\ast\mathord{\Tr_{h}}\ast g^\ast}\\
              & &  g_{!}g^{\ast} \ar[d]^{\Tr_g}\\
              & & \Id .
}
$$

\item[\textsuperscript{(Var 4)}] （正规化条件）。若 $f$ 有限且常秩为 $n$，则复合同态
$$
\Ff \rightarrow
f_{\ast}f^{\ast}\Ff =f_{!}f^{\ast}\Ff \xrightarrow{\Tr_{f}}\Ff
$$
为乘以 $n$ 的同态。
\end{itemize}
\end{enonce*}

\sisi{\emph{首先注意，条件 \emph{(Var 2)} 至 \emph{(Var 4)} 蕴含下述引理。}}{
首先注意，条件 (Var 2) 至 (Var 4) 蕴含下述引理。}

\begin{enonce*}{引理 6.2.3.1}\etiquette[6.2.3.1]{XVII.lem6.2.3.1}
若 $X=X'\amalg X''$\sisi{}{，}令 $f'=f|X'$ 和 $f''=f|X''$，则态射 $\Tr_{f}$：
$$
f_{!}f^{\ast}\Ff =f'_{!}{f'}^{\ast}\Ff \oplus
f''_{!}{f''}^{\ast}\Ff \longrightarrow \Ff
$$
是 $\Tr_{f'}$ 和 $\Tr_{f''}$ 之和。
$$
\xymatrix@C=2.25cm@R=.75cm{
  X' \ar@{^{(}->}[rd]^{j'}&           &\\
   & X   \ar[r]^{f}  \ar[r]^{f}    & S\\
X'' \ar@{^{(}->}[ru]_{j''}& &.
}
$$
\end{enonce*}

若\pageoriginale 对 $j'$ 应用关于基变换 $j'$ 的 (Var 2)，并对恒等映射应用 (Var 4)，可知
$$
\Tr_{j'}:j'_{!}{j'}^{\ast}f^{\ast}\Ff \longrightarrow
f^{\ast}\Ff
$$
是由 $f^{\ast}\Ff $ 的直和分解
$$
f^{\ast}\Ff =j'_{!}{j'}^{\ast}f^{\ast}\Ff \oplus
j''_{!}{j''}^{\ast}f^{\ast}\Ff {\quoi}.
$$
所定义的典范映射。

若对 $f'=f\circ j'$ 应用 (Var 3)，可知下图交换
\[
\xymatrix@R=1.25cm{
f'_{!} f'^{\ast}\Ff  \ar[rr] \ar[rd]_{\Tr_{f'}}& &f_{!}f^{\ast}\Ff \ar[ld]^{\Tr_f}\\
                                   & \Ff    &    ,
}
\]
结合 $f''$ 的类似陈述，这正是 \hyperref[XVII.lem6.2.3.1]{6.2.3.1}。

设 $f:X\rightarrow S$ 是分离、平坦、拟有限的有限表现态射。
对每个 $U\in \Ob S_{\text{\rm et}}$，设 $\Psi'_{f}(U)$ 是
$\Psi_{f}(U)$（\sisi{\hyperref[XVII.prop6.1.1]{6.1.1}}{\hyperref[XVII.subsubsec6.2.1]{6.2.1}}）
的满子范畴，由满足下述条件的分解 $(V_{i})_{i\in I}$ 组成：
当 $i\neq 0$ 时，$V_{i}$ 在 $U$ 上局部自由且常秩。
诸范畴\pageoriginale $\Psi'_{f}(U)$ 组成 $\Psi_{f}$ 的一个纤维化子范畴，
在 $\Psi_{f}$ 中共尾。

设 $\Ff $ 和 $\mathcal{G}$ 是 $S$ 上的阿贝尔层。
由 \hyperref[XVII.prop6.2.2]{6.2.2} (ii)，给定一个态射
$g:f_{!}f^{\ast}\Ff \rightarrow \mathcal{G}$，等价于对每个分解
$\varphi=(V_{i})_{i\in I}\in\Psi'_{f}(U)$ 给定一个态射
$g_{\varphi}:\Ff ^{I^{\ast}}|U\rightarrow \mathcal{G}|U$，
与 \eqhyperref[XVII.eq6.2.1.2]{6.2.1.2} 的态射 ${}^{t}\sigma$ 相容。
若 $f_{i}$ 是从 $V_{i}$ 到 $U$ 的投影，则 $g_{\varphi}$ 是坐标箭头为下述复合的箭头：
\[
g^{i}_{\varphi}:\Ff |U\rightarrow
f_{i\ast}f^{\ast}_{i}\Ff |U\rightarrow
f_{!}f^{\ast}\Ff |U\xrightarrow{g} \mathcal{G}|U{\quoi}.
\]

由此，存在唯一的箭头 $\Tr_{f}$，它与平展局部化相容，且满足 (Var 4) 和
\hyperref[XVII.lem6.2.3.1]{6.2.3.1}。显然，它满足 (Var 1) 和 (Var~2)。

为验证 (Var 3)，通过基变换可归结到 $S$ 为代数闭域的谱这一简单情形。

\setcounter{subsubsection}{3}
\subsubsection{}\etiquette[6.2.4]{XVII.subsubsec6.2.4}
分离拟有限态射 $f:X\rightarrow S$ 的一个\emph{赋权}（weighting），
是对 $X$ 的每个几何点 $x$ 给定一个整数 $n(x)$，它仅依赖于 $x$ 所在的 $X$ 的底层点，且满足：
\begin{itemize}
\item[$(\ast)$] \emph{对于每个 $U\in \Ob S_{\text{\rm et}}$ 以及
$f^{-1}(U)$ 中在 $U$ 上有限的每个开闭子空间 $X_{1}$，
将 $U$ 的几何\sisi{点}{点} $s$ 映到数 $\sum_{x\in X_{1}(s)}n(x)$ 的函数是局部常值的。}
\end{itemize}

若 $f=gh$ 是两个已赋权态射的复合，则令
$n_{f}(x)=n_{h}(x)\sisi{.}{\cdot}n_{g}(h(x))$，以此给 $f$ 赋权。

若\pageoriginale 态射 $f'$ 由已赋权态射 $f$ 经基变换 $g$ 得到，
则令 $n_{f'}(x)=n_{f}(g'(x))$，以此给 $f'$ 赋权。

上述论证实际上证明了：

\begin{enonce*}{命题 6.2.5}\etiquette[6.2.5]{XVII.prop6.2.5}
存在唯一的方式，对每个已赋权的分离拟有限态射 $f:X\rightarrow S$ 以及
$S$ 上的每个阿贝尔层 $\Ff $，定义迹态射
$\Tr_{f}:f_{!}f^{\ast}\Ff \rightarrow \Ff $，
它对 $\Ff $ 具有函子性，与基变换及复合相容，并满足：
\begin{itemize}
\item[\textsuperscript{(Var $\sisi{\leftexp{\prime}{4}}{4'}$)}]
若 $f$ 有限，并存在整数 $n$，使对 $S$ 的每个几何点 $s$ 都有
$\sum_{x\in X(s)}n(x)=n$，则复合箭头
$$
\Ff \longrightarrow
f_{\ast}f^{\ast}\Ff =f_{!}f^{\ast}\Ff \xrightarrow{\Tr_{f}}\Ff
$$
为乘以 $n$ 的箭头。
\end{itemize}
\end{enonce*}

\begin{enonce*}[remark]{例 6.2.6}\etiquette[6.2.6]{XVII.exem6.2.6}
设 $S$ 是\sisi{}{约化的} Noether 几何单支概形（例如正规概形），
$f:X\rightarrow S$ 是分离拟有限态射。对 $X$ 的每个几何点 $x$，
其在 $S$ 中的像为 $s$，设 $\Oo_{x}$（\resp $\Oo_{s}$）是
$X$ 在 $x$ 处（\resp $S$ 在 $s$ 处）的严格局部环。
概形 $\Spec(\Oo_{s})$ 约化且有唯一的极大点 $\eta$。
此外，$\Oo_{x}$ 在 $\Oo_{s}$ 上有限。令 $n(x)=\rg_{\eta}(\Oo_{x})$。
容易验证，$n(x)$ 是 $f$ 的赋权，因而给出迹态射 $\Tr_{f}$。
\end{enonce*}

\setcounter{subsubsection}{6}
\subsubsection{}\etiquette[6.2.7]{XVII.subsubsec6.2.7}
设\pageoriginale $u:X\rightarrow Y$ 是有限表现的分离拟有限平坦态射。
因为迹态射（\hyperref[XVII.thm6.2.3]{6.2.3}）具有函子性，
它对 $Y$ 上的每个阿贝尔层复形 $K$ 定义一个复形态射
$$\etiquette[6.2.7.1]{XVII.eq6.2.7.1}
\Tr_{u}:u_{!} u^{\ast}K\longrightarrow K{\quoi}.\leqno(6.2.7.1)
$$
对于通过应用函子从 \eqhyperref[XVII.eq6.2.7.1]{6.2.7.1} 或
\hyperref[XVII.thm6.2.3]{6.2.3} 得到的任意态射，我们将不加区别地称为
\emph{迹态射}、\emph{Gysin 态射}或\emph{态射} $u_{\ast}$；例如下面这些态射。

考虑交换图
$$
\xymatrix@C=1.5cm@R=1.75cm{
X  \ar[rr]^{u} \ar[rd]_{f}&   & Y\ar[ld]^{g}\\
                          & S &
}
$$
其中 $g$ 分离且局部有限型。对 $Y$ 上每个阿贝尔层 $\Ff $，迹态射定义一个态射
$$\etiquette[6.2.7.2]{XVII.eq6.2.7.2}
\Tr_{u}\quad\mbox{\rm 或}\quad u_{\ast}:f_{!}u^{\ast}\Ff \simeq
g_{!}u_{!}u^{\ast}\Ff \xrightarrow{g_{!}(\Tr_{u})} g_{!}\Ff {\quoi}.\leqno(6.2.7.2)
$$

若 $g$ 可紧化，$K$ 是上同调为挠的下有界复形，或为 $g^{\ast}\Aa$-模复形，
其中 $\Aa$ 是 $S$ 上的挠环层，则迹态射 \eqhyperref[XVII.eq6.2.7.1]{6.2.7.1}
同样诱导
\begin{align*}
\Tr_{u}\quad\mbox{\rm 或}\quad u_{\ast} &: Rf_{!}u^{\ast}K\simeq
Rg_{!}u_{!}u^{\ast}K\xrightarrow{Rg_{!}(\Tr_{u})}Rg_{!}K{\quoi},\tag{6.2.7.3}\\
\Tr_{u}\quad\mbox{\rm 或}\quad u_{\ast} &: R^{q}f_{!}u^{\ast}K\simeq
R^{q}g_{!}u_{!}u^{\ast}K\xrightarrow{R^{q}g_{!}(\Tr_{u})}
R^{q}g_{!}K{\quoi}.\tag{6.2.7.4}
\end{align*}\etiquette[6.2.7.3]{XVII.eq6.2.7.3}
\etiquette[6.2.7.4]{XVII.eq6.2.7.4}%The vertical space following this paragraph is too large.

当\pageoriginale $S$ 是代数闭域的谱时，\eqhyperref[XVII.eq6.2.7.4]{6.2.7.4} 是态射
$$\etiquette[6.2.7.5]{XVII.eq6.2.7.5}
u_{\ast}:\HH^{q}_{c}(X,u^{\ast}K)\longrightarrow \HH^{q}_{c}(Y,K){\quoi}.\leqno(6.2.7.5)
$$

由此可见，\emph{固有支集上同调关于拟有限态射是协变的}。

\subsubsection{}\etiquette[6.2.8]{XVII.subsubsec6.2.8}
设 $u:X\rightarrow Y$ 如 \hyperref[XVII.subsubsec6.2.7]{6.2.7}，
$\Delta_{n}=[0,n]$ 是含有 $\sisi{}{(}n+1\sisi{}{)}$ 个元素的标准集合，
$X_{n}$ 是 $X$ 在 $Y$ 上的 $(n+1)$ 重迭代纤维积。
各 $X_{n}$ 组成一个增广到 $Y$ 的单纯概形（\cf \sisi{VI${}_{\text{B}}$}{VBIS}）。
以 $u_{n}:X_{n}\rightarrow Y$ 表示增广态射。
若 $\Ff $ 是 $Y$ 上的阿贝尔层，则由 \eqhyperref[XVII.eq6.2.7.2]{6.2.7.2}
（作记号变换 $S\sisi{\rightarrow}{\mapsto} Y$、
$X\sisi{\rightarrow}{\mapsto} X_{n}$、$Y\sisi{\rightarrow}{\mapsto} Y$ 或 $X_{m}$ 后应用），
各层 $\Ff _{n}=u_{n!}u_{n}^{\ast}\Ff $ 组成增广到 $\Ff $ 的\emph{单纯}层。
按定义，此单纯层的“\emph{非退化链}”微分复形的第 $n$ 个分量是
$\Ff _{n}$ 的一个商，即以 $\Ff _{n}$ 为目标的退化算子的直和，
以及对称算子 $\sigma$ 的态射 $\sigma-\epsilon(\sigma)$ 的余核。

\begin{enonce*}{命题 6.2.9}\etiquette[6.2.9]{XVII.prop6.2.9}
在上述假设下，若 $u$ 是分离、平展、有限型且满的态射，
则对于 $Y$ 上每个阿贝尔层 $\Ff $，有
\begin{enumerate}
\renewcommand{\theenumi}{\roman{enumi}}
\renewcommand{\labelenumi}{(\theenumi)}
\item 各层 $\Ff _{n}$ 组成 $\Ff $ 的一个（左）解消。

\item 单纯层 $(\Ff _{n})$ 的非退化链微分复形是 $\Ff $ 的一个（左）解消。
若 $u$ 的所有几何纤维至多有 $d$ 个点，则此解消的长度 $\leq d$。
\end{enumerate}
\end{enonce*}

由\pageoriginale 基变换，只须在 $Y$ 是代数闭域的谱时验证此陈述；
此时 \hyperref[XVII.subsubsec6.2.8]{6.2.8} 归结为计算一个单纯形的单纯同调。

设 $(U_{i})_{i\in I}$ 是概形 $Y$ 的有限开覆盖，$\Ff $ 是 $Y$ 上的阿贝尔层；
对 $I$ 的每个子集 $P$，设 $j_{P}$ 为 $U_{P}=\cap_{i\in P}U_{i}$ 到
$\sisi{X}{Y}$ 的包含。假设 $I$ 上已选定全序。
若 $X=\coprod_{i\in I}U_{i}$，$u$ 是从 $X$ 到 $Y$ 的投影，
则 $\Ff $ 的解消 \hyperref[XVII.prop6.2.9]{6.2.9} (ii) 为
$$\etiquette[6.2.9.1]{XVII.eq6.2.9.1}
(\Ff _{n})_{n\,\text{deg}}=\bigoplus_{|P|=n+1}j_{P!}j^{\ast}_{P}\Ff {\quoi}.\leqno(6.2.9.1)
$$

当 $I$ 无限时，这个解消仍有意义。

\setcounter{subsubsection}{9}
\subsubsection{}\etiquette[6.2.10]{XVII.subsubsec6.2.10}
在 \hyperref[XVII.prop6.2.9]{6.2.9} 的假设下，设 $\Aa$ 是 $Y$ 上的挠环层，
$K\in \Ob^{-}(Y,\Aa)$。将 $X$ 看作赋环 $u^{\ast}\Aa$ 的概形。
因为解消 \hyperref[XVII.prop6.2.9]{6.2.9} (i)
（\resp \hyperref[XVII.prop6.2.9]{6.2.9} (ii)）对 $\Ff $ 具有函子性，
它给 $K$ 定义一个由双复形 $K_{1}$ 构成的解消。
此解消称为 $K$ 的非常 \emph{Čech 解消}
（\resp \emph{非常交错} Čech 解消）。

考虑交换图
$$
\xymatrix@C=1.5cm@R=1.75cm{
X  \ar[rr]^{u} \ar[rd]_{f}&   & Y\ar[ld]^{g}\\
                          & S & ,
}
$$
其中 $g$ 可紧化。沿用 \hyperref[XVII.subsubsec6.2.8]{6.2.8} 的记号，
若 $f_{n}$ 是从 $X_{n}$ 到 $S$ 的投影，则有 $Rf_{n!}=Rg_{!}\sisi{U}{u}_{n!}$。
因而，按\pageoriginale 半单纯次数过滤的复形 $K_{1}$ 定义一个谱序列
$$\etiquette[6.2.10.1]{XVII.eq6.2.10.1}
E^{-p,q}_{1}=R^{q}f_{p!}(u^{\ast}_{p}K)\Longrightarrow R^{-p+q} g_{!}
K{\quoi}.\leqno(6.2.10.1)
$$
微分 $d_{1}$ 来自单纯箭头，因而
$$\etiquette[6.2.10.2]{XVII.eq6.2.10.2}
E^{-p,q}_{2}=\check{\Hh}^{p}(R^{q}f_{p!}(u^{\ast}_{p}K)\sisi{}{)}\Longrightarrow
R^{-p+q}g_{!}K{\quoi}.\leqno(6.2.10.2)
$$

在 \eqhyperref[XVII.eq6.2.9.1]{6.2.9.1} 的假设下，类似地得到谱序列
$$\etiquette[6.2.10.3]{XVII.eq6.2.10.3}
E^{-p,q}_{1}=\bigoplus_{|P|=p+1}R^{q}f_{p!}(j^{\ast}_{p}K)\Longrightarrow R^{-p+q}g_{!}K{\quoi}.\leqno(6.2.10.3)
$$

谱序列 \eqhyperref[XVII.eq6.2.10.1]{6.2.10.1} 和
\eqhyperref[XVII.eq6.2.10.3]{6.2.10.3} 称为\emph{上局部化谱序列}。

\begin{enonce*}{命题 6.2.11}\etiquette[6.2.11]{XVII.prop6.2.11}
若 $f$ 是有限型分离平展态射，则函子 $f_{!}$ 是函子 $f^{\ast}$ 的左伴随，
其伴随箭头是迹态射。
\end{enonce*}

已知，若 $X$ 是景 $S$ 的一个对象，则将阿贝尔层限制到 $X$ 的函子 $f^{\ast}$
有左伴随 $f_{!}$。若 $F$ 是 $X$ 上的阿贝尔层，则 $f_{!}F$ 是由预层
$$
V\longrightarrow \bigoplus_{\xi \in \Hom(V,X)}F(\xi){\quoi}.
$$
生成的层。在平展情形，暂将此函子记为 $f\sisi{''}{^0}_{!}$。
其伴随性质立即蕴含其构造与任意基变换交换。
迹态射定义一个从 $f_{!}\Ff $ 到 $f^{0}_{!}\Ff $ 的态射 $t$，
与\pageoriginale 任意基变换相容。只须逐纤维说明 $t$ 是同构，
\ie 当 $S$ 为可分闭域的谱时验证，此时结论平凡。

\setcounter{subsubsection}{11}
\paragraph{}\etiquette[6.2.11.1]{XVII.par6.2.11.1}
须注意，$f^{0}_{!}$ 的这种解释专属于平展景的情形，
例如不能推广到预概形的 fpqf 景。

\subsection{连续系数的迹的变式}\etiquette[6.3]{XVII.subsec6.3}

本节证明 \textsc{Grothendieck} \cite[第 6 节]{sga4.exp.XVII.bib.Gr} 的断言，
并将这些结果用于为适当的 fppf 层定义“迹态射”。


\subsubsection{}\etiquette[6.3.1]{XVII.subsubsec6.3.1}
设 $A$ 是交换环。若 $L$ 和 $M$ 是两个有限生成投射 $A$-模，则从
$\Hom(L,M)$ 到 $\Hom(\bigwedge\limits^{m}L,\bigwedge\limits^{m}M)$
的映射 $u\mapsto \bigwedge\limits^{m}u$ 与任意基变换相容，并且是
$m$\sisi{次齐次}{\textsuperscript{次齐次}}的
（\sisi{\hyperref[XVII.subsubsec5.5.2]{5.5.2}}{\hyperref[XVII.par5.5.2.2]{5.5.2.2}}）。
由 \sisi{\hyperref[XVII.subsubsec5.5.2]{5.5.2}}{\hyperref[XVII.par5.5.2.3]{5.5.2.3}}，
它定义态射
$$\etiquette[6.3.1.1]{XVII.eq6.3.1.1}
\bigwedge^{m}:\Gamma^{m}\Hom(L,M)\longrightarrow
\Hom(\bigwedge^{m} L, \bigwedge^{m}M)\leqno(6.3.1.1)
$$
使得每次基变换之后，对 $u\in \Hom(L,M)$ 都有
$\bigwedge\limits^{m}(\gamma^{m}(u))=\bigwedge\limits^{m}u$。
因为 $\Hom(L,M)$ 平坦，此态射等同于
（\sisi{\hyperref[XVII.subsubsec5.5.2]{5.5.2}}{\hyperref[XVII.par5.5.2.5]{5.5.2.5}}）一个态射
$$\etiquette[6.3.1.2]{XVII.eq6.3.1.2}
\bigwedge^{m}:TS^{m}\Hom(L,M)\longrightarrow
\Hom(\bigwedge^{m}L, \bigwedge^{m}M){\quoi}.\leqno(6.3.1.2)
$$
注意，只要 $2$ 在 $A$ 中可逆，也可以定义这样的同态：
$\bigotimes\limits^{m}\Hom(L,M)$ 的一个对称元素定义从
$\bigotimes\limits^{m}L$ 到 $\bigotimes\limits^{m}M$ 的一个对称同态，
它下降到商上，定义从 $\bigwedge\limits^{m}L$ 到 $\bigwedge\limits^{m}M$ 的同态。

给定\pageoriginale 三个有限生成投射 $A$-模，图
$$
\vcenter{\etiquette[6.3.1.3]{XVII.eq6.3.1.3}
\xymatrix{
\Gamma^{m} \Hom(L,M) \otimes \Gamma^{m} \Hom(M,N)  \ar[r] \ar[d]^{\bigwedge\limits^{m}\otimes \bigwedge\limits^{m}}& \Gamma^{m}
\Hom(L,N) \ar[d]^{\bigwedge\limits^{m}}\\
\Hom(\bigwedge\limits^{m} L,
\bigwedge\limits^{m} M) \otimes
\Hom(\bigwedge\limits^{m} M,
\bigwedge\limits^{m} N)         \ar[r]     &
\Hom(\bigwedge\limits^{m} L,
\bigwedge\limits^{m} N),
}}\leqno{(6.3.1.3)}
$$
交换，其中水平箭头来自复合映射（第一个还来自（\hyperref[XVII.par5.5.2.8]{5.5.2.8}））。
这表达了公式
$\bigwedge\limits^{m}(vu)=\bigwedge\limits^{m}v\circ \bigwedge\limits^{m}u$。

若 $L$ 是秩为 $n$ 的自由 $A$-模，\eqhyperref[XVII.eq6.3.1.2]{6.3.1.2}
定义一个典范态射
$$\etiquette[6.3.1.4]{XVII.eq6.3.1.4}
\det:\Gamma^{n}\End_{A}(L)=TS^{n}\End_{A}(L)\longrightarrow
A{\quoi},\leqno(6.3.1.4)
$$
由 \eqhyperref[XVII.eq6.3.1.3]{6.3.1.3}，它是代数同态。按定义有
$$\etiquette[6.3.1.5]{XVII.eq6.3.1.5}
\det (u^{\otimes n})=\det (u\text{; }L){\quoi}.\leqno(6.3.1.5)
$$

若进一步 $L$ 是一个 $A$-代数 $B$ 上的模，则 \eqhyperref[XVII.eq6.3.1.4]{6.3.1.4}
通过复合定义代数同态
$$\etiquette[6.3.1.6]{XVII.eq6.3.1.6}
{\det}_{L}:TS^{n}_{A}(B)\longrightarrow A{\quoi}.\leqno(6.3.1.6)
$$

\begin{enonce*}{引理 6.3.2}\etiquette[6.3.2]{XVII.lem6.3.2}
假设 $L$ 出现在正合列
\[
\sisi{F}{E}:0\longrightarrow L'\longrightarrow L\longrightarrow
L''\longrightarrow 0
\]
中，其中 $L'$（\resp $L''$）局部自由，秩为 $n_{1}$（\resp $n_{2}$）。
令 $n=n_{1}+n_{2}$，设 $\alpha$ 是从扩张 $E$ 的自同态\pageoriginale 群
$\End(E)$ 到 $\End(L')\oplus \End(L'')$ 的“伴随分次”态射；
对两个平坦 $A$-模 $X$ 和 $Y$，设 $\beta$ 是由同构
（\hyperref[XVII.par5.5.2.6]{5.5.2.6}）
\[
TS^{n}(X\oplus Y)\simeq \sisi{\sum}{\bigoplus}_{n'+n''=n} TS^{n'\sisi{_{1}}{}}(X)\otimes
TS^{n''}(Y){\quoi}.
\]
得到的从 $TS^{n}(X\oplus Y)$ 到 $TS^{n_{1}}(X)\otimes TS^{n_{2}}(Y)$ 的态射。
则下图交换：{\footnotesize
\[
\xymatrix@C=1cm@R=1.5cm{
  TS^{n}(\End\sisi{)}{} (E)) \ar[d] \ar[r]^(.4){TS^{n}(\alpha)} & TS^{n}(\End(L')\oplus
\End(L''))  \ar[r]^(.45){\beta} & TS^{n_1} (\End(L')) \otimes TS^{n_2}
(\End(L'')) \ar[d]^{\det \otimes \det}\\
TS^{n}(\sisi{L}{\End(L)})         \ar[rr]^{\det} &                 & A.
}
\]
}
\end{enonce*}

由 $TS^{n}$ 的泛性质 \hyperref[XVII.subsubsec5.5.2]{5.5.2}，
只须验证对 $u\in \End(E)$，$u^{\otimes n}$ 的两个像相等。
若 $\alpha{u}=(u',u'')$，这两个像分别为 $\det(u)$ 和 $\det(u').\det(u'')$，
这就证明了断言。

\begin{enonce*}{引理 6.3.3}\etiquette[6.3.3]{XVII.lem6.3.3}
设 $A$ 是\sisi{}{交换}环，$B$ 是\sisi{}{交换} $A$-代数，$C$ 是未必交换的
$B$-代数，$E$ 是在 $A$ 上自由且秩为 $n$ 的 $B$-模，
$F$ 是在 $B$ 上自由且秩为 $m$ 的 $C$-模。假设 $C$ 在 $A$ 上平坦。
由 $\Gamma^{nm}$ 的泛性质 \hyperref[XVII.par5.5.2.2]{5.5.2.2}
和 \hyperref[XVII.par5.5.2.3]{5.5.2.3}，存在唯一的箭头
$\sigma:TS^{nm}_{A}(C)\rightarrow TS^{n}_{A}TS^{\sisi{n}{m}}_{B}(C)$，
使每次环扩张 $A\rightarrow A'$ 后，都有
$\sigma(c^{\otimes nm})=(c^{\otimes m})^{\otimes n}$。于是下图交换：
$$
\xymatrix@C=2cm@R=1.5cm{
TS^n_A (B) \ar[d]^{\det_E}  &  & TS^{n}_A TS^{m}_B (C) \ar[ll]_{TS^n_A
  (\det_F)}\\
A    & & TS^{nm}_A (C) \ar[ll]_{\det_{E\otimes_B F}} \ar[u]_{\sigma}.
}
$$
\end{enonce*}

\begin{enonce*}[remark]{证明}
由\pageoriginale \hyperref[XVII.par5.5.2.2]{5.5.2.2} 和
\hyperref[XVII.par5.5.2.3]{5.5.2.3}，只须说明 $TS^{nm}_{A}(C)$ 的元素
$u^{\otimes nm}$ 的两个像相等。这两个像是 $\det_{E}(\det_{F}(u))$ 和
$\det_{E\otimes_{B}F}(u)$；它们可由 $A$、$B$、$E$、$F$ 及
$u\in \End_{B}(F)$ 表示。要验证它们相等，可以先用它在 $\End(E)$ 中的
（交换）像 $B_{1}$ 代替 $B$，并用 $F\otimes_{B}B_{1}$ 代替 $F$。
此时待证公式就是分块行列式公式（\textsc{Bourbaki}，
\sisi{Alg.}{\emph{Alg.}} \sisi{第}{\chap}~8 章 \S~12，引理 2，
\sisi{第 140 页}{第~140 页}）。
\end{enonce*}

当不再要求 $C$ 在 $A$ 上平坦时，同样的证明更一般地给出代数的交换图
$$\vcenter{\etiquette[6.3.3.1]{XVII.eq6.3.3.1}
\xymatrix@C=3cm@R=1.5cm{
\Gamma^{n}_A (B)  \ar[d]^{\det_E}& \Gamma^{n}_A
\Gamma^m_{B}(C) \ar[l]_{\Gamma_A^{n}(\det_F)}\\
A            & \ar[l]_{\det_{E\otimes_B F}}  \Gamma^{nm}_A (C) \ar[u]_{\sigma}.
}}\leqno{(6.3.3.1)}
$$

\setcounter{subsubsection}{3}
\subsubsection{}\etiquette[6.3.4]{XVII.subsubsec6.3.4}
若 $B$ 是交换 $A$-代数，则 $\Spec(TS^{n}(B))$ 恰为 $\Spec(B)$ 在 $\Spec(A)$
上的第 $n$ 个对称幂，\ie $(\Spec(B)/\Spec(A))^{n}$ 模对称群的商
（\SGA 1 \sisi{IV 3}{V 1}）。

若\pageoriginale $X$ 是\sisi{}{在 $S$ 上}分离的概形，
则 $X$ 上的拟凝聚 $\Oo_{X}$-模 $\Ff $，若其支集在 $S$ 上有限，
且其直像在 $S$ 上局部自由、秩有限（\resp 秩为 $n$），
则称它\emph{在} $S$ \emph{上有限局部自由}（\resp 秩为 $n$）。
这里 $\Ff $ 的\emph{支集}仍指由 $\Ff $ 的零化子定义的 $X$ 的\emph{子概形}
（该零化子是拟凝聚理想，因为 $\Ff $ 是有限型拟凝聚层）。

若 $\Ff $ 在 $S$ 上有限局部自由，秩为 $n$，支集为 $E$，则 $E$ 在 $S$ 上仿射。
在 $S$ 上整体化，$\Ff $ 借助 \eqhyperref[XVII.eq6.3.1.6]{6.3.1.6} 定义
$\Sym^{n}_{S}(E)$ 的截面 $\sigma_{E}(\Ff )$，从而由函子性，定义代数空间
\footnote[1]{\MakeLinkTarget*{sga4.exp.XVII.footnote.6.3.4.artin-1}\label{XVII.noteArtin1}
M. \textsc{Artin}，Algebraization of formal moduli I\sisi{}{
，载于 \emph{Global Analysis: Papers in Honor of K.~Kodaira}，Princeton University Press，1969，21--71}。
我们只在 $\Sym^{n}_{S}(X)$ 作为概形存在时使用上述构造；当 $X$ 在 $S$ 上拟射影时即如此。}
$\Sym^{n}_{S}(X)$ 的一个截面 $\sigma(\Ff )$，
$$\etiquette[6.3.4.1]{XVII.eq6.3.4.1}
\sigma(\Ff )\in \Gamma(\Sym^{n}_{S}(X)/S){\quoi}.\leqno(6.3.4.1)
$$

特别地，若 $f:S'\rightarrow S$ 有限局部自由，秩为 $n$，则取 $\Ff =\Oo_{S'}$，
得到典范截面
$$\etiquette[6.3.4.2]{XVII.eq6.3.4.2}
t\in\Gamma(\Sym^{n}_{S}(S')/S){\quoi}.\leqno(6.3.4.2)
$$

\begin{enonce*}{命题 6.3.5}\etiquette[6.3.5]{XVII.prop6.3.5}
设 $X$ 是在 $S$ 上分离的概形，$\Ff _{1}$ 和 $\Ff _{2}$ 是两个
在 $S$ 上有限局部自由的 $\Oo_{X}$-模，秩分别为 $n_{1}$ 和 $n_{2}$。
令 $n=n_{1}+n_{2}$。设 $d$ 是典范箭头
$d:\Sym^{n_{1}}_{S}(X)\times \Sym^{n_{2}}_{S}(X)\rightarrow \Sym^{n}_{S}(X)$。
若 $\Ff $ 是 $\Ff _{1}$ 被 $\Ff _{2}$ 的扩张，则
$$
\sigma(\Ff )=d(\sigma(\Ff _{1}),\sigma(\Ff _{2})){\quoi}.
$$
\end{enonce*}

可以\pageoriginale 用 $\Ff _{1}\oplus \Ff _{2}$ 的支集（它在 $S$ 上有限）代替 $X$，
并假设 $S$ 仿射。将命题译成环论语言，它就由 \hyperref[XVII.lem6.3.2]{6.3.2} 得到。

\begin{enonce*}{引理 6.3.6}\etiquette[6.3.6]{XVII.lem6.3.6}
若 $X$ 是在 $S$ 上分离的概形，$\Ff $ 是在 $S$ 上有限局部自由、秩为 $n$ 的
$\Oo_{X}$-模，$\Ll$ 是 $X$ 上的可逆层，则
$$
\sigma(\Ff )=\sigma(\Ff \otimes\Ll){\quoi}.
$$
\end{enonce*}

用 $\Ff $ 的支集（\hyperref[XVII.subsubsec6.3.4]{6.3.4}）代替 $X$\sisi{，}{。}
在 $S$ 上局部地，$\Ll$ 同构于 $\Oo_{X}$，断言便显然成立。

读者可验证：

\begin{enonce*}{命题 6.3.7}\etiquette[6.3.7]{XVII.prop6.3.7}
对每个在 $S$ 上分离的 $S$-概形图
\[
\xymatrix{
X        \ar[rr]^{u} \ar[rd]_{f} &        &  Y  \ar[ld]^{g}\\
                                & S      &,
}
\]
以及每个在 $S$ 上有限局部自由、秩为 $n$ 的 $\Oo_{X}$-模 $\Ff $，
$\Oo_{Y}$-模 $u_{\ast}\Ff $ 在 $S$ 上有限局部自由，秩为 $n$，且
\[
\sisi{u}{(\Sym^n_S u)}(\sigma(\Ff ))=\sigma(u_{\ast}\Ff ){\quoi}.
\]
\end{enonce*}

\setcounter{subsubsection}{7}
\paragraph{}\etiquette[6.3.7.1]{XVII.par6.3.7.1}
最后注意，态射 $\sigma$ \eqhyperref[XVII.eq6.3.4.1]{6.3.4.1} 和
$t$ \eqhyperref[XVII.eq6.3.4.2]{6.3.4.2} 的构造\emph{与任意基变换相容}；
由 \hyperref[XVII.par5.5.2.7]{5.5.2.7}，这一说法有意义。

\subsection*{应用 1．光滑曲线的情形}\pageoriginale

\subsubsection{}\etiquette[6.3.8]{XVII.subsubsec6.3.8}
设 $f:X\rightarrow S$ 是 $S$ 上的光滑拟射影曲线，\ie 纯相对维数为 $1$
的光滑拟射影态射。Hilbert 概形理论
\sisi{[}{}\textsc{Grothendieck} \cite{sga4.exp.XVII.bib.Gr}）表明，
fpqc 层 $\sheaf{\Div}^{n}_{X/S}$ 可表；它给每个 $S$-概形 $S'$ 赋予
$X'=\fprod{X}{S'}{S}$ 上在 $S$ 上\emph{有限}且次数为 $n$ 的相对除子所成的集合。

对 $X$ 上在 $S$ 上有限且次数为 $n$ 的每个相对除子 $E$，有典范元素
\sisi{(6.3.4.3)}{\eqhyperref[XVII.eq6.3.4.2]{6.3.4.2}}
$t(E)\in \Sym^{n}_{S}(E)$，其在 $\Sym^{n}_{S}(X)$ 中的像为 $\sigma(E)$。
因为这一构造与任意基变换相容，它定义
$$\etiquette[6.3.8.1]{XVII.eq6.3.8.1}
\sigma:\sheaf{\Div}^{n}_{X/S}\longrightarrow
\Sym^{n}_{S}(X):E\longrightarrow \sigma(\Oo_{E}){\quoi}.\leqno(6.3.8.1)
$$

$f$ 的每个截面 $s$ 都定义一个在 $S$ 上有限且次数为 $1$ 的相对除子，
即除子 $s(S)$。态射
\[
(X/S)^{n}\longrightarrow \sheaf{\Div}^{n}_{X/S}:(s_{i})_{0<i\leq
n}\longmapsto \sum s_{i}(S)
\]
对称，因而诱导态射
$$\etiquette[6.3.8.2]{XVII.eq6.3.8.2}
\alpha:\Sym^{n}_{S}(X)\longrightarrow \sheaf{\Div}^{n}_{X/S}{\quoi}.\leqno(6.3.8.2)
$$

\begin{enonce*}{命题 6.3.9}\etiquette[6.3.9]{XVII.prop6.3.9}
\hyperref[XVII.subsubsec6.3.8]{6.3.8} 的态射 $\sigma$ 和 $\alpha$ 是互逆同构。
\end{enonce*}

\begin{enonce*}[remark]{证明}
\begin{enumerate}
\renewcommand{\theenumi}{\alph{enumi}}
\renewcommand{\labelenumi}{\theenumi)}
\item $\sigma\alpha=\Id$。只须证明复合箭头
$$
(X/S)^{n}\xrightarrow{\pi}\Sym^{n}_{S}(X)\xrightarrow{\alpha}\sheaf{\Div}^{n}_{X/S}\xrightarrow{\sigma}\Sym^{n}_{S}(X)
$$
是\pageoriginale $\pi$。当 $n=1$ 时断言显然；一般情形由下面的引理得到。

\begin{enonce*}{引理 6.3.9.1}\etiquette[6.3.9.1]{XVII.lem6.3.9.1}
沿用 \hyperref[XVII.prop6.3.5]{6.3.5} 和 \eqhyperref[XVII.eq6.3.8.1]{6.3.8.1} 的记号，有
\[
\sigma(E+F)=d(\sigma(E),\sigma(F)){\quoi}.
\]
\end{enonce*}

此引理由 \hyperref[XVII.prop6.3.5]{6.3.5}、\hyperref[XVII.lem6.3.6]{6.3.6}
以及正合列
\[
0\longrightarrow \Oo(-F)\otimes \Oo_{E}\longrightarrow
\Oo_{E+F}\longrightarrow \Oo_{F}\longrightarrow 0{\quoi}.
\]
得到。

\item 已知 $\sheaf{\Div}^{n}_{X/S}$ 在 $S$ 上光滑
（曲线上除子的无穷小提升没有障碍（\cf \SGA 3 XI 1））。
态射 $\alpha\pi:(X/S)^{n}\rightarrow \sheaf{\Div}^{n}_{X/S}$
是光滑 $S$-概形\sisi{}{之间}的拟有限满态射，其源的纤维等维。
因而它平坦（\EGA ${\rm 0_{IV}}$ 17.3.5），所以 $\alpha\pi$ 是 fppf 层的满态射，
从而 $\alpha$ 也是。结合 a)，便完成 \hyperref[XVII.prop6.3.9]{6.3.9} 的证明。
\end{enumerate}
\end{enonce*}

\begin{enonce*}{引理 6.3.10}\etiquette[6.3.10]{XVII.lem6.3.10}
设 $k$ 是交换环，$M$ 是自由 $k$-模，$u$ 是 $M$ 的自同态。
以 $N_{\det(u-T)=0}$ 表示从 $k[T]/\det(u-T)$ 到 $k$ 的范数。
则对 $P\in k[T]$，有
\[
N_{\det(u-T)=0} (P(T))=\det(P(u),M){\quoi}.
\]
\end{enonce*}

这归结为验证 $u$ 的矩阵各元素与 $P$ 的系数之间的一个代数恒等式。
因此，为证明 \hyperref[XVII.lem6.3.10]{6.3.10}，只须
（由熟知的论证）处理 $k$ 代数闭且 $u$ 半单、特征值为 $a_{i}$ 的情形。

此时\pageoriginale 有
\begin{align*}
& \det(u-T)=\Pi(a_{i}-T)\\
& N_{\det(u-T)=0}(P(T))=\Pi P(a_{i})\quad\mbox{\rm 且}\\
& \det(P(u),M)=\Pi P(a_{i}){\quoi}.
\end{align*}

\setcounter{subsubsection}{11}
\setcounter{paragraph}{-1}
\paragraph{}\etiquette[6.3.11.0]{XVII.par6.3.11.0}
设 $f:X\rightarrow S$ 是有限表现的拟射影态射，平坦且纤维是纯一维
Cohen--Macaulay 概形；因此，当 $X$ 是 $S$ 上的光滑拟射影曲线时，这些条件满足。
设 $\Ff $ 是\emph{在} $S$ \emph{上}有限局部自由的 $\Oo_{X}$-模。
因为 $\Ff $ 在 $X$ 的一个逐纤维稠密开子集上为零，存在典范箭头
$\psi:\Oo_{X}\rightarrow \det^{\ast}(\Ff )$
（\SGA 6 XI 或 \textsc{Mumford}
\sisi{[1 第~5 章 \S~3]}{\cite[\chap~5 \S~3]{sga4.exp.XVII.bib.Mu}}）。
以 $\delta(\Ff )$ 表示由 $\coker(\psi)$ 的零化理想定义的除子（同上）；
此除子在 $S$ 上有限。

\begin{enonce*}{命题 6.3.11.1}\etiquette[6.3.11.1]{XVII.prop6.3.11.1}
在 \hyperref[XVII.par6.3.11.0]{6.3.11.0} 的假设下，沿用
\hyperref[XVII.subsubsec6.3.8]{6.3.8} 和 \eqhyperref[XVII.eq6.3.4.1]{6.3.4.1} 的记号，
若 $\Ff $ 的秩为 $n$，则 $\delta(\Ff )$ 的次数为 $n$，且
\[
\sigma(\Ff )=\sigma(\delta(\Ff )){\quoi}.
\]
\end{enonce*}

\begin{enonce*}[remark]{证明}
\begin{enumerate}
\renewcommand{\theenumi}{\alph{enumi}}
\renewcommand{\labelenumi}{\theenumi)}
\item \emph{$f$ 光滑的情形。}由标准归约，可以假设 $S$ 是剩余域代数闭的 Artin 局部概形。
$\sigma$ 的可加性（\hyperref[XVII.prop6.3.5]{6.3.5}）使我们可以假设
$\Ff $ 的支集集中于 $X$ 的一个有理点 $x$。
立即可验证，若 $T$ 是 $x$ 处的局部一致化参数，定义
$p:X\rightarrow \sisi{A}{\AAA}^{1}_{S}$，则
$p_{\ast}(\Oo_{\sisi{S}{\delta}(\Ff )})\simeq \Oo_{\delta(p_{\ast}\Ff )}$
（因为所考虑的构造与完备化交换）。结合 \hyperref[XVII.prop6.3.7]{6.3.7}，
便归结到 $X=\sisi{A'}{\AAA^1}_{S}$ 的情形。

若\pageoriginale $X=\AAA\sisi{'}{^1}_{S}$ 且 $S$ 仿射，$S=\Spec(k)$，
则 $\Ff $ 可等同于赋有自同态 $u$（即 $T$ 的作用）的局部自由 $k$-模 $M$。
作为 $k[T]$-模，$M$ 有解消
\[
0\rightarrow
k[T]\otimes_{k}M\xrightarrow{\varphi}k[T]\otimes_{k}M\xrightarrow{\epsilon}
M\rightarrow 0{\quoi},
\]
其中 $\epsilon(P\otimes m)=P(u)(m)$，
$\varphi(1\otimes m)=1\otimes um-T\otimes m$。
由公式 $\varphi=u-T$，可知除子 $\delta(\Ff )$ 的方程为 $\det(u-T,M)=0$。

为验证 \sisi{6.3.11}{\hyperref[XVII.prop6.3.11.1]{6.3.11.1}}，
只须说明在每次基变换之后，对 $\Oo_{X}$ 的每个截面 $\varphi$，有
\[
{\det}_{k}(\varphi,M)={\det}_{k}(\varphi,\Oo_{\delta(\Ff )})(=N_{\delta(\Ff )/S}(\varphi)){\quoi}.
\]
当 $\sisi{f}{\varphi}=P(T)$ 时，这正是 \hyperref[XVII.lem6.3.10]{6.3.10} 的断言。


\item 一般情形。我们先以下面更具体的形式，明确陈述所要证明的断言：

\begin{enonce*}{推论 6.3.11.2}\etiquette[6.3.11.2]{XVII.cor6.3.11.2}
在关于 $f:X\rightarrow S$ 的假设 \hyperref[XVII.par6.3.11.0]{6.3.11.0} 下，
设 $M$ 是 $X$ 上的有限局部自由模，$u$ 是 $M$ 的自同态，使 $\det u$ 是
$\Oo_{X}$ 的一个在每个纤维 $X_{s}$ 上都正则的截面
（从而 $\Coker u=M/u(M)$ 和 $\Coker \det u=\Oo_{X}/\det u$
都在 $S$ 上平坦且有限表现，且有\sisi{拟有限纤维}{拟有限支集}）。
假设这些层的共同支集\nde{$M/uM$ 在 \hyperref[XVII.subsubsec6.3.4]{6.3.4}
意义下的支集，是 $\Oo/\det u$ 的支集的一个由幂零理想定义的子概形。}
在 $S$ 上有限，从而它们定义在 $S$ 上有限局部自由的层。
则对 $\Oo_{X}$ 的每个截面 $\varphi$，$\varphi$ 作用于 $\Coker u$ 与
$\Coker \det u$ 的、在 $\Oo_{S}$ 上的行列式相等：
$$\etiquette[6.3.11.3]{XVII.eq6.3.11.3}
{\det}_{S}(\varphi,M/u(M))={\det}_{S}(\varphi,\Oo_{X}/\det
u){\quoi}.\leqno(6.3.11.3)
$$
\end{enonce*}


我们\pageoriginale 注意到，由齐次性，此陈述蕴含所考虑的两个局部自由模在 $S$ 上的秩相等
（将 $\varphi$ 换成 $T\varphi$，其中 $T$ 是 $\Oo_{S}$ 的一个“未定”截面）。
此外，%\c
将 $\varphi$ 换成 $1+T\varphi$，其中 $\sisi{j}{T}$ 是这样的未定元
（更确切地说，作基变换 $S'=S[T]\rightarrow S$），
便从 \eqhyperref[XVII.eq6.3.11.3]{6.3.11.3} 推出关于特征多项式的表面上更一般的公式
$$\etiquette[6.3.11.4]{XVII.eq6.3.11.4}
\Pol_{S}(\varphi,M/uM)(T)=\Pol_{S}(\varphi,\Oo_{X}/\det u)(T)\leqno(6.3.11.4)
$$
成立。

为证明 \eqhyperref[XVII.eq6.3.11.3]{6.3.11.3}，将归结到已经证明的光滑情形。
通过在 $S$ 上平展局部化，归结到 $S$ 严格局部、闭点为 $s$ 的情形；
再由可加性，归结到 $M/u(M)$ 的支集在 $X_{s}$ 中只有一个点 $x$ 的情形。
对底作平坦扩张之后，可以假设剩余域扩张 $k(x)/k(s)$ 平凡；
再将 $X$ 换成 $x$ 的充分小邻域，并给 $\varphi$ 加上一个在 $M/u(M)$ 及
$\Oo_{X}/ \det u$ 上诱导零运算的函数（这不改变待证公式），
容易看出可以假设 $\varphi$ 诱导 $\Oo_{X_{s}}$ 的一个在 $x$ 处正则的截面，
从而（缩小 $X$ 后）定义一个拟有限平坦态射
$$
p:X\longrightarrow X'=S[T]{\quoi}.
$$
所以，通过在 $X$ 和 $X'$ 上平展局部化，可以假设有如上的\emph{有限平坦}态射 $p$，
其中 $X'$ 在 $S$ 上光滑，且 $\varphi=p^{\ast}(\varphi')$，
$\varphi'$ 是 $\Oo_{X'}$ 的一个截面。令
\begin{align*}
& M'=p_{\ast}(M)\quav u'=p_{\ast}(u){\quoi},\\
& P=\Coker u\quav Q=\Coker \det u{\quoi},\\
& P'=\Coker u'\quav Q'=\Coker \det u'{\quoi}.
\end{align*}
我们\pageoriginale 知道
\[
\det u'=N_{\sisi{X'/X}{X/X'}}\det u,\quad\mbox{\rm 故}\quad Q'=\Coker
N_{\sisi{X'/X}{X/X'}}(\det u){\quoi}.
\]
显然有
\[
{\det}_{S}(\varphi,P)={\det}_{S}(\varphi',P\sisi{}{'}){\quoi},
\]
而由已在光滑情形证明的公式 \eqhyperref[XVII.eq6.3.11.3]{6.3.11.3}，有
\[
{\det}_{S}(\varphi',P')={\det}_{S}(\varphi',Q'){\quoi},
\]
所以尚须证明公式
\[
{\det}_{S}(\varphi',Q')={\det}_{S}(\varphi,Q){\quoi},
\]
它\emph{等价于}所要建立的公式 \eqhyperref[XVII.eq6.3.11.3]{6.3.11.3}。
由此可见，后者是否成立仅依赖于 $\Oo_{X}$ 的截面 $\det u$；
等价地，它等价于将 $M$ 换成 $\Oo_{X}$、将 $u$ 换成 $\det u$ 时的同一公式。
但当 $M=\Oo_{X}$ 时，公式 \eqhyperref[XVII.eq6.3.11.3]{6.3.11.3} 是重言式！
证明至此完成。
\end{enumerate}
\end{enonce*}

\begin{enonce*}[remark]{注 6.3.11.5}\etiquette[6.3.11.5]{XVII.rem6.3.11.5}
在\pageoriginale $S$ 为 Artin 概形的情形，$\Coker u$ 与 $\Coker \det u$
（沿用 \hyperref[XVII.cor6.3.11.2]{6.3.11.2} 的记号）在 $S$ 上秩相等，
是 \EGA IV 21.10.17.3 的特例；\eqhyperref[XVII.eq6.3.11.3]{6.3.11.3}
可视为它的一个“相对”版本。
\end{enonce*}

\begin{enonce*}{命题 6.3.12}\etiquette[6.3.12]{XVII.prop6.3.12}
设 $X$ 和 $Y$ 是 $S$ 上的两条有限表现的拟射影曲线，平坦且纤维为
Cohen--Macaulay 概形，并设 $u:X\rightarrow Y$ 是拟有限平坦态射。

则下图
\[
\xymatrix@C=4em{
\sheaf{\Div}^{n}_{X/S}  \ar[r]^{u_{\ast}} \ar[d]_{\sigma_X} &
\sheaf{\Div}^{n}_{Y/S} \ar[d]_{\sigma_Y}\\
\Sym^{n}_S(X)     \ar[r]^{\Sym^{n}(u)} & \Sym^{n}_S(Y)
}
\]
交换\footnote[1]{\MakeLinkTarget*{sga4.exp.XVII.footnote.6.3.12.1}\label{XVII.note192}
因为我们只考虑在 $S$ 上\emph{有限}的除子，态射 $u_{\ast}$ 有定义。
例如，可取 $u_{\ast}(D)=\delta(u_{\ast}\Oo_{D})$ 作为定义；
当 $u$ 有限时，这一公式与 \EGA IV 21.5.5 一致。}。
\end{enonce*}

\sisi{\noindent}{}%
因为 $u_{\ast}(D)=\delta(u_{\ast}\Oo_{D})$，\hyperref[XVII.prop6.3.12]{6.3.12}
由 \hyperref[XVII.prop6.3.7]{6.3.7} 和 \hyperref[XVII.prop6.3.11.1]{6.3.11.1} 得到。

\subsection*{应用 2．迹态射}

\setcounter{subsubsection}{12}
\subsubsection{}\etiquette[6.3.13]{XVII.subsubsec6.3.13}
设 $G$ 是概形 $S$ 的大 fppf 景上的阿贝尔层（\hyperref[XVII.subsec0.10]{0.10}）。
假设 $G$ 满足下述条件。

\paragraph{}\etiquette[6.3.13.1]{XVII.par6.3.13.1}
\emph{对 $S$-概形的每个有限表现仿射平坦态射 $f:X\rightarrow Y$，
以及每个整数 $n\geq 1$，态射}
\[
\Hom_{S}(\Sym^{n}_{Y}(X),G)\longrightarrow \Hom_{S}((X/Y)^{n},G)
\]
\emph{是\pageoriginale 单射，且其像为从 $(X/Y)^{n}$ 到 $G$ 的对称态射所成的集合。}

此条件在基变换下稳定。

设 $f:T\rightarrow S$ 是有限局部自由态射，秩为 $n$。
设 $\pr_{i}$ 是从 $(T/S)^{n}$ 到 $T$ 的第 $i$ 个投影，$t$ 是
$\Sym^{n}_{S}(T)$ 的截面 \sisi{(6.3.4.3)}{\eqhyperref[XVII.eq6.3.4.2]{6.3.4.2}}。
若 $u\in G(T)$，则 $\sum \pr^{\ast}_{i}(u)$ 是
$G\sisi{}{(}(T/S)^{n}$) 的对称元素，因而定义
\[
u_{n}\in \sisi{}{G(}\Sym^{n}_{S}(T)\sisi{}{)}{\quoi}.
\]
令
\[
\Tr_{f}(u)=t^{\ast}(u_{n}){\quoi}.
\]

当 $f$ 有限局部自由时，在 $S$ 上局部地用上述公式定义 $\Tr_{f}(u)$。
这一构造与任意基变换相容，因而定义\emph{迹态射}
$$\etiquette[6.3.13.2]{XVII.eq6.3.13.2}
\Tr_{f}:f_{\ast}f^{\ast}G\longrightarrow G{\quoi}.\leqno(6.3.13.2)
$$

\subsubsection{}\etiquette[6.3.14]{XVII.subsubsec6.3.14}
下列情形满足条件（\hyperref[XVII.par6.3.13.1]{6.3.13.1}）：
\begin{enumerate}
\renewcommand{\theenumi}{\alph{enumi}}
\renewcommand{\labelenumi}{\theenumi)}
\item 若 $G$ 可表，由对称幂的定义即知。

\item $G$ 是 $S$ 的小平展景（\hyperref[XVII.subsec0.10]{0.10}）上一个层
在 $S_{\text{\rm fppf}}$ 上的逆像。

\item $G$ 由拟凝聚 $\Oo_{S}$-模 $G_{0}$ 定义。
\end{enumerate}

断言 c) 归结到（\hyperref[XVII.par5.5.2.7]{5.5.2.7}）。

\begin{enonce*}{命题 6.3.15}\etiquette[6.3.15]{XVII.prop6.3.15}
对\pageoriginale 满足（\hyperref[XVII.par6.3.13.1]{6.3.13.1}）的 $G$，有：
\begin{enumerate}
\renewcommand{\theenumi}{\roman{enumi}}
\renewcommand{\labelenumi}{(\theenumi)}
\item 态射 $\Tr_{f}$ 对 $G$ 具有函子性，因而可加。

\item 其构造与任意基变换 $S'\rightarrow S$ 相容。

\item 若 $T=\coprod_{i}T_{i}$，从而 $f_{\ast}G\simeq \sum f_{i\ast}G$，
则 $\Tr_{f}=\sum \Tr_{f_{i}}$。

\item 若 $T$ 在 $S$ 上常秩为 $n$，则复合态射
\[
G\longrightarrow f_{\ast}f^{\ast}G\xrightarrow{\Tr_{f}} G
\]
为乘以 $n$ 的态射。

\item 迹态射与复合相容\emph{（\cf \hyperref[XVII.thm6.2.3]{6.2.3} Var 3）。}
\end{enumerate}
\end{enonce*}

断言 (i) 和 (ii) 显然，(iii) 是 \hyperref[XVII.prop6.3.5]{6.3.5} 的简单推论。
沿用 \hyperref[XVII.subsubsec6.3.13]{6.3.13} 的记号，若 $u\in G(S)$，
则 $u_{n}$ 是 $u$ 的拉回的 $n$ 倍，故 $\Tr_{f}(f^{\ast}u)=t^{\ast}(u_{n})=nu$。

为证明 (v)，在 $S$ 上平展局部化并利用 (iii)，便归结为证明：对
\[
f=gh:S''\xrightarrow{h} S'\xrightarrow{g}S{\quoi},
\]
其中 $\sisi{f}{g}$ 和 $\sisi{g}{h}$ 有限局部自由，\emph{常秩}分别为 $n$ 和 $m$，图
$$\vcenter{\etiquette[6.3.15.1]{XVII.eq6.3.15.1}
\xymatrix{
G(S'')    \ar[rr]^{\Tr_h}\ar[rd]_{\Tr_f} &        &G(S') \ar[ld]^{\Tr_g}\\
                &  G(S)         &
}}\leqno{(6.3.15.1)}
$$
交换。对 $u\in G(S'')$，考虑图
$$\vcenter{\etiquette[6.3.15.2]{XVII.eq6.3.15.2}
\xymatrix@C=2.5cm@R=1.75cm{
\Sym^{n}_S(S')  \ar[rr]^{\Sym^{n}(t)} \ar[rd]^{\Tr_{\sisi{g}{h}}(u)_n} & & \Sym^{n}_S
\Sym^{m}_{S'} (S'')  \ar[dd] \ar[ld]^{(u_m)_n}\\
                                          & G &\\
          S   \ar@<3pt>[ru]^{\Tr_{\sisi{f}{g}}\Tr_{\sisi{g}{h}}(u)} \ar@<-3pt>[ru]_{\Tr_{\sisi{fg}{gh}}(u)}  \ar[rr]^{t}  \ar[uu]_{t}   &        & \Sym^{nm}_S(S'') \ar[lu]_{u_{nm}}.
}}\leqno{(6.3.15.2)}
$$
右侧\pageoriginale 的竖直箭头有定义，因为 $S\sisi{'}{''}$ 在 $S$ 上平坦。
每个三角形都交换，因而 (v) 将由外框的交换性推出。

译成环论语言，这由 \hyperref[XVII.lem6.3.3]{6.3.3} 的特例 $B=E$、$C=F$ 得到。

\setcounter{subsubsection}{15}
\subsubsection{}\etiquette[6.3.16]{XVII.subsubsec6.3.16}
现在设 $f:T\rightarrow S$ 是\sisi{分离的}{}有限表现拟有限平坦态射。
若 $ℱ$ 是 $T$ 上的 fppf 层，定义它在 $S$ 上的固有支集直像
$f_{!}\sisi{T}{ℱ}$ 为下述预层：它给 $S'/S$ 赋予 $\sisi{F}{ℱ}$ 在 $T'$ 上
支集在 $S'$ 上固有的截面所成的集合。由 \hyperref[XVII.prop6.1.1]{6.1.1}，此预层是层。
沿用 \hyperref[XVII.prop6.1.1]{6.1.1} 的记号，也可将它表示成一个局部正向极限
（\sisi{IV 14.\quad}{\hyperrefext[V][V.subsec8.2]{8.2}}）：
\[
f_{!}\sisi{F}{ℱ(U)}=\varinjlim_{V\in\Phi (U)}f_{\ast}(\sisi{F}{ℱ}|V){\quoi},
\]
此公式\pageoriginale 解释了一个相当显然的事实：$f_{!}$ 的性质归结为 $f'$ 有限时
$f'_{\ast}$ 的性质。由此得到：

\begin{enonce*}{命题 6.3.17}\etiquette[6.3.17]{XVII.prop6.3.17}
存在唯一的方式，对每个\sisi{}{分离的}有限表现拟有限平坦态射 $f:T\rightarrow S$
和 $S$ 上满足（\hyperref[XVII.par6.3.13.1]{6.3.13.1}）的每个 fppf 层 $G$，定义迹态射
$$
\Tr_{f}:f_{!}f^{\ast}G\longrightarrow G
$$
满足：
\begin{enumerate}
\renewcommand{\theenumi}{\roman{enumi}}
\renewcommand{\labelenumi}{(\theenumi)}
\item 当 $f$ 有限时，得到的正是态射 \eqhyperref[XVII.eq6.3.13.2]{6.3.13.2}。

\item 经适当调整后，\hyperref[XVII.thm6.2.3]{6.2.3} 的条件
\emph{(Var 1)} 至 \emph{(Var 4)} 满足。
\end{enumerate}
\end{enonce*}

拟有限态射 $f$ 诱导景的交换图（\cf（\hyperref[XVII.subsec0.11]{0.11}）），
其中 $\alpha_{\ast}$ 是“限制”函子：
$$
\xymatrix{
T_{\text{\rm fppf}} \ar[rr]^{\alpha} \ar[d]^{f}& & T_{\text{et}} \ar[d]^{f_{\et}}\\
S_{\text{\rm fppf}} \ar[rr]^{\alpha}& & S_{\et}.
}
$$
对于 $T_{\text{\rm fppf}}$ 上的 $G$，有
$$
f_{!}\alpha_{\ast}\simeq \alpha_{\ast}f_{!}{\quoi}.
$$
还\pageoriginale 有“基变换”态射
$$
c:f^{\ast}_{\text{\rm et}}\alpha_{\ast}\longrightarrow
\alpha_{\ast}f^{\ast}{\quoi}.
$$
下面迹态射 \hyperref[XVII.thm6.2.3]{6.2.3} 与 \hyperref[XVII.prop6.3.17]{6.3.17}
的相容性的验证，留给读者。对 $S_{\text{\rm fppf}}$ 上满足
（\hyperref[XVII.par6.3.13.1]{6.3.13.1}）的 $G$，下图交换：
$$\MakeLinkTarget*{sga4.exp.XVII.equation.6.3.17.1}\etiquette[6.3.17.1]{XVII.eq6.3.17.1}
\vcenter{\xymatrix{
f_{\et!}f^{\ast}_{\et}\alpha_{\ast} G \ar[r]^{f_{!}c} \ar[d]_{\Tr_f}^{\hyperref[XVII.thm6.2.3]{(6.2.3)}}&
f_{\et!}\alpha_{\ast}f^{\ast}G  \ar@{<->}[r]^{\sim}&
\alpha_{\ast}f_!f^{\ast}G \ar[d]_{\alpha_{\ast}\Tr_f}^{\hyperref[XVII.prop6.3.17]{(6.3.17)}}\\
\alpha_{\ast}G                   \ar@{=}[rr] &        & \alpha_{\ast}G.
}}\leqno(6.3.17.1)
$$

更具体地，这就是说：若 $G$ 是 $S$ 上的 fppf 层，其在平展景上的限制定义平展层
$\alpha_{\ast}G$，若 $u$ 是逆像 $f^{\ast}_{\text{\rm et}}(\alpha_{\ast}G)$
的固有支集截面，其在 $f^{\ast}G$ 中的像为 $u_{1}$，
则元素 $\Tr_{f}(u)$（按 \hyperref[XVII.thm6.2.3]{6.2.3} 的意义）
与 $\Tr_{f}(u_{1})$（按 \hyperref[XVII.prop6.3.17]{6.3.17} 的意义）在 $G(S)$ 中相等。

\begin{enonce*}[remark]{例 6.3.18}\etiquette[6.3.18]{XVII.exem6.3.18}
设 $f:S'\rightarrow S$ 是有限局部自由态射，秩为 $n$，且
$u\in \Gamma(S',\Oo_{S'})=\Hom_{S}(S',\EE^{1}_{S})$。
\nde{这里 $\EE^1_S$ 表示 $S$ 上的仿射直线。}态射 $u$ 定义一个态射
\[
\Sym(u):\Sym^{n}_{S}(S')\longrightarrow
\Sym^{n}_{S}(\EE^{1}_{S})=\Spec(S[\sigma_{1}\ldots
\sigma_{n}]\sisi{)}{\quoi},
\]
其中各 $\sigma_{i}$ 是基本对称函数（\cf \textsc{Bourbaki}，\emph{Alg.} 5，
附录 I 第~2 节\nde{此引用待核对。新版中可能是 IV \S~6 第~1 节。}）。
于是有（简单计算留给读者）：
$$\etiquette[6.3.18.1]{XVII.eq6.3.18.1}
[\Sym(u)\circ t]^{\ast}(\sigma_{i})=\sigma_{i}(u){\quoi},\leqno(6.3.18.1)
$$
其中\pageoriginale $t\in \Gamma$ $(\Sym^{n}_{S}(S')/S)$ 是典范元素
\sisi{(6.3.4.3)}{\eqhyperref[XVII.eq6.3.4.2]{6.3.4.2}}，而 $\sigma_{i}(u)$ 表示
$f_{\ast}\Oo_{S'}$ 上乘以 $u$ 的特征多项式的第 $i$ 个系数。
由此，“迹”映射
\[
\begin{cases}%raghu
\GG_{a}(S')\longrightarrow \GG_{a}(S)\\
\GG_{m}(S')\longrightarrow \GG_{m}(S)
\end{cases}
\]
分别是通常的\emph{迹}与\emph{范数}\sisi{}{}；事实上，由 $\EE^{1}_{S}$ 上的
加法（\resp 乘法）诱导的态射
$\Sym^{n}_{S}(\EE^{1}_{S})\rightarrow \EE^{1}_{S}$，由
$\Sym^{n}_{S}(\EE^{1}_{S})\simeq \Spec(S[\sigma_{1},\ldots,\sigma_{n}])$
上的函数 $\sigma_{1}$（\resp $\sigma_{n}$）给出。

\sisi{\noindent}{}%
于是由 \hyperref[XVII.eq6.3.17.1]{(6.3.17.1)}，当 $n$ 在 $S$ 上可逆时，平展层图
$$\vcenter{\etiquette[6.3.18.2]{XVII.eq6.3.18.2}
\xymatrix{
0 \ar[r]& f_{\ast}\mu_n \ar[r] \ar[d]_{\Tr_f}^{\hyperref[XVII.thm6.2.3]{(6.2.3)}}& f_{\ast}\sisi{G}{\GG}_m
\ar[r] \ar[d]^{N_f}
& f_{\ast}\sisi{G}{\GG}_m \ar[d]^{N_f} \ar[r]& 0\\
0 \ar[r]& \mu_n \ar[r]        & \relax\sisi{G}{\GG}_m    \ar[r]     & \relax\sisi{G}{\GG}_m  \ar[r]       & 0
}}\leqno{(6.3.18.2)}
$$
交换。
\end{enonce*}

\begin{enonce*}{命题 6.3.19}\etiquette[6.3.19]{XVII.prop6.3.19}
考虑交换图
\[
\xymatrix{
X \ar[rr]^{u} \ar[rd]_{p}&   & Y\ar[ld]^{q}\\
  & S &
}
\]
其中\pageoriginale $X$ 与 $Y$ 是 $S$ 上的光滑拟射影曲线，$u$ 是拟有限态射。
\begin{enumerate}
\renewcommand{\theenumi}{\roman{enumi}}
\renewcommand{\labelenumi}{(\theenumi)}
\item 若 $D$ 是 $X$ 上在 $S$ 上有限的除子（见注\emph{第~\pageref{XVII.note192} 页}），
$g$ 是 $G_{Y}$ 在 $Y$ 上的截面，则
\[
\Tr_{D/S}(u^{\ast}g)=\Tr_{(u_{\ast}D)/S}(g){\quoi}.
\]

\item 若 $E$ 是 $Y$ 上在 $S$ 上有限的除子，$g$ 是 $G_{X}$ 在 $X$ 上的截面，
且 $u$ 有限，则
\[
\Tr_{E/S}(\Tr_{X/Y}(g))=\Tr_{(u\sisi{_{\ast}}{^\ast}E)/S}(\sisi{G}{g}){\quoi}.
\]
\end{enumerate}
\end{enonce*}

\begin{enonce*}[remark]{证明}
\begin{enumerate}
\renewcommand{\theenumi}{\roman{enumi}}
\renewcommand{\labelenumi}{(\theenumi)}
\item 设 $g_{n}$ 是由 $g$ 得到的 $G$ 在 $\Sym^{n}_{S}(Y)$ 上的截面
（\hyperref[XVII.subsubsec6.3.13]{6.3.13}）。由 \hyperref[XVII.prop6.3.12]{6.3.12}，有
\begin{align*}
\Tr_{D/S}(u^{\ast}g)\overset{\dfn}{=}
\sigma(D)^{\ast}(u^{\ast}g)_{n} & =\sigma(D)^{\ast}(\sisi{u}{\Sym^n_S(u)}^{\ast}g_{n})=\sisi{u}{(\Sym^n_S(u)}\circ
\sigma(D))^\ast
(g_{n})\\
& =\sigma(u_{\ast}D)^\ast(g_{n})=\Tr_{u_{\ast}D}(g_{n}){\quoi}.
\end{align*}

\item 公式 (ii) 由 \hyperref[XVII.prop6.3.15]{6.3.15} (ii) 和 (v) 得到：
\[
\Tr_{E/S}(\Tr_{X/Y}(g))=\Tr_{E/S}(\Tr_{u^{\ast}E/E}(g))=\Tr_{u^{\ast}E/S}(g){\quoi}.
\]
\end{enumerate}
\end{enonce*}

\subsection*{应用 3．挠子的迹}

\setcounter{subsubsection}{19}
\subsubsection{}\etiquette[6.3.20]{XVII.subsubsec6.3.20}
设 $G$ 是 $S$ 的大 fppf 景上的层，满足（\hyperref[XVII.par6.3.13.1]{6.3.13.1}）。
设 $u:X\rightarrow Y$ 是 $S$-概形的有限局部自由态射，$K$ 是 $X$ 上的
$G_{X}$-挠子（torsor）。假设在 $Y$ 上就 fppf 拓扑而言局部地，
$X$ 上的挠子 $K$ 平凡。当 $K$ 在平展拓扑下局部平凡时，即为这种情形。

\begin{enonce*}{命题 6.3.21}\etiquette[6.3.21]{XVII.prop6.3.21}
在\pageoriginale \hyperref[XVII.subsubsec6.3.20]{6.3.20} 的假设下，
在唯一同构意义下恰有一个 $Y$ 上的 $G_{Y}$-挠子 $\Tr_{u}(K)$，
赋有态射 $\Tr_{u}$ 或 $\Tr^{K}_{u}:u_{\ast}K\rightarrow \Tr_{u}(K)$，
对 $Y$ 上 $G$ 的局部截面 $g$（按 fppf 意义）及 $u_{\ast}K$ 的局部截面 $k$ 满足
$$
\Tr^{K}_{u}(gk)=\Tr_{u}(g)\cdot\Tr_{u}\sisi{}{^K}(k){\quoi}.
$$
\end{enonce*}

问题在 $Y$ 上是局部的，因而可以假设 $K$ 平凡：$K=G_{X}$。
取 $K=G_{Y}$，并取态射 \eqhyperref[XVII.eq6.3.13.2]{6.3.13.2} 为迹态射，
便得到问题的一个解。立即可验证，它是唯一的解。

特别地，当 $G$ 满足下述条件时，迹函子
\sisi{\hyperref[XVII.subsubsec6.3.20]{6.3.20}}{\hyperref[XVII.prop6.3.21]{6.3.21}} 有定义。

\setcounter{subsubsection}{21}
\paragraph{对于每个 $S$-概形 $X$，$X$ 上的每个 $G_{X}$-挠子都在平展拓扑下局部平凡}
\etiquette[6.3.21.1]{XVII.par6.3.21.1}

当 $G$ 是光滑群、$G$ 是 $S$ 的小平展景上一个层的逆像，
或 $G$ 由 $S$ 上的拟凝聚层定义时，这个条件满足。

由构造有
$$\etiquette[6.3.21.2]{XVII.eq6.3.21.2}
\Tr_{u}(K)=u_{\ast}(K) \times_{u_{\ast}G_{X}}G_{Y}{\quoi}.\leqno(6.3.21.2)
$$

\subsubsection{}\etiquette[6.3.22]{XVII.subsubsec6.3.22}
为在更一般的假设下定义迹函子，关键是：对于 $S$-概形的有限表现拟射影平坦态射
$f:X\rightarrow T$，定义一个合理的函子 $K\mapsto K_{n}$，
从 $X$ 上的 $G_{X}$-挠子到 $Z=\Sym^{n}_{T}(X)$ 上的 $G_{Z}$-挠子。
我们\pageoriginale 先有如下引理，其中 $f_{n}$ 表示从 $\Sym^{n}_{T}(X)$ 到 $T$ 的投影；
其证明与 \hyperref[XVII.prop6.3.21]{6.3.21} 相同：

\hypertarget{sga4.exp.XVII.lemma.6.3.22.1}{}\begin{enonce*}{引理 6.3.22.1}\etiquette[6.3.22.1]{XVII.lem6.3.22.1}
沿用上述记号，若挠子 $K$ 在 $T$ 上局部平凡，则在唯一同构意义下，存在唯一的
$\Sym^{n}_{T}(X)$ 上的挠子 $K_{n}$，赋有态射
$\alpha:f_{\ast}K\rightarrow f_{n{\ast}}(K_{n})$，
使对 $f_{\ast}G_{X}$ 的局部\emph{（fppf）}截面 $g$ 及
$\sisi{u}{f}_{\ast}K$ 的局部截面 $k$，有
$\alpha(gk)=g_{n}\sisi{\circ}{\cdot} \alpha(k)$
（$g_{n}$ 的记号见 \hyperref[XVII.subsubsec6.3.13]{6.3.13}；
这里 $g_{n}$ 的定义使用平坦性假设\etc）。
\end{enonce*}

一般情形将用一个略显刻意的下降论证处理。首先有：

\hypertarget{sga4.exp.XVII.proposition.6.3.23}{}\begin{enonce*}{命题 6.3.23}\etiquette[6.3.23]{XVII.prop6.3.23}
考虑交换图
$$\vcenter{\etiquette[6.3.23.1]{XVII.eq6.3.23.1}
\xymatrix{
X \ar[rr]^{u} \ar[rd]& & Y\ar[ld]\\
  & T &
}}\leqno{(6.3.23.1)}
$$
其中各对象是在 $T$ 上平坦、有限表现且拟射影的概形，$u$ 平展且满。
则概形的半单纯系统 $\Sym^{n}_{T}((X/Y)^{k+1})$，
是概形 $\Sym^{n}_{T}(Y)$ 的平展拓扑的一个超覆盖。
\end{enonce*}

\begin{enonce*}[remark]{证明}
\begin{enumerate}
\renewcommand{\theenumi}{\alph{enumi}}
\renewcommand{\labelenumi}{\theenumi)}
\item 我们将使用以下辅助结果。

\setcounter{subsubsection}{23}\setcounter{paragraph}{1}
\paragraph{}\etiquette[6.3.23.2]{XVII.par6.3.23.2}
\emph{设 $(X,Y,T)$ 如上，$t\in T$，$Z$ 是一个 $T$-概形，
它是有限多个 $T$-概形 $\Spec(k_{i})$
的不交并，其中 $k_{i}$ 是 $k_{t}$ 的有限扩张\nde{这里 $k_t$ 表示 $T$ 在 $t$ 处的剩余域。}，
并\pageoriginale 考虑 $T$-概形的交换图}
$$\vcenter{\etiquette[6.3.23.3]{XVII.eq6.3.23.3}
\xymatrix{
              & Z \ar@{^{(}->}[ld]_{x} \ar@{^{(}->}[rd]^{y}     \\
X_t \ar[rr]^{u}&    & Y_t,
}}\leqno{(6.3.23.3)}
$$
\emph{其中 $x$ 和 $y$ 是闭浸入。设 $z\in \Sym^{n}_{T}(Z)$，
$x(\sisi{Z}{z})$ 和 $y(\sisi{Z}{z})$ 是它在 $\Sym^{n}_{T}(X)$
与 $\Sym^{n}_{T}(Y)$ 中的像。则 $\Sym^{n}_{T}(u)$ 在 $x(\sisi{Z}{z})$ 处平展。}

概形 $\Sym^{n}_{T}(X)$ 的构造与任意基变换
$\sisi{S'}{T'}\rightarrow \sisi{S}{T}$ 相容（\hyperref[XVII.par5.5.2.7]{5.5.2.7}）。
归结到 Noether 情形，可知 $\Sym^{n}_{T}(X)$ 在 $\sisi{S}{T}$ 上有限表现；
此外，它在 $\sisi{S}{T}$ 上平坦（\hyperref[XVII.par5.5.2.4]{5.5.2.4}）。
由逐纤维平坦性判别法，为证明（\hyperref[XVII.par6.3.23.2]{6.3.23.2}），
可归结到 $\sisi{S}{T}$ 为域的谱的情形。
此时，$\Sym^{n}_{T}(X)$ 在 $\sisi{Z}{z}$ 处的完备化的仿射代数，
是概形 $\Sym^{n}_{T}(Z_{k})$ 在 $x(z)$ 处局部化的仿射代数的逆向极限，
其中 $Z_{k}$ 遍历 $Z$ 在 $X$ 中的无穷小邻域。
由假设，$u$ 在 $Z$ 于 $X$ 中的无穷小邻域与于 $Y$ 中的无穷小邻域之间诱导同构。
因此，$\Sym^{n}_{T}(u)$ 在 $x(z)$ 和 $y(z)$ 处的完备化是同构，
这就证明了（\hyperref[XVII.par6.3.23.2]{6.3.23.2}）。

由 \hyperref[XVII.par6.3.23.2]{6.3.23.2}，将 $\Sym^{n}_{T}(Y)$ 在 $y(z)$ 处的
Hensel 化对 $z\in \Sym^{n}_{T}(Z)$ 取不交并所得的概形，
仅依赖于 $Z$ 以及 $T$-概形 $X^{h}_{Z}$；后者是 $X$ 在 $Z$ 的各点处的 Hensel 化的不交并。
我们将此概形记为 $\Sym^{h\,n}_{T}(X^{h}_{Z})$；它对 $X^{h}_{Z}$ 具有函子性。

\item 通过\pageoriginale 取极限，为证明 \hyperref[XVII.prop6.3.23]{6.3.23}，
只须说明：对于在其纤维中闭的 $y\in \Sym^{n}_{T}(Y)$，
将单纯概形 $\Sym^{n}_{T}((X/Y)^{k+1})$ 拉回到 $\Sym^{n}_{T}(Y)$ 在此点处的
Hensel 化后，得到平展拓扑的一个超覆盖。
这个问题在 $T$ 上对平展拓扑是局部的；因此归结到存在图
\eqhyperref[XVII.eq6.3.23.3]{6.3.23.3}，且 $y=x(z)$ 的情形
（先用 $k_{t}$ 的可分扩张消去剩余域扩张）。
此时，所考虑的半单纯系统由系统
$$
\Sym^{h\,n}_{T}((X^{h}_{u^{-1}(Z)}/Y^{h}_{Z})^{k+1}),\quad\mbox{\rm 增广到}\quad \Sym^{h\,n}_{T}(Y^{h}_{Z})
$$
经基变换得到，所以只须考虑后者。
半单纯系统 $(X^{h}_{u^{-1}(Z)}/Y^{h}_{Z})^{k+1}$ 与常值系统 $Y^{h}_{Z}$
半单纯同伦，因为 $X^{h}_{u^{-1}(Z)}/Y^{h}_{Z}$ 有截面。
所以系统 $\Sym^{h\,n}_{T}((X^{h}_{u^{-1}(Z)}/Y^{h}_{Z})^{k+1})$
也同伦于 $\Sym^{h\, n}_{T} (Y^{h}_{Z})$，特别地，是它的一个超覆盖。
这就证明了 \sisi{6.3.2.3}{\hyperref[XVII.prop6.3.23]{6.3.23}}。
\end{enumerate}
\end{enonce*}

\subsubsection{}\etiquette[6.3.24]{XVII.subsubsec6.3.24}
设 $f:Y\rightarrow T$ 是 $S$-概形之间的有限表现拟射影平坦态射，
$K$ 是 $Y$ 上的 $G_{Y}$-挠子。考虑下述问题：

\setcounter{paragraph}{-1}
\paragraph{}\etiquette[6.3.24.0]{XVII.par6.3.24.0}
对每个态射 $u:X\rightarrow Y$，其中 $X$ 在 $T$ 上平坦、拟射影且有限表现，
定义 $\Sym^{n}_{T}(X)$ 上的一个挠子 $(u^{\ast}K)_{n}$；
并对每个 $v:X'\rightarrow X$ 定义同构
$c_{v}:(v^{\ast}u^{\ast}K)_{n}\simeq \Sym^{n}_{T}(v)^{\ast}(u^{\ast}K)_{n}$，使得
\begin{enumerate}
\renewcommand{\theenumi}{\roman{enumi}}
\renewcommand{\labelenumi}{(\theenumi)}
\item 同构\pageoriginale $c_{v}$ 与复合相容；

\item 当 $u^{\ast}K$ 在 $T$ 上局部平凡时，得到的正是
\hyperref[XVII.lem6.3.22.1]{6.3.22.1} 的挠子及其函子性。
\end{enumerate}

由 \hyperref[XVII.prop6.3.23]{6.3.23}，对每个从 $K$ 通过
\hyperref[XVII.par6.3.24.0]{6.3.24.0} 所考虑类型的逆像得到的挠子，
问题 \hyperref[XVII.par6.3.24.0]{6.3.24.0} 的解的存在性与唯一性，
在 $Y$ 上对平展拓扑是局部问题。
特别地，若 $G$ 满足（\hyperref[XVII.par6.3.21.1]{6.3.21.1}），
则问题 \hyperref[XVII.par6.3.24.0]{6.3.24.0} 有唯一解。

现在假设 $G$ 满足：

\paragraph*{(6.3.24.1)}\MakeLinkTarget*{sga4.exp.XVII.paragraph.6.3.24.1}\etiquette[6.3.24.1]{XVII.par6.3.24.1}
\emph{对每个整数 $n$ 和每个仿射 $S$-概形图
\[
\xymatrix{
X \ar[rr]^{u} \ar[rd]_{f}& & Y\ar[ld]^{g}\\
  &T&,
}
\]
其中 $f$ 和 $g$ 平坦且有限表现，$u$ 有限、局部自由且满，
增广到 $\Sym^{n}_{T}(Y)$ 的截断半单纯概形
$(\Sym^{n}_{T}((X/Y)^{k+1}))_{0\leq k\leq 2}$，
是 $G$-挠子范畴的一个下降图\emph{（\textsc{Giraud} \cite{sga4.exp.XVII.bib.Gi}）}。}

\paragraph*{6.3.24.1 bis}\etiquette[6.3.24.1 bis]{XVII.par6.3.24.1bis}
由 \hyperref[XVII.prop6.3.23]{6.3.23}，用标准方式可知，
条件（\ref{XVII.par6.3.24.1}）蕴含如下类似条件：只假设 $f$ 和 $g$
平坦、拟射影且有限表现，而 $u$ 平坦且满。
事实上，在平展拓扑下局部地，$u$ 有有限局部自由的拟截面；
将 \hyperref[XVII.prop6.3.23]{6.3.23} 和 \textsc{Giraud} \cite{sga4.exp.XVII.bib.Gi}
应用于组成下降数据方形的图，便得结论。

给定\pageoriginale 在 $T$ 上平坦、拟射影且有限表现的概形 $Y$ 上的一个
$G_{Y}$-挠子 $K$，按定义存在忠实平坦的有限表现态射 $u:X\rightarrow Y$，
使 $u^{\ast}K$ 平凡。\nde{可以假设 $T$ 仿射。
此时存在所述的 $u$，甚至可以假设它仿射，以便随后应用前一段。}
若 $u_{i}$ 表示从 $(X/Y)^{i+1}$ 到 $Y$ 的投影 $(i\geq 0)$，
则概形 $\Sym^{n}_{T}((X/Y)^{i+1})$ 上的挠子 $(u^{\ast}_{i}K)_{n}$
（\hyperref[XVII.lem6.3.22.1]{6.3.22.1}），在这一增广到 $\Sym^{n}_{T}(Y)$
的半单纯概形上组成一个下降数据 $(i=0,1,2)$。假设 $G$ 满足条件：

\paragraph*{(6.3.24.2)}\MakeLinkTarget*{sga4.exp.XVII.paragraph.6.3.24.2}\etiquette[6.3.24.2]{XVII.par6.3.24.2}
\emph{上述构造的下降数据均有效}（\cf \ref{XVII.par6.3.24.1}）。

此时，以 $K_{n}$ 表示该下降数据所定义的 $\Sym^{n}_{T}(Y)$ 上的挠子
（在唯一同构意义下确定）。

\subsubsection{}\etiquette[6.3.25]{XVII.subsubsec6.3.25}
假设 $G$ 满足（\hyperref[XVII.par6.3.21.1]{6.3.21.1}），或满足条件
（\ref{XVII.par6.3.24.1}）和（\ref{XVII.par6.3.24.2}）。
下面的形式规则的明确含义及其验证，留给读者：

\paragraph{}\etiquette[6.3.25.1]{XVII.par6.3.25.1}
\emph{基变换}。设
\[
\xymatrix{
X' \ar[rr]^{g'} \ar[d]_{f'}& & X \ar[d]^{f}\\
T' \ar[rr]_{g}& & T.
}
\]
是交换方形。以 $g'_{n}$ 表示从 $\Sym^{n}_{T'}(X')$ 到 $\Sym^{n}_{T}(X)$ 的箭头。
对于 $X$ 上的一个 $G$-挠子 $K$，有
\[
{g'_n}^{\ast} K_{n}\simeq ({g'}^{\ast}K)_{n}{\quoi}.
\]

现在设\pageoriginale $f:X\rightarrow T$ 是态射，$K$ 是 $X$ 上的 $G_{X}$-挠子。

\paragraph{}\etiquette[6.3.25.2]{XVII.par6.3.25.2}
\emph{群的变换。}设 $r:G_{1}\rightarrow G_{2}$ 是群之间的态射，两者都满足
（\hyperref[XVII.par6.3.21.1]{6.3.21.1}），或
（\ref{XVII.par6.3.24.1}）、（\ref{XVII.par6.3.24.2}）。有
\[
r(K_{n})\simeq r(K)_{n}{\quoi}.
\]

\paragraph{}\etiquette[6.3.25.3]{XVII.par6.3.25.3}
\emph{对 $K$ 的可加性}\sisi{：}{。}有
\[
(K'+K'')_{n}\simeq K'_{n}+K''_{n}{\quoi}.
\]

\paragraph{}\etiquette[6.3.25.4]{XVII.par6.3.25.4}
\emph{对 $n$ 的可加性}。设 $n_{1}$ 和 $n_{2}$ 是两个整数，$n=n_{1}+n_{2}$，并设
\[
\alpha:\fprod{\Sym_{T}{}^{n_{1}}(X)}{\Sym_{T}{}^{n_{2}}(X)}{T}\longrightarrow
\Sym^{n}_{T}(X)
\]
为典范箭头；则
\[
\alpha^{\ast}K_{n}\simeq \pr^{\ast}_{1}
K_{n_{1}}+\pr^{\ast}_{2}K_{n_{2}}{\quoi}.
\]

\paragraph{}\etiquette[6.3.25.5]{XVII.par6.3.25.5}
\emph{对 $n$ 的乘法性。}设 $n$ 和 $m$ 是两个整数，$\alpha$ 是从
$\Sym^{n}_{T}(\Sym^{m}_{T}(X))$ 到 $\Sym^{nm}_{T}(X)$ 的典范箭头
（\cf \hyperref[XVII.lem6.3.3]{6.3.3}）。有
\[
\alpha^{\ast}(K_{nm})\simeq (K_{m})_{n}{\quoi}.
\]

\subsubsection{}\etiquette[6.3.26]{XVII.subsubsec6.3.26}
假设 $G$ 满足（\hyperref[XVII.par6.3.21.1]{6.3.21.1}），或
（\ref{XVII.par6.3.24.1}）、（\ref{XVII.par6.3.24.2}）。
若 $f:X\rightarrow T$ 是 $S$-概形之间的秩为 $n$ 的有限局部自由态射，
$t$ 是 $\Sym^{n}_{T}(X)$ 的截面 \hyperref[XVII.eq6.3.4.2]{(6.3.4.3)}，
且 $K$ 是 $X$ 上的 $G_{X}$-挠子，则用公式
$$\etiquette[6.3.26.1]{XVII.eq6.3.26.1}
\Tr_{f}(K)=t^{\ast}(K_{n}){\quoi}.\leqno(6.3.26.1)
$$
定义 $T$ 上的一个 $G_{T}$-挠子 $\Tr_{f}(K)$。

若\pageoriginale $f$ 局部自由，则在 $T$ 上局部地，用同一公式定义 $\Tr_{f}(K)$。
在共同适用的范围内，通过典范且具函子性的同构，此定义与
\hyperref[XVII.prop6.3.21]{6.3.21} 的定义等价
（定义 \hyperref[XVII.prop6.3.21]{6.3.21} 给出了所提出的下降问题的解）。
由（\hyperref[XVII.par6.3.25.1]{6.3.25.1}）、（\hyperref[XVII.par6.3.25.2]{6.3.25.2}）
及（\hyperref[XVII.par6.3.25.3]{6.3.25.3}）立即可知，函子 $\Tr_{f}$
\emph{与基变换 $T'\rightarrow T$ 相容，与群的变换 $r:G_{1}\rightarrow G_{2}$ 相容，
并且与挠子的加法相容。}

图 \eqhyperref[XVII.eq6.3.15.2]{6.3.15.2} 的外框的交换性，结合
（\hyperref[XVII.par6.3.25.1]{6.3.25.1}）和
（\hyperref[XVII.par6.3.25.4]{6.3.25.4}），给出\emph{与态射复合的相容性}：
$$
\Tr_{fg}(K)\simeq \Tr_{f}(\Tr_{g}(K)\sisi{}{)}{\quoi}.
$$

\subsubsection{}\etiquette[6.3.27]{XVII.subsubsec6.3.27}
在相对曲线的特殊情形中，还有：

\begin{enonce*}{构造 6.3.27.1}\etiquette[6.3.27.1]{XVII.const6.3.27.1}
\emph{若} $f:X\rightarrow S$ 是 $S$ 上的光滑拟射影曲线，$D_{1}$ 和 $D_{2}$
是在 $S$ 上有限的两个相对除子，$K$ 是 $X$ 上的一个 $G$-挠子，
则令 $\Tr_{D/S}(K)=\Tr_{D/S}(K|D)$，有典范同构
$$
\Tr_{D_{1}+D_{2}}(K)\simeq \Tr_{D_{1}}(K)+\Tr_{D_{2}}(K){\quoi}.
$$
\end{enonce*}

\sisi{\noindent}{}%
此同构立即由同构（\hyperref[XVII.par6.3.25.4]{6.3.25.4}）及
\hyperref[XVII.lem6.3.9.1]{6.3.9.1} 得到。最后，\hyperref[XVII.prop6.3.19]{6.3.19}
的论证还给出：

\begin{enonce*}{构造 6.3.27.2}\etiquette[6.3.27.2]{XVII.const6.3.27.2}
考虑交换图
$$
\xymatrix{
X   \ar[rr]^{u} \ar[rd]_{p}&     &  Y \ar[ld]^{q}\\
     &   S &,
}
$$
其中\pageoriginale $X$ 和 $Y$ 是 $S$ 上的光滑拟射影曲线，$u$ 是拟有限态射。
\begin{enumerate}
\renewcommand{\theenumi}{\roman{enumi}}
\renewcommand{\labelenumi}{(\theenumi)}
\item 若 $D$ 是 $X$ 上在 $S$ 上有限的除子，$K$ 是 $Y$ 上的 $G$-挠子，
则有\emph{（\cf 第~\pageref{XVII.note192} 页的注）}：
$$
\Tr_{D/S}(u^{\ast}K)\simeq \Tr_{(u_{\ast}D)/S}(K){\quoi}.
$$

\item 若 $E$ 是 $Y$ 上在 $S$ 上有限的除子，$K$ 是 $X$ 上的 $G$-挠子，
且 $u$ 有限，则有
$$
\Tr_{E/S}(\Tr_{X/Y}(K))\simeq \Tr_{u^{\ast}E/S}(K){\quoi}.
$$
\end{enumerate}
\end{enonce*}

\hypertarget{sga4.exp.XVII.proposition.6.3.28}{}\begin{enonce*}{命题 6.3.28}\etiquette[6.3.28]{XVII.prop6.3.28}
设 $A$ 和 $B$ 是两个（交换）$R$-代数，$u:A\rightarrow B$ 是同态，
使 $B$ 成为忠实、有限型、局部自由的 $A$-模。
则代数的\sisi{}{余}半单纯复形
$(TS^{n}_{R}(\bigotimes\limits^{k+1}_{A}B))_{k\geq 0}$，
作为微分\sisi{}{余}半单纯 $TS^{n}_{R}(A)$-模，是 $TS^{n}_{R}(A)$ 的分裂解消。
\end{enonce*}

\begin{enonce*}[remark]{证明}
与\sisi{}{余}半单纯模 $TS^{n}_{R}(\bigotimes\limits^{k+1}_{A}B)$ 伴随的模复形，
仅依赖于赋有模同态 $u:A\rightarrow B$ 的 $A$-模 $B$。
作为 $A$-模，$B$ 可写成 $B=A\oplus L$，由此可显式写出一个同伦算子。
\end{enonce*}

回顾：

\hypertarget{sga4.exp.XVII.proposition.6.3.29}{}\begin{enonce*}{命题 6.3.29}\etiquette[6.3.29]{XVII.prop6.3.29}
设 $A$ 是环，$B^{\ast}$ 是余半单纯 $A$-代数，$c:A\rightarrow B^{\ast}$ 是结构态射。
要使模的逆像函子（$=$ 标量扩张）$\epsilon^{\ast}$ 忠实（\resp 全忠实），
只须 $A$-模的序列
$$
0\longrightarrow A\longrightarrow B^{0}
$$
（\resp\pageoriginale $0\rightarrow A\rightarrow B^{0}\xrightarrow{j_{0}-j_{1}}B^{1}$）
正合，且与任意 $A$-模作张量积后仍正合。
\end{enonce*}

证明立即可得，留给读者。

由（\hyperref[XVII.prop6.3.28]{6.3.28}）和
（\hyperref[XVII.prop6.3.29]{6.3.29}），用标准方式可得：

\hypertarget{sga4.exp.XVII.corollary.6.3.30}{}\begin{enonce*}{推论 6.3.30}\etiquette[6.3.30]{XVII.cor6.3.30}
每个仿射群 $G$ 都满足条件（\ref{XVII.par6.3.24.1}）。
\end{enonce*}

（利用 $G$-挠子可由一个在 $T$ 上仿射的概形表示这一事实。）

将 \hyperref[XVII.cor6.3.30]{6.3.30} 推广到每个可由代数空间表示的群 $G$，
应该不难\footnote[1]{\MakeLinkTarget*{sga4.exp.XVII.footnote.6.3.30.artin-2}\label{XVII.noteArtin2}
M.~\textsc{Artin} 的 Algebraic Spaces。Yale University，Whittemore Lectures，1969 年 5 月。}。
我猜想，每个可由平坦代数空间表示的群都满足条件（\ref{XVII.par6.3.24.2}）。
我们有一个略带诙谐的陈述：

\hypertarget{sga4.exp.XVII.proposition.6.3.31}{}\begin{enonce*}{命题 6.3.31}\etiquette[6.3.31]{XVII.prop6.3.31}
设 $S$ 上下列序列正合：
$$
0\rightarrow G\xrightarrow{r}G'\xrightarrow{u}H\rightarrow 0{\quoi}.
$$
假设 $G'$ 和 $H$（从而 $G$）满足（\hyperref[XVII.par6.3.13.1]{6.3.13.1}），
$G'$ 满足（\hyperref[XVII.par6.3.21.1]{6.3.21.1}），
且 $G$ 满足（\ref{XVII.par6.3.24.1}）。则 $G$ 满足（\ref{XVII.par6.3.24.2}）。
\end{enonce*}

例如，这适用于作为可表光滑群之间满态射的核的仿射群。

\begin{enonce*}{6.3.31 的证明}
\emph{设\pageoriginale $K$ 是 $X$ 上的一个 $G$-挠子。
由 \hyperref[XVII.prop6.3.23]{6.3.23}，问题在 $X$ 上对平展拓扑是局部的，
所以只须处理 $r(K)$ 平凡的情形。
满足 $r(K)$ 平凡的 $G$-挠子 $K$，可等同于赋有同构 $k:H\isomto u(K')$
的平凡 $G'$-挠子 $K'$。令 $Z=\Sym^{n}_{T}(X)$。
平凡挠子之间的态射 $k$ 诱导 $k_{n}:H_{Z}\isomto u(K'_{n})$，
而 $\Sym^{n}\sisi{}{_T}(X)$ 上对应的 $G$-挠子表示所考虑的下降数据。}
\end{enonce*}

% END native input: chapters/exp-17-part5-zh.tex
% BEGIN native input: chapters/exp-17-part6-zh.tex
\section{附录}\etiquette[7]{XVII.sec7}\pageoriginale

\setcounter{subsection}{-1}
\subsection{预备知识}\etiquette[7.0]{XVII.subsec7.0}

本附录使用 \sisi{${\rm VI_{B}}$}{VBIS} 的结果，并沿用其中的记号。

\subsubsection{}\etiquette[7.0.1]{XVII.subsubsec7.0.1}
设 $\Sch_{!}$ 为以拟紧拟分离概形为对象、以有限型分离态射为态射的范畴。以 $\Ee$ 表示由概形上的平展层定义的、$\Sch_{!}$ 上的双纤维化范畴，其纤维范畴为拓扑斯的反范畴 \sisi{${\rm VI_{B}}$ (4.3.0)}{（\hyperrefext[Vbis][Vbis.subsubsec4.3.0]{4.3.0}）}。固定 $\Sch_{!}$ 的一个对象 $R$，以 $\Ee_{R}$ 表示由 $\Ee$ 导出的、$\Sch_{!R}$ 上的双纤维化范畴。最后，设 $\mathcal{A}$ 是拓扑斯 $\sheafGamma(\Ee^{\circ}_{R})$ 中的环，满足以下条件：
\begin{enumerate}
\renewcommand{\theenumi}{\roman{enumi}}
\renewcommand{\labelenumi}{(\theenumi)}
\item $\mathcal{A}$ 是 $\Ee^{\circ}_{R}$ 在 $\Sch^{\circ}_{!R}$ 上的\emph{余笛卡儿}（cocartesian）截面。

\item 对 $\Sch_{!R}$ 中的每个概形 $X$，$\mathcal{A}_{X}$ 是\emph{挠}层。
\end{enumerate}

在这一约定下，我们在范畴 $\sheaf{\Mod}(\sisi{}{\sheafGamma(}\Ee^{\circ}_{R}\sisi{}{)},\mathcal{A})$ 中工作（\refvb[VI B ]{Vbis.subsubsec1.3.5}）。

目标是对 $\Sch_{!R}$ 中的每个态射 $f:X\rightarrow Y$，定义函子 $Rf_{!}:D(X,\mathcal{A}_X)\rightarrow
D(Y,\mathcal{A}_{Y})$，使它在 $f$ 可紧化时与前面定义的函子一致，并使随 $X$ 变动的范畴 $D(X,\mathcal{A}_{X})$ 构成 $\Sch_{!R}$ 上的余纤维化范畴。\nde{\MakeLinkTarget*{sga4.exp.XVII.editor-note.7.0.1.compactification}\label{XVII.ndeapp}由 Nagata 紧化定理，$\Sch_{!}$ 中的每个态射都可紧化（见 \nderef{XVII.ndeConrad}）。本附录说明，在定义 $Rf_!$ 时，如何依靠 Chow 引理（\hyperref[XVII.lem7.0.6]{7.0.6}）来避免使用紧化定理。}

\subsubsection{}\etiquette[7.0.2]{XVII.subsubsec7.0.2}
本附录的全部结果虽然在 \hyperref[XVII.subsubsec7.0.1]{7.0.1} 指定的框架下陈述，但作显然的修改后，也可移到以下框架：取 $\sheafGamma(\Ee^{\circ}_{R})$ 中以常值层 $\ZZ$ 为值的常值截面，并且对于每个概形 $X$，限制到导出范畴 $D^{\ast}_{\tors}(X)$，它由上同调层为挠层的阿贝尔群层复形定义（\refvb[VI B ]{Vbis.subsubsec4.3.1}）。

\begin{enonce*}{引理 7.0.3}\etiquette[7.0.3]{XVII.lem7.0.3}
设\pageoriginale 给定概形的交换图
$$\vcenter{\etiquette[7.0.3.1]{XVII.eq7.0.3.1}
\xymatrix{
X' \ar[rr]^{j'} \ar[d]_{f'}& & X \ar[d]^{f}\\
Y' \ar[rr]_{j}& & Y.
}}\leqno{(7.0.3.1)}
$$
其中 $f$ 和 $f'$ 固有，$j$ 和 $j'$ 为浸入。那么，对于 $D^{+}_{\tors}(X')$ 中的每个复形 $K^{\bigdot}$，典范箭头\nde{经过一个简单的逐步分解，此箭头的定义与 \eqhyperref[XVII.eq5.1.5.2]{5.1.5.2} 中的定义类似。}
$$
j_{!}\circ R^{+}f'\sisi{}{_*}(K^{\bigdot})\longrightarrow R^{+}f\sisi{}{_*}\circ
j_{!}(K^{\bigdot})
$$
是同构。
\end{enonce*}

由“远端有界函子（way-out functor）引理”（\cite{sga4.exp.XVII.bib.RD} I \sisi{(}{}7.1），只需证明，对于每个整数 $i$ 和每个挠层 $F$，箭头
$$\etiquette[7.0.3.2]{XVII.eq7.0.3.2}
j_{!}\circ R^{i}f'\sisi{}{_*}(F)\longrightarrow R^{i}f\sisi{}{_*}\circ j_{!}(F)\leqno(7.0.3.2)
$$
是同构。

用 $j$ 和 $j'$ 的闭像代替它们，可归结为 $j$ 和 $j'$ 为像处处稠密的开浸入的情形。于是 $f$ 和 $f'$ 的固有性蕴含图 \eqhyperref[XVII.eq7.0.3.1]{7.0.3.1} 为笛卡儿图。在这些条件下，由固有基变换定理的 \hyperrefext[XII][XII.cor5.2]{5.2} 形式，对于 $Y$ 的每个几何点 $\xi$，箭头 $(j_{!}\circ R^{i}f'\sisi{}{_*}(F))_{\xi}\rightarrow (R^{i}f\sisi{}{_*}\circ
j_{!}(F))_{\xi}$ 是同构。

\begin{enonce*}{构造 7.0.4}\etiquette[7.0.4]{XVII.const7.0.4}
\emph{设 $D$ 为范畴，$X$ 和 $Y$ 为 $\Sch_{!R}$ 中的两个 $D$-对象，并且对 $D$ 中的每个箭头 $\alpha\rightarrow \beta$，态射 $X_{\beta}\rightarrow X_{\alpha}$ 和 $Y_{\beta}\rightarrow Y_{\alpha}$ 都固有。设 $j:X\rightarrow
Y$ 为态射，并且对 $D$ 的每个对象 $\alpha$，$j_{\alpha}:X_{\alpha}\rightarrow Y_{\alpha}$ 为浸入。对于 $i=0$ 的同构 \eqhyperref[XVII.eq7.0.3.2]{7.0.3.2} 满足通常的相容性条件，由此可以构造一个\emph{正合}函子：}
\[
j_{!}:\Mod(\sheafGamma(\bar{X}),\mathcal{A})\longrightarrow
\Mod(\sheafGamma(\bar{Y}),\mathcal{A})
\]
\emph{（记号见 \sisi{(}{}\refvb{Vbis.subsubsec1.2.5}）\sisi{。}{}它\pageoriginale 诱导一个仍记为 $j_{!}$ 的三角函子：}
\[
j_{!}:D(\sheafGamma(\bar{X}),\mathcal{A})\longrightarrow
D(\sheafGamma(\bar{Y},\mathcal{A}))
\]

\emph{于是，对于 $D$ 的每个对象 $\alpha$，有交换图，其中 $j_{\alpha !}$ 表示通常的“$0$ 延拓”函子：}
\[
\xymatrix{
D(\sheafGamma(\bar{X}), \sisi{}{\Aa}) \ar[rr]^{j_!} \ar[d]_{R e^{\sisi{+}{*}}_{\alpha}}&
&D(\sheafGamma(\bar{Y}),\Aa) \ar[d]^{R e^{\sisi{+}{*}}_{\alpha}}\\
D(X_{\alpha}, \Aa_{X_{\alpha}})  \ar[rr]_{j_{\alpha !}} & &
  D(Y_{\alpha},\Aa_{Y_\alpha}).
}
\]
\end{enonce*}

\begin{enonce*}{构造 7.0.5}\etiquette[7.0.5]{XVII.constru7.0.5}
设给定 $\Sch_{!R}$ 中的交换图
$$\vcenter{\etiquette[7.0.5.1]{XVII.eq7.0.5.1}
\xymatrix{
X'   \ar[rr]^{i'} \ar[d]_{f'}&    & X \ar[d]^{f}\\
Y'   \ar[rr]^{i} &    & Y
}}\leqno{(7.0.5.1)}
$$
其中 $i$ 和 $i'$ 为浸入，$f$ 是固有态射，$f'$ 是固有\emph{满射}态射。它定义 $\Sch_{!R}$ 中单纯对象的交换图
$$\vcenter{\etiquette[7.0.5.2]{XVII.eq7.0.5.2}
\xymatrix{
[X'|Y']      \ar[d]_{u_{f'}}  \ar[r]^{j} & [X|Y] \ar[d]^{u_f}\\
C^{\Delta}_{Y'}  \ar[r]_{i} & C^{\Delta}_{Y}
}}\leqno{(7.0.5.2)}
$$
，其中对每个 $n$，$j_{n}:[X'|Y']_{n}\rightarrow [X|Y]_{n}$ 为浸入。沿用此记号，有典范同构
\[
i_{!}\isomto  R^{+}\bar{u}_{f\sisi{+}{*}}\circ j_{!}\circ
L^{+}\bar{u}_{f'}\sisi{}{^*}{\quoi}.
\]
\end{enonce*}

用 $i$ 和 $i'$ 的闭像代替它们，可设 $i$ 和 $i'$ 为像处处稠密的开浸入。于是图 \eqhyperref[XVII.eq7.0.5.1]{7.0.5.1} 是笛卡儿图，且 $f$ 是\emph{满射}。此外，对于每个整数 $n$，$j_{n}$ 是开浸入\pageoriginale，而图
\[
\xymatrix{
[X'|Y']_n        \ar[r]^{j_n} \ar[d]& [X|Y]_n \ar[d]\\
Y'  \ar[r]_{i} & Y
}
\]
是笛卡儿图。

因此，$L^{+}\bar{u}_{f}^{\sisi{+}{*}}\circ i_{!}$ 同构于 $j_{!}\circ
L^{+}\bar{u}_{f'}^{\ast}$。利用 $u_{f}$ 是上同调下降增广这一事实（\sisi{${\rm VI_{B}}$ (\refnoexpose{Vbis.prop4.3.2})}{\refvb{Vbis.prop4.3.2}}），得到一串同构
\[
i_{!}\isomto R^{+}\bar{u}_{f\sisi{+}{*}}\circ
L^{+}(\bar{u}^{\sisi{+}{*}}_{f})\circ i_{!}\isomto
R^{+}\bar{u}_{f\ast}\circ j_{!}\circ L^{+}\bar{u}_{f'}^{\ast}
\]
，其复合即为所求同构。\nde{更直接地，不假设 $f'$ 满射，仍有一串典范箭头
\[i_!\to i_! R^+\bar{u}_{f'*}L^+\bar{u}^*_{f'}\isomto  R^+\bar{u}_{f*} j_! L^+u^*_{f'},\]
，其中第二个箭头是同构（\hyperref[XVII.lem7.0.3]{7.0.3}）。当 $f'$ 满射时，由上同调下降，第一个箭头也是同构，而复合与正文中构造的同构一致。
}

最后，要用到 Chow 引理的以下变式，其证明将出现于 \EGA II 第 2\sisi{版}{版}\nde{见 \nderef{XVII.ndeConrad} 中引用的 [C, 2.6]。}，它改进了（\sisi{\hyperref[XVII.subsec7.1]{7.1}}{\hyperrefext[XII][XII.lem7.1]{7.1}}）：

\begin{enonce*}{引理 7.0.6}\etiquette[7.0.6]{XVII.lem7.0.6}
设 $S$ 为拟紧拟分离概形，$f:X\rightarrow S$ 为有限型分离态射。那么存在一个满足以下条件的交换图
\[
\xymatrix{
X' \ar[rr]^{p} \ar[d]_{f'}& & X \ar[lld]^{f}\\
S                         & &
}
\]
：
\begin{enumerate}
\renewcommand{\theenumi}{\roman{enumi}}
\renewcommand{\labelenumi}{(\theenumi)}
\item $p$ 固有且\emph{满射}；

\item $X$ 有一个稠密开子集 $U$，使 $p$ 从 $p^{-1}(U)$ 到 $U$ 诱导同构；

\item $f'$ 拟射影。
\end{enumerate}
\end{enonce*}

\subsection{$R^{+}f_{!}$ 的定义}\etiquette[7.1]{XVII.subsec7.1}\pageoriginale

\begin{enonce*}{定义 7.1.1}\etiquette[7.1.1]{XVII.def7.1.1}
设 $f:X\rightarrow S$ 是概形的态射。$f$ 的一个\emph{伪紧化}（pseudo-compactification）是指一个交换图
$$\vcenter{
\xymatrix{
X_1 \ar[r]^{\sisi{}{j_1}} \ar[d]_{p_1}& T_1\ar[ldd]^{f_1}\\
X   \ar[d]_{f}& \\
S   &
}}\leqno{(7.1.1.0)}
$$\etiquette[7.1.1.0]{XVII.eq7.1.1.0}
，其中 $p_{1}$ 固有且满射，$j_{1}$ 为浸入，$f_{1}$ 固有。
\end{enonce*}

\setcounter{subsubsection}{1}
\subsubsection{}\etiquette[7.1.2]{XVII.subsubsec7.1.2}
设 $f:X\rightarrow S$ 是 $\Sch_{!R}$ 中的态射。由 Chow 引理（\hyperref[XVII.lem7.0.6]{7.0.6}），$f$ 有一个伪紧化 $(p_{1},j_{1},f_{1})$。图 \eqhyperref[XVII.eq7.1.1.0]{7.1.1.0} 定义态射
\[
[X_{1}|X]\xrightarrow{j_{1}}[T_{1}|S]
\]
，仍记为 $j_{1}$，使得对每个 $n$，$j_{1n}$ 为浸入。令 $R^{+(1)}f_{!}=\sisi{L^{+}\bar{u}_{p_{1}}\cdot
j_{1!}\circ R^{+}\bar{u}_{f_{1}}}{R^+\bar{u}_{f_1*}\circ j_{1!} \circ L^+\bar{u}_{p_1}^*}$，便得到从 $D^{+}(X,\mathcal{A}_{X})$ 到 $D^{+}(S,\mathcal{A}_{S})$ 的\emph{三角}函子。

\begin{enonce*}{构造 7.1.3}\etiquette[7.1.3]{XVII.constru7.1.3}
对于 $f$ 的每对伪紧化 $((p_{1},j_{1},f_{1}),(p_{2},j_{2},f_{2}))$，存在同构 $\alpha_{12}:R^{+(1)}f_{!}\rightarrow R^{+(2)}f_{!}$，使得对于任意另一个伪紧化 $(p_{3}, j_{3},f_{3})$，有 $\alpha_{23}\circ \alpha_{12}=\alpha_{13}$。
\end{enonce*}

\begin{enumerate}
\renewcommand{\theenumi}{\alph{enumi}}
\renewcommand{\labelenumi}{\theenumi)}
\item 令
\[
[[X_{1},X_{2}|X]]_{[n]\times [m]}=[X_{1}|X]_{n}\times
[X_{2}|X]_{m}
\]
以定义双单纯系统 $[[X_{1},X_{2}|X]]$。类似地，令
\[
[[X_{1}|X]]_{[n]\times[m]}=[X_{1}|X]_{n}{\quoi},
\]
，其中由态射 $[m']\rightarrow [m]$ 诱导的所有态射 $[[X_{1}|X]]_{[n]\times [m]}\rightarrow
[[X_{1}|X]]_{[n]\times [m']}$ 均为恒等态射。以相同方式定义双单纯系统 $[[T_{1},T_{2}|S]]$ 和 $[[T_{1}|S]]$，由此有交换图
\[
\xymatrix@C=1em{
  & & [[T_1,T_2|S]] \ar[rrrr]\ar'[dr][rrrddd]_(.2){u_{f_1f_2}}|(.38)\hole& &   & &  [[T_1|S]]\ar[lddd]^{u_{f_1}}\\
  [[X_1,X_2|X]]
  \ar[rrrr]\ar[rru]^{j_{12}}\ar[rrrddd]_{u_{p_1 p_2}} & &  & & [[X_1|X]]\ar[rru]_{j_1}\ar[lddd]_(.75){u_{p_1}}  \\
  & & & & &   &\\
  & & & & & C^{\Delta\times\Delta}_S\\
  &&& C^{\Delta\times\Delta}_X \ar[rru]
}
\]
，其中\pageoriginale $j_{12}$ 是逐次的浸入。\sisi{VI b (\refnoexpose{Vbis.subsec2.3})}{\hyperrefext[Vbis][Vbis.subsec2.3]{2.3}} 中描述的计算过程使我们可以构造态射
\[
R^{+(1)} f_{!}\xrightarrow{\alpha_{1}}R^{+}\bar{u_{f_{1}f_{2}}}\sisi{}{_*}\circ
j_{12!} \circ L^{+}\bar{u_{p_{1}p_{2}}}^{\sisi{+}{\ast}}=R^{{+}(1,2)}f_{!}
\]
，它由 $S$ 上层的三重复形的一个态射得到。

\item $\alpha_{1}$ \emph{是同构}：

为证明这一点，可在前述三重复形的态射中，固定来自 $\bigl[[T_{1},T_{2}]\bigr]_{[n]\times[m]}$ 的第一个指标 $n$，并考虑所得双复形的态射。可验证，这个态射恰为引理 \hyperref[XVII.constru7.0.5]{7.0.5} 对图
\[
\xymatrix{
**[r][X_1|X]_n \times X_2  \ar[r] \ar[d] & **[r][T_1|S]_n \times
T_2 \ar[d]\\
**[r][X_1|X]_n             \ar[r] & **[r][T_1|S]_n
}
\]
定义的同构。

\item 以相同方式构造同构
\[
R^{+(2)}f_{!}\xrightarrow{\alpha_{2}}R^{+}\bar{u_{f_{1}f_{2}}}_{\sisi{+}{*}}\circ
j_{12!} \circ L^{+}\bar{u_{p_{1}p_{2}}}\sisi{}{^*}
\]
，并令 $\alpha_{12}=\alpha^{-1}_{2}\circ \alpha_{1}$。
\end{enumerate}

为验证关系 $\alpha_{23}\circ \alpha_{12}=\alpha_{13}$，使用一个三重单纯对象。所得图的交换性留给读者验证：\pageoriginale
$$
\xymatrix{
  R^{+(2,3)}{f_!} \ar[rd]& R^{+(2)}f_! \ar[r] \ar[l]& R^{+(1,2)}f_! \ar[ld]\\
                &         R^{+(1,2,3)}f_!  &                  \\
R^{+(\sisi{2}{3})}f_! \ar[uu]\sisi{}{\ar[rd]}     &                         &  R^{+(1)}f_! \ar[uu]\sisi{}{\ar[ld]}\\
                &    R^{+(1,3)}f_! \ar[uu]        &
}
$$

\begin{enonce*}{构造 7.1.4}\etiquette[7.1.4]{XVII.constru7.1.4}
设 $Z\xrightarrow{g}Y$ 和 $Y\xrightarrow{f}X$ 是 $\Sch_{!R}$ 中的两个态射。对于每个三元组 $((p_{1},j_{1},f_{1}),
(p_{2},j_{2}g_{2}), (p_{3},j_{3},(fg)_{3}))$，其中 $(p_{1},j_{1},f_{1})$ 是 $f$ 的伪紧化，$(p_{2},j_{2},g_{2})$ 是 $g$ 的伪紧化，$(p_{3},j_{3},(fg)_{3})$ 是 $fg$ 的伪紧化，存在一个与前面定义的同构 $\alpha_{ij}$ 相容的同构 $\sisi{R^{+(3)}fg_{!}\isomfrom  R^{+(2)}g_{!}\circ
R^{+(1)}f_{!}}{R^{+(1)}f_!\circ R^{+(2)}g_!\isomto R^{+(3)}(fg)_!}$。
\end{enonce*}

分几个步骤进行。

\setcounter{subsubsection}{4}
\paragraph{}\etiquette[7.1.4.1]{XVII.par7.1.4.1}
设 $(p_{1},j_{1},f_{1})$（\resp $(p_{2},j_{2},f_{2})$）是 $f:Y\rightarrow X$（\resp $g:Z\rightarrow Y$）的一个伪紧化：
$$\vcenter{\etiquette[7.1.4.1.1]{XVII.eq7.1.4.1.1}
\xymatrix{
  &     & Z\ar[d]^{g}  & Z_2\ar[l]_{p_2}\ar@{^{(}->}[r]^{j_2}  & T_2 \ar[lld]^{g_2}\\
T_1  \ar[rrd]_{f_1}& Y_1 \ar[r]^{p_1} \ar@{_{(}->}[l]_{j_1}& Y \ar[d]^{f} &      &    \\
    &     & X        &      &
}}\leqno{(7.1.4.1.1)}
$$

由 $g$ 经基变换得到的态射 $Z'=\fprod{Z}{Y_{1}}{Y}\rightarrow Y_{1}$ 赋有伪紧化：
$$\vcenter{\etiquette[7.1.4.1.2]{XVII.eq7.1.4.1.2}
\xymatrix{
  T'_2 = \fprod{T_2}{Y_1}{Y}   \ar[rdd]_{g'_2}& \fprod{Z_2}{Y_1}{Y} =Z'_2 \ar[d]^{p'_2}
\ar@{_{(}->}[l]_{j'_2}\\
&  \fprod{Z}{Y_1}{Y} = Z' \ar[d]^{g'}\\
&    Y_1
}}\leqno{(7.1.4.1.2)}
$$

现在\pageoriginale 选择一个\sisi{}{伪}紧化，其对象为
\sisi{$j_{1}\cdot g'_{2}$。
$$\etiquette[7.1.4.1.3]{XVII.eq7.1.4.1.3}
\vcenter{
\xymatrix{
T_{12}        \ar[d]_{g''_2} & \fprod{T_2}{Y_1}{Y}=T'_2
\ar@{_{(}->}[l]_(.6){j'_1} \ar[d]^{g'_2}\\
T_1            &   Y_1 \ar@{_{(}->}[l]
}}\leqno(7.1.4.1.3)
$$%
}{$j_1\circ g'_2$\nde{原文选择了一个紧化，借助紧化定理这是可以做到的。由于附录的目的是避免这一定理（\cf \nderef{XVII.ndeapp}），编辑冒昧地将其换为伪紧化，并对 \hyperref[XVII.par7.1.4.2]{7.1.4.2} 作了略微修改。另一种解决方法是只考虑使 $f_1$ 为\emph{射影}态射的伪紧化 $(p_1,j_1,f_1)$。}
$$\etiquette[7.1.4.1.3]{XVII.eq7.1.4.1.3}
\vcenter{
\xymatrix{
T_{12}        \ar[dd]_{g''_2} &T_2''\ar[d]\ar[d]^p\ar@{_{(}->}[l]_{j'_1}\\
& \fprod{T_2}{Y_1}{Y}=T'_2\ar[d]^{g'_2}\\
T_1            &   Y_1 \ar@{_{(}->}[l]_{j_1}
}}\leqno(7.1.4.1.3)
$$}%
由此定义复合 $fg$ 的一个伪紧化；相应的函子记为 $R^{+(12)}\sisi{\int g}{(fg)}_{!}$。

\paragraph{}\etiquette[7.1.4.2]{XVII.par7.1.4.2}
\emph{前面的资料使我们可以定义同构} $R^{+(12)}\sisi{\int g}{(fg)}_{!}\isomfrom  R^{+(1)}f_{!}\sisi{\cdot}{\circ} R^{+(2)}g_{!}$：

存在一个交换图，其中 $Z_2''=\fprod{T''_2}{Z'_2}{T'_2}$，而所有箭头都自然地由 \eqhyperref[XVII.eq7.1.4.1.1]{7.1.4.1.1}、\eqhyperref[XVII.eq7.1.4.1.2]{7.1.4.1.2} 和 \eqhyperref[XVII.eq7.1.4.1.3]{7.1.4.1.3} 导出。
$$\etiquette[7.1.4.2.1]{XVII.eq7.1.4.2.1}
\vcenter{
\xymatrix{
[T_{12}|X] \ar[rd]_{g''_2} &   &[T\sisi{'}{''}_2|Y] \ar[rd]_{g'_2\sisi{}{p}} \ar[ll]_{j'_1}&
[Z\sisi{'}{''}_{2}|Z] \ar[r]^{r}\ar[l]_{j\sisi{'}{''}_2} \ar[d]_{l}&
C^{\Delta}_Z \ar[d]^{g}\\
            &    [T_1|X] \ar[rrrd]_{u_{f_1}} &  & [Y_1|Y]\ar[r]_{u_{p_1}}\ar[ll]_{j_1} & C_Y^{\Delta} \ar[d]^{f}\\
            &       &  &   &  C_X^{\Delta}
}}\leqno(7.1.4.2.1)
$$

我们\pageoriginale 将定义一个自然箭头：
\[
\beta:R^{+(1)}f_{!}\circ R^{+(2)}g_{!}\longrightarrow R^{+(12)}\sisi{fg}{(fg)}_{!}
\]
，并将证明它是同构。

为此使用两个引理：

\begin{enonce*}{引理 7.1.4.2.2}\etiquette[7.1.4.2.2]{XVII.lem7.1.4.2.2}
设给定交换图：
\[
\xymatrix{
X_1    \ar[d]_{p_1}     \ar[r]^{j_1} & T_1 \ar[d]^{h'}\\
X    \ar[d]_{f}       \ar[r]_{j}           & T \ar[ld]^{h}\\
    S       &
}
\]
其中 $h$ 和 $h'$ 固有，$p_{1}$ 固有且满射，$j$ 和 $j_{1}$ 为浸入。那么典范箭头
\[
R^{+}h\sisi{}{_*}\circ j_{!}\longrightarrow R^{+(1)}\sisi{(hh')}{f}_{!}
\]
是同构。
\end{enonce*}

事实上，由（\hyperref[XVII.constru7.0.5]{7.0.5}），箭头
\[
j_{!}\longrightarrow R^{+}\bar{u}_{h'\sisi{}{*}}\circ j_{1!}\circ
L^{+}\bar{u}_{p_{1}}\sisi{}{^*}
\]
是同构。应用 $R^{+}h\sisi{}{_*}$ 即得结论。

\begin{enonce*}{引理 7.1.4.2.3}\etiquette[7.1.4.2.3]{XVII.lem7.1.4.2.3}
沿用图\emph{\eqhyperref[XVII.eq7.1.4.2.1]{7.1.4.2.1}} 的记号，典范箭头
\begin{gather*}
R^{+(1)}f_{!}\circ R^{+}\bar{u}_{p_{1}\sisi{}{*}}\longrightarrow
R^{+}\bar{u}_{f_{1}\sisi{}{*}}\circ j_{1!}\\
R^{+(2)}g_{!}\circ R^{+}\bar{r}\sisi{}{_*}\longrightarrow
R^{+}\bar{u}_{p_{1}\sisi{}{*}}\circ R^{+}\sisi{g'_{2}}{(g_2'p)_*}\circ j\sisi{'}{''}_{2!}
\end{gather*}
都是同构。
\end{enonce*}

只证明第一个同构，第二个的推导\pageoriginale 留给读者。

采用显然的记号\sisi{，令 $C=\fprod{[Y_{1}|Y]}{X_{1}}{X}$ 和 $\bar{C}=\fprod{[T_{1}|S]}{T_{1}}{S}$，于是有交换图：
\[
\xymatrix{
                                & \bar{C} \ar[rdd]|\hole \ar[rr]  & & T_1 \ar[ldd]^{f_1}\\
C                 \ar[rr]  \ar[ru]^{i}  \ar[d]^{q}    &    &  Y_1  \ar[ru]_{j_1} \ar[d]^{p_1}&      \\
[Y_1|Y] \ar[rr]_{u_{p_1}}&  &      Y         &
}
\]
其中 $i$ 为浸入，$q$ 固有且满射。接着构造图
\[
\xymatrix{
\bigl[\bigl[C \underset{q}{\bigm|}[Y_1|Y]\bigr]\bigr] \ar[r]^{i} \ar[d]^{u_q}&
[[\bar{C}|C^{\Delta}_Y]] \ar[ldd]^{k}\\
|X_1 | X |                             \ar[d]         &\\
Y                       &                        &
}
\]
，由（\hyperref[XVII.lem7.1.4.2.2]{7.1.4.2.2}），函子 $R\bar{u}j_{1}\circ
j_{1}$ 等同于函子 $R^{+}\bar{k}\circ i_{!}\circ
L^{+}\bar{u}_{q}$；再由（\hyperref[XVII.lem7.0.3]{7.0.3}）及基变换定理，后一个函子等同于 $R^{+(1)}f_{!}\circ
R^{+}\bar{u}_{p_{1}}$。}{并沿用 \hyperref[XVII.constru7.1.3]{7.1.3} a) 的记号，存在交换图
\[\xymatrix{&[Y_1|Y]\ar[rr]^{j_1}\ar[dd]^(.7){\bar{u}_{p_1}}|\hole &&[T_1|X]\ar[dd]^{\bar{u}_{f_1}}\\
[[Y_1,Y_1|Y]]\ar[ur]^{k}\ar[dd]_{q}\ar[rr]^(.7)i && [[T_1,T_1|X]]\ar[dd]\ar[ur]\\
&Y\ar[rr]^(.3){f}|\hole && X\\
[Y_1|Y]\ar[ur]^{\bar{u}_{p_1}}\ar[rr]^{j_1} &&[T_1|X]\ar[ur]_{\bar{u}_{f_1}} }\]
其中 $q$ 和 $k$ 是从 $[[Y_1,Y_1|Y]]$ 到 $[Y_1|Y]$ 的两个投影。由上同调下降，$R^+\bar{u}_{f_1 *} \circ j_{1!}$ 等同于 $R^+\bar{u}_{f_1 *} \circ j_{1!} \circ R^+ q_* L^+ q^*$；对上述立方体的前面和顶面应用 \hyperref[XVII.lem7.0.3]{7.0.3}，后者等同于 $R^+\bar{u}_{f_1*}\circ j_{1!}\circ R^+k_* \circ L^+q^*$。由基变换定理，这最后一个函子等同于 $R^+\bar{u}_{f_1 *} \circ j_{1!} \circ L^+\bar{u}_{p_1*} \circ R^+\bar{u}_{p_1 *}$，按定义即为 $R^{+(1)}f_{!}\circ
R^{+}\bar{u}_{p_{1}*}$。\nde{原文中的证明略有不同，%and not very clear
并使用 \hyperref[XVII.lem7.1.4.2.2]{7.1.4.2.2}。}

}%
引理（\hyperref[XVII.lem7.1.4.2.3]{7.1.4.2.3}）已经证明，因此有一串同构
\begin{align*}
R^{+(\sisi{3}{12})} \sisi{fg}{(fg)}_{!} &= R^{+}\bar{u}_{f_{1}\sisi{}{*}}\circ R^{+}g''_{2}\circ
j'_{1!}\circ j\sisi{'}{''}_{2!}\circ L^{+}\bar{r}^{\sisi{+}{*}}\\
&\isomfrom  R^{+}\bar{u}_{f_{1}\sisi{+}{*}}\circ j_{1!}\circ R^{+}
\sisi{g'_{2}}{(g'_2p)_*}\circ j'_{2!}\circ L^{+}\bar{r}^{\sisi{+}{*}}\\
&\isomfrom  R^{+(1)}f_{!}\circ R^{+}\bar{u}_{p_{1}\sisi{+}{*}}\circ
R^{+}\sisi{g'_{2+}}{(g'_2p)_*}\circ j'_{2!}\circ L^{+}\bar{r}\sisi{}{^*}\\
&\isomfrom  R^{+(1)}f_{!}\circ R^{+(2)}g_{!}\circ
R^{+}\bar{r}\sisi{}{_*}\circ L^{+}\bar{r}^{\sisi{+}{*}}\\
&\isomfrom  R^{+(1)}f_{!}\circ R^{+(2)}g_{!}
\end{align*}
，其\pageoriginale 复合给出 $\beta$。

\paragraph{}\etiquette[7.1.4.3]{XVII.par7.1.4.3}
若 $(p_{3},j_{3},(fg)_{3})$ 是 $fg$ 的伪紧化，则将 $\alpha_{(12)3}$ 与 $\beta$ 复合，便定义同构
\[
R^{+(1)}f_{!}\circ
R^{+(2)}g_{!}\xrightarrow{\gamma_{123}}R^{+(3)}fg_{!}
\]
。

\paragraph{}\etiquette[7.1.4.4]{XVII.par7.1.4.4}
现在设 $(p_{1'},j_{1'},f_{1'})$ 和 $(p_{2'},j_{2'},g_{2'})$ 是 $f$ 和 $g$ 的另外两个伪紧化，并配有图 \sisi{(\hyperref[XVII.par7.1.4.3]{7.1.4.3})}{\eqhyperref[XVII.eq7.1.4.1.3]{7.1.4.1.3}}；考虑一个 \sisi{(\hyperref[XVII.par7.1.4.4]{7.1.4.4})}{\eqhyperref[XVII.eq7.1.4.2.1]{7.1.4.2.1}} 型的图，将其中的单纯对象换为双单纯对象，就可以验证下图交换：
\[
\xymatrix{
R^{+(12)}\sisi{fg}{(fg)}_1 \ar[d]_{\alpha_{(12)(1'2')}} & R^{+(1)}f_1 \circ
R^{+(2)}g_! \ar[d]^{\alpha_{11'}\ast \alpha_{22'}} \ar[l]_(.6){\beta}\\
R^{+(1'2')}\sisi{fg}{(fg)}_!  & R^{+(1')}f_! \circ R^{+(2')}g_! \ar[l]^(.6){\beta'}
}
\]

由此，对于 $fg$ 的每个伪紧化 $(p_{3'},j_{3'},(fg)_{3'})$，下图交换：
\[
\xymatrix{
R^{+(3)}\sisi{fg}{(fg)}_! \ar[d]_{\alpha_{33'}}& R^{+(1)}f_! \circ R^{+(2)}g_!
\ar[l]_(.6){\gamma_{123}} \ar[d]^{\alpha_{11'}\ast
  \alpha_{22'}}\\
R^{+(3')}\sisi{fg}{(fg)}_!  & R^{+(1')}f_! \circ R^{+(2')}g_! \ar[l]_(.6){\gamma_{1'2'3'}}
}
\]
这就完成了 \hyperref[XVII.constru7.1.4]{7.1.4} 的证明。

\subsubsection{}\etiquette[7.1.5]{XVII.subsubsec7.1.5}
对 $\Sch_{!R}$ 中的每个态射 $f$ 选择一个伪紧化，并使恒等态射对应平凡的伪紧化，由此对每个态射 $f:X\rightarrow Y$ 定义三角函子
\[
R^{+}f_{!}:D^{+}(X,\mathcal{A}_{X})\longrightarrow
D^{+}(Y,\mathcal{A}_{Y})
\]
，并且\pageoriginale 对每对态射 $(f,g)$ 定义同构
\[
R^{+}fg_{!}\xrightarrow{\substack{c_{fg}\\\sim}}
R^{+}f_{!}\circ R^{+}g_{!}
\]
此外，若 $f$ 为恒等态射，则 $R^{+}f_{!}$ 为恒等函子，$c_{fg}$ 为恒等态射。

\begin{enonce*}{命题 7.1.6}\etiquette[7.1.6]{XVII.prop7.1.6}
对应 $f\rightarrow R^{+}f_{!}$ 定义从 $\Sch_{!R}$ 到范畴的 $2$-范畴的正规化伪函子。换言之，令 $X$ 遍历 $\Sch_{!R}$ 的对象，范畴 $D^{+}(X,\mathcal{A}_{X})$ 可以组织为 $\Sch_{!R}$ 上的余纤维化范畴：对于 $\Sch_{!R}$ 中的每个态射 $f:X\rightarrow Y$，令
\[
\Hom_{f}(K,L)=\Hom(R^{+}f_{!}K,L)
\]
，其中 $K\in D^{+}(X,\mathcal{A}_{X})$，$L\in D^{+}(Y,\mathcal{A}_{Y})$。
\end{enonce*}

还须说明等同 $R^{+}fg_{!}\xrightarrow{\substack{c_{fg}\\\sim}}
R^{+}f_{!}\circ R^{+}g_{!}$ 与三重复合相容。设 $h:T\rightarrow Z$、$g:Z\rightarrow Y$ 和 $f:Y\rightarrow X$ 为 $\Sch_{!R}$ 中的三个态射；由前面的论证，为验证这一相容性，只需在任意选定 $h$、$g$ 和 $f$ 的伪紧化后验证。

有图：
$$\vcenter{\etiquette[7.1.6.1]{XVII.eq7.1.6.1}
  \xymatrix{
 T_5 \ar[rd]& T_4 \ar[rd] \ar@{_{(}->}[l]& T_3 \ar[rd]\ar@{_{(}->}[l]
 & T'\ar@{_{(}->}[l] \ar[r]^{p_3} \ar[d]& T \ar[d]^{h} \\
        & T_3 \ar[rd]  & T_2 \ar@{_{(}->}[l] \ar[rd]  & Z'\ar[r]^{p_2}\ar@{_{(}->}[l] \ar[d]& Z
  \ar[d]^{g}\\
        &    & T_1 \ar[rrd]& Y_1\ar@{_{(}->}[l] \ar[r]^{p_1}& Y \ar[d]^{f}\\
        &    &     &     & X
}}\leqno{(7.1.6.1)}
$$
其中所有斜箭头都固有，而 $p_{1}$、$p_{2}$ 和 $p_{3}$ 固有且满射。

由 \eqhyperref[XVII.eq7.1.6.1]{7.1.6.1} 构造一个（\hyperref[XVII.par7.1.4.4]{7.1.4.4}）型的图，即得结论。

\begin{enonce*}{命题 7.1.7}\etiquette[7.1.7]{XVII.prop7.1.7}
若\pageoriginale $f:X\rightarrow S$ 为可紧化态射，则函子 $Rf_{!}$ 与第\emph{XVII}讲中定义的函子一致。
\end{enonce*}

事实上，存在一个交换图：
\[
\xymatrix{
X_1         \ar[r]^{j_1} \ar[d]_{p_1}& T_1 \ar[d]^{h'}\\
X           \ar[r]_{j}  \ar[d]_{f} & T \ar[ld]^{h}\\
S
}
\]
，对它应用（\hyperref[XVII.lem7.1.4.2.2]{7.1.4.2.2}）即可。

\begin{enonce*}{定义 7.1.8}\etiquette[7.1.8]{XVII.def7.1.8}
若 $f:X\rightarrow Y$ 是 $\Sch_{!}$ 中的态射，则刚定义的函子 $R^{+}f_{!}$ 称为\emph{具有固有支集的正像函子}（导出范畴意义下）。函子 $H^{q}\circ R^{+}f_{!}$，记为 $R^{q}f_{!}$，称为第 $q$ 阶\emph{具有固有支集的正像函子}。
\end{enonce*}

沿用此记号，下面的构造留给读者验证。

\begin{enonce*}{构造 7.1.9}\etiquette[7.1.9]{XVII.constru7.1.9}
设 $(\pi_{1}, j_{1}, f_{1})$ 为 $f:X\rightarrow Y$ 的伪紧化，以 $\pi^{n}_{1}$ 表示典范态射 $[X_{1}|X]_{n}\rightarrow X$。那么对于 $D^{+}(X,\mathcal{A}_{X})$ 中的每个复形 $K^{\bigdot}$，有两个双正则谱序列：
\begin{align*}
E^{pq}_{1} &= R^{q}f_{!}(K^{p})\Longrightarrow R^{p+q}
f_{!}(K^{\bigdot})\tag{7.1.9.1}\\
E^{pq}_{1} &= R^{q}(f\circ
\pi^{p}_{1})_{!}(\pi^{p\sisi{+}{*}}_{1}(K^{\bigdot}))\Longrightarrow
R^{p+q}f_{!}(K^{\bigdot})\tag{7.1.9.2}
\end{align*}\etiquette[7.1.9.1]{XVII.eq7.1.9.1}
\etiquette[7.1.9.2]{XVII.eq7.1.9.2}%The vertical space following this paragraph is too large.
\end{enonce*}

（例如，取 $j_{1!}\circ
L^{+}\bar{u}_{p_{1}}^{\sisi{+}{*}}(K^{\bigdot})$ 的典范松弛消解（\cf XVII （\hyperref[XVII.subsec4.2]{4.2}）），再取其在 $u_{f_{1}}$ 下的正像。）

\subsection[$R^{+}f_{!}$ 的若干性质]{\mathversion{bold}$R^{+}f_{!}$ 的若干性质}\etiquette[7.2]{XVII.subsec7.2}
\pageoriginale

\begin{enonce*}{命题 7.2.1}\etiquette[7.2.1]{XVII.prop7.2.1}
设 $f:X\rightarrow Y$ 是 $\Sch_{!R}$ 中相对维数 $\leq d$ 的态射。那么函子 $R^{+}f_{!}$ 的上同调维数 $\leq 2d$。
\end{enonce*}

只需证明，对于 $X$ 上的每个 $\mathcal{A}_{X}$-模层 $F$，当 $i>2d$ 时，有 $R^{i}f_{!}(F)=0$。（于是由标准论证可得：对于 $D^{+}(X,\mathcal{A}_{X})$ 中任意复形 $K^{\bigdot}$，若当 $i>k$ 时 $H^{i}(K^{\bigdot})=0$，则当 $i>k+2d$ 时 $R^{i}f_{!}(K^{\bigdot})=0$。）

设 $\mathcal{U}=(U_{i})_{i\in I}$ 是 $X$ 的一个由仿射开子集组成的有限覆盖。令
\[
K^{n}_{\mathcal{U}}=\bigoplus_{|P|=\sisi{}{-}n+1} j_{P!}j^{*}_{P}(F)
\]
即定义（\hyperref[XVII.prop6.2.9]{6.2.9}）$F$ 的一个左消解 $K^{\bigdot}_{\mathcal{U}}\rightarrow F$，其中对于 $I$ 的每个子集 $P$，$j_{P}$ 为 $\cap_{i\in P} U_{i}$ 到 $X$ 的包含态射。

由 \eqhyperref[XVII.eq7.1.9.1]{7.1.9.1} 得到谱序列
\[
E^{pq}_{1}=\bigoplus_{\substack{|P|=\sisi{}{-}p+1\\ P\sisi{\in}{\subset} I\sisi{}{\text{ 非空}}}} R^{q}(fj_{P})_{!}
j^{\ast}_{P}(F)\Longrightarrow R^{p+q}f_{!}(F)
\]
，而 $fj_{P}$ 可紧化这一事实（第 XVII 讲已处理的情形）表明，当 $p>0$ 或 $q>2d$ 时，$E^{pq}_{1}=0$；由此，当 $n>2d$ 时 $R^{n}f_{!}(F)=0$，这正是所要证明的。

\begin{enonce*}{命题 7.2.2}\etiquette[7.2.2]{XVII.prop7.2.2}
设 $f:X\rightarrow Y$ 为 $\Sch_{!R}$ 中的态射，$K^{\bigdot}$ 为 $D^{b}(X,\mathcal{A}_{X})$ 中有限 \tTor-维数的对象。那么 $R^{+}f_{!}(K^{\bigdot})$ 具有有限 \tTor-维数。
\end{enonce*}

如上，设 $\mathcal{U}=(U_{i})_{i\in I}$ 是 $X$ 的一个由仿射开子集组成的有限覆盖。由于有限 \tTor-维数可以逐点验证（\hyperref[XVII.def4.1.9]{4.1.9}），$K^{n}_{\mathcal{U}}$ 的函子性使得可以用有限 \tTor-维数的复形构造 $K^{\bigdot}$ 的有界左消解。以 $K^{\bigdot\bigdot}_{\mathcal{U}}$ 表示所得双复形，其中\emph{第二个}指标是与 $K^{\bigdot}$ 对应的指标。按照（\sisi{VI b.(\refnoexpose{Vbis.subsubsec2.3.2})}{\hyperrefext[Vbis][Vbis.subsubsec2.3.2]{2.3.2}}）的约定，对每个整数 $n$，有以下正合列，其中运算 $\sigma_{\leq n}$ 在 $C^{+}(C^{+}(\Mod(X,\mathcal{A}_{X}))$ 中进行：
\[
\sisi{0\longrightarrow (\sigma_{\leq
n}(K^{\bigdot\bigdot}_{\mathcal{U}}))_{s}\longrightarrow (\sigma_{\leq
n+1}(K^{\bigdot\bigdot}_{\mathcal{U}}))_{s}\longrightarrow (-1)^{n+1}
K^{n+1,\bigdot}_{\mathcal{U}}\longrightarrow 0}{0\to (-1)^{n+1}K^{n+1,\bigdot}_{\mathcal{U}}\to(\sigma_{\leq
n+1}(K^{\bigdot\bigdot}_{\mathcal{U}}))_{s}\to(\sigma_{\leq
n}(K^{\bigdot\bigdot}_{\mathcal{U}}))_{s}\to 0}
\]
因而\pageoriginale 在 $D^{+}(X,\mathcal{A}_{X})$ 中有特殊三角：
\[
\xymatrix@C=1em{
  & K^{n+1,\bigdot}_{\Uu} \sisi{\ar[ld]}{\ar[rd]}&    \\
  (\sigma_{\leq n}(K^{\bigdot\bigdot}_{\Uu}))_s \sisi{\ar[rr]}{\ar[ur]}&  &(\sigma_{\leq
  n+1}(K^{\bigdot\bigdot}_{\Uu}))_s \sisi{\ar[lu]}{\ar[ll]}\\
}
\]

利用对每个 $n$，$R^{+}f_{!}(K^{n+1,\bigdot}_{\mathcal{U}})$ 具有有限 \tTor-维数这一事实（\hyperref[XVII.thm5.2.10]{5.2.10}），便得到 $R^{+}f_{!}((\sigma_{\leq
0}(K^{\bigdot\bigdot}_{\mathcal{U}}))_{s})=R^{+}f_{!}(K^{\bigdot})$ 具有有限 \tTor-维数。

\begin{enonce*}[remark]{注 7.2.3}\etiquette[7.2.3]{XVII.rem7.2.3}
前面的论证并不能表明，只要 $K^{\bigdot}$ 的 \tTor-维数 $\leq k$，$R^{+}f_{!}(K^{\bigdot})$ 的 \tTor-维数就 $\leq k$（\cf \sisi{(3.7)}{\hyperref[XVII.subsubsec7.3.7]{7.3.7}}）。
\end{enonce*}

\begin{enonce*}{命题 7.2.4}\etiquette[7.2.4]{XVII.prop7.2.4}
设给定 $\Sch_{!R}$ 中的笛卡儿图：
\[
\xymatrix{
X'   \ar[r]^{g'}\ar[d]_{f'}& X \ar[d]^{f}\\
S'  \ar[r]_{g} & S
}
\]
那么对于 $D^{+}(X,\mathcal{A}_{X})$ 中的每个复形 $K^{\bigdot}$，典范态射
\[
g^{\ast}\circ
R^{+}f_{!}(K^{\bigdot})\xrightarrow{\mathrm{ch}_{K^{\bigdot}}}R^{+}f_{!}\sisi{}{'}\circ
{g'}^{\ast}(K^{\bigdot})
\]
是同构。
\end{enonce*}

设 $(\pi_{1}, j_{1}, f_{1})$ 为 $f$ 的伪紧化；经基变换，它给出 $f'$ 的伪紧化 $(\pi'_{1}, j'_{1}, f'_{1})$ 以及交换图：
\[
\xymatrix@R=2ex{
                 && T'_1 \ar'[ldddd]^(.8){f'_1}|(.55)\hole [lldddddddd]|(.25)\hole \ar[rrrr] &&&& T_1 \ar[lldddddddd]^{f_1}\\
                 &\\
X'_1 \ar[rruu]^{j'_1}\ar[ddd]_{\pi'_1} \ar[rrrr] &&&& X_1 \ar[rruu]_{j_1} \ar[ddd]^{\pi_1}&\\
                     & & &\\
                     & *{} & &\\
X'   \ar[rrrr]_{g'} \ar[ddd]_{f'} &&&& X  \ar[ddd]_{f}  &\\
                     & & &\\
                     & & &\\
S'   \ar[rrrr]_{g} &&&&  S         &
}
\]
，由此\pageoriginale 可以定义箭头 $\mathrm{ch}_{K}$。

由 \eqhyperref[XVII.eq7.1.9.2]{7.1.9.2}，得到双正则谱序列的一个态射：
\[
\xymatrix{
g^{\ast}\circ R^{q} (f \circ \pi^{p}_1)_! (\pi^{p\sisi{+}{\ast}}_1
(K^{\bigdot})) \ar[d]^{u^{pq}_1} \ar[r]&  g^{\ast} \circ R^{p+q}f_{!}(K^{\bigdot}) \ar[d]^{\mathrm{ch}_K}\\
R^{q}(f'\circ \pi'^{p}_1)_! (\pi'^{p\ast}_1 (g'^{\ast}(K^{\bigdot})))  \ar[r] &
  R^{+}f'_! \circ g'^{\ast}(K\sisi{}{^{\bigdot}})
}
\]
而 $\mathrm{ch}_{K^{\bigdot}}$ 为同构，因为由（\hyperref[XVII.thm5.2.6]{5.2.6}），所有 $u^{pq}_{1}$ 都是同构。

\subsection{$Rf_{!}$ 的定义}\etiquette[7.3]{XVII.subsec7.3}

现在拟对 $\Sch_{!R}$ 中的每个态射 $f$，把函子 $R^{+}f_{!}$“以合理的方式”延伸为函子 $Rf_{!}:D(X,\mathcal{A}_{X})\rightarrow D(Y,\mathcal{A}_{Y})$。

\subsubsection{}\etiquette[7.3.1]{XVII.subsubsec7.3.1}
设 $(p_{1}, j_{1}, f_{1})$ 为 $f$ 的伪紧化，$F$ 为 $\mathcal{A}_{X}$-模。取 $j_{1!}\circ L\bar{u}_{p_{1}}^{\ast}(F)$ 的典范松弛消解（\cf \hyperref[XVII.subsec4.2]{4.2}），再取其在 $f_{1}$ 下的正像，最后取所得双复形的伴随单复形，就得到 $\mathcal{A}_{Y}$-模的复形：
$$\etiquette[7.3.1.1]{XVII.eq7.3.1.1}
0\longrightarrow
\Xx^{0}(F)\xrightarrow{S^{0}(F)}\Xx^{1}(F)\ldots
\Xx^{n}(F)\xrightarrow{S^{n}(F)}
\Xx^{n+1}(F)\ldots.\leqno(7.3.1.1)
$$
，它关于 $F$ 是函子性的\pageoriginale，并在 $D^{+}(Y,\mathcal{A}_{Y})$ 中同构于 $R^{+}f_{!}(F)$。

由 \hyperref[XVII.prop7.2.1]{7.2.1}，当 $i>2d$ 时 $H^{i}(\Xx^{\bigdot}(F))=0$，其中 $d$ 为使 $f$ 的相对维数 $\leq d$ 的整数；因此，复形
$$\etiquette[7.3.1.2]{XVII.eq7.3.1.2}
0\to \Xx^{0}(F)\to
\Xx^{1}(F)\to \ldots \sisi{}{\to}
\Xx^{2d-1}(F)\to \Ker S^{2d}(F)\to
0\leqno(7.3.1.2)
$$
在 $D^{+}(Y,\mathcal{A}_{Y})$ 中典范同构于复形 \eqhyperref[XVII.eq7.3.1.1]{7.3.1.1}。此外：

\begin{enonce*}{引理 7.3.2}\etiquette[7.3.2]{XVII.lem7.3.2}
函子 $\Xx^{0}$、$\Xx^{1},\ldots,\Xx^{2d-1}$ 和 $\Ker S^{2d}$ 关于 $F$ 正合。
\end{enonce*}

函子 $\Xx^{n}$ 的正合性由典范松弛消解的正合性，以及松弛层短正合列在态射下的正像仍为短正合列这一事实得出。

为证明 $\Ker S^{2d}$ 的正合性，对于 $X$ 上层的短正合列 $0\rightarrow F'\rightarrow F\rightarrow F''\rightarrow 0$，考虑图
\[
\xymatrix{
\Ker S^{2d}(F')  \ar[r] \ar[d]& \Xx^{2d+1}(F') \ar[r] \ar[d]&
\Xx^{2d+2}(F') \ar[d] \ldots\\
\ker S^{2d}(F)   \ar[r] \ar[d]& \Xx^{2d+1}{(F)} \ar[r] \ar[d]&
\Xx^{2d+2}(F)\ar[d] \ldots \\
\ker S^{2d}(F'')   \ar[r] & \Xx^{2d+1}{(F'')} \ar[r] &
\Xx^{2d+2}{(F'')} \ldots \\
}
\]
，并利用正合列
\[
0\longrightarrow \sigma_{\geq 2d+1}
(\Xx^{\bigdot}(F'))\longrightarrow \sigma_{\geq 2d+1}
(\Xx^{\bigdot}(F))\longrightarrow \sigma_{\geq
2d+1}(\Xx^{\bigdot}(F''))\to0.
\]
相伴的同调正合列，即得结论。

现在有以下一般引理：

\begin{enonce*}{引理 7.3.3}\etiquette[7.3.3]{XVII.lem7.3.3}
设 $A$ 和 $B$ 为两个阿贝尔范畴，$\mathrm{Je}^{\bigdot}$ 为 $C^{b}(\Ffex(A,B))$ 的一个对象，其中 $\Ffex(A,B)$ 表示从 $A$ 到 $B$ 的正合函子范畴。将 $C(A)$ 中的每个复形 $K^{\bigdot}$ 对应到双复形 $\mathrm{Je}^{\bigdot}(K^{\bigdot})$ 的伴随单复形，便得到一个保持拟同构的三角函子 $K(A)\rightarrow K(B)$。
\end{enonce*}

该函子可写成\pageoriginale 复合
\[
C(A)\longrightarrow C^{b}(C(B))\xrightarrow{(\quad)_{s}}C(B),
\]
，其中 $(\quad)_{s}$ 为取伴随单复形的函子；应用（\sisi{VI b (\refnoexpose{Vbis.lem2.3.2.2})}{\hyperrefext[Vbis][Vbis.lem2.3.2.2]{2.3.2.2}} 的注）即可；细节留给读者。

\begin{enonce*}[remark]{习题 7.3.4}\etiquette[7.3.4]{XVII.exer7.3.4}
证明，对于 $D^{+}(X,\mathcal{A}_{X})$ 中的每个复形 $K^{\bigdot}$，典范箭头
\[
\tau_{\leq 2d} \Xx^{\bigdot}(K^{\bigdot})\longrightarrow
\Xx^{\bigdot}(K^{\bigdot})
\]
在相伴的单复形上诱导拟同构。
\end{enonce*}

为此，先注意典范箭头
\[
\varinjlim_{k\geq 2d} (\tau_{\leq k}
\Xx^{\bigdot}(K^{\bigdot}))_{s}\longrightarrow
(\Xx^{\bigdot}(K^{\bigdot}))_{s}
\]
是同构，再注意复形的上同调层与滤过正向极限交换。

\setcounter{subsubsection}{4}
\subsubsection{}\etiquette[7.3.5]{XVII.subsubsec7.3.5}
设 $R^{(1)}f_{!}:D(X,\mathcal{A}_{X})\rightarrow
D(Y,\mathcal{A}_{Y})$ 为利用 \sisi{\hyperref[XVII.rem7.2.3]{7.2.3}}{\hyperref[XVII.lem7.3.3]{7.3.3}} 由 \eqhyperref[XVII.eq7.3.1.2]{7.3.1.2} 定义的函子。鉴于 \sisi{\hyperref[XVII.prop7.2.4]{7.2.4}}{\hyperref[XVII.exer7.3.4]{7.3.4}}，它可直接由 \eqhyperref[XVII.eq7.3.1.1]{7.3.1.1} 计算，且可验证它“不依赖”所选的伪紧化。接着还应去掉上标 $+$，验证陈述 \hyperref[XVII.constru7.1.3]{7.1.3}、\hyperref[XVII.constru7.1.4]{7.1.4} 和 \hyperref[XVII.prop7.1.6]{7.1.6}，但撰稿人没有勇气这样做。

\begin{enonce*}{定义 7.3.6}\etiquette[7.3.6]{XVII.def7.3.6}
在 \hyperref[XVII.subsubsec7.3.5]{7.3.5} 中构造的函子 $Rf_{!}:D(X,\mathcal{A}_{X})\rightarrow
D(Y,\mathcal{A}_{Y})$ 仍称为\emph{具有固有支集的正像函子}。
\end{enonce*}

\setcounter{subsubsection}{6}
\subsubsection{}\etiquette[7.3.7]{XVII.subsubsec7.3.7}
于是可证明函子 $Rf_{!}$ 与基变换交换，并且陈述（\hyperref[XVII.prop5.2.9]{5.2.9}）和（\hyperref[XVII.thm5.2.10]{5.2.10}）对 $\Sch_{!R}$ 中的每个态射仍然成立，从而说明刚才的构造是合理的。

\begin{thebibliography}{\sisi{TF}{99}}
\sisi{\bibitem[1]{sga4.exp.XVII.bib.DP}}{\bibitem{sga4.exp.XVII.bib.DP}} A. Dold\pageoriginale 和 D. Puppe. Homologie nicht\sisi{ }{-}additiver
Funktoren. Anwendungen. Ann. Inst. Fourier \sisi{XI 1961}{\textbf{11} (1961), 201--312}.

\sisi{\bibitem[1]{sga4.exp.XVII.bib.GZ}}{\bibitem{sga4.exp.XVII.bib.GZ}} Gabriel-Zisman. Calculus of Fractions and Homotopy
Theory, Ergebnisse, Bd 35, Springer 1967.

\sisi{\bibitem[1]{sga4.exp.XVII.bib.Gi}}{\bibitem{sga4.exp.XVII.bib.Gi}} J. Giraud. Méthode de la descente, Mémoire nᵒ 2 de la Société Mathématique de France (1964).

\sisi{\bibitem[TF]{sga4.exp.XVII.bib.TF}}{\bibitem{sga4.exp.XVII.bib.TF}} R. Godement. Théorie des faisceaux, Hermann,
1958. Act. Scient. Ind. nᵒ 1252, (Paris)\sisi{（引作 [TF]）}{}.

\sisi{\bibitem[1]{sga4.exp.XVII.bib.Gr}}{\bibitem{sga4.exp.XVII.bib.Gr}} A. Grothendieck. TDTE IV: les schémas de
Hilbert. Séminaire Bourbaki 221. 1961 年 5 月。重印于 FGA。

\sisi{\bibitem[RD]{sga4.exp.XVII.bib.RD}}{\bibitem{sga4.exp.XVII.bib.RD}} R. Hartshorne. Residues and Duality. Lecture Notes
in Mathematics nᵒ 20. Springer 1966.

\sisi{\bibitem[1]{sga4.exp.XVII.bib.La}}{\bibitem{sga4.exp.XVII.bib.La}} D. Lazard. 学位论文 --- Paris 1968.

\sisi{}{\bibitem{sga4.exp.XVII.bib.Mu} D. Mumford, \emph{Geometric Invariant Theory}, Ergebnisse der Mathematik und ihrer Grenzgebiete 34, Springer-Verlag, 1965.}

\sisi{}{\bibitem{sga4.exp.XVII.bib.Na} M. Nagata, ``A generalization of the imbedding problem of an abstract variety in a complete variety'', \emph{J. Math. Kyoto Univ.} \textbf{3} (1963), 89--102}

\sisi{\bibitem[1]{sga4.exp.XVII.bib.Ro}}{\bibitem{sga4.exp.XVII.bib.Ro}} N. Roby. Lois polynômes et lois formelles en théorie
des modules. Ann. E.N.S. \textbf{80} \sisi{1963, p.}{(1963), }231--\sisi{2}{3}48.

\sisi{\bibitem[CD]{sga4.exp.XVII.bib.CD}}{\bibitem{sga4.exp.XVII.bib.CD}} J.L. Verdier. Catégories dérivées (état 0),
由 I.H.E.S. 油印。

\sisi{\bibitem[2]{sga4.exp.XVII.bib.V2}}{\bibitem{sga4.exp.XVII.bib.V2}} J.L. Verdier. Catégories dérivées（将刊于
Grundlehren der Math. Wiss, Springer）。

\sisi{\bibitem[3]{sga4.exp.XVII.bib.V3}}{\bibitem{sga4.exp.XVII.bib.V3}} J.L. Verdier. Dualité dans la cohomologie des
espaces localement compacts. Séminaire Bourbaki 300. 1965 年 11 月。
\end{thebibliography}


%0.2 no ref. 1957, 1329, 5575, 5594, 5821
%2.1.2.2 no ref. 1386
%2.1.2.3 no ref  1412
%2.1.12.1 no ref 1641
%3.2.4.3 no ref  2680
%5.1.8.2 no ref 3925
%XII 4.1 no ref 4690, 6830
%6.3.4.3 no ref 7561, 7876, 8060,8409
%6.3.11 no ref 7736, 7833
%6.2.17 no ref 8042
%6.3.17.1 no ref 8078
%6.3.2.3 no ref 8271
%4.3.0 no ref 8569
%3.7 no ref 9229


\ifx\danslelivre\undefined
\let\SGAXVIIenddocument\enddocument
\else
\let\SGAXVIIenddocument\relax
\fi
\SGAXVIIenddocument


% END native input: chapters/exp-17-part6-zh.tex
% BEGIN native input: chapters/exp-18-part1-zh.tex
\chapter{整体对偶公式}
%\addtocontents{toc}{P. Deligne}
\begin{center}
P. Deligne
\end{center}
\phantomsection\MakeLinkTarget*{sga4.exp.XVIII}\label{sga4.exp.XVIII}

% The edition is fixed to the corrected (second) reading. This expandable
% selector is output-equivalent and avoids current Xy-pic conditional failures.
\renewcommand{\sisi}[2]{#2}
\let\Arrowvert\Vert

\setcounter{pageoriginal}{480}
\newcommand{\acorriger}[1]{\nde{#1}}

\setcounter{section}{-1}
\section{引言}\etiquette[0]{XVIII.sec0}\pageoriginale

\subsection{}\etiquette[0.1]{XVIII.subsec0.1}
本讲讨论 Poincaré 对偶。我们仿照 \textsc{Verdier} [1] 的方法构造这一对偶；他处理的是局部紧 Hausdorff 拓扑空间的情形。我们先验地证明第 XVII 讲的函子 $Rf_{!}$ 有一个右伴随 $Rf^{!}$，并建立其许多性质，然后才在光滑态射的情形下更明确地计算它。

这一构造之所以可能，似乎基于函子 $Rf_{!}$ 的下列性质：
\begin{enumerate}
\renewcommand{\theenumi}{\alph{enumi}}
\renewcommand{\labelenumi}{\theenumi)}
\item 对于伴随 $Rf^{!}$ 的存在性及其形式性质：
\begin{itemize}
\item[(I)] 函子 $Rf_{!}$ 与基变换可交换；

\item[(II)] 函子 $R^{i}f_{!}$ 与滤过归纳极限可交换；当 $i$ 充分大时，它们为零；

\item[(III)] 函子 $Rf_{!}$ 可以“在复形层次上”计算，即从复形范畴之间的一个函子出发来计算。
\end{itemize}

\item 对于 $f$ 光滑时这一伴随的计算：
\begin{itemize}
\item[(IV)] 计算“小” $U$ 的 $Rf_{!}((\ZZ/n\ZZ)_{U})$。
\end{itemize}
\end{enumerate}

\subsection{}\etiquette[0.2]{XVIII.subsec0.2}
在独立于第 \hyperref[XVIII.sec1]{\S 1}、2 节的第 \hyperref[XVIII.sec3]{\S 3} 节第 1 小节中，我们构造函子 $Rf^{!}$；那里用到了适合这一特定情形的构造，而不只是 (I)、(II)、(III)。这使该小节颇不令人愉快；作者承认，他并非总能完全理解其中发生了什么。

第 \hyperref[XVIII.sec2]{2} 节讨论 (IV)；它完全由\pageoriginale{}研究曲线上同调的第 \hyperref[XVIII.sec1]{\S 1} 节形式地推出。

在第 \hyperref[XVIII.sec1]{1} 节中，我们还研究连续系数的情形，所证明的内容远多于第 \hyperref[XVIII.sec2]{2} 节所需。该节深受 \textsc{Serre} \cite{sga4.exp.XVIII.bib.6} 的影响。对后文有用的要点是：
\begin{enumerate}
\renewcommand{\theenumi}{\alph{enumi}}
\renewcommand{\labelenumi}{\theenumi)}
\item 第~\hyperref[XVIII.subsec1.1]{1.1} 小节讨论曲线情形的迹态射，并在 \hyperref[XVIII.thm2.9]{2.9} 中加以推广；

\item 消没定理 \hyperref[XVIII.lemfond1.6.9]{1.6.9}。愿意接受超越方法论证的读者，可以从 \hyperref[XVIII.lem1.6.6]{1.6.6}（第二种证明）开始阅读第~\hyperref[XVIII.subsec1.6]{1.6} 小节，而不依赖第~\hyperref[XVIII.subsec1.2]{1.2} 至 \hyperref[XVIII.subsec1.5]{1.5} 小节\nde{对于初读者，我们极力推荐这条阅读路线：在 $\S$\hyperref[XVIII.subsec1.1]{1.1} 之后，直接转到 $\S$\hyperref[XVIII.sec2]{2} 和 $\S$\hyperref[XVIII.sec3]{3}。}。\hyperref[XVIII.lemfond1.6.9]{1.6.9} 的关键成分是代数闭域上光滑曲线的任一种形式的 Poincaré 对偶，再加上无环性引理 \hyperrefext[XV][XV.lem2.6]{2.6}（它是光滑基变换定理 \hyperrefext[XVI][XVI.thm1.1]{1.1} 的一个预备结果）。其结论在 \hyperref[XVIII.thm2.14]{2.14} 中得到推广。
\end{enumerate}

第 \hyperref[XVIII.subsec3.2]{3.2} 小节汇集第 \hyperref[XVIII.sec2]{2} 节的推论。

\subsection{}\etiquette[0.3]{XVIII.subsec0.3}
\textsc{Verdier} 在 Poincaré 对偶中采用的方法，使我们能够对可紧化的 $f$ 定义函子 $Rf^{!}$。这种一般性的代价是，必须验证的相容性并不熟悉，有时即使是最平凡的相容性，如 \hyperref[XVIII.lem3.2.3]{3.2.3}，也需要十分奇特的证明。作者承认，他证明的这类相容性没有应有的那么多。

有兴趣的读者不难利用第 XVII 讲的附录，将 $Rf^{!}$ 的定义推广到有限型分离态射，只要其终点为凝聚概形（即拟紧拟分离概形）。

\subsection{}\etiquette[0.4]{XVIII.subsec0.4}
主要结果 \hyperref[XVIII.thm3.2.5]{3.2.5} 用具体的语言陈述，使我们可以形式地回到讨论班口头报告所采用的观点：在那里，函子 $Rf^{!}$ 仅对\pageoriginale{}可光滑化态射定义，所用方法是将态射分解为一个光滑态射与一个浸入的复合。口头报告中证明的概要可见 \textsc{Verdier} \cite{sga4.exp.XVIII.bib.8}。

\section{曲线的上同调}\etiquette[1]{XVIII.sec1}\pageoriginale

\subsection{迹态射}\etiquette[1.1]{XVIII.subsec1.1}

\subsubsection{}\etiquette[1.1.1]{XVIII.subsubsec1.1.1}
设 $n$ 是一个 $\geq 1$ 的整数，$X$ 是一个其上 $n$ 可逆的概形。于是 $X$ 上的层 $𝛍_n$（IX \hyperref[IX.eq3.1]{3.1}）是常值环层 $\ZZ/n$ 上的秩为 1 的局部自由模层。对 $i\in \ZZ$，以 $\ZZ/n(i)$ 表示它的 $i^{\text{\rm th}}$ 张量幂：
\begin{equation*}
\ZZ/n(i)=\Hhom(\ZZ/n,\GG_{m})^{\otimes
    i}{\quoi}.\tag{1.1.1.1}
\end{equation*}\etiquette[1.1.1.1]{XVIII.eq1.1.1.1}

若 $n=dn'$，取 $d$ 次幂给出一个同构
$$
𝛍_n\otimes_{\ZZ/n}\ZZ/n'\xrightarrow{\sim}𝛍_{n'},
$$
从而给出称为典范同构的同构：
\begin{equation*}
\ZZ/n(i)\otimes_{\ZZ/n}\ZZ/n'\xrightarrow{\sim} \ZZ/n'(i){\quoi}.\tag{1.1.1.2}
\end{equation*}\etiquette[1.1.1.2]{XVIII.eq1.1.1.2}

若 $F$ 是满足 $nF=0$ 的阿贝尔层，置
\begin{equation*}
F(i)=F\otimes_{\ZZ/n}\ZZ/n(i){\quoi}.\tag{1.1.1.3}
\end{equation*}\etiquette[1.1.1.3]{XVIII.eq1.1.1.3}%
若已有 $n'F=0$，同构 \eqhyperref[XVIII.eq1.1.1.2]{1.1.1.2} 使我们能够将 $F\otimes_{\ZZ/n}\ZZ/n(i)$ 与 ${F\otimes_{\ZZ/n'}\ZZ/n'(i)}$ 等同，因而 $F(i)$ 只依赖于 $F$，不依赖于 $n$ 的选择。更一般地，若 $F$ 是与 $X$ 的剩余特征互素的挠层（\hyperrefext[XVII][XVII.subsec0.13]{0.13}），${}_{m}F$ 是 $F$ 上乘以 $m$ 的核，则定义 $F(i)$ 为以下归纳极限，其中 $m$ 在 $X$ 上可逆的整数中按乘法趋于无穷：
\begin{equation*}
F(i)=\varinjlim_{m}{}_{m}F(i){\quoi}.\tag{1.1.1.4}
\end{equation*}\etiquette[1.1.1.4]{XVIII.eq1.1.1.4}

函子 $F\mapsto F(i)$ 是正合的。若 $f:Y\rightarrow X$ 是一个概形态射（\resp{}拟紧拟分离概形态射），且 $F$ 是 $X$ 上（\resp $Y$ 上）与 $X$ 的剩余特征互素\pageoriginale{}的挠层，则 $f^{\ast}(F(i))\simeq(f^{\ast}F)(i)$（\resp $(f_{\ast}F)(i)\simeq f_{\ast}(F(i))$）。若 $\mathcal{A}$ 是 $X$ 上被某个在 $X$ 上可逆的整数 $n$ 零化的环层，这就推广为从 $D(X,\mathcal{A})$ 到 $D(X,\mathcal{A})$ 的一个正合函子 $K\mapsto K(i)$。该函子与逆像函子和正像函子可交换。

函子 $F\mapsto F(i)$ 有时称为“\emph{Tate 扭曲}（Tate twist）”函子。

\begin{enonce*}{定义 1.1.2}\etiquette[1.1.2]{XVIII.def1.1.2}
概形 $S$ 上的\emph{平坦曲线}（flat curve）是一个有限表现平坦分离态射 $f:X\rightarrow S$，其纤维均为纯一维。
\end{enonce*}

“纯一维”这一说法并不排除空概形。与 \EGA II 7.4.2 不同，这里允许域上的曲线为空，但要求它是分离的。

当 $S$ 是一个域的谱时，我们只称之为 $S$ 上的曲线；域上的曲线是拟射影的\nde{在 \EGA II 7.4.10 中，这一断言仅在 $X$ 正规的附加假设下得到证明，而其中曾预告（\EGA II 7.4.12），一般情形的证明将出现在尚未出版的 \EGA V 中。读者可以将填补这一空缺作为练习。无论如何，这并不影响迹态射的构造，因为可以先在局部（即拟射影开子概形上）进行构造，再如 \hyperref[XVIII.prop1.1.6]{1.1.6} 的证明中那样将迹态射黏合。}。

\setcounter{subsubsection}{2}
\subsubsection{}\etiquette[1.1.3]{XVIII.subsubsec1.1.3}
设 $X$ 是代数闭域 $k$ 上的曲线，其特征指数 $p$ 与整数 $n\geq 1$ 互素。先假定 $X$ 约化，且设 $\bar{X}$ 为 $k$ 上的一条完备曲线，$X$ 是其中的稠密\emph{开子概形}。

补集 $Y=\bar{X}-X$ 的维数为 $0$。因而上同调正合列（\hyperrefext[XVII][XVII.eq5.1.16.3]{5.1.16.3}）诱导一个同构
\begin{equation*}
H^{2}_{c}(X,\ZZ/n(1))\xrightarrow{\sim} H^{2}(\bar{X},\ZZ/n(1)){\quoi}.\tag{1.1.3.1}
\end{equation*}\etiquette[1.1.3.1]{XVIII.eq1.1.3.1}%
我们还有一个同构 \hyperrefext[IX][IX.cor4.7]{4.7}（由 Kummer 理论给出）
\begin{equation*}
H^{2}(\bar{X},\ZZ/n(1))=\Pic(\bar{X})/n=(\ZZ/n)^{c}\tag{1.1.3.2}
\end{equation*}\etiquette[1.1.3.2]{XVIII.eq1.1.3.2}%
其中 $c$ 表示 $X$ 的不可约分支的集合，或等价地，表示 $\bar{X}$ 的不可约分支的集合。

若\pageoriginale{}不再要求 $X$ 约化，且 $c$ 是其不可约分支的集合，则对 $i\in c$，以 $n_{i}$ 表示 $X$ 在其第 $i^{\text{\rm th}}$ 个不可约分支的一般点 $\eta_{i}$ 的重数：
$$
n_{i}=\lg(O_{X,\eta_{i}}){\quoi}.
$$

我们知道，从 $H^{2}_{c}(X,\ZZ/n(1))$ 到 $H^{2}_{c}(X_{\red}, \ZZ/n(1))$ 的限制映射是一个同构。因此，将 \eqhyperref[XVIII.eq1.1.3.1]{1.1.3.1} 和 \eqhyperref[XVIII.eq1.1.3.2]{1.1.3.2} 应用于 $X_{\text{\rm red}}$，得到 $H^{2}_{c}(X,\ZZ/n(1))$ 与 $(\ZZ/n)^{c}$ 之间的一个典范同构。以 $\Tr_{X}$，或简记为 $\Tr$，表示复合映射
\begin{equation*}
\Tr:H^{2}_{c}(X,\ZZ/n(1))\xrightarrow{\sim}(\ZZ/n)^{c}\xrightarrow{t}\ZZ/n{\quoi},\tag{1.1.3.3}
\end{equation*}\etiquette[1.1.3.3]{XVIII.eq1.1.3.3}%
其中
$$
t((a_{i})_{i\in c})=\sum n_{i}a_{i}{\quoi}.
$$

若 $F$ 是被 $n$ 零化的阿贝尔群，则有同构
$$
H^{2}_{c}(X,F(1))\mathop{\longleftrightarrow}^{\sim}F\otimes
H^{2}_{c}(X,\ZZ/n(1))
$$
从而由 \eqhyperref[XVIII.eq1.1.3.3]{1.1.3.3} 得到一个\emph{迹态射}（trace morphism）
\begin{equation*}
\Tr:H^{2}_{c}(X,F(1))\longrightarrow F{\quoi}.\tag{1.1.3.4}
\end{equation*}\etiquette[1.1.3.4]{XVIII.eq1.1.3.4}

若 $n=n'd$，则图
$$\vcenter{\etiquette[1.1.3.5]{XVIII.eq1.1.3.5}
\xymatrix@R=1.25cm{
0  \ar[r] & 𝛍_n  \ar[r]\ar[d]_{x^d} & \GG_m \ar[d]_{x^d} \ar[r]^{x^n} & {\GG}_m 
\ar[r] \ar@{=}[d]& 0\\
0  \ar[r] & 𝛍_{n'}  \ar[r] & \GG_m  \ar[r]^{x^{n'}} & {\GG}_m 
\ar[r] & 0
}}\leqno{(1.1.3.5)}
$$
交换，因而根据\pageoriginale \hyperrefext[IX][IX.cor4.7]{4.7} 中同构 \eqhyperref[XVIII.eq1.1.3.2]{1.1.3.2} 的定义，图
$$\vcenter{\etiquette[1.1.3.6]{XVIII.eq1.1.3.6}
  \xymatrix@C=2cm@R=1.25cm{
H^{2}_c(X,\ZZ/n(1)) \ar[d] \ar[r]^-{\sim} & (\ZZ/n)^{c} \ar[d]\\
H^{2}_c(X,\ZZ/n'(1)) \ar[r]^-{\sim}  & (\ZZ/n')^{c}
}}\leqno{(1.1.3.6)}
$$
也交换。

因此，映射 \eqhyperref[XVIII.eq1.1.3.4]{1.1.3.4} 只依赖于 $F$，而不依赖于与 $p$ 互素、满足 $nF=0$ 的整数 $n$ 的选择。通过取极限，我们对每个与 $p$ 互素的挠阿贝尔群 $F$ 定义该映射。

下面的引理由定义立即得到。

\begin{enonce*}{引理 1.1.4}\etiquette[1.1.4]{XVIII.lem1.1.4}
在上述假设下，若对 $i\in c$，$U_{i}$ 是包含在 $X$ 的第 $i^{\text{\rm th}}$ 个不可约分支中的一个非空开子概形，则图
$$
\xymatrix@R=1.5cm@C=1.25cm{
\displaystyle{\mathop{\oplus}_{i\in c}} H^2_c(U_i, F(1))
\ar[rr]^{\sim} \ar[rd]_{\displaystyle{\mathop{\oplus}_{i\in c}}
  Tr_{U_i}}& & H^2_c(X,F(1))\ar[ld]^{Tr_X}\\
 & F &
}
$$
交换。
\end{enonce*}

\begin{enonce*}{引理 1.1.5}\etiquette[1.1.5]{XVIII.lem1.1.5}
设 $X$、$Y$ 是代数闭域 $k$ 上的两条曲线，$u:X\rightarrow Y$ 是拟有限平坦态射，$n$ 是与 $k$ 的特征指数 $p$ 互素的整数。

那么，下图交换，其中 $\Tr_{u}$ 是映射\emph{（\hyperrefext[XVII][XVII.thm6.2.3]{6.2.3}）：}
$$
\xymatrix@R=1.5cm{
  H^2_c(X,\ZZ/n(1)) \ar[rr]^{Tr_u} \ar[rd]_{Tr_X}&              & H^2_c(Y,\ZZ/n(1)) \ar[ld]^{Tr_Y}\\
                         & \ZZ/n &. 
}
$$
\end{enonce*}

由引理 \hyperref[XVIII.lem1.1.4]{1.1.4}，可以归结为 $X$、$Y$ 不可约且 $X_{\text{\rm red}}$、$Y_{\text{\rm red}}$ 光滑的情形。若 $e$（\resp $f$）是 $X$（\resp $Y$）在一般点的重数，且 $e=df$，则下图中标有 $+$ 的三角形交换：
$$
\vcenter{\etiquette[1.1.5.1]{XVIII.eq1.1.5.1}
\xymatrix@R=1.5cm{
H^2_c(X,\ZZ/n(1))
\ar[rr]^{\sim}_(.5){{\displaystyle\mathop{\phantom{a}}_{\displaystyle\mathop{{\phantom{a}}}_{{\displaystyle\mathop{+}}}}}}  \ar[dd]_{\Tr_u} \ar[rd]_{\Tr_X}& &
H^2_c(X_{\red},\ZZ/n(1)) \ar[dd]^{d.\Tr_{u_{\red}}} \ar[ld]^{e.\Tr_{X_{\red}}}\\
                               & \ZZ/n &\\
H^2_c(Y,\ZZ/n(1)) \ar[rr]^{\displaystyle\mathop{\sim}^{\displaystyle\mathop{\phantom{a}}^{{\displaystyle\mathop{\phantom{a}}^{\textstyle{+}}}}}} \ar[ru]^{\Tr_Y}&  & H^2_c(Y_{\red},\ZZ/n(1))\ar[lu]_{f.\Tr_{Y_{\red}}}, 
}}\sisi{}{\leqno{(1.1.5.1)}}
$$

它的\pageoriginale{}外框交换，这是下图交换的结果：
$$
\xymatrix@C=2cm@R=1.5cm{
u_!\ZZ/n_X(1)           \ar[r]^{\sim} \ar[d]^{\Tr_u} &
u_{{\red}_!}\ZZ/n_{X_{\red}}(1) \ar[d]^{d.\Tr_{u_{\red}}}\\
\ZZ/n_Y(1)  \ar[r]^{\sim}& \ZZ/n_{Y_{\red}}(1).}
$$
只需在 $Y$ 的一般点验证后一相容性，这显然成立。

\sisi{}{因此，为了证明 \eqhyperref[XVIII.eq1.1.5.1]{1.1.5.1} 中左侧三角形交换，只需证明右侧三角形交换。所以，现在还可以假定 $X$、$Y$ 约化\nde{本段为本版所增，下两段也作了修改。}。}

设 $\bar{X}$、$\bar{Y}$ 是分别以 \sisi{$X_{\red}$}{$X$}、\sisi{$Y_{\red}$}{$Y$} 为稠密开子概形的完备\sisi{非光滑}{光滑}曲线（\EGA II 7.4.11）。态射 \sisi{$u_{\red}$}{$u$} 延拓为一个平坦态射 $\bar{u}:\bar{X}\rightarrow \bar{Y}$（\EGA II 7.4.9），而由 \hyperref[XVIII.lem1.1.4]{1.1.4}，只需对 $\bar{u}$ 验证 \hyperref[XVIII.lem1.1.5]{1.1.5}。

由（XVII \hyperref[XVII.eq6.3.18.2]{6.3.18.2}），下面这个 $\bar{Y}$ 上层的图交换，其中 $N$ 表示范数：
$$
\xymatrix@R=1.25cm{
0  \ar[r] & \bar{u}_{\ast}𝛍_n  \ar[r] \ar[d]^{\Tr_{\bar{u}}}&
\bar{u}_{\ast}\GG_m  \ar[r]^{x^n} \ar[d]^{N}&
\bar{u}_{\ast} \GG_m  \ar[r] \ar[d]^{N}& 0 \\
0  \ar[r] & 𝛍_n  \ar[r] & \GG_m  \ar[r]^{x^n} & \GG_m 
\ar[r] & 0 ,
}
$$
因而图
$$
\xymatrix@C=1.5cm@R=1.25cm{
\Pic(\bar{X})/n \ar[r]^(.45){\sim}\ar[d]^{N}& H^2_c(\bar{X},\ZZ/n(1))
\ar[d]^{\Tr_{\bar{u}}}\\ 
\Pic(\bar{Y})/n \ar[r]^(.45){\sim}& H^2_c(\bar{Y},\ZZ/n(1)),\\
}
$$
也交换；\pageoriginale{}剩下只需注意，若 $\Ll$ 是 $\bar{X}$ 上次数为 1 的可逆层，则其范数的次数仍为 $1$（考虑到除子范数的定义 \EGA IV 21.5.5，这一事实由相容性 \sisi{\EGA IV 21.10.7}{\EGA IV 21.10.17} 得到）。

\begin{enonce*}{命题 1.1.6}\etiquette[1.1.6]{XVIII.prop1.1.6}
对于每条可紧化平坦曲线 $f:X\rightarrow S$ 和 $S$ 上每个与 $S$ 的剩余特征互素的挠层 $F$，存在唯一一种方式定义迹态射
$$
\Tr_{f}:R^{2}f_{!}(f^{\ast}F(1))\longrightarrow F{\quoi},
$$
使之关于 $F$ 具有函子性，其形成与任意基变换相容\emph{（\cf \hyperrefext[XVII][XVII.thm6.2.3]{6.2.3}，(VAR 2)）}，而当 $S$ 为代数闭域的谱时，它与迹态射 \eqhyperref[XVIII.eq1.1.3.4]{1.1.3.4} 一致。
\end{enonce*}

态射 $\Tr_{f}$ 是逐纤维指定的；因此其唯一性是显然的。

设 $p:\PP^{1}_{S}\rightarrow S$ 为 $S$ 上的射影直线，$n$ 是在 $S$ 上可逆的整数。有 $\sheafPic_{S}(\PP^{1}_{S})=\ZZ_{S}$，因而 Kummer 理论（\hyperrefext[IX][IX.Thr3.2]{3.2}）给出一个态射 $\ZZ/n\rightarrow R^{2}p_{\ast}\ZZ/n(1)$，逐纤维验证（\hyperrefext[XII][XII.cor5.2]{5.2}）可知它是同构。其逆 $\Tr_{p}$ 逐纤维诱导迹态射 \eqhyperref[XVIII.eq1.1.3.4]{1.1.3.4}。

设 $f:X\rightarrow S$ 是 $S$ 上的平坦曲线，使得存在从 $X$ 到 $\PP^1_S$ 的拟有限平坦 $S$-态射 $u$。由 \hyperref[XVIII.lem1.1.5]{1.1.5}，复合态射 $\Tr_{p}\circ \Tr_{u}$（\hyperrefext[XVII][XVII.thm6.2.3]{6.2.3}）在 $S$ 的每个几何点的纤维都是迹态射 \eqhyperref[XVIII.eq1.1.3.4]{1.1.3.4}。特别地，它不依赖于 $u$ 的选择。

假设\pageoriginale $X$ 是开子概形 $U_{i}$ 的并，且每个 $U_{i}$ 都存在到 $\PP^{1}_S$ 的拟有限平坦态射。若 $\alpha_{i}$（\resp $\alpha_{ij}$）是 $U_{i}$（\resp $U_{ij}=U_{i}\cap U_{j}$）到 $X$ 的包含，则序列
$$
\sum R^{2}(f\alpha_{ij})_{!}\ZZ/n(1)\rightrightarrows
\sum R^2(f\alpha_{i})_{!}\ZZ/n(1)\rightarrow
R^{2}f_{!}\ZZ/n(1)\rightarrow 0{\quoi}.
$$
由（\hyperrefext[XVII][XVII.prop6.2.9]{6.2.9}）及 $R^{2}f_{!}$ 的右正合性（因为 $R^{3}f_{!}=0$）是正合的。利用 \hyperref[XVIII.lem1.1.5]{1.1.5} 逐纤维验证，可知迹态射之和 $\underset{i}\sum R^{2}(f\alpha_{i})_{!}\ZZ/n(1)\rightarrow \ZZ/n$ 通过一个从 $R^{2}f_{!}\ZZ/n(1)$ 出发的迹态射 $\Tr_{f}$ 分解，此态射逐纤维诱导 \eqhyperref[XVIII.eq1.1.3.4]{1.1.3.4}。


对于一般的 $X$，设 $j:X'\rightarrow X$ 是 $X$ 的最大开子集，使 $f|X'$ 为 Cohen--Macaulay 态射（\cf \EGA IV 12.1.1 (vi)）。$X'$ 的每一点 $x$ 在 $X'$ 中都有一个邻域 $U_{x}$，使得存在从 $U_{x}$ 到 $\PP^{1}_{S}$ 的拟有限 $S$-态射 $u_{x}$。该态射自动平坦（\EGA IV 11.3.10 和 15.4.2），所以 $X'$ 满足前面的假设。此外，$X'$ 在 $X$ 的每个纤维中稠密；因而从 $R^{2}(fj)_{!}\ZZ/n(1)$ 到 $R^{2}f_{!}\ZZ/n(1)$ 的态射是同构，这可以逐纤维验证，而复合箭头
$$
\Tr_{f}:R^{2}f_{!}\ZZ/n(1)\xleftarrow{\sim}
R^{2}(fj)_{!}\ZZ/n(1)\xrightarrow{\Tr_{fj}}\ZZ/n
$$
在每个纤维上诱导 \eqhyperref[XVIII.eq1.1.3.4]{1.1.3.4}。

若 $F$ 满足 $nF=0$，则有（\hyperrefext[XVII][XVII.thm5.2.6]{5.2.6}）
\begin{equation*}
R^{2}f_{!}\ZZ/n(1)\otimes F\xrightarrow{\sim}
R^{2}f_{!}(f^{\ast}F(1))\tag{1.1.6.1}
\end{equation*}\etiquette[1.1.6.1]{XVIII.eq1.1.6.1}%
并定义 $\Tr_{f}$ 为复合箭头
$$
\Tr_{f}:R^{2}f_{!}(f^{\ast}F(1))\xleftarrow{\sim}
R^{2}f_{!}\ZZ/n(1) \otimes F\rightarrow F{\quoi}.
$$
一般情形下，将 $\Tr_{f}$ 定义为关于各层 ${}_nF$ 的迹态射的滤过余极限，其中 $n$ 在 $S$ 上可逆。这些箭头在每个纤维上诱导 \eqhyperref[XVIII.eq1.1.3.4]{1.1.3.4}，因而与任意基变换可交换。\cqfd

利用 \hyperref[XVIII.lem1.1.4]{1.1.4}、\hyperref[XVIII.lem1.1.5]{1.1.5} 及 \eqhyperref[XVIII.eq1.1.6.1]{1.1.6.1} 逐纤维\pageoriginale{}验证，可得以下结果：

\begin{enonce*}{引理 1.1.7}\etiquette[1.1.7]{XVIII.lem1.1.7}
设 $f:Y\rightarrow S$ 是 $S$ 上的可紧化平坦曲线，$u:X\rightarrow Y$ 是有限表现拟有限平坦态射，$F$ 是 $S$ 上的挠层，其挠性与 $S$ 的剩余特征互素。下图交换：
$$
\xymatrix@C=1.5cm@R=1.25cm{
R^2(fu)_! (fu)^{\ast}F(1) \ar[d]^{\Tr_{fu}}\ar@{=}[r]& R^2f_!(u_!
u^{\ast})f^{\ast}F(1) \ar[r]^{\Tr_u}& R^2f_!f^{\ast}F(1)\ar[d]^{\Tr_f} \\
F     \ar@{=}[rr] &&  F,
}
$$
\end{enonce*}

\begin{enonce*}{引理 1.1.8}\etiquette[1.1.8]{XVIII.lem1.1.8}
设 $u:Y\rightarrow S$ 是有限表现拟有限平坦态射，$f:X\rightarrow Y$ 是 $S$-可紧化的平坦曲线\emph{（\hyperrefext[XVII][XVII.def3.2.1]{3.2.1}）}，且 $F$ 如 \hyperref[XVIII.lem1.1.7]{1.1.7} 中所述。下图交换\nde{这一相容性值得作一些论证。为验证它，可以先作基变换，归结为 $S$ 是代数闭域 $k$ 的谱的情形。再可假定 $Y$ 连通且非空，此时 $Y\to S$ 有一个（唯一的）截面 $S\to Y$，它是由一个幂零理想定义的闭浸入。\hyperref[XVIII.prop1.1.6]{1.1.6} 证明中的论证使我们可以假定 $X=\mathbf{P}^1_Y$。剩下只需注意，$Y$ 的重数等于 $X$ 在其一般点的重数。另见 \eqhyperref[XVIII.eq2.9.3]{2.9.3} 交换性的证明。}：
$$
\xymatrix@C=1.5cm@R=1.25cm{
R^{2}(uf)_! (uf)^{\ast}F(1)  \ar@{=}[r] \ar[d]^{\Tr_{uf}}&
u_!R^2f_!f^{\ast} u^{\ast}F(1) \ar[r]^{\Tr_f} &
u_!u^{\ast}F \ar[d]^{\Tr_u}\\
F        \ar@{=}[rr] & &  F
}
$$
\end{enonce*}

同样，由 \eqhyperref[XVIII.eq1.1.3.2]{1.1.3.2} 可验证：

\begin{enonce*}{引理 1.1.9}\etiquette[1.1.9]{XVIII.lem1.1.9}
若 $f:X\rightarrow S$ 是可紧化光滑曲线，且 $f$ 的几何纤维不可约，则迹态射（\hyperref[XVIII.prop1.1.6]{1.1.6}）是同构。
\end{enonce*}

\subsection{射影空间的 $1$-无环性}\etiquette[1.2]{XVIII.subsec1.2}

本节原则上始终使用 fppf 拓扑（\hyperrefext[XVII][XVII.subsec0.10]{0.10}）。

\subsubsection{}\etiquette[1.2.1]{XVIII.subsubsec1.2.1}

设 $\epsilon$ 是概形 $S$ 上处处秩\pageoriginale $\geq 2$ 的局部自由模，$p:\PP(\epsilon)\rightarrow S$ 为相应的射影丛，$G$ 是 $S$ 的大 fppf 景上的阿贝尔层（\hyperrefext[XVII][XVII.subsec0.10]{0.10}）。每一对 $(K_{0},\varphi)$，其中 $K_{0}$ 为 $G$ 下的一个挠子（torsor），$\varphi:\GG_{mS}\rightarrow G$ 为同态，都在 $\PP(\epsilon)$ 上定义一个 $p^{\ast}G$-挠子 $K(K_{0},\varphi)$，它是 $K_{0}$ 与 $\GG_{m}$-挠子 $\Oo(1)$ 在 $\varphi$ 下的像之和\nde{本节中，我们默认借助如下对应等同线丛与 $\GG_{m}$-挠子：将一个线丛对应到其可逆截面层，而 $\GG_{m}$ 通过乘法作用于该层。注重细节的读者或许会希望知道，依照 Grothendieck 的约定，射影概形 $\PP(\epsilon)$ 是 $\epsilon$ 的超平面丛；同样，去零截面的向量丛 $\VV(\epsilon)^{\ast}\to \PP(\epsilon)$ 是 $\VV(\epsilon)$ 的一个开子集，后者为 $\epsilon$ 的对称代数的谱。}：
\begin{equation*}
K(K_{0},\varphi)=K_{0}+\varphi \Oo(1){\quoi}.\tag{1.2.1.1}
\end{equation*}\etiquette[1.2.1.1]{XVIII.eq1.2.1.1}%
规定
$$
\Hom((K_{0},\varphi),(K'_{0},\varphi'))=
\begin{cases}
∅\quad\mbox{\rm 若}\quad \varphi\neq \varphi'\\
\Hom(K_{0},K'_{0})\quad\mbox{\rm 若}\quad \varphi=\varphi'{\quoi}.
\end{cases}
$$
便将这些对 $(K_{0},\varphi)$ 组成一个范畴。于是，挠子 $K(K_{0},\varphi)$ 关于 $(K_{0},\varphi)$ 具有函子性。它与局部化相容，因而定义了从 $S$ 上由各对 $(K_{0},\varphi)$ 组成的叠，到 $S$ 上由 $\PP(\epsilon)$ 上的 $p^{\ast}G$-挠子组成的叠的一个函子（后者的叠在 $U/S$ 上的截面就是 $\fprod{U}{\PP}{S}(\epsilon)$ 上的 $p^{\ast}G$-挠子）。这些叠都是群胚叠（即在一个 $U/S$ 上的每个态射都是同构）。

\begin{enonce*}{定理 1.2.2}\etiquette[1.2.2]{XVIII.thm1.2.2}
在 \hyperref[XVIII.subsubsec1.2.1]{1.2.1} 的假设下，若 $G$ 是 $S$ 上有限表现的平坦交换群概形，则 \hyperref[XVIII.subsubsec1.2.1]{1.2.1} 中的函子 $K$ 是等价。
\end{enonce*}

\setcounter{subsubsection}{2}
\paragraph{}\etiquette[1.2.2.1]{XVIII.par1.2.2.1}

这一断言等价于以下两个条件同时成立：
\begin{gather*}
G\xrightarrow{\sim} p_{\ast}p^{\ast}G\quad\mbox{\rm 且}\\
\Hhom(\GG_{m},G)\xrightarrow{\sim} R^{1}p_{\ast}p^{\ast}G{\quoi}.
\end{gather*}
事实上，群胚叠的一个态射是等价，当且仅当它在由对象同构类预层生成的层上诱导同构（第二个公式），并在任意局部对象的自同构层上诱导同构（第一个公式）。

我们将给出全忠实性的两个证明（\hyperref[XVIII.lem1.2.4]{1.2.4} 和 \hyperref[XVIII.subsubsec1.2.5]{1.2.5}）\nde{对于本质满性，读者应当注意，证明将用到 \EGA 引言中解释的一般证明技巧。}。

\begin{enonce*}{引理 1.2.3}\etiquette[1.2.3]{XVIII.lem1.2.3}
设\pageoriginale $f:X\rightarrow S$ 是固有平坦态射。假设普遍地有 $f_{\ast}\Oo_{X}=\Oo_{S}$，且普遍地有 $R^{1}f_{\ast}\Oo_{X}=0$。则从 $X$ 到一个有限型 $S$-群概形 $G$ 的每个 $S$-态射都通过 $f$ 和 $G$ 的一个截面分解；更确切地说，有 $G\xrightarrow{\sim} f_{\ast}f^{\ast}G$。
\end{enonce*}

只需证明第一个断言。

由刚性引理（\textsc{Mumford} \cite[6.1. p.~115]{sga4.exp.XVIII.bib.4}；注意 \textsc{Mumford} 并未用到对 $S$ 所加的 Noether 假设），只需处理 $S$ 是代数闭域的谱的情形。先假定 $G$ 是阿贝尔概形，设 $G^{\ast}$ 为其对偶阿贝尔概形。由于 $G=G^{{\ast}{\ast}}$，给定 $\varphi:X\rightarrow G$ 等价于给定 $X\times G^{\ast}$ 上的一个可逆层 $\Ll$，它沿 $X\times \{e\}$ 被平凡化，且在从 $X\times G^{\ast}$ 到 $X$ 的投影 $\pr_{1}$ 的纤维上代数等价于零。若 $x:S\to X$ 是 $X$ 的一个有理点，则 $\Ll\sim \pr^{\ast}_{2}(x\times G^\ast)^{\ast}\Ll\otimes \Mm$，其中 $\Mm$ 沿 $x\times G^{\ast}$ 和 $X\times \{e\}$ 被平凡化。给定 $\Mm$ 等价于给定从 $G^{\ast}$ 到 $\sheafPic_{X/S}$ 的一个将 $e$ 映为 $e$ 的概形态射。由于 $H^{1}(X,\Oo_{X})=0$，有 $\sheafPic^0_{X/S}=0$\nde{这是由典范同构 $\Lie(\sheafPic_{X/S})=H^{1}(X,\Oo_{X})=0$ 得到的，而该同构来自将这两个群分别与 $\Pic(X\times_k \Spec (k[\epsilon]))\to \Pic(X)$ 的核比较（其中 $S=\Spec(k)$）。关于 Picard 概形的性质，更多细节见 \cite{sga4.exp.XVIII.bib.2bis}。}，$\Mm\sim\Oo$，而且 $\Ll\sim \pr^{\ast}_{2}\Ll'$，其中 $\Ll'$ 在 $G^{\ast}$ 上\sisi{等价}{代数等价}于零，因此它定义了一个 $g\in G(S)$，使得 $\varphi=g\circ f$。

在一般情形，$G$ 是一个阿贝尔概形 $A$ 由一个仿射群 $G_0$ 的扩张。前面的结果使我们归结为 $A=0$ 的情形，此时
$$
\Hom_{S}(\Spec (H^0(X,\Oo_{X})),G)\xrightarrow{\sim}
\Hom_{S}(X,G){\quoi}.
$$

\begin{enonce*}{引理 1.2.4}\etiquette[1.2.4]{XVIII.lem1.2.4}
在 \hyperref[XVIII.subsubsec1.2.1]{1.2.1} 的假设下，若 $G$ 是 $S$ 上有限表现的群概形，或者由 $S$ 上的一个拟凝聚层定义，则函子 \eqhyperref[XVIII.eq1.2.1.1]{1.2.1.1} 是全忠实的。
\end{enonce*}

\begin{itemize}
\item[A.] 我们证明：若 $K(K_{0},\varphi)$ 与 $K(K'_0,\varphi')$ 同构，则 $\varphi=\varphi'$。对于拟凝聚的 $G$，有 $\varphi=\varphi'=0$，无须证明。对于有限表现的 $G$，由（\SGA 3 IX 5.1），只需\pageoriginale{}在 $S$ 为代数闭域 $k$ 的谱时证明这一断言。设 $T$ 是 $G$ 的最大子环面，置 $H=G/T$。上同调正合列给出
$$
0\rightarrow H^{0}(\PP(\epsilon),T)\rightarrow
H^{0}(\PP(\epsilon),G) \xrightarrow{\alpha}
H^{0}(\PP(\epsilon),H)
 \xrightarrow{\delta}
H^{1}(\PP(\epsilon),T)\xrightarrow{\beta}H^{1}(\PP(\epsilon),G){\quoi}.
$$
由 \hyperref[XVIII.lem1.2.3]{1.2.3}，$H^0(\PP(\epsilon),G)\sim G(k)$，且 $H^0(\PP(\epsilon),H)\sim H(k)$，所以 $\alpha$ 为满射，$\delta$ 为零。若 $Y(T)=\Hom(\GG_{m},T)$，则 $T\sim G_{m}\otimes Y(T)$ 且 $H^{1}(\PP(\epsilon),T)\sim Y(T)$。所以，$\beta$ 的单射性表示由 $\Oo(1)$ 定义的典范箭头 $\Hom(\GG_{m},G)\simeq \Hom(\GG_{m},T)\rightarrow H^{1}(\PP(\epsilon),G)$ 的单射性，断言得证。

\item[B.] 证明函子 \eqhyperref[XVIII.eq1.2.1.1]{1.2.1.1} 全忠实是关于 $S$ 的局部问题。因此，可以假定 $\PP(\epsilon)/S$ 有一个截面 $s$\sisi{}{，并选定一个同构 $s^\ast\Oo(1)\sim \Oo$\nde{原版缺少的这一必要附加数据，已在本版中补入。}}。若 $K(K_{0},\varphi)$ 与 $K(K'_{0},\varphi')$ 同构，则由 A 有 $\varphi=\varphi'$，且 $K_{0}\sim s^{\ast}K(K_{0},\varphi)\sim s^{\ast}K(K'_{0},\varphi')\sim K'_{0}$。根据公式
\begin{equation*}
\Hom (K(K_{0},\varphi), K(K'_{0},\varphi'))\sim \Hom(K(K_{0}-K'_{0},\varphi-\varphi'),G_{\PP(\epsilon)}){\quoi},\tag{1.2.4.1}
\end{equation*}\etiquette[1.2.4.1]{XVIII.eq1.2.4.1}%
只需再证明
$$
\Hom_{\tors.}(G,G)\xrightarrow{\sim}
\Hom_{\tors.}(p^{\ast}G,p^{\ast}G){\quoi}.
$$
这由 \hyperref[XVIII.lem1.2.3]{1.2.3} 得到。
\end{itemize}

\setcounter{subsubsection}{4}
\subsubsection{}\etiquette[1.2.5]{XVIII.subsubsec1.2.5}

这里给出全忠实性的另一个证明，它更有 \sisi{Kunneth}{Künneth} 公式的风格。问题关于 $S$ 是局部的；因而由 \eqhyperref[XVIII.eq1.2.4.1]{1.2.4.1}，只需验证
$$
\Hom_{\PP(\epsilon)}(G,K(G,\varphi))
\begin{cases}
=∅ &\text{\rm 若}\quad \varphi\neq 0\\
\xleftarrow{\sim} G(S)\ &\text{\rm 若}\quad \varphi= 0{\quoi}.
\end{cases}
$$
这正是下面引理的内容\nde{这一转述值得作一些解释。$\PP(\epsilon)$ 上的线丛 $\Oo(1)$ 沿典范投影 $\VV(\epsilon)^{\ast}\to \PP(\epsilon)$ 拉回后，被典范地平凡化。因此，$G$-挠子 $K(G,\varphi)$ 拉回到 $\VV(\epsilon)^{\ast}$ 后，也被典范地平凡化。\emph{通过}这一等同，同构 $G\to K(G,\varphi)$ 因而诱导 $\VV(\epsilon)^{\ast}$ 上平凡 $G$-挠子的一个自同构，即一个 $S$-态射 $\VV(\epsilon)^{\ast}\to G$。为了从 $\PP(\epsilon)$ 下降而来，该态射必须满足类似所述条件的相容性，而反向结论由忠实平坦下降得到。严格说来，审阅者认为正确的条件应为 $\Psi(\lambda v)=\varphi(\lambda)^{-1}\Psi(v)$，但这没有影响，因为可以将引理应用于 $\Psi^{-1}$ 而非 $\Psi$。原版中的条件 $\Psi(\lambda v)=\varphi(\lambda)\Psi(v)$ 也给出正确的陈述（因为可以将 $\varphi$ 换成 $\varphi^{-1}$），但此处作了修改，以便不必修改证明。}：

\begin{enonce*}{引理 1.2.6}\etiquette[1.2.6]{XVIII.lem1.2.6}
设 $\epsilon$ 是 $S$ 上处处秩 $\geq 2$ 的局部自由模，$G$ 是 $S$ 上有限表现的群概形，$\varphi:\GG_{m}\rightarrow G$ 是群同态，$\VV(\epsilon)$ 是由 $\epsilon$ 定义的向量丛，$\VV(\epsilon)^{\ast}$\pageoriginale{}为相应的去零截面向量丛，$\Psi:\VV(\epsilon)^{\ast}\rightarrow G$ 是一个 $S$-态射，满足以下从 $\GG_{m}\times \VV(\epsilon)^{\ast}$ 到 $G$ 的 $S$-态射之间的恒等式：
$$
\Psi(\lambda v)=\sisi{\varphi (\lambda)\Psi(v)}{\Psi(v)\varphi(\lambda)}{\quoi}.
$$
则 $\varphi=e$，且 $\Psi$ 通过从 $\PP(\epsilon)$ 到 $S$ 的投影 $p$ 分解。
\end{enonce*}

只需证明第二个断言。立即归结为 $S$ 是 Noether 概形的情形，再由（\SGA 3 IX 5.1）及应用于 $\PP(\epsilon)$ 的刚性引理（\textsc{Mumford} \cite[6.1 p.~115]{sga4.exp.XVIII.bib.4}），归结为 $S$ 是域的谱的情形。设 $n$ 是一个整数，$W_{n}$、$G_{n}$ 分别为 $\VV(\epsilon)$、$G$ 中零截面的第 $n^{\text{\rm th}}$ 个无穷小邻域，$\mathfrak{g}_{n}$ 是 $G_{n}$ 的仿射代数，$P_{n}(\Psi):\VV(\epsilon)^{\ast}\rightarrow \Hhom_{S}(W_{n},G)$ 是 $\Psi$ 的 $n^{\text{\rm th}}$ 阶主部\nde{更具体地，$P_n(\Psi)$ 是通过伴随对应于如下态射的态射：$\VV(\epsilon)^{\ast}\times_S W_n\to G$（我们也将其记作 $P_n$），它是将 $\Psi$ 与 $\VV(\epsilon)$ 上的加法在这些子概形上诱导的态射 $\VV(\epsilon)^{\ast}\times_{S} W_n\to \VV(\epsilon)^\ast$ 复合得到的。用这些 $S$-概形的点函子来表述，可更具体地解释为：若 $v$ 是 $\VV(\epsilon)^{\ast}$ 的一个截面，$t$ 是 $W_n$ 的一个截面，则 $P_n(\Psi)(v,t)=\Psi(v+t)$。}。
$\Hhom_{S}(W_{n},G)$ 的每个截面，通过按零截面的像的逆作右平移，都定义一个从 $W_{n}$ 到 $G_{n}$ 的将 $0$ 映到 $e$ 的态射；由此定义
$$
\Psi_{1}=P_{n}(\Psi).\Psi^{-1}:\VV(\epsilon)^{\ast}\longrightarrow
\Hhom_{S}(W_{n},G_{n}){\quoi}.
$$
假设蕴含 $\Psi_{1}$ 与位似可交换，其中位似在右侧的作用，是由其在 $W_{n}\subset \VV(\epsilon)$ 上的作用通过结构搬运得到的\nde{用上面的记号，$\Psi_1$ 可由关系 $\Psi_1(v,t).\Psi(v)=\Psi(v+t)$ 刻画。立即得到，若 $\lambda$ 是 $\GG_{m}$ 的一个截面，则 $\Psi_1(\lambda v,\lambda t)=\Psi_1(v,t)$；只要 $\GG_{m}$ 在 $\Hhom_S(W_n,G_n)$ 上正确地作用，即 $\GG_{m}$ 的截面 $\lambda$ 通过在右侧复合 $W_n$ 上乘以 $\lambda$ 的\emph{逆}的映射而作用，这就是所断言的与位似的可交换性。}。$W_{n}$ 的仿射代数是 $\sum^{n}_{i=0}\Sym^{i}(\epsilon)$。因此，由齐次性，$\Psi_{1}$ 由以下 $\Oo_{\PP(\epsilon)}$-模在 $\PP(\epsilon)$ 上的一个截面定义：
$$
\sum^{n}_{i=0}p^{\ast}\Hhom(\mathfrak{g}_{n},\Sym^{i}(\epsilon))(-i){\quoi}.\nde{为了详细验证这一断言，可以先定义该截面在 $\VV(\epsilon)^{\ast}$ 上的拉回；于是由定义，得到一个平凡丛的截面 $s=(s_0,\dots,s_n)$，其中该平凡丛是将结构层与向量空间 $\sum^{n}_{i=0}\Hhom(\mathfrak{g}_{n},\Sym^{i}(\epsilon))$ 作张量积得到的。它满足关系 $s_i(v)=\lambda^i s_i(\lambda v)$。这确实是给出 $\PP(\epsilon)$ 上所需截面的下降数据。也可以论证，这些截面 $s_i$ 延拓到整个 $\VV(\epsilon)$（因为原点是 $\VV(\epsilon)$ 中余维 $\geq 2$ 的闭子集）。由 $\GG_m$ 在 $\AAA^1$ 中的稠密性，延拓后的截面 $s_i$ 对 $\lambda\in\AAA^1$ 满足关系 $s_i(v)=\lambda^i s_i(\lambda v)$。取 $\lambda=0$，得到当 $i>0$ 时 $s_i=0$。}
$$
这个模中指标 $i>0$ 的各分量除了零截面以外没有其他截面，由此推出
$$
\Psi_{1}=0.
$$
由于这对每个 $n$ 都成立，它蕴含 $\Psi$ 通过 $G$ 的一个截面分解。

第二个证明可以推广，用以证明下面的 \hyperref[XVIII.lem1.2.4]{1.2.4} 的解析变体。当 $G$ 仿射时，\hyperref[XVIII.thm1.2.2]{1.2.2} 的证明可以不依赖于这一变体。

\setcounter{subsubsection}{6}
\subsubsection{}\etiquette[1.2.7]{XVIII.subsubsec1.2.7}

设\pageoriginale $L$ 是一个完备离散赋值环 $\Oo$ 的分式域，$U_{m}$ 是 $\GG_{m}$ 的形式完备化的刚性解析“一般纤维”群，$\Oo(1)^{\wedge}$ 是在 $\PP^{r}_{L}$ 上通过对 $\Oo(1)$ 作形式完备化得到的 $U_{m}$-挠子。若 $G$ 是 $L$ 上的刚性解析群，则函子 $(K_{0},\varphi)\mapsto p^{\ast}K_{0}+\varphi\Oo(1)^{\wedge}$（\cf \hyperref[XVIII.subsubsec1.2.1]{1.2.1}）是从以下范畴到 $\PP^{r}_{L}$ 上的 $p^{\ast}G$-挠子范畴的一个全忠实函子：该范畴的对象为对 $(K_{0},\varphi)$，其中第一个分量为 $\Spec(L)$ 上的 $G$-挠子，$\varphi\in \Hom(U_m,G)$\acorriger{读者可在 \cite{sga4.exp.XVIII.bib.5bis} 中找到刚性解析几何语言的入门介绍。在这一情形中，当 $n<0$ 时层 $\Oo(n)$ 不存在非平凡截面，例如可由刚性版本的 GAGA 得到；见 \cite{sga4.exp.XVIII.bib.3ter}。}。

\begin{enonce*}{引理 1.2.8}\etiquette[1.2.8]{XVIII.lem1.2.8}
设 $k$ 是特征指数为 $p$ 的域，$f:X\rightarrow \Spec(k)$ 是 $k$ 上的固有概形，满足 $H^0(X,\Oo_{X})=k$ 和 $\sheafPic^{\tau}_{X}=0$\nde{\hyperref[XVIII.def1.6.4]{1.6.4} 中给出了更一般情形下 $\sheafPic^{\tau}$ 的定义。}。对于 $k$ 上每个不含环面的有限型交换仿射群概形 $G$，有 $R^{1}f_{\ast}(f^{\ast}G)=0$。
\end{enonce*}


只需证明，在每个基变换 $c:S\rightarrow \Spec(k)$ 之后，从 $S$ 上的 $G$-挠子范畴到 $X_{S}$ 上的 $G$-挠子范畴的函子 $f^{\ast}$ 都是等价。该问题关于 $\Spec (k)$ 是局部的，所以为了解决它，可以归结为如下情形：$X$ 有有理点，$G$ 是各群 $\GG_{a}$、$\alpha_{p}$、$𝛍_{p}$、$\ZZ/p$ 和 $\ZZ/\ell_{i}$ 的逐次扩张，其中 $\ell_{i}$ 与 $p$ 互素，且 $\ZZ/\ell_{i}\sim 𝛍_{\ell_i}$。

我们知道 $\Lie(\sheafPic_{X})=H^{1}(X,\Oo_{X})=0$。由 \hyperref[XVIII.lem1.2.3]{1.2.3} 及截面的存在性，所讨论的函子 $f^{\ast}$ 是全忠实的；因此，只需证明 $R^{1}f_{\ast}(f^{\ast}G)=0$，或者由分步化简，只需证明
$$
R^{1}f_{\ast}\GG_{a}=R^{1}f_{\ast}\alpha_{p}=R^{1}f_{\ast}𝛍_{p}=R^{1}f_{\ast}\ZZ/p=R^{1}f_{\ast}\ZZ/\ell_{i}=0{\quoi}.
$$

对于 $\GG_{a}$，这由 $H^{1}(X,\Oo_{X})=0$ 得到。$\alpha_{p}$（\resp $\ZZ/p$）的情形则通过 \hyperref[XVIII.lem1.2.3]{1.2.3}，由正合列
\begin{align*}
& 0\longrightarrow \alpha_{p}\longrightarrow
  \GG_{a}\xrightarrow{x^{p}} \GG_{a}\longrightarrow 0\\
& 0\longrightarrow \ZZ/p\longrightarrow
  \GG_{a}\xrightarrow{x-x^{p}}\GG_{a}\longrightarrow
  0\qquad (p\neq 1)){\quoi}.\tag*{(\resp}
\end{align*}
定义的上同调长正合列推出。$\ZZ/\ell_{i}$（\resp $𝛍_{p}$）的情形，则由同构 $\ZZ/\ell_{i}\sim 𝛍_{\ell_{i}}$、\pageoriginale \hyperref[XVIII.lem1.2.3]{1.2.3} 以及由 Kummer 正合列
$$
0\longrightarrow 𝛍_{n}\longrightarrow \GG_{m}\longrightarrow
\GG_{m}\longrightarrow 0{\quoi}.
$$
推出的长正合列得到。

\begin{enonce*}{引理 1.2.9}\etiquette[1.2.9]{XVIII.lem1.2.9}
设 $S$ 是局部 Artin 概形，$\bar{s}$ 是 $S$ 的一个几何点，$f:X\rightarrow S$ 是固有平坦态射，$G$ 是 $S$ 上有限型的平坦交换群概形。假设 $f$ 光滑，或者 $G$ 仿射。若 $H^0(X_{\bar{s}},\Oo)=k(\bar{s})$，$\sheafPic^{\tau}_{X_{\bar{s}}}=0$，且 $G_{\bar{s}}$ 不含子群 $\GG_{m}$，则从 $S$ 上的 $G$-挠子范畴到 $X$ 上的 $f^{\ast}G$-挠子范畴的函子 $f^{\ast}$ 是等价。
\end{enonce*}

\begin{itemize}
\item[a.] 当把有限平坦态射 $u:S'\rightarrow S$ 作为局部化态射时，所讨论的问题关于 $S$ 是局部的。这使我们可以假定 $f$ 有一个截面 $x$\nde{在作这样的基变换后截面的存在性，可以由 \hyperref[XVIII.lem2.3]{2.3} 和 \hyperref[XVIII.lem2.4]{2.4} 推出。}。于是，由 \hyperref[XVIII.lem1.2.3]{1.2.3} 及 $x$ 的存在性，函子 $f^{\ast}$ 全忠实，而只需再证明，每个沿 $x$ 被平凡化的挠子都是平凡的。

\item[b.] 假定 $S$ 是域的谱。若 $G$ 仿射，则断言由 \hyperref[XVIII.lem1.2.8]{1.2.8} 得到。若 $X$ 光滑，而 $G$ 是一个阿贝尔簇 $A$ 由一个仿射群 $G_{0}$ 的扩张，以 $G^{n}$ 表示 ${}_{n}A$ 在 $G$ 中的逆像。群 $H^{1}(X,A)$ 是挠群（\textsc{Raynaud} \cite[XIII 2.6]{sga4.exp.XVIII.bib.5}），因而是各 $H^{1}(X,{}_{n}A)$ 的像的并。所以群 $H^{1}(X,G)$ 是各 $H^{1}(X,G^{n})$ 的像的并。由于群 $G^{n}$ 仿射，结论仍由 \hyperref[XVIII.lem1.2.8]{1.2.8} 得到。

\item[c.] 对 $S$ 的长度作归纳，证明一般情形（假定截面 $x$ 存在）。设 $i:S_{0}\hookrightarrow S$ 是由一个平方为零、长度为 1 的理想 $I$ 定义的闭子概形。
\end{itemize}

由归纳假设，使 $x^{\ast}P$ 平凡的 $X$ 上的 $G$-挠子 $P$，都是 $X\times S_{0}$ 上平凡 $G$-挠子 $P_{0}$ 的形变。若 $G$ 光滑，且 $L$ 是 $G_{s}$ 的 $\Lie$ 代数，则这些形变由
$$
H^{1}(X_{s},f^{\ast}(L\otimes I))=0{\quoi},
$$
分类，从而\pageoriginale{}在这种情形下证明了 \hyperref[XVIII.lem1.2.9]{1.2.9}。

对于 $S$ 上平坦的 $G$，这一论证可借助 $G_{s}$ 的 \emph{Lie 复形} $L$ 推广。
% SGA-SRC-001128 — reconstructed by mathematical inference; scan inconclusive
若 $\mathrm{L}\Lambda(G_{s}/k(s))$ 表示 $G_{s}$ 的相对余切复形，则此复形 $L\in D^{b}(k(s))$ 是 $\mathrm{L}e^{\ast}(\mathrm{L}\Lambda(G_{s}/k(s)))$ 的对偶。有 $H^{i}(L)=0$，当 $i\neq 0$、$1$ 时\sisi{}{（\textsc{Illusie} \cite[VII~2.4.1]{sga4.exp.XVIII.bib.3}）}。

$S_{0}$ 上（\resp $X_{0}$ 上）的平凡 $G$-挠子在 $S$ 上（\resp $X$ 上）的形变，由 $H^{1}(L\otimes I)$（\resp $H^{1}(X_{s},f^{\ast}(L\otimes I))$（\textsc{Illusie} \cite[VII~2.4.4~(ii)]{sga4.exp.XVIII.bib.3}）分类。
按假设，有 $H^{1}(X_{s},ℋ^0(f^{\ast}(L\otimes I)))=0$（因为 $H^{1}(X_{s},\Oo)=0$），且 $H^{1}(L\otimes I)\sim H^0(X_{s}, ℋ^{1}(f^{\ast}L\otimes I))$）。于是得到
$$
H^{1}(L\otimes I)\xrightarrow{\sim} H^{1}(X_{s},f^{\ast}(L\otimes
I)){\quoi},
$$
这意味着，$X_{0}$ 上的平凡 $G$-挠子在 $X$ 上的每个形变，都是 $S_{0}$ 上的平凡 $G$-挠子在 $S$ 上的唯一一个形变的逆像。因此，若此形变沿一个截面平凡，则它就是平凡的，\hyperref[XVIII.lem1.2.9]{1.2.9} 得证。

\setcounter{subsubsection}{9}
\subsubsection{}\etiquette[1.2.10]{XVIII.subsubsec1.2.10}
在 $S$ 为 Artin 概形时证明 \hyperref[XVIII.thm1.2.2]{1.2.2}。设 $T$ 为 $G$ 的最大子环面；由于已经知道全忠实性（\hyperref[XVIII.lem1.2.4]{1.2.4}），问题关于 $S$ 对平展拓扑是局部的：可以假定 $T$ 可对角化，$T\sim \GG^{n}_{m}$。由 \hyperref[XVIII.lem1.2.3]{1.2.3} 及 \hyperref[XVIII.lem1.2.9]{1.2.9}，上同调正合列给出
$$
R^{1}f_{\ast}T\xrightarrow{\sim} R^{1}f_{\ast}G{\quoi},
$$
于是断言成立，因为它在 $G=T$ 时是平凡的，且 $\Hom(\GG_{m},T)\xrightarrow{\sim}\Hom(\GG_{m},G)$（\cf \hyperref[XVIII.par1.2.2.1]{1.2.2.1}）。

\subsubsection{}\etiquette[1.2.11]{XVIII.subsubsec1.2.11}
在 $S$ 为完备 Noether 局部概形时证明 \hyperref[XVIII.thm1.2.2]{1.2.2}。设 $s$ 是~$p$ 的一个截面\nde{如上面一样，还必须选定一个同构 $s^\ast\Oo(1)\sim \Oo$。}。

\noindent
若 $K$ 是 $\PP(\epsilon)$ 上的一个 $G$-挠子，则由 \hyperref[XVIII.subsubsec1.2.10]{1.2.10}，存在唯一一对 $(\widehat{\varphi},\widehat{\Psi})$，其中 $\widehat{\varphi}:\GG_{m}\rightarrow G$ 为形式同态，$\widehat{\Psi}:\widehat{p^{\ast}s^{\ast}K}+\widehat{\varphi}\widehat{\Oo(1)}\xrightarrow{\sim}\widehat{K}$ 为形式同构\pageoriginale\nde{必须要求这一同构在应用 $s^\ast$ 后成为显然的同构。}\nde{这里的意思是，若 $A$ 是 $S$ 的环，$I$ 为其极大理想，$S_n=\Spec(A/I^n)$，则一个形式态射由每个 $n$ 对应的 $S_n$ 上对象之间的一族相容态射组成。代数性问题就是：这一相容系统是否由 $S$ 上类似的数据取逆像而得。}。若 $G$ 仿射，则 $\widehat{\varphi}$ 可代数化（\SGA 3 IX 7.1）。暂且承认 $\widehat{\varphi}$ 总能代数化为 $\varphi:\GG_{m}\rightarrow G$，这一点将在（\hyperref[XVIII.subsubsec1.2.13]{1.2.13}）中证明。

若 $K_{0}=s^{\ast}K$，且 $K_{1}=p^{\ast}K_{0}+\varphi\Oo(1)$，则有形式同构 $\widehat{\Psi}:\widehat{K_{1}}=\widehat{K(K_{0},\varphi)}\xrightarrow{\sim}\widehat{K}$，即 $K-K_{1}$ 的一个形式平凡化。由形式 GAGA，这个截面是代数的，\hyperref[XVIII.subsubsec1.2.11]{1.2.11} 得证。更确切地说，若 $G$ 拟仿射，则 $K-K_{1}$ 可表，并使用 \EGA III 5.1.4。在一般情形，对代数空间 $K-K_{1}$ 以同样的方式进行论证。

\subsubsection{$\widehat{\varphi}$ 的代数性：在离散赋值环的谱上的预备验证}\etiquette[1.2.12]{XVIII.subsubsec1.2.12}

假设 $S$ 是一个完备离散赋值环的谱（trait），$S=\Spec(V)$，设 $L$ 为~$V$ 的分式域。
\begin{itemize}
\item[a)] 在 $L$ 上，存在唯一三元组 $(K_{0L},\varphi_{L},\Psi_{L})$：
$$
\Psi_{L}:K(K_{0L},\varphi_{L})\xrightarrow{\sim} K|\Spec (L){\quoi}.
$$

\item[b)] 由（\hyperref[XVIII.subsubsec1.2.10]{1.2.10}）（\hyperref[XVIII.subsubsec1.2.11]{1.2.11}），采用 \hyperref[XVIII.subsubsec1.2.7]{1.2.7} 的记号，有刚性解析映射 $\widehat{\varphi_{L}}:U_{m}\rightarrow \widehat{G}$、一个 $\widehat{G}$-挠子 $\widehat{K_{0}}$，以及一个刚性解析同构
$$
\widehat{\Psi}:p^{\ast}\widehat{K_0}+\widehat{\varphi_{L}}
\widehat{\Oo(1)}\xrightarrow{\sim}
\widehat{K}{\quoi}.
$$
由唯一性断言 \hyperref[XVIII.subsubsec1.2.7]{1.2.7}，$\widehat{\varphi_{L}}$ 与从 $\varphi_{L}$ 得到的映射等同。因此，\hyperref[XVIII.subsubsec1.2.11]{1.2.11} 的形式态射 $\widehat{\varphi}$ 在 $𝛍_{n}$ 上诱导一个映射，它与从代数映射 $\varphi_{L}:\GG_{m}\rightarrow G$ 得到的映射一致。
\end{itemize}

\subsubsection{$\widehat{\varphi}$ 的代数性}\etiquette[1.2.13]{XVIII.subsubsec1.2.13}

对在 $S$ 上可逆的 $\ell$，\hyperref[XVIII.subsubsec1.2.11]{1.2.11} 中的形式映射 $\widehat{\varphi}:\GG_{m}\rightarrow G$ 诱导映射
$$
\varphi_{\ell}:T_{\ell}(\widehat{\GG_{m}})=T_{\ell}(\GG_{m})\longrightarrow
T_{\ell}(G){\quoi}.
$$
由 \hyperref[XVIII.subsubsec1.2.12]{(1.2.12)} 及（\EGA II 7.1.9），对于 $S$ 的每一点 $t$，都存在一个\pageoriginale{}代数态射 $\varphi_{t}:\GG_{m_{t}}\rightarrow G_{t}$，使得 $T_{\ell}(\varphi_{t})=T_{\ell}(\widehat{\varphi})_{t}$。因此，由（\SGA 3 XV 4.1 bis），$\widehat{\varphi}$ 可代数化。

\subsubsection{}\etiquette[1.2.14]{XVIII.subsubsec1.2.14}
现在证明 \hyperref[XVIII.thm1.2.2]{1.2.2}。通过取极限，归结为 $S$ 是 Noether 概形的情形，进而归结为 $S$ 是仿射坐标环为 $A$ 的 Noether 局部概形的情形。设 $s$ 为 $p$ 的一个截面。设 $K$ 是 $\PP(\epsilon)$ 上的一个 $G$-挠子。在 $S'=\Spec(\widehat{A})$ 上，存在唯一一个同态 $\varphi':\GG_{mS'}\rightarrow G_{S'}$，使得 $K$ 在 $\fprod{\PP(\epsilon)}{S'}{S}$ 上的逆像 $K'$ 形如 $K(K'_{0},\varphi')$。这一同态 $\varphi'$ 下降为 $\varphi:\GG_{m}\rightarrow G$。为了证明函子 \hyperref[XVIII.subsubsec1.2.1]{1.2.1} 的本质满性，可以将 $K$ 换成 $K-\varphi\Oo(1)$，即假定 $\varphi=0$。于是在 $S'$ 上，存在唯一同构 $\Psi:{p'}^{\ast}s^{\ast}K'\xrightarrow{\sim}K'$，它在截面 $s$ 上诱导恒等。这一态射下降到 $S$（fppf 挠子态射的 fpqc 下降），从而证明了 \hyperref[XVIII.thm1.2.2]{1.2.2}。

\begin{enonce*}{变体 1.2.15}\etiquette[1.2.15]{XVIII.var1.2.15}
\hyperref[XVIII.thm1.2.2]{1.2.2} 的结论在下列情形仍成立。
\begin{enumerate}
\renewcommand{\theenumi}{\roman{enumi}}
\renewcommand{\labelenumi}{(\theenumi)}
\item $G$ 由 $S$ 上的一个拟凝聚模定义。

\item $G$ 是 $S$ 的小平展景上某个挠层的逆像。
\end{enumerate}
\end{enonce*}

在情形 (i) 中，考虑到 $G$ 的 fppf 上同调与它的 Zariski 拓扑上同调一致，箭头（\hyperref[XVIII.par1.2.2.1]{1.2.2.1}）为同构是一个标准事实。

在情形 (ii) 中，由固有态射的基变换定理及代数闭域 $k$ 上 $\PP^{r}_{k}$ 的单连通性，箭头（\hyperref[XVIII.par1.2.2.1]{1.2.2.1}）为同构。

% END native input: chapters/exp-18-part1-zh.tex
% BEGIN native input: chapters/exp-18-part2-zh.tex
\subsection{挠子的积分}\etiquette[1.3]{XVIII.subsec1.3}

\subsubsection{}\etiquette[1.3.1]{XVIII.subsubsec1.3.1} 设 $S$ 为概形，$G$ 为 $S$ 的大 fppf 景上的阿贝尔层，满足以下三个条件。

\paragraph{}\etiquette[1.3.1.1]{XVIII.par1.3.1.1} \begin{itemize} \item[(i)] \emph{对\pageoriginale 每个满足 $r\geq 1$ 的射影空间 $p:\PP^r_{S}\rightarrow S$，在 $S$ 的大 fppf 景上}（\hyperrefext[XVII][XVII.subsec0.10]{0.10}）\emph{有} \begin{gather*}
G\xrightarrow{\sim} p_{\ast}p^{\ast}G\quad\mbox{且}\\
\Hhom(\GG_{m},G)\xrightarrow{\sim} R^{1}p_{\ast}p^{\ast}G
\end{gather*}

\item[(ii)] \emph{对每个 $n$ 及每个光滑仿射 $S$-态射 $f:X\rightarrow T$，若 $K_{1}$ 和 $K_{2}$ 为 $\Sym^{n}_{T}(X)$ 上的两个 $G$-挠子，则集合 $\Hom(K_{1},K_{2})$ 可与从 $K_{1}$ 在 $(X/T)^{n}$ 上的逆像到 $K_{2}$ 在 $(X/T)^{n}$ 上的逆像的对称态射所成的集合等同。}

此条件等价于以下两个条件同时成立：

\item[${\rm (ii')}$] $\Hom_{S}(\Sym^{n}_{T}(X),G)\xrightarrow{\sim}\Hom_{S}((X/T)^{n},
G)^{\sigma_n}$，以及

\item[${\rm (ii'')}$] 若 $\Sym^{n}_{T}(X)$ 上的挠子 $K$ 在 $(X/T)^{n}$ 上的逆像有对称截面，则 $K$ 平凡。

\item[(iii)] $G$ \emph{或者满足条件}（\hyperrefext[XVII][XVII.par6.3.21.1]{6.3.21.1}），\emph{或者同时满足条件}（\ref{XVII.par6.3.24.1}）\emph{及}（\ref{XVII.par6.3.24.2}）。 \end{itemize}

于是 \hyperref[XVIII.thm1.2.2]{1.2.2} 的结论成立（\cf \hyperref[XVIII.par1.2.2.1]{1.2.2.1}），\hyperrefext[XVII][XVII.subsubsec6.3.26]{6.3.26} 的公式也成立。

条件 (i) 已在 \hyperref[XVIII.thm1.2.2]{1.2.2} 及 \hyperref[XVIII.var1.2.15]{1.2.15} 中讨论。当 $G$ 可表示且为有限表现，或为 $S$ 的小平展景上的层的逆像，或为平坦拟凝聚层时，条件 (ii) 均成立。因此，以下各情形均满足条件 \hyperref[XVIII.par1.3.1.1]{1.3.1.1}。

\paragraph{}\etiquette[1.3.1.2]{XVIII.par1.3.1.2} \begin{enumerate} \renewcommand{\theenumi}{\alph{enumi}} \renewcommand{\labelenumi}{(\theenumi)} \item $G$ 光滑且为有限表现（\hyperref[XVIII.thm1.2.2]{1.2.2} 及 \hyperrefext[XVII][XVII.par6.3.21.1]{6.3.21.1}）。

\item $G$ 为 $S$ 的小平展景上的挠层的逆像（\hyperref[XVIII.var1.2.15]{1.2.15} (ii) 及 \hyperrefext[XVII][XVII.par6.3.21.1]{6.3.21.1}）。

\item $G$ 仿射，且为光滑群之间的满态射的核（\hyperref[XVIII.thm1.2.2]{1.2.2} 及 XVII \hyperref[XVII.eq6.3.3.1]{6.3.3.1}）。

\item $G$ 由 $S$ 上的平坦拟凝聚层定义（\hyperref[XVIII.var1.2.15]{1.2.15} (i) 及 \hyperrefext[XVII][XVII.par6.3.21.1]{6.3.21.1}）。 \end{enumerate}

\subsubsection{}\etiquette[1.3.2]{XVIII.subsubsec1.3.2} 设\pageoriginale $f:X\rightarrow S$ 为概形 $S$ 上的光滑射影曲线，纤维非空且几何连通。设 $\Ll$ 为 $X$ 上的可逆层，它在以下意义下“充分”相对丰沛。

\paragraph{}\etiquette[1.3.2.1]{XVIII.par1.3.2.1} 对 $S$ 的每个几何点 $s$，若 $X_{s}$ 的亏格为 $g$，则 \begin{gather*}
\deg_{X_{s}}(\Ll)>2g-2,\\
\deg_{X_{s}}(\Ll)\geq 1\quad\mbox{\rm 若}\quad g=0,\quaet
\deg_{X_{s}}(\Ll)\geq 2\quad\mbox{\rm 若}\quad g=1
\end{gather*}

此时有 $R^{1}f_{\ast}\Ll=0$，而 $f_{\ast}\Ll$ 局部自由，其构造与每个基变换可交换，且处处的秩 $\geq 2$。若 $f_{\ast}\Ll$ 的一个截面 $s$ 在任何点都不为零，将 $s$ 看作 $X$ 上的 $\Ll$ 的截面时，其零点概形 $D_{s}$ 是一个相对除子（\EGA IV 21.15.2）；由于 $X/S$ 的相对维数为 1，它在 $S$ 上有限局部自由。有典范同构 $$
s:\Oo(D_{s})\xrightarrow{\sim}\Ll{\quoi}.
$$ 除子 $D_{s}$ 对 $s$ 的依赖仅差乘以 $\Oo^{\ast}_{S}$ 的一个截面。因此，上述构造在由 $X$ 经基变换 $$
p:\PP(f_{\ast}(\Ll)^{\vee})\longrightarrow S{\quoi}.
$$ 所得的概形上给出一个相对除子 $Z_{\Ll}$。

设 $K$ 为 $G_{X}$ 下的挠子，$K|Z_{\Ll}$ 为它在 $Z_{\Ll}$ 上的逆像。此挠子从 $Z_{\Ll}$ 到 $\PP(f_{\ast}(\Ll)^{\vee})$ 的迹（\hyperrefext[XVII][XVII.subsubsec6.3.26]{6.3.26}），是 $\PP(f_{\ast}(\Ll)^{\vee})$ 上的一个挠子 $K_{\Ll}$： $$
K_{\Ll}=\Tr_{Z_{\Ll}/\PP}(K|Z_{\Ll}){\quoi}.
$$

由 \hyperref[XVIII.thm1.2.2]{1.2.2}，存在一个在唯一同构意义下唯一的三元组，由 $S$ 上的一个 $G$-挠子 $\langle \Ll$, $K$]、一个态射 $\varphi_{\Ll,K}:\GG_{m}\rightarrow G$ 及一个同构 $$
K_{\Ll}\simeq
p^{\ast}\langle \Ll,K]+\varphi_{\Ll,K}\Oo(1){\quoi}.
$$ 组成。

\subsubsection{}\etiquette[1.3.3]{XVIII.subsubsec1.3.3} 若\pageoriginale 将挠子的合成用加法记号表示，则有典范同构及等式 \begin{align*}
& \langle \Ll,K_{1}+K_{2}] \simeq
\langle \Ll,K_{1}]+\langle \Ll,K_{2}]\tag{1.3.3.1}\\
& \varphi_{\Ll, K_{1}+K_{2}}=\varphi_{\Ll,K_{1}}+\varphi_{\Ll,K_{2}}{\quoi}.\tag{1.3.3.2}
\end{align*}\etiquette[1.3.3.1]{XVIII.eq1.3.3.1} \etiquette[1.3.3.2]{XVIII.eq1.3.3.2}

设 $\Ll_{1}$ 和 $\Ll_{2}$ 为 $X$ 上的两个满足（\hyperref[XVIII.par1.3.2.1]{1.3.2.1}）的可逆层，并设 $\Ll=\Ll_{1}\otimes \Ll_{2}$。截面的张量积定义一个态射 $$
\xymatrix@R=1.5cm@C=.5cm{
 & **[r] \PP
  (f_{\ast}(\Ll_1)^{\vee}) \ar[ld]^{\pr_1}& \times_S \ar[dd]^{p_{12}}& **[l]\PP
  (f_{\ast}(\Ll_2)^{\vee})\ar[rr]^{\pi}\ar[rd]^{\pr_2} &     &
 \PP(f_{\ast}(\Ll)^{\vee}) \ar@{-}[dd]^{p}\\
**[r]\PP(f_*(\Ll_1)^{\vee})\ar[rrd]_{p_1}&          &                &                 &  **[l]\PP(f_{\ast}(\Ll_2)^{\vee}) \ar[lld]^{p_2}  &\\
 &          &  S     \ar@{=}[rrr]        &                 &      & S.
}
$$ 有 $\pi^{\ast}Z_{\Ll}\simeq Z_{\Ll_{1}}+Z_{\Ll_{2}}$，从而得到一个

同构（\hyperrefext[XVII][XVII.const6.3.27.1]{6.3.27.1}） \begin{equation*}
\pi^{\ast}K_{\Ll}=\pr^{\ast}_{1}K_{\Ll_{1}}+\pr^{\ast}_{2}K_{\Ll_{2}}{\quoi}.\tag{1.3.3.3}
\end{equation*}\etiquette[1.3.3.3]{XVIII.eq1.3.3.3}

有典范同构 $$
\pi^{\ast}\Oo(1)\simeq \pr^{\ast}_{1}\Oo(1)\otimes
\pr^{\ast}_{2}\Oo(1){\quoi},
$$ 因而可将 \eqhyperref[XVIII.eq1.3.3.3]{1.3.3.3} 改写为 \begin{gather*}
\pi^{\ast}(p^{\ast}\langle \Ll,K]+\varphi_{\Ll,K}\Oo(1))=
p^{\ast}_{12}\langle \Ll,K]+\pr^{\ast}_{1}\varphi_{\Ll,K}\Oo(1)+\pr^{\ast}_{2}\varphi_{\Ll,K}\Oo(1) \\
= p^{\ast}_{12}\langle \Ll_{1},K]+p^{\ast}_{12}\langle \Ll_{2},K]+\pr^{\ast}_{1}\varphi_{\Ll_{1},K}\Oo(1)+\pr^{\ast}_{2}\varphi_{\Ll_{2},K}\Oo(1){\quoi}.
\end{gather*}

两次应用 \hyperref[XVIII.thm1.2.2]{1.2.2}，得到一个典范同构 \begin{gather*}
\langle \Ll_{1}\otimes \Ll_{2}, K]=\langle \Ll_{1},K]+\langle \Ll_{2},K]\tag{1.3.3.4}
\end{gather*}\etiquette[1.3.3.4]{XVIII.eq1.3.3.4} 及\pageoriginale 一个恒等式 $$
\varphi_{\Ll_{1}\otimes\Ll_{2},K}=\varphi_{\Ll_{1},K}=\varphi_{\Ll_{2},K}{\quoi}.
$$

由此可见 $\varphi_{\Ll,K}$ 仅依赖于 $K$，于是置 \begin{equation*}
\deg(K)=\varphi_{\Ll,K}{\quoi}.\tag{1.3.3.5}
\end{equation*}\etiquette[1.3.3.5]{XVIII.eq1.3.3.5} 因此， \begin{align*}
& K_{\Ll}\simeq
p^{\ast}\langle \Ll,K]+\deg(K)\Oo(1){\quoi}.\tag{1.3.3.6}\\
& \deg(K_{1}+K_{2})=\deg(K_{1})+\deg(K_{2}){\quoi}.\tag{1.3.3.7}
\end{align*}\etiquette[1.3.3.6]{XVIII.eq1.3.3.6} \etiquette[1.3.3.7]{XVIII.eq1.3.3.7}

设 $\mathcal{U}$ 为一个可逆层，它局部为 $f_{\ast}\Ll$ 的直和因子，定义 $\PP((f_{\ast}\Ll)^{\vee})$ 的一个截面 $u$。设 $D_{u}$ 为 $\mathcal{U}$ 的局部可逆截面定义的相对除子。 $$
u^{\ast}\Oo(1)\simeq \mathcal{U}^{\vee}
$$ 因此，由 \eqhyperref[XVIII.eq1.3.3.6]{1.3.3.6} 得到 \begin{equation*}
u^{\ast}K_{\Ll}=\Tr_{D_u/S}(K)\simeq
\langle \Ll,K]-\deg(K)\mathcal{U}{\quoi}\nde{记号 $\Tr_{D_u/S}(K)$ 将在下文 \eqhyperref[XVIII.eq1.3.5.0]{1.3.5.0} 中引入。}.\tag{1.3.3.8}
\end{equation*}\etiquette[1.3.3.8]{XVIII.eq1.3.3.8}

以下相容性的验证留给读者：

\begin{enonce*}{命题 1.3.4}\etiquette[1.3.4]{XVIII.prop1.3.4} 同构 \eqhyperref[XVIII.eq1.3.3.1]{1.3.3.1} 及 \eqhyperref[XVIII.eq1.3.3.4]{1.3.3.4} 的以下图式交换： $$
\xymatrix@C=.1cm@R=.1cm{
\langle \Ll_{1}\otimes \Ll_{2}, K_{1}+K_{2}] & \simeq & \langle
\Ll_{1}\otimes \Ll_{2}, K_{1}]+\langle \Ll_{1} \otimes \Ll_{2}, K_{2}]\\
 & & \rotatebox{90}{$\simeq$} \\
\rotatebox{90}{$\simeq$} & & \langle \Ll_{1}, K_{1}]+\langle\Ll_{2}, K_{1}]+\langle
    \Ll_{1}, K_{2}]+\langle \Ll_{2}, K_{2}]\\
 & & \rotatebox{-50}{$\Arrowvert$}\\
\langle\Ll_{1}, K_{1}+K_{2}]+\langle\Ll_{2},K_{1}+K_{2}] & \simeq  &
    \langle\Ll_{1}, K_{1}]+\langle \Ll_{1},K_{2}]+\langle\Ll_{2},K_{1}]+\langle\Ll_{2},K_{2}]. 
}
$$

%$$
%\xymatrix{
%\langle \Ll_1\otimes \Ll_2, K_1 + K_2] \ar@{-}[dd]&\simeq& \langle \Ll_1\otimes \Ll_2, K_1] + \langle \Ll_! \otimes
%  \Ll_2, K_2]\ar@{-}[d]\\
%  & & \langle \Ll_1\otimes \Ll_2, K_1] + \langle \Ll_! \otimes
%  \Ll_2, K_2] \ar@{=}[d]\\
%\langle \Ll_1\otimes \Ll_2, K_1] + \langle \Ll_! \otimes
%  \Ll_2, K_2] & \simeq& \langle \Ll_1\otimes \Ll_2, K_1] + \langle \Ll_! \otimes
%  \Ll_2, K_2].
%}
%$$


\end{enonce*}

此外，同构 \eqhyperref[XVIII.eq1.3.3.1]{1.3.3.1} 及 \eqhyperref[XVIII.eq1.3.3.4]{1.3.3.4} 与结合及交换同构相容，精确表述留给读者。于是，当 $\Ll$ 不\pageoriginale 必满足条件（\hyperref[XVIII.par1.3.2.1]{1.3.2.1}）时，可以用公式 \eqhyperref[XVIII.eq1.3.3.4]{1.3.3.4} 通过“线性”延拓 $\langle \Ll,K]$ 的定义。

\setcounter{subsubsection}{4} \subsubsection{}\etiquette[1.3.5]{XVIII.subsubsec1.3.5} 对 $X$ 上的每个有效相对除子 $D$，置 \begin{equation*}
\Tr_{D/S}(K)=\Tr_{D/S}(K|D){\quoi}.\tag{1.3.5.0}
\end{equation*}\etiquette[1.3.5.0]{XVIII.eq1.3.5.0}

若 $\Oo(D)$ 满足（\hyperref[XVIII.par1.3.2.1]{1.3.2.1}），则 $\Oo(D)$ 的截面 \hyperref[XVIII.sec1]{1} 是 $f_{\ast}\Oo(D)$ 的一个截面，作为这样的截面，它生成一个局部为 $f_{\ast}\Oo(D)$ 的直和因子的可逆层，并定义该层的一个平凡化。因此，同构 \eqhyperref[XVIII.eq1.3.3.8]{1.3.3.8} 给出一个同构 \begin{equation*}
\langle \Oo(D),K]\simeq \Tr_{D/S}(K){\quoi}.\tag{1.3.5.1}
\end{equation*}\etiquette[1.3.5.1]{XVIII.eq1.3.5.1}

立即验证，图式 $$\vcenter{
\xymatrix{
  \langle \Oo(D) , K_1 + K_2] \ar[r]^{\sim} \ar@{-}[d]_{\text{\rotatebox{90}{$\sim$}}}& \Tr_{D/S}(K_1 + K_2)\ar@{-}[d]^{\text{\rotatebox{90}{$\sim$}}}\\
\langle \Oo(D),K_1] + \langle \Oo(D) ,K_2] \ar[r]^{\sim}& \Tr_{D/S}(K_1) + \Tr_{D/S}(K_2)
}}\leqno{(1.3.5.2)}
$$\etiquette[1.3.5.2]{XVIII.eq1.3.5.2} 及 $$\vcenter{
\xymatrix{
\langle \Oo(D_1 + D_2),K] \ar[r]^{\sim}
\ar@{-}[d]^{\text{\rotatebox{90}{$\sim$}}}& \Tr_{(D_1 +
  D_{2})/S}(K)\ar@{-}[d]^{\text{\rotatebox{90}{$\sim$}}\quad (XVII\
  \hyperref[XVII.const6.3.27.1]{6.3.27.1})}\\ 
\langle \Oo(D_1),K] + \langle \Oo(D_2),K] \ar[r]^{\sim}& \Tr_{D_1 /S}(K) +
    \Tr_{D_{2} /S} (K)
}}\leqno{(1.3.5.3)}
$$\etiquette[1.3.5.3]{XVIII.eq1.3.5.3} 交换。由线性，可对每个相对除子定义同构 \eqhyperref[XVIII.eq1.3.5.1]{1.3.5.1}。

\subsubsection{}\etiquette[1.3.6]{XVIII.subsubsec1.3.6} 设 $D$ 和 $E$ 为两个有效相对 Cartier 除子，$f$ 为满足 $$
\text{\rm div\,}(f)=D-E{\quoi},
$$ 的有理函数，即 $\Oo(D)$ 与 $\Oo(E)$ 之间的一个同构。$\Oo(D)$ 与 $\Oo(E)$ 之间的同构 $f$ 通过 \eqhyperref[XVIII.eq1.3.5.1]{1.3.5.1} 定义一个同构 \begin{equation*}
\langle f,K]:\Tr_{D/S}(K)\xrightarrow{\sim} \Tr_{E/S}(K){\quoi}.\tag{1.3.6.1}
\end{equation*}\etiquette[1.3.6.1]{XVIII.eq1.3.6.1} 有\pageoriginale 以下公式，其中第二个由 \eqhyperref[XVIII.eq1.3.5.2]{1.3.5.2} 得到： \begin{align*}
& \langle fg,K]=\langle f,K].\langle g,K]\\
& \langle f,K_{1}+K_{2}]=\langle f,K_{1}]+\langle f,K_{2}]{\quoi}.
\end{align*}

若 $\lambda$ 为 $\Oo^{\ast}_{S}$ 的一个截面，则对 $X$ 上的每个可逆层 $\Ll$，乘以 $\lambda$ 是 $\Ll$ 的一个自同构。通过结构转移，此自同构在 \eqhyperref[XVIII.eq1.3.3.8]{1.3.3.8} 的两边诱导自同构，且它们在同构 \eqhyperref[XVIII.eq1.3.3.8]{1.3.3.8} 下相互对应。于是 $$
\langle \lambda, K]=\deg (K)(\lambda){\quoi}\;;
$$ 因而若 $\deg(K)\neq 0$，则 $\Tr_{D/S}(K)$ 与 $\Tr_{E/S}(K)$ 之间的同构 $\langle f,K]$ 依赖于 $f$ 的选取。

设 $\lambda\in \Gamma(S,\Oo^{\ast}_{S})$、$g\in \Gamma(S,G)$，设 $\Ll$ 为 $X$ 上的可逆层，$K$ 为 $X$ 上的 $G$-挠子。若 $\langle \lambda,g]$ 为由 $\Ll$ 的自同构 $\lambda$ 及 $K$ 的自同构 $g$ 通过结构转移得到的 $\langle \Ll,K]$ 的自同构，则将 $\langle \lambda,g]$ 与 $G$ 的一个截面等同后，它由 $\langle \lambda,g]=\langle \lambda,K]+\langle \Ll,g]$ 给出；由于 $\langle \Oo(D),g]=\Tr_{D/S}(g)$，前一公式给出 \begin{equation*}
\langle \lambda,g]=\deg(K)(\lambda)+\deg(\Ll).g{\quoi}.\tag{1.3.6.2}
\end{equation*}\etiquette[1.3.6.2]{XVIII.eq1.3.6.2}

\subsubsection{}\etiquette[1.3.7]{XVIII.subsubsec1.3.7} 设 $f:X\rightarrow S$ 为 $S$ 上的光滑射影曲线，$G$ 为 $S$ 上满足 \hyperref[XVIII.par1.3.1.1]{1.3.1.1} 的群。若 $X\xrightarrow{f_{0}}T\xrightarrow{\pi}S$ 为 $f$ 的 Stein 分解，则 $\pi$ 为 $S$ 的平展覆盖，且层 $f_{\ast}G_{m}$ 为环面 $\prod_{T/S}\GG_{m}$。在 $S$ 上局部地（对于平展拓扑），$X$ 为 $S$ 上一族具有非空几何连通纤维的曲线 $(X_{i})_{i\in I}$ 的不交并。若 $K$ 为 $G$ 下的挠子，$\Ll$ 为 $X$ 上的可逆层，则置 $$
\langle \Ll,K]=\sum_{i}\langle \Ll|X_{i},K|X_{i}]
$$ 并定义 $K:f_{\ast}\GG_{m}=\GG^{I}_{m}\rightarrow G$ 的次数，其坐标为 $\deg(K|X_{i}):\GG_{m}\rightarrow G$。局部构造可整体化，从而给出\pageoriginale 函子及态射 \begin{align*}
& \langle \Ll,K]_{X/S}=\Tr_{T/S}(\langle \Ll,K]_{X/T})\tag{1.3.7.1}\\
& \deg(K):f_{\ast}\GG_{m}=\prod_{T/S}\GG_{m}\longrightarrow G:\deg_{X/S}(K)=\prod_{T/S}\deg_{X/T}(K){\quoi}.\tag{1.3.7.2}
\end{align*}\etiquette[1.3.7.1]{XVIII.eq1.3.7.1} \etiquette[1.3.7.2]{XVIII.eq1.3.7.2}

对于 $\Ll$ 的由 $\Gamma(T,\Oo^{\ast}_{T})$ 的一个元素 $\lambda_{0}$ 定义的自同构 $\lambda$，有 \begin{equation*}
\langle \lambda,K]=\deg(K)(\lambda){\quoi}.\tag{1.3.7.3}
\end{equation*}\etiquette[1.3.7.3]{XVIII.eq1.3.7.3}

$K$ 的总次数定义为复合 $$
\deg \tot(K):\GG_{m}\longrightarrow
f_{\ast}\GG_{m}\xrightarrow{\deg(K)}G{\quoi}.
$$ 由前述结果得到：

\begin{enonce*}[remark]{公式汇集 1.3.8}\etiquette[1.3.8]{XVIII.for1.3.8} 对每个光滑射影曲线 $f:X\rightarrow S$，有一个函子 $\langle~,~]$，其构造与每个基变换相容，它将 $X$ 上的可逆层 $\Ll$ 及 $f^{\ast}G$ 下的挠子 $K$ 对应到 $S$ 上 $G$ 下的挠子 $\langle \Ll,K]$。此函子赋有以下附加结构，并满足以下条件： \end{enonce*}

\setcounter{subsubsection}{8} \paragraph{}\etiquette[1.3.8.1]{XVIII.par1.3.8.1} $\langle \Ll,K]$ 关于 $\Ll$ 和 $K$ 双可加；双可加同构与结合及交换同构和基变换相容，满足相容性（\hyperref[XVIII.prop1.3.4]{1.3.4}），并且与下述同构（\hyperref[XVIII.par1.3.8.2]{1.3.8.2}）相容（通过 \hyperrefext[XVII][XVII.par6.3.25.3]{6.3.25.3} 及 \hyperrefext[XVII][XVII.const6.3.27.1]{6.3.27.1}）。

\paragraph{}\etiquette[1.3.8.2]{XVIII.par1.3.8.2} 对 $X/S$ 上的每个有效相对除子 $E$，有典范同构 $$
\langle \Oo(E),K]\simeq \Tr_{E/S}(K|E){\quoi}.
$$

\paragraph{}\etiquette[1.3.8.3]{XVIII.par1.3.8.3} 对 $\Ll$ 的自同构 $\lambda\in f_{\ast}\GG_{m}$，次数 \eqhyperref[XVIII.eq1.3.7.2]{1.3.7.2} 满足 $$
\langle \lambda,K]=\deg(K)(\lambda){\quoi}.
$$ 对 $\lambda\in \GG_{m}(S)$，总次数满足 $$
\langle \lambda,K]=\deg\tot(K)(\lambda){\quoi}.
$$ 同样，\pageoriginale 对 $g\in G(S)$ 有 $\langle \Ll,g]=\deg(\Ll).g$。

由此推出典范同构： \begin{equation*}
\langle f^{\ast}\Ll,K]\simeq \deg\tot(K)(\Ll){\quoi}.\tag{1.3.8.4}
\end{equation*}\etiquette[1.3.8.4]{XVIII.eq1.3.8.4}

由定义立即可知，对每个同态 $\varphi:G\rightarrow G'$，有一个与上述资料相容的典范同构： \begin{equation*}
\langle \Ll,\varphi(K)]\simeq \varphi(\langle \Ll,K]){\quoi}.\tag{1.3.8.5}
\end{equation*}\etiquette[1.3.8.5]{XVIII.eq1.3.8.5}

\setcounter{paragraph}{5} \paragraph{}\etiquette[1.3.8.6]{XVIII.par1.3.8.6} 设 $u:X\rightarrow Y$ 为 $S$ 上光滑射影曲线之间的平坦态射。则有与双可加性及基变换相容的同构： \begin{itemize} \item[--] 对 $X$ 上的可逆层 $\Ll$ 及 $Y$ 上的 $G_{Y}$-挠子 $K$，有 $$
\langle \Ll, u^{\ast}K]\xrightarrow{\sim}\langle N_{X/Y}(\Ll),K]{\quoi}.
$$

\item[--] 对 $Y$ 上的可逆层 $\Ll$ 及 $X$ 上的 $G_{X}$-挠子 $K$，有 $$
\langle u^{\ast}\Ll,K]\xrightarrow{\sim}
\langle \Ll,\Tr_{X/Y}K]{\quoi}.
$$ \end{itemize} 此外，这些同构对曲线态射的复合 $uv$ 满足显然的相容性；对 $\Ll=\Oo(E)$，它们与同构（\hyperrefext[XVII][XVII.const6.3.27.2]{6.3.27.2}）等同： \begin{align*}
\Tr_{E/S}(u^{\ast}K) &\xrightarrow{\sim} \Tr_{u_{\ast}E/S}(K)\\
\Tr_{u^{\ast}E/S}(K) &\xrightarrow{\sim} \Tr_{E/S}\Tr_{X/Y}K{\quoi}.
\end{align*} 还须证明 \hyperref[XVIII.par1.3.8.6]{1.3.8.6}。归结为 $X$ 和 $Y$ 具有非空几何连通纤维的情形，并注意到：对 $\Ll$ 的每个作为 $\Ll$ 的直接像的截面不为零的截面 $s$，上述显式公式给出一个所需类型的同构 $\Psi_{s}$，它关于 $s$ 齐次，因而当 $\Ll$ 满足（\hyperref[XVIII.par1.3.2.1]{1.3.2.1}）时与 $s$ 无关。由此即得一般情形。

\subsubsection{}\etiquette[1.3.9]{XVIII.subsubsec1.3.9} 在 \hyperref[XVIII.subsubsec1.3.7]{1.3.7} 的假设下，已在（\hyperref[XVIII.for1.3.8]{1.3.8}）中看到，$X$ 上的每个 $G$-挠子 $K$ 定义一个函子 $\langle \ast,K]$，以下记为 $F(\ast)$，它赋有\pageoriginale 以下附加资料，并满足以下条件： \begin{enumerate} \renewcommand{\theenumi}{\roman{enumi}} \renewcommand{\labelenumi}{(\theenumi)} \item $F$ 为 $S$ 上的 fppf 叠的态射（\hyperrefext[XVII][XVII.subsec0.10]{0.10} 及 \textsc{Giraud} \cite{sga4.exp.XVIII.bib.1}） \begin{itemize} \item[--] 其源为如下叠：对一个 $S$-概形 $T$，其 $T$-对象为 $X_{T}$ 上的可逆层；

\item[--] 其目标为 $G$-挠子的叠。 \end{itemize}

\item 函子 $F$ 赋有可加同构 $F(\Ll_{1}\otimes \Ll_{2})\xrightarrow{\sim}F(\Ll_{1})+F(\Ll_{2})$，这些同构具有函子性，与结合及交换同构相容，并且与基变换相容。 \end{enumerate}

\begin{enonce*}{定理 1.3.10}\etiquette[1.3.10]{XVIII.thm1.3.10} 在 \hyperref[XVIII.subsubsec1.3.7]{1.3.7} 的假设下，函子 $\Phi:K\mapsto \langle \ast,K]$ 是 $X$ 上的 $G$-挠子范畴与（\hyperref[XVIII.subsubsec1.3.9]{1.3.9}）中所考虑的叠态射范畴之间的等价。 \end{enonce*}

容易归结为 $X$ 具有非空连通纤维的情形。

设 $F$ 满足 \hyperref[XVIII.subsubsec1.3.9]{1.3.9} (i) 及 (ii)。对每个 $S$-概形 $T$ 及 $X_{T}/T$ 的每个截面 $t$，也用 $t$ 表示相对除子 $t(T)$。则 $F(\Oo(t))$ 为 $T$ 上的 $G_{T}$-挠子。当 $t$ 为由 $X$ 到 $X$ 的恒等映射定义的 $X$ 的取值于 $X$ 的“泛点”时，用 $\Psi(F)$ 表示这样得到的 $X$ 上的挠子。对上述每个截面 $t$，有 \begin{equation*}
F(\Oo(t))\simeq t^{\ast}\Psi(F){\quoi}.\tag{1.3.10.1}
\end{equation*}\etiquette[1.3.10.1]{XVIII.eq1.3.10.1}

同构 \hyperref[XVIII.par1.3.8.2]{1.3.8.2} 给出函子的同构 $\Psi\circ\Phi \simeq \Id$。来构造一个同构 $\Phi\circ\Psi\simeq \Id$。

对 $X_{T}$ 在 $T$ 上的每个截面 $t$，同构 \eqhyperref[XVIII.eq1.3.10.1]{1.3.10.1} 是一个同构 $$
F(\Oo(t))\xrightarrow[\alpha]{\sim}\langle \Oo(t),\Psi(F)]{\quoi}.
$$

由可加性，$X_{T}$ 在 $T$ 上的每族截面 $t_{1},\ldots, t_{n}$ 定义一个同构 $$
F(\Oo(\sum t_{i}))\xrightarrow[\alpha]{\sim}\langle \Oo(\sum
t_{i}),\Psi(F)]{\quoi},
$$ 由 \hyperref[XVIII.subsubsec1.3.9]{1.3.9} (ii)，此同构不依赖于各 $t_{i}$ 的顺序。

取\pageoriginale $T=(X/S)^{n}$，并取 $t_{1}\ldots t_{n}$ 为截面的泛 $n^{\text{\rm 元组}}$。若 $T'=\Sym^{n}_{S}(X)$，则在 $X_{T'}$ 上有一个除子 $\sum t_{i}$，其逆像为 $X_{T}$ 上的除子 $\sum t_{i}$。在 $T$ 上，$\alpha$ 为一个对称同构 $$
\alpha:F(\Oo(\sum t_{i}))\xrightarrow{\sim} \langle \Oo(\sum
t_{i}),\Psi(F)]
$$，它联系的挠子均为 $T'$ 上挠子的逆像。由假设（\hyperref[XVIII.par1.3.1.1]{1.3.1.1}）(ii)，此同构来自 $T'$ 上唯一的同构： $$
\alpha:F(\Oo(\sum t_{i}))\xrightarrow{\sim} \langle \Oo (\sum
t_{i}),\Psi (F)]\qquad (\text{\rm 在 } T'){\quoi}.
$$

由 \hyperrefext[XVII][XVII.prop6.3.9]{6.3.9}，对 $X$ 上的每个相对除子 $D$，有一个同构 \begin{equation*}
\alpha_{D}:F(\Oo(D))\xrightarrow{\sim}
\langle \Oo(D),\Psi(F)]{\quoi},\tag{1.3.10.2}
\end{equation*}\etiquette[1.3.10.2]{XVIII.eq1.3.10.2} 它关于 $D$ 可加，且其构造与每个基变换相容。

设 $\Ll$ 为 $X$ 上的在 \hyperref[XVIII.par1.3.2.1]{1.3.2.1} 意义下充分丰沛的可逆层。$f_{\ast}\Ll$ 的每个处处非零截面 $s$，采用 \hyperref[XVIII.subsubsec1.3.2]{1.3.2} 的记号，定义一个同构 $$
s:\Oo(D_{s})\xrightarrow{\sim}\Ll{\quoi},
$$ 从而由 \eqhyperref[XVIII.eq1.3.10.2]{1.3.10.2} 定义一个同构 \begin{equation*}
\alpha_{s}:F(\Ll)\xrightarrow{\sim}\langle \Ll,\Psi(F)]{\quoi}.\tag{1.3.10.3}
\end{equation*}\etiquette[1.3.10.3]{XVIII.eq1.3.10.3}

证明此同构不依赖于 $s$。事实上： \begin{enumerate} \renewcommand{\theenumi}{\roman{enumi}} \renewcommand{\labelenumi}{(\theenumi)} \item $\alpha_{s}$ 的构造与每个基变换可交换。

\item 存在同态 $\varphi:\GG_{m}\rightarrow G$，使 $$
\alpha_{\lambda s}=\varphi(\lambda).\alpha_{s}
$$ （由 \eqhyperref[XVIII.eq1.3.10.3]{1.3.10.3} 两边的函子性）。 \end{enumerate}

必有 $\varphi=0$，否则由拉回 $F(\Ll)-\langle \Ll,\Psi(F)]$ 得到的 $\PP((f_{\ast}\Ll)^{\vee})$ 上的挠子将具有次数 $\neq 0$。因此，各 $\alpha_{s}$ 定义一个从 $\PP((f_{\ast}\Ll)^{\vee})$ 到 $G$ 的态射，而由假设\pageoriginale（\hyperref[XVIII.par1.3.1.1]{1.3.1.1}）(i)，此态射为常值。因此可置 $\alpha_{\Ll}=\alpha_{s}$。

各同构 $\alpha_{\Ll}$ 因此构成函子 $F$ 与 $\langle *,\Psi(F)]$ 在充分丰沛层上的限制之间的一个函子同构。此同构与可加同构相容，因而可通过线性延拓到所有可逆层，从而完成同构 $$
\alpha:\Phi\circ\Psi\simeq \Id{\quoi}.
$$ 的构造。

\begin{enonce*}[remark]{例 1.3.11}\etiquette[1.3.11]{XVIII.exem1.3.11} $G=\GG_{m}$。在 \hyperref[XVIII.subsubsec1.3.7]{1.3.7} 的假设下，特别当 $G=\GG_{m}$ 时，使用记号 $\langle~,~\rangle $ 代替 $\langle~,~]$；若 $\Ll$ 和 $\Mm$ 为 $X$ 上的两个可逆层，则 $\langle \Ll$, $\Mm\rangle $ 为 $S$ 上的可逆层。此外，有同构（\hyperref[XVIII.par1.3.8.2]{1.3.8.2}） $$
\langle \Oo(D),\Mm\rangle \simeq N_{D/S}(\Mm){\quoi}.
$$

若 $D$ 和 $E$ 为不相交的有效相对除子，则 $\Oo(E)|D$ 典范同构于 $\Oo$，从而有同构 \begin{equation}
\sigma_{D,E}:\langle \Oo(D),\Oo(E)\rangle \simeq N_{D/S}(\Oo(E))\simeq
N_{D/S}(\Oo)\simeq \Oo{\quoi}.\tag{1.3.11.1}
\end{equation}\etiquette[1.3.11.1]{XVIII.eq1.3.11.1} \end{enonce*}

\begin{enonce*}{命题 1.3.12}\etiquette[1.3.12]{XVIII.prop1.3.12} 设 $(D_{i})_{i=1,2}$ 和 $E$ 为有效相对除子，且 $D_{i}$ 与 $E$ 不相交 $(i=1,2)$。设 $f$ 为满足 $\text{\rm div\,}(f)=D_{1}-D_{2}$ 的有理函数，即一个同构 $f:\Oo(D_{1})\xrightarrow{\sim}\Oo(D_{2})$。以下同构的图式交换： \begin{equation*}
\vcenter{
\xymatrix@C=2.5cm@R=1.25cm{
\langle \Oo(D_1), \Oo(E)\rangle   \ar[d]^{\langle f,\Oo(E)\rangle } \ar[r]^(.65){\sigma_{D_{1},E}} & \Oo\ar[d]^{N_{E/S}(f)}\\
\langle \Oo(D_2),\Oo(E)\rangle    \ar[r]^(.65){\sigma_{D_{2},E}} & \Oo
}}\tag{1.3.12.1}
\end{equation*}\etiquette[1.3.12.1]{XVIII.eq1.3.12.1} \begin{equation*}
\vcenter{
\xymatrix@C=2.5cm@R=1.25cm{
\langle \Oo(E), \Oo(D_1)\rangle   \ar[d]^{\langle \Oo(E),f\rangle } \ar[r]^(.65){\sigma_{E,D_{1}}} & \Oo\ar[d]^{N_{E/S}(f)}\\
\langle \Oo(E),\Oo(D_2)\rangle    \ar[r]^(.65){\sigma_{E,D_{2}}} & \Oo
}}\tag{1.3.12.2}
\end{equation*}\etiquette[1.3.12.2]{XVIII.eq1.3.12.2} \end{enonce*}

公式\pageoriginale \eqhyperref[XVIII.eq1.3.12.2]{1.3.12.2} 仅仅表示图式 $$
\xymatrix@C=2.5cm@R=1.25cm{
N_{E/S}(\Oo(D_1))  \ar[r]^(.55){\sim} \ar[d]^{N_{E/S}(f)} & \Oo \ar[d]^{N_{E/S}(f)} \\         
N_{E/S}(\Oo(D_2))  \ar[r]^(.55){\sim} & \Oo.
}
$$ 的交换性。

对每个 $S$-概形 $T$ 及 $X_{T}$ 上的每个可逆层 $\Ll$，置 $F(\Ll)=N_{E_{T}/T}(\Ll)$。此函子满足 \hyperref[XVIII.thm1.3.10]{1.3.10} 的假设，因而存在唯一的可逆层 $\Mm$，赋有与双可加性相容的函子同构： $$
N_{E_{T}/T}(\Ll)\xrightarrow{\sim}\langle \Ll,\Mm\rangle {\quoi}.
$$

此外，若 $d:T\rightarrow X_{T}$ 为 $X$ 在 $T$ 上的一个截面，则有典范同构 $$
d^{\ast}\Mm\xrightarrow{\sim} N_{E_{T}/T}(\Oo(d)){\quoi}.
$$ 取 $d$ 为 $X$ 的泛截面 $\Delta:X\rightarrow
X_{X}$，得到 $$
\Mm\xrightarrow{\sim} N_{E_{X}/X}(\Oo(\Delta)){\quoi}.
$$

在 $S$ 上局部地，$E$ 是若干截面之和 $E=\sum t_{i}$（\hyperrefext[XVII][XVII.subsec6.3]{6.3}）；对这样的 $E$，有 $$
N_{E_{X}/X}(\Oo(\Delta))=\otimes
t^{\ast}_{i}\Oo(\Delta)=\sum\Oo(t_{i})=\Oo(E){\quoi},
$$ 而此同构不依赖于各 $t_{i}$ 的顺序，因而定义一个典范同构\nde{此论证似乎假定，局部地，各 $t_i$ 在相差一个置换的意义下良定义，但这一点并不成立。例如，$S=\Spec k[\epsilon]$ 为对偶数代数的谱，$X=\PP^1_S$，而 $E$ 是 $\AAA^1_S=\Spec \Oo_S[T]\subset X$ 中由方程 $(T-\varepsilon)(T+\varepsilon)=T^2$ 定义的闭子概形时，便可看出这一点。看来应先应用“消去下降”条件 \hyperref[XVIII.par1.3.1.1]{1.3.1.1}~(ii)，在经 $\Sym^n_S X\to S$ 基变换后 $X$ 上的泛除子的情形下得到结果，再对除子 $E$ 推出结果。} \begin{equation*}
N_{E_{X}/X}(\Oo(\Delta))\sim \Oo(E){\quoi},\tag{1.3.12.3}
\end{equation*}\etiquette[1.3.12.3]{XVIII.eq1.3.12.3} 从而 \begin{equation*}
N_{E_{T}/T}(\Ll)\xrightarrow{\sim}\langle \Ll,\Oo(E)\rangle{\quoi}.%
\tag{1.3.12.4}
\end{equation*}\etiquette[1.3.12.4]{XVIII.eq1.3.12.4}

同构 \eqhyperref[XVIII.eq1.3.12.3]{1.3.12.3} 与双可加性及基变换相容，并且还由以下性质刻画：对 $X_{T}$ 在 $T$ 上的每个\pageoriginale 截面 $d$，同构 $$
N_{E/T}(\Oo(d))\xrightarrow{\sim}\langle \Oo(d),\Oo(E)\rangle
\xrightarrow{\sim} d^{\ast}\Oo(E)
$$ 为泛同构 \eqhyperref[XVIII.eq1.3.12.3]{1.3.12.3} 的拉回。特别地，当 $d$ 与 $E$ 不相交时，图式 $$
\xymatrix{
N_{E/T}(\Oo(d))  \ar[r]^{\sim} \ar@{=}[d]& \langle \Oo(d), \Oo(E)\rangle 
\ar[r]^{\sim} & d^\ast\Oo(E) \ar@{=}[d]\\
\Oo              \ar@{=}[rr] &          & \Oo  
}
$$ 交换。由可加性\nde{上述同一异议在此也适用。可用同样方式，考虑经 $\Sym^n_S (X-E)\to S$ 基变换后定义的“与 $E$ 不相交的泛除子”。}，可知当 $D$ 与 $E$ 不相交时，图式 \begin{equation*}
\vcenter{
\xymatrix{
N_{E/T}(\Oo(D))  \ar[r]^{\sim} \ar@{=}[d]& \langle \Oo(D), \Oo(E)\rangle 
\ar[r]^{\sim} & N_{D/T}(\Oo(E)) \ar@{=}[d]\\
\Oo              \ar@{=}[rr] &          & \Oo  
}}\tag{1.3.12.5}
\end{equation*}\etiquette[1.3.12.5]{XVIII.eq1.3.12.5} 交换。因此，箭头 $\sigma_{E,D}$ 为 \hyperref[XVIII.eq1.3.12.4]{1.3.12.4} 的逆与一个显然的箭头的复合： $$
\langle \Oo(D),\Oo(E)\rangle \xleftarrow{\sim} N_{E/S}(\Oo(D))\longrightarrow
\Oo
$$ 由于 \hyperref[XVIII.eq1.3.12.4]{1.3.12.4} 具有函子性，立即得到 \eqhyperref[XVIII.eq1.3.12.1]{1.3.12.1}。

由 \hyperref[XVIII.prop1.3.12]{1.3.12} 推出（\cf \textsc{Serre} \cite[Ch.~III, prop.~7, p.~46]{sga4.exp.XVIII.bib.6}）

\begin{enonce*}{推论 1.3.13}\etiquette[1.3.13]{XVIII.cor1.3.13} 设 $(D_{i})_{i=1,2}$ 及 $(E_{i})_{i=1,2}$ 为有效相对除子，且 $D_{i}$ 与 $E_{j}$ 不相交；设 $f$ 为满足 $\text{\rm div\,}(f)=D_{1}-D_{2}$ 的有理函数，$g$ 为满足 $\text{\rm div\,}(g)=E_{1}-E_{2}$ 的有理函数。则 $$
N_{E_{1}/S}(f).N_{E_{2}/S}(f)^{-1}=N_{D_{1}/S}(g).N_{D_{2}/S}(g)^{-1}
$$ 简写为 \begin{equation*}
f(\text{\rm div\,}(g))=g(\text{\rm div\,}(f)){\quoi}.\tag{1.3.13.1}
\end{equation*}\etiquette[1.3.13.1]{XVIII.eq1.3.13.1} \end{enonce*}

将此公式写成 $$
N_{D_{2}/S}(g)N_{E_{1}/S}(f)=N_{E_{2}/S}(f)N_{D_{1}/S}(g){\quoi},
$$，则经由 \eqhyperref[XVIII.eq1.3.12.1]{1.3.12.1} 及 \eqhyperref[XVIII.eq1.3.12.2]{1.3.12.2}，它表示图式 $$
\xymatrix@C=2.25cm@R=1.5cm{
\langle \Oo(D_1), \Oo(E_1)\rangle   \ar[r]^{\langle \Oo(D_1), g\rangle } \ar[d]^{\langle f,\Oo(E_1)\rangle }&
\langle \Oo(D_1),\Oo(E_2)\rangle  \ar[d]^{\langle f,\Oo(E_2)\rangle }\\
\langle \Oo(D_2), \Oo(E_1)\rangle   \ar[r]^{\langle \Oo(D_2),g\rangle } & \langle \Oo(D_2), \Oo(E_2)\rangle .
}
$$ 的交换性。

\begin{enonce*}{推论 1.3.14}\etiquette[1.3.14]{XVIII.cor1.3.14} 当\pageoriginale $G=\GG_{m}$ 时，次数（\hyperref[XVIII.subsubsec1.3.7]{1.3.7}）与可逆层通常的次数相同。 \end{enonce*}

只须在 $f$ 为 $f_0\in \Gamma(S,\Oo^{\ast}_{S})$ 的拉回时，比较公式 \eqhyperref[XVIII.eq1.3.7.3]{1.3.7.3} 及 \eqhyperref[XVIII.eq1.3.12.1]{1.3.12.1}。

\setcounter{subsubsection}{14} \subsubsection{}\etiquette[1.3.15]{XVIII.subsubsec1.3.15} 对任意充分丰沛的 $\Ll$ 和 $\Mm$，在 $S$ 上局部地存在不相交的除子 $D$ 和 $E$ 及同构 $$
\alpha:\Ll\longrightarrow \Oo(D)\quaet
\beta:\Mm\longrightarrow \Oo(E){\quoi}.
$$。用 $\tau$ 表示使图式 $$
\xymatrix@C=1.5cm@R=1.25cm{
\langle \Ll,\Mm\rangle  \ar[r]^{\langle \alpha,\beta\rangle }
\ar[d]^{\tau}_{\text{\rotatebox{90}{$\backsim$}}}& \langle \Oo(D), \Oo(E)\rangle  \ar[r]^(.7){\sigma_{D,E}}& \Oo\ar@{=}[d]\\
\langle \Mm,\Ll\rangle  \ar[r]^{\langle \beta,\alpha\rangle }& \langle \Oo(E), \Oo(D)\rangle  \ar[r]^(.7){\sigma_{E,D}}& \Oo.
}
$$ 交换的同构。由 \hyperref[XVIII.prop1.3.12]{1.3.12}，此同构不依赖于 $D$、$E$、$\alpha$ 和 $\beta$ 的具体选择；因而可整体化为一个典范同构 \begin{equation*}
\tau:\langle \Ll,\Mm\rangle \xrightarrow{\sim}
\langle \Mm,\Ll\rangle {\quoi}.\tag{1.3.15.1}
\end{equation*}\etiquette[1.3.15.1]{XVIII.eq1.3.15.1}

此同构满足 $\tau^{2}=1$；由于 \eqhyperref[XVIII.eq1.3.11.1]{1.3.11.1} 与双可加性相容，它也与双可加性相容，并通过线性延拓到任意可逆层。

由 \eqhyperref[XVIII.eq1.3.12.5]{1.3.12.5}，态射 \eqhyperref[XVIII.eq1.3.12.4]{1.3.12.4} 可与复合 $$
N_{E_{T}/T}(\Ll)\xleftarrow{\sim}
\langle \Oo(E),\Ll\rangle \xrightarrow{\sim\tau}\langle \Ll,\Oo(E)\rangle 
$$ 等同（在 $\Ll=\Oo(D)$ 且 $D$ 与 $E$ 不相交时验证）。

\begin{enonce*}{公式汇集 1.3.16}\etiquette[1.3.16]{XVIII.for1.3.16} 对\pageoriginale 每个光滑射影曲线 $f:X\rightarrow
S$，存在一个函子 $\langle~,~\rangle $，其构造与每个基变换可交换，它将 $X$ 上的两个可逆层 $\Ll$ 和 $\Mm$ 对应到 $S$ 上的可逆层 $\langle \Ll,\Mm\rangle $。此函子赋有以下附加结构，并满足以下条件： \end{enonce*}

\setcounter{subsubsection}{16} \paragraph{}\etiquette[1.3.16.1]{XVIII.par1.3.16.1} 采用 \hyperref[XVIII.for1.3.8]{1.3.8} 的记号，当 $G=\GG_{m}$ 时，有 $$
\langle \Ll,\Mm\rangle =\langle \Ll,\Mm]{\quoi},
$$ 特别地，$\langle \Ll,\Mm\rangle $ 在 \hyperref[XVIII.par1.3.8.1]{1.3.8.1} 所指的意义下关于 $\Ll$、$\Mm$ 双可加，且此双可加性与如下典范同构的右边的双可加性相容（同构关于 $\Mm$ 具有函子性）： $$
\langle \Oo(E),\Mm\rangle \xrightarrow{\sim} N_{E/S}(\Mm){\quoi}.
$$

\paragraph{}\etiquette[1.3.16.2]{XVIII.par1.3.16.2} 有一个对称同构 $$
\tau:\langle \Ll,\Mm\rangle \xrightarrow{\sim}\langle \Mm,\Ll\rangle {\quoi},
$$，它关于 $\Ll$ 和 $\Mm$ 具有函子性，与基变换及双可加同构相容，并满足 $\tau^{2}=1$。对不相交的除子 $D$ 和 $E$，它使以下图式交换： $$
\xymatrix{
N_{D/S}(\Oo(E))\ar@{-}[d]^{\text{\rotatebox{90}{$\backsim$}}}& \langle \Oo(D), \Oo(E)\rangle \ar[l]\ar[r]^{\tau} &
\langle \Oo(E), \Oo(D)\rangle  \ar[r]&
N_{E/S}(\Oo(D))\ar@{-}[d]^{\text{\rotatebox{90}{$\backsim$}}}\\
\Oo \ar@{=}[rrr]&    & & \Oo.
}
$$

这些公式推出公式 \eqhyperref[XVIII.eq1.3.12.1]{1.3.12.1}、\eqhyperref[XVIII.eq1.3.12.2]{1.3.12.2}、\eqhyperref[XVIII.eq1.3.13.1]{1.3.13.1}，以及

\paragraph{}\etiquette[1.3.16.3]{XVIII.par1.3.16.3} 对 $a$、$b\in \Gamma(S,\Oo^{\ast}_{S})$，用 $\langle a,b\rangle $ 表示由 $\Ll$ 和 $\Mm$ 的自同构 $a$ 及 $b$，通过函子性得到的 $\langle \Ll,\Mm\rangle $ 的自同构，并将其与 $\GG_{m}$ 的一个截面等同，则有 $$
\langle a,b\rangle =a^{\deg(\Mm)} b^{\deg(\Ll)}{\quoi}.
$$

\paragraph{}\etiquette[1.3.16.4]{XVIII.par1.3.16.4} 若 $\Ll_0$（\resp $\Mm_0$）为 $S$ 上的可逆模，则有典范同构 \begin{align*}
\langle f^{\ast}\Ll_0,\Mm\rangle  &\xrightarrow{\sim}
\Ll_0^{\otimes\deg(\Mm)}\\
\langle \Ll,f^{\ast}\Mm_0\rangle  &\xrightarrow{\sim} \Mm_0^{\otimes\deg(\Ll)}\qquad){\quoi}.\tag*{(\resp} 
\end{align*}

\paragraph{}\etiquette[1.3.16.5]{XVIII.par1.3.16.5} 若\pageoriginale $u:X\rightarrow Y$ 为 $S$ 上光滑射影曲线之间的有限平坦态射，则对 $X$ 上的可逆层 $\Ll$ 及 $Y$ 上的可逆层 $\Mm$，同构 \hyperref[XVIII.par1.3.8.6]{1.3.8.6} \begin{align*}
\langle \Ll,u^{\ast}\Mm\rangle  &\sim
\langle N_{X/Y}\Ll,\Mm\rangle \\
\langle u^{\ast}\Mm,\Ll\rangle  &\sim \langle \Mm,N_{X/Y}\Ll\rangle 
\end{align*} 通过对称性 \eqhyperref[XVIII.par1.3.16.2]{1.3.16.2} 相互推出（在 $\Ll=\Oo(E)$ 及 $\Mm=\Oo(F)$ 且 $E$ 与 $F$ 不相交时验证）。

\paragraph{}\etiquette[1.3.16.6]{XVIII.par1.3.16.6} 对 $X$ 上的可逆层 $\Ll$，对称箭头 $$
\tau:\langle \Ll,\Ll\rangle \longrightarrow
\langle \Ll,\Ll\rangle 
$$ 为乘以 $(-1)^{\deg(\Ll)}$。

\subsubsection{}\etiquette[1.3.17]{XVIII.subsubsec1.3.17} 还须证明 \hyperref[XVIII.par1.3.16.6]{1.3.16.6}。给出两个证明，第一个使用行列式理论（\SGA 6 XI）\nde{如 \cite{sga4.exp.XVIII.bib.0ter} 所述，该讲并未写成；该处还提醒：行列式对于区别三角形的乘法性质须谨慎处理，否则会有反例。关于此问题可参阅 \cite{sga4.exp.XVIII.bib.3bis}。}。

\sisi{}{在给出各证明前，可先注意到：若 $\Ll_1$ 和 $\Ll_2$ 为 $X$ 上的两个可逆层，且 $\Ll=\Ll_1\otimes\Ll_2$，则自同构 $\tau_{\Ll}\colon \langle
\Ll,\Ll\rangle \longrightarrow \langle \Ll,\Ll\rangle$（与 $S$ 上的一个可逆函数等同）是自同构 $\tau_{\Ll_1}$ 与 $\tau_{\Ll_2}$ 的乘积。因此，当 $S$ 为域的谱（且 $X$ 几何连通）时，可假设 $\Ll=\Oo(D)$，其中 $D$ 为非零有效除子。由于 $\tau^2=\Id$，必要时将 $\Ll$ 换为一个充分大的奇数次幂（或写成两个丰沛丛之商），还可使 $\Ll$ 任意丰沛。

下述论证还需归结为 $S$ 为域的谱的情形。为此，只须归结为 $S$ 约化的情形；这可通过证明在 $S$ 上 fppf 局部地，资料 $(X,\Ll)$ 来自约化基底上的类似资料来实现。在作此归约前，可先归结为 $X\to S$ 赋有一个截面 $s\colon S\to
X$，且 $X\to S$ 的几何纤维连通并具有某个亏格 $g$ 的情形。利用 \cite[\S 5.2]{sga4.exp.XVIII.bib.4}，可得到一个亏格为 $g$ 的光滑射影曲线 $X_0\to S_0$，其中 $S_0$ 在 $\Spec(\ZZ)$ 上光滑且有限型，使\emph{局部地}，上述每个相对曲线均由它经基变换得到。利用截面 $s_0\colon S_0\to X_0$，可将 Picard 概形 $S_1=\sheafPic_{X_0/S_0}\to S_0$ 刚性化，从而在 $X_1=X_0\times_{S_0}S_1$ 上得到一个“泛”线丛 $\Ll_1$。局部地，每个上述赋有线丛的相对曲线，均来自 $S_1$ 上的 $(X_1,\Ll_1)$，而该基底确实约化\nde{前两段是本版增补的。}。}

\emph{第一种证明}：设 $u:T\rightarrow S$ 为有限局部自由态射。若 $\Ll$ 为 $T$ 上的可逆层，则有同构 \begin{equation*}
N_{T/S}(\Ll)\simeq \det(u_{\ast}\Ll)\det(u_{\ast}\Oo_{T})^{-1}\tag{1.3.17.1}
\end{equation*}\etiquette[1.3.17.1]{XVIII.eq1.3.17.1} 它由以下三个条件唯一确定： \begin{enumerate} \renewcommand{\theenumi}{\alph{enumi}} \renewcommand{\labelenumi}{\theenumi)} \item 其构造与基变换可交换；

\item 当 $\Ll=\Oo_{T}$ 时，它为 $\Oo_{S}$ 的恒等态射；

\item 它关于 $\Ll$ 具有函子性。 \end{enumerate}

设 $D$ 和 $E$ 为 $X$ 上的两个有效相对除子，$i$ 为 $D$ 到 $X$ 的包含，$u$ 为 $D$ 到 $S$ 的投射。短正合列
$$
0\longrightarrow \Oo (-E)\longrightarrow \Oo\longrightarrow
\Oo_{E}\longrightarrow 0
$$
给出同构
$$
\langle \Oo(D),\Oo(-E)\rangle \simeq N_{D/S}(\Oo(-E))\simeq
\det(u_{\ast}i^{\ast}\Oo(-E))\det (u_{\ast}i^{\ast}\Oo)^{-1}
$$
作为\pageoriginale 一个同构
$$
\langle \Oo(D),\Oo(-E)\rangle =\det Ru_{\ast}(Li^{\ast}\Oo_{E})^{-1}\simeq
\det Rf_{\ast}(\Oo_{D}\mathop{\otimes}^{\LL}
\Oo_{E})^{-1}{\quoi}.
$$
或者
\begin{equation*}
\langle \Oo(D),\Oo(E)\rangle \simeq \det Rf_{\ast}(\Oo_{D}\mathop{\otimes}^{\LL}\Oo_{E}){\quoi}.\tag{1.3.17.2}
\end{equation*}\etiquette[1.3.17.2]{XVIII.eq1.3.17.2}

对不相交的 $D$ 和 $E$，有 $\Oo_{D}\displaystyle{\mathop{\otimes}^{\LL}}\Oo_{E}=0$，由此得到的 $\langle \Oo(D)$, $\Oo(E)\rangle $ 与 $\Oo_{S}$ 的同构，便是已考虑过的同构。所以，对不相交的 $D$ 和 $E$，对称性 $\tau$ 与对称同构 $\Oo_{D}\displaystyle{\mathop{\otimes}^{\LL}}\Oo_{E}\xrightarrow{\sim}\Oo_{E}\displaystyle{\mathop{\otimes}^{\LL}}\Oo_{D}$ 通过 \eqhyperref[XVIII.eq1.3.17.2]{1.3.17.2} 相容。由特殊化\nde{此处似乎须先按前述方法，归结为 $S$ 为域的谱的情形。此时，可让 $D$ 和 $E$ 出现在由整概形参数化的两个除子族中，使这两个族中的除子在泛点上方不相交。例如，可考虑这样的泛情形：在整概形 $\Sym^n_S
X\times \Sym^m_S X$ 的一个稠密开子集上，条件 $D$ 与 $E$ 不相交成立。}，这对任意 $D$ 和 $E$ 仍然成立。

当 $D=E$ 由理想 $I$ 定义时，有 \begin{alignat*}{2}
&
ℋ^{0}(\Oo_{D}\mathop{\otimes}^{\LL}\Oo_{D})=\Oo_{D}
&\qquad& \text{\rm 对称 } = +\sisi{1}{\Id}\\
& ℋ^{-1}(\Oo_{D}\mathop{\otimes}^{\LL}\Oo_{D})=I/I^{2} &&
\text{\rm 对称 } =-\sisi{1}{\Id}\\
& ℋ^{i}(\Oo_{D}\mathop{\otimes}^{\LL}\Oo_{D})=0 &&
\text{\rm 当 } i\neq 0, -1{\quoi}.
\end{alignat*} 因此，映射 $$
\tau:\langle \Oo(D), \Oo(D)\rangle \xrightarrow{\sim}\langle \Oo(D), \Oo(D)\rangle 
$$ 为 $(-1)^{\rk I/I^{2}}=(-1)^{\deg(\Oo(D))}$，这就证明了 \hyperref[XVIII.par1.3.16.6]{1.3.16.6}。

\smallskip \noindent \emph{第二种证明}：归结为“泛”情形后，可假设 $S$ 约化。此情形又归结为 $S$ 为代数闭域 $k$ 的谱的情形。此时，设 $X_{1}$、$X_{2}$、$Y_{1}$ 和 $Y_{2}$ 为 $X$ 上的四个除子，$f$ 和 $g$ 为两个有理函数，并假设 \begin{align*}
&\text{\rm div\,}(f)=X_{2}-X_{1}\\
&\text{\rm div\,}(g)=Y_{2}-Y_{1}\\
&(X_{1}\cup Y_{2})\cap (X_{2}\cup Y_{1})=∅{\quoi}.
\end{align*} 于是，映射 $f$ 和 $g$ 定义\nde{与原版相比，从此处至本段末，许多地方的“符号”已有修改。} $$
\langle f,g\rangle :\langle \Oo(X_{2}), \Oo(Y_{2})\rangle \longrightarrow
\langle \Oo(X_{1}),\Oo(Y_{1})\rangle {\quoi}.
$$ 两\pageoriginale 边均典范同构于 $k$ \eqhyperref[XVIII.eq1.3.11.1]{1.3.11.1}，因而 $\langle f,g\rangle $ 可与 $k^{\ast}$ 的一个元素等同。

\hypertarget{sga4.exp.XVIII.lemma.1.3.18}{}\begin{enonce*}{引理 1.3.18}\etiquette[1.3.18]{XVIII.lem1.3.18} 有 $$
\langle f,g\rangle =\prod_{x\in X_{1}\cup Y_{2}}(-1)^{v_{x}(f)v_{x}(g)}
(f^{v_{x}(g)}/g^{v_{x}(f)})(x){\quoi}.
$$ \end{enonce*}

证明（\hyperref[XVIII.lem1.3.18]{1.3.18}）推出（\hyperref[XVIII.par1.3.16.6]{1.3.16.6}）。可假设 $\Ll\sim \Oo(X)\sim \Oo(Y)$，其中 $X$ 和 $Y$ 有效，且 $X\cap Y=∅$\nde{前面已看到，可以假设 $\Ll$ 非常丰沛。}。设 $f$ 为满足 $\text{\rm div\,}(f)=Y-X$ 的有理函数\nde{注意勿将 $X$ 与同样记为 $X$ 的整体曲线混淆。}。交换图式 $$
\xymatrix@C=2.5cm@R=1.25cm{
\langle \Ll, \Ll\rangle  \ar[r]^{\tau} \ar@{=}[d]&
\langle \Ll,\Ll\rangle  \ar@{=}[d]\\
\langle \Oo(X), \Oo(Y)\rangle   \ar[r]^{\langle f^{-1},f\rangle } & \langle
\Oo(Y), \Oo(X)\rangle 
}
$$ 表明 $$
\tau=\langle f^{-1},f\rangle =\prod_{x\in
Y}(-1)^{-v_{x}(f)^{2}}=(-1)^{\deg(\Ll)}{\quoi}.
$$ 来证明 \hyperref[XVIII.lem1.3.18]{1.3.18}。归结为假设 $X_{i}$ 和 $Y_{i}$ 非常丰沛。此时存在分解 $f=f_{1}f_{2}$；$g=g_{1}g_{2}$，满足 \begin{alignat*}{2}
&\text{\rm div\,}(f_{1}) = X'-X_{1} &\qquad& \text{\rm div\,}(g_{1})
= Y'-Y_{1}\\
& \text{\rm div\,}(f_{2}) = X_{2}-X' &\qquad& \text{\rm div\,}(g_{2})
= Y_{2}-Y'{\quoi},
\end{alignat*}，其中除子 $X'$、$Y'$ 及 $X_{1}\cup X_{2}\cup Y_{1}\cup Y_{2}$ 两两不相交。

用 $\langle~,~\rangle^{\ast}$ 表示（\hyperref[XVIII.lem1.3.18]{1.3.18}）的右边，并证明 $$ \langle f,g\rangle ^{\ast}=\langle f_{1},g_{1}\rangle ^{\ast}\cdot
\langle f_{2},g_{2}\rangle ^{\ast}{\quoi}.
$$ 由 \textsc{Serre} \cite[Ch.~III, prop.~6, p.~44]{sga4.exp.XVIII.bib.6}，有 \begin{equation*}
\prod_{x}(-1)^{v_{x}(f_{1})v_{x}(g_{2})}
(f_{1}^{v_{x}(g_{2})}/g_{2}^{v_{x}(f_{1})}) (x)=1{\quoi}.\tag{1.3.18.1}
\end{equation*}\etiquette[1.3.18.1]{XVIII.eq1.3.18.1} 于是，
% SGA-SRC-001129 — exact continuation restored from the original scan
\begin{eqnarray*}
\langle f_{1},g_{1}\rangle ^{\ast}\cdot \langle f_{2},g_{2}\rangle
^{\ast} &= &g_{1}(X_{1})f_{1}(Y')g_{2}(X')f_{2}(Y_{2})\cdot\\
& & \prod_{X_{1}\cup X'\cup Y'\cup Y_{2}}(-1)^{v_{x}(f_{1})v_{x}(g_{2})}
(f_{1}^{v_{x}(g_{2})}\cdot g_{2}^{-v_{x}(f_{1})})(x)\\
&=&\prod_{x\in X'} g_{2}(x)^{-v_{x}(f_{2})}\cdot
g_{2}(x)^{-v_{x}(f_{1})}\cdot \prod_{x\in
Y'}f_{1}(x)^{v_{x}(g_{1})}\cdot f_{1}(x)^{v_{x}(g_{2})}\cdot\\
&&\prod_{x\in X_{1}\cup Y_{2}}(-1)^{v_{x}(f)v_{x}(g)}
((f_{1}f_{2})^{v_{x}(g_{2})}\cdot
(g_{1}g_{2})^{-v_{x}(f_{1})})(x)\\
&&=\langle f_{1}f_{2},g_{1}g_{2}\rangle
^{\ast}{\quoi}. 
\end{eqnarray*} 由 $\langle~,~\rangle $ 的函子性\pageoriginale 有 $\langle f,g\rangle =\langle f_{1},g_{1}\rangle \langle
f_{2},g_{2}\rangle $，且容易验证 $\langle
f_{i},g_{i}\rangle =\langle f_{i},g_{i}\rangle ^{\ast}$\nde{这是另行假设 $X_1\cap Y_2=∅$ 时的特殊情形；它由 \hyperref[XVIII.eq1.3.12.1]{1.3.12.1} 及 \hyperref[XVIII.eq1.3.12.2]{1.3.12.2} 得到。}。引理得证。

\subsection{严格交换的 Picard 范畴}\etiquette[1.4]{XVIII.subsec1.4}

为用一个“泛系数公式”表述 \hyperref[XVIII.subsec1.3]{1.3} 的结果，需要“普通拓扑”中的一些结果。

本节结果由 A.~\textsc{Grothendieck} 致 J.L. \textsc{Verdier} 的一封信（1966 年）提出。

双函子的结合性与交换性概念（见下文）由 \textsc{Mac Lane} 引入。

\subsubsection{}\etiquette[1.4.1]{XVIII.subsubsec1.4.1} 设 $C$ 为范畴，$F:C\times C\rightarrow C$ 为函子，$$
\sigma:F(F(X,Y),Z)\xrightarrow{\sim} F(X,F(Y,Z))
$$ 为三函子的一个同构。若下述条件成立，则称对 $(F,\sigma)$ 为一个\emph{结合函子}，或称 $\sigma$ 为 $F$ 上的一个\emph{结合性资料}：

\begin{itemize} \item[(Ass)] 对 $C$ 的每族对象 $(X_{i})_{i\in I}$，以 $e:I\rightarrow M(I)$ 表示从 $I$ 到由 $I$ 生成的自由幺半群（不含单位元）的典范映射，则存在映射 $ℱ:M(I)\rightarrow \Ob C$、同构 $a_{i}:ℱ(e_{i})\xrightarrow{\sim} X_{i}$ 及同构 $a_{g,h}:ℱ(gh)\xrightarrow{\sim}F(ℱ(g),ℱ(h))$，使图式 \begin{equation*}
\vcenter{
\xymatrix@C=1.25cm@R=1.25cm{
ℱ(f(gh))  \ar@{=}[d]\ar[r]^-{\sim}_-{a_{f,gh}} &
F(ℱ(f),ℱ(gh))
\ar[r]^-{\sim}_-{a_{g,h}}&
F(ℱ(f),F(ℱ(g),ℱ(h)))\\
ℱ((fg)h)  \ar[r]^-{\sim}_-{a_{fg,h}}& F(ℱ(fg),ℱ(h))
\ar[r]^-{\sim}_-{a_{f,g}}& F(F(ℱ(f),ℱ(g)),ℱ(h))\ar[u]_{\sigma}^{\text{\rotatebox{90}{$\backsim$}}}
}}\tag{1.4.1.1}
\end{equation*}\etiquette[1.4.1.1]{XVIII.eq1.4.1.1} 交换。

我们不需要以下事实：公理 (Ass) 等价于“五边形公理”，即以下图式交换： $$
\xymatrix@C=-.5cm{
                                  & F(F(X,Y),F(Z,T))&  \\
F(X,F(Y,F(Z,T))) \ar@{<->}[ru]    &                 & F(F(F(X,Y),Z),T)\ar@{<->}[lu] \\
F(X,F(F(Y,Z),T)) \ar@{<->}[rr]\ar@{<->}[u]    &    & F(F(X,F(Y,Z)),T). \ar@{<->}[u] \\
}
$$

现假设\pageoriginale 给定具有函子性的同构 \begin{align*}
\sigma_{X,Y,Z} &: F(F(X,Y),Z)\xrightarrow{\sim}F(X,F(Y,Z))\\
\tau_{X,Y} &: F(X,Y)\xrightarrow{\sim} F(Y,X){\quoi}.
\end{align*} 若下述条件成立，则称 $\sigma$ 及 $\tau$ 使 $F$ 成为一个\emph{结合且严格交换}的函子。

\item[(Ass.s.Com)] 对 $C$ 的每族对象 $(X_{i})_{i\in I}$，以 $e:I\rightarrow N(I)$ 表示从 $I$ 到由 $I$ 生成的自由阿贝尔幺半群（不含单位元）的典范映射，则存在映射 $ℱ:N(I)\rightarrow \Ob C$、同构 $a_{i}:ℱ(e_{i})\xrightarrow{\sim} X_{i}$ 及同构 $a_{g,h}:ℱ(gh)\xrightarrow{\sim}
F(ℱ(g),ℱ(h))$，使图式 \eqhyperref[XVIII.eq1.4.1.1]{1.4.1.1} 以及图式 \begin{equation*}
\vcenter{
\xymatrix@C=2.25cm{
ℱ(gh) \ar@{=}[d] \ar[r]^(.4){a_{g,h}}& F(ℱ(g),ℱ(h))\ar[d]^{\tau}\\
ℱ(hg)    \ar[r]^(.4){a_{h,g}}& F(ℱ(h),ℱ(g)).        
}}\tag{1.4.1.2}
\end{equation*}\etiquette[1.4.1.2]{XVIII.eq1.4.1.2} 均交换。 \end{itemize}

我们不需要以下事实：此公理等价于以下条件同时成立： \begin{enumerate} \renewcommand{\labelenumi}{\theenumi)} \item 五边形公理；

\item $\tau_{X,X}:X+X\rightarrow X+X$ 为恒等；

\item $\tau_{Y,X}\circ\tau_{X,Y}:X+Y\rightarrow Y+X\rightarrow X+Y$ 为恒等；

\item 六边形公理：图式 $$
\xymatrix{
          &     X + (Y + Z)                 &\\
 (X+Y)+Z  \ar@{-}[ru]_{\sigma} \ar@{-}[d]^{\tau}& & X+(Z+Y)\ar@{-}[lu]^{\tau}\\
Z+(X+Y) \ar@{-}[rd]_{\sigma}  &  & (X+Z) + Y \ar@{-}[u]_{\sigma}\\
          &    (Z+X)+Y    \ar@{-}[ru]_{\tau}                  &
}
$$ \pageoriginale 交换。 \end{enumerate}

注意，条件 2) 在实际中往往不满足（因此使用\emph{严格}交换的术语）。

\hypertarget{sga4.exp.XVIII.definition.1.4.2}{}\begin{enonce*}{定义 1.4.2}\etiquette[1.4.2]{XVIII.def1.4.2} 一个\emph{严格交换的 Picard 范畴 $\Pp$}，是指一个非空范畴，其中所有态射均为同构，赋有一个函子 $+:\Pp\times
\Pp\rightarrow \Pp:(X,Y)\mapsto X+Y$ 及具有函子性的同构 \begin{align*}
\sigma &: (X+Y)+Z\xrightarrow{\sim} X+(Y+Z)\\
\tau &: X+Y\longrightarrow Y+X
\end{align*}，使 $+$ 为结合且严格交换的函子，并且还要求：对每个 $X\in \Ob \Pp$，函子 $Y\mapsto X+Y$ 为范畴等价。 \end{enonce*}

读者可验证：

\hypertarget{sga4.exp.XVIII.lemma.1.4.3}{}\begin{enonce*}{引理 1.4.3}\etiquette[1.4.3]{XVIII.lem1.4.3} 若 $(X_{i})_{i\in I}$ 为一个严格交换的 Picard 范畴 $\Pp$ 的一族对象，则存在映射 $\Sigma:\ZZ^{(I)}\rightarrow \Ob \Pp$、同构 $a_{i}:\Sigma(e_{i})\xrightarrow{\sim} X_{i}$ 及 $a_{\underline{n},\underline{m}}:\Sigma(\underline{n}+\underline{m})
\xrightarrow{\sim}\Sigma(\underline{n})+\Sigma(\underline{m})$，使 \eqhyperref[XVIII.eq1.4.1.1]{1.4.1.1} 及 \eqhyperref[XVIII.eq1.4.1.2]{1.4.1.2} 类型的图式交换。系统 $(\Sigma,(a_{i}),(a_{\underline{n},\underline{m}}))$ 在唯一同构意义下唯一，且关于 $(X_{i})_{i\in I}$ 具有函子性。 \end{enonce*}

\setcounter{subsubsection}{3} \subsubsection{}\etiquette[1.4.4]{XVIII.subsubsec1.4.4} 严格交换的 Picard 范畴在唯一同构意义下恰有一个\emph{中性对象}，在此可将其定义为一个对 $(e,\varphi)$，由对象 $e$ 及同构 $\varphi:e+e\xrightarrow{\sim}e$ 组成。若 $(e,\varphi)$ 为中性对象，则存在唯一函子同构 $$
\alpha_{s}:e+X\xrightarrow{\sim} X
$$，使图式 $$
\xymatrix@C=2cm@R=1.25cm{
e + (e + X) \ar@{<->}[r] \ar[d]^{e+\alpha_s} & (e + e) + X \ar[d]^{\varphi}\\
e + X \ar@{=}[r]                       & e + X .
}
$$ 交换。

类似地\pageoriginale 定义 $\alpha_{d}:X+e\xrightarrow{\sim}X$，而 $\varphi$ 同时是 $\alpha_{s}$ 和 $\alpha_{d}$ 的特殊情形；$\tau$ 将 $\alpha_{s}$ 与 $\alpha_{d}$ 互换。群 $\Aut(e)$ 是阿贝尔群，且对每个 $X\in \Ob\Pp$，函子 $X+$ 及 $+X$ 在 $\Aut(e)$ 与 $\Aut(X)$ 之间诱导同一个同构。

\hypertarget{sga4.exp.XVIII.definition.1.4.5}{}\begin{enonce*}{定义 1.4.5}\etiquette[1.4.5]{XVIII.def1.4.5} 景 $\Ss$ \emph{上的严格交换 Picard 叠 $\Pp$}，是指 $\Ss$ 上的一个群胚叠 $\Pp$（\textsc{Giraud} \cite{sga4.exp.XVIII.bib.1}），赋有函子 $+:\fprod{\Pp}{\Pp}{\Ss}\rightarrow \Pp$ 及函子同构 \begin{align*}
\tau_{x,y} &: x+y\longrightarrow y+x\\
\sigma_{x,y,z} &: (x+y)+z\longrightarrow x+(y+z)
\end{align*}，它们对每个 $U\in \Ob \Ss$ 均使 $\Pp(U)$ 成为严格交换的 Picard 范畴。 \end{enonce*}

下文只说 \emph{Picard 叠}，而不指明“严格交换”。由 \hyperref[XVIII.subsubsec1.4.4]{1.4.4}，每个 Picard 叠 $\Pp$ 均有一个整体的“中性对象” $e$，而 $\sheafAut(e)$ 为阿贝尔层。

\setcounter{subsubsection}{5} \subsubsection{}\etiquette[1.4.6]{XVIII.subsubsec1.4.6} $\Ss$ 上的 Picard 叠之间的一个\emph{可加函子} $F:\Pp_{1}\rightarrow
\Pp_{2}$，是一个 $\Ss$-函子（必为笛卡尔的），赋有一个函子同构 $$
F(x+y)\xrightarrow{\sim} F(x)+F(y)
$$，使图式 $$
\xymatrix@C=2.25cm@R=1.5cm{
F(x+y)  \ar[r] \ar[d]^{F(\tau)}& F(x) + F(y) \ar[d]^{\tau}\\
F(y+x)         \ar[r] & F(y) + F(x)
}
$$ 及 $$
\xymatrix@R=1.5cm{
  F((x+y)+z) \ar[r]\ar[d]^{F(\sigma)} & F(x+y)+F(z) \ar[r]&(F(x) + F(y)) + F(z)       \ar[d]^{\sigma}  &
\\
F(x + (y + z))    \ar[r]        & F(x) + F(y +z)  \ar[r] & F(x) + (F(y)
+ F(z)). 
}
$$ 交换。

一个\pageoriginale\emph{可加函子的态射} $u:F\rightarrow G$，是一个使图式 $$\vcenter{
\xymatrix@C=2cm{
F(x+y)       \ar[r]^{u_{x+y}} \ar[d]       & G(x+y) \ar[d]\\
F(x) + F(y)  \ar[r]^{u_x + u_y}     & G(x) + G(y). 
}}
$$ 交换的 $\Ss$-函子的态射（自动为同构）。

\subsubsection{}\etiquette[1.4.7]{XVIII.subsubsec1.4.7} 若 $\Pp_{1}$ 和 $\Pp_{2}$ 为 $\Ss$ 上的两个 Picard 叠，则 Picard 叠 $\HOM(\Pp_{1},\Pp_{2})$ 定义如下： \begin{enumerate} \renewcommand{\theenumi}{\alph{enumi}} \renewcommand{\labelenumi}{\theenumi)} \item 若 $U\in \Ob \Ss$，其在 $U$ 上的对象为从 $\Pp_{1}|U$ 到 $\Pp_{2}|U$ 的可加函子；其态射为可加函子的态射；

\item 两个可加函子 $F_{1}$ 与 $F_{2}$ 的和由公式 $(F_{1}+F_{2})(x)=F_{1}(x)+F_{2}(x)$ 定义；其结构同构为使图式 $$
\xymatrix@C=1cm{
\scriptstyle{(F_1 + F_2)(x+y)} \ar[r]^{\rm{\scriptstyle{\text{结构}\qquad}}}\ar@{=}[dd]& \scriptstyle{(F_1 + F_2)(x) + (F_1 + F_2)(y)}\ar@{=}[rd]& \\
                                 & &\scriptstyle{F_1(x) + F_2(x) +  F_1(y)+F_2(y)}\\
\scriptstyle{F_1(x+y)+F_2(x+y)}\ar[rd]                    & &\\
&\scriptstyle{F_1(x) + F_1(y) + F_2(x) + F_2(y)}\ar@{<->}[ruu]_{\tau}&
}
$$ 交换的同构；

\item 结合及交换同构通过 $\Pp_{2}$ 中的类似同构定义。 \end{enumerate}

\subsubsection{}\etiquette[1.4.8]{XVIII.subsubsec1.4.8} 若 $\Pp_{1}$、$\Pp_{2}$ 和 $\Pp$ 为三个 Picard 叠，则从 $\Pp_{1}\times
\Pp_{2}$ 到 $\Pp$ 的一个双可加函子，是一个从 $\fprod{\Pp_{1}}{\Pp_{2}}{\Ss}$ 到 $\Pp$ 的 $\Ss$-函子 $F$，赋有函子同构 \begin{align*}
F(x_{1}+y_{1},x_{2}) &\xrightarrow{\sim} F(x_{1},
x_{2})+F(y_{1},x_{2})\\
F(x_{1}, x_{2}+y_{2}) &\xrightarrow{\sim} F(x_{1}, x_{2})+F(x_{1},
y_{2}){\quoi},
\end{align*}

使得

\begin{enumerate} \renewcommand{\theenumi}{\alph{enumi}} \renewcommand{\labelenumi}{\theenumi)} \item 固定 $x_{1}$（\resp $x_{2}$）时，$F(x_{1},\ast)$（\resp $F(\ast,x_{2})$）满足 \hyperref[XVIII.subsubsec1.4.6]{1.4.6} 的相容性条件。

\item 对\pageoriginale $U\in \Ob \Ss$、$x_{1}$、$y_{1}\in
\Ob \Pp_{1}(U)$ 及 $x_{2}$、$y_{2}\in \Ob \Pp_{2}(U)$，图式 $$
\xymatrix@C=0cm@R=.5cm{
\scriptscriptstyle{F(x_1+y_1, x_2+y_2)}\ar[dd] \ar[r] &\scriptscriptstyle{F(x_1+y_1, x_2) + F(x_1+y_1,y_2)}\ar[rd]\\
 & &\scriptscriptstyle{F(x_1,x_2) + F(y_1,x_2) +F(x_1,y_2)+F(y_1,y_2)}\\
\scriptscriptstyle{F(x_1,x_2+y_2) + F(y_1,x_2+y_2)}\ar[rd] \\
& \scriptscriptstyle{F(x_1,x_2)+F(x_1,y_2) + F(y_1,x_2)+F(y_1,y_2)}\ar@{<->}[ruu]_{\tau}
}
$$ 交换。 \end{enumerate}

例如，典范函子 $$
\HOM(\Pp_{1},\Pp_{2})\times
\Pp_{1}\longrightarrow \Pp_{2}:(F,x)\longmapsto F(x)
$$ 是双可加的。

如同 \hyperref[XVIII.subsubsec1.4.7]{1.4.7}，可证明从 $\Pp_{1}\times \Pp_{2}$ 到 $\Pp$ 的双可加函子构成一个 Picard 叠 $\HOM(\Pp_{1}, \Pp_{2};\Pp)$。

将在（\hyperref[XVIII.const1.4.20]{1.4.20}）中看到，存在一个“$2$-泛”的双可加函子：$\Pp_{1}\times \Pp_{2}\rightarrow
\Pp_{1}\otimes \Pp_{2}$；更精确地说，存在一个 Picard 叠 $\Pp_{1}\otimes \Pp_{2}$ 及双可加函子 $\otimes:\Pp_{1}\times
\Pp_{2}\rightarrow \Pp_{1}\otimes \Pp_{2}$，使对每个 Picard 叠 $\Pp$，由 $\otimes$ 定义的函子 \begin{equation*}
\HOM(\Pp_{1}\otimes \Pp_{2},\Pp)\longrightarrow \HOM (\Pp_{1},\Pp_{2};\Pp)\tag{1.4.8.1}
\end{equation*}\etiquette[1.4.8.1]{XVIII.eq1.4.8.1} 为等价。此叠在唯一等价意义下唯一（该等价本身又在唯一同构意义下唯一）。

\subsubsection{}\etiquette[1.4.9]{XVIII.subsubsec1.4.9} 设 $u:\Ss_{1}\rightarrow \Ss_{2}$ 为景的态射。$\Ss_{1}$ 上的 Picard 叠 $\Pp$ 的直接像 $u_{\ast}\Pp$ 为 $\Ss_{2}$ 上由 $$
u_{\ast}\Pp(V)=\Pp(u^{\ast}V){\quoi}.
$$ 定义的 Picard 叠。其在 $V\in \Ob \Ss_{2}$ 上的对象为 $\Pp$ 在 $u^{\ast}V$ 上的对象；其态射、加法及相容性同构均取自 $\Pp$。直接像 $u_{\ast}$ 为一个 $2$-函子。

\subsubsection{}\etiquette[1.4.10]{XVIII.subsubsec1.4.10} 景 $\Ss$ 上的一个\pageoriginale\emph{Picard 预叠} $\Pp$，是 $\Ss$ 上的一个预叠 $\Pp$（\textsc{Giraud} \cite{sga4.exp.XVIII.bib.1}），赋有函子 $+:\fprod{\Pp}{\Pp}{\Ss}\rightarrow \Pp$ 及结合与交换同构 $\sigma$ 和 $\tau$，使对每个 $U\in \Ob \Ss$，$\Pp(U)$ 为一个（严格交换的）Picard 范畴。若 $\Pp_{1}$ 和 $\Pp_{2}$ 为两个 Picard 预叠，则以显然的方式（\cf \hyperref[XVIII.subsubsec1.4.6]{1.4.6}，\hyperref[XVIII.subsubsec1.4.7]{1.4.7}）定义从 $\Pp_{1}$ 到 $\Pp_{2}$ 的\emph{可加函子}，以及它们构成的 Picard 预叠 $\HOM(\Pp_{1}, \Pp_{2})$。若 $\Pp$ 为 Picard 预叠，$j:\Pp\rightarrow
a\Pp$ 为其生成的叠，则在本质唯一的等价意义下，存在唯一的一对结构，由 $a\Pp$ 上的 Picard 叠结构及 $j$ 上的可加函子结构组成；赋予这些资料后，称对 $(j,a\Pp)$ 为 $\Pp$ \emph{生成的 Picard 叠}。对每个 Picard 叠 $\Pp_{1}$，有等价 \begin{equation*}
\HOM(a\Pp,\Pp_{1})\xrightarrow{\sim} \HOM(\Pp,\Pp_{1}){\quoi}.\tag{1.4.10.1}
\end{equation*}\etiquette[1.4.10.1]{XVIII.eq1.4.10.1}

\subsubsection{}\etiquette[1.4.11]{XVIII.subsubsec1.4.11} 用 $C^{[-1,0]}(\Ss)$ 表示 $\Ss$ 上的阿贝尔层复形 $K$ 所成的范畴，其中当 $i\not\in [-1,0]$ 时 $K^{i}=0$。对每个复形 $K\in \Ob C^{[-1,0]}(\Ss)$ $$
K:d:K^{-1}\longrightarrow K^{0}
$$，对应以下 Picard 预叠 $\pch(K)$：

\begin{enumerate} \renewcommand{\theenumi}{\Roman{enumi}} \renewcommand{\labelenumi}{(\theenumi)} \item 对 $U\in \Ob \Ss$，有 $\Ob \pch(K)(U)=K^{0}(U)$。

\item 若 $x$、$y\in K^{0}(U)$，则从 $x$ 到 y 的箭头是 $K^{-1}(U)$ 的一个元素 $f$，满足 $df=y-x$。

\item 箭头的合成律为 $K^{-1}(U)$ 中的加法。

\item 函子 $+$ 由 $K^{0}(U)$ 和 $K^{-1}(U)$ 中的加法给出。

\item 结合及交换同构由 $K^{-1}(U)$ 的零元给出。 \end{enumerate}

用 $\ch(K)$ 表示 Picard 预叠 $\pch(K)$ 生成的 Picard 叠（\hyperref[XVIII.subsubsec1.4.10]{1.4.10}）。层 $ℋ^{0}(K)$ 解释为以下伴随层 \begin{align*}
ℋ^{0}(K) & \simeq a\; \text{\rm ($U\longmapsto$ 范畴}\tag{1.4.11.1}\\
&\qquad\qquad \text{\rm $\ch(K)(U)$ 的对象同构类的群)}{\quoi}.
\end{align*}\etiquette[1.4.11.1]{XVIII.eq1.4.11.1}

另\pageoriginale 一方面， \begin{equation*}
ℋ^{-1}(K)\simeq \sheafAut(e)\tag{1.4.11.2}
\end{equation*}\etiquette[1.4.11.2]{XVIII.eq1.4.11.2}

\subsubsection{}\etiquette[1.4.12]{XVIII.subsubsec1.4.12} 复形的每个态射 $f:K\rightarrow L$ 诱导一个可加函子 $\pch(f):\pch(K)\rightarrow \pch(L)$，从而诱导一个可加函子 $$
\ch(f):\ch(K)\longrightarrow \ch(L){\quoi}.
$$

由 \eqhyperref[XVIII.eq1.4.11.1]{1.4.11.1} 及 \eqhyperref[XVIII.eq1.4.11.2]{1.4.11.2} 可知，$\ch(f)$ 为等价当且仅当 $f$ 为拟同构。由 \eqhyperref[XVIII.eq1.4.10.1]{1.4.10.1} 可推出，对两个态射 $f$、$g:K\rightarrow L$，有 $$
\Hom(\ch(f),\ch(g))\xrightarrow{\sim} \Hom(\pch(f),\pch(g)){\quoi}.
$$

一个\emph{函子态射} $h:\pch(f)\rightarrow \pch(g)$ 是一个层态射 $h:K^{0}\rightarrow L^{-1}$，满足 \begin{equation*}
g(x)-f(x)=dh(x)\tag{1.4.12.1}
\end{equation*}\etiquette[1.4.12.1]{XVIII.eq1.4.12.1} 并且对每个满足 $y-x=du$ 的三元组 $(x,y,u)$，有 \begin{equation*}
h(y)+f(u)=g(u)+h(x){\quoi}\tag{1.4.12.2}
\end{equation*}\etiquette[1.4.12.2]{XVIII.eq1.4.12.2}

$h$ 为可加函子的态射，意味着 \begin{equation*}
h(x+y)=h(x)+h(y){\quoi}.\tag{1.4.12.3}
\end{equation*}\etiquette[1.4.12.3]{XVIII.eq1.4.12.3} 条件 \eqhyperref[XVIII.eq1.4.12.3]{1.4.12.3} 使得可将 \eqhyperref[XVIII.eq1.4.12.2]{1.4.12.2} 改写为 \begin{equation*}
g(u)-f(u)=h\sisi{du}{(du)} ;\tag{$1.4.12.2'$}
\NextLinkTarget{sga4.exp.XVIII.equation.1.4.12.2-prime}\label{1.4.12.2'}
\end{equation*} 因此，条件 \eqhyperref[XVIII.eq1.4.12.1]{1.4.12.1} 至 \eqhyperref[XVIII.eq1.4.12.3]{1.4.12.3} 表明 $h$ 为从 $f$ 到 $g$ 的一个同伦，且 \begin{equation*}
\Hom(\ch(f),\ch(g))\simeq \{H:\;\text{\rm 同伦 } K\longrightarrow
L| g-f=dH+Hd\}\tag{1.4.12.4}
\end{equation*}\etiquette[1.4.12.4]{XVIII.eq1.4.12.4}

\begin{enonce*}{引理 1.4.13}\etiquette[1.4.13]{XVIII.lem1.4.13} \begin{enumerate} \renewcommand{\theenumi}{\Roman{enumi}} \renewcommand{\labelenumi}{(\theenumi)} \item 对每个 Picard 叠 $\Pp$，存在一个复形 $K$，使 $\Pp=\ch(K)$。

\item 对每个可加函子 $F:\ch(K)\rightarrow \ch(L)$，存在一个拟同构 $k:K'\rightarrow K$ 及一个态射 $\ell:K'\rightarrow L$，使 $F\simeq \ch(\ell)\ch(k)^{-1}$。 \end{enumerate} \end{enonce*}

来证明\pageoriginale (I)。设 $(k_{i},U_{i})_{i\in I}$ 为一个族，使 \begin{enumerate} \renewcommand{\theenumi}{\alph{enumi}} \renewcommand{\labelenumi}{(\theenumi)} \item $k_{i}\in \Ob \Pp(U_{i})$；

\item $\Pp$ 的每个对象均局部同构于某个 $k_{i}$ 的逆像。

置 $K^{0}=\displaystyle{\mathop{\oplus}_{i\in
I}}\ZZ_{U_{i}}$。由 \hyperref[XVIII.lem1.4.3]{1.4.3}，存在一个从 $\ch(0\rightarrow K^{0})$ 到 $\Pp$ 的可加函子 $F$，将 $\ZZ_{U_{i}}$ 的基底 $e_{i}$ 映到 $k_{i}$。设 $K^{-1}$ 为由 $K^{0}$ 的一个局部截面 $x$ 及同构 $t:F(0)\xrightarrow{\sim} F(x)$ 组成的对的层。$K^{-1}$ 上的加法定义为 $(x_{1},
t_{1})+(x_{2},t_{2})=(x_{1}+x_{2},t)$，其中 $t$ 使图式 $$
\xymatrix@C=1.5cm{
F(0) + F(0)    \ar[r]^{t_1+t_2} \ar@{<->}[d]   & F(x_1) + F(x_2)\ar@{<->}[d]\\
F(0+0)          \ar[r]^{t}         & F(x_1+x_2).
}
$$ 交换。用 $d(x,t)=x$ 定义 $d:K^{-1}\rightarrow K^{0}$；$d$ 是可加的。

若 $y-x=dt$，则存在唯一态射 $F(t)$ 使图式 $$
\xymatrix@C=1.75cm{
F(x)           \ar[r]^{F(t)}   \ar@{<->}[d]& F(y) \ar@{<->}[d]\\
F(0)+F(x)      \ar[r]^{t+F(x)} & F(y-x)+F(x)
}
$$ 交换；并可验证此构造定义一个等价 $$
F:\ch(K)\longrightarrow \Pp{\quoi}.
$$

来证明 (II)。设 $F:\ch(K)\rightarrow \ch(L)$ 为一个可加函子。则存在一个族 $(k_{i},\ell_{i}, \alpha_{i},
U_{i})_{i\in I}$，使得

\item $k_{i}\in K^{0}(U_{i})$、$\ell_{i}\in L^{0}(U_{i})$，且 $\alpha_{i}$ 为从 $F(k_{i})$ 到 $\ell_{i}$ 的态射；

\item 坐标为 $k_{i}$ 的同态 $k^{0}:\displaystyle{\mathop{\oplus}_{i\in
I}}\ZZ_{U_{i}}\rightarrow K^{0}$ 为满态射。

置 ${K'}^{0}=\displaystyle{\mathop{\oplus}_{i\in
I}}\ZZ_{U_{i}}$，并设 $\ell^{0}:{K'}^{0}\rightarrow L^{0}$ 为坐标为 $\ell_{i}$ 的同态。由于函子 $F$ 可加，存在唯一的一族同构 $\alpha_{x}:Fk^{0}(x)\xrightarrow{\sim}\ell^{0}(x)$（$x$ 为 ${K'}^{0}$ 的局部截面），使得

\item 对 $\ZZ_{U_{i}}$ 的 \hyperref[XVIII.sec1]{截面 1} $e_{i}$，有 $\alpha_{e_{i}}=\alpha_{i}$；

\item 以下\pageoriginale 图式 $$
\xymatrix@C=2cm{
Fk^0(x+y)  \ar[r]^{\alpha_{x+y}} \ar@{<->}[d]& \ell^0(x+y)
\ar@{=}[d]\\
Fk^0(x)+ Fk^0(y)    \ar[r]^{\alpha_{x} + \alpha_{y}}     & \ell^0(x)+\ell^0(y)
}
$$ 交换。

设 ${K'}^{-1}=\fprod{K^{-1}}{{K'}^{0}}{K^{0}}$；由 (d)，显然的箭头为一个复形的拟同构 $k:K'\rightarrow K$。若 $x$、$y\in {K'}^{0}(U)$，且 $t\in
{K'}^{-1}(U)$ 满足 $y-x=dt$，则存在唯一的 $L^{-1}(U)$ 的截面 $\ell^{-1}(t;x,y)$，定义使图式

\item $\begin{array}[t]{c}
\xymatrix@C=2.75cm@R=1.75cm{
&F(k^0(x))  \ar[r]^{F(k^{-1}(t))} \ar[d]^{\alpha_x}&
  F(k^0(y))\ar[d]^{\alpha_y} \\ 
&\ell^0(x)     \ar[r]^{\ell^{-1}(t;x,y)}& \ell^0(y).
}
\end{array}$

交换的态射。$F$ 为函子这一事实给出

\item $\ell^{-1}(t;x,y)+\ell^{-1}(t';y,z)=\ell^{-1}(t+t';x,z)$。

$F$ 的可加同构的函子性给出

\item $\ell^{-1}(t_{1}+t_{2};x_{1}+x_{2},
y_{1}+y_{2})=\ell^{-1}(t_{1};x_{1},y_{1})+\ell^{-1}(t_{2};x_{2},y_{2})$。

由 (h) 得到 \begin{align*}
&\ell^{-1}(0;x,x)+\ell^{-1}(0;x,x)=\ell^{-1}(0;x,x)\qquad \text{\ie}\\
& \ell^{-1}(0;x,x)=0{\quoi}.
\end{align*} 再由 (i) 可知 $\ell^{-1}(t;x,y)$ 仅依赖于 $t$：$$
\ell^{-1}(t;x,y)=\ell^{-1}(t)
$$，且 $\ell^{-1}(t)$ 关于 $t$ 可加。因此，$\ell=(\ell^{-1}, \ell^{0})$ 为复形的态射，而由 (f) 和 (g)，各箭头 $\alpha$ 构成一个可加函子的同构 $F\circ \ch(k)=\ch(\ell)$。 \end{enumerate}

\setcounter{subsubsection}{13} \subsubsection{}\etiquette[1.4.14]{XVIII.subsubsec1.4.14} 设 $\Ch^{\flat}(\Ss)$ 为如下范畴：对象为 $\Ss$ 上的小 Picard 叠，箭头为可加\pageoriginale 函子的同构类。构造 $\ch$ 定义一个从 $C^{[-1,0]}(\Ss)$（\hyperref[XVIII.subsubsec1.4.11]{1.4.11}）到 $\Ch^{\flat}(\Ss)$ 的函子。设 $D^{[-1,0]}(\Ss)$ 为 $\Ss$ 的导出范畴中由复形 $K$ 组成的子范畴，其中当 $i\neq 0$ 或 $-1$ 时有 $H^{i}(K)=0$。由 \hyperref[XVIII.subsubsec1.4.12]{1.4.12} 及 \hyperref[XVIII.lem1.4.13]{1.4.13} 可知

\begin{enonce*}{命题 1.4.15}\etiquette[1.4.15]{XVIII.prop1.4.15} 函子 $\ch$ 诱导一个范畴等价 $$
\ch:D^{[-1,0]}(\Ss)\xrightarrow{\sim}
\Ch^{\flat}(\Ss){\quoi}.
$$ \end{enonce*}

用 $\flat$ 表示 $\ch$ 的逆等价。对 $\Ss$ 上的每个 Picard 叠 $\Pp$，$\Pp^{\flat}\in \Ob D^{[-1,0]}(\Ss)$ 在相差一个等价的同构类的意义下确定 $\Pp$。

\begin{enonce*}{引理 1.4.16}\etiquette[1.4.16]{XVIII.lem1.4.16} 设 $L\in \Ob C^{[-1,0]}(\Ss)$ 满足 $L^{-1}$ 内射。 \begin{enumerate} \renewcommand{\theenumi}{\Roman{enumi}} \renewcommand{\labelenumi}{(\theenumi)} \item $\pch(L)$ 已经是叠。

\item 对 $K\in \Ob C^{[-1,0]}(\Ss)$，每个可加函子 $F:\ch(K)\rightarrow \ch(L)$ 均同构于某个函子 $\ch(f)$，其中 $f$ 为从 $K$ 到 $L$ 的态射。 \end{enumerate} \end{enonce*}

\begin{enumerate} \renewcommand{\theenumi}{\Roman{enumi}} \renewcommand{\labelenumi}{(\theenumi)} \item 对任意 $U\in \Ob \Ss$ 和 $x\in \Ob \ch(L)(U)$，由 $L^{0}|U$ 的一个局部截面 $\ell$ 及 $x$ 与 $\ell$ 之间的一个同构组成的对的层，为 $L^{-1}$ 下的主齐性空间；因而有一个截面，且 $j:\pch(L)\rightarrow \ch(L)$ 在对象的同构类上满射，从而为等价。

\item 设 $L'$ 为内射对象组成的复形，赋有拟同构 $\varphi:L{\xrightarrow[\approx]{}}L'$，使 $i\leq -1$ 时 $\varphi^{i}$ 为同构。由 \hyperref[XVIII.lem1.4.13]{1.4.13} 及 \cite{sga4.exp.XVIII.bib.9}，存在一个在相差同伦的意义下交换的图式 $$
\xymatrix@R=1cm@C=1.5cm{
L  \ar[r]^{\varphi}_{\approx} & L'\\
K' \ar[u]^{f_0}  \ar[r]^{\Psi}_{\approx}   & K\ar[u]_{f_1}
}
$$，使 $F\simeq \ch(f_{0})\ch(\Psi)^{-1}$。此外，$K=\tau_{\leq
0}K$ 且 $L=\tau_{\leq 0}L'$，故 $f_{1}$ 通过 $f:K\rightarrow L$ 分解，且 $F\simeq \ch(f)$。 \end{enumerate}

此引理通过 \hyperref[XVIII.subsubsec1.4.12]{1.4.12} 的以下推论细化了 \hyperref[XVIII.prop1.4.15]{1.4.15}。

\begin{enonce*}{推论 1.4.17}\etiquette[1.4.17]{XVIII.cor1.4.17} 构造\pageoriginale $\ch$ 定义以下两者之间的 $2$-范畴等价： \begin{enumerate} \renewcommand{\theenumi}{\alph{enumi}} \renewcommand{\labelenumi}{\theenumi)} \item $\Ss$ 上的 Picard 叠的 $2$-范畴。

\item 如下 $2$-范畴：其对象和 $1$-箭头为 $C^{[-1,0]}(\Ss)$ 的由满足 $L^{-1}$ 内射的复形 $L$ 组成的满子范畴的对象和箭头，其 $2$-箭头为箭头之间的同伦：$\Hom(f,g)=\{H|g-f=\text{\rm dH}+\text{\rm Hd}\}$。 \end{enumerate} \end{enonce*}

\begin{enonce*}[remark]{构造 1.4.18}\etiquette[1.4.18]{XVIII.const1.4.18} 有 $$
\HOM(\Pp_{1}, \Pp_{2})^{\flat}\simeq \tau_{\leq
0}R\Hhom (\Pp^{\flat}_{1}, \Pp^{\flat}_{2}){\quoi}.
$$ \end{enonce*}

设 $K$、$L\in \Ob C^{[-1,0]}(\Ss)$，且 $L^{-1}$ 内射。由 \hyperref[XVIII.lem1.4.16]{1.4.16} (II) 及 \hyperref[XVIII.subsubsec1.4.12]{1.4.12}，有 \begin{equation*}
\ch(\tau_{\leq
0}\Hhom(K,L))\xrightarrow{\sim}\HOM(\ch(K),\ch(L))\tag{1.4.18.1}
\end{equation*}\etiquette[1.4.18.1]{XVIII.eq1.4.18.1} 由此以函子的方式推出构造 \hyperref[XVIII.const1.4.18]{1.4.18}。

\begin{enonce*}[remark]{构造 1.4.19}\etiquette[1.4.19]{XVIII.const1.4.19} 对景的态射 $f:\Ss_{1}\to\Ss_{2}$，有 $$
(f_{\ast}\Pp)^{\flat}\simeq \tau_{\leq
0}\,Rf_{\ast}(\Pp^{\flat}){\quoi}.
$$ \end{enonce*}

设 $L\in \Ob C^{[-1,0]}$，且 $L^{-1}$ 内射。由 \hyperref[XVIII.lem1.4.16]{1.4.16} (I)，有

\begin{equation*}
\ch(f_{\ast}L)\xrightarrow{\sim} f_{\ast}\ch(L)\tag{1.4.19.1}
\end{equation*}\etiquette[1.4.19.1]{XVIII.eq1.4.19.1} 由此推出具有函子性的构造 \hyperref[XVIII.const1.4.19]{1.4.19}。

我们不需要使用以下结果

\begin{enonce*}[remark]{构造 1.4.20}\etiquette[1.4.20]{XVIII.const1.4.20} 在 \hyperref[XVIII.subsubsec1.4.8]{1.4.8} 中所允诺的张量积存在，且 $$
(\Pp_{1}\otimes \Pp_{2})^{\flat}\simeq \tau_{\geq
-1}\Pp^{\flat}_{1}\mathop{\otimes}^{\LL}\Pp^{\flat}_{2}{\quoi}.
$$ \end{enonce*}

如同 \hyperref[XVIII.subsubsec1.4.12]{1.4.12} 及 \hyperref[XVIII.lem1.4.13]{1.4.13} (II)，可验证：对 $C^{[-1,0]}(\Ss)$ 中所有 $K_{1}$、$K_{2}$ 和 $L$，有 \begin{enumerate} \renewcommand{\theenumi}{\alph{enumi}} \renewcommand{\labelenumi}{\theenumi)} \item 每个态射 $f\in \Hom(K_{1}\otimes K_{2},L)\simeq
\Hom(\tau_{\geq-1}(K_{1}\otimes K_{2}),L)$ 定义一个双可加函子 $\ch(f)\in \Hom(\ch(K_{1})$、$\ch(K_{2})$；$\ch(L)$）；此外，函子的态射可与同伦等同；

\item 对每个双可加函子 $F:\ch(K_{1})$、$\ch(K_{2})\rightarrow \ch(L)$，存在拟同构 $k_{i}:K'_{i}\rightarrow K_{i}$ 和 $f:K'_{1}\otimes
K'_{2}\rightarrow L$，使 $F\simeq \ch(f)\circ
(\ch(k_{1})^{-1},\ch(k_{2})^{-1})$。 \end{enumerate}

若\pageoriginale 此时 $K^0_{1}$ 或 $K^0_{2}$ 平坦，且 ${K'}^0_{1}$ 或 ${K'}^0_{2}$ 也平坦，则 $k_{1}\otimes k_{2}:\tau_{\geq -1}{K'}_{1}\otimes {K'}_{2}\rightarrow
\tau_{\geq -1}K_{1}\otimes K_{2}$ 为拟同构，从而 $\ch(\tau_{\geq -1}(K_{1}\otimes K_{2}))$ 满足泛性质 \eqhyperref[XVIII.eq1.4.8.1]{1.4.8.1}，于是得到 $\otimes$ 的存在性及具有函子性的构造 \hyperref[XVIII.const1.4.20]{1.4.20}。

\setcounter{subsubsection}{20} \subsubsection{}\etiquette[1.4.21]{XVIII.subsubsec1.4.21} 设 $G$ 为阿贝尔层，$G[1]$ 为将 $G$ 置于次数 $-1$ 的复形。预叠 $\pch(G[1])$ 可与 $G$ 下的平凡挠子的预叠等同，因此 $\ch(G[1])$ “正是” $G$ 下的挠子的叠，赋予其通常的加法。

\subsubsection{}\etiquette[1.4.22]{XVIII.subsubsec1.4.22} 设 $K\in C^{[-1,0]}(\Ss)$，并设 $G$ 为阿贝尔层。$K$ 被 $G$ 扩张的各扩张 $E$，在
% SGA-SRC-001136
Baer 加法下构成一个 Picard 叠 $\EXT(K,G)$。 $$
E:0\rightarrow
G[0]\xrightarrow{\alpha}E\xrightarrow{\beta}K\rightarrow 0{\quoi}.
$$

对每个扩张 $E$，对应一个从 $\ch(K)$ 到 $\ch(G[1])$ 的可加函子，它将 $K^{0}$ 的局部截面 $x$ 对应到挠子 $\beta^{-1}(x)$。读者可验证下述结果。

\begin{enonce*}{命题 1.4.23}\etiquette[1.4.23]{XVIII.prop1.4.23} 上述构造为一个 Picard 叠的等价 $$
\EXT(K,G)\xrightarrow{\sim} \HOM(\ch(K),\ch(G[1])){\quoi}.
$$ \end{enonce*}

\setcounter{subsubsection}{23} \subsubsection{}\etiquette[1.4.24]{XVIII.subsubsec1.4.24} 设 $F$ 为景 $\Ss$ 的一个对象。也用 $F$ 表示由 $F$ 表示的层，用 $\ZZ^{(F)}$ 表示它生成的阿贝尔层。设 $C$ 为 $\Ss$ 上的 Picard 叠。容易验证（\cf \hyperref[XVIII.lem1.4.3]{1.4.3}），给出以下任何一种资料，均等价于给出其他资料： \begin{enumerate} \renewcommand{\theenumi}{\alph{enumi}} \renewcommand{\labelenumi}{\theenumi)} \item 一个可加函子 $H:\ch(\ZZ^{(F)})\rightarrow C$；

\item 一个\emph{叠的态射}：

（$F$ 的局部截面的叠，仅有恒等态射）$\rightarrow C$；

\item $C(F)$ 的一个对象。 \end{enumerate}

假设 $\Ss$ 中存在纤维积（以便应用定义 \hyperref[XVIII.subsubsec1.4.9]{1.4.9}），并设 $f$ 为景的典范态射 $f:\Ss/F\rightarrow \Ss$。上述构造给出一个等价 \begin{equation*}
\HOM(\ch(\ZZ^{(F)}),C)\xrightarrow{\sim}
f_{\ast}f^{\ast}C{\quoi}.\tag{1.4.24.1}
\end{equation*}\etiquette[1.4.24.1]{XVIII.eq1.4.24.1}

由 \eqhyperref[XVIII.eq1.4.24.1]{1.4.24.1} 通过构造 $$
\tau_{\leq 0} R\Hhom(\ZZ^{(F)},C^{\flat})\xrightarrow{\sim}
\tau_{\leq 0} Rf_{\ast} f^{\ast}C^{\flat}{\quoi},
$$ 得到的同构\pageoriginale，也可由更一般的同构 \begin{equation*}
R\Hhom(\ZZ^{(F)},K)\xrightarrow{\sim}
Rf_{\ast}f^{\ast}K\quad (K\in \Ob D^{+}(\Ss)){\quoi}.\tag{1.4.24.2}
\end{equation*}\etiquette[1.4.24.2]{XVIII.eq1.4.24.2}
推出。


% END native input: chapters/exp-18-part2-zh.tex
% BEGIN native input: chapters/exp-18-part3-zh.tex
\subsection{泛系数公式}\etiquette[1.5]{XVIII.subsec1.5}

\subsubsection{}\etiquette[1.5.1]{XVIII.subsubsec1.5.1}
设 $f:X\rightarrow S$ 是 $S$ 上的光滑射影曲线。以 $\PIC(X/S)$ 表示 $S$ 的大 fppf 景（\hyperrefext[XVII][XVII.subsec0.10]{0.10}）上的 Picard 叠，即 $X$ 上可逆层叠的正像（\hyperref[XVIII.subsubsec1.4.9]{1.4.9}）。有
\begin{align*}
\PIC(X/S)&= f_{\ast}\ch\GG_{m}[1]\tag{1.5.1.1}\\
\PIC(X/S)^{\flat} &=\tau_{\leq 0}
Rf_{\ast}(\GG_{m}[1]){\quoi}.\tag{1.5.1.2}
\end{align*}\etiquette[1.5.1.1]{XVIII.eq1.5.1.1}
\etiquette[1.5.1.2]{XVIII.eq1.5.1.2}

$X$ 的每个截面 $t$ 都在 $X$ 上定义一个可逆层 $\Oo(t)$。采用 \hyperref[XVIII.subsubsec1.4.24]{1.4.24} 的记号（\cf \hyperref[XVIII.subsubsec1.4.24]{1.4.24} $a\leftrightarrow b$），由于这一构造与任意基变换相容，它定义了 $S$ 上的一个 Picard 叠态射
\begin{equation*}
\ch \ZZ^{(X)}\longrightarrow \PIC(X/S){\quoi}.\tag{1.5.1.3}
\end{equation*}\etiquette[1.5.1.3]{XVIII.eq1.5.1.3}

若 $C$ 是 $S$ 上的任意 Picard 叠，则态射 \eqhyperref[XVIII.eq1.5.1.3]{1.5.1.3} 和 \eqhyperref[XVIII.eq1.4.24.1]{1.4.24.1} 定义一个加法函子
\begin{equation*}
\HOM(\PIC(X/S),C)\longrightarrow\HOM(\ch(\ZZ^{(X)}),
C)\xrightarrow{\sim} f_{\ast}f^{\ast}C{\quoi}.\tag{1.5.1.4}
\end{equation*}\etiquette[1.5.1.4]{XVIII.eq1.5.1.4}

应用构造 $\flat$ 后，函子 \eqhyperref[XVIII.eq1.5.1.3]{1.5.1.3} 定义一个态射\nde{这里应当理解，平移函子 $[1]$ 的运算优先级高于截断 $\tau_{\leq 0}$。}
\begin{equation*}
\ZZ^{(X)}\longrightarrow \tau_{\leq 0}
Rf_{\ast}\GG_{m}[1]\longrightarrow
Rf_{\ast}\GG_{m}[1]{\quoi},\tag{1.5.1.5}
\end{equation*}\etiquette[1.5.1.5]{XVIII.eq1.5.1.5}%
而与函子 \eqhyperref[XVIII.eq1.5.1.4]{1.5.1.4} 类似的对象，由 \eqhyperref[XVIII.eq1.4.24.2]{1.4.24.2}，是一个态射
\begin{equation*}
R\Hhom(\tau_{\leq 0} Rf_{\ast}(\GG_{m}[1]),K)\longrightarrow
Rf_{\ast} f^{\ast}K{\quoi}.\tag{1.5.1.6}
\end{equation*}\etiquette[1.5.1.6]{XVIII.eq1.5.1.6}

若 $C$ 是满足（\hyperref[XVIII.par1.3.1.1]{1.3.1.1}）的一个群 $G$ 下的挠子叠 $\ch(G[1])$，定理 \hyperref[XVIII.thm1.3.10]{1.3.10} 断言，函子 \eqhyperref[XVIII.eq1.5.1.4]{1.5.1.4} 是等价，\pageoriginale{}其逆等价将 $G_{X}$ 下的挠子 $K$ 对应到加法函子 $\langle \ast,K]$。\eqhyperref[XVIII.eq1.5.1.4]{1.5.1.4} 为等价这一断言，等价于说从 \eqhyperref[XVIII.eq1.5.1.4]{1.5.1.4} 或 \eqhyperref[XVIII.eq1.5.1.6]{1.5.1.6} 得到的态射
\begin{equation*}
\tau_{\leq 1}R\Hhom(\tau_{\leq
0}Rf_{\ast}(\GG_{m}[1]),G)\longrightarrow \tau_{\leq 1}
Rf_{\ast}f^{\ast}G\tag{1.5.1.7}
\end{equation*}\etiquette[1.5.1.7]{XVIII.eq1.5.1.7}%
是同构。

\begin{enonce*}{定理 1.5.2}\etiquette[1.5.2]{XVIII.thm1.5.2}
\emph{（泛系数公式，universal coefficient formula）。}

设 $f:X\rightarrow S$ 是 $S$ 上的光滑射影曲线，$C$ 是 $S_{\text{\rm fppf}}$ 上的 Picard 叠，并设 $K=C^{\flat}$（\hyperref[XVIII.prop1.4.15]{1.4.15}）。假设 $K$ 局部同构（在 $D^{+}(S_{\text{\rm fppf}})$ 中）于形如 $G_{-1}\rightarrow G_{0}$ 的复形，其中层 $G_{i}$ 满足（\hyperref[XVIII.par1.3.1.1]{1.3.1.1}）。则：
\begin{enumerate}
\renewcommand{\theenumi}{\Roman{enumi}}
\renewcommand{\labelenumi}{(\theenumi)}
\item 函子 \eqhyperref[XVIII.eq1.5.1.4]{1.5.1.4}
\begin{equation*}
\HOM(\PIC(X/S),C)\longrightarrow f_{\ast}f^{\ast}C\tag{1.5.2.1}
\end{equation*}\etiquette[1.5.2.1]{XVIII.eq1.5.2.1}%
是 Picard 叠的等价。

\item 从 \eqhyperref[XVIII.eq1.5.1.6]{1.5.1.6} 或 \eqhyperref[XVIII.eq1.5.1.4]{1.5.1.4} 得到的态射
\begin{equation*}
\tau_{\leq 0}R\Hhom(\tau_{\leq 0}
Rf_{\ast}(\GG_{m}[1]),K)\longrightarrow \tau_{\leq 0}
Rf_{\ast}f^{\ast}K\tag{1.5.2.2}
\end{equation*}\etiquette[1.5.2.2]{XVIII.eq1.5.2.2}%
是同构。
\end{enumerate}
\end{enonce*}

这一定理使 $\tau_{\leq 0}Rf_{\ast}(\GG_{m}[1])$ 扮演经典泛系数公式中同调所扮演的角色。

显然，对每个叠 $C$ 都有 $(I)\leftrightarrow (II)$，而且 $(I)$、$(II)$ 关于 $S$ 都是局部性质。若 $C$ 形如 $\ch(G[1])$，则如上所述，定理由 \hyperref[XVIII.thm1.3.10]{1.3.10} 得到。还需证明：若 $K$ 形如
$$
K:d:G_{-1}\longrightarrow G_{0}{\quoi},
$$
且 \eqhyperref[XVIII.eq1.5.2.2]{1.5.2.2} 对 $G_{-1}[1]$ 和 $G_{0}[1]$ 为同构，则 \eqhyperref[XVIII.eq1.5.2.2]{1.5.2.2} 对 $K$ 也是同构。为验证这一点，只需将上同调长正合列
$$
0\longrightarrow R^{-1}f_{\ast}K\longrightarrow
f_{\ast}G_{-1}\longrightarrow f_{\ast}G_{0}\longrightarrow
R^{0}f_{\ast}K\longrightarrow R^{1} f_{\ast}G_{-1}\longrightarrow
R^{1}f_{\ast}G_{0}
$$
与\pageoriginale \eqhyperref[XVIII.eq1.5.2.2]{1.5.2.2} 左侧的类似长正合列比较，并应用五引理。

\setcounter{subsubsection}{2}
\subsubsection{}\etiquette[1.5.3]{XVIII.subsubsec1.5.3}
设 $u$ 是 $S$ 上光滑射影曲线之间的一个平坦态射
$$
\xymatrix@=1.25cm{
X              \ar[rr]^{u} \ar[rd]_{f_1}&       &Y \ar[ld]^{f_2}\\
                                       &  S    &           .
}
$$
于是，加法函子图
$$\vcenter{\etiquette[1.5.3.1]{XVIII.eq1.5.3.1}
\xymatrix{
\ZZ^{(X)}  \ar[r]\ar[d]^{u} & \PIC(X/S) \ar[d]^{N_{X/Y}}\\
\ZZ^{(Y)}  \ar[r] & \PIC(Y/S)
}}\leqno{(1.5.3.1)}
$$
本质交换。

因此，对于 $S$ 上的每个 Picard 叠 $C$，图
$$\vcenter{\etiquette[1.5.3.2]{XVIII.eq1.5.3.2}
\xymatrix@R=1.75cm@C=2cm{
\HOM(\PIC(Y/S),C)         \ar[r] \ar[d]^{\HOM(N_{X/Y},C)}&
f_{2\ast}f^{\ast}_2 C \ar[d]^{u^{\ast}}\\
\HOM(\PIC(X/S), C)        \ar[r] & f_{1\ast}f^{\ast}_1 C
}}\leqno{(1.5.3.2)}
$$
本质交换；对于 $K\in \Ob D^{+}(S)$，图
$$\vcenter{\etiquette[1.5.3.3]{XVIII.eq1.5.3.3}
\xymatrix@C=2cm@R=1.5cm{
R\Hhom(\tau_{\leq 0} Rf_{2\ast}\GG_m[1],K) \ar[r]\ar[d]^{R\Hhom(N_{X/Y},K)}&
Rf_{2\ast}f^{\ast}_2 K  \ar[d]^{u^{\ast}}\\
R\Hhom(\tau_{\leq 0} Rf_{1\ast}\GG_m[1],K) \ar[r]&
Rf_{1\ast}f^{\ast}_{1} K
}}\leqno{(1.5.3.3)}
$$
交换。

再设 $G$ 是 $S$ 上满足（\hyperref[XVIII.par1.3.1.1]{1.3.1.1}）的阿贝尔层。迹 \hyperrefext[XVII][XVII.subsec6.3]{6.3} 给出一个加法函子
\begin{equation*}
\Tr_{u}:f_{1\ast}f^{\ast}_{1}\ch(G[1])\longrightarrow f_{2\ast}f^{\ast}_{2}\ch(G[1])\tag{1.5.3.4}
\end{equation*}\etiquette[1.5.3.4]{XVIII.eq1.5.3.4}%
而\pageoriginale \hyperref[XVIII.par1.3.8.6]{1.3.8.6} 给出一个使下图本质交换的同构
$$\vcenter{\etiquette[1.5.3.5]{XVIII.eq1.5.3.5}
\xymatrix@R=1.75cm{
\HOM(\PIC(X/S), \ch(G[1])) \ar[d]^{\HOM(u^{\ast}, \ch(G[1]))}&  f_{1\ast}f^{\ast}_1
\ch(G[1]) \ar[l] \ar[d]^{\Tr_u}\\
\HOM(\PIC(Y/S), \ch(G[1])) & f_{2\ast}f^{\ast}_2 \ch(G[1]).\ar[l]
}}\leqno{(1.5.3.5)}
$$

因此，图
$$\vcenter{\etiquette[1.5.3.6]{XVIII.eq1.5.3.6}
\xymatrix@R=1.75cm{
\tau_{\leq 1} R\Hhom(\tau_{\leq 0} Rf_{1\ast} \GG_{m}[1],
G)  \ar[d]_{R\Hhom(u^{\ast},G)}  & \tau_{\leq 1} Rf_{1\ast}
f^{\ast}_1G  \ar[l]_-{\sim} \ar[d]^{Tr_u}\\
\tau_{\leq 1} R\Hhom(\tau_{\leq 0} Rf_{2\ast} \GG_m[1], G)&
  \tau_{\leq 1} Rf_{2\ast}f^{\ast}_2 G \ar[l]_-{\sim}
}}\leqno{(1.5.3.6)}
$$
交换。

\subsubsection{}\etiquette[1.5.4]{XVIII.subsubsec1.5.4}
采用 \hyperref[XVIII.thm1.5.2]{1.5.2} 的记号，设 $G$ 是 $S$ 上满足 \hyperref[XVIII.par1.3.1.1]{1.3.1.1} 的阿贝尔层。关于 $S$ 对平展拓扑局部地，有（非典范地）
$$
\tau_{\leq 0}
Rf_{\ast}\GG_{m}[1]\simeq
f_{\ast}\GG_{m}[1]+R^{1}f_{\ast}\GG_{m}{\quoi}.
$$

事实上，在 Picard 叠的层次上，可将可逆层的每个同构类对应到属于该类、且沿 $X/S$ 的适当截面刚化的可逆层，由此定义一个截面 $\ch(R^{1}f_{\ast}\GG_{m})\rightarrow f_{\ast}\ch(\GG_{m}[1])$\nde{回顾此叠等同于 $\PIC(X/S)$，\cf \eqhyperref[XVIII.eq1.5.1.1]{1.5.1.1}。}。

因此，同构 \eqhyperref[XVIII.eq1.5.2.2]{1.5.2.2} 给出一个同构
\begin{equation*}
\Hhom(R^{1}f_{\ast}\GG_{m},G)\xrightarrow{\sim}\Hhom_{S}(X,G){\quoi}\tag{1.5.4.1}
\end{equation*}\etiquette[1.5.4.1]{XVIII.eq1.5.4.1}%
和一个正合列
\begin{equation*}
0\longrightarrow
\Eext^{1}(R^{1}f_{\ast}\GG_{m},G)\longrightarrow
R^1f_{\ast}G\longrightarrow
\Hhom(f_{\ast}\GG_{m},G)\longrightarrow 0{\quoi}.\tag{1.5.4.2}
\end{equation*}\etiquette[1.5.4.2]{XVIII.eq1.5.4.2}

\eqhyperref[XVIII.eq1.5.4.1]{1.5.4.1} 和 \eqhyperref[XVIII.eq1.5.4.2]{1.5.4.2} 中的态射解释如下：
\begin{enumerate}
\renewcommand{\theenumi}{\alph{enumi}}
\renewcommand{\labelenumi}{\theenumi)}
\item 态射\pageoriginale \eqhyperref[XVIII.eq1.5.4.1]{1.5.4.1} 和 \eqhyperref[XVIII.eq1.5.4.2]{1.5.4.2} 中的第一个态射，由从 $X$ 到 $\sheafPic_{X/S}$ 的典范映射定义。

\item \eqhyperref[XVIII.eq1.5.4.2]{1.5.4.2} 中的第二个态射是次数 \eqhyperref[XVIII.eq1.3.7.2]{1.3.7.2}。
\end{enumerate}

关于 $S$ 对平展拓扑局部地，有（非典范地）$\sheafPic_{X/S}\sim \sheafPic^{0}_{X/S}\times \ZZ^{k}$；所以
$$
\Eext^{1}(R^{1}f_{\ast}\GG_{m},G)\xrightarrow{\sim}\Eext^{1}(\sheafPic^{0}_{X/S}, G)
$$
而 \eqhyperref[XVIII.eq1.5.4.2]{1.5.4.2} 也可以写成
\begin{equation*}
0 \longrightarrow \Eext^{1}(\sheafPic^{0}_{X/S},G)\longrightarrow
R^{1}f_{\ast}G\longrightarrow
\Hhom(f_{\ast}\GG_{m},G)\longrightarrow 0{\quoi}.\tag{1.5.4.3}
\end{equation*}\etiquette[1.5.4.3]{XVIII.eq1.5.4.3}

\subsubsection{}\etiquette[1.5.5]{XVIII.subsubsec1.5.5}
在 \eqhyperref[XVIII.eq1.5.2.2]{1.5.2.2} 中取 $K=\GG_{m}[1]$ 时，定理 \hyperref[XVIII.thm1.5.2]{1.5.2} 就成为 $\tau_{\leq 0}(Rf_{\ast}\GG_{m}[1])$ 的一个自对偶定理。代表 $\tau_{\leq 0} Rf_{\ast} \GG_{m}[1]$ 的任意复形，都有一个典范的三步滤过，其逐次商在唯一拟同构意义下为
$$
\sheafPic_{X/S}/\sheafPic^{0}_{X/S}\quav \sheafPic^{0}_{X/S}\quav
(f_{\ast}\GG_{m})[1]{\quoi}.
$$

前述自对偶对应于 $\sheafPic_{X/S}/\sheafPic^{0}_{X/S}$ 与 $f_{\ast}\GG_{m}$ 之间的一个对偶（取值于 $\GG_{m}$），以及 $\sheafPic^{0}_{X/S}$ 上的一个自对偶（取值于 $\GG_{m}[1]$）（Jacobian 的自对偶）。

\subsubsection{}\etiquette[1.5.6]{XVIII.subsubsec1.5.6}
对于不假定在 $S$ 上固有的光滑曲线 $S:f:X\rightarrow S$，情形明显更坏。加法群的表现似乎无法控制。不过，当 $G$ 是挠性与 $S$ 的剩余特征互素的平展挠层时，前面的论证仍可推广到如下光滑曲线：它们由 $S$ 上的固有曲线去掉一个在 $S$ 上有限的部分得到。可以期望，平展 $p$-挠层也有合理的表现；见 \textsc{Serre} \cite[Ch.~VI nᵒ 11 p.~126]{sga4.exp.XVIII.bib.6}。

这里，定理 \hyperref[XVIII.thm1.2.2]{1.2.2} 的作用将由无环性定理\pageoriginale \hyperrefext[XV][XV.cor2.2]{2.2} 取代；后者有如下推论。

\begin{enonce*}{引理 1.5.7}\etiquette[1.5.7]{XVIII.lem1.5.7}
设 $p:X\rightarrow S$ 是 $S$ 上标准仿射空间的投影，$G$ 是 $S_{\text{\rm et}}$ 上的阿贝尔挠层，其挠性与 $S$ 的剩余特征互素。仍以 $G$ 表示 $G$ 在景 $S_{\text{\rm fppf}}$ 上的逆像。则函子 $p^{\ast}$ 在 $S_{\text{\rm et}}$ 上（或 $S_{\text{\rm fppf}}$ 上，由 {\textsc{Grothendieck} \cite[Corollaire~11.9]{sga4.exp.XVIII.bib.2}}，二者等价）的 $G$ 下的挠子范畴与 $X$ 上的 $G_{X}$ 下的挠子范畴之间给出范畴等价。
\end{enonce*}

\setcounter{subsubsection}{7}
\subsubsection{}\etiquette[1.5.8]{XVIII.subsubsec1.5.8}
设 $\bar{f}:\bar{X}\rightarrow S$ 是有限表现的固有态射，其纤维均为纯一维；$Y$ 是 $\bar{X}$ 的一个在 $S$ 上有限的闭子概形，由理想 $\mathfrak{m}$ 定义，并假定 $X=\bar{X}-Y$ 是 $S$ 上的光滑曲线。另设 $G$ 是 $S_{\text{\rm et}}$ 上的挠层，其挠性与 $S$ 的剩余特征互素。我们仍以 $G$ 表示作为 $G$ 的逆像得到的 $S$ 的大 fppf 景（\hyperrefext[XVII][XVII.subsec0.10]{0.10}）上的层。

设 $\GG_{m}(\mathfrak{m})$ 是 $\GG_{m\bar{X}}$ 的子层，由在 $Y$ 上等于 1 的截面组成。$\GG_{m}(\mathfrak{m})$ 下的一个挠子，等同于 $\bar{X}$ 上的一个在 $Y$ 上被平凡化的可逆层。$X$ 上在 $S$ 上有限的每个相对除子，

都定义这样的一个可逆层。关于 $S$ 对平展拓扑局部地，存在这样的除子，给出相对丰沛可逆层。

\begin{enonce*}{引理 1.5.9}\etiquette[1.5.9]{XVIII.lem1.5.9}
在 \hyperref[XVIII.subsubsec1.5.8]{1.5.8} 的假设下，设 $S$ 仿射\nde{实际上，必须假定 $S$ 是闭点为 $s$ 的局部概形。若要保留原来的假设，则在结论中应要求 $m$ 只在 $s$ 在 $S$ 中的一个开邻域的逆像上有定义，如下文 \hyperref[XVIII.par1.5.10.1]{1.5.10.1} 中的情形。}；设 $\Oo_{\bar{X}}(1)$ 是 $\bar{X}$ 上沿 $Y$ 被平凡化的相对丰沛可逆层，$\Mm$ 是 $\bar{X}$ 上沿 $Y$ 被平凡化的可逆层。对于每一点 $s\in S$，存在整数 $n_{0}$，使得对每个整数 $n\geq n_{0}$ 和纤维 $\bar{X}_{s}$ 上 $\Mm(n)\otimes k(s)$ 的每个在 $Y_{s}$ 上等于 $1$ 的截面 $m_{s}$，都存在 $\Mm(n)$ 的一个在 $Y$ 上等于 $1$ 的截面 $m$，提升 $m_{s}$。
\end{enonce*}

归结为 $S$ 是 Noether 概形的情形（因为一族满态射在基变换之后仍是一族满态射\nde{去掉这一 Noether 假设，比乍看起来要微妙一些，因为问题是仿射的，而非线性的。若 $\bar{X}$、$Y$、$\Mm$ 是由 Noether 概形 $S_0$ 上的相应数据 $\bar{X}_0$、$Y_0$、$\Mm_0$ 通过基变换 $S\to S_0$ 得到的（假定这是局部概形之间的局部态射）（\cf \EGA IV 8），且已知所需断言在这一 Noether 情形成立，则可通过一次令人不快的图追踪得到 $S$ 上所需的结果；在追踪过程中，可能需要考虑所需断言的一个变体，要求截面在 $Y_s$ 或 $Y$ 上等于 $0$，而不是 $1$。}）。设 $\Nn$ 为 $\Mm$ 中在\pageoriginale $\bar{X}_{s}$ 和 $Y$ 上消失的局部截面组成的凝聚层。对于每个 $n$，问题的局部解组成 $\Nn(n)$ 下的一个主齐性空间，因此提升 $m_{s}$ 的障碍属于 $H^{1}(\bar{X},\Nn(n))=H^{0}(S,R^{1}\bar{f}_{\ast}\Nn(n))$。但当 $n$ 充分大时，$R^{1}\bar{f}_{\ast}\Nn(n)=0$（\EGA III 2.2.1）。

\setcounter{subsubsection}{9}
\subsubsection{}\etiquette[1.5.10]{XVIII.subsubsec1.5.10}
设 $\Ll$ 是 $\bar{X}$ 上沿 $Y$ 被平凡化的可逆层，并设 $s\in S$。考虑以下条件

\paragraph{}\etiquette[1.5.10.1]{XVIII.par1.5.10.1}
$\bar{X}_{s}$ 上 $\Ll\otimes k(s)$ 的每个在 $Y_{s}$ 上等于 1 的截面，都可提升为 $s$ 的一个邻域的逆像 $U$ 上 $\Ll$ 的一个截面，它在 $\fprod{Y}{U}{S}$ 上等于 1。

\paragraph{}\etiquette[1.5.10.2]{XVIII.par1.5.10.2}
存在 $\bar{X}_{s}$ 上 $\Ll\otimes k(s)$ 的四个截面 $m_{i}$，它们在 $Y_s$ 上等于 1，且在 $X_{s}$ 的每个不可约分支的一般点线性无关。

考虑下面的条件，它们分别涉及 $\Ll$ 的一个截面 $m$、一对截面 $(m_{1},m_{2})$ 及 $\Ll$ 的三个截面 $(m_{1},m_{2},m_{3})$。

\paragraph{}\etiquette[1.5.10.3]{XVIII.par1.5.10.3}
$m$ 在 $Y$ 上等于 1，且由 $m=0$ 定义的 $X$ 的子概形 $Z(m)$ 是 $S$ 上的一个相对除子（这意味着它在 $S$ 上拟有限，因而在 $S$ 上有限，因为 $X=\bar{X}-Y$ 在 $S$ 上光滑）。

\paragraph{}\etiquette[1.5.10.4]{XVIII.par1.5.10.4}
设 $\lambda$ 为仿射直线 $g:\EE^{1}_{S}\rightarrow S$ 上的典范坐标。${g'}^{\ast}\Ll$ 的截面 $\lambda m_{1}+(1-\lambda)m_{2}$ 满足（\hyperref[XVIII.par1.5.10.3]{1.5.10.3}）。

\paragraph{}\etiquette[1.5.10.5]{XVIII.par1.5.10.5}
类似地，$\lambda m_{1}+\mu m_{2}+(1-\lambda -\mu)m_{3}$ 满足（\hyperref[XVIII.par1.5.10.3]{1.5.10.3}）。

下面的引理留给读者用一般位置论证来验证\nde{这里应注意，考虑到（\hyperref[XVIII.par1.5.10.1]{1.5.10.1}）以及在局部情形只考察 $s$ 上的纤维就能验证条件（\hyperref[XVIII.par1.5.10.3]{1.5.10.3}）的第二部分这一事实，引理现在只涉及 $\Ll$ 在纤维 $X_s$ 上的限制的截面。例如，对于 (I)，条件（\hyperref[XVIII.par1.5.10.2]{1.5.10.2}）给出一个维数为 $3$ 的\emph{仿射}空间 $V$，其中的元素是 $X_s$ 上 $\Ll\otimes k(s)$ 的在 $Y_s$ 上等于 $1$ 的截面，而且属于这个空间 $V$ 的截面在 $X_s$ 的任何一般点都不为零。我们要确定 $k\in V$，使对每个 $i$，$(m_i,k)$ 满足（\hyperref[XVIII.par1.5.10.4]{1.5.10.4}），即使连接点 $m_i$ 与 $k$ 的仿射直线不经过原点（在原点与 $X_s$ 的不同不可约分支相联系的各个表现中均如此）；换言之，只需 $k$ 不属于连接 $m_i$ 与原点的直线。于是得到一组同时要求的条件，它们分别在 $V$ 的非空开子集上满足。为了完全严格，并考虑到不可约分支未必几何不可约这一事实，当必须在一个扩张中取 $V$ 的点时（若 $k(s)$ 为有限域），实际上需要引入一个有限 Galois 扩张 $s'\to s$，使得 $X_{s'}$ 的不可约分支几何不可约。在 $X_{s'}$ 上作上述论证得到的 $V_{s'}$ 的稠密开子集，是 Galois 不变的，因而定义了所需的 $V$ 的稠密开子集。}。

\begin{enonce*}{引理 1.5.11}\etiquette[1.5.11]{XVIII.lem1.5.11}
设 $\Ll$ 是 $\bar{X}$ 上沿 $Y$ 被平凡化的可逆层，$\bar{s}$ 是 $S$ 的一个几何点，其在 $S$ 中的像为 $s$。若 $\Ll$ 在 $s$ 满足（\hyperref[XVIII.par1.5.10.1]{1.5.10.1}）和（\hyperref[XVIII.par1.5.10.2]{1.5.10.2}），则：
\begin{enumerate}
\renewcommand{\theenumi}{\Roman{enumi}}
\renewcommand{\labelenumi}{(\theenumi)}
\item 若 $\Ll$ 的 $n$ 个截面 $m_{i}$ 满足（\hyperref[XVIII.par1.5.10.3]{1.5.10.3}），则存在 $\bar{s}$ 的一个平展邻域 $U$ 和 $\Ll|\bar{f}^{-1}(U)$ 的一个截面 $k$，使各对 $(m_{i},k)$ 满足（\hyperref[XVIII.par1.5.10.4]{1.5.10.4}）。

\item 设\pageoriginale $m_{1}$、$m_{2}$、$k_{1}$、$k_{2}$ 是 $\Ll$ 的四个截面，使各对 $(m_{i},k_{j})$ 满足（\hyperref[XVIII.par1.5.10.4]{1.5.10.4}）。则存在 $\bar{s}$ 的一个平展邻域 $U$ 和 $\Ll|\bar{f}^{-1}(U)$ 的一个截面 $k$，使各三元组 $(m_{i}, k_{j}, k)$ 满足（\hyperref[XVIII.par1.5.10.5]{1.5.10.5}）。
\end{enumerate}
\end{enonce*}

当 $k(s)$ 是无限域时，只需取 $U$ 为一个 Zariski 邻域即可。

\setcounter{subsubsection}{11}
\subsubsection{}\etiquette[1.5.12]{XVIII.subsubsec1.5.12}
再设 $\Ll$ 是 $\bar{X}$ 上沿 $Y$ 刚化的可逆层，$K$ 是 $X$ 上的一个 $G$-挠子，其中 $G$ 如 \hyperref[XVIII.lem1.5.7]{1.5.7} 中所述。我们要在 $S$ 上定义一个 $G$-挠子 $\langle \Ll,K]$，将 \hyperref[XVIII.subsec1.3]{1.3} 中所研究的挠子推广，并使之关于 $\Ll$ 具有加法性。

设 $(m_{1},m_{2})$ 是 $\Ll$ 的一对满足（\hyperref[XVIII.par1.5.10.4]{1.5.10.4}）的截面。采用 \hyperref[XVIII.par1.5.10.3]{1.5.10.3} 及 \hyperref[XVIII.par1.5.10.4]{1.5.10.4} 的记号，将挠子 ${g'}^{\ast}K$ 从 $\lambda m_{1}+(1-\lambda)m_{2}=0$ 取迹到 $\EE^{1}_{S}$，得到 $\EE^{1}_{S}$ 上的一个挠子 $K_{1,2}$。它沿截面 $\lambda =0$（\resp $\lambda = 1$）的拉回，是挠子 $K_{2}=\Tr_{Z(m_{2})/S}(K)$（\resp $K_{1}=\Tr_{Z(m_{1})/S}(K)$）。由 \hyperref[XVIII.lem1.5.7]{1.5.7}，挠子 $K_{1,2}$ 典范地是 $S$ 上某个 $G$ 下的挠子的拉回；因此 $K_{1}$ 与 $K_{2}$ 之间有一个典范同构。

若 $m_{1}$、$m_{2}$、$k$ 是 $\Ll$ 的三个截面，使各对 $(m_{i},k)$ 满足（\hyperref[XVIII.par1.5.10.4]{1.5.10.4}），则通过复合仍可得到同构
\begin{equation*}
\Psi_{k}:K_{1}=\Tr_{Z(m_{1})/S}(K)\xrightarrow{\sim}
K_{2}=\Tr_{Z(m_{2})/S}(K){\quoi}.\tag{1.5.12.1}
\end{equation*}\etiquette[1.5.12.1]{XVIII.eq1.5.12.1}

若 $E$ 是 $X$ 上在 $S$ 上有限的相对除子，则 $\Ll\otimes \Oo(E)$ 的截面 $m'_{i}=m_{i}\otimes 1$ 仍满足（\hyperref[XVIII.par1.5.10.3]{1.5.10.3}），且有（对 $i=1,2$）
\begin{equation*}
K'_{i}=\Tr_{Z(m'_{i})/S}(K)\simeq K_{i}+\Tr_{E/S}(K)\tag{1.5.12.2}
\end{equation*}\etiquette[1.5.12.2]{XVIII.eq1.5.12.2}%
而同构 $\Psi_{k\otimes 1}:K'_{1}\rightarrow K'_{2}$ 从 $\Psi_{k}$ 得到：
\begin{equation*}
\Psi_{k\otimes 1}=\Psi_{k}+\id_{\Tr_{E/S}}(K){\quoi}.\tag{1.5.12.3}
\end{equation*}\etiquette[1.5.12.3]{XVIII.eq1.5.12.3}

为验证 $\Psi_{k}$ 不依赖于 $k$，只需在 $S$ 的每一点 $s$ 的一个平展邻域中验证；而由 \hyperref[XVIII.eq1.5.12.3]{1.5.12.3} 及 \hyperref[XVIII.lem1.5.9]{1.5.9}，我们\pageoriginale{}归结为条件（\hyperref[XVIII.par1.5.10.1]{1.5.10.1}）和（\hyperref[XVIII.par1.5.10.2]{1.5.10.2}）成立的情形。采用 \hyperref[XVIII.lem1.5.11]{1.5.11} (II) 的记号，$\Psi_{\lambda k_{i}+(1-\lambda)k}-\Psi_{k_{i}}$ 是从 $\EE^{1}_{S}$ 到 $G$ 的一个态射，在 $\lambda =1$ 处为零，因而恒为零；所以 $\Psi_{k_{i}}=\Psi_{k}$，进而 $\Psi_{k_{1}}=\Psi_{k_{2}}$。

设 $m_{1}$ 和 $m_{2}$ 是 $\Ll$ 的两个满足 \hyperref[XVIII.par1.5.10.3]{1.5.10.3} 的截面，且设 $K_{i}=\Tr_{Z(m_{i})/S}(K)$。为定义同构 $\Psi_{m_{2},m_{1}}:K_{1}\rightarrow K_{2}$，只需在 $S$ 的每一点 $s$ 的一个平展邻域内，以与基变换相容的方式定义它。若 $E$ 是 $X$ 上在 $S$ 上有限的相对除子，使得 $\Ll\otimes \Oo(E)$ 满足（\hyperref[XVIII.par1.5.10.1]{1.5.10.1}）和（\hyperref[XVIII.par1.5.10.2]{1.5.10.2}）（见 \hyperref[XVIII.lem1.5.9]{1.5.9}），且各对 $(m_{i}\otimes 1,k)$ 满足（\hyperref[XVIII.par1.5.10.4]{1.5.10.4}），则置（\cf \eqhyperref[XVIII.eq1.5.12.3]{1.5.12.3}）
$$
\Psi_{m_{2},m_{1}}+\Tr_{E/S}(K)=\Psi_{k}{\quoi}.
$$
这就定义了 $\Psi_{m_{2},m_{1}}$，而利用（\hyperref[XVIII.lem1.5.11]{1.5.11}. (I)）可验证 $\Psi_{m_{3},m_{2}}\Psi_{m_{2},m_{1}}=\Psi_{m_{3},m_{1}}$。这些构造与基变换相容。若存在 $\Ll$ 的一个满足（\hyperref[XVIII.par1.5.10.3]{1.5.10.3}）的截面 $m$，则置
$$
\langle \Ll,K]=\Tr_{Z(m)/S}(K)
$$
由前面的论证，这一挠子在典范同构意义下不依赖于 $m$ 的选择。

在一般情形，注意为定义 $\langle \Ll,K]$，只需在 $S$ 的每一点 $s$ 的一个平展邻域中，以与基变换相容的方式定义它。若如上所述，$\Ll\otimes \Oo(E)$ 满足（\hyperref[XVIII.par1.5.10.1]{1.5.10.1}）及（\hyperref[XVIII.par1.5.10.2]{1.5.10.2}），则 $\Ll\otimes \Oo(E)$ 有满足（\hyperref[XVIII.par1.5.10.3]{1.5.10.3}）的截面，并置
$$
\langle \Ll,K]=\langle \Ll\otimes \Oo(E), K]-\Tr_{E/S}(K){\quoi}.
$$

\subsubsection{}\etiquette[1.5.13]{XVIII.subsubsec1.5.13}
在 \hyperref[XVIII.subsubsec1.5.8]{1.5.8} 的假设下，若 $x$、$y$ 是 $\bar{f}_{\ast}\GG_{m}(\mathfrak{m})$ 的两个截面，则当 $\lambda$ 在 $S$ 上仿射直线的某个具有非空纤维的开子集中取值时，$\lambda x+(1-\lambda)y$ 仍是一个截面。

这可以表述为

\paragraph{}\etiquette[1.5.13.1]{XVIII.par1.5.13.1}
“$\bar{f}_{\ast}\GG_{m}(\mathfrak{m})$ 是道路连通的。”

从仿射直线的一个具有非空纤维的开子集到 $G$ 的每个\pageoriginale{}态射都是常值的。因此，不存在从 $\bar{f}_{\ast}\GG_{m}(\mathfrak{m})$ 到 $G$ 的非平凡同态，而且 $\Ll$ 的自同构在 $\langle \Ll,K]$ 上平凡地作用。所以，挠子 $\langle \Ll,K]$ 只依赖于 $K$ 和由 $\Ll$ 定义的 $R^{1}\bar{f}_{\ast}\GG_{m}(\mathfrak{m})$ 的截面。通过整体化，由此可以对 $R^{1}\bar{f}_{\ast}\GG_{m}(\mathfrak{m})$ 的截面 $\lambda$ 定义符号 $\langle \lambda,K]$。挠子 $\langle \lambda,K]$ 关于 $\lambda$ 具有加法性，因此对应于 $R^{1}\bar{f}_{\ast}\GG_{m}(\mathfrak{m})$ 由 $G$ 的一个扩张 $e(K)$（见 \hyperref[XVIII.prop1.4.23]{1.4.23}，或 \SGA 7 VII 1.1.6 和 1.2）。

$e(K)$ 的形成与任意基变换及挠子的加法相容，因而给出一个加法函子
$$
e:\bar{f}_{\ast}\ch(G[1])\longrightarrow
\EXT(R^{1}\bar{f}_{\ast}\GG_{m}(\mathfrak{m}),G){\quoi}.
$$

设 $j$ 是从 $X$ 到 $R^{1}\bar{f}_{\ast}\GG_{m}(\mathfrak{m})$ 的典范映射，$j:t\rightarrow$（$\Oo(t)$ 的类）。沿 $j$ 的拉回是一个加法函子
$$
j^{\ast}:\EXT(R^{1}\bar{f}_{\ast}\GG_{m}(\mathfrak{m}),G)\longrightarrow
\bar{f}_{\ast}\ch(G[1])
$$
显然 $j^{\ast}e\sim \Id$。如 \hyperref[XVIII.thm1.3.10]{1.3.10} 和 \hyperref[XVIII.thm1.5.2]{1.5.2} 中一样，可验证对于 fppf 拓扑，有

\begin{enonce*}{命题 1.5.14}\etiquette[1.5.14]{XVIII.prop1.5.14}
在 \hyperref[XVIII.subsubsec1.5.8]{1.5.8} 的假设下，\hyperref[XVIII.subsubsec1.5.13]{1.5.13} 中的函子 $e$、$j^{\ast}$ 是互为拟逆的等价。由 \hyperref[XVIII.const1.4.19]{1.4.19} 和 \hyperref[XVIII.prop1.4.23]{1.4.23}，得到一个同构
$$
\Ext^{1}(R^{1}\bar{f}_{\ast}\GG_{m}(\mathfrak{m}),G)\xrightarrow{\sim}
H^{1}(X,G){\quoi}.
$$
\end{enonce*}

\setcounter{subsubsection}{14}
\subsubsection{}\etiquette[1.5.15]{XVIII.subsubsec1.5.15}
设 $f_{i}:X_{i}\rightarrow S$ $(i=1,2)$ 是两条如 \hyperref[XVIII.subsubsec1.5.8]{1.5.8} 中所述的已紧化曲线。

设 $\bar{u}:\bar{X}_{1}\rightarrow \bar{X}_{2}$ 是一个态射，使得 $Y_{1}$ 是 $Y_{2}$ 的逆像的子概形，且它诱导一个从 $X_{1}$ 到 $X_{2}$ 的有限平坦态射。则对 $\bar{X}_{2}$ 上沿 $Y_2$ 刚化的可逆层 $\Ll$ 及 $X_{1}$ 上的挠子 $K$，有
$$
\langle \Ll,\tr_{X_{1}/X_{2}}(K)]\xrightarrow{\sim}
\langle \bar{u}^{\ast}\Ll,K]{\quoi},
$$
由此\pageoriginale{}推出下图交换
$$\vcenter{\etiquette[1.5.15.1]{XVIII.eq1.5.15.1}
\xymatrix@R=1.75cm{
\Ext^{1}(R^1\bar{f}_{1\ast} \GG_m (\mathfrak{m}_1), G) 
\ar[r]^-{\sim}\ar[d]^{\Ext^{1}(\bar{u}^{\ast},\Id)} &   H^{1}(X_1,G) \ar[d]^{\Tr_{X_1 /X_2}}\\
\Ext^1(R^1\bar{f}_{2\ast}\GG_m (\mathfrak{m}_2), G)  \ar[r]^-{\sim}&
H^1(X_2,G).
}}\leqno{(1.5.15.1)}
$$

若 $\bar{u}:\bar{X}_{1}\rightarrow \bar{X}_{2}$ 有限平坦，且满足 $Y_{1}=\bar{u}^{-1}Y_{2}$（作为概形），则对 $\bar{X}_{1}$ 上沿 $Y_1$ 刚化的可逆层 $\Ll$ 及 $X_{2}$ 上的挠子 $K$，有
$$
\langle \Ll,u^{\ast}K]\xrightarrow{\sim}
\langle N_{\bar{X}_{1}/\bar{X}_{2}}\Ll,K]
$$
由此推出下图交换
\begin{equation*}
\vcenter{
\xymatrix@R=1.75cm@C=1.5cm{
\Ext^{1}(R^{1}\bar{f}_{2\ast} \GG_m (\mathfrak{m}_2) , G)
\ar[d]^{\Ext^{1}(N_{\bar{X}_1 /\bar{X}_2}, \Id)}\ar[r]^-{\sim} &
H^{1}(X_2,G) \ar[d]^{u^{\ast}}\\
\Ext^{1}(R^{1}\bar{f}_{1\ast} \GG_m (\mathfrak{m}_1), G)  \ar[r]^-{\sim} &
H^{1}(X_1, G).
}}\tag{1.5.15.2}
\end{equation*}\etiquette[1.5.15.2]{XVIII.eq1.5.15.2}

当然，可以用 Picard 叠的语言更精确地陈述这一结果。

\subsection{一个消没定理}\etiquette[1.6]{XVIII.subsec1.6}

本节从 \hyperref[XVIII.lem1.6.6]{1.6.6}（第二种证明）到末尾，可以独立于第 \hyperref[XVIII.subsec1.2]{1.2} 至 \hyperref[XVIII.subsec1.5]{1.5} 小节阅读。

\begin{enonce*}{引理 1.6.1}\etiquette[1.6.1]{XVIII.lem1.6.1}
设 $X$ 是概形，$Y$ 是 $X$ 的一个由理想 $\mathfrak{m}$ 定义的闭子概形，其补集为 $U$。
$$
U\mathop{\hookrightarrow}^{j} X\mathop{\hookleftarrow}^{i} Y{\quoi}.
$$
设 $\GG_{m}(\mathfrak{m})$ 是如下 \emph{fpqc} 层：它在 $X$-概形 $X_{1}$ 上的截面，是 $\Oo^{\ast}_{X_{1}}$ 中在 $Y$ 的逆像上等于 $1$ 的截面。对在 $X$ 上可逆的 $n\geq 1$，序列
$$
0\rightarrow j_{!}𝛍_{n}\rightarrow
\GG_{m}(\mathfrak{m})\xrightarrow{n} \GG_{m}(\mathfrak{m})\rightarrow
0
$$
是 $X$ 的大平展景\emph{（\hyperrefext[XVII][XVII.subsec0.10]{0.10}）}上的层的正合列。
\end{enonce*}

设\pageoriginale $i$ 为由 $i$ 诱导的大平展景态射 $i:Y\rightarrow X$：
$$
i^{\ast}(X_{1}/X)=\fprod{X_{1}}{Y}{X}{\quoi}.
$$
函子 $i_{\ast}$ 是正合的。在图
$$\vcenter{\etiquette[1.6.1.1]{XVIII.eq1.6.1.1}
\xymatrix{
             &    0  \ar[d]           &  0  \ar[d]                   &    0\ar[d]    \\
0  \ar[r]    &j_!𝛍_n  \ar[r]\ar[d]  & \GG_m(\mathfrak{m})\ar[r]\ar[d]  &
             \GG_m(\mathfrak{m}) \ar[r] \ar[d] &0\\
0  \ar[r]     &𝛍_n  \ar[r]  \ar[d]  & \GG_m  \ar[r] \ar[d]  &
  \GG_m  \ar[r] \ar[d]  & 0\\
0  \ar[r]    & i_{\ast}𝛍_n  \ar[d] \ar[r] &
  i_{\ast}\GG_m  \ar[r]\ar[d]  & i_{\ast}\GG_m   \ar[r]\ar[d] & 0\\
             &    0             &  0                    &    0   \\
}}\leqno{(1.6.1.1)}
$$
中，各列均正合，最后两行也正合（\hyperrefext[IX][IX.Thr3.2]{3.2}）。因此第一行正合。

\setcounter{subsubsection}{1}
\subsubsection{}\etiquette[1.6.2]{XVIII.subsubsec1.6.2}
设 $k$ 是一个域，
$$
G'_{1}\longrightarrow G'_{2}\longrightarrow G\longrightarrow
G''_{2}\longrightarrow G''_{1}
$$
是 $\Spec (k)$ 的大 fppf 景上的阿贝尔层的正合列。回顾，若 $G'_{1}$、$G'_{2}$ 可表且有限型，而 $G''_{2}$、$G''_{1}$ 可表且局部有限型，则 $G$ 可表且局部有限型\acorriger{当 $G'_1=G'_2=0$ 时，这是显然的。当 $G''_1=G''_2=0$ 时，利用 \SGA 3 V 10.1.1（新版），可归结为 $G'_1\to G'_2$ 为单态射的情形，而此情形见（\SGA 3 VI${}_{\text{A}}$ 3.2）。于是我们归结为 $G'_1=G''_1=0$ 的情形。这时可将层 $G$ 看作概形 $G''_2$ 上 $G'_2$ 下的挠子。若假定 $G'_2$ 仿射（如果在进行这些化简之前假定 $G'_1$、$G'_2$ 仿射，则由（\SGA 3 VI${}_{\text{B}}$ 11.17）可归结到这种情形），则 $G$ 的可表性立即由忠实平坦下降得到。这个特殊情形足以用于本段中的应用。若 $G'_2$ 拟射影，则可利用（\SGA 1 VIII 7.7）及类似于 \cite[引理~3.3.1]{sga4.exp.XVIII.bib.4bis} 中的论证得到 $G$ 的可表性。事实上，域上的每个有限型群概形 $G$ 都拟射影：可以假定 $G$ 连通；当 $G$ 光滑时，这就是 \cite[推论~1.2]{sga4.exp.XVIII.bib.0bis}，而一般情形可以归结到光滑情形（M. Raynaud），因为由（\SGA 3 VII${}_{\text{A}}$ 8.3），存在从 $G$ 到一个光滑群概形的有限态射。}。

若 $P$ 是 $k$ 上的有限型连通交换代数群，且 $n$ 是在 $k$ 中可逆的整数，则序列
\begin{equation*}
0\rightarrow {}_{n}P\rightarrow P\xrightarrow{n} P\rightarrow 0\tag{1.6.2.1}
\end{equation*}\etiquette[1.6.2.1]{XVIII.eq1.6.2.1}%
正合。若 $G$ 是 $k$ 上被 $n$ 零化的有限平展交换群，则正合列
$$
\Hom(P,G)\rightarrow \Hom({}_{n}P,G)\xrightarrow{j}
\Ext^{1}(P,G)\xrightarrow{n} \Ext^{1}(P,G)
$$
中的第一个群和最后一个箭头为零，从而有同构
\begin{equation*}
j:\Hom({}_{n}P,G)\xrightarrow{\sim} \Ext^{1}(P,G){\quoi}.\tag{1.6.2.2}
\end{equation*}\etiquette[1.6.2.2]{XVIII.eq1.6.2.2}

\subsubsection{}\etiquette[1.6.3]{XVIII.subsubsec1.6.3}
在\pageoriginale \hyperref[XVIII.lem1.6.1]{1.6.1} 的假设下，假定 $X$ 在域 $k$ 的谱上固有
% SGA-SRC-001141
$$
p:X\longrightarrow \Spec(k){\quoi}.
$$
正合列
$$
0\longrightarrow \GG_{m}(\mathfrak{m})\longrightarrow
\GG_{m}\longrightarrow i_{\ast}\GG_{m}\longrightarrow 0
$$
定义了 $\Spec (k)$ 上 fppf 层上同调的一个正合列，容易直接将它解释为
\begin{equation*}
 0\rightarrow p_{\ast}\GG_{m}(\mathfrak{m})\rightarrow
p_{\ast}\GG_{m_{X}}\rightarrow
(pi)_{\ast}\GG_{m_{Y}}\xrightarrow{j}\tag{1.6.3.1}
 R^{1}p_{\ast}\GG_{m}(\mathfrak{m})\rightarrow
\sheafPic_{X/k}\rightarrow \sheafPic_{Y/k}{\quoi}.
\end{equation*}\etiquette[1.6.3.1]{XVIII.eq1.6.3.1}

层 $p_{\ast}\GG_{m_{X}}$ 和 $(pi)_{\ast}\GG_{m_{Y}}$ 可由有限型\sisi{}{且仿射}的群概形表出。层 $\sheafPic_{X/k}$ 和 $\sheafPic_{Y/k}$ 可由局部有限型群概形表出。由 \hyperref[XVIII.subsubsec1.6.2]{1.6.2}，层 $p_{\ast}\GG_{m}(\mathfrak{m})$ 和 $R^{1}p_{\ast}\GG_{m}(\mathfrak{m})$ 可表。由 \eqhyperref[XVIII.par1.5.13.1]{1.5.13.1}，前者由一个连通（有限型）群概形表出；后者局部有限型。

\begin{enonce*}{定义 1.6.4}\etiquette[1.6.4]{XVIII.def1.6.4}
在上述假设下，以 $\sheafPic_{\mathfrak{m},X/k}$ 表示表出 $R^{1}p_{\ast}\GG_{m}(\mathfrak{m})$ 的群概形。以 $\sheafPic^{0}_{\mathfrak{m},X/k}$ 表示其单位分支，并以 $\sheafPic^{\tau}_{\mathfrak{m},X/k}$ 表示 $\sheafPic_{\mathfrak{m},X/k}/\sheafPic^{0}_{\mathfrak{m},X/k}$ 的挠子群在 $\sheafPic_{\mathfrak{m},X/k}$ 中的逆像。
\end{enonce*}

概形 $\sheafPic_{\mathfrak{m},X/k}$ 表出由如下预层生成的 fppf 层：它将 $k$ 上的每个概形 $S$ 对应到 $\fprod{S}{X}{k}$ 上沿 $\fprod{S}{Y}{k}$ 被平凡化的可逆模的同构类集合（\hyperref[XVIII.subsubsec1.5.8]{1.5.8}）。

当 $k$ 代数闭且 $n$ 在 $k$ 中可逆时，Kummer 正合列 \hyperref[XVIII.lem1.6.1]{1.6.1} 给出一个同构
\begin{equation*}
H^{1}_{c}(U,𝛍_{n})\xrightarrow{\sim}
{}_{n}\sheafPic^{\tau}_{\mathfrak{m},X/k}(k){\quoi}.\tag{1.6.4.1}
\end{equation*}\etiquette[1.6.4.1]{XVIII.eq1.6.4.1}

\setcounter{subsubsection}{4}
\subsubsection{}\etiquette[1.6.5]{XVIII.subsubsec1.6.5}
设\pageoriginale $X$ 是代数闭域 $k$ 上的光滑曲线，$\bar{X}$ 是以 $X$ 为稠密开子概形的光滑射影曲线，$Y$ 是 $\bar{X}$ 的一个闭子概形，使得 $X=\bar{X}-Y$，$\mathfrak{m}$ 是定义 $Y$ 的理想，$n$ 是在 $k$ 中可逆的整数，$G$ 是被 $n$ 零化的有限阿贝尔群。

群 $\sheafPic_{\mathfrak{m},\bar{X}/k}$ 非典范地是子群 $\sheafPic^{0}_{\mathfrak{m},\bar{X}/k}$ 与一个自由阿贝尔群的积。因此，由 \eqhyperref[XVIII.eq1.6.4.1]{1.6.4.1}、\eqhyperref[XVIII.eq1.6.2.2]{1.6.2.2} 及 \hyperref[XVIII.prop1.5.14]{1.5.14}，有一个典范的复合同构
\begin{gather*}
\Hom(H^{1}_{c}(X,𝛍_{n}),G)\xleftarrow{\sim}
\Hom({}_{n}\sheafPic^{\tau}_{\mathfrak{m},\bar{X}/k}(k),G)=\Hom({}_{n}\sheafPic^{0}_{\mathfrak{m},\bar{X}/k},G)\\
\xrightarrow{\sim}
\Ext^{1}(\sheafPic^{0}_{\mathfrak{m},\bar{X}/k},G)\xleftarrow{\sim}
\Ext^{1}(\sheafPic_{\mathfrak{m},\bar{X}/k},G)\xrightarrow{\sim} H^{1}(X,G)
\end{gather*}

若 $u:X\rightarrow Y$ 是 $k$ 上光滑曲线之间的有限态射，则由 \hyperrefext[XVII][XVII.subsec6.3]{6.3}. 和 \eqhyperref[XVIII.eq1.5.15.2]{1.5.15.2} 得到下图交换
\begin{equation*}
\vcenter{
\xymatrix@C=2cm@R=1.5cm{
\Hom(H^{1}_c(Y,𝛍_n),G)  \ar[r]^(.55){\sim} \ar[d]^{\Hom(\Tr_u,G)} &
H^{1}(Y,G)               \ar[d]^{u^{\ast}}\\
\Hom(H^1_c(X,𝛍_n),G)    \ar[r]^(.55){\sim} & H^{1}(X,G).
}}\tag{1.6.5.1}
\end{equation*}\etiquette[1.6.5.1]{XVIII.eq1.6.5.1}

\begin{enonce*}{引理 1.6.6}\etiquette[1.6.6]{XVIII.lem1.6.6}
设 $X$ 是代数闭域 $k$ 上的光滑曲线，$n$ 是在 $k$ 中可逆的整数，$u:X'\rightarrow X$ 是一个驯分歧的平展覆盖（tame étale covering）。则态射
$$
\Tr_{u}:H^{1}_{c}(X',\ZZ/n)\longrightarrow
H^{1}_{c}(X,\ZZ/n)
$$
与态射
$$
u^{\ast}:H^{1}(X,\ZZ/n)\longrightarrow
H^{1}(X',\ZZ/n){\quoi}.
$$
的转置（非典范地\nde{若已选定一个同构 $𝛍_n\simeq \mathbf{Z}/n$，则此同构是典范的。}）同构。
\emph{第一种证明。}由于 $𝛍_{n}$ 与 $\ZZ/n$ 同构，\hyperref[XVIII.lem1.6.6]{1.6.6} 是 \eqhyperref[XVIII.eq1.6.5.1]{1.6.5.1} 中取 $G=\ZZ/n$ 的特殊情形。
\end{enonce*}

\noindent
\emph{第二种证明。}
若 $k$ 是复数域 $\CC$，则由比较定理，态射 $\Tr_{u}$ 和 $u^{\ast}$ 等同于由拓扑空间的连续映射 $u:X'(\CC)\rightarrow X(\CC)$ 所定义的类似态射：
\begin{align*}
\Tr_{u} &: H^{1}_{c}(X'(\CC),\ZZ/n) \longrightarrow
H^{1}_{c}(X(\CC), \ZZ/n)\quad \text{和}\\
u^{\ast} &: H^{1}(X(\CC),\ZZ/n)\longrightarrow H^{1}(X'(\CC),\ZZ/n){\quoi}.
\end{align*}
这\pageoriginale{}两个同态在 Poincaré 对偶下互为转置。

由 Lefschetz 原理，当 $k$ 的特征为 $0$ 时，断言 \hyperref[XVIII.lem1.6.6]{1.6.6} 仍成立（事实上，所讨论的上同调群在代数闭基域扩张下不变）。

若 $k$ 的特征为 $p>0$，设 $W(k)$ 为 $k$ 上的 Witt 向量环。设 $\bar{X}$ 是 $X$ 的光滑射影完备化。由（\SGA 1 III 7.4），曲线 $\bar{X}$ 提升为 $W(k)$ 上的一条光滑射影曲线 $\bar{X}_{1}$。由于 $\bar{X}$ 光滑，每一点 $s_{i}\in S=\bar{X}-X$ 都提升为 $\bar{X}_{1}$ 在 $W(k)$ 上的一个截面 $s'_{i}$；置 $S_{1}=⋃_{i}s'_{i}(\Spec(W(k)))$，$X_{1}=\bar{X}_{1}-S_{1}$。

回顾（\SGA 1 XIII 2.4），$X$ 的覆盖 $X'$ 提升为 $X_{1}$ 的一个平展覆盖 $X'_{1}$，而且 $X'_{1}$ 由 $W(k)$ 上的一条光滑射影曲线 $\bar{X}'_{1}$ 去掉有限个两两不交截面的并 $S'_{1}$ 得到。
$$
\xymatrix@=.5cm{
                  &     X_1 \ar[ddddd]^{f_1}       &                 &
                  X\ar[ll]\ar[ddddd]\\
 & & &\\
X'_1 \ar[ruu]_{u_1}\ar[rddd]_{f'_1}&                & X' \ar[ll]\ar[ruu]_{u} \ar[rddd] &   \\
& & &\\
& & &\\
                  &    \Spec(W(k)) &                 &\Spec(k).\ar[ll]
}
$$

平展层 $R^{1}f_{1\ast}(\ZZ/n)$、$R^{1}f'_{1\ast}(\ZZ/n)$、$R^{1}f_{1!}(\ZZ/n)$ 和 $R^{1}f'_{1!}(\ZZ/n)$ 局部常值，且其形成与基变换相容，\pageoriginale{}这由联系 $\bar{X}_{1}$、$X_{1}$、$S_{1}$ 或 $\bar{X}'_{1}$、$X'_{1}$、$S'_{1}$ 的上同调的正合列得到。

因此，态射
$$
\Tr_{u_{1}}:R^{1}f'_{1!}\ZZ/n\longrightarrow
R^{1}f_{1!}\ZZ/n
$$
或
$$
u^{\ast}_{1}:R^{1}f_{1\ast}\ZZ/n\longrightarrow
R^{1}f'_{1\ast}\ZZ/n
$$
的特殊几何纤维与一般几何纤维同构，于是归结为已经处理过的 $k$ 特征为零的情形。

\begin{enonce*}{引理 1.6.7}\etiquette[1.6.7]{XVIII.lem1.6.7}
设 $X$ 是特征指数为 $p$ 的代数闭域 $k$ 上的连通光滑曲线，$n$ 是与 $p$ 互素的整数。存在一个主覆盖\emph{（即有限平展 Galois（满射）覆盖）} $u:X'\rightarrow X$，其次数整除 $n$ 的某个幂，使得迹态射
$$
\Tr_{u}:H^{1}_{c}(X', \ZZ/n)\longrightarrow
H^{1}_{c}(X,\ZZ/n)
$$
为零。
\end{enonce*}

回顾群 $H^{1}(X,\ZZ/n)$ 是有限的（\hyperrefext[IX][IX.thm4.6]{4.6} 或 \hyperrefext[XVI][XVI.thm5.2]{5.2}）。设 $u:X'\rightarrow X$ 是 $X$ 的最大连通阿贝尔主覆盖，其 Galois 群被 $n$ 零化（其 Galois 群记作 $H_{1}(X,\ZZ/n)$）。由构造，同态
$$
u^{\ast}:H^{1}(X,\ZZ/n)\longrightarrow H^{1}(X', \ZZ/n)
$$
为零。因此，由 \hyperref[XVIII.lem1.6.6]{1.6.6}，同态
$$
\Tr_{u}:H^{1}_{c}(X',\ZZ/n)\longrightarrow
H^{1}_{c}(X,\ZZ/n)
$$
为零。

回顾下面的引理（\EGA IV 15.6.5）：

\begin{enonce*}{引理 1.6.8}\etiquette[1.6.8]{XVIII.lem1.6.8}
设 $f:X\rightarrow S$ 是光滑态射，$x$ 是 $f$ 的一个截面。则存在 $X$ 的一个包含 $x$ 的 Zariski 开子集 $U$，使得 $f|U$ 的几何纤维连通。
\end{enonce*}

\begin{enonce*}{基本引理 1.6.9}\etiquette[1.6.9]{XVIII.lemfond1.6.9}
设\pageoriginale $f:X\rightarrow S$ 是可紧化的光滑曲线（\hyperref[XVIII.def1.1.2]{1.1.2}）\emph{（\hyperrefext[XVII][XVII.subsec3.2]{3.2}）}，$x$ 是 $X$ 的一个几何点，$s$ 是其在 $S$ 中的像，$n\geq 1$ 是在 $S$ 上可逆的整数。存在 $s$ 在 $S$ 中的一个平展邻域 $V$，以及 $x$ 在 $f^{-1}(V)=\fprod{V}{X}{S}$ 中的一个平展邻域 $U$
\begin{equation*}
\vcenter{
\xymatrix@R=1.5cm@C=1cm{
x \ar[r]\ar@{-}[d]& U\ar[rd]_{f'}\ar[r]^{u} & f^{-1}(V)\ar[r] \ar[d]^{f_v}& X\ar[d]^{f}\\
s \ar[rr]&   & V\ar[r]         & S,
}}\tag{1.6.9.1}
\end{equation*}\etiquette[1.6.9.1]{XVIII.eq1.6.9.1}%
使得 $R^{0}f'_{!}\sheaf{\ZZ}/n=0$，箭头\emph{ \hyperrefext[XVII][XVII.subsec6.2]{6.2}}
$$
\Tr_{u}:R^{1}f'_{!}\sheaf{\ZZ}/n\longrightarrow
R^{1}f_{v!}\sheaf{\ZZ}/n
$$
为零，且迹态射
$$
\Tr_{f'}:R^{2}f'_{!}\sheaf{\ZZ}/n(1)\longrightarrow
\sheaf{\ZZ}/n
$$
为同构。
\end{enonce*}

为使 $R^{0}f'_{!}\sheaf{\ZZ}/n=0$，只需 $f'$ 拟仿射；为使 $\Tr_{f'}$ 为同构，只需 $f'$ 的几何纤维连通（\hyperref[XVIII.lem1.1.9]{1.1.9}）；由 \hyperref[XVIII.lem1.6.8]{1.6.8}，关于 $S$ 对平展拓扑局部地，容易获得 $f'$ 的这些性质。还需对具有几何连通纤维的 $X$，证明存在图 \eqhyperref[XVIII.eq1.6.9.1]{1.6.9.1}，使所讨论的箭头 $\Tr_{u}$ 为零：之后只需缩小 $U$、$V$，就能使 $f'$ 满足另外两个所需的性质。

问题关于 $S$ 在 $s$ 的邻域中是局部的；由 \hyperrefext[XVII][XVII.thm5.2.6]{5.2.6}（基变换），立即归结为 $S$ 是 Noether 概形的情形。通过取极限，并利用 \hyperrefext[XVII][XVII.thm5.3.6]{5.3.6}（可构造性），归结为 $S$ 是严格局部 Noether 概形，且闭点为 $s$ 的像的情形。

设 $t$ 为 $S$ 的一个几何点。由 \hyperref[XVIII.lem1.6.7]{1.6.7}，存在一个次数与 $k(t)$ 的特征指数
% SGA-SRC-001143
$p$ 互素的主覆盖 $u:X^{1}_{t}\rightarrow X_{t}$，使得箭头
$$
\Tr_{u}:H^{1}_{c}(X^{1}_{t},\ZZ/n)\longrightarrow
H^{1}_{c}(X_{t},\ZZ/n)
$$
为\pageoriginale{}零。由无环性定理 \hyperrefext[XV][XV.thm2.1]{2.1}（或 \hyperrefext[XV][XV.lem2.6]{2.6}），存在 $x$ 在 $X$ 中的一个平展邻域 $v:U\rightarrow X$，使得 $X^{1}_{t}$ 在 $U_{t}$ 上分裂，因而 $v_{t}:U_{t}\rightarrow X_{t}$ 通过 $X^{1}_{t}$ 分解。所以仍有
$$
\Tr_{v_{t}}=0:H^{1}_{c}(U_{t},\ZZ/n)\longrightarrow
H^{1}_{c}(X_{t},\ZZ/n){\quoi}.
$$

因此，对于 $S$ 的每个几何点 $t$，都存在 $x$ 在 $X$ 中的一个平展邻域
$$
\xymatrix{
U \ar[r]^{v} \ar[rd]_{f'}& X \ar[d]^{f}\\
  & S
}
$$
使得
\begin{equation*}
\Tr_{v}:R^{1}f'_{!}\ZZ/n\longrightarrow
R^{1}f_{!}\ZZ/n\tag{1.6.9.2}
\end{equation*}\etiquette[1.6.9.2]{XVIII.eq1.6.9.2}%
在 $t$ 的纤维为零（\hyperrefext[XVII][XVII.thm5.2.6]{5.2.6}）。

当 $U$ 变化时，箭头 \eqhyperref[XVIII.eq1.6.9.2]{1.6.9.2} 的像组成可构造层 $R^{1}f_{!}\ZZ/n$ 的一个滤过递减子层系统，它在 $S$ 的任意给定几何点最终变为零。利用 Noether 归纳法可验证，这样的系统总是稳定的，且稳定值为零。因此存在 $U$，使 \eqhyperref[XVIII.eq1.6.9.2]{1.6.9.2} 为零，\hyperref[XVIII.lemfond1.6.9]{1.6.9} 得证。

% END native input: chapters/exp-18-part3-zh.tex
% BEGIN native input: chapters/exp-18-part4-zh.tex
\section{迹态射}\etiquette[2]{XVIII.sec2}\pageoriginale

本 \S 致力于证明定理 \hyperref[XVIII.thm2.9]{2.9} 及 \hyperref[XVIII.prop2.12]{2.12}. \hyperref[XVIII.thm2.9]{2.9} 的证明仅使用第 \hyperref[XVIII.sec1]{1} \S 中的第 1 小节.

\begin{enonce*}{引理 2.1}\etiquette[2.1]{XVIII.lem2.1}
设 $f:X\rightarrow S$ 为相对维数 $\leq d$ 的可紧化态射 \emph{(\hyperrefext[XVII][XVII.def3.2.1]{3.2.1})}, $j:U\hookrightarrow X$ 为 $X$ 的有限型开子概形, 其补在每个纤维中的维数 $< d$ (例如, $U$ 在 $X$ 的每个纤维中稠密). 若 $F$ 为 $X$ 上的挠层, 则典范箭头
$$
R^{2d}(fj)_{!}j^{\ast}F\xrightarrow{\sim} R^{2d}f_{!}F
$$
为同构.
\end{enonce*}

设 $i$ 为 $Y=X-U$ 到 $X$ 的嵌入. 正合列 (XVII \hyperref[XVII.eq5.1.16.2]{5.1.16.2}) 给出正合列
$$
R^{2d-1}(fi)_{!}i^{\ast}F\longrightarrow
R^{2d}(fj)_{!}j^{\ast}F\longrightarrow R^{2d}f_{!}F\longrightarrow
R^{2d}(fi)_{!}i^{\ast}F
$$
由 (\hyperrefext[XVII][XVII.cor5.2.8.1]{5.2.8.1}), 其两端的项为零.

\begin{enonce*}{引理 2.2}\etiquette[2.2]{XVIII.lem2.2}
设 $f:X\rightarrow S$ 为相对维数 $\leq d$ 的可紧化态射, $u_{i}:X_{i}\rightarrow X$ 为有限型分离平展态射族; 置 $u_{ij}:X_{ij}=\fprod{X_{i}}{X_{j}}{X}\rightarrow X$. 对 $X$ 上的每个挠层 $F$, 箭头 \emph{(XVII \hyperref[XVII.eq6.2.7.2]{6.2.7.2})} 所成的序列
$$
\bigoplus_{i,j} R^{2d}(fu_{ij})_{!}u^{\ast}_{ij}F\rightrightarrows
\bigoplus_{i} R^{2d}(fu_{i})_{!}u^{\ast}_{i}F\rightarrow
R^{2d}f_{!}F\rightarrow 0
$$
正合.
\end{enonce*}

这立即由函子 $R^{2d}f_{!}$ 的右正合性及 \hyperrefext[XVII][XVII.prop6.2.9]{6.2.9} 得到.

\begin{enonce*}{引理 2.3}\etiquette[2.3]{XVIII.lem2.3}
设 $f:X\rightarrow S$ 为有限表示的平坦态射. $X$ 中 $f$ 为 Cohen--Macaulay 的开子集 \emph{(\EGA IV 6.8.1 及 12.2.1 (vii))} 相对有限表示, 且在每个纤维中稠密.
\end{enonce*}

此\pageoriginale 问题在 $S$ 上是局部的, 归约后可假设其为 Noether 概形, 从而得到第一个断言. 第二个断言逐纤维验证, 因而可取 $S$ 为域的谱. $X$ 在 $X$ 的极大点处的局部环是 Artin 环, 因而为 Cohen--Macaulay 环; 由于“Cohen--Macaulay”是“开”性质, 断言得证.

\begin{enonce*}{引理 2.4}\etiquette[2.4]{XVIII.lem2.4}
设 $f:X\rightarrow S$ 为有限表示、平坦且 Cohen--Macaulay 的态射, $S$ 拟分离. $X$ 的每个点都有一个 (Zariski) 开邻域 $U$, 使存在从 $U$ 到标准向量丛 $\EE^{d}_{S}$ 的拟有限、平坦且有限表示的 $S$ 态射.
\end{enonce*}

这是 \EGA IV 15.4.3 的特殊情形.

\begin{enonce*}{引理 2.5}\etiquette[2.5]{XVIII.lem2.5}
设 $\underline{u}=(u_{1},\ldots,u_{d})$ 与 $\underline{v}=(v_{1},\ldots,v_{d})$ 为维数为 $d$、极大理想为 $\mathfrak{m}$ 的 Cohen--Macaulay 局部环 $A$ 的两组参数系. 存在 $A$ 的参数系序列 $\underline{w}_{0},\ldots,\underline{w}_{n}$, 使 $\underline{w}_{0}=\underline{u}$、$\underline{w}_{n}=\underline{v}$, 且每个 $\underline{w}_{i+1}$ 只改变 $\underline{w}_{i}$ 中的一个参数而得.
\end{enonce*}

对 $d$ 作归纳, 当 $d<2$ 时断言为空泛的; 因而假设 $d\geq 2$. 元素 $x\in \mathfrak{m}$ 为 $A/u_{1}$ 正则 (\resp $A/v_{1}$ 正则) 的充要条件是 $x$ 不属于 $\Ass(A/u_{1})$ (\resp $\Ass(A/v_{1})$) 的任何成员; 由于素理想 $\neq \mathfrak{m}$ 的任何有限并都不等于 $\mathfrak{m}$, 存在同时为 $A/u_{1}$ 正则与 $A/v_{1}$ 正则的 $x\in \mathfrak{m}$, 因而也存在参数系 $\underline{a}=(u_{1},x,\ldots)$ 与 $\underline{b}=(v_{1},x,\ldots)$. 在 $A/u_{1}$ 中对 $\underline{u}-\{u_{1}\}$ 与 $\underline{a}-\{u_{1}\}$、在 $A/x$ 中对 $\underline{a}-\{x\}$ 与 $\underline{b}-\{x\}$、在 $A/v_{1}$ 中对 $\underline{b}-\{v_{1}\}$ 与 $\underline{v}-\{v_{1}\}$ 应用归纳假设, 即得结论.

\begin{enonce*}{引理 2.6}\etiquette[2.6]{XVIII.lem2.6}
设 $X$ 为代数闭域 $k$ 上的有限型 Cohen--Macaulay 概形, $u$、$v:X\rightarrow \EE^{d}_{k}$ 为从 $X$ 到标准向量丛的两个拟有限平坦态射. $X$ 的每个闭点 $x$ 都有一个 (Zariski) 开邻域 $U$, 使存在态射序列\pageoriginale $w_{0},\ldots,w_{n}:U\rightarrow \EE^{d}_{k}$, 满足 $w_{0}=u|U$、$w_{n}=v|U$, 且 $w_{i+1}$ 与 $w_{i}$ 仅有一个坐标不同.
\end{enonce*}

先假设 $u(x)=v(x)=0$. 此时, 在 $x$ 处从 $X$ 到 $\AAA^{d}_{k}$ 的拟有限平坦态射的芽,

若将 $x$ 映到 $0$, 就等同于局部环 $\Oo_{X,x}$ 的参数系, 而 \hyperref[XVIII.lem2.6]{2.6} 由 \hyperref[XVIII.lem2.5]{2.5} 得到. 在一般情形下, 先用态射链将 $u$ 接到 $u-u(x)$ (\resp 将 $v$ 接到 $v-v(x)$):
$$
u-(u_{1}(x),\ldots,u_{i}(x),0,\ldots,0)(\text{\resp}
v-(v_{1}(x),\ldots,v_{i}(x),0,\ldots,0)){\quoi}.
$$

\setcounter{subsection}{6}
\subsection{}\etiquette[2.7]{XVIII.subsec2.7}
设有一对可复合的 $S$ 可紧化态射 $X\xrightarrow{g} Y\xrightarrow{f}S$. 假设 $f$ 的相对维数 $\leq d$, $g$ 的相对维数 $\leq
e$. 则对 $X$ 上的每个挠层 $F$, 复合谱序列
\begin{equation*}
R^{p}f_{!}R^{q}g_{!}F\Longrightarrow R^{p+q}(fg)_{!}F\tag{2.7.1}
\end{equation*}\etiquette[2.7.1]{XVIII.eq2.7.1}%
给出同构 (利用最高次循环项的论证)
\begin{equation*}
R^{2d}f_{!}R^{2e}g_{!}F\simeq R^{2(d+e)}(fg)_{!}F{\quoi}.\tag{2.7.2}
\end{equation*}\etiquette[2.7.2]{XVIII.eq2.7.2}

\subsection{}\etiquette[2.8]{XVIII.subsec2.8}
设
$$
a^{d}:\EE^{d}_{S}\longrightarrow S
$$
为凝聚 (即拟紧、拟分离) 概形 $S$ 上的标准向量丛. 若 $F$ 为 $S$ 上与 $S$ 的剩余特征互素的挠层, 则如下对 $d$ 作归纳, 定义迹态射
\begin{equation*}
\Tr_{a^d}:R^{2d}a^{d}_{!}F(d)\xrightarrow{\sim} F{\quoi}.\tag{2.8.1}
\end{equation*}\etiquette[2.8.1]{XVIII.eq2.8.1}

当 $d=0$ 时, $\Tr_{a^0}$ 为恒等态射. 当 $d=1$ 时, $\Tr_{a^1}$ 为同构 (\hyperref[XVIII.prop1.1.6]{1.1.6}) (\cf \hyperref[XVIII.lem1.1.9]{1.1.9}). 最后, 有典范同构
$$
\EE^{d+1}_{S}=\fprod{\EE^{1}_{S}}{\EE^{d}_{S}}{S}=\EE^{1}_{\EE_{S}^d}
$$
将\pageoriginale $\Tr_{a^{d+1}}$ 定义为同构 $\Tr_{a^{d}_{S}}$ 与
$Ra^{d}_{S!}(\Tr_{a^1_{\EE^d_S}})$
的复合 (\cf \hyperref[XVIII.thm2.9]{2.9} 的图表 Var 3).

\eqhyperref[XVIII.eq2.8.1]{2.8.1} 的源 (及靶) 的形成与任意基变换相容. 因而每个在 $\EE^{d}_{S}$ 上作用的 $S$ 上的代数群, 都作用于这些层; 连通代数群的作用必为平凡作用. 因此, 包含于仿射群中的坐标置换群在 \eqhyperref[XVIII.eq2.8.1]{2.8.1} 的源 (及靶) 上的作用是平凡的.

\setcounter{subsubsection}{1}
\subsubsection{}\etiquette[2.8.2]{XVIII.subsubsec2.8.2}
同构 \eqhyperref[XVIII.eq2.8.1]{2.8.1} 在坐标置换下不变.

也可用杯积解释 \eqhyperref[XVIII.eq2.8.1]{2.8.1}, 从而验证这一点 (XVII \hyperref[XVII.eq5.4.2.2]{5.4.2.2} 及 \hyperrefext[XVII][XVII.lem5.4.3.5]{5.4.3.5}).

\begin{enonce*}{定理 2.9}\etiquette[2.9]{XVIII.thm2.9}
考虑三元组 $(f,d,F)$, 由可紧化态射 \emph{(\hyperrefext[XVII][XVII.subsec6.3]{6.3})} $f:X\rightarrow Y$、整数 $d$ 及 $Y$ 上与 $Y$ 的剩余特征互素的挠层 $F$ 组成, 其中态射 $f$ 满足条件 $(\ast)_{d}$: 存在 $X$ 的开子集 $U$, 使 $f$ 在 $U$ 上的限制为有限表示的平坦态射, 其纤维维数为 $\leq d$, 并使 $X-U$ 的纤维维数 $< d$.

有且仅有一种方式, 对每个上述三元组 $(f,d,F)$ 关联迹态射
$$
\Tr_{f}:R^{2d}f_{!}f^{\ast}F(d)\longrightarrow F
$$
使下列条件满足:
\begin{itemize}
\item[(Var 1)] (函子性). 迹态射对 $F$ 具有函子性.

\item[(Var 2)] (与基变换相容). 对每个笛卡尔方块
$$
\xymatrix{
X' \ar[r]_{g'}\ar[d]^{f'}& X\ar[d]^{f}\\
Y' \ar[r]_{g} & Y
}
$$
其中\pageoriginale $Y'$ 拟紧且拟分离, 图表
$$
\xymatrix@R=1.25cm{
R^{2d}f'_!f'^{\ast}g^{\ast}F(d) \ar@{=}[r]\ar[d]^{Tr_{f'}}& R^{2d}f'_!
g'^{\ast} f^{\ast} F(d) & g^{\ast}
R^{2d}f_!f^{\ast} F(d) \ar[l]_{\sim} \ar[d]^{g^{\ast}(Tr_f)}\\
g^{\ast}F \ar@{=}[rr]&            & g^{\ast}F
}
$$
交换, 其中水平同构为基变换箭头 \emph{\eqhyperrefext[XVII][XVII.eq5.2.2.1]{5.2.2.1}}.

\item[(Var 3)] (与复合相容). 对每一对可复合的 $Z$ 可紧化态射
$$
X\xrightarrow{g}Y\xrightarrow{f}Z{\quoi},
$$
以及使 $f$ 满足 $(\ast)_{d}$、$g$ 满足 $(\ast)_{e}$ 的整数 $d$ 与 $e$, 和 $Z$ 上的挠层 $F$, 复合态射 $fg$ 满足 $(\ast)_{d+e}$, 且图表
$$
\xymatrix@C=2.25cm{
R^{2d}f_! R^{2e} g_! g^{\ast} f^{\ast}F(d)(e) \ar@{<->}[d]_{\text{\rotatebox{90}{$\backsim$}}} \ar[r]^(.6){R^{2d}f_!(\Tr_g)} &
R^{2d}f_! f^{\ast}F(d) \ar[d]^{\Tr_f}\\
R^{2(d+e)} (fg)_! (fg)^{\ast} F(d+e) \ar[r]^(.65){\Tr_{fg}} & F
}
$$
交换, 其中竖直同构为箭头 \eqhyperref[XVIII.eq2.7.2]{2.7.2}.

\item[(Var 4)] (规范化条件).
\begin{itemize}
\item[(I)] 若 $f$ 为秩为 $n$ 的有限局部自由态射, 且 $d=0$, 则复合态射
$$
F\longrightarrow
f_{\ast}f^{\ast}F=f_{!}f^{\ast}F\xrightarrow{\tr_{f}}F
$$
为乘以 $n$.

\item[(II)] 若\pageoriginale $f$ 为标准仿射直线 $f:\EE^{1}_{Y}\rightarrow Y$, 且 $d=1$, 则迹态射
$$
\Tr_{f}: R^{2}f_{!}f^{\ast}F(1)\longrightarrow F
$$
为同构 \eqhyperref[XVIII.eq2.8.1]{2.8.1}.
\end{itemize}
\end{itemize}
\end{enonce*}

\begin{enumerate}
\renewcommand{\theenumi}{\alph{enumi}}
\renewcommand{\labelenumi}{\theenumi)}
\item 对于 $d=0$, 条件 $(\ast)_{d}$ 表示 $f$ 拟有限、平坦且有限表示. 由 (\hyperrefext[XVII][XVII.thm6.2.3]{6.2.3}), 对 $d=0$, 迹态射必与 \loccit{} 中所构造的态射一致.

当 $f$ 为标准仿射空间 $\AAA^{d}_{S}$ 到 $S$ 的投影时, 由 (Var~4)~(II) 及 (Var~3) 可知, $\Tr_{f}$ 必须为同构 \eqhyperref[XVIII.eq2.8.1]{2.8.1}.

\item 设 $f:X\rightarrow Y$ 为满足条件 $(\ast)_{d}$ 的可紧化态射: 存在从 $X$ 到标准向量丛 $a^{d}:\EE^{d}_{Y}\rightarrow Y$ 的拟有限、平坦且有限表示的 $Y$ 态射 $u$.

设 $t(f,u)$ 为迹态射 \eqhyperref[XVIII.eq2.8.1]{2.8.1} 与 (\SGA \hyperrefext[XVII][XVII.thm6.2.3]{6.2.3}) 的“复合”态射 (\cf (Var 3))
$$
t(f,u)=\Tr_{a^d}\circ R^{2d}a^{d}_{!}(\Tr_{u}){\quoi}.
$$

此态射 $t(f,u)$ 满足 (Var 1) 及 (Var 2). 通过基变换, 为验证 $t(f,u)$ 只依赖于 $f$ 及 $d$, 而不依赖于 $u$, 只须在 $Y$ 为代数闭域 $k$ 的谱时验证之. 由 \hyperref[XVIII.lem2.2]{2.2}, 此问题在 $X$ 上是局部的, 所以由 \hyperref[XVIII.lem2.6]{2.6}, 只须证明
$t(f,u)=t(f,v)$
其中 $u$ 与 $v$ 只有一个坐标不同; 由 \hyperref[XVIII.subsubsec2.8.2]{2.8.2}, 可以假设此坐标为第一个. 置 $E=\EE^{d-1}_{Y}$, 于是有交换图表
$$
\xymatrix{
                              &    &\EE^1_{E} \ar[rdd]_(.3){a} &  &  \\
X \ar[rrrr]_{v}\ar[rru]^{u} \ar[rdd]_{f} \ar[rrrd]_{x}&    &                 &
                              &\EE^{1}_E \ar[ld]_{a}\\
                              &    &                 &E \ar[lld]^w& \\
                              & Y  &                 &  &
}
$$
根据\pageoriginale 构造 (\hyperref[XVIII.subsec2.8]{2.8}),
$$
\Tr_{wa}=\Tr_{w}\circ R^{2(d-1)}w_!(\Tr_{a}){\quoi}\;;
$$
推得 $t(f,u)$ 为 $\Tr_{w}$ 与 $t(x,u)$ 的“复合”:
\begin{align*}
t(f,u) &= \Tr_{wa}\circ R^{2d}(wa)_{!}(\Tr_{u})\\
       &= \Tr_{w}\circ R^{2(d-1)}w_{!}(\Tr_{a})\circ
       R^{2d}(wa)_{!}(\Tr_{u})\\
       &= \Tr_{w}\circ R^{2(d-1)}w_{!}(\Tr_{a}\circ
       R^{2}a_{!}(\Tr_{u})){\quoi}.
\end{align*}
此外 (\hyperref[XVIII.lem1.1.7]{1.1.7}), $t(x,u)$ 恰为 \hyperref[XVIII.prop1.1.6]{1.1.6} 的态射 $\Tr_{x}$; 因而 $t(x,u)=t(x,v)=\Tr_{x}$ 且 $t(f,u)=t(f,v)$.

以上面构造的态射 $t(f,u)$ 记为 $\Tr_{f}$. 由 \hyperrefext[XVII][XVII.thm6.2.3]{6.2.3}, 此迹态射满足 (Var 1)、(Var 2)、$e=0$ 时的 (Var 3) 及 (Var 4). 只要它有定义, 就必与所求迹态射一致, 如果后者存在.

\item 设 $f:X\rightarrow Y$ 为可紧化的 Cohen--Macaulay 态射, 平坦且有限表示, 纯相对维数为 $d$. 设 $(U_{i})_{i\in I}$ 为 $X$ 的拟紧开子集覆盖, 使存在从 $U_{i}$ 到 $\EE^{d}_{Y}$ 的拟有限平坦 $Y$ 态射 (\hyperref[XVIII.lem2.4]{2.4}). 由 \hyperref[XVIII.lem2.2]{2.2} 及 b), 存在唯一态射 $\Tr_{f}$ 使下图交换, 这里采用 \hyperref[XVIII.lem2.2]{2.2} 的记号:
$$
\xymatrix{
\bigoplus_{i,j}
R^{2d}(fu_{ij})_!(fu_{ij})^{\ast} F(d) \ar@<2pt>[r]
\ar@<-2pt>[r]&\bigoplus_{i}
R^{2d}(fu_i)_!(fu_i)^{\ast}F(d)
\ar[rdd]_{\bigoplus_{i}\Tr_{fu_i}}\ar[r]
&
R^{2d}f_!f^{\ast}F(d) \ar[dd]^{\Tr_f}\\
& & \\
& & F
}
$$

此迹态射不依赖于所选取的开覆盖 (用并覆盖比较两个覆盖); 它满足 (Var 1)、(Var 2)、$e=0$ 时的 (Var 3) 及 (Var 4), 这立即由 b) 得到.

\item 设 $f:X\rightarrow Y$ 为满足 $(\ast)_{d}$ 的可紧化态射, $j:U\hookrightarrow X$\pageoriginale 为 $(\ast)_{d}$ 所保证的开子集. 根据 \hyperref[XVIII.lem2.3]{2.3}, 将 $U$ 换成较小开子集后, 可以假设 $fj$ 为 Cohen--Macaulay 且纯维数为 $d$.

\noindent
由 \hyperref[XVIII.lem2.1]{2.1}, 典范态射
$$
R^{2d}(fj)_{!}(fj)^{\ast}F(d)\xrightarrow{\sim}
R^{2d}f_{!}f^{\ast}F(d)
$$
为同构, 将 $\Tr_{f}$ 定义为复合
$$
\Tr_{f}:R^{2d}f_{!}f^{\ast}F(d)\xleftarrow{\sim}
R^{2d}(fj)_{!}(fj)^{\ast}F(d)\xrightarrow{\Tr_{fj}}F{\quoi}.
$$
由 c) 中建立的 $e=0$ 时的 (Var 3), 立即可知 $\Tr_{f}$ 与 $U$ 的选择无关. 此迹态射满足 (Var 1)、(Var 2)、$e=0$ 时的 (Var 3) 及 (Var 4), 且是唯一能满足所加条件的态射. 还须证明它满足 (Var 3), 为此, 上述论证将我们归结为只考虑分别满足 $(\ast)_{d}$ 与 $(\ast)_{e}$ 的对 $(f,d)$、$(g,e)$.

希望细致整理以下论证的读者, 可先验证下述引理:

\begin{enonce*}{引理 2.9.1}\etiquette[2.9.1]{XVIII.lem2.9.1}
设 $(f,g,h)$ 为三个可复合的 $T$ 可紧化态射, 满足 $(\ast)_{k}$、$(\ast)_{\ell}$ 及 $(\ast)_{m}$.
$$
X\xrightarrow{h}Y\xrightarrow{g}Z\xrightarrow{f}T{\quoi}.
$$
若对 $(f,g)$ 与 $(g,h)$ 满足 {\rm (Var 3)}, 则 $(f,gh)$ 满足 {\rm (Var 3)} 当且仅当 $(fg,h)$ 满足 {\rm (Var 3)}.
\end{enonce*}

结论中的条件也表示, $\Tr_{fgh}$ 为 $\Tr_{f}$、$\Tr_{g}$ 及 $\Tr_{h}$ 的“复合”.

\item 设给定交换图表:
$$
\xymatrix@R=1.25cm@C=1.5cm{
X \ar[r]^{v} \ar[rd]_{g}& \EE^{e}_Y \ar[r]^{u'}\ar[d]^{c} & \EE_Z^{d+e}\ar[d]^{b} \\
            &     Y   \ar[r]^{u} \ar[rd]_{f}                       & \EE_Z^{d}\ar[d]^{a}  \\
            &                                  &     Z
}
$$
相对于\pageoriginale $fg$ 的迹态射为相对于 $ab$ 与 $u'v$ 的迹态射的“复合”; 根据 $ab$ 时的构造以及 $u'v$ 时的 \hyperrefext[XVII][XVII.thm6.2.3]{6.2.3}, 它也是相对于 $a$、$b$、$u'$ 及 $v$ 的迹态射的“复合”. 另一方面, 相对于 $f$ 与 $g$ 的迹态射的“复合”, 是相对于 $a$、$u$、$c$ 及 $v$ 的迹态射的“复合”. 因而必须验证, 对每个笛卡尔图表
\begin{equation*}
\vcenter{
\xymatrix@C=1.25cm@R=1cm{
\EE^{e}_S  \ar[d]^{b'} \ar[r]^{u'}   & \EE^{e}_T \ar[d]^{b}\\
    S      \ar[r]^{u}     &      T
}}\tag{2.9.2}
\end{equation*}\etiquette[2.9.2]{XVIII.eq2.9.2}%
其中 $u$ 拟有限、平坦且有限表示, 迹态射 \eqhyperref[XVIII.eq2.8.1]{2.8.1} 及 \hyperrefext[XVII][XVII.subsec6.2]{6.2} 的图表
\begin{equation*}
\vcenter{
\xymatrix@C=2.25cm{
R^{2e}b_!u'_! u'^{\ast}b {}^{\ast} F(e) \ar[r]^(.55){R^{2e}b_!
  (\Tr_{u'})}\ar@{=}[d]& R^{2e}b_! b^{\ast} F(e)\ar[r]^{\Tr_b}
&  F \ar@{=}[d]\\
u_! R^{2e}b'_! b'^{\ast} u^{\ast} F(e)\ar[r]^(.55){u_!(\Tr_{b'})} & u_!
u^{\ast} F \ar[r]^{\Tr_u}& F
}}\tag{2.9.3}
\end{equation*}\etiquette[2.9.3]{XVIII.eq2.9.3}%
交换.
\end{enumerate}


通过基变换, 只须在 $T$ 为代数闭域 $k$ 的谱时验证之. 若 $S$ 为若干概形 $S_{i}$ 的和, 只须对以 $u_{i}:S_{i}\rightarrow T$ 为基的图表 \eqhyperref[XVIII.eq2.9.2]{2.9.2} 验证 \eqhyperref[XVIII.eq2.9.3]{2.9.3} 的交换性. 因而可以假设 $S$ 为有限次数 $n=[A:k]$ 的局部 Artin $k$ 代数 $A$ 的谱. 此时若通过函子 $u^{\ast}$ (\resp ${u'}^{\ast}$) (VIII \hyperref[VIII.thm1.1]{1.1}) 将 $T$ (\resp $\EE^{e}_{T}$) 上的层等同于 $S$ (\resp $\EE^{e}_{S}$) 上的层, 则态射 $\Tr_{b}$ 等同于态射 $\Tr_{b'}$ (\eqhyperref[XVIII.eq2.8.1]{2.8.1} 与基变换的相容性), 而 $\Tr_{u}$ 与 $\Tr_{u'}$ 等同于乘以 $n$, 断言得证\nde{也可将此相容性归结为情形 $e=1$, 它由结合 \hyperref[XVIII.lem1.1.7]{1.1.7} 及 \hyperref[XVIII.lem1.1.8]{1.1.8} 得到. 此外, 这里所给的论证与 \hyperref[XVIII.lem1.1.8]{1.1.8} 中默然略过的论证类似.}. 至此完成了 \hyperref[XVIII.thm2.9]{2.9} 的证明.

在证明过程中已经\pageoriginale 看到:

\begin{enonce*}{命题 2.10}\etiquette[2.10]{XVIII.prop2.10}
迹态射 \emph{\hyperrefext[XVII][XVII.thm6.2.3]{6.2.3}、XVIII~\hyperref[XVIII.prop1.1.6]{1.1.6}} 及 \eqhyperref[XVIII.eq2.8.1]{2.8.1} 均为迹态射 \hyperref[XVIII.thm2.9]{2.9} 的特殊情形.
\end{enonce*}

\begin{enonce*}[remark]{注 2.10.1}\etiquette[2.10.1]{XVIII.rem2.10.1}
不难证明, \hyperref[XVIII.thm2.9]{2.9} 的迹态射为同构, 当且仅当 $f$ 的几何纤维恰有一个维数为 $d$ 的不可约分支, 且其“重数”与 $n$ 互素. 这里不给出这一事实的直接证明, 它将由第 \hyperref[XVIII.sec3]{3} \S 的“整体对偶定理”立即得到.
\end{enonce*}

\setcounter{subsection}{10}
\subsection{}\etiquette[2.11]{XVIII.subsec2.11}
设 $f:X\rightarrow S$ 与 $g:Y\rightarrow S$ 为两个 $S$ 可紧化态射, 分别满足 $(\ast)_{d}$ 及 $(\ast)_{e}$.
$$
\xymatrix{
     & \fprod{X}{Y}{S} \ar[dd]^{f\times g}\ar[ld]_{\pr_1}\ar[rd]^{\pr_2} & \\
X\ar[rd]_{f} &                                        & Y\ar[ld]^{g} \\
             &    S                                   &
}
$$

Künneth 同构 (\hyperrefext[XVII][XVII.thm5.4.3]{5.4.3})
$$
Rf_{!}(\ZZ/n){\mathop{\otimes}^{\LL}}_{\ZZ/n}Rg_{!}(\ZZ/n)\xrightarrow{\sim}
R(f\times g)_{!}(\ZZ/n)
$$
诱导同构
\begin{equation*}
R^{2d}f_{!}\ZZ/n\otimes
R^{2e}g_{!}\ZZ/n\xrightarrow{\sim} R^{2(d+e)}(f\times
g)_{!}(\ZZ/n)\tag{2.11.1}
\end{equation*}\etiquette[2.11.1]{XVIII.eq2.11.1}

显然 $f\times g$ 满足 $(\ast)_{d+e}$; 有如下相容性:

\begin{enonce*}{命题 2.12}\etiquette[2.12]{XVIII.prop2.12}
下图
$$\vcenter{\etiquette[2.12.1]{XVIII.eq2.12.1}
\xymatrix@R=1.25cm{
R^{2d}f_!\ZZ/n(d) \otimes R^{2e}g_!\ZZ/n (e)
\ar[r]^{\sim} \ar[d]^{\Tr_f\otimes \Tr_g} & R^{2(d+e)}(f\times g)_!
\ZZ/n(d+e) \ar[d]^{\Tr_{f\times g}}\\
\ZZ/n  \ar@{=}[r] & \ZZ/n,
}}\leqno{(2.12.1)}
$$
交换, 其中第一个水平箭头是 \eqhyperref[XVIII.eq2.11.1]{2.11.1} 的扭转形式.\pageoriginale
\end{enonce*}

为验证这一点, 只须利用 (Var 3) 以及在 \hyperrefext[XVII][XVII.thm5.4.3]{5.4.3} 证明中使用的 Künneth 箭头的不对称解释 (\hyperrefext[XVII][XVII.lem5.4.3.5]{5.4.3.5}).

\setcounter{subsection}{12}
\subsection{}\etiquette[2.13]{XVIII.subsec2.13}
设 $f:X\rightarrow S$ 为满足 $(\ast)_{d}$ 的 $S$ 可紧化态射, $\mathcal{A}$ 为 $S$ 上的环层, $n$ 为 $S$ 上可逆且满足 $n\mathcal{A}=0$ 的整数.

对每个 $K\in D(S,\mathcal{A})$, 有 (\hyperrefext[XVII][XVII.thm5.2.6]{5.2.6})
\begin{equation*}
K{\mathop{\otimes}^{\LL}}_{\ZZ/n}Rf_{!}\ZZ/n(d)\xrightarrow{\sim} Rf_{!}(f^{\ast}K(d)){\quoi}.\tag{2.13.1}
\end{equation*}\etiquette[2.13.1]{XVIII.eq2.13.1}

迹态射
$$
\Tr_{f}:R^{2d}f_{!}\ZZ/n(d)\longrightarrow \ZZ/n
$$
也可解释为态射
$$
\Tr_{f}:Rf_{!}\ZZ/n(d)[2d]\longrightarrow \ZZ/n{\quoi}.
$$
将此态射与 $K$ 作张量积, 经由 \eqhyperref[XVIII.eq2.13.1]{2.13.1}, 定义迹态射
\begin{equation*}
\Tr_{f}:Rf_{!}(f^{\ast}K(d)[2d])\longrightarrow K{\quoi}.\tag{2.13.2}
\end{equation*}\etiquette[2.13.2]{XVIII.eq2.13.2}

如将在 \hyperref[XVIII.thm3.2.5]{3.2.5} 中看到的, 下述定理实质上等价于平展上同调中的整体对偶定理.

\begin{enonce*}{定理 2.14}\etiquette[2.14]{XVIII.thm2.14}
设 $f:X\rightarrow S$ 为纯相对维数 $d$ 的光滑可紧化态射, $n$ 为 $S$ 上可逆的整数 $\geq 1$. 对 $X$ 的每个几何点 $x$, 其在 $S$ 中的像为 $s$, 存在 $S$ 中 $s$ 的平展邻域 $V$, 以及 $f^{-1}(V)$ 中 $x$ 的平展邻域 $U$
\begin{equation*}
\vcenter{
\xymatrix@C=1.5cm@R=1.25cm{
U \ar[rd]_{f'_V} \ar[r]^{j} & f^{-1}(V)   \ar[r] \ar[d]^{f_V}& X\ar[d]^{f}\\
                                             &     V           \ar[r]  & S
}}\tag{2.14.1}
\end{equation*}\etiquette[2.14.1]{XVIII.eq2.14.1}%
使得\pageoriginale 态射 \sisi{\emph{(XVII 6.2.7.2)}}{\eqhyperrefext[XVII][XVII.eq6.2.7.4]{6.2.7.4}}
$$
R^{i}f'_{V!}\ZZ/n\longrightarrow R^{i}f_{V!}\ZZ/n
$$
对 $i<2d$ 为零, 且态射
$$
\Tr_{f'_{V}}:R^{2d}f'_{V!}\ZZ/n(d)\longrightarrow \ZZ/n
$$
为同构.
\end{enonce*}

先证明下述引理.

\begin{enonce*}{引理 2.14.2}\etiquette[2.14.2]{XVIII.lem2.14.2}
设 $\mathcal{A}$ 为阿贝尔范畴, $k$ 为整数 $\geq 0$, $(K_{i})_{0\leq i\leq 2k}$ 为 $D^{b}(\mathcal{A})$ 中的对象, 满足对 $p\not\in [0,k]$ 有 $H^{p}(K_{i})=0$, $f_{i}:K_{i}\rightarrow
K_{i+1}$ $(0\leq i\leq 2k-1)$ 为态射, $f$ 为其复合. 若对 $p<k$ 有 $H^{p}(f_{i})=0$, 则在 $D^{b}(\mathcal{A})$ 中存在从复形 $H^{k}(K_{0})[-k]$ (由置于次数 $k$ 的 $H^{k}(K_{0})$ 构成) 到 $K_{2k}$ 的态射 $\varphi$, 使图表
$$
\xymatrix{
K_{0}  \ar[d]\ar[r]^{f} & K_{2k}\\
\sigma_{\geq k}(K_{0})  \ar@{=}[r] & H^{k}(K_{0})[-k] \ar[u]_{\varphi}.
}
$$
交换.
\end{enonce*}

当 $k=0$ 时引理平凡. 对 $k$ 作归纳证明. 对复形 $\sigma_{\geq
1}(K_{i})[1]$ 应用归纳假设, 得到态射 $\varphi':H^{k}(K_{0})[-k]\rightarrow \sigma_{\geq 1}K_{2k-2}$ 使下图交换, 其中 $f'=f_{2k-3}\ldots f_{0}$:
$$
\xymatrix{
K_{0}  \ar[d]\ar[r]^{f'} & \sigma_{\geq 1}K_{2k-2}            \\
\sigma_{\geq k}(K_{0})  \ar@{=}[r] & H^{k}(K_{0})[-k]. \ar[u]_{\varphi'}
}
$$

对每个复形 $L$, 特殊三角
$$
\xymatrix{
                            &\sigma_{\geq 1}(L)\ar[ld]_{+1}  &       \\
  \sigma_{\leq 0}(L) \ar[rr] &                    & L\ar[lu]
}
$$
对每个复形 $M$ 给出\pageoriginale 正合列
$$
\ldots\rightarrow \Hom(M,\sigma_{\leq 0}(L))\rightarrow \Hom
(M,L)\rightarrow \Hom (M,\sigma_{\geq
1}L)\xrightarrow{j}\Ext^{1}(M,\sigma_{\leq 0}L){\quoi}.
$$

特别地, 有图表
$${\def\objectstyle{\scriptscriptstyle}
\xymatrix@C=0.4cm{
\Hom(\sigma_{\geq k}K_{0}, \sigma_{\leq 0} K_{2k-2})\ar[r] \ar[d]&
\overset{\text{\ding{176}}}{\underset{\phantom{\text{\ding{176}}}}
{\Hom}}(\sigma_{\geq k}K_{0},
    K_{2k-2}) \ar[r]\ar[d] &
\overset{\text{\ding{174}}}{\underset{\phantom{\text{\ding{174}}}}
{\Hom}}(\sigma_{\geq
      k}K_{0}, \sigma_{\geq 1} K_{2k-2})\ar[r] \ar[d]&
\overset{\text{\ding{175}}}{\underset{\phantom{\text{\ding{175}}}}
{\Ext}^1}(\sigma_{\geq
      k}K_{0}, \sigma_{\leq 0} K_{2k-2})\ar[d]\\
\overset{\text{\ding{178}}}{\underset{\phantom{\text{\ding{178}}}}
{\Hom}}(K_0,\sigma_{\leq 0} K_{2k-2})\ar[r] &
\overset{\text{\ding{172}}}{\underset{\phantom{\text{\ding{172}}}}
{\Hom}}(K_0,K_{2k-2}) \ar[r] &
\overset{\text{\ding{173}\ding{177}}}{\underset{\phantom{\text{\ding{173}\ding{177}}}}
{\Hom}}(K_0, \sigma_{\geq
  1}K_{2k-2})\ar[r]& \Ext^{1}(K_0, \sigma_{\leq 0} K_{2k-2})
}}\leqno{(2.14.3)}\etiquette[2.14.3]{XVIII.eq2.14.3}
$$
它映到 $K_{2k-1}$ 与 $K_{2k}$ 的类似图表. 根据假设, 位于 \ding{172} 的态射 $f'$ 在 \ding{173} 处的像可提升到 \ding{174}. 将元素 \ding{174} 提升到 \ding{176} 的障碍 \ding{175}, 在 $K_{2k-1}$ 的图表 \eqhyperref[XVIII.eq2.14.3]{2.14.3} 中的像为零. 因而在此新图表中, \ding{174} 可提升到 \ding{176}, 而 \ding{172} 与 \ding{176} 的像之差在 \ding{177} 处的像为零, 因而是某个元素 \ding{178} 的像. 此元素 \ding{178} 在 $K_{2k}$ 的图表 \hyperref[XVIII.eq2.14.3]{2.14.3} 中的像为零; 因而在此新图表中, $f$ 为 \ding{176} 的像, 问题得解.

证明对于给定的 $d$ 值, 断言 \hyperref[XVIII.thm2.14]{2.14} 蕴含下述结果.

\begin{enonce*}{推论 2.14.4}\etiquette[2.14.4]{XVIII.cor2.14.4}
在 \hyperref[XVIII.thm2.14]{2.14} 的假设下, 存在图表 \eqhyperref[XVIII.eq2.14.1]{2.14.1}, 使 \emph{\hyperrefext[XVII][XVII.subsubsec6.2.4]{6.2.4}} 中的态射在 $D^{b}(V,\sheaf{\ZZ}/n)$ 中有分解
$$
\xymatrix@C=.75cm@R=1cm{
Rf'_{V!}\sheaf{\ZZ}/n(d)  \ar[rr] \ar[rd]_{t}&  & Rf_{V!}\sheaf{\ZZ}/n(d)\\
                          &\sheaf{\ZZ}/n[-2d], \ar[ru]_{\varphi}
}
$$
其中 $t$ 为箭头
% SGA-SRC-001151
$$
t:Rf'_{V!}\ZZ/n(d)\rightarrow \sigma_{\geq
2d}Rf'_{V!}\ZZ/n(d)=H^{2d}(Rf'_{V!}\ZZ/n(d))[-2d]\xrightarrow{\Tr}\sheaf{\ZZ}/n[-2d]{\quoi}.
$$
\end{enonce*}

置\pageoriginale $f'_{0}=f$, 归纳构造 $4d$ 个满足 \hyperref[XVIII.thm2.14]{2.14} 的图表 \eqhyperref[XVIII.eq2.14.1]{2.14.1}
$$
\xymatrix@C=2cm@R=1.5cm{
U^{'}_{i+1} \ar[r]\ar[rd]_{f'_{i+1}}& f'^{-1}_{i}(V'_{i+1}) \ar[d]\ar[r]& U'_i \ar[d]^{f'_{i}}\\
            & V'_{i+1}  \ar[r]            & V'_i
}
$$
置 $V=V'_{4d}$ 及 $U_{i}=\fprod{U'_{i}}{V}{V'_{i}}$. 得到图表
$$
\xymatrix@R=1.5cm@C=1cm{
U_{4d}\ar[r] \ar[rrrd]^{f_{4d}}& \cdots \ar[r]& U_1 \ar[r]\ar[rd]^{f_1}& U_0\ar[r]\ar[d]^{f^0} & X \ar[d]^{f}\\
            &              &           & V \ar@{-}[r]  & S ,
}
$$
其中 $U_{0}=f^{-1}(V)$, 使箭头
$$
R^{i}f_{k+1!}\ZZ/n(d)\longrightarrow R^{i}f_{k!}\ZZ/n(d)
$$
为零 (对 $i<2d$), 且箭头
$$
\Tr:R^{2d}f_{4d!}\sheaf{\ZZ}/n(d)\longrightarrow \sheaf{\ZZ}/n
$$
为同构. 由 \hyperref[XVIII.lem2.14.2]{2.14.2}, 推论对 $U=U_{4d}$ 成立.

现在已\sisi{准备就绪}{准备就绪}, 对 $d$ 作归纳证明 \hyperref[XVIII.thm2.14]{2.14}. 情形 $d=0$ 平凡, 情形 $d=1$ 恰为 \hyperref[XVIII.lemfond1.6.9]{1.6.9}. 因而可以假设 $d\geq 2$.

由于问题在 $X$ 与 $S$ 上是局部的, 可以假设 $f$ 有分解 $f=gh$, 其中 $g$ 与 $h$ 光滑且可紧化, 纯相对维数为 $d'$ 及 $d''$, 并且 $d'$、$d''<d=d'+d''$. 对 $g$ 以及由 $h$ 经基变换所得的态射应用 \hyperref[XVIII.cor2.14.4]{2.14.4}, 得到图表
$$
\xymatrix@C=2.2cm@R=1.5cm{
  U  \ar[r]^{u} \ar[rd]^{h'_0} & f^{-1}W \ar[d]^{h_0}       \ar[r] & X \ar[d]^{h} \\
V  \ar[r]^{v} \ar[rd]^{g'_0} & g^{-1}W \ar[r] \ar[d]^{g_0}       & Y \ar[d]^{g}\\
               &    W   \ar[r]               & S,
}
$$
使得\pageoriginale 态射 (XVII \hyperref[XVII.eq6.2.7.3]{6.2.7.3}) 有分解
\begin{align*}
 Rh'_{0!}\ZZ/n(d'') & \xrightarrow{\Tr}
\ZZ/n[-2d'']\xrightarrow{\varphi} Rh_{0!}\ZZ/n(d'')\\
Rg'_{0!}\ZZ/n(d') & \xrightarrow{\Tr}\ZZ/n[-2d']\xrightarrow{\Psi} Rg_{0!}\ZZ/n(d')
\end{align*}
且 $\Tr_{h'_0}$ 与 $\Tr_{g'_0}$ 为同构.

设 $h_{1}$ 与 $h'_{1}$ 为由 $h_{0}$ 与 $h'_0$ 经基变换到 $V$ 所得的态射, 并设 $f'=g'_{0}h'_{1}$、$f_{W}=g_{0}h_{0}$.

下图交换:
$$
\xymatrix{
  Rf'_{!}\ZZ/n(d) \ar[rr]\ar@{=}[d] & & Rf_{W!}\ZZ/n(d) \ar@{=}[d]\\
Rg'_{0!} Rh'_1\ZZ/n(d) \ar[r]& Rg_{0!} Rh'_{0!}\ZZ/n(d) \ar[r]& Rg_{0!}
Rh_{0!}\ZZ/n(d),
}
$$
图表
$$
\xymatrix@C=2.1cm@R=1cm{
  Rg'_{0!} Rh'_{1!} \ZZ/n(d)   \ar@{=}[d] \ar@<-5ex>@/_2pc/[ddd]& &\\
Rg'_{0!} v^{\ast} Rh'_{0!} \ZZ/n(d)
 \ar[dd] \ar[r]_{Rg'_{0!}v^{\ast}(t)(d')}&
Rg'_{0!}\ZZ/n(d')[-2d''] \ar[d]^{t[-2d'']}& \\
   & \ZZ/n[-2d] \ar[d]^{\Psi[-2d'']}\\
Rg_{0!} Rh'_{0!} \ZZ/n(d) \ar[r]_{Rg_{0!}(t)(d')} \ar@/_2.2pc/[rr]& Rg_{0!}\ZZ/n(d')[-2d''] \ar[r]_{Rg_{0!}(\varphi)(d')}& Rg_{0!}Rh_{0!}\ZZ/n(d)
}
$$
也交换. 因而对于\pageoriginale 图表
$$
\xymatrix@C=1.5cm@R=1.25cm{
U_1 \ar[r]\ar[rd]^{f'} & f^{-1}W \ar[r]\ar[d]^{f_W}& X\ar[d]^{f}\\
                       &  W \ar@{-}[r]     & S,
}
$$
从 $Rf'_{!}\ZZ/n(d)$ 到 $Rf_{W!}\ZZ/n(d)$ 的典范箭头经由 $\ZZ/n[-2d]$ 分解; 因而从 $R^{i}f'_{!}\ZZ/n(d)$ 到 $R^{i}f_{W!}\ZZ/n(d)$ 的箭头对 $i<2d$ 为零. 最后, $\Tr_{f'}$ 为同构, 因为它是同构 $\Tr_{g'_{0}}$ 与 $R^{2d'}g'_{0!}\Tr_{h'_{1}}$ 的复合.

至此完成了 \hyperref[XVIII.thm2.14]{2.14} 的证明.



% END native input: chapters/exp-18-part4-zh.tex
% BEGIN native input: chapters/exp-18-part5-zh.tex
\section{整体对偶定理}\etiquette[3]{XVIII.sec3}\pageoriginale

\subsection{函子 $Rf^{!}$}\etiquette[3.1]{XVIII.sec3.1}

我们需要典范松软消解的一个变体.

\begin{enonce*}{引理 3.1.1}\etiquette[3.1.1]{XVIII.lem3.1.1}
设 $X$ 为凝聚概形 (\ie 拟紧且拟分离). 函子 $\varinjlim$ 在可构造集合层 (\resp 群层, \resp 阿贝尔群层) 范畴的 Ind-对象范畴与集合层 (\resp Ind-有限群层, \resp 挠阿贝尔群层) 范畴之间诱导一个等价.
\end{enonce*}

这来自 \hyperrefext[IX][IX.cor2.7.2]{2.7.2} 及 \hyperrefext[IX][IX.cor2.7.3]{2.7.3}.

\setcounter{subsubsection}{1}
\subsubsection{}\etiquette[3.1.2]{XVIII.subsubsec3.1.2}
设 $\Ff$ 为概形 $X$ 上的挠阿贝尔层, 它是可构造层 $(\Ff_{i})_{i\in I}$ 的滤过归纳极限, 并设 $X_{0}$ 为 $X$ 的一个保守点集. $\mathcal{F}$ 的\emph{修正典范松软消解}（modified canonical flasque resolution）由公式
$$
C^{\ast}_{\ell}(\Ff)=\varinjlim_{i}C^{\ast}(\Ff_{i}){\quoi},
$$ 定义, 其中 $C^{\ast}(\Ff_{i})$ 为 \hyperrefext[XVII][XVII.def4.2.2]{4.2.2} 的典范松软消解. 由 \hyperref[XVIII.lem3.1.1]{3.1.1}, 当 $X$ 凝聚时, 这一定义与 $\Ff_{i}$ 的选取无关. 它还与局部化可交换; 因而当 $\Ff$ 为任意概形 $X$ (未必凝聚) 上的挠层时, 可通过整体化定义 $C^{\ast}_{\ell}(\Ff)$.

\emph{复形 $C^{\ast}_{\ell}(\Ff)$ 是 $\Ff$ 的一个消解, 关于 $\Ff$ 具有函子性, 其构造与局部化及滤过归纳极限可交换; 函子 $C^{n}_{\ell}$ 正合.}

由此, 对 $X$ 上的环层 $\mathcal{A}$, 函子 $C^{n}_{\ell}$ 将 $\mathcal{A}$-模映为 $\mathcal{A}$-模; 它们也将挠层映为挠层. 众所周知,

\begin{enonce*}{引理 3.1.3}\etiquette[3.1.3]{XVIII.lem3.1.3}
设\pageoriginale $\mathcal{A}$ 为阿贝尔 $\sheaf{\mathcal{U}}$-范畴, 具有一个小生成集, 并且由小滤过有序集索引的归纳极限在 $\mathcal{A}$ 中可表且正合. 从 $\mathcal{A}^{\circ}$ 到 $(\Ens)$ 的函子 $F$ 可表, 当且仅当它将 (任意) 小归纳极限映为投射极限.
\end{enonce*}

\Cf 例如 \SGA 4 I \hyperref[I.prop8.12.5]{8.12.5}, 其中既不假设 $\mathcal{A}$ 阿贝尔, 也不假设滤过的 $\varinjlim$ 正合.

我们将此引理应用于赋环景 $(X,\mathcal{A})$ 上的模层范畴 $\Mod(X,\mathcal{A})$. 在这一特例中, 若函子 $F$ 由层 $ℱ$ 表示, 则
\begin{equation*}
ℱ(U)=F(\mathcal{A}_{U}){\quoi}.\tag{3.1.3.1}
\end{equation*}\etiquette[3.1.3.1]{XVIII.eq3.1.3.1}

对 $ℱ$ 的这一描述直接验证了当前情形中的 \hyperref[XVIII.lem3.1.3]{3.1.3}.

回顾, 若 $f:X\rightarrow S$ 为可紧化态射 (\hyperrefext[XVII][XVII.subsec3.2]{3.2}.), 则 $S$ 拟紧且拟分离 (原文如此), 并且 $f$ 的纤维维数有界.

\begin{enonce*}{定理 3.1.4}\etiquette[3.1.4]{XVIII.thm3.1.4}
设 $f:X\rightarrow S$ 为可紧化态射, $\mathcal{A}$ 为 $S$ 上的挠环层. 则函子
\begin{equation*}
Rf_{!}:D(X,f^{\ast}\mathcal{A})\longrightarrow D(S,\mathcal{A})\tag{3.1.4.1}
\end{equation*}\etiquette[3.1.4.1]{XVIII.eq3.1.4.1}%
容许部分右伴随
\begin{equation*}
Rf^{!}:D^{+}(S,\mathcal{A})\longrightarrow D^{+}(X,f^{\ast}\mathcal{A}):\tag{3.1.4.2}
\end{equation*}\etiquette[3.1.4.2]{XVIII.eq3.1.4.2}%
对 $K\in \Ob D(X,f^{\ast}\mathcal{A})$ 及 $L\in \Ob
D^{+}(S,\mathcal{A})$, 有函子性同构
\begin{equation*}
\Hom(Rf_{!}K,L)\simeq \Hom(K,Rf^{!}L){\quoi}.\tag{3.1.4.3}
\end{equation*}\etiquette[3.1.4.3]{XVIII.eq3.1.4.3}%
赋予使图
\begin{equation*}
\vcenter{
\xymatrix{
  \Hom(Rf_!(K[1]),L)
  \ar@{<->}[d]_{\text{\rotatebox{90}{$\backsim$}}}\ar@{<->}[r]^-{\sim}&
  \Hom(K[1],Rf^{!}L)\ar@{-}[r]^-\sim&
\Hom(K,(Rf^{!}L)[-1])\ar@{<->}[d]^{\text{\rotatebox{90}{$\backsim$}}}\\
\Hom((Rf_!K)[1],L)\ar@{-}[r]^-\sim & \Hom(Rf_!K,L[-1])
\ar@{<->}[r]^-{\sim}_-{\txt{\rm 伴随}} & \Hom(K,Rf^{!}(L[-1])),
}}\tag{3.1.4.4}
\end{equation*}\etiquette[3.1.4.4]{XVIII.eq3.1.4.4}%
交换的平移态射后,\pageoriginale 这一函子 $Rf^{!}$ 为三角函子.
\end{enonce*}

注意. 记号 $Rf^{!}$ 有滥用之处, 因为 $Rf^{!}$ 一般不是某个函子 $f^{!}$ 的导出函子.

\noindent
\emph{证明}. 取整数 $d$, 使 $f$ 的每个纤维维数均为 $<d$, 并选择紧化
$$
\xymatrix@R=1.25cm@C=1.5cm{
X \ar@{^{(}->}[r]^{j} \ar[d]_{f}& \bar{X}\ar[ld]^{\bar{f}}\\
S.
}
$$

若 $K$ 为 $X$ 上的 $f^{\ast}\mathcal{A}$-模复形, 则
$$
Rf_{!}K\simeq R\bar{f}_{\ast}(j_{!}K){\quoi}.
$$

对 $X$ 上的每个层 $F$, 复形 $\tau_{\leq 2d}C^\ast_{\ell}j_{!}F$ 的各分量 $F^{i}$ 满足
\begin{equation*}
R^{k}\bar{f}_{\ast}F^{i}=0\quad \text{\rm 当}\quad k>0{\quoi}.\tag{3.1.4.5}
\end{equation*}\etiquette[3.1.4.5]{XVIII.eq3.1.4.5}%
事实上, 当 $i\neq 2d$ 时, $F^{i}$ 为松软层的归纳极限, 而当 $i=2d$ 且 $k>0$ 时, 有 (\hyperrefext[XVII][XVII.cor5.2.8.1]{5.2.8.1})
$$
R^{k}\bar{f}_{\ast}F^{2d}=R^{k+2d}f_{!}F=0{\quoi}.
$$

$K$ 的截断修正典范松软消解双复形所伴随的单复形, 是 $K$ 的一个消解. 因而由 \eqhyperref[XVIII.eq3.1.4.5]{3.1.4.5}, 有 (记号见 \hyperrefext[XVII][XVII.subsubsec1.1.15]{1.1.15})
\begin{equation*}
R\bar{f}_{\ast}(j_{!}K)\xleftarrow{\sim} \bar{f}_{\ast}\tau''_{\leq
2d} C^{\ast}_{\ell}j_{!}K{\quoi}.\tag{3.1.4.6}
\end{equation*}\etiquette[3.1.4.6]{XVIII.eq3.1.4.6}

以 $f^{\bigdot}_{!}$ 表示由下式定义的从 $X$ 上的模层到 $S$ 上的层复形的\pageoriginale 函子:
\begin{equation*}
f^{\bigdot}_{!}(F)=\bar{f}_{\ast}\tau_{\leq
2d}C^{\ast}_{\ell}j_{!}F{\quoi}.\tag{3.1.4.7}
\end{equation*}\etiquette[3.1.4.7]{XVIII.eq3.1.4.7}

\begin{enonce*}{引理 3.1.4.8}\etiquette[3.1.4.8]{XVIII.lem3.1.4.8}
\begin{enumerate}
\renewcommand{\theenumi}{\roman{enumi}}
\renewcommand{\labelenumi}{(\theenumi)}
\item 函子 $f^{i}_{!}$ 正合且与滤过归纳极限可交换; 当 $i\not\in[0,2d]$ 时有 $f^{i}_{!}=0$.

\item 由于函子 $f^{\bigdot}_{!}$ 有界且各分量正合, 其导出函子\emph{(\hyperrefext[XVII][XVII.prop1.2.10]{1.2.10})} 平凡地得到, 记为 $Rf^{\bigdot}_{!}$, 或简记为 $f^{\bigdot}_{!}:D(X,f^{\ast}\mathcal{A})\rightarrow
D(S,\mathcal{A})$. 箭头 \eqhyperref[XVIII.eq3.1.4.6]{3.1.4.6} 为函子的同构
$$
Rf^{\bigdot}_{!}\xrightarrow{\sim} Rf_{!}{\quoi}.
$$
\end{enumerate}
\end{enonce*}

函子 $C^{i}_{\ell}$ 正合且与滤过归纳极限可交换 (\hyperref[XVIII.subsubsec3.1.2]{3.1.2}). 它们将每个层映为无上同调层. 函子 $\bar{f}_{\ast}$ 与 $j_{!}$ 同滤过归纳极限可交换. 从而
\begin{enumerate}
\renewcommand{\theenumi}{\alph{enumi}}
\renewcommand{\labelenumi}{\theenumi)}
\item 函子 $\bar{f}_{\ast}C^{i}_{\ell}j_{!}$, 从而函子 $f^{i}_{!}$, 都与滤过归纳极限可交换;

\item 当 $i<2d$ 时, 函子 $f^{i}_{!}=\bar{f}_{\ast}C^{i}_{\ell}j_{!}$ 正合. 当 $i=2d$ 时, 由 \eqhyperref[XVIII.eq3.1.4.5]{3.1.4.5} 可验证 $f^{i}_{!}$ 正合.
\end{enumerate}

\hyperref[XVIII.lem3.1.4.8]{3.1.4.8} 的其余断言是平凡的.

\hypertarget{sga4.exp.XVIII.remark.3.1.4.9}{}\begin{enonce*}[remark]{注 3.1.4.9}\etiquette[3.1.4.9]{XVIII.rem3.1.4.9}
$f^{\bigdot}_{!}$ 的定义与环层 $\mathcal{A}$ 无关, 并且对任意阿贝尔层仍有意义. 断言 \hyperref[XVIII.lem3.1.4.8]{3.1.4.8} (i) 在 $X$ 与 $S$ 上的挠阿贝尔层范畴中仍然成立.
\end{enonce*}

由 \hyperref[XVIII.lem3.1.4.8]{3.1.4.8} 及 \hyperref[XVIII.lem3.1.3]{3.1.3}, 函子
$$
f^{i}_{!}:\Mod(X,f^{\ast}\mathcal{A})\longrightarrow
\Mod(S,\mathcal{A})
$$
有右伴随 $f^{!}_{i}$; 由于 $f^{i}_{!}$ 正合, $f^{!}_{i}$ 将内射对象映为内射对象.

采用 (\hyperrefext[XVII][XVII.subsubsec1.1.12]{1.1.12}) 的符号约定, 函子 $f^{!}_{i}$ 构成左正合函子的一个复形, 在上同调次数 $n\not\in [-2d,0]$ 处为零:
\begin{equation*}
f^{!}_{\bigdot}:\Mod(S,\mathcal{A})\longrightarrow
C^{b}(X,f^{\ast}\mathcal{A}){\quoi}.\tag{3.1.4.10}
\end{equation*}\etiquette[3.1.4.10]{XVIII.eq3.1.4.10}

函子\pageoriginale $f^{!}_{\bigdot}$ 导出 (\hyperrefext[XVII][XVII.prop1.2.10]{1.2.10}) 为一个三角函子
\begin{equation*}
Rf^{!}:D^{+}(S,\mathcal{A})\longrightarrow
D^{+}(X,f^{\ast}\mathcal{A}){\quoi},\tag{3.1.4.11}
\end{equation*}\etiquette[3.1.4.11]{XVIII.eq3.1.4.11}%
它右伴随于函子 $Rf_{!}$ $(=Rf^{\bigdot}_{!}=Lf^{\bigdot}_{!})$ (“平凡对偶公式”). 更准确地, 若 $I\in \Ob C^{+}(S,\mathcal{A})$ 为内射 $\mathcal{A}$-模层的复形, 则 $\mathcal{A}$-模 $f^{!}_{i} I^{n}$ 内射, 并有三重复形的同构 (\hyperrefext[XVII][XVII.subsubsec1.1.12]{1.1.12})
\begin{equation*}
\Hom^{\bigdot}(f^{\bigdot}_{!}K, I)\xrightarrow{\sim}
\Hom^{\bigdot}(K,f^{!}_{\bigdot}I){\quoi}.\tag{3.1.4.12}
\end{equation*}\etiquette[3.1.4.12]{XVIII.eq3.1.4.12}%
过渡到伴随的单复形 (用乘积计算) 的 $H^{0}$, 即得伴随性陈述. 相容性 \eqhyperref[XVIII.eq3.1.4.4]{3.1.4.4} 的验证留给读者.

\hypertarget{sga4.exp.XVIII.remark.3.1.5}{}\begin{enonce*}[remark]{注 3.1.5}\etiquette[3.1.5]{XVIII.rem3.1.5}
假设函子 $Rf^{!}$ 有有限上同调振幅. 由 \hyperrefext[XVII][XVII.prop1.2.10]{1.2.10}, 函子 $f^{!}_{\bigdot}$ 于是容许导出函子
\begin{equation*}
Rf^{!}:D(S,\mathcal{A})\longrightarrow D(X,f^{\ast}\mathcal{A}){\quoi}.\tag{3.1.5.1}
\end{equation*}\etiquette[3.1.5.1]{XVIII.eq3.1.5.1}%
这一导出函子仍右伴随于函子 $Rf_{!}$ \eqhyperref[XVIII.eq3.1.4.1]{3.1.4.1}. 事实上, 有
\begin{align*}
& \Hom_{D(S)}(Rf_{!}K,L)\simeq
\varinjlim_{L\xrightarrow{\sim}L'}\Hom_{K(S)}(f^{\bigdot}_{!}K,L')\simeq
\varinjlim_{\substack{K'\xrightarrow{\sim}K\\ L\xrightarrow{\sim}
L'}}\Hom_{K(S)} (f^{\bigdot}_{!}K', L')\\
& \Hom_{D(X)}(K,Rf^{!}L)\simeq
\varinjlim_{K'\xrightarrow{\sim}K}\Hom_{K(X)}(K',Rf^{!}L)\simeq
\varinjlim_{\substack{K'\xrightarrow{\sim}K\\
L\xrightarrow{\sim}L'}}\Hom_{K(X)}(K', f^{!}_{\bigdot}L')
\end{align*}
并且由伴随性,
% SGA-SRC-001156
$$
\begin{gathered}
\Hom_{K(S)}(f^{\bigdot}_{!}K',L')=H^{0}\Hom(f^{\bigdot}_{!}K',L')\\
\simeq H^{0}\Hom(K',f^{!}_{\bigdot}L')=\Hom_{K(X)}(K',f^{!}_{\bigdot}L'){\quoi}.
\end{gathered}
$$
\end{enonce*}

\hypertarget{sga4.exp.XVIII.definition.3.1.6}{}\begin{enonce*}{定义 3.1.6}\etiquette[3.1.6]{XVIII.def3.1.6}
右伴随于函子 $Rf_{!}$ \eqhyperref[XVIII.eq3.1.4.1]{3.1.4.1} 的三角函子
$$
Rf^{!}:D^{+}(S,\mathcal{A})\longrightarrow
D^{+}(X,f^{\ast}\mathcal{A})
$$
(或,\pageoriginale 在 \hyperref[XVIII.rem3.1.5]{3.1.5} 的假设下,
$$
Rf^{!}:D(S,\mathcal{A})\longrightarrow
D(X,f^{\ast}\mathcal{A})){\quoi},
$$
称为\emph{例外逆像函子}（exceptional inverse-image functor）.
\end{enonce*}

\hypertarget{sga4.exp.XVIII.proposition.3.1.7}{}\begin{enonce*}{命题 3.1.7}\etiquette[3.1.7]{XVIII.prop3.1.7}
假设可紧化态射 $f:X\rightarrow S$ 的相对维数为 $\leq d$.
\begin{enumerate}
\renewcommand{\theenumi}{\roman{enumi}}
\renewcommand{\labelenumi}{(\theenumi)}
\item 若 $L\in \Ob D^{+}(S,\mathcal{A})$ 在 $i\leq k$ 时满足 $H^{i}(L)=0$, 则在 $i\leq k-2d$ 时有 $H^{i}Rf^{!}L=0$.

\item 对 $L\in \Ob D^{b}(S,\mathcal{A})$, 有
$$
\dim \inj (Rf^{!}L)\leq \dim \inj (L){\quoi}.
$$
\end{enumerate}
\end{enonce*}

函子 $Rf^{!}$ 是 $f^{!}_{\bigdot}:\Mod(S,\mathcal{A})\rightarrow
C^{b}\Mod(X,f^{\ast}\mathcal{A})$ 的右导出函子. 断言 (i) 来自当 $i<-2d$ 时 $(f^{!}_{\bigdot})^{i}=0$ 这一事实; 断言 (ii) 则来自当 $i>0$ 时 $(f^{!}_{\bigdot})^{i}=0$, 以及函子 $f^{!}_{i}$ 将内射对象映为内射对象.

\hypertarget{sga4.exp.XVIII.proposition.3.1.8}{}\begin{enonce*}{命题 3.1.8}\etiquette[3.1.8]{XVIII.prop3.1.8}
设 $f:X\rightarrow S$ 为拟有限可紧化态射.
\begin{enumerate}
\renewcommand{\theenumi}{\roman{enumi}}
\renewcommand{\labelenumi}{(\theenumi)}
\item 函子 $f_{!}:\Mod(X,f^{\ast}\mathcal{A})\rightarrow
\Mod(S,\mathcal{A})$ 容许右伴随 $f^{!}:\Mod(S,\mathcal{A})\rightarrow \Mod(X,f^{\ast}\mathcal{A})$, 且 $Rf^{!}$ 为 $f^{!}$ 的导出函子.

\item 若 $f$ 为闭浸入, 则 $f^{!}$ 是“支集在 $X$ 中的截面”函子.

\item 若 $f$ 平展, 则 $f^{!}$ 等同于函子 $f^{\ast}$, 并以
$$
\Tr_{f}:f_{!}f^{\ast}F\longrightarrow F
$$
为伴随态射.
\end{enumerate}
\end{enonce*}

断言 (i) 直接来自 \hyperref[XVIII.thm3.1.4]{3.1.4} 证明中 $d=0$ 的特例: 有 $f^{\bigdot}_{!}=f_{!}$ 及 $f^{!}_{\bigdot}=f^{!}$. 等价地, 伴随函子 $f^{!}$ 的存在性来自 \hyperref[XVIII.lem3.1.3]{3.1.3}, 以及 $f_{!}$ 正合\pageoriginale 且与滤过归纳极限可交换这一事实; 伴随性 $(f_{!},Rf^{!})$ 是一个“平凡对偶公式”.

断言 (ii) 显然, (iii) 即为 \hyperrefext[XVII][XVII.prop6.2.11]{6.2.11}.

\setcounter{subsubsection}{8}
\subsubsection{}\etiquette[3.1.9]{XVIII.subsubsec3.1.9}
设 $f:X\rightarrow S$ 为可紧化态射, $\mathcal{A}$ 为 $S$ 上的挠环层, 并以 $\mathcal{A}$ 的逆像环层给 $X$ 赋环. 设 $K\in \Ob
D^{-}(X,f^{\ast}\mathcal{A})$, $L\in \Ob D^{+}(X,f^{\ast}\mathcal{A})$, 并以一个下有界内射层复形表示 $L$, 从而
$$
R\Hhom(K,L)\simeq \Hhom(K,L){\quoi}.
$$
复形 $\Hhom(K,L)$ 的各分量松软, 因而
\begin{equation*}
Rf_{\ast}R\Hhom(K,L)\simeq f_{\ast}\Hhom(K,L){\quoi}.\tag{3.1.9.1}
\end{equation*}\etiquette[3.1.9.1]{XVIII.eq3.1.9.1}%
如 \hyperref[XVIII.thm3.1.4]{3.1.4} 的证明中那样, 选取 $f$ 的紧化及整数 $d$, 得到函子 $f^{\bigdot}_{!}$ \eqhyperref[XVIII.eq3.1.4.7]{3.1.4.7}.
\begin{equation*}
\vcenter{
\xymatrix@C=1.5cm@R=.5cm{
X_V \ar[rr]^{k_X} \ar@{^{(}->}[rd]_{j_V} \ar[ddd]^{f_V}&      & X\ar[ddd]^{f} \ar@{^{(}->}[rd]^{j}         \\
    & \bar{X}_V \ar[ldd]^{\bar{f}_V} \ar[rr]_(.3){k_{\bar{X}}}   &   & \bar{X} \ar[ldd]_{\bar{f}}\\
    &         &        &\\
V  \ar[rr]^{k} &      & S.
}}\tag{3.1.9.2}
\end{equation*}\etiquette[3.1.9.2]{XVIII.eq3.1.9.2}

对每个 $V\in \Ob S_{\text{\rm et}}$, 都有同构\nde{这一态射由 \hyperrefext[XVII][XVII.rem4.2.6]{4.2.6} 的构造给出, 但除非先在所考察的拓扑斯上相容地选取点, 否则没有理由认为它是同构. 为使这一条件成立, \hyperrefext[XVII][XVII.rem4.2.6]{4.2.6} 中出现的方块须在某种意义下为笛卡儿方块. 更准确地, 此处应固定拓扑斯 $S_{\text{\rm et}}$ 的一族点 $(p)_{p\in P}$, 并对每个 $V\in \Ob S_{\text{\rm et}}$, 采用由 \hyperrefext[IV][IV.subsec6.7]{6.7} 的对应与满足 $p\in P$ 且 $\xi\in V_p$ 的对 $(p,\xi)$ 相联系的点族, 构造修正典范松软消解.}
\begin{equation*}
k^{\ast}f^{\bigdot}_{!}\xrightarrow{\sim}
f^{\bigdot}_{V!}k^{\ast}_{X}{\quoi},\tag{3.1.9.3}
\end{equation*}\etiquette[3.1.9.3]{XVIII.eq3.1.9.3}%
从而
\begin{equation*}
\Hhom^{\bigdot}(f^{\bigdot}_{!}K,f^{\bigdot}_{!}L)(V)\simeq
\Hhom^{\bigdot}(f^{\bigdot}_{V!}(k^{\ast}_{X}K),f^{\bigdot}_{V!}(k^{\ast}_{X}L)){\quoi}.
\end{equation*}%
由 $f^{\bigdot}_{V!}$ 的函子性, 得到
$$
\Hhom^{\bigdot}(K,L)(X_{V})=\Hom^{\bigdot}(k^{\ast}_{X}K,k^{\ast}_{X}L)\longrightarrow
\Hom^{\bigdot}(f^{\bigdot}_{V!}(k^{\ast}_{X}K),f^{\bigdot}_{V!}(k^{\ast}_{X}L)){\quoi},
$$
从而在复形层面得到箭头
\begin{equation*}
f_{\ast}\Hhom^{\bigdot}(K,L)\longrightarrow
\Hhom^{\bigdot}(f^{\bigdot}_{!}K,f^{\bigdot}_{!}L){\quoi}.\tag{3.1.9.4}
\end{equation*}\etiquette[3.1.9.4]{XVIII.eq3.1.9.4}

由\pageoriginale \eqhyperref[XVIII.eq3.1.9.1]{3.1.9.1}, 这一箭头导出为
\begin{equation*}
Rf_{\ast}R\Hhom(K,L)\longrightarrow
R\Hhom(Rf_{!}K,Rf_{!}L){\quoi}.\tag{3.1.9.5}
\end{equation*}\etiquette[3.1.9.5]{XVIII.eq3.1.9.5}

现设 $K\in \Ob D^{-}(X,\mathcal{A})$ 且 $L\in \Ob
D^{+}(S,\mathcal{A})$. 从 $Rf_{!}Rf^{!}L$ 到 $L$ 的伴随箭头, 连同 \eqhyperref[XVIII.eq3.1.9.4]{3.1.9.4}, 给出复合箭头
$$
Rf_{\ast}R\Hhom(K,Rf^{!}L)\longrightarrow R\Hhom(Rf_{!}K,
Rf_{!}Rf^{!}L)\longrightarrow R\Hhom(Rf_{!}K,L){\quoi},
$$
它是伴随同构在 $S$ 上的局部形式:
\begin{equation*}
Rf_{\ast}R\Hhom(K,Rf^{!}L)\longrightarrow
R\Hhom(Rf_{!}K,L){\quoi}.\tag{3.1.9.6}
\end{equation*}\etiquette[3.1.9.6]{XVIII.eq3.1.9.6}

\hypertarget{sga4.exp.XVIII.proposition.3.1.10}{}\begin{enonce*}{命题 3.1.10}\etiquette[3.1.10]{XVIII.prop3.1.10}
箭头 \hyperref[XVIII.eq3.1.9.6]{3.1.9.6} 为同构.
\end{enonce*}

设 $V\in \Ob S_{\text{\rm et}}$, 并沿用 \eqhyperref[XVIII.eq3.1.9.2]{3.1.9.2} 的记号. (“传递性”) 同构
$$
k_{!}Rf_{V!}\simeq Rf_{!}k_{X!}{\quoi},
$$
的存在性是函子 $Rf_{!}$ 理论的基础; 它转置为\emph{局部化同构} (\hyperref[XVIII.prop3.1.8]{3.1.8}. (iii))
\begin{equation*}
k^{\ast}_{X}Rf^{!}\simeq Rf^{!}_{V}k^{\ast}{\quoi},\tag{3.1.10.1}
\end{equation*}\etiquette[3.1.10.1]{XVIII.eq3.1.10.1}%
在复形层面由态射
$$
k^{\ast}_{X}f^{!}_{\bigdot}\longrightarrow f^{!}_{V\bigdot}k^{\ast}{\quoi}.
$$
表示.

要验证 \eqhyperref[XVIII.eq3.1.9.6]{3.1.9.6} 为同构, 只须验证对每个 $V\in \Ob S_{\text{\rm et}}$, \eqhyperref[XVIII.eq3.1.9.6]{3.1.9.6} 都诱导同构
\begin{align*}
\Hom &(k^{\ast}_{X}K,k^{\ast}_{X}Rf^{!}L)\simeq
 H^{0}(V,Rf_{\ast}R\Hhom (K,Rf^{!}L))\tag{3.1.10.2}\\
& \longrightarrow H^{0}(V,R\Hom(Rf_{!}K,L))\simeq
 \Hhom(k^{\ast}Rf_{!}K,k^{\ast}L)
\end{align*}\etiquette[3.1.10.2]{XVIII.eq3.1.10.2}

当 $V=S$ 时, 这恰是 \eqhyperref[XVIII.eq3.1.4.3]{3.1.4.3}. 一般情形来自下述相容性, \emph{编者尚未验证这一相容性}:

\hypertarget{sga4.exp.XVIII.lemma.3.1.10.3}{}\begin{enonce*}{引理 3.1.10.3}\etiquette[3.1.10.3]{XVIII.lem3.1.10.3}
图\pageoriginale
% SGA-SRC-001158
$$
\xymatrix@C=2.5cm@R=1.5cm{
\Hom(k^{\ast}_XK, k^{\ast}_XRf^{!}L)
\ar[r]^{\scriptstyle{\hyperref[XVIII.eq3.1.10.2]{(3.1.10.2)}}} \ar[d]^{\scriptstyle{\hyperref[XVIII.eq3.1.10.1]{(3.1.10.1)}}}_{\text{\rotatebox{90}{$\backsim$}}}&
\Hom(k^{\ast}Rf_!K,k^{\ast}L) \ar[d]^{\txt{\rm （局部化）}}_{\text{\rotatebox{90}{$\backsim$}}}\\
\Hom(k^{\ast}_XK,Rf^{!}_Vk^{\ast}L)
\ar@{<->}[r]^{\scriptstyle{\rm{\hyperref[XVIII.eq3.1.4.3]{(3.1.4.3\;\text{在}\; V)}}}}_{\sim}& \Hom(Rf_{V!}k^{\ast}_XK,k^{\ast}L)
}
$$
交换.
\end{enonce*}

\setcounter{subsubsection}{10}
\subsubsection{}\etiquette[3.1.11]{XVIII.subsubsec3.1.11}
考虑笛卡儿方块
\begin{equation*}
\vcenter{
\xymatrix@C=2cm@R=1.5cm{
X' \ar[r]^{g'}\ar[d]^{f'}& X\ar[d]^{f}\\
S' \ar[r]^{g} & S,
}}\tag{3.1.11.1}
\end{equation*}\etiquette[3.1.11.1]{XVIII.eq3.1.11.1}%
其中 $f$ 可紧化, $\mathcal{A}$ 为 $S$ 上的挠环层, $\mathcal{B}$ 为 $S'$ 上的挠环层, $M$ 为 $(g^{\ast}\mathcal{A},\mathcal{B})$-双模的上有界复形.
% SGA-SRC-001159
在复形层面, 于是有关于 $K$ $(K\in \Ob
C(X,f^{\ast}\mathcal{A}))$ 具有函子性的态射
\begin{equation*}
M\otimes_{\mathcal{A}}g^{\ast}f^{\bigdot}_{!}K\longrightarrow
f'_{!}{}^{\bigdot}
({f'}^{\ast}M\otimes_{\mathcal{A}}{g'}^{\ast}K){\quoi},\tag{3.1.11.2}
\end{equation*}\etiquette[3.1.11.2]{XVIII.eq3.1.11.2}%
其转置为关于 $L$ $(L\in \Ob
C(S',\mathcal{B}))$ 具有函子性的态射
\begin{equation*}
g'_{\ast}\Hhom_{\mathcal{B}}({f'}^{\ast}M,f'_{\bigdot}{}^{!}L)\longrightarrow
f^{!}_{\bigdot}g_\ast\Hhom_{\mathcal{B}}(M,L){\quoi}.\tag{3.1.11.3}
\end{equation*}\etiquette[3.1.11.3]{XVIII.eq3.1.11.3}

过渡到导出函子, 得到一个关于 $L$ 具有函子性的从 $D^{+}$ $(S',\mathcal{B})$ 到 $D^{+}(X,f^{\ast}\mathcal{A})$ 的态射:
\begin{equation*}
Rg'_{\ast}R\Hhom_{\mathcal{B}}({f'}^{\ast}M,Rf'{}^{!}L)\longrightarrow
Rf^{!}Rg_{\ast}R\Hom_{\mathcal{B}}(M,L){\quoi},\tag{3.1.11.4}
\end{equation*}\etiquette[3.1.11.4]{XVIII.eq3.1.11.4}%
使下图交换, 其中 $K\in \Ob
D^{-}(X,f^{\ast}\mathcal{A})$ 且 $L\in \Ob D^{+}(S',\mathcal{B})$:
$$
\xymatrix{
\scriptscriptstyle{\Hom(K,Rg'_{\ast} R\Hhom_{\Bb}(f'^{\ast}M,Rf'^{!}L))}
\ar@{<->}[d] \ar[r]^{\text{\ding{172}}} &
\scriptscriptstyle{\Hom(K,Rf^{!}Rg_{\ast}R\Hhom_{\Bb}(M,L))}
\ar@{<->}[r]_{{\scriptscriptstyle{\text{伴随}}}}& \scriptscriptstyle{\Hom(Rf_!K, Rg_{\ast}
R\Hhom_{\Bb}(M,L))}\ar@{<->}[d]\\
\scriptscriptstyle{\Hom(g'^{\ast}K,R\Hhom_{\Bb}(f'^{\ast} M, Rf'^{!}L))}\ar@{<->}[d]
& &  \scriptscriptstyle{\Hom(g^{\ast}Rf_! K, R\Hhom_{\Bb} (M,L))} \ar@{<->}[d]\\
\scriptscriptstyle{\Hom(f'^{\ast}M\otimes_{\Aa}g'^{\ast}K,
Rf'^{!}L)} \ar@{<->}[r]_{\scriptscriptstyle{\text{伴随}}}& \scriptscriptstyle{\Hom(Rf'^{\cdot}_!(f'^{\ast}M\otimes^{L}_\Aa
g'^{\ast}K),L)}
\ar[r]^{\text{\ding{173}}} &
\scriptscriptstyle{\Hom_{\Bb}(M\mathop{\otimes}^{\LL}_\Aa g^{\ast}Rf_!K, L)}.
}
$$

在\pageoriginale 这一图中, \ding{173} (由箭头 \hyperrefext[XVII][XVII.eq5.2.2.1]{5.2.2.1} 导出) 为同构 (\hyperrefext[XVII][XVII.thm5.2.6]{5.2.6}). 因而箭头 \ding{172} 对每个 $K$ 都是同构, 由此得出 \eqhyperref[XVIII.eq3.1.11.4]{3.1.11.4} 为同构.

固有基变换定理于是转置为下述结论.

\hypertarget{sga4.exp.XVIII.proposition.3.1.12}{}\begin{enonce*}{命题 3.1.12}\etiquette[3.1.12]{XVIII.prop3.1.12}
箭头 \eqhyperref[XVIII.eq3.1.11.4]{3.1.11.4} 为同构.
\end{enonce*}

这一同构称为\emph{诱导同构}（induction isomorphism）.

有如下三个特例.

\hypertarget{sga4.exp.XVIII.corollary.3.1.12.1}{}\begin{enonce*}{推论 3.1.12.1}\etiquette[3.1.12.1]{XVIII.cor3.1.12.1}
设 $f:X\rightarrow S$ 为可紧化态射, $\mathcal{A}$ 与 $\mathcal{B}$ 为 $S$ 上的两个挠环层, 并设 $\varphi:\mathcal{A}\rightarrow \mathcal{B}$ 为同态. 以 $\rho$ 表示从 $\mathcal{B}$ 到 $\mathcal{A}$ 的标量限制函子. 则对 $K\in \Ob
D^{+}(S,\mathcal{B})$, 有
$$
\rho Rf^{!}K\xrightarrow{\sim} Rf^{!}(\rho K){\quoi}.
$$
\end{enonce*}

只须在 \eqhyperref[XVIII.eq3.1.11.3]{3.1.11.3} 中取 $S=S'$, $X=X'$, $\mathcal{A}=\mathcal{A}$, $\mathcal{B}=\mathcal{B}$, $M=\mathcal{B}$.

因而环层在 $Rf^{!}$ 的构造中不起本质作用.

\hypertarget{sga4.exp.XVIII.corollary.3.1.12.2}{}\begin{enonce*}{推论 3.1.12.2}\etiquette[3.1.12.2]{XVIII.cor3.1.12.2}
设 $f:X\rightarrow S$ 为可紧化态射, $\mathcal{A}$ 为 $S$ 上的挠环层, 并设 $K\in \Ob D^{-}$ $(S,\mathcal{A})$, $L\in \Ob D^{+}(S,\mathcal{A})$. 有\emph{诱导公式}
\begin{gather*}
R\Hhom(f^{\ast}K,Rf^{!}L)\xrightarrow{\sim} Rf^{!}R\Hhom(K,L)\\
(\text{\rm 取 } S=S', X=X'){\quoi}.
\end{gather*}
\end{enonce*}

\hypertarget{sga4.exp.XVIII.corollary.3.1.12.3}{}\begin{enonce*}{推论 3.1.12.3}\etiquette[3.1.12.3]{XVIII.cor3.1.12.3}
% SGA-SRC-001161
对每个方块 \hyperref[XVIII.subsubsec3.1.11]{3.1.11} 以及每个 $L\in \Ob D^{+}(S',\mathcal{B})$, 有
\begin{gather*}
Rg'_{\ast}R{f'}^{!}L\xrightarrow{\sim} Rf^{!}Rg_{\ast}L\\
(\text{\rm 取 } \mathcal{B}=g^{\ast}\mathcal{A}=M){\quoi}.
\end{gather*}
\end{enonce*}

\setcounter{subsubsection}{12}
\subsubsection{}\etiquette[3.1.13]{XVIII.subsubsec3.1.13}
对\pageoriginale $S$-可紧化态射的复合 $gh$
$$
X\xrightarrow{h}Y\xrightarrow{g}Z
$$

复合同构
$$
c_{g,h}:Rg_{!}Rh_{!}\xrightarrow{\sim} R(gh)_{!}
$$
转置为复合态射\nde{转置态射实际上沿另一方向, 但由于它是同构, 这不造成问题.}
\begin{equation*}
c^{!}_{g,h}:Rh^{!}Rg^{!}\xrightarrow{\sim} R(gh)^{!}\tag{3.1.13.1}
\end{equation*}\etiquette[3.1.13.1]{XVIII.eq3.1.13.1}%
对三重复合满足通常的上循环条件. 换言之, \emph{范畴 $D^{+}$ $(X,f^{\ast}\mathcal{A})$ 构成在以 $S$-概形} $f:X\rightarrow S$ \emph{为对象、以 $S$-可紧化态射为箭头的范畴上的一个纤维且余纤维范畴的各纤维, 正像} (\resp \emph{逆像}) \emph{函子为函子} $Rf_{!}$ (\resp $Rf^{!}$). 因而对可紧化态射的每个交换方块,
$$
\xymatrix@=1.75cm{
X' \ar[r]_{g} \ar[d]^{f'}& X \ar[d]^{f}\\
S' \ar[r]_{g'}& S,
}
$$
都有余基变换态射 (\hyperrefext[XVII][XVII.prop2.1.3]{2.1.3})
\begin{equation*}
ch^{!}:Rf'_{!}Rg^{!}\longrightarrow Rg^{!} Rf_{!}{\quoi}.\tag{3.1.13.2}
\end{equation*}\etiquette[3.1.13.2]{XVIII.eq3.1.13.2}

\subsubsection{}\etiquette[3.1.14]{XVIII.subsubsec3.1.14}
考虑笛卡儿方块
$$
\xymatrix@=1.75cm{
X' \ar[r]^{g'} \ar[d]^{f'}& X \ar[d]^{f}\\
S' \ar[r]^{g}& S,
}
$$
其中 $f$ 可紧化, 并设 $\mathcal{A}$ 为 $S$ 上的挠环层. 以 $\mathcal{A}$ 的逆像作为环层给 $S'$, $X$ 及 $X'$ 赋环, 滥用记号仍\pageoriginale 将这些环层记为 $\mathcal{A}$. 基变换同构
\begin{equation*}
g^{\ast}Rf_{!}\xrightarrow{\sim} Rf'_{!}{g'}^{\ast}\tag{3.1.14.1}
\end{equation*}\etiquette[3.1.14.1]{XVIII.eq3.1.14.1}%
使我们可以将范畴 $D(S,\mathcal{A})$, $D(X,\mathcal{A})$, $D(S',\mathcal{A})$ 及 $D(X',\mathcal{A})$ 看成以交换图
$$
\xymatrix@C=2cm@R=1cm{
X' \ar[d]& X\ar[l]\ar[d]\\
S' & S\ar[l],
}
$$
定义的范畴上的一个余纤维范畴的纤维, 正像函子为函子 $Rf_{!}$, $Rf'_{!}$, $g^{\ast}$ 及 ${g'}^{\ast}$ (原文如此). 若限于范畴 $D^{+}$, 甚至得到这一图上的一个纤维且余纤维范畴, 正像函子为函子 $Rf^{!}$, $R{f'}^{!}$, $Rg_{\ast}$ 及 $Rg'_{\ast}$ (原文如此). 由此得到
\begin{enumerate}
\renewcommand{\theenumi}{\alph{enumi}}
\renewcommand{\labelenumi}{\theenumi)}
\item “逆像”函子的一个“传递性同构” $(\sic)$
% SGA-SRC-001163
$$
Rg'_{\ast} R{f'}^{!}\xrightarrow{\sim} Rf^{!}Rg_{\ast}
$$
它已在 \hyperref[XVIII.cor3.1.12.3]{3.1.12.3} 中得到 (同样通过转置 \eqhyperref[XVIII.eq3.1.14.1]{3.1.14.1} 来定义).

\item 一个“基变换态射” (\hyperrefext[XVII][XVII.prop2.1.3]{2.1.3})
\begin{equation*}
{g'}^{\ast}Rf^{!}\longrightarrow R{f'}^{!}g^{\ast}\tag{3.1.14.2}
\end{equation*}\etiquette[3.1.14.2]{XVIII.eq3.1.14.2}%
由 \loccit, 可用下述两种方式描述它:
\end{enumerate}

\begin{enumerate}
\renewcommand{\theenumi}{\alph{enumi}}
\renewcommand{\labelenumi}{\theenumi)}
\item 由伴随性, 给出态射 \eqhyperref[XVIII.eq3.1.14.2]{3.1.14.2} 相当于给出态射
$$
Rf'_{!}{g'}^{\ast}Rf^{!}\longrightarrow g^{\ast}{\quoi}.
$$
基变换定理将这一态射给为复合态射
$$
Rf'_{!}{g'}^{\ast}Rf^{!}\xleftarrow{\sim}
g^{\ast}Rf_{!}Rf^{!}\xrightarrow{g^{\ast}\ast \text{伴随}} g^{\ast}{\quoi}.
$$

\item 由\pageoriginale 伴随性, 给出态射 \eqhyperref[XVIII.eq3.1.14.2]{3.1.14.2} 相当于给出态射
$$
Rf^{!}\longrightarrow Rg'_{\ast}R{f'}^{!}g^{\ast}{\quoi}.
$$
同构 (\hyperref[XVIII.cor3.1.12.3]{3.1.12.3}) 将这一态射给为复合态射
$$
Rf^{!}\xrightarrow[Rf^{!}\ast \text{伴随}]{} Rf^{!}Rg_{\ast}
g^{\ast}\xleftarrow{\sim} Rg'_{\ast}R{f'}^{!}g^{\ast}{\quoi}.
$$
\end{enumerate}

由它们在纤维且余纤维范畴中的解释, 显然, \emph{当 $f$ 或 $g$ 为两个态射的复合时, 态射} (\hyperref[XVIII.cor3.1.12.3]{3.1.12.3}), \eqhyperref[XVIII.eq3.1.13.2]{3.1.13.2} \emph{及} \eqhyperref[XVIII.eq3.1.14.2]{3.1.14.2} \emph{各自满足 \hyperrefext[XII][XII.prop4.4]{4.4} 类型的两项相容性\nde{这些构造也可在交换结构与交叉函子的背景下解释, \cf \cite{sga4.exp.XVIII.bib.0}.}.}

以下一般拓扑的结果将在下一小节使用.

\subsubsection{}\etiquette[3.1.15]{XVIII.subsubsec3.1.15}
设 $f:X\rightarrow S$ 为景态射. 以 $\Gamma(f)$ 表示下述景:
\begin{enumerate}
\renewcommand{\theenumi}{\alph{enumi}}
\renewcommand{\labelenumi}{\theenumi)}
\item 范畴 $\Gamma(f)$ 为三元组 $(U,V,\varphi)$ 的范畴, 其中
$$
U\in \Ob (X), V\in \Ob S,\quad \varphi\in \Hom
(U,f^{\ast}V){\quoi}.
$$

\item 一族态射 $(u_{i},v_{i}):(U_{i},V_{i},\varphi_{i})\rightarrow (U,V,\varphi)$ 为覆盖族, 当且仅当在 $X$ 中, 态射 $u_{i}:U_{i}\rightarrow U$ 覆盖 $U$.
\end{enumerate}

于是有:

\paragraph{}\etiquette[3.1.15.1]{XVIII.par3.1.15.1}
由
$$
(U,V,\varphi)\longmapsto U
$$
给出的从 $\Gamma(f)$ 到 $X$ 的函子为景的等价 $\gamma:X\rightarrow \Gamma(f)$.

\paragraph{}\etiquette[3.1.15.2]{XVIII.par3.1.15.2}
从 $S$ 到 $\Gamma(f)$ 的函子 $f^{\ast}_{\Gamma}$
$$
V\longmapsto (f^{\ast}V,V,\Id_{f^{\ast}V})
$$
是一个景态射, 记为 $f_{\Gamma}:\Gamma(f)\rightarrow S$, 或滥用记号简记 $f$. 它满足 $f_{\Gamma}\gamma\simeq f$. 函子 $f^{\ast}_{\Gamma}$ 有左伴随 $f_{\Gamma\bigdot}$:
$$
f_{\Gamma\bigdot}:(U,V,\varphi)\longmapsto V{\quoi}.
$$
因而函子\pageoriginale $f_{\Gamma\bigdot}$ 为从 $\Gamma(f)$ 到 $S$ 的景余态射（comorphism of sites）\nde{由于这一术语并不常见, 此处应说明, 所指的是 \hyperrefext[III][III.sec2]{2} 及 \hyperrefext[IV][IV.subsec4.7]{4.7} 中出现的概念 (应记住, 原版中第二个引用之处的排印错误, 曾使人无法理解箭头究竟沿哪个方向).}, 并与景态射 $f_{\Gamma}$ 定义相同的拓扑斯态射.
\begin{equation*}
\vcenter{
\xymatrix@R=2.2cm@C=1.75cm{
X &         &\Gamma(f) \ar[ll]^{\approx}
\ar@<4pt>[ld]^{f_{\Gamma}}& \\
  &    S \ar[lu]^{f^{\ast}}  \ar@<4pt>[ru]^{f^{\ast}_{\Gamma}} & .
}}\tag{3.1.15.3}
\end{equation*}\etiquette[3.1.15.3]{XVIII.eq3.1.15.3}

由此得到计算层的逆像的如下规则\nde{此处使用了 \hyperrefext[III][III.prop2.3]{2.3} 或 \hyperrefext[IV][IV.eq4.7.3]{4.7.3}.}

\setcounter{paragraph}{3}
\paragraph{}\etiquette[3.1.15.4]{XVIII.par3.1.15.4}
设 $\Ff$ 为 $S$ 上的层, $f^{\bigdot}_{\Gamma}\Ff$ 为 $\Gamma(f)$ 上由
$$
f^{\bigdot}_{\Gamma}\Ff(U,V,\varphi)=\Ff(V){\quoi},
$$
给出的预层, $af^{\bigdot}_{\Gamma}\Ff$ 为其伴随层. 则
$$
f^{\ast}\Ff(U)\simeq
af^{\bigdot}_{\Gamma}\Ff(U,V,\varphi){\quoi}.
$$

\subsubsection{}\etiquette[3.1.16]{XVIII.subsubsec3.1.16}
设 $f:X\rightarrow S$ 为景态射, $\mathcal{A}$ 为 $S$ 上的环层, $\mathcal{B}$ 为 $X$ 上的环层, $d$ 为整数. 设 $K$ 为下述类型的对象:

\paragraph{}\etiquette[3.1.16.1]{XVIII.par3.1.16.1}
$K$ 将 $X$ 的每个对象 $U$ 映为 $S$ 上的一个 $(\mathcal{A},f_{|U\ast}(\mathcal{B}|U))$-模复形 $K(U)$, 在次数 $>d$ 处为零, 且 $K(U)$ 为 $U$ 的\emph{协变}函子.

对 $S$ 上的每个 $\mathcal{A}$-模层 $ℱ$, 以 $K^{!}ℱ$ 表示 $X$ 上由 $\mathcal{B}$-模预层的复形
$$
U\longmapsto \Hom_{\mathcal{A}}(K(U),ℱ){\quoi}.
$$
生成的层复形.

以 $RK^{!}$ 表示 $K^{!}$ 的导出函子
$$
RK^{!}:D^{+}(S,\mathcal{A})\longrightarrow
D^{+}(X,\mathcal{B}){\quoi}.
$$
若 $U:K_{1}\rightarrow K_{2}$ 为 \hyperref[XVIII.par3.1.16.1]{3.1.16.1} 的对象之间的态射, 则 $U$ 诱导\pageoriginale 函子态射
$$
{}^{t}u:K^{!}_{2}\longrightarrow K^!_{1}\quaet RK^{!}_{2}\longrightarrow
RK^{!}_{1}{\quoi}.
$$

在上述假设下, 有:

\begin{enonce*}{命题 3.1.17}\etiquette[3.1.17]{XVIII.prop3.1.17}
\begin{enumerate}
\renewcommand{\theenumi}{\roman{enumi}}
\renewcommand{\labelenumi}{(\theenumi)}
\item 对 $X$ 的每个点 $x$ 以及每个 $L\in
\Ob D^{+}(S,\mathcal{A})$, 以 $V(x)$ 表示 $x$ 的邻域范畴, 则有
\begin{equation*}
ℋ^{n}(RK^{!}L)_{x}\simeq \varinjlim_{U\in \Ob V(x)}
\Ext^{n}(K(U),L){\quoi}.\tag{3.1.17.1}
\end{equation*}\etiquette[3.1.17.1]{XVIII.eq3.1.17.1}

\item 因而有谱序列
\begin{equation*}
% SGA-SRC-001166
E^{pq}_{2}=\varinjlim_{U\in \Ob V(x)}
\Ext^{p}(ℋ^{-q}(K(U)),L)\Longrightarrow
ℋ^{p+q}(RK^{!}L)_{x}{\quoi}.\tag{3.1.17.2}
\end{equation*}\etiquette[3.1.17.2]{XVIII.eq3.1.17.2}

\item 若 $X^{0}$ 为 $X$ 的一个保守点集, 且 $u:K_{1}\rightarrow K_{2}$ 为对象 \hyperref[XVIII.par3.1.16.1]{3.1.16.1} 的一个态射, 则要使
$$
{}^{t}u:RK^{!}_{2}\longrightarrow RK^{!}_{1}
$$
为同构, 充分必要条件是, 对每个 $n$ 以及每个 $x\in X^{0}$, $\Mod(S,\mathcal{A})$ 中的 pro-对象态射
\begin{equation*}
u_{x}:\mathop{``\varprojlim\mbox{''}}_{U\in \Ob V(x)}
ℋ^{n}K_{1}(U)\longrightarrow \mathop{``\varprojlim\mbox{''}}_{U\in
\Ob V(x)} ℋ^{n}K_{2}(U)\tag{3.1.17.3}
\end{equation*}\etiquette[3.1.17.3]{XVIII.eq3.1.17.3}%
为同构.
\end{enumerate}
\end{enonce*}

为证明 (i), 取 $L$ 为内射复形, 则 \eqhyperref[XVIII.eq3.1.17.1]{3.1.17.1} 是直接的; 显然 (i) $\Rightarrow$ (ii).

(iii) 中条件 \eqhyperref[XVIII.eq3.1.17.3]{3.1.17.3} 的充分性直接来自 \eqhyperref[XVIII.eq3.1.17.2]{3.1.17.2}. 我们不需要其必要性. 若 ${}^{t}u$ 为同构, 则由 (\hyperref[XVIII.par3.1.16.1]{3.1.16.1}), 对 $(S,\mathcal{A})$ 上的每个内射层 $I$, 有
$$
\Hom(\mathop{``\varprojlim\mbox{''}}_{V(x)}ℋ^{n}K_{2}(U),I)\xrightarrow{\sim}\Hom(\mathop{``\varprojlim\mbox{''}}_{V(x)}
ℋ^{n}K_{1}(U),I){\quoi}, 
$$
最后注意, 在每个具有足够多内射对象的阿贝尔范畴 (这里为 $\Mod(S,\mathcal{A})$) 中, 当 $I$ 为内射 $\mathcal{A}$-模时, 函子 $\Hom(\ast,I)$ 构成 $\Pro(\mathcal{A})$ 上的一个保守函子族, 即得结论.

\begin{enonce*}[remark]{例 3.1.18}\etiquette[3.1.18]{XVIII.exem3.1.18}
设\pageoriginale $f:X\rightarrow S$ 为可紧化概形态射, $\mathcal{A}$ 为 $S$ 上的挠环层. 以 $f^{\ast}\mathcal{A}$ 为环层给 $X$ 赋环. 选择一个整数 $d$, 作为 $f$ 的纤维维数的上界, 以及 $f$ 的一个紧化, 并令 (\eqhyperref[XVIII.eq3.1.4.7]{3.1.4.7})
\begin{equation*}
K(U)=f^{\bigdot}_{!}((f^{\ast}\mathcal{A})_{U}){\quoi}.\tag{3.1.18.1}
\end{equation*}\etiquette[3.1.18.1]{XVIII.eq3.1.18.1}
\end{enonce*}

由 $f^{\bigdot}_{!}$ 的伴随的公式 \eqhyperref[XVIII.eq3.1.3.1]{3.1.3.1}, 有同构
\begin{equation*}
f^{!}_{\bigdot}\simeq K^{!}:\Mod(S,\mathcal{A})\longrightarrow C^{b}\Mod(X,f^{\ast}\mathcal{A}){\quoi}.\tag{3.1.18.2}
\end{equation*}\etiquette[3.1.18.2]{XVIII.eq3.1.18.2}%
因而
\begin{equation*}
Rf^{!}=RK^{!}:D^{+}(S,\mathcal{A})\longrightarrow
D^{+}(X,f^{\ast}\mathcal{A}){\quoi}.\tag{3.1.18.3}
\end{equation*}\etiquette[3.1.18.3]{XVIII.eq3.1.18.3}

\begin{enonce*}[remark]{例 3.1.19}\etiquette[3.1.19]{XVIII.exem3.1.19}
设 $f:X\rightarrow S$, $\mathcal{A}$, $\mathcal{B}$ 及 $d$ 如 \hyperref[XVIII.subsubsec3.1.16]{3.1.16}. 设 $K$ 为\sisi{下述}{下述}类型的对象.
\end{enonce*}

\setcounter{subsubsection}{19}
\setcounter{paragraph}{0}
\paragraph{}\etiquette[3.1.19.1]{XVIII.par3.1.19.1}
对每个 $W=(U,V,\varphi)\in \Ob \Gamma(f)$ (\hyperref[XVIII.subsubsec3.1.15]{3.1.15})\nde{按定义, $\varphi$ 是一个 $X$-态射 $U\to V\times_S X$. 要正确理解这一例子, 必须采用等价的观点, 将 $\varphi$ 解释为一个 $S$-态射 $\varphi\colon U\to V$.}, $K(W)$ 为 $V$ 上的 $(\mathcal{A}_{V},\varphi_{\ast}\mathcal{B}_{|U})$-模复形, 关于 $W$ 具有函子性: 对每个态射 $(u,v):
W_{1}=(U_{1},V_{1},\varphi_{1})\to W_{2}=(U_{2},V_{2},\varphi_{2})$, 由
\begin{align*}
K(u,v) &: K(W_{1})\longrightarrow v^{\ast}K(W_{2}),\quad\mbox{\ie}\\
K(u,v) &: v_{!}K(W_{1})\longrightarrow K(W_{2}){\quoi}.
\end{align*}
给出.

与每个 (\hyperref[XVIII.par3.1.19.1]{3.1.19.1}) 类型的对象 $K$ 相伴, 有一个关于 $f_{\Gamma}:\Gamma(f)\rightarrow S$ 的 (\hyperref[XVIII.par3.1.16.1]{3.1.16.1}) 类型的对象 $K'$, 由
\begin{equation*}
K'(U,V,\varphi)=j_{!}K(U,V,\varphi)\tag{3.1.19.2}
\end{equation*}\etiquette[3.1.19.2]{XVIII.eq3.1.19.2}%
给出, 其中 $j$ 为 $V$ 到 $S$ 的“包含态射”. 令
$$
K^{!}\mathop{=}_{\defn} {K'}^{!}\quaet RK^{!}\mathop{=}_{\defn}
R{K'}^{!}{\quoi}.
$$

\begin{enonce*}[remark]{例 3.1.20}\etiquette[3.1.20]{XVIII.exem3.1.20}
用 \hyperref[XVIII.exem3.1.19]{3.1.19} 的记号, 取 $\mathcal{B}=f^{\ast}\mathcal{A}$, 并令
\begin{equation*}
K(U,V,\varphi)=\mathcal{A}_{V}\quad\text{\rm(置于次数 0)}\tag{3.1.20.1}
\end{equation*}\etiquette[3.1.20.1]{XVIII.eq3.1.20.1}
\end{enonce*}

由\pageoriginale \hyperref[XVIII.par3.1.15.4]{3.1.15.4}, 有函子的同构
\begin{gather*}
K^{!}\simeq f^{\ast}\tag{3.1.20.2}\\
RK^{!}\simeq f^{\ast}{\quoi}.\tag{3.1.20.3}
\end{gather*}\etiquette[3.1.20.2]{XVIII.eq3.1.20.2}
\etiquette[3.1.20.3]{XVIII.eq3.1.20.3}

\subsection{Poincaré 对偶}\etiquette[3.2]{XVIII.subsec3.2}

\subsubsection{}\etiquette[3.2.1]{XVIII.subsubsec3.2.1}
设 $f$ 为可紧化的平坦态射, 纯相对维数为 $d$ (或更一般地, 满足 \hyperref[XVIII.thm2.9]{2.9} 的条件 $(\ast)_{d}$):
$$
\xymatrix@C=2cm@R=1.5cm@M=.2cm{
X \ar@{^{(}->}[r]^{j} \ar[d]_{f}& \bar{X}\ar[ld]_{\bar{f}} &\\
S &.
}
$$
设 $n$ 为在 $S$ 上可逆的整数 $\geq 1$, 并假设 $S$ 以环层 $\mathcal{A}$ 赋环, 使得 $n\mathcal{A}=0$. 由 \hyperref[XVIII.exem3.1.18]{3.1.18}, 函子 $Rf^{!}:D^{+}(S,\mathcal{A})\rightarrow D^{+}(X,f^{\ast}\mathcal{A})$ 可用 \hyperref[XVIII.subsubsec3.1.16]{3.1.16} 的方法计算, 其中 $K$ 由 \eqhyperref[XVIII.eq3.1.18.1]{3.1.18.1} 给出.

设 $K'$ 为 (\hyperref[XVIII.par3.1.19.1]{3.1.19.1}) 类型的函子, 它将 $(U,V,\varphi)\in \Ob \Gamma(f)$ 映为 $V$ 上的层复形 $\ZZ/n(-d)[-2d]$, 由 $\ZZ/n(-d)$ 放在次数 $2d$ 处组成; 并设 $K''$ 为由 \eqhyperref[XVIII.eq3.1.19.2]{3.1.19.2} 对应的关于 $f:\Gamma(f)\rightarrow S$ 的 (\hyperref[XVIII.par3.1.16.1]{3.1.16.1}) 类型的对象.

可以将 $K$ 与 $K''$ 视为关于 $\Gamma(f)\rightarrow S$ 的 \hyperref[XVIII.par3.1.16.1]{3.1.16.1} 的对象. 此外, 对每个 $(U,V,\varphi)\in \Ob
\Gamma(f)$, 迹态射 \hyperref[XVIII.thm2.9]{2.9} 定义 $V$ 上的层态射
$$
\Tr:R^{2d}\varphi_{!}\ZZ/n\longrightarrow
\ZZ/n(-d){\quoi},
$$
从而以 $j$ 表示从 $V$ 到 $S$ 的态射, 得到函子态射
\begin{gather*}
K(U)=\bar{f}_{\ast}\tau_{\leq
2d}C^{\ast}_{\ell}\ZZ/n_{U}\longrightarrow
ℋ^{2d}(K(U))[-2d]\simeq\tag{3.2.1.1}\\
R^{2d}f_{!}(\ZZ/n)_{U}\simeq j_{!}R^{2d}\varphi_{!} \ZZ/n
\xrightarrow{\Tr_{\varphi}} j_{!}
\ZZ/n(-d)=K''(U,V,\varphi){\quoi}.
\end{gather*}\etiquette[3.2.1.1]{XVIII.eq3.2.1.1}

\hyperref[XVIII.exem3.1.20]{3.1.20} 的\pageoriginale 一个扭转形式表明
$$
{RK''}^{!}=Rf^{\ast}(d)[2d]{\quoi}.
$$
因而由 \hyperref[XVIII.subsubsec3.1.16]{3.1.16}, 有复合态射
\begin{equation*}
t_{f}:Rf^{\ast}(d)[2d]\simeq {RK''}^{!}\longrightarrow RK^{!}\simeq
Rf^{!}{\quoi}.\tag{3.2.1.2}
\end{equation*}\etiquette[3.2.1.2]{XVIII.eq3.2.1.2}

\begin{enonce*}[remark]{注 3.2.2}\etiquette[3.2.2]{XVIII.rem3.2.2}
若函子 $Rf^{!}$ 有有限上同调维数, 则 $t_f$ 仍有定义, 为从 $D(S,\mathcal{A})$ 到 $D(X,f^{\ast}\mathcal{A})$ 的函子之间的态射 (定义见 \hyperref[XVIII.rem3.1.5]{3.1.5}).
\end{enonce*}

还有另一种构造箭头 \eqhyperref[XVIII.eq3.2.1.2]{3.2.1.2} 的方法, 即由 \eqhyperref[XVIII.eq2.13.2]{2.13.2} 通过伴随性得到
$$
\Tr:Rf_{!}(f^{\ast}L(d)[2d])\longrightarrow L{\quoi}.
$$
这两条箭头一致:

\begin{enonce*}{引理 3.2.3}\etiquette[3.2.3]{XVIII.lem3.2.3}
对 $L\in \Ob D^{+}(S,\mathcal{A})$ (\resp $D(S,\mathcal{A})$, \emph{若函子 $Rf^{!}$ 有有限上同调维数}), 图
$$
\xymatrix@C=2.5cm@R=1.75cm{
Rf_!(f^{\ast}L(d)[2d]) \ar[d]^{\Tr_f} \ar[r]^{Rf_!(t_f)}  & Rf_!
Rf^{!}L \ar[d]^{\txt{\rm 伴随}}\\
L                              \ar@{=}[r] & L
}
$$
交换.
\end{enonce*}

先用函子 $f^!_\bigdot$ \eqhyperref[XVIII.eq3.1.4.10]{3.1.4.10} 给出 $t_f$ 定义的一个变体. 设 $U\in \Ob X_{\text{\rm et}}$, $V\in \Ob S_{\text{\rm et}}$ 且 $\varphi\in \Hom_{f}(U,V)$. 按定义,
$$
f^{!}_{\bigdot}(L)(U)=\Hom(f^{\bigdot}_{!}\mathcal{A}_{U},L){\quoi}\;;
$$
迹态射定义一个态射
$$
\Tr_{f}:f^{\bigdot}_{!}\mathcal{A}_{U}\longrightarrow
\mathcal{A}_{V}(-d)[-2d]
$$
其转置为
$$
t_{\varphi}:f^{!}_{\bigdot}(L)(U)\longleftarrow L(d)[2d](V){\quoi}.
$$
态射\pageoriginale $t_{\varphi}$ 定义一个态射
$$
t_{f}:f^{\ast}L(d)[2d]\longrightarrow f^{!}_{\bigdot}L
$$
将验证它在导出范畴中与 \eqhyperref[XVIII.eq3.2.1.2]{3.2.1.2} 一致的工作留给读者. 这一描述表明, 对模复形 $L$, 下图在复形层面交换:
$$
\xymatrix@R=1.25cm{
\Hom^{\bigdot}(\Aa_V,L) \ar[r] \ar[d]^{\Tr_f}& \Hom^{\bigdot}(f^{\ast}\Aa_V(d) [2d], f^{\ast}
L(d) [2d])\ar[d]^{t_f}\\
\Hom(f^{\bigdot}_!f^{\ast}\Aa_V(d)[2d],L) \ar@{-}[r]^{\sim}& \Hom^{\bigdot}(f^{\ast}\Aa_V(d)[2d], f^{!}_{\bigdot}L).
}
$$
由这一图中的箭头关于 $\mathcal{A}_{V}$ 的函子性, 可得对每个层或层复形 $K$, 下图交换:
$$
\xymatrix{
\Hom^{\bigdot}(K,L) \ar[r] \ar[d]^{\Tr_f}& \Hom^{\bigdot}(f^{\ast}K(d)[2d],
f^{\ast}L(d)[2d]) \ar[d]\\
\Hom^{\bigdot}(f^{\bigdot}_!f^{\ast}K(d)[2d],L) \ar@{<->}[r]^{\sim} &\Hom^{\bigdot}(f^{\ast}K(d)[2d], f^{!}_{\bigdot}L).
}
$$
取 $K=L$ 及 $\Id:L\rightarrow L$, 就在复形层面得到 \hyperref[XVIII.lem3.2.3]{3.2.3}.

若有人理解了这一证明并愿向我解释, 我将不胜感激.

\setcounter{subsubsection}{3}
\subsubsection{}\etiquette[3.2.4]{XVIII.subsubsec3.2.4}
由此引理以及通过伴随性定义的复合同构 (\hyperref[XVIII.subsubsec3.1.13]{3.1.13}), 立即可得: 若 $f$ 为两个平坦、有限表示、可紧化态射的复合, 它们的纯相对维数分别为 $d_{1}$ 与 $d_{2}:f=gh$, 且 $d=d_{1}+d_{2}$, 则图
$$
\xymatrix{
h^{\ast}(g^{\ast}L(d_2)[2d_2])(d_1)[2d_1]
\ar@{-}[r]^(.6){\sim}
\ar[d]^{t_h\ast t_g}&(gh)^{\ast}L(d)[2d] \ar[d]^{t_{gh}}\\
Rh^{!} \; Rg^{!} \; L  \ar@{-}[r]^{\sim}&  R(gh)^{!}L
}
$$
\pageoriginale 交换.

\begin{enonce*}{定理 3.2.5}[Poincaré 对偶]%
\etiquette[3.2.5]{XVIII.thm3.2.5}
设 $f:X\rightarrow S$ 为可紧化光滑态射\emph{(从而 $S$ 拟紧且拟分离 (原文如此))}, 设 $d$ 为 $X$ 上的局部常值函数“$f$ 的相对维数”, $n\geq 1$ 为在 $S$ 上可逆的整数, $\mathcal{A}$ 为 $S$ 上满足 $n\mathcal{A}=0$ 的环层. 则函子
$$
Rf_{!}:D(X,f^{\ast}\mathcal{A})\longrightarrow D(S,\mathcal{A})
$$
容许如下从 $D(S,\mathcal{A})$ 到 $D(X,f^{\ast}\mathcal{A})$ 的函子为右伴随:
$$
K\longmapsto f^{\ast}K(d)[2d]
$$
并以态射 \eqhyperref[XVIII.eq2.13.2]{2.13.2}
$$
\Tr_{f}:Rf_{!}f^{\ast}K(d)[2d]\longrightarrow K
$$
为伴随箭头:
\hypertarget{sga4.exp.XVIII.equation.3.2.5.1}{}$$
\Hom(K,f^{\ast}L(d)[2d])\xrightarrow{\sim} \Hom(Rf_{!}K,L){\quoi}.
$$
\end{enonce*}

\noindent
\emph{证明}. 化为 $d$ 常值的情形. 由箭头 $t_{f}$ \eqhyperref[XVIII.eq3.2.1.2]{3.2.1.2} 的描述 \hyperref[XVIII.subsubsec3.2.1]{3.2.1}, 定理 \hyperref[XVIII.thm2.14]{2.14} 连同判据 \hyperref[XVIII.prop3.1.17]{3.1.17} (iii), 蕴含箭头 $t_{f}$ 为同构 (对 $L\in D^{+}(S,\mathcal{A})$):
$$
t_{f}:f^{\ast}L(d)[2d]\xrightarrow{\sim} Rf^{!}L{\quoi}.
$$

特别地, 函子 $Rf^{!}$ 有有限上同调维数, 因而在无界导出范畴上有定义. 对 $L\in D(S,\mathcal{A})$, 态射 $t_{f}$ 仍为同构\nde{更准确地, $Rf^!$ 有有限上同调维数这一事实, 使我们可以证明能够在无界导出范畴上定义函子 $Rf^!$, 使得对内射层复形 $I^\bigdot$, $Rf^!I^\bigdot$ 为双复形 $f^!_\bigdot I^\bigdot$ 所伴随的全复形. 于是, 对每个 $n\in\ZZ$, $t_f$ 在 $I^\bigdot$ 处取值后在第 $n$ 个上同调层上诱导的态射, 与对于足够小的 $k$ 所取的愚蠢截断 $I^{\geq k}$ 的同一态射等同, 从而证明断言.}.

按\pageoriginale 定义, $Rf^{!}$ 为函子 $Rf_{!}$ 的右伴随; 由 \hyperref[XVIII.lem3.2.3]{3.2.3}, 从 $Rf_{!}Rf^{!}K$ 到 $K$ 的伴随箭头, 经 $t_{f}$ 等同于 $\Tr_{f}$, 这就完成了 \hyperref[XVIII.thm3.2.5]{3.2.5} 的证明.

从现在起, 利用 $t_{f}$ 将函子 $f^{\ast}L(d)[2d]$ 与 $Rf^{!}L$ 等同.

\setcounter{subsubsection}{5}
\subsubsection{}\etiquette[3.2.6]{XVIII.subsubsec3.2.6}
设 $S$ 为代数闭域 $k$ 的谱, $X$ 为可紧化概形, 在 $k$ 上光滑且纯相对维数为 $d$; 设 $n$ 为 $k$ 中可逆的整数, $F$ 为局部常值可构造 $\ZZ/n$-模层, \ie 被 $n$ 零化的“系数系统”. 以 ${}^\vee$ 表示 $\ZZ/n$-对偶, 有
$$
R\Hhom(F,\ZZ/n)=F^{\vee}\quad\text{\rm (置于次数 0)}{\quoi}.
$$
这一公式来自 $\ZZ/n$ 为内射 $\ZZ/n$-模这一事实.
% SGA-SRC-001173
因而 \hyperref[XVIII.thm3.2.5]{3.2.5} 的伴随同构, 或更准确地, 从中导出的同构
$$
R\Gamma\;R\Hhom(F,\ZZ/n(d)[2d])\xrightarrow{\sim}
R\Hom(R\Gamma_{c}(F),\ZZ/n){\quoi},
$$
在此写成
\begin{equation*}
R\Gamma (F^{\vee}(d)[2d])\xrightarrow{\sim} R\Hom
(R\Gamma_{c}(F),\ZZ/n){\quoi}.\tag{3.2.6.1}
\end{equation*}\etiquette[3.2.6.1]{XVIII.eq3.2.6.1}%
过渡到上同调群, 再次利用 $\ZZ/n$ 为内射 $\ZZ/n$-模这一事实, 得到
\begin{equation*}
H^{2d-i}(X,F^{\vee}(d))\xrightarrow{\sim}H^{i}_{c}(X,F)^{\vee}{\quoi},\tag{3.2.6.2}
\end{equation*}\etiquette[3.2.6.2]{XVIII.eq3.2.6.2}%
这就是 Poincaré 对偶的通常形式.

\begin{thebibliography}{99}\pageoriginale%

\sisi{}{%
\bibitem{sga4.exp.XVIII.bib.0} J. \textsc{Ayoub}, Les six opérations de Grothendieck et le
formalisme des cycles évanescents dans le monde motivique. I,
Astérisque 314 (2007).

\bibitem{sga4.exp.XVIII.bib.0bis} B. \textsc{Conrad}, A modern proof of Chevalley's theorem on
algebraic groups, Journal of the Ramanujan Mathematical Society 17 (2002),
p.~1--18.

\bibitem{sga4.exp.XVIII.bib.0ter} D. \textsc{Ferrand}, On the non additivity of the trace
in derived categories, \url{http://arxiv.org/abs/math/0506589}

}

\bibitem{sga4.exp.XVIII.bib.1} J. \textsc{Giraud}, Cohomologie non abélienne. Columbia
  University.

\bibitem{sga4.exp.XVIII.bib.2} A. \textsc{Grothendieck}, Le groupe de Brauer
  III. In: Dix exposés sur
la cohomologie des schémas. North Holl. Publ. Co.

\sisi{}%
{\bibitem{sga4.exp.XVIII.bib.2bis} A. \textsc{Grothendieck}, Technique de descente et théorème
d'existence en géométrie algébrique. VI. Les schémas de Picard : propriétés
générales, Séminaire Bourbaki 1961-1962, exposé nᵒ 236,
pages 221-243.}

\sisi{\bibitem{sga4.exp.XVIII.bib.3} L. \textsc{Illusie}, Thèse}%
{\bibitem{sga4.exp.XVIII.bib.3} L. \textsc{Illusie}, Complexe cotangent et déformations II.
Lecture Notes in Math. 283. Springer Verlag.}

\sisi{}{%
\bibitem{sga4.exp.XVIII.bib.3bis} F. \textsc{Knudsen}, D. \textsc{Mumford}, The
projectivity of the moduli space of stable curves. I : Preliminaries on
``det'' and ``Div'', Mathematica Scandinavica 39 (1976), p.~19--55.

\bibitem{sga4.exp.XVIII.bib.3ter} U. \textsc{Köpf},
Über eigentliche Familien algebraischer Varietäten
über affinoiden Räumen, Schriftenreihe des Mathematischen Instituts
der Universität Münster. 2. Serie, Heft 7 (1974).}

\bibitem{sga4.exp.XVIII.bib.4} D. \textsc{Mumford}, Geometric invariant
  theory. Springer-Verlag 1965.

\sisi{}{%
\bibitem{sga4.exp.XVIII.bib.4bis} F. \textsc{Orgogozo}, Isomotifs de dimension inférieure ou
égale à un, Manuscripta Mathematica 115 (2004), p.~339--360.}

\bibitem{sga4.exp.XVIII.bib.5} M. \textsc{Raynaud}, Faisceaux amples sur les schémas
  en groupes et les espaces homogènes. Lecture Notes in Math. 119. Springer Verlag.

\sisi{}{%
\bibitem{sga4.exp.XVIII.bib.5bis} M. \textsc{Raynaud}, Géométrie analytique rigide d'après
Tate, Kiehl. Mémoires de la Société Mathématique de France, tome 39--40
(1974), p.~319--327.}

\bibitem{sga4.exp.XVIII.bib.6} J.P. \textsc{Serre}, Groupes algébriques et corps de
classe. Publ. Inst. Math. Nancago. Hermann 1959.

\bibitem{sga4.exp.XVIII.bib.7} J.L. \textsc{Verdier}, Dualité dans la cohomologie des espaces
localement compacts. Sém. Bourbaki 300. Nov. 1965.

\bibitem{sga4.exp.XVIII.bib.8} J.L. \textsc{Verdier}, A duality theorem in the etale
  cohomology of schemes. Woods-hole seminar on alg. geometry. July 1964.

\bibitem{sga4.exp.XVIII.bib.9} Catégories dérivées, J.L. \textsc{Verdier}. IHÉS 油印讲义.

\end{thebibliography}

%2.17  - 106
%IX 3.1 - 138
%XVII  6.3.18.2 - 345
%1.2.0 - 855
%1.2.7.3 - 865
%1.2.2.11 - 892
%1.2.12 - 920
%6.3.3.1 - 1026
%1.3.12.4 - 1631, 1636, 1704
%1.3.2.7 - 1669
%1.2.9.7 - 1778
%1.2.16.2 - 1784
%2.12.3 - 4257
%3.1.9.10 -  5337
%3.2.5.1 - 5461

\ifx\danslelivre\undefined
\let\SGAXVIIIenddocument\enddocument
\else
\let\SGAXVIIIenddocument\relax
\fi
\SGAXVIIIenddocument



% END native input: chapters/exp-18-part5-zh.tex
% BEGIN native input: chapters/exp-19-zh.tex
\chapter{等特征优预概形的上同调}
%\addtocontents{toc}{M. Artin}
\begin{center}
M. Artin
\end{center}
\phantomsection\MakeLinkTarget*{sga4.exp.XIX}\label{sga4.exp.XIX}

% The edition is fixed to the corrected (second) reading. This expandable
% selector is output-equivalent and avoids current Xy-pic conditional failures.
\renewcommand{\sisi}[2]{#2}

\newcommand{\UU}{\mathbf{U}}

本\pageoriginale 讲讨论等特征（equal characteristic）优预概形（excellent prescheme）
（\EGA IV 7.8.1）的平展上同调，\emph{并假设这类预概形的
奇点消解\pageoriginale 定理成立}，具体形式如下：

设 $X=\Spec A$ 为等特征优局部概形，
$U\subset X$ 为正则开子集。则存在一个固有双有理态射
$f:X'\rightarrow X$ 及 $X$ 上的开浸入
$U\rightarrow X'$，使 $X'$ 正则，而
$Y=X'-U$ 处处为余维 1，且具有正规交叉。

\emph{因此，这些证明“目前”只在特征零时有效}
\cite{sga4.exp.XIX.bib.4}。也可以由这些证明推出
特征 $p>0$ 的低维预概形的一些结果
\cite{sga4.exp.XIX.bib.1}。

在奇点消解假设下，对于等特征预概形及与特征互素的系数，
我们将得到或多或少令人满意的结果（\cf \hyperref[XIX.thm3.2]{3.2}，\hyperref[XIX.thm4.1]{4.1}，
\hyperref[XIX.thm5.1]{5.1}，\hyperref[XIX.thm6.1]{6.1}）。另一方面，对于一般混合特征预概形的上同调，
几乎一无所知；即使在已有奇点消解的
2 维情形也是如此 \cite{sga4.exp.XIX.bib.2}。
例如，对任何维数 $>1$ 的混合特征完备环，
纯性定理 \hyperref[XIX.thm2.1]{2.1} 仍是未知的。
这一结果在形式幂级数环 $k[[x_{1},\ldots, x_{n}]]$ 上的情形
（\hyperref[XIX.thm1.2]{1.2}），是该理论的主要工具之一。
还应指出，\hyperref[XIX.thm1.2]{1.2} 的证明并不使用消解定理。
\emph{因此，此定理已在任意特征下得到证明。}

我们将经常使用 \EGA IV 7.8 中建立的优\pageoriginale 概形的性质，
而不明确指出，例如优性在有限型扩张下的稳定性
（\EGA IV 7.8.3(ii)），以及在严格局部化下的稳定性。
后一个事实由奇点消解及 \EGA IV 7.9.5 得到。

\section{环 $k[[x_{1},\ldots,x_{n}]]$ 的纯性}\etiquette[1]{XIX.sec1}

回顾 \hyperrefext[XVI][XVI.thm3.7]{3.7} 的以下推论：

\begin{enonce*}{推论 1.1}\etiquette[1.1]{XIX.cor1.1}
设 $A$ 为严格局部环，并用 $A\{x\}$ 表示 $A[x]$
在由 $\rad A$ 和 $x$ 生成的素理想处的严格局部化。
设 $U=\Spec A\{x\}[1/x]$。则对与 $A$ 的剩余特征互素的 $n$，有
$$
H^{q}(U,𝛍_{n})\simeq
\begin{cases}
𝛍_{n}(A) & \text{\rm 若}\quad q=0\\
\ZZ/n & \text{\rm 若}\quad q=1\\
0            & \text{\rm 若}\quad q>1{\quoi}.
\end{cases}
$$
\end{enonce*}

事实上，设 $X=\Spec A[x]$，并设 $Y=V(x)$ 为 $x$
在 $X$ 中的零点集。则 $(Y,X)$ 是
$\Spec A$ 上的一个光滑对（\hyperrefext[XVI][XVI.def3.1]{3.1}）。
由 (V 4, VIII 5)，$H^{q}(U,𝛍_{n})$
可与层 $ℋ^{q+1}_{Y}$ $(X,𝛍_{n})$
在位于 $X$ 的点 $(\rad A,x)$ 上方的一个几何点处的纤维等同，
因而此推论立即由（\hyperrefext[XVI][XVI.thm3.7]{3.7}）得到。
注意，$𝛍_{n}$ 和 $\ZZ/n$ 都是 $n$ 阶循环群。
此处使用它们是为了给出典范\pageoriginale 公式（\cf \hyperref[XIX.thm3.4]{3.4}）。

\begin{enonce*}{定理 1.2}\etiquette[1.2]{XIX.thm1.2}
设 $k$ 为可分闭域，$U=\Spec
k[[x_{1},\ldots,x_{r}]]$ $[1/x_{r}]$。则对与
$k$ 的特征互素的 $n$，有
$$
H^{q}(U,𝛍_{n})=
\begin{cases}
𝛍_{n}(A) &\text{\rm 若}\quad q=0\\
\ZZ/n &\text{\rm 若}\quad q=1\\
0 &\text{\rm 若}\quad q>1{\quoi}.
\end{cases}
$$
\end{enonce*}

\begin{enonce*}[remark]{注}
当 $k=\CC$ 时，这一证明同样适用于收敛幂级数环。
此外，只须将结果的陈述稍加斟酌，便容易去掉 $k$ 可分闭的假设。
不过，为了建立一般的纯性定理 \hyperref[XIX.thm3.2]{3.2}
（当然是在奇点消解假设下），我们只需要
$k$ 可分闭时的结果，而这些其他断言都是
定理 \hyperref[XIX.thm3.2]{3.2} 的推论。
\end{enonce*}

在维数为 0 和 1 时，可对每个正则严格局部环证明该结果：

\begin{enonce*}{引理 1.3}\etiquette[1.3]{XIX.lem1.3}
设 $A$ 为正则严格局部环，$x\in \rad A$ 为一个
局部参数。设 $U=\Spec A[1/x]$。则
$H^{0}(U,𝛍_{n})=𝛍_{n}(A)$，且有典范同构
$$
\ZZ/n\simeq H^{0}(U,\GG_{m})/(H^{0}(U,\GG_{m}))^{n}\simeq
H^{1}(U,𝛍_{n}){\quoi},
$$
其中\pageoriginale $\ZZ/n$ 的生成元为 $x\in H^{0}$ $(U,\GG_{m})$ 的“剩余类”。
\end{enonce*}

\medskip
\noindent
\emph{证明}。对于 $q=0$，断言就是 $U$ 连通且非空。
这确实成立。对于 $q=1$，应用 Kummer 正合列
（\hyperrefext[IX][IX.rems3.7]{3.7}）。由于 $A$ 正则且为析因环，
而 $U$ 是 $X=\Spec A$ 的开子集，有 $\Pic U=0$。
于是有正合列
$$
H^{0}(U,\GG_{m})\xrightarrow{n}H^{0}(U,\GG_{m})\rightarrow
H^{1}(U,𝛍_{n})\rightarrow 0{\quoi},
$$
还须证明 $H^{0}(U,\GG_{m})/(H^{0}(U,\GG_{m}))^{n}$
同构于 $\ZZ/n$，且由 $x$ 的剩余类生成。
然而，有显然的正合列
$$
\xymatrix@R=.1cm{
0 \ar[r]& H^0(X,\GG_m)\ar[r] & H^{0}(U,\GG_m)\ar[r] &
\ZZ \ar[r]& 0\\
& \text{\rotatebox{-90}{$\simeq$}} & \text{\rotatebox{-90}{$\simeq$}} &\\
& A^{\ast}  & A[1/x]^{\ast}   
}
$$
其中截面 $a\in H^{0}$ $(U,\GG_{m})$ 在
$\ZZ$ 中的像为 $a$ 沿 $\{x=0\}$ 的消失阶数。
群 $A^{\ast}$ 是 $n$-可除的，因为 $A$ 严格局部，
且 $n$ 与剩余特征互素。事实上，若
$u\in H^{0}$ $(X,\GG_{m})$，则方程 $Y^{n}-u=0$
定义 $A$ 的一个有限平展扩张，故该扩张完全分裂。
于是 $H^{1}$ $(U,𝛍_{n})\simeq \ZZ/n$，引理得证。

\medskip
\noindent
\emph{\hyperref[XIX.thm1.2]{1.2} 的证明}。还须处理
$q>1$ 的情形。回顾 (V，附录) 的术语。有
$$
H^{q}(U,𝛍_{n})=\varinjlim_{\UU}
H^{q}(\UU,𝛍_{n})
$$ 
其中\pageoriginale $\xymatrix{\UU = \bigg\{U_0 & U_1
\ar@<3pt>[l] \ar@<-3pt>[l] & \ldots \ar@<5.5pt>[l]
\ar[l]\ar@<-5.5pt>[l]\bigg\}}$ 遍历 $U$ 的\emph{超覆盖}（hypercovering）
所成的范畴（在同伦意义下滤过）。设 $\UU$ 为这样一个超覆盖，
并取各 $U_{i}$ 为分离且有限型的。
由于要证明 $H^{q}$ $(U,𝛍_{n})=0$，只须找出一个
超覆盖 $\VV$ 及态射 $\VV\rightarrow \UU$，使诱导态射
$$
H^{q}(\UU,𝛍_{n})\longrightarrow
H^{q}(\VV,𝛍_{n})
$$
为零态射。

设 $\{W_{1},\ldots,W_{m}\}$ 为 $i\leq q+1$ 时所有
$U_{i}$ 的连通分支所成的集合，并设 $B_{j}$ 为
$A=k[[x_{1},\ldots,x_{r}]]$ 在 $W_{j}$ 的有理函数域
$R(W_{j})$ 中的正规化。则由于 $W_{j}\rightarrow U$ 平展，
概形 $W_{j}$ 是 $\Spec B_{j}$ 的一个开子集，
可设它为 $\Spec C_{j}$ 的开补集，其中 $C_{j}$
为 $B_{j}$ 的适当的约化商环。

另取 $B_{j}$ 的元素 $b_{j}$，使 $b_{j}$ 生成
$k((x_{1},\ldots,x_{r}))$ 的可分扩张 $R(W_{j})$。
设 $f_{j}(Y)\in A[Y]$ 为 $b_{j}$ 在 $A$ 上的
首一不可约多项式，则 $B_{j}$ 同构于环 $A[Y]/(f_{j})$ 的正规化。
最后，设 $d_{j}\in A$ 为 $f_{j}(Y)$ 的判别式，
它是 $A$ 的一个非零元素。

现在可以将 $i<r$ 的各 $x_{i}$ 换为 $A$ 的其他元素，
唯一的条件是 $\{x_{1},\ldots,x_{r}\}$ 构成 $A$ 的一个参数系。

容易看出，必要时更换这些元素后，
可\pageoriginale 假设满足以下条件：

\setcounter{subsection}{3}
\subsection{}\etiquette[1.4]{XIX.subsec1.4}
\begin{itemize}
\item [(i)] 对每个使 $d_{j}$ 不是单位的 $j$，
$x_{1},\ldots,x_{r-1}$ 与 $d_{j}$ 生成一个根为
$\rad A$ 的准素理想。

\item[(ii)] 对每个 $j$，$x_{1},\ldots,x_{r-1}$ 在
$C_{j}$ 中的像生成一个根为 $\rad C_{j}$ 的准素理想。
\end{itemize}
于是有以下引理：

\begin{enonce*}{引理 1.5}\etiquette[1.5]{XIX.lem1.5}
设 $A^{0}=k[[x_{1},\ldots,x_{r-1}]]$ $\{x_{r}\}$ 为
$k[[x_{1},\ldots,x_{r-1}]]$ $[x_{r}]$ 在点
$(x_{1},\ldots,x_{r})$ 处的 Hensel 化（henselization），
并设 $X^{0}=\Spec A^{0}$，$U^{0}=\Spec
A^{0}$ $[1/x_{r}]$。在上述 $\hyperref[XIX.subsec1.4]{1.4}$ 的条件下，
存在 $U^{0}$ 的一个（$q+1$ 阶）超覆盖 $\UU^{0}$，以及
同构 $\fprod{\UU^{0}}{X}{X^{0}}\simeq\UU$。
\end{enonce*}

先承认此引理。环 $A^{0}$ 是优环（\EGA IV 7.8.3
(ii), 7.9.5），因此 $U^{0}_{i}$ 和 $U_{i}$ 的连通分支
双射对应（\EGA IV 7.8.3 (vii)）。同构
$\fprod{\UU^{0}}{X}{X^{0}}\simeq
\UU$ 因而给出一个同构
$$
H^{q}(\UU^{0},𝛍_{n})\simeq
H^{q}(\UU,𝛍_{n}){\quoi}.
$$
由于 $k$ 可分闭，$A^{0}$ 为严格局部环，
可将 \hyperref[XIX.cor1.1]{1.1} 应用于环 $A^{0}$。
于是 $H^{q}$ $(U^{0},𝛍_{n})=0$。然而
$$
H^{q}(U^{0},𝛍_{n})\simeq
\varinjlim_{\VV^{0}}H^{q}(\VV^{0},𝛍_{n}){\quoi},
$$
其中 $\VV^{0}$ 遍历 $U^{0}$ 的 $q+1$ 阶超覆盖所成的范畴，\pageoriginale 且
$H^{q}$ $(\UU^{0},𝛍_{n})$ 为有限群。
因此存在超覆盖态射 $\VV^{0}\rightarrow \UU^{0}$，使诱导态射
$$
H^{q}(\UU^{0},𝛍_{n})\longrightarrow
H^{q}(\VV^{0},𝛍_{n})
$$
为零态射。置 $\VV=\fprod{\VV^{0}}{X}{X^{0}}$。
则 $H^{q}$ $(\VV^{0},𝛍_{n})\simeq
H^{q}(\VV,𝛍_{n})$（用与 $\UU^{0}$ 相同的论证），所以
$$
H^{q}(\UU,𝛍_{n})\longrightarrow
H^{q}(\VV,𝛍_{n})
$$
也为零态射，\cqfd

\medskip
\noindent
\emph{\hyperref[XIX.lem1.5]{1.5} 的证明}。只须找出有限
$A^{0}$-代数 $B^{0}_{j}$ 及 $C^{0}_{j}$，以及 $A$-同构
\begin{align*}
B^{0}_{j}\otimes_{A^{0}}A &\xrightarrow{\sim} B_{j}\\
C^{0}_{j}\otimes_{A^{0}}A &\xrightarrow{\sim} C_{j}{\quoi}.
\end{align*}
事实上，由（\cite{sga4.exp.XIX.bib.3}~1.4），
态射 $B_{j}\rightarrow C_{j}$ 由一个态射
$B^{0}_{j}\rightarrow C^{0}_{j}$ 诱导，故开子集 $W_{j}\subset
\Spec B_{j}$ 由某个开子集 $W^{0}_{j}\subset \Spec
B^{0}_{j}$ 诱导。此外，各 $W_{j}$ 之间的 $U$-态射
由各 $A$-代数 $B_{j}$ 之间的态射诱导，从而
由各 $B^{0}_{j}$ 之间的态射诱导
（\cite{sga4.exp.XIX.bib.3}~1.4），
因而也由各 $W^{0}_{j}$ 之间的 $U^{0}$-态射诱导。
所以，这些 $W^{0}_{j}$ 构成 $U^{0}$ 上的一个单纯形对象
$\UU^{0}$ 的连通分支集，且它诱导出 $\UU$；
立即可见，它是 $U^{0}$ 的一个超覆盖。

由 \hyperref[XIX.subsec1.4]{1.4} (ii) 确实可知 $C_{j}$ 由
一个 $A^{0}$-代数\pageoriginale $C^{0}_{j}$ 诱导，这很容易验证
\footnote[1]{事实上，由 \hyperref[XIX.subsec1.4]{1.4} (ii)，$C_{j}$ 在
$k[[X_{1},\ldots,X_{r-1}]]$ 上有限，因而在 $A^{0}$ 上更是有限，
只须取 $C^{0}_{j}=C_{j}$，赋予由限制标量得到的
$A^{0}$-模结构。}。对于 $B_{j}$，只须找出
首一多项式 $f^{0}_{j}(Y)\in A^{0}[Y]$，使两个
$A$-代数
$$
A[Y]/(f_{j})\quav
A[Y]/(f^{0}_{j})=A\otimes_{A^{0}}A^{0}[Y]/(f^{0}_{j})
$$
同构。事实上，此时 $A^{0}[Y]/(f^{0}_{j})$ 约化，
取 $B^{0}_{j}$ 为 $A^{0}[Y]/(f^{0}_{j})$ 的正规化
（\EGA IV 7.8.3 (vi)）。环 $B^{0}_{j}\otimes_{A^{0}}A$
仍然正规（\EGA IV 7.8.3 (vii)），并且与 $B_{j}$ 双有理，
因而等于该环。

由 \hyperref[XIX.subsec1.4]{1.4} (i)，由 $d_{j}$ 生成的理想 $I_{j}\subset A$
是由 $k[[x_{1},\ldots,x_{r-1}]]$ $[x_{r}]$ 的一个理想诱导的
（Weierstrass 预备定理 \cite{sga4.exp.XIX.bib.5}\cite{sga4.exp.XIX.bib.6}），
从而由 $A^{0}$ 的一个理想 $I^{0}_{j}$ 诱导。
此外，$A/I_{j}$ 是一个有限 $k[[x_{1},\ldots,x_{r-1}]]$-代数，
因而同构于 $A^{0}/I^{0}_{j}$。故后一环完备，
从而 $A^{0}$ 的 $I^{0}_{j}$-进完备化是完备的，且同构于 $A$。
于是可找到一个首一多项式
$f^{0}_{j}\in A^{0}$ $[Y]$，使在 $A$ $[Y]$ 中有
$$
f_{j}\equiv f^{0}_{j}(\text{\rm 模 } d_{j}{}^{2}\rad A){\quoi}.
$$
由（\cite{sga4.exp.XIX.bib.3}，1.3），$A[Y]/(f_{j})\simeq A[Y]/(f^{0}_{j})$，
引理得证。

\section{严格局部环的情形}\etiquette[2]{XIX.sec2}

定理如下：

\begin{enonce*}{定理 2.1}\etiquette[2.1]{XIX.thm2.1}
\footnote[1]{依赖于奇点消解（\cf 导言）。}。\pageoriginale 设 $A$ 为等特征的优、正则、
严格局部环，$x$ 为 $A$ 的一个“局部参数”
\footnote[2]{此处指：属于某个\emph{正则}参数系的元素。}。
设 $U=\Spec A[1/x]$。则对与 $A$ 的特征互素的 $n$，有
$$
H^{q}(U,𝛍_{n})=
\begin{cases}
𝛍_{n}(A) &\text{\rm 若}\quad q=0\\
\ZZ/n &\text{\rm 若}\quad q=1\\
0 &\text{\rm 若}\quad q>1{\quoi}.
\end{cases}
$$
\end{enonce*}

证明对 $q$ 作归纳。用 $P(N)$ 表示
\hyperref[XIX.thm2.1]{2.1} 对各 $q\leq N$ 的断言。
由 \hyperref[XIX.lem1.3]{1.3}，断言 $P(1)$ 成立。

\begin{enonce*}{引理 2.2}\etiquette[2.2]{XIX.lem2.2}
假设 $P(N)$ 成立。设 $A\rightarrow A'$ 为等特征的优、正则、
严格局部环之间的一个局部态射。
设 $\{x_{1},\ldots,x_{r}\}$ 为 $A$ 的某个正则参数系的一部分，
且它在 $A'$ 中的像集仍为某个正则参数系的一部分。设
\begin{align*}
U &= \Spec A[1/x_{1},\ldots,1/x_{r}]\\
U'&= \Spec A'[1/x'_{1},\ldots,1/x'_{r}]{\quoi}.
\end{align*}
则典范态射
$$
H^{q}(U,𝛍_{n})\longrightarrow
H^{q}(U',𝛍_{n})
$$
对每个 $q\leq N$ 均为\pageoriginale 双射。
\end{enonce*}

\medskip
\noindent
\emph{证明}。\footnote[3]{有一个比下述证明更短、但使用杯积形式体系的证明：
在假设 $P(N)$ 下证明，上同调
$H^{\ast}(U,\ZZ/n\ZZ)$ 在维数 $\leq N$ 内为
$H^{1}(U,\ZZ/n\ZZ)\simeq (\ZZ/n\ZZ)^{r}$
上的（$\ZZ/n\ZZ)$ 上的）外代数，
从而只须证明 $q=1$ 时的 \hyperref[XIX.lem2.2]{2.2}，而这是熟知的情形。
例如，\SGA 7 中明确给出了在假设 $P(N)$ 下的上述计算。}
对 $r$ 作归纳：$r=0$ 时断言平凡。
设 $r>0$，并假设定理对 $r-1$ 已获证明。
置 $B=A/(x_{r})$，$V=\Spec A[1/x_{1},\ldots,1/x_{r-1}]$，
$Y=\Spec B[1/x_{1},\ldots,1/x_{r-1}]$
（$x_{i}$ 在 $B$ 中的像仍用同一符号表示）。
有开浸入 $i:U\rightarrow V$，而 $Y$ 为由元素
$x_{r}$ 定义的互补闭子集。

将 $P(N)$ 应用于 $V$ 在 $Y$ 的各几何点处的严格局部化，得到
\begin{equation*}
R^{q}i_{\ast}(𝛍_{n})_{U}=
\begin{cases}
(𝛍_{n})_{V} &\text{\rm 若}\quad q=0\\
(\ZZ/n)_{Y} &\text{\rm 若}\quad q=1\\
0 &\text{\rm 若}\quad 1<q\leq N{\quoi},
\end{cases}\tag{2.3}
\end{equation*}\etiquette[2.3]{XIX.eq2.3}
因为这些严格局部化是等特征优正则环，
而元素 $x_{r}$ 在 $Y$ 的每一点都是局部参数。

用“撇号”表示由 $A'$ 得到的类似对象。得到一个谱序列态射
$$
\xymatrix@C=.4cm{
E^{pq}_2 \ar[d]^{\varphi^{pq}_2} & = \text{\phantom{ab}} H^{p}(V,R^{q}i_{\ast}
𝛍_n)\ar[d]^{\varphi} \ar@{=>}[r] & H^{p+q}(U,𝛍_n) \ar[d]^{\varphi}&\\
E'^{pq}_2 &= \text{\phantom{ab}} H^{p}(V',R^{q}i'_{\ast}𝛍_n) \ar@{=>}[r]& H^{p+q}(U',𝛍_n).
}
$$
由 \hyperref[XIX.lem1.3]{1.3} 及 \eqhyperref[XIX.eq2.3]{2.3} 显然\pageoriginale 可知，
基变换 $V\xleftarrow{g}V'$ 对 $q\leq N$ 诱导同构
$$
g^{\ast}R^{q}i_{\ast}𝛍_{n}\xrightarrow{\sim}
R^{q}i'_{\ast}𝛍_{n}
$$
（只有在 $q=0$、$1$ 时才需要验证）。
因此，将归纳假设应用于态射
$$
H^{q}(V,𝛍_{n})\longrightarrow
H^{q}(V',𝛍_{n}){\quoi},
$$
可知当 $q=0$ 且 $p\leq N$ 时，$\varphi^{pq}_{2}$ 为同构。

同样，可将归纳假设应用于环态射 $B\rightarrow B'$；
这两个环都是等特征的正则、优、严格局部环，
而元素 $x_{1},\ldots,x_{r-1}$ 构成某个参数系的一部分。
于是
$$
H^{q}(Y,𝛍_{n})\xrightarrow{\sim} H^{q}(Y',
𝛍_{n})\quad\text{\rm 若}\quad q\leq N{\quoi}.
$$
由于 $(𝛍_{n})\simeq (\ZZ/n)_{Y}$，可知当 $p\leq N$ 且 $q=1$ 时，
$\varphi^{pq}_{2}$ 也为双射。

当 $1<q\leq N$ 时，
$R^{q}i_{\ast}𝛍_{n}=R^{q}i'_{\ast}𝛍_{n}=0$，
故当 $1<q\leq N$ 时有 $E^{pq}_{2}={E'}^{pq}_{2}=0$；
因此对 $1\leq m\leq N$，谱序列态射归结为正合列态射
$$
\xymatrix@C=.4cm{
H^{m-2}(Y,\ZZ/n) \ar[r] \ar[d]^{a}& H^{m}(V,𝛍_n) \ar[r]\ar[d]^{b}&
H^m(U,𝛍_n) \ar[r]^(.45){e} \ar[d]^{c}&
H^{m-1}(Y,\ZZ/n)  \ar[r] \ar[d]^{d}&
H^{m+1}(V,𝛍_n) \ar[d]\\
H^{m-2}(Y',\ZZ/n) \ar[r] & H^{m}(V',𝛍_n) \ar[r]&
H^m(U',𝛍_n) \ar[r]^(.45){e} & H^{m-1}(Y',\ZZ/n)  \ar[r]&
H^{m+1}(V',𝛍_n).
}
$$
已经\pageoriginale 看到 $a$、$b$ 及 $d$ 为双射，
所以 $c$ 为\emph{单射}。为证明 $c$ \emph{满射}，
只须证明当 $m\leq N$ 时，态射
$$
H^{m}(U,𝛍_{n})\xrightarrow{e}H^{m-1}(Y,\ZZ/n)
$$
为满射。

% SGA-SRC-001196
设 $k\subset A$ 为任意可分闭子域，并设
$A^{0}=k\{x_{1},\ldots, x_{r}\}$ 为
$k[x_{1},\ldots,x_{r}]$ 在原点的 Hensel 化。
则态射 $A^{0}\rightarrow A$ 及元素 $x_{i}$ 满足
\hyperref[XIX.thm2.1]{2.1} 的假设。采用显然的记号，有交换方块
$$
\xymatrix@R=1.25cm{
H^{m}(U^{0},𝛍_n)  \ar[r]^(.45){e^0} \ar[d]^{c^0}& H^{m-1}(Y^0,
\ZZ/n) \ar[d]^{d^{0}}\\
H^{m}(U,𝛍_n)  \ar[r]^(.45){e} & H^{m-1}(Y,
\ZZ/n)
}
$$
且已知箭头 $d^{0}$ 为双射。因此，只须证明态射
$e^{0}$ 满射。于是，略去上标 ${}^{0}$，
只须验证 $A=k\{x_{1},\ldots,x_{r}\}$ 时的满射性。

对于这个环，它是 $k$ 上光滑概形的极限，
我们已在 \hyperrefext[XVI][XVI.thm3.7]{3.7} 中证明了纯性定理。
故\pageoriginale 对每个 $q>1$ 有 $R^{q}i_{\ast}𝛍_{n}=0$，
从而对每个 $m\geq 1$ 有正合列
$$
H^{m}(U,𝛍_{n})\xrightarrow{e}H^{m-1}(Y,\ZZ/n)\rightarrow
H^{m+1}(V,𝛍_{n})\xrightarrow{f}H^{m+1}(U,𝛍_{n})
$$
。因此，$e$ 的满射性等价于 $f$ 的单射性。

设 $Z=\Spec k\{x_{1},\ldots,x_{r-1}\}$
$[1/x_{1},\ldots,1/x_{r-1}]$。由环的包含给出的态射
$\Spec k\{x_{1},\ldots,x_{r}\}\rightarrow \Spec k\{x_{1},\ldots,x_{r-1}\}$
是无环的。这就是光滑态射的局部无环性
（\hyperrefext[XVI][XVI.thm1.1]{1.1}）。
由于 $\Spec k\{x_{1},\ldots, x_{r}\}$ 的开子集 V 是
$Z$ 的逆像，态射
$$
H^{q}(Z,𝛍_{n})\longrightarrow
H^{q}(V,𝛍_{n})
$$
对每个 $q$ 均为双射（\hyperrefext[XV][XV.def1.3]{1.3}）。
因此，只须证明态射
$$
H^{q}(Z,𝛍_{n})\longrightarrow
H^{q}(U,𝛍_{n})
$$
对每个 $q$ 均为单射。

环 $k\{x_{1},\ldots,x_{r}\}$ 是各环
$A_{\alpha}$ 的归纳极限，这些环在
$k\{x_{1},\ldots,x_{r-1}\}$ $[x_{r}]$ 上平展，
且各自赋有点 $x_{1}=x_{2}=\ldots=x_{r}=0$ 到该环的一个提升。
因此，$U$ 是各开子集
$U_{\alpha}=\Spec A_{\alpha}[1/x_{1},\ldots,1/x_{r}]$ 的投射极限，且有
$$
H^{q}(U,𝛍_{n})=\varinjlim_{\alpha}H^{q}(U_{\alpha},
𝛍_{n}){\quoi}.
$$
% SGA-SRC-001198
由于 $A_{\alpha}[1/x_{r}]$ 在 Hensel 环
$k\{x_{1},\ldots,x_{r-1}\}$ 上光滑，\pageoriginale 且
$A_{\alpha}[1/x_{r}]$ 的闭纤维非空，立即可知
$k\{x_{1},\ldots, x_{r-1}\}$-代数 $A_{\alpha}[1/x_{r}]$ 有截面。
因而 $U_{\alpha}$ 在 $Z$ 上有截面，于是
$$
H^{q}(Z,𝛍_{n})\longrightarrow
H^{q}(U_{\alpha},𝛍_{n})
$$
对每个 $q$ 均为单射，因而
$$
H^{q}(Z,𝛍_{n})\longrightarrow
H^{q}(U,𝛍_{n})
$$
也为单射，引理得证。

\hypertarget{sga4.exp.XIX.lemma.2.4}{}\begin{enonce*}{引理 2.4}\etiquette[2.4]{XIX.lem2.4}
假设 $P(N)$ 成立。设 $A\rightarrow A'$ 为等特征的优严格局部环
（不必正则）之间的一个正则态射\emph{（\EGA IV 6.8.1）}，
$U\subset \Spec A$ 为正则开子集。设
$U'=\fprod{U}{\Spec A'}{\Spec A}$。则对与 $A$ 的剩余特征互素的 $n$，
态射
$$
H^{q}(U,𝛍_{n})\longrightarrow
H^{q}(U',𝛍_{n})
$$
在 $q\leq N$ 时为双射。
\end{enonce*}

\medskip
\noindent
\emph{证明。}设 $f:X\rightarrow \Spec A$ 为 $\Spec A$ 的一个奇点消解，
使 $\Spec A$ 的开子集 $U$ 提升为开浸入
$i:U\rightarrow X$，且闭子集 $Y=X-U$ 纯为余维 1 并具有正规交叉。
用撇号表示由 $A\rightarrow A'$ 诱导的基变换。
态射 $g:X'\rightarrow X$ 正则（\EGA IV 6.8.3），故
$f':X'\rightarrow \Spec A'$ 为 $\Spec A'$ 的一个消解，
$i':U'\rightarrow X'$ 为开浸入，且\pageoriginale
$Y'=X'-U'$ 纯为余维 1 并具有正规交叉。
这些概形都是优概形。

设 $p'$ 为 $X'$ 的一点，$p$ 为其在 $X$ 中的像。
由于 $Y$ 在 $X$ 中具有正规交叉，存在元素
$x_{1},\ldots,x_{r}\in O_{X,p}$，使 $\{x_{1},\ldots,x_{r}\}$
为 $O_{X,p}$ 的一个参数系的一部分，且 $\Spec
O_{X,p}\cap Y=V(x_{1}\cdots x_{r})$。
此外，由 $g:X'\rightarrow X$ 正则立即可知，
各 $x_{i}$ 在 $O_{X',p'}$ 中的像 $x'_{i}$ 构成该环的
一个参数系的一部分，且 $\Spec O_{X',p'}\cap
Y'=V(x'_{1}\cdots x'_{r})$。

设 $B'$ 为 $O_{X',p'}$ 在 $p'$ 上方的一个几何点
$\bar{p}$ 处的严格局部化，$B$ 为
$O_{X,p}$ 在 $X$ 的相应几何点处的严格局部化
（这些环是优的（\EGA IV 7.8.3 (ii), 7.9.5））。设
$$
U_{\bar{p}}=\Spec B[1/x_{1},\ldots,1/x_{r}]\quav U'_{\bar{p}}=\Spec
B'[1/x'_{1},\ldots,1/x'_{r}]{\quoi}.
$$
引理 \hyperref[XIX.lem2.2]{2.2} 适用，从而态射
\begin{equation*}
H^{q}(U_{\bar{p}},𝛍_{n})\longrightarrow
H^{q}(U'_{\bar{p}},𝛍_{n})\tag{$\ast$}
\end{equation*}
在 $q\leq N$ 时为双射。

计算开浸入 $i:U\rightarrow X$ 及 $i':U'\rightarrow X'$ 的层
$R^{q}i_{\ast}𝛍_{n}$ 及 $R^{q}i'_{\ast}𝛍_{n}$：
在几何点 $\bar{p}$ 处的茎
$(R^{q}i'_{\ast}𝛍_{n})_{\bar{p}}$ 恰为
$H^{q}(U'_{\bar{p}},𝛍_{n})$（VIII 5）。同样，
$(R^{q}i_{\ast}𝛍_{n})_{\bar{p}}=H^{q}(U_{\bar{p}},𝛍_{n})$。
因此，$(\ast)$ 表明基变换态射
\begin{equation*}
g^{\ast}R^{q}i_{\ast}𝛍_{n}\longrightarrow R^{q}i'_{\ast}𝛍_{n}\tag{${\ast}{\ast}$}
\end{equation*}
在 $q\leq N$ 时\pageoriginale 为双射。

由固有态射的基变换定理
（\hyperrefext[XII][XII.cor5.5]{5.5}），对每个挠层 $F$
及每个 $q$，有（因为 $X$ 和 $X'$ 在严格局部环的谱上固有，
且 $X$ 和 $X'$ 的闭纤维相等）
$$
H^{q}(X,F)\xrightarrow{\sim} H^{q}(X',g^{\ast}F)
$$
。因此，由 $({\ast}{\ast})$，谱序列态射
$$
\xymatrix@C=.5cm{
E^{pq}_2 \ar[d]& =& H^{p}(X,R^{q}i_{\ast}𝛍_n) \ar@{=>}[r]\ar[d]^{\varphi^{pq}_2}&
H^{p+q}(U,𝛍_n) \ar[d]^{\varphi}\\
E'^{pq}_2   &=&   H^{p}(X',R^{q}i'_{\ast}𝛍_n)\ar@{=>}[r]
& H^{p+q}(U',𝛍_n),
}
$$
中的箭头 $\varphi^{pq}_{2}$ 对每个 $q\leq N$ 均为双射。
容易看出，这意味着当 $q\leq N$ 且 $p+q\leq N+1$ 时，
$\varphi^{pq}_{\infty}$ 为双射。因此，收敛项之间的态射
$$
H^{m}(U,𝛍_{n})\longrightarrow
H^{m}(U',𝛍_{n})
$$
在 $m\leq N$ 时为双射，引理得证。

以下引理将完成 \hyperref[XIX.thm2.1]{2.1} 的证明：

\hypertarget{sga4.exp.XIX.lemma.2.5}{}\begin{enonce*}{引理 2.5}\etiquette[2.5]{XIX.lem2.5}
$P(N)\Rightarrow P(N+1)$。
\end{enonce*}

\medskip
\noindent
\emph{证明}。设 $A$ 为等\pageoriginale 特征的严格局部、正则、优环，
$x\in \rad A$ 为局部参数，$U=\Spec A[1/x]$。
设 $\hat{A}$ 为 $A$ 的完备化，$\hat{x}$ 为 $x$ 的像，
$\hat{U}=\Spec \hat{A}[1/x]$，$g:\hat{U}\rightarrow U$ 为典范态射。
只须证明
\begin{equation*}
\begin{split}
g_{\ast}(𝛍_{n})_{\hat{U}}\simeq
(𝛍_{n})_{U}\quad\text{\rm 且}\\
R^{q}g_{\ast}𝛍_{n}=0\quad\text{\rm 若}\quad 1\leq
q\leq N{\quoi}.
\end{split}\tag{2.6}
\end{equation*}\etiquette[2.6]{XIX.eq2.6}
事实上，Leray 谱序列
$$
E^{pq}_{2}=H^{p}(U,R^{q}g_{\ast}𝛍_{n})\Longrightarrow
H^{p+q}(\hat{U},𝛍_{n})
$$
将在 $q\leq N$ 时给出同构
$$
H^{q}(U,𝛍_{n})\xrightarrow{\sim}
H^{q}(\hat{U},𝛍_{n})
$$
以及正合列
$$
0\longrightarrow H^{N+1}(U,𝛍_{n})\longrightarrow
H^{N+1}(\hat{U},𝛍_{n})\longrightarrow
H^{0}(U,R^{N+1}g_{\ast}𝛍_{n}){\quoi}.
$$
然而，已经知道 $P(1)$ 成立（\hyperref[XIX.lem1.3]{1.3}），
因此只须证明当 $1<q\leq N+1$ 时 $H^{q}(U,𝛍_{n})=0$。
于是可将 $A$ 换为 $\hat{A}$，后者同构于形式幂级数环
$\hat{A}\simeq k[[x_{1},\ldots,x_{r}]]$，其中 $k$ 可分闭，
且 $\hat{x}=x_{r}$。这一情形已在 \hyperref[XIX.thm1.2]{1.2} 中处理。

现在验证 \eqhyperref[XIX.eq2.6]{2.6}：
$g_{\ast}(𝛍_{n})_{\hat{U}}\simeq
(𝛍_{n})_{U}$ 这一事实，正是 \eqhyperref[XIX.eq2.3]{2.3}
在 $q=0$ 时的断言。回顾 $R^{q}g_{\ast}𝛍_{n}$ 为
预层 $\Rr^{q}$ 的伴随层，其中
$$
\Rr^{q}(V)=H^{q}(\hat{V},𝛍_{n})
$$
对\pageoriginale 平展的 $V\rightarrow U$ 成立。
只须取 $V$ 分离且有限型；此时 $V$ 是若干整概形的不交并，
而只须考虑这样的 $V$。设 $B$ 为 $A$ 在
$\Gamma(V,O_{V})$ 的分式域中的正规化。
则 $B$ 为有限 $A$-代数（\EGA IV 7.8.3 (vi)），
立即可见 V 为 $\Spec B$ 的正则开子集，
而 $\hat{V}$ 为 $\Spec \hat{B}$ 的相应开子集，
其中 $\hat{B}=B\otimes_{A}\hat{A}$。
于是开子集 $V\subset \Spec B$ 及态射
$B\rightarrow \hat{B}$ 满足 \hyperref[XIX.lem2.4]{2.4} 的假设。故
$$
\Rr^{q}(V)\simeq H^{q}(V, 𝛍_{n})\simeq
H^{q}(\hat{V}, 𝛍_{n})
$$
在 $q\leq N$ 时成立。但预层
$V\mapsto H^{q}(V,𝛍_{n})$ 的伴随层显然为零，引理得证。

\section{纯性}\etiquette[3]{XIX.sec3}

采用与 \hyperrefext[XVI][XVI.sec3]{3} 类似的术语：

\hypertarget{sga4.exp.XIX.definition.3.1}{}\begin{enonce*}{定义 3.1}\etiquette[3.1]{XIX.def3.1}
一个\emph{余维 $c$ 的正则对 $(Y,X)$}（regular pair），是指局部 Noether 正则概形之间的闭浸入
$i:Y\rightarrow X$，使对每个 $y\in Y$ 均有 $\codim_{y}(Y,X)=c$
（此时存在 $y$ 在 $X$ 中的一个开邻域 $X'$，以及生成
$Y$ 的理想的截面 $x_{1},\ldots,x_{c}\in \Gamma(X',O_{X})$，
使它们在 $Y\cap X'$ 的每一点均构成 $X$ 的一个正则参数系的一部分）。
记 $U=X-Y$，并记 $j:U\rightarrow X$ 为相应的开浸入。
\end{enonce*}

余维 $c$ 的正则对的\emph{态射}\pageoriginale
$(Y',X')\rightarrow (Y,X)$，是一个使 $Y'=\fprod{Y}{X'}{X}$ 的态射
$X'\rightarrow X$。正则三元组、正则三元组的态射等概念，
也以类似的显然方式定义。

\hypertarget{sga4.exp.XIX.theorem.3.2}{}\begin{enonce*}{定理 3.2}\etiquette[3.2]{XIX.thm3.2}
\footnote[1]{依赖于奇点消解（\cf 导言）。}。
\emph{（上同调纯性）}设 $(Y,X)$ 为余维 $c>0$ 的正则对，
其中 $X$ 为等特征优预概形。设 $F$ 为 $X$ 上的局部常值层，
其值为阶数 $n$ 与 $X$ 的特征互素的循环群。
则采用 \emph{(V 4, VIII 5)} 的记号，有
\begin{align*}
ℋ^{q}_{Y} (X,F) &= (R^{q}i^{!})F=0\quad\text{\rm 若}\quad q\neq
2c{\quoi},\\
(R^{q}j_{\ast})j^{\ast}F &= 0\quad\text{\rm 若}\quad q\neq 0,
2c-1,\quad \text{且}\\
ℋ^{2c}_{Y}(X,F) &= i^{\ast}(R^{2c-1}j_{\ast})j^{\ast}F
\end{align*}
\noindent
为 $Y$ 上的 $n$ 阶循环群的局部常值层。
\end{enonce*}

\medskip
\noindent
\emph{证明}。（\cf \hyperrefext[XVI][XVI.thm3.7]{3.7}）。立即可见
$$
ℋ^{q}_{Y}(X,F)=0\quad\text{\rm 若}\quad q=0,1{\quoi}.
$$
由 (V \hyperref[V.cor4.5]{4.5})，
$$
ℋ^{q}_{Y}(X,F)=(R^{q}i^{!})F\simeq
i^{\ast}(R^{q-1}j_{\ast})j^{\ast}F\quad\text{\rm 若}\quad q>1{\quoi},
$$
故关于这三个层的断言等价。由于这些断言对于平展拓扑在
$X$ 上是局部的，可\pageoriginale 假设 $F=𝛍_{n}$。

当 $c=1$ 时，计算 $(R^{q}j_{\ast})j^{\ast}𝛍_{n}$ 的茎
（VIII 5），便归结为 \hyperref[XIX.thm2.1]{2.1} 中所考虑的情形，从而
$$
((R^{q}j_{\ast})j^{\ast}𝛍_{n})_{\bar{y}}=
\begin{cases}
\ZZ/n &\text{\rm 若}\quad q=1\\
0 &\text{\rm 若}\quad q>1
\end{cases}
$$
对 $Y$ 的每个几何点 $\bar{y}$ 成立。
无论如何，这些层在 $Y$ 外均为零，因此这证明了 $q>1$ 时的断言。
对 $q=1$ $(=2c-1)$，应用
（\hyperrefext[IX][IX.prop2.13]{2.13} (i)）的判别法及
\hyperref[XIX.lem1.3]{1.3} 的显式计算。

$c>1$ 的情形用归纳法处理（\cf
\hyperrefext[XVI][XVI.thm3.7]{3.7}）：设 $(Y,X)$ 为余维 $c>1$ 的正则对。
由定义显然可知，在 $X$ 上局部地可找到正则三元组
$(Y,Z,X)$，其中 $(Z,X)$ 的余维为 1，而 $(Y,Z)$ 的余维为 $c-1$。设
$$
\xymatrix{
  Y \ar[rr]^{u} \ar[rd]_{i}   &   & Z\ar[ld]^{v}\\
      & X & 
}
$$
为这些闭浸入。有谱序列
$$
E^{pq}_{2}=(R^{p}u^{!})(R^{q}v^{!})F\Longrightarrow
(R^{p+q}i^{!})F{\quoi}.
$$
由归纳假设，$q\neq 2$ 时 $(R^{q}v^{!})F=0$，
而 $q=2$ 时它局部常值，等等。
于是再次将归纳假设应用于态射 $u$，
当 $(p,q)\neq (2c-2,2)$ 时\pageoriginale
有 $(R^{p}u^{!})(R^{2}v^{!})F=0$，
而当 $(p,q)=(2c-2,2)$ 时它局部常值，等等；
由此便得关于收敛项 $(R^{p+q}i^{!})F$ 的结果，\cqfd

\setcounter{subsection}{2}
\subsection{}\etiquette[3.3]{XIX.subsec3.3}
还须在\emph{典范}同构意义下确定
（\hyperref[XIX.thm3.2]{3.2}）中的局部常值层。
对每个预概形 $S$，引入局部常值层
$(𝛍_{n}^{\otimes r})_{S}=$ 层
$(𝛍_{n})_{S}$ 在常值环层 $\ZZ_{S}$ 上的 $r$ 重张量积。
另置 $(𝛍_{n}^{\otimes 0})_{S}=(\ZZ/n)_{S}$。
立即可见，若 $n$ 在 $S$ 上可逆，
则 $(𝛍_{n}^{\otimes r})_{S}$ 为 $n$ 阶循环群的局部常值层。
有形如 $(𝛍_{n}^{\otimes r})\otimes
(𝛍_{n}^{\otimes
s})=(𝛍_{n}^{\otimes r+s})$ 的公式，且若 $S'\rightarrow
S$，则 $(𝛍_{n}^{\otimes r})_{S'}$ 典范同构于
$(𝛍_{n}^{\otimes r})_{S}$ 在 $S'$ 上的逆像。
在不会引起混淆时，常略去符号 $S$。

设 $(Y,X)$ 为余维 1 的正则对，并假设 $Y$
由方程 $x=0$ 定义。有层的正合列
$$
0\longrightarrow (\GG_{m})_{X}\longrightarrow
j_{\ast}(\GG_{m})_{U}\longrightarrow \ZZ_{Y}\longrightarrow
0{\quoi},
$$
于是应用（\hyperref[XIX.lem1.3]{1.3}），得到典范同构
\begin{equation*}
(R^{1}j_{\ast})𝛍_{n}\simeq
j_{\ast}(\GG_{m})_{U}/(j_{\ast}(\GG_{m})_{U})^{n}\simeq
(\ZZ/n)_{Y}{\quoi},\tag{3.3.1}
\end{equation*}\etiquette[3.3.1]{XIX.eq3.3.1}
这些同构显然不依赖于 $x$ 的选取。因此，有典范同构
\begin{equation*}
\varphi:\ZZ/n\xrightarrow{\sim}
(R^{2}i^{!})𝛍_{n}\simeq i^{\ast}(R^{1}j_{\ast})𝛍_{n}{\quoi}.\tag{3.3.2}
\end{equation*}\etiquette[3.3.2]{XIX.eq3.3.2}

\begin{enonce*}{定理 3.4}\etiquette[3.4]{XIX.thm3.4}

\footnote[1]{依赖于奇点消解（\cf 导言）。此定理也可看成
\hyperref[XIX.thm3.2]{3.2} 与第 XVIII \Exp 中借助“局部基本类”建立的变性理论的结合；
在那里，当 $Y$ 在 $X$ 中局部为完全交时，研究了
同态 $\varphi(Y,X)$（它们不必为双射）。}。\pageoriginale
设 $C$ 为等特征优概形的正则对所成的族。
存在唯一函数 $\varphi$，它将每个余维 $c$ 的正则对
$(Y,X)$ 对应到一个层同构（称为典范同构）
$$
\varphi(Y,X)=\varphi:\ZZ/n\xrightarrow{\sim}
ℋ^{2c}_{Y}(X,𝛍_{n}^{\otimes
c})=(R^{2c}i^{!})𝛍_{n}^{\otimes c}
$$
并使 $\varphi$ 满足以下条件：

\begin{enumerate}
\renewcommand{\theenumi}{\roman{enumi}}
\renewcommand{\labelenumi}{(\theenumi)}
\item 若 $c=1$，且 $Y$ 在 $X$ 中由方程 $x=0$ 定义，
% SGA-SRC-001210
则 $\varphi$ 为态射 \eqhyperref[XIX.eq3.3.2]{3.3.2}。

\item $\varphi$ 与正则对的光滑态射相容；即若
$(Y',X')\rightarrow (Y,X)$ 为一个态射，设 $g:Y'\rightarrow Y$ 为
诱导的态射。则 $Y'$ 上的层的图式
$$
\xymatrix{
\ZZ/n  \ar[rr]^{\varphi(Y',X')} & &
(R^{2c}i^{!})𝛍_n^{{\otimes}c}\\
g^{\ast}(\ZZ/n) \ar[u]\ar[rr]^{g^{\ast}\varphi(Y,X)}& &
g^{\ast}(R^{2c}i^{!})𝛍_n^{{\otimes}c} \ar[u]
}
$$
交换（特别地，这意味着右边的竖箭头为同构）。

\item \emph{传递性：}设 $(Z,Y,X)$ 为正则三元组，即
正则预概形之间的闭浸入的交换图式，等等，
$$
\xymatrix{
Z \ar[rr]^{v}\ar[rd]_{w}&    &Y\ar[ld]^{u}\\
            & X  & .
}
$$
分别\pageoriginale 设 $a$、$b$、$c$ 为 $u$、$v$、$w$ 的余维
\emph{（故 $c=a+b$）}。有同构
\begin{align*}
\varphi(Y,X)\otimes 𝛍_{n} {}^{\otimes
b}:𝛍_{n} {}^{\otimes b} \longrightarrow &
 (R^{2a}u^{!})(𝛍_{n}^{\otimes a}) \otimes
𝛍_{n} {}^{\otimes b}\simeq\\
& \simeq (R^{2a}u^{!})𝛍_{n}{}^{\otimes c}{\quoi}.
\end{align*}
谱序列 $(R^{p}v^{!})(R^{q}u^{!})F\Rightarrow
(R^{p+q}w^{!})F$ 给出同构
$$
\xi : (R^{2b}v^{!})(R^{2a}u^{!})𝛍_{n}{}^{\otimes
c}\xrightarrow{\sim} (R^{2c}w^{!})𝛍_{n}{}^{\otimes
c}{\quoi}.
$$
于是得到同构的图式
% SGA-SRC-001212 — tensor typography checked against the scan; missing \varphi reconstructed from transitivity
$$\vcenter{\etiquette[3.4.1]{XIX.eq3.4.1}
\xymatrix{
\ZZ/n \ar[r]^{\varphi(Z,Y)} \ar[d]^{\varphi(Z,X)}&
(R^{2b}v^{!})𝛍_{n}^{\otimes b} \ar[d]^{\eta =
  (R^{2b}v^{!})(\varphi(Y,X)\otimes 𝛍_{n}^{\otimes b})}\\
(R^{2c}w^{!})𝛍_{n}^{\otimes c} &
(R^{2b}v^{!})(R^{2a}u^{!})𝛍_{n}^{\otimes c}\ar[l]
}}\leqno{(3.4.1)}
$$
传递性断言就是这一图式交换。
\end{enumerate}
\end{enonce*}

\medskip
\noindent
\emph{证明。}存在性与唯一性对余维 $c$ 作归纳。
当 $c=1$ 时，由 \eqhyperref[XIX.eq3.3.1]{3.3.1} 局部地定义
态射 $\varphi$。由于此态射不依赖于参数 $x$ 的选择，
通过粘合，对每个余维 1 的 $(Y,X)$ 得到一个同构 $\varphi$；
由该态射的定义及 \hyperref[XIX.lem1.3]{1.3}，
显然这样构造的 $\varphi$ 满足 (ii)。$c=1$ 时不涉及传递性。

现\pageoriginale 假设余维 $<c$ 时存在性及唯一性已获证明，
且 $c>1$。应用归纳假设，立即可知在
\hyperref[XIX.thm3.4]{3.4} (iii) 的情形下，
\eqhyperref[XIX.eq3.4.1]{3.4.1} 中已定义的三个箭头，
即 $\varphi (Z,Y)$、$\xi$、$\eta$，
与正则三元组的态射相容。在此情形下置
$$
\Psi (Z,Y,X)=\xi \eta \varphi (Z,Y){\quoi},
$$
于是 $\Psi$ 也与三元组的态射相容。
若 $(Z,X)$ 为余维 $c>1$ 的正则对，
则在 $X$ 上局部地可找到一个满足 $0<\codim(Y,X)<c$ 的三元组
$(Z,Y,X)$。因此，只须证明箭头 $\Psi(Z,Y,X)$ 不依赖于 $Y$。
事实上，这将证明余维 $c$ 时 $\varphi$ 的唯一性，
而存在性随即由粘合得到。

若 $(Z,Y,X)$ 任意且 $\codim(Y,X)>1$，
则在 $X$ 上局部地可找到正则三元组 $(Y,W,X)$，
使 $(W,X)$ 的余维为 1。留给读者的一个关于传递性的图追踪论证，
就将问题归结为在 $(Y,X)$ 的余维为 1 时验证与 $Y$ 无关。

设 $(Z,Y_{0},X)$ 及 $(Z,Y_{1},X)$ 为两个这样的三元组。
只须逐纤维验证该断言，因此可归结为 $X$ 严格局部，
而各 $Y_{i}$ 由一个方程定义的情形，设方程为 $f_{i}=0$
$(i=0,1)$。设 $\bar{X}=\sheaf{\Spec}O_{X}[t]$，
$\bar{Z}=\fprod{Z}{\bar{X}}{X}$，并设 $\bar{Y}$ 为
$\bar{X}$ 中由方程
$$
(f_{1}-f_{0})t+f_{0}=0{\quoi}.
$$
定义的闭子概形。则\pageoriginale $\bar{Y}$ 在位于
$\bar{X}/X$ 的截面 $\epsilon_{0}:\{t=0\}$ 及 $\epsilon_{1}:\{t=1\}$
上方的点处显然正则。设 $C$ 为 $\bar{Y}$ 中由
$\bar{Y}$ 不正则的点组成的闭子集（\EGA IV 7.8.3 (iv)），
并分别将 $\bar{X}$、$\bar{Y}$、$\bar{Z}$ 换为
$\bar{X}-C$、$\bar{Y}-C$、$\bar{Z}-C\cap \bar{Z}$。
则 $(\bar{Z},\bar{Y},\bar{X})$ 为正则三元组，
且各截面 $\epsilon_{i}$ 给出三元组的态射
$$
(Z,Y_{i},X)\longrightarrow (\bar{Z},\bar{Y},\bar{X})\quad
(i=0,1){\quoi}.
$$

设 $F=(R^{2c}w^{!})𝛍_{n}{}^{\otimes c}$ 为 $Z$ 上的层。
它是一个同构于 $\ZZ/n$ 的局部常值层
（取任一个同构 $\Psi (Z,Y_{i},X)$ 即可），
立即可见，$\bar{Z}$ 上的层
$\bar{F}=(R^{2c}\bar{w}^{!})𝛍_{n}{}^{\otimes c}$
是 $F$ 在 $\bar{Z}$ 上的逆像。由于 $\bar{Z}\rightarrow Z$ 的纤维
连通且非空，而 $\bar{Z}\rightarrow Z$ 光滑，因而是开映射，
可知诱导态射
$$
H^{0}(Z,F)\longrightarrow H^{0}(\bar{Z},\bar{F})
$$
为\emph{双射}；因而以下各态射
$$\vcenter{
\xymatrix@R=.5cm{
                 &                                & H^{0}(Z,F)\\
H^{0}(Z,F) \ar[r]& H^{0}(\bar{Z},\bar{F})\ar[ru]^{\epsilon^{\ast}_0} \ar[rd]_{\epsilon^{\ast}_1} &           \\
                 &                                & H^{0}(Z,F)
}}\leqno{(*)}
$$
均为双射，故 $\epsilon^{\ast}_{0}=\epsilon^{\ast}_{1}$。

给出层 $F$ 与 $\ZZ/n$ 的一个同构，等价于给出
一个为生成元的整体截面（即 1 的像）。设 $\bar{\alpha}$、
$\alpha_{i}$ 分别为同构\pageoriginale
$\Psi(\bar{Z},\bar{Y},\bar{X})$、$\Psi(Z,Y_{i},X)$
给出的 $\bar{F}$、$F$ 的整体截面。
由于 $\Psi(.,.,.)$ 与态射相容，图式 $(\ast)$ 给出
$$
\alpha_{0}=\epsilon^{\ast}_{0}(\bar{\alpha})=\epsilon^{\ast}_{1}(\bar{\alpha})=\alpha_{1}{\quoi},
$$
所以 $\Psi(Z,Y_{0},X)=\Psi(Z,Y_{1},X)$，\cqfd

\section{正则态射的局部无环性}\etiquette[4]{XIX.sec4}

\begin{enonce*}{定理 4.1}\etiquette[4.1]{XIX.thm4.1}
\footnote[1]{依赖于奇点消解（\cf 导言）。}
\emph{（局部无环性）。}设
$g:X'\rightarrow X$ 为等特征优概形之间的正则态射
\emph{（\EGA IV 6.8.1）}。则对与特征互素的 $n$，
$g$ 为普遍局部无环的。
\end{enonce*}

\medskip
\noindent
\emph{证明}。由于假设在有限型基变换下稳定，
由 \hyperrefext[XV][XV.prop1.13]{1.13} ii) 立即可知，
只须证明 $g$ 对 $n$ 局部无环。
应用 \hyperrefext[XV][XV.cor1.17]{1.17} 的判别法：
设 $\bar{x}'$ 为 $X'$ 的一个几何点，$\bar{x}$ 为 $X$ 的相应几何点。
由正则态射的定义立即可知，$g$ 在这些点处诱导的
严格局部化之间的态射仍然正则。于是归结为
$X=\Spec A$ 和 $X'=\Spec A'$ 严格局部，且 $g$ 为局部态射的情形。
由 \hyperrefext[XV][XV.cor1.17]{1.17}，只须证明
$g$ 的\pageoriginale 几何纤维对 $n$ 无环。
将 $X$ 换为一个整闭子概形，$X'$ 换为相应的闭子概形，
便归结为几何泛纤维 $X'_{\bar{x}}$ 的情形。
取 $\bar{x}=\Spec \bar{K}$，其中 $\bar{K}$ 为
$X$ 的有理函数域的可分闭包。则
$$
\bar{x}=\varprojlim_{\alpha}U_{\alpha}
$$
其中 $U_{\alpha}$ 遍历一个连通平展 $X$-概形的滤过范畴，
并可假设它们\emph{正则}（\EGA IV 7.8.3 (v)）。
设 $B_{\alpha}$ 为 $A$ 在有理函数域
$R(U_{\alpha})$ 中的正规化。
置 $U'_{\alpha}=\fprod{X'}{U_{\alpha}}{X}$，
$B'_{\alpha}=A'\otimes_{A}B_{\alpha}$。
则 $U_{\alpha}$ 为 $\Spec B_{\alpha}$ 的一个正则开子概形，
且 $B_{\alpha}\rightarrow B'_{\alpha}$
% SGA-SRC-001217
正则，故由 \hyperref[XIX.lem2.4]{引理 2.4}，对每个 $q$ 有
$$
H^{q}(U_{\alpha},𝛍_{n})\xrightarrow{\sim}
H^{q}(U'_{\alpha},𝛍_{n})
$$
。纤维 $X'_{\bar{x}}$ 是各 $U'_{\alpha}$ 的极限，因而
\begin{align*}
H^{q}(X'_{\bar{x}},𝛍_{n}) &= \varinjlim_{\alpha}
H^{q}(U'_{\alpha},𝛍_{n})\\
&= \varinjlim_{\alpha} H^{q}(U_{\alpha}, 𝛍_{n})\\
 &= H^{q}(\bar{x},𝛍_{n})=0\quad\mbox{\rm 若}\quad q>0,\cqfd.
\end{align*}

\begin{enonce*}{推论 4.2}\etiquette[4.2]{XIX.cor4.2}
（沿正则态射的基变换）。设
$$
\xymatrix@C=2.25cm@R=1cm{
X \ar[d]^{f}& X'\ar[l]_{g'} \ar[d]^{f'}\\
Y & Y'\ar[l]_{g}
}
$$
为笛卡尔图式，其中 $g:Y'\rightarrow Y$ 为等特征优预概形之间的正则态射，
且 $f$ 拟紧且拟分离。\pageoriginale 则对每个挠性与特征互素的
阿贝尔挠层 $F$，基变换态射
\emph{（XII \hyperref[XII.eq4.2]{4.2}）}
$$
g^{\ast}(R^{q}f_{\ast})F\longrightarrow (R^{q}f'_{\ast}){g'}^{\ast}F
$$
对每个 $q\geq 0$ 均为双射。
\end{enonce*}

这由 \hyperref[XIX.thm4.1]{4.1} 及
\hyperrefext[XVI][XVI.thm1.1]{1.1} 得到。

由 \hyperref[XIX.thm4.1]{4.1} 及 \hyperrefext[XV][XV.cor1.17]{1.17}
立即推出

\begin{enonce*}{推论 4.3}\etiquette[4.3]{XIX.cor4.3}
设 $A$ 为等特征优严格局部环，
$f:X\rightarrow \Spec A$ 为一个态射。设 $A\rightarrow A'$ 为正则态射，
其中 $A'$ 也是等特征优严格局部环。
用 $'$ 表示由 $A\rightarrow A'$ 诱导的基变换。
对 $X$ 上的每个挠性与 $A$ 的特征互素的挠层 $F$，有
$$
H^{q}(X,F)\xrightarrow{\sim} H^{q}(X',F')
$$
对每个 $q\geq 0$ 成立。
\end{enonce*}

特别地，当 $A'$ 为 $A$ 的完备化 $\hat{A}$ 时，此推论适用。
（在这一情形下，假设“$A$ 严格局部”还可换为
“$A$ 为 Hensel 局部环”，这是立即可见的。）

\section{有限性定理}\etiquette[5]{XIX.sec5}\pageoriginale

\begin{enonce*}{定理 5.1}\etiquette[5.1]{XIX.thm5.1}
\footnote[1]{\hypertarget{sga4.exp.XIX.note.resolution}{}依赖于奇点消解（\cf 导言）。}
\emph{（有限性）。}设 $f:X\rightarrow S$ 为等特征优概形之间的有限型态射，
$F$ 为 $X$ 上的可构造阿贝尔挠层，其挠性与特征互素。
则对每个 $q$，$(R^{q}f_{\ast})F$ 同样可构造。
\end{enonce*}

此结果推广了 \hyperrefext[XVI][XVI.thm5.1]{5.1} 中的定理。
证明几乎相同：应用纯性结果
\hyperref[XIX.thm3.2]{3.2}，代替 \hyperrefext[XVI][XVI.thm3.7]{3.7}。

\begin{enonce*}{推论 5.2}\etiquette[5.2]{XIX.cor5.2}
设 $A$ 为等特征优严格局部环，
$f:X\rightarrow \Spec A$ 为有限型态射。
设 $F$ 为 $X$ 上的可构造阿贝尔层，
其挠性与 $A$ 的特征互素。则对每个 $q$，群 $H^{q}(X,F)$ 均为有限群。
\end{enonce*}

\section{仿射态射的上同调维数}\etiquette[6]{XIX.sec6}

设 $f:X\rightarrow Y$ 为优预概形之间的有限型态射，
$F$ 为 $X$ 上的层。回顾，当 $Y$ 局部时，
已定义了一个整数 $\delta(F)$（\hyperrefext[XIV][XIV.def2.2]{2.2}）。
因此，对任意 $Y$，可将\pageoriginale $\delta$ 看成一个函数 $\delta(F,f)$，置
\begin{equation*}
\delta(F,f,y)=\delta(F'){\quoi},\tag{6.0}
\end{equation*}\etiquette[6.0]{XIX.eq6.0}
其中 $Y'=\Spec O_{Y,y}$，$X'\rightarrow Y'$ 为
由 $f$ 经基变换 $Y'\rightarrow Y$ 所得的态射，
$F'$ 为 $F$ 在 $X'$ 上的逆像层。

有以下结果，它推广了 \hyperrefext[XVI][XVI.def3.1]{3.1}：

\begin{enonce*}{定理 6.1}\etiquette[6.1]{XIX.thm6.1}
\sgarepeatedfootnotemark{sga4.exp.XIX.note.resolution}。设 $f:X\rightarrow Y$ 为有限型仿射态射，
其中 $Y$ 为等特征优概形，$F$ 为 $X$ 上的
阿贝尔挠层。则在每个满足 $\delta(F,f,y)<q$ 的点
$y\in Y$ 处，$R^{q}f_{\ast}F$ 均为零。
\end{enonce*}

以下是局部的变形，它显然等价于 \hyperref[XIX.thm6.1]{6.1}：

\begin{enonce*}{定理 6.1 bis}\etiquette[6.1 bis]{XIX.thmbis6.1}
\sgarepeatedfootnotemark{sga4.exp.XIX.note.resolution}。设 $A$ 为等特征优严格局部环，
$f:X\rightarrow Y=\Spec A$ 为有限型仿射态射。
设 $F$ 为 $X$ 上的挠层。则
$$
H^{q}(X,F)=0\quad\text{\rm 若}\quad q>\delta(F){\quoi}.
$$
\end{enonce*}

回顾（\hyperrefext[X][X.thm5.1]{5.1}），当 $F$ 为 $p$-挠层，
其中 $p$ 为特征时，这类断言没有特别的意义。\pageoriginale
因此可假设 $F$ 的挠性与特征互素。

有以下推论：

\begin{enonce*}{推论 6.2}\etiquette[6.2]{XIX.cor6.2}
设 $A$ 为等特征优 Hensel 环，剩余域为 $k$，
$U\subset \Spec A$ 为仿射开子概形。
设 $\ell \in \PP$。则
$$
\cd_{\ell}U\leq \cd_{\ell}k+\dim A{\quoi}.
$$
\end{enonce*}

事实上，设 $\widetilde{A}$ 为 $A$ 的严格局部化，
它是群为 $G=G(\bar{k}/k)$ 的无限 Galois 平展覆盖，
并设 $\widetilde{U}$ 为 $U$ 在 $\Spec \widetilde{A}$ 中诱导的仿射开子概形。
应用 Hochschild--Serre 谱序列（VIII \hyperref[VIII.prop8.4]{8.4}）
$$
H^{p}(G,H^{q}(\widetilde{U},.))\Longrightarrow H^{p+q}(U,.)
$$ 
便归结为 $A=\widetilde{A}$ 的情形，而这就是
\hyperref[XIX.thm6.1]{6.1} bis 的特殊情形。

采用以上记号，设 $R(A)$ 为 $A$ 上的有理函数环。
有 $\Spec R(A)=\varprojlim_{\alpha}U_{\alpha}$，其中 $U_{\alpha}$
为 $\Spec A$ 的仿射开子概形。因此，取极限理论给出

\begin{enonce*}{推论 6.3}\etiquette[6.3]{XIX.cor6.3}
设 $A$ 为等特征优 Hensel 环，剩余域为 $k$。
设 $R(A)$ 为 $\Spec A$ 上的有理函数域。
则对 $\ell \in \PP$，有
$$
\cd_{\ell}R(A)\leq \cd_{\ell}k+\dim A
$$
且\pageoriginale 当 $\ell$ 在 $k$ 中可逆，并且另有
$\ell\neq 2$ 或 $k$ 不可序时，等号成立。
\end{enonce*}

后一情形下的等号由 \hyperrefext[X][X.cor2.4]{2.4} 得到。

最后，应用 \hyperrefext[X][X.cor4.2]{4.2}、
\hyperrefext[X][X.cor4.4]{4.4}，得到

\begin{enonce*}{推论 6.4}\etiquette[6.4]{XIX.cor6.4}
设 $X$ 为等特征 Noether 优概形，$R(X)$ 为
$X$ 上的有理函数环。设 $\ell\in \PP$ 在 $X$ 上可逆，
并假设 $\ell\neq 2$ 或 $X$ 的剩余域均不可序。则
$$
\cd_{\ell}(X)\leq \cd_{\ell}R(X)+\dim X\leq 2\cd_{\ell}R(X){\quoi}.
$$
\end{enonce*}

\medskip
\noindent
\emph{\hyperref[XIX.thm6.1]{6.1} 的证明}。首先作与
\hyperrefext[XIV][XIV.sec3]{3} 类似的归约。对 $\delta$ 作归纳：
显然，$\delta\leq d$ 时的 \hyperref[XIX.thm6.1]{6.1} 等价于
$\delta\leq d$ 时的（\hyperref[XIX.thm6.1]{6.1} bis）。
$\delta=0$ 时结论平凡。以 \hyperref[XIX.thm6.1]{6.1} bis 的形式
处理 $\delta\leq 1$ 的情形。由于 $F$ 为其各可构造子层的
滤过归纳极限（\hyperrefext[IX][IX.prop2.9]{2.9} (iii)），
可假设 $F$ 可构造；此时 $F$ 的支集包含在某个闭子集
$Z\subset X$ 中，且对每个 $z\in Z$ 有 $\delta(z)\leq 1$
（\hyperrefext[XIV][XIV.def2.2]{2.2}）。于是，$Z$ 的每个约化不可约分支
$Z_{i}$，或者是 $Y$ 的闭点 $y$ 上的维数 $\leq 1$ 的有限型仿射概形，
或者是在 $Y$ 的某个维数为 1 的约化闭子概形 $C$ 上拟有限的概形。
在后一情形中，$C=\Spec B$，其中 $B$ 为维数 1 的严格局部整环。
这样的概形 $Z_{i}$ 是维数 $\leq 1$ 的严格局部整环的谱，
或者是 $B$ 的分式域的某个有限扩张域的谱。
因此，由 \hyperrefext[X][X.thm2.3]{2.3} 及
\hyperrefext[IX][IX.cor5.7]{5.7}，无论何种情形，
$Z_{i}$ 的上同调维数至多为\pageoriginale 1。有正合列
$$
0\longrightarrow F\longrightarrow \prod_{i} F_{i}\longrightarrow
\epsilon \longrightarrow 0
$$
其中 $F_{i}$ 为 $F$ 在 $Z_{i}$ 上的限制，由此即得 $Z$ 的结果。

现假设定理对各 $\delta<d$ 已获证明，
来证明 $\delta(F)\leq d$ 时的结论，并假设 $d\geq 2$。

\hypertarget{sga4.exp.XIX.lemma.6.5}{}\begin{enonce*}{引理 6.5}\etiquette[6.5]{XIX.lem6.5}
只须处理 $X=\EE^{1}_{Y}=\sheaf{\Spec}O_{Y}[t]$
为 $1$ 维仿射空间，且 $Y$ 严格局部的情形。
\end{enonce*}

证明与 \hyperrefext[XIV][XIV.lem4.2]{4.2} 相同：
显然可假设 $Y$ 严格局部，且 $\delta(F)\leq d$。
须证明 $q>d$ 时 $H^{q}(X,F)=0$。可将 $X$ 嵌入
$\EE^{N}_{Y}$，因此只须处理 $X=\EE^{N}_{Y}$ 的情形。
对 $N$ 作归纳：考虑图式
$$
\xymatrix@C=1.2cm@R=1cm{
\EE^{N}_Y \ar[rr]^{g} \ar[rd]_{f}&   & \EE^{N-1}_Y \ar[ld]^{h}\\
                     & Y   &  
}
$$
于是 $\EE^{N}_{Y}$ 是 $\EE^{N-1}_{Y}$ 上的 1 维仿射空间。
应用归纳假设及 Leray 谱序列
$$
H^{p}(\EE^{N-1},R^{q}g_{\ast}F)\Longrightarrow
H^{p+q}(\EE^{N},F)
$$
可知只须证明
$$
\delta(R^{q}g_{\ast}F)\leq d-q{\quoi}.
$$
这\pageoriginale 就是说，若 $z\in \EE^{N-1}$ 满足
$\delta(z)>d-q$，则 $R^{q}g_{\ast}F=0$ 在点 $z$ 处成立。
然而，关于 $N$ 的归纳假设适用于态射 $g$，
故若 $\delta(F,g,z)<q$，则 $R^{q}g_{\ast}F=0$ 在点 $z$ 处成立。
所以只须证明不等式
$$
\delta(F,g,z)+\delta(z)\leq \delta(F){\quoi}.
$$
这一乐趣留给读者。

\hypertarget{sga4.exp.XIX.lemma.6.6}{}\begin{enonce*}{引理 6.6}\etiquette[6.6]{XIX.lem6.6}
为证明各 $\delta\leq d$ 时的 $\hyperref[XIX.thm6.1]{6.1}$，只须证明：
设 $B$ 为等特征、维数 $d$ 的完备正规局部环。
设 $U\subset \Spec B$ 为仿射开子集。
假设存在元素 $x\in \rad B$ 在 $U$ 上可逆。则
$$
H^{q}(U,\ZZ/\ell)=0\quad\text{\rm 若}\quad q>d{\quoi}.
$$
\end{enonce*}

此外，此引理是推论 \hyperref[XIX.cor6.2]{6.2} 的特殊情形。

\medskip
\noindent
\emph{证明}（\cf \hyperrefext[XIV][XIV.lem4.4]{4.4}）。由 \hyperref[XIX.lem6.5]{6.5}，
只须处理 $Y$ 严格局部且 $X=\EE^{1}_{Y}$ 的情形。考虑图式
$$
\xymatrix@C=1cm@R=.8cm{
\EE^{1}_Y \ar[rr]^{i} \ar[rd]_{f}&   & \PP^{1}_Y \ar[ld]^{g}\\
          & Y &,
}
$$
其中 $\PP^{1}_{Y}$ 为射影直线，$\EE^{1}_{Y}$ 是由
$\PP^{1}$ 去掉无穷远截面 $Y_{\infty}$ 得到的。
设 $F$ 为 $\EE^{1}$ 上的挠层，且 $\delta(F)\leq d$。有谱序列
$$
H^{p}(\PP^{1},R^{q}i_{\ast}F)\Longrightarrow
H^{p+q}(\EE^{1},F){\quoi}.
$$
当 $q>0$ 时，\pageoriginale $R^{q}i_{\ast}F$ 集中于
同构于 $Y$ 的截面 $Y_{\infty}$ 上。
由于 $Y$ 严格局部，当 $p>0$ 且 $q>0$ 时，有
$H^{p}(\PP^{1},R^{q}i_{\ast}F)=0$。
此外，由 \hyperrefext[XII][XII.cor5.3]{5.3}，当 $p>2$ 时有
$H^{p}(\PP^{1}, i_{\ast}F)=0$。
由于要证明当 $n>d$ 时 $H^{n}(\EE^{1},F)=0$，且 $d\geq 2$，只须证明
$$
H^{0}(\PP^{1},
R^{q}i_{\ast}F)=H^{0}(Y_{\infty},R^{q}i_{\ast}F)=0\quad\text{\rm
若}\quad q>d{\quoi}.
$$
此群同构于 $R^{q}i_{\ast}F$ 在闭纤维中
$\PP^{1}$ 的无穷远点 $y_{\infty}$ 处的茎。

现在只须（VII \hyperref[VII.prop3.3]{3.3} 及 \hyperrefext[IX][IX.prop2.9]{2.9} (iii)）
处理 $F$ 可构造的情形。此时，由于 $\delta(F)\leq d$，
存在 $\PP^{1}$ 的一个闭子概形 $X$，满足 $\dim X\leq d$，
且 $X$ 包含 $F$ 的支集。设 $\Spec A=\widetilde{X}$ 为
$X$ 在点 $y_{\infty}$ 处的严格局部化，
$\widetilde{U}=\widetilde{X}-\fprod{\widetilde{X}}{Y_{\infty}}{\PP^{1}}$，
并设 $\widetilde{F}$ 为 $F$ 在 $\widetilde{U}$ 上诱导的层。
注意，$\widetilde{U}$ 由 $\widetilde{X}$ 对元素
$x\in \Gamma(X,O_{X})$ 局部化得到。有
$$
(R^{q}i_{\ast}F)_{y_{\infty}}\simeq
H^{q}(\widetilde{U},\widetilde{F}){\quoi}.
$$
因此，只须证明对 $\widetilde{U}$ 上的每个挠层 $\widetilde{F}$，
若 $q>d$，则 $H^{q}(\widetilde{U},\widetilde{F})=0$。

现在可应用 \hyperrefext[IX][IX.prop5.6]{5.6}，取 $\varphi=\delta$，
可知只须证明
$$
H^{q}(V,\ZZ/\ell)=0\quad\text{\rm 若}\quad q>\delta(V)
$$
对\pageoriginale 每个在 $\widetilde{U}$ 上有限且整的 $V$ 以及可逆的 $\ell$ 成立。
设 $x:\bar{V}\rightarrow V$ 为 $V$ 的正规化。有正合列
$$
0\longrightarrow (\ZZ/\ell)_{V}\longrightarrow
x_{\ast}(\ZZ/\ell)_{\bar{V}}\longrightarrow C\longrightarrow 0
$$
且 $\delta(C)<\delta(V)$；由此（对 $\delta$ 作归纳）
可知只须取 $V$ \emph{正规}。

设 $B$ 为 $A$ 在 $V$ 上的有理函数域中的正规化。
则 $B$ 为等特征优严格局部环，$V$ 是 $\Spec B$ 的一个仿射开子集，
由对 $x$ 局部化得到。设 $\hat{B}$ 为 $B$ 的完备化，
$\hat{V}$ 为诱导的开子集（它由 $\Spec \hat{B}$
对 $x$ 在 $\hat{B}$ 中的像局部化得到）。由 \hyperref[XIX.cor4.3]{4.3}，有
$$
H^{q}(V,\ZZ/\ell)\simeq H^{q}(\hat{V},\ZZ/\ell){\quoi}.
$$
因此，归结为证明
$$
H^{q}(\hat{V},\ZZ/\ell)=0\quad\text{\rm 若}\quad q>\delta
(V)=\dim B=\dim \hat{B}{\quoi}.
$$
由归纳假设，若 $\delta(V)<d$，则定理成立，
故定理 \hyperref[XIX.thm6.1]{6.1} bis 适用于此情形。
因此只须处理 $\dim \hat{B}=d$ 的情形；
由于 $\hat{B}$ 正规，此时正是引理的情形，于是得到结果。

\setcounter{subsection}{7}
\setcounter{subsubsection}{-1}
\subsubsection{}\etiquette[6.7.0]{XIX.subsubsec6.7.0}
首先在附加假设下证明 \hyperref[XIX.lem6.6]{6.6}，其中
$B$ 的特征为 $p>0$。设 $A$ 为等特征、等维的完备环，
$k\subset A$ 为使 $A/\rad A$ 为 $k$ 的有限扩张的子域。
称\pageoriginale $k$-代数 $A$ \emph{形式可分}（formally separable），
若存在 $A$ 的一个形如 $k[[x_{1},\ldots,x_{n}]]$ 的 $k$-子代数，
使态射 $\Spec A\rightarrow \Spec k[[x]]$ 有限，
且在 $\Spec k[[x]]$ 的泛点处平展
$(x=(x_{1},\ldots,x_{n}))$。

\hypertarget{sga4.exp.XIX.lemma.6.7}{}\begin{enonce*}{引理 6.7}\etiquette[6.7]{XIX.lem6.7}
设 $k\subset B$ 为使 $B/\rad B$ 在其上有限的子域。
在上述关于 $d$ 的归纳假设下，若 $\hyperref[XIX.lem6.6]{6.6}$ 的假设满足，
且另有以下任一条件满足，则 \hyperref[XIX.lem6.6]{6.6} 的结论成立：
\begin{enumerate}
\renewcommand{\theenumi}{\roman{enumi}}
\renewcommand{\labelenumi}{(\theenumi)}
\item $k$ 的特征为 $0$。

\item 存在元素 $x\in \rad B$，它在 $U$ 上可逆，
且 $\bar{B}=B/(x)$ 为形式可分的 $k$-代数。
\end{enumerate}
\end{enonce*}

\medskip
\noindent
\emph{证明}。选取 $\rad B$ 中的元素 $x_{1}=x$、
$x_{2},\ldots, x_{d}$，使 $\{x_{1},\ldots,x_{d}\}$
生成一个根为 $\rad B$ 的准素理想。则 $B$ 为有限 $A$
-代数，其中 $A=k[[x_{1},\ldots,x_{d}]]$，
而 $\bar{B}$ 为 $k[[x_{2},\ldots, x_{d}]]=\bar{A}$-代数 $B/x_{1}B$。
在情形 (ii) 中，由于 $\bar{B}$ 形式可分，
可选取元素 $x_{2},\ldots,x_{d}$，使
$\bar{B}$ 成为泛点处平展的 $k[[x_{2},\ldots,x_{d}]]$-代数。
立即可知，$B$ 也是泛点处平展的 $A$-代数，
而在情形 (i) 中这当然也成立。

设 $b\in B$ 为生成 $R(B)$ 在 $k((x))$ 上的域扩张的元素。
设 $\epsilon \in k[[x]]$ 为 $b$ 在 $k[[x]]$ 上的
首一不可约方程的判别式。在情形 (ii) 中，
可选择 $b$，使 $x_{1}$ 不整除 $\epsilon$。
另写\pageoriginale $\Spec B-U=V(x_{1})\cup Y$，
其中 $Y$ 不包含 $V(x_{1})$ 的任何不可约分支。
设 $C$ 为环 $\Gamma(Y,O_{Y})$，这里赋予 $Y$ 诱导的约化结构。
必要时更换 $x_{2},\ldots,x_{d}$，可假设
$V(x_{d})$ 不包含在 $Y$ 在 $\Spec A$ 中的像内，
并且在情形 (ii) 中，$V(x_{d})$ 也不包含在 $V(\epsilon)$ 中。

\hypertarget{sga4.exp.XIX.lemma.6.8}{}\begin{enonce*}{引理 6.8}\etiquette[6.8]{XIX.lem6.8}
设 $A^{0}=k[[x_{1},\ldots,x_{d-1}]]$ $\{x_{d}\}$ 为
$k[[x_{1},\ldots,x_{d-1}]]$ $[x_{d}]$ 在点
$x_{1}=x_{2}=\ldots=x_{d}=0$ 处的严格局部化。
在上述条件下，存在一个 $A^{0}$-代数 $B^{0}$、
一个仿射开子集 $U^{0}\subset \Spec B^{0}$，以及
一个同构 $A\otimes_{A^{0}}B^{0}\simeq B$，
使 $U^{0}$ 在 $\Spec B$ 中诱导的开子集为 $U$。
\end{enonce*}

此引理的证明与 \hyperref[XIX.lem1.5]{1.5} 类似：
只须找出 $A^{0}$-代数 $B^{0}$、$C^{0}$ 及同构
$A\otimes_{A^{0}}B^{0}\simeq B$、$A\otimes_{A^{0}}C^{0}\simeq C$，
因为此时存在唯一的（满射）态射 $B^{0}\rightarrow C^{0}$，
诱导出态射 $B\rightarrow C$
（\cite{sga4.exp.XIX.bib.3}~1.4）。
因此取 $U^{0}=\Spec B^{0}-\{(\Spec C^{0})\cup V(x_{1})\}$，
立即验证 $U$ 由 $U^{0}$ 诱导。
$C^{0}$ 的存在性，以及情形 (ii) 中 $B^{0}$ 的存在性，
均与 \hyperref[XIX.lem1.5]{1.5} 中一样证明。
情形 (i) 中 $B^{0}$ 的存在性由
（\cite{sga4.exp.XIX.bib.3} 5.1）得到。
（也可以将 Abhyankar 引理应用于沿 $\{x_{1}=0\}$ 的分歧，
再使用下降，直接证明。）

现在可以完成 \hyperref[XIX.lem6.7]{6.7} 的证明。
再次应用 \hyperref[XIX.cor4.3]{4.3}，得到
$$
H^{q}(U,\ZZ/\ell)\simeq H^{q}(U^{0},\ZZ/\ell){\quoi},
$$
由此\pageoriginale 可知，只须证明当 $q>d$ 时后一个群为零。
$\Spec A^{0}$ 是在 $k[[x_{1},\ldots,x_{d-1}]]$ $[x_{d}]$ 上
仿射且平展的概形的极限，故 $\Spec B^{0}$ 是
$\Spec k[[x_{1},\ldots,x_{d-1}]]$ 上的有限型概形的极限，
而这些概形的各纤维维数均 $\leq 1$。
由于 $U^{0}\subset \Spec B^{0}[1/x_{1}]$，有
$$
U^{0}=\varprojlim U^{0}_{\alpha}
$$
其中 $U^{0}_{\alpha}$ 在概形
$V^{0}=\Spec k[[x_{1},\ldots, x_{d-1}]]$ $[1/x_{1}]$ 上仿射且有限型，
并且 $U^{0}_{\alpha}/V^{0}$ 的各纤维维数均 $\leq 1$。

设 $f^{0}:U^{0}\rightarrow V^{0}$ 为结构态射。
由于 $f^{0}_{\alpha}$ 的各纤维维数均 $\leq 1$，对
$v\in V^{0}$ 有
$$
\delta(\ZZ/\ell, f^{0}_{\alpha},v)\leq \dim O_{V^{0},v}+1\leq
d-1
$$
（因为 $\dim O_{V^{0},v}\leq d-2$）。故关于
% SGA-SRC-001224
$d$ 的归纳假设适用于态射 $f^{0}_{\alpha}$；
应用 \hyperref[XIX.thm6.1]{6.1}，得到
$$
R^{q}f^{0}_{\alpha\ast}(\ZZ/\ell)=0\quad\text{\rm 在点}\quad v\quad\text{\rm 若}\quad q>\dim O_{V^{0},v}+1{\quoi},
$$
即若 $q>d-\dim\bar{v}$。
因此，$R^{q}f^{0}_{\alpha\ast}(\ZZ/\ell)\neq 0$ 在点 $v$ 处成立
意味着 $\dim\bar{v}\leq d-q$。
于是，对于包含 $i:V^{0}\rightarrow \Spec
k[[x_{1},\ldots, x_{d-1}]]$（\cf \eqhyperref[XIX.eq6.0]{6.0}），有
$$
\delta (R^{q}f^{0}_{\alpha\ast}(\ZZ/\ell))\leq d-q{\quoi},
$$
将 \hyperref[XIX.thm6.1]{6.1} bis 形式的归纳假设
应用于态射\pageoriginale $i$，得到
$$
H^{p}(V^{0},R^{q}f^{0}_{\alpha\ast}(\ZZ/\ell))=0\quad\text{\rm
若}\quad p+q>d{\quoi}.
$$
因此 Leray 谱序列
$$
E^{pq}_{2}=H^{p}(V^{0},R^{q}f^{0}_{\alpha\ast}(\ZZ/\ell))\Longrightarrow
H^{p+q}(U^{0}_{\alpha},\ZZ/\ell)
$$
表明 $n>d$ 时 $H^{n}(U^{0}_{\alpha},\ZZ/\ell)=0$；
由于 $U^{0}$ 是各 $U^{0}_{\alpha}$ 的极限，同样得到
$$
H^{n}(U^{0},\ZZ/\ell)=0\quad\text{\rm 若}\quad n>d{\quoi}.
$$

这就完成了 \hyperref[XIX.lem6.7]{6.7} 的证明，
从而完成了特征 0 时定理 \hyperref[XIX.thm6.1]{6.1} 的证明。

\section{仿射态射——证明的完成}\etiquette[7]{XIX.sec7}

回顾，为完成特征 $p>0$ 时的 \hyperref[XIX.thm6.1]{6.1} 的证明
（仍在奇点消解假设下），已由 \hyperref[XIX.lem6.6]{6.6}
归结为对 $d\geq 2$ 证明以下断言：

\hypertarget{sga4.exp.XIX.lemma.7.1}{}\begin{enonce*}{引理 7.1}\etiquette[7.1]{XIX.lem7.1}
设 $B$ 为等特征完备正规局部环，剩余域可分闭，维数为 $d$。
设 $U\subset \Spec B$ 为仿射开子集，且存在 $x\in \rad B$
在 $U$ 上可逆。则
$$
H^{q}(U,\ZZ/\ell)=0\quad\text{\rm 若}\quad q>d{\quoi}.
$$

此外，\pageoriginale 由归纳假设，可假设
$\hyperref[XIX.thm6.1]{6.1}$ 对各 $\delta<d$ 成立，且
$\hyperref[XIX.lem7.1]{7.1}$ 已在特殊情形
$\hyperref[XIX.lem6.7]{6.7}$ {\rm (ii)} 下得到证明。
\end{enonce*}

首先注意，$B$ 正规的条件并不重要。
事实上，设 $B$ 任意，$\pi:\bar{B}\rightarrow B$
为 $B$ 的正规化，$\bar{U}\subset\Spec \bar{B}$ 为 $U$ 的逆像。
有正合列
$$
0\longrightarrow (\ZZ/\ell)_{U}\longrightarrow
\pi_{\ast}(\ZZ/\ell)_{\bar{U}}\longrightarrow C\longrightarrow
0{\quoi},
$$
其中 $C$ 的支集维数 $<d$，故当 $q>d-1$ 时有 $H^{q}(U,C)=0$。
由于 $H^{q}(U,\pi_{\ast}\bigdot)\simeq
H^{q}(\bar{U},\bigdot)$（VIII \hyperref[VIII.cor5.6]{5.6}），可知
\begin{equation*}
\begin{split}
& H^{q}(U,\ZZ/\ell)=0\quad \text{\rm 若}\quad q>d\quad
\text{\rm 等价于}\\
& H^{q}(\bar{U},\ZZ/\ell)=0\quad\text{\rm 若}\quad q>d{\quoi}.
\end{split}\tag{7.2}
\end{equation*}\etiquette[7.2]{XIX.eq7.2}

\setcounter{subsection}{3}
\setcounter{subsubsection}{-1}
\subsubsection{}\etiquette[7.3.0]{XIX.subsubsec7.3.0}
简要回顾形式可分概念所需的性质，不给证明。
读者可参阅（\EGA IV 18.11.10, 18.11.11）：
设 $A$ 为完备 Noether 局部 $k$-代数，
其剩余域为 $k$ 的有限扩张，并用 $A'$ 表示局部环
$A\otimes_{k}k'$ 的完备化（其中 $k'=k^{1/p}$）。
则扩张 $A\rightarrow A'$ 是\emph{纯不可分}（radicial）的
（可能为无限扩张），这很容易看出——这是 Nagata 的著名例子。
因此，为计算平展上同调，可将 $A$ 换为 $A'$ 而不影响结果（VIII 1）。

现在，一个维数\pageoriginale $d$ 的等维代数 $A$
形式\emph{可分}的条件，即它为 $k[[x_{1},\ldots, x_{d}]]$
的有限且泛点处平展的扩张这一条件，可以用模
$\widehat{\Omega}^{1}_{A/k}$（\EGA IV 18.11.10），
或 Jacobian 判别法（\EGA $0_{IV}$ 22.6, IV 7.1）
\footnote[1]{本讨论班的共同编者未能弄清这一引用的意义，
因为所引之处的情形在他看来颇不相同；
他也未能弄清下一个断言的理由。A. G.}
表述；由此可见，若 $A$ 不形式可分，则 $k'$-代数 $A'$ 不约化。
例如，为证明此事，可以归结为 $A$ 由一个方程 $f=0$ 给定的情形，
其中 $f\in k[[x_{1},\ldots,x_{d+1}]]$，
并应用 Jacobian 判别法及以下引理，
其有趣的证明留给读者：

\hypertarget{sga4.exp.XIX.lemma.7.3}{}\begin{enonce*}{引理 7.3}\etiquette[7.3]{XIX.lem7.3}
设 $f(x_{1},\ldots,x_{n})\in k[[x_{1},\ldots,x_{n}]]$ 为形式幂级数。
假设对每个 $i=1,\ldots,n$，存在可逆元素
$u_{i}\in k[[x_{1},\ldots,x_{n}]]$，使 $u_{i}f$
为关于 $x_{i}^{p}$ 及
$x_{1},\ldots,x_{i-1},x_{i+1},\ldots,x_{n}$ 的级数。
则存在可逆元素 $v\in k[[x_{1},\ldots,x_{n}]]$，使
$$
vf\in k[[x_{1}^{p},\ldots,x_{n}^{p}]]{\quoi}.
$$
\end{enonce*}

容易得到，有限次应用
% SGA-SRC-001226
将 $(A,k)$ 换为 $(\bar{A}',k')$ 的步骤，
其中 $\bar{A}'$ 为 $A'$ 的正规化，
便可归结为 $A$ 确实形式可分的情形。
同样可证明，对给定的 $x\in\rad A$，
对代数 $A$ 应用同样步骤后，
环 $(A/(x))_{\text{\rm réd}}$ 将变为形式可分代数。

应用 VIII \hyperref[VIII.thm1.1]{1.1} 及 \hyperref[VIII.def7.2]{7.2}，得到

\hypertarget{sga4.exp.XIX.corollary.7.4}{}\begin{enonce*}{推论 7.4}\etiquette[7.4]{XIX.cor7.4}
只须在如下附加条件下证明 $\hyperref[XIX.lem7.1]{(7.1)}$：
$B$ 包含一个域 $k$，使 $B/\rad B$ 在其上有限，
且 $k$-代数 $(B/(x))_{\text{\rm réd}}$ 形式可分。\pageoriginale
\end{enonce*}
 
这一情形下 \hyperref[XIX.lem7.1]{7.1} 的证明，
是将它归结为 \hyperref[XIX.lem6.7]{6.7} (ii) 中处理的特殊情形：

采用 \hyperref[XIX.lem7.1]{7.1}、\hyperref[XIX.cor7.4]{7.4} 的记号，设
$$
D(x)=\sum_{i}r_{i}X_{i}
$$
为 $x$ 的除子，故每个 $X_{i}$ 是 $\Spec B$ 的
一个余维 1 的闭不可约子概形，且在 $k$ 上形式可分
（即 $B/\mathfrak{J}(X_{i})$ 为形式可分 $k$-代数，
而这又等价于 $(B/(x))_{\text{\rm réd}}$ 形式可分）。
选择 $y\in\rad B$，使有
$$
D(y)=\sum_{i}(r_{i}-1)X_{i}+\Gamma
$$
其中 $\Gamma$ 不包含任何 $X_{i}$，并考虑理想 $(x,y)$ 的爆破（blow-up）
$$
f:Z\longrightarrow \Spec B
$$
。环 $B[y/x]$ 和 $B[x/y]$ 的谱构成 $Z$ 的一个仿射开覆盖。

设 $B_{1}$ 为 $B[x/y]$ 在由 $(\rad B,x/y)$ 生成的极大理想处的严格局部化。
对 $\Spec B[y/x]\rightarrow \Spec B$ 的闭纤维的每个闭点 $Q$，
记 $A_{Q}$ 为 $B[y/x]$ 在 $Q$ 处的严格局部化，
$U_{Q}$ 为 $\Spec A_{Q}$ 中作为 $U$ 的逆像的仿射开子集。

\hypertarget{sga4.exp.XIX.lemma.7.5}{}\begin{enonce*}{引理 7.5}\etiquette[7.5]{XIX.lem7.5}
为\pageoriginale 对环 $B$ 证明 $\hyperref[XIX.lem7.1]{7.1}$，
只须证明当 $q>d(\geq 2)$ 时有
\begin{equation*}
H^{q}(\Spec B_{1}[y/x],\ZZ/\ell)=0\tag{i}
\end{equation*}
及
\begin{equation*}
H^{q}(U_{Q},\ZZ/\ell)=0\quad\text{对每个 $Q$}{\quoi}.\tag{ii}
\end{equation*}
\end{enonce*}

\medskip
\noindent
\emph{证明。}有开浸入
$$
U\hookrightarrow \Spec B[1/x]\hookrightarrow \Spec B[y/x]{\quoi},
$$
从而有图式
$$
\xymatrix{
U \ar[rr]^{i} \ar[rd]_{k}&         & \Spec B[y/x] \ar[ld]^j\\
                        &   Z     &      . 
}
$$
若 $Q$ 为 $\Spec B[y/x]\rightarrow \Spec B$ 的闭纤维的一个闭点，
则由假设，当 $q>d$ 时有 $(R^{q}i_{\ast}\ZZ/\ell)_{Q}=0$。
对 $\Spec B[y/x]$ 的每个不是
$\Spec B[y/x]\rightarrow \Spec B$ 的闭纤维中的闭点的点 $P$，
有 $\delta(F,i,P)<d$，其中 $F=(\ZZ/\ell)_{U}$
（\eqhyperref[XIX.eq6.0]{6.0} 的记号）。
因此，将归纳假设及 \hyperref[XIX.thm6.1]{6.1} 应用于态射 $i$，
可知对于态射 $\Spec B[y/x]\rightarrow \Spec B$ 有
$$
\delta(R^{q}i_{\ast}\ZZ/\ell)\leq d-q{\quoi}.
$$
所以，将 \hyperref[XIX.thm6.1]{6.1} bis 形式的归纳假设
应用于后一个态射，得到
$$
H^{p}(\Spec B[y/x],R^{q}i_{\ast}\ZZ/\ell)=0
$$
若\pageoriginale $q>0$ 且 $p>d-q$。利用 Leray 谱序列
$$
E^{pq}_{2}=H^{p}(\Spec
B[y/x],R^{q}i_{\ast}\ZZ/\ell)\Longrightarrow
H^{p+q}(U,\ZZ / \ell){\quoi}.
$$
证明 \hyperref[XIX.lem7.1]{7.1} 的问题归结为证明
$$
H^{p}(\Spec B[y/x],i_{\ast}\ZZ/\ell)=0\quad\text{\rm 若}\quad
p>d{\quoi}.
$$

有正合列
\begin{equation*}
0\longrightarrow \ZZ/\ell\longrightarrow
i_{\ast}(\ZZ/\ell)_{U}\longrightarrow C\longrightarrow 0\tag{7.6}
\end{equation*}\etiquette[7.6]{XIX.eq7.6}
其中 $\delta(C)<d$。因此（再次由归纳法），只须证明
\begin{equation*}
H^{p}(\Spec B[y/x],\ZZ/\ell)=0\quad\text{\rm 若}\quad
p>d{\quoi}.\tag{7.7}
\end{equation*}\etiquette[7.7]{XIX.eq7.7}

开浸入 $\Spec B[y/x]\xrightarrow{j}Z$ 由从 $Z$ 中去掉
“无穷远处”的闭子集 $C$ 得到，而在 $Z$ 的另一个开子集
$\Spec B[x/y]$ 中，它与闭子集 $C=V(x/y)$ 等同。
各 $R^{q}j_{\ast}\ZZ/\ell$ 集中在 $C$ 上，
且态射 $C\rightarrow \Spec B$ 为闭浸入。
因此，当 $p>0$ 且 $q>0$ 时有 $H^{p}(Z,R^{q}j_{\ast}\ZZ/\ell)=0$。
由于 $Z\rightarrow \Spec B$ 的闭纤维维数为 1，且该态射固有，
由 \hyperrefext[XII][XII.cor5.3]{5.3} bis 可知
$p>2$ 时有 $H^{p}(Z,j_{\ast}\ZZ/\ell)=0$。所以谱序列
$$
E^{pq}_{2}=H^{p}(Z,R^{q}j_{\ast}\ZZ/\ell)\Longrightarrow
H^{p+q}(\Spec B[y/x],\ZZ/\ell)
$$
给出同构
$$
H^{q}(\Spec B[y/x],\ZZ/\ell)\simeq
H^{0}(Z,R^{q}j_{\ast}\ZZ/\ell)\simeq
(R^{q}j_{\ast}\ZZ/\ell)_{P}
$$
若\pageoriginale $q>2$，其中 $P$ 为 $Z$ 的“无穷远点”。
这些茎正是（VIII 5）
% SGA-SRC-001231
$$
(R^{q}j_{\ast}\ZZ/\ell)_{P}\simeq H^{q}(\Spec
B_{1}[y/x],\ZZ/\ell){\quoi}.
$$

由 \eqhyperref[XIX.eq7.7]{7.7}，只须证明后面的这些群在 $q>d$ $(\geq 2)$ 时为零。
有一个与 \eqhyperref[XIX.eq7.6]{7.6} 类似的正合列，
结合归纳假设，它表明只须有
$$
H^{q}(\Spec B_{1}[y/x],\ZZ/\ell)=0\quad\mbox{\rm 若}\quad
q>d{\quoi},
$$
引理得证。

\medskip
\noindent
\emph{\hyperref[XIX.lem7.5]{7.5} (i) 的证明。}

考虑 $\Spec B[x/y]$ 的闭子集 $C=V(x/y)$，赋予它诱导的约化结构。
立即可见，态射 $Z\rightarrow \Spec B$ 诱导的
$C$ 到 $\Spec B$ 的闭浸入，将 $C$ 与
$\Spec B$ 的闭子概形 $X=\bigcup X_{i}$ 等同。
由于
\begin{align*}
D(x) &= \sum r_{i}X_{i}\\
D(y) &= \sum (r_{i}-1)X_{i}+\Gamma{\quoi},
\end{align*}
元素 $x/y$ 在 $X_{i}$ 的泛点附近局部生成理想
$\mathfrak{J}(X_{i})$；在该处，由于 $B$ 正规，局部环是离散赋值环。
因此，态射 $\Spec B[x/y]\rightarrow \Spec B$ 在这样的点的
一个邻域上为同构，而元素 $x/y$ 在 $C$ 的每个
不可约\pageoriginale 分支 $C_{i}$（对应于 $X_{i}$）上的消失阶数为 1。

由于闭子集 $C$ 可提升到 $\Spec B_{1}$，
将环 $B[x/y]$ 换为 $B_{1}$ 后上述结论仍成立，
故 $B_{1}/(x/y)$ 是一个\emph{完备}环，且为形式可分的 $k$-代数。

为证明 \hyperref[XIX.lem7.5]{7.5} (i)，首先可应用 \hyperref[XIX.cor4.3]{4.3}，
由此可将 $B_{1}$ 换为其完备化 $\hat{B}_{1}$，
而不改变环 $B_{1}/(x/y)$。由 \eqhyperref[XIX.eq7.2]{7.2}，
还可将 $\hat{B}_{1}$ 换为其正规化 $\bar{B}_{1}$。
由于 $B_{1}$（从而 $\hat{B}_{1}$）在各 $C_{i}$ 的泛点处已经正则，
态射 $\Spec \bar{B}_{1}\rightarrow \Spec \hat{B}_{1}$ 在这些点处为同构，
故 $\bar{B}_{1}/(x/y)$ 仍为形式可分的 $k$-代数。
于是归结为特殊情形 \hyperref[XIX.lem6.7]{6.7} (ii)，结果得证。

\medskip
\noindent
\emph{\hyperref[XIX.lem7.5]{7.5} (ii) 的证明。}

选取 $\Spec B[y/x]\rightarrow \Spec B$ 的闭纤维中的点 $Q$。
严格局部化 $A_{Q}$ 可写为极限
$$
A_{Q}=\varinjlim A
$$
其中 $A$ 遍历在 $B[y/x]$ 上平展且有限型的环的一个滤过范畴。
设 $U_{A}$ 为 $\Spec A$ 中作为 $U$ 的逆像的仿射开子集。
由于 $x$ 在 $U_{A}$ 上可逆，态射 $U_{A}\rightarrow \Spec B$ 平展，
且 $U_{A}$ 正规。因此，若 $C$ 表示 $B$ 在
$A$ 的有理函数域 $R(A)$ 中的正规化，则有开浸入
$$
U_{A}\longrightarrow \Spec C{\quoi}.
$$
只\pageoriginale 须（VII \hyperref[VII.thm5.7]{5.7}）证明，
对每个 $A$，若 $q>d$，则 $H^{q}(U_{A},\ZZ/\ell)$
在 $H^{q}(U_{Q},\ZZ/\ell)$ 中的像为零。
由于 $x$ 在 $U_{Q}$ 上可逆，若 $f\in B$ 为任一满足
$f$ 在 $B[y/x]$ 中被 $x$ 整除，且 $f/x$ 在点 $Q$ 处不为零的元素，
则态射 $U_{Q}\rightarrow \Spec A$ 通过 $\Spec A[1/f]$ 分解。
所以只须对某个适当的这样的 $f$ 证明
$$
H^{q}(U_{A,f},\ZZ/\ell)=0\quad\text{\rm 若}\quad q>d{\quoi},
$$
其中 $U_{A,f}=U_{A}-V(f)$ 为 $\Spec A[1/f]$ 的仿射开子集，
同时也是 $\Spec C$ 的仿射开子集。

设 $c\in k$ 使 $y/x+c$ 在 $Q$ 处不为零，
并设 $J$ 为 $B$ 中闭子集 $\Gamma\cap X$ 的理想。则若
\begin{equation*}
f\equiv y+cx(\text{\rm 模 }J^{N})\quad N>0{\quoi},\tag{7.8}
\end{equation*}\etiquette[7.8]{XIX.eq7.8}
元素 $f$ 便满足上述条件：事实上，立即可见在
$\Spec B[y/x]$ 中有 $V(J)=V(x)$。
故 \eqhyperref[XIX.eq7.8]{7.8} 确实意味着 $f$ 被 $x$ 整除，从而
$$
f/x\equiv y/x+c(\text{\rm 模 }J){\quoi}.
$$
由于 $J$ 在 $Q$ 处为零，而 $y/x+c$ 在 $Q$ 处非零，断言得证。

现在，$\Gamma\cap X$ 为 $\Spec B$ 的一个余维 2 的闭子集。
因此，可选取一个满足 \eqhyperref[XIX.eq7.8]{7.8} 的 $f$，
记为 $z$，使 $V(z)\cap V(y)$ 的余维也为\pageoriginale 2。
另取元素 $x_{3},\ldots,x_{d}$，使
$y$、$z$、$x_{3},\ldots,x_{d}$ 构成环 $B$ 的一个参数系。
则 $B$，从而 $C$，为有限 $k[[y,z,x_{3},\ldots,x_{d}]]$-代数。

设 $C^{0}$ 为 $k[[y,z,x_{3},\ldots,x_{d}]]$ 在
$k$ $((y,z,x))$ 于有理函数域 $R(C)$ 中的可分闭包内的正规化，
于是 $C$ 为 $C^{0}$ 的纯不可分扩张，
而 $C^{0}$ 为 $k$ $((y,z,x))$ 的泛点处平展的扩张
（其中 $x=x_{3},\ldots,x_{d}$）。

选择一个线性组合
$$
f=az+by\quad (a,b\in k){\quoi},
$$
使态射
$$
\Spec C^{0}\longrightarrow \Spec k[[z,y,x]]
$$
在 $V(f)$ 的泛点上方平展，且 $f/x$ 在
$\Spec B[y/x]$ 的 $Q$ 处不为零。只须证明
$$
H^{q}(U_{A,f},\ZZ/\ell)=0\quad\text{\rm 若}\quad q>d{\quoi},
$$
因而（VIII \hyperref[VIII.thm1.1]{1.1}）只须证明
$$
H^{q}(U^{0}_{A,f},\ZZ/\ell)=0\quad\text{\rm 若}\quad
q>d{\quoi},
$$
其中 $U^{0}_{A,f}$ 为 $\Spec C^{0}$ 中对应于
$\Spec C$ 的开子集 $U_{A,f}$ 的开子集；
由下降可知它是仿射的（\EGA IV 2.7.1 (xiii)）。
但由构造，$C^{0}$ 的元素 $f$ 在 $U^{0}_{A,f}$ 上可逆，
且 $C^{0}/(f)$ 为形式可分的 $k$-代数。
所以这正是 \hyperref[XIX.lem6.7]{6.7} (ii) 的情形，
从而完成了 \hyperref[XIX.lem7.5]{7.5} (ii) 的证明，
也就完成了 \hyperref[XIX.thm6.1]{6.1} 的证明。

\begin{thebibliography}{99}
\bibitem{sga4.exp.XIX.bib.1} \textsc{Abhyankar}, S.\pageoriginale Local uniformization on
algebraic surfaces over ground fields of characteristic $p\neq
0$. Annals of Math. vol. 63 (1963), and Resolution of singularities of
embedded algebraic surfaces, Academic Press (1966).

\bibitem{sga4.exp.XIX.bib.2} \textsc{Abhyankar}, S. Resolution of singularities of
arithmetical surfaces. Arithmetical Algebraic Geometry, New York
(1966).

\bibitem{sga4.exp.XIX.bib.3} \textsc{Artin}, M. Etale coverings of schemes over hensel
rings, Amer, Journal Math. vol 88 ${\rm n^{\circ}}$ 4 p 915-934
(1966).

\bibitem{sga4.exp.XIX.bib.4} \textsc{Hironaka}, H. Resolution of singularities of an
algebraic variety over a field of characteristic zero I, II. Annals of
Math. vol. 79, ${\rm n^{\circ}}$ 1, 2 (1964).

\bibitem{sga4.exp.XIX.bib.5} \textsc{Weierstrass}, K. Einige auf die analytischen
Functionen mehrerer Veränderlichen sich beziehende
Sätze. Mathematische Werke von K. Weierstrass, Bd 2, Berlin (1895).

\bibitem{sga4.exp.XIX.bib.6} \textsc{Bourbaki}, N. Commutative Algebra, Chap. 3.
\end{thebibliography}


%XII 4.2 no ref. 1072

\ifx\danslelivre\undefined
\let\SGAXIXenddocument\enddocument
\else
\let\SGAXIXenddocument\relax
\fi
\SGAXIXenddocument


% END native input: chapters/exp-19-zh.tex
\backmatter
\end{document}
